% 纲要标题：# 第二卷　代数拓扑与微分拓扑
% 纲要位置：几何学计划纲要.md，第 992 行。
% 正文状态：未撰写；本文件目前是框架。

% 纲要内容开始
% 
% > 围绕空间分解、不变量的构造、局部到整体及光滑实现组织。代数拓扑由基本拓扑和群论独立进入；微分拓扑另用第一卷流形、流和定向。谱序列、障碍与运算不再挤在同一层的提要中。
% 
% 纲要内容结束

% !TeX TXS-program:compile = cmd /c build-Book2.bat -NoPause
% !TeX TXS-program:view = txs:///view-pdf-internal --embedded "Book2.pdf"

\documentclass[UTF8,zihao=-4,openany,fontset=fandol]{ctexbook}
\newcommand{\BookTitle}{代数拓扑与微分拓扑}
\newcommand{\BookVolumeName}{第二卷}
\newcommand{\BookNumber}{2}
\newcommand{\BookDraftDate}{2026年10月}
% 《几何学》五卷共用的基础排版与语义宏。
% 各卷主文件须先定义 \BookTitle、\BookVolumeName 和 \BookDraftDate。
\newcommand{\BookAuthor}{R. EverStarine}
% 稿件状态由各卷元数据指定；尚未开写的卷沿用框架状态。
\providecommand{\BookStatus}{框架预览·正文未撰写}

\usepackage[a4paper,inner=28mm,outer=24mm,top=27mm,bottom=28mm,
  headheight=16pt,headsep=8mm,footskip=12mm]{geometry}

\usepackage{amsmath}
\usepackage{mathtools}
\usepackage{amsthm}
\usepackage{aliascnt}
\usepackage{unicode-math}
% 数学字母使用默认数学字体内的对应字形；禁止 mathrm 等默认继承正文族。
% 运算符（sin、log、operatorname 等）随 mathrm 使用数学正体。
\unimathsetup{mathrm=sym,mathit=sym,mathbf=sym,mathsf=sym,mathtt=sym}
\setmainfont{texgyrepagella-regular.otf}[
  BoldFont=texgyrepagella-bold.otf,
  ItalicFont=texgyrepagella-italic.otf,
  BoldItalicFont=texgyrepagella-bolditalic.otf
]
\setsansfont{texgyreheros-regular.otf}[
  BoldFont=texgyreheros-bold.otf,
  ItalicFont=texgyreheros-italic.otf,
  BoldItalicFont=texgyreheros-bolditalic.otf
]
\setmonofont{lmmonolt10-regular.otf}[
  BoldFont=lmmonolt10-bold.otf,
  ItalicFont=lmmonolt10-oblique.otf,
  BoldItalicFont=lmmonolt10-boldoblique.otf
]
% 人物注释中的古希腊文、斯拉夫文字、日文异体字与希伯来文原名使用有相应字形的字体。
\newfontfamily\GeometryOriginalNameFont{Times New Roman}
\NewDocumentCommand{\OriginalName}{m}{{\GeometryOriginalNameFont #1}}
\newCJKfontfamily\GeometryOriginalHanFont{SimSun}
\NewDocumentCommand{\OriginalHanName}{m}{{\GeometryOriginalHanFont #1}}
\newfontfamily\GeometryHebrewFont{Times New Roman}[Script=Hebrew]
\TeXXeTstate=1
\NewDocumentCommand{\OriginalHebrewName}{m}{{\GeometryHebrewFont\beginR#1\endR}}
% 行内、行间及标题内公式均使用 unicode-math 默认 Latin Modern Math。
% 公式之外的正文西文仍由 setmainfont 使用 Pagella；黑板粗体单独采用 AMS。
\let\mathbb\relax
\DeclareMathAlphabet{\mathbb}{U}{msb}{m}{n}
% unicode-math 的 Latin Modern Math 将 \varnothing 映为瘦长椭圆；空集单独
% 采用传统 AMS MSBM 的圆形斜杠字形，不改变其余数学符号的字体。
\DeclareSymbolFont{GeometryAMSb}{U}{msb}{m}{n}
\DeclareMathSymbol{\GeometryRoundEmptySet}{\mathord}{GeometryAMSb}{"3F}
% 当前 Latin Modern Math 缺少横截符号 U+22D4；仅此关系符使用 AMSa 的现有字形。
\DeclareSymbolFont{GeometryAMSa}{U}{msa}{m}{n}
\DeclareMathSymbol{\GeometryTransverse}{\mathrel}{GeometryAMSa}{"74}
\AtBeginDocument{%
  \let\varnothing\GeometryRoundEmptySet
  \let\emptyset\GeometryRoundEmptySet
  \let\pitchfork\GeometryTransverse
}

\usepackage{microtype}
\usepackage{xcolor}
\usepackage{graphicx}
\usepackage{booktabs}
\usepackage{longtable}
\usepackage{array}
\usepackage{tabularx}
\usepackage{algorithm}
\usepackage{algpseudocode}
\usepackage{enumitem}
\usepackage{fancyhdr}
% 双面排版补齐奇数页时，空白页不残留上一章的页眉或页码。
\usepackage{emptypage}
\usepackage[most]{tcolorbox}
\usepackage{tikz-cd}
\usepackage{csquotes}
\usepackage[
  backend=biber,
  style=numeric-comp,
  sorting=nyt,
  giveninits=true,
  maxbibnames=99,
  doi=true,
  url=true,
  isbn=true
]{biblatex}
\addbibresource{../Shared/References.bib}
% 骨架没有引文时保留书目入口；写入真实条目后由完整构建自动排入书目。
\newif\ifGeometryHasBibliography
\AtDataInput{\global\GeometryHasBibliographytrue}
\NewDocumentCommand{\PrintBookBibliography}{}{%
  \ifGeometryHasBibliography
    \printbibliography[heading=bibintoc,title={参考文献}]%
  \else
    \chapter*{参考文献}%
    \addcontentsline{toc}{chapter}{参考文献}%
    \noindent 待撰写\par
  \fi
}
% 访问日期保留在书目库供内部核验，成品参考文献不显示 visited on。
\renewbibmacro*{url+urldate}{\usebibmacro{url}}
% biblatex 在书目内重绑定格式宏；在其作用域中将强调统一为正体，
% 同时避免书名中的中文及标点仍继承旧的斜体样式。
\AtBeginBibliography{\let\emph\textup\let\textit\textup}
% 中文条目的责任者连接符：不可套用西文的 “and”，按 GB/T 7714—2015 的著录
% 惯例改用逗号分隔。中文条目在 Shared/References.bib 中以 langid = {chinese}
% 标注（含译本，因其著录信息为中文形式），西文条目保持 “and” 不变。
% 按条目切换须在每条书目项处理时判定；\DeclareDelimFormat 的可选参数不是
% 语言选择器（实测无效），故用 \AtEveryBibitem 配合 \iffieldequalstr。
\AtEveryBibitem{%
  \iffieldequalstr{langid}{chinese}%
    {\DeclareDelimFormat{finalnamedelim}{\addcomma\space}%
     \DeclareDelimFormat{multinamedelim}{\addcomma\space}}%
    {\DeclareDelimFormat{finalnamedelim}{\addspace\bibstring{and}\space}%
     \DeclareDelimFormat{multinamedelim}{\addcomma\space}}%
}
\usepackage{etoolbox}
\usepackage{xparse}
\usepackage[noautomatic]{imakeidx}

\definecolor{GeometryBlue}{HTML}{244A64}
\definecolor{GeometryBlueLight}{HTML}{EEF4F7}
\definecolor{GeometryGray}{HTML}{4F5962}
\definecolor{GeometryGrayLight}{HTML}{F4F5F6}

\ctexset{
  part = {
    % 编题页分两行：第一行“第X编”，与封面“第一卷”同款（大号无衬线灰色）；
    % 第二行编名，与封面“几何学”同款（大号无衬线蓝色粗体），字号保持原样。
    % 第一行与行距由 number／aftername 给出，format 只负责编名的字体，
    % 以免字号被重置或被重复输出。
    format = {\centering\fontsize{30}{38}\selectfont\sffamily\bfseries\color{GeometryBlue}},
    name = {,},
    number = {\centering\normalfont\normalsize{\fontsize{24}{30}\selectfont\sffamily\color{GeometryGray}第\chinese{part}编}},
    aftername = {\par\vspace{0.7cm}},
    beforeskip = 0pt,
    afterskip = 2.5\baselineskip,
    pagestyle = plain
  },
  chapter = {
    format = \raggedright\Huge\bfseries\color{GeometryBlue},
    name = {第,章},
    number = \chinese{chapter},
    aftername = \quad,
    beforeskip = 0pt,
    afterskip = 30pt,
    pagestyle = plain
  },
  section = {
    format = \Large\bfseries\color{GeometryBlue},
    beforeskip = 2.4ex plus .6ex minus .2ex,
    afterskip = 1.1ex plus .2ex
  },
  subsection = {
    format = \large\bfseries\color{GeometryGray},
    beforeskip = 2ex plus .4ex minus .2ex,
    afterskip = .8ex plus .2ex
  }
}

\setlength{\parindent}{2em}
\setlength{\parskip}{0pt}
\linespread{1.18}
% 仅用于两个连续的普通正文段落；列表、公式、定理、标题等结构边界不用。
% 段间额外留白统一由这一长度控制，避免改动全局 \parskip 后扰乱各类环境。
\newlength{\GeometryProseParagraphGap}
\AtBeginDocument{%
  \setlength{\GeometryProseParagraphGap}{0.40\baselineskip}%
}
\NewDocumentCommand{\ProseParagraphBreak}{}{%
  \par\addvspace{\GeometryProseParagraphGap}%
}
% 与代数学保持一致：小节编号并进入目录，便于按具体主题定位。
\setcounter{secnumdepth}{3}
\setcounter{tocdepth}{2}
\allowdisplaybreaks[2]
\emergencystretch=2em
\clubpenalty=10000
\widowpenalty=10000
\displaywidowpenalty=10000

% 固定列表内部段距，避免 flushbottom 将默认的弹性 \parsep 拉大到超过正文段距。
\setlist[itemize]{leftmargin=2.2em,itemsep=.35ex,parsep=0pt,topsep=.6ex}
\setlist[enumerate]{leftmargin=2.5em,itemsep=.35ex,parsep=0pt,topsep=.6ex,
  font=\normalfont\upshape}
\setlist[enumerate,1]{label=\arabic*.,ref=\arabic*}
\setlist[enumerate,2]{label=\arabic*),ref=\arabic*}
\setlist[enumerate,3]{label=\alph*),ref=\alph*}
\setlist[description]{leftmargin=2.5em,labelindent=0pt,itemsep=.35ex,parsep=0pt,topsep=.6ex}

% “下列命题等价”专用列表；普通罗列仍使用 enumerate。
\newlist{equivalences}{enumerate}{2}
\setlist[equivalences,1]{leftmargin=2.8em,itemsep=.35ex,parsep=0pt,topsep=.6ex,
  label=\Roman*.,ref=\Roman*,font=\normalfont\upshape}
\setlist[equivalences,2]{leftmargin=2.8em,itemsep=.35ex,parsep=0pt,topsep=.6ex,
  label=\roman*),ref=\roman*,font=\normalfont\upshape}

\pagestyle{fancy}
\fancyhf{}
\fancyhead[LE]{\small\color{GeometryGray}\thepage\quad\leftmark}
% 奇数页采用页前最后有效的节标记，避免页中开始新节时提前更换页眉。
% 新节或习题的标记从后续页生效；章首页使用 plain，不显示上一章的标记。
\fancyhead[RO]{\small\color{GeometryGray}\TopMark{2e-right}\quad\thepage}
\renewcommand{\headrulewidth}{0.35pt}
% 保留数学标题的原始大小写，避免 book 类的 \MakeUppercase
% 与 unicode-math 命令在页眉标记中发生冲突。
% 偶数页页眉显示“章序＋章标题”（如“第一章 代数对象、映射与基本构造”）；
% 章序与章标题之间以 \quad 分隔，间距与奇数页节眉“节号＋节标题”之间一致；
% 附录章的章序由 \CTEXthechapter 按本卷附录字母（如 A）给出。
\renewcommand{\chaptermark}[1]{\markboth{\CTEXthechapter\quad #1}{}}
\renewcommand{\sectionmark}[1]{\markright{\thesection\quad #1}}
% 章末习题不编号，但进入习题页后应覆盖上一节留下的奇数页页眉；
% 同时写入目录项。PDF 书签由 hyperref 对 \section*{习题} 自动生成，
% 不再另行 \pdfbookmark，避免“习题”下再嵌套一个“习题”。
\newcommand{\BookExercises}{%
  \section*{习题}%
  \markright{习题}%
  \addcontentsline{toc}{section}{习题}%
}
\fancypagestyle{plain}{
  \fancyhf{}
  \fancyfoot[C]{\small\color{GeometryGray}\thepage}
  \renewcommand{\headrulewidth}{0pt}
}

\usepackage[unicode,colorlinks=true,linkcolor=GeometryBlue,
  citecolor=GeometryBlue,urlcolor=GeometryBlue,pdfborder={0 0 0}]{hyperref}
\usepackage{bookmark}
% PDF 书签沿用目录条目中的编号：正文章为“第几章”，附录章为“附录 A”，
% 节与小节分别保留 x.x、x.x.x；无编号的导读、习题和卷末单元仍不加号。
\bookmarksetup{numbered}

% unicode-math 会调整原始 Unicode 希腊字母的类别码。若把这些字母直接写进
% \texorpdfstring 的书签分支，辅助文件回读时类别码不同，LaTeX 会在每一轮都
% 误报标签发生变化。书签分支统一写成下列可展开宏，既保留希腊字母显示，
% 又使写入和回读的标签文本稳定。
\newcommand{\GeometryPdfAlpha}{α}
\newcommand{\GeometryPdfGamma}{Γ}
\newcommand{\GeometryPdfLambda}{λ}
\newcommand{\GeometryPdfMu}{μ}
\newcommand{\GeometryPdfPi}{π}
\newcommand{\GeometryPdfSigma}{σ}
\newcommand{\GeometryPdfTheta}{θ}
\newcommand{\GeometryPdfXi}{ξ}
\newcommand{\GeometryPdfZeta}{ζ}
\newcommand{\GeometryPdfCapitalLambda}{Λ}
\newcommand{\GeometryPdfOmega}{Ω}
\newcommand{\GeometryPdfCapitalSigma}{Σ}
\newcommand{\GeometryPdfDelta}{δ}
\newcommand{\GeometryPdfEpsilon}{ε}
\newcommand{\GeometryPdfIota}{ι}
\newcommand{\GeometryPdfSmallOmega}{ω}
\newcommand{\GeometryPdfEll}{ℓ}

\usepackage[nameinlink,noabbrev]{cleveref}
\newcommand{\crefpairconjunction}{和}
\newcommand{\crefmiddleconjunction}{、}
\newcommand{\creflastconjunction}{和}

% CTeX 2.5.10 的标题前缀宏带一个参数，而 TeX Live 2025 中 cleveref
% 仍按普通前缀宏展开它们，会把节号写成含 \endcsname 的损坏元数据。
% 改为恒等包装后，CTeX 传入的标题编号保持不变，\ref 与 \cref 均可正常工作。
\makeatletter
\def\p@part#1{#1}
\def\p@chapter#1{#1}
\def\p@section#1{#1}
\def\p@subsection#1{#1}
\def\p@subsubsection#1{#1}
\def\p@paragraph#1{#1}
\def\p@subparagraph#1{#1}
\makeatother
% 非数学西文字母统一正体；不改变括号类别、标点或数学字体。
% 2026-09-11：非数学模式的西文字母统一使用 Pagella 正体，保留粗细区别。
% 保留字母级边界，不重分类数字、括号、标点及中文。
% 此设置也覆盖列表序号字母和术语原文；数学模式保持原样。
\newfontfamily\GeometryLetterFont{texgyrepagella-regular.otf}[
  BoldFont=texgyrepagella-bold.otf,
  ItalicFont=texgyrepagella-italic.otf,
  BoldItalicFont=texgyrepagella-bolditalic.otf
]
\ExplSyntaxOn
\bool_new:N \l_geometry_text_letters_bool
\bool_set_true:N \l_geometry_text_letters_bool
\cs_new_protected:Npn \GeometryPreserveMathText
  { \bool_set_false:N \l_geometry_text_letters_bool }
\cs_new_protected:Npn \geometry_letter_group_begin:
  {
    \group_begin:
    \bool_if:NT \l_geometry_text_letters_bool
      { \GeometryLetterFont \upshape }
  }
% 继承 Default 的全部 xeCJK 转移，保留中西文间距及标点处理；
% 不能用只有开关分组的转移覆盖 Boundary 或 CJK 的原处理。
\xeCJK_new_class:n { GeometryLatinLetter }
\seq_map_inline:Nn \g__xeCJK_class_seq
  {
    \str_if_eq:nnF {#1} { GeometryLatinLetter }
      {
        \xeCJK_copy_inter_class_toks:nnnn
          {#1} {GeometryLatinLetter} {#1} {Default}
        \xeCJK_copy_inter_class_toks:nnnn
          {GeometryLatinLetter} {#1} {Default} {#1}
        \xeCJK_app_inter_class_toks:nnn {#1} {GeometryLatinLetter}
          { \geometry_letter_group_begin: }
        \xeCJK_pre_inter_class_toks:nnn {GeometryLatinLetter} {#1}
          { \group_end: }
      }
  }
% ASCII 字母及常用拉丁扩展字母；明确跳过乘号、除号，不重分类标点。
\xeCJKDeclareCharClass { GeometryLatinLetter }
  { "0041->"005A, "0061->"007A, "00C0->"00D6, "00D8->"00F6,
    "00F8->"024F, "1E00->"1EFF }
\ExplSyntaxOff
% 数学组内局部关闭字母切换，连公式中 \text{...} 的字体也保持原样。
% 只追加，不覆盖 amsmath、unicode-math、xeCJK 已有的数学初始化。
\AtBeginDocument{%
  \everymath\expandafter{\the\everymath\GeometryPreserveMathText}%
  \everydisplay\expandafter{\the\everydisplay\GeometryPreserveMathText}%
}


% 定理类环境（定理、定义、命题、引理、推论、断言、猜想、开放问题、注记、
% 习题）的正文：中文用楷体，西文保持正体，数学字体不受影响。
% 注意 \normalfont 只重置西文字族，中文字族仍随默认（宋体），故须显式 \kaishu；
% 不能退回 \itshape——那会把西文一并变斜。
\newtheoremstyle{geometryplain}
  {1.2ex}{1.2ex}{\normalfont\upshape\kaishu}{}
  {\color{GeometryBlue}\bfseries}{.}{0.6em}{}
\newtheoremstyle{geometrydefinition}
  {1.2ex}{1.2ex}{\normalfont\kaishu}{}
  {\color{GeometryBlue}\bfseries}{.}{0.6em}{}
\newtheoremstyle{geometryremark}
  {1.0ex}{1.0ex}{\normalfont\kaishu}{}
  {\color{GeometryGray}\bfseries}{.}{0.6em}{}

\theoremstyle{geometryplain}
\newtheorem{theorem}{定理}[section]
% 正文默认按“章.节.流水号”编号；未分节的附录省略中间的 0。
\renewcommand{\thetheorem}{%
  \ifnum\value{section}=0
    \thechapter.\arabic{theorem}%
  \else
    \thesection.\arabic{theorem}%
  \fi
}
\newaliascnt{lemma}{theorem}
\newtheorem{lemma}[lemma]{引理}
\aliascntresetthe{lemma}
\newaliascnt{proposition}{theorem}
\newtheorem{proposition}[proposition]{命题}
\aliascntresetthe{proposition}
\newaliascnt{corollary}{theorem}
\newtheorem{corollary}[corollary]{推论}
\aliascntresetthe{corollary}
\newaliascnt{claim}{theorem}
\newtheorem{claim}[claim]{断言}
\aliascntresetthe{claim}
\newaliascnt{conjecture}{theorem}
\newtheorem{conjecture}[conjecture]{猜想}
\aliascntresetthe{conjecture}
\newaliascnt{openproblem}{theorem}
\newtheorem{openproblem}[openproblem]{开放问题}
\aliascntresetthe{openproblem}

\theoremstyle{geometrydefinition}
\newaliascnt{definition}{theorem}
\newtheorem{definition}[definition]{定义}
\aliascntresetthe{definition}
\newaliascnt{example}{theorem}
\newtheorem{example}[example]{例}
\aliascntresetthe{example}
% 习题按章单独编号，不占用定理类环境的流水号。
\newtheorem{exercise}{习题}[chapter]

\theoremstyle{geometryremark}
\newaliascnt{remark}{theorem}
\newtheorem{remark}[remark]{注记}
\aliascntresetthe{remark}
\renewcommand{\proofname}{证明}
% 仅在证明紧接定理类陈述时收紧间距，并按陈述结尾是否为行间公式调整。
% 证明标题（证明／解／证明概要）的中文用楷体，
% 西文保持正体；正文不施加楷体，仍为默认宋体。不能退回 \itshape。
\newif\ifGeometryDirectProof
\newif\ifGeometryCoupledProof
\newif\ifGeometryStatementEndsInDisplay
\newskip\GeometryProofTrailingSkip
% 行间公式结束后，TeX 留下空的水平列表；在定理环境结束前辨认这一状态。
% 证明跟在公式之后时不再减去定理尾距，否则公式下方会比上方更紧。
\newcommand{\GeometryRecordStatementTail}{%
  \ifhmode
    \ifnum\lastnodetype=-1\relax
      \global\GeometryStatementEndsInDisplaytrue
    \else
      \global\GeometryStatementEndsInDisplayfalse
    \fi
  \else
    \global\GeometryStatementEndsInDisplayfalse
  \fi
}
\AddToHook{para/begin}{\global\GeometryDirectProoffalse\global\GeometryStatementEndsInDisplayfalse}
\AddToHook{env/theorem/end}{\GeometryRecordStatementTail}
\AddToHook{env/lemma/end}{\GeometryRecordStatementTail}
\AddToHook{env/proposition/end}{\GeometryRecordStatementTail}
\AddToHook{env/corollary/end}{\GeometryRecordStatementTail}
\AddToHook{env/claim/end}{\GeometryRecordStatementTail}
\AddToHook{env/exercise/end}{\GeometryRecordStatementTail}
\AddToHook{env/theorem/after}{\global\GeometryDirectProoftrue}
\AddToHook{env/lemma/after}{\global\GeometryDirectProoftrue}
\AddToHook{env/proposition/after}{\global\GeometryDirectProoftrue}
\AddToHook{env/corollary/after}{\global\GeometryDirectProoftrue}
\AddToHook{env/claim/after}{\global\GeometryDirectProoftrue}
\AddToHook{env/exercise/after}{\global\GeometryDirectProoftrue}
\AddToHook{env/figure/before}{\global\GeometryDirectProoffalse}
\AddToHook{env/table/before}{\global\GeometryDirectProoffalse}
\makeatletter
\renewenvironment{proof}[1][\proofname]{\par
  \pushQED{\qed}%
  \normalfont \topsep6\p@\@plus6\p@\relax
  \GeometryCoupledProoffalse
  \ifGeometryDirectProof
    \GeometryCoupledProoftrue
    \GeometryProofTrailingSkip=\topsep
    \advance\GeometryProofTrailingSkip by \partopsep
    \ifGeometryStatementEndsInDisplay\else
      \kern-5\p@% 只在陈述以普通文字结尾时抵消部分尾距。
    \fi
    \topsep2\p@\partopsep\z@\relax
  \fi
  \global\GeometryDirectProoffalse
  \global\GeometryStatementEndsInDisplayfalse
  \trivlist
  \item[\hskip\labelsep\upshape\kaishu #1\@addpunct{.}]\ignorespaces
}{%
  \popQED
  \ifGeometryCoupledProof\@topsepadd\GeometryProofTrailingSkip\fi
  \endtrivlist\@endpefalse
}
\makeatother
\newenvironment{solution}
  {\begin{proof}[解]}
  {\end{proof}}
\newenvironment{proofsketch}
  {\begin{proof}[证明概要]}
  {\end{proof}}

% 正文选读内容的星号属于编号而不属于标题；第二参数统一接收标签。
\newcommand{\GeometryOptionalMark}{\texorpdfstring{\textsuperscript{*}}{*}}
% 选做习题沿用同一编号星号。可选参数仍只表示习题自身的简短题名，
% 不得再写“选读”或“选做”；星号由本环境自动附在题号右侧。
\NewDocumentEnvironment{optionalexercise}{o}{%
  \begingroup
  \let\GeometrySavedTheExercise\theexercise
  \renewcommand{\theexercise}{\GeometrySavedTheExercise\GeometryOptionalMark}%
  \IfNoValueTF{#1}{\begin{exercise}}{\begin{exercise}[#1]}%
}{%
  \end{exercise}%
  \endgroup
}
% 可选参数只供目录与 PDF 书签使用；正文题名可借此按语义手动换行，
% 又不把换行命令写入目录或书签。
\NewDocumentCommand{\OptionalChapter}{O{}mm}{%
  \begingroup
  \let\GeometrySavedTheChapter\thechapter
  \renewcommand{\thechapter}{\GeometrySavedTheChapter\GeometryOptionalMark}%
  \ctexset{chapter={name={第,章\GeometryOptionalMark},
    number=\chinese{chapter},aftername=\quad}}%
  \ifstrempty{#1}{\chapter{#2}}{\chapter[#1]{#2}}\label{#3}%
  \endgroup
}
\NewDocumentCommand{\OptionalSection}{mm}{%
  \begingroup
  \let\GeometrySavedTheSection\thesection
  \renewcommand{\thesection}{\GeometrySavedTheSection\GeometryOptionalMark}%
  \ctexset{section={number=\thesection}}%
  \section{#1}\label{#2}%
  \endgroup
}
\NewDocumentCommand{\OptionalSubsection}{mm}{%
  \begingroup
  \let\GeometrySavedTheSubsection\thesubsection
  \renewcommand{\thesubsection}{\GeometrySavedTheSubsection\GeometryOptionalMark}%
  \ctexset{subsection={number=\thesubsection}}%
  \subsection{#1}\label{#2}%
  \endgroup
}

% 编题页必须从奇数页开始，题页仅在页脚显示页码，其后保留一张空白页。
\makeatletter
\NewDocumentCommand{\BookPart}{mm}{%
  \cleardoublepage
  \begingroup
  \def\GeometryPartLabel{#2}%
  % \part 会在返回前翻页，故须在 \@endpart 开始处记录题页标签。
  \pretocmd{\@endpart}{\label{\GeometryPartLabel}}{}%
    {\PackageError{geometry}{无法为编题页写入标签}{请检查当前文档类的 \string\@endpart 定义。}}%
  \part{#1}%
  \endgroup
  \thispagestyle{empty}%
  \null
  \clearpage
}
% 各编的“本编导读”是无编号章，在总目录和 PDF 书签中与所属编的各章并列。
% addcontentsline 同时建立目录项和书签，不再另加同名 pdfbookmark。
\NewDocumentCommand{\BookPartGuide}{O{本编导读}}{%
  \chapter*{#1}%
  \addcontentsline{toc}{chapter}{#1}%
}

% “附录”单独排无编号编题页，使用与正文编名相同的字号、字体和颜色。
% 题页从奇数页开始，其后保留一张空白页，再进入第一个附录章。
% startatroot 使根级书签与书后顶层节点分离；phantomsection 须先于
% addcontentsline，目录中的“附录”行才能获得独立锚点而非沿用前一节的
% section*.N。根级“附录”书签连同其子章由 addcontentsline 的 part 行自动
% 建立，不得再手写 \pdfbookmark（实测会产生重复的“附录”节点）。
\NewDocumentCommand{\BookAppendixPart}{m}{%
  \cleardoublepage
  \begingroup
  \def\GeometryPartLabel{#1}%
  % 在 part* 的题页尚未翻走时建立唯一的目录与书签目标。
  \pretocmd{\@endpart}{%
    \bookmarksetup{startatroot}%
    \phantomsection
    \addcontentsline{toc}{part}{附录}%
    \label{\GeometryPartLabel}%
  }{}{\PackageError{geometry}{无法为附录编题页写入目录与标签}{请检查当前文档类的 \string\@endpart 定义。}}%
  \part*{附录}%
  \endgroup
  \thispagestyle{empty}%
  \null
  \clearpage
}
\makeatother

% 列表项的括号与编号整体承担链接，避免只让中间的编号可点击。
\NewDocumentCommand{\itemref}{m}{\hyperref[#1]{(\ref*{#1})}}

% 章、节、小节的引用：整个“第 X 章／第 X 节／第 X 小节”都是跳转链接。
% 链接必须走 hyperref 的标签解析机制（目标名 chapter.N／section.N.M）；
% 早先用 \hyperlink{裸标签名} 指向不存在的目标，链接全部失效。
%
% 章号用汉字（“第五章”）且必须取“被引用那一章”的编号；\thechapter 是当前
% 计数器，在编导读、前言或跨章引用时会给出错误章号（表现为“第零章”）。
% hyperref 下 \ref 受保护、\getrefnumber 在 \edef 中又失真，故改由 zref 记录：
% 各章在 \chapter 之后调用 \BookChapterNumber，把汉字号存为 zref 属性；
% \zref@extract 可展开，能在 \edef 中安全取出。
\usepackage[user,hyperref]{zref}
\makeatletter
\zref@newprop{geometrychapternum}{}
% 在 \chapter 之后调用，参数即该章标签；把本章汉字号存为 zref 属性
\NewDocumentCommand{\BookChapterNumber}{m}{%
  \zref@setcurrent{geometrychapternum}{\chinese{chapter}}%
  \zref@labelbyprops{zch@#1}{geometrychapternum}%
}
% 取出某章标签对应的汉字号（\zref@extract 可展开，能在 \edef 中安全使用）
\NewDocumentCommand{\GeometryChineseChapterOf}{m}{%
  \zref@extract{zch@#1}{geometrychapternum}}
\makeatother
% \ref* 的分组会遮住编号末端的西文字符类别；编号前的间距可由 xeCJK
% 自动产生，编号后的同类间距须在完整链接内部显式补回。
\NewDocumentCommand{\secref}{m}{%
  \hyperref[#1]{第\ref*{#1}\CJKecglue 节}}
\NewDocumentCommand{\subsecref}{m}{%
  \hyperref[#1]{第\ref*{#1}\CJKecglue 小节}}
\NewDocumentCommand{\chapref}{m}{%
  \begingroup
  \xeCJKsetup{CJKecglue={}}%
  \edef\GeometryChapterDisp{第\GeometryChineseChapterOf{#1}章}%
  \hyperref[#1]{\GeometryChapterDisp}%
  \endgroup}

% 算法采用带编号的步骤，并明确标出输入与输出。
\floatname{algorithm}{算法}
\renewcommand{\listalgorithmname}{算法目录}
\algrenewcommand\algorithmicrequire{\textbf{输入：}}
\algrenewcommand\algorithmicensure{\textbf{输出：}}
\numberwithin{algorithm}{chapter}
% 完整中文编号进入题注、各自目录和引用；附录自动使用字母章号。
% 内部 \nobreakspace 与 cleveref 默认的定理引用一致；引用整串置于同一链接，
% 链接后的中文间距由下方 \cref 包装补齐。
\renewcommand{\thealgorithm}{算法\nobreakspace\thechapter-\arabic{algorithm}}
\renewcommand{\thefigure}{图\nobreakspace\thechapter-\arabic{figure}}
\renewcommand{\thetable}{表\nobreakspace\thechapter-\arabic{table}}
\renewcommand{\figurename}{}
\renewcommand{\tablename}{}
\makeatletter
\renewcommand{\fnum@figure}{\thefigure}
\renewcommand{\fnum@table}{\thetable}
\renewcommand{\fnum@algorithm}{\thealgorithm}
\makeatother
\NewDocumentEnvironment{algorithmsteps}{mm}
  {\begin{algorithmic}[1]\Require #1\Ensure #2}
  {\end{algorithmic}}

\newtcolorbox{learninggoals}{
  enhanced,breakable,colback=GeometryBlueLight,colframe=GeometryBlue,
  boxrule=.45pt,arc=1mm,left=1.2mm,right=1.2mm,top=1mm,bottom=1mm,
  title=本章目标,fonttitle=\bfseries
}
\newtcolorbox{partguide}{
  enhanced,breakable,colback=GeometryGrayLight,colframe=GeometryGray,
  boxrule=.35pt,arc=1mm,left=1.2mm,right=1.2mm,top=1mm,bottom=1mm,
  title=本编导读,fonttitle=\bfseries
}
\newtcolorbox{chaptersummary}{
  enhanced,breakable,colback=GeometryBlueLight,colframe=GeometryBlue,
  boxrule=.45pt,arc=1mm,left=1.2mm,right=1.2mm,top=1mm,bottom=1mm,
  title=本章小结,fonttitle=\bfseries,
  % 与阅读提示框一致：仅本框内容为多段散文，须恢复首行缩进。
  before upper={\parindent=2em\parskip=0pt}
}
% 仅用于尚待高阶数学复核的核心断言；成稿前必须处理或移除。
\newtcolorbox{coreverification}{
  enhanced,breakable,colback=yellow!10,colframe=orange!85!black,
  boxrule=.55pt,arc=1mm,left=1.2mm,right=1.2mm,top=1mm,bottom=1mm,
  title=待数学复核,fonttitle=\bfseries
}
% 兼容旧稿的占位命令：只输出普通文字，不另加方框。
\NewDocumentCommand{\GeometryPlaceholder}{m}{待撰写}
\NewDocumentEnvironment{chapterabstract}{}
  {\small\list{}{\rightmargin\leftmargin
    \listparindent=2\ccwd\itemindent=\listparindent}%
    \item\relax\color{GeometryGray}}
  {\endlist}
% 引诗与短引文统一居中，中文使用定理类环境的楷体；章摘要已在上方独立定义，
% 不继承这里的字形与对齐方式。
\RenewDocumentEnvironment{quote}{}
  {\list{}{\rightmargin\leftmargin}\item\relax\centering\kaishu}
  {\endlist}
% 引诗环境保留为语义清楚的兼容入口，版式与 quote 完全一致。
\NewDocumentEnvironment{poemquote}{}
  {\begin{quote}}
  {\end{quote}}
% 引诗末尾的作者或出处在引诗区域内右对齐。
\NewDocumentCommand{\quoteattribution}{m}{%
  \par\smallskip
  {\raggedleft #1\par}%
}

\makeindex[name=chinese,title=中文索引,columns=2,intoc,options=-q]
\makeindex[name=foreign,title=外文索引,columns=2,intoc,options=-q]
\makeindex[name=symbols,title=符号索引,columns=2,intoc,options=-q]
% 索引采用稍小字号与自然栏底，字母组间仅留少量空白；末页平衡不得越过正文高度。
\indexsetup{othercode={\small\raggedcolumns\maxbalancingoverflow=0pt
  \renewcommand{\indexspace}{\par\addvspace{0.2\baselineskip}}}}

% 有条目时读入 makeindex 生成的真实索引；尚无条目时仍保留书后结构与目录入口。
% 一键脚本会在 .idx 为空时删除旧 .ind，故这里不会误排上一次构建的陈旧索引。
\NewDocumentCommand{\PrintBookIndex}{m m}{%
  \IfFileExists{#1.ind}{%
    % 双栏输出可能丢失页眉标记；在索引范围内直接绑定标题并完成分页。
    \cleardoublepage
    \begingroup
      \fancyhead[LE]{\small\color{GeometryGray}\thepage\quad #2}%
      \fancyhead[RO]{\small\color{GeometryGray}#2\quad\thepage}%
      \printindex[#1]%
      \clearpage
    \endgroup
  }{%
    \chapter*{#2}%
    \addcontentsline{toc}{chapter}{#2}%
    \noindent 待撰写\par
  }%
}

% 所有术语在每卷首次正式定义或引入时都显示完整原文，并登记中外文索引。
% 星号形式与缺失原文均视为错误，防止旧的隐藏原文规则被重新带回源码。
% 定理类正文处于楷体；Fandol 楷体没有粗体字形，直接 \textbf 会退回普通楷体。
% 术语主称因此局部恢复正文的宋体字族，再使用与普通正文一致的粗体。
\NewDocumentCommand{\term}{s O{} m o}{%
  \IfBooleanT{#1}{%
    \PackageError{geometry}{\string\term* 已取消}{请改用 \string\term，并保留末尾的原文注解。}%
  }%
  {\songti\textbf{\textup{#3}}}%
  \IfNoValueTF{#4}{%
    \PackageError{geometry}{术语缺少原文注解}{每个术语首次正式定义或引入时都必须写成 \string\term[排序键]{中文主称}[原文]。}%
  }{%
    \textnormal{(\textup{#4})}%
  }%
  \ifstrempty{#2}{\index[chinese]{#3}}{\index[chinese]{#2@#3}}%
  \IfNoValueF{#4}{\index[foreign]{#4}}%
}
\NewDocumentCommand{\termidx}{O{}m}{%
  \ifstrempty{#1}{\index[chinese]{#2}}{\index[chinese]{#1@#2}}%
}
% 别称与主称一样用 \termidx 登记实际页码，不生成“见某某”的转引条目。
\NewDocumentCommand{\foreignidx}{O{}m}{%
  \ifstrempty{#1}{\index[foreign]{#2}}{\index[foreign]{#1@#2}}%
}
\newcounter{geometrySymbolIndexOrder}
\NewDocumentCommand{\symidx}{m}{%
  \stepcounter{geometrySymbolIndexOrder}%
  \edef\GeometrySymbolSortKey{%
    \number\numexpr10000+\value{geometrySymbolIndexOrder}\relax}%
  \index[symbols]{\GeometrySymbolSortKey@#1}%
}

% 一般强调与术语粗体分离。
\NewDocumentCommand{\keypoint}{m}{\CJKunderdot{#1}}
\NewDocumentCommand{\foreignkeypoint}{m}{\underline{\textup{#1}}}


\crefname{chapter}{章}{章}
\crefname{part}{编}{编}
\crefname{section}{节}{节}
\crefname{subsection}{小节}{小节}
\crefname{theorem}{定理}{定理}
\crefname{lemma}{引理}{引理}
\crefname{proposition}{命题}{命题}
\crefname{corollary}{推论}{推论}
\crefname{claim}{断言}{断言}
\crefname{conjecture}{猜想}{猜想}
\crefname{openproblem}{开放问题}{开放问题}
\crefname{definition}{定义}{定义}
\crefname{example}{例}{例}
\crefname{remark}{注记}{注记}
\crefname{exercise}{习题}{习题}
\crefname{equation}{式}{式}
\crefname{figure}{}{}
\crefname{table}{}{}
\crefname{algorithm}{}{}
\crefformat{figure}{#2#1#3}
\Crefformat{figure}{#2#1#3}
\crefformat{table}{#2#1#3}
\Crefformat{table}{#2#1#3}
\crefformat{algorithm}{#2#1#3}
\Crefformat{algorithm}{#2#1#3}

% cleveref 会在自己的分组中排出“定义 3.1.5”一类完整链接，类别名与编号之间
% 使用 \nobreakspace。分组结束后，xeCJK 无法再从末端识别出西文数字，
% 因而不会在紧随其后的汉字前插入 \CJKecglue。这里只在下一个源码记号确为
% CJK 正文字时补回同一种间距；空格、控制序列以及中英文标点均不处理，
% 从而不在“定义3.1.5，”或“定义3.1.5.”前制造额外空白。
\NewCommandCopy{\GeometryOriginalCref}{\cref}
\NewCommandCopy{\GeometryOriginalCrefCapital}{\Cref}
\ExplSyntaxOn
\cs_new_protected:Npn \geometry_cref_maybe_cjk_glue:
  { \peek_after:Nw \geometry_cref_maybe_cjk_glue_aux: }
\cs_new_protected:Npn \geometry_cref_maybe_cjk_glue_aux:
  {
    \token_if_space:NF \l_peek_token
      {
        \token_if_cs:NF \l_peek_token
          {
            \xeCJK_if_same_class:NNT \l_peek_token 中
              { \CJKecglue }
          }
      }
  }
\RenewDocumentCommand{\cref}{s m}
  {
    \IfBooleanTF{#1}
      { \GeometryOriginalCref*{#2} }
      { \GeometryOriginalCref{#2} }
    \geometry_cref_maybe_cjk_glue:
  }
\RenewDocumentCommand{\Cref}{s m}
  {
    \IfBooleanTF{#1}
      { \GeometryOriginalCrefCapital*{#2} }
      { \GeometryOriginalCrefCapital{#2} }
    \geometry_cref_maybe_cjk_glue:
  }
% 附录标签沿用字母编号；名称与编号整体进入同一链接，并与 \cref
% 采用相同的链接后中西文间距规则。
\NewDocumentCommand{\appref}{m}
  {
    \hyperref[#1]{附录\CJKecglue\ref*{#1}}
    \geometry_cref_maybe_cjk_glue:
  }
\ExplSyntaxOff

\newcommand{\MakeBookTitle}[1][封面]{%
  % 书名页在 PDF 侧栏中作为顶级入口；卷级文件可传入更具辨识性的书签名。
  \bookmarksetup{startatroot}%
  \pdfbookmark[0]{#1}{GeometryCover}%
  \begin{titlepage}
    \centering
    \vspace*{2.2cm}
    {\fontsize{36}{46}\selectfont\sffamily 几何学\par}
    \vspace{0.7cm}
    {\Large\sffamily\color{GeometryGray}\BookVolumeName\par}
    \vspace{1.0cm}
    {\fontsize{24}{32}\selectfont\bfseries\color{GeometryBlue}\BookTitle\par}
    \vfill
    {\large \BookAuthor\ 编写\par}
    \vspace{0.5cm}
    {\normalsize\BookDraftDate\par}
    \vspace{0.7cm}
    {\large\sffamily\bfseries\color{GeometryGray}\BookStatus\par}
    \vspace*{1.8cm}
  \end{titlepage}
}

% 目录也应有独立的顶级 PDF 书签。使用公共命令避免五卷分别遗漏。
\newcommand{\MakeBookContents}{%
  \cleardoublepage
  \bookmarksetup{startatroot}%
  \pdfbookmark[0]{目录}{GeometryContents}%
  \tableofcontents
}

% 图、表、算法目录不进入总目录，也不建立 PDF 书签；各自页面使用原有章级标题。
% 只另起一页而不补单双页。
\makeatletter
\newcommand{\GeometryListStart}{%
  \clearpage
}
% 各列表命令自带 \chapter*，局部消除其中的再次换页，保留原有标题与列表排版。
\newcommand{\GeometryListOnCurrentPage}[1]{%
  \begingroup
    \let\clearpage\relax
    \let\cleardoublepage\relax
    #1%
  \endgroup
}
\newcommand{\MakeBookListOfFigures}{%
  \GeometryListStart
  \GeometryListOnCurrentPage{\listoffigures}%
}
\newcommand{\MakeBookListOfTables}{%
  \begingroup
    \renewcommand{\listtablename}{表目录}%
    \GeometryListStart
    \GeometryListOnCurrentPage{\listoftables}%
  \endgroup
}
\newcommand{\MakeBookListOfAlgorithms}{%
  \begingroup
    \GeometryListStart
    % float.sty 在 listof 内重设编号栏；在标题回调中再加宽它。
    \let\GeometryOriginalFloatListHead\float@listhead
    \renewcommand{\float@listhead}[1]{%
      \renewcommand*\l@algorithm{\@dottedtocline{1}{1.5em}{5em}}%
      \GeometryOriginalFloatListHead{##1}%
    }%
    \GeometryListOnCurrentPage{\listofalgorithms}%
  \endgroup
}
\makeatother

% 无编号章（序言、目录、编导读、后记、参考文献、索引等）同样写入页眉标记，
% 使后续偶数页页眉显示该单元标题。置于 hyperref 之后，保证保留 hyperref 的章锚点。
\makeatletter
\let\Geometry@orig@schapter\@schapter
\def\@schapter#1{%
  \markboth{#1}{#1}%
  \Geometry@orig@schapter{#1}%
}
% 目录按层级着色：编为黑色、小节为深灰；章与节保持默认链接蓝。
\pretocmd{\l@part}{\hypersetup{linkcolor=black}}{}{%
  \PackageWarning{geometry}{无法为目录编条目着色}}
\pretocmd{\l@chapter}{\hypersetup{linkcolor=GeometryBlue}}{}{%
  \PackageWarning{geometry}{无法为目录章条目着色}}
\pretocmd{\l@section}{\hypersetup{linkcolor=GeometryBlue}}{}{%
  \PackageWarning{geometry}{无法为目录节条目着色}}
\pretocmd{\l@subsection}{\hypersetup{linkcolor=GeometryGray}}{}{%
  \PackageWarning{geometry}{无法为目录小节条目着色}}
% 编的目录条目与 PDF 书签须用纯文本“第X编”。ctex 的 \CTEX@part@tocline
% 取 \CTEXthepart，而该宏带上编题页的颜色与字体设置，会把格式代码写进目录，
% 并使书签出现 “=GeometryGray第一编” 之类乱码。此覆盖须置于 \ctexset 之后，
% 否则会被 ctex 重新定义；\CTEXifname 是 ctex 的内部开关，此处沿用。
\def\CTEX@part@tocline#1#2{%
  \CTEXifname
    {第\chinese{part}编\hspace{1em}}%
    {}%
  #2}
\makeatother

\hypersetup{pdftitle={几何学·\BookVolumeName：\BookTitle},pdfauthor={\BookAuthor},pdfsubject={几何学书系框架预览}}

\begin{document}
\frontmatter
% !TeX root = ../Book2.tex
\MakeBookTitle[\BookTitle]
\clearpage
\thispagestyle{empty}
\null
\clearpage
% !TeX root = ../../Book2.tex
\chapter*{序言}
\addcontentsline{toc}{chapter}{序言}
本卷为《几何学》\BookVolumeName《\BookTitle》的框架预览。目录依据《几何学计划纲要》建立，选读标记与原纲要一致。

待撰写

\MakeBookContents
% 有实际图、表或算法后再启用相应目录。
% \MakeBookListOfFigures
% \MakeBookListOfTables
% \MakeBookListOfAlgorithms


\mainmatter
% !TeX root = ../Book2.tex
% 纲要标题：# 第一编　空间分解、同调与对偶
% 纲要位置：几何学计划纲要.md，第 996 行。
% 正文状态：未撰写；本文件目前是框架。

% 纲要内容开始
% 
% 纲要内容结束

\BookPart{空间分解、同调与对偶}{g2:part:01}
\BookPartGuide
待撰写

% !TeX root = ../../Book2.tex
% 纲要标题：## 第1章　同伦、覆盖与胞腔空间
% 纲要位置：几何学计划纲要.md，第 998 行。
% 正文状态：未撰写；本文件目前是框架。

% 纲要内容开始
% 
% > 本章目标与证明范围：把独立入口和 CW 构造合成一章。基本群、覆盖在本卷有完整的最低证明；较深任务是余纤维化、相对胞腔构造与胞腔逼近，而不是重复第一卷的全部曲面分类。
% 
% 纲要内容结束

\chapter{同伦、覆盖与胞腔空间}
\label{g2:ch:01}
\BookChapterNumber{g2:ch:01}

\begin{chapterabstract}
待撰写
\end{chapterabstract}

\begin{learninggoals}
待撰写
\end{learninggoals}

待撰写

% !TeX root = ../../Book2.tex
% 纲要标题：### 1.1 同伦与基本群的工作语言
% 纲要位置：几何学计划纲要.md，第 1052 行。
% 正文状态：初稿；尚未正式审核。

% 纲要内容开始
% 
% * 映射同伦、基点变化和形变收缩
% * van Kampen 及图的群展示
% * 覆盖的提升与分类
% * 基本群作用和单值
% 
% 纲要内容结束

\section{同伦与基本群的工作语言}
\label{g2:sec:0101}

把一个圆沿圆周走一遍，与在圆内走一条闭合折线，看起来都只是回到出发点。区别在于：在圆盘中，闭合道路可以连续缩到一点；若只能沿圆周移动，绕行的次数便无法消去。基本群记录这种在指定空间中变形的障碍。计算时，一方面把空间拆成容易处理的部分，另一方面把道路提升到较简单的覆盖空间；这两种方法分别导向 van Kampen 定理和覆盖分类。

\ProseParagraphBreak

本节从连续映射和群论开始建立所需结论。拓扑空间不默认道路连通，覆盖也允许无限叶数；出现局部道路连通或半局部单连通条件时，将在相应结果中写明。道路拼接始终按实际行进次序书写：\(\alpha*\beta\) 先走 \(\alpha\)，后走 \(\beta\)；函数复合 \(g\circ f\) 则仍先作 \(f\)，后作 \(g\)。这一区别决定了纤维上的单值是右作用，而覆盖变换使用左作用。

\subsection{映射同伦、基点变化和形变收缩}
\label{g2:subsec:0101:01}
若两张地图把一个空间放进另一个空间，逐点移动像的位置并不自动给出连续变形：空间变量与变形时间必须一起连续。以后粘合空间时，还常常要求某一部分全程保持不动，因此先把相对条件纳入定义。

\begin{definition}[映射同伦与相对同伦]\label{g2:def:homotopy-relative}
	设 \(X,Y\) 为拓扑空间，\(A\subset X\)，\(f,g:X\rightarrow Y\) 为连续映射，记 \(I=[0,1]\)。若存在连续映射
	\[
		H:X\times I\longrightarrow Y,\qquad H(x,0)=f(x),\quad H(x,1)=g(x),
	\]
	则称 \(f,g\)\term[tonglun]{同伦}[homotopic]，记 \(f\simeq g\)。如果还满足 \(H(a,s)=f(a)=g(a)\) 对一切 \(a\in A,s\in I\) 成立，则称它们\term[xiangduitonglun]{相对于 \(A\) 同伦}[homotopic relative to \(A\)]，记 \(f\simeq g\;(\operatorname{rel}A)\)。带基点映射 \(f,g:(X,x_0)\rightarrow(Y,y_0)\) 的\term[dingjidiantonglun]{定基点同伦}[based homotopy]，指相对于 \(\{x_0\}\) 的同伦。
\end{definition}

同伦是映射集合上的等价关系：恒定变形给出自反性，倒转时间给出对称性，把两个变形各放在时间的一半便给出传递性。后一构造在接缝处连续，因为两式在闭集 \(X\times[0,1/2]\) 和 \(X\times[1/2,1]\) 上的取值相同。下文所用的有限闭集粘合均依据这一事实：有限个闭子集覆盖定义域时，在各闭子集上连续且在交集上一致的映射能够粘成连续映射。

\begin{definition}[道路与基本群]\label{g2:def:fundamental-group}
	设 \(X\) 为拓扑空间，\(x,y,z\in X\)，\(I=[0,1]\)。从 \(x\) 到 \(y\) 的\term[daolu]{道路}[path]是满足 \(\alpha(0)=x,\alpha(1)=y\) 的连续映射 \(\alpha:I\rightarrow X\)。道路的定端点同伦是相对于 \(\partial I=\{0,1\}\) 的同伦。同一起终点的道路类记为 \([\alpha]\)。若 \(\alpha\) 从 \(x\) 到 \(y\)，\(\beta\) 从 \(y\) 到 \(z\)，定义
	\[
		(\alpha*\beta)(t)=
		\begin{cases}
			\alpha(2t),&0\leq t\leq1/2,\\
			\beta(2t-1),&1/2\leq t\leq1,
		\end{cases}
		\qquad \bar\alpha(t)=\alpha(1-t).
	\]
	点 \(x\) 处的常值道路记为 \(e_x\)。起终点均为 \(x\) 的道路称为\term[huanlu]{环路}[loop]。所有这类环路的定端点同伦类，在乘法 \([\alpha][\beta]=[\alpha*\beta]\) 下组成的群称为 \(X\) 在 \(x\) 处的\term[jibenqun]{基本群}[fundamental group]，记为 \(\pi_1(X,x)\)。\symidx{\(\pi_1(X,x)\)：基本群}环路与 \(e_x\) 定端点同伦时，称它\term[linglun]{零伦}[null-homotopic]。
\end{definition}

这里需要验证道路类确实组成群。两个定端点同伦按道路参数各占一半，便给出拼接的良定义性。若 \(\rho:I\rightarrow I\) 连续且固定两端，则
\[
	(t,s)\longmapsto\alpha\bigl((1-s)t+s\rho(t)\bigr)
\]
说明 \(\alpha\circ\rho\) 与 \(\alpha\) 定端点同伦。三条道路的两种括号仅把行进时间分配成 \((1/4,1/4,1/2)\) 或 \((1/2,1/4,1/4)\)，都是同一条三段道路的这种重参数化，因此得到结合律。停顿的时间可以同样消去，故 \([e_x]\) 是单位。最后，公式
\[
	H(t,s)=\alpha\bigl((1-s)\min\{2t,2-2t\}\bigr)
\]
把 \(\alpha*\bar\alpha\) 收缩到 \(e_x\)，并全程固定两端；用于 \(\bar\alpha\) 又给出另一侧逆。因此 \([\bar\alpha]=[\alpha]^{-1}\)。同一计算还说明：从 \(x\) 到 \(y\) 的道路 \(\alpha,\beta\) 定端点同伦，当且仅当 \(\alpha*\bar\beta\) 零伦。

\ProseParagraphBreak

道路可以先后复合，但起终点必须相接。这些道路类合起来构成\term[jibenqunpai]{基本群胚}[fundamental groupoid]：对象是空间的点，从 \(x\) 到 \(y\) 的态射是相应道路类，每个态射都有逆。基本群正是只保留某个对象及其自态射所得的群。群胚的说法能够同时处理不同基点；本节主要使用下面明确的基点变换公式。

\begin{proposition}[诱导同态与基点变化]\label{g2:prop:basepoint-change}
	设 \(X,Y\) 为拓扑空间，\(f:X\rightarrow Y\) 连续，\(x\in X\)。则
	\[
		f_*:\pi_1(X,x)\longrightarrow\pi_1(Y,f(x)),\qquad [\alpha]\longmapsto[f\circ\alpha]
	\]
	是群同态，保持恒等映射与函数复合。若 \(y\in X\)，且 \(\gamma\) 是 \(X\) 中从 \(x\) 到 \(y\) 的道路，则
	\[
		c_\gamma:\pi_1(X,x)\longrightarrow\pi_1(X,y),\qquad
		c_\gamma[\alpha]=[\bar\gamma*\alpha*\gamma]
	\]
	是同构，其逆为 \(c_{\bar\gamma}\)；对 \(z\in X\) 及从 \(y\) 到 \(z\) 的道路 \(\delta\)，有 \(c_{\gamma*\delta}=c_\delta\circ c_\gamma\)。若 \(g:X\rightarrow Y\) 连续，\(H:X\times[0,1]\rightarrow Y\) 是从 \(f\) 到 \(g\) 的同伦，并令 \(\lambda(s)=H(x,s)\)，则
	\[
		g_*=c_\lambda\circ f_*:\pi_1(X,x)\longrightarrow\pi_1(Y,g(x)).
	\]
	特别地，定基点同伦的映射诱导相同的同态。
\end{proposition}
\begin{proof}
	把 \(f\) 复合在道路或同伦的后面，便得到相应道路或同伦；它与拼接的两段公式相容，所以 \(f_*\) 良定义且保持乘法。复合及恒等式逐点成立。对于 \(c_\gamma\)，两次使用相邻的 \(\gamma*\bar\gamma\simeq e_x\) 可得
	\[
		[\bar\gamma*\alpha*\gamma]\,[\bar\gamma*\beta*\gamma]
		=[\bar\gamma*(\alpha*\beta)*\gamma].
	\]
	其逆与复合公式由同样的消去得到。

	最后，把 \(H\) 限制到 \(\alpha(t)\) 上，得到方形 \((t,s)\mapsto H(\alpha(t),s)\)。其下边是 \(f\circ\alpha\)，上边是 \(g\circ\alpha\)，左右两边均为 \(\lambda\)。沿方形两条边界路线从左下角到右上角，得到 \((f\circ\alpha)*\lambda\) 与 \(\lambda*(g\circ\alpha)\)；方形内的直线变形把这两条路线定端点同伦。因此
	\[
		[g\circ\alpha]=[\bar\lambda*(f\circ\alpha)*\lambda]=c_\lambda f_*[\alpha],
	\]
	这也说明为什么普通同伦不能不加说明地给出 \(f_*=g_*\)：两者的目标基点可能不同。
\end{proof}

若换用另一条从 \(x\) 到 \(y\) 的道路 \(\gamma'\)，则 \(c_{\gamma'}c_\gamma^{-1}\) 是 \(\pi_1(X,y)\) 的一个内自同构。因而道路连通空间在不同基点的基本群互相同构，但通常没有与道路选择无关的指定同构。若该群为交换群，内自同构均为恒等，这一歧义才消失。

\begin{definition}[同伦等价与形变收缩]\label{g2:def:homotopy-equivalence-retract}
	设 \(X,Y\) 为拓扑空间，记 \(I=[0,1]\)。若有连续映射 \(f:X\rightarrow Y\)、\(h:Y\rightarrow X\)，满足 \(hf\simeq\operatorname{id}_X\) 和 \(fh\simeq\operatorname{id}_Y\)，则称 \(f\) 为\term[tonglundengjia]{同伦等价}[homotopy equivalence]，称 \(X,Y\) 具有相同同伦型。非空空间与一点同伦等价时，称它\term[kesuo]{可缩}[contractible]。

	设 \(A\subset X\) 带子空间拓扑，\(i:A\hookrightarrow X\) 为包含。满足 \(r\circ i=\operatorname{id}_A\) 的连续映射 \(r:X\rightarrow A\) 称为\term[suohui]{缩回}[retraction]。若还存在同伦
	\[
		H:X\times I\longrightarrow X,\qquad H(x,0)=x,\quad H(x,1)=ir(x),\quad H(a,s)=a
	\]
	对所有 \(x\in X,a\in A,s\in I\) 成立，则称 \(A\) 为 \(X\) 的\term[xingbiansuosuo]{形变收缩核}[deformation retract]，称 \(H\) 为到 \(A\) 的\term[xingbianshousuo]{形变收缩}[deformation retraction]。本书使用的形变收缩固定 \(A\) 的每一点，文献也称之为强形变收缩。
\end{definition}

\begin{proposition}[基本群的同伦不变性]\label{g2:prop:pi-one-homotopy-invariance}
	设 \(f:X\rightarrow Y\) 是拓扑空间间的同伦等价，\(x\in X\)。则 \(f_*:\pi_1(X,x)\rightarrow\pi_1(Y,f(x))\) 是同构。特别地，若 \(A\subset X\) 是形变收缩核，\(a\in A\)，则包含和缩回诱导互逆同构 \(\pi_1(A,a)\cong\pi_1(X,a)\)。可缩空间的所有基本群均平凡。
\end{proposition}
\begin{proof}
	取同伦逆 \(h\)。由\cref{g2:prop:basepoint-change}，\(h_*f_*=(hf)_*\) 与恒等同态只差一个基点变化同构，故为同构，因而 \(f_*\) 单射。同一论证用于 \(h\)，说明以 \(f(x)\) 为基点时 \(h_*\) 也是单射。给定 \(b\in\pi_1(Y,f(x))\)，因 \(h_*f_*\) 满射，可取 \(a\) 使 \(h_*f_*(a)=h_*(b)\)；\(h_*\) 的单射性便给出 \(f_*(a)=b\)。形变收缩中的同伦固定 \(a\)，所以两复合直接诱导恒等同态。一点空间的基本群平凡，最后一项随之成立。
\end{proof}

\begin{example}[径向收缩]\label{g2:exm:radial-retraction}
	设 \(n\geq1\)，\(\mathbb R^n\) 配标准范数，\(S^{n-1}=\{x:\lVert x\rVert=1\}\)。\(\mathbb R^n\setminus\{0\}\) 到 \(S^{n-1}\) 的形变收缩为
	\[
		H(x,s)=\left(1-s+\frac{s}{\lVert x\rVert}\right)x.
	\]
	系数始终为正，所以变形不经过原点；在单位球面上它恒等。同样，若 \(0<r<1<R\)，开圆环 \(\{z\in\mathbb C:r<|z|<R\}\) 形变收缩到单位圆。它与填满的圆盘不同：圆盘可沿 \((z,s)\mapsto(1-s)z\) 缩到中心，而这一变形会经过穿孔平面所删去的点。
\end{example}

\begin{definition}[单连通空间]\label{g2:def:simply-connected}
	设 \(X\) 为非空拓扑空间。若 \(X\) 道路连通，且在某个 \(x\in X\) 处有 \(\pi_1(X,x)=\{1\}\)，则称 \(X\)\term[danliantong]{单连通}[simply connected]。
\end{definition}

由基点变化，单连通定义中的平凡性在任一基点处都成立。这一条件等价于任意两条同起终点的道路定端点同伦。可缩蕴含单连通，但反向结论并不属于基本群能决定的内容；后面的高维同伦与同调将记录基本群看不到的障碍。

\subsection{van Kampen 及图的群展示}
\label{g2:subsec:0101:02}
要把空间分解变成群的计算，必须分别知道局部产生哪些环路、交叠部分给出哪些关系。先计算最基本的局部环路：圆上的绕行。这里直接用实数角度证明，不预先调用覆盖分类。

\begin{proposition}[圆的基本群]\label{g2:prop:circle-fundamental-group}
	设 \(S^1=\{z\in\mathbb C:|z|=1\}\)，以 \(1\) 为基点，记 \(I=[0,1]\)。对每条基点环路 \(\alpha:I\rightarrow S^1\)，存在唯一连续函数 \(\theta:I\rightarrow\mathbb R\)，满足 \(\theta(0)=0\) 和 \(\alpha(t)=\exp(2\pi\mathrm{i}\theta(t))\)。其终值 \(\theta(1)\) 为整数，且
	\[
		\pi_1(S^1,1)\longrightarrow\mathbb Z,\qquad[\alpha]\longmapsto\theta(1)
	\]
	是同构，正向环路 \(t\mapsto\exp(2\pi\mathrm{i}t)\) 对应 \(1\)。整数 \(\theta(1)\) 称为环路的\term[raoshu]{绕数}[winding number]。
\end{proposition}
\begin{proof}
	每个圆点都有一段真开圆弧，其上存在连续实角。用这些圆弧的逆像覆盖紧区间，再把区间细分成有限小段，使每段的像落在一条圆弧内。在第一段选择起值为零的角度，随后各段加上唯一整数，使接点处角度相同，便得到 \(\theta\)。两种选择的差为连续整数值函数，故在连通区间上恒定；起值相同使差恒为零。由于 \(\alpha(1)=1\)，终值是整数。

	为证明同伦不变性，设 \(H:I\times I\rightarrow S^1\) 是基点环路同伦。若两个时间 \(s,s'\) 足够接近，则由紧方形上的一致连续性，\(H(t,s')/H(t,s)\) 对全部 \(t\) 都避开 \(-1\)。它有取值于 \((-1/2,1/2)\) 的连续实角 \(\eta(t)\)，且两端均为零。于是两条环路的实角提升相差 \(\eta\)，终值相同。把变形时间作有限细分，即得整个同伦中绕数不变。

	若两个实角提升 \(\theta_0,\theta_1\) 具有同一终值 \(k\)，则
	\[
		(t,s)\longmapsto\exp\bigl(2\pi\mathrm{i}((1-s)\theta_0(t)+s\theta_1(t))\bigr)
	\]
	是固定基点的同伦。因此绕数也完全决定道路类。拼接时，第二段的实角从第一段终值接着增加，故绕数相加；环路 \(t\mapsto\exp(2\pi\mathrm{i}kt)\) 实现每个整数 \(k\)。
\end{proof}

两个群 \(G_1,G_2\) 的\term[ziyouji]{自由积}[free product] \(G_1*G_2\)，可以用交替的非单位字母组成的有限约化字构造；相邻同属一个群的字母相乘，乘积为单位时删去。它的泛性质是：给出到群 \(K\) 的同态 \(G_1\rightarrow K,G_2\rightarrow K\)，恰好给出唯一同态 \(G_1*G_2\rightarrow K\)。相交部分的环路在两个群中各出现一次，下一定理说明应当把这两份识别。

\begin{theorem}[van Kampen 定理]\label{g2:thm:van-kampen}
	设拓扑空间 \(X=U\cup V\)，其中 \(U,V\) 是道路连通开集，\(W=U\cap V\) 非空且道路连通，固定 \(x_0\in W\)。记包含诱导的同态为
	\[
		i_U:\pi_1(W,x_0)\longrightarrow\pi_1(U,x_0),\qquad
		i_V:\pi_1(W,x_0)\longrightarrow\pi_1(V,x_0).
	\]
	则包含给出的满同态
	\[
		\Phi:\pi_1(U,x_0)*\pi_1(V,x_0)\longrightarrow\pi_1(X,x_0)
	\]
	的核为所有 \(i_U(w)i_V(w)^{-1}\) 所生成的正规子群 \(N\)，其中 \(w\in\pi_1(W,x_0)\)。因而
	\[
		\pi_1(X,x_0)\cong
		\frac{\pi_1(U,x_0)*\pi_1(V,x_0)}{\langle\!\langle i_U(w)i_V(w)^{-1}:w\in\pi_1(W,x_0)\rangle\!\rangle}.
	\]
	等价地，这个群是两同态 \(i_U,i_V\) 的推出。
\end{theorem}

van Kampen 定理用交集中的环路给出两组生成元之间的全部关系。\footnote{Egbertus Rudolf van Kampen（1908--1942），荷兰数学家，研究拓扑与拓扑群。}定理及图的计算可参看 \cite[第1.2、1.A节]{Hatcher2002}。后面的证明把生成性与关系的完备性分开：细分环路得到生成元，细分环路之间的同伦则确定全部关系。推出的含义由下图表示；对任意群 \(K\)，若同态 \(\varphi_U:\pi_1(U,x_0)\rightarrow K\)、\(\varphi_V:\pi_1(V,x_0)\rightarrow K\) 在 \(\pi_1(W,x_0)\) 上一致，就存在唯一经过 \(\pi_1(X,x_0)\) 的同态 \(\varphi\)。
\[
	\begin{tikzcd}[column sep=large,row sep=large]
		\pi_1(W,x_0) \arrow[r,"i_U"] \arrow[d,"i_V"'] & \pi_1(U,x_0) \arrow[d] \arrow[ddr,bend left=20,"\varphi_U"] & \\
		\pi_1(V,x_0) \arrow[r] \arrow[drr,bend right=20,"\varphi_V"'] & \pi_1(X,x_0) \arrow[dr,dashed,"\exists!\,\varphi"] & \\
		& & K
	\end{tikzcd}
\]
\begin{proof}
	先证明 \(\Phi\) 满射。对基点环路 \(\alpha\)，把参数区间细分，使各小段的像落在 \(U\) 或 \(V\) 内。这样的细分存在：\(\alpha^{-1}(U),\alpha^{-1}(V)\) 是紧区间的开覆盖，可以取比其 Lebesgue 数更小的网格。合并相邻的同色段以后，每个换色接点 \(x_j\) 都在 \(W\)。在 \(W\) 中选从 \(x_0\) 到 \(x_j\) 的道路 \(h_j\)，并在首尾取常值道路。若第 \(j\) 段为 \(\alpha_j\)，则
	\[
		\ell_j=h_{j-1}*\alpha_j*\bar h_j
	\]
	是完全位于 \(U\) 或 \(V\) 的基点环路。相邻的 \(\bar h_j*h_j\) 可消去，所以 \([\alpha]=[\ell_1]\cdots[\ell_m]\)，得到满射。

	所有规定的关系显然在 \(X\) 中成立。记商群为 \(Q=(\pi_1(U,x_0)*\pi_1(V,x_0))/N\)，还需证明两个在 \(X\) 中同伦的分解在 \(Q\) 中相等。设它们分别由有限个 \(U\) 或 \(V\) 中的环路拼成，\(F:I^2\rightarrow X\) 是拼接环路之间固定基点的同伦。将方形细分成有限矩形网格，使每个闭小矩形都映入 \(U\) 或 \(V\)，并让下、上边的网格包含原环路分段的接点。存在这种网格仍由 \(F^{-1}(U),F^{-1}(V)\) 的 Lebesgue 数保证；先取细网格，再加入有限条所需分界线即可。

	给每个小矩形指定一个容纳其像的开集，称作它的颜色。对每个顶点 \(v\)，考察所有邻接矩形的颜色；在方形上、下边还把原环路段的颜色列入。若只有颜色 \(U\)，就在 \(U\) 中选道路 \(h_v:x_0\rightarrow F(v)\)；只有 \(V\) 时在 \(V\) 中选；两色都有时，\(F(v)\in W\)，就在 \(W\) 中选。位于左右边及原环路接点上的顶点满足 \(F(v)=x_0\)，一律取 \(h_v=e_{x_0}\)。道路连通条件正是在这里保证这些选择存在。

	对从 \(v\) 到 \(w\) 的有向网格边 \(e\)，用
	\[
		L(e)=[h_v*(F|_e)*\bar h_w]\in Q
	\]
	记录它。这里将边线性参数化为区间；可以使用任何邻接矩形的颜色来理解括号中的环路。若边的两侧颜色不同，则边像及两端所选道路都在 \(W\) 中，定义 \(N\) 的关系恰好保证两种理解在 \(Q\) 中相同。反向边给出 \(L(\bar e)=L(e)^{-1}\)。

	每个小矩形的有向边界在其指定开集中零伦，因为 \(F\) 本身给出了填充。其四边附带的 \(h_v\) 依次消去，故这四个 \(L(e)\) 的乘积为单位。现在从方形下边的左到右路线开始，逐个跨过小矩形，把路线推到上边；每次用一个小矩形的另两边替换原两边，都因刚才的边界关系而不改变 \(Q\) 中的乘积。路线中可能出现的相邻反向边也相消。左右外边的像恒为基点，所贡献的元素为单位，因此下边与上边给出相同元素。

	沿下边的每一原始环路段，所有 \(h_v\) 都位于该段的指定开集，且段首尾的 \(h_v\) 为常值。于是这一段各网格边的乘积正是原环路在对应因子群中的类；上边亦然。因此两个原分解在 \(Q\) 中相同，\(Q\rightarrow\pi_1(X,x_0)\) 单射。商群构造同时给出所述推出泛性质。
\end{proof}

交集的道路连通不能省去。例如，用两条真开圆弧覆盖 \(S^1\)，每条圆弧的基本群都平凡，但交集有两个道路分支。若仍套用上式便会错得 \(\pi_1(S^1)=1\)。此时需要保留多个基点的群胚版本，或者改用满足本定理条件的分解。

\begin{definition}[图与生成树]\label{g2:def:graph-tree}
	由离散顶点集与若干闭区间的互不相交并，在每个区间两端指定顶点并粘合，得到的空间称为\term[tu]{图}[graph]；区间内部称为开边，允许自环及重边。采用弱拓扑：子集闭当且仅当其在顶点集的原像闭，且在每条闭区间的原像闭。只有有限个顶点和边时称为有限图。由若干顶点及完整闭边组成的子空间称为子图。连通且没有有限简单回路的子图称为\term[shu]{树}[tree]；这里简单回路允许单条自环，除起终点外不重复顶点。包含连通图全部顶点的树称为\term[shengchengshu]{生成树}[spanning tree]。
\end{definition}

图的两个顶点能否由有限边道路连接，给出了顶点集的等价关系。把每一等价类连同以其中顶点为端点的所有开边放在一起，在每条闭边上的原像不是整个区间就是空集，因此这一部分及其补集在弱拓扑中都开。连通图于是只有一个这样的分支，任意两顶点之间都有有限边道路。一棵有限树中，删去道路的往返部分后，两顶点间的不重复边道路唯一：若有两条，就能从首次分离至再次相遇处找出回路。不断删去不含指定根顶点的末端边，说明有限树可缩。也可给每条边长度 \(1\)，把每一点沿到根的唯一弧移动其距离的 \(s\) 倍，得到固定根的连续收缩。

\begin{theorem}[有限图的自由群展示]\label{g2:thm:graph-free-group}
	设 \(\Gamma\) 是非空有限连通图，顶点数为 \(v\)，边数为 \(e\)，选根顶点 \(x_0\) 和生成树 \(T\)。给每条不属于 \(T\) 的边 \(a\) 选定方向，记 \(p_u,p_w\) 分别为树中从 \(x_0\) 到 \(a\) 的起点 \(u\)、终点 \(w\) 的唯一不重复边道路。则
	\[
		\pi_1(\Gamma,x_0)
		=\left\langle [p_u*a*\bar p_w]\ (a\notin T)\ \middle|\ \varnothing\right\rangle
		\cong F_{e-v+1},
	\]
	其中 \(F_r\) 为 \(r\) 个生成元的自由群，\(F_0=\{1\}\)。此外，\(\Gamma/T\) 是 \(e-v+1\) 个圆的一点并，商映射 \(\Gamma\rightarrow\Gamma/T\) 是同伦等价。
\end{theorem}
\begin{proof}
	从一个顶点开始，只要还有未加入的顶点，就沿一条连接现有顶点与新顶点的边扩张。有限连通性使这一过程到达全部顶点，且不会形成回路，所以得到生成树；每加一顶点恰加一边，树有 \(v-1\) 条边。

	记 \(q:\Gamma\rightarrow\Gamma/T\)。树被压成一点后，每条其余边的两端都成为该点，所以商是 \(r=e-v+1\) 个圆的一点并 \(R_r\)。定义 \(s:R_r\rightarrow\Gamma\)，使第 \(a\) 个圆依次沿 \(p_u,a,\bar p_w\) 行进，各占三分之一时间；公共点映到 \(x_0\)。各段在端点一致，故映射连续。\(qs\) 在每个圆上先停顿、再绕行、再停顿，因而重参数化便给出 \(qs\simeq\operatorname{id}_{R_r}\)。

	还要说明压缩树没有改变 \(\Gamma\) 的同伦型。取上述沿唯一弧到 \(x_0\) 的收缩 \(h_t:T\rightarrow T\)，并将树内道路参数化为 \(p_v(\tau)=h_{1-\tau}(v)\)。对非树边 \(a:[0,1]\rightarrow\Gamma\)，在时刻 \(t\) 把它改走成三段：从 \(h_t(u)\) 沿树回到 \(u\)，沿 \(a\) 到 \(w\)，再沿树到 \(h_t(w)\)。具体定义
	\[
		H(a(x),t)=
		\begin{cases}
			h_{t-3x}(u),&0\leq x\leq t/3,\\
			a\bigl((x-t/3)/(1-2t/3)\bigr),&t/3\leq x\leq1-t/3,\\
			h_{3x-3+t}(w),&1-t/3\leq x\leq1.
		\end{cases}
	\]
	中段分母至少为 \(1/3\)，三式在各自的闭区域上连续且在接缝处相同；这也涵盖 \(t=0\) 时两条树内尾段消失的情形。在树上定义 \(H(v,t)=h_t(v)\)，与所有边端点的公式一致。有限图是有限个闭边与顶点的并，故这些公式粘成连续同伦；它在 \(t=0\) 为恒等，在 \(t=1\) 为 \(sq\)，并始终固定 \(x_0\)。

	最后计算 \(R_r\)。当 \(r=0\) 时它是一点；当 \(r=1\) 时由\cref{g2:prop:circle-fundamental-group}得到 \(\mathbb Z\)。对于较多圆，取第一个圆连同其他圆在公共点附近的小开弧为 \(U\)，取其他圆连同第一个圆的小开弧为 \(V\)。\(U\) 收缩到第一个圆，\(V\) 收缩到其余圆的一点并，交集为可缩的有限星状邻域。\cref{g2:thm:van-kampen}给出 \(\pi_1(R_r)\cong\mathbb Z*\pi_1(R_{r-1})\)，归纳得到自由群。\(q\) 把题中每条环路送到对应圆的正向生成元，因此也得到了所写的具体展示。
\end{proof}

\begin{example}[三个顶点间的两种计数]\label{g2:exm:graph-computation}
	三角形图有 \(v=e=3\)，删去任意一边便得到生成树，故基本群为 \(F_1=\mathbb Z\)。若在其中两顶点间再加一条边，则 \(e=4,v=3\)，基本群为 \(F_2\)，两条非树边分别给出自由生成元。另一方面，立方体的棱图有 \(e=12,v=8\)，所以基本群为 \(F_5\)；这里仅取棱，不填入六个方形面。填面会加入关系，不能只按图的边数继续计算。
\end{example}

\begin{lemma}[图中紧子集的有限控制]\label{g2:lem:graph-compact-control}
	设 \(\Gamma\) 为采用\cref{g2:def:graph-tree}中弱拓扑的图，\(K\subset\Gamma\) 是紧子集。则 \(K\) 包含在某个有限子图中。特别地，环路与两个环路之间的同伦，其像均包含在有限子图中。
\end{lemma}
\begin{proof}
	若 \(K\) 与无限多条开边相交，从其中可数条不同开边各选一点 \(x_j\in K\)。集合 \(D=\{x_j\}\) 的任意子集，在每条闭边的原像至多为一个内点，在顶点集的原像为空，因而依弱拓扑均为闭集。这说明 \(D\) 是 \(\Gamma\) 的闭离散子集。于是它也是紧空间 \(K\) 的闭子集，故紧；但无限离散空间的单点开覆盖没有有限子覆盖，矛盾。若 \(K\) 含无限多顶点，取可数多个顶点也会得到相同矛盾，因为它们的任意子集在每条闭边的原像至多为两端点。因此 \(K\) 只遇到有限条开边和有限个顶点，把这些边的端点一并加入便得到所需有限子图。\(I\) 和 \(I^2\) 紧，连续像紧，最后一项随之成立。
\end{proof}

\begin{proposition}[一般连通图的自由基]\label{g2:prop:general-graph-free-group}
	设 \(\Gamma\) 是非空连通图，采用由离散顶点集和各闭边粘合所得的弱拓扑，不要求有限或局部有限。则 \(\Gamma\) 有生成树 \(T\)。固定顶点 \(x_0\)，给每条非树边 \(a:u\rightarrow w\) 选定方向，并记 \(p_v\) 为 \(T\) 中从 \(x_0\) 到顶点 \(v\) 的唯一不重复边道路，则 \(\pi_1(\Gamma,x_0)\) 是以所有 \([p_u*a*\bar p_w]\) 为自由基的自由群。
\end{proposition}
\begin{proof}
	在包含全部顶点的无回路子图中，按包含关系应用 Zorn 引理。全序链的并仍无回路，因为每个回路只有有限条边，若它在并中出现，就已在链的某一项中出现。因此存在极大无回路子图 \(T\)。若它不连通，在 \(\Gamma\) 中取一条连接两个分支的有限边道路，其中有一条边的两端分属不同分支。把这条边加入 \(T\) 不产生回路，与极大性矛盾，故 \(T\) 是生成树。

	给定有限子图 \(K\)，把根 \(x_0\) 及从根到 \(K\) 各顶点的有限条树内道路一并加入，得到有限连通子图 \(K'\)。这些树内道路连同根组成 \(T'\)，且 \(K'\cap T=T'\)：若 \(K\) 中一条边属于 \(T\)，两端至根道路的并已包含这条唯一的树内连接边。因此 \(T'\) 是 \(K'\) 的生成树，有限图定理给出的自由基恰是 \(K'\) 中的非树边所对应的那部分指定环路。

	由\cref{g2:lem:graph-compact-control}，每条基点环路的像都落在某个有限子图中，将它扩成上述 \(K'\)，便知指定环路生成全部基本群。若这些生成元之间存在一个非空约化字关系，取其所代表环路的一次零伦；该同伦的紧像也在有限子图中，再加入字中各条边及其树内道路，仍可扩成上述 \(K'\)。这就在有限图 \(K'\) 的自由基之间产生非空约化字关系，与\cref{g2:thm:graph-free-group}矛盾。故没有额外关系。自由群元素仍由有限字表示，并不允许无限乘积。
\end{proof}

\subsection{覆盖的提升与分类}
\label{g2:subsec:0101:03}
在圆上，一条道路可以不断积累角度；若把角度记在实线上，绕行就成为实线上普通的移动。覆盖把这一构造推广到其他空间：局部看来，每一层都与底空间相同，只有沿较长道路继续移动以后，才可能落到另一层。

\begin{definition}[覆盖与提升]\label{g2:def:covering-space}
	设 \(E,X\) 为拓扑空间。若连续满射 \(p:E\rightarrow X\) 满足：每个 \(x\in X\) 都有开邻域 \(U\)，使
	\[
		p^{-1}(U)=\coprod_{\lambda\in\Lambda}U_\lambda,
	\]
	其中 \(U_\lambda\) 是两两不交的开集，且每个 \(p|_{U_\lambda}:U_\lambda\rightarrow U\) 都是同胚，则称 \(p\) 为\term[fugai]{覆盖}[covering map]，称 \(U\) 为均匀覆盖邻域，\(U_\lambda\) 为其上的叶。点 \(x\) 上的\term[xianwei]{纤维}[fiber]是 \(p^{-1}(x)\)；叶数指纤维基数，允许无限。\symidx{\(p^{-1}(x)\)：覆盖纤维}给定拓扑空间 \(Y\) 及连续映射 \(f:Y\rightarrow X\)，满足 \(p\widetilde f=f\) 的连续映射 \(\widetilde f:Y\rightarrow E\) 称为 \(f\) 的\term[tisheng]{提升}[lift]。

	对同一底空间上的两个覆盖 \(p_i:E_i\rightarrow X\)（\(i=1,2\)），满足 \(p_2\circ r=p_1\) 的同胚 \(r:E_1\rightarrow E_2\) 称为\term[fugaitonggou]{覆盖同构}[isomorphism of coverings]。若指定 \(x_0\in X\) 及 \(e_i\in p_i^{-1}(x_0)\)，保持基点还要求 \(r(e_1)=e_2\)。
\end{definition}

\begin{theorem}[道路提升与同伦提升]\label{g2:thm:covering-lifts}
	设 \(p:E\rightarrow X\) 为覆盖，\(I=[0,1]\)。每条道路 \(\alpha:I\rightarrow X\) 以及每个 \(e_0\in p^{-1}(\alpha(0))\)，唯一确定从 \(e_0\) 出发的提升道路 \(\widetilde\alpha\)。更一般地，对任意拓扑空间 \(Y\)、连续同伦 \(H:Y\times I\rightarrow X\) 和连续初始提升 \(\widetilde f_0:Y\rightarrow E\)，其中 \(p\widetilde f_0(y)=H(y,0)\)，存在唯一连续提升
	\[
		\widetilde H:Y\times I\longrightarrow E,\qquad
		p\widetilde H=H,\quad \widetilde H(y,0)=\widetilde f_0(y).
	\]
	尤其，定端点同伦的道路在指定同一起始提升点后，其提升终点相同，且提升道路也定端点同伦。
\end{theorem}
\begin{proof}
	对一条道路，用均匀覆盖邻域的逆像覆盖 \(I\)，取有限细分，使每段像位于一个这样的邻域。从 \(e_0\) 所在叶开始，逐段用 \(p\) 在相应叶上的逆映射提升；每段终点指定下一段的叶，闭区间粘合保证连续。若两条提升在某时刻相同，在一个均匀覆盖邻域中就只能同时走同一叶；若不相同，便位于不同叶，在附近仍不同。因此相等的时间集合及其补集都开，由 \(I\) 连通和起点相同，得到唯一性。

	对一般 \(Y\)，先对每个 \(y\) 沿竖直道路 \(t\mapsto H(y,t)\) 作上述提升，得到唯一候选 \(\widetilde H\)。关键是证明它随 \(y\) 连续。固定 \(y_0\)，选 \(0=t_0<\cdots<t_m=1\) 和均匀覆盖邻域 \(U_j\)，使 \(H(\{y_0\}\times[t_{j-1},t_j])\subset U_j\)。因为这是一条紧线段，对每个 \(j\) 可取 \(y_0\) 的开邻域 \(W_j\)，使 \(H(W_j\times[t_{j-1},t_j])\subset U_j\)。具体地，先在紧线段上取有限个乘积开邻域，再交其 \(Y\) 方向的邻域即可。

	取这些 \(W_j\) 的交，并缩小它，使 \(\widetilde f_0\) 落在 \(U_1\) 上指定的一叶中，于是第一条带上的提升等于该叶的逆映射与 \(H\) 的复合，因而连续。在 \(t_1\) 处再缩小邻域，使已得终值落在 \(U_2\) 上指定的一叶中，便可连续延长到第二条带。有限次重复后，得到某个 \(W\times I\) 上的连续提升；在接缝处由初值一致而粘合。它与逐点定义的候选由道路唯一性相同。由于 \(y_0\) 任意，\(\widetilde H\) 连续；唯一性仍逐条竖直道路成立。

	道路同伦的两端提升都是纤维中的连续道路。纤维在子空间拓扑下离散，区间连通，所以这两条端点道路恒定，得到最后的结论。
\end{proof}

提升一般映射时，需要在定义域的小邻域内连接点；构造万有覆盖时，还要控制底空间的小环路。两种条件分别如下。

\begin{definition}[局部道路连通与半局部单连通]\label{g2:def:local-covering-conditions}
	设 \(X\) 为拓扑空间。如果每一点都有由道路连通开集构成的邻域基，则称 \(X\)\term[jubudaoluliantong]{局部道路连通}[locally path-connected]。如果每个 \(x\in X\) 都有开邻域 \(U\)，使包含诱导的 \(\pi_1(U,x)\rightarrow\pi_1(X,x)\) 为零同态，则称 \(X\)\term[banjubudanliantong]{半局部单连通}[semilocally simply connected]。若空间还局部道路连通，可把此处的 \(U\) 选成道路连通开邻域；这时零同态条件对 \(U\) 中任意基点都成立。
\end{definition}

\begin{proposition}[覆盖中的基本群与提升判据]\label{g2:prop:covering-lift-criterion}
	设 \(p:E\rightarrow X\) 为覆盖，\(e_0\in E\)，\(x_0=p(e_0)\)。则 \(p_*:\pi_1(E,e_0)\rightarrow\pi_1(X,x_0)\) 单射，其像由那些从 \(e_0\) 提升后仍闭合的基点环路类组成。

	进一步，设 \(Y\) 道路连通且局部道路连通，\(y_0\in Y\)，\(f:(Y,y_0)\rightarrow(X,x_0)\) 连续。存在保持所给基点的提升 \(\widetilde f:(Y,y_0)\rightarrow(E,e_0)\)，当且仅当
	\[
		f_*\pi_1(Y,y_0)\subset p_*\pi_1(E,e_0).
	\]
	若存在，提升唯一。更一般地，对任意连通拓扑空间 \(Z\) 及连续映射 \(g:Z\rightarrow X\)，若其两个提升 \(\widetilde g_1,\widetilde g_2:Z\rightarrow E\) 在一点相同，就处处相同。
\end{proposition}
\begin{proof}
	若 \(E\) 中的基点环路投影后零伦，提升该零伦便把它变到常值环路，故 \(p_*\) 单射。闭合提升的投影显然属于像；反之，若一条底环路与某个闭合提升的投影定端点同伦，同伦提升保证它的提升也闭合，故像的描述准确。

	若 \(\widetilde f\) 存在，\(f_*=p_*\widetilde f_*\) 立即给出必要性。反过来，对每个 \(y\in Y\)，选从 \(y_0\) 到 \(y\) 的道路 \(\eta\)，定义 \(\widetilde f(y)\) 为 \(f\eta\) 从 \(e_0\) 提升后的终点。若换成 \(\eta'\)，则 \(f(\eta*\bar\eta')\) 的类属于 \(p_*\pi_1(E,e_0)\)，所以其提升闭合。逆向走提升道路的唯一性说明 \(f\eta,f\eta'\) 的提升终点相同，定义因而与选择无关。

	为证明连续性，在 \(f(y)\) 附近取均匀覆盖邻域 \(U\)。局部道路连通性使 \(y\) 有道路连通开邻域 \(W\subset f^{-1}(U)\)。从 \(y_0\) 先到 \(y\)，再在 \(W\) 内去到任意 \(z\)，说明 \(\widetilde f|_W=(p|_{\widetilde U})^{-1}f|_W\)，其中 \(\widetilde U\) 为包含 \(\widetilde f(y)\) 的叶，因此连续。

	最后，对同一映射的两个提升，在每个均匀覆盖邻域内，它们若落在同一叶便相等，若在不同叶便不同。所以相等点集及其补集均开。在连通定义域中，只要相等点集非空，它就等于全空间。
\end{proof}

提升判据的图形如下。虚线存在的障碍是环路能否闭合提升，而非逐点能否选到某个原像；局部道路连通性则保证由道路终点构造出来的选择连续。
\[
	\begin{tikzcd}[column sep=large,row sep=large]
		& (E,e_0) \arrow[d,"p"] \\
		(Y,y_0) \arrow[r,"f"'] \arrow[ur,dashed,"\widetilde f"] & (X,x_0)
	\end{tikzcd}
\]

要构造一个让所有环路都能展开的覆盖，需要保证足够小的底环路在整个空间中已经零伦。这比小邻域本身单连通弱。

\begin{theorem}[道路类万有覆盖]\label{g2:thm:universal-cover}
	设 \(X\) 非空、道路连通、局部道路连通且半局部单连通，固定 \(x_0\in X\)。令
	\[
		\widetilde X=\{[\alpha]:\alpha\text{ 是从 }x_0\text{ 出发的道路}\},\qquad
		p([\alpha])=\alpha(1),
	\]
	其中仅把同终点且定端点同伦的道路识别。对道路连通开集 \(U\subset X\)，要求对每个 \(u\in U\)，包含诱导的 \(\pi_1(U,u)\rightarrow\pi_1(X,u)\) 为零；若 \(\alpha(1)\in U\)，取
	\[
		B(\alpha,U)=\{[\alpha*\eta]:\eta\text{ 位于 }U,\ \eta(0)=\alpha(1)\}.
	\]
	这些集合构成 \(\widetilde X\) 的一个拓扑基，使 \(p:\widetilde X\rightarrow X\) 成为单连通覆盖，称为 \(X\) 的\term[wanyoufugai]{万有覆盖}[universal cover]。其基点取 \(\widetilde x_0=[e_{x_0}]\)。对任意带基点连通覆盖 \(q:(E,e_0)\rightarrow(X,x_0)\)，存在唯一保持基点的覆盖映射 \(\widetilde X\rightarrow E\) 与到底空间的投影相容。
\end{theorem}
\begin{proof}
	先说明小开集足够多。给定 \(x\) 的任意开邻域 \(O\)，先取半局部单连通条件中的邻域，再由局部道路连通性在它与 \(O\) 的交内取道路连通开邻域 \(U\)。其基本群到 \(X\) 的同态因子化经过零同态，故也为零；换基点只作共轭，不改变这一性质。因此定理中的 \(U\) 构成 \(X\) 的开基。

	固定这样的 \(U\)。\(p:B(\alpha,U)\rightarrow U\) 满射，因为 \(U\) 道路连通；它也单射，因为 \(U\) 中同起终点的两条道路，其差环路在 \(X\) 中零伦。若 \([\beta]\in B(\alpha,U)\)，则附加或消去一条 \(U\) 内的道路可得 \(B(\beta,U)=B(\alpha,U)\)。若 \([\beta]\) 同时属于两个基本集合，在 \(\beta(1)\) 处取包含于两开集交的小开集 \(W\)，就有 \(B(\beta,W)\) 包含在二者交内，故这些集合确为拓扑基。该基还说明 \(p|_{B(\alpha,U)}\) 与其逆均连续；固定 \(U\) 时各个 \(B(\alpha,U)\) 要么相同要么不交，且铺满 \(p^{-1}(U)\)。因此 \(p\) 是覆盖。\(X\) 道路连通也保证 \(p\) 满射。

	对从 \(x_0\) 出发的道路 \(\alpha\)，令 \(\alpha_t(s)=\alpha(ts)\)。则 \(t\mapsto[\alpha_t]\) 是 \(\alpha\) 从 \(\widetilde x_0\) 出发的提升：在任意小时间段内，\(\alpha\) 的尾段位于一个上述 \(U\)，因此道路类始终落在同一基本集合，连续性由该集合与 \(U\) 的同胚给出。这条提升连接 \(\widetilde x_0\) 与 \([\alpha]\)，故 \(\widetilde X\) 道路连通。

	若 \(\alpha\) 为 \(x_0\) 处环路，上述提升终点是 \([\alpha]\)；它闭合当且仅当 \([\alpha]=[e_{x_0}]\)。由\cref{g2:prop:covering-lift-criterion}，\(p_*\) 单射，且像只含单位，因此 \(\widetilde X\) 单连通。

	对覆盖 \(q:E\rightarrow X\)，底空间局部道路连通使 \(E\) 也局部道路连通；连通且局部道路连通的空间道路连通，因为其道路分支是开集。提升判据把 \(p\) 唯一提升为保持基点的 \(r:\widetilde X\rightarrow E\)。它满射：对任意 \(e\in E\)，从 \(e_0\) 到 \(e\) 选道路，再把其投影提升到 \(\widetilde X\)，由提升唯一性可知终点经 \(r\) 映到 \(e\)。在 \(X\) 中取同时被 \(p,q\) 均匀覆盖的道路连通小邻域；\(r\) 把 \(p\) 的每一叶同胚映到 \(q\) 的一叶，因此 \(r\) 本身也是覆盖。
\end{proof}

半局部单连通性是存在单连通覆盖的必要条件：若有这样的覆盖，在任意均匀覆盖邻域 \(U\) 中的基点环路都可在一叶内闭合提升，而在单连通覆盖中零伦；投影后，它就在 \(X\) 中零伦。例如，平面中以 \((1/n,0)\) 为圆心、\(1/n\) 为半径的圆的并，在共同点 \(0\) 的任意邻域内都含完整的小圆。把其他圆压到 \(0\)、保留这一小圆，得到到该圆的连续缩回：在 \(0\) 处它不增大到 \(0\) 的距离，在保留圆的其他点附近则为恒等。因此小圆环路在整个空间中不零伦。这一空间不是半局部单连通，不能用上述构造得到单连通覆盖。

\begin{theorem}[连通覆盖的子群分类]\label{g2:thm:covering-classification}
	设 \(X\) 非空、道路连通、局部道路连通且半局部单连通，\(x_0\in X\)，\(G=\pi_1(X,x_0)\)。带基点连通覆盖 \(p:(E,e_0)\rightarrow(X,x_0)\) 的保持基点同构类，与 \(G\) 的子群一一对应，所对应的子群为
	\[
		H=p_*\pi_1(E,e_0)\subset G.
	\]
	任意子群 \(H\leq G\) 可由以下道路类空间实现。令 \(\widetilde X\) 为以 \(x_0\) 为起点的定端点同伦道路类所构成的万有覆盖；对从 \(x_0\) 出发的道路，规定
	\[
		\alpha\sim_H\beta
		\quad\Longleftrightarrow\quad
		\alpha(1)=\beta(1),\qquad[\alpha*\bar\beta]\in H.
	\]
	在 \(X_H=\widetilde X/{\sim_H}\) 上采用商拓扑，\(p_H([\alpha]_H)=\alpha(1)\) 是连通覆盖，其指定基点为 \([e_{x_0}]_H\)。不指定覆盖基点时，同构类对应于 \(G\) 的子群的共轭类。进一步，对两个带基点连通覆盖 \(p_i:(E_i,e_i)\rightarrow(X,x_0)\)，令 \(H_i=(p_i)_*\pi_1(E_i,e_i)\)（\(i=1,2\)），则存在连续映射 \(r:E_1\rightarrow E_2\) 满足 \(r(e_1)=e_2\) 及 \(p_2r=p_1\)，当且仅当 \(H_1\subset H_2\)；该映射若存在则唯一且为覆盖。
\end{theorem}
\begin{proof}
	\(\sim_H\) 的自反、对称、传递分别来自 \(H\) 含单位、对取逆封闭和对乘法封闭；传递时使用 \([\alpha*\bar\beta][\beta*\bar\gamma]=[\alpha*\bar\gamma]\)。它也不依赖道路代表。对万有覆盖上的一个基本集合 \(B(\alpha,U)\)，投影到 \(X_H\) 后仍与 \(U\) 一一对应。两个在同一 \(U\) 上的此类集合若有一点被识别，先在该底点重选两条代表道路 \(\alpha,\beta\)，便可在每个同一底点处同时被识别，因为
	\[
		[(\alpha*\eta)*\overline{\beta*\eta}]=[\alpha*\bar\beta]
	\]
	对共同尾段 \(\eta\) 成立。它们因而整叶相同或不交。商映射是开映射：基本开集的饱和集是同类基本开集的并。于是其像是开叶，且每一叶上的投影为同胚，证明 \(p_H\) 是覆盖。\(X_H\) 是道路连通空间 \(\widetilde X\) 的连续像，所以道路连通。

	基点环路 \(\alpha\) 的提升终点为 \([\alpha]_H\)，它闭合恰好等价于 \([\alpha]\in H\)。由提升判据中像的刻画，\((p_H)_*\pi_1(X_H,[e_{x_0}]_H)=H\)，因此每个子群都实现。

	对已有两个覆盖，\(E_1\) 道路连通且局部道路连通；把 \(p_1\) 看成待提升映射，\cref{g2:prop:covering-lift-criterion}恰给出条件 \(H_1\subset H_2\)。提升唯一，并由同时均匀覆盖的小邻域看出它是局部的逐叶同胚；其满射性与万有覆盖证明末段相同，所以是覆盖。若两子群相等，就有两个方向的保持基点提升，它们的复合与恒等映射在基点一致，提升唯一性使复合为恒等。因此两个覆盖保持基点同构；反向结论从投影相容性立即成立。

	最后，在同一纤维中把基点从 \(e_0\) 换到 \(e_1\)，选提升空间中的道路 \(\widetilde\gamma:e_0\rightarrow e_1\)，记投影环路的类为 \(g=[p\widetilde\gamma]\)。基点变化公式给出新子群 \(g^{-1}Hg\)。反之，任意 \(g\in G\) 都有从 \(e_0\) 出发的提升，取其终点便实现这个共轭子群。结合带基点分类，得到不指定基点的分类。
\end{proof}

上述存在性构造需要半局部单连通性；提升判据和两个已有覆盖之间的唯一性则不需要它。区分这两种用途，才能在不具备万有覆盖的空间中继续正确使用道路提升。关于这些构造可参看 \cite[第1.3节]{Hatcher2002}；这里道路类商空间的等价关系也将在下一小节解释为基本群的作用。

\subsection{基本群作用和单值}
\label{g2:subsec:0101:04}
覆盖分类把一个空间换成子群，单值则把同一信息放在纤维点怎样被环路置换上。它还说明，覆盖变换与沿环路移动纤维点是两种不同操作；左右次序只有在写出提升公式后才清楚。

\begin{definition}[覆盖的单值作用与覆盖变换]\label{g2:def:monodromy-deck}
	设 \(p:E\rightarrow X\) 为覆盖，\(x_0\in X\)，\(G=\pi_1(X,x_0)\)，\(F=p^{-1}(x_0)\)。对 \(e\in F\) 与 \(g=[\alpha]\in G\)，令 \(e\cdot g\) 为环路 \(\alpha\) 从 \(e\) 出发的唯一提升的终点，称此为纤维上的\term[danzhizuoyong]{单值作用}[monodromy action]。与投影相容的自同胚 \(D:E\rightarrow E\)，即满足 \(pD=p\) 的自同胚，称为\term[fugaibianhuan]{覆盖变换}[deck transformation]；这些变换在通常函数复合下组成的群记为 \(\operatorname{Deck}(E/X)\)。\symidx{\(\operatorname{Deck}(E/X)\)：覆盖变换群}
\end{definition}

\begin{proposition}[右单值与左覆盖变换]\label{g2:prop:monodromy-deck}
	设 \(X\) 道路连通且局部道路连通，\(p:E\rightarrow X\) 为连通覆盖，固定 \(x_0\in X\) 和 \(e_0\in p^{-1}(x_0)\)，并记 \(G=\pi_1(X,x_0)\)、\(H=p_*\pi_1(E,e_0)\)、\(F=p^{-1}(x_0)\)。对 \(e\in F\) 与 \(g=[\alpha]\in G\)，令 \(e\cdot g\) 为 \(\alpha\) 从 \(e\) 出发的提升的终点。则这是传递右作用，对任意 \(e\in F\)、\(g,k\in G\) 有
	\[
		(e\cdot g)\cdot k=e\cdot(gk),\qquad e\cdot1=e.
	\]
	\(e_0\) 的稳定子为 \(H\)，并有右 \(G\) 集合同构
	\[
		H\backslash G\longrightarrow p^{-1}(x_0),\qquad Hg\longmapsto e_0\cdot g,
	\]
	其中 \(H\backslash G=\{Hg:g\in G\}\)，右作用为 \((Hg)\cdot k=Hgk\)。因此覆盖叶数等于指数 \([G:H]\)。覆盖变换与单值交换，且
	\[
		\operatorname{Deck}(E/X)\cong N_G(H)/H,
		\qquad N_G(H)=\{a\in G:aHa^{-1}=H\}.
	\]
	若再假设 \(X\) 半局部单连通，并使用道路类模型 \(E=X_H\)，其中 \([\gamma]_H\) 表示从 \(x_0\) 出发的道路按同终点且 \([\gamma*\bar\eta]\in H\) 的关系识别所得的类，则上述同构具体为
	\[
		aH\longmapsto D_a,\qquad D_a([\gamma]_H)=[a*\gamma]_H
		\quad(a\in N_G(H)),
	\]
	式中 \(a\) 任选一条代表环路。特别地，在万有覆盖 \(\widetilde X\) 上有左作用
	\[
		D_a([\gamma])=[a*\gamma],\qquad D_a\circ D_b=D_{ab},
	\]
	且 \(\operatorname{Deck}(\widetilde X/X)\cong G\)。
\end{proposition}
\begin{proof}
	同伦提升保证 \(e\cdot[\alpha]\) 不依赖代表；拼接提升先提升 \(\alpha\)，再从所得终点提升下一条环路，因此恰为右作用。常值道路提升仍恒定。空间 \(E\) 连通且局部道路连通，因而道路连通；任意两纤维点之间的道路投影为底环路，故作用传递。提升判据中对闭合提升的刻画，说明 \(e_0\) 的稳定子正是 \(H\)。于是 \(e_0\cdot g=e_0\cdot k\) 当且仅当 \(gk^{-1}\in H\)，也就是 \(Hg=Hk\)，给出所述双射及右作用相容性。

	若 \(D\) 是覆盖变换，则 \(D\widetilde\alpha\) 正是 \(\alpha\) 从 \(D(e)\) 出发的提升，所以
	\[
		D(e\cdot g)=D(e)\cdot g.
	\]
	而且 \(D\) 只要固定一个点，就由提升唯一性成为恒等。任意覆盖变换因此由 \(D(e_0)\) 决定。写 \(D(e_0)=e_0\cdot a\)，这一点的稳定子为 \(a^{-1}Ha\)；\(D\) 可逆且与单值交换，故两点的稳定子相同，得到 \(a\in N_G(H)\)。反之，对 \(a\in N_G(H)\)，\(e_0\) 与 \(e_0\cdot a\) 对应同一子群，已有覆盖间的提升判据给出唯一覆盖变换 \(D_a\)，把前一点送到后一点。这一论证不需要半局部单连通性。

	在纤维的陪集模型上，\(D_a(Hg)=Hag\)。因此
	\[
		D_aD_b(Hg)=Habg=D_{ab}(Hg),
	\]
	再由一个纤维点决定覆盖变换，得到 \(D_a\circ D_b=D_{ab}\)；核为使 \(Ha=H\) 的元素，即 \(H\)。由前段的满射性，得到 \(N_G(H)/H\) 的群同构。

	在道路类模型中，若 \(\gamma\sim_H\eta\)，则
	\[
		[(a*\gamma)*\overline{a*\eta}]
		=a[\gamma*\bar\eta]a^{-1}\in H,
	\]
	所以前置环路公式良定义。它把基本开叶送到基本开叶，逆为 \(D_{a^{-1}}\)，故为覆盖变换，且把基点送到 \([a]_H\)，与刚才的 \(D_a\) 相同。取 \(H=\{1\}\) 就得到万有覆盖上的左作用。对任意 \(H\)，等价关系 \(\sim_H\) 也恰是此左作用中 \(H\) 的轨道关系，因此 \(X_H=H\backslash\widetilde X\)。
\end{proof}

若把右单值写成置换 \(R_g(e)=e\cdot g\)，通常函数复合满足 \(R_{gk}=R_k\circ R_g\)；它与左覆盖变换的 \(D_{gk}=D_g\circ D_k\) 不同。即使在万有覆盖的同一纤维上，前者是 \(h\mapsto hg\)，后者是 \(h\mapsto gh\)。这解释了为什么不能把“沿环路移动终点”直接当作同一个左群表示。

\begin{corollary}[正规覆盖]\label{g2:cor:regular-covering}
	设 \(X\) 道路连通且局部道路连通，\(p:E\rightarrow X\) 为连通覆盖，\(x_0\in X\)，\(e_0\in p^{-1}(x_0)\)，\(H=p_*\pi_1(E,e_0)\leq G=\pi_1(X,x_0)\)。覆盖变换群在每个纤维上传递，当且仅当 \(H\trianglelefteq G\)。满足此条件的覆盖称为\term[zhengguifugai]{正规覆盖}[regular covering]；此时 \(\operatorname{Deck}(E/X)\cong G/H\)。
\end{corollary}
\begin{proof}
	在 \(H\backslash G\) 中，覆盖变换把 \(H\) 送到 \(Ha\)，其中恰好允许 \(a\in N_G(H)\)。这能够到达全部陪集，当且仅当 \(N_G(H)=G\)，即 \(H\) 正规。若在一个纤维上传递，沿连接两个底点的道路提升便把该纤维与另一纤维双射，并与覆盖变换交换，故在每个纤维上传递。群同构由\cref{g2:prop:monodromy-deck}给出。
\end{proof}

\begin{example}[圆的覆盖与循环单值]\label{g2:exm:circle-coverings}
	实指数映射 \(p:\mathbb R\rightarrow S^1\)，\(p(t)=\exp(2\pi\mathrm{i}t)\)，是万有覆盖：小圆弧的原像是互不相交的开区间，且 \(\mathbb R\) 可缩。正向生成环路从整数 \(j\) 提升为 \(t\mapsto j+t\)，把纤维点 \(j\) 送到 \(j+1\)。覆盖变换为 \(t\mapsto t+k\)，\(k\in\mathbb Z\)。

	对整数 \(n\geq1\)，覆盖 \(p_n:S^1\rightarrow S^1\)，\(p_n(z)=z^n\)，有纤维 \(\{\zeta^j:0\leq j<n\}\)，其中 \(\zeta=\exp(2\pi\mathrm{i}/n)\)。正向基环路从 \(\zeta^j\) 提升为 \(\zeta^j\exp(2\pi\mathrm{i}t/n)\)，单值因而是 \(j\mapsto j+1\pmod n\)。对应子群为 \(n\mathbb Z\)，覆盖变换为 \(z\mapsto\zeta^kz\)，其群为 \(\mathbb Z/n\mathbb Z\)。\(\mathbb Z\) 的子群只有 \(n\mathbb Z\)（\(n\geq1\)）与零子群，所以这些有限覆盖和实线覆盖已列出圆的全部连通覆盖，且全部正规。
\end{example}

\begin{example}[三叶覆盖可以没有非恒等覆盖变换]\label{g2:exm:nonregular-three-cover}
	设 \(R_2=S^1\vee S^1\)，两个定向圆分别标记为 \(a,b\)。取三个顶点 \(1,2,3\)，从每个顶点 \(j\) 画一条标记为 \(a\) 的有向边，终点按置换 \((12)\) 给出；再画一条标记为 \(b\) 的有向边，终点按置换 \((123)\) 给出。把每条标记边同向映到 \(R_2\) 的相应圆，得到三叶覆盖：每个顶点处每一标记恰有一个入边端和一个出边端，所以小星状邻域逐叶同胚映到底图的小星状邻域。\(b\) 边已连接三个顶点，故覆盖连通。

	单值由 \((12)\)、\((123)\) 生成，在三个点上传递。覆盖变换在纤维上的置换必须同时与这两者交换；与三循环交换的置换只能是该三循环的幂，其中只有恒等置换又与 \((12)\) 交换。因此覆盖变换群平凡。特别地，传递单值并不蕴含覆盖变换传递。这一覆盖对应 \(F(a,b)\) 中一个指数为 \(3\) 的非正规子群；覆盖图有 \(6\) 条边、\(3\) 个顶点，所以该子群由覆盖诱导单射识别为自由群 \(F_4\)。
\end{example}
% 纲要条目排为编号小节标题，并进入目录与书签；说明段仍保留于纲要注释中。

% !TeX root = ../../Book2.tex
% 纲要标题：### 1.2 粘合的同伦性质
% 纲要位置：几何学计划纲要.md，第 1009 行。
% 正文状态：未撰写；本文件目前是框架。

% 纲要内容开始
% 
% * 映射柱、映射锥与相对空间
% * 同伦延拓性质
% * 余纤维化与良好空间对
% * 直接商与同伦余极限的区别
% 
% 纲要内容结束

\section{粘合的同伦性质}
\label{g2:sec:0102}

待撰写

\subsection{映射柱、映射锥与相对空间}
\label{g2:subsec:0102:01}
待撰写

\subsection{同伦延拓性质}
\label{g2:subsec:0102:02}
待撰写

\subsection{余纤维化与良好空间对}
\label{g2:subsec:0102:03}
待撰写

\subsection{直接商与同伦余极限的区别}
\label{g2:subsec:0102:04}
待撰写
% 纲要条目排为编号小节标题，并进入目录与书签；说明段仍保留于纲要注释中。

% !TeX root = ../../Book2.tex
% 纲要标题：### 1.3 CW 复形的构造
% 纲要位置：几何学计划纲要.md，第 1016 行。
% 正文状态：未撰写；本文件目前是框架。

% 纲要内容开始
% 
% * 胞腔附着、骨架与弱拓扑
% * 闭有限性和紧子集的控制
% * 相对 CW 复形和逐骨架构造
% * 球面、射影空间与 Grassmann 型模型
% 
% 纲要内容结束

\section{CW 复形的构造}
\label{g2:sec:0103}

待撰写

\subsection{胞腔附着、骨架与弱拓扑}
\label{g2:subsec:0103:01}
待撰写

\subsection{闭有限性和紧子集的控制}
\label{g2:subsec:0103:02}
待撰写

\subsection{相对 CW 复形和逐骨架构造}
\label{g2:subsec:0103:03}
待撰写

\subsection{球面、射影空间与 Grassmann 型模型}
\label{g2:subsec:0103:04}
待撰写
% 纲要条目排为编号小节标题，并进入目录与书签；说明段仍保留于纲要注释中。

% !TeX root = ../../Book2.tex
% 纲要标题：### 1.4 胞腔逼近定理
% 纲要位置：几何学计划纲要.md，第 1023 行。
% 正文状态：未撰写；本文件目前是框架。

% 纲要内容开始
% 
% * 球面映射的降维引理
% * 绝对和相对胞腔逼近
% * 同伦类和附着映射的简化
% * 有限复形与一般 CW 空间的区别
% 
% ---
% 
% 纲要内容结束

\section{胞腔逼近定理}
\label{g2:sec:0104}

待撰写

\subsection{球面映射的降维引理}
\label{g2:subsec:0104:01}
待撰写

\subsection{绝对和相对胞腔逼近}
\label{g2:subsec:0104:02}
待撰写

\subsection{同伦类和附着映射的简化}
\label{g2:subsec:0104:03}
待撰写

\subsection{有限复形与一般 CW 空间的区别}
\label{g2:subsec:0104:04}
待撰写
% 纲要条目排为编号小节标题，并进入目录与书签；说明段仍保留于纲要注释中。


\begin{chaptersummary}
待撰写
\end{chaptersummary}

\BookExercises
待撰写

% !TeX root = ../../Book2.tex
% 纲要标题：## 第2章　奇异同调、切除与连接映射
% 纲要位置：几何学计划纲要.md，第 1032 行。
% 正文状态：未撰写；本文件目前是框架。

% 纲要内容开始
% 
% > 本章目标与证明范围：在链层面建立同调理论；棱柱算子、重心细分和小链复形均有实际构造。长正合列不是定理名称串联，而由短正合复形和具体连接映射产生。
% 
% 纲要内容结束

\chapter{奇异同调、切除与连接映射}
\label{g2:ch:02}
\BookChapterNumber{g2:ch:02}

\begin{chapterabstract}
待撰写
\end{chapterabstract}

\begin{learninggoals}
待撰写
\end{learninggoals}

待撰写

% !TeX root = ../../Book2.tex
% 纲要标题：### 2.1 链复形及其代数
% 纲要位置：几何学计划纲要.md，第 1036 行。
% 正文状态：未撰写；本文件目前是框架。

% 纲要内容开始
% 
% * 链、圈、边界和增广
% * 链映射、链同伦与映射锥
% * 短正合复形和连接同态
% * 长正合列的正合性与自然性
% 
% 纲要内容结束

\section{链复形及其代数}
\label{g2:sec:0201}

待撰写

\subsection{链、圈、边界和增广}
\label{g2:subsec:0201:01}
待撰写

\subsection{链映射、链同伦与映射锥}
\label{g2:subsec:0201:02}
待撰写

\subsection{短正合复形和连接同态}
\label{g2:subsec:0201:03}
待撰写

\subsection{长正合列的正合性与自然性}
\label{g2:subsec:0201:04}
待撰写
% 纲要条目排为编号小节标题，并进入目录与书签；说明段仍保留于纲要注释中。

% !TeX root = ../../Book2.tex
% 纲要标题：### 2.2 奇异链与同伦不变性
% 纲要位置：几何学计划纲要.md，第 1093 行。
% 正文状态：初稿；尚未正式审核。

% 纲要内容开始
% 
% * 奇异单形和边界符号
% * 棱柱算子的显式公式
% * 相对链复形和约化同调
% * 对及三元组的长正合列
% 
% 纲要内容结束

\section{奇异链与同伦不变性}
\label{g2:sec:0202}

把空间分成三角形再计算边界，是理解同调的一种直接方式；但一般空间未必已经给定三角剖分，连续映射也未必把某个三角形送到另一个三角形。奇异链采用另一种做法：把所有从标准单形到空间的连续映射都作为生成元。这样，连续映射只需与这些参数映射复合，就能作用于链。链群虽然很大，边界公式却仍由单形的有限个面决定。

本节始终使用整数系数。先从单形各面的次序确定符号，再把同伦所扫过的棱柱拆成单形，以说明同伦为什么在同调上消失。指定子空间以后，同样的链模型还能记录“边界允许落在该子空间中”的链；空间对的连接同态，正是把这条尚未消失的边界取出来。相关链构造可参阅 \cite[第2章，§2.1]{Hatcher2002}。

\subsection{奇异单形和边界符号}
\label{g2:subsec:0202:01}
单形的面需要保留参数次序，而不能只记其像。取
\[
	\Delta^n=\left\{(t_0,\ldots,t_n)\in\mathbb R^{n+1}:
	t_i\ge0,\ \sum_{i=0}^nt_i=1\right\},
\]
并按次序记其顶点为 \(v_0,\ldots,v_n\)。当 \(n\ge1\) 时，仿射面映射
\[
	\delta_i^n:\Delta^{n-1}\longrightarrow\Delta^n,\qquad
	\delta_i^n(t_0,\ldots,t_{n-1})
	=(t_0,\ldots,t_{i-1},0,t_i,\ldots,t_{n-1})
\]
将标准单形送到略去第 \(i\) 个顶点的面。上标表示靶单形的维数。

\begin{definition}[奇异链与奇异同调]\label{g2:def:singular-chain}
	设 \(X\) 为拓扑空间，\(n\ge0\) 为整数，取标准单形
	\[
		\Delta^n=\{(t_0,\ldots,t_n)\in\mathbb R^{n+1}:
		t_i\ge0,\ \textstyle\sum_{i=0}^nt_i=1\},
	\]
	其顶点按标准基的顺序编号。从 \(\Delta^n\) 到 \(X\) 的连续映射 \(\sigma:\Delta^n\to X\) 称为一个\term[qiyidanxing]{奇异单形}[singular simplex]。以所有这样的映射为基的自由阿贝尔群记作
	\[
		C_n(X)=\bigoplus_{\sigma:\Delta^n\to X\ \text{连续}}\mathbb Z\,\sigma.
	\]
	它的元素是有限形式和 \(c=\sum_\sigma a_\sigma\sigma\)，其中 \(a_\sigma\in\mathbb Z\)，称为整数系数的\term[qiyilian]{奇异链}[singular chain]。规定 \(C_n(X)=0\)（\(n<0\)），并定义
	\[
		\partial_n\sigma=\sum_{i=0}^n(-1)^i\sigma\circ\delta_i^n
		\quad(n\ge1),\qquad \partial_0=0,
	\]
	其中 \(\delta_i^n:\Delta^{n-1}\to\Delta^n\) 是在第 \(i\) 个坐标插入 \(0\) 的保序仿射面映射，坐标从 \(0\) 起编号。将边界按加法延拓后，它满足 \(\partial_{n-1}\partial_n=0\)，所得链复形记作 \(C_\bullet(X)\)。其同调
	\[
		H_n(X;\mathbb Z)=H_n(X)
		=\frac{\ker(\partial_n:C_n(X)\to C_{n-1}(X))}
		{\operatorname{im}(\partial_{n+1}:C_{n+1}(X)\to C_n(X))}
	\]
	称为 \(X\) 的整数系数\term[qiyitongdiao]{奇异同调}[singular homology]；当 \(n<0\) 时规定 \(H_n(X)=0\)。\symidx{\(C_n(X)\)：整数系数奇异链群}\symidx{\(H_n(X)\)：整数系数奇异同调群}
\end{definition}

这里没有要求 \(\sigma\) 单射，也不要求它的像具有 \(n\) 维。常值映射同样是生成元；像相同而参数映射不同的两个单形，也仍是不同的生成元。特别地，改变顶点次序得到的奇异单形，并不在链群的定义中直接等同于原单形的正号或负号。所有抵消都必须落实为同一参数映射的系数相加。

定义中所用的 \(\partial^2=0\) 来自每个余维二的面恰好出现两次。若 \(0\le i<j\le n\)，则
\[
	\delta_j^n\circ\delta_i^{n-1}
	=\delta_i^n\circ\delta_{j-1}^{n-1};
\]
两边都按原有次序嵌入略去 \(v_i,v_j\) 的面。因此，当 \(n\ge2\) 时，
\[
	\partial_{n-1}\partial_n\sigma
	=\sum_{0\le i<j\le n}
	\left((-1)^{j+i}+(-1)^{i+j-1}\right)
	\sigma\circ\delta_j^n\circ\delta_i^{n-1}
	=0.
\]
当 \(n=1\) 时，结论由 \(\partial_0=0\) 得到。这既核验了符号，也证明了定义中的商确有意义。

例如，把一个奇异一维单形看成道路 \(\gamma:[0,1]\to X\)，其中参数 \(t=t_1\)，便有
\[
	\partial\gamma=\gamma(1)-\gamma(0).
\]
这里的点代表相应的零维奇异单形。对于一个三角形参数映射 \(\sigma\)，有
\[
	\partial\sigma
	=\sigma|[v_1,v_2]-\sigma|[v_0,v_2]+\sigma|[v_0,v_1],
\]
其中每个限制都通过保序仿射映射重新参数化为定义在 \(\Delta^1\) 上的映射。再次取边界时，三个顶点的系数分别抵消。

\begin{proposition}[连续映射诱导的链映射]\label{g2:prop:singular-functoriality}
	设 \(f:X\to Y\) 为拓扑空间之间的连续映射，\(C_\bullet(X)\)、\(C_\bullet(Y)\) 为整数系数奇异链复形。对每个生成元定义 \(f_\#(\sigma)=f\circ\sigma\)，并按加法延拓，则
	\[
		\partial f_\#=f_\#\partial.
	\]
	因此 \(f_\#\) 诱导同态 \(f_*:H_n(X)\to H_n(Y)\)，满足 \(f_*[c]=[f_\#c]\)。对于连续映射 \(g:Y\to Z\)，有
	\[
		(g\circ f)_\#=g_\#f_\#,\qquad
		(g\circ f)_*=g_*f_*,\qquad
		(\operatorname{id}_X)_*=\operatorname{id}_{H_n(X)}.
	\]
\end{proposition}
\begin{proof}
	对 \(n\ge1\) 的生成元，复合映射的结合律给出
	\[
		\partial(f\circ\sigma)
		=\sum_{i=0}^n(-1)^if\circ\sigma\circ\delta_i^n
		=f_\#(\partial\sigma).
	\]
	零次等式由 \(\partial_0=0\) 得到。于是 \(f_\#\) 把圈送到圈，并把形如 \(\partial b\) 的边界送到 \(\partial(f_\#b)\)，故在同调商群上良定义。其余等式先在生成元上由复合映射的结合律和恒等映射性质得到，再下降到同调。
\end{proof}

\begin{proposition}[点和道路分支的同调]\label{g2:prop:homology-point-components}
	设 \(X\) 为任意拓扑空间，\(\mathcal P(X)\) 为其道路分支的集合，\(H_n\) 表示整数系数奇异同调。则对所有 \(n\ge0\)，各分支的包含映射给出自然同构
	\[
		H_n(X)\cong\bigoplus_{Q\in\mathcal P(X)}H_n(Q).
	\]
	若 \(*\) 是单点空间，则 \(H_0(*)\cong\mathbb Z\)、\(H_n(*)=0\)（\(n>0\)）。此外，
	\[
		H_0(X)\cong\mathbb Z[\mathcal P(X)]
		=\bigoplus_{Q\in\mathcal P(X)}\mathbb Z[Q],
	\]
	其中零维圈 \(\sum_i a_i x_i\) 被送到 \(\sum_i a_i[Q(x_i)]\)，\(Q(x_i)\) 是含 \(x_i\) 的道路分支。空集的全部普通同调群均为零。
\end{proposition}
\begin{proof}
	标准单形道路连通，所以每个奇异单形的像完全落在一个道路分支中。由子空间拓扑，它也给出到该分支的连续映射。因而
	\[
		C_n(X)=\bigoplus_{Q\in\mathcal P(X)}C_n(Q),
	\]
	且边界算子保持各直和项。链只有有限项，所以一个链为圈当且仅当各分支中的分量都是圈，一个链为边界当且仅当各分量都是边界。这证明了同调的直和分解，不需要各道路分支为开集。

	单点空间在每个非负次数恰有一个奇异单形 \(s_n\)，所以 \(C_n(*)=\mathbb Zs_n\)，并且
	\[
		\partial s_n=
		\begin{cases}
			s_{n-1},&n>0\ \text{为偶数},\\
			0,&n\ \text{为奇数}.
		\end{cases}
	\]
	相邻的非零边界是恒等同态，故正次数的同调全部为零；由于 \(\partial_0=\partial_1=0\)，零次同调是 \(\mathbb Z\)。

	最后，对每个零维生成元 \(x\) 定义 \(\Phi(x)=[Q(x)]\)。道路的两端在同一道路分支中，故 \(\Phi\partial_1=0\)。映射 \(\Phi:C_0(X)\to\mathbb Z[\mathcal P(X)]\) 满射。若 \(c=\sum_i a_ix_i\) 属于其核，则在每个实际出现的分支中，系数之和为零。在这样的分支中选一点 \(x_Q\)，并为每个 \(x_i\in Q\) 选一条从 \(x_Q\) 到 \(x_i\) 的道路 \(\gamma_i\)。该分支中的链满足
	\[
		\sum_{x_i\in Q}a_ix_i
		=\sum_{x_i\in Q}a_i(x_i-x_Q)
		=\partial\left(\sum_{x_i\in Q}a_i\gamma_i\right).
	\]
	求和只涉及有限个点和有限个分支，故 \(c\in\operatorname{im}\partial_1\)。由此得到所述 \(H_0\) 同构。空集没有非负维奇异单形，所以各链群均为零。
\end{proof}

这个计算辨别的是道路分支，而不只是连通分支；两点之差成为边界的原因，是存在连接它们的道路。它也说明了奇异链的大小与同调的大小可以相差很大：即使单点空间在每个非负次数都有非零链群，其正次数同调仍全部消失。

\subsection{棱柱算子的显式公式}
\label{g2:subsec:0202:02}
若 \(f,g:X\to Y\) 由同伦 \(H:X\times I\to Y\) 连接，一条奇异单形 \(\sigma:\Delta^n\to X\) 在同伦中扫过的对象是 \(\Delta^n\times I\)。它的上下底面分别给出 \(f\circ\sigma\) 和 \(g\circ\sigma\)，侧面则来自 \(\sigma\) 的边界。将这个几何关系写成链等式，需要给棱柱中的每块单形确定次序和系数。

对 Euclid 空间中的点 \(q_0,\ldots,q_m\)，以下用 \([q_0,\ldots,q_m]\) 表示仿射映射
\[
	\Delta^m\longrightarrow\mathbb R^N,\qquad
	(t_0,\ldots,t_m)\longmapsto\sum_{j=0}^m t_jq_j.
\]
当这些点位于某个凸集时，映射的像也位于该凸集。这一记号保留有序参数，不预先附加顶点置换关系。

\begin{theorem}[奇异同调的同伦不变性]\label{g2:thm:singular-homotopy-invariance}
	设 \(X,Y\) 为拓扑空间，\(I=[0,1]\)，连续映射 \(H:X\times I\to Y\) 满足 \(H(x,0)=f(x)\)、\(H(x,1)=g(x)\)。在整数系数奇异链上，对 \(\sigma:\Delta^n\to X\)、\(n\ge0\)，以 \(v_0,\ldots,v_n\) 表示 \(\Delta^n\) 的有序标准顶点，记
	\[
		v_j^\varepsilon=(v_j,\varepsilon)\in\Delta^n\times I
		\quad(0\le j\le n,\ \varepsilon=0,1),\qquad
		F_\sigma=H\circ(\sigma\times\operatorname{id}_I).
	\]
	定义\term[lengzhu suanzi]{棱柱算子}[prism operator]
	\[
		P_n(\sigma)=\sum_{i=0}^n(-1)^i
		F_\sigma\circ
		[v_0^0,\ldots,v_i^0,v_i^1,\ldots,v_n^1],
	\]
	其中方括号表示由所列有序顶点确定的仿射单形。按加法延拓为 \(P_n:C_n(X)\to C_{n+1}(Y)\)，并规定 \(P_{-1}=0\)，则
	\begin{equation}\label{g2:eq:prism-identity}
		\partial_{n+1}P_n+P_{n-1}\partial_n=g_\#-f_\#
		\qquad(n\ge0).
	\end{equation}
	因此 \(f_*=g_*:H_n(X)\to H_n(Y)\) 对所有次数成立。
\end{theorem}
\begin{proof}
	先说明公式中的单形如何填满棱柱，再核对每类面的符号。对点 \(((t_0,\ldots,t_n),s)\in\Delta^n\times I\)，取满足
	\[
		\sum_{j=i+1}^nt_j\le s\le\sum_{j=i}^nt_j
	\]
	的指标 \(i\)。这样的指标总存在。在第 \(i\) 块单形中，顶点 \(v_j^0\)（\(j<i\)）的系数为 \(t_j\)，顶点 \(v_j^1\)（\(j>i\)）的系数也为 \(t_j\)，而两个同指标顶点的系数分别为
	\[
		\sum_{j=i}^nt_j-s,\qquad
		s-\sum_{j=i+1}^nt_j.
	\]
	这些数非负，总和为 \(1\)，并给出原点的坐标。因此这 \(n+1\) 块单形覆盖棱柱；两个不等式取等号的位置正是相邻块共享的面。

	核对边界时，简记 \(u_j=v_j^0\)、\(w_j=v_j^1\)，并在各式中省略共同的后复合 \(F_\sigma\)。对第 \(i\) 块
	\[
		T_i=[u_0,\ldots,u_i,w_i,\ldots,w_n],
	\]
	略去 \(u_j\)（\(j\le i\)）的边界系数是 \((-1)^{i+j}\)，略去 \(w_j\)（\(j\ge i\)）的系数是 \((-1)^{i+j+1}\)，因为 \(w_j\) 在 \(T_i\) 中的位置为 \(j+1\)。

	先看 \(j=i\) 的两类面。略去 \(u_i\) 得到
	\[
		[u_0,\ldots,u_{i-1},w_i,\ldots,w_n],
	\]
	系数为 \(+1\)；当 \(i\ge1\) 时，它恰好等于从 \(T_{i-1}\) 中略去 \(w_{i-1}\) 的面，后者系数为 \(-1\)。这些内部分界面逐对抵消。剩下的两项是从 \(T_0\) 略去 \(u_0\) 得到的上底面 \([w_0,\ldots,w_n]\)，以及从 \(T_n\) 略去 \(w_n\) 得到的下底面 \(-[u_0,\ldots,u_n]\)。复合 \(F_\sigma\) 后，它们分别是 \(g_\#\sigma\) 和 \(-f_\#\sigma\)。

	余下的项来自侧面。若 \(j<i\)，从 \(T_i\) 略去 \(u_j\) 后，所得有序单形属于第 \(j\) 个底面上的棱柱；略去该底面顶点后，转折指标由 \(i\) 变成 \(i-1\)。它在 \(P_{n-1}\partial_n\sigma\) 中的系数为
	\[
		(-1)^j(-1)^{i-1}=-(-1)^{i+j},
	\]
	恰与它在 \(\partial_{n+1}P_n\sigma\) 中的系数相反。若 \(j>i\)，从 \(T_i\) 略去 \(w_j\) 后，底面棱柱的转折指标仍为 \(i\)，故它在 \(P_{n-1}\partial_n\sigma\) 中的系数为
	\[
		(-1)^j(-1)^i=-(-1)^{i+j+1}.
	\]
	这两类已经包括全部侧面。因此，内部面抵消、侧面与 \(P\partial\) 抵消以后，只留下上底面减下底面，得到 \eqref{g2:eq:prism-identity}。当 \(n=0\) 时，\(P_0(x)\) 就是道路 \(t\mapsto H(x,t)\)，其边界直接等于 \(g(x)-f(x)\)，同一公式仍成立。

	若 \(c\) 是奇异 \(n\) 圈，则
	\[
		g_\#c-f_\#c=\partial P_nc,
	\]
	所以二者代表同一个同调类。这证明了同伦不变性。
\end{proof}

在 \(n=1\) 时，棱柱是一个正方形，公式把它沿连接 \(v_0^0\) 与 \(v_1^1\) 的对角线分成两个三角形：
\[
	P_1(\sigma)=F_\sigma\circ[v_0^0,v_0^1,v_1^1]
	-F_\sigma\circ[v_0^0,v_1^0,v_1^1].
\]
共同对角线在边界中出现两次，系数相反；左右两边则是两个端点的同伦道路。这是 \(\partial P+P\partial=g_\#-f_\#\) 在最低非平凡维数中的具体形状。

\begin{corollary}[同伦等价与可缩空间]\label{g2:cor:contractible-homology}
	设 \(f:X\to Y\) 为拓扑空间之间的同伦等价，\(H_n\) 为整数系数奇异同调。则 \(f_*:H_n(X)\to H_n(Y)\) 对每个 \(n\ge0\) 都是同构。特别地，若非空空间 \(X\) 可缩，即 \(\operatorname{id}_X\) 同伦于一个常值映射，则
	\[
		H_n(X)\cong
		\begin{cases}
			\mathbb Z,&n=0,\\
			0,&n>0.
		\end{cases}
	\]
\end{corollary}
\begin{proof}
	取 \(f\) 的同伦逆 \(u:Y\to X\)。函子性及同伦不变性给出 \(u_*f_*=\operatorname{id}\)、\(f_*u_*=\operatorname{id}\)。非空可缩空间到单点空间的映射，与所选收缩点的包含映射互为同伦逆，再应用 \cref{g2:prop:homology-point-components} 即得所述计算。
\end{proof}

棱柱公式还保留了映射的所在范围：\(P_n\sigma\) 的每一项都落在 \(H(\sigma(\Delta^n)\times I)\) 中。这一点使同一个构造能够用于空间对，而不必另选一种同伦算子。

\subsection{相对链复形和约化同调}
\label{g2:subsec:0202:03}
设一条链的边界没有在整个空间中消失，却完全落在指定子空间 \(A\) 内。如果只研究 \(X\) 相对于 \(A\) 的变化，应允许这样的链成为圈。实现这一点的代数操作，是将 \(A\) 中的全部链置为零。由于边界仍留在 \(A\) 中，这个操作与边界算子相容。

\begin{definition}[相对奇异同调]\label{g2:def:relative-singular-homology}
	设 \(X\) 为拓扑空间，\(A\subseteq X\) 赋予子空间拓扑，\(C_\bullet(X)\)、\(C_\bullet(A)\) 为整数系数奇异链复形。通过包含映射将 \(C_n(A)\) 视为 \(C_n(X)\) 的子群，定义
	\[
		C_n(X,A)=C_n(X)/C_n(A),\qquad
		\partial(c+C_n(A))=\partial c+C_{n-1}(A).
	\]
	由于 \(\partial C_n(A)\subseteq C_{n-1}(A)\)，此边界良定义并满足 \(\partial^2=0\)。商链复形 \(C_\bullet(X,A)\) 的同调
	\[
		H_n(X,A;\mathbb Z)=H_n(X,A)
		=H_n(C_\bullet(X,A))
	\]
	称为空间对 \((X,A)\) 的整数系数\term[xiangdui qiyi tongdiao]{相对奇异同调}[relative singular homology]。负次数的相对链群及相对同调均为零。\symidx{\(C_n(X,A)\)：相对奇异链群}\symidx{\(H_n(X,A)\)：相对奇异同调群}
\end{definition}

上述包含在链上确实是单射：两个到 \(A\) 的连续映射，与包含映射复合后相等，当且仅当它们本身相等；而一个到 \(X\) 的奇异单形只要像在 \(A\) 中，就由子空间拓扑给出到 \(A\) 的连续映射。因此 \(C_n(A)\) 恰由那些像完全落在 \(A\) 中的生成元张成。

记 \(\overline c=c+C_n(A)\)。它是相对圈，当且仅当
\[
	\partial c\in C_{n-1}(A).
\]
两个这样的链 \(c,c'\) 给出同一个相对同调类，当且仅当存在 \(b\in C_{n+1}(X)\)、\(a\in C_n(A)\)，使
\[
	c-c'=\partial b+a.
\]
特别地，相对同调中的零类允许由“一个边界加一条 \(A\) 中的链”表示。这是商复形的圈和边界条件展开后的原式。

若只看阿贝尔群，\(C_n(X,A)\) 可以用像不完全落在 \(A\) 中的奇异单形作为一组基；但这些生成元的某些面可能落在 \(A\) 中，所以它们在 \(C_n(X)\) 中张成的子群未必组成子复形。相对边界是在商群中计算的。两个极端情形为
\[
	C_\bullet(X,\varnothing)=C_\bullet(X),\qquad
	C_\bullet(X,X)=0.
\]

\begin{proposition}[相对同调的函子性与同伦不变性]\label{g2:prop:relative-homotopy-invariance}
	设 \((X,A)\)、\((Y,B)\) 为空间对，\(f:(X,A)\to(Y,B)\) 为连续空间对映射，即 \(f:X\to Y\) 连续且 \(f(A)\subseteq B\)。则
	\[
		f_\#(\overline c)=\overline{f_\#c}
	\]
	给出整数系数相对链复形之间的链映射，并诱导 \(f_*:H_n(X,A)\to H_n(Y,B)\)。这些映射保持恒等映射和复合映射。若空间对映射 \(f,g\) 之间存在连续同伦 \(H:X\times I\to Y\)，满足
	\[
		H(x,0)=f(x),\quad H(x,1)=g(x),\quad H(A\times I)\subseteq B,
	\]
	则 \(f_*=g_*\)。
\end{proposition}
\begin{proof}
	条件 \(f(A)\subseteq B\) 保证 \(f_\#C_n(A)\subseteq C_n(B)\)，所以映射在商上良定义。它与边界交换，以及保持复合与恒等映射，都由绝对链上的相应等式下降得到。

	对于所给同伦，\cref{g2:thm:singular-homotopy-invariance} 的棱柱算子满足
	\[
		P_n(C_n(A))\subseteq C_{n+1}(B),
	\]
	因为这些棱柱的像包含在 \(H(A\times I)\) 中。故 \(P_n\) 下降为相对链算子，且等式 \(\partial P+P\partial=g_\#-f_\#\) 在商上仍成立。将其用于相对圈，就得到两个诱导同态相等。
\end{proof}

另一个常用调整发生在零次。一个非空道路连通空间总有 \(H_0\cong\mathbb Z\)，即使它可缩也不例外。若希望可缩空间在每个次数的同调都为零，可以把“零维链的总系数”也当作一次边界计算。

\begin{definition}[约化奇异同调]\label{g2:def:reduced-homology}
	设 \(X\) 为任意拓扑空间。对整数系数奇异链复形定义增广
	\[
		\varepsilon_X:C_0(X)\longrightarrow\mathbb Z,\qquad
		\varepsilon_X\left(\sum_i a_ix_i\right)=\sum_i a_i.
	\]
	定义增广链复形 \(\widetilde C_\bullet(X)\) 为
	\[
		\widetilde C_n(X)=
		\begin{cases}
			C_n(X),&n\ge0,\\
			\mathbb Z,&n=-1,\\
			0,&n<-1,
		\end{cases}
		\qquad
		\widetilde\partial_n=
		\begin{cases}
			\partial_n,&n\ge1,\\
			\varepsilon_X,&n=0,\\
			0,&n<0.
		\end{cases}
	\]
	道路两端的总系数之差为零，所以 \(\varepsilon_X\partial_1=0\)。其同调 \(\widetilde H_n(X)=H_n(\widetilde C_\bullet(X))\) 称为整数系数\term[yuehua tongdiao]{约化同调}[reduced homology]。对于连续映射 \(f:X\to Y\)，增广链映射在非负次数取 \(f_\#\)，在 \(-1\) 次取 \(\operatorname{id}_{\mathbb Z}\)。\symidx{\(\widetilde H_n(X)\)：整数系数约化同调}
\end{definition}

这里明确保留了空集：\(\varepsilon_\varnothing:0\to\mathbb Z\) 仍是一个同态，只是不满射。这个约定可以使含空集的公式继续写在增广链层面，但必须记住它在 \(-1\) 次留下的同调。

\begin{proposition}[约化同调与普通同调的关系]\label{g2:prop:reduced-homology-degrees}
	设 \(X\) 为拓扑空间，所有同调均取整数系数。对于 \(n>0\)，有自然同构 \(\widetilde H_n(X)=H_n(X)\)。在零次及负次数，有自然正合列
	\[
		0\longrightarrow\widetilde H_0(X)
		\longrightarrow H_0(X)
		\xrightarrow{\bar\varepsilon_X}\mathbb Z
		\longrightarrow\widetilde H_{-1}(X)
		\longrightarrow0,
	\]
	且 \(\widetilde H_n(X)=0\)（\(n<-1\)）。若 \(X\ne\varnothing\)，则 \(\widetilde H_{-1}(X)=0\)；选定一点 \(x_0\in X\) 后，有分解
	\[
		H_0(X)\cong\widetilde H_0(X)\oplus\mathbb Z,
	\]
	其中第二项由 \([x_0]\) 生成。若 \(X=\varnothing\)，则只有 \(\widetilde H_{-1}(\varnothing)=\mathbb Z\) 非零。约化同调具有同伦不变性，因而非空可缩空间的全部约化同调均为零。
\end{proposition}
\begin{proof}
	正次数的圈和边界与普通链复形完全相同。在零次，
	\[
		\widetilde H_0(X)=\ker\varepsilon_X/\operatorname{im}\partial_1,
		\qquad
		H_0(X)=C_0(X)/\operatorname{im}\partial_1.
	\]
	由于 \(\varepsilon_X\partial_1=0\)，增广下降为 \(\bar\varepsilon_X:H_0(X)\to\mathbb Z\)，其核正是 \(\widetilde H_0(X)\)。在 \(-1\) 次，全部 \(\mathbb Z\) 都是圈，而边界是 \(\operatorname{im}\varepsilon_X\)，因此 \(\widetilde H_{-1}(X)\) 是该增广的余核。这证明了正合列。

	当 \(X\ne\varnothing\) 时，\(\varepsilon_X(x_0)=1\)，所以增广满射；同态 \(m\mapsto m[x_0]\) 给出其一个截面。具体分解是
	\[
		[c]\longmapsto
		\bigl([c]-\bar\varepsilon_X([c])[x_0],\,
		\bar\varepsilon_X([c])\bigr).
	\]
	这个分裂使用了 \(x_0\) 的选择，不能不加说明地把第二个直和项当作无选择的自然分裂。空集的结论则直接由其增广复形只有 \(-1\) 次的 \(\mathbb Z\) 非零得到。

	连续映射满足 \(\varepsilon_Yf_\#=\varepsilon_X\)，所以定义中的增广链映射确与边界交换。对两个同伦的映射，沿用棱柱算子，并取 \(P_{-1}=0\)。在零次，\(\partial P_0+P_{-1}\varepsilon_X=g_\#-f_\#\) 仍成立；在 \(-1\) 次，两条链映射都是 \(\operatorname{id}_{\mathbb Z}\)，差为零，其余次数与原证明相同。故它们在约化同调上诱导同一映射。单点空间的增广还使零次同调消失，结合 \cref{g2:cor:contractible-homology} 的同伦逆构造，即得非空可缩空间的结论。
\end{proof}

结合道路分支的计算，非空空间的 \(\widetilde H_0(X)\) 就是 \(\mathbb Z[\mathcal P(X)]\) 中总系数为零的部分。选定含 \(x_0\) 的分支 \(Q_0\) 后，各差 \([Q]-[Q_0]\)（\(Q\ne Q_0\)）构成它的一组基。若恰有 \(r\ge1\) 个道路分支，就得到 \(\widetilde H_0(X)\cong\mathbb Z^{r-1}\)。

\subsection{对及三元组的长正合列}
\label{g2:subsec:0202:04}
相对圈的边界落在 \(A\) 中，所以它本身给出 \(A\) 中的一个低一维圈。连接同态把这条边界作为输出。要保证输出与代表链的选择无关，并知道这些同态形成正合列，可以把前一节的短正合复形定理直接用于奇异链；在这个模型中，抽象的“提升后取边界”有具体的几何含义。

\begin{theorem}[空间对的长正合列]\label{g2:thm:pair-LES}
	设 \(X\) 为拓扑空间，\(A\subseteq X\) 为任意子空间，所有奇异链和同调均取整数系数。包含映射 \(i:A\hookrightarrow X\) 与商链映射 \(j:C_\bullet(X)\to C_\bullet(X,A)\) 给出自然长正合列
	\[
		\cdots\longrightarrow H_n(A)\xrightarrow{i_*}H_n(X)
		\xrightarrow{j_*}H_n(X,A)
		\xrightarrow{\delta_n}H_{n-1}(A)\longrightarrow\cdots,
	\]
	其低次末端为
	\[
		H_1(X,A)\xrightarrow{\delta_1}H_0(A)
		\xrightarrow{i_*}H_0(X)\xrightarrow{j_*}H_0(X,A)
		\longrightarrow0.
	\]
	若 \(c\in C_n(X)\) 满足 \(\partial c\in C_{n-1}(A)\)，则连接同态为
	\[
		\delta_n[\overline c]=[\partial c]\in H_{n-1}(A).
	\]
	对于任意连续空间对映射 \(f:(X,A)\to(Y,B)\)，它诱导两条长正合列之间的交换图，特别是
	\[
		(f|_A)_*\delta_n=\delta_n f_*.
	\]
	这里 \(f|_A\) 视为从 \(A\) 到 \(B\) 的映射；两侧的连接同态分别属于所对应的空间对。
\end{theorem}
\begin{proof}
	包含映射在奇异链上单射，商映射满射，且其核正是 \(C_n(A)\)。它们均与边界交换，故得到短正合复形
	\begin{equation}\label{g2:eq:pair-short-exact}
		0\longrightarrow C_\bullet(A)
		\xrightarrow{i_\#}C_\bullet(X)
		\xrightarrow{j}C_\bullet(X,A)\longrightarrow0.
	\end{equation}
	\cref{g2:thm:homology-LES} 适用于这个逐次短正合列，给出所述长正合列。它的连接构造先将商复形中的圈 \(\overline c\) 提升为 \(c\)，再取 \(\partial c\)。由于 \(j(\partial c)=\partial\overline c=0\)，这个边界在 \(C_{n-1}(A)\) 中；其在该子群中的原像就是它本身，所以连接公式没有额外负号。

	也可直接核验公式中的选择无关性。若换成同一相对同调类的代表 \(c'=c+\partial b+a\)，其中 \(a\in C_n(A)\)，则
	\[
		\partial c'=\partial c+\partial a,
	\]
	二者在 \(A\) 中只差一个边界。又因 \(\partial^2c=0\)，输出确是 \(A\) 中的同调类。

	给定空间对映射 \(f\)，它把 \eqref{g2:eq:pair-short-exact} 映到 \((Y,B)\) 的对应短正合复形。前一节的自然性结论因此适用。对连接项，具体等式是
	\[
		\delta_n[f_\#c]=[\partial f_\#c]
		=[f_\#\partial c]=(f|_A)_*[\partial c].
	\]
	它对应下面的交换方块，其余方块由包含或商映射的定义交换：
	\[
		\begin{tikzcd}[column sep=large,row sep=large]
			H_n(X,A) \arrow[r,"\delta_n"] \arrow[d,"f_*"']
				& H_{n-1}(A) \arrow[d,"(f|_A)_*"] \\
			H_n(Y,B) \arrow[r,"\delta_n"']
				& H_{n-1}(B).
		\end{tikzcd}
	\]
\end{proof}

例如，在相对项处的正合性有一个直接的解释。若相对圈 \(c\) 的边界在 \(A\) 中又是边界，即 \(\partial c=\partial a\)、\(a\in C_n(A)\)，则 \(c-a\) 是 \(X\) 中的绝对圈，并且与 \(c\) 给出相同的相对类。因此 \(\ker\delta_n=\operatorname{im}j_*\)。这说明连接同态恰好检测一个相对类能否改用绝对圈表示。

\begin{corollary}[零次相对同调]\label{g2:cor:relative-zero-homology}
	设 \(A\subseteq X\) 为拓扑空间的任意子空间，使用整数系数奇异同调。则 \(H_0(X,A)\) 是以那些不与 \(A\) 相交的道路分支为基的自由阿贝尔群。特别地，若 \(A\) 与 \(X\) 的每个道路分支都相交，则 \(H_0(X,A)=0\)。
\end{corollary}
\begin{proof}
	空间对长正合列的末端说明 \(H_0(X,A)\) 是 \(H_0(A)\to H_0(X)\) 的余核。根据 \cref{g2:prop:homology-point-components}，一个 \(A\) 中的点被送到它在 \(X\) 中所属的道路分支。因此像恰是与 \(A\) 相交的分支所张成的自由直和项，取商后留下所述其余基元。
\end{proof}

\begin{proposition}[约化形式的空间对长正合列]\label{g2:prop:reduced-pair-LES}
	设 \(A\subseteq X\) 为拓扑空间的非空子空间，\(\widetilde H_n\) 为增广次数 \(-1\) 取 \(\mathbb Z\) 的整数系数约化同调，\(H_n(X,A)\) 为普通相对奇异同调。则有自然长正合列
	\[
		\cdots\longrightarrow\widetilde H_n(A)
		\longrightarrow\widetilde H_n(X)
		\longrightarrow H_n(X,A)
		\xrightarrow{\widetilde\delta_n}\widetilde H_{n-1}(A)
		\longrightarrow\cdots,
	\]
	末端为
	\[
		H_1(X,A)\xrightarrow{\widetilde\delta_1}\widetilde H_0(A)
		\longrightarrow\widetilde H_0(X)
		\longrightarrow H_0(X,A)\longrightarrow0.
	\]
	连接同态仍由 \(\widetilde\delta_n[\overline c]=[\partial c]\) 给出；在 \(n=1\) 时，\(\partial c\) 作为总系数为零的零维圈表示 \(\widetilde H_0(A)\) 中的类。
\end{proposition}
\begin{proof}
	把 \eqref{g2:eq:pair-short-exact} 中前两个复形都增广，并让包含映射在 \(-1\) 次为 \(\operatorname{id}_{\mathbb Z}\)，就得到短正合复形
	\[
		0\longrightarrow\widetilde C_\bullet(A)
		\longrightarrow\widetilde C_\bullet(X)
		\longrightarrow C_\bullet(X,A)\longrightarrow0.
	\]
	非负次数正合性与原来相同；在 \(-1\) 次，它是
	\[
		0\longrightarrow\mathbb Z
		\xrightarrow{\operatorname{id}}\mathbb Z
		\longrightarrow0\longrightarrow0.
	\]
	因此两个增广复形之商恰是普通相对链复形，而不是一个额外带有 \(-1\) 次 \(\mathbb Z\) 的复形。这解释了长正合列中相对项仍写作 \(H_n(X,A)\) 的原因。

	应用 \cref{g2:thm:homology-LES} 得到结论及自然性。连接同态仍提升到 \(C_n(X)\) 后取边界；在 \(n=1\) 时，\(\varepsilon_A(\partial c)=0\)，所以输出确实落在约化零次同调中。非空性保证 \(\widetilde H_{-1}(A)=\widetilde H_{-1}(X)=0\)，于是末端具有陈述中的形式。
\end{proof}

上面这个增广短正合复形在 \(A=\varnothing\) 时仍然存在，但其长正合列不能在 \(H_0(X,\varnothing)\) 后直接截断：下一项是 \(\widetilde H_{-1}(\varnothing)=\mathbb Z\)。非空条件的作用就在这里。通常所谓非空空间对的“约化相对同调”，仍由同一个普通相对链复形计算，不必再定义一套不同的相对群。

\begin{corollary}[选定一点后的相对同调]\label{g2:cor:reduced-relative-point}
	设 \(X\) 为拓扑空间，\(x_0\in X\)，使用整数系数奇异同调。由链商映射诱导的映射给出自然同构
	\[
		\widetilde H_n(X)\xrightarrow{\ \cong\ }H_n(X,\{x_0\})
		\qquad(n\ge0).
	\]
	这里的自然性针对保持所选点的连续映射。
\end{corollary}
\begin{proof}
	单点空间的全部约化同调均为零。在 \cref{g2:prop:reduced-pair-LES} 中取 \(A=\{x_0\}\)，所述映射的左右两项都为零，故它是同构。自然性来自该长正合列对空间对映射的自然性。
\end{proof}

若指定两个嵌套的子空间，就能在已经忽略较小子空间的条件下，再忽略较大子空间。商群的第三同构定理在链复形中的应用，产生下一条序列。

\begin{theorem}[空间三元组的长正合列]\label{g2:thm:triple-LES}
	设 \(X\) 为拓扑空间，\(B\subseteq A\subseteq X\) 均赋予子空间拓扑，所有链和同调均取整数系数。自然映射给出长正合列
	\[
		\cdots\longrightarrow H_n(A,B)\longrightarrow H_n(X,B)
		\longrightarrow H_n(X,A)
		\xrightarrow{\delta_n}H_{n-1}(A,B)\longrightarrow\cdots,
	\]
	并终止于 \(H_0(X,A)\to0\)。若 \(c\in C_n(X)\) 满足 \(\partial c\in C_{n-1}(A)\)，则
	\[
		\delta_n[c+C_n(A)]
		=[\partial c+C_{n-1}(B)].
	\]
	这个序列对满足 \(f(A)\subseteq A'\)、\(f(B)\subseteq B'\) 的连续映射 \(f:X\to X'\) 自然，其中 \(B'\subseteq A'\subseteq X'\) 为另一空间三元组。
\end{theorem}
\begin{proof}
	在每个次数，自然映射组成短正合列
	\[
		0\longrightarrow C_n(A)/C_n(B)
		\longrightarrow C_n(X)/C_n(B)
		\longrightarrow C_n(X)/C_n(A)\longrightarrow0.
	\]
	第一条映射单射，因为一个 \(A\) 中的链在中间商中为零，当且仅当它已在 \(C_n(B)\) 中。最后一条映射满射，其核由 \(C_n(A)\) 中的链模去 \(C_n(B)\) 构成，正好是第一条映射的像。边界保持所有这些子群，故这是短正合复形
	\[
		0\longrightarrow C_\bullet(A,B)
		\longrightarrow C_\bullet(X,B)
		\longrightarrow C_\bullet(X,A)\longrightarrow0.
	\]
	应用 \cref{g2:thm:homology-LES} 得到长正合列。具体地，相对圈 \(c+C_n(A)\) 先提升为 \(c+C_n(B)\)，其边界是 \(\partial c+C_{n-1}(B)\)，位于 \(C_{n-1}(A,B)\) 中，这就给出陈述的公式。若换代表 \(c'=c+\partial b+a\)、\(a\in C_n(A)\)，输出只改变 \(C_\bullet(A,B)\) 中的边界 \(\partial a+C_{n-1}(B)\)，所以公式与代表无关。

	满足陈述中条件的 \(f\) 保持三条子链复形，并与边界及商映射交换，从而诱导这两条短正合复形之间的链映射。一般连接同态的自然性给出整条序列的自然性；也可在连接项上直接用 \(\partial f_\#c=f_\#\partial c\) 核验。
\end{proof}

取 \(B=\varnothing\)，三元组序列恢复空间对序列；取一点 \(B=\{x_0\}\subseteq A\)，再用 \cref{g2:cor:reduced-relative-point} 识别前两个相对项，就得到空间对序列的约化形式。

\begin{example}[区间的相对基本类]\label{g2:exm:interval-relative-class}
	设 \(I=[0,1]\)、\(A=\{0,1\}\)，使用整数系数奇异同调。则
	\[
		H_n(I,A)\cong
		\begin{cases}
			\mathbb Z,&n=1,\\
			0,&n\ne1.
		\end{cases}
	\]
	奇异一维单形 \(\sigma:\Delta^1\to I\)、\(\sigma(t_0,t_1)=t_1\)，给出 \(H_1(I,A)\) 的一个生成元，连接同态将它送到 \([1]-[0]\)。
\end{example}
\begin{proof}
	区间可缩，而 \(A\) 是两个点的不交并。由 \cref{g2:prop:homology-point-components,g2:cor:contractible-homology}，两者的正次数绝对同调都为零，\(H_0(A)\cong\mathbb Z^2\)、\(H_0(I)\cong\mathbb Z\)。空间对序列因此在低次化为
	\[
		0\longrightarrow H_1(I,A)\xrightarrow{\delta_1}
		\mathbb Z^2\xrightarrow{(a,b)\mapsto a+b}\mathbb Z
		\longrightarrow H_0(I,A)\longrightarrow0.
	\]
	求和同态满射，其核由 \((-1,1)\) 生成，故 \(\delta_1\) 将 \(H_1(I,A)\) 同构到该核。链 \(\sigma\) 满足 \(\partial\sigma=1-0\)，所以它正好给出这个生成元。更高次数的消失也由同一长正合列得到，负次数按定义为零。
\end{proof}

反射 \(r:I\to I\)、\(r(t)=1-t\) 保持端点集合，连接同态的自然性给出
\[
	\delta_1 r_*[\overline\sigma]=[0]-[1]
	=-\delta_1[\overline\sigma].
\]
因为 \(\delta_1\) 单射，所以 \(r_*\) 在 \(H_1(I,A)\cong\mathbb Z\) 上为乘以 \(-1\)。虽然 \(r\) 与恒等映射作为 \(I\to I\) 的映射可由 \(H(t,s)=(1-s)t+s(1-t)\) 连接，它们却不能通过始终保持端点集合的空间对映射同伦，否则相对同调上的诱导映射必须相同。相对同伦不变性中要求 \(H(A\times I)\subseteq A\)，正是为了保留这一信息。
% 纲要条目排为编号小节标题，并进入目录与书签；说明段仍保留于纲要注释中。

% !TeX root = ../../Book2.tex
% 纲要标题：### 2.3 细分与切除的证明
% 纲要位置：几何学计划纲要.md，第 1050 行。
% 正文状态：未撰写；本文件目前是框架。

% 纲要内容开始
% 
% * 重心细分及细分链算子
% * 细分与恒等映射的链同伦
% * 从开覆盖得到小链复形
% * 切除定理与闭包条件
% * Mayer–Vietoris 序列的导出
% 
% 纲要内容结束

\section{细分与切除的证明}
\label{g2:sec:0203}

待撰写

\subsection{重心细分及细分链算子}
\label{g2:subsec:0203:01}
待撰写

\subsection{细分与恒等映射的链同伦}
\label{g2:subsec:0203:02}
待撰写

\subsection{从开覆盖得到小链复形}
\label{g2:subsec:0203:03}
待撰写

\subsection{切除定理与闭包条件}
\label{g2:subsec:0203:04}
待撰写

\subsection{Mayer–Vietoris 序列的导出}
\label{g2:subsec:0203:05}
待撰写
% 纲要条目排为编号小节标题，并进入目录与书签；说明段仍保留于纲要注释中。

% !TeX root = ../../Book2.tex
% 纲要标题：### 2.4 局部同调和初步应用
% 纲要位置：几何学计划纲要.md，第 1058 行。
% 正文状态：未撰写；本文件目前是框架。

% 纲要内容开始
% 
% * 球面、球及穿孔欧氏空间
% * 流形点的局部同调
% * 维数不变性的同调证明
% * 约化同调与悬挂
% 
% ---
% 
% 纲要内容结束

\section{局部同调和初步应用}
\label{g2:sec:0204}

待撰写

\subsection{球面、球及穿孔欧氏空间}
\label{g2:subsec:0204:01}
待撰写

\subsection{流形点的局部同调}
\label{g2:subsec:0204:02}
待撰写

\subsection{维数不变性的同调证明}
\label{g2:subsec:0204:03}
待撰写

\subsection{约化同调与悬挂}
\label{g2:subsec:0204:04}
待撰写
% 纲要条目排为编号小节标题，并进入目录与书签；说明段仍保留于纲要注释中。


\begin{chaptersummary}
待撰写
\end{chaptersummary}

\BookExercises
待撰写

% !TeX root = ../../Book2.tex
% 纲要标题：## 第3章　胞腔同调、映射度与实际计算
% 纲要位置：几何学计划纲要.md，第 1067 行。
% 正文状态：未撰写；本文件目前是框架。

% 纲要内容开始
% 
% > 本章目标与证明范围：完整比较胞腔链与奇异链，使用整数矩阵保留挠信息。将附着映射、映射度和边界矩阵放在同一证明流程中。
% 
% 纲要内容结束

\chapter{胞腔同调、映射度与实际计算}
\label{g2:ch:03}
\BookChapterNumber{g2:ch:03}

\begin{chapterabstract}
待撰写
\end{chapterabstract}

\begin{learninggoals}
待撰写
\end{learninggoals}

待撰写

% !TeX root = ../../Book2.tex
% 纲要标题：### 3.1 骨架滤过给出的链复形
% 纲要位置：几何学计划纲要.md，第 1121 行。
% 正文状态：初稿；尚未正式审核。

% 纲要内容开始
% 
% * 骨架对的相对同调
% * 连接映射与胞腔微分
% * 胞腔微分平方为零
% * 与奇异同调的自然比较
% 
% 纲要内容结束

\section{骨架滤过给出的链复形}
\label{g2:sec:0301}

奇异链复形允许所有连续单形，因而适用于任意空间；但即使空间只有一个点，链群也在每个非负次数都有生成元。第一章给出的胞腔分解则只记录空间实际怎样由圆盘拼成。我们希望每个定向的 \(n\) 胞腔只贡献一个 \(n\) 维生成元，同时保留奇异同调中的全部信息。

\ProseParagraphBreak

困难在于胞腔的边界不必嵌入旧骨架，甚至可能多次经过同一个低维胞腔。因此不能把它简单当成若干几何面的集合。本节先用相对同调把每个圆盘的内部单独取出，再由连接同态恢复其附着边界。全节使用整数系数，约定 \(X^{-1}=\varnothing\)，普通同调在负次数为零。

\subsection{骨架对的相对同调}
\label{g2:subsec:0301:01}
相对同调 \(H_q(X^n,X^{n-1})\) 把旧骨架中的链视为零。每个新圆盘的边界因而消失，留下的应当是一个球面的顶维类。下面的切除论证说明，这个直观在胞腔数无限、附着映射不单射时仍然成立。

\begin{theorem}[骨架对的相对同调]\label{g2:thm:skeletal-relative-homology}
	设 \(X\) 为 CW 复形，\(E_n\) 为其 \(n\) 胞腔的指标集，\(\phi_\alpha:(D^n_\alpha,S^{n-1}_\alpha)\to(X^n,X^{n-1})\) 为特征映射。对所有整数 \(q\ge0\)，这些映射诱导同构
	\[
		\bigoplus_{\alpha\in E_n}H_q(D^n_\alpha,S^{n-1}_\alpha)
		\xrightarrow{\ \cong\ }H_q(X^n,X^{n-1}).
	\]
	当 \(n=0\) 时，将圆盘边界理解为空集。因此右边在 \(q\ne n\) 时为零，在 \(q=n\) 时是以定向 \(n\) 胞腔为基的自由阿贝尔群。同一结论适用于通过附着一族 \(n\) 胞腔得到的空间对 \((Y,A)\)。
\end{theorem}
\begin{proof}
	先处理 \(n\ge1\)。在附着空间 \(Y\) 中，取旧空间 \(A\) 连同各圆盘半径大于 \(1/2\) 的外环所形成的集合 \(U\)。在每个特征圆盘上的原像是 \(\{u:\|u\|>1/2\}\)，在 \(A\) 上的原像是全部 \(A\)，故商拓扑保证 \(U\) 是包含 \(A\) 的开集。外环中的 \(u=r\theta\) 沿
	\[
		r\theta\longmapsto\bigl((1-t)r+t\bigr)\theta
	\]
	推向圆盘边界，并在 \(A\) 上保持不动。这给出 \(U\) 到 \(A\) 的强形变收缩。各圆盘上的公式在边界一致；限制到开集的商映射仍是商映射，而与区间相乘仍为商映射（\cref{g2:lem:quotient-interval}），所以此同伦确实连续。

	由于 \(H_q(U,A)=0\)，三元组序列（\cref{g2:thm:triple-LES}）给出 \(H_q(Y,A)\cong H_q(Y,U)\)。附着构造中 \(A\) 闭，且 \(A\subset U=\operatorname{int}_Y U\)，所以切除定理的闭包条件成立。切去 \(A\) 后，各胞腔内部成为互不相交的开球，得到
	\[
		H_q(Y,U)\cong
		H_q\left(\coprod_{\alpha\in E_n}\mathring D^n_\alpha,
		\coprod_{\alpha\in E_n}\{u:1/2<\|u\|<1\}\right)
		\cong\bigoplus_{\alpha\in E_n}H_q(\mathring D^n,\{1/2<\|u\|<1\}).
	\]
	最后一个同构在链上已经成立：每个奇异单形的参数域连通，像落在一个不交分支中；每条链又只含有限个单形。因此这里是直和，不是允许无限支撑的直积。

	为核对同构恰由特征映射诱导，在每个圆盘中也取外环 \(V=\{u:1/2<\|u\|\le1\}\)。边界是 \(V\) 的强形变收缩核，所以 \(H_q(D^n,S^{n-1})\to H_q(D^n,V)\) 为同构；切除边界又将后者识别为 \(H_q(\mathring D^n,V\setminus S^{n-1})\)。这些映射与 \(\phi_\alpha\) 及前述切除图交换，于是其直和正是陈述中的映射。由\cref{g2:prop:ball-puncture-homology}，每个圆盘对只在次数 \(n\) 贡献一个 \(\mathbb Z\)。

	当 \(n=0\) 时，\(Y\) 是 \(A\) 与新增离散点的不交并，相对链只留下这些点的链复形。单点同调计算直接给出所述结论。取 \(Y=X^n\)、\(A=X^{n-1}\)，即得骨架对的情形。
\end{proof}

这里给胞腔选定向，代数上就是在 \(H_n(D^n,S^{n-1})\cong\mathbb Z\) 中选一个生成元。记它为 \([D^n,S^{n-1}]\)，并把其特征映射下的像记作 \(e^n_\alpha\)。这不会把一个可能自我识别的闭胞腔当成嵌入圆盘；生成元来自参数圆盘的相对同调。

\begin{proposition}[骨架的维数与稳定性]\label{g2:prop:skeletal-homology-stability}
	设 \(X\) 为 CW 复形，\(m,q\ge0\)。则 \(H_q(X^m)=0\) 当 \(q>m\)；包含映射 \(X^m\hookrightarrow X\) 在 \(q<m\) 时诱导同构，在 \(q=m\) 时诱导满射。不要求 \(X\) 有限维或局部有限。
\end{proposition}
\begin{proof}
	在 \((X^r,X^{r-1})\) 的长正合列中考察
	\[
		H_{q+1}(X^r,X^{r-1})\longrightarrow H_q(X^{r-1})
		\longrightarrow H_q(X^r)\longrightarrow H_q(X^r,X^{r-1}).
	\]
	两端只可能分别在 \(r=q+1\)、\(r=q\) 时非零。因此加入 \(r\) 胞腔只可能增加 \(r\) 次同调、消去 \(r-1\) 次同调，其他次数的包含映射均为同构。从离散的 \(X^0\) 开始归纳，便得维数消失性；固定 \(q\)，从 \(X^q\) 到 \(X^{q+1}\) 的映射满射，此后各步均为同构。

	还须说明无限个骨架的并不会出现额外的类或关系。任一奇异链仅含有限个单形，其支撑是紧集，因而由\cref{g2:thm:finite-subcomplex}落在有限子复形中，特别地落在某个有限维骨架中。于是 \(X\) 中的每个 \(q\) 圈可用刚才的有限步结论移到 \(X^q\)；如果 \(X^m\) 中的一个圈在 \(X\) 中成为边界，那么实现这个关系的有限链也落在某个骨架 \(X^r\) 中。当 \(q<m\) 时，\(H_q(X^m)\to H_q(X^r)\) 已是同构，故原圈在 \(X^m\) 中就是边界。这分别给出满射与单射。
\end{proof}

\subsection{连接映射与胞腔微分}
\label{g2:subsec:0301:02}
现在每个维数都有所需的自由群，但各维群之间还没有边界映射。一个 \(n\) 胞腔的边界首先给出 \(X^{n-1}\) 中的绝对同调类；再忽略更低的 \(X^{n-2}\)，才得到 \((n-1)\) 胞腔的线性组合。这两步不能颠倒，也不能把其中的相对化步骤省去。

\begin{definition}[胞腔链复形]\label{g2:def:cellular-chain-complex}
	设 \(X\) 为 CW 复形。定义其整数\term[bao qiang lian qun]{胞腔链群}[cellular chain group]为
	\[
		C_n^{\mathrm{cell}}(X)=H_n(X^n,X^{n-1})\quad(n\ge0),\qquad
		C_n^{\mathrm{cell}}(X)=0\quad(n<0).
	\]
	记空间对的连接同态为 \(\delta_n:H_n(X^n,X^{n-1})\to H_{n-1}(X^{n-1})\)，相对化同态为 \(j_r:H_r(X^r)\to H_r(X^r,X^{r-1})\)。定义\term[bao qiang wei fen]{胞腔微分}[cellular differential]
	\begin{equation}\label{g2:eq:cellular-differential}
		\mathrm{d}_n=j_{n-1}\delta_n:C_n^{\mathrm{cell}}(X)\longrightarrow C_{n-1}^{\mathrm{cell}}(X)\quad(n\ge1),\qquad \mathrm{d}_0=0.
	\end{equation}
	下文证明这些同态组成链复形，称为\term[bao qiang lian fu xing]{胞腔链复形}[cellular chain complex]，其同调记为 \(H_n^{\mathrm{cell}}(X)\)。
\end{definition}

\symidx{\(C_n^{\mathrm{cell}}(X)\)：胞腔链群}
\symidx{\(H_n^{\mathrm{cell}}(X)\)：胞腔同调群}
这一定义中的 \(C_n^{\mathrm{cell}}\) 本身已经是一个相对同调群；在它们之间再取核与像，才得到胞腔同调。相比之下，奇异链群 \(C_n(X)\) 是由全部连续单形生成的自由群。两种“链”处在不同构造层次，不能直接把一个胞腔符号当成某个唯一的奇异单形。

\ProseParagraphBreak

若 \(c\) 是代表某个胞腔链的奇异相对圈，则 \(\partial c\) 落在 \(X^{n-1}\) 中。\cref{g2:thm:pair-LES}约定连接同态为 \([c]\mapsto[\partial c]\)，没有另加负号；再模去 \(X^{n-2}\) 中的链，就得到 \(\mathrm{d}_n[c]\)。三元组 \((X^n,X^{n-1},X^{n-2})\) 的连接同态正是这个复合。

\begin{example}[有向边的边界]\label{g2:exm:cellular-edge-boundary}
	设一条一胞腔由特征路径 \(\phi:[0,1]\to X^1\) 给出，起点与终点分别为顶点 \(v_-\)、\(v_+\)，以参数从 \(0\) 到 \(1\) 为正方向。由\cref{g2:exm:interval-relative-class}，连接同态给出
	\[
		\mathrm{d}_1 e=[v_+]-[v_-].
	\]
	若两端附着到同一个顶点，结果为零，而不是两倍这个顶点。取两个顶点 \(v,w\)，用三条均从 \(v\) 指向 \(w\) 的边 \(a,b,c\) 相连，则
	\[
		\mathrm{d}_1(xa+yb+zc)=(x+y+z)(w-v).
	\]
	所以一维圈群由 \(b-a,c-a\) 自由生成。图没有二维胞腔，这些圈不会成为胞腔边界；它的胞腔同调为 \(H_1^{\mathrm{cell}}\cong\mathbb Z^2\)、\(H_0^{\mathrm{cell}}\cong\mathbb Z\)。下一小节后的比较定理将说明，这也就是奇异同调。
\end{example}

\subsection{胞腔微分平方为零}
\label{g2:subsec:0301:03}
\begin{proposition}[胞腔微分的复合]\label{g2:prop:cellular-d-square}
	对任意 CW 复形 \(X\) 及整数 \(n\)，胞腔微分满足 \(\mathrm{d}_{n-1}\mathrm{d}_n=0\)。若 \(f:X\to Y\) 是胞腔映射，则空间对映射 \((X^n,X^{n-1})\to(Y^n,Y^{n-1})\) 诱导的同态 \(f_n^{\mathrm{cell}}\) 组成链映射。
\end{proposition}
\begin{proof}
	对于 \(n\ge2\)，由定义有
	\[
		\mathrm{d}_{n-1}\mathrm{d}_n
		=j_{n-2}\delta_{n-1}j_{n-1}\delta_n.
	\]
	空间对 \((X^{n-1},X^{n-2})\) 的长正合列中，\(j_{n-1}\) 与 \(\delta_{n-1}\) 是相邻的两条映射，故 \(\delta_{n-1}j_{n-1}=0\)。在代表链上，这也表示绝对圈已经没有边界。\(n\le1\) 时结论由 \(\mathrm{d}_0=0\) 及负次数链群为零得到。

	胞腔映射保持每层骨架；空间对长正合列的自然性给出
	\[
		(f|_{X^{n-1}})_*\delta_n^X=\delta_n^Y f_n^{\mathrm{cell}},\qquad
		f_{n-1}^{\mathrm{cell}}j_{n-1}^X=j_{n-1}^Y(f|_{X^{n-1}})_*.
	\]
	复合两式，便得 \(\mathrm{d}_n^Yf_n^{\mathrm{cell}}=f_{n-1}^{\mathrm{cell}}\mathrm{d}_n^X\)，这就是链映射条件。
\end{proof}

\(\mathrm{d}^2=0\) 因而不是对附着图逐个猜测所得的巧合，而是长正合列的结构保证。实际计算出整数边界矩阵后，矩阵乘积为零仍值得核对：理论保证正确输入满足它，计算中漏掉一次经过或弄反一个方向，却会破坏这个条件。

\ProseParagraphBreak

反转一个胞腔的定向只把相应生成元乘以 \(-1\)。如果在第 \(n\) 维用对角矩阵 \(P_n\) 记录这些符号变化，边界矩阵按 \(A_n\mapsto P_{n-1}^{-1}A_nP_n\) 改变，故新旧链复形同构。胞腔的定向是写矩阵的选择，同调群并不依赖这一选择。

\subsection{与奇异同调的自然比较}
\label{g2:subsec:0301:04}
关键问题现在变成：胞腔圈对应哪些奇异圈，胞腔边界对应哪些奇异边界？先看 \(n\) 骨架上的绝对同调，再只加入下一层胞腔，就足以回答这两个问题。

\begin{theorem}[胞腔同调与奇异同调]\label{g2:thm:cellular-singular-comparison}
	设 \(X\) 为任意 CW 复形，使用整数系数。对每个 \(n\ge0\)，存在同构
	\[
		\Theta_{X,n}:H_n^{\mathrm{cell}}(X)\xrightarrow{\ \cong\ }H_n(X),
	\]
	并且对胞腔映射自然：若 \(f:X\to Y\) 为胞腔映射，则
	\[
		\Theta_{Y,n}H_n(f^{\mathrm{cell}})=f_*\Theta_{X,n}.
	\]
	若给任意连续映射取胞腔逼近，这个公式也确定其在胞腔同调上的同态，且与逼近的选择无关。
\end{theorem}
\begin{proof}
	当 \(n\ge1\) 时，\(H_n(X^{n-1})=H_{n-1}(X^{n-2})=0\)，这里 \(n=1\) 的第二项为 \(H_0(\varnothing)=0\)。两个空间对的长正合列于是说明
	\[
		j_n:H_n(X^n)\hookrightarrow C_n^{\mathrm{cell}}(X),\qquad
		j_{n-1}:H_{n-1}(X^{n-1})\hookrightarrow C_{n-1}^{\mathrm{cell}}(X)
	\]
	都是单射。因 \(\mathrm{d}_n=j_{n-1}\delta_n\)，第二条单射保证 \(\ker\mathrm{d}_n=\ker\delta_n\)；第一条正合列再给出
	\begin{equation}\label{g2:eq:cellular-cycles-skeleton}
		j_n:H_n(X^n)\xrightarrow{\ \cong\ }\ker\mathrm{d}_n.
	\end{equation}
	这表明胞腔圈恰是 \(n\) 骨架上的奇异同调类。并不是每个胞腔链都能提升为绝对圈：能否提升，正由 \(\delta_n\) 检测。

	接着考察加入 \((n+1)\) 胞腔后发生的变化。相对项 \(H_n(X^{n+1},X^n)\) 为零，所以空间对序列给出
	\[
		C_{n+1}^{\mathrm{cell}}(X)\xrightarrow{\delta_{n+1}}H_n(X^n)
		\longrightarrow H_n(X^{n+1})\longrightarrow0.
	\]
	由\cref{g2:prop:skeletal-homology-stability}，\(H_n(X^{n+1})\to H_n(X)\) 是同构。因此
	\[
		H_n(X)\cong H_n(X^n)/\operatorname{im}\delta_{n+1}.
	\]
	在 \(j_n\) 下，分母的像恰为 \(\operatorname{im}(j_n\delta_{n+1})=\operatorname{im}\mathrm{d}_{n+1}\)，分子由 \eqref{g2:eq:cellular-cycles-skeleton} 恰对应 \(\ker\mathrm{d}_n\)。取商即得所需同构。\(n=0\) 时，\(j_0:H_0(X^0)\to H_0(X^0,\varnothing)\) 就是恒等识别，\(\ker\mathrm{d}_0=C_0^{\mathrm{cell}}(X)\)；同一段关于 \((X^1,X^0)\) 的正合列和稳定性证明仍然适用。

	为看清同构的选择无关性，给定胞腔圈 \(z\)，先取由单射 \(j_n\) 唯一确定的类 \(\xi\in H_n(X^n)\)，满足 \(j_n\xi=z\)，再将 \(\xi\) 映入 \(H_n(X)\)。若 \(z\) 改变一个胞腔边界，\(\xi\) 恰改变一个 \(\operatorname{im}\delta_{n+1}\) 中的类，它在 \(X\) 中为零。这就是 \(\Theta_{X,n}\) 的具体构造。

	对胞腔映射，包含、相对化与连接映射全部自然，故刚才的核、像和取商识别与 \(f\) 交换。可将结果画成
	\[
		\begin{tikzcd}[column sep=large,row sep=large]
			H_n^{\mathrm{cell}}(X) \arrow[r,"\Theta_{X,n}","\cong"'] \arrow[d,"H_n(f^{\mathrm{cell}})"']
				& H_n(X) \arrow[d,"f_*"] \\
			H_n^{\mathrm{cell}}(Y) \arrow[r,"\Theta_{Y,n}","\cong"']
				& H_n(Y).
		\end{tikzcd}
	\]
	最后，由\cref{g2:thm:cellular-approximation}，任意连续映射 \(f\) 都有同伦的胞腔代表 \(f'\)。奇异同调的同伦不变性给出 \(f'_*=f_*\)，所以对应的胞腔同态必为 \(\Theta_{Y,n}^{-1}f_*\Theta_{X,n}\)，与 \(f'\) 的选择无关。其恒等与复合性质也由这个公式得到。
\end{proof}

这个比较是同调群上的自然比较。它没有声称每个胞腔天然指定一条唯一的奇异链，也没有声称同伦映射在胞腔链群上总是给出相同矩阵；稳定的是取同调后的映射。不同 CW 分解也通过同一个奇异同调得到同构，这为本章最后的计算核对提供依据。上述构造可与\cite[第2.2节，Cellular Homology]{Hatcher2002}对照。

\begin{theorem}[CW 对的胞腔同调]\label{g2:thm:relative-cellular-comparison}
	设 \(X\) 为 CW 复形，\(A\subset X\) 为子复形。包含诱导胞腔子链复形 \(C_\bullet^{\mathrm{cell}}(A)\subset C_\bullet^{\mathrm{cell}}(X)\)。令
	\[
		C_\bullet^{\mathrm{cell}}(X,A)=C_\bullet^{\mathrm{cell}}(X)/C_\bullet^{\mathrm{cell}}(A).
	\]
	其第 \(n\) 群以 \(X\) 中不属于 \(A\) 的定向 \(n\) 胞腔为基，边界由原边界删除属于 \(A\) 的项得到，并有对胞腔空间对映射自然的同构
	\[
		H_n(C_\bullet^{\mathrm{cell}}(X,A))\cong H_n(X,A;\mathbb Z)\qquad(n\ge0).
	\]
\end{theorem}
\begin{proof}
	包含映射在各相对骨架同调上把一个胞腔生成元映到同名生成元，因此逐次单射；它与胞腔微分交换，故确实给出子链复形。商群的基与边界公式随之成立。

	为把这个商复形与奇异相对同调比较，令 \(F_{-1}=A\)、\(F_r=A\cup X^r\)。\cref{g2:prop:cw-pair-relative}说明 \(F_r\) 由 \(F_{r-1}\) 附着那些不在 \(A\) 中的 \(r\) 胞腔得到，拓扑也与 \(X\) 中的子空间拓扑一致。骨架相对同调定理的附着版本给出
	\[
		H_q(F_r,F_{r-1})=0\quad(q\ne r),\qquad
		H_r(F_r,F_{r-1})\cong C_r^{\mathrm{cell}}(X,A).
	\]
	由三元组 \((F_r,F_{r-1},A)\) 的正合列，从 \(H_q(F_{-1},A)=0\) 开始归纳，得到 \(H_q(F_r,A)=0\) 当 \(q>r\)；固定 \(q\)，包含诱导的群在 \(r\ge q+1\) 后稳定。与绝对情形相同，奇异相对圈及实现其边界关系的链都有紧支撑，所以稳定值正是 \(H_q(X,A)\)。

	用三元组序列替换上一证明中的空间对序列，便有单射
	\[
		H_n(F_n,A)\hookrightarrow H_n(F_n,F_{n-1})=C_n^{\mathrm{cell}}(X,A),
	\]
	其像为相对胞腔圈群；下一层的连接映射给出该单射下的全部相对胞腔边界，并且
	\[
		H_n(X,A)\cong H_n(F_n,A)/\operatorname{im}\bigl(H_{n+1}(F_{n+1},F_n)\to H_n(F_n,A)\bigr).
	\]
	这里的微分与商复形的微分相同，因为两者均由代表链取边界，再忽略 \(A\) 及低一层骨架得到。取商于是给出陈述的同构；自然性仍来自三元组序列的自然性。整个论证在 \(n=0\) 也适用，因为 \(H_0(F_{-1},A)=0\)。
\end{proof}

例如，若 \(Y=X\cup_\alpha D^m\)，其中 \(X\) 为 CW 复形且 \(\alpha:S^{m-1}\to X^{m-1}\)，那么相对胞腔链只有次数 \(m\) 的一个 \(\mathbb Z\)，故 \(H_m(Y,X)=\mathbb Z\)，其余次数为零。这个相对群只记录新增胞腔；它的边界究竟消去旧空间的哪个类，仍须回到连接同态及附着映射去计算。
% 纲要条目排为编号小节标题，并进入目录与书签；说明段仍保留于纲要注释中。

% !TeX root = ../../Book2.tex
% 纲要标题：### 3.2 附着映射和边界系数
% 纲要位置：几何学计划纲要.md，第 1128 行。
% 正文状态：初稿；尚未正式审核。

% 纲要内容开始
% 
% * 球面映射度的同调定义
% * 局部度、复合与悬挂
% * 边界系数的度公式
% * 射影空间附着映射的核验
% 
% 纲要内容结束

\section{附着映射和边界系数}
\label{g2:sec:0302}

上一节把每个定向胞腔变成链群的一名基元，尚须回答：一个 \(n\) 胞腔的边界，在某个 \(n-1\) 胞腔上究竟出现多少次？这里的“多少次”必须带符号；两个覆盖片可能相加，也可能抵消。先研究球面映射对唯一的顶维同调生成元做了什么，再将附着映射压到一颗目标球面，就能逐项读出边界矩阵。

\subsection{球面映射度的同调定义}
\label{g2:subsec:0302:01}

由\cref{g2:thm:sphere-homology}，\(n\ge1\) 时 \(H_n(S^n;\mathbb Z)\cong\mathbb Z\)。选择一个生成元 \([S^n]\) 相当于在同调意义下给球面选定向，另一个选择是 \(-[S^n]\)。任何这个群的自同态都只能是乘一个整数，这正是连续映射可以留下的顶维信息。

\begin{definition}[球面映射度]\label{g2:def:sphere-degree}
	设 \(n\ge1\)，\(f:S^n\rightarrow S^n\) 连续。唯一满足
	\[
		f_*[S^n]=(\deg f)[S^n]
	\]
	的整数称为 \(f\) 的\term[yingshe du]{映射度}[degree]。自映射的源和目标使用同一个生成元，故同时更换生成元的符号不改变 \(\deg f\)。若源、目标是两颗分别定向的 \(n\) 维球面，则以各自的正生成元定义映射度；只反转其中一端的定向会改变映射度的符号。\symidx{\(\deg f\)：定向球面之间的映射度}
\end{definition}

当 \(n=0\) 时，可以在 \(\widetilde H_0(S^0)=\mathbb Z([+]-[-])\) 上作同样定义。恒等映射、交换两点的映射与常值映射的度分别为 \(1,-1,0\)。普通 \(H_0(S^0)\) 的秩是二，不能直接把它的任意自同态称为“乘一个整数”；这里使用约化同调是实质性的区别。

\begin{proposition}[映射度的基本性质]\label{g2:prop:degree-basic}
	设 \(n\ge1\)，\(f,g:S^n\rightarrow S^n\) 连续。恒等映射的度为 \(1\)，同伦的映射具有相同的度，并且
	\[
		\deg(g\circ f)=(\deg g)(\deg f).
	\]
	若 \(f\) 不是满射，则 \(\deg f=0\)；若 \(f\) 是同伦等价，则 \(\deg f\in\{1,-1\}\)。这些结论对用约化零次同调定义的 \(S^0\) 映射度也成立。
\end{proposition}
\begin{proof}
	恒等映射与复合公式直接来自同调的函子性。例如
	\[
		(g\circ f)_*[S^n]=g_*\bigl((\deg f)[S^n]\bigr)
		=(\deg f)(\deg g)[S^n].
	\]
	同伦不变性来自\cref{g2:thm:singular-homotopy-invariance}。若 \(y\notin f(S^n)\)，则 \(f\) 分解经过 \(S^n\setminus\{y\}\)，后者同胚于 \(\mathbb R^n\)，因而可缩，其 \(n\) 次同调为零。故 \(f_*=0\)。若 \(h\) 是 \(f\) 的同伦逆，则 \((\deg h)(\deg f)=\deg(hf)=1\)，两个整数只能同为 \(1\) 或同为 \(-1\)。\(S^0\) 的结论可直接检查三类映射，也可使用相同的约化同调论证。
\end{proof}

特别地，非零映射度强迫满射，而零伦映射的度为零。这两条判断的反向含义不能混为一谈：一个映射可以覆盖整个球面，却让不同部分的同调贡献互相抵消。下面的局部公式将把这种抵消写成真正的整数加法。

\subsection{局部度、复合与悬挂}
\label{g2:subsec:0302:02}

先解决最基本的正负号。将球面上下两个半球交换，赤道上的点保持不动；由于同调类的共同边界在两个半球中以相反符号出现，这个交换应当反转顶维生成元。

\begin{proposition}[反射与对径映射]\label{g2:prop:reflection-antipodal-degree}
	设 \(n\ge0\)。球面 \(S^n\subset\mathbb R^{n+1}\) 关于任一过原点超平面的反射，映射度为 \(-1\)。对径映射 \(A_n(x)=-x\) 的映射度为
	\[
		\deg A_n=(-1)^{n+1}.
	\]
	\(n=0\) 时度取约化零次同调的定义。
\end{proposition}
\begin{proof}
	\(n=0\) 时反射交换两个点，将 \([+]-[-]\) 送到其负元。设 \(n\ge1\)，先取反射 \(r(x_1,\ldots,x_n,t)=(x_1,\ldots,x_n,-t)\)。沿用\cref{g2:thm:sphere-homology}的开覆盖 \(U=S^n\setminus\{S\}\)、\(V=S^n\setminus\{N\}\)，连接同构为
	\[
		\delta:\widetilde H_n(S^n)\xrightarrow{\cong}
		\widetilde H_{n-1}(U\cap V)\cong\widetilde H_{n-1}(S^{n-1}).
	\]
	把代表圈细分后写为 \(z=a+b\)，其中 \(a\) 位于 \(U\)、\(b\) 位于 \(V\)。按\cref{g2:eq:mayer-vietoris-connecting}的约定，\(\delta[z]=[\partial a]\)。反射交换 \(U,V\)，所以 \(r_\#z\) 的 \(U\) 部分是 \(r_\#b\)，从而
	\[
		\delta r_*[z]=[\partial r_\#b]
		=-r_*[\partial a].
	\]
	赤道投影 \(U\cap V\rightarrow S^{n-1}\) 与 \(r\) 复合后不变，因此 \(r\) 在交集的同调上诱导恒等。这就给出 \(\delta r_*=-\delta\)，而 \(\delta\) 是同构，故 \(\deg r=-1\)。任一过原点超平面的反射都是 \(h r h^{-1}\)，其中 \(h\) 为正交线性映射限制在球面上的同胚；复合公式使 \(\deg h\) 与 \(\deg h^{-1}\) 抵消，仍得到 \(-1\)。最后，\(A_n\) 是分别改变 \(n+1\) 个坐标符号的反射之复合，所以其度为 \((-1)^{n+1}\)。
\end{proof}

要把整体的度拆成各原像点的贡献，先把球面的生成元限制到一点附近。对 \(x\in S^n\)、\(n\ge1\)，由于 \(S^n\setminus\{x\}\) 非空可缩，约化对的长正合列给出同构
\[
	H_n(S^n)\xrightarrow{\cong}H_n(S^n,S^n\setminus\{x\}).
\]
记 \([S^n]_x\) 为 \([S^n]\) 的像。\cref{g2:prop:local-homology-neighborhood}又将它识别为任意开邻域对 \((U,U\setminus\{x\})\) 的生成元。这样定义的局部正号来自同一个整体选择，各点之间不必另外任意指定符号。

\begin{definition}[局部度]\label{g2:def:local-degree}
	设 \(n\ge1\)，\(f:S^n\rightarrow S^n\) 连续，\(f(x)=y\)，且 \(x\) 是 \(f^{-1}(y)\) 的孤立点。取开邻域 \(U\ni x\)，使 \(U\cap f^{-1}(y)=\{x\}\)。空间对映射
	\[
		f:(U,U\setminus\{x\})\longrightarrow(S^n,S^n\setminus\{y\})
	\]
	在 \(n\) 次同调上将局部正生成元送到 \(\deg_x(f)\,[S^n]_y\)；整数 \(\deg_x(f)\) 称为 \(f\) 在 \(x\) 处的\term[ju bu du]{局部度}[local degree]。缩小 \(U\) 所诱导的局部同调同构与 \(f_*\) 交换，所以该整数与邻域选择无关。\symidx{\(\deg_x(f)\)：映射在孤立原像点处的局部度}
\end{definition}

局部度的定义只要求连续性和孤立原像，并未要求可微或局部可逆。若 \(f\) 将 \(x\) 的某个邻域同胚地映到 \(y\) 的邻域，局部同调同态是同构，故局部度只能为 \(\pm1\)；这时称局部正生成元被保持或反转，分别就是符号 \(+1\) 与 \(-1\) 的含义。

\begin{theorem}[局部度求和]\label{g2:thm:local-degree-sum}
	设 \(n\ge1\)，\(f:S^n\rightarrow S^n\) 连续。如果某点 \(y\) 的原像是有限集 \(F=\{x_1,\ldots,x_r\}\)，则
	\begin{equation}\label{g2:eq:local-degree-sum}
		\deg f=\sum_{i=1}^r\deg_{x_i}(f).
	\end{equation}
	允许 \(F\) 为空，此时右端为空和 \(0\)。特别地，同胚在每一点的局部度都等于其整体映射度。
\end{theorem}
\begin{proof}
	空原像的情形已经包含在\cref{g2:prop:degree-basic}中。设 \(r>0\)，取两两不交的开邻域 \(U_i\ni x_i\)。记 \(U=\bigcup_iU_i\)。在空间对 \((S^n,S^n\setminus F)\) 中切去闭集 \(S^n\setminus U\)；它包含于开集 \(S^n\setminus F\)，满足\cref{g2:thm:singular-excision}的闭包条件。由切除及不交并的链群分解，得到
	\[
		H_n(S^n,S^n\setminus F)
		\cong\bigoplus_{i=1}^r H_n(U_i,U_i\setminus\{x_i\}).
	\]
	限制到单个点 \(x_i\) 对应于投影到第 \(i\) 项：其余 \(U_j\) 在该相对群中已被全部除去。于是整体生成元 \([S^n]\) 在这个直和中的像正是各局部正生成元组成的向量 \((1,\ldots,1)\)，而不是任意选出的 \(r\) 个符号。

	由空间对映射的自然性，以下方块交换，其中上方箭头包含上述切除同构的逆，右方箭头是各 \(f|_{U_i}\) 所诱导同态之和：
	\[
		\begin{tikzcd}[column sep=large,row sep=large]
			H_n(S^n) \arrow[r] \arrow[d,"f_*"'] & \displaystyle\bigoplus_i H_n(U_i,U_i\setminus\{x_i\}) \arrow[d,"\sum_i(f|_{U_i})_*"] \\
			H_n(S^n) \arrow[r,"\cong"'] & H_n(S^n,S^n\setminus\{y\}).
		\end{tikzcd}
	\]
	沿右边走，\((1,\ldots,1)\) 变成 \(\sum_i\deg_{x_i}(f)\)；沿左边走，生成元变成 \(\deg f\) 倍。因此交换性给出公式。同胚的每个目标点只有一个原像，遂得到最后一项。
\end{proof}

在连续映射中，直接按原像点的个数计数一般会错。局部同胚也可能反转符号；若原像虽孤立但映射并不局部可逆，单点的局部度甚至可以是任意整数。公式中每一项都要由相应的局部同调同态决定。

\begin{proposition}[局部复合与悬挂]\label{g2:prop:degree-composition-suspension}
	设 \(f,g:S^n\rightarrow S^n\) 连续，\(n\ge1\)，\(f(x)=y\)、\(g(y)=z\)。若 \(x\) 是 \(f^{-1}(y)\) 的孤立点，\(y\) 是 \(g^{-1}(z)\) 的孤立点，则 \(x\) 是 \((g\circ f)^{-1}(z)\) 的孤立点，且
	\[
		\deg_x(g\circ f)=\deg_x(f)\deg_y(g).
	\]
	此外，对任意 \(n\ge0\) 的球面映射 \(f:S^n\rightarrow S^n\)，其悬挂 \(Sf:S^{n+1}\rightarrow S^{n+1}\) 满足 \(\deg(Sf)=\deg f\)；\(n=0\) 的右端取约化同调定义。
\end{proposition}
\begin{proof}
	先取开邻域 \(V\ni y\)，使 \(g^{-1}(z)\cap V=\{y\}\)，再缩小 \(U\ni x\)，使 \(f(U)\subset V\) 且 \(f^{-1}(y)\cap U=\{x\}\)。于是 \((gf)^{-1}(z)\cap U=\{x\}\)，并有空间对映射的复合
	\[
		(U,U\setminus\{x\})\xrightarrow{f}
		(V,V\setminus\{y\})\xrightarrow{g}
		(S^n,S^n\setminus\{z\}).
	\]
	局部同调的函子性使两个整数相乘，得到第一式。

	对于悬挂，\cref{g2:thm:suspension-homology}提供自然同构
	\[
		\delta:\widetilde H_{n+1}(SS^n)\xrightarrow{\cong}\widetilde H_n(S^n),
		\qquad \delta(Sf)_*=f_*\delta.
	\]
	通过该节给出的同胚 \(SS^n\cong S^{n+1}\) 识别源、目标，取两个球面上经 \(\delta\) 对应的生成元。交换等式说明两端的自同态乘同一个整数，因此悬挂保持映射度。两个生成元同时反号不改变结论。
\end{proof}

\begin{example}[任意整数的实现]\label{g2:exm:all-integer-degrees}
	把 \(S^1\) 看成复平面内的单位圆。对任意整数 \(k\)，映射 \(p_k(z)=z^k\) 的度为 \(k\)。当 \(k=0\) 时它为常值。设 \(k>0\)，任意 \(y\in S^1\) 有 \(k\) 个原像；在各原像附近的角坐标中，映射是 \(t\mapsto kt\)，因此局部度为 \(+1\)。具体地，旋转可把源点与目标点移到 \(1\)，旋转同伦于恒等，局部度为 \(1\)；而正伸缩 \(t\mapsto kt\) 与恒等之间的同伦 \(t\mapsto((1-s)k+s)t\) 始终把非零点送到非零点，在足够小的坐标区间内给出空间对同伦。因此局部同调同态不变，局部度求和给出 \(k\)。若 \(k<0\)，则 \(p_k\) 是 \(p_{|k|}\) 与反射 \(z\mapsto\bar z\) 的复合，故其度为 \(-|k|\)。对 \(p_k\) 作 \(n-1\) 次悬挂，就在每个 \(n\ge1\) 的维数得到度为 \(k\) 的 \(S^n\) 自映射。
\end{example}

\subsection{边界系数的度公式}
\label{g2:subsec:0302:03}

固定一个 CW 复形 \(X\)，并给各胞腔选定向。记 \(e^n_\alpha\) 同时表示开胞腔及其在 \(C_n^{\mathrm{cell}}(X)=H_n(X^n,X^{n-1})\) 中的正基元，实际含义由所在位置区分。对一个 \(n-1\) 胞腔 \(e^{n-1}_\beta\)，把 \(X^{n-1}\) 中该开胞腔以外的全部点压成一点，得到商映射
\[
	q_\beta:X^{n-1}\longrightarrow
	X^{n-1}/(X^{n-1}\setminus e^{n-1}_\beta)
	\cong D^{n-1}_\beta/\partial D^{n-1}_\beta\cong S^{n-1}_\beta.
\]
这个商只保留所问胞腔的内部，故应当对应于从链群中提取它的系数。即使 \(X\) 有无限多个胞腔，商映射也连续：它在各特征圆盘上分别是通常的圆盘商映射或常值映射，CW 弱拓扑保证这些限制给出连续映射。

\ProseParagraphBreak

给商球面定向以前，还须核对圆盘的相对生成元确实映成球面的生成元。设 \(k\ge1\)，\(q:D^k\to D^k/S^{k-1}=Q\)，\(*\) 为压缩边界所得的点，取圆盘中心 \(x\)，记 \(\bar x=q(x)\)。径向形变收缩说明 \(S^{k-1}\hookrightarrow D^k\setminus\{x\}\) 是同伦等价；比较约化对的长正合列，得到下面左侧竖箭头为同构。\(Q\) 同胚于 \(S^k\)，且 \(Q\setminus\{\bar x\}\) 可缩，所以右侧竖箭头也为同构。
\[
	\begin{tikzcd}[column sep=large,row sep=large]
		H_k(D^k,S^{k-1}) \arrow[r,"q_*"] \arrow[d,"\cong"'] & H_k(Q,*) \arrow[d,"\cong"] \\
		H_k(D^k,D^k\setminus\{x\}) \arrow[r,"q_*"'] & H_k(Q,Q\setminus\{\bar x\}).
	\end{tikzcd}
\]
下侧箭头同样是同构：在 \(x\) 的一个小开邻域上，\(q\) 是同胚，再用\cref{g2:prop:local-homology-neighborhood}即可。方块由空间对映射诱导而交换，因此上侧 \(q_*\) 是同构。这也覆盖 \(k=1\)，其左侧比较使用的是边界两点的约化零次同调。由此，将相对圆盘的正生成元的商像选为 \(H_k(Q,*)\) 的正生成元，再经 \(H_k(Q)\cong H_k(Q,*)\) 给商球面定向。上面的 \(S^{n-1}_\beta\) 就采用由 \(e^{n-1}_\beta\) 这样确定的定向。

\begin{theorem}[胞腔边界的度公式]\label{g2:thm:cellular-boundary-degree}
	设 \(X\) 为 CW 复形，各胞腔已经选定向，\(n\ge2\)。对 \(n\) 胞腔的特征映射 \(\Phi_\alpha:D^n_\alpha\rightarrow X^n\)，记附着映射为 \(a_\alpha=\Phi_\alpha|_{\partial D^n_\alpha}\)，并给 \(\partial D^n_\alpha\) 取由连接同态传来的边界定向。则胞腔微分为
	\begin{equation}\label{g2:eq:cellular-degree-formula}
		\mathrm{d}_n(e^n_\alpha)=\sum_\beta
		\deg(q_\beta\circ a_\alpha)\,e^{n-1}_\beta.
	\end{equation}
	和式只有有限个非零项。对 \(n=1\)，将定向区间的起点、终点像分别记为 \(v_-,v_+\)，则 \(\mathrm{d}_1(e^1_\alpha)=v_+-v_-\)；若两端映到同一顶点，结果为零。
\end{theorem}
\begin{proof}
	取与 \(e^n_\alpha\) 对应的相对球生成元 \(u_\alpha\in H_n(D^n_\alpha,\partial D^n_\alpha)\)，其连接像是边界球面的正生成元 \([\partial D^n_\alpha]\)。\cref{g2:thm:skeletal-relative-homology}将 \((\Phi_\alpha)_*u_\alpha\) 识别为 \(e^n_\alpha\)。连接同态的自然性以及\cref{g2:def:cellular-chain-complex}中的 \(\mathrm{d}_n=j\delta\) 给出
	\[
		\mathrm{d}_n(e^n_\alpha)
		=j\,(a_\alpha)_*[\partial D^n_\alpha]
		\in H_{n-1}(X^{n-1},X^{n-2}).
	\]
	这里 \(j\) 把绝对同调类送到相对同调类，故目前的公式已经说明边界完全由附着映射决定。

	映射 \(q_\beta\) 把 \(X^{n-2}\) 送到基点 \(*\)，因而诱导
	\[
		(q_\beta)_*:H_{n-1}(X^{n-1},X^{n-2})
		\longrightarrow H_{n-1}(S^{n-1}_\beta,*)\cong\mathbb Z.
	\]
	它把 \(e^{n-1}_\beta\) 送到 \(1\)，把其他胞腔基元送到 \(0\)：前者来自该特征圆盘上的边界压缩及我们选定的目标定向，后者来自常值映射。因此它正是第 \(\beta\) 个坐标投影。在上式两边作用这个投影，所得整数就是 \(q_\beta a_\alpha\) 在球面顶维同调上乘的整数，即其度。由于 \(n-1\ge1\)，这里 \(H_{n-1}(S^{n-1}_\beta)\rightarrow H_{n-1}(S^{n-1}_\beta,*)\) 是同构，不会引入额外的零次分支项。

	附着球面的像是紧集；由\cref{g2:thm:finite-subcomplex}，它包含在有限子复形内。不在该子复形中的 \(e^{n-1}_\beta\) 对应常值复合 \(q_\beta a_\alpha\)，故其系数为零。最后，\(n=1\) 时相对区间生成元可由正向道路代表，连接同态直接给出终点减起点；这正是所述公式，无须将任意顶点集误当作一个零维球面。
\end{proof}

度公式中的符号来自两项选择：源胞腔的定向规定附着球面的边界定向，目标胞腔的定向规定压缩球面的正生成元。反转源胞腔会改变这列边界系数的全部符号，反转目标胞腔则改变对应行的全部符号。链群基随之改变，同调群不变；计算中应保存这些基选择，不能在某个矩阵中单独改一个看起来不顺眼的负号。

\begin{proposition}[附着字的指数和]\label{g2:prop:word-boundary-degrees}
	设 \(X\) 有一个零胞腔、有限个定向一胞腔 \(a_1,\ldots,a_r\)，且某个二胞腔 \(e\) 的附着映射沿有限字 \(w\) 依次遍历这些圆周，字母 \(a_i^{-1}\) 表示反向遍历。若 \(\nu_i(w)\) 为 \(a_i\) 的出现次数减去 \(a_i^{-1}\) 的出现次数，则
	\[
		\mathrm{d}_2(e)=\sum_{i=1}^r\nu_i(w)a_i.
	\]
	只要一般附着映射与上述逐段遍历的映射同伦，同一公式仍成立。
\end{proposition}
\begin{proof}
	压掉除 \(a_i\) 以外的圆周，在剩余开边内部选一点 \(y\)。复合 \(q_i a\) 的原像每出现一个 \(a_i\) 或 \(a_i^{-1}\) 就有一点，其他字母的线段全部落在基点，不贡献原像。正向遍历的局部角坐标是正伸缩，局部度为 \(+1\)；反向遍历再复合一次反射，局部度为 \(-1\)。由局部度求和，复合的度就是 \(\nu_i(w)\)，再代入胞腔边界公式即可。同伦的附着映射经 \(q_i\) 复合后仍同伦，故各映射度不变。
\end{proof}

例如，一个二胞腔沿 \(aba^{-1}b^{-1}\) 附着时，两项指数和都是零，所以它的胞腔边界为零；沿 \(a^2b^2\) 附着时则得到 \(2a+2b\)。两幅粘合图都可能把同一条边用两次，区别在于两次所带的方向。这正是只数边数不能替代边界矩阵的原因。

\subsection{射影空间附着映射的核验}
\label{g2:subsec:0302:04}

\cref{g2:exm:projective-cw}给出了实射影空间每维一个胞腔的模型：在 \(\mathbb RP^n\) 中新添一个 \(n\) 胞腔，其附着映射是
\[
	\pi:S^{n-1}\longrightarrow\mathbb RP^{n-1},\qquad x\longmapsto[x].
\]
这确实是二重覆盖，但边界系数不能据此立刻写成 \(2\)：度公式要求先把低骨架 \(\mathbb RP^{n-2}\) 压掉，并比较两个覆盖片的符号。对径映射的奇偶性正是在这一步进入。

\begin{proposition}[实射影空间的胞腔微分]\label{g2:prop:rp-cellular-boundary}
	对 \(\mathbb RP^N\) 的标准 CW 结构，可以逐维选择胞腔正生成元 \(e^0,\ldots,e^N\)，使每个 \(1\le n\le N\) 满足
	\begin{equation}\label{g2:eq:rp-cellular-boundary}
		\mathrm{d}_n(e^n)=\bigl(1+(-1)^n\bigr)e^{n-1}
		=\begin{cases}2e^{n-1},&n\text{ 为偶数},\\0,&n\text{ 为奇数}.\end{cases}
	\end{equation}
	有限维数 \(N\) 的取值不影响已经选定的低维胞腔和微分。
\end{proposition}
\begin{proof}
	\(n=1\) 时两个端点都附着到唯一的零胞腔，故 \(\mathrm{d}_1=0\)。现在设 \(n\ge2\)，记
	\[
		F_n:S^{n-1}\xrightarrow{\pi}\mathbb RP^{n-1}
		\xrightarrow{q}\mathbb RP^{n-1}/\mathbb RP^{n-2}\cong S^{n-1}.
	\]
	源球面中的赤道 \(x_n=0\) 映入 \(\mathbb RP^{n-2}\)，两个开半球 \(x_n>0\) 与 \(x_n<0\) 则分别同胚地映到目标球面除基点以外的部分。因为每条不在低维射影子空间中的直线恰有两个单位代表，这两个半球上的映射没有其他原像片。

	从已选定的 \(e^{n-1}\) 出发，在北半球取一点 \(x\)。因为 \(F_n\) 在 \(x\) 附近是局部同胚，其局部度为 \(\pm1\)；选择 \(e^n\) 的定向，使这个局部度为 \(+1\)。反转 \(e^n\) 的定向即反转其边界定向，所以总能做到。此时 \(F_n^{-1}(F_n(x))=\{x,-x\}\)。对径映射 \(A_{n-1}\) 满足 \(F_nA_{n-1}=F_n\)，局部复合公式给出
	\[
		\deg_{-x}(F_n)
		=\deg_{-x}(F_nA_{n-1})
		=\deg_x(F_n)\deg_{-x}(A_{n-1})
		=(-1)^n.
	\]
	最后一步使用：对径映射是同胚，其局部度等于整体度，而\cref{g2:prop:reflection-antipodal-degree}在维数 \(n-1\) 给出的符号是 \((-1)^n\)。因此局部度求和给出 \(\deg F_n=1+(-1)^n\)，胞腔边界的度公式便得到结论。按维数递增作定向选择，不会改变此前已经选定的低维生成元，故所有等式可以同时成立。
\end{proof}

\(\mathbb RP^2\) 的附着圆周有两片同向贡献，所以 \(\mathrm{d}_2=2\)；\(\mathbb RP^3\) 新添胞腔的边界球面上两片贡献反向，所以 \(\mathrm{d}_3=0\)。“覆盖了两次”在相邻维数给出了完全不同的整数边界。下一节将把这些整数代回链复形，直接看见 \(\mathbb Z/2\) 挠项从何而来。映射度、局部度和胞腔边界公式的相关论述可参阅\cite[第2.2节，Degree 与 Cellular Homology]{Hatcher2002}。
% 纲要条目排为编号小节标题，并进入目录与书签；说明段仍保留于纲要注释中。

% !TeX root = ../../Book2.tex
% 纲要标题：### 3.3 典型空间的整数同调
% 纲要位置：几何学计划纲要.md，第 1135 行。
% 正文状态：初稿；尚未正式审核。

% 纲要内容开始
% 
% * 实射影空间和复射影空间
% * 可定向及不可定向曲面
% * 透镜空间与扭项
% * 附着一个高维胞腔改变什么
% 
% 纲要内容结束

\section{典型空间的整数同调}
\label{g2:sec:0303}

胞腔链群告诉我们各维有多少个待检查的生成元，附着映射才决定哪些生成元成为圈、哪些圈又成为边界。下面的空间都有很少的胞腔，却已经展示了三种不同现象：边界为零时保留自由生成元，非零整数倍的边界产生有限阶同调类，而新增胞腔也可能使原有关系之间出现新的关系。本节的同调均取整数系数。

\subsection{实射影空间和复射影空间}
\label{g2:subsec:0303:01}

\cref{g2:exm:projective-cw}给出的标准分解中，\(\mathbb RP^N\) 在实维数 \(0,1,\ldots,N\) 各有一个胞腔，\(\mathbb CP^N\) 只在实维数 \(0,2,\ldots,2N\) 各有一个胞腔。两个空间都可以用“逐次加入一条新的坐标轴”构造，但它们相邻维数的链群很不相同：实情形允许非零边界，复情形的每一个边界映射都面对一个零群。

\begin{proposition}[射影空间的同调]\label{g2:prop:projective-integer-homology}
	对整数 \(N\geq1\)，实射影空间的整数同调为
	\[
		H_k(\mathbb RP^N)\cong
		\begin{cases}
			\mathbb Z,&k=0,\text{ 或 }k=N\text{ 且 }N\text{ 为奇数},\\
			\mathbb Z/2,&0<k<N\text{ 且 }k\text{ 为奇数},\\
			0,&\text{其余情形}.
		\end{cases}
	\]
	对整数 \(N\geq0\)，复射影空间的整数同调为
	\[
		H_k(\mathbb CP^N)\cong
		\begin{cases}
			\mathbb Z,&k=0,2,\ldots,2N,\\
			0,&\text{其余情形}.
		\end{cases}
	\]
	\(\mathbb RP^0\) 是单点，只有零次同调为 \(\mathbb Z\)。
\end{proposition}
\begin{proof}
	选用\eqref{g2:eq:rp-cellular-boundary}的相容胞腔定向，实射影空间的链复形在 \(0\) 至 \(N\) 次各为 \(\mathbb Z\)，并且
	\[
		\mathrm{d}_k=
		\begin{cases}
			2,&k\text{ 为正偶数},\\
			0,&k\text{ 为奇数}.
		\end{cases}
	\]
	若 \(0<k<N\) 为奇数，所有 \(k\) 链都是圈，而边界是它们的偶数倍，所以 \(H_k=\mathbb Z/2\)；若 \(k>0\) 为偶数，\(\mathrm{d}_k\) 单射，故该次同调为零。最高次须单独看：没有来自 \(N+1\) 次的边界，所以当 \(N\) 为奇数时 \(H_N=\ker\mathrm{d}_N=\mathbb Z\)，不是 \(\mathbb Z/2\)。零次的唯一生成元也没有非零边界。

	在复情形，奇数次链群为零，因此所有胞腔边界均为零，每个偶数次胞腔都给出一个无限阶同调生成元。两种计算通过\cref{g2:thm:cellular-singular-comparison}与奇异同调一致。
\end{proof}

例如，\(\mathbb RP^2\) 的一胞腔本身不是二维胞腔的边界，但它的两倍是边界。这就是 \(\mathbb Z/2\) 所记录的几何事实。对于阿贝尔群 \(G\)，满足 \(ng=0\) 的非零元素 \(g\)（其中 \(n\) 为某个正整数）称为\term[naoyuan]{挠元}[torsion element]；零元与全部挠元组成\term[naoziqun]{挠子群}[torsion subgroup]。本节标题中的“扭项”指这种有限阶信息，它由整数边界系数 \(2\) 保留下来。

\ProseParagraphBreak

无限射影空间要连同拓扑一起说明。令
\[
	\mathbb RP^\infty=\bigcup_{N\geq0}\mathbb RP^N,
	\qquad
	\mathbb CP^\infty=\bigcup_{N\geq0}\mathbb CP^N,
\]
其中包含映射在齐次坐标末尾补零，并赋予上述胞腔分解的弱拓扑。它们的每条胞腔链仍是有限和。由胞腔比较定理，直接使用无限链复形便得到
\[
	H_k(\mathbb RP^\infty)\cong
	\begin{cases}
		\mathbb Z,&k=0,\\
		\mathbb Z/2,&k>0\text{ 为奇数},\\
		0,&k>0\text{ 为偶数},
	\end{cases}
	\qquad
	H_k(\mathbb CP^\infty)\cong
	\begin{cases}
		\mathbb Z,&k\geq0\text{ 为偶数},\\
		0,&k\geq0\text{ 为奇数}.
	\end{cases}
\]
这里没有最高维胞腔，因而没有实射影奇数维顶部的那个例外。若从有限空间理解同一个计算，\cref{g2:thm:finite-subcomplex}保证每条奇异链的紧像落在某个有限子复形中；一个圈成为边界时，其填充链也只用有限多个胞腔。固定次数的结果因此由足够高的有限阶段决定。

\subsection{可定向及不可定向曲面}
\label{g2:subsec:0303:02}

我们使用两类熟悉的闭曲面模型。\(\Sigma_g\) 表示 \(g\) 个环面的连通和，\(g\geq1\)，并令 \(\Sigma_0=S^2\)；\(N_h\) 表示 \(h\) 个实射影平面的连通和，\(h\geq1\)。连通和是从两个曲面各删去一个开圆盘，再将留下的边界圆周粘合。前一类可定向，后一类不可定向；下面的计算只需它们的具体多边形模型。

\ProseParagraphBreak

环面由正方形按边界词 \(aba^{-1}b^{-1}\) 粘合，\(\mathbb RP^2\) 则由圆盘边界的两段按同向的 \(aa\) 粘合。对两个这样的多边形模型作连通和，可以在各自二维胞腔中挖去小圆盘，再切开一条连接内外边界的弧；粘合后，这两条切缝合成一条可重新合起的边，剩余外边界依次读出原来的两个词。因此逐次连通和给出一个顶点、若干条边和一个二维胞腔，其附着词分别为
\begin{equation}\label{g2:eq:surface-attaching-words}
	\Sigma_g:\quad\prod_{i=1}^g a_i b_i a_i^{-1}b_i^{-1},
	\qquad
	N_h:\quad a_1^2a_2^2\cdots a_h^2.
\end{equation}
在这些模型中，每条多边形边都成对粘合；边的内点有圆盘邻域，顶点周围的小扇形也首尾接成一个圆盘。因而商确为无边界曲面。可定向模型的成对边在多边形边界上反向出现，使二维胞腔的定向在粘合处相容；\(a_i^2\) 则在相应闭路周围反转局部定向。

\begin{proposition}[闭曲面模型的同调]\label{g2:prop:surface-integer-homology}
	设 \(\Sigma_g\) 为 \(g\ge0\) 个柄的闭可定向曲面模型（\(\Sigma_0=S^2\)），\(N_h\) 为 \(h\ge1\) 个实射影平面的连通和，取 \eqref{g2:eq:surface-attaching-words} 的胞腔分解。它们的整数同调由下表给出，所有高于二次的同调均为零。
	\begin{center}
		\begin{tabular}{cccc}
			\toprule
			空间 & \(H_0\) & \(H_1\) & \(H_2\)\\
			\midrule
			\(\Sigma_g\)（\(g\geq0\)） & \(\mathbb Z\) & \(\mathbb Z^{2g}\) & \(\mathbb Z\)\\
			\(N_h\)（\(h\geq1\)） & \(\mathbb Z\) & \(\mathbb Z^{h-1}\oplus\mathbb Z/2\) & \(0\)\\
			\bottomrule
		\end{tabular}
	\end{center}
\end{proposition}
\begin{proof}
	先设 \(g\geq1\)。每条边的两端都是唯一的顶点，所以 \(\mathrm{d}_1=0\)。\cref{g2:prop:word-boundary-degrees}把二维胞腔的边界系数化为附着词中各字母的指数和。在 \(\Sigma_g\) 的词中，每个 \(a_i,b_i\) 都以正负指数各出现一次，故 \(\mathrm{d}_2=0\)。其链复形为
	\[
		0\longrightarrow\mathbb Z\xrightarrow{0}\mathbb Z^{2g}
		\xrightarrow{0}\mathbb Z\longrightarrow0,
	\]
	于是得到第一行。\(g=0\) 时用球面的一零胞腔、一二胞腔分解，结论相同。

	在 \(N_h\) 中，每个字母的指数和是 \(2\)，所以
	\[
		\mathrm{d}_2(1)=2(a_1+\cdots+a_h),\qquad\mathrm{d}_1=0.
	\]
	这个 \(\mathrm{d}_2:\mathbb Z\to\mathbb Z^h\) 单射，故 \(H_2=0\)。令 \(v=a_1+\cdots+a_h\)，则 \(v,a_2,\ldots,a_h\) 是 \(\mathbb Z^h\) 的一组整数基，因为 \(a_1=v-a_2-\cdots-a_h\)。在这组基下只剩一个关系 \(2v=0\)，因此
	\[
		H_1(N_h)=\mathbb Z^h/\langle2v\rangle
		\cong\mathbb Z/2\oplus\mathbb Z^{h-1}.
	\]
	两种模型的一维骨架都连通，零次同调为 \(\mathbb Z\)，没有更高维胞腔。
\end{proof}

这里可以看见“边界抵消”与曲面定向的关系。对 \(\Sigma_g\)，二维胞腔的各对边在整数链中互相抵消，整个胞腔成为一个二维圈。对 \(N_h\)，同向出现的边留下系数 \(2\)，不能得到非零二维圈。这些模型中的顶维圈由此直接从多边形边界的符号算出。

\subsection{透镜空间与扭项}
\label{g2:subsec:0303:03}

实射影三维空间的 \(\mathbb Z/2\) 可以推广成任意 \(\mathbb Z/p\)。设 \(p\geq2\)，\(q\in\mathbb Z\) 与 \(p\) 互素，并记 \(\zeta=e^{2\pi\mathrm{i}/p}\)。在
\[
	S^3=\{(z_1,z_2)\in\mathbb C^2:|z_1|^2+|z_2|^2=1\}
\]
上考虑周期变换
\[
	\rho(z_1,z_2)=(\zeta z_1,\zeta^q z_2).
\]
定义三维\term[toujingkongjian]{透镜空间}[lens space]为
\[
	L(p,q)\coloneqq S^3/\langle\rho\rangle.
\]
\symidx{\(L(p,q)\)：三维透镜空间}这里固定的是商作用的参数约定；有些文献改用透镜面上的旋转角作第二参数，得到的公式会出现 \(q\) 的模 \(p\) 逆元。若 \(\rho^j\) 固定一个点，则至少一个坐标非零；由 \(p\mid j\) 或 \(p\mid qj\) 及互素条件，都得到 \(p\mid j\)。因此该作用自由。有限多个不同像可以用互不相交的小邻域分开，所以商投影是 \(p\) 叶覆盖，商空间是闭三维流形。

\begin{figure}[tbp]
	\centering
	\begin{tikzpicture}[scale=1.05,>=stealth]
		\draw[thick] (0,0) circle (1.25);
		\draw[dashed] (-1.25,0) arc[start angle=180,end angle=0,x radius=1.25,y radius=.28];
		\draw[thick] (-1.25,0) arc[start angle=180,end angle=360,x radius=1.25,y radius=.28];
		\node at (0,.65) {\(F_1\)};
		\node at (0,-.65) {\(F_0\)};
		\node[left] at (-1.25,0) {\(K\)};
		\draw[->,thick] (1.5,-.65) to[bend right=22] node[right] {\(\tau:z\mapsto\zeta^r z\)} (1.5,.65);
	\end{tikzpicture}
	\caption{基本区域三球 \(B\) 的边界示意。两张盘面 \(F_0,F_1\) 沿圆周 \(K\) 相接，并由 \(\tau\) 识别；\(K\) 本身还按 \(p\) 重旋转识别。盘面继承的边界定向相反，边缘的多重绕行则保留在二维边界中。}
	\label{g2:fig:lens-boundary-faces}
\end{figure}

\begin{proposition}[透镜空间的胞腔模型]\label{g2:prop:lens-cellular-model}
	设 \(p\geq2\) 且 \(\gcd(p,q)=1\)。按上述约定定义的 \(L(p,q)\) 有一个在维数 \(0,1,2,3\) 各含一个胞腔的 CW 分解。可以选取胞腔定向，使其整数链复形为
	\begin{equation}\label{g2:eq:lens-cellular-complex}
		0\longrightarrow\mathbb Z\xrightarrow{0}\mathbb Z
		\xrightarrow{p}\mathbb Z\xrightarrow{0}\mathbb Z
		\longrightarrow0.
	\end{equation}
\end{proposition}
\begin{proof}
	我们先找一个基本区域，再读出它的面和边的粘合，\cref{g2:fig:lens-boundary-faces}显示其边界两张盘面的关系。取整数 \(r\) 使 \(rq\equiv1\pmod p\)，并令 \(\tau=\rho^r\)。它仍生成同一个周期群，但在第二坐标上恰旋转 \(2\pi/p\)。在 \(S^3\) 中取扇区
	\[
		B=\left\{\left(z,\sqrt{1-|z|^2}\,e^{\mathrm{i}\theta}\right):
		|z|\leq1,\ 0\leq\theta\leq2\pi/p\right\}.
	\]
	当 \(|z|=1\) 时，第二坐标为零，所有 \(\theta\) 给出同一点。这个集合是闭三球，而不是边界带棱的环面。一个明确的参数说明如下：从普通闭三球 \(\{(z,s):|z|^2+s^2\leq1\}\subset\mathbb C\times\mathbb R\) 出发，在 \(|z|<1\) 时令
	\[
		\theta=\frac{\pi}{p}\left(1+\frac{s}{\sqrt{1-|z|^2}}\right),
	\]
	再代入 \(B\) 的公式；在 \(|z|=1\) 时取 \((z,0)\)。因为趋近边缘时第二坐标的模趋于零，该映射连续；它是紧空间到 Hausdorff 空间的连续双射，故是同胚。

	三球边界分成两个闭圆盘
	\[
		F_j=\left\{\left(z,\sqrt{1-|z|^2}\,e^{2\pi\mathrm{i}j/p}\right):|z|\leq1\right\},
		\qquad j=0,1,
	\]
	它们沿公共圆周 \(K=\{(z,0):|z|=1\}\) 相交。\(\tau\) 把 \(F_0\) 送到 \(F_1\)，在盘坐标中是旋转 \(z\mapsto\zeta^r z\)。若 \(z_2\ne0\) 且其辐角不是 \(2\pi/p\) 的整数倍，该轨道在扇区内部恰有一个代表；若辐角为 \(2\pi/p\) 的整数倍，则在两个盘面内部各有一个代表，它们由上述面粘合识别。剩下的点都位于 \(K\)，其识别为 \(z\sim\zeta z\)。因此从 \(B\) 出发作这些识别，恰得到 \(L(p,q)\)。

	边缘商 \(K/\langle\zeta\rangle\) 是圆周，商坐标可取 \(z\mapsto z^p\)，给它选一个顶点和一条开边。盘面内部的公共像是一开二胞腔，\(B\) 的内部是一开三胞腔；其余点都在刚才的圆周内。这给出所述四个胞腔。它们的特征映射来自圆盘和三球到商空间的映射；空间紧且只有有限多个胞腔，因此附着拓扑与所得商拓扑一致。

	给二胞腔选取 \(F_0\) 的复平面定向，其边界绕 \(K\) 一周；在商圆周中则由 \(z\mapsto z^p\) 绕行 \(p\) 周。因此\cref{g2:thm:cellular-boundary-degree}给出 \(\mathrm{d}_2=p\)。唯一一胞腔的两端相同，故 \(\mathrm{d}_1=0\)。

	最后计算 \(\mathrm{d}_3\)。把三球边界的两个半球都用同一个盘坐标 \(z\) 表示时，它们从三球继承的边界定向相反。两面的识别 \(z\mapsto\zeta^r z\) 又保持盘定向。在二胞腔内部取一点，它在这两个面中各有一个原像，局部度一正一负；压掉一维骨架以后，其度之和为 \(1-1=0\)。这就是 \(\mathrm{d}_3=0\) 的来源。
\end{proof}

两个盘面的定向抵消只说明三胞腔的边界系数为零；二胞腔边缘上的 \(p\) 重绕行仍然存在。由\eqref{g2:eq:lens-cellular-complex}立即得到
\begin{equation}\label{g2:eq:lens-integer-homology}
	H_k(L(p,q))\cong
	\begin{cases}
		\mathbb Z,&k=0,3,\\
		\mathbb Z/p,&k=1,\\
		0,&\text{其余情形}.
	\end{cases}
\end{equation}
一胞腔给出的圈有阶恰为 \(p\)：它的 \(p\) 倍被二胞腔填充，任何较小正倍数都不在 \(p\mathbb Z\) 内。这一计算保留了覆盖重数 \(p\)，却没有保留盘面旋转中的全部参数 \(q\)；因此相同的这组同调群本身不能回答两个透镜空间是否同胚。特别地，\(L(2,1)=\mathbb RP^3\)，与前面的结果吻合。上述胞腔模型也可对照\cite[第2.2节，Cellular Homology]{Hatcher2002}。

\subsection{附着一个高维胞腔改变什么}
\label{g2:subsec:0303:04}

附着球面在旧空间中的同调类，可以是无限阶、有限阶或零。空间对的长正合列分别给出这三种情形中新生成的类与新施加的关系，并说明一次附着只可能影响两个相邻次数的同调。

\begin{proposition}[单胞腔附着的同调变化]\label{g2:prop:single-cell-homology-change}
	设 \(X\) 是非空 CW 复形，\(m\geq2\)，\(\alpha:S^{m-1}\to X^{m-1}\) 连续，并令 \(Y=X\cup_\alpha D^m\)。给圆盘选定向，球面取相应边界定向，记 \(u=\alpha_*[S^{m-1}]\in H_{m-1}(X)\)。则包含 \(X\hookrightarrow Y\) 在 \(k\notin\{m,m-1\}\) 时诱导同构 \(H_k(X)\cong H_k(Y)\)，且有正合列
	\begin{equation}\label{g2:eq:single-cell-six-term}
		\begin{gathered}
			0\longrightarrow H_m(X)\longrightarrow H_m(Y)\longrightarrow\mathbb Z\\
			\xrightarrow{\,1\mapsto u\,}H_{m-1}(X)\longrightarrow H_{m-1}(Y)\longrightarrow0.
		\end{gathered}
	\end{equation}
	特别地，\(H_{m-1}(Y)\cong H_{m-1}(X)/\langle u\rangle\)。若 \(u\) 为无限阶，则 \(H_m(Y)\cong H_m(X)\)；若 \(u\) 有有限阶（包括 \(u=0\)），则存在群同构 \(H_m(Y)\cong H_m(X)\oplus\mathbb Z\)，但后一分解一般没有指定的自然选择。
\end{proposition}
\begin{proof}
	相对胞腔复形 \(C^{\mathrm{cell}}_\bullet(Y,X)\) 只剩新增的一个 \(m\) 次生成元，所有微分为零。由\cref{g2:thm:relative-cellular-comparison}，相对奇异同调为
	\[
		H_k(Y,X)\cong
		\begin{cases}
			\mathbb Z,&k=m,\\
			0,&k\ne m.
		\end{cases}
	\]
	代入\cref{g2:thm:pair-LES}便得到不受影响的次数和\eqref{g2:eq:single-cell-six-term}。还要认清其中的连接同态：特征映射是空间对映射 \((D^m,S^{m-1})\to(Y,X)\)，其相对定向类给出上面 \(\mathbb Z\) 的生成元。连接同态的自然性，以及“取代表链的边界”这一约定，说明它的像恰为 \(\alpha_*[S^{m-1}]=u\)。

	由正合性，低一次群是所述商群，而最高一段可以写成
	\[
		0\longrightarrow H_m(X)\longrightarrow H_m(Y)
		\longrightarrow K\longrightarrow0,
		\qquad K=\{a\in\mathbb Z:au=0\}.
	\]
	若 \(u\) 为无限阶，则 \(K=0\)。若 \(u\) 的阶为 \(s<\infty\)，约定零元的阶为 \(1\)，则 \(K=s\mathbb Z\cong\mathbb Z\)。在 \(H_m(Y)\) 中选一个映到 \(s\) 的元素，它的整数倍给出该短正合列的分裂，因而得到直和分解。更换这个提升可以加上 \(H_m(X)\) 中的元素，所以分解不由原数据唯一指定。
\end{proof}

若 \(u\ne0\) 的阶为 \(s\)，新增胞腔本身的边界还没有消失，但它的 \(s\) 倍的边界能在 \(X\) 中填充。用 \(s\) 份新胞腔减去这条旧填充链，才得到一个新的 \(m\) 维圈。于是一次附着可以同时消去低一次的挠类，并产生高一次的无限阶类。

\begin{example}[在球面上附着一个胞腔]\label{g2:exm:degree-cell-attachment}
	设 \(m\geq2\)，\(f:S^{m-1}\to S^{m-1}\) 的度为 \(a\in\mathbb Z\)，并令 \(Y=S^{m-1}\cup_fD^m\)。取球面的一零胞腔、一 \((m-1)\) 胞腔分解，则其唯一可能非零的高次边界为 \(\mathrm{d}_m=a\)。若 \(a\ne0\)，则
	\[
		\widetilde H_{m-1}(Y)=\mathbb Z/a\mathbb Z,
		\qquad \widetilde H_k(Y)=0\quad(k\ne m-1).
	\]
	若 \(a=0\)，则 \(\widetilde H_m(Y)=\widetilde H_{m-1}(Y)=\mathbb Z\)，其余约化同调为零。这也可逐项代入\eqref{g2:eq:single-cell-six-term}得到。对于 \(|a|=1\)，附着后约化同调全部消失；对于 \(|a|>1\)，留下阶为 \(|a|\) 的类。同调计算本身没有进一步证明这个空间在前一种情形可缩，也没有把后一种情形的整个同伦型归结为这些群。
\end{example}

这些例子把胞腔数、附着的整数系数和最终同调分开了。胞腔数相同可以有不同的同调；同调相同也只表示这里计算的阿贝尔群信息相同。要可靠地处理胞腔较多的空间，还需在整数基中同时控制核与像，这正是下一节边界矩阵计算的任务。
% 纲要条目排为编号小节标题，并进入目录与书签；说明段仍保留于纲要注释中。

% !TeX root = ../../Book2.tex
% 纲要标题：### 3.4 有限复形的整数同调计算
% 纲要位置：几何学计划纲要.md，第 1142 行。
% 正文状态：初稿；尚未正式审核。

% 纲要内容开始
% 
% * 链群基与整数边界矩阵
% * Smith 标准形及核像计算
% * Euler 示性数的同调公式
% * 不同 CW 分解之间的结果核对
% 
% ---
% 
% 纲要内容结束

\section{有限复形的整数同调计算}
\label{g2:sec:0304}

已知每个胞腔边界中的整数系数以后，同调计算便可以交给矩阵运算。但工作环必须保留为 \(\mathbb Z\)：把关系 \(6z=0\) 除以 \(6\)，就会把一个可能非零的六阶元素消去。本节给出既能手算也能由程序执行的整数算法，并说明怎样用换基矩阵核验计算结果。计算需用圈群的一组整数基表示边界群，再取商。

\subsection{链群基与整数边界矩阵}
\label{g2:subsec:0304:01}

设 \(X\) 是有限 CW 复形，\(c_n\) 是其 \(n\) 维胞腔数。固定各维胞腔的次序和定向，\cref{g2:def:cellular-chain-complex}给出的胞腔链群便成为带基的自由阿贝尔群
\[
	C_n^{\mathrm{cell}}(X)\cong\mathbb Z^{c_n}.
\]
将一个胞腔的边界系数按列排列，得到 \(c_{n-1}\times c_n\) 的整数矩阵 \(A_n\)。若链 \(x\) 的坐标写为列向量，则 \(\mathrm{d}_nx\) 的坐标是 \(A_nx\)。由边界平方为零，必须有
\[
	A_nA_{n+1}=0.
\]
这一等式保证边界坐标落在相应的核中，因而可以在圈群内计算边界群。若乘积非零，应回查定向、附着映射或抄录的系数。

\begin{definition}[幺模矩阵]\label{g2:def:unimodular-matrix}
	方阵 \(P\in M_r(\mathbb Z)\) 若满足 \(\det P=\pm1\)，称为\term[yaomo juzhen]{幺模矩阵}[unimodular matrix]。等价地，\(P\) 的逆矩阵仍有整数系数。
\end{definition}

等价性的一个方向由伴随矩阵公式 \(P^{-1}=(\det P)^{-1}\operatorname{adj}P\) 得到；另一方向由 \(\det P\det P^{-1}=1\) 及两因子为整数得到。交换两个基元、把一个基元乘以 \(-1\)、把一个基元加上另一个基元的整数倍，都由幺模矩阵实现。把某个基元乘以 \(2\) 则不是换基：所得元素只生成原自由群的一个真子群。

\ProseParagraphBreak

若第 \(n\) 次的新基用旧基表示的列组成 \(P_n\)，则旧坐标与新坐标满足 \(x=P_nx'\)，因而新边界矩阵为
\begin{equation}\label{g2:eq:boundary-matrix-change}
	A_n'=P_{n-1}^{-1}A_nP_n.
\end{equation}
同一群 \(C_n\) 同时是 \(\mathrm{d}_n\) 的源和 \(\mathrm{d}_{n+1}\) 的靶，故它的换基会同时改变这两张矩阵。对 \(A_n\) 作列变换以后，必须对 \(A_{n+1}\) 作相应的逆行变换；这使两张矩阵中的坐标继续属于同一组基，从而可以计算 \(\ker\mathrm{d}_n/\operatorname{im}\mathrm{d}_{n+1}\)。

\subsection{Smith 标准形及核像计算}
\label{g2:subsec:0304:02}

整数换基用带余除法逐步减小主元，直到一个正主元整除其后的所有系数；这一步替代了域上将主元归一为 \(1\) 的操作。这个整除关系正好保留了有限阶同调类。

\begin{theorem}[Smith 标准形]\label{g2:thm:smith-normal-form}
	设 \(A\) 是 \(p\times q\) 整数矩阵。存在幺模矩阵 \(U\in M_p(\mathbb Z)\)、\(V\in M_q(\mathbb Z)\)，使
	\[
		UAV=S=\operatorname{diag}(s_1,\ldots,s_r,0),\qquad
		0<s_1\mid s_2\mid\cdots\mid s_r.
	\]
	这里的对角形允许长方形矩阵，未列出的元素均为零，\(r=\operatorname{rank}_{\mathbb Q}A\)。整数 \(s_i\) 由 \(A\) 唯一确定，称为其\term[bubian yinzi]{不变因子}[invariant factors]；矩阵 \(S\) 称为 \(A\) 的\term[Smith biaozhunxing]{Smith 标准形}[Smith normal form]。\footnote{Henry John Stephen Smith（1826--1883），爱尔兰数学家，研究数论和初等因子。}矩阵 \(U,V,S\) 可以由有限次整数初等行、列变换求出。
\end{theorem}
\begin{proof}
	先证明存在性，并把构造写成可执行的过程。零矩阵无须处理。若 \(A\ne0\)，把一个非零元素移到左上角，再变号使主元 \(a=A_{11}>0\)。对第一列中的元素 \(b\) 作带余除法 \(b=ka+r\)，\(0\le r<a\)，用相应行减去第一行的 \(k\) 倍。如果 \(r\ne0\)，交换该行与第一行，得到严格较小的正主元，重新开始。若余数均为零，则可以把主元以下的第一列清零。以相同方法对第一行作列变换；一旦出现非零余数，交换至主元位置并重新开始。

	每次需要重新开始时，正主元严格变小，所以这种下降只能发生有限次。两次下降之间只须处理有限个行、列元素，故最终得到
	\[
		\begin{pmatrix}a&0\\0&B\end{pmatrix},\qquad a>0.
	\]
	此时还不能立即结束：必须检查 \(a\) 是否整除 \(B\) 中的每个元素。若某个 \(B_{ij}=b\) 不被 \(a\) 整除，把它所在的行加到第一行；第一列的主元仍是 \(a\)，第一行相应位置变为 \(b\)。再对该列减去第一列的整数倍，出现非零余数 \(0<r<a\)，交换两列后主元再次严格变小。因此，有限次下降以后可以同时做到第一行、列清零以及 \(a\) 整除整个余下矩阵。

	只在 \(B\) 内重复上述过程。由于 \(B\) 的各项原来均为 \(a\) 的倍数，任何整数行、列变换所得的项仍为 \(a\) 的倍数，所以第二个正主元必被第一个整除。按矩阵较小维数归纳，得到要求的对角形。每一步都是交换、变号或整数倍相加，因而累计的左右乘矩阵 \(U,V\) 均幺模。

	再证明唯一性。令 \(\Delta_0(A)=1\)，对 \(1\le k\le\min\{p,q\}\)，令 \(\Delta_k(A)\) 为全部 \(k\) 阶子式的非负最大公因数；若这些子式全为零，则取 \(\Delta_k(A)=0\)。一次整数初等行变换后，每个新子式都由旧子式的整数线性组合得到：对变换涉及的行展开行列式即可，出现重复行的项为零。因此旧子式的公因数整除全部新子式；应用逆变换得到反向整除。列变换同理。对于一般的整数左、右乘矩阵，也可逐行、逐列作多线性展开，得到同样的整数线性组合；若乘矩阵幺模，再用其整数逆矩阵就得到反向整除。因此 \(\Delta_k(UAV)=\Delta_k(A)\)。

	对于所得对角形，\(k\le r\) 时每个非零 \(k\) 阶子式是 \(k\) 个不同 \(s_i\) 的乘积。由 \(s_i\mid s_{i+1}\)，这些乘积均被 \(s_1\cdots s_k\) 整除，而这个乘积本身就是一个子式。因此
	\[
		\Delta_k(A)=s_1\cdots s_k\quad(1\le k\le r),
		\qquad s_k=\frac{\Delta_k(A)}{\Delta_{k-1}(A)}.
	\]
	当 \(k>r\) 时全部子式为零。这确定了 \(r\) 及每个 \(s_i\)，也表明 \(r\) 就是在 \(\mathbb Q\) 上的矩阵秩。
\end{proof}

从标准形可以直接读出整数核与像的一组基。对于同态 \(A:\mathbb Z^q\to\mathbb Z^p\)，由 \(Ax=0\iff S(V^{-1}x)=0\) 可知，\(V\) 的后 \(q-r\) 列构成 \(\ker A\) 的自由基。而 \(\operatorname{im}A\) 的一组基为
\[
	U^{-1}s_1e_1,\ldots,U^{-1}s_re_r,
\]
其中 \(e_i\) 是 \(\mathbb Z^p\) 的标准基。于是
\[
	\mathbb Z^p/\operatorname{im}A
	\cong\mathbb Z^{p-r}\oplus\bigoplus_{i=1}^r\mathbb Z/s_i\mathbb Z.
\]
当 \(s_i=1\) 时，相应的商群为零，应从最终结果中略去。像的基一般不能扩充为靶群的一组整数基，例如 \(2\mathbb Z\subset\mathbb Z\)；这种不能扩充的现象正是商群出现挠的原因。

\begin{proposition}[在圈群中计算边界]\label{g2:prop:integer-homology-matrices}
	设
	\[
		\mathbb Z^b\xrightarrow{B}\mathbb Z^c\xrightarrow{A}\mathbb Z^a,
		\qquad AB=0,
	\]
	且 \(UAV=\operatorname{diag}(s_1,\ldots,s_r,0)\) 为 Smith 标准形。则 \(V^{-1}B\) 的前 \(r\) 行为零。将其余 \(c-r\) 行组成的矩阵记为 \(K\)，并设
	\[
		PKQ=\operatorname{diag}(t_1,\ldots,t_\ell,0),
		\qquad 0<t_1\mid\cdots\mid t_\ell,
	\]
	其中 \(P,Q\) 幺模。则
	\begin{equation}\label{g2:eq:integer-homology-result}
		\ker A/\operatorname{im}B
		\cong\mathbb Z^{c-r-\ell}\oplus
		\bigoplus_{i=1}^\ell\mathbb Z/t_i\mathbb Z.
	\end{equation}
	若 \(J:\mathbb Z^{c-r}\to\mathbb Z^c\) 把坐标填入最后 \(c-r\) 个位置，则这些商群生成元可由圈 \(VJP^{-1}e_i\) 表示。
\end{proposition}
\begin{proof}
	由 \(AB=0\) 得
	\[
		S(V^{-1}B)=UAB=0.
	\]
	前 \(r\) 个非零对角元是非零整数；整数群中没有被非零整数消去的非零元素，所以 \(V^{-1}B\) 的前 \(r\) 行均为零。映射 \(y\mapsto VJy\) 是从 \(\mathbb Z^{c-r}\) 到 \(\ker A\) 的同构，并把 \(\operatorname{im}K\) 送到 \(\operatorname{im}B\)。于是所求群等于 \(\mathbb Z^{c-r}/\operatorname{im}K\)。再用 \(P\) 变换坐标，像变成由 \(t_ie_i\)（\(1\le i\le\ell\)）生成的子群，逐坐标取商就得到\eqref{g2:eq:integer-homology-result}及所述代表圈。
\end{proof}

\begin{algorithm}[tbp]
	\caption{整数同调}
	\label{g2:alg:integer-homology}
	\begin{algorithmsteps}{有限整数矩阵列 \(A_n:\mathbb Z^{c_n}\to\mathbb Z^{c_{n-1}}\)。}{各次同调的自由秩、挠不变因子及代表圈；用于核验的幺模换基矩阵。}
		\State 检查各矩阵尺寸相容，且 \(A_nA_{n+1}=0\)；不满足时停止并回查输入。
		\State 对所求次数 \(n\)，计算 \(UA_nV=S\)，保存 \(U,V\)，记非零对角元数为 \(r\)。
		\State 计算 \(V^{-1}A_{n+1}\)，检查前 \(r\) 行为零，取余下行组成 \(K\)。
		\State 计算 \(PKQ=T\)，保存 \(P,Q\)，记非零对角元数为 \(\ell\)。
		\State 输出自由秩 \(c_n-r-\ell\)、\(T\) 中大于 \(1\) 的对角元，以及相应的圈 \(VJP^{-1}e_i\)。
		\State 用整数运算复核两组矩阵等式、换基矩阵的行列式为 \(\pm1\)、对角元为正且依次整除。
	\end{algorithmsteps}
\end{algorithm}

核对\cref{g2:alg:integer-homology}中的矩阵等式、换基矩阵的幺模性及对角元的整除关系，即可独立验证结果。零矩阵以及零行、零列的情形也包含在内：相应的秩为零，空直和为零群。若只想求某一次同调，可以独立使用原来的相邻矩阵 \(A_n,A_{n+1}\)；若把多次变换原地写回整条链复形，则须按\eqref{g2:eq:boundary-matrix-change}同步更新相邻矩阵。

\begin{example}[两组有限阶关系]\label{g2:exm:integer-matrix-certificate}
	取两个顶点 \(v_0,v_1\)，两条从 \(v_0\) 指向 \(v_1\) 的边 \(a,b\)，以及一条以 \(v_0\) 为基点的环边 \(c\)。沿环路 \((ba^{-1})^2\) 和 \((ba^{-1})^2c^6\) 分别附着一个二维胞腔。这里的字从左向右行走，故 \(ba^{-1}\) 确实从 \(v_0\) 出发又返回 \(v_0\)。按基 \((a,b,c)\) 写出的边界矩阵为
	\[
		A_1=\begin{pmatrix}-1&-1&0\\1&1&0\end{pmatrix},
		\qquad
		A_2=\begin{pmatrix}-2&-2\\2&2\\0&6\end{pmatrix}.
	\]
	直接相乘得到 \(A_1A_2=0\)。取
	\[
		U=\begin{pmatrix}-1&0\\1&1\end{pmatrix},
		\qquad
		V=\begin{pmatrix}1&-1&0\\0&1&0\\0&0&1\end{pmatrix}.
	\]
	其行列式分别为 \(-1,1\)，且
	\[
		UA_1V=\begin{pmatrix}1&0&0\\0&0&0\end{pmatrix},
		\qquad
		V^{-1}A_2=\begin{pmatrix}0&0\\2&2\\0&6\end{pmatrix}.
	\]
	因此圈群的基是 \(u=b-a\)、\(v=c\)，两条二维边界在这组基中的矩阵为
	\[
		K=\begin{pmatrix}2&2\\0&6\end{pmatrix}.
	\]
	把第二列减去第一列，即取
	\[
		Q=\begin{pmatrix}1&-1\\0&1\end{pmatrix},\qquad
		KQ=\begin{pmatrix}2&0\\0&6\end{pmatrix},
	\]
	得到 \(H_1(X)\cong\mathbb Z/2\mathbb Z\oplus\mathbb Z/6\mathbb Z\)，其生成元分别是 \([u],[v]\)。零次同调由两个顶点的差关系给出 \(H_0(X)\cong\mathbb Z\)。矩阵 \(A_2\) 的两列线性无关，故 \(H_2(X)=0\)，更高次也为零。只在有理数上消元，会看到一次同调的自由秩为零，却无法恢复这两组有限阶关系。
\end{example}

\subsection{Euler 示性数的同调公式}
\label{g2:subsec:0304:03}

完整的整数同调保留了每个有限阶关系，但有时我们先需要一个快速的总量核对。例如，增添一对相邻维数的胞腔可能改变边界矩阵的大小，却不改变交错计数。把各次矩阵的秩写出来，就能说明这种抵消从何而来。

\begin{definition}[Euler 示性数]\label{g2:def:euler-characteristic}
	设 \(X\) 是有限 CW 复形，\(c_n\) 是其 \(n\) 维胞腔数。称
	\[
		\chi(X)\coloneqq\sum_{n\ge0}(-1)^nc_n
	\]
	为这个胞腔分解的\term[Euler shixingshu]{Euler 示性数}[Euler characteristic]。\footnote{Leonhard Euler（1707--1783），瑞士数学家，在分析、几何和数论等领域作出广泛贡献。}由下一定理，它实际上只依赖于 \(X\) 的同伦型。\symidx{\(\chi(X)\)：Euler 示性数}
\end{definition}

\begin{theorem}[Euler 示性数的同调公式]\label{g2:thm:euler-homology}
	设 \(X\) 是有限 CW 复形，\(A_n\) 是其整数胞腔边界矩阵，令 \(r_n=\operatorname{rank}_{\mathbb Q}A_n\)，并令 \(b_n\) 为 \(H_n(X;\mathbb Z)\) 的自由秩。则
	\[
		b_n=c_n-r_n-r_{n+1},\qquad
		\chi(X)=\sum_{n\ge0}(-1)^nb_n.
	\]
	特别地，两个同伦等价的有限 CW 复形具有相同的 Euler 示性数。
\end{theorem}
\begin{proof}
	由\cref{g2:prop:integer-homology-matrices}，\(\ker A_n\) 的秩是 \(c_n-r_n\)。在它的整数基中表示的 \(A_{n+1}\) 具有秩 \(r_{n+1}\)，因为换基和删去零行均不改变有理矩阵秩。因此商群的自由秩为 \(c_n-r_n-r_{n+1}\)。\cref{g2:thm:cellular-singular-comparison}把这个商群与 \(H_n(X;\mathbb Z)\) 识别起来，得到第一式。

	设 \(X\) 的维数为 \(N\)，在两端取 \(r_0=r_{N+1}=0\)。对第一式作交错求和，两个秩和逐项相消：
	\[
		\sum_{n=0}^N(-1)^nb_n
		=\sum_{n=0}^N(-1)^nc_n
		-\sum_{n=0}^N(-1)^nr_n
		-\sum_{n=0}^N(-1)^nr_{n+1}
		=\sum_{n=0}^N(-1)^nc_n.
	\]
	同伦等价诱导奇异同调群同构，故各次自由秩相等，第二个结论随之成立。
\end{proof}

这一公式也见\cite[第2.2节，Cellular Homology]{Hatcher2002}。有限性保证交错和只有有限项；对于具有无限多个胞腔的空间，不能直接把发散的胞腔计数作抵消。挠群的自由秩为零，因此 Euler 示性数本身不能检查挠信息是否计算正确。例如\cref{g2:exm:integer-matrix-certificate}中 \(\chi(X)=2-3+2=1\)，与 \(b_0=1\)、其余 \(b_n=0\) 相符；把一次同调误写成零仍会通过这一项核对。

\subsection{不同 CW 分解之间的结果核对}
\label{g2:subsec:0304:04}

同一个空间可以有大小差别很大的边界矩阵，胞腔数与代表圈也随分解改变；这些矩阵算出的同调群则彼此同构。\cref{g2:thm:cellular-singular-comparison}提供了比较依据：每种胞腔分解的同调都自然同构于空间本身的奇异同调。对两个分解之间的恒等映射作胞腔逼近，还可以把这项比较落实为胞腔链映射。

\begin{example}[圆周的两个分解]\label{g2:exm:circle-cell-decompositions}
	圆周可以只有一个顶点 \(v\) 和一个边胞腔 \(e\)，此时 \(\mathrm{d}_1e=v-v=0\)，链复形是 \(\mathbb Z\xrightarrow{0}\mathbb Z\)。也可以在圆周上顺次取三个顶点 \(v_0,v_1,v_2\)，沿正向取三条边 \(e_{01},e_{12},e_{20}\)。后一分解的边界矩阵为
	\[
		A_1=\begin{pmatrix}-1&0&1\\1&-1&0\\0&1&-1\end{pmatrix}.
	\]
	解 \(A_1(x,y,z)^{\mathsf T}=0\)，得到 \(x=y=z\)，故一次圈群由 \(e_{01}+e_{12}+e_{20}\) 生成。零次边界群由 \(v_1-v_0,v_2-v_1\) 生成，恰为坐标和为零的子群，所以 \(H_0\cong\mathbb Z\)。两种分解都给出 \(H_1\cong\mathbb Z\)，其余正次数同调为零。相应 Euler 示性数分别为 \(1-1=0\) 和 \(3-3=0\)。

	还可以直接看出两套计算为何相容。从粗分解到细分解定义链映射 \(F(v)=v_0\)、\(F(e)=e_{01}+e_{12}+e_{20}\)。反向定义 \(G(v_i)=v\)，并取 \(G(e_{01})=G(e_{12})=0\)、\(G(e_{20})=e\)。边界公式直接给出二者均为链映射，且 \(GF=\operatorname{id}\)。令次数上升一的同态为
	\[
		h(v_0)=0,\qquad h(v_1)=e_{01},\qquad
		h(v_2)=e_{01}+e_{12},\qquad h|_{C_1}=0.
	\]
	在三个顶点上，\(\mathrm{d}h=\operatorname{id}-FG\)；在前两条边上，\(h\mathrm{d}\) 分别等于该边，在 \(e_{20}\) 上则等于 \(-e_{01}-e_{12}\)，也正是 \((\operatorname{id}-FG)(e_{20})\)。因此 \(\operatorname{id}-FG=\mathrm{d}h+h\mathrm{d}\)。按\cref{g2:prop:chain-homotopy-homology}，两套链复形诱导互逆的同调同构；这里细分增加的自由生成元在链同伦中被消去。
\end{example}

在实际核对中，首先检查两边的各次自由秩及非平凡挠不变因子，再检查 Euler 示性数；后者是较弱但很便宜的整体核对。若还要比较某个指定的同调类或映射，就必须保留链映射或代表圈的坐标，只比较抽象群的名称不够。例如两个计算都得到 \(\mathbb Z/6\mathbb Z\)，仍须说明某个指定圈对应生成元、其两倍还是零。

\ProseParagraphBreak

这里非平凡挠因子按整除次序排列后确实由群本身决定。具体地，对有限挠群 \(F=\bigoplus_i\mathbb Z/t_i\mathbb Z\) 和素数 \(p\)，记 \(F[p^k]=\{x:p^kx=0\}\)，则
\[
	|F[p^k]|=p^{\sum_i\min\{k,\nu_p(t_i)\}},
\]
其中 \(\nu_p(t_i)\) 是 \(t_i\) 中 \(p\) 的指数。相邻两个指数和的差给出满足 \(\nu_p(t_i)\ge k\) 的因子数，因而恢复各个 \(p\) 的指数列；按整除次序从右端对齐这些列，就恢复全部 \(t_i>1\)。自由秩则是除去有限阶元素后的自由阿贝尔群秩。这样，两种胞腔分解算出的整数同调同构，既意味着自由秩一致，也意味着有限阶关系一致。

% 纲要条目排为编号小节标题，并进入目录与书签；说明段仍保留于纲要注释中。


\begin{chaptersummary}
待撰写
\end{chaptersummary}

\BookExercises
待撰写

% !TeX root = ../../Book2.tex
% 纲要标题：## 第4章　系数变换、乘积与上同调环
% 纲要位置：几何学计划纲要.md，第 1151 行。
% 正文状态：初稿；尚未正式审核。

% 纲要内容开始
% 
% > 本章目标与证明范围：在确实需要系数变化时建立自由消解、Ext 与 Tor。随后从链级乘积推出环结构和 Künneth；不把自然正合列的非典范分裂写成自然同构。
% 
% 纲要内容结束

\chapter{系数变换、乘积与上同调环}
\label{g2:ch:04}
\BookChapterNumber{g2:ch:04}

\begin{chapterabstract}
	本章研究系数与乘法怎样丰富空间的不变量。自由消解把取同态、取张量时的正合性障碍组织为 Ext 与 Tor，万有系数定理据此从自由链复形的同调恢复任意系数同调及上同调。Alexander--Whitney 映射与 shuffle 映射比较乘积空间的链与张量复形，导出 Künneth 短正合列。随后通过对角映射定义杯积，证明自然性、结合性、单位与分次交换性，并建立相对乘积和取值配对。环面、闭曲面及复射影空间的计算说明，即使加性同调相同，乘法仍能区分不同的同伦类型。
\end{chapterabstract}

\begin{learninggoals}
	\begin{itemize}
		\item 构造自由消解及其比较链映射，计算循环模的 Ext、Tor，并说明两者分别测量哪种正合性障碍。
		\item 使用两种万有系数正合列计算系数同调与上同调，区别自然映射和依赖选择的分裂。
		\item 写出 Alexander--Whitney 与 shuffle 映射，核对总微分及交换因子的符号，用 Künneth 定理计算积空间中的自由类与挠类。
		\item 从上链定义杯积，证明其环性质及分次交换性，判断相对杯积的目标及所需切除条件。
		\item 计算环面、可定向曲面、不可定向曲面的模二环及复射影空间的上同调环，用乘法关系约束映射和区分同伦类型。
	\end{itemize}
\end{learninggoals}

上一章中，射影平面的胞腔链含有一条乘以 \(2\) 的边界。用整数计算，它产生一次同调中的二阶挠类，却不产生二次同调；把系数改成 \(\mathbb F_2\) 后，这条边界成为零，二次圈反而出现了。系数变化不仅改变已有群中元素的写法，还可能把原来有边界的链变成新的圈。因此，需要在链复形上处理系数变化，再计算其同调。

\ProseParagraphBreak

另一个问题来自 \(\mathbb CP^2\) 与 \(S^2\vee S^4\)。两者的整数同调都在次数 \(0,2,4\) 各有一个 \(\mathbb Z\)，单看这些群无法区分它们。可是 \(\mathbb CP^2\) 中两条坐标射影直线交于一点，这暗示两个二次信息同时出现时可能留下一个四次信息；在球面的楔和中，两个维数的部分则只在基点接触。杯积会把这种差别变成可以证明、可以计算的乘法关系。

\ProseParagraphBreak

系数变化带来的新圈和新上同调类，将由自由消解、Ext 与 Tor 描述；这些构造从模、核、商、自由基和线性映射开始建立。整数与域上的计算是主要例子，采用主理想整环的统一表述，是因为其自由模的子模仍然自由，足以把所需障碍放进两项消解。转向乘积后，链的张量构造把两组维数合为总次数；Künneth 定理据此计算乘积同调，杯积再沿对角映射给出同一空间上的乘法。杯积只要求系数为含幺交换环，Künneth 的短正合列则使用主理想整环条件。

\ProseParagraphBreak

几何部分使用\chapref{g2:ch:02}的奇异链、链同伦、空间对正合列和切除：链同伦用于比较乘积的不同链模型，正合列与切除用于构造相对乘积。\chapref{g2:ch:03}的胞腔计算则提供换系数前的整数同调。关于上同调与万有系数、乘积与杯积的阅读对照，可分别查\cite[第3.1节、第3.A节；第3.2节、第3.B节]{Hatcher2002}。

% !TeX root = ../../Book2.tex
% 纲要标题：### 4.1 自由消解与系数障碍
% 纲要位置：几何学计划纲要.md，第 1105 行。
% 正文状态：未撰写；本文件目前是框架。

% 纲要内容开始
% 
% * 自由消解和比较同伦
% * Ext、Tor 的定义与整数循环模计算
% * 维数平移的最低形式
% * 张量复形与符号约定
% 
% 纲要内容结束

\section{自由消解与系数障碍}
\label{g2:sec:0401}

待撰写

\subsection{自由消解和比较同伦}
\label{g2:subsec:0401:01}
待撰写

\subsection{Ext、Tor 的定义与整数循环模计算}
\label{g2:subsec:0401:02}
待撰写

\subsection{维数平移的最低形式}
\label{g2:subsec:0401:03}
待撰写

\subsection{张量复形与符号约定}
\label{g2:subsec:0401:04}
待撰写
% 纲要条目排为编号小节标题，并进入目录与书签；说明段仍保留于纲要注释中。

% !TeX root = ../../Book2.tex
% 纲要标题：### 4.2 万有系数定理
% 纲要位置：几何学计划纲要.md，第 1112 行。
% 正文状态：未撰写；本文件目前是框架。

% 纲要内容开始
% 
% * 自由链复形的圈边界分解
% * 上同调的自然短正合列
% * 同调的系数变换形式
% * 非典范分裂及函子性的限制
% 
% 纲要内容结束

\section{万有系数定理}
\label{g2:sec:0402}

待撰写

\subsection{自由链复形的圈边界分解}
\label{g2:subsec:0402:01}
待撰写

\subsection{上同调的自然短正合列}
\label{g2:subsec:0402:02}
待撰写

\subsection{同调的系数变换形式}
\label{g2:subsec:0402:03}
待撰写

\subsection{非典范分裂及函子性的限制}
\label{g2:subsec:0402:04}
待撰写
% 纲要条目排为编号小节标题，并进入目录与书签；说明段仍保留于纲要注释中。

% !TeX root = ../../Book2.tex
% 纲要标题：### 4.3 空间乘积与 Eilenberg–Zilber
% 纲要位置：几何学计划纲要.md，第 1119 行。
% 正文状态：未撰写；本文件目前是框架。

% 纲要内容开始
% 
% * Alexander–Whitney 与 shuffle 映射
% * 链同伦等价和符号
% * 域及主理想整环上的 Künneth 定理
% * Tor 项和积空间算例
% 
% 纲要内容结束

\section{空间乘积与 Eilenberg–Zilber}
\label{g2:sec:0403}

待撰写

\subsection{Alexander–Whitney 与 shuffle 映射}
\label{g2:subsec:0403:01}
待撰写

\subsection{链同伦等价和符号}
\label{g2:subsec:0403:02}
待撰写

\subsection{域及主理想整环上的 Künneth 定理}
\label{g2:subsec:0403:03}
待撰写

\subsection{Tor 项和积空间算例}
\label{g2:subsec:0403:04}
待撰写
% 纲要条目排为编号小节标题，并进入目录与书签；说明段仍保留于纲要注释中。

% !TeX root = ../../Book2.tex
% 纲要标题：### 4.4 杯积与环结构
% 纲要位置：几何学计划纲要.md，第 1176 行。
% 正文状态：初稿；尚未正式审核。

% 纲要内容开始
% 
% * 对角映射及杯积
% * 自然性、结合性与单位
% * 分次交换性的链同伦证明
% * 相对杯积及配对
% 
% 纲要内容结束

\section{杯积与环结构}
\label{g2:sec:0404}

同调外积把两个空间的类放到乘积空间中。如果两个类原来就来自同一个空间，还需要一条把它们放回该空间的自然映射。对角映射 \(x\mapsto(x,x)\) 正好提供这条映射，但它的方向是从 \(X\) 到 \(X\times X\)。因此，适合沿它拉回的对象是上同调。杯积把“同一点的两份信息同时取值”转成乘法，使上同调能够记录加性群看不见的关系。本节使用任意含幺交换环 \(R\)，上边界始终取 \(\delta\varphi=\varphi\partial\)。有关杯积及其相对形式可参阅 \cite[第3.2节]{Hatcher2002}。

\subsection{对角映射及杯积}
\label{g2:subsec:0404:01}

先把两个上链作用在张量链上的含义说清。若 \(\varphi\in C^p(X;R)\)、\(\psi\in C^q(Y;R)\)，定义次数 \(p+q\) 的线性泛函
\[
	(\varphi\otimes\psi)(a\otimes b)=\varphi(a)\psi(b)
	\quad(|a|=p,\ |b|=q),
\]
并在其他双次数上取零。这只是一个明确的泛函定义，不要求
\(\operatorname{Hom}(C_p,R)\otimes\operatorname{Hom}(C_q,R)\to
\operatorname{Hom}(C_p\otimes C_q,R)\) 为同构。按上一节的总微分，
\begin{equation}\label{g2:eq:tensor-cochain-coboundary}
	\delta(\varphi\otimes\psi)
	=\delta\varphi\otimes\psi+(-1)^p\varphi\otimes\delta\psi.
\end{equation}
所需符号由张量总微分给出。

\begin{definition}[杯积]\label{g2:def:cup-product}
	设 \(X\) 为拓扑空间，\(R\) 为含幺交换环。令 \(\Delta_X:X\to X\times X\) 为对角映射，记
	\[
		D_X=\operatorname{AW}\circ(\Delta_X)_\#
		:C_*(X;R)\longrightarrow C_*(X;R)\otimes_R C_*(X;R).
	\]
	对 \(\varphi\in C^p(X;R)\)、\(\psi\in C^q(X;R)\)，定义其\term[beiji]{杯积}[cup product]为
	\[
		\varphi\smile\psi=(\varphi\otimes\psi)D_X.
	\]
	等价地，对 \(\sigma:\Delta^{p+q}\to X\)，
	\begin{equation}\label{g2:eq:cup-front-back}
		(\varphi\smile\psi)(\sigma)
		=\varphi\bigl(\sigma[0,\ldots,p]\bigr)
			\psi\bigl(\sigma[p,\ldots,p+q]\bigr).
	\end{equation}
\end{definition}

切分点的位置由两个上链的次数决定：前 \(p\) 维面给第一个上链，后 \(q\) 维面给第二个。三角形上的一次上链杯积只使用边 \(01\) 与边 \(12\)；边 \(02\) 不直接出现在这条取值公式中，却通过上链的闭性参与不同代表之间的比较。这种不对称属于链层面的选择，稍后会证明它在上同调中只留下分次交换符号。\symidx{\(\smile\)：上链与上同调的杯积}

\begin{proposition}[杯积的上边界]\label{g2:prop:cup-leibniz}
	对拓扑空间 \(X\)、含幺交换环 \(R\) 以及 \(\varphi\in C^p(X;R)\)、\(\psi\in C^q(X;R)\)，有
	\[
		\delta(\varphi\smile\psi)
		=\delta\varphi\smile\psi+(-1)^p\varphi\smile\delta\psi.
	\]
	因此杯积诱导双线性映射
	\[
		H^p(X;R)\otimes_R H^q(X;R)\longrightarrow H^{p+q}(X;R),
		\qquad
		[\varphi]\otimes[\psi]\longmapsto[\varphi\smile\psi].
	\]
\end{proposition}
\begin{proof}
	由\cref{g2:prop:aw-shuffle-chain-maps}，\(D_X\) 是链映射。将\eqref{g2:eq:tensor-cochain-coboundary}与 \(D_X\) 复合，即得所述等式。如果 \(\delta\varphi=\delta\psi=0\)，其杯积也闭；若将 \(\varphi\) 改成 \(\varphi+\delta u\)，杯积的变化为 \(\delta(u\smile\psi)\)。若将 \(\psi\) 改成 \(\psi+\delta v\)，变化为 \((-1)^p\delta(\varphi\smile v)\)。所以所得上同调类只依赖两个原类。
\end{proof}

\subsection{自然性、结合性与单位}
\label{g2:subsec:0404:02}

杯积的边界公式使它下降到上同调。接下来比较拉回、三次相乘和常值零次上链与这项运算的关系：拉回保持单形的各个面，三次相乘按两个切分点分割单形，常值 \(1\) 则在零维端面上给出单位元。

\begin{theorem}[杯积的环性质]\label{g2:thm:cup-properties}
	设 \(R\) 为含幺交换环，\(X\) 为拓扑空间。杯积使
	\[
		H^*(X;R)\coloneqq\bigoplus_{n\ge0}H^n(X;R)
	\]
	成为含幺分次 \(R\) 代数，乘法把次数 \(p,q\) 相加。其单位是把每个点取值为 \(1\) 的零次上圈所代表的类。对连续映射 \(f:X\to Y\)，
	\[
		f^*(\alpha\smile\beta)=f^*\alpha\smile f^*\beta,
		\qquad f^*(1_Y)=1_X.
	\]
	特别地，同伦等价诱导上同调环同构。
\end{theorem}
\begin{proof}
	双线性已经证明。若 \(\varphi,\psi,\eta\) 的次数分别为 \(p,q,r\)，则两种结合方式在任意 \((p+q+r)\) 单形 \(\sigma\) 上都给出
	\[
		\varphi(\sigma[0,\ldots,p])\,
		\psi(\sigma[p,\ldots,p+q])\,
		\eta(\sigma[p+q,\ldots,p+q+r]).
	\]
	因此结合律在上链上就严格成立。令 \(\mathbf1\) 在每个零单形上取值 \(1\)，则 \((\delta\mathbf1)(\gamma)=1-1=0\)；前后面公式分别取切分位置 \(0\) 和 \(p\)，直接得到
	\[
		\mathbf1\smile\varphi=\varphi=\varphi\smile\mathbf1.
	\]
	这给出单位。空空间的上同调为零，按其零环的通常约定处理。

	连续映射与取面交换，故
	\[
		\bigl(f^*(\varphi\smile\psi)\bigr)(\sigma)
		=\varphi(f\sigma[0,\ldots,p])\psi(f\sigma[p,\ldots,p+q])
		=\bigl(f^*\varphi\smile f^*\psi\bigr)(\sigma).
	\]
	单位也保持。最后，\cref{g2:thm:singular-homotopy-invariance}的棱柱链同伦经 \(\operatorname{Hom}(-,R)\) 给出上链同伦：若 \(g_\#-f_\#=\partial P+P\partial\)，则对上圈 \(\varphi\) 有 \(g^*\varphi-f^*\varphi=\delta(\varphi P)\)。故同伦映射诱导相同的上同调映射；把它用于同伦等价及其逆，即得环同构。
\end{proof}

我们称这个分次代数为 \(X\) 的\term[shangtongdiao huan]{上同调环}[cohomology ring]。零次上同调可以理解为每个道路分支上的常值函数，杯积在零次就是逐点相乘。若 \(X\) 道路连通，则 \(H^0(X;R)=R\)，其中的元素在其他次数上按通常的标量乘法作用。

\subsection{分次交换性的链同伦证明}
\label{g2:subsec:0404:03}

一般上链的杯积并不交换：在一个三角形上，先看边 \(01\) 再看边 \(12\)，与把两条边的任务互换，通常给出不同的数。要证明上同调上的交换律，必须找到两种取值之差的原上链，而不能直接交换前后面。

\begin{proposition}[分次交换性]\label{g2:prop:graded-commutativity}
	设 \(X\) 为拓扑空间，\(R\) 为含幺交换环。对 \(\alpha\in H^p(X;R)\)、\(\beta\in H^q(X;R)\)，
	\[
		\alpha\smile\beta=(-1)^{pq}\beta\smile\alpha.
	\]
	更具体地，存在自然算子 \(K:C_n(X;R)\to(C_*(X;R)\otimes C_*(X;R))_{n+1}\)，使
	\begin{equation}\label{g2:eq:diagonal-switch-homotopy}
		\tau D_X-D_X=\mathrm{d}K+K\partial,
		\qquad
		\tau(a\otimes b)=(-1)^{|a||b|}b\otimes a.
	\end{equation}
\end{proposition}
\begin{proof}
	先核对 \(\tau\) 是链映射。对 \(|a|=i\)、\(|b|=j\)，\(\mathrm{d}\tau(a\otimes b)\) 的两项系数分别为 \((-1)^{ij}\)、\((-1)^{ij+j}\)；\(\tau\mathrm{d}(a\otimes b)\) 中对应的 \(b\otimes\partial a\)、\(\partial b\otimes a\) 项系数分别为 \((-1)^{(i-1)j}\)、\((-1)^{i+i(j-1)}\)。前者与 \(ij+j\)、后者与 \(ij\) 模二相同，故两式一致。

	于是 \(D_X\) 与 \(\tau D_X\) 是两条自然链映射。它们在零次均把一个点送到该点与自身的张量。源的通用基元是 \(\operatorname{id}_{\Delta^n}\)，目标模型 \(C_*(\Delta^n;R)\otimes C_*(\Delta^n;R)\) 的增广收缩已在\cref{g2:lem:product-model-homotopy}后给出。应用该引理便得到 \(K\)。具体地，已知低维 \(K\) 后，在模型中填充圈
	\[
		(\tau D-D)\operatorname{id}_{\Delta^n}
			-K\partial\operatorname{id}_{\Delta^n},
	\]
	再沿 \(\sigma\otimes\sigma\) 推到任意奇异单形 \(\sigma\)。这给出所需的自然递归构造，而且 \(K\sigma\) 的两个因子始终留在 \(\sigma(\Delta^n)\) 内。

	取代表上圈 \(\varphi,\psi\)，并令 \(\lambda=\varphi\otimes\psi\)。由\eqref{g2:eq:tensor-cochain-coboundary}，\(\lambda\mathrm{d}=0\)。把 \(\lambda\) 作用在\eqref{g2:eq:diagonal-switch-homotopy}两边，得到
	\[
		(-1)^{pq}\psi\smile\varphi-\varphi\smile\psi
		=(\lambda K)\partial=\delta(\lambda K).
	\]
	故差为上边界，所述上同调等式成立。若 \(p=q=0\)，右侧没有负次上链，此时两种取值已由 \(R\) 的交换律严格相同。
\end{proof}

这样的环称为\term[fenci jiaohuan huan]{分次交换环}[graded-commutative ring]。若 \(\alpha\) 为奇次数类，只能立即推出 \(2\alpha^2=0\)，不能在整数系数下直接除以 \(2\)。当系数域的特征不为 \(2\) 时，才可据此断言所有奇次数类的平方为零。特征 \(2\) 时符号消失，奇次数平方则可能非零；下一节的非定向曲面会给出实例。

\subsection{相对杯积及配对}
\label{g2:subsec:0404:04}

相对上链 \(\varphi\in C^p(X,A;R)\) 可以看成在 \(A\) 中每条链上都取零的 \(X\) 上链。若另一个上链 \(\psi\) 在 \(B\) 上为零，则它们的杯积在完全落入 \(A\) 的单形上为零，也在完全落入 \(B\) 的单形上为零。不过，一条落在 \(A\cup B\) 的单形可能同时穿过两部分。这是相对杯积必须处理的额外问题。

\begin{proposition}[相对杯积]\label{g2:prop:relative-cup-product}
	设 \(X\) 为拓扑空间，\(R\) 为含幺交换环，\(A,B\subset X\) 满足
	\begin{equation}\label{g2:eq:relative-cup-excision-condition}
		A\cup B
		=\operatorname{int}_{A\cup B}A\ \cup\ \operatorname{int}_{A\cup B}B.
	\end{equation}
	则有自然的双线性乘积
	\[
		H^p(X,A;R)\otimes_R H^q(X,B;R)
		\longrightarrow H^{p+q}(X,A\cup B;R).
	\]
	它与绝对杯积及相对群到绝对群的映射相容，并满足分次交换律。特别地，\(A,B\) 都开时满足此条件。取 \(A=B\) 或其中之一为空，总可得到相对上同调的乘法以及绝对上同调对它的作用。
\end{proposition}
\begin{proof}
	先用整数链，令
	\[
		S_*=C_*(A;\mathbb Z)+C_*(B;\mathbb Z)
		\subset C_*(A\cup B;\mathbb Z).
	\]
	由\eqref{g2:eq:relative-cup-excision-condition}和\cref{g2:thm:small-chain-isomorphism}，包含 \(S_*\to C_*(A\cup B;\mathbb Z)\) 诱导同调同构。商复形
	\[
		Q_*=C_*(A\cup B;\mathbb Z)/S_*
	\]
	因此无同调。每个 \(Q_n\) 仍为自由阿贝尔群，以不完全落入 \(A\) 或 \(B\) 的单形为基；其圈群和边界群也自由。\cref{g2:lem:free-complex-splitting}在同调全为零时给出的两项分解，说明 \(Q_*\) 是可缩复形。于是张量 \(R\) 后仍可缩，再取 \(\operatorname{Hom}_R(-,R)\) 也可缩。

	自然满射
	\[
		q:C_*(X;R)/(C_*(A;R)+C_*(B;R))
		\longrightarrow C_*(X,A\cup B;R)
	\]
	的核就是 \(Q_*\otimes_{\mathbb Z}R\)。两侧链模都有由单形得到的自由基，故对应短正合列逐次分裂；对其取 \(\operatorname{Hom}_R(-,R)\) 后仍正合，核复形的可缩性说明 \(q^*\) 诱导上同调同构。

	现在取相对上圈 \(\varphi,\psi\)。前后面公式说明 \(\varphi\smile\psi\) 在 \(C_*(A)+C_*(B)\) 上取零，故给出 \(q\) 的源上的上圈类；经上述 \(q^*\) 的逆同构，得到 \(H^{p+q}(X,A\cup B;R)\) 中的类。更换代表时，\cref{g2:prop:cup-leibniz}中的上边界公式仍在这个商复形上成立，所以乘积良定义。连续映射保持相应子空间时，杯积和 \(q\) 都与之交换；同构的逆也自然，因此得到相对自然性。

	\cref{g2:prop:graded-commutativity}构造的 \(K\sigma\) 的两个因子都留在 \(\sigma\) 的像内。若 \(\sigma\) 落在 \(A\)，\(\varphi\) 使其取值为零；若落在 \(B\)，\(\psi\) 使其取值为零。因此 \((\varphi\otimes\psi)K\) 也在 \(S_*\) 上为零，原来的交换同伦下降到同一个商复形，证明相对分次交换律。与绝对杯积的相容性由各商映射与前后面公式直接得到。
\end{proof}

当 \(A,B,C\) 都为开集时，可以继续相乘并得到 \(H^*(X,A\cup B\cup C;R)\) 中的结合律：三重上链乘积在 \(C_*(A)+C_*(B)+C_*(C)\) 上取零，小链比较与刚才相同，而两种结合方式在这个商上已经相等。这个说法也适用于任意有限个开子空间。不满足小链比较条件的一般子空间对，不能仅凭前后面公式就把值域写成 \(H^*(X,A\cup B;R)\)。

\begin{definition}[取值配对与上同调外积]\label{g2:def:kronecker-cross-cohomology}
	设 \(R\) 为含幺交换环。相对上圈 \(\varphi\) 与同次数相对圈 \(z\) 的\term[quzhi peidui]{取值配对}[evaluation pairing]定义为
	\[
		\langle[\varphi],[z]\rangle=\varphi(z)\in R.
	\]
	对 \(\alpha\in H^p(X;R)\)、\(\beta\in H^q(Y;R)\)，定义\term[shangtongdiao waiji]{上同调外积}[cohomology cross product]为
	\[
		\alpha\times\beta
		=\operatorname{pr}_X^*\alpha\smile\operatorname{pr}_Y^*\beta
		\in H^{p+q}(X\times Y;R).
	\]
	若 \(A\subset X\)、\(B\subset Y\) 为开子空间，同一公式利用相对杯积给出
	\[
		H^p(X,A;R)\otimes_R H^q(Y,B;R)
		\longrightarrow
		H^{p+q}(X\times Y,A\times Y\cup X\times B;R).
	\]
\end{definition}

配对与代表无关：上边界在圈上为零，上圈在边界上为零；相对情形中，剩余链落在子空间，仍被相对上链消去。配对不总能看见所有上同调类。例如 \(H^2(\mathbb{RP}^2;\mathbb Z)=\mathbb Z/2\) 而 \(H_2(\mathbb{RP}^2;\mathbb Z)=0\)，所以这次配对恒为零。万有系数定理中的 Ext 项正是其中的一种不可见信息。\symidx{\(\langle\alpha,z\rangle\)：上同调类与同调类的取值配对}

\begin{proposition}[外积的取值]\label{g2:prop:cross-cup-evaluation}
	设 \(R\) 为含幺交换环，\(\alpha\in H^p(X;R)\)、\(\beta\in H^q(Y;R)\)、\(a\in H_p(X;R)\)、\(b\in H_q(Y;R)\)。则
	\[
		\langle\alpha\times\beta,a\times b\rangle
		=\langle\alpha,a\rangle\langle\beta,b\rangle.
	\]
	若因子配对的次数不匹配而总次数相等，则取值为零。对开子空间 \(A\subset X\)、\(B\subset Y\)，同一结论适用于相对外积和相对配对。另有
	\[
		\Delta_X^*(\alpha\times\beta)=\alpha\smile\beta,
	\]
	以及外积的乘法规则
	\[
		(\alpha\times\beta)\smile(\alpha'\times\beta')
		=(-1)^{|\beta||\alpha'|}
			(\alpha\smile\alpha')\times(\beta\smile\beta').
	\]
\end{proposition}
\begin{proof}
	取代表上圈 \(\varphi,\psi\)。上同调外积的上链代表恰为 \((\varphi\otimes\psi)\operatorname{AW}\)，这可直接对乘积中的单形应用前后面公式。于是
	\[
		\langle\alpha\times\beta,a\times b\rangle
		=(\varphi\otimes\psi)\operatorname{AW}\operatorname{sh}(z\otimes w).
	\]
	由\cref{g2:thm:eilenberg-zilber}，\(\operatorname{AW}\operatorname{sh}\) 与恒等链同伦；\(z\otimes w\) 是圈，\(\varphi\otimes\psi\) 是上圈，所以右侧等于 \(\varphi(z)\psi(w)\)。双次数不匹配时，泛函的定义使它为零。

	相对情形还要核对模型。令 \(Z=A\times Y\cup X\times B\)。两条链映射及其自然链同伦分别保持来自 \(A\times Y\)、\(X\times B\) 的链，因而下降为
	\[
		\frac{C_*(X\times Y;R)}{C_*(A\times Y;R)+C_*(X\times B;R)}
		\ \simeq_{\mathrm{ch}}\
		C_*(X,A;R)\otimes_R C_*(Y,B;R).
	\]
	右侧的商描述由各奇异单形基逐项得到。由于 \(A,B\) 开，左侧分母是 \(Z\) 的开覆盖小链复形。\cref{g2:prop:relative-cup-product}证明中的可缩商论证，给出左侧与 \(C_*(X\times Y,Z;R)\) 的自然同调、上同调比较同构；通过它定义相对同调外积。杯积定义的相对上同调外积在小链商上仍由 \((\varphi\otimes\psi)\operatorname{AW}\) 表示，因此完全相同的链同伦取值证明适用。

	对角公式来自 \(\operatorname{pr}_X\Delta_X=\operatorname{pr}_Y\Delta_X=\operatorname{id}_X\) 及自然性。最后，把四个拉回类按 \(\alpha,\beta,\alpha',\beta'\) 的顺序相乘，只需把 \(\beta\) 与 \(\alpha'\) 交换一次；\cref{g2:prop:graded-commutativity}给出系数 \((-1)^{|\beta||\alpha'|}\)，其余由结合律完成。
\end{proof}

这个配对公式还给出常用的上同调 Künneth 形式。若 \(k\) 为域，且对每个 \(p\)，\(H_p(X;k)\) 有限维，则
\[
	\bigoplus_{p+q=n}H^p(X;k)\otimes_k H^q(Y;k)
	\xrightarrow{\ \times\ }H^n(X\times Y;k)
\]
为自然同构。理由是：域上的万有系数同构把上同调识别为同调的线性对偶，\cref{g2:thm:kunneth}计算乘积同调，而有限维条件保证
\[
	\operatorname{Hom}_k(U\otimes V,k)
	\cong\operatorname{Hom}_k(U,k)\otimes\operatorname{Hom}_k(V,k)
	\quad(\dim U<\infty).
\]
选 \(U\) 的有限基及对偶基，就可把任意双线性泛函写成这类张量的有限和；这也证明所用的代数事实。外积取值公式确认所得同构正是上面显示的映射。若 \(R\) 为主理想整环，两个空间各次同调都为有限秩自由 \(R\) 模，同一论证也适用；Ext、Tor 项在此都为零。其分次环的乘法必须采用命题中的符号规则。

\begin{proposition}[局部生成元的乘积]\label{g2:prop:local-cross-normalization}
	设 \(p,q\ge1\)，在 \(\mathbb R^p,\mathbb R^q\) 上取标准坐标定向，在 \(\mathbb R^{p+q}=\mathbb R^p\times\mathbb R^q\) 上先列第一因子坐标。令
	\[
		u_p\in H^p(\mathbb R^p,\mathbb R^p\setminus\{0\};\mathbb Z),
		\qquad
		u_q\in H^q(\mathbb R^q,\mathbb R^q\setminus\{0\};\mathbb Z)
	\]
	分别在正向局部同调生成元上取值 \(1\)。则相对外积 \(u_p\times u_q\) 是
	\[
		H^{p+q}(\mathbb R^{p+q},\mathbb R^{p+q}\setminus\{0\};\mathbb Z)
	\]
	的正向生成元。这个结论可迭代到任意有限个因子，也适用于原点的球形或立方体邻域。
\end{proposition}
\begin{proof}
	\cref{g2:thm:manifold-local-homology}把这两个局部同调群分别识别为次数 \(p,q\) 的 \(\mathbb Z\)，其他次数为零。万有系数定理因而使相应上同调与同调的对偶一致，故在正向生成元上取值 \(1\) 唯一确定 \(u_p,u_q\)。

	各取一个以原点为内点的小闭立方体，按标准坐标定向三角剖分。把各最高维单形带定向求和，得到链 \(c_p,c_q\)。在“立方体模边界”的相对胞腔链中，这个和代表正向最高维胞腔的生成元；径向投影把穿孔立方体形变收缩到边界，所以由空间对的长正合列和局部切除，它也代表正向局部类。其 shuffle 链是两个剖分的乘积剖分：每个单形乘积按格路径切分，\eqref{g2:eq:shuffle-map}中的符号使所有块带标准乘积定向，内部面在相邻块之间抵消。故 \(\operatorname{sh}(c_p\otimes c_q)\) 正是乘积立方体的正向基本链，代表局部生成元 \(c_{p+q}\)。特别地，这里给出的是正号，而不只是一个未定号的生成元。

	两个穿孔子空间都开，且
	\[
		(\mathbb R^p\setminus\{0\})\times\mathbb R^q
		\ \cup\ \mathbb R^p\times(\mathbb R^q\setminus\{0\})
		=\mathbb R^{p+q}\setminus\{0\}.
	\]
	因此可用\cref{g2:prop:cross-cup-evaluation}的相对形式，得到
	\[
		\langle u_p\times u_q,c_{p+q}\rangle
		=\langle u_p,c_p\rangle\langle u_q,c_q\rangle=1.
	\]
	这证明了归一化。反复应用即得多因子结论；限制到包含原点的球或立方体邻域时，切除给出同一局部群，正向基本链的选取不变。
\end{proof}

例如在 \(\mathbb C^n\) 中采用 \((x_1,y_1,\ldots,x_n,y_n)\) 的实坐标次序，每个坐标复平面的正向局部二次类相乘，给出原点的正向 \(2n\) 次类。相对杯积因此能够把若干坐标超平面的局部条件集中到它们的交点；下一节计算复射影空间的乘法时，正是用这个可直接核验的局部事实来确定系数。

% !TeX root = ../../Book2.tex
% 纲要标题：### 4.5 上同调环的计算与比较
% 纲要位置：几何学计划纲要.md，第 1183 行。
% 正文状态：初稿；尚未正式审核。

% 纲要内容开始
% 
% * 环面和闭曲面的上同调环
% * 复射影空间的生成元和关系
% * 相同加性同调、不同乘法结构的模型
% 
% ---
% 
% 纲要内容结束

\section{上同调环的计算与比较}
\label{g2:sec:0405}

上一章已经算出环面、闭曲面和射影空间的同调群。现在的问题是：两个低次上同调类能否相乘成为高次的非零类？加性群只告诉我们每个次数有哪些元素，杯积还记录这些元素怎样共同作用于空间。下面先从坐标投影计算环面，再用连通和的压缩映射处理闭曲面；复射影空间的关键则是把若干个相对上同调类集中到同一个坐标邻域中。本节写 \(\mathbb F_2=\mathbb Z/2\mathbb Z\)，所有环同构均保留次数和单位元。

\subsection{环面和闭曲面的上同调环}
\label{g2:subsec:0405:01}
取圆周的一次整数上同调生成元 \(\xi\)，使它在正向绕行的圈上取值 \(1\)。在环面 \(T^m=(S^1)^m\) 上，第 \(i\) 个坐标投影 \(p_i\) 给出类 \(x_i=p_i^*\xi\)。这些类分别来自各个圆周因子，它们的乘积可用前两节的外积及其取值公式计算：各因子的一次生成元相乘，检测相应坐标子环面的基本圈。

\begin{proposition}[环面的上同调环]\label{g2:prop:torus-cohomology-ring}
	设 \(m\geq1\)。整数系数下，\(T^m=(S^1)^m\) 的上同调环为
	\[
		H^*(T^m;\mathbb Z)\cong\Lambda_{\mathbb Z}(x_1,\ldots,x_m),\qquad |x_i|=1.
	\]
	右端是由次数为一的生成元及关系 \(x_i^2=0\)、\(x_ix_j=-x_jx_i\) 构成的\term[waidaishu]{外代数}[exterior algebra]；每个次数 \(k\) 的加性基由 \(x_{i_1}\cdots x_{i_k}\)、\(i_1<\cdots<i_k\) 给出。对任意交换含幺环 \(R\)，相同陈述成立，只需将 \(\mathbb Z\) 换成 \(R\)。
\end{proposition}
\begin{proof}
	先取整数系数。圆周的整数同调仅在次数 \(0,1\) 为 \(\mathbb Z\)。由\cref{g2:thm:eilenberg-zilber,g2:thm:kunneth}，反复对圆周作乘积时，外积给出的类在每个次数构成基；这里所有参与的群均为自由群，没有额外的 \(\operatorname{Tor}\) 项。\cref{g2:prop:cross-cup-evaluation}将相应上同调外积识别为上述坐标类的杯积，并证明它们与对应乘积圈互为对偶基。

	因为 \(\xi^2\in H^2(S^1;\mathbb Z)=0\)，自然性给出 \(x_i^2=p_i^*(\xi^2)=0\)。对两个不同的坐标，\cref{g2:prop:graded-commutativity}给出 \(x_ix_j=-x_jx_i\)。这些关系把每个单项式化为一个严格递增指标的单项式，或化为零；而外积已经证明所有递增单项式线性无关，所以没有其他关系。

	对一般 \(R\)，把它先视为阿贝尔群。由于环面的各次整数同调有限生成且自由，\cref{g2:thm:uct-cohomology}说明整数上同调基的系数变换给出 \(H^k(T^m;R)\) 的 \(R\) 基。系数变换与杯积相容，所以刚才的生成元及全部乘法关系随之成立。在特征 \(2\) 时，负号虽然消失，\(x_i^2=0\) 仍是独立的关系，不能只从交换性中推出。
\end{proof}

对于 \(T^2\)，因此有 \(x^2=y^2=0\)，而 \(xy\) 是二次生成元。两个一次类分别测量两个方向的绕行，它们的乘积能够检测整个二维胞腔。若改成两个圆周与一个球面的楔和，虽然各次加性群一样，两个一次类却不能这样相乘；本节末尾会明确比较这两个空间。

\ProseParagraphBreak

为了把环面的计算搬到多柄曲面上，先说明楔和与连通和之间所用的映射。若 \(X_1,\ldots,X_s\) 是连通 CW 复形，各以一个零胞腔为基点，则有限楔和 \(W=\bigvee_iX_i\) 的约化胞腔链分解为各因子的直和。经胞腔比较及自然的万有系数正合列，限制映射给出
\begin{equation}\label{g2:eq:wedge-positive-cohomology}
	\widetilde H^k(W;R)\xrightarrow{\ \cong\ }\bigoplus_i\widetilde H^k(X_i;R),\qquad k>0.
\end{equation}
这里的乘法按分量进行：来自不同因子的两个正次类，其乘积限制到每个因子上都为零，由该同构的单射性，乘积本身就是零。零次则只有一个共同的单位元，而不是每个因子各保留一个单位元。

\begin{lemma}[连通和的压缩映射]\label{g2:lem:surface-pinch-cohomology}
	设 \(M=M_1\mathbin{\#}\cdots\mathbin{\#}M_s\) 是 \(s\geq1\) 个连通闭曲面模型的连通和。若各 \(M_i\) 可定向，给它们选择与连通和相容的定向并取 \(R=\mathbb Z\)；若不要求可定向，则取 \(R=\mathbb F_2\)。存在连续映射 \(q:M\to\bigvee_iM_i\)，使一次拉回
	\[
		q^*:H^1\Bigl(\bigvee_iM_i;R\Bigr)\longrightarrow H^1(M;R)
	\]
	为同构。若 \(\omega_i\in H^2(M_i;R)\) 是对应定向的生成元，模 \(2\) 时取唯一非零类，则各 \(\omega_i\) 在楔和中的延拓均被 \(q^*\) 送到同一个二次生成元 \(\omega\in H^2(M;R)\)。
\end{lemma}
\begin{proof}
	先考虑两个曲面的连通和。将粘合用的分离圆周压成一点，左右两块各自成为 \((M_i\setminus\operatorname{int}D^2)/\partial D^2\)。这与 \(M_i\) 同胚：边界的圆环领圈在内边界压缩后变成一个圆盘，外边界上的粘合保持不变。于是得到到 \(M_1\vee M_2\) 的映射。多次作此操作，得到各因子由若干基点连成树的空间；再压掉连接这些基点的树，可把它们归并为一个楔点。等价地，可在原曲面的各连接圆周之间选互不交叉的连接弧，一并压缩。这个构造在每个因子内保留一组标准边环路。

	由\eqref{g2:eq:surface-attaching-words}，可定向曲面的整数一维同调以各柄的两条边为基；模 \(2\) 时，无论可否定向，各条标准边也独立生成一维同调，因为二维胞腔的边界系数均为偶数。\(q_*\) 将这些边送到相应因子的同名边，故在一次同调上是同构。\cref{g2:thm:uct-cohomology}在一次给出自然同构 \(H^1(-;R)\cong\operatorname{Hom}_{\mathbb Z}(H_1(-;\mathbb Z),R)\)，因为零次同调是自由群；这就证明了关于 \(q^*\) 的断言。

	还要确定二维生成元的倍数。将多边形模型细分为与连接圆周及连接弧相容的有限三角剖分；可先切分各边和各多边形，再作重心细分，使粘合后的三角形有一致的顶点次序。可定向时按曲面定向给各三角形取符号，模 \(2\) 时全部取系数 \(1\)。这些三角形之和 \(z\) 是奇异二圈：内部边出现两次并抵消，原多边形的成对边也按所取系数抵消。相对化到原二维胞腔时，这些小三角片各出现一次并铺满特征圆盘，所以其胞腔系数为 \(1\)；胞腔比较定理因而将 \([z]\) 识别为二维生成元。把 \(q\) 再与到第 \(i\) 个因子的收缩复合，记为 \(q_i:M\to M_i\)。在该因子内选一个避开压缩集合的三角形内点 \(y\)，其原像只有一个点，且 \(q_i\) 在此处保持定向地为局部同胚；模 \(2\) 时无须选择符号。

	映到局部同调 \(H_2(M_i,M_i\setminus\{y\};R)\) 后，\((q_i)_*[z]\) 的系数因此为 \(1\)。同一个三角形和构造说明 \(M_i\) 的二维生成元也映到局部生成元：切除到该三角形内部，只剩包含 \(y\) 的那一片，边界定向与圆盘相对生成元一致。由于两边的二次群均为秩一的自由 \(R\) 模，\((q_i)_*[z]\) 恰为所选生成元。最后用取值配对的自然性得到 \(\langle q_i^*\omega_i,[z]\rangle=1\)。可定向情形下一维整数同调自由，模 \(2\) 情形下系数是域，因此相应的上同调取值同构保证 \(q_i^*\omega_i=\omega\)。
\end{proof}

\begin{theorem}[可定向曲面的上同调环]\label{g2:thm:orientable-surface-ring}
	设 \(\Sigma_g\) 是 \(g\geq1\) 个环面的连通和，并选定定向。存在一次类 \(a_1,b_1,\ldots,a_g,b_g\) 和二次类 \(\omega\)，使
	\[
		H^0(\Sigma_g;\mathbb Z)=\mathbb Z,
		\quad H^1(\Sigma_g;\mathbb Z)=\bigoplus_{i=1}^g(\mathbb Za_i\oplus\mathbb Zb_i),
		\quad H^2(\Sigma_g;\mathbb Z)=\mathbb Z\omega.
	\]
	所有更高次群为零，正次生成元的乘法由
	\begin{equation}\label{g2:eq:orientable-surface-products}
		a_i a_j=b_i b_j=0,\qquad a_i b_j=\delta_{ij}\omega,\qquad b_j a_i=-\delta_{ij}\omega
	\end{equation}
	完全决定。\(\Sigma_0=S^2\) 的环为 \(\mathbb Z[\omega]/(\omega^2)\)，其中 \(|\omega|=2\)。
\end{theorem}
\begin{proof}
	取\cref{g2:lem:surface-pinch-cohomology}的映射 \(q:\Sigma_g\to\bigvee_iT^2\)。在每个环面上选择\cref{g2:prop:torus-cohomology-ring}的一次基 \(x_i,y_i\)，使 \(x_i y_i=\omega_i\)，并通过收缩其他因子把这些类延拓到楔和。令 \(a_i=q^*x_i\)、\(b_i=q^*y_i\)。同因子中的平方为零，异因子之间的全部正次乘积为零；而 \(q^*\omega_i=\omega\)。杯积自然性于是给出全部关系。加性群已由\cref{g2:prop:surface-integer-homology}和万有系数定理确定，所以这些关系没有遗漏其他乘法。\(S^2\) 的结论直接来自它的上同调只在次数 \(0,2\) 非零。
\end{proof}

不可定向曲面宜先用模 \(2\) 系数观察乘法。这样即使一条边沿同一个方向出现两次，它在边界中仍会抵消。与可定向曲面的主要区别不是交换律中的负号，而是某个一次类的平方可以非零。我们先把这个现象在 \(\mathbb RP^2\) 上算清楚。

\begin{lemma}[射影平面的一次平方]\label{g2:lem:projective-plane-square}
	设 \(t\) 是 \(H^1(\mathbb RP^2;\mathbb F_2)\) 的非零类，则 \(t^2\) 是 \(H^2(\mathbb RP^2;\mathbb F_2)\) 的非零类。因此
	\[
		H^*(\mathbb RP^2;\mathbb F_2)\cong\mathbb F_2[t]/(t^3),\qquad |t|=1.
	\]
\end{lemma}
\begin{proof}
	考虑系数短正合列
	\[
		0\longrightarrow\mathbb F_2\xrightarrow{\ \iota\ }\mathbb Z/4\xrightarrow{\ \rho\ }\mathbb F_2\longrightarrow0,
		\qquad\iota(1)=2,
	\]
	其中 \(\rho\) 为模 \(2\) 映射。奇异链群是自由阿贝尔群，所以对每个奇异单形逐一提升上链的值，得到上链复形的短正合列及连接同态 \(\beta:H^1(X;\mathbb F_2)\to H^2(X;\mathbb F_2)\)。它的具体算法是：把模 \(2\) 闭上链 \(c\) 的每个值提升为 \(0\) 或 \(1\)，得到模 \(4\) 上链 \(\widetilde c\)；其上边界取值于 \(\{0,2\}\)，除以 \(2\) 再模 \(2\)，所得闭上链代表 \(\beta[c]\)。

	在一个奇异二单形 \([v_0v_1v_2]\) 上，记 \(c[v_0v_1]=r\)、\(c[v_1v_2]=s\)，其中 \(r,s\in\{0,1\}\)。闭性条件给出 \(c[v_0v_2]\equiv r+s\pmod2\)，因此
	\[
		\frac{(\delta\widetilde c)[v_0v_1v_2]}2
		=\frac{s-\bigl((r+s)\bmod2\bigr)+r}{2}
		\equiv rs\pmod2.
	\]
	右端正是\cref{g2:def:cup-product}的 \((c\smile c)[v_0v_1v_2]\)。于是对任意空间，\(\beta[c]=[c]^2\) 在一次都成立。

	对 \(X=\mathbb RP^2\)，\cref{g2:prop:projective-integer-homology}给出 \(H_1(X;\mathbb Z)=\mathbb Z/2\)。一次万有系数同构及其自然性说明 \(H^1(X;\mathbb Z/4)\to H^1(X;\mathbb F_2)\) 对应 \(\operatorname{Hom}(\mathbb Z/2,\mathbb Z/4)\to\operatorname{Hom}(\mathbb Z/2,\mathbb F_2)\)。前一群的同态只能把生成元送到 \(0\) 或 \(2\)，经 \(\rho\) 都变成零。系数长正合列因而说明 \(\beta\) 在一次单射。由\cref{g2:cor:cellular-coefficients}，模 \(2\) 胞腔微分全为零；域系数的上同调取值同构因而说明 \(H^1(X;\mathbb F_2)\) 与 \(H^2(X;\mathbb F_2)\) 都是一维；于是 \(t^2=\beta(t)\ne0\)。高于二次的群为零，得到所述环。
\end{proof}

\begin{proposition}[不可定向曲面的模二上同调环]\label{g2:prop:nonorientable-surface-ring}
	设 \(N_h\) 是 \(h\geq1\) 个实射影平面的连通和。存在一次基 \(c_1,\ldots,c_h\) 及二次生成元 \(\omega\)，使
	\[
		H^0(N_h;\mathbb F_2)=\mathbb F_2,\qquad
		H^1(N_h;\mathbb F_2)=\mathbb F_2^h,\qquad
		H^2(N_h;\mathbb F_2)=\mathbb F_2\omega,
	\]
	更高次群为零，而乘法为
	\begin{equation}\label{g2:eq:nonorientable-surface-products}
		c_i c_j=0\quad(i\ne j),\qquad c_i^2=\omega\quad(1\leq i\leq h).
	\end{equation}
\end{proposition}
\begin{proof}
	\cref{g2:lem:surface-pinch-cohomology}给出 \(q:N_h\to\bigvee_i\mathbb RP^2\)，且一次拉回同构，二次拉回将各因子的非零类送到共同的 \(\omega\)。把各因子的一次类记为 \(t_i\)，令 \(c_i=q^*t_i\)。由\cref{g2:lem:projective-plane-square}，\(c_i^2=q^*(t_i^2)=\omega\)；异因子的乘积由\eqref{g2:eq:wedge-positive-cohomology}为零。这给出了全部一次乘法，而更高正次乘积因次数超过 \(2\) 为零。
\end{proof}

这里的系数必须明确。整数系数下，\(H^2(N_h;\mathbb Z)=\mathbb Z/2\) 来自万有系数正合列中的 \(\operatorname{Ext}\) 项；模 \(2\) 系数下，顶次类又能通过三角形和来评价。它们不是可以不加说明地互换的同一个计算。特别地，模 \(2\) 交换律并不迫使奇次平方为零，\(c_i^2=\omega\) 正说明了这一点。

\subsection{复射影空间的生成元和关系}
\label{g2:subsec:0405:02}
\cref{g2:prop:projective-integer-homology}已经给出 \(\mathbb CP^n\) 的整数同调：每个偶数次数 \(0,2,\ldots,2n\) 各有一个 \(\mathbb Z\)，其余为零。由于这些群自由，万有系数定理给出完全相同的上同调加性群。但是，仅凭这个结果还不能断言二次生成元 \(u\) 的平方是四次生成元：它可能是零，也可能是该生成元的某个非单位整数倍。下面的证明正是为了确定这个整数倍。

\begin{theorem}[复射影空间的上同调环]\label{g2:thm:complex-projective-ring}
	对整数 \(n\geq1\)，若 \(u\) 是 \(H^2(\mathbb CP^n;\mathbb Z)\) 的任一生成元，则
	\[
		H^*(\mathbb CP^n;\mathbb Z)\cong\mathbb Z[u]/(u^{n+1}),\qquad |u|=2.
	\]
	标准包含 \(\mathbb CP^k\hookrightarrow\mathbb CP^n\) 在 \(0\leq k\leq n\) 时将该环限制为截断到次数 \(2k\) 的环。赋予标准 CW 弱拓扑的无限复射影空间满足
	\[
		H^*(\mathbb CP^\infty;\mathbb Z)\cong\mathbb Z[u],\qquad |u|=2.
	\]
\end{theorem}
\begin{proof}
	先证有限维情形。记 \(X=\mathbb CP^n\)，在齐次坐标 \([z_0:\cdots:z_n]\) 中取点 \(p=[1:0:\cdots:0]\) 及 \(n\) 个超平面
	\[
		H_j=\{z_j=0\},\qquad A_j=X\setminus H_j=\{z_j\ne0\},\qquad 1\leq j\leq n.
	\]
	每个 \(A_j\) 都是仿射坐标图 \(\mathbb C^n\)，因而可缩。\cref{g2:prop:relative-cohomology-exactness}的空间对长正合列说明忘去相对条件的映射 \(H^2(X,A_j;\mathbb Z)\to H^2(X;\mathbb Z)\) 是同构，因为它两侧的 \(H^1(A_j;\mathbb Z)\)、\(H^2(A_j;\mathbb Z)\) 都为零。于是 \(u\) 有唯一相对提升 \(\widetilde u_j\)。这些提升可以理解为：计算 \(u\) 时，把超平面 \(H_j\) 之外的部分作为可以忽略的区域。

	相对杯积的目标由被忽略区域的并集决定。所有 \(A_j\) 都开，且
	\[
		A_1\cup\cdots\cup A_n=X\setminus\{p\}.
	\]
	由\cref{g2:prop:relative-cup-product}，迭代相对杯积得到
	\[
		v=\widetilde u_1\smile\cdots\smile\widetilde u_n
		\in H^{2n}(X,X\setminus\{p\};\mathbb Z),
	\]
	它在绝对上同调中的像为 \(u^n\)。证明现在分成两件事：\(v\) 是局部群的生成元，而且这个局部群到绝对顶次上同调的映射是同构。

	取含 \(p\) 的仿射图 \(U=\{z_0\ne0\}\)，以 \(w_j=z_j/z_0\) 识别为 \(\mathbb C^n\)。在这里，\(A_j\cap U=\{w_j\ne0\}\)。要知道 \(\widetilde u_j\) 限制后的整数倍数，考虑坐标直线
	\[
		L_j=\{[z_0:0:\cdots:z_j:\cdots:0]\}\cong\mathbb CP^1.
	\]
	标准胞腔包含在二次同调上为同构，故 \(H^2(X;\mathbb Z)\to H^2(L_j;\mathbb Z)\) 也是同构。对于非第一条坐标直线，交换坐标给出同样的结论；这种交换是 \(X\) 的同胚，并把第一条坐标直线送到 \(L_j\)。又因 \(A_j\cap L_j=L_j\setminus\{p\}\cong\mathbb C\)，下图两条水平箭头和右侧箭头均为同构：
	\[
		\begin{tikzcd}[column sep=large]
			H^2(X,A_j;\mathbb Z) \arrow[r] \arrow[d] & H^2(X;\mathbb Z) \arrow[d] \\
			H^2(L_j,L_j\setminus\{p\};\mathbb Z) \arrow[r] & H^2(L_j;\mathbb Z).
		\end{tikzcd}
	\]
	因此左侧箭头也为同构，\(\widetilde u_j\) 限制到这条直线的局部类是生成元。再切除到 \(L_j\cap U\)，得到 \(H^2(\mathbb C,\mathbb C\setminus\{0\};\mathbb Z)\) 的生成元。另一方面，空间对 \((U,A_j\cap U)\) 沿其他 \(n-1\) 个坐标同时收缩到这条坐标轴，故它的相对二次上同调也由该限制完全确定。若 \(\eta_j\) 表示 \(\mathbb C\) 在原点处的一个相对生成元，\(\pi_j:U\to\mathbb C\) 为坐标投影，则 \(\widetilde u_j|_U=\pm\pi_j^*\eta_j\)。

	于是 \(v|_U\) 就是 \(n\) 个平面局部生成元的外积，至多相差符号。\cref{g2:prop:local-cross-normalization}说明，该外积与乘积圆盘的相对基本链配对为 \(\pm1\)，所以它生成
	\[
		H^{2n}(\mathbb C^n,\mathbb C^n\setminus\{0\};\mathbb Z)\cong\mathbb Z.
	\]
	这一步也可直接把每个平面换成原点周围的小闭圆盘：各相对二次类在圆盘上取值 \(\pm1\)，乘积链的评价是这些数的乘积，而乘积圆盘的边界恰对应“至少一个坐标落在边界上”。\cref{g2:prop:relative-cohomology-exactness}的切除结论把这个局部群与 \(H^{2n}(X,X\setminus\{p\};\mathbb Z)\) 识别，故 \(v\) 是后者的生成元。

	最后，\(X\setminus\{p\}\) 变形收缩到超平面 \(\{z_0=0\}\cong\mathbb CP^{n-1}\)。显式收缩为
	\[
		[z_0:z_1:\cdots:z_n]\longmapsto[(1-t)z_0:z_1:\cdots:z_n],\qquad 0\leq t\leq1.
	\]
	在所用定义域中，\((z_1,\ldots,z_n)\ne0\)，所以此公式始终有定义并固定目标超平面。因而补空间在次数 \(2n-1,2n\) 的上同调均为零，空间对长正合列给出
	\[
		H^{2n}(X,X\setminus\{p\};\mathbb Z)\xrightarrow{\ \cong\ }H^{2n}(X;\mathbb Z).
	\]
	\(v\) 的像 \(u^n\) 因此是顶次生成元。

	对每个 \(1\leq k\leq n\)，标准包含 \(\mathbb CP^k\hookrightarrow\mathbb CP^n\) 在二次和 \(2k\) 次上同调上都是同构；这由相容的偶维胞腔基及万有系数同构得到。将上述顶次计算用于 \(\mathbb CP^k\)，可知 \(u^k\) 生成 \(H^{2k}(\mathbb CP^n;\mathbb Z)\)。这些幂连同单位元已经穷尽全部加性基，且 \(2n\) 次以上没有群，故唯一所需截断关系是 \(u^{n+1}=0\)。

	对于 \(\mathbb CP^\infty\)，每个固定次数的胞腔链和对应同调在足够大的有限阶段已经稳定；相应同调群自由，万有系数同构于是保证 \(H^{2k}(\mathbb CP^\infty;\mathbb Z)\to H^{2k}(\mathbb CP^n;\mathbb Z)\) 在 \(n\geq k\) 时为同构。选择二次生成元 \(u\)，再限制到各有限阶段，就知每个 \(u^k\) 都生成该次群。分次上同调环按各次群的直和定义，一个元素只有有限多个非零齐次分量，因此得到多项式环 \(\mathbb Z[u]\)，而不是形式幂级数环。
\end{proof}

证明中起决定作用的是超平面的共同交点。在 \(\mathbb CP^2\) 中，两个坐标射影直线交于 \(p\)，它们对应的相对二次类在 \(p\) 附近分别检测两个复坐标方向；相乘后就检测一个四维邻域。这个图景解释了 \(u^2\) 为什么有机会非零，而局部外积配对为 \(\pm1\) 则进一步排除了它只是顶次生成元的两倍或三倍。这里的坐标与相对杯积证明可与\cite[第3.2节，Spaces with Polynomial Cohomology]{Hatcher2002}对照。

\subsection{相同加性同调、不同乘法结构的模型}
\label{g2:subsec:0405:03}
前面的计算给出两种区分空间的办法：环面的一次类可以相乘得到顶次类，复射影平面的二次类可以平方得到四次生成元。下面分别与加性同调相同的楔和比较，考察这两种乘积在楔和中为何消失。

\begin{example}[环面与球面圆周的楔和]\label{g2:exm:torus-wedge-ring}
	令 \(X=T^2\)、\(Y=S^1\vee S^1\vee S^2\)。两者的整数同调群在次数 \(0,1,2\) 分别是 \(\mathbb Z,\mathbb Z^2,\mathbb Z\)，其他次数为零。但 \(H^*(X;\mathbb Z)\) 中存在一次类 \(x,y\) 使 \(xy\ne0\)，而 \(H^*(Y;\mathbb Z)\) 中任意两个一次类的乘积均为零。后一断言可由\eqref{g2:eq:wedge-positive-cohomology}逐因子核对：在圆周因子上二次群为零，在球面因子上一次群为零。于是两环不可能同构。因为连续映射诱导保持杯积的拉回，同伦等价必诱导上同调环同构，所以 \(X,Y\) 不同伦等价。
\end{example}

\begin{example}[复射影平面与两个球面的楔和]\label{g2:exm:projective-plane-wedge-ring}
	令 \(X=\mathbb CP^2\)、\(Y=S^2\vee S^4\)。两者的整数同调群都仅在次数 \(0,2,4\) 为 \(\mathbb Z\)。然而 \(X\) 的二次生成元 \(u\) 满足 \(u^2\ne0\)，而 \(Y\) 的每个二次类平方为零：它在 \(S^2\) 上的平方因次数为零，在 \(S^4\) 上的限制本身就为零。故两者不同伦等价。这里没有一维同调参与，差异完全发生在偶数次的乘法中。
\end{example}

系数变化还可能消去加性群的区别，同时保留乘法上的差异。对 \(g\geq1\)，\(\Sigma_g\) 与 \(N_{2g}\) 的模 \(2\) 加性同调都在次数 \(0,1,2\) 为 \(\mathbb F_2,\mathbb F_2^{2g},\mathbb F_2\)。前者任意一次类的平方为零：\eqref{g2:eq:orientable-surface-products}模 \(2\) 后，每个基元平方为零，交叉项成对出现。后者则有 \(c_i^2=\omega\ne0\)，所以两者的模 \(2\) 上同调环不同。这项比较的系数始终为 \(\mathbb F_2\)。


\begin{chaptersummary}
	自由消解把一个模写成自由生成元及其关系，并通过比较链映射和比较同伦保证计算不依赖所选表示。取同态后的一次上同调是 Ext，取张量后的一次同调是 Tor；在主理想整环上，自由模的子模仍自由，所以所需消解可以缩短到两项。

	\ProseParagraphBreak

	自由链复形的圈、边界及链群之间有两条不同的短正合列：链群到低一次边界的列可以分裂，边界到圈的列却可能留下挠同调。把两者结合便得到万有系数定理。上同调的取值映射会漏掉低一次同调的 Ext，系数同调中的旧圈张量则会漏掉低一次同调的 Tor。正合列中的映射自然，选出的直和分解一般不自然。

	\ProseParagraphBreak

	Alexander--Whitney 映射沿前后面切开单形，shuffle 映射把单形乘积按格路径剖分。标准模型中的逐次填圈给出自然链同伦，因而空间乘积可以用张量复形计算。Künneth 公式中的张量项来自圈的外积，Tor 项还使用挠圈整数倍的填充，两者在总次数上相差一。

	\ProseParagraphBreak

	杯积是上同调外积沿对角映射的拉回。前后面公式在上链上给出结合律、单位和自然性，而交换律需要一个链同伦，并留下 \((-1)^{pq}\) 的符号。相对杯积还要检查小链比较条件；局部外积配对使各个坐标方向的生成元相乘后仍为乘积邻域的生成元。

	\ProseParagraphBreak

	环面的环是一次生成元的外代数，可定向曲面的乘法记录成对的柄方向，不可定向曲面的模二一次平方可以非零。复射影空间中，坐标超平面的相对类相乘集中到共同交点，证明二次生成元的每个允许幂都是生成元。这些乘法关系既能区分加性同调相同的空间，也能限制连续映射在上同调上的作用。
\end{chaptersummary}

\BookExercises

\begin{exercise}\label{g2:ex:cyclic-coefficient-maps}
	令 \(u:\mathbb Z/12\mathbb Z\to\mathbb Z/4\mathbb Z\) 把 \(1\) 的类送到 \(1\) 的类，令 \(G=\mathbb Z/6\mathbb Z\)。写出 \(u\) 在两项自由消解上的一个提升，并计算
	\[
		u_*:\operatorname{Tor}^{\mathbb Z}_1(\mathbb Z/12\mathbb Z,G)\longrightarrow\operatorname{Tor}^{\mathbb Z}_1(\mathbb Z/4\mathbb Z,G)
	\]
	及
	\[
		u^*:\operatorname{Ext}_{\mathbb Z}^1(\mathbb Z/4\mathbb Z,G)\longrightarrow\operatorname{Ext}_{\mathbb Z}^1(\mathbb Z/12\mathbb Z,G).
	\]
	说明为何仅知道两端群的同构类型，还不足以给出这两个映射。
\end{exercise}
\begin{solution}
	源、靶的消解分别为 \(\mathbb Z\xrightarrow{12}\mathbb Z\)、\(\mathbb Z\xrightarrow{4}\mathbb Z\)，右侧项为零次。取零次提升 \(f_0=1\)，一次提升 \(f_1=3\)，因为 \(4f_1=f_0\,12\)，且零次诱导的商映射正是 \(u\)。

	取张量 \(G\) 后，Tor 分别是 \(G[12]=G\) 与 \(G[4]=\{0,3\}\)。诱导映射由一次提升给出，故 \(u_*(g)=3g\)。它满射，核为 \(\{0,2,4\}\)，而不是把 \(G[12]\) 不经计算地包含到 \(G[4]\) 中。

	取同态后箭头反向。一次上链映射仍由预复合乘 \(3\) 给出，所以
	\[
		u^*:G/4G\longrightarrow G/12G,\qquad[g]\longmapsto[3g].
	\]
	这里 \(4G=\{0,2,4\}\)、\(12G=0\)，故源为二阶群，靶为 \(G\)，该映射将非零类送到 \(3\)，是单射。两处的乘 \(3\) 来自同一个消解提升，但 Tor 在核之间协变作用，Ext 在余核之间反变作用。只有群的阶或同构类型不能确定箭头方向及具体倍数。
\end{solution}

\begin{exercise}\label{g2:ex:uct-needs-free-chains}
	固定 \(m\ge2\)，令整数链复形 \(C\) 仅有 \(C_0=\mathbb Z/m\mathbb Z\) 非零。分别计算 \(H^1(C;\mathbb Z)\) 与 \(H_1(C;\mathbb Z/m\mathbb Z)\)。证明：如果删去两种万有系数定理中“链群自由”的条件，两条短正合列都会在这个例子中失败。
\end{exercise}
\begin{solution}
	原复形只在零次有 \(H_0(C)=\mathbb Z/m\mathbb Z\)。对 \(\operatorname{Hom}(C,\mathbb Z)\)，零次项也为零，因为有限阶元素只能映到 \(0\)；故 \(H^1(C;\mathbb Z)=0\)。如果仍套用上同调万有系数列的 \(n=1\) 情形，就会要求非零群
	\[
		\operatorname{Ext}_{\mathbb Z}^1(H_0(C),\mathbb Z)=\mathbb Z/m\mathbb Z
	\]
	单射到零群中，矛盾。

	张量 \(G=\mathbb Z/m\mathbb Z\) 后，复形仍只在零次非零，所以 \(H_1(C;G)=0\)。如果套用同调万有系数列的 \(n=1\) 情形，就会要求零群满射到
	\[
		\operatorname{Tor}^{\mathbb Z}_1(H_0(C),G)=G[m]=G\ne0,
	\]
	也矛盾。这里直接对非自由复形作张量或取同态，与先将模替换成自由消解再作运算并不是同一件事；失败发生在定理的输入条件，而不是 Ext 或 Tor 的循环模公式。
\end{solution}

\begin{exercise}\label{g2:ex:coefficient-boundary-matrix}
	令自由整数链复形 \(C\) 的非零项为 \(C_2=\mathbb Z\langle c\rangle\)、\(C_1=\mathbb Z\langle x,y\rangle\)、\(C_0=\mathbb Z\langle v\rangle\)，规定 \(\mathrm{d}_2c=4x+6y\)、\(\mathrm{d}_1=0\)。计算整数同调，并在 \(G=\mathbb Z/6\mathbb Z\) 时计算全部 \(H_n(C;G)\)、\(H^n(C;G)\)。给出代表元或坐标描述，并指出万有系数定理中的每个非零项。
\end{exercise}
\begin{solution}
	向量 \(a=2x+3y\)、\(b=x+y\) 构成 \(C_1\) 的一组整数基，因为坐标矩阵的行列式为 \(2-3=-1\)。在此基下 \(\mathrm{d}_2c=2a\)，故
	\[
		H_0(C)=\mathbb Z,\qquad H_1(C)=\mathbb Z\langle[b]\rangle\oplus(\mathbb Z/2\mathbb Z)\langle[a]\rangle,\qquad H_n(C)=0\quad(n\ge2).
	\]
	在 \(G\) 中，\(\mathrm{d}_2(tc)=2t\,a\)，所以二次圈由 \(3c\) 生成，得到 \(H_2(C;G)\cong\mathbb Z/2\mathbb Z\)。一次同调为
	\[
		H_1(C;G)=(G/2G)\langle[a]\rangle\oplus G\langle[b]\rangle\cong\mathbb Z/2\mathbb Z\oplus\mathbb Z/6\mathbb Z,
	\]
	零次同调为 \(G\)，其余次数为零。UCT 中一次项全部来自 \(H_1(C)\otimes G\)，而二次项全部来自 \(\operatorname{Tor}_1(H_1(C),G)=G[2]\)。

	令 \(a^*,b^*\) 表示这组基的 \(G\) 值坐标函数。一次上链 \(r a^*+s b^*\) 的上边界为 \(2r\) 倍的二次坐标函数，所以它闭当且仅当 \(r\in\{0,3\}\)，而 \(s\) 任意。由于零次上链微分为零，
	\[
		H^0(C;G)=G,\qquad H^1(C;G)=G[2]\langle a^*\rangle\oplus G\langle b^*\rangle,\qquad H^2(C;G)=G/2G.
	\]
	在 \(H^1\) 的第一项中，\(G[2]\langle a^*\rangle\) 指允许的系数 \(r\) 属于 \(G[2]\)，其非零类具体为 \(3a^*\)。其他次数为零。UCT 中 \(H^1\) 全部来自 \(\operatorname{Hom}(H_1(C),G)\)，因为 \(H_0(C)\) 自由；\(H^2\) 全部来自 \(\operatorname{Ext}^1(H_1(C),G)=G/2G\)。低一次同调中的同一个二阶挠因子，因而分别在二次同调和二次上同调中通过核与余核出现。
\end{solution}

\begin{exercise}\label{g2:ex:mod-p-betti-count}
	设 \(C\) 为非负自由整数链复形，\(n\ge0\)，并且 \(H_n(C)\)、\(H_{n-1}(C)\) 都有限生成。将其写为自由部分与有限循环群之直和，记自由秩为 \(r_j\)，记该分解中阶被素数 \(p\) 整除的非零循环因子数为 \(t_j(p)\)，约定 \(r_{-1}=t_{-1}(p)=0\)。证明
	\[
		\dim_{\mathbb F_p}H_n(C;\mathbb F_p)=r_n+t_n(p)+t_{n-1}(p),\qquad
		\dim_{\mathbb Q}H_n(C;\mathbb Q)=r_n.
	\]
	再证明同样的两个维数公式对 \(H^n(C;\mathbb F_p)\)、\(H^n(C;\mathbb Q)\) 成立。最后说明这些维数不能区分 \(p\) 阶与 \(p^2\) 阶的循环挠因子。
\end{exercise}
\begin{solution}
	一个自由因子 \(\mathbb Z\) 张量 \(\mathbb F_p\) 贡献一个 \(\mathbb F_p\)，而其一次 Tor 为零。对循环因子 \(\mathbb Z/a\mathbb Z\)，\cref{g2:prop:cyclic-ext-tor}给出张量为 \(\mathbb F_p/a\mathbb F_p\)、Tor 为 \(\mathbb F_p[a]\)。若 \(p\mid a\)，乘 \(a\) 为零，两项各为一维；若 \(p\nmid a\)，乘 \(a\) 可逆，两项均为零。有限直和逐项计算，再由同调万有系数短正合列两端维数相加，就得到第一式。对于 \(\mathbb Q\)，每个非零整数的乘法可逆，所以所有循环因子的张量与 Tor 均为零，只剩自由秩 \(r_n\)。

	上同调公式中，一个自由因子的 Hom 贡献一维，Ext 为零。循环因子的 Hom 为 \(\mathbb F_p[a]\)，Ext 为 \(\mathbb F_p/a\mathbb F_p\)，所以同样各由是否 \(p\mid a\) 决定。上同调万有系数列给出相同的维数。取有理系数时，有限阶群到 \(\mathbb Q\) 的同态为零，而 \(\mathbb Q/a\mathbb Q=0\)，故仍只剩自由秩。

	以上维数只检测循环因子的阶是否被 \(p\) 整除，不检测被 \(p\) 整除的次数。例如分别以乘 \(p\)、乘 \(p^2\) 为微分的两项自由复形，整数零次同调分别为 \(\mathbb Z/p\mathbb Z\)、\(\mathbb Z/p^2\mathbb Z\)，但张量 \(\mathbb F_p\) 后两者微分都为零，所有模 \(p\) 同调维数相同；有理系数也都不能区分它们。
\end{solution}

\begin{exercise}\label{g2:ex:lens-product-groups}
	设 \(m,n\ge2\)，\(g=\gcd(m,n)\)。计算 \(L(m,1)\times L(n,1)\) 的全部整数同调群，并指出哪些项来自 Künneth 公式中的 Tor。所写的直和分解是否都由外积自然给出？
\end{exercise}
\begin{solution}
	由\cref{g2:prop:lens-cellular-model}，\(L(m,1)\) 的整数同调在次数 \(0,1,3\) 分别为 \(\mathbb Z,\mathbb Z/m,\mathbb Z\)，其余为零；另一个因子把 \(m\) 换成 \(n\)。自由因子不贡献 Tor，两个挠群的张量与 Tor 均为 \(\mathbb Z/g\)。因此
	\[
		H_j(L(m,1)\times L(n,1);\mathbb Z)\cong
		\begin{cases}
			\mathbb Z,&j=0,6,\\
			\mathbb Z/m\oplus\mathbb Z/n,&j=1,4,\\
			\mathbb Z/g,&j=2,\\
			\mathbb Z^2\oplus\mathbb Z/g,&j=3,\\
			0,&j=5\text{ 或 }j>6.
		\end{cases}
	\]
	其中次数二的 \(\mathbb Z/g\) 来自 \(H_1\otimes H_1\)；次数三的 \(\mathbb Z/g\) 才来自 \(\operatorname{Tor}_1(H_1,H_1)\)。次数四的两个挠项来自 \(H_1\otimes H_3\)、\(H_3\otimes H_1\)，不是 Tor。除次数三以外，所有显示的群都由外积自然识别；次数三应保留自然短正合列
	\[
		0\longrightarrow\mathbb Z^2\longrightarrow H_3
		\longrightarrow\mathbb Z/g\longrightarrow0.
	\]
	次数三的挠类不是两个同调类的外积。不过，本题还有一个比一般分裂定理更强的特点：左端自由、右端有限，而正合列已知分裂，因此 \(H_3\) 的挠子群经右箭头同构到 \(\mathbb Z/g\)。其逆给出这条列的唯一分裂，因为两个分裂之差须为 \(\mathbb Z/g\to\mathbb Z^2\) 的同态，只能为零。所以本题的第三次直和可由外积子群与挠子群共同确定。若 \(g=1\)，相应循环群为零，所有次数都由外积给出。
\end{solution}

\begin{exercise}\label{g2:ex:explicit-mixed-tor-cycle}
	令 \(E\) 的基为 \(x\in E_1\)、\(a\in E_0\)，微分 \(\mathrm{d}x=ma\)；令 \(F\) 的基为 \(y\in F_1\)、\(b\in F_0\)，微分 \(\mathrm{d}y=nb\)，其中 \(m,n\ge2\)，其他次数为零。直接计算 \(E\otimes F\) 的同调，写出次数一的生成圈，并计算交换链映射 \(\tau:E\otimes F\to F\otimes E\) 对所选生成圈的作用。
\end{exercise}
\begin{solution}
	置 \(g=\gcd(m,n)\)。次数零以 \(a\otimes b\) 为基，来自次数一的边界系数是 \(m,n\)，故 \(H_0=\mathbb Z/g\)。次数一的链
	\[
		u(x\otimes b)+v(a\otimes y)
	\]
	为圈当且仅当 \(mu+nv=0\)。除以 \(g\) 后 \(m/g,n/g\) 互素，因此所有整数解恰为
	\[
		u=k\,n/g,\qquad v=-k\,m/g,\qquad k\in\mathbb Z.
	\]
	于是圈群由
	\[
		z=\frac ng\,x\otimes b-\frac mg\,a\otimes y
	\]
	生成。而
	\[
		\mathrm{d}(x\otimes y)=m\,a\otimes y-n\,x\otimes b=-gz,
	\]
	故 \(H_1=\mathbb Z/g\)。这个边界同态从次数二的 \(\mathbb Z\) 单射，所以 \(H_2=0\)，更高次数也为零。

	在 \(F\otimes E\) 中按同一规则选
	\[
		z'=\frac mg\,y\otimes a-\frac ng\,b\otimes x.
	\]
	由于参与 \(z\) 的每个张量都有一个零次因子，交换符号均为 \(+1\)，直接计算得到 \(\tau(z)=-z'\)。所以在这两组指定的循环群生成元下，次数一的交换同态为乘以 \(-1\)。这不否定 Tor 的对称同构，而是说明把其抽象同构写成具体整数生成元时，必须同时说明生成圈的符号。
\end{solution}

\begin{exercise}\label{g2:ex:torus-map-determinant-cup}
	令 \(u,v\in H^1(T^2;\mathbb Z)\) 分别来自两个正向圆周，并令 \(f:T^2\to T^2\) 为连续映射。若
	\[
		f^*u=au+bv,\qquad f^*v=cu+dv,
		\qquad a,b,c,d\in\mathbb Z,
	\]
	计算 \(f^*\) 在 \(H^2(T^2;\mathbb Z)\) 上的作用。由此判断：若一个环面自映射在 \(H^1\) 上可逆，它在 \(H^2\) 上是否可逆？
\end{exercise}
\begin{solution}
	圆周的二次上同调为零，故由拉回的乘法相容性 \(u^2=v^2=0\)。分次交换律给出 \(vu=-uv\)，而\cref{g2:prop:cross-cup-evaluation}说明 \(uv\) 在两个正向圆周的外积上取值 \(1\)，所以它是 \(H^2(T^2;\mathbb Z)\cong\mathbb Z\) 的生成元。于是
	\[
		f^*(uv)=(au+bv)(cu+dv)
		=ad\,uv+bc\,vu=(ad-bc)\,uv.
	\]
	若 \(f^*\) 在 \(H^1\cong\mathbb Z^2\) 上可逆，其整数矩阵的行列式必为 \(\pm1\)，故在 \(H^2\) 上也是同构。特别地，交换两个圆周方向使矩阵行列式为 \(-1\)，因此把二次生成元变号；这一符号与乘积定向的交换一致。
\end{solution}

\begin{exercise}\label{g2:ex:cup-obstruction-two-open-cover}
	设 \(X=U\cup V\)，其中 \(U,V\) 为开子空间，\(R\) 为含幺交换环。证明：若 \(\alpha\in H^p(X;R)\) 在 \(U\) 上限制为零，\(\beta\in H^q(X;R)\) 在 \(V\) 上限制为零，则 \(\alpha\smile\beta=0\)。据此证明，对 \(p,q\ge1\)，\(S^p\times S^q\) 不存在这样的两开集覆盖：两个包含映射 \(U\to S^p\times S^q\)、\(V\to S^p\times S^q\) 都同伦于常值映射。
\end{exercise}
\begin{solution}
	先设 \(p>0\)，取代表上圈 \(\varphi\)。限制类为零意味着 \(\varphi|_{C_*(U)}=\delta h\)，其中 \(h\) 为 \(U\) 上的 \(p-1\) 次上链。把 \(h\) 在 \(U\) 之外的奇异单形基元上取值为零，就得到 \(X\) 上的延拓 \(\widehat h\)。于是
	\[
		\varphi'=\varphi-\delta\widehat h
	\]
	仍代表 \(\alpha\)，并且在 \(U\) 中的全部链上为零；因此 \(\alpha\) 来自 \(H^p(X,U;R)\)。若 \(p=0\)，零次没有上边界，限制类为零已经说明代表上圈本身在 \(U\) 上为零，仍得同一结论。同理，\(\beta\) 来自 \(H^q(X,V;R)\)。由\cref{g2:prop:relative-cup-product}，两相对提升的杯积属于
	\[
		H^{p+q}(X,U\cup V;R)=H^{p+q}(X,X;R)=0.
	\]
	映到绝对上同调即得 \(\alpha\smile\beta=0\)。

	现在取整数系数，令 \(\alpha,\beta\) 分别为两个球面正向生成元的拉回。它们的杯积在球面基本类的外积上取值 \(1\)，所以非零。若 \(U,V\) 的包含映射都同伦于常值，则任意正次上同调类在它们上的限制都为零：常值映射经一点分解，而一点的正次上同调为零。上一部分将迫使 \(\alpha\smile\beta=0\)，矛盾。
\end{solution}

\begin{exercise}\label{g2:ex:complex-local-orientation}
	在 \(\mathbb C^n\) 上采用实坐标次序 \((x_1,y_1,\ldots,x_n,y_n)\)，令 \(u\) 为 \(H^{2n}(\mathbb C^n,\mathbb C^n\setminus\{0\};\mathbb Z)\) 的正向生成元。计算下列映射对 \(u\) 的拉回：
	\[
		c_r(z_1,\ldots,z_n)
		=(\overline z_1,\ldots,\overline z_r,z_{r+1},\ldots,z_n),
		\qquad 0\le r\le n,
	\]
	以及任意复坐标置换。要求用局部生成元的外积说明符号。
\end{exercise}
\begin{solution}
	先在实轴的局部对 \((\mathbb R,\mathbb R\setminus\{0\})\) 中考察反射。正向小区间的连接边界是“正侧点减负侧点”；反射交换两侧，所以把局部一维同调生成元变号，也把其取值对偶变号。复共轭在一个坐标复平面中保持 \(x\) 不变、把 \(y\) 变号。\cref{g2:prop:local-cross-normalization}把局部二次生成元写成 \(x,y\) 的一维局部生成元的外积，故复共轭把它乘以 \(-1\)。

	再次使用该命题，\(\mathbb C^n\) 的 \(u\) 是 \(n\) 个局部二次生成元按给定坐标次序的外积。外积自然性于是给出
	\[
		c_r^*u=(-1)^r u.
	\]
	交换两个复坐标只交换两个次数为 \(2\) 的局部类，系数为 \((-1)^{2\cdot2}=1\)。任何复坐标置换都可写成相邻交换的复合，所以都保持 \(u\)。这说明一个复坐标块的共轭与两个完整复坐标块的交换，对实定向有不同作用。
\end{solution}

\begin{exercise}\label{g2:ex:surface-ring-genus-bound}
	设 \(g,h\geq1\)，\(f:\Sigma_g\to\Sigma_h\) 连续。选定两曲面的定向，并用二次生成元定义整数 \(d\)：\(f^*\omega_h=d\omega_g\)。证明：若 \(d\ne0\)，则 \(g\geq h\)。此外，若 \(g=h\)，设一次拉回在正文的有序基 \((a_1,b_1,\ldots,a_g,b_g)\) 下的矩阵为 \(P\)，证明 \((\det P)^2=d^{2g}\)。
\end{exercise}
\begin{solution}
	在一次有理上同调上定义双线性配对 \(\lambda_g\)：若 \(uv=r\omega_g\)，令 \(\lambda_g(u,v)=r\)。由\eqref{g2:eq:orientable-surface-products}，它在所列基下的矩阵是 \(g\) 个块 \(\left(\begin{smallmatrix}0&1\\-1&0\end{smallmatrix}\right)\) 的对角和，记为 \(J_g\)。该矩阵可逆。杯积自然性给出
	\[
		\lambda_g(f^*u,f^*v)=d\lambda_h(u,v),\qquad P^{\mathsf T}J_gP=dJ_h.
	\]
	这里 \(P\) 一般为 \(2g\times2h\) 矩阵。若 \(f^*u=0\)，上述等式说明 \(d\lambda_h(u,v)=0\) 对所有 \(v\) 成立。因为 \(d\ne0\) 且 \(J_h\) 可逆，必有 \(u=0\)。所以 \(f^*\) 在一次有理上同调上单射，\(2h\leq2g\)。若 \(g=h\)，则 \(P\) 是方阵，对矩阵恒等式取行列式，利用 \(\det J_g=1\)，得到 \((\det P)^2=d^{2g}\)。
\end{solution}

\begin{exercise}\label{g2:ex:projective-product-even-degree}
	令 \(X=\mathbb CP^1\times\mathbb CP^1\)，两因子的二次生成元拉回后分别记为 \(x,y\)，\(\mathbb CP^2\) 的二次生成元记为 \(u\)。设 \(f:X\to\mathbb CP^2\) 连续且 \(f^*u=ax+by\)，其中 \(a,b\in\mathbb Z\)。计算 \(f^*(u^2)\)，并证明不存在使 \(f^*:H^4(\mathbb CP^2;\mathbb Z)\to H^4(X;\mathbb Z)\) 为同构的连续映射。然后证明任意连续映射 \(r:\mathbb CP^2\times\mathbb CP^2\to\mathbb CP^1\) 在二次整数上同调上的拉回都是零。
\end{exercise}
\begin{solution}
	\cref{g2:thm:complex-projective-ring}给出 \(x^2=y^2=0\)，而外积说明 \(xy\) 是 \(H^4(X;\mathbb Z)\) 的生成元。因 \(x,y\) 的次数均为 \(2\)，交换时没有负号，故
	\[
		f^*(u^2)=(ax+by)^2=2ab\,xy.
	\]
	这是乘以偶数的同态，不能为乘以 \(1\) 或 \(-1\)，所以不可能是整数群的同构。

	对于第二问，把源的两个二次生成元记为 \(v,w\)，靶的二次生成元记为 \(s\)。若 \(r^*s=cv+ew\)，由 \(s^2=0\) 得
	\[
		0=(cv+ew)^2=c^2v^2+2ce\,vw+e^2w^2.
	\]
	整数 Künneth 同构及\cref{g2:thm:complex-projective-ring}说明 \(v^2,vw,w^2\) 是四次上同调的一组自由基。因此 \(c^2=e^2=0\)，从而 \(c=e=0\)。这个障碍来自源与靶的平方关系不同；加性群中的任意同态未必能来自连续映射。
\end{solution}

\begin{exercise}\label{g2:ex:nonorientable-square-kernel}
	在 \(N_h\) 的一次模 \(2\) 上同调中令
	\[
		V=\{v\in H^1(N_h;\mathbb F_2):v^2=0\}.
	\]
	证明 \(V\) 是维数 \(h-1\) 的子空间。若用 \(vw=\lambda(v,w)\omega\) 定义双线性配对，计算它限制在 \(V\) 上的根空间 \(\{v\in V:\lambda(v,w)=0\text{ 对所有 }w\in V\}\)，并区分 \(h\) 的奇偶性。
\end{exercise}
\begin{solution}
	写 \(v=\sum_i a_ic_i\)，\(a_i\in\mathbb F_2\)。由\eqref{g2:eq:nonorientable-surface-products}，有
	\[
		v^2=\Bigl(\sum_i a_i\Bigr)\omega,\qquad
		\lambda\Bigl(\sum_i a_ic_i,\sum_i b_ic_i\Bigr)=\sum_i a_ib_i.
	\]
	故 \(V\) 是系数和为零的超平面，维数为 \(h-1\)。令 \(e=c_1+\cdots+c_h\)，则 \(V=e^\perp\)，其中环境配对就是标准点积，因而非退化。直接验证也可确定 \(V^\perp\)：当 \(h\geq2\) 时，向量 \(c_i+c_j\) 都属于 \(V\)，与它们正交要求各坐标相等，所以 \(V^\perp=\mathbb F_2e\)；当 \(h=1\) 时此结论同样成立。因此限制配对的根空间为 \(V\cap\mathbb F_2e\)。若 \(h\) 为奇数，\(e\) 的系数和为 \(1\)，根空间为零；若 \(h\) 为偶数，\(e\in V\)，根空间就是一维空间 \(\mathbb F_2e\)。
\end{solution}

\begin{exercise}\label{g2:ex:projective-product-ring-automorphisms}
	设 \(1\leq n\leq m\)，令 \(X=\mathbb CP^n\times\mathbb CP^m\)，把两个因子的二次生成元在 \(X\) 上的拉回记为 \(u,v\)。求 \(H^*(X;\mathbb Z)\) 的全部保次数、保单位元的环自同构，区分 \(n<m\) 与 \(n=m\)。证明每个这样的环自同构都由 \(X\) 的某个同胚诱导。
\end{exercise}
\begin{solution}
	由\cref{g2:thm:complex-projective-ring}及无挠的上同调外积，
	\[
		H^*(X;\mathbb Z)=\mathbb Z[u,v]/(u^{n+1},v^{m+1}),\qquad |u|=|v|=2.
	\]
	其加性基为 \(u^iv^j\)、\(0\leq i\leq n\)、\(0\leq j\leq m\)。一个分次环自同构 \(\Phi\) 必须满足
	\[
		\Phi(u)=au+bv,\qquad\Phi(v)=cu+dv,\qquad ad-bc=\pm1.
	\]
	若 \(n<m\)，考察关系 \(\Phi(u)^{n+1}=0\) 中的 \(v^{n+1}\) 系数，它等于 \(b^{n+1}\)。该单项式属于加性基，故 \(b=0\)。行列式条件随即给出 \(a,d\in\{1,-1\}\)。再看 \(\Phi(v)^{m+1}=0\) 中的 \(uv^m\) 系数，得到 \((m+1)cd^m=0\)，所以 \(c=0\)。因此自同构恰是两个生成元各自独立变号。

	若 \(n=m\)，关系 \(\Phi(u)^{n+1}=0\) 的 \(u^nv\) 系数给出 \((n+1)a^nb=0\)，所以 \(a=0\) 或 \(b=0\)。同理，\(c=0\) 或 \(d=0\)。结合行列式为 \(\pm1\)，矩阵只能是两个对角元为 \(\pm1\) 的对角矩阵，或两个反对角元为 \(\pm1\) 的反对角矩阵。也就是说，除了独立变号，还可以交换 \(u,v\)。这些候选确实都保持截断关系，因而都是环自同构。

	最后说明几何实现。复共轭 \([z_0:\cdots:z_r]\mapsto[\overline z_0:\cdots:\overline z_r]\) 是 \(\mathbb CP^r\) 的同胚。在标准 \(\mathbb CP^1\cong S^2\) 上，它在仿射图 \(z=x+\mathrm{i}y\) 中是 \((x,y)\mapsto(x,-y)\)，经球面的立体投影就是一个反射，故由\cref{g2:prop:reflection-antipodal-degree}诱导二次生成元变号。标准包含在二次上同调上为同构，所以复共轭在任意 \(\mathbb CP^r\) 上也使二次生成元变号。对两个因子分别选恒等或复共轭，即实现全部独立变号；当 \(n=m\) 时，再用交换两个因子的同胚实现交换。因而上面列出的每个环自同构都有所要求的实现。
\end{solution}

% !TeX root = ../../Book2.tex
% 纲要标题：## 第5章　局部系统与扭系数同调
% 纲要位置：几何学计划纲要.md，第 1191 行。
% 正文状态：初稿；尚未正式审核。

% 纲要内容开始
% 
% > 本章目标与证明范围：从覆盖的单值进入非平凡系数，完整构造链模型和函子性。定向局部系统、纤维化中的系数作用及障碍系数在此有共同的实际基础。 分类主版本固定连通、局部道路连通和半局部单连通空间；更一般群胚表述另述。
% 
% 纲要内容结束

\chapter{局部系统与扭系数同调}
\label{g2:ch:05}
\BookChapterNumber{g2:ch:05}

\begin{chapterabstract}
	本章把固定系数推广为随位置变化、可沿道路运输的系数。先建立路径群胚、局部常值层与基本群表示之间的联系，再用万有覆盖上的群环链复形构造扭系数同调及上同调，并给出首顶点系数的几何公式。定向局部系统说明符号单值如何来自流形自身；圆周和不可定向曲面的计算展示核、余核及顶次整数类怎样随单值改变。最后把纤维同调组织为基底上的局部系统，证明映射环面的正合列，为后续纤维化、谱序列与障碍理论准备系数数据。
\end{chapterabstract}

\begin{learninggoals}
	\begin{itemize}
		\item 用道路运输定义局部系统，在明确的局部条件下证明它与局部常值层、基本群模的等价，说明换基点时需要保留的比较映射。
		\item 区分左右群环作用，构造局部系数链与上链，证明函子性、正合列、同伦不变性、切除及胞腔计算的依据。
		\item 从流形局部同调构造定向系统，计算圆周与不可定向曲面的扭系数同调及上同调，解释不变元和余不变元的差别。
		\item 构造纤维同调与上同调的运输，计算映射环面，并区分局部系统的不变类、几何截面和更高延拓问题。
	\end{itemize}
\end{learninggoals}

在前两章的链中，一个单形的系数是整数，取边界时把同一个整数放在各个面前即可。这个习惯默认了不同位置的系数已经处在同一个群中。可是，在不可定向曲面上移动一小片有定向的圆盘，绕某条闭路回来时，其局部定向可能反转。若系数用来记录这片圆盘的定向，那么“绕一圈后仍取原系数”就不再符合几何。

\ProseParagraphBreak

另一种情形来自一族空间。各个纤维的同调群可能彼此同构，但沿基底环路比较它们以后，这个群会经历一个非平凡自同构。只记录“每个纤维的群都是 \(A\)”会漏掉这个自同构；局部系统把群和比较它们的道路运输一起保留下来。上一章改变的是所有位置共用的系数，本章还允许绕行改变系数的识别。

\ProseParagraphBreak

\chapref{g2:ch:01}的覆盖与单值给出运输的几何模型，\chapref{g2:ch:02}的奇异链和正合列用于构造带运输的边界，\chapref{g2:ch:03}的胞腔方法则将计算缩到少量生成元。上链及其乘积沿用\chapref{g2:ch:04}的构造。层在这里记录可以局部选择并粘合的系数，其定义和与运输的联系在本章建立。路径群胚上的运输可在任意空间上定义；它与基本群表示、局部常值层之间的分类对应，在连通、局部道路连通、半局部单连通的条件下证明。

\ProseParagraphBreak

本章的一个持续约定是：\(\alpha*\beta\) 表示先走 \(\alpha\)、再走 \(\beta\)，故运输满足 \(T_{\alpha*\beta}=T_\beta T_\alpha\)。它决定后面的右群作用、首顶点边界和上链中的逆映射。局部系数的阅读对照主要见\cite[第3.H节]{Hatcher2002}；不同书的左右作用与路径复合约定可能不同，比较公式时须同时比较这些约定。

% !TeX root = ../../Book2.tex
% 纲要标题：### 5.1 局部系统的等价描述
% 纲要位置：几何学计划纲要.md，第 1195 行。
% 正文状态：初稿；尚未正式审核。

% 纲要内容开始
% 
% * 路径群胚上的运输函子
% * 局部常值层与纤维
% * 连通局部良好空间上的基本群表示
% * 换基点和同构的相容
% 
% 纲要内容结束

\section{局部系统的等价描述}
\label{g2:sec:0501}

在通常的整数链中，单形上的系数可以直接相加，因为它们都属于同一个指定的群 \(\mathbb Z\)。如果每一点只给出一个同构于 \(\mathbb Z\) 的群，情形就不同了：把一个系数送到另一点需要选择道路，绕一圈回来还可能把它变成相反数。定向的改变正会产生这种现象。本节先把“沿道路比较系数”本身定义清楚，再说明在适当的局部条件下，它也可以表现为一个局部常值的层，或一个基本群作用。

\ProseParagraphBreak

以下 \(R\) 为含幺交换环，系数纤维为 \(R\) 模，\(R=\mathbb Z\) 时就是阿贝尔群。这里不要求模自由或有限生成。沿用\cref{g2:def:fundamental-group}的顺序，\(\alpha*\beta\) 表示先走 \(\alpha\)、后走 \(\beta\)，\(\bar\alpha\) 为反向道路；因此运输的复合写成 \(T_{\alpha*\beta}=T_\beta\circ T_\alpha\)。

\subsection{路径群胚上的运输函子}
\label{g2:subsec:0501:01}
同一端点的环路组成基本群；若同时保留所有点以及不同点之间的道路，就得到一个能直接描述运输的对象。

\begin{definition}[局部系统]\label{g2:def:local-system}
	设 \(X\) 为拓扑空间。它的\term[lujingqunpei]{路径群胚}[fundamental groupoid] \(\Pi_1(X)\) 以 \(X\) 的点为对象，以定端点同伦道路类 \([\alpha]:x\to y\) 为态射，复合为
	\[
		[\beta]\circ[\alpha]=[\alpha*\beta].
	\]
	常值道路给出恒等态射，反向道路给出逆态射；这些运算良定义由\cref{g2:def:fundamental-group}之后对重参数化和往返道路的论证保证。\symidx{\(\Pi_1(X)\)：路径群胚}

	一个 \(R\) 模\term[jubuxitong]{局部系统}[local system] \(\mathcal L\) 是从 \(\Pi_1(X)\) 到 \(R\) 模范畴的函子。具体说，它给每点 \(x\) 一个模 \(\mathcal L_x\)，给每条道路类一个线性同构 \(T_\alpha:\mathcal L_x\to\mathcal L_y\)，并满足
	\begin{equation}\label{g2:eq:local-transport-composition}
		T_{e_x}=\operatorname{id},\qquad T_{\alpha*\beta}=T_\beta T_\alpha.
	\end{equation}
	\(\mathcal L_x\) 称为在 \(x\) 处的纤维，\(T_\alpha\) 称为沿 \(\alpha\) 的运输。逆道路恒等式给出 \(T_{\bar\alpha}=T_\alpha^{-1}\)。局部系统之间的态射 \(u:\mathcal L\to\mathcal K\) 是一族线性映射 \(u_x:\mathcal L_x\to\mathcal K_x\)，满足每条 \(\alpha:x\to y\) 上的自然性
	\[
		\begin{tikzcd}[column sep=large,row sep=large]
			\mathcal L_x \arrow[r,"T^{\mathcal L}_\alpha"] \arrow[d,"u_x"'] & \mathcal L_y \arrow[d,"u_y"] \\
			\mathcal K_x \arrow[r,"T^{\mathcal K}_\alpha"'] & \mathcal K_y.
		\end{tikzcd}
	\]
	\symidx{\(\mathcal L_x,T_\alpha\)：局部系统的纤维与道路运输}
\end{definition}

这个定义在任意拓扑空间上有意义；它只使用道路及其同伦，没有预先给纤维的并集指定拓扑。两个局部系统同构，指存在各纤维上均为同构的态射，其逆映射族自动满足同一自然性方形。固定一个 \(R\) 模 \(M\)，在每点都取 \(M\)、每条道路都取恒等映射，得到\term[changzhixitong]{常值系统}[constant system]，记为 \(\underline M\)。局部系统平凡，是指它同构于某个这样的系统；这并不表示已经选定那个同构。

\begin{proposition}[拉回与同伦运输]\label{g2:prop:local-system-pullback}
	设 \(f:Y\to X\) 为连续映射，\(\mathcal L\) 为 \(X\) 上的 \(R\) 模局部系统。规定
	\[
		(f^*\mathcal L)_y=\mathcal L_{f(y)},\qquad T^{f^*\mathcal L}_\alpha=T^{\mathcal L}_{f\circ\alpha},
	\]
	便得到 \(Y\) 上的局部系统，称为 \(\mathcal L\) 的\term[lahuijubuxitong]{拉回局部系统}[pullback local system]。它也逐纤维拉回态射，并满足 \((f\circ g)^*=g^*f^*\) 及恒等映射的相容性。若 \(H:Y\times I\to X\) 是从 \(f_0\) 到 \(f_1\) 的同伦，令 \(\lambda_y(t)=H(y,t)\)，则
	\[
		T_{\lambda_y}:(f_0^*\mathcal L)_y\longrightarrow(f_1^*\mathcal L)_y
	\]
	组成局部系统同构。此同构由所给同伦确定；不同的同伦可能给出不同的同构。
\end{proposition}
\begin{proof}
	连续映射把定端点同伦送到定端点同伦，并保持道路拼接及常值道路，所以拉回满足\eqref{g2:eq:local-transport-composition}。态射、恒等和复合的公式都是对相应点或道路逐项求值。

	对道路 \(\alpha:y\to z\)，将 \(H\) 限制在 \(\alpha\) 上得到方形 \((s,t)\mapsto H(\alpha(s),t)\)。方形内两条从左下到右上的边界路线定端点同伦，故
	\[
		(f_0\alpha)*\lambda_z\simeq\lambda_y*(f_1\alpha).
	\]
	应用运输公式即得 \(T_{\lambda_z}T_{f_0\alpha}=T_{f_1\alpha}T_{\lambda_y}\)，这正是态射的自然性。每个 \(T_{\lambda_y}\) 可逆，故得到同构。取 \(Y\) 为一点、\(f_0=f_1\) 时，同伦就是基点环路，所得同构就是该环路的单值；因此它确实可能依赖同伦。
\end{proof}

拉回的定义解释了系数为何需要跟着空间映射一起变化：若 \(y\) 映到 \(f(y)\)，在 \(y\) 使用的系数正应取自 \(\mathcal L_{f(y)}\)。稍后定义链映射时，这会成为系数所在位置的相容条件。

\subsection{局部常值层与纤维}
\label{g2:subsec:0501:02}
道路运输给出不同纤维的比较，层则记录一个开集上能同时选取哪些彼此相容的系数。在运输平凡的小邻域内，选定一个值后可把它搬到各点，得到一份局部常值截面；层的限制与粘合正是为了组织这样的局部选择。

\begin{definition}[局部常值层]\label{g2:def:locally-constant-sheaf}
	拓扑空间 \(X\) 上的 \(R\) 模\term[ceng]{层}[sheaf] \(\mathscr F\)，给每个开集 \(U\) 一个 \(R\) 模 \(\mathscr F(U)\)，并给每个包含 \(V\subset U\) 一个线性限制映射 \(s\mapsto s|_V\)。限制保持恒等与复合，且满足粘合条件：若 \(U=\bigcup_iU_i\)，而 \(s_i\in\mathscr F(U_i)\) 在每个交集上相等，则存在唯一 \(s\in\mathscr F(U)\) 使 \(s|_{U_i}=s_i\)。\(\mathscr F(U)\) 的元素称为 \(U\) 上的\term[jiemian]{截面}[section]。

	在 \(x\) 处，把定义在 \(x\) 的某个开邻域上的截面按“在更小邻域上一致”识别，所得模称为\term[jing]{茎}[stalk] \(\mathscr F_x\)；截面 \(s\) 的类记为 \(s_x\)，称为它在 \(x\) 处的\term[ya]{芽}[germ]。加法和数乘在两个代表的公共邻域内进行，进一步缩小邻域不改变结果，故定义良好。

	给模 \(A\) 离散拓扑，\(U\mapsto\{\text{局部常值函数 }U\to A\}\) 构成常值层 \(\underline A\)。若每个 \(x\in X\) 都有开邻域 \(U\)，使限制在 \(U\) 上的层 \(\mathscr F|_U\) 同构于 \(\underline A|_U\)，其中 \(A\) 可随邻域变化，则称 \(\mathscr F\) 为\term[jubuchangzhiceng]{局部常值层}[locally constant sheaf]。层的态射是一族与限制相容的线性映射。
\end{definition}

这里的常值层允许一个不连通开集的不同部分取不同值。若仅给每个非空开集同一个模 \(A\)，并把所有限制设为恒等，就通常不满足粘合条件：两个不交开集上选的不同常值无法粘成一个共同常值。另一方面，茎只看一点附近，所以 \(\underline A_x\cong A\)，同构把一个局部常值函数的芽送到它在 \(x\) 的值。

\begin{proposition}[层的纤维空间]\label{g2:prop:local-sheaf-fiber-space}
	设 \(X\) 为拓扑空间，\(R\) 为含幺交换环。局部常值 \(R\) 模层等价于下列数据：一个覆盖 \(p:E\to X\)，各纤维具有 \(R\) 模结构，并在每点附近有保持纤维线性结构的同胚 \(p^{-1}(U)\cong U\times A\)，其中 \(A\) 取离散拓扑。态射分别取层态射和与投影相容、在各纤维上线性的连续映射。这样的 \(E\) 通过道路提升确定一个 \(R\) 模局部系统。
\end{proposition}
\begin{proof}
	先由层构造空间。令 \(E(\mathscr F)=\coprod_{x\in X}\mathscr F_x\)，投影把芽送到它所在的点。对 \(s\in\mathscr F(U)\)，定义 \(\widehat s(U)=\{s_x:x\in U\}\)。以这些集合为开集的基：若 \(s_x=t_x\)，芽的定义给出 \(x\) 的邻域 \(W\)，使 \(s|_W=t|_W\)，所以 \(\widehat s(W)=\widehat t(W)\) 包含在原来两个集合的交内。每个芽有截面代表，故这些集合也覆盖 \(E(\mathscr F)\)。投影在 \(\widehat s(U)\) 上是到 \(U\) 的双射，其逆 \(x\mapsto s_x\) 连续；后一结论再次来自“两个芽相同就在某个邻域上一致”。投影本身连续，因为开集的逆像是这类基本开集的并。因此投影在每个 \(\widehat s(U)\) 上是同胚。

	若 \(\mathscr F|_U\cong\underline A|_U\)，茎的对应给出 \(E(\mathscr F)|_U\cong U\times A\)。在这个对应中，基本开集就是局部常值函数的图像；它们与离散 \(A\) 的乘积拓扑一致。因此投影不仅是局部同胚，而且在同一 \(U\) 上把全部纤维同时平凡化，确为覆盖。一个层态射在茎上诱导线性映射；它把截面芽族送到像截面的芽族，所以得到连续的纤维线性映射。

	反过来，对题述 \(E\) 定义 \(\mathscr S_E(U)\) 为所有连续截面 \(s:U\to E\)。加法和数乘逐点进行，在乘积邻域中就是对局部常值函数进行运算，因而仍连续。相容截面逐点粘合后连续，且唯一，故 \(\mathscr S_E\) 是层；在每个乘积邻域上，它正是常值层。映射 \(\mathscr F(U)\to\mathscr S_{E(\mathscr F)}(U)\)、\(s\mapsto(x\mapsto s_x)\) 单射，因为芽处处相同的截面局部相同，再由粘合唯一性相同。它满射，因为一个连续截面在任一点附近落入某个 \(\widehat t(V)\)，于是局部来自截面 \(t\)，这些代表在交集上具有相同的芽，因而能够唯一粘合。另一方面，\(E(\mathscr S_E)\to E\) 把截面的芽送到该点的值；在乘积邻域中它是恒等的纤维对应，故为纤维线性同胚。这两个对应也保持态射。

	最后，\cref{g2:thm:covering-lifts}已证明：覆盖中从指定纤维点出发的道路提升唯一，提升终点只依赖定端点同伦类。把提升终点定义为 \(T_\alpha\)，拼接提升给出\eqref{g2:eq:local-transport-composition}。将道路有限细分到乘积邻域，在每一段内运输都是常值坐标的识别，所以线性，复合后仍线性。连续纤维线性映射把道路提升送到道路提升，由唯一性便与运输交换。
\end{proof}

纤维空间使用离散的系数拓扑。即使 \(R=\mathbb R\)、纤维为 \(\mathbb R^r\)，这里也不是把通常的向量空间拓扑放在纤维上；离散拓扑使“沿道路保持局部常值”成为一项确定的数据。后文谈平坦丛时，可以把这项运输结构与其对应的光滑向量丛联系起来。

\begin{proposition}[运输的局部实现]\label{g2:prop:local-system-realization}
	设 \(X\) 局部道路连通且半局部单连通。每个 \(R\) 模局部系统 \(\mathcal L\) 的纤维并集 \(\coprod_x\mathcal L_x\) 都有自然的覆盖拓扑及局部模平凡化，使其道路提升恢复 \(\mathcal L\) 的运输。它与\cref{g2:prop:local-sheaf-fiber-space}的构造一起，给出局部系统与局部常值 \(R\) 模层的等价。
\end{proposition}
\begin{proof}
	由\cref{g2:def:local-covering-conditions}及\cref{g2:thm:universal-cover}证明开头的取邻域论证，可以取由道路连通开集 \(U\) 构成的开基，且 \(U\) 中的每条环路在 \(X\) 中零伦。固定 \(u\in U\)。若 \(y\in U\)，从 \(u\) 到 \(y\) 的两条 \(U\) 内道路的差环路在 \(X\) 中零伦，所以它们在 \(X\) 中定端点同伦，给出相同的运输。对每个 \(m\in\mathcal L_u\) 因而有确定的族
	\[
		s_m(y)=T_\eta(m),\qquad \eta:u\to y\text{ 位于 }U.
	\]
	对任意 \(y\)，\(m\mapsto s_m(y)\) 是线性同构，故这些族互不相交且填满 \(U\) 上的纤维。

	把全部集合 \(s_m(U)\) 宣布为开集的基。若两个这样的族在 \(y\) 处相同，在两个底开集的交内取含 \(y\) 的上述小邻域 \(W\)。两族在 \(W\) 上均由共同的值沿 \(W\) 内道路运输得到，故在整个 \(W\) 上相同，证明了基的交集条件。投影在每个 \(s_m(U)\) 上与 \(U\) 同胚，并有保持模结构的局部同胚
	\[
		U\times\mathcal L_u\longrightarrow p^{-1}(U),\qquad(y,m)\longmapsto s_m(y).
	\]
	这给出所求覆盖。道路可有限细分到这些开集；每一段的提升正沿某个 \(s_m\) 走，其终点就是原来的运输，拼接后便恢复整条道路的运输。

	局部系统态射在这样的邻域中由一点的线性映射决定，因此在所构造的覆盖上连续。若一开始已有局部模平凡化覆盖，再取同时包含在平凡化邻域内的上述小开集，则刚刚构造的 \(s_m\) 恰好是原覆盖的叶，故原拓扑被恢复。结合\cref{g2:prop:local-sheaf-fiber-space}对层与覆盖双向恢复的证明，得到连同态射在内的等价。
\end{proof}

半局部单连通性使小开集内的环路在整个空间中零伦，其运输因而为恒等。配合局部道路连通性，一个纤维中的值便能沿小开集内的道路唯一地展开为局部常值截面。这给出了路径运输与局部常值层在上述局部条件下的对应。

\subsection{连通局部良好空间上的基本群表示}
\label{g2:subsec:0501:03}
现在假设 \(X\) 非空、连通、局部道路连通且半局部单连通。连通和局部道路连通使 \(X\) 道路连通：道路分支都是开集，多于一个分支便会给出分离。任选 \(x_0\in X\)，从这里到每一点都能走道路，所以一个纤维加上环路运输理应包含整个局部系统的信息。需要说明的是，如何从这份信息重建其他纤维，并把道路选择的歧义准确地除掉。

\begin{theorem}[局部系统的分类]\label{g2:thm:local-system-classification}
	设 \(X\) 非空、连通、局部道路连通且半局部单连通，\(R\) 为含幺交换环，\(x_0\in X\)，\(\pi=\pi_1(X,x_0)\)。取基点纤维及其单值
	\begin{equation}\label{g2:eq:local-right-action}
		M=\mathcal L_{x_0},\qquad m\cdot g=T_g(m)
	\end{equation}
	给出以下三类数据之间的等价：\(X\) 上的 \(R\) 模局部系统、\(X\) 上的局部常值 \(R\) 模层、以及带 \(R\) 线性右 \(\pi\) 作用的 \(R\) 模。第三类数据的态射为保持右作用的 \(R\) 线性映射。

	具体地，对给定的右 \(\pi\) 模 \(M\)，取道路类万有覆盖 \(p:\widetilde X\to X\)，其中 \(g[\gamma]=[g*\gamma]\) 为左覆盖变换作用。相应的系数纤维空间为
	\begin{equation}\label{g2:eq:associated-local-system}
		E_M=(M\times\widetilde X)/{\sim},\qquad(m\cdot g,\widetilde x)\sim(m,g\widetilde x),
	\end{equation}
	\(M\) 取离散拓扑，\(E_M\) 取商拓扑，投影为 \([m,\widetilde x]\mapsto p(\widetilde x)\)。
\end{theorem}
\begin{proof}
	先检查右作用。由\eqref{g2:eq:local-transport-composition}，\((m\cdot g)\cdot h=T_hT_g(m)=T_{gh}(m)=m\cdot(gh)\)，而单位环路作用为恒等。局部系统态射在 \(x_0\) 处与每条环路的运输交换，所以得到保持此右作用的线性映射。

	反向构造用\cref{g2:thm:universal-cover,g2:prop:monodromy-deck}已证明的万有覆盖及左作用。式\eqref{g2:eq:associated-local-system}的等价关系正是左作用
	\[
		g\cdot(m,\widetilde x)=(m\cdot g^{-1},g\widetilde x)
	\]
	的轨道关系：先作用 \(h\) 再作用 \(g\)，第一坐标成为 \(m\cdot h^{-1}g^{-1}=m\cdot(gh)^{-1}\)，第二坐标成为 \(gh\widetilde x\)。每个作用映射是同胚，故轨道商映射是开映射，其开集饱和为这些同胚像的并。

	选万有覆盖的一个道路连通均匀覆盖邻域 \(U\) 及其上一叶，记相应的逆映射为 \(s:U\to\widetilde X\)。\(\pi\) 在每个万有覆盖纤维上自由且传递：对于同终点道路类 \([\alpha],[\beta]\)，唯一的 \(g=[\beta*\bar\alpha]\) 满足 \(g[\alpha]=[\beta]\)。覆盖变换因此把所选叶送到全部其他叶。每个 \(E_M\) 中底点位于 \(U\) 的元素，都唯一写成 \([m,s(y)]\)，因为 \([m,gs(y)]=[m\cdot g,s(y)]\)。于是
	\[
		U\times M\longrightarrow E_M|_U,\qquad(y,m)\longmapsto[m,s(y)]
	\]
	为双射。它连续，并且把每个 \(V\times\{m\}\)、\(V\subset U\) 开，送到开集：该集合在商前的饱和集是 \(\{m\cdot g^{-1}\}\times gs(V)\) 对全部 \(g\) 的并。故这是同胚。换一叶只把纤维坐标改为 \(m\mapsto m\cdot g\)，这是线性同构，所以每个纤维的模结构与选择无关。

	用 \(m\mapsto[m,\widetilde x_0]\)、\(\widetilde x_0=[e_{x_0}]\)，识别基点纤维与 \(M\)。代表 \(g\) 的环路从 \(\widetilde x_0\) 提升后终点为 \(g\widetilde x_0\)，因此 \(E_M\) 中相应运输的终值为
	\[
		[m,g\widetilde x_0]=[m\cdot g,\widetilde x_0],
	\]
	恰好恢复所给右作用。

	若 \(M\) 来自已有局部系统 \(\mathcal L\)，定义
	\[
		\Phi:E_M\longrightarrow\coprod_x\mathcal L_x,\qquad
		\Phi([m,[\gamma]])=T_\gamma(m).
	\]
	它尊重商关系，因为 \(T_{g*\gamma}(m)=T_\gamma(m\cdot g)\)。固定某条从 \(x_0\) 到 \(x\) 的道路，\(T_\gamma\) 给出每个纤维上的线性双射。在\cref{g2:prop:local-system-realization}的局部平凡化中，\(\Phi\) 是由固定线性同构给出的乘积映射，故连续，且逆连续；沿任一道路接上尾段也说明它保持运输。因此恢复的是原局部系统。

	最后，保持右作用的线性映射 \(u:M\to N\) 诱导 \([m,\widetilde x]\mapsto[u(m),\widetilde x]\)，由商关系可知良定义，并在局部平凡化中连续且线性。任意局部系统态射都由基点分量唯一确定，因为沿 \(\gamma:x_0\to x\) 的自然性要求
	\[
		u_x=T^{\mathcal K}_\gamma\,u_{x_0}\,(T^{\mathcal L}_\gamma)^{-1}.
	\]
	上述构造给出每一个保持右作用的基点映射，所以不只对象的同构类对应，态射也一一对应。局部常值层的等价已由\cref{g2:prop:local-system-realization}证明。
\end{proof}

按已经固定的道路拼接顺序，\(g\mapsto T_g\) 满足 \(T_{gh}=T_hT_g\)，因而给出右作用。取逆便将复合顺序转换为通常的左表示：
\begin{equation}\label{g2:eq:local-left-representation}
	\rho:\pi\longrightarrow\operatorname{Aut}_R(M),\qquad\rho(g)=T_g^{-1}.
\end{equation}
此时 \(\rho(gh)=\rho(g)\rho(h)\)。当 \(M=R^r\) 并选定基时，这成为到 \(\operatorname{GL}_r(R)\) 的表示；换纤维基使表示同时共轭，而非产生一个新的局部系统同构类。群丛、群环模与局部系数的关联还可参见 \cite[第3.H节]{Hatcher2002}；比较其中的公式时须一并转换左右作用。

\begin{example}[圆周上的符号系统]\label{g2:exm:sign-local-system}
	取 \(X=S^1\)、\(R=\mathbb Z\)、\(M=\mathbb Z\)。由\cref{g2:prop:circle-fundamental-group}，\(\pi_1(S^1,1)=\mathbb Z\)，正向生成元记为 \(t\)。令 \(m\cdot t=-m\)，得到一个局部系统 \(\mathbb Z_-\)。在 \(\widetilde X=\mathbb R\)、\(p(s)=\exp(2\pi\mathrm{i}s)\) 的模型中，它的纤维空间为
	\[
		(\mathbb Z\times\mathbb R)/\bigl((m,s+1)\sim(-m,s)\bigr).
	\]
	沿一小段圆弧，整数系数保持不变；绕圆一周以后，同一起点纤维内的整数变成相反数。它不能同构于 \(\underline{\mathbb Z}\)：若基点同构为 \(u\)，自然性将要求 \(u(-m)=u(m)\)，取 \(m=1\) 即与 \(u\) 单射矛盾。它的每个纤维仍然都是 \(\mathbb Z\)，可见只列出纤维的同构类型不够。
\end{example}

\subsection{换基点和同构的相容}
\label{g2:subsec:0501:04}
分类中选择基点，是为了用一个模保存全部纤维的信息。改用另一个基点不会改变局部系统，但比较两个模时必须给出连接道路。\cref{g2:prop:basepoint-change}已用往返道路的消去证明基本群的换基点公式，下面将它与系数运输同时写出。

\begin{proposition}[换基点与单值]\label{g2:prop:basepoint-local-system}
	设 \(\mathcal L\) 为拓扑空间 \(X\) 上的 \(R\) 模局部系统，\(c:x_0\to x_1\) 为道路，\(M_i=\mathcal L_{x_i}\)。定义
	\[
		\psi_c:\pi_1(X,x_1)\longrightarrow\pi_1(X,x_0),\qquad
		\psi_c([h])=[c*h*\bar c].
	\]
	这是群同构，且 \(T_c:M_0\to M_1\) 满足
	\begin{equation}\label{g2:eq:basepoint-module-transport}
		T_c(m\cdot\psi_c(h))=T_c(m)\cdot h.
	\end{equation}
	若 \(c':x_0\to x_1\) 为另一条道路，令 \(g=[c'*\bar c]\in\pi_1(X,x_0)\)，则
	\[
		T_{c'}=T_cT_g,\qquad\psi_{c'}(h)=g\psi_c(h)g^{-1}.
	\]
	局部系统态射 \(u:\mathcal L\to\mathcal K\) 满足 \(u_{x_1}T_c^{\mathcal L}=T_c^{\mathcal K}u_{x_0}\)，所以换基点同时保持态射以及同构的对应。
\end{proposition}
\begin{proof}
	\(\psi_c\) 是\cref{g2:prop:basepoint-change}中从 \(x_0\) 到 \(x_1\) 的换基点同构的逆。直接展开可见
	\[
		T_{c*h*\bar c}=T_c^{-1}T_hT_c.
	\]
	两边左复合 \(T_c\) 就得到\eqref{g2:eq:basepoint-module-transport}。又因 \(c'\simeq(c'*\bar c)*c\)，运输公式给出 \(T_{c'}=T_cT_g\)。将 \(c'\simeq g*c\) 及其反向道路代入 \(c'*h*\bar c'\)，得到 \(\psi_{c'}(h)=g\psi_c(h)g^{-1}\)。最后的等式就是态射沿 \(c\) 的自然性方形。因此道路选择造成的两项变化彼此匹配，不会改变所描述的局部系统。
\end{proof}

若再接一条 \(d:x_1\to x_2\)，则 \(T_{c*d}=T_dT_c\)，同时 \(\psi_{c*d}=\psi_c\psi_d\)。这说明连续两次换基点与一次沿拼接道路换基点相容。纤维间的比较随连接道路而定；即使基本群交换，非恒等的环路运输仍会改变这个比较。

\begin{proposition}[整体截面与平凡性]\label{g2:prop:parallel-sections}
	设 \(X\) 非空、连通、局部道路连通且半局部单连通，\(\mathcal L\) 为 \(R\) 模局部系统，\(\mathscr F\) 为相应的局部常值层。固定 \(x_0\)，记 \(M=\mathcal L_{x_0}\)、\(\pi=\pi_1(X,x_0)\)。在 \(x_0\) 求值给出自然同构
	\[
		\mathscr F(X)\cong M^\pi,
		\qquad M^\pi=\{m\in M:m\cdot g=m\text{ 对每个 }g\in\pi\}.
	\]
	局部系统 \(\mathcal L\) 平凡，当且仅当每个基点环路作用均为恒等。
\end{proposition}
\begin{proof}
	把整体截面看成纤维覆盖的连续截面。它沿任何道路的限制都是从给定初值出发的提升，由提升唯一性，其终值必为初值的运输；所以基点值确定整个截面，而且被每条环路固定。反过来，给定 \(m\in M^\pi\)，对 \(x\in X\) 选 \(\gamma:x_0\to x\)，令 \(s(x)=T_\gamma m\)。若改用 \(\eta\)，则 \(g=[\gamma*\bar\eta]\) 满足 \([\gamma]=[g*\eta]\)，于是 \(T_\gamma m=T_\eta(m\cdot g)=T_\eta m\)，定义与选择无关。在局部平凡化邻域内，它正是沿小道路运输同一个值所得的截面，故连续。两构造互逆且线性。

	常值系统的环路运输均为恒等，同构保持这一性质。若全部环路作用为恒等，上述方法可以同时用于每个 \(m\in M\)，给出 \(\underline M\to\mathcal L\) 的纤维线性同构，并与道路运输交换，故 \(\mathcal L\) 平凡。
\end{proof}

例如，\(\mathbb Z_-\) 的整体截面只有零，因为 \(m=-m\) 在 \(\mathbb Z\) 中只允许 \(m=0\)。在单连通的局部良好空间上，每个局部系统则平凡；选定一点的纤维后，沿道路运输给出不依赖道路选择的比较。这也是把系数拉到万有覆盖上以后能够使用一个固定模的原因。
% 纲要条目排为编号小节标题，并进入目录与书签；说明段仍保留于纲要注释中。

% !TeX root = ../../Book2.tex
% 纲要标题：### 5.2 扭系数链与上链
% 纲要位置：几何学计划纲要.md，第 1202 行。
% 正文状态：未撰写；本文件目前是框架。

% 纲要内容开始
% 
% * 万有覆盖链复形的群环作用
% * 左右作用及张量、Hom 的约定
% * 局部系数边界和连接映射
% * 同伦不变性与切除
% 
% 纲要内容结束

\section{扭系数链与上链}
\label{g2:sec:0502}

待撰写

\subsection{万有覆盖链复形的群环作用}
\label{g2:subsec:0502:01}
待撰写

\subsection{左右作用及张量、Hom 的约定}
\label{g2:subsec:0502:02}
待撰写

\subsection{局部系数边界和连接映射}
\label{g2:subsec:0502:03}
待撰写

\subsection{同伦不变性与切除}
\label{g2:subsec:0502:04}
待撰写
% 纲要条目排为编号小节标题，并进入目录与书签；说明段仍保留于纲要注释中。

% !TeX root = ../../Book2.tex
% 纲要标题：### 5.3 定向系数和循环模型
% 纲要位置：几何学计划纲要.md，第 1209 行。
% 正文状态：未撰写；本文件目前是框架。

% 纲要内容开始
% 
% * 流形局部同调给出的定向系统
% * 圆周单值下的核、余核计算
% * Möbius 型和不可定向曲面的计算
% * 平凡化系数和遗失的单值信息
% 
% 纲要内容结束

\section{定向系数和循环模型}
\label{g2:sec:0503}

待撰写

\subsection{流形局部同调给出的定向系统}
\label{g2:subsec:0503:01}
待撰写

\subsection{圆周单值下的核、余核计算}
\label{g2:subsec:0503:02}
待撰写

\subsection{Möbius 型和不可定向曲面的计算}
\label{g2:subsec:0503:03}
待撰写

\subsection{平凡化系数和遗失的单值信息}
\label{g2:subsec:0503:04}
待撰写
% 纲要条目排为编号小节标题，并进入目录与书签；说明段仍保留于纲要注释中。

% !TeX root = ../../Book2.tex
% 纲要标题：### 5.4 纤维上同调如何变化
% 纲要位置：几何学计划纲要.md，第 1216 行。
% 正文状态：未撰写；本文件目前是框架。

% 纲要内容开始
% 
% * 纤维运输在同调上的作用
% * 映射环面的扭曲
% * 为 Serre 谱序列和障碍理论准备的系数函子
% 
% ---
% 
% 纲要内容结束

\section{纤维上同调如何变化}
\label{g2:sec:0504}

待撰写

\subsection{纤维运输在同调上的作用}
\label{g2:subsec:0504:01}
待撰写

\subsection{映射环面的扭曲}
\label{g2:subsec:0504:02}
待撰写

\subsection{为 Serre 谱序列和障碍理论准备的系数函子}
\label{g2:subsec:0504:03}
待撰写
% 纲要条目排为编号小节标题，并进入目录与书签；说明段仍保留于纲要注释中。


\begin{chaptersummary}
	局部系统给每点一个模，给每条道路类一个运输同构。在连通、局部道路连通且半局部单连通的空间上，选定基点后，全部数据等价于一个带基本群右作用的模；覆盖纤维空间和局部常值层给出另外两种描述。换基点通过一条连接道路比较纤维，换道路时单值和基本群同构同时改变。

	\ProseParagraphBreak

	万有覆盖的链是左群环自由模。将它与右系数模取张量得到扭系数链；取等变同态时，须先把系数通过逆作用改写为左模。几何上，系数放在单形首顶点：删掉首顶点的面要向前运输系数，上链则要把面上的值逆向送回。单形内的运输相容性保证边界平方为零，棱柱和小链同伦则给出同伦不变性与切除。

	\ProseParagraphBreak

	圆周把全部单值集中在一个自同构 \(T\) 上。一次同调和零次上同调取 \(\ker(T-1)\)，零次同调和一次上同调取 \(\operatorname{coker}(T-1)\)。前者要求一个值绕行后不变，后者把绕行前后的差值置零。无边界流形的局部顶次同调组成定向系统；它在不可定向闭曲面上补偿局部方向的反转，使顶次整数类重新出现。

	\ProseParagraphBreak

	局部平凡纤维丛沿道路给出纤维的同伦运输。同调向前推送，上同调则通过逆运输的拉回向前比较。映射环面的正合列把这些系数系统在圆周上的核、余核送入总空间计算。局部系统保留单值，但不包含丛的全部粘合信息；同调中有不变元也不保证存在几何截面。后续谱序列与障碍理论将在这些明确的运输数据上继续组织更高层信息。
\end{chaptersummary}

\BookExercises
除另行说明外取整数系数，路径拼接与左右作用沿用正文约定。计算扭系数时先写出单值及相应复形；由短正合列得到抽象群的直和分解时，说明分裂的依据。

\begin{exercise}\label{g2:ex:constant-presheaf-gluing}
	设 \(A\) 为非零 \(R\) 模，\(X=U\amalg V\)，其中 \(U,V\) 为非空连通开集。给每个非空开集 \(W\subset X\) 指定 \(P(W)=A\)，并把非空开集之间的限制映射设为恒等，另令 \(P(\varnothing)=0\)。证明这些数据不构成层；计算真正的常值层 \(\underline A(X)\)。再证明 \(P\) 与 \(\underline A\) 在每一点的茎同构，解释为何单独知道全部茎的同构类型不能判定粘合性质。
\end{exercise}
\begin{solution}
	取两个不同的元素 \(a,b\in A\)，分别作为 \(U,V\) 上的截面。它们在空交集上的限制同为零，满足相容条件；若能由 \(c\in P(X)=A\) 粘合，则恒等限制同时要求 \(c=a\) 与 \(c=b\)，矛盾。所以 \(P\) 不满足层的粘合条件。

	一个到离散模 \(A\) 的局部常值函数在连通集上必为常值：任意取值的逆像及其补集均开，若两者都非空就分离了该连通集。因此 \(\underline A(X)=A\times A\)，两坐标为在 \(U,V\) 上的值，任意一对都能唯一粘成局部常值函数。

	在一点 \(x\) 附近，\(P\) 的每个代表只是一个 \(a\in A\)，其进一步限制不改变 \(a\)，故 \(P_x\cong A\)。一个局部常值函数在 \(x\) 的某个邻域上恒等于其值，所以 \(\underline A_x\cong A\)。把 \(a\) 送到恒值函数还给出自然映射 \(P\to\underline A\)，它在每个茎上为同构，在 \(X\) 上却只是对角映射 \(A\to A\times A\)。茎记录已有截面的局部相等关系，不能自行保证不同开集上的相容截面已有一个整体代表；后一项正是这里缺少的粘合公理。
\end{solution}

\begin{exercise}\label{g2:ex:unipotent-circle-system}
	设 \(\mathcal L\) 为圆周上的 \(\mathbb Q\) 模局部系统，基点纤维为列向量空间 \(\mathbb Q^2\)，正向生成环路的运输矩阵为
	\[
		J=\begin{pmatrix}1&1\\0&1\end{pmatrix}.
	\]
	对整数 \(k\) 定义 \(f_k:S^1\to S^1\)、\(f_k(z)=z^k\)，其中 \(f_0\) 为常值映射。求 \(f_k^*\mathcal L\) 的单值及整体截面，判断哪些 \(k\) 给出平凡系统，并证明当 \(k\ne0\) 时 \(f_k^*\mathcal L\cong\mathcal L\)。
\end{exercise}
\begin{solution}
	正向环路经 \(f_k\) 绕行 \(k\) 周，故拉回系统的生成运输为 \(J^k\)。写 \(J=I+N\)，有 \(N^2=0\) 及 \((I+N)^{-1}=I-N\)，所以对所有整数 \(k\) 都有
	\[
		J^k=I+kN=\begin{pmatrix}1&k\\0&1\end{pmatrix}.
	\]
	由\cref{g2:prop:parallel-sections}，整体截面对应于 \(\ker(J^k-I)\)。当 \(k=0\) 时这是全部 \(\mathbb Q^2\)，系统平凡；当 \(k\ne0\) 时这是 \(\mathbb Q(1,0)^{\mathsf T}\)，且 \(J^k\ne I\)，故系统不平凡。

	对 \(k\ne0\)，令 \(S=\operatorname{diag}(k,1)\)，则 \(S\) 可逆，且 \(SJ=J^kS\)。这说明 \(S\) 保持原系统与拉回系统的右生成元作用，因此由\cref{g2:thm:local-system-classification}延拓为局部系统同构。拉回改变了指定基下的单值矩阵，但并不一定改变局部系统的同构类。
\end{solution}

\begin{exercise}\label{g2:ex:projective-rank-one-system}
	分类 \(\mathbb{RP}^2\) 上纤维同构于 \(\mathbb Z\) 的整数局部系统。对反极点覆盖 \(p:S^2\to\mathbb{RP}^2\)，计算这些系统的拉回，并比较它们在 \(\mathbb{RP}^2\) 上的整体截面。
\end{exercise}
\begin{solution}
	\(S^2\) 单连通，而 \(p\) 是二叶覆盖，所以它是万有覆盖，覆盖变换为恒等和反极点映射。由\cref{g2:prop:monodromy-deck}，\(\pi_1(\mathbb{RP}^2)\cong\mathbb Z/2\mathbb Z\)；同一计算也见\cref{g2:ex:projective-pi1}。两个空间均是连通流形，具备分类定理要求的局部条件。整数模 \(\mathbb Z\) 的自同构只有乘以 \(1\) 和乘以 \(-1\)，两者的平方均为恒等，所以恰有两种右作用。它们不共轭，因为 \(\operatorname{Aut}_{\mathbb Z}(\mathbb Z)=\{1,-1\}\) 是交换群。因此有且仅有两个同构类，分别由平凡单值与符号单值给出。

	两个系统拉回到 \(S^2\) 后均平凡，因为 \(\pi_1(S^2)=0\)，没有非平凡环路作用。原空间上，平凡系统的整体截面为 \(\mathbb Z\)，符号系统的整体截面为 \(\{m\in\mathbb Z:m=-m\}=0\)。所以在万有覆盖上平凡，并不意味着原局部系统平凡；下降回底空间时仍需记住覆盖变换怎样作用于系数。
\end{solution}

\begin{exercise}\label{g2:ex:noncommuting-transport}
	令 \(X=S^1\vee S^1\)，取两条正向环路 \(a,b\) 为自由生成元。设 \(\mathbb Q^2\) 上的生成运输为
	\[
		T_a=A=\begin{pmatrix}1&1\\0&1\end{pmatrix},\qquad
		T_b=B=\begin{pmatrix}1&0\\1&1\end{pmatrix}.
	\]
	说明这些数据确定一个局部系统，计算 \(T_{a*b}\) 和 \(T_{a*b*\bar a*\bar b}\)。若换基点的两条道路 \(c,c'\) 满足 \([c'*\bar c]=[a*b]\)，写出 \(T_{c'}\) 与 \(T_c\) 的关系，并说明为何不能把 \(T_aT_b\) 直接当成 \(T_{a*b}\)。
\end{exercise}
\begin{solution}
	\cref{g2:thm:graph-free-group}给出 \(\pi_1(X)=\langle a,b\rangle\)，没有关系需要检验。可先定义左表示 \(\rho(a)=A^{-1}\)、\(\rho(b)=B^{-1}\)，由自由群性质延拓到整个基本群，再令 \(T_g=\rho(g)^{-1}\)。它满足 \(T_{gh}=T_hT_g\)，所以由\cref{g2:thm:local-system-classification}确定所求局部系统。

	按行进顺序复合运输，得到
	\[
		T_{a*b}=BA=\begin{pmatrix}1&1\\1&2\end{pmatrix},\qquad
		T_{a*b*\bar a*\bar b}=B^{-1}A^{-1}BA
		=\begin{pmatrix}0&-1\\1&3\end{pmatrix}.
	\]
	这里可先计算 \(BA\)，左乘 \(A^{-1}\) 得 \(\left(\begin{smallmatrix}0&-1\\1&2\end{smallmatrix}\right)\)，再左乘 \(B^{-1}\) 得所示矩阵。换基点公式给出 \(T_{c'}=T_cT_{a*b}=T_cBA\)。函数复合先应用右边矩阵，故先沿 \(a\) 运输再沿 \(b\) 运输必须是 \(BA\)。由于 \(AB=\left(\begin{smallmatrix}2&1\\1&1\end{smallmatrix}\right)\ne BA\)，在这个例子中倒转顺序会实质改变答案，而不只是改变记号。
\end{solution}

\begin{exercise}\label{g2:ex:twisted-interval-pair}
	设 \(I=[0,1]\)，\(A=\{0,1\}\)，\(\mathcal L\) 为 \(I\) 上的 \(R\) 模局部系统，\(M=\mathcal L_0\)，\(T:\mathcal L_0\to\mathcal L_1\) 为沿正向区间的运输。计算 \(H_n(I,A;\mathcal L)\) 和 \(H^n(I,A;\mathcal L)\)，并在不把两个端点纤维预先视为同一个模的情况下写出同调连接映射。
\end{exercise}
\begin{solution}
	沿区间内唯一的同端点道路类运输，可以把整个系统平凡化为 \(M\)。区间和端点都是 CW 空间，所以\cref{g2:thm:twisted-cellular-comparison}适用。相对胞腔复形只保留一条一胞腔，系数为 \(M\)，边界为零，故
	\[
		H_n(I,A;\mathcal L)\cong H^n(I,A;\mathcal L)\cong
		\begin{cases}M,&n=1,\\0,&n\ne1.\end{cases}
	\]
	这里上同调的一次识别也使用从首顶点出发的平凡化；相对上链在端点上为零，故只有一维项。若用正向一胞腔把一次相对同调识别为 \(M\)，\eqref{g2:eq:twisted-boundary}给出
	\[
		\partial_*:M\longrightarrow H_0(A;\mathcal L)=\mathcal L_0\oplus\mathcal L_1,
		\qquad m\longmapsto(-m,Tm).
	\]
	两个端点的系数不能直接相减；它们是落在两个直和项中的元素。映入 \(H_0(I;\mathcal L)\) 后，正向路径的边界恰把 \(m\) 与 \(Tm\) 识别起来，所以这对元素的像为零，符合长正合列。
\end{solution}

\begin{exercise}\label{g2:ex:two-loop-local-matrices}
	设 \(X=S^1\vee S^1\)，两条正向环路记为 \(a,b\)。在基点纤维 \(M=\mathbb Z^2\) 上规定运输
	\[
		T_a(x,y)=(-x,y),\qquad T_b(x,y)=(x,-y).
	\]
	计算所得局部系统 \(\mathcal L\) 的全部同调与上同调群，并分别给出 \(H_1\) 的生成元和 \(H^1\) 的自由部分、挠部分的代表。
\end{exercise}
\begin{solution}
	楔和有一个零胞腔、两条一胞腔，没有更高胞腔。两条提升边的边界分别为 \((a-1)\widetilde v\)、\((b-1)\widetilde v\)。同调复形为
	\[
		0\longrightarrow M\oplus M\xrightarrow{\mathrm d}M\longrightarrow0,
		\qquad \mathrm d((x_1,x_2),(y_1,y_2))=(-2x_1,-2y_2).
	\]
	因此核由“第一条边的系数 \((0,1)\)”和“第二条边的系数 \((1,0)\)”生成，像为 \(2\mathbb Z\oplus2\mathbb Z\)，得到
	\[
		H_1(X;\mathcal L)=\mathbb Z^2,\qquad
		H_0(X;\mathcal L)=(\mathbb Z/2)^2,
		\qquad H_n(X;\mathcal L)=0\quad(n\ge2).
	\]
	由于两个运输都是自身的逆，上链复形为
	\[
		0\longrightarrow M\xrightarrow{\delta}M\oplus M\longrightarrow0,
		\qquad\delta(u,v)=((-2u,0),(0,-2v)).
	\]
	其核为零；商群中第一条边的第二坐标、第二条边的第一坐标给出两个自由生成元，而第一条边的第一坐标、第二条边的第二坐标分别只按偶数关系取商，给出两个二阶生成元。因此
	\[
		H^0(X;\mathcal L)=0,\qquad
		H^1(X;\mathcal L)=\mathbb Z^2\oplus(\mathbb Z/2)^2,
		\qquad H^n(X;\mathcal L)=0\quad(n\ge2).
	\]
	这里零次同调与零次上同调已不同：前者允许把运输差值置零，后者要求一个值同时被两个运输固定。
\end{solution}

\begin{exercise}\label{g2:ex:regular-local-coefficients-circle}
	令 \(\Gamma=\mathbb Z[t,t^{-1}]\)，在圆周上取基点纤维为右正则模 \(M=\Gamma\)、正向单值为 \(m\mapsto mt\) 的局部系统 \(\mathcal R\)。计算 \(H_*(S^1;\mathcal R)\) 与 \(H^*(S^1;\mathcal R)\)，并说明为什么后一组群不能直接解释为万有覆盖 \(\mathbb R\) 的普通整数上同调。
\end{exercise}
\begin{solution}
	胞腔链只有一次和零次，边界为乘 \(t-1\)。Laurent 多项式环是整环，故这个乘法单射；取值 \(t=1\) 的同态满射到 \(\mathbb Z\)，核为 \((t-1)\Gamma\)。核的断言可将 Laurent 多项式乘以足够大的 \(t\) 次幂，化为普通多项式的因式分解；\(t\) 在 \(\Gamma\) 中可逆，故不会改变该理想。因此
	\[
		H_0(S^1;\mathcal R)=\mathbb Z,\qquad H_n(S^1;\mathcal R)=0\quad(n>0).
	\]
	上链余边界为乘 \(t^{-1}-1=-t^{-1}(t-1)\)。因 \(-t^{-1}\) 是单位，这个映射仍单射且像仍为 \((t-1)\Gamma\)，所以
	\[
		H^1(S^1;\mathcal R)=\mathbb Z,\qquad H^n(S^1;\mathcal R)=0\quad(n\ne1).
	\]
	同调模型中 \(\Gamma\otimes_\Gamma C_*(\mathbb R;\mathbb Z)\cong C_*(\mathbb R;\mathbb Z)\)，所以同调确实恢复万有覆盖的普通整数同调。上链模型却是 \(\operatorname{Hom}_\Gamma(C_*(\mathbb R;\mathbb Z),\Gamma^-)\)，要求等变性且取值于群环；普通上链为 \(\operatorname{Hom}_{\mathbb Z}(C_*(\mathbb R;\mathbb Z),\mathbb Z)\)，两者的对象、取值和等变条件都不同。\(\mathbb R\) 可缩，其普通整数上同调仅在零次为 \(\mathbb Z\)，与上述结果并不矛盾。
\end{solution}

\begin{exercise}\label{g2:ex:sign-coefficient-connecting}
	在圆周上令 \(\mathcal Z_-\) 为单值 \(-1\) 的整数局部系统，\(\mathbb F_2\) 为常值系统。模 \(2\) 把 \(-1\) 化为 \(1\)，故有局部系统短正合列
	\[
		0\longrightarrow\mathcal Z_-\xrightarrow{\times2}\mathcal Z_-\longrightarrow\mathbb F_2\longrightarrow0.
	\]
	证明它给出同调及上同调的系数长正合列，并计算连接同态
	\[
		H_1(S^1;\mathbb F_2)\longrightarrow H_0(S^1;\mathcal Z_-),\qquad
		H^0(S^1;\mathbb F_2)\longrightarrow H^1(S^1;\mathcal Z_-).
	\]
\end{exercise}
\begin{solution}
	局部系数链逐单形取纤维的直和，所以上述系数序列给出短正合链复形。上链逐单形取纤维的直积；每个模 \(2\) 值可选整数代表，因而末项仍满射，也给出短正合上链复形。运输相容性使系数映射与边界、余边界交换。于是\cref{g2:thm:homology-LES}的提升构造及其上链形式适用：先把模 \(2\) 的圈或上圈提升为整数链或上链，再取边界，其值必为偶数，最后除以 \(2\) 得到连接类。

	用一顶点一边的胞腔模型，符号系数的链边界和上链余边界均为乘 \(-2\)，模 \(2\) 系数的两种微分均为零。因此
	\[
		H_1(S^1;\mathbb F_2)=H^0(S^1;\mathbb F_2)=\mathbb F_2,
		\qquad H_0(S^1;\mathcal Z_-)=H^1(S^1;\mathcal Z_-)=\mathbb Z/2.
	\]
	一次模 \(2\) 链生成元可提升为系数 \(1\) 的整数边，其边界为 \(-2\)，除以 \(2\) 得到 \(-1\)，代表 \(\mathbb Z/2\) 的非零元。零次模 \(2\) 上链生成元提升为顶点值 \(1\)，余边界同样为 \(-2\)，所以第二个连接映射也把 \(1\) 送到 \([-1]\)。两映射都是同构。胞腔比较可在选取 \(J\) 后与所有系数映射同时使用，故胞腔中算出的连接类就是奇异系数正合列中的连接类。
\end{solution}

\begin{exercise}\label{g2:ex:integer-unipotent-circle}
	取整数 \(q\geq2\)，令 \(A=\mathbb Z^2\)，单值自同构为
	\[
		T(x,y)=(x+qy,y).
	\]
	计算圆周上的 \(H_k(S^1;\mathcal L_T)\) 和 \(H^k(S^1;\mathcal L_T)\)，并给出核、余核中所用的生成元。再将系数逐纤维改成 \(\mathbb Q\)，比较结果。说明这次出现的挠元怎样来自单值。
\end{exercise}
\begin{solution}
	记标准基为 \(e_1,e_2\)，则 \((T-1)(x,y)=(qy,0)\)，所以
	\[
		\ker(T-1)=\mathbb Ze_1,\qquad
		(T-1)A=q\mathbb Ze_1,\qquad
		A/(T-1)A=(\mathbb Z/q)[e_1]\oplus\mathbb Z[e_2].
	\]
	由\cref{g2:prop:circle-local-coefficients}，\(H_1=H^0=\mathbb Ze_1\)，\(H_0=H^1=(\mathbb Z/q)[e_1]\oplus\mathbb Z[e_2]\)，更高次数为零。上链微分用 \(T^{-1}-1\)，其公式为 \((x,y)\mapsto(-qy,0)\)，确实具有相同的核与像。

	改为有理数系数后，像为整个 \(\mathbb Qe_1\)，核为 \(\mathbb Qe_1\)，余核为 \(\mathbb Q[e_2]\)。因此零次、一次同调与上同调都同构于 \(\mathbb Q\)，更高次数为零。整数挠元 \([e_1]\) 的阶恰为 \(q\)：绕行把 \(e_2\) 变成 \(e_2+qe_1\)，零次同调要求两者相等，留下关系 \(qe_1=0\)，却没有关系 \(e_1=0\)。这说明扭系数计算中的挠元不仅能来自空间本身的整数同调，也能来自单值矩阵的整数像。
\end{solution}

\begin{exercise}\label{g2:ex:surface-other-sign-systems}
	设 \(h\geq2\)，在 \(N_h\) 的标准边 \(a_1,\ldots,a_h\) 上指定符号 \(\varepsilon_i\in\{1,-1\}\)，且至少一个符号为 \(1\)，至少一个为 \( -1\)。证明这些符号定义一个秩一整数局部系统 \(\mathcal L_\varepsilon\)，并计算所有 \(H_k(N_h;\mathcal L_\varepsilon)\)。将所得结果与常值系统及定向系统比较。
\end{exercise}
\begin{solution}
	关系 \(a_1^2\cdots a_h^2=1\) 在符号作用下变成 \(\varepsilon_1^2\cdots\varepsilon_h^2=1\)，所以各生成元的指定作用满足群关系，确实给出一个右 \(\mathbb Z[\pi_1(N_h)]\) 模，进而由\cref{g2:thm:local-system-classification}得到局部系统。

	令 \(I=\{i:\varepsilon_i=-1\}\)、\(J=\{i:\varepsilon_i=1\}\)，两集合均非空。提升多边形公式\eqref{g2:eq:lifted-word-boundary}中，各前缀 \(Q_i=a_1^2\cdots a_{i-1}^2\) 的作用均为 \(1\)，故胞腔复形为
	\[
		0\longrightarrow\mathbb Z
		\xrightarrow{\ 1\mapsto2\sum_{j\in J}e_j\ }\mathbb Z^h
		\xrightarrow{\ (x_i)\mapsto-2\sum_{i\in I}x_i\ }\mathbb Z
		\longrightarrow0.
	\]
	第一映射因 \(J\ne\varnothing\) 单射，所以 \(H_2=0\)；末端映射的像因 \(I\ne\varnothing\) 恰为 \(2\mathbb Z\)，所以 \(H_0=\mathbb Z/2\)。

	固定 \(i_0\in I\)，中间核有一组整数基
	\[
		e_i-e_{i_0}\quad(i\in I\setminus\{i_0\}),\qquad e_j\quad(j\in J).
	\]
	令 \(u_J=\sum_{j\in J}e_j\)。在后一组基中，选定一个 \(j_0\in J\)，以 \(u_J\) 代替 \(e_{j_0}\) 仍为整数基，因为 \(e_{j_0}=u_J-\sum_{j\ne j_0}e_j\)。二维边界仅给出关系 \(2u_J=0\)，于是
	\[
		H_1(N_h;\mathcal L_\varepsilon)\cong\mathbb Z^{h-2}\oplus\mathbb Z/2.
	\]
	高于二次没有胞腔，故其群为零。所有符号为 \(1\) 时是常值系统，结果为 \(H_0=\mathbb Z\)、\(H_1=\mathbb Z^{h-1}\oplus\mathbb Z/2\)、\(H_2=0\)；所有符号为 \( -1\) 时是定向系统，结果为 \(H_0=\mathbb Z/2\)、\(H_1=\mathbb Z^{h-1}\)、\(H_2=\mathbb Z\)。因此不是任意非平凡秩一系统都能恢复顶次整数类，定向系统的符号与几何附着相匹配才产生这一效果。
\end{solution}

\begin{exercise}\label{g2:ex:moebius-relative-orientation}
	令 \(B=(\mathbb R\times[-1,1])/((s,u)\sim(s+1,-u))\) 为闭 Möbius 带，\(r:B\to S^1\) 为核心收缩，取 \(\mathcal L=r^*\mathbb Z_{-}\)。证明 \(\mathcal L|_{\partial B}\) 为常值系统，并计算 \(H_k(B,\partial B;\mathcal L)\)。在长正合列中明确写出 \(H_0(\partial B;\mathcal L)\to H_0(B;\mathcal L)\)。
\end{exercise}
\begin{solution}
	边界是一个圆周：从 \([0,1]\) 出发沿上边走到 \([1,1]=[0,-1]\)，再沿下边走到 \([1,-1]=[0,1]\)，恰绕完边界一次。投到核心时，这条路径绕行两周。因此边界生成元的单值为 \((-1)^2=1\)，边界上的系统为常值整数系统；选一个边界基点处的生成元，就得到其平凡化。

	由核心收缩和\cref{g2:prop:circle-local-coefficients}，
	\[
		H_1(B;\mathcal L)=0,\qquad H_0(B;\mathcal L)=\mathbb Z/2,\qquad
		H_1(\partial B;\mathcal L)=H_0(\partial B;\mathcal L)=\mathbb Z.
	\]
	沿一条从边界基点到核心基点的路径识别纤维，包含诱导的零次映射将整数 \(m\) 送到其在 \(\mathbb Z/2\) 中的类。换一条识别路径至多反号，在模二余核中不改变该映射。故空间对的长正合列在低次成为
	\[
		\begin{aligned}
			0\longrightarrow H_2(B,\partial B;\mathcal L)\longrightarrow\mathbb Z\longrightarrow0,\\
			0\longrightarrow H_1(B,\partial B;\mathcal L)\longrightarrow\mathbb Z\xrightarrow{\bmod\,2}\mathbb Z/2
			\longrightarrow H_0(B,\partial B;\mathcal L)\longrightarrow0.
		\end{aligned}
	\]
	正合性给出 \(H_2(B,\partial B;\mathcal L)\cong\mathbb Z\)、\(H_1(B,\partial B;\mathcal L)\cong2\mathbb Z\cong\mathbb Z\)、\(H_0(B,\partial B;\mathcal L)=0\)。高次由两端的消失同样为零。这里相对二维类的边界生成 \(H_1(\partial B;\mathcal L)\)，而相对一次类到边界零次群的像是 \(2\mathbb Z\)；两条连接映射的整数倍数不同。
\end{solution}

\begin{exercise}\label{g2:ex:product-orientation-system}
	设 \(M,N\) 分别为连通、Hausdorff、第二可数、无边界的 \(m,n\) 维拓扑流形，\(m,n\geq1\)。固定乘积中先列 \(M\) 坐标、后列 \(N\) 坐标的次序。证明相对外积给出局部系统同构
	\[
		\mathcal O_{M\times N}\cong
		\operatorname{pr}_M^*\mathcal O_M\otimes_{\mathbb Z}\operatorname{pr}_N^*\mathcal O_N.
	\]
	推出 \(M\times N\) 可定向当且仅当 \(M,N\) 都可定向，并解释为何两个不可定向流形相乘不会自动使定向反转抵消。
\end{exercise}
\begin{solution}
	在点 \((x,y)\) 处，两个被删去的单点在 Hausdorff 流形中闭，故其补空间开，且
	\[
		(M\setminus\{x\})\times N\ \cup\ M\times(N\setminus\{y\})
		=(M\times N)\setminus\{(x,y)\}.
	\]
	因此\cref{g2:prop:cross-cup-evaluation}证明中构造的相对同调外积适用，得到
	\[
		\mathcal O_{M,x}\otimes_{\mathbb Z}\mathcal O_{N,y}
		\longrightarrow\mathcal O_{M\times N,(x,y)}.
	\]
	切除到两个坐标球后，\cref{g2:prop:local-cross-normalization}的乘积基本链证明说明：两个局部生成元的外积仍为乘积邻域的生成元，且坐标顺序固定其正号。两端均为无限循环群，因此此映射是同构。

	这些同构还必须保持运输。分别在两个坐标球内用共同相对群延拓生成元，其外积在乘积邻域内仍为同一个相对类在各点的像；相对外积的自然性保证这些延拓与各点同构交换。分段任意乘积路径并在每段使用这种乘积邻域，就得运输相容性。于是逐点同构组成局部系统同构。

	以 \((x_0,y_0)\) 为基点，乘积闭路 \((\alpha,\beta)\) 的作用为
	\[
		w_{M\times N}([\alpha],[\beta])=w_M([\alpha])\,w_N([\beta]).
	\]
	若两个因子都可定向，右侧恒为 \(1\)，乘积可定向。反之，只在第一因子走闭路、第二因子保持不动，可知乘积可定向迫使 \(w_M\equiv1\)；交换两因子的角色得到 \(w_N\equiv1\)。两个反定向闭路同时走时，符号确实相乘为 \(1\)，但乘积空间还包含只在一个因子中绕行的闭路，后者的负号没有抵消。
\end{solution}

\begin{exercise}\label{g2:ex:torus-unipotent-mapping-torus}
	令 \(F=\mathbb R^2/\mathbb Z^2\)，矩阵
	\[
		A=\begin{pmatrix}1&2\\0&1\end{pmatrix}
	\]
	诱导同胚 \(f:F\to F\)。计算映射环面 \(E_f\) 的全部整数同调与上同调群。说明为何纤维和基底的整数同调都无挠，而总空间仍能出现挠元，并判断此丛是否具有截面。
\end{exercise}
\begin{solution}
	矩阵 \(A\) 及其逆都是整数矩阵，故确实在商环面上诱导互逆的连续映射。环面同调为 \(H_0=H_2=\mathbb Z\)、\(H_1=\mathbb Z^2\)。\(f_*\) 在一次同调上的矩阵是 \(A\)，在零次为恒等；在二次为乘 \(\det A=1\)，可由\cref{g2:prop:torus-cohomology-ring}的两个一次类相乘计算得到。于是
	\[
		(A-I)(x,y)=(2y,0),\qquad
		\ker(A-I)=\mathbb Ze_1,\qquad
		\operatorname{coker}(A-I)=(\mathbb Z/2)[e_1]\oplus\mathbb Z[e_2].
	\]
	将这些群放入\cref{g2:thm:mapping-torus-sequences}，各短正合列的商群均为自由阿贝尔群，故可选择基的提升使之分裂。得到
	\[
		H_k(E_f;\mathbb Z)\cong
		\begin{cases}
			\mathbb Z,&k=0,3,\\
			\mathbb Z^2\oplus\mathbb Z/2,&k=1,\\
			\mathbb Z^2,&k=2,\\
			0,&k>3.
		\end{cases}
	\]
	上同调中，\(f^*\) 在对偶一次基上的矩阵为 \(A^{\mathsf T}\)，所以 \((A^{\mathsf T}-I)(u,v)=(0,2u)\)。其核为 \(\mathbb Ze_2^*\)，余核为 \(\mathbb Z[e_1^*]\oplus(\mathbb Z/2)[e_2^*]\)；在零次和二次仍为恒等。上同调短正合列的商群同样自由，故
	\[
		H^k(E_f;\mathbb Z)\cong
		\begin{cases}
			\mathbb Z,&k=0,3,\\
			\mathbb Z^2,&k=1,\\
			\mathbb Z^2\oplus\mathbb Z/2,&k=2,\\
			0,&k>3.
		\end{cases}
	\]
	这里同调挠元来自单值关系 \(f_*e_2-e_2=2e_1\)，上同调挠元来自对偶关系；纤维群本身自由，并不保证运输差值的像是直和项。原点是 \(f\) 的固定点，故 \(e^{2\pi\mathrm{i}t}\mapsto[0,t]\) 给出截面。截面只提供一个基底方向，不会消除纤维单值造成的挠元。
\end{solution}

\begin{exercise}\label{g2:ex:rotation-mapping-torus}
	取 \(\theta\in\mathbb R\setminus\mathbb Z\)，令 \(f:S^1\to S^1\)、\(f(z)=e^{2\pi\mathrm{i}\theta}z\)。证明 \(f\) 没有固定点，但 \(E_f\to S^1\) 有截面。进一步构造一个覆盖基底恒等映射的同胚 \(E_f\cong S^1\times S^1\)，并写出它给出的截面在柱坐标中的道路。
\end{exercise}
\begin{solution}
	若 \(f(z)=z\)，因 \(z\ne0\) 得 \(e^{2\pi\mathrm{i}\theta}=1\)，与 \(\theta\notin\mathbb Z\) 矛盾。定义
	\[
		\Phi:E_f\longrightarrow S^1\times S^1,
		\qquad \Phi([z,t])=(e^{2\pi\mathrm{i}\theta t}z,e^{2\pi\mathrm{i}t}).
	\]
	在 \(t=1\) 处的值为 \((e^{2\pi\mathrm{i}\theta}z,1)\)，与 \([f(z),0]\) 的值一致，故此连续公式下降到商空间。给定 \((w,e^{2\pi\mathrm{i}t})\)，取 \(0\leq t<1\)，其原像为 \([e^{-2\pi\mathrm{i}\theta t}w,t]\)，因此 \(\Phi\) 双射。定义域是紧空间的商，靶为 Hausdorff 空间，所以连续双射为同胚；第二坐标正是原投影。

	乘积中的截面 \(b\mapsto(1,b)\) 对应于道路 \(c(t)=e^{-2\pi\mathrm{i}\theta t}\)。它满足 \(c(0)=1=f(c(1))\)，因此 \(t\mapsto[c(t),t]\) 下降为截面。这个例子把“单值同胚有固定点”和“丛有截面”区分开：截面可以沿纤维移动，不必由某个固定纤维点给出。
\end{solution}

\begin{exercise}\label{g2:ex:permutation-mapping-torus}
	设 \(F\) 是含 \(d\geq1\) 个点的离散空间，\(f\) 是其一个置换，循环分解含 \(r\) 个循环，长度为 \(\ell_1,\ldots,\ell_r\)。计算 \(E_f\) 的同调与上同调，给出其道路分支，并判断投影 \(E_f\to S^1\) 何时有截面。解释为何 \(H_0(F;\mathbb Z)\) 总有非零单值不变元，却不总有截面。
\end{exercise}
\begin{solution}
	每个点贡献一条区间，末端按置换接到下一点的区间首端。同一循环的 \(\ell_j\) 条区间依次粘成一个圆周，不同循环之间不相交。因此 \(E_f\) 是 \(r\) 个圆周的不交并，第 \(j\) 个分支到基底的投影绕行 \(\ell_j\) 周。直接得到
	\[
		H_0(E_f;\mathbb Z)=H_1(E_f;\mathbb Z)=\mathbb Z^r,
		\qquad H^0(E_f;\mathbb Z)=H^1(E_f;\mathbb Z)=\mathbb Z^r,
	\]
	其余非负次数的群为零。也可从正合列计算：\(H_0(F)=\mathbb Z^d\)，\(f_*\) 置换点基；\(\ker(f_*-1)\) 由各循环内全部点的和生成，\(\operatorname{coker}(f_*-1)\) 则将同一循环中的点基识别，二者均为 \(\mathbb Z^r\)。高次纤维同调为零，故没有其他贡献；上同调用对偶置换得到相同结果。

	由\cref{g2:prop:mapping-torus-section}，截面存在恰当某个纤维道路分支被 \(f\) 固定。离散空间的道路分支是单点，所以条件正是某个 \(\ell_j=1\)。循环中的点之和总是非零不变元，但当循环长度大于一时，它是多个点的形式和，不是一个能单独返回自身的纤维点。这正是同调不变元与几何截面的差别。
\end{solution}

% !TeX root = ../../Book2.tex
% 纲要标题：## 第6章　流形对偶、支集与相交配对
% 纲要位置：几何学计划纲要.md，第 1174 行。
% 正文状态：未撰写；本文件目前是框架。

% 纲要内容开始
% 
% > 本章目标与证明范围：以局部基本类和支集为证明主线，完整建立所用 Poincaré–Lefschetz 版本。闭、非紧、带边界和非定向不能仅在公式旁追加一句说明。
% 
% 纲要内容结束

\chapter{流形对偶、支集与相交配对}
\label{g2:ch:06}
\BookChapterNumber{g2:ch:06}

\begin{chapterabstract}
待撰写
\end{chapterabstract}

\begin{learninggoals}
待撰写
\end{learninggoals}

待撰写

% !TeX root = ../../Book2.tex
% 纲要标题：### 6.1 局部基本类的粘合
% 纲要位置：几何学计划纲要.md，第 1228 行。
% 正文状态：初稿；尚未正式审核。

% 纲要内容开始
% 
% * 欧氏局部同调生成元
% * 定向系统与紧集上的基本类
% * 从局部一致性到整体基本类
% 
% 纲要内容结束

\section{局部基本类的粘合}
\label{g2:sec:0601}

在一个小球内，最高次局部同调记录这一小块空间的定向。对偶定理需要把这些局部数据同时放到流形上：它们能否来自同一个同调类？闭流形与非紧流形在这里已有差别。有限奇异链能够覆盖整个紧流形，却不能覆盖整个非紧流形；因此先在任意紧集附近构造基本类，比一开始只讨论一个整体类更合适。

\ProseParagraphBreak

本节的流形均为 Hausdorff、第二可数的无边界拓扑 \(n\) 维流形，\(R\) 为含幺交换环。第五章已经构造整数定向局部系统 \(\mathcal O_M\)，并证明局部系数的切除与同伦不变性。在每个坐标球内，这些工具给出局部生成元；紧集的有限覆盖和相对 Mayer--Vietoris 序列将把它们粘成一个紧集附近的相对类。

\subsection{欧氏局部同调生成元}
\label{g2:subsec:0601:01}
回忆\cref{g2:thm:manifold-local-homology}：对 \(x\in M\)，整数局部同调在次数 \(n\) 为 \(\mathbb Z\)，其他次数为零。这个无限循环群并没有首选生成元；一张有序坐标图给出的两种方向，对应两个互为负数的生成元。按\cref{g2:def:orientation-local-system}，这些群连同坐标球内的运输组成 \(\mathcal O_M\)。现把系数推广为 \(R\)，记
\[
	\mathcal O_{R}\coloneqq\mathcal O_M\otimes_{\mathbb Z}R.
\]
\symidx{\(\mathcal O_R\)：以 \(R\) 为系数的定向局部系统}
若 \(\mathcal L\) 是 \(M\) 上任意 \(R\) 模局部系统，则有自然识别
\begin{equation}\label{g2:eq:local-homology-with-coefficients}
	H_i(M,M\setminus\{x\};\mathcal L)\cong
	\begin{cases}
		\mathcal O_{R,x}\otimes_R\mathcal L_x,&i=n,\\
		0,&i\ne n.
	\end{cases}
\end{equation}
公式中的系数纤维 \(\mathcal L_x\) 来自局部平凡化：先切除到 \(x\) 的坐标球，在球内把各纤维运输到 \(x\)，再计算常系数局部同调。球面计算中的悬挂连接同构对系数模成立，故最高次得到 \(\mathcal O_{M,x}\otimes_{\mathbb Z}\mathcal L_x\)；这与式中张量积规范相同。也可在球面的自由胞腔链上直接核对：顶次只有一个系数模，其余约化链为零。球内换一组定向，局部同调生成元和 \(\mathcal O_{R,x}\) 的基同时反号，所以识别不依赖所选坐标方向。

\begin{proposition}[定向系统的平方]\label{g2:prop:orientation-square}
	设 \(M\) 是 Hausdorff、第二可数、无边界的拓扑 \(n\) 维流形，\(R\) 为含幺交换环。存在规范的局部系统同构
	\begin{equation}\label{g2:eq:orientation-square}
		b:\mathcal O_R\otimes_R\mathcal O_R\xrightarrow{\ \cong\ }\underline R,
		\qquad b((\mu\otimes a)\otimes(\mu\otimes c))=ac,
	\end{equation}
	其中 \(\mu\) 是整数局部同调的任一生成元，\(\underline R\) 表示常值 \(R\) 系统。因此每点有一个规范元素
	\[
		\mathbf 1_x=(\mu\otimes1)\otimes(\mu\otimes1)
		\in H_n(M,M\setminus\{x\};\mathcal O_R),
	\]
	它对应于 \(1\in R\)，并随 \(x\) 平行变化。
\end{proposition}
\begin{proof}
	整数生成元仅有 \(\mu\) 和 \(-\mu\) 两个。同时更换两因子的基，积中的符号为 \((-1)^2=1\)，故 \(b\) 无歧义。在一个坐标球内固定整数生成元族，这个映射就是 \(R\otimes_RR\to R\) 的乘法同构。路径运输也同时作用于两因子，闭路造成的两个定向符号彼此抵消，所以局部同构相容地给出全局同构。其逆把 \(1\) 送到陈述中的元素；再用\eqref{g2:eq:local-homology-with-coefficients}，便得到所述局部类。
\end{proof}

这个平方解释了为什么不可定向流形仍能有定向系数基本类：几何上的反号与系数上的反号同时发生，局部类本身不变。配对的规范归一化取整数生成元的像；换成任意 \(R\) 基 \(\theta=u\mu\) 时，\(b(\theta,\theta)=u^2\)，归一化值随之变化。

\ProseParagraphBreak

下面还要用到一点之外的 Euclid 模型。设 \(C\subset\mathbb R^n\) 是非空凸紧集，\(x\in C\)。选以 \(x\) 为中心、包含 \(C\) 于内部的大球，其边界为 \(S\)。沿从 \(x\) 发出的射线，\(\mathbb R^n\setminus C\) 和 \(\mathbb R^n\setminus\{x\}\) 都形变收缩到 \(S\)：一条射线与凸集的交是从起点出发的一个闭区间，因此一个已在 \(C\) 外的点移动到 \(S\) 时不会再次进入 \(C\)。由对的长正合列，映射
\begin{equation}\label{g2:eq:convex-local-class}
	H_i(\mathbb R^n,\mathbb R^n\setminus C;A)
	\longrightarrow H_i(\mathbb R^n,\mathbb R^n\setminus\{x\};A)
\end{equation}
对任意 \(R\) 模 \(A\) 都是同构。这里允许 \(x\) 位于 \(C\) 的边界，也允许 \(C\) 没有内部；只需 \(C\) 非空且凸。若 \(n=0\)，非空凸集只有一个点，同一结论由零次链直接得到。

\subsection{定向系统与紧集上的基本类}
\label{g2:subsec:0601:02}
在紧集 \(K\) 上指定局部值，需要先说清“相容”的含义。令 \(\mathcal E\) 为一个局部系统。记 \(\Gamma(K;\mathcal E|_K)\) 为其纤维空间在 \(K\) 上的连续截面：在任一局部平凡化中，截面的系数是 \(K\) 上的局部常值函数。即使 \(K\) 不是局部道路连通，也使用这个定义，不能只要求沿 \(K\) 内的道路保持不变。例如在一列点及其极限点组成的紧集上，每条道路都常值，但连续截面仍须在足够靠近极限点时取同一个系数。

\begin{lemma}[紧集上的局部值]\label{g2:lem:compact-localization}
	设 \(M\) 是 Hausdorff、第二可数、无边界的拓扑 \(n\) 维流形，\(R\) 为含幺交换环，\(\mathcal L\) 为 \(M\) 上的 \(R\) 模局部系统，\(K\subset M\) 紧。则
	\[
		H_i(M,M\setminus K;\mathcal L)=0\qquad(i>n),
	\]
	且限制到各点给出同构
	\begin{equation}\label{g2:eq:compact-localization}
		\rho_K:H_n(M,M\setminus K;\mathcal L)
		\xrightarrow{\ \cong\ }\Gamma(K;(\mathcal O_R\otimes_R\mathcal L)|_K).
	\end{equation}
	特别地，一个顶次相对类若在 \(K\) 的每一点都有零局部值，则它为零。对紧集包含 \(K\subset L\)，相对化映射 \(H_n(M,M\setminus L;\mathcal L)\to H_n(M,M\setminus K;\mathcal L)\) 对应截面限制；对系数同态以及开邻域的切除，以上识别同样自然。
\end{lemma}
\begin{proof}
	我们先说明相对类为何确实给出连续截面，再证明存在性和唯一性。令相对圈 \(z\) 表示一个顶次类，则有限链 \(\partial z\) 的像构成一个与 \(K\) 不交的紧集。每个 \(x\in K\) 都有一个闭坐标球 \(B\)，使 \(x\in\operatorname{int}B\) 且 \(B\) 避开这个紧集。于是 \(z\) 还定义 \(H_n(M,M\setminus B;\mathcal L)\) 中的类。由\eqref{g2:eq:convex-local-class}及局部系数切除，它在 \(B\) 内的各局部值属于同一个平行族。因此原来的局部值在 \(K\) 上连续，\(\rho_K\) 有定义。所有限制均来自空间对或系数映射，故自然性一旦有同构便随之成立。

	粘合所需的正合性可以直接写出。暂记 \(G_i(K)=H_i(M,M\setminus K;\mathcal L)\)。对紧集 \(A,B\subset M\)，局部系数的相对 Mayer--Vietoris 序列包含
	\begin{equation}\label{g2:eq:compact-relative-mv}
		\begin{aligned}
			\cdots\longrightarrow G_{i+1}(A\cap B)&\longrightarrow G_i(A\cup B)\\
			&\xrightarrow{(r_A,r_B)}G_i(A)\oplus G_i(B)\\
			&\xrightarrow{\mathrm{res}_A-\mathrm{res}_B}G_i(A\cap B)\longrightarrow\cdots.
		\end{aligned}
	\end{equation}
	为核对它的来源，令 \(C=C_*(M;\mathcal L)\)、\(P=C_*(M\setminus A;\mathcal L)\)、\(Q=C_*(M\setminus B;\mathcal L)\)。链短正合列为
	\[
		0\longrightarrow C/(P\cap Q)\xrightarrow{([z],[z])}C/P\oplus C/Q
		\xrightarrow{[u]-[v]}C/(P+Q)\longrightarrow0.
	\]
	这里 \(P\cap Q=C_*(M\setminus(A\cup B);\mathcal L)\)。另一方面，\(M\setminus A\) 与 \(M\setminus B\) 是 \(M\setminus(A\cap B)\) 的开覆盖；\cref{g2:lem:twisted-small-chain-retraction}说明 \(P+Q\) 包含到这个补空间的全部链是链同伦等价。比较以 \(C\) 为中间项的两个商短正合列，其子复形映射与中间映射都在同调上为同构，长正合列的图追踪便给出末项同调与 \(H_i(M,M\setminus(A\cap B);\mathcal L)\) 的同构。因而取长正合列就得到\eqref{g2:eq:compact-relative-mv}。箭头上的两个限制之差作用为 \((a,b)\mapsto a|_{A\cap B}-b|_{A\cap B}\)。

	若结论对 \(A,B,A\cap B\) 已成立，就对 \(A\cup B\) 成立。高于 \(n\) 次的群在上述正合列中夹在两个零群之间；次数 \(n\) 时，因为交集的 \(n+1\) 次群为零，\((r_A,r_B)\) 单射，其像正是交上限制相同的两个类。相容截面也恰由 \(A,B\) 上限制相同的截面组成：这是有限闭覆盖上的连续映射粘合，\(A,B\) 在其并集中闭。故这一段正合列同时给出顶次局部值的存在与唯一性。

	先在一张坐标图 \(U\cong\mathbb R^n\) 内证明结论；\(\mathcal L|_U\) 平凡化为一个固定模 \(A_0\)，定向系统也平凡。若紧集凸，\eqref{g2:eq:convex-local-class}给出全部结论，因为凸集连通，其局部常值截面都是常值。若紧集是有限个凸紧集的并，则利用上一段归纳。归纳中，一个新凸集与此前各凸集的交仍为凸紧集或空集；所以交集是更少个同类集合的并，归纳假设确实适用。空集的相对群与截面模均为零。

	现在令 \(K\subset U\) 是任意紧集。对 \(\alpha\in H_i(M,M\setminus K;\mathcal L)\)，先用\cref{g2:thm:twisted-excision}将它表示为 \(U\) 中的有限相对圈 \(z\)。链 \(\partial z\) 的像与 \(K\) 是两个不交的紧集，在坐标中有正距离；若 \(\partial z=0\)，则任选一个正距离界。以 \(K\) 中的点为中心，取有限个足够小的闭球，使其内部覆盖 \(K\)，且其并 \(N\) 避开 \(\partial z\) 的像。于是 \(z\) 已在 \(H_i(M,M\setminus N;\mathcal L)\) 定义一个类 \(\alpha_N\)，并映到 \(\alpha\)。当 \(i>n\) 时，有限凸集并的情形给出 \(\alpha_N=0\)。当 \(i=n\) 且 \(\alpha\) 的全部局部值为零时，\(\alpha_N\) 限制到每个小球的中心为零；凸集情形说明它在这个球上的相对类为零。因此 \(\alpha_N\) 在 \(N\) 上的所有局部值为零，有限并的唯一性又给出 \(\alpha_N=0\)。这证明了一般 \(K\subset U\) 的高次消失和顶次单射。

	还须实现 \(K\) 上给定的截面。在上述平凡化中，它是连续映射 \(s:K\to A_0\)，其中 \(A_0\) 取离散拓扑。紧性说明它只取有限多个值 \(a_1,\ldots,a_t\)；每个非空纤维 \(K_j=s^{-1}(a_j)\) 在 \(K\) 中既开又闭，因此是紧集。不同 \(K_j\) 的坐标距离为正。可分别选由有限个小闭球组成的紧邻域 \(N_j\supset K_j\)，使这些 \(N_j\) 两两不交。在每个 \(N_j\) 上，常值 \(a_j\) 由有限凸集并的情形给出唯一顶次相对类。对不交并使用已证的粘合，得到 \(\bigcup_jN_j\) 上的类，再限制到 \(K\)，其局部值恰为 \(s\)。这就完成了单张坐标图内的证明。

	最后处理一般 \(K\subset M\)。每个点有闭包紧且包含在坐标图内的小坐标球，\(\mathcal L\) 在该图内平凡。用有限个这样的球内部覆盖 \(K\)，令 \(K_1,\ldots,K_m\) 为 \(K\) 与各闭球的交；它们紧、覆盖 \(K\)，且各在一张图内。对覆盖项数 \(m\) 归纳。令 \(A=K_1\cup\cdots\cup K_{m-1}\)、\(B=K_m\)，则
	\[
		A\cap B=(K_1\cap K_m)\cup\cdots\cup(K_{m-1}\cap K_m)
	\]
	仍由至多 \(m-1\) 个各位于坐标图内的紧集组成。归纳假设适用于 \(A,B,A\cap B\)，粘合步骤遂得到 \(K\) 的结论。若 \(n=0\)，流形离散，紧集有限，结论也可逐点直接计算。
\end{proof}

引理的关键是顶次的单射性：只要局部值相同，来自不同球、不同有限覆盖或不同代表链的构造自动给出同一个相对类。高一次群的消失使相对 Mayer--Vietoris 序列提供唯一粘合。常系数的上述论证可对照\cite[第3.3节，Orientations and Homology]{Hatcher2002}；这里把同一局部到整体过程写成了局部系数与紧集上连续截面的形式。

\begin{definition}[紧集上的基本类]\label{g2:def:compact-fundamental-class}
	设 \(M\) 是 Hausdorff、第二可数、无边界的拓扑 \(n\) 维流形，\(R\) 为含幺交换环，\(K\subset M\) 紧。在\eqref{g2:eq:compact-localization}中取 \(\mathcal L=\mathcal O_R\)，并用\eqref{g2:eq:orientation-square}，常值截面 \(1\) 唯一对应的类
	\[
		[M]_K\in H_n(M,M\setminus K;\mathcal O_R)
	\]
	称为 \(M\) 在 \(K\) 上的\term[jinjishangdejibenlei]{基本类}[fundamental class over a compact set]。其在每个 \(x\in K\) 的局部值为 \(\mathbf1_x\)。
\end{definition}

一个有向区间就能显示这种相对类如何工作。取 \(M=\mathbb R\)、\(K=[a,b]\)，用标准方向识别 \(\mathcal O_{\mathbb Z}\) 与 \(\mathbb Z\)。从 \(a-1\) 到 \(b+1\) 的有向线段是相对圈，因为两个端点都在 \(M\setminus K\) 中；它在每个 \(x\in K\) 的局部值为 \(1\)，所以代表 \([\mathbb R]_K\)。改换更远的端点不改此类。若 \(K\) 有两个不交的紧区间，顶次相对群则为 \(\mathbb Z^2\)，允许分别指定两个局部整数；基本类选的是 \((1,1)\)。

\subsection{从局部一致性到整体基本类}
\label{g2:subsec:0601:03}
若紧集 \(K\subset L\)，从 \(L\) 附近的相对类限制到 \(K\) 附近，仍保留 \(K\) 上规定的局部值。局部值引理的唯一性于是使各个紧集上的基本类自动相容；紧流形的整体基本类就是取 \(K=M\) 所得的类。

\begin{theorem}[基本类及其限制]\label{g2:thm:fundamental-class}
	设 \(M\) 是 Hausdorff、第二可数、无边界的拓扑 \(n\) 维流形，\(R\) 为含幺交换环。对所有紧集 \(K\subset M\)，\([M]_K\) 存在且唯一；若 \(K\subset L\) 均紧，则
	\begin{equation}\label{g2:eq:fundamental-class-restriction}
		[M]_L\longmapsto[M]_K.
	\end{equation}
	若 \(U\subset M\) 开且包含 \(K\)，在 \(\mathcal O_{R,U}\cong\mathcal O_{R,M}|_U\) 的自然识别下，切除同构将 \([U]_K\) 送到 \([M]_K\)。若 \(M\) 紧，则
	\[
		[M]\coloneqq[M]_M\in H_n(M;\mathcal O_R)
	\]
	是规范的整体基本类，且 \(H_n(M;\mathcal O_R)\cong\Gamma(M;\underline R)\)。特别地，\(M\) 连通且紧时，\(H_n(M;\mathcal O_R)\cong R\)，\([M]\) 对应 \(1\)。
\end{theorem}
\begin{proof}
	存在唯一性正是\cref{g2:lem:compact-localization}与规范平方同构的结论。紧集限制和开集切除都将每个 \(x\in K\) 处的规范局部值送到自身，由唯一性便有陈述中的相容性。紧时取 \(K=M\)，相对群成为绝对群。连通空间上的局部常值函数为常值，故最后的截面模为 \(R\)。
\end{proof}

常系数基本类需要一个定向选择。一个 \(R\)\term[huandingxiang]{定向}[orientation over a ring]是 \(\mathcal O_R\) 的平行截面 \(\theta\)，其在每点都生成该秩一自由 \(R\) 模；等价地，在整数局部生成元下，其系数是单位。\cref{g2:lem:compact-localization}在 \(\mathcal L=\underline R\) 时给出唯一类
\[
	[M]^{\theta}_K\in H_n(M,M\setminus K;R),\qquad
	\rho_K([M]^{\theta}_K)=\theta|_K.
\]
若要从规范扭系数类直接得到它，应使用系数同构 \(b(\theta,-):\mathcal O_R\to\underline R\)：在局部写 \(\theta=u\mu\)，它把 \(\mu\) 送到 \(u\)，所以 \(\mu\otimes\mu\) 的像恰为 \(u\mu=\theta\)。当 \(\theta\) 来自整数定向时，这就是熟悉的按所选正方向将系数识别为 \(R\)。整数定向反转，常系数基本类变号；规范扭系数基本类则不变。

\begin{example}[闭流形的顶次群]\label{g2:exm:closed-manifold-top-homology}
	设 \(M\) 是连通闭 \(n\) 维拓扑流形，\(R\) 为含幺交换环。若 \(M\) 可定向，则选定整数定向后 \(H_n(M;R)\cong R\)。若不可定向，则
	\[
		H_n(M;R)\cong\{a\in R:2a=0\}.
	\]
	事实上，\cref{g2:lem:compact-localization}把常系数顶次类识别为 \(\mathcal O_R\) 的平行截面。固定一点的整数局部生成元后，截面由该点的一个系数 \(a\) 决定；定向保持环路不加条件，定向反转环路要求 \(a=-a\)。因此整数系数下不可定向闭流形的顶次群为零，而模二系数下仍有基本类。取 \(R=\mathbb Z/4\) 时，不可定向闭流形的顶次群是 \(\{0,2\}\cong\mathbb Z/2\)，但非零元素 \(2\) 不是单位，所以它不构成 \(R\) 定向的基本类。这区别了“顶次群非零”与“存在处处为生成元的类”。
\end{example}

在非紧流形上，各紧集附近的基本类仍然存在，并随紧集的包含关系相容。对于连通非紧流形，普通有限链的顶次同调为零，定向系数下也如此；下面的证明将这一消失归结为局部值在连通空间上的一致性。

\begin{proposition}[非紧流形的顶次消失]\label{g2:prop:noncompact-top-homology}
	设 \(M\) 是连通、非紧、Hausdorff、第二可数、无边界的拓扑 \(n\) 维流形，\(R\) 为含幺交换环，\(\mathcal L\) 为任意 \(R\) 模局部系统。则普通奇异同调满足
	\[
		H_i(M;\mathcal L)=0\qquad(i\geq n).
	\]
\end{proposition}
\begin{proof}
	取有限圈 \(z\)，其单形像的并是紧集。选开集 \(U\) 包含此并，使 \(\overline U\) 紧；具体地，用有限个闭包紧的小坐标球覆盖该紧集，再取这些球内部的并。令 \(V=M\setminus\overline U\)。两个开集 \(U,V\) 不交，且
	\[
		M\setminus V=\overline U,\qquad
		M\setminus(U\cup V)=\partial U
	\]
	均紧。三空间 \(V\subset U\cup V\subset M\) 的相对链短正合列给出
	\[
		H_{i+1}(M,U\cup V;\mathcal L)\longrightarrow H_i(U\cup V,V;\mathcal L)
		\longrightarrow H_i(M,V;\mathcal L).
	\]
	由于不交并，中央项就是 \(H_i(U;\mathcal L)\)。当 \(i>n\) 时，两端由\cref{g2:lem:compact-localization}的高次消失同时为零，故 \(z\) 已在 \(U\) 内为边界。

	当 \(i=n\) 时，左端仍为零，所以 \(H_n(U;\mathcal L)\to H_n(M,V;\mathcal L)\) 单射。类 \([z]\in H_n(M;\mathcal L)\) 在各点的局部值构成 \(\mathcal O_R\otimes\mathcal L\) 的全局平行截面；这一点仍由避开 \(\partial z=0\) 的小球论证得到。因 \(M\) 非紧，有限链的紧像之外存在一点，那里局部值为零。连通流形道路连通，局部系统运输是同构，所以这一截面处处为零。于是 \(z\) 在 \(H_n(M,V;\mathcal L)=H_n(M,M\setminus\overline U;\mathcal L)\) 中的类，由紧集 \(\overline U\) 上局部值的单射性而为零。中央项的单射又说明 \(z\) 在 \(U\) 内为边界，从而在 \(M\) 内为边界。
\end{proof}

证明先把有限链的支撑放进一个紧集，再通过绝对群到相对群的单射使用局部值的唯一性，从而由各点局部值为零得到原类为零。\(\mathbb R^n\) 展示了这一差别：每个紧集都有相容的 \([\mathbb R^n]_K\)，普通 \(H_n(\mathbb R^n)\) 却为零（\(n\geq1\)）。下一节引入紧支上链和局部有限链，使这些紧集附近的类能够参与非紧流形的对偶。
% 纲要条目排为编号小节标题，并进入目录与书签；说明段仍保留于纲要注释中。

% !TeX root = ../../Book2.tex
% 纲要标题：### 6.2 帽积、支集与自然性
% 纲要位置：几何学计划纲要.md，第 1184 行。
% 正文状态：未撰写；本文件目前是框架。

% 纲要内容开始
% 
% * 帽积的链级定义
% * 边界公式及符号
% * 紧支上同调和局部有限链
% * 普通、紧支和局部有限理论的配对
% 
% 纲要内容结束

\section{帽积、支集与自然性}
\label{g2:sec:0602}

待撰写

\subsection{帽积的链级定义}
\label{g2:subsec:0602:01}
待撰写

\subsection{边界公式及符号}
\label{g2:subsec:0602:02}
待撰写

\subsection{紧支上同调和局部有限链}
\label{g2:subsec:0602:03}
待撰写

\subsection{普通、紧支和局部有限理论的配对}
\label{g2:subsec:0602:04}
待撰写
% 纲要条目排为编号小节标题，并进入目录与书签；说明段仍保留于纲要注释中。

% !TeX root = ../../Book2.tex
% 纲要标题：### 6.3 Poincaré 对偶的证明
% 纲要位置：几何学计划纲要.md，第 1241 行。
% 正文状态：初稿；尚未正式审核。

% 纲要内容开始
% 
% * 欧氏空间上的对偶
% * Mayer–Vietoris 与五引理
% * 可数穷竭及极限步骤
% * 闭流形的 Poincaré 对偶
% 
% 纲要内容结束

\section{Poincaré 对偶的证明}
\label{g2:sec:0603}

帽积已经把一个上同调类变成维数互补的链，基本类则给出了进行这项运算的几何背景。现在要证明，这个过程既不丢失信息，也能得到每一个互补维数的同调类。证明从 Euclid 空间开始，通过开覆盖的正合列粘合，最后处理穷尽中的极限。闭流形只是其中紧支条件自动消失的一种情形。

\ProseParagraphBreak

本节均取 Hausdorff、第二可数、无边界的拓扑 \(n\) 维流形，\(R\) 为含幺交换环，\(\mathcal L\) 为 \(R\) 模局部系统。写 \(\mathcal O_R\) 表示定向系统。先证明
\[
	H_c^k(M;\mathcal L)\cong H_{n-k}(M;\mathcal O_R\otimes_R\mathcal L),
\]
再证明普通上同调与局部有限同调的对应。后一版本将有限链改为局部有限链，其极限比较同时处理逆极限及其修正项。

\subsection{欧氏空间上的对偶}
\label{g2:subsec:0603:01}
对紧集 \(K\subset M\)，\cref{g2:thm:fundamental-class}给出规范类 \([M]_K\in H_n(M,M\setminus K;\mathcal O_R)\)。相对帽积定义
\[
	D_K:H^k(M,M\setminus K;\mathcal L)\longrightarrow H_{n-k}(M;\mathcal O_R\otimes_R\mathcal L),
	\qquad \alpha\longmapsto\alpha\frown[M]_K.
\]
为何输出是绝对同调？若上圈 \(\varphi\) 在 \(M\setminus K\) 上为零，相对基本圈 \(c\) 的边界在那里，则边界公式中的 \(\varphi\frown\partial c\) 和 \(\delta\varphi\frown c\) 都为零。改变相对基本圈或上圈代表时，同一公式说明结果只改变绝对边界。

\ProseParagraphBreak

对 \(K\subset L\)，基本类的相对化把 \([M]_L\) 送到 \([M]_K\)，故 \(D_K\) 与上同调的支集扩大相容。取直接极限得到\term[jinzhi duiyou yingshe]{紧支对偶映射}[compactly supported duality map]
\begin{equation}\label{g2:eq:compact-duality-map}
	D_{c,M}:H_c^k(M;\mathcal L)\longrightarrow H_{n-k}(M;\mathcal O_R\otimes_R\mathcal L).
\end{equation}
若 \(j:U\hookrightarrow M\) 开，则 \(D_{c,M}j_!=j_*D_{c,U}\)：把紧支上圈及对应基本圈切除到 \(U\) 后，两边使用相同的帽积。这里的延拓 \(j_!\) 由\cref{g2:prop:compact-support-MV}中的切除同构给出，因而与基本类的限制相容。

\begin{lemma}[Euclid 模型的对偶]\label{g2:lem:euclidean-duality}
	设 \(R\) 为含幺交换环，\(A\) 为任意 \(R\) 模，\(U\) 为同胚于 \(\mathbb R^n\) 的空间。对常值系统 \(\underline A\)，映射 \(D_{c,U}\) 在各次均为同构。选定一组整数定向后，其唯一可能非零的部分是
	\[
		H_c^n(U;A)\cong A\xrightarrow{\ D_{c,U}\ }H_0(U;A)\cong A,
	\]
	此映射为恒等。
\end{lemma}
\begin{proof}
	先取 \(U=\mathbb R^n\)、\(n>0\)。用越来越大的闭立方体 \(K_j\) 穷尽。补空间形变收缩到 \(S^{n-1}\)，对的长正合列给出 \(H^k(U,U\setminus K_j;A)=0\)（\(k\ne n\)），以及 \(H^n(U,U\setminus K_j;A)=A\)。这些识别由定向及连接映射确定，扩大立方体的映射保持它们，故直接极限仍为 \(A\)。普通同调因 \(U\) 可缩而只在零次为 \(A\)。

	还要核对对偶映射在这两个 \(A\) 之间没有引入额外倍数。取一个包含 \(K_j\) 于内部的有向 \(n\) 单形 \(\Delta\)。径向移动把相应空间对与 \((\Delta,\partial\Delta)\) 的生成元对应起来。相对上圈 \(\varphi\) 若在 \(\Delta\) 上取值 \(a\)，前后面公式便给出 \(\varphi\frown\Delta=a[v_n]\)，其零次同调类就是 \(a\)。因此所求映射为恒等。换同胚或反转所选整数定向，会同时改变基本链和定向系统的基，规范的扭系数映射保持不变。\(n=0\) 时空间只有一点，帽积直接是 \(A\otimes_RR\cong A\)。
\end{proof}

任意局部系统在一个坐标球内都平凡：球单连通，\cref{g2:prop:parallel-sections}保证从一点运输到其他点没有道路歧义。因此上述球内计算适用于任意 \(\mathcal L\)，球与球之间的系数比较仍由原局部系统的运输给出。

\subsection{Mayer–Vietoris 与五引理}
\label{g2:subsec:0603:02}
设 \(M=U\cup V\)。紧支上同调的延拓方向与普通同调的包含方向相同，所以两条 Mayer--Vietoris 序列可以按互补次数比较。要比较连接映射，须在帽积边界公式中保留相对基本圈位于补空间的边界，再利用上链的支集条件处理这些项。

\begin{lemma}[对偶与开集粘合]\label{g2:lem:duality-MV}
	设 \(M\) 为 Hausdorff、第二可数、无边界的拓扑 \(n\) 维流形，\(R\) 为含幺交换环，\(\mathcal L\) 为 \(R\) 模局部系统，\(M=U\cup V\) 为开覆盖。紧支上同调序列与系数为 \(\mathcal O_R\otimes_R\mathcal L\) 的普通同调序列，在对偶映射下除连接方块外均交换；连接方块满足
	\begin{equation}\label{g2:eq:duality-MV-sign}
		\partial_{\mathrm{MV}}D_{c,M}=(-1)^{k+1}D_{c,U\cap V}\delta_{\mathrm{MV}}
		\quad\text{在 }H_c^k(M;\mathcal L)\text{ 上}.
	\end{equation}
	若 \(D_{c,U},D_{c,V},D_{c,U\cap V}\) 在各次为同构，则 \(D_{c,M}\) 在各次为同构。
\end{lemma}
\begin{proof}
	普通箭头由帽积自然性与开集延拓的定义交换。为计算连接方块，取紧集 \(K\subset U\)、\(L\subset V\)，记 \(A=M\setminus K\)、\(B=M\setminus L\)。代表上圈 \(\varphi\) 在 \(A\cap B\) 上为零。由\cref{g2:prop:compact-support-MV}证明中的上链短正合列，可写
	\[
		\varphi=\varphi_A-\varphi_B,\qquad
		\varphi_A|_A=0,\quad\varphi_B|_B=0,\quad
		\delta\varphi_A=\delta\varphi_B.
	\]
	连接像由 \(\delta\varphi_A\) 表示，它在 \(A\) 内及 \(B\) 内的单形上都为零；通过小链比较，这正是支集为 \(K\cap L\) 的类。

	选代表 \([M]_{K\cup L}\) 的有限链 \(c\)，按开覆盖 \(U\setminus L,U\cap V,V\setminus K\) 作足够细分，写成 \(c=a+b+d\)，三个部分依次落在这三个开集中。该覆盖确实覆盖 \(M\)：在 \(U\) 外的点在 \(V\setminus K\)，在 \(V\) 外的点在 \(U\setminus L\)。\(b\) 表示交集上的基本类，\(a+b\) 表示 \(K\) 上的基本类，故
	\[
		\partial(a+b)\in C_{n-1}(U\setminus K;\mathcal O_R),\qquad
		\varphi_B\frown\partial a=0.
	\]
	第一个包含也可逐单形看出：\(\partial c\) 避开 \(K\cup L\)，而 \(d\) 及其面避开 \(K\)，所以 \(\partial(a+b)\) 避开 \(K\)；它仍是 \(U\) 中的链。

	沿上同调连接箭头再作对偶，得到 \((\delta\varphi_A)\frown b\)。由帽积边界公式，它在 \(U\cap V\) 中与 \(\varphi_A\frown\partial b\) 同调。另一方向先得到绝对圈 \(\varphi\frown c\)，把它分成 \(U\) 中的 \(\varphi\frown a\) 和 \(V\) 中的 \(\varphi\frown(b+d)\)。普通 Mayer--Vietoris 连接映射取第一部分的边界，而
	\[
		\begin{aligned}
			\partial(\varphi\frown a)
			&=(-1)^k\varphi\frown\partial a\\
			&=(-1)^k\varphi_A\frown\partial a
			=(-1)^{k+1}\varphi_A\frown\partial b.
		\end{aligned}
	\]
	最后一步使用 \(\varphi_A\frown\partial(a+b)=0\)。这是所需符号。局部系数在每个单形中平凡化，所有公式保持运输；细分与相对小链比较均有\cref{g2:lem:twisted-small-chain-retraction}的链同伦，所以这些代表计算确实给出原序列的映射。

	最后取任意五个连续位置的正合列，将必要竖箭头乘以 \(-1\)，使方块交换；改变符号不改变是否同构。中间箭头两侧的四个箭头由假设为同构，五引理给出中间箭头同构。其图追踪含义是：满射性先用右边箭头找候选原像，再用左边箭头修正差值；单射性先将一个核元移到左侧，再用更左一项消去它。故 \(D_{c,M}\) 在所有次数均为同构。
\end{proof}

这个引理将已知开集及其交集上的对偶推广到它们的并。坐标图的交通常只是 Euclid 开集，因此下一步先从球和有限并出发证明任意 Euclid 开集上的对偶，再用于流形的坐标覆盖。

\subsection{可数穷竭及极限步骤}
\label{g2:subsec:0603:03}
先处理直接极限。若 \(W_1\subset W_2\subset\cdots\) 为开集且 \(M=\bigcup_jW_j\)，则
\begin{equation}\label{g2:eq:duality-direct-continuity}
	\varinjlim_jH_c^k(W_j;\mathcal L)\cong H_c^k(M;\mathcal L),\qquad
	\varinjlim_jH_q(W_j;\mathcal E)\cong H_q(M;\mathcal E).
\end{equation}
第一个式子由紧支集给出：每个紧集包含于某个 \(W_j\)，这是紧集对递增开覆盖取有限子覆盖的结果；两个表示相等，也已在某个较大的紧支集处相等。第二个式子由有限链给出：一条圈及证明它为边界的链都只含有限个单形，其像为紧集，所以都已包含在某个 \(W_j\) 内。这些性质分别给出两式的满射性与单射性。

\begin{theorem}[紧支 Poincaré 对偶]\label{g2:thm:poincare-duality-compact}
	设 \(M\) 是 Hausdorff、第二可数、无边界的拓扑 \(n\) 维流形，\(R\) 为含幺交换环，\(\mathcal L\) 为任意 \(R\) 模局部系统。帽积给出自然同构
	\begin{equation}\label{g2:eq:poincare-compact}
		D_{c,M}:H_c^k(M;\mathcal L)\xrightarrow{\ \cong\ }H_{n-k}(M;\mathcal O_R\otimes_R\mathcal L)
	\end{equation}
	对所有整数 \(k\) 成立；负次数群取零。这里的自然性包括系数同态以及开嵌入的紧支延拓与普通同调包含。
\end{theorem}
\begin{proof}
	在 \(\mathbb R^n\) 的一个开集 \(W\) 中，取可数个闭包包含于 \(W\) 的开矩形覆盖 \(W\)，例如端点有理的坐标矩形。每个非空有限交仍是开矩形，因而同胚于 \(\mathbb R^n\)。\cref{g2:lem:euclidean-duality}给出一个矩形的结论；对有限个矩形的并归纳，新增矩形与此前并集的交是更少个矩形的并，故\cref{g2:lem:duality-MV}适用。再对前 \(j\) 个矩形的并使用\eqref{g2:eq:duality-direct-continuity}，得到 \(W\) 的结论。局部系统只需在每个矩形内平凡，不要求在整个 \(W\) 上平凡。

	一般流形因第二可数而有可数个坐标球 \(U_1,U_2,\ldots\) 覆盖。令 \(V_j=U_1\cup\cdots\cup U_j\)。\(U_j\cap V_{j-1}\) 是 \(U_j\) 中的开集，在坐标下是 Euclid 空间的开集，上一段已证明其对偶。用开集粘合引理归纳得到所有 \(V_j\) 的结论，再用直接极限得到 \(M\) 的结论。各步使用的映射始终是\eqref{g2:eq:compact-duality-map}的帽积映射，所以得到的同构具有已核对的自然性。对于 \(n=0\)，坐标球是单点，相同的有限并与可数并论证仍成立。
\end{proof}

上述局部到整体证明可对照\cite[第3.3节，Duality for Noncompact Manifolds；第3.H节]{Hatcher2002}。现在转向普通上同调。要保留一个非紧流形的全部局部定向，基本链必须允许局部有限的无限和。

\begin{proposition}[局部有限基本类]\label{g2:prop:lf-fundamental-class}
	设 \(M\) 是 Hausdorff、第二可数、无边界的拓扑 \(n\) 维流形，\(R\) 为含幺交换环。存在唯一
	\[
		[M]_{\mathrm{lf}}\in H_n^{\mathrm{lf}}(M;\mathcal O_R)
	\]
	使其在每个紧集 \(K\) 上的相对化均为 \([M]_K\)。其各点局部值均为 \(\mathbf1_x\)。当 \(M\) 紧时，它就是普通基本类。
\end{proposition}
\begin{proof}
	取紧穷尽 \(K_j\subset\operatorname{int}K_{j+1}\)。\cref{g2:lem:compact-localization}给出 \(H_{n+1}(M,M\setminus K_j;\mathcal O_R)=0\)。故\eqref{g2:eq:lf-limit-sequence}中的修正项为零，而相容的族 \(([M]_{K_j})_j\) 唯一提升为所求局部有限同调类。任意紧集包含于某个 \(K_j\)，所以这个类具有陈述中的全部相对化性质，且与所选穷尽无关。紧空间的局部有限链本来就是有限链。
\end{proof}

为了从紧支对偶推出局部有限对偶，需要一个中间结论。对紧集 \(K\subset M\)，定义
\[
	\mathcal N^k(K;\mathcal L)=\varinjlim_{V\supset K\text{ 开}}H^k(V;\mathcal L|_V),
\]
其中过渡映射是缩小邻域的限制。这记录一个在 \(K\) 的某个邻域上定义的类，并把在更小邻域上一致的类识别。任意紧集可能非常不规则，邻域极限与 \(H^k(K;\mathcal L|_K)\) 的比较需要另加条件。以下直接使用这个邻域定义。

\begin{lemma}[紧集的邻域对偶]\label{g2:lem:neighborhood-duality}
	设 \(M,R,\mathcal L\) 满足\cref{g2:thm:poincare-duality-compact}的条件，\(K\subset M\) 紧。在邻域 \(V\) 中以 \([M]_K\) 的切除类取帽积，得到自然同构
	\begin{equation}\label{g2:eq:neighborhood-duality}
		\Phi_K:\mathcal N^k(K;\mathcal L)\xrightarrow{\ \cong\ }
		H_{n-k}(M,M\setminus K;\mathcal O_R\otimes_R\mathcal L).
	\end{equation}
	它与紧集包含所诱导的两侧限制相容。
\end{lemma}
\begin{proof}
	映射本身由相对帽积定义：\(H^k(V;\mathcal L)\) 帽 \(H_n(V,V\setminus K;\mathcal O_R)\) 后，再映入 \(M\) 的相对同调。缩小邻域时，基本类的切除与帽积自然性保证相容。

	构造证明同构所需的正合列。取紧集 \(C\) 和开邻域 \(V\)，使 \(K\subset V\subset\overline V\subset\operatorname{int}C\)。对三元组
	\[
		M\supset V\cup(M\setminus C)\supset M\setminus C
	\]
	使用上同调长正合列。左端的相对群是 \(H^*(M,M\setminus(C\setminus V);\mathcal L)\)，中间是 \(H^*(M,M\setminus C;\mathcal L)\)，第三项由不交开并识别为 \(H^*(V;\mathcal L)\)。让 \(C\) 增大、\(V\) 缩小，便得到
	\[
		\cdots\to H_c^k(M\setminus K;\mathcal L)\to H_c^k(M;\mathcal L)
		\to\mathcal N^k(K;\mathcal L)\to H_c^{k+1}(M\setminus K;\mathcal L)\to\cdots.
	\]
	其共尾性需分别核对：每个 \(M\) 内的紧支集可包含于足够大的 \(C\)；每个 \(M\setminus K\) 内的紧集可先用邻域 \(V\) 与 \(K\) 分开，再包含于 \(C\setminus V\)；任意 \(K\) 的邻域都包含一个闭包紧的更小邻域，且该闭包可放入某个 \(\operatorname{int}C\)。局部紧 Hausdorff 性和 \(K\) 的紧性保证这些选择。直接极限的正合性已在\cref{g2:prop:compact-support-MV}的证明中核对。

	把这条序列与空间对 \((M,M\setminus K)\) 的同调长正合列比较，其他位置使用已证的紧支对偶，中间使用 \(\Phi_K\)。普通箭头由相对帽积自然性交换。连接箭头也可直接检查：在上述三元组中，把 \(V\) 上的上圈 \(\varphi\) 延拓为一个在 \(M\setminus C\) 上为零的上链 \(\widehat\varphi\)，则 \(\delta\widehat\varphi\) 的支集位于 \(C\setminus V\)。取 \([M]_C\) 的代表 \(c\)，用开覆盖 \(V,M\setminus K\) 作相容的小链比较，\(\widehat\varphi\frown c\) 在模去 \(M\setminus K\) 后代表 \(\Phi_K[\varphi]\)。由于 \(\widehat\varphi\frown\partial c=0\)，有
	\[
		\partial(\widehat\varphi\frown c)=(-1)^{k+1}(\delta\widehat\varphi)\frown c.
	\]
	这正是两条连接映射相差的符号。上链的延拓只需逐单形选值；其余边界承担连接映射，不要求这个延拓本身是上链映射。小链比较采用\cref{g2:lem:twisted-small-chain-retraction}，也不会改变任一相对类。调整竖箭头的符号后应用五引理，得到 \(\Phi_K\) 同构。对更小紧集的自然性在构造时已经成立。
\end{proof}

邻域极限为何有用，可以从一个相容的序列看出来。若 \(K_j\subset U_j\subset K_{j+1}\)，其中 \(K_j\) 紧、\(U_j\) 开且递增穷尽，则对每个次数 \(k\) 有
\[
	H^k(U_{j+1};\mathcal L)\longrightarrow\mathcal N^k(K_{j+1};\mathcal L)
	\longrightarrow H^k(U_j;\mathcal L)\longrightarrow\mathcal N^k(K_j;\mathcal L).
\]
第二箭头存在，是因为 \(K_{j+1}\) 的每个邻域都包含 \(U_j\)。相邻两个箭头的复合就是原来各群塔的过渡映射。因此普通邻域群与 \(\mathcal N^k(K_j)\) 是一个交错群塔的两个共尾子塔。

\ProseParagraphBreak

交错不会改变 \(\varprojlim\) 和 \(\varprojlim{}^1\)。第一项由相容序列唯一填回被省去坐标可见；第二项也可从\cref{g2:prop:locally-finite-limit}的余核定义验证。若群塔箭头为 \(q_i:G_{i+1}\to G_i\)，每隔一个坐标保留一次时，把乘积元素 \((b_i)\) 送到 \((b_{2i-1}+q_{2i-1}b_{2i})_i\)。它把 \(\Delta\) 的像送到子塔的 \(\Delta\) 像，并在余核上满射；若所得类为零，先解出奇数坐标，再令 \(x_{2i}=b_{2i}+q_{2i}x_{2i+1}\)，即可把原类写成 \(\Delta(x_i)\)，证明单射。这说明两个共尾子塔在这两个极限项上确实一致。

\begin{theorem}[局部有限 Poincaré 对偶]\label{g2:thm:poincare-duality-locally-finite}
	设 \(M\) 是 Hausdorff、第二可数、无边界的拓扑 \(n\) 维流形，\(R\) 为含幺交换环，\(\mathcal L\) 为 \(R\) 模局部系统。帽积给出同构
	\begin{equation}\label{g2:eq:poincare-locally-finite}
		D_{\mathrm{lf},M}:H^k(M;\mathcal L)\xrightarrow{\ \cong\ }
		H_{n-k}^{\mathrm{lf}}(M;\mathcal O_R\otimes_R\mathcal L),
		\qquad\alpha\longmapsto\alpha\frown[M]_{\mathrm{lf}}.
	\end{equation}
	它与系数同态及开集限制相容。将紧支上同调类忘记支集后取此映射，等于先取 \(D_{c,M}\)，再把有限链视为局部有限链。
\end{theorem}
\begin{proof}
	帽积及支集相容性由\cref{g2:prop:support-cap-pairings}给出。关键是证明它为同构，同时保留逆极限的修正项。选一个代表 \([M]_{\mathrm{lf}}\) 的局部有限圈 \(c\)。递归选择紧集 \(K_j\) 和相对紧开集 \(U_j\)，使
	\[
		K_j\subset U_j\subset K_{j+1},\qquad
		\overline U_j\subset\operatorname{int}K_{j+1},\qquad \bigcup_jK_j=M,
	\]
	并使 \(U_j\) 包含 \(c\) 中所有遇到 \(K_j\) 的单形的像。这样的单形只有有限条，其像的并紧；把这个并及 \(K_j\) 放入相对紧开集，再把闭包及预先选定穷尽的下一紧集放入 \(\operatorname{int}K_{j+1}\)，便可完成递归。

	令 \(c_j\) 为 \(c\) 中遇到 \(K_j\) 的那些项之和。它是 \(U_j\) 中的有限链，且 \(\partial c_j\) 完全位于 \(M\setminus K_j\)：一个遇到 \(K_j\) 的面，其在 \(\partial c\) 中的全部贡献都来自已经保留的单形，所以这些贡献仍抵消。因而 \(c_j\) 表示 \([M]_{K_j}\)。同时 \(c_{j+1}-c_j\) 的每个单形都避开 \(K_j\)。

	取上链复形塔 \(E_j^*=C^*(U_j;\mathcal L)\)，其限制映射逐次满射，因为可在不完全位于 \(U_j\) 的单形上任意补值。每条奇异单形的像紧，故最终位于某个 \(U_j\)，于是 \(\varprojlim E_j^*=C^*(M;\mathcal L)\)。另一方面令
	\[
		Q_{j,*}=C_*(M,M\setminus K_j;\mathcal O_R\otimes_R\mathcal L),
		F_j^k=Q_{j,n-k},\qquad \mathrm d_{F_j}^k=(-1)^{k+1}\partial.
	\]
	最后这个符号只是把降次链复形重排为升次上链复形。帽积给出上链映射
	\[
		E_j^*\longrightarrow F_j^*,\qquad \varphi\longmapsto\varphi\frown c_j.
	\]
	边界公式以及 \(\partial c_j\) 避开 \(K_j\) 验证与微分相容；\(c_{j+1}-c_j\) 避开 \(K_j\) 则验证这些映射与群塔过渡严格相容。在逆极限上，它就是帽积 \(\varphi\frown c\)。

	在第 \(j\) 项的上同调上，此映射是
	\[
		H^k(U_j;\mathcal L)\longrightarrow\mathcal N^k(K_j;\mathcal L)
		\xrightarrow{\Phi_{K_j}}H_{n-k}(Q_{j,*}).
	\]
	第二箭头由邻域对偶为同构，第一箭头和前述交错系统一起，在 \(\varprojlim\) 及 \(\varprojlim{}^1\) 上都诱导同构。对两个上链复形塔分别应用\cref{g2:cor:cochain-limit-sequence}，便得到两条自然短正合列；左端比较的是次数 \(k-1\) 的 \(\varprojlim{}^1\)，右端比较的是次数 \(k\) 的 \(\varprojlim\)，两者均为同构。中间箭头因此为同构，而中间两群正是 \(H^k(M;\mathcal L)\) 与 \(H_{n-k}^{\mathrm{lf}}(M;\mathcal O_R\otimes_R\mathcal L)\)。这一比较同时覆盖修正项非零的情形。

	系数自然性来自链公式。对于开嵌入，局部有限基本类的限制具有同一组局部值，由唯一性等于开子流形的基本类；按\cref{g2:prop:support-functoriality}中相对商塔的限制构造，帽积与这些商及切除相容，故得到开集自然性。
\end{proof}

\subsection{闭流形的 Poincaré 对偶}
\label{g2:subsec:0603:04}
紧空间上，每条局部有限链有限，每个上链也自动具有紧支集。因此刚才两种对偶在闭流形上重合。若再选择一个 \(R\) 定向，将 \(\mathcal O_R\) 识别为常值系统，便得到通常不显式写出支集和扭系数的公式。

\begin{corollary}[闭流形的 Poincaré 对偶]\label{g2:cor:closed-poincare-duality}
	设 \(M\) 为闭的 Hausdorff、第二可数拓扑 \(n\) 维流形，\(R\) 为含幺交换环，\(\mathcal L\) 为 \(R\) 模局部系统。则
	\[
		H^k(M;\mathcal L)\xrightarrow{\ \alpha\mapsto\alpha\frown[M]\ }
		H_{n-k}(M;\mathcal O_R\otimes_R\mathcal L)
	\]
	为同构。若 \(M\) 已选定 \(R\) 定向，则化为 \(H^k(M;\mathcal L)\cong H_{n-k}(M;\mathcal L)\)；常值系数时特别得到 \(H^k(M;R)\cong H_{n-k}(M;R)\)。若整体定向反转，对偶映射乘以 \(-1\)。
\end{corollary}
\begin{proof}
	紧性使 \(H_c^k=H^k\)，所以第一结论是\cref{g2:thm:poincare-duality-compact}的直接应用。设所选 \(R\) 定向为 \(\theta\)，用\cref{g2:eq:orientation-square}的规范配对 \(b(\theta,-)\colon\mathcal O_R\to\underline R\) 平凡化定向系统，规范扭基本类就变为该定向的普通基本类。整体反向把这个类变成相反数，帽积对它线性，故对偶映射也反号。
\end{proof}

\begin{example}[圆周上的两个支集版本]\label{g2:exm:duality-line-circle}
	对已定向的 \(S^1\)，单位上同调类 \(1\in H^0(S^1;\mathbb Z)\) 帽基本圈得到该圈，正向一次生成元帽基本圈得到零次的点类。因此 \(H^0\cong H_1\)、\(H^1\cong H_0\)。把圆周换成 \(\mathbb R\)，普通 \(H_1(\mathbb R)=0\)，原来的第一条公式就不成立。正确的两组结果是
	\[
		H_c^1(\mathbb R;\mathbb Z)\cong H_0(\mathbb R;\mathbb Z)=\mathbb Z,
		\qquad H^0(\mathbb R;\mathbb Z)\cong H_1^{\mathrm{lf}}(\mathbb R;\mathbb Z)=\mathbb Z.
	\]
	第二个右端的生成元可由全部整数区间的有向和 \(\sum_{j\in\mathbb Z}[j,j+1]\) 表示，相邻端点抵消，且每个紧集只遇有限项。它不是普通有限链；这一点正对应左端允许常值上圈具有非紧支集。
\end{example}

流形对偶将互补次数的上同调与同调比较；同一次数中，上同调与同调的线性对偶之间还由万有系数定理联系。整数系数下，\cref{g2:thm:uct-cohomology}保留了低一次同调的 Ext 项，挠信息由此参与上同调。下一节加入边界条件，最后一节再将对偶与取值配对结合，分别构造自由部分的交形式和挠部分的链接配对。
% 纲要条目排为编号小节标题，并进入目录与书签；说明段仍保留于纲要注释中。

% !TeX root = ../../Book2.tex
% 纲要标题：### 6.4 边界与扭系数的修正
% 纲要位置：几何学计划纲要.md，第 1198 行。
% 正文状态：未撰写；本文件目前是框架。

% 纲要内容开始
% 
% * 相对基本类和边界基本类
% * Poincaré–Lefschetz 两种配对
% * 非定向情形的定向局部系统
% * 非紧情形为何必须改变支集
% 
% 纲要内容结束

\section{边界与扭系数的修正}
\label{g2:sec:0604}

待撰写

\subsection{相对基本类和边界基本类}
\label{g2:subsec:0604:01}
待撰写

\subsection{Poincaré–Lefschetz 两种配对}
\label{g2:subsec:0604:02}
待撰写

\subsection{非定向情形的定向局部系统}
\label{g2:subsec:0604:03}
待撰写

\subsection{非紧情形为何必须改变支集}
\label{g2:subsec:0604:04}
待撰写
% 纲要条目排为编号小节标题，并进入目录与书签；说明段仍保留于纲要注释中。

% !TeX root = ../../Book2.tex
% 纲要标题：### 6.5 交形式及拓扑约束
% 纲要位置：几何学计划纲要.md，第 1255 行。
% 正文状态：未撰写；本文件目前是框架。

% 纲要内容开始
% 
% * 中维相交形式的对称性
% * 曲面和四维基本模型
% * 扭部分及链接配对的入口
% * Betti 数与 Euler 示性数的限制
% 
% ---
% 
% 纲要内容结束

\section{交形式及拓扑约束}
\label{g2:sec:0605}

待撰写

\subsection{中维相交形式的对称性}
\label{g2:subsec:0605:01}
待撰写

\subsection{曲面和四维基本模型}
\label{g2:subsec:0605:02}
待撰写

\subsection{扭部分及链接配对的入口}
\label{g2:subsec:0605:03}
待撰写

\subsection{Betti 数与 Euler 示性数的限制}
\label{g2:subsec:0605:04}
待撰写
% 纲要条目排为编号小节标题，并进入目录与书签；说明段仍保留于纲要注释中。


\begin{chaptersummary}
待撰写
\end{chaptersummary}

\BookExercises
待撰写

% !TeX root = ../../Book2.tex
% 纲要标题：## 第7章　de Rham、Čech 与层上同调比较
% 纲要位置：几何学计划纲要.md，第 1264 行。
% 正文状态：初稿；尚未正式审核。

% 纲要内容开始
% 
% > 本章目标与证明范围：把三个计算模型及其比较映射真正建立起来。de Rham 部分明确要求光滑流形和形式；本章不使用后置 Hodge 定理来证明 de Rham 后又反向引用。 积分比较所需的单形及单形同伦的相对光滑逼近在本章就地证明，不借后置的一般横截性。
% 
% 纲要内容结束

\chapter{de Rham、Čech 与层上同调比较}
\label{g2:ch:07}
\BookChapterNumber{g2:ch:07}

\begin{chapterabstract}
	局部存在的函数、形式或截面未必能够整体粘合。本章从层的茎和局部正合性出发，以典范消解定义层上同调，证明松弛层、细层的消失及无上同调消解比较。Čech 复形把粘合条件写在开集的交叠上；细化同伦、覆盖比较与覆盖极限说明其与层模型的关系。对光滑流形，就地证明单形及单形同伦的相对光滑逼近，再以积分和双复形建立 de Rham 同构，核对乘法、相对与紧支版本。最后通过周期、射影空间及平坦局部系数，说明实形式记录了哪些拓扑信息、又遗漏了哪些整数信息。
\end{chapterabstract}

\begin{learninggoals}
	\begin{itemize}
		\item 区分预层、层、层化与茎，说明层上的局部正合为什么不保证整体截面正合。
		\item 构造层上同调，证明松弛与细层的消失，并使用无上同调消解计算整体障碍。
		\item 写出 Čech 微分及细化映射，证明指标选择不影响上同调，辨认单份覆盖与覆盖极限的比较条件。
		\item 保持公共面地光滑化奇异单形，证明积分给出自然的 de Rham 同构，并核对乘法和支集。
		\item 计算环面、球面与射影空间的形式类，区分实周期、整周期及挠信息，将局部系统实现为平坦丛的形式复形。
	\end{itemize}
\end{learninggoals}

在圆周的一小段开弧上，总能选择一个连续角度；绕完整个圆周后，角度却可能增加 \(2\pi\)。如果用几段开弧覆盖圆周，每一段的角度都没有问题，问题出现在交叠处：两份角度的差是常数，这些差能否通过修改各片的常数同时消掉？闭一次形式的局部原函数也面临同一个问题。局部求解与整体粘合之间的距离，正是本章要记录的上同调。

\ProseParagraphBreak

第4章把上同调表示为奇异上链，并建立杯积；第5章用局部系统记录绕行运输，第6章用基本类把上同调与同调联系起来。现在我们增加可以直接微分和积分的表示。de Rham 理论\footnote{Georges de Rham（1903--1990），瑞士数学家，研究微分拓扑与代数拓扑，建立了微分形式与拓扑上同调之间的比较。}以闭形式模恰当形式为对象，Čech 理论\footnote{Eduard Čech（1893--1960），捷克数学家，研究射影微分几何与拓扑，发展了以开覆盖组织空间拓扑信息的方法。}以覆盖交叠上的数据为对象，层上同调则允许把同一种局部到整体的问题用于不同类型的系数。对象看起来不同，比较映射必须说明它们为何测量同一障碍。

\begin{table}[htbp]
	\centering
	\caption{本章各模型保存的数据及比较办法}
	\label{g2:tab:cohomology-models}
	\begin{tabular}{lll}
		\toprule
		模型 & 直接保存的数据 & 本章使用的比较 \\
		\midrule
		微分形式 & 局部可微的交替场 & 在光滑单形上积分 \\
		奇异上链 & 每条单形的数值 & 小链与光滑链比较 \\
		Čech 复形 & 各级交叠上的截面 & 细化及覆盖双复形 \\
		层上同调 & 取整体截面后的障碍 & 无上同调消解 \\
		\bottomrule
	\end{tabular}
\end{table}

层态射的满射性给出局部提升，消解把这些提升之间的粘合障碍组织为一个复形；覆盖则将复形写成交叠上的数据。对微分形式，第一卷第14.2节的同伦公式提供局部原形式，第12.3节的单位分解用来组合各片的形式，第14.4节的带角 Stokes 公式使积分与边界相容。第三节先构造保持公共面的光滑逼近，再将这几项局部工具接成积分比较的整体同构。

\ProseParagraphBreak

本章的层论构造先允许一般拓扑空间，覆盖极限比较明确限于 Hausdorff、第二可数、局部紧空间。形式部分取无边界光滑流形和实系数；整数信息需要另由整数上同调比较，不能从微分形式的实向量空间结构中自动恢复。紧支条件也一直单独保留：一个闭紧支形式可能有普通原形式，却没有紧支原形式。

% !TeX root = ../../Book2.tex
% 纲要标题：### 7.1 层与局部正合性
% 纲要位置：几何学计划纲要.md，第 1268 行。
% 正文状态：初稿。

% 纲要内容开始
% 
% * 预层、层化与茎
% * 层的核余核及茎上的正合检测
% * 整体截面与层上同调
% * 松弛层、细层及无上同调机制
% 
% 纲要内容结束

\section{层与局部正合性}
\label{g2:sec:0701}

一个闭一形式在每个小球上都有原函数，原函数却未必能拼成整体函数。此时应同时保存各开集上的解、缩小定义域的限制，以及重叠处的差。层把这些局部资料组织在一起；层上同调则记录局部解无法整体拼合的障碍。本节先把“局部正合”与“整体正合”分开，再构造能够测量两者差别的群。

以下层均为含幺交换环 \(R\) 上的模层；只研究加法时取 \(R=\mathbb Z\)。底空间在需要局部有限覆盖时假设为 Hausdorff、第二可数、局部紧空间，特别包含本书的流形。预层、层化、正合性和松弛层的论证本身不需要这些拓扑条件。

\subsection{预层、层化与茎}
\label{g2:subsec:0701:01}
局部资料先按开集组织：定义域缩小时有一个限制映射，连续缩小两次应与直接限制一致。预层记录这些映射，层再要求相容的局部资料能够唯一粘合。

\begin{definition}\label{g2:def:presheaf-sheaf}
	拓扑空间 \(X\) 上的一个 \(R\) 模\term[yuceng]{预层}[presheaf] \(\mathcal F\)，给每个开集 \(U\) 指定一个模 \(\mathcal F(U)\)，给每个包含 \(V\subseteq U\) 指定限制同态 \(s\mapsto s|_V\)，满足恒等限制及逐次限制律。若对任意开覆盖 \(U=\bigcup_iU_i\)，彼此相容的截面 \(s_i\in\mathcal F(U_i)\) 都能唯一粘成 \(s\in\mathcal F(U)\)，便称 \(\mathcal F\) 为\term[ceng]{层}[sheaf]。空覆盖约定给出 \(\mathcal F(\varnothing)=0\)。与所有限制映射交换的同态族称为预层或层的态射。
\end{definition}

连续函数、光滑函数及光滑微分形式都是层：正则性是局部条件，相同的局部函数又能按其值逐点粘合。相反，“定义域上有界的连续函数”只形成预层；在 \(\mathbb R\) 上用有界开区间覆盖，各处的函数 \(x\mapsto x\) 相容，粘成的函数却不再有界。层公理要求保留任意局部可粘资料，所以有界性这种整体限制通常不应放进层的定义。

\begin{definition}\label{g2:def:stalk}
	对 \(x\in X\)，把所有 \(x\) 的开邻域上的截面作如下识别：\((U,s)\) 与 \((V,t)\) 等价，当且仅当存在 \(x\in W\subseteq U\cap V\)，使 \(s|_W=t|_W\)。所得模
	\[
		\mathcal F_x=\varinjlim_{x\in U}\mathcal F(U)
	\]
	称为 \(\mathcal F\) 在 \(x\) 的\term[jing]{茎}[stalk]，截面 \(s\) 的等价类 \(s_x\) 称为它的芽。加法和数乘在共同邻域上进行，再取等价类。
\end{definition}

芽保留的是某个未指定小邻域内的全部信息，通常远多于点值。例如光滑函数芽 \(f_x\) 决定所有阶导数，\(f(x)=0\) 却不意味着 \(f_x=0\)。后一个等式要求 \(f\) 在某邻域上恒为零。对层而言，两个截面相等当且仅当它们在每点的芽相等：芽相等给出局部相等的覆盖，再用粘合的唯一性。

\begin{proposition}[层化]\label{g2:prop:sheafification}
	任意 \(R\) 模预层 \(\mathcal F\) 有层 \(\mathcal F^{\#}\) 及自然态射 \(\eta:\mathcal F\to\mathcal F^{\#}\)，满足：\(\eta_x\) 在每点都是同构；对任意层 \(\mathcal G\)，每个预层态射 \(\mathcal F\to\mathcal G\) 都唯一经 \(\eta\) 分解。
\end{proposition}
\begin{proof}
	定义 \(\mathcal F^{\#}(U)\) 为芽的族 \((a_x)_{x\in U}\)，要求每个 \(x\in U\) 都有邻域 \(V\subseteq U\) 及 \(s\in\mathcal F(V)\)，使 \(a_y=s_y\) 对所有 \(y\in V\) 成立。限制就是删去定义域外的芽。局部代表可以逐点选择，故相容的芽族唯一粘合；相加时取共同邻域即可，所以它是模层。令 \(\eta(s)=(s_x)_x\)。
	在 \(x\) 附近，一个层化截面按定义来自原预层截面，所以 \(\eta_x\) 满射。若两个原截面的像有同一芽，则在某邻域内它们的芽族一致；尤其在 \(x\) 处原来的芽相等，故 \(\eta_x\) 单射。给定 \(f:\mathcal F\to\mathcal G\)，把局部代表 \(s\) 送至 \(f(s)\)；不同代表在每点有同一芽，因 \(\mathcal G\) 为层，其像在重叠处一致，遂能唯一粘合。局部代表也迫使此延拓唯一。
	因此所得分解由下图表达，其中右下方的箭头由局部代表确定：
	\[
		\begin{tikzcd}[column sep=large]
			\mathcal F \arrow[r,"\eta"] \arrow[dr,"f"'] & \mathcal F^{\#} \arrow[d,dashed,"\exists!\,\bar f"] \\
			& \mathcal G.
		\end{tikzcd}
	\]
\end{proof}

把每个非空开集都赋予同一个模 \(A\)、限制均取恒等，得到常值预层；它的层化记为 \(\underline A\)，称为\term[changzhiceng]{常值层}[constant sheaf]。其截面恰是取值于离散模 \(A\) 的局部常值函数。一个不连通开集上的截面可以在不同分支取不同值，所以“常值层”并不表示每个开集只有一个共同常数。第5章的局部系统也可视为层：在每个平凡化邻域取局部常值截面，以运输给出的过渡映射粘合。

\subsection{层的核余核及茎上的正合检测}
\label{g2:subsec:0701:02}
局部方程 \(f(s)=0\) 可以直接逐开集求解，商却更微妙：一个局部商类可能只能在不同小邻域选不同代表。因此核可以逐开集取，余核则通常要层化。

\begin{proposition}[核、余核与正合性]\label{g2:prop:stalk-exactness}
	设 \(f:\mathcal F\to\mathcal G\) 为拓扑空间 \(X\) 上的 \(R\) 模层态射。核由 \(U\mapsto\ker f_U\) 给出，余核为预层 \(U\mapsto\mathcal G(U)/f\mathcal F(U)\) 的层化；像为预层像的层化。它们在点 \(x\) 的茎分别是 \(\ker f_x\)、\(\operatorname{coker}f_x\) 和 \(\operatorname{im}f_x\)。因而层的复形正合，当且仅当每点的茎复形正合；层态射为同构，当且仅当所有茎映射为同构。
\end{proposition}
\begin{proof}
	核截面的相容族在 \(\mathcal F\) 中粘合，像在 \(\mathcal G\) 中局部为零，故整体也为零；这证明核已是层。商的泛性质先在预层中成立，再用\cref{g2:prop:sheafification}得到层范畴中的余核泛性质。
	关键是取芽保持正合。若 \(s_x\) 在某映射下为零，就能缩小代表 \(s\) 的定义域，使其像作为截面已经为零。若一个芽属于某像，选择该原像芽的代表，再缩小邻域便得到截面层面的原像。有限次等式和有限个代表只需取有限个邻域的交；这证明有向极限保持模的核、像及余核。层化不改茎，故上述茎公式成立。
	若每点核与像相同，则核截面在每点附近来自原像，正是像层的定义；反向可直接取茎。最后，茎同构意味着核层与余核层的所有茎为零，而所有茎为零的层只能是零层；也可以把局部逆像利用唯一性直接粘成逆态射。
\end{proof}

特别地，满的层态射只保证截面\emph{局部}有原像，不保证每个开集上的截面映射满射。后文写
\[
	0\longrightarrow\mathcal F\longrightarrow\mathcal G\longrightarrow\mathcal H\longrightarrow0
\]
时，正合性均按\cref{g2:prop:stalk-exactness}理解。层化使商对象补足局部代表，因而核与余核满足通常的模运算规则；涉及核、像的图追逐可以逐茎进行，但追逐所得的局部原像仍可能没有整体选择。

\subsection{整体截面与层上同调}
\label{g2:subsec:0701:03}
记 \(\Gamma(U,\mathcal F)=\mathcal F(U)\)，特别称 \(\Gamma(X,\mathcal F)\) 为整体截面模。对上述短正合列，
\[
	0\longrightarrow\Gamma(X,\mathcal F)\longrightarrow\Gamma(X,\mathcal G)\longrightarrow\Gamma(X,\mathcal H)
\]
总是正合：映到零的整体截面在局部属于核，核本身已经是层，故各局部原像唯一粘合。这叫整体截面的\term[zuozhenghe]{左正合性}[left exactness]，最后一箭却可能不满。

\begin{example}\label{g2:exm:sheaf-antiderivative-obstruction}
	在 \(S^1\) 上，令 \(\mathcal C^{\infty}\) 为实光滑函数层，\(\Omega^1\) 为实一形式层。局部区间上的积分说明
	\[
		0\longrightarrow\underline{\mathbb R}\longrightarrow\mathcal C^{\infty}\xrightarrow{\mathrm d}\Omega^1\longrightarrow0
	\]
	在茎上正合。可是若把 \(S^1\) 写成 \(\mathbb R/\mathbb Z\)，形式 \(\mathrm dt\) 沿一周的积分为 \(1\)，而任何 \(\mathrm df\) 沿一周的积分为 \(0\)。所以整体截面上的最后一箭不满。局部原函数在重叠处只相差常数；这些常数无法全被重新选择的局部常数消去，才是整体障碍。
\end{example}

为记录局部提升的粘合障碍，先允许在每点独立指定一个芽，再取商记录这些任意芽与真实局部截面之间的差别。反复进行这两项操作，就能构造一条具有自然映射的消解。对层 \(\mathcal F\)，设
\[
	T\mathcal F(U)=\prod_{x\in U}\mathcal F_x.
\]
这里不要求芽随点连续，也不要求它们由同一个局部截面产生；限制只是投影。故 \(T\mathcal F\) 是层，且 \(\mathcal F\to T\mathcal F\)，\(s\mapsto(s_x)_x\)，为单射。取层余核 \(Q\mathcal F=T\mathcal F/\mathcal F\)，再迭代，令
\begin{equation}\label{g2:eq:canonical-sheaf-resolution}
	G^q\mathcal F=TQ^q\mathcal F,\qquad
	0\longrightarrow\mathcal F\longrightarrow G^0\mathcal F\xrightarrow{\mathrm d_G}G^1\mathcal F\xrightarrow{\mathrm d_G}\cdots,
\end{equation}
其中 \(Q^0\mathcal F=\mathcal F\)，微分先取商 \(TQ^q\mathcal F\to Q^{q+1}\mathcal F\)，再嵌入 \(TQ^{q+1}\mathcal F\)。像与核都等于中间的 \(Q^{q+1}\mathcal F\)，所以这是一条正合的层复形，称为本节的典范消解。

\begin{definition}\label{g2:def:sheaf-cohomology}
	拓扑空间 \(X\) 上的 \(R\) 模层 \(\mathcal F\) 的\term[cengshangtongdiao]{层上同调}[sheaf cohomology]定义为
	\[
		H^q(X,\mathcal F)=H^q\bigl(\Gamma(X,G^\bullet\mathcal F)\bigr),\qquad q\ge0.
	\]
	若对所有 \(q>0\) 都有 \(H^q(X,\mathcal F)=0\)，称 \(\mathcal F\) 在 \(X\) 上\term[wushangtongdiao]{无上同调}[acyclic]。
\end{definition}

零次核恰是 \(\Gamma(X,\mathcal F)\)，故 \(H^0(X,\mathcal F)=\Gamma(X,\mathcal F)\)。在开集 \(U\) 上限制时，茎、层商及上述逐点乘积都与限制交换，因此 \(G^\bullet\mathcal F|_U=G^\bullet(\mathcal F|_U)\)。这既定义了上同调的限制，也说明 \(H^q(U,\mathcal F)\) 无歧义。

\begin{proposition}[层上同调的长正合列]\label{g2:prop:sheaf-long-exact}
	拓扑空间 \(X\) 上每条 \(R\) 模层短正合列 \(0\to\mathcal F\to\mathcal G\to\mathcal H\to0\) 都给出自然长正合列
	\[
		0\to H^0(X,\mathcal F)\to H^0(X,\mathcal G)\to H^0(X,\mathcal H)
		\xrightarrow{\partial}H^1(X,\mathcal F)\to\cdots.
	\]
	连接同态 \(\partial s\) 为零当且仅当 \(s\in\Gamma(X,\mathcal H)\) 能整体提升到 \(\mathcal G\)。
\end{proposition}
\begin{proof}
	茎序列正合，而模的任意乘积保持短正合：满射的每个坐标各选一个原像即可。因此 \(T\) 以及 \(\Gamma(X,T(-))\) 都保持短正合。把短正合列与其 \(T\) 像排成两行，竖箭为单射；逐茎追逐得到余核列 \(0\to Q\mathcal F\to Q\mathcal G\to Q\mathcal H\to0\) 正合。更具体地，右侧一个商类先提升到中间 \(T\) 项，差额来自左侧 \(T\) 项；若中间商类落到零，再减去中间原层的局部提升，便回到左侧商类。这也证明 \(Q\) 正合，迭代后 \(Q^q\) 正合。
	故 \(\Gamma(X,G^\bullet(-))\) 将原序列变成逐次短正合的上链复形序列。连接映射的构造与第2章的链版本相同：将右侧闭上链提升至中间，对其取微分，再由左侧单射识别所得闭上链。两种提升相差中间的左侧上链，其微分只改变一个边界。若连接类为零，减去该左侧原像后，中间提升便是闭上链；反之闭提升的微分为零。这证明连接处的正合性，另外两处由同样的提升和取像直接得到。所有步骤与短正合列的态射交换，故自然。零次的结论正是最后一句。
\end{proof}

每个层态射都逐点作用于芽，并诱导层商上的态射，因而自动给出典范消解之间的链映射。接下来先考察使高次上同调消失的延拓和局部组合机制；下一节再据此用更适合计算的消解替换典范消解。

\subsection{松弛层、细层及无上同调机制}
\label{g2:subsec:0701:04}
\begin{definition}\label{g2:def:flabby-fine}
	若层 \(\mathcal F\) 对所有开集包含 \(V\subseteq U\) 的限制 \(\mathcal F(U)\to\mathcal F(V)\) 都满射，称它为\term[songchiceng]{松弛层}[flabby sheaf]。在仿紧 Hausdorff 空间 \(X\) 上，若对每个局部有限开覆盖 \(\{U_i\}\)，都有层端同态 \(\rho_i:\mathcal F\to\mathcal F\) 及局部有限闭集族 \(K_i\subseteq U_i\)，使 \(\rho_i|_{X\setminus K_i}=0\) 且 \(\sum_i\rho_i=1\)，则称 \(\mathcal F\) 为\term[xiceng]{细层}[fine sheaf]。端同态之和在每个足够小的邻域都是有限和。
\end{definition}

松弛性允许任意延拓；细性允许把局部资料乘上有支集控制的权重后相加。两者解决障碍的方式不同。\(T\mathcal F\) 总是松弛，因为一族芽可以在较大开集其余各点取零。光滑函数层通常不松弛，例如 \(1/x\) 在 \((0,1)\) 上无法延拓成 \(\mathbb R\) 上的光滑函数；但它是细层。

\begin{lemma}[松弛核的提升性质]\label{g2:lem:flabby-lifting}
	若拓扑空间 \(X\) 上的层短正合列 \(0\to\mathcal A\to\mathcal B\to\mathcal C\to0\) 的核 \(\mathcal A\) 松弛，则 \(\mathcal B(U)\to\mathcal C(U)\) 对每个开集 \(U\) 满射。若 \(\mathcal B\) 也松弛，则 \(\mathcal C\) 松弛。
\end{lemma}
\begin{proof}
	固定 \(s\in\mathcal C(U)\)，考虑它在开子集 \(V\subseteq U\) 上的提升 \(b_V\)，按延拓排序。全序族可由层公理在开集之并上粘合，故 Zorn 引理给出极大提升。若 \(x\in U\setminus V\)，茎上的满射性给出邻域 \(W\) 上的提升 \(b_W\)。在 \(V\cap W\) 上，两者之差来自 \(\mathcal A(V\cap W)\) 的唯一截面；利用松弛性把这个差延拓至 \(W\)，从 \(b_W\) 中减去后便与 \(b_V\) 相合。因此提升能延至 \(V\cup W\)，违背极大性。所以 \(V=U\)。若 \(\mathcal B\) 也松弛，先把 \(\mathcal C(V)\) 的截面提升到 \(\mathcal B(V)\)，再延至 \(\mathcal B(U)\) 并取像，就得到 \(\mathcal C\) 的松弛性。
\end{proof}

\begin{theorem}[松弛层与细层的消失]\label{g2:thm:flabby-fine-acyclic}
	拓扑空间 \(X\) 上的松弛 \(R\) 模层无上同调，其在每个开子集上的限制也无上同调。仿紧 Hausdorff 空间 \(X\) 上的细 \(R\) 模层在 \(X\) 上无上同调。
\end{theorem}
\begin{proof}
	若 \(\mathcal F\) 松弛，则在 \(0\to\mathcal F\to T\mathcal F\to Q\mathcal F\to0\) 中，前两项都松弛。\cref{g2:lem:flabby-lifting}表明 \(Q\mathcal F\) 松弛，且取整体截面仍正合。逐次迭代，典范消解的整体截面复形在正次数均正合，故上同调消失。松弛性显然保留于开子集上的限制，同一论证遂适用。
	现设 \(\mathcal F\) 细，取 \(q>0\) 及 \(z\in\Gamma(X,G^q\mathcal F)\) 满足 \(\mathrm d_Gz=0\)。层消解的正合性给出一个开覆盖，使每个开集上存在 \(b_i\) 满足 \(\mathrm d_Gb_i=z\)。用仿紧性取局部有限开细化，仍记为 \(\{U_i\}\)，并选细性给出的 \(\rho_i,K_i\)。函子性使 \(\rho_i\) 诱导消解端同态 \(G^\bullet\rho_i\)，与 \(\mathrm d_G\) 交换。由于消解构造与开限制交换，\(G^\bullet\rho_i\) 在 \(X\setminus K_i\) 上也为零。
	因此 \((G^{q-1}\rho_i)b_i\) 可在 \(U_i\) 外补零：在 \(U_i\cap(X\setminus K_i)\) 上两个定义都是零，层公理保证粘合。令 \(b\) 为这些整体截面的局部有限和。每点附近只有有限项，故加法、函子作用及微分均与此和交换，得到
	\[
		\mathrm d_Gb=\sum_i(G^q\rho_i)z
		=G^q\Bigl(\sum_i\rho_i\Bigr)z=z.
	\]
	每个正次闭上链都成为边界，证明细层无上同调。这里没有要求一般微分与任意函数相乘交换；交换的是同一个层端同态在\emph{典范消解}上诱导的映射。
\end{proof}

\begin{corollary}\label{g2:cor:smooth-module-acyclic}
	有限维、Hausdorff、第二可数的光滑流形 \(M\) 上，每个光滑函数层 \(\mathcal C^{\infty}_M\) 的模层都是细层，因而无上同调。特别地，光滑向量丛截面层及微分形式层 \(\Omega^p_M\) 均无上同调；它们在任意开子集上的限制也无上同调。
\end{corollary}
\begin{proof}
	第一卷定理12.3.5说明：对每个开覆盖，可取光滑函数 \(\chi_i\)，支集族局部有限、\(\operatorname{supp}\chi_i\subseteq U_i\)，且 \(\sum_i\chi_i=1\)。在模层上取 \(\rho_i(s)=\chi_i s\)，便满足\cref{g2:def:flabby-fine}。流形的每个开子集仍为 Hausdorff、第二可数光滑流形，故限制后的结论同样成立。
\end{proof}

应用于\cref{g2:exm:sheaf-antiderivative-obstruction}，长正合列与 \(\mathcal C^{\infty}\) 的消失给出
\[
	H^1(S^1,\underline{\mathbb R})
	\cong\Omega^1(S^1)/\mathrm d\mathcal C^{\infty}(S^1)\cong\mathbb R.
\]
最后一个同构由沿一周积分给出。若 \(a(t)\mathrm dt\) 的积分为零，\(f(t)=\int_0^t a(u)\,\mathrm du\) 是周期函数，故在圆上给出原函数；这证明其核恰为恰当形式。于是层上同调确实把局部原函数的粘合失败转成了可以计算的整体不变量。
% 纲要条目排为编号小节标题，并进入目录与书签；说明段仍保留于纲要注释中。

% !TeX root = ../../Book2.tex
% 纲要标题：### 7.2 消解和 Čech 模型
% 纲要位置：几何学计划纲要.md，第 1225 行。
% 正文状态：未撰写；本文件目前是框架。

% 纲要内容开始
% 
% * 无上同调消解及比较定理
% * 覆盖的 Čech 复形
% * 细化映射与覆盖选择
% * 无上同调覆盖的比较定理
% 
% 纲要内容结束

\section{消解和 Čech 模型}
\label{g2:sec:0702}

待撰写

\subsection{无上同调消解及比较定理}
\label{g2:subsec:0702:01}
待撰写

\subsection{覆盖的 Čech 复形}
\label{g2:subsec:0702:02}
待撰写

\subsection{细化映射与覆盖选择}
\label{g2:subsec:0702:03}
待撰写

\subsection{无上同调覆盖的比较定理}
\label{g2:subsec:0702:04}
待撰写
% 纲要条目排为编号小节标题，并进入目录与书签；说明段仍保留于纲要注释中。

% !TeX root = ../../Book2.tex
% 纲要标题：### 7.3 积分与 de Rham 定理
% 纲要位置：几何学计划纲要.md，第 1282 行。
% 正文状态：初稿；尚未正式审核。

% 纲要内容开始
% 
% * 光滑奇异链及光滑逼近
% * 积分映射与 Stokes
% * 好覆盖和局部 Poincaré 引理
% * 双复形的两种收缩
% * 乘法和支集版本的比较
% 
% 纲要内容结束

\section{积分与 de Rham 定理}
\label{g2:sec:0703}

本节把局部形式与拓扑上链联系起来。第一卷已经构造外微分、形式积分及 Stokes 公式；这里先证明连续单形能够在保留公共面的条件下光滑化，再用覆盖双复形证明积分诱导同构。整个比较均在有限维、Hausdorff、第二可数、无边界光滑流形上进行。

\subsection{光滑奇异链及光滑逼近}
\label{g2:subsec:0703:01}
奇异单形只要求连续，而拉回微分形式需要导数。因此，积分以前先要解决一个几何问题：能否把用于表示同调类的有限个单形同时光滑化，并使公共面仍然完全相同？只把每个单形单独逼近，原先抵消的边界可能不再抵消；保留公共面是下面构造的核心。

\begin{definition}[光滑奇异链]\label{g2:def:smooth-singular-chain}
	设 \(M\) 为有限维、Hausdorff、第二可数的无边界光滑流形，\(R\) 为含幺交换环。映射 \(\sigma:\Delta^k\to M\) 若可延拓为标准闭单形在其仿射包中某个开邻域上的光滑映射，就称为\term[guanghua qiyi danxing]{光滑奇异单形}[smooth singular simplex]。以这些映射为基的自由 \(R\)-模记为 \(C_k^\infty(M;R)\)。面映射 \(\delta_i:\Delta^{k-1}\to\Delta^k\) 是仿射映射，故
	\[
		\partial\sigma=\sum_{i=0}^k(-1)^i\sigma\circ\delta_i
	\]
	仍是光滑链。所得同调记为 \(H_k^\infty(M;R)\)，光滑奇异上链复形记为 \(C_\infty^k(M;R)=\operatorname{Hom}_R(C_k^\infty(M;R),R)\)。
\end{definition}

这里要求在闭单形附近延拓，因而边、顶点处也有足够的导数。单形本身有角，并不要求它的像是嵌入子流形；常值单形和微分降秩的单形都允许。光滑链的构造与比较可参阅 \cite[第18章，Smooth Singular Homology]{Lee2012Smooth}。下面把所需的相对逼近直接证明出来。

\begin{lemma}[紧像附近的欧氏模型]\label{g2:lem:compact-image-retraction}
	设 \(M\) 为有限维、Hausdorff、第二可数的无边界光滑流形，\(Q\subset M\) 紧。存在 \(Q\) 的开邻域 \(V\)、光滑嵌入 \(e:V\to\mathbb R^N\)、包含 \(e(Q)\) 的开集 \(W\subset\mathbb R^N\) 和光滑映射 \(r:W\to V\)，使 \(r(e(x))=x\) 对 \(Q\) 的某个开邻域中的 \(x\) 成立。
\end{lemma}
\begin{proof}
	第一卷推论12.3.6的闭集截断结论允许在相对紧坐标图 \((U_i,\varphi_i)\) 内选紧支光滑函数 \(\chi_i\)，使有限个集合 \(\{\chi_i>0\}\) 覆盖 \(Q\)。令 \(V=\{\sum_i\chi_i>0\}\)，定义
	\[
		e(x)=\bigl(\chi_1(x),\chi_1(x)\varphi_1(x),\ldots,\chi_s(x),\chi_s(x)\varphi_s(x)\bigr).
	\]
	每项在相应坐标图之外补零；支集紧含于图内保证补零光滑。如果 \(\chi_i(x)>0\)，从第 \(i\) 组分量的商便恢复 \(\varphi_i(x)\)。这同时证明 \(e\) 单射、逆映射在像上连续，以及 \(\mathrm{d}e_x\) 单射，所以 \(e\) 是嵌入。

	说明为什么在 \(e(Q)\) 附近能回到这张嵌入像上。在每个 \(x\in V\) 处取 \(\mathrm{d}e_x(\mathrm{T}_xV)\) 的欧氏正交补，考虑
	\[
		\Psi(x,v)=e(x)+v,\qquad v\perp\mathrm{d}e_x(\mathrm{T}_xV).
	\]
	这些正交补组成光滑向量丛：在坐标图内，切空间由一个满秩矩阵 \(A(x)\) 的列生成，其正交投影为 \(A(A^{\mathsf T}A)^{-1}A^{\mathsf T}\)，随 \(x\) 光滑变化。在零截面处，\(\mathrm{d}\Psi\) 把切向与法向分量相加，是同构。第一卷定理11.3.3的逆函数定理因而给出每个 \((x,0)\) 附近的光滑逆。

	选 \(Q\subset V_0\Subset V\)。在基点属于 \(V_0\)、\(|v|\) 充分小时，\(\Psi\) 可以同时取为单射。否则存在两组不同的 \((x_j,v_j),(y_j,w_j)\)，基点属于 \(\overline V_0\)，法向长度趋于零，而 \(\Psi(x_j,v_j)=\Psi(y_j,w_j)\)。取收敛子列后有 \(e(x)=e(y)\)，故 \(x=y\)；两组点最终都落入 \((x,0)\) 附近同一逆函数邻域，与其单射性矛盾。类似地，紧性使微分在足够小的这一邻域处处可逆。因此 \(\Psi\) 从这个开邻域到其开像是微分同胚，投影其逆映射的第一分量便得到 \(r\)。
\end{proof}

平滑目标为向量空间时，面间兼容可以用一个直接的延拓公式处理。这个步骤也解释了为什么“每个面上光滑”在单形的角处足够。

\begin{lemma}[单形边界数据的光滑延拓]\label{g2:lem:simplex-face-extension}
	设 \(b:\partial\Delta^k\to\mathbb R^N\) 连续，且在每个闭面上都可光滑延拓。则存在 \(\Delta^k\) 的开邻域上的光滑映射 \(B\)，使 \(B|_{\partial\Delta^k}=b\)。若连续映射 \(F:\Delta^k\to\mathbb R^N\) 满足 \(F|_{\partial\Delta^k}=b\)，则对每个 \(\varepsilon>0\)，可取这样的 \(B\) 使 \(\sup_{\Delta^k}|B-F|<\varepsilon\)。
\end{lemma}
\begin{proof}
	先只做边界延拓。在一个边界点附近，把经过该点的 \(q\) 个面写成 \(x_1=0,\ldots,x_q=0\)，单形位于 \(x_i\geq0\) 一侧；未经过该点的面可通过缩小邻域避开。对非空 \(I\subseteq\{1,\ldots,q\}\)，记 \(p_I\) 为把这些坐标置零的投影，\(b_I\) 为 \(b\) 在相应交面上的限制，并选其光滑延拓 \(\widetilde b_I\)。在足够小的邻域内定义
	\begin{equation}\label{g2:eq:corner-face-extension}
		B_0(x)=\sum_{\varnothing\ne I\subseteq\{1,\ldots,q\}}(-1)^{|I|+1}\widetilde b_I(p_Ix).
	\end{equation}
	若 \(x\) 位于第 \(j\) 个面内，含 \(j\) 与不含 \(j\) 的项两两抵消，只留下 \(b(x)\)。这里所有投影点仍在相应闭面内，所以不同延拓在面外怎样选择并不影响这次抵消。用边界附近的有限光滑单位分解粘合这些 \(B_0\)，再乘一个在边界附近等于一的截断函数并补零，得到整个单形邻域上的光滑 \(B_1\)。

	现在令 \(h=F-B_1\)。它连续且在边界为零。由一致连续性，\(h\) 在足够窄的边界邻域内任意小。取光滑截断 \(\chi\)，支集紧含于单形内部，在这个窄邻域之外等于一，则 \(\chi h\) 在全仿射包上补零仍连续，并满足 \(\|\chi h-h\|_\infty<\varepsilon/2\)。选非负光滑核 \(\rho_\delta\)，支集在半径 \(\delta\) 的球内且积分为一。卷积 \(h_\delta=\rho_\delta*(\chi h)\) 光滑；当 \(\delta\) 足够小时，它的支集仍紧含于单形内部，而且
	\[
		\|h_\delta-\chi h\|_\infty
		\leq\sup_{|y|\leq\delta,\,x}|\chi(x-y)h(x-y)-\chi(x)h(x)|<\varepsilon/2.
	\]
	因此 \(B=B_1+h_\delta\) 满足要求。零维时没有边界，直接取给定点值即可。
\end{proof}

\begin{lemma}[相对单形光滑逼近]\label{g2:lem:relative-simplex-smoothing}
	设 \(K\) 是有限单形复形，\(A\subset K\) 为子复形，\(M\) 为有限维、Hausdorff、第二可数的无边界光滑流形。若 \(f:|K|\to M\) 连续，且 \(f\) 在 \(A\) 的每个闭单形上光滑，则存在逐单形光滑的 \(g:|K|\to M\) 和相对 \(|A|\) 的同伦 \(H:f\simeq g\)。可以要求 \(H\) 任意接近 \(f\)：对每个单形预先指定包含其像的开集，整个单形同伦仍留在该开集内。

	结论也适用于有限个单形按面识别得到的复形，包括同一个单形的不同面之间的识别。若 \(f\) 在某个单形上原已光滑，则 \(H\) 在该单形与 \([0,1]\) 的乘积上可取光滑。给定一条连续单形同伦，若其两端及指定侧面已逐面光滑，则可保持它们不变，将同伦光滑化于棱柱的有限单形剖分上。
\end{lemma}
\begin{proof}
	把紧像 \(f(|K|)\) 放入\cref{g2:lem:compact-image-retraction}的模型中，记 \(F=e\circ f\)。先在 \(A\) 上保持 \(G=F\)，再依维数把 \(G\) 定义到其余单形。假设一个 \(k\) 维单形的边界上已经给定逐面光滑的 \(b\)，而且 \(b\) 接近 \(F\)。边界误差 \(b-F\) 可连续延拓到整个单形而不增加上界：从重心 \(a\) 出发，把 \(x\ne a\) 唯一写成 \(x=(1-t)a+ty\)，其中 \(y\in\partial\Delta^k\)、\(0<t\leq1\)，令
	\[
		E(x)=t\bigl(b(y)-F(y)\bigr),\qquad E(a)=0.
	\]
	这个延拓连续，\(F+E\) 恰好限制为 \(b\)。应用\cref{g2:lem:simplex-face-extension}，得到一个光滑 \(G\)，在边界上严格等于 \(b\)，并可任意接近 \(F+E\)。若第 \(k\) 维步骤另添的误差小于 \(\varepsilon/(\dim K+1)\)，全部步骤结束后的误差小于 \(\varepsilon\)。同维单形内部互不影响，公共面在处理它们以前已经固定，因此这些映射粘合为连续 \(G\)。对于按面识别的有限单形族，给每个已识别的面只选择一次数据，同一构造仍成立。

	取 \(\varepsilon\) 充分小，使连接 \(F(x)\) 与 \(G(x)\) 的所有线段均落在回缩的定义域 \(W\) 内。定义
	\begin{equation}\label{g2:eq:relative-smoothing-homotopy}
		g(x)=r(G(x)),\qquad H(x,t)=r\bigl((1-t)F(x)+tG(x)\bigr).
	\end{equation}
	在 \(A\) 上两端完全相等，故同伦固定 \(A\)。紧性以及 \(r\) 的连续性保证：进一步缩小误差后，每个单形的整条同伦都留在其指定开集内。若 \(f\) 在一个单形上原已光滑，公式中的 \(F\)、\(G\) 都有邻域延拓，因而 \(H\) 在该单形的棱柱上光滑。

	最后，把 \(\Delta^k\times[0,1]\) 按第2章的棱柱剖分写成有限个单形。若一条连续同伦的两端和指定侧面已经逐面光滑，就将这些面作为子复形 \(A\)，再次应用刚证的构造。得到的同伦与原同伦有相同的两端及指定侧面，并在每个棱柱单形上光滑。它在剖分界面上不必具有来自两个单形的同一法向导数；积分链只要求各闭单形光滑且公共面相同，这正是所需的相容条件。
\end{proof}

\begin{theorem}[光滑链与连续链的比较]\label{g2:thm:smooth-singular-comparison}
	设 \(M\) 为有限维、Hausdorff、第二可数的无边界光滑流形，\(R\) 为含幺交换环。包含映射诱导自然同构
	\[
		H_k^\infty(M;R)\xrightarrow{\ \cong\ }H_k(M;R).
	\]
	特别地，限制实上链得到自然同构
	\[
		H^k(M;\mathbb R)\xrightarrow{\ \cong\ }H^k(C_\infty^\bullet(M;\mathbb R)).
	\]
	这里的自然性指光滑映射之间的自然性。
\end{theorem}
\begin{proof}
	每条链只涉及有限个奇异单形。把它们以及全部面列出，相同的参数映射只保留一次，并按面关系识别其定义域。这给出上引理允许的有限单形族；即使两个不同面恰好是同一个奇异单形，它们也得到同一个光滑替代和同伦。以 \(S\) 表示这个有限族上的光滑替代，面相容性给出 \(\partial S=S\partial\)。

	设 \(z\) 是连续奇异圈。对其有限单形族应用上引理，则 \(Sz\) 为光滑圈。用同伦棱柱定义 \(P\)，第2章\cref{g2:eq:prism-identity}的边界计算逐个单形给出
	\[
		\partial P+P\partial=S-\operatorname{id}.
	\]
	原公式只使用相邻面上的同伦相同，故适用于这里逐单形给出的相容同伦。于是 \(Sz-z=\partial Pz\)，证明满射。若光滑圈 \(z\) 在连续链中满足 \(z=\partial b\)，把 \(b\)、\(z\) 及其全部面同时列出，以 \(z\) 的单形及其面为固定子复形。相对逼近保持这些单形不变，故 \(\partial Sb=Sz=z\)，证明单射。这个论证使用的都是有限 \(R\)-线性组合，不要求链的系数为整数。

	为说明实上同调的结论，回顾向量空间复形的直接事实：限制到圈诱导
	\[
		H^k(\operatorname{Hom}_{\mathbb R}(C_*,\mathbb R))
		\cong\operatorname{Hom}_{\mathbb R}(H_k(C_*),\mathbb R).
	\]
	每个同调上的线性泛函先给出圈上的泛函，再经选取基延拓到 \(C_k\)，因其消去边界而成为余圈；反过来，一个余圈若在所有圈上为零，就经 \(\partial:C_k\to B_{k-1}\) 因子化，把该泛函从 \(B_{k-1}\) 延拓到 \(C_{k-1}\) 后，它便是余边界。这证明该识别，且求值公式说明识别自然。将已证同调同构代入，便得实上同调同构。光滑映射把光滑单形送到光滑单形，与包含及求值交换，故自然性成立。
\end{proof}

\begin{corollary}[光滑小链与开集对]\label{g2:cor:smooth-small-chains}
	设 \(M\) 为有限维、Hausdorff、第二可数的无边界光滑流形，\(\mathcal U\) 为其开覆盖。由像落在某个 \(U\in\mathcal U\) 内的光滑单形生成的子复形 \(C_*^{\infty,\mathcal U}(M;R)\) 包含于 \(C_*^\infty(M;R)\)，并诱导同调同构。若 \(U\subset M\) 开，则光滑相对链与连续相对链也计算同一个 \(H_*(M,U;R)\)；实系数相对上同调亦然。
\end{corollary}
\begin{proof}
	\cref{g2:prop:subdivision-chain-map,g2:prop:subdivision-chain-homotopy}中的细分及链同伦，只在参数单形内使用仿射映射与仿射锥。因此，原单形光滑时，产生的每个单形仍光滑，而且像始终留在原单形的像中。有限光滑链经足够多次细分成为小链，细分链同伦给出满射；若小圈在全部光滑链中有填充，把这个填充细分为小链，再利用原小圈与其细分之间仍属小链的链同伦，便得到小链填充。这正是\cref{g2:thm:small-chain-isomorphism}的证明，其中所有链现在都保留光滑性。

	开子集 \(U\) 也是无边界光滑流形。比较 \((M,U)\) 的两条相对同调长正合列，绝对项由上一定理同构，五引理给出相对同构；实上同调再用证明中的线性泛函识别。这里要求 \(U\) 开，是为了让光滑单形在 \(U\) 中的邻域延拓与在 \(M\) 中的定义相符。
\end{proof}

\subsection{积分映射与 Stokes}
\label{g2:subsec:0703:02}
以下 \(M\) 为有限维、Hausdorff、第二可数、无边界的光滑流形，形式取实系数。写 \(C_\infty^k(M;\mathbb R)=\operatorname{Hom}_{\mathbb R}(C_k^\infty(M;\mathbb R),\mathbb R)\)，上边界仍为 \((\delta\varphi)(c)=\varphi(\partial c)\)。一个形式给上链赋值的办法已经很自然：把它在每条光滑单形上积分。

\begin{definition}[积分上链]\label{g2:def:integration-cochain}
	对 \(\omega\in\Omega^k(M)\)，定义
	\[
		I_M\omega(\sigma)=\int_{\Delta^k}\sigma^*\omega.
	\]
	标准单形采用有序顶点 \((0,e_1,\ldots,e_k)\) 及定向 \(\mathrm{d}t^1\wedge\cdots\wedge\mathrm{d}t^k\)；零次积分是函数在点上的值。此定义按线性延拓到链上。
\end{definition}

这里不要求 \(\omega\) 紧支，因为每条积分单形的定义域紧。第一卷定理14.4.5的带角 Stokes 公式\footnote{George Gabriel Stokes（1819--1903），生于爱尔兰的英国数学家、物理学家，研究流体力学与光学。}适用于闭单形：每个面带上边界定向，恰好产生奇异边界中的交替符号。因此对 \((k+1)\) 单形 \(\sigma\)，
\begin{equation}\label{g2:eq:integration-cochain-map}
	I_M(\mathrm{d}\omega)(\sigma)
	=\int_{\Delta^{k+1}}\mathrm{d}(\sigma^*\omega)
	=\sum_{j=0}^{k+1}(-1)^j\int_{\Delta^k}(\sigma\circ d_j)^*\omega
	=(\delta I_M\omega)(\sigma).
\end{equation}
第一卷定理14.5.5已证明这一链级等式；现在\cref{g2:thm:smooth-singular-comparison}使它能够取值于通常的奇异上同调。具体说，光滑链包含诱导上链限制 \(j^*:C^*(M;\mathbb R)\to C_\infty^*(M;\mathbb R)\)，它在上同调上为同构，故积分定义
\[
	\mathcal I_M:H_{\mathrm{dR}}^k(M)\longrightarrow H^k(M;\mathbb R),
	\qquad \mathcal I_M=(H^k(j^*))^{-1}\circ H^k(I_M).
\]
\symidx{\(\mathcal I_M\)：由光滑单形积分定义的 de Rham 比较映射}
上链限制在实数域上保持同调同构，可直接用万有系数定理的域版本说明：向量空间复形可选取边界、同调及其补空间，取线性对偶后仅留下同调的对偶；不需要这些空间有限维。

\begin{proposition}[积分的自然性]\label{g2:prop:integration-naturality}
	若 \(f:M\to N\) 光滑，则 \(I_Mf^*=f^*I_N\)。因此 \(\mathcal I\) 与光滑拉回相容，特别与开子集限制相容。若 \(z\) 是光滑 \(k\) 圈，则
	\[
		\langle\mathcal I_M[\omega],[z]\rangle=\int_z\omega.
	\]
\end{proposition}
\begin{proof}
	在任意光滑单形 \(\sigma\) 上，\(\sigma^*f^*\omega=(f\circ\sigma)^*\omega\)，所以两个上链的值相等。光滑链包含也与 \(f\) 相容，故在其诱导同构两侧取逆仍得到自然性。求值公式正是定义。若改变闭形式代表或圈代表，差分别由\eqref{g2:eq:integration-cochain-map}化为对圈的余边界求值或对边界的闭上链求值，均为零。
\end{proof}

例如圆周上写 \(\theta=(x\,\mathrm{d}y-y\,\mathrm{d}x)/(2\pi)\)，其中 \(x^2+y^2=1\)。沿 \(t\mapsto(\cos2\pi t,\sin2\pi t)\) 的积分为 \(1\)，所以 \(\mathcal I[\theta]\) 检测一次正向绕行。这不仅说明 \(\theta\) 闭而非恰当，还固定了它与整数生成元比较时的归一化。

\subsection{好覆盖和局部 Poincaré 引理}
\label{g2:subsec:0703:03}
局部求原形式与整体求原形式的差别，可以在最简单的零次就看见：每个连通小开集上的闭函数是常数，但若选取了几份不同常数，它们仍须在交集上一致才能粘成一个函数。高次形式也有同样的相容问题，只是交集上的差需要用更低次的原形式继续修正。

第一卷定理14.2.2给出了星形开集上的具体算子。若 \(U\subset\mathbb R^n\) 关于 \(0\) 星形、\(k>0\)，则
\[
	(K\omega)_x(v_1,\ldots,v_{k-1})
	=\int_0^1t^{k-1}\omega_{tx}(x,v_1,\ldots,v_{k-1})\,\mathrm{d}t
\]
满足 \(\mathrm{d}K+K\mathrm{d}=\operatorname{id}\)；零次相应公式为 \(K\mathrm{d}f=f-f(0)\)。这里的 \(t^{k-1}\) 来自收缩映射对其余 \(k-1\) 个切向量各乘一次 \(t\)。因此每个正次数闭形式局部有原形式，闭函数局部为常数。以 \(\underline{\mathbb R}\) 表示局部常值实函数层，得到层的正合序列
\begin{equation}\label{g2:eq:de-rham-resolution}
	0\longrightarrow\underline{\mathbb R}\longrightarrow\Omega^0
	\xrightarrow{\mathrm{d}}\Omega^1\xrightarrow{\mathrm{d}}\cdots
	\xrightarrow{\mathrm{d}}\Omega^n\longrightarrow0.
\end{equation}
正合性是在茎上检查的：给定一点的闭形式芽，先缩小到坐标星形邻域再应用上式，并不要求原形式在最初的整个开集存在。每个 \(\Omega^q\) 是光滑函数层的模，单位分解的乘法使它成为细层；\cref{g2:thm:acyclic-resolution}遂给出
\begin{equation}\label{g2:eq:de-rham-sheaf-comparison}
	H^k(M;\underline{\mathbb R})\cong H_{\mathrm{dR}}^k(M).
\end{equation}
这一式比较了层与形式，接下来还需将它与奇异上链的积分映射对应起来。

\begin{definition}[光滑好覆盖]\label{g2:def:good-cover}
	光滑流形上的开覆盖 \(\mathfrak U=(U_i)\) 若满足：每个非空有限交 \(U_{i_0\cdots i_p}\) 都能在自身内光滑收缩到一点，则称为\term[guanghuahaofugai]{光滑好覆盖}[smooth good cover]。本节用形式作收缩时，所说的好覆盖均指这一版本，它比只要求连续可缩的定义多了光滑性条件。
\end{definition}

对这种交集，第一卷定理14.2.1的同伦公式给出正次数 de Rham 上同调为零，零次为 \(\mathbb R\)；光滑棱柱同伦给光滑奇异链相同的结论。因而局部比较只剩常数函数在点上取值，积分在零次正是恒等。

\ProseParagraphBreak

\(\mathbb R^n\) 的任意开集都有可数的这种覆盖：取闭包仍在该开集内、端点为有理数的全部开矩形。有限交若非空仍为凸开矩形，线段收缩留在其中。圆周也可取三个略长于三分之一周的开弧，使每两个相交于一段开弧、三重交为空。这份覆盖的 Čech 复形是三角形的三条边而没有内部面，故 \(\check H^1\cong\mathbb R\)。若只用两个长开弧覆盖圆周，其交有两段，常值层在交集上有两份独立的常数，这份覆盖不满足本定义。

\ProseParagraphBreak

下面先在有好覆盖的区域证明比较，再处理一般流形。一般情形用坐标开集及其交，不要求这些交都可缩；交集可以有洞，只需已经知道其积分比较为同构。这种安排把局部计算与覆盖的几何形状分开。

\subsection{双复形的两种收缩}
\label{g2:subsec:0703:04}
固定可数开覆盖 \(\mathfrak U=(U_i)\)。用上一节的交替指标约定，考虑两个第一象限双复形
\[
	A^{p,q}=\prod_{i_0<\cdots<i_p}\Omega^q(U_{i_0\cdots i_p}),
	\qquad B^{p,q}=\prod_{i_0<\cdots<i_p}C_\infty^q(U_{i_0\cdots i_p};\mathbb R).
\]
横向微分是覆盖的交替限制 \(\delta_{\mathrm C}\)，纵向分别是 \(\mathrm d\) 与奇异上边界 \(\delta_{\mathrm s}\)。两种原始微分交换；总微分必须写成
\[
	\mathrm D_A=\delta_{\mathrm C}+(-1)^p\mathrm d,
	\qquad \mathrm D_B=\delta_{\mathrm C}+(-1)^p\delta_{\mathrm s}
	\quad\text{在双次数 }(p,q)\text{ 上}.
\]
这样两项交叉合成相消，\(\mathrm D^2=0\)。总次数 \(k\) 只含 \((0,k),\ldots,(k,0)\) 这些有限项；无限覆盖的乘积位于每一项内部，不会产生沿总对角线求无限和的问题。

\begin{lemma}[两种增广行的正合性]\label{g2:lem:de-rham-augmented-rows}
	形式的横行增广为 \(\Omega^q(M)\to A^{0,q}\to A^{1,q}\to\cdots\)。奇异上链的横行增广为
	\[
		C_{\infty,\mathfrak U}^q(M;\mathbb R)\longrightarrow B^{0,q}\longrightarrow B^{1,q}\longrightarrow\cdots,
	\]
	其中 \(C_{\infty,\mathfrak U}^q\) 是由像完全落在某个 \(U_i\) 内的光滑单形生成的小链复形的线性对偶。两条增广行均正合，且两种增广均在总上同调上给出同构。
\end{lemma}
\begin{proof}
	对形式选从属于 \(\mathfrak U\) 的光滑单位分解 \((\rho_j)\)。第一卷定理12.3.5保证乘法支集留在相应开集内，且求和局部有限。对交替指标上链，令
	\[
		(ha)_{i_0\cdots i_{p-1}}=\sum_j\rho_j a_{j i_0\cdots i_{p-1}}.
	\]
	每项先在交集定义，再利用 \(\operatorname{supp}\rho_j\subset U_j\) 向其余部分延为零。在边界附近它原本就为零，故延拓仍光滑。展开交替和时，删去 \(j\) 的项为 \((\sum_j\rho_j)a=a\)，其余项两两抵消，得到 \(\delta_{\mathrm C}h+h\delta_{\mathrm C}=\operatorname{id}\)，包括增广位置。

	对奇异上链，不需要给上链乘单位分解。为每条覆盖小单形 \(\sigma\) 选一个指标 \(j(\sigma)\)，使其整个像落在 \(U_{j(\sigma)}\)。定义
	\[
		(ha)_{i_0\cdots i_{p-1}}(\sigma)
		=a_{j(\sigma)i_0\cdots i_{p-1}}(\sigma).
	\]
	只在 \(\sigma\) 完全落在相应交集时求值，故右端有定义。对固定 \(\sigma\)，仍是插入一个首指标的同一交替和消去，得到横向收缩。零次增广所需的上链也按这个逐单形公式定义，不需任何连续性。

	这两个横向算子一般不与纵向微分交换；我们没有把它们直接称为总复形的收缩。\cref{g2:lem:first-quadrant-comparison}按有限总对角线逐项消去多余分量，正好把横行正合性转化为两个增广的总上同调同构。
\end{proof}

小链与全部光滑链也计算相同的上同调。\cref{g2:prop:subdivision-chain-map,g2:prop:subdivision-chain-homotopy,g2:thm:small-chain-isomorphism}的构造只把单形与仿射细分、仿射锥和棱柱复合，因而保持光滑性。每条有限链细分充分多次后落入覆盖小链；圈及填充链都可如此处理。故小链包含诱导同调同构，实数域上的线性对偶再给出上同调同构。

\begin{theorem}[de Rham 定理]\label{g2:thm:de-rham}
	设 \(M\) 为有限维、Hausdorff、第二可数、无边界的光滑流形。积分比较
	\[
		\mathcal I_M:H_{\mathrm{dR}}^k(M)\xrightarrow{\ \cong\ }H^k(M;\mathbb R)
	\]
	在每个次数为同构，并与光滑映射的拉回相容。它与\eqref{g2:eq:de-rham-sheaf-comparison}一道识别形式、常值实数层与奇异上同调；覆盖极限的 Čech 上同调也自然等于这些群。对好覆盖，单份覆盖的常值 Čech 复形已经计算它们，比较由常数包含和积分给出。
\end{theorem}
\begin{proof}
	证明的目标不只是找一个抽象同构，而是核对刚才的积分映射。逐个交集上的积分组成双复形映射 \(I:A^{\bullet,\bullet}\to B^{\bullet,\bullet}\)，因为它与限制及两种纵向微分相容。两条增广在积分下也交换，得到
	\[
		\begin{tikzcd}[column sep=large]
			\Omega^\bullet(M) \arrow[r] \arrow[d,"I"] & \operatorname{Tot} A \arrow[d,"I"] \\
			C_{\infty,\mathfrak U}^\bullet(M;\mathbb R) \arrow[r] & \operatorname{Tot} B .
		\end{tikzcd}
	\]
	上下两个横箭头已在上同调上为同构。

	若 \(\mathfrak U\) 是好覆盖，各交集上的纵向上同调只在零次为 \(\mathbb R\)。积分把常数函数送到同值的零上圈，因此每一列都诱导同构。第一象限比较引理给出中间两个总复形的同构，再由上图得到积分同构。两种双复形还都接受常值 Čech 复形的包含，这两个包含诱导同构且与积分交换，证明所述 Čech 识别。层的识别与此相容，因为常值层到形式消解的增广正是这里使用的常数包含。

	特别地，\(\mathbb R^n\) 的每个开子集都由前述有理开矩形覆盖，故已经满足结论。现在取一般 \(M\) 的可数坐标开覆盖。任一非空有限交包含于其中一个坐标域，在该坐标下是 \(\mathbb R^n\) 的开子集；上一段已给出它的积分同构。于是逐列积分仍诱导同构，虽然这些列的正次上同调未必消失。再用同一双复形比较及增广图，得到一般 \(M\) 的结论。光滑映射的自然性由\cref{g2:prop:integration-naturality}给出。流形满足 Hausdorff、第二可数、局部紧条件，故\cref{g2:thm:cech-sheaf-comparison}还将覆盖极限的 Čech 上同调识别为常值层上同调；这些比较都由同一个常数增广诱导，因而相容。
\end{proof}

在好覆盖上，可以把这段证明理解为两种计算次序：先用单位分解消去覆盖方向，留下整体形式；或者先在可缩交集上求原形式，留下常数及其 Čech 相容条件。一般坐标覆盖则需要保留交集本身的上同调。积分构造与从 Euclid 开集推广到一般流形的思路可对照\cite[第18章，Smooth Singular Homology 与 The de Rham Theorem]{Lee2012Smooth}。

\subsection{乘法和支集版本的比较}
\label{g2:subsec:0703:05}
楔积与杯积的比较要经过乘积空间：先在乘积链上积分，再沿对角映射拉回，就能证明积分同构保持上同调乘法。在单条单形上，两种公式通常不同：楔积在整个单形上积分，Alexander--Whitney 杯积则在前面、后面分别取值。下面的证明将这种链级差别与上同调上的相容区分开来。

\begin{proposition}[乘法的相容性]\label{g2:prop:de-rham-products}
	在\cref{g2:thm:de-rham}的条件下，\(\mathcal I_M\) 保持单位和乘法：
	\[
		\mathcal I_M([\omega]\wedge[\eta])
		=\mathcal I_M[\omega]\smile\mathcal I_M[\eta].
	\]
	对两个这样的流形，积分也把形式的外积 \(\pi_M^*\omega\wedge\pi_N^*\eta\) 送到上同调外积。
\end{proposition}
\begin{proof}
	设 \(\deg\omega=p\)、\(\deg\eta=q\)，两形式闭。\cref{g2:thm:kunneth}在实数域上没有 \(\operatorname{Tor}\) 项，说明 \(H_{p+q}(M\times N;\mathbb R)\) 由各次数的乘积圈类生成。\cref{g2:thm:uct-cohomology}的域版本又说明，一个上同调类由它在全部同调类上的值确定。故只需在 \([z]\times[w]\) 上比较两类，其中 \(z,w\) 可由光滑比较取为光滑圈。

	若 \(\deg z=p\)、\(\deg w=q\)，第4章的乘积链是把有向 \(\Delta^p\times\Delta^q\) 剖成 shuffle 单形后逐项求和。内部公共面相消，定向取第一因子在前。逐片积分再用有限可加性及连续函数在紧乘积区域上的迭代积分公式，得到
	\[
		\int_{z\times w}\pi_M^*\omega\wedge\pi_N^*\eta
		=\left(\int_z\omega\right)\left(\int_w\eta\right).
	\]
	若 \(\deg z=a\)、\(\deg w=b\)、\(a+b=p+q\) 而 \((a,b)\ne(p,q)\)，则 \(a<p\) 或 \(b<q\)，对应因子上的拉回形式为零。上同调外积的值也分别由\cref{g2:prop:cross-cup-evaluation}给出相同的乘积或零。因此外积比较成立。

	最后取 \(N=M\)，沿光滑对角映射 \(\Delta(x)=(x,x)\) 拉回。形式外积拉回为楔积，奇异上同调外积拉回为杯积，自然性给出结论。零次常函数 \(1\) 的积分是常值零上圈 \(1\)，所以单位也保持。这里使用的是同调的域上 Künneth 同构，没有要求上同调外积的张量映射在无限维情形仍为同构。
\end{proof}

紧支比较还要控制原形式的支集：一个闭紧支形式即使有普通原形式，该原形式也可能在无穷远处保持非零。因此，要在紧支复形中判断恰当性，需同时构造紧支的原形式。定义
\[
	\Omega_c^k(M)=\{\omega\in\Omega^k(M):\operatorname{supp}\omega\text{ 紧}\},
	\qquad H_{c,\mathrm{dR}}^k(M)=H^k(\Omega_c^\bullet(M),\mathrm d).
\]
若 \(\omega\) 支于紧集 \(K\)，则 \(I\omega\) 在完全位于 \(M\setminus K\) 的单形上为零，因而是相对上链。第6章的定义给出自然映射
\[
	\mathcal I_{c,M}:H_{c,\mathrm{dR}}^k(M)\longrightarrow
	H_c^k(M;\mathbb R)=\varinjlim_K H^k(M,M\setminus K;\mathbb R).
\]

\begin{theorem}[相对及紧支比较]\label{g2:thm:de-rham-support}
	设 \(M\) 满足\cref{g2:thm:de-rham}的条件，\(V\subset M\) 开。复形
	\[
		A^k(M,V)=\Omega^k(M)\oplus\Omega^{k-1}(V),\qquad
		\mathrm D(\omega,\eta)=(\mathrm d\omega,\omega|_V-\mathrm d\eta)
	\]
	约定负次形式为零，其上同调经积分自然同构于 \(H^k(M,V;\mathbb R)\)。此外，\(\mathcal I_{c,M}\) 在每个次数为同构；它与忘记支集、开嵌入的扩零及固有光滑映射的拉回相容，并把紧支楔积配对送到相应杯积配对。
\end{theorem}
\begin{proof}
	先解释相对复形。一个圈由闭形式 \(\omega\) 及 \(V\) 上满足 \(\omega|_V=\mathrm d\eta\) 的原形式组成；相对条件记录的是这一份原形式数据。直接计算 \(\mathrm D^2=0\)。由投影到 \(\Omega^\bullet(M)\) 得到的短正合列，产生与空间对完全同形的长正合列。

	把形式换成光滑奇异上链，得到相同的复形。上链限制 \(C_\infty^*(M)\to C_\infty^*(V)\) 逐次满射，因为可给其余单形补零；因此相对上链的核包含 \(\varphi\mapsto(\varphi,0)\) 诱导此复形的上同调同构。也可逐项验证：对一个圈 \((a,b)\)，延拓 \(b\) 为 \(\widetilde b\)，减去余边界 \(\mathrm D(\widetilde b,0)\) 后，第二项为零，第一项落在限制核内；相同操作处理边界。积分与两条长正合列交换，对 \(M,V\) 用普通 de Rham 定理及五引理，得到相对比较。光滑与连续相对模型的比较则由绝对比较和对的长正合列得到。

	现在令 \(V=M\setminus K\)，并使紧集 \(K\) 增大。证明自然包含
	\[
		\Omega_c^\bullet(M)\longrightarrow\varinjlim_K A^\bullet(M,M\setminus K),
		\qquad\omega\longmapsto(\omega,0)
	\]
	在上同调上为同构。给定右侧圈 \((\omega,\eta)\)，取光滑函数 \(\rho\)，使其在 \(K\) 的一个邻域为零，在某个更大紧集 \(L\) 外为 \(1\)。第一卷推论12.3.6的闭集截断可作此选择。把 \(\rho\eta\) 延为 \(M\) 上的形式 \(\widetilde\eta\)；它在 \(K\) 附近为零，所以延拓光滑。在更大紧集 \(L\) 对应的相对复形中，减去 \(\mathrm D(\widetilde\eta,0)\) 后得到
	\[
		(\omega-\mathrm d\widetilde\eta,0),
	\]
	且第一项在 \(M\setminus L\) 为零。这证明上同调满射。零次没有 \(\eta\)，同一结论直接由 \(\omega|_{M\setminus K}=0\) 得到。

	为证单射，设紧支闭形式 \(\omega\) 在极限中为边界。在足够大紧集 \(K\) 的阶段，有 \(\omega=\mathrm d\beta\)、\(\beta|_{M\setminus K}=\mathrm d\gamma\)。用相同截断延拓 \(\rho\gamma\)，则 \(\beta-\mathrm d(\rho\gamma)\) 是紧支形式，其外微分仍为 \(\omega\)。当 \(\deg\omega=1\) 时 \(\gamma\) 的次数为负，条件已经说 \(\beta\) 在补集为零；零次没有非零边界。故各次均得单射。向量空间的滤过直接极限保持正合：圈、边界及两个代表相等都由有限个等式见证，故在某个更晚阶段已经成立。于是可交换这里的上同调与直接极限，结合相对比较得到 \(\mathcal I_{c,M}\) 同构。

	忘记支集与固有拉回在积分链级相容；固有性保证紧集的逆像仍紧。开嵌入扩零则在包含支集的相对对上由切除实现，与形式的扩零相同。乘法证明也保留支集：若第一形式支于 \(K\)，就在空间对 \((M,M\setminus K)\) 上作外积比较。乘积链映射及其同伦保持补集子链，故第4章的域上 Künneth 论证下降到相对商；上述迭代积分计算仍成立。沿对角拉回后，支集回到 \(K\)。若两形式都紧支，所得支集包含在两支集的交中，取直接极限便得到紧支杯积配对。
\end{proof}

取 \(f\in C_c^\infty(\mathbb R)\)、\(\int_{\mathbb R}f(x)\,\mathrm dx=1\)。形式 \(f(x)\,\mathrm dx\) 有普通原函数 \(F(x)=\int_{-\infty}^x f(t)\,\mathrm dt\)，但 \(F\) 在正向远端为 \(1\)，不可能紧支；任何其他原函数只差常数，也不能在两端都为零。因此该形式在 \(H_{c,\mathrm{dR}}^1(\mathbb R)\) 非零，在普通 \(H_{\mathrm{dR}}^1(\mathbb R)\) 却为零。这正对应第6章 \(H_c^1(\mathbb R;\mathbb R)=\mathbb R\) 与 \(H^1(\mathbb R;\mathbb R)=0\) 的区别。

% !TeX root = ../../Book2.tex
% 纲要标题：### 7.4 de Rham 上同调的计算与局部系数
% 纲要位置：几何学计划纲要.md，第 1290 行。
% 正文状态：初稿；尚未正式审核。

% 纲要内容开始
% 
% * 周期与整上同调
% * 环面、球面及射影空间计算
% * 局部系统值形式的平坦联络模型
% 
% ---
% 
% 纲要内容结束

\section{de Rham 上同调的计算与局部系数}
\label{g2:sec:0704}

积分比较把一个闭形式变成对圈的测量，但它只给出实系数信息。一个周期是整数，究竟是对形式本身的限制，还是对它所代表的类的限制？如果某个整数类有有限阶，它还能由一个非零闭形式表示吗？先澄清这些问题，再把第4章的环计算落实为具体微分形式。最后，我们把实数换成沿道路运输的向量，说明第5章的局部系统怎样进入微分方程。

本节的流形均为有限维、Hausdorff、第二可数的无边界光滑流形；紧性只在明确提出时使用。积分所用链为上一节的光滑奇异链，系数为整数或实数。由其与连续奇异链的比较，它们计算同样的同调群。

\subsection{周期与整上同调}
\label{g2:subsec:0704:01}
设 \(\omega\in\Omega^k(M)\) 满足 \(\mathrm d\omega=0\)。对光滑整数 \(k\) 圈 \(z=\sum_i n_i\sigma_i\)，定义其\term[zhouqi]{周期}[period]为
\[
	\operatorname{Per}_{\omega}([z])=\int_z\omega
	\coloneqq\sum_i n_i\int_{\Delta^k}\sigma_i^*\omega.
\]
若 \(z-z'=\partial c\)，Stokes 公式给出 \(\int_{z-z'}\omega=\int_c\mathrm d\omega=0\)；若以 \(\omega+\mathrm d\eta\) 替换 \(\omega\)，在圈上的积分也不变。因此周期定义了群同态 \(H_k(M;\mathbb Z)\to\mathbb R\)，且只依赖 \([\omega]\in H^k_{\mathrm{dR}}(M)\)。当 \(k=0\) 时，闭零形式是局部常函数，它在点所代表的零圈上的周期就是函数值。

整数链上的积分仍取实数值：例如在 \(\mathbb R/\mathbb Z\) 上，\(a\,\mathrm dt\) 的一周积分是任意指定的 \(a\in\mathbb R\)。另一方面，积分在挠类上恒为零：若 \(r[z]=0\)，则 \(r\int_z\omega=0\)，故 \(\int_z\omega=0\)。整周期条件所限制的是这些实数取值是否全部属于 \(\mathbb Z\)。

\begin{proposition}[整周期判据]\label{g2:prop:integral-periods}
	设 \(M\) 为无边界光滑流形，\(k\geq0\)，并且 \(H_k(M;\mathbb Z)\)、\(H_{k-1}(M;\mathbb Z)\) 有限生成，其中 \(H_{-1}=0\)。对闭形式 \(\omega\in\Omega^k(M)\)，下列条件等价：
	\begin{enumerate}
		\item 对每个光滑整数 \(k\) 圈 \(z\)，都有 \(\int_z\omega\in\mathbb Z\)。
		\item 在\cref{g2:thm:de-rham}的积分同构下，\([\omega]\) 属于系数映射 \(H^k(M;\mathbb Z)\to H^k(M;\mathbb R)\) 的像。
	\end{enumerate}
	该像同构于 \(H^k(M;\mathbb Z)/\operatorname{Tor}H^k(M;\mathbb Z)\)，在 \(H^k(M;\mathbb R)\) 中是一个满格。特别地，此判据适用于紧光滑流形。
\end{proposition}
\begin{proof}
	先恢复万有系数定理在这里的作用。\cref{g2:thm:uct-cohomology}的评价映射给出交换图
	\[
		\begin{tikzcd}[column sep=large]
			H^k(M;\mathbb Z) \arrow[r,two heads,"\rho_{\mathbb Z}"] \arrow[d] & \operatorname{Hom}(H_k(M;\mathbb Z),\mathbb Z) \arrow[d] \\
			H^k(M;\mathbb R) \arrow[r,"\rho_{\mathbb R}"'] & \operatorname{Hom}(H_k(M;\mathbb Z),\mathbb R).
		\end{tikzcd}
	\]
	下方箭头是同构：有限生成群分解为自由群与有限循环群的直和，而 \(\operatorname{Ext}^1(\mathbb Z,\mathbb R)=0\)、\(\operatorname{Ext}^1(\mathbb Z/r,\mathbb R)=\mathbb R/r\mathbb R=0\)，后一个公式来自\cref{g2:prop:cyclic-ext-tor}。因此实上同调类完全由其在整数圈上的取值确定。

	积分比较将 \(\rho_{\mathbb R}[\omega]\) 识别为刚定义的周期同态。如果它的值全部为整数，就得到一个同态 \(H_k(M;\mathbb Z)\to\mathbb Z\)；上方箭头满射，故这个同态提升为整数上同调类，交换性和下方箭头的单射性说明其实像正是 \([\omega]\)。反过来，整数上同调类在整数圈上当然取整数值，交换图给出所需结论。

	写 \(H_k(M;\mathbb Z)\cong\mathbb Z^b\oplus T\)，其中 \(T\) 有限。右上角是 \(\mathbb Z^b\)，右下角是 \(\mathbb R^b\)，右侧映射是通常的整数格包含。左上角评价的核为 \(\operatorname{Ext}^1(H_{k-1}(M;\mathbb Z),\mathbb Z)\)，由同样的有限生成分解可知它是有限群；评价的目标无挠，所以这个核恰是全部挠子群。由此得到最后的断言。紧流形的整数同调有限生成已由\cref{g2:lem:compact-manifold-finite-homology}证明。
\end{proof}

具有上述性质的闭形式称为具有\term[zhengzhouqi]{整周期}[integral periods]。这个条件既不要求形式在坐标中的系数为整数，也不指定唯一的整数上同调提升。后者可以相差一个挠类，而所有周期仍完全相同。实上同调是向量空间，整周期类只是其中的格；任意实倍数通常会离开这个格。

还有一个较弱的条件应单独辨认：\(\omega\) 为闭形式只说明积分对同调不变，并不说明任何周期为零。恰当形式的周期全部为零；在上述有限生成条件下，反方向由积分同构及评价同构成立，即所有周期为零恰好等价于 \(\omega\) 恰当。对 \(S^1=\mathbb R/\mathbb Z\)，\(\mathrm dt\) 闭且有整周期，但其一周积分为 \(1\)，所以不是恰当形式。

\subsection{环面、球面及射影空间计算}
\label{g2:subsec:0704:02}
在 \(T^m=\mathbb R^m/\mathbb Z^m\) 上，坐标函数 \(t_i\) 不下降为实值函数，其微分却下降为整体一形式 \(\theta_i\)：整数平移改变函数一个常数，不改变微分。令 \(\gamma_j\) 为正向的第 \(j\) 条坐标圆周，则
\[
	\int_{\gamma_j}\theta_i=\delta_{ij}.
\]
所以积分同构将 \([\theta_i]\) 送到\cref{g2:prop:torus-cohomology-ring}中的第 \(i\) 个一次生成元的实像。该同构保持乘法，因而
\begin{equation}\label{g2:eq:torus-de-rham-ring}
	H^*_{\mathrm{dR}}(T^m)
	=\Lambda_{\mathbb R}([\theta_1],\ldots,[\theta_m]).
\end{equation}
具体说，每个闭 \(k\) 形式都与唯一一个常系数组合 \(\sum_{i_1<\cdots<i_k}a_{i_1\cdots i_k}\theta_{i_1}\wedge\cdots\wedge\theta_{i_k}\) 相差恰当形式；这些系数就是它在相应坐标子环面上的积分。因而该类具有整周期，当且仅当所有 \(a_{i_1\cdots i_k}\) 为整数。这些常系数代表由积分比较和环面的杯积计算共同确定。

对 \(n\geq1\)，在单位球面 \(S^n\subset\mathbb R^{n+1}\) 上取向外定向，并令
\[
	\nu=\frac{1}{\operatorname{vol}(S^n)}
	\left.\sum_{j=0}^n(-1)^j x_j\,\mathrm dx_0\wedge\cdots\wedge\widehat{\mathrm dx_j}\wedge\cdots\wedge\mathrm dx_n\right|_{S^n}.
\]
分母是通常球面面积，故 \(\int_{S^n}\nu=1\)。顶次形式自动闭；它不是恰当形式，否则 Stokes 公式会使这个积分为零。球面的同调计算与积分比较给出
\[
	H^k_{\mathrm{dR}}(S^n)=
	\begin{cases}
		\mathbb R[1],&k=0,\\
		\mathbb R[\nu],&k=n,\\
		0,&\text{其余次数}.
	\end{cases}
\]
全部正次乘积为零。任意顶次形式 \(\omega\) 都有 \([\omega]=(\int_{S^n}\omega)[\nu]\)；这给出一个无需显式求原函数的恰当性判据。

复射影空间需要一个二形式，而不只是顶次体积形式。下面从单位球面构造它，使归一化可以直接检查。记 \(q:S^{2m+1}\to\mathbb CP^m\) 为 \(z\mapsto[z]\)，其中 \(m\geq1\)、\(z_j=x_j+\mathrm i y_j\)，并在球面上取
\[
	\alpha=\frac{1}{2\pi}\sum_{j=0}^m(x_j\,\mathrm dy_j-y_j\,\mathrm dx_j).
\]
圆周作用为 \(z\mapsto e^{\mathrm i\vartheta}z\)，正向参数 \(0\leq\vartheta\leq2\pi\) 的纤维上有 \(\alpha=\mathrm d\vartheta/(2\pi)\)。在仿射图中，把某个非零齐次坐标选为正实数再归一化，得到光滑局部截面 \(s\)。两条这样的截面在足够小的交叠上满足 \(s'=e^{\mathrm i f}s\)，其中 \(f\) 是实值光滑函数；直接代入给出
\[
	(s')^*\alpha=s^*\alpha+\frac{\mathrm df}{2\pi}.
\]
因此 \(\mathrm d(s^*\alpha)\) 在交叠上一致，粘成 \(\mathbb CP^m\) 上的闭二形式 \(\omega\)。这也说明 \(q^*\omega=\mathrm d\alpha\)：在局部乘积坐标中，\(\alpha\) 等于截面上的形式加上纤维角度的微分，而后者的外微分为零。\(\mathrm d\alpha\) 在球面上恰当，并不意味着下降后的 \(\omega\) 也恰当；原始的一形式 \(\alpha\) 在纤维方向不为零，不能下降到射影空间。

在坐标直线 \(\mathbb CP^1\) 的图 \([1:z]\)、\(z=x+\mathrm i y\) 中，可取 \(s(z)=(1,z)/\sqrt{1+|z|^2}\)。于是
\[
	s^*\alpha=\frac{x\,\mathrm dy-y\,\mathrm dx}{2\pi(1+x^2+y^2)},
	\qquad \omega=\frac{\mathrm dx\wedge\mathrm dy}{\pi(1+x^2+y^2)^2}.
\]
取复坐标给出的定向，利用极坐标积分得到
\[
	\int_{\mathbb CP^1}\omega
	=\frac1\pi\int_0^{2\pi}\int_0^\infty\frac{r}{(1+r^2)^2}\,\mathrm dr\,\mathrm d\vartheta=1.
\]
被省略的点没有面积，且形式在另一张图上光滑，故这里的广义积分等于紧流形上的积分。\cref{g2:thm:complex-projective-ring}说明 \(H^2(\mathbb CP^m;\mathbb Z)\to H^2(\mathbb CP^1;\mathbb Z)\) 是同构，所以 \([\omega]\) 是在正向坐标直线上取值 \(1\) 的整数生成元的实像。乘法相容性遂给出
\begin{equation}\label{g2:eq:projective-de-rham-ring}
	H^*_{\mathrm{dR}}(\mathbb CP^m)
	\cong\mathbb R[u]/(u^{m+1}),\qquad u=[\omega],\quad |u|=2.
\end{equation}
特别地，\(1,\omega,\ldots,\omega^m\) 的类分别给出全部偶数次的一维上同调，奇数次为零。非零的楔幂在这里由第4章已经证明的整数杯积计算保证。

实射影空间则清楚显示实系数丢失了什么。\cref{g2:prop:projective-integer-homology}与万有系数定理给出，当 \(m\geq1\) 时
\[
	H^k(\mathbb RP^m;\mathbb Z)=
	\begin{cases}
		\mathbb Z,&k=0,\text{ 或 }k=m\text{ 为奇数},\\
		\mathbb Z/2,&k\text{ 为正偶数且 }k\leq m,\\
		0,&\text{其余次数}.
	\end{cases}
\]
这里偶数次的 \(\mathbb Z/2\) 来自前一次同调的 Ext 项。换成实系数以后，这些项全部消失，故
\[
	H^k_{\mathrm{dR}}(\mathbb RP^m)=
	\begin{cases}
		\mathbb R,&k=0,\text{ 或 }k=m\text{ 为奇数},\\
		0,&\text{其余次数}.
	\end{cases}
\]
例如 \(H^2(\mathbb RP^2;\mathbb Z)=\mathbb Z/2\)，却没有非零的二次 de Rham 类。若 \(m\) 为奇数，球面对径映射保持定向，归一化球面体积形式 \(\nu\) 下降为 \(\bar\nu\)，且二重覆盖给出 \(\int_{\mathbb RP^m}\bar\nu=1/2\)；因此具有整周期的正向顶次生成元由 \(2\bar\nu\) 表示。当 \(m\) 为偶数时，对径映射反转定向，这个体积形式不能作为普通实值形式下降。

\subsection{局部系统值形式的平坦联络模型}
\label{g2:subsec:0704:03}
第5章的局部系统沿道路比较不同位置的系数。若希望用微分形式作同样的计算，首先要知道怎样微分一个取值于不同纤维的对象。设 \(E\to M\) 是有限秩实光滑向量丛，取联络 \(\nabla\)，并记
\[
	\Omega^k(M;E)=\Gamma(\Lambda^k\mathrm T^*M\otimes E).
\]
沿用第一卷第15章的行标架 \(e=(e_1,\ldots,e_r)\) 与列分量约定：若 \(\nabla e=eA\)，则 \(\nabla(eu)=e(\mathrm du+Au)\)。形式上的\term[waixiebianweifen]{外协变微分}[exterior covariant derivative]定义为
\begin{equation}\label{g2:eq:flat-valued-differential}
	\mathrm d_\nabla(e\eta)=e(\mathrm d\eta+A\wedge\eta).
\end{equation}
换标架 \(e'=eg\) 时，\(A'=g^{-1}Ag+g^{-1}\mathrm dg\)、\(\eta'=g^{-1}\eta\)，代入可知右侧表示同一个整体形式。因此这一定义不依赖标架。展开两次求导并使用 \(\mathrm d(A\wedge\eta)=\mathrm dA\wedge\eta-A\wedge\mathrm d\eta\)，得到
\begin{equation}\label{g2:eq:flat-curvature-square}
	\mathrm d_\nabla^2(e\eta)=e(F\wedge\eta),\qquad
	F=\mathrm dA+A\wedge A.
\end{equation}
这正是第一卷定理15.3.5的曲率公式。只有曲率为零的\term[pingtanlianluo]{平坦联络}[flat connection]才能给出所需上链复形；要求平方在零形式上都为零，就已经强迫 \(F=0\)。

平行标架使局部联络矩阵为零，于是向量丛值形式的外协变微分就化为逐分量的普通外微分。第一卷定理15.5.2已经证明，\(F=0\) 等价于每点附近存在这样的标架。在坐标盒中写 \(A=\sum_i A_i\,\mathrm dx^i\)，其构造先解带参数的线性常微分方程
\[
	\partial_1g=-A_1g,\qquad g(0,x^2,\ldots,x^n)=I.
\]
解对参数光滑；同时解 \(\partial_1h=hA_1\)、\(h(0,x')=I\)，可得 \(hg=I\)，故换标架合法。换后 \(B_1=0\)，而零曲率给出 \(0=F_{1j}'=\partial_1B_j\)，所以其余系数只依赖 \(x^2,\ldots,x^n\)。在这些变量上归纳消去剩余系数，就得到 \(A=0\) 的标架。这里用到的参数常微分方程定理及光滑依赖已在第一卷第13.1节证明；每一步可缩小坐标盒，以取得共同定义域。反过来，平行标架中 \(A=0\) 显然使曲率为零。

两个平行标架满足 \(e'=eg\) 时，\(0=g^{-1}\mathrm dg\)，所以转移矩阵局部常值。在这样的标架中保持向量分量不变，便能沿道路运输。把一个道路同伦的参数方形细分到每个小格都落在一个平行标架域，格子内两条边路的运输相同，公共边逐一抵消，证明运输只依赖定端点同伦类；这就是第一卷定理15.5.4的论证。因此平行截面组成有限维实\term[jubu xitong]{局部系统}[local system] \(\mathcal L\)。其纤维可识别为 \(E_x\)，运输即平行移动。反过来，把有限秩实局部系统的局部常转移矩阵用于粘合 \(U_i\times\mathbb R^r\)，便得到光滑向量丛；因为转移矩阵的微分为零，各片的普通微分粘成平坦联络。这两个构造在相同局部平行标架中互相恢复，也保持纤维线性且与运输相容的映射。

\begin{theorem}[局部系统的 de Rham 比较]\label{g2:thm:twisted-de-rham}
	设 \(M\) 为无边界光滑流形，\(E\to M\) 为有限秩实光滑向量丛，\(\nabla\) 为平坦联络，\(\mathcal L\) 为其平行截面局部系统。则运输后的积分诱导自然同构
	\[
		H^k\bigl(\Omega^*(M;E),\mathrm d_\nabla\bigr)
		\cong H^k(M;\mathcal L).
	\]
	对常值系统及其平凡平坦丛，此同构就是\cref{g2:thm:de-rham}的积分同构。
\end{theorem}
\begin{proof}
	先说明积分取值在哪里。对光滑单形 \(\sigma:\Delta^k\to M\)，将 \(\sigma^*E\) 的各纤维沿单形内道路运输到首顶点 \(v_0\)。单形可缩，平坦运输使这个比较与道路无关。于是 \(\sigma^*\eta\) 成为取值于固定向量空间 \(E_{\sigma(v_0)}\) 的形式，可逐分量积分；记所得元素为 \(I_\nabla\eta(\sigma)\)。这种描述没有选择基，因而积分与平行标架的改变相容。

	在首顶点平行标架中，\(\mathrm d_\nabla\) 变成普通外微分，Stokes 公式逐分量成立。边界中去掉 \(v_0\) 的面以 \(v_1\) 为首顶点，必须先把它的积分沿边 \([v_0,v_1]\) 反向运回 \(v_0\)；其余各面首顶点仍是 \(v_0\)。因此
	\begin{align*}
		I_\nabla(\mathrm d_\nabla\eta)(\sigma)
		&=T_{01}^{-1}I_\nabla\eta(\sigma|[v_1,\ldots,v_{k+1}])\\
		&\quad+\sum_{i=1}^{k+1}(-1)^iI_\nabla\eta(\sigma|[v_0,\ldots,\widehat v_i,\ldots,v_{k+1}]).
	\end{align*}
	这恰好是\cref{g2:def:twisted-cochains}的余边界公式，故积分是上链映射。上一节的相对光滑逼近比较也适用于这个局部系统：每个单形同伦的运输把新旧首顶点纤维对应起来，棱柱公式中的公共面仍以相反符号和相同运输抵消。因此光滑局部系数上链与连续局部系数上链给出同一上同调。

	在平行标架坐标小球内，形式复形在正次数的局部恰当性逐分量成立，零次核就是 \(\mathcal L\)。因此层复形
	\[
		0\longrightarrow\mathcal L\longrightarrow\Omega^0(E)
		\xrightarrow{\mathrm d_\nabla}\Omega^1(E)
		\xrightarrow{\mathrm d_\nabla}\cdots
	\]
	正合。形式层可以乘光滑单位分解，故由本章的细层无上同调结论及\cref{g2:thm:acyclic-resolution}，其整体截面计算 \(\mathcal L\) 的层上同调。还须把这项计算与奇异上同调通过指定的积分联系起来。

	取由平行标架坐标域 \(U_i\) 组成的开覆盖。每个非空有限交 \(U_{i_0\cdots i_p}\) 是第一张坐标图中的 Euclid 开集，且第一张图的平行标架在交集上仍然平凡化 \(\mathcal L\)。该交集未必可缩，但\cref{g2:thm:de-rham}已经对任意 Euclid 开集证明普通积分同构，逐分量应用便得到这个交集上带 \(\mathcal L\) 系数的积分同构。现在将上一节的覆盖双复形换成 \(E\) 值形式与 \(\mathcal L\) 值小单形上链。形式增广行仍由光滑单位分解收缩；上链增广行仍可对每个小单形选择一个容纳它的覆盖指标作收缩，限制不改变其首顶点纤维，因而\cref{g2:lem:de-rham-augmented-rows}的证明保持有效。刚才在有限交上的积分同构说明每列的比较是同构，\cref{g2:lem:first-quadrant-comparison}遂给出总复形同构，再沿两个增广映射得到所述整体积分同构。

	光滑映射的拉回与单形上的运输、积分都交换；保持联络的丛映射在平行标架中是常矩阵，也与积分交换。这给出所述自然性。
\end{proof}

\begin{example}[圆周上的平坦系数]\label{g2:exm:flat-circle-de-rham}
	设 \(E\to S^1=\mathbb R/\mathbb Z\) 的实际正向环路运输为 \(T\in\operatorname{GL}(V)\)。在区间 \([0,1]\) 上沿路径平行平凡化，截面的列函数满足 \(f(1)=T^{-1}f(0)\)，而不是 \(f(1)=Tf(0)\)：终点的平行标架本身已相对于起点乘了 \(T\)。拉回到实线后，形式写作 \(f(t)\) 或 \(g(t)\,\mathrm dt\)，并满足所有导数的相应准周期条件。

	此时微分为普通导数。零次核是满足 \(T^{-1}v=v\) 的常函数，故为 \(\ker(T-1)\)。对一次形式，令 \(b=\int_0^1g(t)\,\mathrm dt\)。若它是 \(\mathrm df\)，则
	\[
		b=f(1)-f(0)=(T^{-1}-1)f(0).
	\]
	反过来，若 \(b=(T^{-1}-1)v\)，取 \(f(t)=v+\int_0^t g(s)\,\mathrm ds\)。端点值满足准周期关系；一阶以上导数的关系来自 \(g\)，所以 \(f\) 下降为光滑截面。任给 \(b\in V\)，取支集位于 \((0,1)\)、积分为 \(1\) 的光滑函数 \(\chi\)，令 \(g=\chi b\) 并作准周期延拓，就实现这个积分。于是
	\[
		H^0_{\mathrm d_\nabla}(S^1;E)=\ker(T-1),\qquad
		H^1_{\mathrm d_\nabla}(S^1;E)=V/(T^{-1}-1)V=V/(T-1)V.
	\]
	这与\cref{g2:prop:circle-local-coefficients}的胞腔计算相合，并解释了一次余边界为什么含逆运输。若 \(V=\mathbb R\)、\(T=-1\)，两个群均为零；但整数符号系统的一次上同调为 \(\mathbb Z/2\)。把系数几何化为实平坦丛依然不能恢复已经被实数消掉的挠信息。
\end{example}

第5章把纤维的实同调或上同调组织为基底上的局部系统；当基底光滑且这些纤维向量空间有限维时，刚才的构造将它们实现为平坦丛。于是沿基底变化的纤维上同调类可以写成平坦丛截面，\(\mathrm d_\nabla\) 测量它是否局部保持平行。这里要求有限维，且只得到所给局部系统的微分形式模型；一个局部系统仍不包含原几何族的全部粘合信息，这与\secref{g2:sec:0504}的区别一致。
% 纲要条目排为编号小节标题，并进入目录与书签；说明段仍保留于纲要注释中。


\begin{chaptersummary}
	层把开集上的截面及其限制组织在一起；层化允许局部可表示的数据粘合，茎则只保留一点附近的信息。核、余核与正合性可以逐茎检测，整体截面却只自动保持左正合。逐点茎乘积的典范消解把整体提升失败定义为高次层上同调；松弛性允许延拓，细性允许带支集地加权粘合，两者各自给出消失机制。

	\ProseParagraphBreak

	Čech 微分交替删去指标，记录局部数据在各级交叠上的相容关系。细化指标的选择通过链同伦消失。无上同调覆盖可直接计算层上同调；在 Hausdorff、第二可数、局部紧空间中，即使不能一次找到这样的覆盖，也可以通过细化逐次消掉固定总次数中的局部障碍，得到覆盖极限的比较同构。

	\ProseParagraphBreak

	光滑逼近必须保持公共面，否则链边界的抵消会丢失。紧像附近的欧氏模型、角面延拓与相对平滑给出光滑链比较。Stokes 公式随后使积分成为上链映射；覆盖双复形一边通过单位分解恢复整体形式，另一边通过局部积分比较接到奇异上同调。楔积与杯积在上链上一般不相等，但乘积积分及对角拉回证明其在上同调上相容。

	\ProseParagraphBreak

	相对形式复形同时保留一个闭形式及其在子空间上的原形式；紧支比较还须通过截断消去远处的原形式。实形式的整周期可检测整数上同调在实上同调中的像，却不能区分其中的挠类。平坦联络的局部平行标架把局部系统变为向量丛值形式，曲率为零正是外协变微分平方为零的条件；圆周上的准周期函数又把这一模型还原为第5章的单值核与余核计算。
\end{chaptersummary}

\BookExercises

\begin{exercise}\label{g2:ex:skyscraper-sheaf}
	设 \(X\) 为 Hausdorff 空间，\(x\in X\)，\(A\) 为 \(R\) 模。定义 \(\mathcal S_x(A)(U)=A\)（\(x\in U\)）及 \(0\)（\(x\notin U\)），限制在两端均含 \(x\) 时取恒等，其余取零。证明它是层，计算各点的茎及全部层上同调。进一步证明，对任意 \(R\) 模层 \(\mathcal F\)，有自然同构
	\[
		\operatorname{Hom}(\mathcal F,\mathcal S_x(A))\cong\operatorname{Hom}_R(\mathcal F_x,A).
	\]
\end{exercise}
\begin{solution}
	若覆盖的开集之并不含 \(x\)，所有截面都是零。若并含 \(x\)，至少一个覆盖成员含 \(x\)，而所有含 \(x\) 的成员两两交也含 \(x\)，所以相容族必须给出同一个 \(a\in A\)。它唯一粘成并集上的 \(a\)，故这是层。在 \(x\) 的每个邻域上模都为 \(A\)，限制为恒等，故茎为 \(A\)。若 \(y\ne x\)，Hausdorff 性给出不含 \(x\) 的 \(y\) 邻域，其上的模为零，故 \(y\) 处的茎为零。
	所有限制均满射，因此该层松弛。由\cref{g2:thm:flabby-fine-acyclic}，\(H^0(X,\mathcal S_x(A))=A\)，而 \(H^q(X,\mathcal S_x(A))=0\)（\(q>0\)）。
	给定层态射 \(f\)，它在含 \(x\) 的开集上的值与限制相容，遂定义 \(\mathcal F_x\to A\)，\(s_x\mapsto f_U(s)\)。若代表改变，在某个共同小邻域上相等，所得值不变。反之，给定 \(u:\mathcal F_x\to A\)，在含 \(x\) 的开集上设 \(f_U(s)=u(s_x)\)，在其余开集上设为零；它显然与限制交换。这两个构造互逆，并与两边的模或层态射相容，故自然。这样的层称为集中于 \(x\) 的摩天楼层。
\end{solution}

\begin{exercise}\label{g2:ex:exponential-sheaf-circle}
	在 \(S^1\) 上令 \(\mathcal C_{\mathbb R}\) 为连续实函数层，\(\mathcal C_{S^1}\) 为连续圆值函数层，后者采用逐点乘法。证明
	\[
		0\longrightarrow\underline{\mathbb Z}\longrightarrow\mathcal C_{\mathbb R}\xrightarrow{\ f\mapsto e^{2\pi\mathrm i f}\ }\mathcal C_{S^1}\longrightarrow0
	\]
	为阿贝尔群层短正合列。由此计算 \(H^1(S^1,\underline{\mathbb Z})\)，并说明连接同态如何记录连续映射 \(g:S^1\to S^1\) 的绕数。
\end{exercise}
\begin{solution}
	指数映射的核是连续整数值函数，也就是局部常值整数函数。对任意 \(g\) 和任意 \(x\)，可缩小 \(x\) 的邻域，使 \(g\) 的像落在圆的一个真开弧内；这个开弧上有连续角度函数，故指数映射在茎上满射。注意这没有断言每个整体圆值函数都有整体角度。
	连续单位分解把 \(\mathcal C_{\mathbb R}\) 变成细层：端同态为乘以连续权函数。由第一卷第4章的单位分解定理与\cref{g2:thm:flabby-fine-acyclic}，它在圆上无上同调。因此长正合列给出
	\[
		H^1(S^1,\underline{\mathbb Z})\cong
		 C(S^1,S^1)/\{e^{2\pi\mathrm i f}:f\in C(S^1,\mathbb R)\}.
	\]
	把圆参数化为 \(t\mapsto e^{2\pi\mathrm it}\)，\(0\le t\le1\)。第一卷第5章的道路提升定理对覆盖 \(\mathbb R\to S^1\) 给出 \(g(e^{2\pi\mathrm it})=e^{2\pi\mathrm i\theta(t)}\) 的连续提升。端点投影相同，故 \(k=\theta(1)-\theta(0)\in\mathbb Z\)。改变初值只使整个提升加同一个整数，故 \(k\) 不变；相乘的提升可相加，所以绕数是群同态。
	若 \(g=e^{2\pi\mathrm if}\) 有整体角度，则取 \(\theta(t)=f(e^{2\pi\mathrm it})\)，两端相等，绕数为零。反之绕数为零时 \(\theta(0)=\theta(1)\)，所以 \(\theta\) 下降为圆上的连续函数 \(f\)。每个整数 \(k\) 都由 \(g(z)=z^k\) 实现，故上述商群恰为 \(\mathbb Z\)。以该识别作约定，连接类 \(\partial g\) 对应绕数 \(k\)；零类恰好等价于整体角度存在。
\end{solution}

\begin{exercise}\label{g2:ex:twisted-cech-circle}
	设 \(A\) 为 \(R\) 模，\(T:A\to A\) 为自同构。在圆的两弧覆盖 \(U,V\) 上给定局部常值层：两个交集分支的过渡映射分别为 \(1\) 与 \(T\)。在共同坐标中把 Čech 微分写成 \(\delta(a,b)=(b-a,b-Ta)\)。计算该覆盖的零次、一次 Čech 上同调，并分别讨论 \(A=\mathbb Z,T=-1\) 和 \(A=\mathbb R,T=-1\)。
\end{exercise}
\begin{solution}
	零次核满足 \(b=a=Ta\)，所以
	\[
		\check H^0\cong\ker(T-1),\qquad
		\check H^1\cong A/(T-1)A.
	\]
	第二个同构可由 \((u,v)\mapsto[v-u]\) 给出：微分的像映到 \([(1-T)a]=0\)。它满射，因为 \((0,w)\) 映到 \([w]\)。若 \(v-u=(1-T)a\)，取 \(b=u+a\)，则 \(b-a=u\)、\(b-Ta=u+(1-T)a=v\)，故该映射的核恰是微分的像。不存在三重交的项，所以高次 Čech 群为零。
	当 \(A=\mathbb Z,T=-1\) 时，不变元只有零，余不变元为 \(\mathbb Z/2\)。当 \(A=\mathbb R,T=-1\) 时，乘 \(2\) 可逆，零次与一次群都为零。相同的反号运输可以在整数系数下留下挠类，在实系数下完全消失；这不是由覆盖形状单独决定的。本题计算的是指定覆盖复形，认同于层上同调时还须核验有限交上的无上同调条件。
\end{solution}

\begin{exercise}\label{g2:ex:smooth-jet-cokernel}
	设 \(m\ge1\)，在 \(\mathbb R\) 的光滑函数层 \(\mathcal C^{\infty}\) 上考虑乘法态射 \(f\mapsto x^m f\)。证明其核为零，余核同构于集中在 \(0\) 的摩天楼层 \(\mathcal S_0(\mathbb R^m)\)，其中商映射记录 \((f(0),f'(0),\ldots,f^{(m-1)}(0))\)。解释为什么这个余核并不是一个在每点都有 \(m\) 维茎的常值层。
\end{exercise}
\begin{solution}
	若 \(x^m f(x)=0\)，则在所有非零点 \(f=0\)，连续性又给出 \(f(0)=0\)，故核层为零。定义层态射 \(J:\mathcal C^{\infty}\to\mathcal S_0(\mathbb R^m)\)：含 \(0\) 的开集上取所列导数，不含 \(0\) 时取零。任意给定的 \(m\) 个导数值都由一个次数小于 \(m\) 的多项式实现，故 \(J\) 在茎上满射。
	在非零点附近，\(x^m\) 可逆，像层就是全部光滑函数层。在 \(0\) 附近，若 \(f^{(j)}(0)=0\)（\(0\le j<m\)），Taylor 公式的积分余项给出
	\[
		f(x)=x^m g(x),\qquad
		g(x)=\frac1{(m-1)!}\int_0^1(1-t)^{m-1}f^{(m)}(tx)\,\mathrm dt.
	\]
	积分区间固定且被积函数光滑，故 \(g\) 在 \(0\) 的一个邻域光滑。因此 \(\ker J\) 的每个茎都等于乘 \(x^m\) 的像，\cref{g2:prop:stalk-exactness}便给出所需短正合列与余核同构。余核在 \(0\) 处的茎为 \(\mathbb R^m\)，在其余点为零；常值层 \(\underline{\mathbb R^m}\) 的每点茎都为 \(\mathbb R^m\)，所以两者不同。这个例子把“局部无法除以 \(x^m\)”精确压缩为有限阶导数资料。
\end{solution}

\begin{exercise}\label{g2:ex:cech-cycle-refinement}
	设圆有一个 \(N\ge3\) 个开弧组成的覆盖，只有相邻两弧相交，所有非空交连通，没有三重交。对常值层 \(\underline R\)，按绕圆的方向给各相邻交定向，把一次上链写成 \((c_0,\ldots,c_{N-1})\)。证明一次 Čech 类由 \(\sum_i c_i\) 唯一确定。再设一个 \(2N\) 弧覆盖细化它，指标选择按模 \(N\) 记为 \(r(2i)=i\)、\(r(2i+1)=i+1\)；计算细化对上述和的影响，并说明反转圆的方向产生什么变化。
\end{exercise}
\begin{solution}
	零次上链为 \((b_i)\in R^N\)，一次微分在定向的相邻交上为 \((\delta b)_i=b_{i+1}-b_i\)，指标按模 \(N\) 取值。这仍是交错 Čech 约定，只是最后一条边采用从 \(N-1\) 指向 \(0\) 的方向，而不是指标递增方向。所有余边界的分量和为零。反之若 \(\sum_i c_i=0\)，取 \(b_0=0\)，并依次令 \(b_{i+1}=b_i+c_i\)；总和为零保证最后回到 \(b_N=b_0\)。于是 \(c=\delta b\)。故求和诱导 \(\check H^1\cong R\)，同时零次闭元为共同常数。
	在细化覆盖上，从 \(2i\) 到 \(2i+1\) 的边映到旧覆盖的 \(i\to i+1\) 边，故取值为 \(c_i\)；从 \(2i+1\) 到 \(2i+2\) 的边，两端旧指标同为 \(i+1\)，按重复指标约定取值为零。因此细化后的分量为 \((c_0,0,c_1,0,\ldots,c_{N-1},0)\)，总和不变，细化诱导此 \(R\) 上的恒等映射。若反转绕行方向，各边取负且次序反转，总和变号。这里识别一次群所用的方向是额外选择，而群及细化映射本身并未依赖该选择。
\end{solution}

\begin{exercise}\label{g2:ex:transverse-face-extension}
	在坐标轴的并集上给定 \(b(x,0)=x^2\)、\(b(0,y)=y^3\)，写出一个光滑延拓。再考虑直线 \(L=\{(x,0)\}\) 与抛物线 \(P=\{(x,x^2)\}\)，在其并集上给定 \(c(x,0)=0\)、\(c(x,x^2)=x\)。说明 \(c\) 连续且在每条曲线上光滑，却不能延拓为原点附近的 \(C^1\) 函数。解释这与\cref{g2:lem:simplex-face-extension}并不矛盾。
\end{exercise}
\begin{solution}
	第一组数据由 \(B(x,y)=x^2+y^3\) 延拓。两条坐标轴相交于原点，减去被重复计算的公共值后，正好得到\cref{g2:eq:corner-face-extension}的二维形式。

	第二组数据在两曲线的唯一交点都等于零。若并集中的点趋于原点，它的横坐标趋于零，两种公式的值都趋于零；在其他点附近只有一条曲线需要考虑，所以 \(c\) 连续。参数 \(x\mapsto(x,0)\) 与 \(x\mapsto(x,x^2)\) 都是光滑嵌入，沿它们的限制分别为 \(0\) 和 \(x\)，故逐曲线光滑。假设有 \(C^1\) 延拓 \(C\)，其 \(\partial_yC\) 在原点附近有界，于是
	\[
		x=C(x,x^2)-C(x,0)=\int_0^{x^2}\partial_yC(x,t)\,\mathrm{d}t=O(x^2),
	\]
	这在 \(x\to0\)、\(x\ne0\) 时不可能。单形的相邻面可以同时写成不同坐标超平面，法向线性独立；这里的两条曲线却相切，不能用同一坐标变换化为两条坐标轴。因此，逐片光滑的数据能否延拓还取决于这些片的相交方式。
\end{solution}

\begin{exercise}\label{g2:ex:close-circle-loops}
	把 \(S^1\) 看作复平面内的单位圆。设光滑闭路 \(\gamma,\eta:[0,1]\to S^1\) 满足 \(\gamma(0)=\gamma(1)=\eta(0)=\eta(1)=1\)，并且 \(\sup_t|\gamma(t)-\eta(t)|<2\)。构造保持端点的逐棱柱光滑同伦，并证明两闭路的绕数相同。再说明常数 \(2\) 不能放宽为允许等号。
\end{exercise}
\begin{solution}
	两个单位复数之间的线段经过原点，当且仅当它们互为相反数；此时两者距离恰为 \(2\)。所以题设使
	\[
		H(t,s)=\frac{(1-s)\gamma(t)+s\eta(t)}{|(1-s)\gamma(t)+s\eta(t)|}
	\]
	处处有定义。分母在紧正方形上有正下界，因而这个公式在正方形附近有光滑延拓，尤其在任意棱柱剖分的各单形上光滑；\(t=0,1\) 时恒为 \(1\)，故端点固定。取单位圆上的闭形式 \(\alpha=(x\,\mathrm{d}y-y\,\mathrm{d}x)/(2\pi)\)。给正方形取 \(\mathrm{d}s\wedge\mathrm{d}t\) 定向，剖分为两个有向三角形并用 Stokes 公式，内部边抵消，两个端点侧边为常值，得到
	\[
		\int_\eta\alpha-\int_\gamma\alpha=\int_{[0,1]^2}H^*(\mathrm{d}\alpha)=0.
	\]
	这两个积分就是绕数。若取 \(\gamma(t)=e^{2\pi\mathrm{i}t}\)、\(\eta(t)=1\)，最大距离为 \(2\)，而绕数分别为 \(1\) 与 \(0\)。它们不可能相对端点同伦：第一卷定理5.1.13证明，以 \(1\) 为基点的圆周环路的定端点同伦类由绕数完全分类。这个例子说明，欧氏空间中的线性插值必须留在回缩的定义域内；仅仅写出插值公式还没有得到目标流形内的同伦。
\end{solution}

\begin{exercise}\label{g2:ex:integration-not-strict-product}
	在 \(\mathbb R^2\) 上取 \(\omega=\mathrm dx\)、\(\eta=\mathrm dy\)，并令 \(\sigma:\Delta^2\to\mathbb R^2\) 是顶点为 \((0,0),(1,0),(0,1)\) 的标准仿射单形。计算 \(I(\omega\wedge\eta)(\sigma)\) 与 \((I\omega\smile I\eta)(\sigma)\)。再构造一个一次光滑奇异上链 \(T\)，使 \(I\omega\smile I\eta-I(\omega\wedge\eta)=\delta T\)，并说明这与正文的乘法定理相容。
\end{exercise}
\begin{solution}
	三角形有向面积为 \(1/2\)，故第一值为 \(1/2\)。前后面杯积给出 \(\int_{01}\mathrm dx\cdot\int_{12}\mathrm dy=1\cdot1=1\)。这说明积分不是严格的微分分次代数同态。

	令 \(a=I(x)\) 为零上链，即在点上取其 \(x\) 坐标，定义
	\[
		T=a\smile I(\mathrm dy)-I(x\,\mathrm dy).
	\]
	积分与微分相容，杯积满足 Leibniz 公式，且 \(\delta I(\mathrm dy)=0\)，所以
	\[
		\delta T=I(\mathrm dx)\smile I(\mathrm dy)-I(\mathrm dx\wedge\mathrm dy).
	\]
	具体地，对一条光滑道路 \(\tau\)，\(T(\tau)=x(\tau(0))(y(\tau(1))-y(\tau(0)))-\int_\tau x\,\mathrm dy\)。在边 \(01,12,02\) 上依次为 \(0,1/2,0\)，故 \((\delta T)(\sigma)=1/2\)，与前述差相符。两种上链不相等，却相差余边界，因而上同调乘法仍相同。
\end{solution}

\begin{exercise}\label{g2:ex:relative-cone-line}
	取 \(M=\mathbb R\)、\(V=\mathbb R\setminus\{0\}\)。在相对形式复形 \(A^*(M,V)\) 中，令 \(\eta\) 在负半轴为 \(0\)、在正半轴为 \(1\)，考虑一次圈 \((0,\eta)\)。证明其类非零，求它在从 \(-1\) 到 \(1\) 的有向相对区间类上的值，并用紧支形式表示它在 \(\varinjlim_KH^1(A^*(\mathbb R,\mathbb R\setminus K))\) 中的像。
\end{exercise}
\begin{solution}
	\(\eta\) 在 \(V\) 上局部常值，故 \(\mathrm D(0,\eta)=(0,-\mathrm d\eta)=0\)。若它是边界，则存在光滑函数 \(f\) 满足 \(\mathrm df=0\)、\(f|_V=\eta\)。前者迫使 \(f\) 在连通实轴上为一个常数，后者却要求两个不同常数，矛盾。

	把 \(\eta\) 的零上链任意延拓为实轴上的零上链 \(b\)。相对比较将 \((0,I\eta)\) 与 \((-\delta b,0)\) 识别，所以在区间上的值是 \(-b(1)+b(-1)=-1\)。负号来自本章相对复形的约定 \(\mathrm D(\omega,\eta)=(\mathrm d\omega,\omega|-\mathrm d\eta)\)，不能只凭两端值之差猜它的顺序。

	取光滑函数 \(\chi\)，在 \(( -\infty,-1]\) 为 \(0\)，在 \([1,\infty)\) 为 \(1\)。当紧集扩大到包含 \([-1,1]\) 时，\(\chi|=\eta\)，故减去 \(\mathrm D(\chi,0)\) 把原圈变为 \((-\mathrm d\chi,0)\)。其第一项紧支，且 \(\int_{\mathbb R}-\mathrm d\chi=-1\)，表示所求非零紧支类。
\end{solution}

\begin{exercise}\label{g2:ex:proper-pullback-compact-forms}
	用积分将 \(H_{c,\mathrm{dR}}^1(\mathbb R)\) 识别为 \(\mathbb R\)。计算固有光滑映射 \(f(x)=x^2\)、\(g(x)=-x\) 在该群上的拉回。再说明覆盖映射 \(p:\mathbb R\to S^1\)、\(p(t)=(\cos2\pi t,\sin2\pi t)\) 为什么不能以通常拉回来定义从全部紧支形式到紧支形式的映射。
\end{exercise}
\begin{solution}
	紧支一次形式写成 \(h(y)\,\mathrm dy\)。若积分为零，其原函数 \(F(y)=\int_{-\infty}^yh(s)\,\mathrm ds\) 在两端均为零，故紧支；任意实数又可由紧支函数的积分实现。因此积分确实给出所述识别。

	\(f^*(h\,\mathrm dy)=2x h(x^2)\,\mathrm dx\)，这是紧支奇函数系数的一次形式，其全轴积分为零，故 \(f^*=0\)。还可直接取 \(G(y)=-\int_y^\infty h(s)\,\mathrm ds\)，则 \(G(x^2)\) 在 \(|x|\) 足够大时为零，并满足 \(\mathrm d(G(x^2))=2x h(x^2)\,\mathrm dx\)。对 \(g\)，拉回为 \(-h(-x)\,\mathrm dx\)，其积分为 \(-\int h\)，故 \(g^*\) 是乘以 \(-1\)。

	圆周紧，所以其每个形式都紧支，但归一化角形式 \(\theta=(x\,\mathrm dy-y\,\mathrm dx)/(2\pi)\) 的拉回为 \(\mathrm dt\)，支集是整条实轴。\(p\) 的确不是固有映射：紧集 \(S^1\) 的逆像不紧。因此普通 de Rham 拉回存在，紧支拉回却不在所述复形之间定义。
\end{solution}

\begin{exercise}\label{g2:ex:integral-one-form-circle-map}
	设 \(M\) 为有限维、Hausdorff、第二可数的连通无边界光滑流形，\(\omega\in\Omega^1(M)\) 为闭实一形式。记 \(\vartheta\) 为 \(\mathbb R/\mathbb Z\) 上由 \(\mathrm dt\) 下降的一形式。证明：\(\omega\) 在每条分片光滑闭路上的积分均为整数，当且仅当存在光滑映射 \(f:M\to\mathbb R/\mathbb Z\)，使 \(\omega=f^*\vartheta\)。说明在固定一点的 \(f\) 值后，这个映射是否唯一。
\end{exercise}
\begin{solution}
	先设周期为整数。固定 \(p\in M\) 及 \(c\in\mathbb R/\mathbb Z\)，对从 \(p\) 到 \(x\) 的分片光滑道路 \(\gamma\)，定义
	\[
		f(x)=c+\left[\int_\gamma\omega\right]\pmod{\mathbb Z}.
	\]
	连通光滑流形可由坐标球连接任意两点，故这样的道路存在。两条道路的积分差为闭路 \(\gamma*\bar\gamma'\) 的积分，按假设为整数，因此 \(f\) 良定义。在小坐标球 \(U\) 上，Poincaré 引理给出 \(\omega|_U=\mathrm dh\)。取 \(x_0\in U\)，并把到 \(x\) 的道路分成固定的 \(p\) 到 \(x_0\) 一段及 \(U\) 内一段，便有 \(f(x)=f(x_0)+[h(x)-h(x_0)]\)。所以 \(f\) 局部具有光滑实值提升，并且 \(f^*\vartheta=\mathrm dh=\omega\)。

	反过来，若 \(\omega=f^*\vartheta\)，对一条闭路 \(\gamma\)，将 \(f\circ\gamma\) 提升到实线，记提升为 \(\widetilde\gamma\)。端点之差是整数，且
	\[
		\int_\gamma\omega=\int_{f\circ\gamma}\vartheta
		=\widetilde\gamma(1)-\widetilde\gamma(0)\in\mathbb Z.
	\]
	若 \(f_1,f_2\) 的拉回均为 \(\omega\)，则群值差 \(f_1-f_2\) 的每个局部实提升都具有零微分，故局部常值。\(M\) 连通使它整体为常值；固定 \(p\) 处的值后，差为零，因而唯一。这个题目表明整周期的一形式可能没有实值原函数，却有圆周值原函数。
\end{solution}

\begin{exercise}\label{g2:ex:dense-period-group}
	在 \(T^2=\mathbb R^2/\mathbb Z^2\) 上取 \(\omega=\mathrm dx+\sqrt2\,\mathrm dy\)。求其周期群 \(\operatorname{Per}_{\omega}(H_1(T^2;\mathbb Z))\)，证明它在 \(\mathbb R\) 中稠密，并证明任何非零实数倍 \(a\omega\) 都不具有整周期。
\end{exercise}
\begin{solution}
	两条坐标圆周生成一次整数同调，周期分别为 \(1,\sqrt2\)，所以周期群为 \(G=\mathbb Z+\sqrt2\mathbb Z\)。要证稠密，给定整数 \(N\geq1\)，考察 \(0,\sqrt2,\ldots,N\sqrt2\) 的小数部分。把 \([0,1)\) 分成 \(N\) 个等长区间，必有两个小数部分落在同一区间内。因此存在非零整数 \(q\)、整数 \(p\)，使 \(0<|q\sqrt2-p|<1/N\)；严格正性来自 \(\sqrt2\) 无理。适当变号，得到 \(G\) 中一个 \(0<\varepsilon<1/N\)。任给 \(r\in\mathbb R\)，取最接近 \(r/\varepsilon\) 的整数 \(j\)，则 \(|r-j\varepsilon|\leq\varepsilon/2\)。令 \(N\) 任意大，证明 \(G\) 稠密。

	若 \(a\omega\) 具有整周期，两条坐标圆周要求 \(a\in\mathbb Z\) 且 \(a\sqrt2\in\mathbb Z\)。当 \(a\ne0\) 时，两者之比会使 \(\sqrt2\) 为有理数，矛盾。周期群可以稠密，这并不与\cref{g2:prop:integral-periods}的格结论冲突：后者说的是具有整周期的上同调类在整个上同调向量空间中组成一个格，而不是说任一实上同调类的周期值总组成离散群。
\end{solution}

\begin{exercise}\label{g2:ex:projective-power-map-forms}
	对整数 \(d\geq1\)、\(m\geq1\)，定义
	\[
		F_d:\mathbb CP^m\longrightarrow\mathbb CP^m,
		\qquad [z_0:\cdots:z_m]\longmapsto[z_0^d:\cdots:z_m^d].
	\]
	证明此映射良定义且光滑。若 \(\omega\) 是\eqref{g2:eq:projective-de-rham-ring}的归一化闭二形式，求 \([F_d^*\omega]\) 及 \(F_d\) 在最高次上同调上的作用，从而求其映射度。
\end{exercise}
\begin{solution}
	同乘一个非零复数 \(\lambda\) 后，各坐标的 \(d\) 次幂同乘 \(\lambda^d\)，故射影点不变；至少一个输入坐标非零，输出也不会全为零。在第 \(j\) 个输入坐标非零的仿射图上，输出的同一个坐标非零，局部表达为各仿射坐标逐个取 \(d\) 次幂，所以光滑。

	坐标直线 \(\mathbb CP^1\) 在此映射下不变，其仿射表达为 \(z\mapsto z^d\)。值 \(1\) 有恰好 \(d\) 个原像；每个原像处复导数 \(dz^{d-1}\) 非零，相应实 Jacobian 的行列式为 \(|dz^{d-1}|^2>0\)。由\cref{g2:thm:local-degree-sum}，这个球面映射的度为 \(d\)。因此 \(F_d^*\omega\) 在正向坐标直线上的积分为 \(d\)，而 \(\omega\) 的积分为 \(1\)。二次上同调是一维，故
	\[
		[F_d^*\omega]=d[\omega].
	\]
	拉回保持楔积，所以 \(F_d^*[\omega^j]=d^j[\omega^j]\)。特别地，\([\omega^m]\ne0\) 生成顶次实上同调，\(F_d^*\) 在该处为乘 \(d^m\)。用与顶次基本类的自然评价，顶次上同调上的乘数就是映射度，于是 \(\deg F_d=d^m\)。这里仅用了闭形式的类；一般没有 \(F_d^*\omega=d\omega\) 的逐点等式。
\end{solution}

\begin{exercise}\label{g2:ex:curvature-obstructs-complex}
	在 \(\mathbb R^2\) 的平凡实线丛上取联络 \(\nabla f=\mathrm df+x f\,\mathrm dy\)。计算 \(\mathrm d_\nabla^2f\)，判断 \(\operatorname{im}(\nabla:\Omega^0\to\Omega^1)\) 是否包含在 \(\ker(\mathrm d_\nabla:\Omega^1\to\Omega^2)\) 中，并求任一非空连通开集上的全部平行截面。
\end{exercise}
\begin{solution}
	联络矩阵为 \(A=x\,\mathrm dy\)，故 \(F=\mathrm dA+A\wedge A=\mathrm dx\wedge\mathrm dy\)，从而
	\[
		\mathrm d_\nabla^2f=f\,\mathrm dx\wedge\mathrm dy.
	\]
	例如 \(\nabla1=x\,\mathrm dy\)，其外协变微分为 \(\mathrm dx\wedge\mathrm dy\ne0\)，所以像不包含于核，不能用这两个空间的商定义上同调。若 \(\nabla f=0\)，再作用一次 \(\mathrm d_\nabla\) 得 \(f\,\mathrm dx\wedge\mathrm dy=0\)；面积形式处处非零，因此 \(f=0\)。这在每个非空开集上都成立，说明非平坦联络甚至可能没有局部非零平行截面，不能把它的平行截面当成与纤维同秩的局部系统。
\end{solution}

\begin{exercise}\label{g2:ex:flat-torus-contracting-homotopy}
	在 \(T^2=\mathbb R^2/\mathbb Z^2\) 的平凡实线丛上取
	\[
		\mathrm d_\nabla=\mathrm d+(a\,\mathrm dx+b\,\mathrm dy)\wedge,
		\qquad (a,b)\in\mathbb R^2\setminus\{(0,0)\}.
	\]
	证明它给出一个平坦联络，并构造该形式复形的收缩同伦，从而证明所有次数的上同调均为零。
\end{exercise}
\begin{solution}
	一形式 \(A=a\,\mathrm dx+b\,\mathrm dy\) 闭，且在线丛中 \(A\wedge A=0\)，故曲率为零。交换坐标后可设 \(a\ne0\)。在光滑双周期函数上令 \(D_x=\partial_x+a\)、\(D_y=\partial_y+b\)。对给定 \(f\)，方程 \(D_xu=f\) 的周期解唯一，且为
	\[
		u(x,y)=e^{-ax}\left(\frac{\int_0^1e^{at}f(t,y)\,\mathrm dt}{e^a-1}
		+\int_0^x e^{at}f(t,y)\,\mathrm dt\right).
	\]
	由于 \(a\) 为非零实数，分母不为零。公式给出 \(u(1,y)=u(0,y)\)，微分方程及 \(f\) 的周期性进一步给出各阶导数的端点一致性；对 \(y\) 的光滑性和周期性可直接在公式中检查。齐次解 \(ce^{-ax}\) 只有在 \(c=0\) 时周期，所以解唯一。记此逆算子为 \(B=D_x^{-1}\)，并逐分量作用于形式。它与 \(D_x,D_y\) 交换：对 \(D_x\) 由互逆关系成立，对 \(D_y\) 则由公式中参数微分与积分交换成立。

	在形式上有 \(\mathrm d_\nabla=\mathrm dx\wedge D_x+\mathrm dy\wedge D_y\)。定义降一次算子 \(K=B\iota_{\partial_x}\)。收缩算子满足
	\[
		\iota_{\partial_x}(\mathrm dx\wedge\eta)+\mathrm dx\wedge\iota_{\partial_x}\eta=\eta,
		\qquad
		\iota_{\partial_x}(\mathrm dy\wedge\eta)+\mathrm dy\wedge\iota_{\partial_x}\eta=0.
	\]
	结合 \(BD_x=1\) 及所述交换性，得到 \(\mathrm d_\nabla K+K\mathrm d_\nabla=1\)。因此每个闭形式 \(\eta\) 都满足 \(\eta=\mathrm d_\nabla K\eta\)，全部上同调为零。这个平坦丛的底层线丛虽平凡，指定联络沿两条坐标圆周的运输却为 \(e^{-a}\)、\(e^{-b}\)；至少一个不是 \(1\)，所以其微分形式上同调可以与普通环面上同调完全不同。
\end{solution}


% !TeX root = ../Book2.tex
% 纲要标题：# 第二编　同伦、谱序列与逐层障碍
% 纲要位置：几何学计划纲要.md，第 1298 行。
% 正文状态：未撰写；本文件目前是框架。

% 纲要内容开始
% 
% 纲要内容结束

\BookPart{同伦、谱序列与逐层障碍}{g2:part:02}
\BookPartGuide
待撰写

% !TeX root = ../../Book2.tex
% 纲要标题：## 第8章　高阶同伦、纤维化与长正合列
% 纲要位置：几何学计划纲要.md，第 1300 行。
% 正文状态：初稿；尚未正式审核。

% 纲要内容开始
% 
% > 本章目标与证明范围：以提升问题连接高阶同伦和纤维化；同时处理相对群、基本群作用及低阶集合项。覆盖只是离散纤维的特例，不将纤维丛和一般纤维化等同。
% 
% 纲要内容结束

\chapter{高阶同伦、纤维化与长正合列}
\label{g2:ch:08}
\BookChapterNumber{g2:ch:08}

\begin{chapterabstract}
	球面映射能否收缩、底空间中的变形能否提升，是本章的两个基本问题。高阶同伦群以立方体拼接组织球面映射，相对同伦群则允许一个边界面在指定子空间中移动。由此构造空间对的长正合列，并说明基点变化、基本群作用和低阶尖点集的含义。Serre 与 Hurewicz 纤维化以同伦提升性质刻画映射；路径空间、拉回及纤维化替换提供具体模型。纤维化的提升把底空间同伦群识别为总空间的相对同伦群，进而给出自然长正合列。复与四元数 Hopf 丛、经典群的球面丛将这些构造用于低阶计算，并给出同调上为零而在同伦上具有无穷阶的球面映射。
\end{chapterabstract}

\begin{learninggoals}
	\begin{itemize}
		\item 在球面、立方体与空间对模型之间转换，构造拼接运算并解释高阶交换性。
		\item 用道路追踪基点变化，辨认相对同伦序列中的群、尖点集及其作用。
		\item 写出两种同伦提升性质，构造路径纤维化、拉回与映射的纤维化替换。
		\item 由提升的顶面构造连接映射，证明长正合列的正合性、自然性及比较结论。
		\item 建立两种 Hopf 丛与经典群的局部平凡化，计算低阶同伦群，并检测非平凡球面映射。
	\end{itemize}
\end{learninggoals}

在单位三维球面 \(S^3\subset\mathbb C^2\) 上，一个点 \((z_0,z_1)\) 张成一条复直线。忘掉共同相位，就把这个点送到复射影直线 \(\mathbb{CP}^1\cong S^2\)。每条纤维是一个圆周；在小范围内可以连续选取相位，沿整个球面作这种选择却会遇到障碍。这个投影能否连续变成常值映射？同调群暂时不能回答：\(S^3\) 的三次同调为 \(\mathbb Z\)，而 \(S^2\) 的三次同调为零，所以投影在这一次数必定诱导零映射。

\ProseParagraphBreak

要直接研究投影本身的变形，我们把它看成一个三维球面在目标空间中的同伦类。一般地，\(\pi_n(X)\) 保存 \(S^n\to X\) 的带基点同伦类。它能提出同调群未能直接解决的填充问题，但运算与计算也需要新的工具。对上面的投影，如果将底空间中的球面从一面开始向上提升，最后留下的边界会落在圆周纤维中。边界的绕行恰好把二维球面映射与一维环路联系起来，成为长正合列中的连接映射。

\ProseParagraphBreak

这种提升方法依靠对参数和边界的控制。第一节先在立方体的各个面上规定哪些值固定、哪些值可以移动，从而建立相对同伦序列；第二节证明纤维丛和路径空间具有所需的提升能力；第三节把两项构造接起来，逐项解释正合性及低阶分支作用。第四节给出可直接计算的局部平凡化，并用长正合列证明两种 Hopf 映射都有无穷阶。

\ProseParagraphBreak

本章使用第1章的道路与同伦、基本群、覆盖提升及相对 CW 构造。用胞腔逼近证明低维球面到高维球面的映射可缩时，会恢复固定基点这一条件；检测球面恒等映射的整数倍时，会用第3章的映射度及其同伦不变性。矩阵群的例子还用有限维线性代数与正交化，四元数的乘法约定在使用前给出。进一步阅读可参看\cite[第4.1节；第4.2节，Fiber Bundles]{Hatcher2002}。

% !TeX root = ../../Book2.tex
% 纲要标题：### 8.1 高阶与相对同伦群
% 纲要位置：几何学计划纲要.md，第 1254 行。
% 正文状态：未撰写；本文件目前是框架。

% 纲要内容开始
% 
% * 球面、立方体和空间对模型
% * 拼接运算及高阶交换性
% * 基点变化与基本群作用
% * 相对同伦序列的低阶解释
% 
% 纲要内容结束

\section{高阶与相对同伦群}
\label{g2:sec:0801}

待撰写

\subsection{球面、立方体和空间对模型}
\label{g2:subsec:0801:01}
待撰写

\subsection{拼接运算及高阶交换性}
\label{g2:subsec:0801:02}
待撰写

\subsection{基点变化与基本群作用}
\label{g2:subsec:0801:03}
待撰写

\subsection{相对同伦序列的低阶解释}
\label{g2:subsec:0801:04}
待撰写
% 纲要条目排为编号小节标题，并进入目录与书签；说明段仍保留于纲要注释中。

% !TeX root = ../../Book2.tex
% 纲要标题：### 8.2 纤维化和替换构造
% 纲要位置：几何学计划纲要.md，第 1311 行。
% 正文状态：初稿；尚未正式审核。

% 纲要内容开始
% 
% * Serre 与 Hurewicz 提升性质
% * 路径空间和同伦纤维
% * 纤维丛的同伦提升
% * 拉回和映射的纤维化替换
% 
% 纲要内容结束

\section{纤维化和替换构造}
\label{g2:sec:0802}

设一个点在底空间 \(B\) 中移动，而我们已经在它的起点上方选定了 \(E\) 中的一点。覆盖的道路提升给出一种跟随它移动的方法。如果同时移动的是一族点，就需要提升整个同伦。这个要求能够把底空间中的零伦搬到总空间，从而把一个同伦类是否消失的问题转成纤维中的边界问题。纤维化正是用这种提升能力描述映射。

\subsection{Serre 与 Hurewicz 提升性质}
\label{g2:subsec:0802:01}

\begin{definition}[同伦提升与纤维化]\label{g2:def:serre-fibration}
	设 \(p:E\rightarrow B\) 连续，\(Z\) 为拓扑空间。若对任意连续映射 \(a:Z\rightarrow E\) 与 \(H:Z\times I\rightarrow B\)，只要 \(H(z,0)=p(a(z))\)，就存在连续映射 \(\widetilde H:Z\times I\rightarrow E\)，满足
	\[
		p\widetilde H=H,\qquad \widetilde H(z,0)=a(z),
	\]
	则称 \(p\) 对 \(Z\) 具有\term[tongluntishengxingzhi]{同伦提升性质}[homotopy lifting property]。若它对每个闭圆盘 \(D^m\)（\(m\ge0\)）都有此性质，称 \(p\) 为\term[Serrexianweihua]{Serre 纤维化}[Serre fibration]；若它对每个拓扑空间 \(Z\) 都有此性质，称 \(p\) 为\term[Hurewiczxianweihua]{Hurewicz 纤维化}[Hurewicz fibration]。\footnote{Witold Hurewicz（1904--1956），波兰数学家，后在美国工作，系统发展了高阶同伦群，并在维数理论中作出重要贡献。}这里不另外要求 \(p\) 满射。
\end{definition}

提升数据组成下图；虚线同时满足两个三角形中的等式。
\[
	\begin{tikzcd}[column sep=large,row sep=large]
		Z\times\{0\} \arrow[r,"a"] \arrow[d,hook] & E \arrow[d,"p"] \\
		Z\times I \arrow[r,"H"'] \arrow[ur,dashed,"\widetilde H"] & B
	\end{tikzcd}
\]
取 \(Z=D^0\) 得到道路提升，取 \(Z=I\) 得到一族道路的提升。Hurewicz 纤维化当然是 Serre 纤维化。后者只要求闭圆盘上的提升，却足以处理球面映射及其同伦，因为这些问题都由有限个圆盘及其边界组成。

\begin{example}[乘积投影与覆盖]\label{g2:exm:product-cover-fibrations}
	对任意拓扑空间 \(B,F\)，投影 \(p:B\times F\rightarrow B\) 是 Hurewicz 纤维化。把初始提升写成 \(a(z)=(H(z,0),v(z))\)，则
	\[
		\widetilde H(z,t)=(H(z,t),v(z))
	\]
	就是所需提升。覆盖也是 Hurewicz 纤维化：\cref{g2:thm:covering-lifts}已经对任意参数空间证明了连续同伦的提升。覆盖提升由初值唯一决定；乘积投影的提升通常有很多，例如还可以在 \(F\) 内连续移动第二个坐标。
\end{example}

在证明长正合列时，常常还要指定同伦在参数空间边界上的提升。下面的相对形式说明，这些边界条件可以一起保持。

\begin{lemma}[相对 CW 对上的提升]\label{g2:lem:relative-homotopy-lift}
	设 \(p:E\rightarrow B\) 为 Serre 纤维化，\((K,L)\) 为相对 CW 复形。给定连续映射 \(H:K\times I\rightarrow B\)，以及定义在
	\[
		A=(K\times\{0\})\cup(L\times I)
	\]
	上的连续提升 \(a:A\rightarrow E\)，满足 \(pa=H|_A\)，则 \(a\) 可以延拓成 \(H\) 在整个 \(K\times I\) 上的提升。特别地，对一个闭立方体，若已在除某一个面之外的全部边界面上给出相容提升，它就能延拓到整个立方体。
\end{lemma}
\begin{proof}
	先处理一个新胞腔。需要提升的空间是 \(C=D^m\times I\)，已经指定提升的部分为
	\[
		A_m=(D^m\times\{0\})\cup(\partial D^m\times I).
	\]
	空间对 \((C,A_m)\) 同胚于 \((D^m\times I,D^m\times\{0\})\)。可以在圆柱边界上看见这个同胚：沿每条从下底中心经过下底边缘、侧面、上底边缘而到上底中心的折线，保持角方向，把“下底加侧面”所占的参数区间线性缩为下底所占的区间，把余下区间线性伸长。这样得到边界的自同胚，将 \(A_m\) 送到下底；再沿从圆柱中心出发的径向线段延拓，就得到整个圆柱的自同胚。\(m=0\) 时两对原本相同。把给定的底映射与已知提升都通过这个同胚搬运，Serre 提升性质便给出 \(C\) 上的提升，再搬回去。这个论证只用面上的相容性，对角点附近的映射没有额外要求。

	现在从 \(L\times I\) 和整个 \(K\times\{0\}\) 出发，按相对胞腔的维数递增构造。处理一个特征映射 \(\phi:D^m\rightarrow K\) 时，下底的提升由初值给定，侧面的提升已在前一骨架上构造；两者在交集一致。刚才的圆柱结论将它们延到 \(D^m\times I\)。同一维的不同胞腔在边界上都沿用旧提升，因而能粘合。\cref{g2:lem:weak-product}保证，只要在 \(L\times I\) 及每个特征圆盘乘 \(I\) 上连续，所得映射就在 \(K\times I\) 上连续；因此胞腔可以无限多个。

	立方体的一个面可取为上底，其余面正是下底与各侧面的并。圆盘与立方体的同胚把单胞腔结论转为最后的断言；任意指定的面都可通过交换、反转坐标移到上底。
\end{proof}

\subsection{路径空间和同伦纤维}
\label{g2:subsec:0802:02}

路径本身也可以作为空间中的点。赋予 \(C(I,B)\) 紧开拓扑，其子基由“某一紧集 \(K\subset I\) 的像落在给定开集 \(U\subset B\) 内”的路径组成。在\cref{g2:lem:quotient-interval}的证明中已经建立了这里需要的指数律：对任意拓扑空间 \(Z\)，映射 \(Z\rightarrow C(I,B)\) 连续，当且仅当其伴随映射 \(Z\times I\rightarrow B\) 连续。证明的两个关键点是 \(I\) 的局部紧 Hausdorff 性和紧集的有限覆盖性质。因此，下面的路径族可以通过含一个额外区间参数的公式直接检验连续性。

\begin{definition}[基点道路空间与回路空间]\label{g2:def:path-loop-space}
	设 \((B,b_0)\) 为带基点拓扑空间。在紧开拓扑下令
	\[
		P_{b_0}B\coloneqq\{\gamma:I\rightarrow B:\gamma(0)=b_0\},\qquad
		\Omega_{b_0}B\coloneqq\{\gamma\in P_{b_0}B:\gamma(1)=b_0\}.
	\]
	两者分别称为\term[jidiandaolukongjian]{基点道路空间}[based path space]与\term[huilukongjian]{回路空间}[loop space]。它们都以常值道路为基点。记终点求值映射为 \(\operatorname{ev}_1(\gamma)=\gamma(1)\)。\symidx{\(P_{b_0}B\)：从 \(b_0\) 出发的道路空间}\symidx{\(\Omega_{b_0}B\)：以 \(b_0\) 为基点的回路空间}
\end{definition}

\begin{proposition}[道路空间的纤维化]\label{g2:prop:path-fibration}
	对任意带基点拓扑空间 \((B,b_0)\)，终点求值映射
	\[
		\operatorname{ev}_1:P_{b_0}B\longrightarrow B
	\]
	是 Hurewicz 纤维化，其在 \(b_0\) 上的纤维为 \(\Omega_{b_0}B\)，而总空间 \(P_{b_0}B\) 可缩。
\end{proposition}
\begin{proof}
	设 \(Z\) 为任意拓扑空间，初始提升给出连续道路族 \(\alpha_z\)，底同伦为 \(H:Z\times I\rightarrow B\)，满足 \(H(z,0)=\alpha_z(1)\)。移动到时刻 \(s\) 时，把原道路与已经走过的底道路接在一起；同时让后一段所占时间从零连续增长。具体定义
	\[
		\lambda_{z,s}(t)=
		\begin{cases}
			\alpha_z((1+s)t),&0\le t\le(1+s)^{-1},\\
			H(z,(1+s)t-1),&(1+s)^{-1}\le t\le1.
		\end{cases}
	\]
	在接缝处，两式均等于 \(\alpha_z(1)=H(z,0)\)。它们分别定义在两个闭集上并且连续，所以 \((z,s,t)\mapsto\lambda_{z,s}(t)\) 连续。指数律给出连续的道路族 \(\widetilde H(z,s)=\lambda_{z,s}\)。其起点恒为 \(b_0\)，终点为 \(H(z,s)\)，且 \(s=0\) 时恰为原道路 \(\alpha_z\)，故这是所需提升。

	定义收缩
	\[
		K(\gamma,s)(t)=\gamma((1-s)t).
	\]
	指数律再次保证其连续。它从恒等映射变到常值道路，始终固定常值道路，故 \(P_{b_0}B\) 可缩。纤维的描述直接来自终点条件。
\end{proof}

这个例子把“球面映射怎样移动”转成了“回路怎样成族”：总空间中的道路可以缩回起点，底空间的同伦信息于是集中在回路空间中。下一节的长正合列将把这种观察写成 \(\pi_n(B)\cong\pi_{n-1}(\Omega_{b_0}B)\)。

\begin{definition}[同伦纤维]\label{g2:def:homotopy-fiber}
	设 \(f:X\rightarrow B\) 连续，\(b_0\in B\)。赋予下列空间以 \(X\times C(I,B)\) 的子空间拓扑：
	\[
		\operatorname{hofib}_{b_0}(f)\coloneqq
		\{(x,\gamma):\gamma(0)=f(x),\ \gamma(1)=b_0\}.
	\]
	称它为 \(f\) 在 \(b_0\) 上的\term[tonglunxianwei]{同伦纤维}[homotopy fiber]。若已选 \(x_0\in X\) 且 \(f(x_0)=b_0\)，它的基点取 \((x_0,c_{b_0})\)，其中 \(c_{b_0}\) 是常值道路。\symidx{\(\operatorname{hofib}_{b_0}(f)\)：映射在 \(b_0\) 上的同伦纤维}
\end{definition}

通常纤维要求 \(f(x)=b_0\)，同伦纤维则把从 \(f(x)\) 到 \(b_0\) 的道路也保留下来。这样，底点沿道路移动时，已经选择的路径可以随之拼接。比如 \(X\) 只有一个点且 \(f(*)=b_0\) 时，同伦纤维就是 \(\Omega_{b_0}B\)。

\subsection{纤维丛的同伦提升}
\label{g2:subsec:0802:03}

局部乘积中的提升只要保持纤维坐标不变。要把这种办法用于纤维丛，需要让不同局部乘积中的选择在交界处接合。闭圆盘的紧性使我们可以把整个提升问题分成有限多个落在单个平凡化邻域内的小问题。

回顾\cref{g2:def:topological-fiber-bundle}：连续满射 \(p:E\to B\) 局部平凡，指每个 \(b\in B\) 都有开邻域 \(U\)、一个纤维模型 \(F_U\) 及同胚
\[
	h_U:p^{-1}(U)\longrightarrow U\times F_U,
	\qquad \operatorname{pr}_U\circ h_U=p|_{p^{-1}(U)}.
\]
当这些模型都可选为同一个非空空间 \(F\) 时，称其纤维为 \(F\)。以下证明用到的是这样的局部乘积图及其连续性。

\begin{theorem}[纤维丛的提升]\label{g2:thm:bundle-serre}
	任意局部平凡纤维丛 \(p:E\rightarrow B\) 都是 Serre 纤维化，并且具有\cref{g2:lem:relative-homotopy-lift}中的相对 CW 同伦提升性质。
\end{theorem}
\begin{proof}
	闭圆盘同胚于闭立方体，所以只需处理 \(H:I^m\times I\rightarrow B\) 及其初始提升。对一组平凡化邻域 \(\{U_\alpha\}\)，逆像 \(\{H^{-1}(U_\alpha)\}\) 是紧度量空间 \(I^m\times I\) 的开覆盖。取一个 Lebesgue 数，再把空间坐标与时间坐标都作足够细的有限等距细分，使每个小棱柱 \(C\times[t_j,t_{j+1}]\) 的像包含在某个 \(U_\alpha\) 内。Lebesgue 数的存在可由有限子覆盖直接证明：若不存在，就能选取直径趋零却不落在任何一个覆盖成员内的子集，选点取收敛子列后，与极限点的一个覆盖邻域矛盾。

	按时间片顺序提升。在一片内部，先提升空间网格的顶点，再提升边、二维面，直到所有 \(m\) 维小立方体。处理一个 \(r\) 维网格胞腔 \(C\) 时，在
	\[
		A_C=(C\times\{t_j\})\cup(\partial C\times[t_j,t_{j+1}])
	\]
	上已经有提升：下底来自前一时间片，侧面来自本片较低维胞腔。取包含整个 \(H(C\times[t_j,t_{j+1}])\) 的平凡化邻域 \(U_\alpha\)。把已有提升写成
	\[
		h_\alpha\bigl(\widetilde H(z,t)\bigr)
		=\bigl(H(z,t),v(z,t)\bigr),\qquad (z,t)\in A_C.
	\]
	由\cref{g2:prop:disk-hep}，经缩放和圆盘、立方体间的同胚，有回缩 \(r_C:C\times[t_j,t_{j+1}]\rightarrow A_C\)。定义新提升为
	\[
		\widetilde H(z,t)=h_\alpha^{-1}\bigl(H(z,t),v(r_C(z,t))\bigr).
	\]
	它的投影严格等于当前 \(H(z,t)\)，并在 \(A_C\) 上等于原提升。对于零维胞腔，回缩就是退回该时间片的起点。因此整个逐维构造可以开始，也不会改变已完成的边界。有限闭集粘合保证每片连续，再逐时间片粘合，得到全部同伦的提升。由\cref{g2:lem:relative-homotopy-lift}即得相对 CW 形式。
\end{proof}

覆盖对应离散纤维的局部乘积，其提升唯一性使各局部选择自动一致。一般纤维丛需要上面保持边界的延拓步骤。纤维丛的这一提升证明也可对照 \cite[第4.2节，Fiber Bundles]{Hatcher2002}。

\begin{example}[纤维化中纤维的形状可以改变]\label{g2:exm:fibration-not-bundle}
	令
	\[
		E=\{(b,v)\in\mathbb R^2:|v|\le|b|\},
		\qquad p(b,v)=b.
	\]
	这是一个 Hurewicz 纤维化。事实上，给定任意 \(H:Z\times I\rightarrow\mathbb R\) 和初始提升 \(a(z)=(H(z,0),v_0(z))\)，令
	\[
		v(z,t)=\max\{-|H(z,t)|,\min\{v_0(z),|H(z,t)|\}\}.
	\]
	最大、最小运算连续，\((H(z,t),v(z,t))\) 始终落在 \(E\)，且因 \(|v_0(z)|\le|H(z,0)|\)，在时刻零恢复给定初值。\(p^{-1}(0)\) 是一点，\(p^{-1}(b)\) 在 \(b\ne0\) 时是非退化闭区间，所以 \(p\) 在零点附近没有局部平凡化。提升性质控制连续变形，局部平凡化还要求邻近纤维具有相同的拓扑形状。
\end{example}

\subsection{拉回和映射的纤维化替换}
\label{g2:subsec:0802:04}

先看怎样限制底空间。若 \(u:C\rightarrow B\) 连续，将 \(E\) 中位于 \(u(c)\) 上方的点与 \(c\) 配对，得到
\[
	u^*E\coloneqq\{(c,e)\in C\times E:u(c)=p(e)\},
	\qquad p_C(c,e)=c.
\]
其中取乘积空间的子空间拓扑。这个\term[lahuixianweihua]{拉回}[pullback]给出交换图
\[
	\begin{tikzcd}[column sep=large,row sep=large]
		u^*E \arrow[r,"\operatorname{pr}_E"] \arrow[d,"p_C"'] & E \arrow[d,"p"] \\
		C \arrow[r,"u"'] & B.
	\end{tikzcd}
\]
图中左上空间只由“右边投影与下边映射相等”的条件确定。

\begin{proposition}[拉回保持提升性质]\label{g2:prop:fibration-pullback}
	设 \(p:E\rightarrow B\) 与 \(u:C\rightarrow B\) 连续。若 \(p\) 为 Serre 纤维化，则 \(p_C:u^*E\rightarrow C\) 也是；若 \(p\) 为 Hurewicz 纤维化，则 \(p_C\) 也为 Hurewicz 纤维化。若 \(p\) 是以 \(F\) 为纤维的纤维丛，则 \(p_C\) 也是以 \(F\) 为纤维的纤维丛。
\end{proposition}
\begin{proof}
	设 \(H:Z\times I\rightarrow C\) 及初始提升 \(a(z)=(H(z,0),e(z))\) 已给定。对底同伦 \(uH\)，用 \(p\) 的提升性质取以 \(e(z)\) 为初值的提升 \(\widetilde H_E\)。则
	\[
		\widetilde H_C(z,t)=\bigl(H(z,t),\widetilde H_E(z,t)\bigr)
	\]
	落在 \(u^*E\)，且满足所需初值与投影等式。允许的参数空间 \(Z\) 与原提升性质相同，因此两种纤维化均得到保持。对纤维丛，在 \(u^{-1}(U)\) 上用原平凡化的第二个坐标，就得到 \(p_C^{-1}(u^{-1}(U))\cong u^{-1}(U)\times F\)。
\end{proof}

任意映射也可以加入一条路径，变成具有提升性质的映射。加入的路径可以随底同伦延长，而把它缩回初点又能恢复原来的源空间。

\begin{proposition}[映射的纤维化替换]\label{g2:prop:fibration-replacement}
	设 \(f:X\rightarrow B\) 为拓扑空间之间的连续映射。令
	\[
		E_f\coloneqq\{(x,\gamma)\in X\times C(I,B):\gamma(0)=f(x)\},
		\qquad p_f(x,\gamma)=\gamma(1),
	\]
	并取 \(j(x)=(x,c_{f(x)})\)、\(q(x,\gamma)=x\)。则 \(p_f\) 为 Hurewicz 纤维化，\(p_fj=f\)，\(qj=\operatorname{id}_X\)，且 \(jq\) 与 \(\operatorname{id}_{E_f}\) 相对 \(j(X)\) 同伦。因此 \(j:X\rightarrow E_f\) 是同伦等价。\(p_f\) 在 \(b_0\in B\) 上的实际纤维就是 \(\operatorname{hofib}_{b_0}(f)\)。\symidx{\(E_f\)：映射的道路空间替换}
\end{proposition}
\begin{proof}
	给定任意参数空间 \(Z\)、初始提升 \(a(z)=(x(z),\alpha_z)\) 及底同伦 \(H:Z\times I\rightarrow B\)，其中 \(H(z,0)=\alpha_z(1)\)。保持 \(x(z)\) 不动，并采用\cref{g2:prop:path-fibration}中的道路 \(\lambda_{z,s}\)，定义
	\[
		\widetilde H(z,s)=\bigl(x(z),\lambda_{z,s}\bigr).
	\]
	每条 \(\lambda_{z,s}\) 的起点仍是 \(\alpha_z(0)=f(x(z))\)，所以它落在 \(E_f\)；终点与初值也分别给出 \(p_f\widetilde H=H\) 和 \(\widetilde H(z,0)=a(z)\)。上述公式的连续性已由指数律证明，故 \(p_f\) 为 Hurewicz 纤维化。

	公式
	\[
		K((x,\gamma),s)=\bigl(x,t\mapsto\gamma((1-s)t)\bigr)
	\]
	给出 \(\operatorname{id}_{E_f}\) 到 \(jq\) 的同伦，并固定所有常值道路；\(qj=\operatorname{id}_X\) 与 \(p_fj=f\) 可直接求值得到。最后，\(p_f^{-1}(b_0)\) 的两个路径端点条件与同伦纤维的定义完全相同。
\end{proof}

纤维化替换把 \(f\) 分解为
\[
	X\xrightarrow[\simeq]{\ j\ }E_f\xrightarrow{\ p_f\ }B.
\]
于是可以对 \(p_f\) 使用长正合列，并用同伦等价 \(j\) 将总空间中的同伦群换回 \(X\)。这正是同伦纤维进入一般映射问题的用途。

% !TeX root = ../../Book2.tex
% 纲要标题：### 8.3 纤维化的同伦长正合列
% 纲要位置：几何学计划纲要.md，第 1318 行。
% 正文状态：初稿；尚未正式审核。

% 纲要内容开始
% 
% * 连接映射的提升构造
% * 逐项验证正合性
% * 基点、连通分支和群作用
% * 自然性与比较
% 
% 纲要内容结束

\section{纤维化的同伦长正合列}
\label{g2:sec:0803}

设 \(p:E\to B\) 是 Serre 纤维化，选定 \(e_0\in E\)，记 \(b_0=p(e_0)\)、\(F=p^{-1}(b_0)\)。底空间中的一个球面映射未必能提升为总空间中的球面映射：沿参数逐步提升以后，最后一个面可能停在 \(F\) 中的一张非平凡球面上。连接映射记录的正是这个未能收拢的面。先允许这一个面在纤维中移动，提升就总能完成；这把纤维化问题与已经建立的相对同伦联系起来。

\subsection{连接映射的提升构造}
\label{g2:subsec:0803:01}

沿用\cref{g2:def:relative-homotopy}的立方体模型，令 \(I^n=I^{n-1}\times I\)，最后一个坐标记为 \(t\)，并令
\[
	J^{n-1}=(\partial I^{n-1}\times I)\cup(I^{n-1}\times\{0\}).
\]
顶面 \(t=1\) 可以映入 \(F\)，其余面必须映到 \(e_0\)。投影到 \(B\) 后，整个边界都成为 \(b_0\)，于是有映射
\[
	P_n:\pi_n(E,F,e_0)\longrightarrow\pi_n(B,b_0),\qquad [v]\longmapsto[pv].
\]

\begin{proposition}
	\label{g2:prop:fibration-relative-comparison}
	若 \(p:E\to B\) 是 Serre 纤维化，\(e_0\in E\)、\(b_0=p(e_0)\)、\(F=p^{-1}(b_0)\)，则 \(P_n\) 在 \(n\ge2\) 时是群同构，在 \(n=1\) 时是保持基点的集合双射。
\end{proposition}

\begin{proof}
	给定 \(u:(I^n,\partial I^n)\to(B,b_0)\)，先在 \(J^{n-1}\) 上指定常值提升 \(e_0\)。\cref{g2:lem:relative-homotopy-lift}把这个提升延拓到整个立方体，得到 \(\widetilde u\)。因为 \(u\) 在顶面也等于 \(b_0\)，\(\widetilde u\) 的顶面落在 \(F\)，所以它定义相对类，且 \(P_n[\widetilde u]=[u]\)。这证明满射。

	再设 \(v_0,v_1\) 的投影由保持边界基点的同伦 \(H:I^n\times I_s\to B\) 连接。在 \(s=0,1\) 上分别指定提升 \(v_0,v_1\)，在 \(J^{n-1}\times I_s\) 上指定 \(e_0\)。它们在交叠处一致。把 \(t\) 看成提升的时间坐标，其余 \(n\) 个坐标为 \((x,s)\)，这些已知数据恰好位于底面和全部侧面；唯有 \(t=1\) 的面尚未指定。再次应用相对提升引理，得到 \(\widetilde H\)。在顶面上，\(p\widetilde H=H=b_0\)，所以顶面始终留在 \(F\)；这正是 \(v_0\) 与 \(v_1\) 的相对同伦。这个论证也适用于 \(n=1\)，直接证明了尖点集上的单射。

	当 \(n\ge2\) 时，投影保持沿第一坐标的拼接，故 \(P_n\) 是群同态。已经得到的双射因而是群同构。
\end{proof}

\begin{definition}
	\label{g2:def:fibration-connecting-map}
	在上述条件下，纤维化的\term[lian jie ying she]{连接映射}[connecting map]定义为
	\[
		\partial_n=\partial\circ P_n^{-1}:\pi_n(B,b_0)\longrightarrow\pi_{n-1}(F,e_0).
	\]
	这里右侧的 \(\partial\) 是空间对 \((E,F)\) 的顶面限制。因此，对上述提升 \(\widetilde u\)，有
	\[
		\partial_n[u]=[\widetilde u(-,1)].
	\]
	当 \(n=1\) 时，右侧指提升道路终点所在的道路连通分支。
\end{definition}

\cref{g2:prop:fibration-relative-comparison}同时解释了构造为何与选择无关：同一底空间同伦类的所有提升给出同一相对类，顶面因而给出同一纤维同伦类。\(n\ge2\) 时连接映射是群同态；\(n=1\) 时其值域为 \(\pi_0(F)\)，通常只具有带基点集合的结构。

\subsection{逐项验证正合性}
\label{g2:subsec:0803:02}

连接映射的意义可以表述为一个可检验的问题：顶面能否在纤维中缩掉？若能，就可以把提升补成一个球面。反过来，如果已经有一个总空间中的球面，它的投影便不会留下这个障碍。下面把这两个方向与另外两处核、像关系一起证明。

\begin{theorem}
	\label{g2:thm:fibration-homotopy-sequence}
	设 \(p:E\to B\) 是 Serre 纤维化，\(e_0\in E\)，\(b_0=p(e_0)\)，\(F=p^{-1}(b_0)\)，\(i:F\hookrightarrow E\) 为包含。存在自然的长正合列
	\[
		\cdots\longrightarrow\pi_n(F,e_0)\xrightarrow{i_*}\pi_n(E,e_0)
		\xrightarrow{p_*}\pi_n(B,b_0)\xrightarrow{\partial_n}\pi_{n-1}(F,e_0)
		\longrightarrow\cdots
	\]
	其低阶部分为
	\[
		\begin{aligned}
			\pi_2(B,b_0)&\xrightarrow{\partial_2}\pi_1(F,e_0)
			\xrightarrow{i_*}\pi_1(E,e_0)\xrightarrow{p_*}\pi_1(B,b_0)\\
			&\xrightarrow{\partial_1}\pi_0(F,e_0)\xrightarrow{i_*}\pi_0(E,e_0)
			\xrightarrow{p_*}\pi_0(B,b_0).
		\end{aligned}
	\]
	群项上的正合指像等于核；尖点集项上的正合指前一映射的像等于后一映射所取基点的原像。若 \(B\) 道路连通，则 \(\pi_0(F)\to\pi_0(E)\) 满射，可以在其后接一个单点集。
\end{theorem}

\begin{proof}
	\cref{g2:thm:pair-homotopy-sequence}给出 \((E,F,e_0)\) 的相对同伦序列。以\cref{g2:prop:fibration-relative-comparison}中的 \(P_n\) 替换相对项，就得到定理中的映射：\(\pi_n(E)\) 先相对化再由 \(P_n\) 投影，正是通常的 \(p_*\)。下面也直接辨认每一种正合位置，以说明各个核代表什么几何事实。

	在 \(\pi_n(E)\) 处，纤维中的映射投影为常值，所以 \(\operatorname{im}i_*\subseteq\ker p_*\)。若带基点球面映射 \(a\) 的投影有带基点零伦，则用球面及其基点组成的 CW 对，提升这个零伦并令基点始终为 \(e_0\)。初始提升为 \(a\)，最终提升整个落在 \(F\)。因此 \([a]\) 来自 \(\pi_n(F)\)。这一论证对 \(n\ge1\) 都成立。

	在 \(\pi_n(B)\) 处，\(n\ge2\) 时，先取 \(u\) 的提升 \(v\)，其顶面为 \(a\)。若 \([a]=0\in\pi_{n-1}(F)\)，取 \(a\) 在 \(F\) 中的带基点零伦 \(A(x,s)\)。在 \(v\) 顶端接上这个零伦，得到
	\[
		w(x,t)=
		\begin{cases}
			v(x,2t),&0\le t\le\tfrac12,\\
			A(x,2t-1),&\tfrac12\le t\le1.
		\end{cases}
	\]
	\(w\) 的整个边界都等于 \(e_0\)，而 \(pw\) 等于先运行 \(u\)、再停在 \(b_0\) 的重参数映射，故 \([u]=p_*[w]\)。若 \([u]\) 一开始来自 \(\pi_n(E)\)，便可选整个边界为 \(e_0\) 的提升，连接类为零。\(n=1\) 时，同一个构造把“顶面零伦”换成“在 \(F\) 内从提升终点走回 \(e_0\) 的道路”，证明 \(\partial_1^{-1}([e_0])=\operatorname{im}p_*\)。

	在 \(\pi_{n-1}(F)\) 处，\(n\ge2\) 时，连接映射给出的顶面连同提升立方体，已经提供它在 \(E\) 中的零伦。反过来，若 \(a:I^{n-1}/\partial I^{n-1}\to F\) 在 \(E\) 中有带基点零伦，反向运行这个零伦就得到 \(v:I^n\to E\)，底面和侧面为 \(e_0\)，顶面为 \(a\)。此时 \(pv\) 的整个边界为 \(b_0\)，而 \(\partial_n[pv]=[a]\)。这也包括可能非交换的 \(\pi_1(F)\)。

	在 \(\pi_0(F)\) 处，连接映射得到的点由提升道路与 \(e_0\) 相连，所以它落在 \(\pi_0(E)\) 的基点原像中。反过来，若 \(y\in F\) 在 \(E\) 中与 \(e_0\) 由道路 \(c\) 相连，则 \(pc\) 是 \(b_0\) 处的环路，且 \(\partial_1[pc]=[y]\)。在 \(\pi_0(E)\) 处，若 \(e\in E\) 的投影在 \(b_0\) 的分支中，提升从 \(p(e)\) 到 \(b_0\) 的道路，得到从 \(e\) 到 \(F\) 的道路；反方向由 \(p(F)=\{b_0\}\) 立即得到。若 \(B\) 道路连通，这就对 \(E\) 的所有分支成立。
\end{proof}

\begin{example}
	\label{g2:exm:loop-space-shift}
	令 \(B\) 是带基点空间，\(P_{b_0}B\) 为从 \(b_0\) 出发的道路空间，终点映射的纤维为回路空间 \(\Omega B\)。\cref{g2:prop:path-fibration}给出这个 Hurewicz 纤维化。\(P_{b_0}B\) 由 \(\gamma(t)\mapsto\gamma((1-s)t)\) 保持常值道路地收缩，所以长正合列给出
	\[
		\pi_n(B,b_0)\cong\pi_{n-1}(\Omega B,c_{b_0})\quad(n\ge2),
		\qquad \pi_1(B,b_0)\cong\pi_0(\Omega B).
	\]
	后一个式子的集合双射在回路拼接下也是群同构。对 \(u(x,t)\)，可以取提升 \(\widetilde u(x,t)(r)=u(x,tr)\)；其顶面正是 \(x\mapsto(t\mapsto u(x,t))\)。这说明指标降低一次来自把一个参数变成回路参数。
\end{example}

\subsection{基点、连通分支和群作用}
\label{g2:subsec:0803:03}

序列末端还能记录总空间有多少个道路连通分支。沿底空间中的环路提升，一个纤维分支可能被带到另一个分支；能够这样互相到达的分支，在总空间中恰好连在一起。这是第1章覆盖单值作用在一般纤维化中的延续。

\begin{proposition}
	\label{g2:prop:fibration-component-action}
	设 \(p:E\to B\) 是 Serre 纤维化，\(B\) 道路连通，且 \(F=p^{-1}(b_0)\ne\varnothing\)。沿环路提升的终点定义右作用
	\[
		\pi_0(F)\times\pi_1(B,b_0)\longrightarrow\pi_0(F),\qquad
		[y]\cdot[\beta]=[\widetilde\beta(1)],\quad\widetilde\beta(0)=y.
	\]
	此作用的轨道集合与 \(\pi_0(E)\) 自然对应。对 \(e_0\in F\)，连接映射为 \(\partial_1[\beta]=[e_0]\cdot[\beta]\)，而 \([e_0]\) 的稳定子是 \(p_*\pi_1(E,e_0)\)。
\end{proposition}

\begin{proof}
	固定一个起点，两条提升的终点在 \(F\) 中相连：把两条提升放在一个方形的两条相对边上，在初始边指定该起点，底空间映射取与第二参数无关的原道路；缺少终点边的三条边组成提升引理中的部分边界。补全提升后，终点边就是所需的纤维内道路。若起点在 \(F\) 中沿一条道路变化，把这条道路作为方形的初始边，同样得到终点分支不变。底空间道路的端点固定同伦也由相同的方形提升处理。因此作用良定义。

	道路乘法 \(\alpha*\beta\) 先走 \(\alpha\)，后走 \(\beta\)。依次提升两段，就有 \(([y]\cdot[\alpha])\cdot[\beta]=[y]\cdot[\alpha*\beta]\)；常值道路作用为恒等，反向道路给出逆变换，所以这是右作用。作用同轨的点通过提升道路在 \(E\) 中相连；若两点在 \(E\) 中相连，投影道路就是实现相应运输的底空间环路。再由 \(B\) 道路连通时每个总空间分支都遇到 \(F\)，得到轨道集合的断言。

	若环路是 \(E\) 中基点环路的投影，就固定 \([e_0]\)。反过来，若其提升终点与 \(e_0\) 在 \(F\) 中相连，接上这条纤维内道路就得到 \(E\) 中的环路，投影与原环路同伦。因而稳定子恰为 \(p_*\pi_1(E,e_0)\)。
\end{proof}

\begin{example}
	\label{g2:exm:covering-sequence-tail}
	对 \(p:S^1\to S^1\)、\(p(z)=z^m\)，\(m\ge2\)，取基点 \(1\)。纤维是 \(m\) 次单位根，低阶序列为
	\[
		0\longrightarrow\mathbb Z\xrightarrow{\times m}\mathbb Z
		\xrightarrow{\partial_1}\{1,\zeta,\ldots,\zeta^{m-1}\}\longrightarrow\{*\},
		\qquad\zeta=e^{2\pi\mathrm i/m},
	\]
	其中 \(\partial_1(k)=\zeta^k\)。整数绕行数通过提升终点记录其模 \(m\) 的余数；纤维分支只有一个轨道，所以总空间连通。在一般连通覆盖中，纤维仍是一个传递的 \(\pi_1(B)\)-集合，稳定子可以不是正规子群。末端的集合因而要由作用理解，不能一概当作商群。
\end{example}

高阶群还需要追踪基点。若 \(\lambda\) 是 \(E\) 中从 \(e_0\) 到 \(e_1\) 的道路，\(\beta=p\lambda\)，则以 \(\lambda\) 规定移动的基点，可以运输纤维的带基点球面：对 \(a:S^r\to p^{-1}(\beta(0))\)，提升底空间同伦 \((z,t)\mapsto\beta(t)\)，同时规定 \((*,t)\mapsto\lambda(t)\)。相对提升引理给出终端球面，并由带第二同伦参数的同一构造保证其类与选择无关。反向道路提供逆映射。若 \(\lambda\) 是环路，依次提升对应于道路拼接，便得到 \(\pi_1(E,e_0)\) 在 \(\pi_r(F,e_0)\) 上的右作用；当 \(\lambda\) 位于 \(F\) 内时，它就是\cref{g2:prop:basepoint-action}中的基点变化作用。

\ProseParagraphBreak

这也解释了为何 \(\pi_0(F)\) 的运输只需底空间环路，而带基点高阶群还要指定基点的移动。选择不同的纤维内回程，最终识别之间相差一个基本群作用。实际计算中，若纤维道路连通且单连通，这个差异消失；对一般纤维，则把移动基点的道路连同群写清楚。沿 \(\lambda\) 改变总空间基点、沿 \(p\lambda\) 改变底空间基点，再用上述运输改变纤维基点，得到的长正合列彼此对应；连接映射的相容性来自对同一立方体及其顶面同时进行提升。

\subsection{自然性与比较}
\label{g2:subsec:0803:04}

纤维化经常通过拉回、坐标变换或映射路径构造与另一个纤维化比较。要把一次计算搬到新模型上，必须确认连接映射也随这些映射一起变化。它的提升定义恰好使这种相容性成为直接的复合关系。

\begin{proposition}
	\label{g2:prop:fibration-sequence-naturality}
	设 \(p:E\to B\)、\(p':E'\to B'\) 为 Serre 纤维化，带基点映射 \(h:E\to E'\)、\(g:B\to B'\) 满足 \(p'h=gp\)。记对应基点纤维为 \(F,F'\)，\(h_F=h|_F\)。则 \(h_F,h,g\) 诱导长正合列之间的态射，特别地
	\[
		\begin{tikzcd}[column sep=large]
			\pi_n(E)\arrow[r,"p_*"]\arrow[d,"h_*"']&\pi_n(B)\arrow[r,"\partial_n"]\arrow[d,"g_*"]&\pi_{n-1}(F)\arrow[d,"(h_F)_*"]\\
			\pi_n(E')\arrow[r,"p'_*"]&\pi_n(B')\arrow[r,"\partial'_n"]&\pi_{n-1}(F')
		\end{tikzcd}
	\]
	交换。\(n=1\) 时右端按带基点集合解释；\(\pi_0(F)\) 的右作用也与 \(g_*\) 和 \((h_F)_*\) 相容。
\end{proposition}

\begin{proof}
	若 \(v\) 是 \(u\) 的提升，则 \(hv\) 是 \(gu\) 的提升，且 \((hv)(-,1)=h_F(v(-,1))\)。这给出连接映射所在方块的交换性。其他方块只涉及包含与投影的复合；道路提升的终点也由 \(h_F\) 送到相应终点，故分支作用相容。
\end{proof}

下面的比较结论在替换空间或替换纤维时尤其有用。若一个连续映射在道路分支上为双射，并在每个基点的所有正次同伦群上为同构，就称它为\term[ruo tong lun deng jia]{弱同伦等价}[weak homotopy equivalence]。

\begin{proposition}
	\label{g2:prop:fibration-comparison}
	设 \(p:E\to B\)、\(p':E'\to B'\) 是 Serre 纤维化，\(B,B'\) 道路连通，连续映射 \(h:E\to E'\)、\(g:B\to B'\) 满足 \(p'h=gp\)。选 \(b_0\in B\)，令 \(F=p^{-1}(b_0)\)、\(F'=(p')^{-1}(g(b_0))\)，并假设 \(F,F'\) 非空。如果 \(g\) 与 \(h_F=h|_F:F\to F'\) 都是弱同伦等价，则 \(h\) 也是弱同伦等价。
\end{proposition}

\begin{proof}
	由\cref{g2:prop:fibration-component-action}，总空间的分支是纤维分支的作用轨道。\(\pi_1(B)\to\pi_1(B')\) 为同构，\(\pi_0(F)\to\pi_0(F')\) 为等变双射，故 \(\pi_0(E)\to\pi_0(E')\) 为双射。每个分支都含有纤维点；基点变化同构允许在这样的点上检查其余同伦群。

	固定 \(e\in F\)。先证明 \(h_*:\pi_n(E,e)\to\pi_n(E',h(e))\) 单射，其中 \(n\ge1\)。若 \(h_*x=0\)，由底空间同伦群上的单射知 \(p_*x=0\)，于是 \(x=i_*a\)。又有 \(i'_*(h_F)_*a=0\)，正合性给出 \((h_F)_*a=\partial'_{n+1}b'\)。由底空间同伦群上的满射，取 \(b\) 使 \(g_*b=b'\)。自然性和 \((h_F)_*\) 的单射给出 \(a=\partial_{n+1}b\)，故 \(x=0\)。在 \(n=1\) 时，零元按群的单位元理解，这个论证仍逐项成立。

	再给定 \(x'\in\pi_n(E')\)。取 \(b\in\pi_n(B)\) 使 \(g_*b=p'_*x'\)。\(n\ge2\) 时，\((h_F)_*\partial_n b=0\)，由纤维群上的单射知 \(\partial_n b=0\)，所以 \(b=p_*x_0\)。\(n=1\) 时用 \(\pi_0(F)\to\pi_0(F')\) 的单射得到同样结论：\(\partial_1b\) 必为基点分支，故也有 \(b=p_*x_0\)。于是 \(x'\) 与 \(h_*x_0\) 的差属于 \(\ker p'_*\)，来自某个纤维类 \(a'\)；选 \(a\) 使 \((h_F)_*a=a'\)，便得到 \(x'\) 的原像。对于 \(n=1\)，具体取 \(a'\) 满足 \(i'_*a'=x'(h_*x_0)^{-1}\)，原像为 \(i_*a\,x_0\)。这证明所有正次同伦群上的满射。
\end{proof}

\cref{g2:prop:fibration-replacement}还把任意带基点映射 \(f:X\to Y\) 换成 \(E_f\to Y\)，其中 \(E_f\simeq X\)。把这个同伦等价诱导的同构代入长正合列，得到
\[
	\cdots\longrightarrow\pi_n(F_f)\longrightarrow\pi_n(X)
	\xrightarrow{f_*}\pi_n(Y)\longrightarrow\pi_{n-1}(F_f)\longrightarrow\cdots.
\]
这里 \(F_f=\operatorname{hofib}_{y_0}(f)\)，其元素同时保存一个 \(x\in X\) 和一条从 \(f(x)\) 到 \(y_0\) 的道路。因而这个序列测量的是映射 \(f\) 的同伦纤维；即使普通原像 \(f^{-1}(y_0)\) 很小，这些道路数据仍能记录 \(f\) 对同伦群的作用。

% !TeX root = ../../Book2.tex
% 纲要标题：### 8.4 Hopf 纤维化与经典群
% 纲要位置：几何学计划纲要.md，第 1275 行。
% 正文状态：未撰写；本文件目前是框架。

% 纲要内容开始
% 
% * 复、四元数 Hopf 映射
% * 球面丛和经典群的纤维序列
% * 低阶同伦群计算
% * 由纤维化构造非平凡球面映射
% 
% ---
% 
% 纲要内容结束

\section{Hopf 纤维化与经典群}
\label{g2:sec:0804}

待撰写

\subsection{复、四元数 Hopf 映射}
\label{g2:subsec:0804:01}
待撰写

\subsection{球面丛和经典群的纤维序列}
\label{g2:subsec:0804:02}
待撰写

\subsection{低阶同伦群计算}
\label{g2:subsec:0804:03}
待撰写

\subsection{由纤维化构造非平凡球面映射}
\label{g2:subsec:0804:04}
待撰写
% 纲要条目排为编号小节标题，并进入目录与书签；说明段仍保留于纲要注释中。


\begin{chaptersummary}
	高阶同伦群把球面映射及其填充问题组织成群；两个独立的拼接方向产生交换性，相对模型则把指定顶面留给子空间。基点运输沿道路给出同构，其环路版本形成基本群作用。纤维化能够保持已有边界数据地提升同伦，因此 \(\pi_n(E,F)\) 与 \(\pi_n(B)\) 对应，顶面限制给出连接映射。长正合列的每一处核、像等式分别表达提升、填充或纤维内收缩；低阶末端通过分支与作用轨道描述总空间的连通性。路径替换使任意映射都能进入这一计算框架。Hopf 丛给出 \(\pi_2(S^2)\cong\mathbb Z\) 及球面之间的高阶比较，经典群的球面丛给出一、二次同伦群的逐阶计算；映射度对同伦加法的相容性进一步保证复与四元数 Hopf 映射具有无穷阶。
\end{chaptersummary}

\BookExercises

\begin{exercise}\label{g2:ex:homotopy-product-retract}
	设 \((X,x_0)\)、\((Y,y_0)\) 为带基点空间，\(n\ge1\)。证明两个投影诱导 \(\pi_n(X\times Y,(x_0,y_0))\cong\pi_n(X,x_0)\times\pi_n(Y,y_0)\)。由此说明 \(S^m\) 的恒等映射作为 \(S^m\times S^1\) 中的球面仍非零伦，其中 \(m\ge1\)。
\end{exercise}
\begin{solution}
	映射 \(I^n\to X\times Y\) 与一对映射 \((f_X,f_Y)\) 一一对应，固定边界条件和定边界同伦也逐分量对应。因此投影在同伦类上给出双射，逆映射为 \(([f_X],[f_Y])\mapsto[(f_X,f_Y)]\)。拼接的两段公式逐分量成立，所以这是群同构。

	令 \(i:S^m\to S^m\times S^1\)、\(i(z)=(z,1)\)，第一投影满足 \(p_1i=\operatorname{id}_{S^m}\)。若 \(i\) 零伦，复合 \(p_1\) 便使球面恒等映射零伦。但恒等映射在 \(H_m(S^m;\mathbb Z)\cong\mathbb Z\) 上诱导恒等，常值映射诱导零映射，与同调的同伦不变性矛盾。因此 \([i]\ne0\)，也相应证明恒等球面代表非零类。
\end{solution}

\begin{exercise}\label{g2:ex:relative-disk-circle}
	以 \(1\in S^1\subset D^2\) 为基点，计算 \(\pi_2(D^2,S^1,1)\) 和 \(\pi_1(D^2,S^1,1)\)。说明边界绕数为 \(k\) 的相对圆盘类如何构造，并判断两个这样的类何时相同。
\end{exercise}
\begin{solution}
	凸圆盘以 \((1-s)z+s\) 收缩到 \(1\)，且始终固定 \(1\)，故它的正次同伦群全为零。空间对序列给出
	\[
		0\longrightarrow\pi_2(D^2,S^1,1)\xrightarrow{\partial}\pi_1(S^1,1)\longrightarrow0.
	\]
	因此 \(\partial\) 是群同构，绕数进一步把它识别为 \(\mathbb Z\)。对于 \(g_k(u)=e^{2\pi\mathrm{i}ku}\)，取
	\[
		f_k(u,t)=(1-t)+t\,g_k(u),\qquad (u,t)\in I^2.
	\]
	凸性保证它进入 \(D^2\)；底边和两条侧边都为 \(1\)，顶边为 \(g_k\)，故 \(\partial[f_k]=k\)。边界同构说明每个相对类等于唯一的 \([f_k]\)，两类相同当且仅当绕数相同。

	低阶序列中 \(\pi_1(D^2)=0\)，而 \(S^1\) 道路连通。每条相对道路的终点都能在 \(S^1\) 内回到 \(1\)，从而相对同伦到圆盘中的环路；后者可收缩。因此 \(\pi_1(D^2,S^1,1)\) 只有指定元素。
\end{solution}

\begin{exercise}\label{g2:ex:relative-circle-two-points}
	令 \(X=S^1\)、\(A=\{1,-1\}\)，基点为 \(1\)。用指数覆盖计算尖点集 \(\pi_1(X,A,1)\)，并写出 \(\pi_1(S^1,1)\to\pi_1(X,A,1)\to\pi_0(A)\) 的两张映射。解释为什么末端 \(\pi_0(A)\) 的两点都能出现。
\end{exercise}
\begin{solution}
	把相对道路 \(\alpha\) 唯一提升成 \(\theta:I\to\mathbb R\)，使 \(\theta(0)=0\)、\(\alpha(t)=e^{2\pi\mathrm{i}\theta(t)}\)。终点条件等价于 \(\theta(1)\in\frac12\mathbb Z\)。在相对同伦中，终点连续地留在离散集 \(A\) 内，因而恒定；提升后的终值也连续地留在离散集 \(\frac12\mathbb Z\) 内，故不变。反过来，两条提升若终值相同，其线性插值固定两端，投影后给出相对同伦。因此 \([\alpha]\mapsto2\theta(1)\) 是到 \(\mathbb Z\) 的尖点集双射，指定元素对应 \(0\)。

	绕数为 \(k\) 的闭环路具有提升终值 \(k\)，所以第一张映射为 \(k\mapsto2k\)；第二张映射把偶数送到 \(\{1\}\)、奇数送到 \(\{-1\}\)。两点在 \(S^1\) 中属于同一道路分支，分别由常值道路与半圆道路到达，所以都出现。这里 \(\mathbb Z\) 用于标记尖点集；原来的相对道路没有自动给定以这个标号相加为定义的拼接运算。
\end{solution}

\begin{exercise}\label{g2:ex:topological-group-simple-action}
	设 \(G\) 为拓扑群，单位元为 \(e\)。证明 \(\pi_1(G,e)\) 在 \(\pi_n(G,e)\) 上的基点运输作用对每个 \(n\ge1\) 都是平凡的，并由 \(n=1\) 的情形推出 \(\pi_1(G,e)\) 交换。
\end{exercise}
\begin{solution}
	设 \(f:(D^n,S^{n-1})\to(G,e)\)，\(\gamma:I\to G\) 是基点环路。点乘公式
	\[
		H(z,s)=\gamma(s)f(z)
	\]
	定义连续同伦，边界值为 \(\gamma(s)\)，初映射和终映射都为 \(f\)。\cref{g2:prop:basepoint-action}证明中的移动边界刻画给出 \(c_\gamma[f]=[H(-,1)]=[f]\)，故作用平凡。对 \(n=1\)，这个作用就是 \(a\cdot g=g^{-1}ag\)，于是所有 \(a,g\in\pi_1(G,e)\) 都满足 \(g^{-1}ag=a\)，即 \(ag=ga\)。证明只使用群乘法连续，无需假设 \(G\) 本身交换。
\end{solution}

\begin{exercise}\label{g2:ex:fibration-composition}
	设 \(p:E\rightarrow B\)、\(q:B\rightarrow C\) 都是 Serre 纤维化。证明 \(qp\) 也是 Serre 纤维化，并证明对 Hurewicz 纤维化也有相同结论。若 \(c\in C\)，说明 \(p:(qp)^{-1}(c)\rightarrow q^{-1}(c)\) 为什么仍为相应类型的纤维化。
\end{exercise}
\begin{solution}
	设允许的参数空间为 \(Z\)，底同伦为 \(H:Z\times I\rightarrow C\)，初始提升为 \(a:Z\rightarrow E\)，满足 \(qpa=H|_{Z\times\{0\}}\)。先由 \(q\) 的提升性质，取以 \(pa\) 为初值的 \(H_B:Z\times I\rightarrow B\)，使 \(qH_B=H\)。再由 \(p\) 的提升性质，取以 \(a\) 为初值的 \(H_E\)，使 \(pH_E=H_B\)。于是 \(qpH_E=H\)，且初值正确。取 \(Z=D^m\) 就得到 Serre 情形，允许任意 \(Z\) 就得到 Hurewicz 情形。

	最后的映射是 \(p\) 沿包含 \(q^{-1}(c)\hookrightarrow B\) 的拉回：拉回空间中的配对 \((b,e)\) 满足 \(b=p(e)\) 和 \(q(b)=c\)，所以与 \((qp)^{-1}(c)\) 同胚。由\cref{g2:prop:fibration-pullback}即得结论。
\end{solution}

\begin{exercise}\label{g2:ex:homotopic-maps-homotopy-fibers}
	设 \(f_0,f_1:X\rightarrow B\) 由连续同伦 \(H:X\times I\rightarrow B\) 连接，固定 \(b_0\in B\)。构造 \(\operatorname{hofib}_{b_0}(f_0)\) 与 \(\operatorname{hofib}_{b_0}(f_1)\) 之间的同伦等价。说明所构造映射怎样作用于路径坐标，并证明两个复合确实同伦于恒等映射。
\end{exercise}
\begin{solution}
	令 \(h_x(s)=H(x,s)\)，它从 \(f_0(x)\) 走到 \(f_1(x)\)。道路拼接采用先左后右的约定，故定义
	\[
		\Phi(x,\gamma)=(x,\bar h_x*\gamma),\qquad
		\Psi(x,\eta)=(x,h_x*\eta).
	\]
	例如 \(\bar h_x*\gamma\) 从 \(f_1(x)\) 经 \(f_0(x)\) 到 \(b_0\)，所以 \(\Phi\) 的靶空间正确。紧开拓扑的指数律与道路的分段拼接公式保证二者连续。

	复合 \(\Psi\Phi\) 的路径坐标是 \(h_x*(\bar h_x*\gamma)\)。先连续移动两个分点，将三段道路的参数分配从 \((1/2,1/4,1/4)\) 改为 \((1/4,1/4,1/2)\)，便得到 \((h_x*\bar h_x)*\gamma\)。这一步只改变走路的速度，端点保持不变。对 \(r\in I\)，在前半段使用往返道路
	\[
		\ell_{x,r}(t)=
		\begin{cases}
			h_x(2(1-r)t),&0\le t\le1/2,\\
			h_x(2(1-r)(1-t)),&1/2\le t\le1.
		\end{cases}
	\]
	当 \(r=0\) 时它是 \(h_x*\bar h_x\)，当 \(r=1\) 时它是 \(f_0(x)\) 处的常值道路。于是 \(\ell_{x,r}*\gamma\) 给出定端点的连续变形；最后将常值段所占时间从 \(1/2\) 缩到零，就得到 \(\gamma\)。每一步都保持 \(x\)，并可由伴随公式验证联合连续。因此 \(\Psi\Phi\simeq\operatorname{id}\)。对 \(\Phi\Psi\)，用同样公式缩去 \(\bar h_x*h_x\) 即得 \(\Phi\Psi\simeq\operatorname{id}\)。
\end{solution}

\begin{exercise}\label{g2:ex:square-map-homotopy-fiber}
	令 \(f:\mathbb R\rightarrow\mathbb R\)、\(f(x)=x^2\)。分别求 \(f^{-1}(-1)\) 与 \(f^{-1}(1)\)，并证明对任意 \(b\in\mathbb R\)，\(\operatorname{hofib}_b(f)\) 都可缩。用一个道路提升问题证明 \(f\) 不是 Serre 纤维化。
\end{exercise}
\begin{solution}
	通常纤维分别为空集和两点集 \(\{-1,1\}\)。同伦纤维中的点则为 \((x,\gamma)\)，其中 \(\gamma(0)=x^2\)、\(\gamma(1)=b\)。先把路径线性变到连接两端的线段：
	\[
		\gamma_s(t)=(1-s)\gamma(t)+s\bigl((1-t)x^2+tb\bigr).
	\]
	这保持端点，且联合连续，因而将同伦纤维强形变收缩到由线段路径组成的子空间。该子空间由 \(x\) 唯一参数化，与 \(\mathbb R\) 同胚。再令 \(x_s=(1-s)x\)，并使用路径 \(t\mapsto(1-t)x_s^2+tb\)，就把它缩到 \(x=0\) 对应的一点。因此同伦纤维总是非空且可缩。

	取底道路 \(H(t)=-t\)，并指定在时刻零提升到 \(x=0\)。若存在提升 \(\widetilde H\)，则 \(\widetilde H(t)^2=-t\)，对任意 \(t>0\) 均不可能。\(f\) 连 \(D^0\) 参数的同伦提升性质都不满足，故不是 Serre 纤维化。
\end{solution}

\begin{exercise}\label{g2:ex:fibration-section-splitting}
	设 \(p:E\to B\) 是 Serre 纤维化，\(B\) 与基点纤维 \(F=p^{-1}(b_0)\) 都道路连通，且有截面 \(s:B\to E\)，取 \(e_0=s(b_0)\)。证明对 \(n\ge2\)，有 \(\pi_n(E,e_0)\cong\pi_n(F,e_0)\oplus\pi_n(B,b_0)\)；对 \(n=1\)，相应结论是分裂群扩张。用 Klein 瓶到圆周的投影说明基本群不必是直积。
\end{exercise}
\begin{solution}
	因为 \(p_*s_*=\operatorname{id}\)，每个 \(p_*\) 都满射。长正合列在底空间群处的正合性给出 \(\partial_n=0\)，于是对每个 \(n\ge1\)，都有分裂短正合列
	\[
		1\longrightarrow\pi_n(F)\xrightarrow{i_*}\pi_n(E)\xrightarrow{p_*}\pi_n(B)\longrightarrow1.
	\]
	在交换群次数也可把两端单位群写成 \(0\)。其中 \(i_*\) 的单射还使用了前一连接同态 \(\partial_{n+1}=0\)。对于 \(n\ge2\)，两个子群均交换，映射 \((a,b)\mapsto i_*a+s_*b\) 给出直和同构；任一 \(x\) 唯一写成 \((x-s_*p_*x)+s_*p_*x\)，前一项来自纤维，这证明满射及唯一性。对于 \(n=1\)，仍有唯一分解 \(x=i_*a\,s_*b\)，但 \(s_*b\) 对 \(i_*\pi_1(F)\) 的共轭作用可以非平凡，故得到半直积。

	具体取
	\[
		K=(S^1\times I)/((z,1)\sim(\bar z,0)),\qquad p([z,t])=e^{2\pi\mathrm i t}.
	\]
	在圆周的开弧上用时间坐标作平凡化，跨越接缝时纤维坐标变为 \(\bar z\)，故 \(p\) 是圆周丛；\(s(e^{2\pi\mathrm i t})=[1,t]\) 是截面。记纤维正向环路为 \(a\)，截面环路为 \(b\)。沿时间方向把纤维环路推到接缝，\(z\mapsto\bar z\) 将正向绕行变为负向绕行，所以
	\[
		b^{-1}ab=a^{-1}.
	\]
	上面的分裂短正合列表明 \(a\) 生成一个嵌入的 \(\mathbb Z\)，每个元素唯一写成 \(a^m b^n\)，故
	\[
		\pi_1(K)\cong\langle a,b\mid b^{-1}ab=a^{-1}\rangle\cong\mathbb Z\rtimes\mathbb Z.
	\]
	这里 \(a\) 有无穷阶，\(a\ne a^{-1}\)，所以基本群非交换，与 \(\mathbb Z^2\) 不同。
\end{solution}

\begin{exercise}\label{g2:ex:weak-equivalence-homotopy-fiber}
	设 \(f:(X,x_0)\to(Y,y_0)\) 为道路连通空间之间的带基点映射。证明 \(f\) 是弱同伦等价，当且仅当 \(F_f=\operatorname{hofib}_{y_0}(f)\) 道路连通且所有正次同伦群均为零。解释证明中为什么必须单独处理 \(\pi_0(F_f)\)。
\end{exercise}
\begin{solution}
	以\cref{g2:prop:fibration-replacement}将 \(f\) 写成 \(X\xrightarrow{j}E_f\xrightarrow{p_f}Y\)。\(j\) 是保持基点的同伦等价，所以 \(E_f\) 道路连通，且 \((p_f)_*\) 与 \(f_*\) 经 \(j_*\) 对应。若 \(f_*\) 在所有正次同伦群上均为同构，对 \(n\ge1\) 考察片段
	\[
		\pi_{n+1}(E_f)\longrightarrow\pi_{n+1}(Y)\longrightarrow\pi_n(F_f)
		\longrightarrow\pi_n(E_f)\longrightarrow\pi_n(Y).
	\]
	第一个映射满射，使进入 \(\pi_n(F_f)\) 的映射为零；最后一个映射单射，使 \(\pi_n(F_f)\to\pi_n(E_f)\) 的像为零。正合性于是迫使 \(\pi_n(F_f)=0\)。

	此外，\(E_f\) 道路连通，故\cref{g2:prop:fibration-component-action}说明 \(\pi_1(Y)\) 在 \(\pi_0(F_f)\) 上传递作用。指定分支的稳定子是 \(\operatorname{im}(p_f)_*\)，而 \(\pi_1(X)\to\pi_1(Y)\) 满射使这个稳定子等于整个群。因此指定分支的轨道只有一点，传递性又说明全部 \(\pi_0(F_f)\) 只有这一点。基点变化随即使纤维内任何基点的正次同伦群都为零。

	反过来，若 \(F_f\) 道路连通且正次同伦群全为零，则对 \(n\ge2\)，长正合列直接给出 \(\pi_n(E_f)\cong\pi_n(Y)\)。在一次处，\(\pi_1(F_f)=0\) 保证单射，\(\pi_0(F_f)\) 只有一点保证满射，故也为同构。源与靶都道路连通，零次条件自动满足，遂得 \(f\) 为弱同伦等价。单独检查零次项是必要的：所有正次群均为零的空间仍可能有多个道路分支。
\end{solution}

\begin{exercise}\label{g2:ex:hopf-pullback-degree}
	设 \(g:(S^2,b)\to(S^2,b)\) 是度为 \(d\in\mathbb Z\) 的带基点连续映射。沿 \(g\) 拉回复 Hopf 丛 \(h:S^3\to S^2\)，得到
	\[
		E_d=\{(x,z)\in S^2\times S^3:g(x)=h(z)\},\qquad p_d(x,z)=x.
	\]
	证明 \(E_d\) 道路连通，计算 \(\pi_1(E_d)\)、\(\pi_2(E_d)\)，并比较 \(\pi_k(E_d)\) 与 \(\pi_k(S^2)\)（\(k\ge3\)）。
\end{exercise}
\begin{solution}
	\cref{g2:prop:fibration-pullback}给出纤维为 \(S^1\) 的局部平凡丛。底空间与纤维均道路连通，先提升底空间道路，再在终点纤维内连接，便说明总空间道路连通。

	先确认度与二次同伦群的关系。由\cref{g2:prop:complex-hopf-homotopy}，\(\pi_2(S^2)\cong\mathbb Z\)；由\cref{g2:lem:degree-homotopy-additivity}，度是到 \(\mathbb Z\) 的群同态，而且恒等映射的度为 \(1\)，所以它满射。无限循环群到自身的满同态是同构，故 \(\deg:\pi_2(S^2)\to\mathbb Z\) 是同构。\cref{g2:prop:degree-basic}的复合公式因而说明 \(g_*\) 在此识别下是乘 \(d\)。

	拉回方块在纤维上为恒等，自然性给出 \(\partial^{(d)}_2=\partial^{\mathrm H}_2g_*\)。Hopf 丛的 \(\partial^{\mathrm H}_2\) 是无限循环群之间的同构，故 \(\partial^{(d)}_2\) 的核、像与乘 \(d\) 的核、像相同。长正合列成为
	\[
		0\longrightarrow\pi_2(E_d)\longrightarrow\mathbb Z\xrightarrow{\ \pm d\ }\mathbb Z
		\longrightarrow\pi_1(E_d)\longrightarrow0.
	\]
	正负号取决于顶面约定与球面、圆周生成元的对应，不影响核与像。因此
	\[
		\pi_1(E_d)\cong\mathbb Z/d\mathbb Z,\qquad
		\pi_2(E_d)\cong
		\begin{cases}\mathbb Z,&d=0,\\0,&d\ne0.\end{cases}
	\]
	这里约定 \(\mathbb Z/0\mathbb Z=\mathbb Z\)。当 \(k\ge3\) 时，\(\pi_k(S^1)=\pi_{k-1}(S^1)=0\)，所以 \((p_d)_*:\pi_k(E_d)\to\pi_k(S^2)\) 为同构。底映射的度改变低阶基本群，却不改变这个高阶比较。
\end{solution}

\begin{exercise}\label{g2:ex:hopf-transitions-sections}
	在\cref{g2:prop:hopf-bundles}的两张坐标图上，令 \(w=q_1q_0^{-1}\)、\(v=q_0q_1^{-1}\)，并沿用局部截面 \(s_0,s_1\)。求交集上的右过渡函数 \(g(w)\)，使 \(s_1(w^{-1})=s_0(w)g(w)\)。再证明复、四元数两种 Hopf 丛均无连续整体截面。
\end{exercise}
\begin{solution}
	交集上 \(w\ne0\)。把分母同乘 \(|w|\)，得到
	\[
		s_1(w^{-1})=\frac{(|w|w^{-1},|w|)}{\sqrt{1+|w|^2}}
		=s_0(w)\frac{\bar w}{|w|}.
	\]
	故 \(g(w)=\bar w/|w|\)。四元数情形中右侧第二坐标的核验是 \(w(\bar w/|w|)=|w|\)，不能任意颠倒系数。复情形在 \(|w|=1\) 上得到 \(g(e^{\mathrm{i}\theta})=e^{-\mathrm{i}\theta}\)。

	若 \(s:S^d\rightarrow S^{2d-1}\) 是 \(h_{\mathbb K}\) 的截面，则 \(h_{\mathbb K}s=\operatorname{id}_{S^d}\)。给总空间选基点 \(s(b)\)，这给出 \(h_{\mathbb K*}s_*\) 为 \(\pi_d(S^d,b)\) 上的恒等同态。但是 \(d<2d-1\)，故 \(\pi_d(S^{2d-1})=0\)。于是 \(\pi_d(S^d)\) 必为零，与恒等映射的度为 \(1\) 矛盾。因此两种丛都无连续整体截面，特别地都不是乘积丛。
\end{solution}

\begin{exercise}\label{g2:ex:projective-higher-homotopy}
	设整数 \(m\ge1\)。从单位向量到其复直线的投影 \(p:S^{2m+1}\rightarrow\mathbb{CP}^m\) 出发，证明
	\[
		\pi_1(\mathbb{CP}^m)=0,\qquad\pi_2(\mathbb{CP}^m)\cong\mathbb Z,
	\]
	以及 \(p_*:\pi_k(S^{2m+1})\rightarrow\pi_k(\mathbb{CP}^m)\) 在 \(k\ge3\) 时为同构。由此确定 \(3\le k\le2m\) 时的群，并说明标准包含 \(\mathbb{CP}^1\hookrightarrow\mathbb{CP}^m\) 在 \(\pi_2\) 上的作用。
\end{exercise}
\begin{solution}
	在第 \(j\) 个齐次坐标非零的图上，把该坐标归一为正实数，得到连续单位代表。于是
	\[
		p^{-1}(U_j)\longrightarrow U_j\times S^1,\qquad
		z\longmapsto([z],z_j/|z_j|)
	\]
	为平凡化，逆映射将所选单位代表乘以第二个坐标。因此 \(p\) 是纤维为 \(S^1\) 的局部平凡丛，且底空间是连通球面的连续像，因而道路连通。

	\(S^{2m+1}\) 的 \(\pi_1,\pi_2\) 都为零，\(S^1\) 道路连通；长正合列给出 \(\pi_1(\mathbb{CP}^m)=0\) 以及连接同构 \(\pi_2(\mathbb{CP}^m)\cong\pi_1(S^1)=\mathbb Z\)。在 \(k\ge3\) 的片段中，\(\pi_k(S^1)=\pi_{k-1}(S^1)=0\)，所以 \(p_*\) 为同构；当 \(3\le k\le2m\) 时，总空间的相应群又由胞腔逼近为零。

	把 \(\mathbb C^2\) 作为前两个坐标放入 \(\mathbb C^{m+1}\)，得到复 Hopf 丛到此圆周丛的映射，在基点纤维上是恒等映射。连接同态的自然性给出交换方块
	\[
		\begin{tikzcd}
			\pi_2(\mathbb{CP}^1) \arrow[r] \arrow[d,"\partial"',"\cong"] & \pi_2(\mathbb{CP}^m) \arrow[d,"\partial","\cong"']\\
			\pi_1(S^1) \arrow[r,"\operatorname{id}"'] & \pi_1(S^1).
		\end{tikzcd}
	\]
	因此标准包含在 \(\pi_2\) 上为同构，并保持由同一正向纤维环路指定的生成元。
\end{solution}

\begin{exercise}\label{g2:ex:two-frames-homotopy}
	设 \(n\ge3\)，令
	\[
		V_2(\mathbb R^n)=\{(u,v)\in S^{n-1}\times S^{n-1}:\langle u,v\rangle=0\}.
	\]
	证明取第一向量给出纤维为 \(S^{n-2}\) 的局部平凡丛，并计算 \(V_2(\mathbb R^n)\) 的 \(\pi_1\) 与 \(\pi_2\)。对 \(n=4\)，可把 \(\mathbb R^4\) 识别为 \(\mathbb H\)，考察 \((u,v)\mapsto(u,\bar u v)\)。
\end{exercise}
\begin{solution}
	在 \(u_0\) 附近连续选取正交矩阵 \(Q(u)\)，使第一列为 \(u\)，则 \(v=Q(u)(0,w)\) 唯一对应 \(w\in S^{n-2}\)。这给出所需平凡化。底与纤维均道路连通，利用道路提升再在终点纤维中连接，得到总空间道路连通。

	\(n=3\) 时，映射 \((u,v)\mapsto(u,v,u\times v)\) 是到 \(\operatorname{SO}(3)\) 的同胚，逆映射取前两列。因此 \(\pi_1(V_2(\mathbb R^3))=\mathbb Z/2\mathbb Z\)，\(\pi_2(V_2(\mathbb R^3))=0\)。

	\(n=4\) 时，\(\operatorname{Re}(\bar u v)=\langle u,v\rangle=0\)，且 \(|\bar u v|=1\)。故
	\[
		V_2(\mathbb R^4)\longrightarrow S^3\times S^2,\qquad(u,v)\longmapsto(u,\bar u v)
	\]
	为同胚，逆映射为 \((u,w)\mapsto(u,uw)\)，其中 \(S^2\) 是单位纯虚四元数球面。映射到积空间及其同伦逐坐标确定，因而 \(\pi_j(X\times Y)\cong\pi_j(X)\times\pi_j(Y)\)。现在 \(\pi_1(S^3)=\pi_1(S^2)=\pi_2(S^3)=0\)，而 \(\pi_2(S^2)=\mathbb Z\)，得到 \(\pi_1(V_2(\mathbb R^4))=0\)、\(\pi_2(V_2(\mathbb R^4))=\mathbb Z\)。

	当 \(n\ge5\) 时，纤维 \(S^{n-2}\) 与底 \(S^{n-1}\) 的一、二次同伦群均为零。长正合列立即给出 \(\pi_1(V_2(\mathbb R^n))=\pi_2(V_2(\mathbb R^n))=0\)。
\end{solution}

\begin{exercise}\label{g2:ex:unitary-root-cover}
	设 \(n\ge1\)、\(m\ge2\) 为整数，定义
	\[
		P_m=\{(A,z)\in\operatorname{U}(n)\times S^1:\det A=z^m\},\qquad p(A,z)=A.
	\]
	证明 \(p\) 是连通的 \(m\) 重覆盖；计算以 \((I,1)\) 为基点的 \(\pi_1(P_m)\)、\(\pi_2(P_m)\)，并说明 \(p_*\pi_1(P_m)\) 在 \(\pi_1(\operatorname{U}(n))\cong\mathbb Z\) 中是哪一个子群。
\end{exercise}
\begin{solution}
	\(p\) 是覆盖 \(S^1\rightarrow S^1,z\mapsto z^m\) 沿 \(\det\) 的拉回。在行列式落入一条真开圆弧的邻域上，选定连续实角 \(\theta(A)\)，其 \(m\) 个互不相同的根为
	\[
		\exp\bigl(2\pi\mathrm{i}(\theta(A)+r)/m\bigr),\qquad 0\le r<m.
	\]
	这直接给出 \(m\) 个覆盖叶。

	\(\operatorname{U}(n)\) 道路连通，故任何 \((A,z)\) 都可沿一条底道路的提升连到 \(I\) 上某点。正文给出的基本群生成环路 \(\gamma(t)=\operatorname{diag}(1,\ldots,1,e^{2\pi\mathrm{i}t})\) 从 \((I,1)\) 提升为 \((\gamma(t),e^{2\pi\mathrm{i}t/m})\)，终点是 \((I,e^{2\pi\mathrm{i}/m})\)。重复提升可连到 \(I\) 上所有 \(m\) 个点，所以 \(P_m\) 道路连通。

	行列式绕数为 \(k\) 的环路，其指定提升的根坐标终值为 \(e^{2\pi\mathrm{i}k/m}\)。因此闭合提升恰对应 \(k\in m\mathbb Z\)。由覆盖基本群的提升判据，\(p_*\) 单射且像为 \(m\mathbb Z\)，所以 \(\pi_1(P_m)\cong m\mathbb Z\cong\mathbb Z\)。离散纤维使覆盖在二次同伦群上诱导同构，故 \(\pi_2(P_m)\cong\pi_2(\operatorname{U}(n))=0\)。
\end{solution}

\begin{exercise}\label{g2:ex:rp2-higher-action}
	利用对径覆盖 \(q:S^2\rightarrow\mathbb{RP}^2\)，证明 \(\pi_2(\mathbb{RP}^2)\cong\mathbb Z\)，并计算 \(\pi_1(\mathbb{RP}^2)\cong\mathbb Z/2\mathbb Z\) 对它的作用。
\end{exercise}
\begin{solution}
	小球面邻域可以选得与其对径邻域不交，故 \(q\) 为二重覆盖；\(S^2\) 单连通。因此底基本群由环路提升的两个可能终点决定，恰为 \(\mathbb Z/2\mathbb Z\)。长正合列及离散纤维给出 \(q_*:\pi_2(S^2)\rightarrow\pi_2(\mathbb{RP}^2)\) 的同构，故目标为 \(\mathbb Z\)。

	先把 \(\pi_2(S^2)\) 的整数具体识别为度。\cref{g2:lem:degree-homotopy-additivity}给出群同态 \(\deg:\pi_2(S^2)\rightarrow\mathbb Z\)，其在恒等映射上取 \(1\)，所以满射；源又已由 Hopf 丛证明为无限循环群，故此同态为同构。

	取球面基点 \(x_0\)。底空间非平凡环路从 \(x_0\) 出发提升后终止于 \(-x_0\)。把定义基本群作用的移动基点同伦提升到球面，就将球面代表 \(f\) 换为对径映射 \(a(x)=-x\) 与 \(f\) 的复合，并沿提升道路把基点识别回来。球面单连通保证这次基点识别只依赖端点；在普通同伦意义下移动基点不改变度。因此该作用在度上是乘以 \(\deg a\)。由\cref{g2:prop:reflection-antipodal-degree}，\(\deg a=(-1)^3=-1\)，所以非单位基本群元素把整数 \(k\) 送到 \(-k\)。
\end{solution}


% !TeX root = ../../Book2.tex
% 纲要标题：## 第9章　Hurewicz、Whitehead 与同伦类型
% 纲要位置：几何学计划纲要.md，第 1334 行。
% 正文状态：未撰写；本文件目前是框架。

% 纲要内容开始
% 
% > 本章目标与证明范围：证明“同伦首先在哪一次被同调看到”及“什么条件下同调能检测等价”。相对论证和 CW 假设单列，避免同调同构的无条件滥用。
% 
% 纲要内容结束

\chapter{Hurewicz、Whitehead 与同伦类型}
\label{g2:ch:09}
\BookChapterNumber{g2:ch:09}

\begin{chapterabstract}
待撰写
\end{chapterabstract}

\begin{learninggoals}
待撰写
\end{learninggoals}

待撰写

% !TeX root = ../../Book2.tex
% 纲要标题：### 9.1 连通度与相对胞腔结构
% 纲要位置：几何学计划纲要.md，第 1338 行。
% 正文状态：初稿；尚未正式审核。

% 纲要内容开始
% 
% * 连通度和胞腔的消去
% * 相对逼近及附着维数
% * 同伦群为零的几何含义
% 
% 纲要内容结束

\section{连通度与相对胞腔结构}
\label{g2:sec:0901}

球面有一个零胞腔和一个顶维胞腔，这使低维球面映射可以被胞腔逼近压成常值。现在把问题反过来：如果一个空间的低阶同伦群都为零，能否给它找一个没有相应低维胞腔的模型？为此需要分清两种操作：把给定映射的像移入子空间，以及通过附着新胞腔改进整个空间的模型。

\subsection{连通度和胞腔的消去}
\label{g2:subsec:0901:01}

\begin{definition}[连通度]\label{g2:def:connectivity}
	设 \(n\ge0\) 为整数。非空空间 \(X\) 称为\term[lian tong du]{\(n\) 连通}[\(n\)-connected]，如果它道路连通，且 \(\pi_i(X,x)=0\) 对 \(1\le i\le n\) 成立。空间对 \((X,A)\) 称为 \(n\) 连通，如果 \(X\) 的每个道路分支都与 \(A\) 相交，并且对每个 \(a\in A\) 有
	\[
		\pi_i(X,A,a)=0\qquad(1\le i\le n).
	\]
	其中相对一阶的等号表示尖点集只有指定元素。连续映射 \(f:A\rightarrow X\) 称为 \(n\) 连通，如果它的映射柱对 \((M_f,A\times\{0\})\) 为 \(n\) 连通。
\end{definition}

零连通就是道路连通，一连通就是单连通。由\cref{g2:prop:basepoint-action}，在道路连通空间内，绝对群的消失与选哪个基点无关。由\cref{g2:thm:pair-homotopy-sequence}，当 \(n\ge1\) 时，映射 \(f\) 为 \(n\) 连通，等价于 \(f_*:\pi_0(A)\rightarrow\pi_0(X)\) 为双射，且在每个基点处满足
\[
	f_*:\pi_i(A,a)\longrightarrow\pi_i(X,f(a))
	\quad\begin{cases}
		\text{为同构},&1\le i<n,\\
		\text{为满射},&i=n.
	\end{cases}
\]
这里映射柱收缩到 \(X\)，所以它的绝对同伦群可用 \(X\) 的同伦群替换。低阶分支条件也有具体含义：相对道路类为零保证同一 \(X\) 分支内的 \(A\) 中两点已能在 \(A\) 内连通。

\begin{lemma}[圆盘压缩]\label{g2:lem:disk-compression}
	设 \(B\subset Y\)、\(m\ge1\)，并且 \(\pi_m(Y,B,b)=0\) 对所有 \(b\in B\) 成立。则每个连续映射 \(u:(D^m,S^{m-1})\rightarrow(Y,B)\) 都可固定整个边界地同伦到一个取值于 \(B\) 的映射。反之，这种压缩性质蕴含上述相对同伦群或同伦集为零。当 \(m=0\) 时，相应的压缩性质等价于 \(Y\) 的每个道路分支与 \(B\) 相交。
\end{lemma}
\begin{proof}
	在 \(S^{m-1}\) 上选一点 \(s_0\)，令 \(b=u(s_0)\)。\cref{g2:def:relative-homotopy}的圆盘模型使 \(u\) 表示 \(\pi_m(Y,B,b)\) 的元素。它为零，故存在 \(H:D^m\times I\rightarrow Y\)，满足 \(H(-,0)=u\)、\(H(-,1)=b\)，且 \(H(S^{m-1}\times I)\subset B\)。这时整个边界可以在 \(B\) 中移动；要把边界固定住，就把这段运动反向放进圆盘外侧的一层领圈。

	写 \(r=\lVert z\rVert\)、\(a_s=1-s/2\)，定义
	\[
		K(z,s)=\begin{cases}
			H(z/a_s,s),&r\le a_s,\\
			H(z/r,2(1-r)),&a_s\le r\le1.
		\end{cases}
	\]
	在 \(r=a_s\) 处，两式的时间参数都为 \(s\)，因而相合。两块是闭集，分母在各自使用的区域内非零，故粘合给出连续同伦。\(K(z,0)=u(z)\)，而在 \(r=1\) 上恒有 \(K(z,s)=H(z,0)=u(z)\)。当 \(s=1\) 时，内半圆盘映到 \(b\)，外半领圈映到 \(B\)，所以整个终映射取值于 \(B\)。

	反过来，压缩后的映射 \(v:D^m\rightarrow B\) 可在 \(B\) 内收缩到 \(v(s_0)\)：取 \(v((1-t)z+ts_0)\)。此同伦固定 \(s_0\)，其边界始终在 \(B\) 中，故为相对零伦。对于 \(m=0\)，一个圆盘就是一点，将它移入 \(B\) 正是选一条通往 \(B\) 的道路。
\end{proof}

领圈保留的是原来指定的边界值。这个细节使压缩可以在许多已经粘合的胞腔上同时进行：公共边界不动，各圆盘上的同伦就能继续粘在一起。

\begin{proposition}[附着维数与连通度]\label{g2:prop:cell-dimension-connectivity}
	设 \((X,A)\) 是 CW 对，\(n\ge0\)，且 \(X\setminus A\) 的全部胞腔维数均大于 \(n\)。则 \((X,A)\) 为 \(n\) 连通。特别地，若 \(n\ge1\)，包含 \(A\hookrightarrow X\) 在 \(\pi_i\) 上对 \(i<n\) 诱导同构，对 \(i=n\) 诱导满射，并在道路分支上诱导双射。
\end{proposition}
\begin{proof}
	先取 \(1\le i\le n\) 的圆盘代表 \((D^i,S^{i-1},s_0)\rightarrow(X,A,a)\)。可先在 \(A\) 内沿道路把 \(a\) 移到其所在胞腔的某个顶点，再用\cref{g2:prop:basepoint-action}运输基点；群是否为零不受影响。给圆盘与边界取以 \(s_0\) 为顶点的 CW 结构，由\cref{g2:cor:cellular-pairs}作空间对的胞腔逼近并固定 \(s_0\)。终映射的像位于 \(X^i\)，而维数假设给出 \(X^i\subset A\)，故该相对类为零。零维同样由胞腔逼近得出：每个点都可沿道路移到 \(X^0\subset A\)。最后应用空间对的同伦长正合列及其分支项即可。
\end{proof}

例如，给 \(A\) 附着四维及更高维胞腔，不改变基本群和第二同伦群，并且使 \(\pi_3(A)\rightarrow\pi_3(X)\) 满射。四胞腔的边界是三球面，所以它已经可能消去 \(\pi_3\) 中的元素。临界维数的“满射”记录的正是这种可能的消去。

\subsection{相对逼近及附着维数}
\label{g2:subsec:0901:02}

一个映射 \(f:A\rightarrow X\) 在同伦群上有两类缺口：\(A\) 中某些球面到了 \(X\) 才能填充，或者 \(X\) 中某些球面尚未由 \(A\) 表示。前一类可在 \(A\) 中附加相应的填充圆盘，后一类可把球面直接添进模型。每轮只处理一个维数，便能保留此前已经得到的同构。

\begin{theorem}[相对 CW 逼近]\label{g2:thm:relative-cw-approximation}
	设 \(A\) 是带有零胞腔 \(a_0\) 的道路连通 CW 复形，\(X\) 是道路连通空间，\(n\ge0\)，且 \(f:(A,a_0)\rightarrow(X,x_0)\) 为 \(n\) 连通映射。则存在包含 \(A\) 为子复形的 CW 复形 \(Z\)，以及延拓 \(q:Z\rightarrow X\)，使得
	\begin{enumerate}
		\item \(Z\setminus A\) 的所有胞腔维数均至少为 \(n+1\)，其附着映射可取保持基点；
		\item \(q|_A=f\)，且 \(q_*:\pi_i(Z,a_0)\rightarrow\pi_i(X,x_0)\) 对每个 \(i\ge1\) 都是同构。
	\end{enumerate}
	这样的 \(q\) 称为一个\term[CW bi jin]{CW 逼近}[CW approximation]。它在道路分支上也是双射，因而是\secref{g2:sec:0903}所定义的弱同伦等价。
\end{theorem}
\begin{proof}
	先设 \(n\ge1\)，从 \(Z_n=A\)、\(q_n=f\) 出发。归纳假定 \(q_k:Z_k\rightarrow X\) 为 \(k\) 连通。对于 \(\ker(q_{k*}:\pi_k(Z_k,a_0)\rightarrow\pi_k(X,x_0))\) 的每个元素，选择一个带基点胞腔代表 \(\alpha:S^k\rightarrow Z_k\)。这里用\cref{g2:thm:cellular-approximation}固定基点，所以 \(\alpha(S^k)\subset Z_k^k\)。沿每个 \(\alpha\) 附着一个 \((k+1)\) 胞腔；由于 \(q_k\alpha\) 零伦，\(q_k\) 可以用这个零伦延拓到所添圆盘。

	再对 \(\pi_{k+1}(X,x_0)\) 的每个元素选一个带基点球面代表，在 \(a_0\) 处以常值边界附着一个 \((k+1)\) 胞腔，得到一个相应的 \((k+1)\) 球面，并把它映成所选代表。记完成两种附着后的空间为 \(Z_{k+1}\)，延拓为 \(q_{k+1}\)。其在 \(\pi_{k+1}\) 上满射，因为所需的球面代表已逐一添入。

	由\cref{g2:prop:cell-dimension-connectivity}，\(Z_k\hookrightarrow Z_{k+1}\) 在 \(\pi_i\)、\(i<k\) 上为同构，在 \(\pi_k\) 上满射。于是 \(q_{k+1}\) 保留了低于 \(k\) 维的同构，也在 \(\pi_k\) 上保持满射。若 \(\beta\in\pi_k(Z_{k+1})\) 被 \(q_{k+1*}\) 送到零，可把它提升为 \(\alpha\in\pi_k(Z_k)\)。这个 \(\alpha\) 属于 \(\ker q_{k*}\)，第一种附着已经给它加上填充圆盘，故 \(\beta=0\)。因此 \(q_{k+1}\) 为 \((k+1)\) 连通。

	当 \(n=0\) 时，先在 \(a_0\) 处添圆周，分别映成 \(\pi_1(X,x_0)\) 的各个元素。所得 \(Z_1\rightarrow X\) 在基本群上满射，且源、靶仍道路连通，故从 \(k=1\) 接入上述归纳。这一步只添一胞腔，符合最低维数的要求。

	令 \(Z=\bigcup_{k\ge n}Z_k\)，在所有胞腔上取 CW 弱拓扑；在 \(n=0\) 情形把刚才的第一步也包含在内。附着映射均为胞腔映射，因而原来 \(A\) 的各维胞腔与新胞腔共同给出 CW 结构。相容的 \(q_k\) 在各特征圆盘上连续，故粘成连续映射 \(q:Z\rightarrow X\)。任一球面映射以及任一球面同伦的像都紧，由\cref{g2:thm:finite-subcomplex}落在有限子复形中，也就落在某个 \(Z_k\) 中。对于固定的 \(i\)，取 \(k>i\)；\(q_k\) 在 \(\pi_i\) 上已经同构，后续只添维数大于 \(i+1\) 的胞腔，保持这个同构。紧像论证遂给出 \(q_*\) 在 \(\pi_i\) 上的双射。所有附着均保持 \(a_0\)，结论成立。
\end{proof}

构造中的两种圆盘作用可用下面两幅交换图表示：左图用圆盘填掉核中的球面，右图把靶中的球面代表添入模型。两图中的 \(Z_{k+1}\) 表示完成本轮附着后的空间。
\[
	\begin{tikzcd}[column sep=large]
		S^k \arrow[r,"\alpha"] \arrow[d,hook] & Z_k \arrow[d,hook] & S^{k+1} \arrow[r,hook] \arrow[d,"\beta"'] & Z_{k+1} \arrow[d,"q_{k+1}"]\\
		D^{k+1} \arrow[r,"\text{特征映射}"'] & Z_{k+1} & X \arrow[r,equal] & X .
	\end{tikzcd}
\]

\begin{corollary}[稀疏胞腔模型]\label{g2:cor:sparse-cw-model}
	设 \(n\ge2\)。若带有零胞腔基点的 CW 复形 \(X\) 为 \((n-1)\) 连通，则存在带基点同伦等价 \(q:Y\rightarrow X\)，使 \(Y^0=\{*\}\)、\(Y^{n-1}=\{*\}\)，且所有正维附着映射保持基点。更一般地，若 \((X,A)\) 为 \((n-1)\) 连通 CW 对且 \(A\) 道路连通，则存在 CW 对 \((Z,A)\) 及相对 \(A\) 的同伦等价 \((Z,A)\rightarrow(X,A)\)，其中 \(Z\setminus A\) 的胞腔维数均至少为 \(n\)。
\end{corollary}
\begin{proof}
	第一项对 \(\{*\}\rightarrow X\) 应用\cref{g2:thm:relative-cw-approximation}，第二项对 \(A\hookrightarrow X\) 应用该定理，均得到所需胞腔维数及弱同伦等价。由\cref{g2:thm:whitehead,g2:cor:relative-whitehead}，它们分别为带基点、相对 \(A\) 的同伦等价。这两个结果将在\secref{g2:sec:0903}由圆盘压缩和逐骨架同伦延拓证明；其证明只用本节的\cref{g2:lem:disk-compression}及第1、8章的结果。
\end{proof}

这里“消去低维胞腔”落实为更换同伦等价的 CW 模型。原来的分解仍可能含有许多一胞腔或二胞腔；新模型把它们承载的可收缩信息吸收进同伦等价中。相对版本则逐点保留 \(A\)，使已经指定的边界和映射仍能继续使用。上述逐维构造与相对模型可参照\cite[第4.1节，CW Approximation]{Hatcher2002}。

\subsection{同伦群为零的几何含义}
\label{g2:subsec:0901:03}

绝对群的消失对应球面的填充，相对群的消失对应圆盘的压缩。这两种图像的维数相差一：\(\pi_i(X)=0\) 让 \(S^i\) 填成 \(D^{i+1}\)，而 \(\pi_i(X,A)=0\) 让一个已经在 \(X\) 中的 \(D^i\) 固定边界地移入 \(A\)。

\begin{proposition}[逐胞腔延拓]\label{g2:prop:connectivity-extension}
	设 \((K,L)\) 为 CW 对，\(Y\) 为非空道路连通空间。若每个 \(K\setminus L\) 中的 \(m\) 胞腔在 \(m\ge2\) 时都满足 \(\pi_{m-1}(Y,y)=0\)，则每个连续映射 \(f:L\rightarrow Y\) 可延拓为 \(K\rightarrow Y\)。特别地，若 \(Y\) 为 \(r\) 连通且 \(K\setminus L\) 的胞腔维数不超过 \(r+1\)，则这种延拓存在。
\end{proposition}
\begin{proof}
	在新增零胞腔上任取 \(Y\) 中的点。对一胞腔，边界是两点，道路连通性提供连接两端像的道路。假设已经定义到相对 \((m-1)\) 骨架；对一个 \(m\) 胞腔，将其附着映射与已定义映射复合，得到 \(S^{m-1}\rightarrow Y\)。当 \(m\ge2\) 时，它在所选边界基点处代表零类，因而有零伦。将零伦的常值末端压成一点，就得到 \(D^m\rightarrow Y\) 的填充。所有圆盘填充与旧映射在边界一致，所以逐层粘合。最终映射在每个特征圆盘上连续，CW 弱拓扑保证整体连续。
\end{proof}

例如，\(S^3\) 为二连通，所以从三维 CW 复形的任意子复形到 \(S^3\) 的映射都能延拓。若再附着四胞腔，则需要填充的是三球面映射；恒等映射 \(S^3\rightarrow S^3\) 的度为 \(1\)，已经不能填成四球内的映射。连通度恰好指出逐胞腔延拓在哪一维首次可能遇到障碍。

% 纲要条目排为编号小节标题，并进入目录与书签；说明段仍保留于纲要注释中。

% !TeX root = ../../Book2.tex
% 纲要标题：### 9.2 Hurewicz 定理
% 纲要位置：几何学计划纲要.md，第 1344 行。
% 正文状态：初稿；尚未正式审核。

% 纲要内容开始
% 
% * Hurewicz 映射的构造
% * 一维交换化与高维首个非零群
% * 绝对定理的证明
% * 相对版本及单连通条件
% 
% 纲要内容结束

\section{Hurewicz 定理}
\label{g2:sec:0902}

球面映射 \(f:S^n\to X\) 一方面给出同伦类，另一方面把球面的基本同调类送成 \(X\) 中的一个 \(n\) 维圈。两种信息何时相同？第8章的 Hopf 映射说明，高阶同伦信息可以在同调中全部消失；而如果所有较低维的球面都已能够填充，首个可能非零的同伦群恰好能由同调计算。Hurewicz 定理把这里的“较低维”写成准确的连通度条件，并给出一张自然的比较映射。

\subsection{Hurewicz 映射的构造}
\label{g2:subsec:0902:01}

固定定向球面 \(S^n\) 的生成元 \(\iota_n\in H_n(S^n;\mathbb Z)\)，其中 \(n\ge1\)。这里的生成元来自\cref{g2:thm:sphere-homology}；改变球面定向会同时改变它的符号。以下同调均取整数系数。

\begin{definition}[Hurewicz 映射]\label{g2:def:hurewicz-map}
	设 \((X,x_0)\) 为带基点空间，\(n\ge1\)。映射
	\[
		h_n:\pi_n(X,x_0)\longrightarrow H_n(X;\mathbb Z),\qquad [f]\longmapsto f_*(\iota_n)
	\]
	称为 \(n\) 次\term[Hurewiczyingshe]{Hurewicz 映射}[Hurewicz map]。对带基点空间对 \((X,A,x_0)\) 及 \(n\ge2\)，选择相对球的生成元 \(u_n\in H_n(D^n,S^{n-1};\mathbb Z)\)，并定义
	\[
		h_n:\pi_n(X,A,x_0)\longrightarrow H_n(X,A;\mathbb Z),\qquad [F]\longmapsto F_*(u_n).
	\]
	\symidx{\(h_n\)：Hurewicz 映射}相对类使用第8章的立方体约定：顶面进入 \(A\)，其余面映到 \(x_0\)。相对球的定向取为，使连接同态把 \(u_n\) 送到这个顶面在压去其边界后所表示的 \(\iota_{n-1}\)。
\end{definition}

同伦不变性保证定义不依赖代表。相对构造也有同样性质，因为相对同伦始终把边界留在 \(A\) 内。这张映射保留了球面所携带的定向圈，却忘掉该圈必须由一个球面参数化这一要求。

\begin{proposition}[自然性与加法]\label{g2:prop:hurewicz-naturality}
	对任意带基点空间 \((X,x_0)\)，\(h_n:\pi_n(X,x_0)\to H_n(X;\mathbb Z)\) 在 \(n\ge1\) 时是群同态；相对 Hurewicz 映射在 \(n\ge2\) 时也是群同态。两者对带基点映射及空间对映射自然，并与空间对长正合列中的映射相容。特别地，采用上述定向时有
	\[
		\begin{tikzcd}[column sep=large,row sep=large]
			\pi_n(X,A,x_0) \arrow[r,"\partial"] \arrow[d,"h_n"'] & \pi_{n-1}(A,x_0) \arrow[d,"h_{n-1}"] \\
			H_n(X,A) \arrow[r,"\delta"'] & H_{n-1}(A).
		\end{tikzcd}
	\]
\end{proposition}
\begin{proof}
	把球面沿赤道压成一点，得到拼接映射 \(c:S^n\to S^n\vee S^n\)。在第8章的加法约定下，\([f]+[g]\) 由 \((f\vee g)c\) 代表。两个球面各自向原球面的投影与 \(c\) 复合，都把一个半球的内部保持定向地送到球面去掉基点的部分，因而次数为 \(1\)。第3章的胞腔同调给出 \(H_n(S^n\vee S^n)\cong\mathbb Z\oplus\mathbb Z\)，所以
	\[
		c_*(\iota_n)=(\iota_n,\iota_n),\qquad h_n([f]+[g])=f_*(\iota_n)+g_*(\iota_n).
	\]
	\(n=1\) 时把左边的加法读作环路拼接，论证仍成立。对相对立方体，沿第一个坐标拼接时，两个半立方体继承相同方向；它们的公共面都映到基点，所代表的相对基本链相加便是整个立方体的基本链。因此相对映射同样保持群运算。

	若 \(a:X\to Y\)，则 \((af)_*(\iota_n)=a_*f_*(\iota_n)\)，这就是自然性；空间对情形完全相同。最后，连接同态在链上通过取边界定义。相对基本链的边界中，除顶面之外的各面均映到基点，取约化同调后为零，而顶面正代表 \(\partial[F]\)。所选定向使这一项带正号，所以 \(\delta h_n[F]=h_{n-1}\partial[F]\)。其余相容方块由包含和取商的自然性得到。
\end{proof}

例如，在 \(X=S^n\) 时，定义直接给出
\[
	h_n([f])=\deg(f)\,\iota_n.
\]
因此，证明球面的 Hurewicz 映射为同构，恰好等于证明球面自映射由次数分类。到目前为止，我们知道次数是同伦不变量且保持加法；它的单射性仍需要证明。

\subsection{一维交换化与高维首个非零群}
\label{g2:subsec:0902:02}

环路的拼接一般不交换，奇异同调群却是阿贝尔群，因此 \(h_1\) 必定忘掉交换子。问题是它是否还忘掉别的信息。

\begin{theorem}[一维 Hurewicz 定理]\label{g2:thm:hurewicz-one}
	设 \(X\) 为道路连通空间，\(x_0\in X\)。则 \(h_1\) 满射，且
	\[
		\ker h_1=[\pi_1(X,x_0),\pi_1(X,x_0)].
	\]
	因而它诱导自然同构
	\[
		\pi_1(X,x_0)^{\mathrm{ab}}\xrightarrow{\ \cong\ }H_1(X;\mathbb Z),
	\]
	其中右上标 \(\mathrm{ab}\) 表示对全部交换子生成的正规子群取商。
\end{theorem}
\begin{proof}
	我们把任意一次圈转成基点环路，并核对二维边界恰好成为交换化后的关系。对每个点 \(p\in X\) 选一条从 \(x_0\) 到 \(p\) 的道路 \(\gamma_p\)，其中 \(\gamma_{x_0}\) 为常值道路。若奇异一维单形 \(\sigma\) 从 \(p\) 到 \(q\)，令
	\[
		L(\sigma)=[\gamma_p*\sigma*\overline\gamma_q]\in\pi_1(X,x_0)^{\mathrm{ab}},
	\]
	并把 \(L\) 加性延拓到奇异一维链群。对于一个奇异二单形 \(\tau\)，记其三条边为 \(\tau_{01},\tau_{12},\tau_{02}\)。三角形内部给出定端点同伦 \(\tau_{01}*\tau_{12}\simeq\tau_{02}\)，故
	\[
		L(\partial\tau)=L(\tau_{12})-L(\tau_{02})+L(\tau_{01})=0.
	\]
	这里沿公共顶点来回走的 \(\overline\gamma_{p_1}*\gamma_{p_1}\) 可以收缩，而交换化允许按链的加法排列三项。因此 \(L\) 在圈群上诱导同态 \(\ell:H_1(X)\to\pi_1(X,x_0)^{\mathrm{ab}}\)。

	为核对它是逆映射，先回忆道路在奇异链中的拼接关系。若 \(\alpha,\beta\) 可以首尾相接，把三角形的三个顶点依次映到参数区间的 \(0,1/2,1\)，在三角形上线性延拓，再复合 \(\alpha*\beta\)，便得到一个二单形，其边界是 \(\beta-(\alpha*\beta)+\alpha\)。因此拼接与链的相加模边界一致；反向道路与原道路的和模边界为零，因为来回道路定端点零伦，而常值一单形本身是常值二单形的边界。

	现在令 \(z=\sum_j a_j\sigma_j\) 为一次圈，\(\sigma_j:p_j\to q_j\)。上述关系说明 \(h_1\ell[z]\) 由
	\[
		\sum_j a_j(\gamma_{p_j}+\sigma_j-\gamma_{q_j})
	\]
	代表。\(\partial z=0\) 恰好说，在每个顶点处出发系数和与到达系数和相等，所以所有 \(\gamma_p\) 项逐点抵消，余下就是 \(z\)。另一方面，若 \(\sigma\) 已是 \(x_0\) 处的环路，则 \(L(\sigma)=[\sigma]\)，故 \(\ell h_1\) 在交换化上为恒等。两者互逆，结论成立；同构的自然性来自 \(h_1\) 的自然性，并不依赖辅助道路的选择。
\end{proof}

在高维，\(\pi_n\) 本来已经交换；需要排除的是较低维的同伦障碍。以下定理的“首个”包括 \(\pi_n\) 仍为零的情况，因而也能用来逐维检测连通度。

\begin{theorem}[Hurewicz 定理]\label{g2:thm:hurewicz-absolute}
	设 \(X\) 为 \((n-1)\) 连通的 CW 复形，\(n\ge2\)，基点取为一个顶点。则
	\[
		\widetilde H_i(X;\mathbb Z)=0\quad(0\le i<n),\qquad
		h_n:\pi_n(X,x_0)\xrightarrow{\ \cong\ }H_n(X;\mathbb Z).
	\]
	同样的结论适用于具有该 CW 同伦型的带基点空间。
\end{theorem}

证明放在下一小节。它将把两种群都写成“\(n\) 维胞腔生成，\((n+1)\) 维胞腔给出关系”的商群，并证明比较映射把相同的生成元与关系对应起来。

\begin{example}[两个次数的区别]\label{g2:exm:hurewicz-first-degree}
	二维球面单连通，定理给出 \(\pi_2(S^2)\cong H_2(S^2)=\mathbb Z\)。对于三次同伦群，\(S^2\) 不满足所需的二连通条件；事实上 \(H_3(S^2)=0\)，而第8章已经证明 Hopf 映射在 \(\pi_3(S^2)\) 中有无穷阶。另一方面，\(S^3\) 二连通，所以 \(\pi_3(S^3)\cong\mathbb Z\)；再结合\cref{g2:prop:complex-hopf-homotopy}中 Hopf 丛诱导的同构 \(\pi_3(S^3)\to\pi_3(S^2)\)，得到 \(\pi_3(S^2)\cong\mathbb Z\)，且 Hopf 映射是生成元。这也完成第8章留下的生成元识别。
\end{example}

\subsection{绝对定理的证明}
\label{g2:subsec:0902:03}

胞腔链能直接计算同调；要证明它也计算首个同伦群，就要控制一个球面经过胞腔附着后何时变得零伦。下面使用的一般同伦切除定理在\cref{g2:thm:homotopy-excision}给出：若 CW 复形 \(Z=U\cup V\) 的交集 \(C=U\cap V\) 非空道路连通，且 \((U,C)\)、\((V,C)\) 分别为 \(r\)、\(s\) 连通，其中 \(r,s\ge0\)，则
\[
	\pi_i(U,C)\longrightarrow\pi_i(Z,V)
\]
在 \(i<r+s\) 时为同构，在 \(i=r+s\) 时为满射。这个一般定理作为外部输入，其精确版本与来源见该处；下面逐项证明本节所需的推论。胞腔证明的组织可对照\cite[第4.2节，Excision for Homotopy Groups]{Hatcher2002}及\cite[第15章，第1节]{May1999}。

\begin{lemma}[球面次数与同伦类]\label{g2:lem:sphere-degree-classification}
	对每个 \(n\ge1\)，次数同态
	\[
		\deg:\pi_n(S^n)\longrightarrow\mathbb Z
	\]
	是同构，且恒等映射的同伦类对应 \(1\)。
\end{lemma}
\begin{proof}
	\(n=1\) 时，\cref{g2:prop:circle-fundamental-group}的绕数给出 \(\pi_1(S^1)\cong\mathbb Z\)，生成元是恒等环路；次数在这个生成元上为 \(1\)，故结论成立。\(n=2\) 时，\cref{g2:prop:complex-hopf-homotopy}用 Hopf 丛的长正合列证明了 \(\pi_2(S^2)\cong\mathbb Z\)。次数是到 \(\mathbb Z\) 的同态，又把恒等映射送到 \(1\)，所以满射；无限循环群到 \(\mathbb Z\) 的满射必为同构。

	设 \(k\ge2\)。把 \(S^{k+1}\) 写成沿赤道 \(S^k\) 粘合的两个闭锥 \(C_+S^k\)、\(C_-S^k\)。由\cref{g2:prop:sphere-connectivity}，\(S^k\) 为 \((k-1)\) 连通；空间对长正合列及锥的可缩性说明两个锥对均为 \(k\) 连通。切除定理在次数 \(k+1<2k\) 给出中间的同构：
	\[
		\pi_k(S^k)\xrightarrow[\partial^{-1}]{\ \cong\ }
		\pi_{k+1}(C_+S^k,S^k)
		\xrightarrow{\ \cong\ }
		\pi_{k+1}(S^{k+1},C_-S^k)
		\xleftarrow{\ \cong\ }\pi_{k+1}(S^{k+1}).
	\]
	外侧同构同样来自空间对序列。把球面映射 \(f\) 分别延为两个锥上的锥映射，便可在这个图中追踪它的相对代表；合起来正是悬挂 \(Sf\)。特别地，恒等映射被送到悬挂球面的恒等映射，定向同时按这两个半球选择。于是若 \(\pi_k(S^k)\) 由恒等类无限循环生成，\(\pi_{k+1}(S^{k+1})\) 也如此。由 \(k=2\) 归纳，所有这些群都由恒等类生成；次数将生成元送到 \(1\)，所以它们的次数同态均为同构。
\end{proof}

\begin{lemma}[球面楔和的首个同伦群]\label{g2:lem:wedge-sphere-first-homotopy}
	设 \(n\ge2\)，\(W=\bigvee_{\alpha\in\Lambda}S^n_\alpha\) 取 CW 拓扑，\(j_\alpha:S^n_\alpha\to W\) 为各包含。则
	\[
		\bigoplus_{\alpha\in\Lambda}\mathbb Z\xrightarrow{\ \cong\ }\pi_n(W),\qquad
		(m_\alpha)_\alpha\longmapsto\sum_\alpha m_\alpha[j_\alpha].
	\]
	若 \(p_\alpha:W\to S^n_\alpha\) 压去其他球面，则 \([f]\) 的第 \(\alpha\) 个坐标为 \(\deg(p_\alpha f)\)。在这些基下，\(h_n:\pi_n(W)\to H_n(W;\mathbb Z)\) 为恒等矩阵。
\end{lemma}
\begin{proof}
	先设 \(\Lambda\) 有限。有限乘积 \(P=\prod_\alpha S^n_\alpha\) 的乘积胞腔分解中，\(W\) 正是 \(n\) 骨架，余下胞腔最早出现在维数 \(2n\)。胞腔逼近说明 \((P,W)\) 为 \((2n-1)\) 连通。由于 \(n+1\le2n-1\)，空间对序列在 \(\pi_n(W)\to\pi_n(P)\) 两旁的相对群均为零，所以这个映射是同构。到乘积的球面映射及其定基点同伦逐坐标对应，拼接也逐坐标进行，故
	\[
		\pi_n(P)\cong\prod_\alpha\pi_n(S^n_\alpha)\cong\bigoplus_\alpha\mathbb Z.
	\]
	上一引理把每个坐标识别为次数，便得到所述基和坐标公式。

	若 \(\Lambda\) 无限，\cref{g2:thm:finite-subcomplex}保证任意球面映射及任意球面同伦的像都落入有限子复形。这里有限子复形就是有限多个球面的楔和。因此有限情形分别给出满射性与单射性，所有坐标也只有有限个非零。最后，胞腔同调把 \((j_\alpha)_*(\iota_n)\) 识别为第 \(\alpha\) 个胞腔基，而 \(h_n[j_\alpha]=(j_\alpha)_*(\iota_n)\)，故比较矩阵为恒等。
\end{proof}

接下来需要比较相对球面与压去子空间后的球面。普通同调对 CW 对可在所有次数这样比较，同伦群则只能在连通度保证的范围内这样做。

\begin{lemma}[压去子复形的同伦比较]\label{g2:lem:relative-quotient-homotopy}
	设 \((X,A)\) 为 CW 对，\(A\ne\varnothing\)，基点取在 \(A\) 的顶点中。若 \((X,A)\) 为 \(r\) 连通、\(A\) 为 \(s\) 连通，其中 \(r,s\ge1\)，则商映射诱导
	\[
		q_*:\pi_i(X,A)\longrightarrow\pi_i(X/A)
	\]
	在 \(1\le i\le r+s\) 时为同构，在 \(i=r+s+1\) 时为满射；次数一的相对项按尖点集理解。
\end{lemma}
\begin{proof}
	沿 \(A\subset X\) 粘上约化锥 \(CA=A\times I/(A\times\{1\}\cup\{x_0\}\times I)\)，得到 \(Z=X\cup_A CA\)。\(CA\) 可保持基点收缩。空间对序列说明 \((CA,A)\) 为 \((s+1)\) 连通，于是同伦切除给出
	\[
		\pi_i(X,A)\longrightarrow\pi_i(Z,CA)
	\]
	在 \(i<r+s+1\) 时同构，在等号处满射。由于 \(CA\) 可缩，\(\pi_i(Z)\to\pi_i(Z,CA)\) 是同构。

	再说明压去 \(CA\) 确实给出同伦等价 \(Z\to Z/CA=X/A\)。由 CW 对的同伦延拓性质（\cref{g2:thm:relative-CW-HEP}），\(CA\) 内固定基点的收缩可以延为 \(Z\) 上从恒等映射出发的同伦 \(H_t\)。终映射在 \(CA\) 上为常值，因而下降成 \(v:Z/CA\to Z\)，满足 \(v\bar q\simeq\operatorname{id}_Z\)，其中 \(\bar q:Z\to Z/CA\)。对每个 \(t\)，\(H_t(CA)\subset CA\)，所以 \(\bar qH_t\) 也下降到商；由\cref{g2:lem:quotient-interval}，这些映射组成连续同伦，给出 \(\bar qv\simeq\operatorname{id}_{Z/CA}\)。这些同伦保持基点。将这两个同构与切除同态复合，所得就是原商映射 \(q\) 的诱导同态，故得到所述范围。
\end{proof}

\begin{lemma}[胞腔给出的同伦关系]\label{g2:lem:cell-attachment-homotopy-presentation}
	设 \(n\ge2\)，\(W=\bigvee_{\alpha\in\Lambda}S^n_\alpha\)，并沿带基点映射 \(\varphi_\beta:S^n\to W\) 附着一族 \((n+1)\) 胞腔，得到 \(V\)。则
	\[
		\pi_n(V)\cong
		\left(\bigoplus_{\alpha\in\Lambda}\mathbb Z[j_\alpha]\right)
		\Big/\left\langle[\varphi_\beta]:\beta\in B\right\rangle.
	\]
	每个关系 \([\varphi_\beta]\) 在第 \(\alpha\) 个坐标的系数为 \(\deg(p_\alpha\varphi_\beta)\)。
\end{lemma}
\begin{proof}
	\((V,W)\) 只附着 \((n+1)\) 胞腔，故为 \(n\) 连通。\(W\) 为 \((n-1)\) 连通，并且
	\[
		V/W=\bigvee_{\beta\in B}S^{n+1}_\beta.
	\]
	上一引理取 \(r=n\)、\(s=n-1\)。由于 \(n+1\le2n-1\)，在所需次数得到
	\[
		\pi_{n+1}(V,W)\cong\pi_{n+1}(V/W)\cong\bigoplus_{\beta\in B}\mathbb Z.
	\]
	各胞腔的特征圆盘给出左边的相对类，压去边界后正是右边各球面的恒等映射。若用第8章的立方体代表，只需把圆盘边界上一块小邻域压到所选基点，再把它对应到立方体的侧面与底面；该压缩与恒等映射同伦，保持相对定向和所代表的类。因此这些特征圆盘类确实组成一组基。

	空间对序列的相关部分为
	\[
		\pi_{n+1}(V,W)\xrightarrow{\partial}\pi_n(W)\longrightarrow\pi_n(V)\longrightarrow0.
	\]
	连接映射取特征圆盘的边界，所以把第 \(\beta\) 个基送到附着类 \([\varphi_\beta]\)。正合性说明这些恰好是全部关系，而前一个楔和引理把它们的坐标识别为 \(\deg(p_\alpha\varphi_\beta)\)。
\end{proof}

\begin{proof}[\cref{g2:thm:hurewicz-absolute}的证明]
	由\cref{g2:cor:sparse-cw-model}，可以用一个基点保持同伦等价 \(Y\to X\) 替换 \(X\)，使 \(Y\) 只有一个零胞腔，没有维数 \(1,\ldots,n-1\) 的胞腔，并使附着映射保持基点。该模型的构造及所用 Whitehead 定理见\cref{g2:sec:0901,g2:sec:0903}。同伦群和同调群在此替换下都同构，\(h_n\) 的自然性使两个比较方块交换，所以只须在 \(Y\) 上证明。

	\(Y\) 的约化胞腔链群在 \(i<n\) 时为零，因而 \(\widetilde H_i(Y)=0\)。令
	\[
		W=Y^n=\bigvee_{\alpha\in\Lambda}S^n_\alpha,\qquad V=Y^{n+1}.
	\]
	胞腔逼近分别作用于球面 \(S^n\) 和同伦 \(S^n\times I\)，给出 \(\pi_n(V)\cong\pi_n(Y)\)：前者把任何球面映射移入 \(n\) 骨架，后者把两个已有代表之间的同伦移入 \((n+1)\) 骨架并固定两端。同调方面，维数高于 \(n+1\) 的胞腔不出现在计算 \(H_n\) 的两张边界映射中，所以 \(H_n(V)\cong H_n(Y)\)。

	记 \(V\) 中各 \((n+1)\) 胞腔的附着映射为 \(\varphi_\beta\)。前一引理把 \(\pi_n(V)\) 写成下列上行的余核；\cref{g2:thm:cellular-boundary-degree}说明胞腔微分的矩阵也恰为 \(\deg(p_\alpha\varphi_\beta)\)，所以得到交换图
	\[
		\begin{tikzcd}[column sep=large,row sep=large]
			\displaystyle\bigoplus_\beta\mathbb Z \arrow[r,"{([\varphi_\beta])}"] \arrow[d,equal] & \pi_n(W) \arrow[r] \arrow[d,"h_n","\cong"'] & \pi_n(V) \arrow[r] \arrow[d,"h_n"] & 0 \\
			\displaystyle\bigoplus_\beta\mathbb Z \arrow[r,"\mathrm d_{n+1}"'] & H_n(W) \arrow[r] & H_n(V) \arrow[r] & 0.
		\end{tikzcd}
	\]
	左方块的交换来自相同的次数系数，也可由特征圆盘的相对 Hurewicz 映射与边界自然性看出；右方块由包含的自然性得到。中间竖箭头是楔和引理给出的同构，而且把上行关系子群送到下行关系子群，因而在余核上诱导同构。交换图保证这个诱导同构就是 \(h_n\)。结合两次骨架比较及 \(Y\to X\) 的同伦等价，定理得证。对具有上述 CW 同伦型的空间，再沿其同伦等价作同一自然性比较即可。
\end{proof}

这个证明也解释了连通度的作用。消去低维胞腔以后，每个 \(n\) 胞腔的边界都已压到一点，所以它本身就是球面；而 \((n+1)\) 胞腔边界上的全部信息，在这个维数恰好由各球面次数记录。若保留低维同伦信息，后一判断便不再成立。

\subsection{相对版本及单连通条件}
\label{g2:subsec:0902:04}

相对群研究的是边界留在 \(A\) 内的圆盘。把 \(A\) 压成一点，会把每个这样的圆盘变成球面；上一小节的商比较告诉我们，\(A\) 的连通度控制这种转换何时丢失信息。

\begin{theorem}[相对 Hurewicz 定理]\label{g2:thm:hurewicz-relative}
	设 \((X,A)\) 为 CW 对，\(X\) 与 \(A\) 均单连通，基点 \(x_0\in A\) 为顶点。若 \(n\ge2\) 且 \((X,A)\) 为 \((n-1)\) 连通，则
	\[
		H_i(X,A;\mathbb Z)=0\quad(0\le i<n),\qquad
		h_n:\pi_n(X,A,x_0)\xrightarrow{\ \cong\ }H_n(X,A;\mathbb Z).
	\]
	该同构对带基点 CW 对映射自然。在 \(n=2\) 时，源群也是阿贝尔群。
\end{theorem}
\begin{proof}
	令 \(q:X\to X/A\) 为商映射。\cref{g2:lem:relative-quotient-homotopy}取 \(r=n-1\)、\(s=1\)，得到
	\[
		q_*:\pi_i(X,A)\xrightarrow{\ \cong\ }\pi_i(X/A)\qquad(1\le i\le n).
	\]
	因此 \(X/A\) 为 \((n-1)\) 连通 CW 复形，绝对 Hurewicz 定理适用。

	同调中的商比较在这里可直接从胞腔链看出。\(X/A\) 有一个代表 \(A\) 的顶点，其余胞腔与 \(X\setminus A\) 的胞腔一一对应。把该顶点的链作为增广的基准，得到链同构
	\[
		C_\bullet^{\mathrm{cell}}(X,A)\cong
		C_\bullet^{\mathrm{cell}}(X/A,\{A\});
	\]
	在零次，把其他顶点送到它与基准顶点之差，就与约化胞腔复形对应。两边的边界均由原附着边界删除落在 \(A\) 中的项得到。由\cref{g2:thm:relative-cellular-comparison}及其自然性，商映射于是诱导
	\[
		H_i(X,A)\xrightarrow{\ \cong\ }\widetilde H_i(X/A)\qquad(i\ge0).
	\]
	绝对定理先给出这些群在 \(i<n\) 时为零，再给出图中的下方同构：
	\[
		\begin{tikzcd}[column sep=large,row sep=large]
			\pi_n(X,A) \arrow[r,"h_n"] \arrow[d,"q_*","\cong"'] & H_n(X,A) \arrow[d,"q_*","\cong"'] \\
			\pi_n(X/A) \arrow[r,"h_n"',"\cong"] & \widetilde H_n(X/A).
		\end{tikzcd}
	\]
	商映射把相对基本类 \(u_n\) 送到 \(D^n/S^{n-1}\) 的球面基本类，所以这个图交换。其余三条箭头为同构，故上方箭头也是同构。自然性已在\cref{g2:prop:hurewicz-naturality}证明。对于 \(n=2\)，空间对序列中的 \(\pi_2(X)\to\pi_2(X,A)\to\pi_1(A)=0\) 还直接说明 \(\pi_2(X,A)\) 是阿贝尔群 \(\pi_2(X)\) 的商。
\end{proof}

\begin{example}[单连通条件失效时]\label{g2:exm:hurewicz-nonsimply-connected}
	令 \(A=S^1\vee S^1\)，\(X=CA\) 为它的锥。\(X\) 可缩，\(A\) 道路连通，所以 \((X,A)\) 为一连通。两个空间对长正合列给出
	\[
		\pi_2(X,A)\cong\pi_1(A)=F(a,b),\qquad
		H_2(X,A)\cong H_1(A)=\mathbb Z^2.
	\]
	\(F(a,b)\) 是两个生成元的自由群。Hurewicz 映射与连接同态相容，所以此处的相对 \(h_2\) 正是自由群的交换化；交换子 \(aba^{-1}b^{-1}\) 被送到零，却是非平凡的相对同伦类。这里 \(X\) 已单连通，失去的是 \(A\) 的单连通性。
\end{example}

从圆盘的角度看，\(A\) 中的基点道路可以绕过非平凡环路，改变带基点的相对类；奇异同调不记录这种绕行。一般的相对 Hurewicz 定理会把基本群作用所给出的关系也取商，相关版本见\cite[第4.2节，The Hurewicz Theorem]{Hatcher2002}。本章使用的单连通条件使这些关系消失，也正好提供商比较在次数 \(n\) 上的同构范围。
% 纲要条目排为编号小节标题，并进入目录与书签；说明段仍保留于纲要注释中。

% !TeX root = ../../Book2.tex
% 纲要标题：### 9.3 Whitehead 定理
% 纲要位置：几何学计划纲要.md，第 1351 行。
% 正文状态：初稿；尚未正式审核。

% 纲要内容开始
% 
% * 弱同伦等价与映射柱
% * 相对逐骨架收缩
% * CW 对上的 Whitehead 定理
% * 非 CW 对象的范围界线
% 
% 纲要内容结束

\section{Whitehead 定理}
\label{g2:sec:0903}

同伦群由球面映射探测空间，CW 复形则由球面附着圆盘构成。这两种球面的来源能够接合：如果一个映射没有留下任何相对同伦类，那么它的映射柱中的每个胞腔都能压回源空间。逐骨架拼合这些压缩，就得到整个空间的同伦逆。Whitehead 定理将这种构造写成一个检测同伦等价的判据。

\subsection{弱同伦等价与映射柱}
\label{g2:subsec:0903:01}

\begin{definition}[弱同伦等价]\label{g2:def:weak-homotopy-equivalence}
	连续映射 \(f:X\rightarrow Y\) 称为\term[ruo tong lun deng jia]{弱同伦等价}[weak homotopy equivalence]，如果 \(f_*:\pi_0(X)\rightarrow\pi_0(Y)\) 是双射，且对每个 \(x\in X\) 和每个 \(i\ge1\)，
	\[
		f_*:\pi_i(X,x)\longrightarrow\pi_i(Y,f(x))
	\]
	都是同构。
\end{definition}

在道路连通空间上，只须在一个基点处验证各阶同构；其他基点可由\cref{g2:prop:basepoint-action}运输过去。相同的抽象同伦群还没有给出这里要求的映射 \(f\)：判据检验的是同一个连续映射在全部维数上的作用。

\ProseParagraphBreak

回顾\cref{g2:def:mapping-cylinder-cone,g2:prop:mapping-cylinder-retract}，映射柱及其分解是
\[
	M_f=(X\times I\amalg Y)/((x,1)\sim f(x)),
	\qquad X\xrightarrow{j}M_f\xrightarrow{p}Y,
\]
其中 \(j(x)=[x,0]\)、\(p([x,t])=f(x)\)、\(p(y)=y\)，且 \(pj=f\)。映射 \(p\) 通过把柱线推向 \(t=1\) 给出到 \(Y\) 的强形变收缩。因此由空间对的同伦序列，\(f\) 是弱同伦等价，恰好等价于 \((M_f,jX)\) 为任意阶连通。这样，把映射变成包含之后，各阶相对群都成为可以直接压缩圆盘的条件。

\begin{lemma}[胞腔映射的映射柱]\label{g2:lem:cellular-mapping-cylinder}
	设 \(f:X\rightarrow Y\) 是 CW 复形之间的胞腔映射，则 \(M_f\) 是 CW 复形，其中 \(X\times\{0\}\) 与 \(Y\) 都是子复形。
\end{lemma}
\begin{proof}
	保留 \(Y\) 的胞腔和 \(X\times\{0\}\) 的胞腔，并对 \(X\) 的每个 \(r\) 胞腔添一个开柱体 \(e^r\times(0,1)\)，把它看成 \((r+1)\) 胞腔。其闭特征圆盘可用 \(D^r\times I\) 表示，边界的三个部分分别进入 \(X\times\{0\}\)、由低维胞腔产生的柱体以及 \(Y^r\)；最后一项使用 \(f(X^r)\subset Y^r\)。所以胞腔边界进入所要求的低维骨架。映射柱的商拓扑在 \(Y\) 和各 \(D^r\times I\) 上检验连续性；\cref{g2:lem:weak-product}说明这正是所列胞腔的弱拓扑。因此得到的是原来的映射柱拓扑。
\end{proof}

\subsection{相对逐骨架收缩}
\label{g2:subsec:0903:02}

圆盘压缩已经解决了一个胞腔的问题。对整个 CW 复形，需要让低维部分先到达指定子空间，并在处理高维胞腔时保留它。第1章的同伦延拓性质负责把局部修正扩展为全空间上的同伦，弱拓扑负责处理无限多个胞腔。

\begin{lemma}[逐骨架压缩]\label{g2:lem:cellular-compression}
	设 \((K,L)\) 为 CW 对，\(B\subset Y\) 非空，\(f:(K,L)\rightarrow(Y,B)\) 连续。对 \(K\setminus L\) 中出现的每个正维胞腔维数 \(m\)，假定 \(\pi_m(Y,B,b)=0\) 对所有 \(b\in B\) 成立；若出现零胞腔，再假定 \(Y\) 的每个道路分支与 \(B\) 相交。则 \(f\) 可以相对 \(L\) 同伦到一个取值于 \(B\) 的映射。
\end{lemma}
\begin{proof}
	按维数依次修正 \(f\)。假定目前的映射 \(f_{m-1}\) 已把 \(L\cup K^{m-1}\) 映入 \(B\)。对每个不属于 \(L\) 的 \(m\) 胞腔，取特征映射 \(\chi:D^m\rightarrow K\)，于是
	\[
		f_{m-1}\chi:(D^m,S^{m-1})\longrightarrow(Y,B).
	\]
	由\cref{g2:lem:disk-compression}，它可以固定整个边界地移入 \(B\)。对同维的全部胞腔同时作这种压缩，并固定 \(L\cup K^{m-1}\)，便得到 \(L\cup K^m\) 上的同伦；各闭圆盘边界相容，\cref{g2:lem:weak-product}保证粘合连续。由\cref{g2:thm:relative-CW-HEP}，把它延拓到整个 \(K\)，仍固定已指定的 \(L\cup K^{m-1}\)。所得终映射 \(f_m\) 把 \(L\cup K^m\) 映入 \(B\)。对于 \(m=0\)，先为每个新增顶点选一条通向 \(B\) 的道路，随后作同样的延拓。

	把第 \(m\) 层修正安排在时间区间
	\[
		[1-2^{-m},\,1-2^{-(m+1)}],\qquad m=0,1,2,\ldots.
	\]
	任一 \(r\) 维特征圆盘在第 \(r\) 层结束后逐点不再移动，故可把 \(t=1\) 的值定义为其稳定值。每个特征圆盘与 \(I\) 的积上，这是一段有限次拼接的连续同伦，之后接常值段；在 \(L\times I\) 上始终为原映射。再由\cref{g2:lem:weak-product}，得到整个 \(K\times I\) 上的连续同伦，包括累积时刻 \(t=1\)。它的终映射取值于 \(B\)，且从始至终固定 \(L\)。
\end{proof}

这段证明还说明，压缩各个胞腔时不必使它们在靶空间中互不相交。需要相容的是公共边界上的映射与同伦；CW 的附着和弱拓扑恰好把这一相容性转化为整体连续性。

\subsection{CW 对上的 Whitehead 定理}
\label{g2:subsec:0903:03}

\begin{theorem}[Whitehead]\label{g2:thm:whitehead}
	设 \(X,Y\) 为 CW 复形，\(f:X\rightarrow Y\) 连续。则 \(f\) 是同伦等价，当且仅当它是弱同伦等价。若 \(X\) 是 \(Y\) 的子复形且 \(f\) 为包含，则弱同伦等价还保证 \(X\) 是 \(Y\) 的强形变收缩核。
\end{theorem}
\begin{proof}
	同伦等价诱导道路分支双射及各阶同伦群同构，必要性成立。证明充分性时，先处理子复形包含 \(X\hookrightarrow Y\)。弱等价和\cref{g2:thm:pair-homotopy-sequence}给出 \(\pi_m(Y,X,x)=0\) 对所有 \(m\ge1\) 和 \(x\in X\) 成立；分支双射还保证每个 \(Y\) 分支遇到 \(X\)。把\cref{g2:lem:cellular-compression}用于恒等映射 \((Y,X)\rightarrow(Y,X)\)，得到
	\[
		H:Y\times I\longrightarrow Y,\qquad
		H(y,0)=y,\quad H(Y,1)\subset X,\quad H(x,t)=x.
	\]
	这正是所需强形变收缩。零维与一维情形已包含在压缩引理的道路条件和相对同伦集条件中。

	对一般 \(f\)，先由\cref{g2:thm:cellular-approximation}取与它同伦的胞腔映射 \(f_c\)。同伦时基点的轨迹通过\cref{g2:prop:basepoint-action}比较诱导映射，因此 \(f_c\) 仍为弱同伦等价。由\cref{g2:lem:cellular-mapping-cylinder}，\((M_{f_c},jX)\) 是 CW 对。映射柱收缩 \(p:M_{f_c}\rightarrow Y\) 为同伦等价，且 \(pj=f_c\)，所以 \(j\) 在道路分支及各阶同伦群上都是同构。刚证的包含情形说明 \(j\) 为同伦等价，因而 \(f_c=pj\) 为同伦等价。与它同伦的 \(f\) 也有同伦逆。
\end{proof}

\begin{corollary}[相对 Whitehead 判据]\label{g2:cor:relative-whitehead}
	设 \((X,A)\)、\((Y,A)\) 为 CW 对，公共子复形 \(A\) 取相同 CW 结构。若 \(f:X\rightarrow Y\) 是弱同伦等价且 \(f|_A=\operatorname{id}_A\)，则 \(f\) 是相对 \(A\) 的同伦等价：存在 \(g:Y\rightarrow X\)，满足 \(g|_A=\operatorname{id}_A\)，且 \(gf\simeq\operatorname{id}_X\)、\(fg\simeq\operatorname{id}_Y\) 的同伦都固定 \(A\)。特别地，保持零胞腔基点的 CW 弱同伦等价具有保持基点的同伦逆和同伦。
\end{corollary}
\begin{proof}
	由胞腔逼近的相对版本，可先把 \(f\) 固定 \(A\) 地变成胞腔映射。取映射柱 \(M_f\)，再把每条 \(\{a\}\times I\)、\(a\in A\)，压成一点，所得空间记为 \(W\)。这使柱的两端共享同一个 \(A\)。它有一个 CW 结构：保留 \(X\) 与 \(Y\) 的胞腔，只把公共 \(A\) 计算一次，对 \(X\setminus A\) 的每个胞腔再添其开柱体。与\cref{g2:lem:cellular-mapping-cylinder}相同的边界与商拓扑论证表明，\(X,Y\) 均为 \(W\) 的子复形。

	柱线向 \(Y\) 推进的强形变收缩仍可下降到 \(W\)，记其终映射为 \(p:W\rightarrow Y\)；它固定 \(A\)，且 \(p|_X=f\)。因此 \(X\hookrightarrow W\) 是弱同伦等价。由\cref{g2:thm:whitehead}的子复形版本，\(W\) 又强形变收缩到 \(X\)，记收缩终映射为 \(r:W\rightarrow X\)。令 \(g=r|_Y\)。将到 \(X\) 的收缩后复合 \(p\) 并限制到 \(Y\)，得到 \(fg\simeq\operatorname{id}_Y\)；将到 \(Y\) 的收缩后复合 \(r\) 并限制到 \(X\)，得到 \(gf\simeq\operatorname{id}_X\)。两种收缩都固定公共的 \(A\)，所以同伦也固定 \(A\)。最后接回最初的相对胞腔逼近同伦，即得原映射的结论。
\end{proof}

于是，\cref{g2:cor:sparse-cw-model}中的模型既有指定的胞腔维数，也保留了所需的基点或子空间。另一个直接应用是：道路连通 CW 复形若所有正阶同伦群均为零，选一个顶点的包含就是弱同伦等价，故该空间可强形变收缩到该顶点。Whitehead 判据与压缩证明可参照\cite[第4.1节，Whitehead's Theorem]{Hatcher2002}。

\subsection{非 CW 对象的范围界线}
\label{g2:subsec:0903:04}

CW 条件在证明中提供了逐胞腔构造同伦的方式。对一般空间，球面探测到的全部信息可能仍不足以控制连续映射的逆。下面的例子把差别集中在道路分支的拓扑上。

\begin{example}\label{g2:exm:rationals-weak-equivalence}
	以 \(\mathbb Q_{\mathrm{disc}}\) 表示带离散拓扑的有理数集，以 \(\mathbb Q\) 表示实数直线的通常子空间。恒等集合映射
	\[
		j:\mathbb Q_{\mathrm{disc}}\longrightarrow\mathbb Q
	\]
	连续，并且是弱同伦等价，却不是同伦等价。

	先说明 \(\mathbb Q\) 的每个连通子集只有一点。若某个子集含有 \(a<b\)，可在其间选无理数 \(c\)；与 \(({-\infty},c)\) 及 \((c,\infty)\) 的交把该子集分成两个非空、不交的相对开集。因此道路以及从任一 \(S^n\)、\(n\ge1\)，到 \(\mathbb Q\) 的映射均为常值。两空间的道路分支都是单点，所有正阶同伦群都为零，故 \(j\) 满足弱等价的全部条件。

	若 \(j\) 有连续同伦逆 \(g:\mathbb Q\rightarrow\mathbb Q_{\mathrm{disc}}\)，同伦 \(jg\simeq\operatorname{id}_{\mathbb Q}\) 在每个点处给出 \(\mathbb Q\) 中的一条道路。这条道路只能恒定，所以 \(g(q)=q\) 对每个 \(q\in\mathbb Q\) 成立。但集合恒等映射 \(\mathbb Q\rightarrow\mathbb Q_{\mathrm{disc}}\) 不连续：离散空间中的开集 \(\{0\}\) 的原像，在通常 \(\mathbb Q\) 中不是开集。矛盾。
\end{example}

这个例子的源是一个零维 CW 复形，靶却不能被上述弱等价替换为同伦等价。若源、靶都已具有 CW 同伦型，Whitehead 判据则仍然适用：用各自的 CW 模型前后复合 \(f\)，所得 CW 映射仍是弱等价，应用定理后再接回两端的同伦等价即可。实际使用判据时，先找到这样的 CW 模型，便能把“各阶同伦群同构”转化为真正的空间同伦等价。

% 纲要条目排为编号小节标题，并进入目录与书签；说明段仍保留于纲要注释中。

% !TeX root = ../../Book2.tex
% 纲要标题：### 9.4 检测与构造的应用
% 纲要位置：几何学计划纲要.md，第 1358 行。
% 正文状态：未撰写；本文件目前是框架。

% 纲要内容开始
% 
% * 单连通 CW 空间的同调判据
% * Moore 空间的构造
% * 相同同调与不同同伦类型
% * Eilenberg–MacLane 空间的初步存在
% 
% 纲要内容结束

\section{检测与构造的应用}
\label{g2:sec:0904}

待撰写

\subsection{单连通 CW 空间的同调判据}
\label{g2:subsec:0904:01}
待撰写

\subsection{Moore 空间的构造}
\label{g2:subsec:0904:02}
待撰写

\subsection{相同同调与不同同伦类型}
\label{g2:subsec:0904:03}
待撰写

\subsection{Eilenberg–MacLane 空间的初步存在}
\label{g2:subsec:0904:04}
待撰写
% 纲要条目排为编号小节标题，并进入目录与书签；说明段仍保留于纲要注释中。

% !TeX root = ../../Book2.tex
% 纲要标题：### 9.5 悬挂、Freudenthal 定理与稳定范围
% 纲要位置：几何学计划纲要.md，第 1365 行。
% 正文状态：初稿；一般同伦切除准确外引，双锥应用与稳定范围在正文证明。

% 纲要内容开始
% 
% 本节承担稳定路线的共同输入，不把悬挂稳定化等同于向量丛加平凡丛。
% 
% * 约化悬挂、回路伴随及自然映射 $X\to\Omega\Sigma X$
% * 同伦切除在两个锥粘合中的使用
% 
%   准确引用同伦切除定理的连通度版本，并完整写出其在双锥中的应用；一般 Blakers–Massey 定理的长证明不增加为本节的隐含任务。[R7]
% 
% * $(n-1)$ 连通带基点 CW 复形的悬挂定理
% 
%   固定 $n\ge2$。对 $(n-1)$ 连通的带基点 CW 复形，悬挂映射 $\pi_i(X)\to\pi_{i+1}(\Sigma X)$ 在 $i\le2n-2$ 时为同构，在 $i=2n-1$ 时为满射；分别核验相对同伦序列中的指标与两个锥的连通度。
% 
% * 球面同伦群在固定差值 $i-n$ 下的稳定化
% * 稳定范围、临界次数和范围以外信息的区别
% * 与带标架配边、谱及 K 理论中不同稳定化的衔接
% 
% ---
% 
% 纲要内容结束

\section{悬挂、Freudenthal 定理与稳定范围}
\label{g2:sec:0905}

把球面上的映射沿一个新方向拉长，再把两端压成一点，就得到维数同时提高一的映射。这个操作称为悬挂。它保留了多少同伦信息，是从逐个同伦群的计算走向稳定同伦的关键问题。本节先把悬挂写成可直接操作的商空间，再利用两个锥的粘合说明：当空间的低维同伦群已经消失时，一定范围内的同伦类在悬挂后既不丢失，也不新增。

\subsection{\texorpdfstring{约化悬挂、回路伴随及自然映射 \(X\to\Omega\Sigma X\)}{约化悬挂、回路伴随及自然映射 X→\GeometryPdfOmega{}\GeometryPdfCapitalSigma{}X}}
\label{g2:subsec:0905:01}
回忆第3章的未约化悬挂：将 \(X\times I\) 的上下两端各压成一点。若 \(X\) 已有基点 \(x_0\)，还可以把整条 \(\{x_0\}\times I\) 压到基点，这样所有构造都保持基点。本节的带基点 CW 复形均取基点为零胞腔；这保证基点包含具有同伦延拓性质。一般地，包含 \(\{x_0\}\hookrightarrow X\) 为余纤维化时，称基点为\term[liangjidian]{良基点}[well-pointed]。

\begin{definition}[约化悬挂与约化锥]\label{g2:def:reduced-suspension}
	设 \((X,x_0)\) 为带基点空间。它的\term[yuehuaxuangua]{约化悬挂}[reduced suspension]和\term[yuehuazhui]{约化锥}[reduced cone]分别为
	\[
		\Sigma X\coloneqq\frac{X\times I}{X\times\{0,1\}\,\cup\,\{x_0\}\times I},\qquad
		C_*X\coloneqq\frac{X\times I}{X\times\{0\}\,\cup\,\{x_0\}\times I}.
	\]
	分母所示子空间各压成一个基点。约化锥中的 \([x,1]\) 给出 \(X\) 的底面，公式 \([x,t]\mapsto[x,(1-s)t]\) 给出固定基点的收缩。带基点映射 \(f:X\rightarrow Y\) 诱导 \(\Sigma f([x,t])=[f(x),t]\)。
\end{definition}

在带基点 CW 模型中，这些商保留 CW 结构，底面包含 \(X\hookrightarrow C_*X\) 是子复形包含。把两份约化锥沿底面相粘，得到
\[
	\Sigma X=C_+X\cup_X C_-X.
\]
约定 \(C_+X\) 由悬挂参数 \(0\le t\le1/2\) 给出，\(C_-X\) 由 \(1/2\le t\le1\) 给出，公共底面是 \(X\times\{1/2\}\)。第一顶锥的锥参数 \(v\) 对应悬挂参数 \(t=v/2\)，因此锥参数从顶点到公共底面递增。两锥的顶点与基点线均已合为基点。由球面的立方体商模型，\(\Sigma S^i\cong I^{i+1}/\partial I^{i+1}\cong S^{i+1}\)，最后一坐标取悬挂参数；\(\Sigma\operatorname{id}_{S^i}\) 对应 \(\operatorname{id}_{S^{i+1}}\)。

\begin{proposition}[悬挂与回路的伴随]\label{g2:prop:suspension-loop-adjunction}
	设 \(X,Y\) 为带基点空间，\(\Omega Y\) 取紧开拓扑。带基点映射由下列公式一一对应，并诱导自然双射
	\[
		[\Sigma X,Y]_*\cong[X,\Omega Y]_*,\qquad
		F\longmapsto\widehat F,\quad \widehat F(x)(t)=F([x,t]).
	\]
	恒等映射 \(\operatorname{id}_{\Sigma X}\) 对应自然映射
	\[
		\eta_X:X\longrightarrow\Omega\Sigma X,\qquad
		\eta_X(x)(t)=[x,t].
	\]
\end{proposition}
\begin{proof}
	先看映射本身。一个从 \(X\times I\) 到 \(Y\) 的连续映射能下降到约化悬挂，当且仅当它在两端及基点线上恒取 \(y_0\)。由区间的紧性及紧开拓扑的指数律，这等价于连续映射 \(x\mapsto(t\mapsto F(x,t))\) 取值于 \(\Omega Y\)，且把 \(x_0\) 送到常值回路。两种公式互为逆映射。

	对带参数的映射使用同一公式，就得到同伦之间的对应；商映射与区间取积仍为商映射，见\cref{g2:lem:quotient-interval}，所以由 \(X\times I\times I\) 下降的同伦连续。对源的预复合和对靶的后复合都与公式交换，故双射自然。把 \(F\) 取为恒等映射便得到 \(\eta_X\)，并且对带基点 \(f:X\rightarrow Y\) 有严格等式 \(\Omega\Sigma f\circ\eta_X=\eta_Y\circ f\)。
\end{proof}

这项伴随把一个几何操作变成了对同一个空间的自然映射。对 \(i\ge1\)，\term[xuangua-tongtai]{悬挂同态}[suspension homomorphism]定义为
\[
	E:\pi_i(X)\longrightarrow\pi_{i+1}(\Sigma X),\qquad [f]\longmapsto[\Sigma f].
\]
球面上的加法由夹断映射定义，悬挂与这一夹断相容，所以 \(E\) 是群同态。经\cref{g2:exm:loop-space-shift}的同构 \(\pi_i(\Omega\Sigma X)\cong\pi_{i+1}(\Sigma X)\)，它正是 \((\eta_X)_*\)：映射 \(\eta_X\circ f\) 的伴随就是 \([u,t]\mapsto[f(u),t]\)。因此，悬挂保留哪些次数的信息，也就是 \(\eta_X\) 在哪些次数诱导同构。

\subsection{同伦切除在两个锥粘合中的使用}
\label{g2:subsec:0905:02}
同调切除允许从粘合中去掉一块公共部分。同伦群记录更细的映射与延拓问题，相应比较只在由连通度控制的范围内成立。下面使用其 CW 版本，证明可查 \cite[第4.2节，Excision for Homotopy Groups]{Hatcher2002}。该定理是本节的外部输入；两个锥的连通度、比较映射及次数换算将在下文逐一落实。

\begin{theorem}[同伦切除]\label{g2:thm:homotopy-excision}
	设 CW 复形 \(Z=A\cup B\)，其中 \(A,B\) 为子复形，\(C=A\cap B\) 非空且道路连通，基点取在 \(C\) 中。设 \(r,s\ge0\)，空间对 \((A,C)\) 为 \(r\) 连通，\((B,C)\) 为 \(s\) 连通。则在次数 \(j\ge1\)，包含诱导的映射
	\[
		\pi_j(A,C)\longrightarrow\pi_j(Z,B)
	\]
	在 \(j<r+s\) 时为同构，在 \(j=r+s\) 时为满射。这里在次数 \(j\ge2\) 使用相对同伦群，在次数 \(j=1\) 使用相应尖点集并将同构理解为双射。
\end{theorem}

其中 \(r\) 连通的含义是相对同伦对象在次数不超过 \(r\) 时消失，零次条件要求每个道路分支与子空间相交。将定理用于双锥，所比较的推出图为
\[
	\begin{tikzcd}
		X \arrow[r,hook] \arrow[d,hook] & C_+X \arrow[d,hook] \\
		C_-X \arrow[r,hook] & \Sigma X.
	\end{tikzcd}
\]
图中两条水平箭头之间的比较，恰是
\[
	\pi_j(C_+X,X)\longrightarrow\pi_j(\Sigma X,C_-X).
\]
两侧都把圆盘边界放进指定子空间；左侧要求整张圆盘落在一顶锥中，右侧允许它使用整个悬挂。切除定理说明，低维时允许使用另一顶锥没有增加新的相对类。

\ProseParagraphBreak
先计算两个锥提供的连通度。假设 \(X\) 为 \((n-1)\) 连通且 \(n\ge2\)。两锥可缩，故\cref{g2:thm:pair-homotopy-sequence}给出
\[
	\partial:\pi_j(C_\pm X,X)\xrightarrow{\ \cong\ }\pi_{j-1}(X)\qquad(j\ge2).
\]
右侧在 \(2\le j\le n\) 为零；次数一的相对尖点集也只有一个元素，因为锥可缩而 \(X\) 道路连通。因此两对 \((C_\pm X,X)\) 都是 \(n\) 连通。切除中的连通度是 \(n,n\)，比 \(X\) 本身的连通度各高一；这一差别决定了最后的稳定范围。

\subsection{\texorpdfstring{\((n-1)\) 连通带基点 CW 复形的悬挂定理}{(n−1) 连通带基点 CW 复形的悬挂定理}}
\label{g2:subsec:0905:03}
悬挂定理把上面的相对比较转回绝对同伦群。Freudenthal 的结果说明，连通度越高，能够无损悬挂的次数越多。\footnote{Hans Freudenthal（1905—1990），生于德国的荷兰数学家，在代数拓扑和数学教育方面作出重要贡献，1937年发表悬挂定理研究。人物资料见 MacTutor 的 \href{https://mathshistory.st-andrews.ac.uk/Biographies/Freudenthal/}{Hans Freudenthal 传记}。}

\begin{theorem}[Freudenthal 悬挂定理]\label{g2:thm:freudenthal}
	设 \(n\ge2\)，\(X\) 为 \((n-1)\) 连通的 CW 复形，基点为一个零胞腔。则悬挂同态
	\[
		E:\pi_i(X)\longrightarrow\pi_{i+1}(\Sigma X)
	\]
	在 \(1\le i\le2n-2\) 时为同构，在 \(i=2n-1\) 时为满射。等价地，自然映射 \(\eta_X:X\rightarrow\Omega\Sigma X\) 在次数小于 \(2n-1\) 时诱导同构，在次数 \(2n-1\) 时诱导满射。
\end{theorem}
\begin{proof}
	在带基点 CW 模型中写 \(\Sigma X=C_+X\cup_XC_-X\)。上一小节已经证明两对 \((C_\pm X,X)\) 是 \(n\) 连通。由\cref{g2:thm:homotopy-excision}，中间的包含映射与两侧的相对同伦序列组成
	\[
		\begin{tikzcd}[column sep=small]
			\pi_i(X) & \pi_{i+1}(C_+X,X) \arrow[l,"\partial"',"\cong"] \arrow[r,"e"] & \pi_{i+1}(\Sigma X,C_-X) & \pi_{i+1}(\Sigma X) \arrow[l,"j_*"',"\cong"] .
		\end{tikzcd}
	\]
	左侧边界同构来自 \(C_+X\) 可缩；右侧同构来自 \(C_-X\) 可缩。后一个断言在 \(i=1\) 也成立：正合列两侧的 \(\pi_2(C_-X)\) 与 \(\pi_1(C_-X)\) 都为零。于是复合 \(j_*^{-1}e\partial^{-1}\) 在 \(i+1<2n\) 时为同构，在 \(i+1=2n\) 时为满射。这分别等价于 \(i\le2n-2\) 与 \(i=2n-1\)。

	还要识别这项复合确实是悬挂。给定 \(f:S^i\rightarrow X\)，将它写成 \(f:(I^i,\partial I^i)\rightarrow(X,x_0)\)。用悬挂参数表示，公式 \(F(u,t)=[f(u),t/2]\) 取值于 \(C_+X\)；它在底面和侧面取基点，在顶面取公共底面的 \(f(u)\)。所以它是\cref{g2:def:relative-homotopy}中的相对圆盘代表，且 \(\partial[F]=[f]\)，即 \([F]=\partial^{-1}[f]\)。

	在球面的立方体商模型中，\(\Sigma f\) 的公式是 \(G(u,t)=[f(u),t]\)。当它作为 \((\Sigma X,C_-X)\) 中的相对映射时，连续变形
	\[
		H_s(u,t)=\bigl[f(u),(1-s/2)t\bigr],\qquad 0\le s\le1,
	\]
	将 \(G\) 变为 \(F\)。底面和侧面始终取基点；顶面的悬挂参数为 \(1-s/2\in[1/2,1]\)，所以顶面始终落在指定子空间 \(C_-X\)。这正是允许的相对同伦，因而 \(j_*[\Sigma f]=e[F]=e\partial^{-1}[f]\)。几何上，它把整段悬挂压到第一顶锥，同时允许顶面沿第二顶锥移动。所用边界就是顶面 \(f\)，与既定边界约定一致。于是 \(j_*^{-1}e\partial^{-1}=E\)，结论成立。
\end{proof}

证明中的三种次数可压缩为一张对应表，使用时先找准被比较的是圆盘维数还是球面维数：
\[
	\begin{array}{c|c|c}
		\text{对象或映射} & \text{同构范围} & \text{满射的临界次数}\\ \hline
		\pi_j(C_+X,X)\rightarrow\pi_j(\Sigma X,C_-X) & j<2n & j=2n\\
		\pi_i(X)\rightarrow\pi_{i+1}(\Sigma X) & i\le2n-2 & i=2n-1
	\end{array}
\]
悬挂使源球面和靶空间同时增加一个悬挂方向，而相对边界把圆盘维数减一；表中的两个换算来自这两项不同操作。

\subsection{\texorpdfstring{球面同伦群在固定差值 \(i-n\) 下的稳定化}{球面同伦群在固定差值 i−n 下的稳定化}}
\label{g2:subsec:0905:04}
球面 \(S^n\) 是 \((n-1)\) 连通的，见\cref{g2:prop:sphere-connectivity}。因此对 \(n\ge2\)，悬挂给出
\[
	\pi_i(S^n)\longrightarrow\pi_{i+1}(S^{n+1}),
	\qquad i\le2n-2\ \text{时为同构}.
\]
每次悬挂都使 \(i,n\) 同时加一，所以应固定差值 \(k=i-n\) 来排列这些群。以 \(k\ge0\) 为例，得到
\[
	\pi_{n+k}(S^n)\xrightarrow{E}\pi_{n+k+1}(S^{n+1})\xrightarrow{E}\pi_{n+k+2}(S^{n+2})\longrightarrow\cdots.
\]
不等式 \(n+k\le2n-2\) 等价于 \(n\ge k+2\)。一旦达到这个维数，后面的每一步都在同构范围内，因为悬挂一次使左侧与右侧之间的余量再增加一。

\begin{definition}[球面的稳定同伦群]\label{g2:def:stable-sphere-homotopy}
	对整数 \(k\)，沿悬挂同态取直接极限，定义球面的第 \(k\) 个\term[wending-tonglunqun]{稳定同伦群}[stable homotopy group]为
	\[
		\pi_k^{\mathrm S}\coloneqq\varinjlim_{n}\pi_{n+k}(S^n).
	\]
	取 \(n\) 充分大以保证 \(n\ge2\)、\(n+k\ge2\)。极限的元素由某一步中的同伦类表示；两项代表相同，当且仅当它们经过有限次悬挂后在同一个群中相等。
\end{definition}

对 \(k\ge0\)，\cref{g2:thm:freudenthal}给出自然同构
\[
	\pi_{n+k}(S^n)\cong\pi_k^{\mathrm S}\qquad(n\ge k+2).
\]
直接极限在这里可以理解为一条最终全由同构组成的链的共同取值。比如差值 \(k=1\) 时，\(\pi_3(S^2)\rightarrow\pi_4(S^3)\) 是临界满射，从 \(\pi_4(S^3)\) 开始得到稳定值；差值 \(k=2\) 时，\(\pi_5(S^3)\rightarrow\pi_6(S^4)\) 是临界满射，从 \(\pi_6(S^4)\) 开始稳定。定理在这里确定比较映射的性质；具体群的数值仍由各自的计算给出。

\ProseParagraphBreak
差值零的稳定群可以直接认出。\cref{g2:prop:complex-hopf-homotopy}通过复 Hopf 纤维化得到 \(\pi_2(S^2)\cong\mathbb Z\)，而\cref{g2:lem:degree-homotopy-additivity}给出的度数同态 \(\deg:\pi_2(S^2)\rightarrow\mathbb Z\) 把恒等映射送到 \(1\)。因此它是满射，也是无限循环群之间的同构。悬挂将恒等映射送到恒等映射，故从 \(n=2\) 开始得到 \(\pi_n(S^n)\cong\mathbb Z\)，并有 \(\pi_0^{\mathrm S}\cong\mathbb Z\)。这里的稳定下标零表示“源与靶的维数差为零”，与通常的道路分支集 \(\pi_0(X)\) 含义不同。

\subsection{稳定范围、临界次数和范围以外信息的区别}
\label{g2:subsec:0905:05}
同构范围内，一个悬挂后的类恰好有一个悬挂前的类；临界次数只保证至少有一个原像，多个原像可能合并。下例把这种合并写成一张明确的附着映射。

\begin{example}[临界悬挂中的非零核]\label{g2:exm:critical-suspension-kernel}
	令 \(W=S^2\vee S^2\)，基点为公共零胞腔。乘积 \(S^2\times S^2\) 由 \(W\) 附着一个四胞腔得到，记其附着映射为 \(a:S^3\rightarrow W\)。则 \([a]\ne0\)，但 \(E[a]=0\)。因此在 \(n=2\)、\(i=3=2n-1\) 的临界次数，悬挂确实可以有核。
\end{example}
\begin{proof}
	把四圆盘取为 \(D^2\times D^2\)。在 \(D^2\times S^1\) 上，\(a\) 是投向第一份 \(D^2/S^1=S^2\)；在 \(S^1\times D^2\) 上，\(a\) 是投向第二份球面；公共部分 \(S^1\times S^1\) 都送到基点。两式在交集相容，给出 \(\partial(D^2\times D^2)\cong S^3\) 上的连续映射。这正是乘积顶胞腔的边界。

	若 \(a\) 零伦，\cref{g2:prop:homotopic-attaching-maps}便给出 \(S^2\times S^2\simeq S^2\vee S^2\vee S^4\)。但在乘积中，两个投影拉回的二次整数上同调生成元 \(u,v\) 满足 \(u\smile v\ne0\)：由\cref{g2:prop:cross-cup-evaluation}，它在乘积基本类上的取值为 \(1\)。在后一楔和中，任意两个二次类的乘积都为零，因为限制到每个 \(S^2\) 因子时四次群为零，限制到 \(S^4\) 因子时二次群为零，而各因子的限制共同检测四次类，见\eqref{g2:eq:wedge-positive-cohomology}。上同调环在同伦等价下保持，所以矛盾，\([a]\ne0\)。

	再设 \(j:W\hookrightarrow S^2\times S^2\) 为包含。构造
	\[
		r:\Sigma(S^2\times S^2)\longrightarrow\Sigma S^2\vee\Sigma S^2\cong\Sigma W
	\]
	如下：先在悬挂参数的中点夹断，得到两个悬挂的楔和，再在前一份使用第一投影的悬挂，在后一份使用第二投影的悬挂。具体地，\(r([(x,y),t])\) 在 \(0\le t\le1/2\) 等于第一份中的 \([x,2t]\)，在 \(1/2\le t\le1\) 等于第二份中的 \([y,2t-1]\)。分界处及被压缩的子空间均映到基点，所以公式连续下降。

	将 \(r\Sigma j\) 限制到第一份 \(\Sigma S^2\)，它先用前半段完成恒等映射，再在后半段停留于基点；限制到第二份时，先在基点停留，再用后半段完成恒等映射。连续缩去停留段，得到固定公共基点的同伦 \(r\Sigma j\simeq\operatorname{id}_{\Sigma W}\)。由于 \(ja\) 延到乘积四胞腔，它零伦，故 \((\Sigma j)(\Sigma a)\) 也零伦。复合左逆 \(r\) 得到 \(\Sigma a\simeq*\)，即 \(E[a]=0\)。
\end{proof}

这个例子还解释了稳定化为何可能简化计算：乘积顶胞腔把两份球面联系起来，联系可以由杯积检测；一次悬挂已经使这张非零附着映射消失。对于 \(i>2n-1\) 的次数，悬挂定理没有给出单射或满射的统一结论，必须继续利用空间的具体结构。若某个空间具有比所选 \(n-1\) 更高的连通度，则应先提高 \(n\) 再应用定理；所谓范围是由已知连通度得到的统一保证。

\subsection{与带标架配边、谱及 K 理论中不同稳定化的衔接}
\label{g2:subsec:0905:06}
球面稳定群可以用映射、几何子流形和一整列空间来研究。这里先列出三种实际操作，说明后续章节怎样与本节相接。

\ProseParagraphBreak
第23章从嵌入 \(M^k\hookrightarrow\mathbb R^{k+n}\) 及法丛的一组标架出发，将管状邻域中的法向圆盘沿标架识别为 \(D^n\)，再把邻域外部压到一点，构造映射 \(S^{k+n}\rightarrow S^n\)。若把环境嵌入再加一个零坐标，并给法丛增加这一坐标的单位向量，源与靶都多出一个方向；相应坍缩映射与本节悬挂的比较将在第23.2节证明。这样，固定维数 \(k\) 的带标架几何对象就与固定差值 \((k+n)-n=k\) 相对应。将同伦与配边建立双向对应，还要在那里处理横截逆像及各项选择的无关性。

\ProseParagraphBreak
附录 C 用\term[pu]{谱}[spectrum]组织一列空间及其结构映射。以序列模型写，数据具有形式 \(E_0,E_1,\ldots\) 及 \(\Sigma E_n\rightarrow E_{n+1}\)；对一个球面映射先悬挂，再复合结构映射，就得到 \(\pi_{n+k}(E_n)\rightarrow\pi_{n+k+1}(E_{n+1})\)。球面谱取 \(E_n=S^n\)，结构映射为 \(\Sigma S^n\cong S^{n+1}\)，因而恢复本节的系统。这种组织方式保留了每一步比较映射，适合把稳定类作为整体来研究。

\ProseParagraphBreak
第24章的向量丛稳定化则从固定底空间 \(X\) 上的丛出发，执行 \(V\mapsto V\oplus\underline{\mathbb C}\)，再考虑稳定同构与虚丛。它增加的是纤维维数，底空间仍是 \(X\)。悬挂稳定化改变空间及映射的维数，这两项操作的联系需要具体构造。例如球面上的复向量丛可由赤道粘接映射 \(S^{m-1}\rightarrow\operatorname{GL}_r(\mathbb C)\) 描述；添加平凡直线对应矩阵包含 \(A\mapsto\operatorname{diag}(A,1)\)，而悬挂球面对应赤道维数的改变。第24章将通过粘接与乘积定理研究这些变化之间的关系。本节给出的稳定范围，是球面映射在悬挂下最终不变的保证。
% 纲要条目排为编号小节标题，并进入目录与书签；说明段仍保留于纲要注释中。


\begin{chaptersummary}
待撰写
\end{chaptersummary}

\BookExercises
待撰写

% !TeX root = ../../Book2.tex
% 纲要标题：## 第10章　滤过复形与谱序列的构造
% 纲要位置：几何学计划纲要.md，第 1334 行。
% 正文状态：未撰写；本文件目前是框架。

% 纲要内容开始
% 
% > 本章目标与证明范围：本章负责代数构造和收敛，不同时承担全部 Serre 应用。必须把高阶微分写成代表元的提升—消去过程，并用隐藏延拓说明极限页尚未恢复的结构。
% 
% 纲要内容结束

\chapter{滤过复形与谱序列的构造}
\label{g2:ch:10}
\BookChapterNumber{g2:ch:10}

\begin{chapterabstract}
待撰写
\end{chapterabstract}

\begin{learninggoals}
待撰写
\end{learninggoals}

待撰写

% !TeX root = ../../Book2.tex
% 纲要标题：### 10.1 几何滤过到代数滤过
% 纲要位置：几何学计划纲要.md，第 1388 行。
% 正文状态：初稿；尚未正式审核。

% 纲要内容开始
% 
% * 骨架滤过与相对链
% * 递增、递减滤过及指标换算
% * 滤过复形与伴随分次
% * 双复形总化及符号
% 
% 纲要内容结束

\section{几何滤过到代数滤过}
\label{g2:sec:1001}

给空间逐层附着胞腔时，每一层都可以用相对同调描述，但这些层如何影响最后的同调，还取决于层与层之间的边界。滤过把“已经加入哪些部分”保存在链复形中；伴随分次先把各层分开，谱序列再逐步恢复层间的联系。本节先从骨架与相对链说明滤过的来源，再固定后面计算所用的指标和符号。

\subsection{骨架滤过与相对链}
\label{g2:subsec:1001:01}

设 \(X\) 是一个 \(N\) 维 CW 复形，记 \(X^p\) 为它的 \(p\) 骨架，并约定 \(X^{-1}=\varnothing\)、\(X^p=X\)（\(p\ge N\)）。以含幺交换环 \(R\) 为系数，奇异链群具有子复形序列
\[
	0\subset C_\bullet(X^0;R)\subset C_\bullet(X^1;R)\subset\cdots\subset C_\bullet(X;R).
\]
这里的包含是把同一个奇异单形看成到更大空间的映射。单形的各面仍落在原来的骨架中，所以边界保持各层。相邻两层的商恰好是相对链复形：
\[
	C_\bullet(X^p;R)/C_\bullet(X^{p-1};R)=C_\bullet(X^p,X^{p-1};R).
\]
\cref{g2:thm:skeletal-relative-homology}及系数变化的\cref{g2:cor:cellular-coefficients}表明，这些相对同调集中在次数 \(p\)，并由 \(p\) 胞腔生成。连接映射追踪胞腔边界如何落入低一层，最终组成胞腔边界。这就是通过分层计算整个空间的一种典型情形。

\ProseParagraphBreak

上链的方向与链相反：限制到较小子空间使信息减少，因此用“在已知部分上为零”定义递减的层。对奇异上链复形 \(C^\bullet(X;R)\)，令
\begin{equation}\label{g2:eq:skeletal-cochain-filtration}
	F^pC^n(X;R)=\ker\bigl(C^n(X;R)\longrightarrow C^n(X^{p-1};R)\bigr).
\end{equation}
于是 \(F^0C^n=C^n\)、\(F^{N+1}C^n=0\)，且 \(F^{p+1}C^n\subset F^pC^n\)。上边界 \(\mathrm d\varphi(\sigma)=\varphi(\partial\sigma)\) 保持这些子模，因为落在 \(X^{p-1}\) 内的单形，其边界也落在那里。

限制上链给出复形同构
\begin{equation}\label{g2:eq:skeletal-associated-graded}
	F^pC^\bullet(X;R)/F^{p+1}C^\bullet(X;R)
	\cong C^\bullet(X^p,X^{p-1};R).
\end{equation}
为核对满射，只需把右侧的上链在不属于 \(X^p\) 的奇异单形上赋值为零。右侧已在 \(X^{p-1}\) 上为零，所得延拓便属于 \(F^p\)；限制映射的核正是 \(F^{p+1}\)。这个赋值用于证明逐次满射，限制映射本身与上边界交换，因此诱导所写的商复形同构。

\ProseParagraphBreak

骨架每次只加入一个维数的胞腔，所以相对群格外集中。更一般的分层，例如把底空间骨架拉回到纤维化的总空间，每一层还携带纤维内部的上同调；这时层间关系会出现在不止一个次数。第11章将从这种几何滤过出发，本章先处理任何滤过复形共有的代数问题。

\subsection{递增、递减滤过及指标换算}
\label{g2:subsec:1001:02}

在模 \(M\) 上，\term[dizeng lvguo]{递增滤过}[increasing filtration]是一族 \(G_pM\subset G_{p+1}M\)；\term[dijian lvguo]{递减滤过}[decreasing filtration]是一族 \(F^{p+1}M\subset F^pM\)。指标取整数，改写 \(F^pM=G_{-p}M\) 就能交换两种记法。这里的指标只表示层数，与链或上链次数是两件事。

\ProseParagraphBreak

本章主要使用上链复形和递减滤过，并在二维表格中记
\[
	\text{滤过次数}=p,\qquad\text{总上链次数}=n=p+q.
\]
因此 \(q=n-p\) 是由前两者决定的辅助指标。总微分把 \(n\) 提高一；若某一步检测到的微分还把滤过次数提高 \(r\)，它的两个指标就必须改变为
\begin{equation}\label{g2:eq:spectral-bidegree-convention}
	(p,q)\longmapsto(p+r,q-r+1).
\end{equation}
箭头向右移动 \(r\) 格，再向下移动 \(r-1\) 格。总次数的增量始终是 \(1\)，而不是页数 \(r\)。

\ProseParagraphBreak

若资料使用递增滤过的链复形 \((C_m,\partial)\)，可先设 \(\widehat C^n=C_{-n}\)，把降一次的边界改写成升一次的微分，再设 \(F^p\widehat C^n=G_{-p}C_{-n}\)。原来的两个同调指标 \((a,b)\)、总次 \(a+b\)，对应上同调指标 \((p,q)=(-a,-b)\)。于是上式变回
\[
	\mathrm d_r:E^r_{a,b}\longrightarrow E^r_{a-r,b+r-1},
\]
它的总同调次数降低一。这解释了同调谱序列和上同调谱序列中箭头方向的差别。比较不同记号时，应先确认总次数以及滤过的增减方向，再换算两个坐标。

\subsection{滤过复形与伴随分次}
\label{g2:subsec:1001:03}

下面把几何例子中真正用到的性质抽出来。系数仍取含幺交换环 \(R\)，所有复形和同态均为 \(R\) 线性的。

\begin{definition}[滤过上链复形]\label{g2:def:filtered-cochain-complex}
	设 \((C^\bullet,\mathrm d)\) 为 \(R\) 模上链复形。一族子模 \(F^pC^n\) 若满足
	\[
		F^{p+1}C^n\subset F^pC^n,\qquad
		\mathrm d(F^pC^n)\subset F^pC^{n+1},
	\]
	就给出一个\term[lvguo shanglian fuxing]{滤过上链复形}[filtered cochain complex]。保持滤过的上链映射 \(f:C\to C'\) 指满足 \(f\mathrm d=\mathrm d'f\) 及 \(f(F^pC^n)\subset F^pC'^n\) 的次数零映射。\symidx{\(F^pC^n\)：上链复形的递减滤过}
\end{definition}

每个 \(F^pC^\bullet\) 都是子复形。要单独观察第 \(p\) 层，将更深的第 \(p+1\) 层压成零即可。

\begin{definition}[伴随分次复形]\label{g2:def:associated-graded-complex}
	滤过上链复形 \((C,F)\) 的\term[bansui fenci]{伴随分次}[associated graded]为
	\[
		\operatorname{gr}_F^pC^n=F^pC^n/F^{p+1}C^n,\qquad
		\operatorname{gr}_F C=\bigoplus_p\operatorname{gr}_F^pC.
	\]
	微分由 \(x+F^{p+1}\mapsto\mathrm dx+F^{p+1}\) 诱导，保持 \(p\)，并使总次数增加一。\symidx{\(\operatorname{gr}_F\)：滤过的伴随分次}
\end{definition}

若换代表 \(x\) 为 \(x+y\)、\(y\in F^{p+1}\)，微分的差 \(\mathrm dy\) 仍在 \(F^{p+1}\)，所以定义良好；平方为零来自 \(\mathrm d^2=0\)。一个保持滤过的上链映射也能在这些商上逐层诱导上链映射。

\begin{example}[被分次隐藏的微分]\label{g2:exm:graded-zero-nonzero-differential}
	取域 \(k\) 上的两项复形 \(C^0=ku\)、\(C^1=kv\)，并令 \(\mathrm du=v\)。赋予 \(u\) 滤过次数 \(0\)、\(v\) 滤过次数 \(2\)，即
	\[
		F^pC^0=\begin{cases}ku,&p\le0,\\0,&p>0,\end{cases}
		\qquad
		F^pC^1=\begin{cases}kv,&p\le2,\\0,&p>2.\end{cases}
	\]
	微分保持递减滤过，但它把 \(u\) 送到了更深的第二层。因此伴随分次中的微分为零，\(u,v\) 都留下非零分次上同调类；原复形的微分却是同构，全部上同调为零。层间跨越两格的这项微分，将在谱序列的第二页被检测到。
\end{example}

由此可见，把 \(C\) 换成 \(\operatorname{gr}_F C\) 会保留逐层数据，却暂时隐去进入更深层的微分。后面的各页正是通过恢复这些部分，逐步接近原复形的上同调。

\subsection{双复形总化及符号}
\label{g2:subsec:1001:04}

第7章的 Čech 与消解计算同时带有两个次数，并用总复形把两种微分合起来。为了直接追踪高阶微分，本章把符号吸收到竖直微分中，使两个方向反交换。

\begin{definition}[双复形与总复形]\label{g2:def:anticommuting-bicomplex}
	一个 \(R\) 模\term[shuang fuxing]{双复形}[bicomplex]由双分次模 \(B^{p,q}\) 及微分
	\[
		\mathrm d_{\mathrm h}:B^{p,q}\longrightarrow B^{p+1,q},\qquad
		\mathrm d_{\mathrm v}:B^{p,q}\longrightarrow B^{p,q+1}
	\]
	组成，满足
	\[
		\mathrm d_{\mathrm h}^{\,2}=\mathrm d_{\mathrm v}^{\,2}=0,\qquad
		\mathrm d_{\mathrm h}\mathrm d_{\mathrm v}+\mathrm d_{\mathrm v}\mathrm d_{\mathrm h}=0.
	\]
	其直和\term[zong fuxing]{总复形}[total complex]定义为
	\[
		\operatorname{Tot}^nB=\bigoplus_{p+q=n}B^{p,q},\qquad
		\mathrm d=\mathrm d_{\mathrm h}+\mathrm d_{\mathrm v}.
	\]
	若 \(p<0\) 或 \(q<0\) 时 \(B^{p,q}=0\)，则称为第一象限双复形。此时每个总次数的直和只有有限项。
\end{definition}

\begin{proposition}\label{g2:prop:total-complex-filtrations}
	双复形的总微分满足 \(\mathrm d^2=0\)。第一象限双复形的总复形具有两种每个次数内有限的递减滤过：
	\begin{equation}\label{g2:eq:column-row-filtrations}
		F^p\operatorname{Tot}^nB=\bigoplus_{a\ge p}B^{a,n-a},\qquad
		G^q\operatorname{Tot}^nB=\bigoplus_{b\ge q}B^{n-b,b}.
	\end{equation}
	第一种滤过的伴随分次微分为 \(\mathrm d_{\mathrm v}\)，第二种为 \(\mathrm d_{\mathrm h}\)。
\end{proposition}
\begin{proof}
	展开平方得到
	\[
		\mathrm d^2=\mathrm d_{\mathrm h}^{\,2}
		+\mathrm d_{\mathrm h}\mathrm d_{\mathrm v}
		+\mathrm d_{\mathrm v}\mathrm d_{\mathrm h}
		+\mathrm d_{\mathrm v}^{\,2}=0.
	\]
	两个微分都不降低水平指标，也都不降低竖直指标，因此保持各自的滤过。对 \(F\) 取相邻商后只保留水平次数 \(p\) 的一列；\(\mathrm d_{\mathrm h}\) 进入下一层而被商去，余下 \(\mathrm d_{\mathrm v}\)。对 \(G\) 同理。第一象限条件还给出 \(F^0\operatorname{Tot}^n=G^0\operatorname{Tot}^n=\operatorname{Tot}^n\) 及 \(F^{n+1}\operatorname{Tot}^n=G^{n+1}\operatorname{Tot}^n=0\)（\(n\ge0\)）。
\end{proof}

\cref{g2:lem:first-quadrant-comparison}使用的是交换微分 \(h,v\) 和公式 \(\mathrm d=h+(-1)^pv\)。在 \(B^{p,q}\) 上设
\[
	\mathrm d_{\mathrm h}=h,\qquad \mathrm d_{\mathrm v}=(-1)^pv,
\]
就得到本节约定：两个混合复合分别是 \((-1)^phv\) 和 \((-1)^{p+1}vh\)，和为零。因此两处的总复形完全相同；区别只是是否把 \((-1)^p\) 预先放进竖直箭头。

\ProseParagraphBreak

计算时，把 \(B^{p,q}\) 画在平面坐标 \((p,q)\) 上，\(\mathrm d_{\mathrm h}\) 向右、\(\mathrm d_{\mathrm v}\) 向上，而一条总次数为 \(n\) 的上链沿斜线 \(p+q=n\) 分布。先对列取上同调，再计算横向诱导微分，是滤过 \(F\) 的前两步；先对行计算则来自 \(G\)。如果一个横向像只是竖直边界，就能用右下方的分量抵消它，抵消过程继续向右下方延伸。这条阶梯将在下一节给出高阶微分的具体代表。
% 纲要条目排为编号小节标题，并进入目录与书签；说明段仍保留于纲要注释中。

% !TeX root = ../../Book2.tex
% 纲要标题：### 10.2 各页与高阶微分
% 纲要位置：几何学计划纲要.md，第 1395 行。
% 正文状态：未撰写；本文件目前是框架。

% 纲要内容开始
% 
% * 零页和第一页的构造
% * 延拓一个代表元到下一层
% * 微分次数和平方为零
% * 页间同调的证明
% * 滤过映射和自然性
% 
% 纲要内容结束

\section{各页与高阶微分}
\label{g2:sec:1002}

待撰写

\subsection{零页和第一页的构造}
\label{g2:subsec:1002:01}
待撰写

\subsection{延拓一个代表元到下一层}
\label{g2:subsec:1002:02}
待撰写

\subsection{微分次数和平方为零}
\label{g2:subsec:1002:03}
待撰写

\subsection{页间同调的证明}
\label{g2:subsec:1002:04}
待撰写

\subsection{滤过映射和自然性}
\label{g2:subsec:1002:05}
待撰写
% 纲要条目排为编号小节标题，并进入目录与书签；说明段仍保留于纲要注释中。

% !TeX root = ../../Book2.tex
% 纲要标题：### 10.3 正合偶与比较方法
% 纲要位置：几何学计划纲要.md，第 1403 行。
% 正文状态：未撰写；本文件目前是框架。

% 纲要内容开始
% 
% * 正合偶及导出偶
% * 与滤过构造的对应
% * 边缘同态和比较定理
% 
% 纲要内容结束

\section{正合偶与比较方法}
\label{g2:sec:1003}

待撰写

\subsection{正合偶及导出偶}
\label{g2:subsec:1003:01}
待撰写

\subsection{与滤过构造的对应}
\label{g2:subsec:1003:02}
待撰写

\subsection{边缘同态和比较定理}
\label{g2:subsec:1003:03}
待撰写
% 纲要条目排为编号小节标题，并进入目录与书签；说明段仍保留于纲要注释中。

% !TeX root = ../../Book2.tex
% 纲要标题：### 10.4 收敛和延拓问题
% 纲要位置：几何学计划纲要.md，第 1409 行。
% 正文状态：初稿；尚未正式审核。

% 纲要内容开始
% 
% * 有界及第一象限收敛
% * 穷尽、分离与完备性的角色
% * 极限页与目标的伴随分次
% * 加性隐藏延拓
% * 乘法谱序列与隐藏延拓
% 
% 纲要内容结束

\section{收敛和延拓问题}
\label{g2:sec:1004}

各页的构造说明了怎样不断修正一个代表元，还需要回答这些修正最终计算什么。最便于把握的情形是：在每个总次数内，滤过只有有限层。此时，足够高阶的“近似闭”已经成为真正的闭，所有可能影响它的边界也已经进入计算。极限页因而给出上同调的逐层商群；把这些商群拼回整个群及其乘法，则留下另一类问题。

\subsection{有界及第一象限收敛}
\label{g2:subsec:1004:01}

这里的有限性按总次数判断。对于固定的 \(n\)，只需存在整数 \(a_n<b_n\)，使 \(F^{a_n}C^n=C^n\)、\(F^{b_n}C^n=0\)；不同次数允许使用不同的界。这比整个复形共用一对上下界更适合第一象限双复形。

\begin{theorem}[有限滤过的收敛]\label{g2:thm:filtered-complex-convergence}
	设 \((C^\bullet,\mathrm d,F)\) 为带递减滤过的 \(R\) 模上链复形，且对每个 \(n\)，\(F^pC^n\) 在 \(p\) 足够小时等于 \(C^n\)，在 \(p\) 足够大时等于零。则对每个 \((p,q)\)，其谱序列 \(E_r^{p,q}\) 在 \(r\) 足够大时稳定，稳定值记为 \(E_\infty^{p,q}\)。在上同调上定义
	\begin{equation}\label{g2:eq:induced-cohomology-filtration}
		F^pH^n(C)\coloneqq\operatorname{im}\bigl(H^n(F^pC)\longrightarrow H^n(C)\bigr).
	\end{equation}
	这是每次数有限的递减滤过，并有自然同构
	\[
		E_\infty^{p,q}\cong F^pH^{p+q}(C)/F^{p+1}H^{p+q}(C).
	\]
	此时记作 \(E_r^{p,q}\Longrightarrow H^{p+q}(C)\)，称该谱序列\term[shoulian]{收敛}[convergence]到 \(H^*(C)\)。
\end{theorem}
\begin{proof}
	固定 \(p,n=p+q\)，记 \(Z^n=\ker(\mathrm d:C^n\to C^{n+1})\)、\(B^n=\mathrm dC^{n-1}\)。我们直接考察\secref{g2:sec:1002}的代表元公式。取 \(r\) 足够大，使
	\[
		F^{p+r}C^{n+1}=0,\qquad F^{p-r+1}C^{n-1}=C^{n-1}.
	\]
	于是分子中“\(\mathrm dx\in F^{p+r}\)”变为“\(\mathrm dx=0\)”，而边界的原像已经允许取遍 \(C^{n-1}\)。具体地，
	\[
		\begin{aligned}
			Z_r^{p,n}&=F^pC^n\cap Z^n,\\
			Z_{r-1}^{p+1,n}&=F^{p+1}C^n\cap Z^n,\\
			\mathrm d Z_{r-1}^{p-r+1,n-1}&=F^pC^n\cap B^n.
		\end{aligned}
	\]
	第三式来自 \(Z_{r-1}^{p-r+1,n-1}=\{y\in C^{n-1}:\mathrm dy\in F^pC^n\}\)。因此所有足够大的页都由同一个商群给出：
	\begin{equation}\label{g2:eq:stable-page-cycles-boundaries}
		E_\infty^{p,n-p}\cong
		\frac{F^pC^n\cap Z^n}{F^{p+1}C^n\cap Z^n+F^pC^n\cap B^n}.
	\end{equation}
	这些识别与页间同调相容，因为代表元已经是真正的圈，高阶微分均取其零微分。

	再把这个商与上同调滤过比较。\eqref{g2:eq:induced-cohomology-filtration}的等价表达是
	\[
		F^pH^n(C)=\frac{(F^pC^n\cap Z^n)+B^n}{B^n}.
	\]
	把 \(x\in F^pC^n\cap Z^n\) 送到 \([x]\) 再模去 \(F^{p+1}H^n(C)\)，得到满射。它的核由那些可以写成 \(x=z+\mathrm dy\) 的元素组成，其中 \(z\in F^{p+1}C^n\cap Z^n\)。因 \(x,z\in F^pC^n\)，必有 \(\mathrm dy\in F^pC^n\cap B^n\)，所以核正是\eqref{g2:eq:stable-page-cycles-boundaries}的分母。这证明所需同构。

	若 \(F^pC^n=C^n\)，则每个 \(n\) 次上同调类都有该层代表，故 \(F^pH^n=H^n\)；若 \(F^pC^n=0\)，则 \(F^pH^n=0\)。因此目标滤过仍然有限。滤过上链映射将上述圈、边界和它们的包含同时送到另一复形中的对应对象，故这些同构自然。
\end{proof}

\begin{corollary}[第一象限双复形]\label{g2:cor:first-quadrant-convergence}
	设 \(C^{p,q}\) 为第一象限反交换双复形，即 \(C^{p,q}=0\) 除非 \(p,q\ge0\)。令 \(\operatorname{Tot}^nC=\bigoplus_{p+q=n}C^{p,q}\)，并按第一指标递减滤过。则相应谱序列收敛到 \(H^*(\operatorname{Tot}C)\)。每格在有限页后稳定，每条总次数对角线上也只有有限个极限项。
\end{corollary}
\begin{proof}
	固定 \(n\ge0\)，对角线上只有 \((0,n),\ldots,(n,0)\) 这 \(n+1\) 个位置。因此 \(F^0\operatorname{Tot}^nC=\operatorname{Tot}^nC\)、\(F^{n+1}\operatorname{Tot}^nC=0\)，可用上一主定理。

	还可直接看出每格的稳定范围。\(\mathrm d_r\) 从 \((p,q)\) 去往 \((p+r,q-r+1)\)，当 \(r>q+1\) 时靶在第一象限外；进入此格的微分来自 \((p-r,q+r-1)\)，当 \(r>p\) 时源在第一象限外。故 \(r>\max\{p,q+1\}\) 后，两端微分都为零，该格不再变化。
\end{proof}

“每格最终不变”与“整个目标由有限个商群恢复”分别使用了箭头的范围和每个总次数上的有限滤过。给一张抽象谱序列画出第一象限，只能先得到前一项；要把它认作某个复形的上同调，还须说明该复形、所用总化及目标滤过。

\subsection{穷尽、分离与完备性的角色}
\label{g2:subsec:1004:02}

当滤过有无限多层时，“一直可以修正”与“存在一个修正完毕的元素”开始分开。下面三个条件分别控制哪些元素被看到、哪些元素能被区分，以及无穷次相容修正能否留下一个元素。

\begin{definition}[滤过的三种性质]\label{g2:def:filtration-exhaustive-separated-complete}
	设 \(M\) 为带递减滤过 \(F^pM\) 的 \(R\) 模。若 \(\bigcup_pF^pM=M\)，称滤过\term[qiongjin-lvguo]{穷尽}[exhaustive]；若 \(\bigcap_pF^pM=0\)，称滤过\term[fenli-lvguo]{分离}[separated]。令
	\[
		\widehat M=\varprojlim_{p\to+\infty}M/F^pM
		=\left\{(\overline x_p)_p:\overline x_{p+1}\bmod F^pM=\overline x_p\right\}.
	\]
	若自然映射 \(M\to\widehat M\)、\(x\mapsto(x\bmod F^pM)_p\) 为同构，称滤过\term[wanbei-lvguo]{完备}[complete]。
\end{definition}

上式中的逆极限就是一列相容的有限精度数据，无需预先掌握逆极限的其他性质。自然映射的核为 \(\bigcap_pF^pM\)，所以完备蕴含分离；其满射性保证每列相容数据来自一个实际元素。对上链复形，这些条件先逐次施加于 \(C^n\)，而上同调上的诱导滤过还需另外考察。

\ProseParagraphBreak

穷尽性防止对象有一部分始终在滤过之外。例如，令 \(C^0=k\)、其他次数为零，却取全部 \(F^pC=0\)，便会得到零谱序列而遗漏 \(H^0(C)=k\)。分离性则防止一个非零元素藏在所有层内：同一个复形若取全部 \(F^pC=C\)，每个相邻商仍为零，整个 \(k\) 都不可见。有限滤过自动排除这两种现象。

\begin{example}[无限修正与完备化]\label{g2:exm:incomplete-filtered-complex}
	设 \(k\) 为域，取两项复形
	\[
		C^0=k[t]\xrightarrow{\,1-t\,}C^1=k[t],
	\]
	并令 \(F^pC^i=t^pk[t]\)（\(p\ge0\)），\(F^pC^i=C^i\)（\(p<0\)）。它的滤过穷尽且分离。在每个相邻商中，乘以 \(1-t\) 的首项都是恒等映射，所以 \(E_1=0\)，所有后续页也为零。但直接计算给出
	\[
		H^0(C)=0,\qquad H^1(C)=k[t]/(1-t)\cong k.
	\]
	这个非零群的诱导滤过完全不分离：因为商中 \(t=1\)，每个 \(F^pH^1\) 都是整个 \(k\)。例如 \(1\) 与 \(t^N\) 之差是边界，故 \([1]\) 在任意深的层中都有代表。

	若想使 \(1\) 成为真正的边界，需要解 \((1-t)x=1\)。有限和 \(x_N=1+t+\cdots+t^{N-1}\) 满足 \((1-t)x_N=1-t^N\)，每提高一次精度都能消去一项，但极限 \(1+t+t^2+\cdots\) 不在 \(k[t]\) 中。完备化把系数模换成 \(k[[t]]\) 后，这个极限存在，乘以 \(1-t\) 成为同构，新的复形才无上同调。这里零谱序列对两个复形都相同，改变的是无限修正是否仍留在原复形内。
\end{example}

完备性解决相容极限的存在，却不单独保证边界像闭合。为看清这一区别，令 \(V=\bigoplus_{i\ge1}ke_i\)，取 \(M=V[[t]]\)，即系数在 \(V\) 中的形式幂级数，每一个系数都是有限线性组合。\(t\) 进滤过使 \(M\) 完备且分离：模 \(t^p\) 的相容数据逐个指定全部系数。定义连续 \(k[[t]]\) 线性映射 \(T:M\to M\)，使 \(T(e_i)=t^ie_i\)。它良定义，因为输出中任一固定 \(t\) 次数只会接收有限个较低次数的输入；按每个 \(e_i\) 的系数观察，\(T\) 就是乘以 \(t^i\)，所以它单射。

\ProseParagraphBreak

在两项复形 \(M\xrightarrow{T}M\) 中，元素 \(y=\sum_{i\ge1}t^ie_i\) 不是边界：其唯一可能原像的常数系数应为 \(\sum_{i\ge1}e_i\)，不属于 \(V\)。但任意有限部分 \(y_N=\sum_{i=1}^Nt^ie_i\) 都是 \(T(\sum_{i=1}^Ne_i)\)，且 \(y-y_N\in F^{N+1}M\)。因此 \([y]\ne0\) 却属于全部 \(F^pH^1\)。这说明即使链级各项穷尽、分离、完备，上同调滤过仍可能不分离。一般无限滤过的收敛需要继续控制这样的边界与极限问题；本章的逐总度有限条件使全部过程在有限步内结束。

\subsection{极限页与目标的伴随分次}
\label{g2:subsec:1004:03}

\cref{g2:thm:filtered-complex-convergence}中的极限同构可以直接理解为“首个非零滤过层”。设 \([x]\in H^n(C)\) 可由 \(x\in F^pC^n\) 中的圈代表。它在 \(E_\infty^{p,n-p}\) 中为零，当且仅当它还可以用更深一层的圈表示，即存在 \(z\in F^{p+1}C^n\cap Z^n\) 和 \(y\in C^{n-1}\)，使
\[
	x=z+\mathrm dy.
\]
因而极限页在第 \(p\) 层记录的是“能在这一层出现、却不能移到下一层”的上同调信息。

\ProseParagraphBreak

这一解释也说明代表元的两种不确定性。加入更深一层的圈，不改变当前的首项；加入任意边界，不改变上同调类。\eqref{g2:eq:stable-page-cycles-boundaries}正好把这两种变化同时取商。于是收敛所给出的完整陈述是
\[
	\operatorname{gr}_F H^n(C)
	=\bigoplus_pF^pH^n(C)/F^{p+1}H^n(C)
	\cong\bigoplus_{p+q=n}E_\infty^{p,q}.
\]
这里的直和发生在伴随分次对象中。把右侧直接写成 \(H^n(C)\) 还需要为各层挑选相容的提升。

\begin{corollary}[维数与单层检测]\label{g2:cor:spectral-sequence-dimension}
	在\cref{g2:thm:filtered-complex-convergence}的条件下，若系数为域 \(k\)，且固定总次数 \(n\) 上的极限项均有限维，则
	\[
		\dim_kH^n(C)=\sum_{p+q=n}\dim_kE_\infty^{p,q}.
	\]
	若某一总次数上只有一个可能非零的极限项 \(E_\infty^{p,n-p}\)，则 \(H^n(C)\) 与该项自然同构；这一结论对一般 \(R\) 模也成立。
\end{corollary}
\begin{proof}
	有限滤过中每个相邻商的维数可加，沿全部层相加即得公式。若只有第 \(p\) 层的商可能非零，则其前面的相邻包含全为等号，故 \(F^pH^n=H^n\)；其后面的相邻包含也全为等号，直到滤过为零，故 \(F^{p+1}H^n=0\)。于是极限同构已经是 \(H^n\cong E_\infty^{p,n-p}\)，没有选择提升的步骤。
\end{proof}

\subsection{加性隐藏延拓}
\label{g2:subsec:1004:04}

知道全部相邻商以后，恢复一个有限滤过群需要逐层求解短正合列。为简化记号，把滤过重编号为
\[
	H=F^0H\supset F^1H\supset\cdots\supset F^{N+1}H=0.
\]
令 \(G^p=F^pH/F^{p+1}H\)，则每一步都有
\begin{equation}\label{g2:eq:hidden-additive-extension}
	0\longrightarrow F^{p+1}H\longrightarrow F^pH\longrightarrow G^p\longrightarrow0.
\end{equation}
已知两端而恢复中间对象的任务称为\term[yincang-yantuo]{隐藏延拓}[hidden extension]问题；这里首先讨论其加性结构。

\begin{proposition}[自由分次块的分裂]\label{g2:prop:free-graded-splitting}
	设一个 \(R\) 模 \(H\) 具有有限递减滤过，且每个相邻商 \(G^p\) 都是自由 \(R\) 模。则存在与滤过相容的模同构
	\[
		H\cong\bigoplus_pG^p.
	\]
	该同构一般依赖各商模的基及其提升，不由滤过唯一指定。特别地，域上的有限滤过总能作为向量空间分裂。
\end{proposition}
\begin{proof}
	在\eqref{g2:eq:hidden-additive-extension}中，对 \(G^p\) 的每个基元任取一个 \(F^pH\) 中的原像，线性延拓得到截面 \(s_p:G^p\to F^pH\)。设商映射为 \(q_p\)，则 \(q_ps_p=\operatorname{id}\)。任意 \(h\in F^pH\) 唯一写成
	\[
		h=\bigl(h-s_pq_p(h)\bigr)+s_pq_p(h),
	\]
	第一项属于 \(F^{p+1}H\)，两项所属子模交为零。因此 \(F^pH=F^{p+1}H\oplus s_p(G^p)\)。从最深非零层向外递推，得到所述同构，并把每个 \(F^pH\) 识别为指标至少为 \(p\) 的各项之和。

	若把一个基元的提升加上 \(F^{p+1}H\) 的元素，它仍为同一基元的提升，却一般改变截面。故论证给出的是依赖选择的分裂。
\end{proof}

对含挠群的商，基元提升可能满足新的倍数关系，因而中间群可能不分裂。即使自由性已经保证加性分裂，选出的截面也未必保持其他映射或乘法。下一节将把这些不同层次的选择落实为可直接计算的群与滤过映射。

\subsection{乘法谱序列与隐藏延拓}
\label{g2:subsec:1004:05}

许多几何应用同时保留杯积。代数上，设 \(C^\bullet\) 具有结合的乘法及单位元 \(1\in C^0\)，其中 \(1\in F^0C^0\)，\(C^nC^m\subset C^{n+m}\)，微分满足
\[
	\mathrm d(xy)=(\mathrm dx)y+(-1)^nx\,\mathrm dy\qquad(x\in C^n),
\]
并且 \(F^pC\cdot F^sC\subset F^{p+s}C\)。这种数据称为\term[chengfa-lvguo]{乘法滤过}[multiplicative filtration]的微分分次代数。乘法把每层的信息送到滤过指标之和所在的层，故也能随各页传递。

\begin{proposition}[谱序列的乘法]\label{g2:prop:multiplicative-spectral-sequence}
	设 \(R\) 为含幺交换环，\(C^\bullet\) 为带乘法滤过的微分分次 \(R\) 代数。则每页都有由代表元乘法诱导的结合乘法
	\[
		E_r^{p,q}\otimes_RE_r^{s,t}\longrightarrow E_r^{p+s,q+t},
	\]
	且 \(\mathrm d_r\) 满足总次数符号的 Leibniz 公式
	\[
		\mathrm d_r(uv)=(\mathrm d_ru)v+(-1)^{p+q}u\,\mathrm d_rv.
	\]
	页间同调同构保持乘法。若滤过逐总度有限，则极限同构保持乘法，因而 \(E_\infty\cong\operatorname{gr}_FH^*(C)\) 为双分次代数同构。若 \(C\) 的乘法分次交换，则各页乘法也分次交换。
\end{proposition}
\begin{proof}
	第零页直接对乘法及微分取相邻商即可。对于 \(r\ge1\)，令 \(x\in Z_r^{p,n}\)、\(y\in Z_r^{s,m}\)。由滤过乘法及 Leibniz 公式，\(xy\in F^{p+s}C^{n+m}\)，且
	\[
		\mathrm d(xy)\in F^{p+r}C^{n+1}F^sC^m+F^pC^nF^{s+r}C^{m+1}
		\subset F^{p+s+r}C^{n+m+1}.
	\]
	故 \(xy\in Z_r^{p+s,n+m}\)。

	为核对选择无关性，先在 \(x\) 上加入 \(u\in Z_{r-1}^{p+1,n}\)。这时 \(uy\in F^{p+s+1}\)，并且 \(\mathrm d(uy)\in F^{p+s+r}\)，所以 \(uy\in Z_{r-1}^{p+s+1,n+m}\)，属于乘积所在页的分母。再在 \(x\) 上加入 \(\mathrm dv\)，其中 \(v\in Z_{r-1}^{p-r+1,n-1}\)。有
	\[
		(\mathrm dv)y=\mathrm d(vy)-(-1)^{n-1}v\,\mathrm dy.
	\]
	这里 \(vy\in F^{p+s-r+1}\)，且 \(\mathrm d(vy)\in F^{p+s}\)，故 \(vy\in Z_{r-1}^{p+s-r+1,n+m-1}\)，第一项是允许的边界。第二项满足 \(v\,\mathrm dy\in F^{p+s+1}\) 及 \(\mathrm d(v\,\mathrm dy)=(\mathrm dv)(\mathrm dy)\in F^{p+s+r}\)，故也是允许的更深层圈。改变 \(y\) 的代表时作同样计算，便知乘法良定义。

	结合律由代表元结合律继承，单位元由 \(1\) 的页类给出；\(\mathrm d_r\) 用 \(\mathrm d\) 表示，所以直接将原 Leibniz 公式取页类，就得到所述符号。两个能提升到下一页的代表元相乘后仍能提升，且其类始终由同一个乘积表示，因此\cref{g2:thm:spectral-sequence-page-transition}的同构保持乘法。

	在收敛条件下，目标滤过的乘法条件来自圈的乘法：\(F^pH^n\cdot F^sH^m\subset F^{p+s}H^{n+m}\)。极限同构把 \(x\) 的页类送到 \([x]\) 的相邻商类，所以也把 \(xy\) 送到 \([x][y]\)，从而保持乘法和单位元。分次交换性同样由代表元的对应等式继承。
\end{proof}

由此得到的仍是伴随分次环。若 \(x\in F^pH\)、\(y\in F^sH\) 的首项乘积为零，只能推出 \(xy\in F^{p+s+1}H\)，其在更深层中的值尚未确定。这就是乘法隐藏延拓的来源。

\begin{example}[同一分次环中的不同平方]\label{g2:exm:hidden-multiplicative-extension}
	设 \(k\) 为域，比较两个三维 \(k\) 代数
	\[
		A=k[x]/(x^3),\qquad B=k[x,y]/(x^2,xy,y^2).
	\]
	在 \(A\) 中记 \(y=x^2\)；两边均以 \(1,x,y\) 为向量空间基。分别取 \(F^0=A\) 或 \(B\)、\(F^1=ky\)、\(F^2=0\)，并向两侧常值延拓。这都是乘法滤过。两者的零次相邻商均为 \(k[x]/(x^2)\)，一次相邻商均为 \(ky\)，\(x\) 在一次商上的作用都为零。因此它们的伴随分次环同构，加性滤过也同构。

	但 \(A\) 中 \(x^2=y\ne0\)、\(x^3=0\)，而 \(B\) 中每个幂零元素都是 \(ax+by\)，平方均为零。因幂零阶在代数同构下保持，两环不同构。将这两个代数分别看成集中在上链次数零、微分为零的复形，会得到从零页起就退化且逐页相同的乘法谱序列；其目标环仍有上述差异。\(A\) 中 \(x^2\) 的首项在 \(F^0/F^1\) 中消失，实际乘积却进入了 \(F^1\)。
\end{example}

在几何计算中，解除这种不确定性通常要重新选取实际上同调类，利用限制映射、已知子空间上的杯积、与基本类的配对，或另一个可比较的滤过。第4章计算复射影空间上同调环时，确定 \(u^2\) 是顶次生成元，正是靠局部外积配对排除了可能的零值与非单位倍数（\cref{g2:thm:complex-projective-ring}）。谱序列提供逐层候选，额外的几何计算则确定这些候选在目标中的具体提升与乘积。
% 纲要条目排为编号小节标题，并进入目录与书签；说明段仍保留于纲要注释中。

% !TeX root = ../../Book2.tex
% 纲要标题：### 10.5 隐藏延拓与比较计算
% 纲要位置：几何学计划纲要.md，第 1417 行。
% 正文状态：未撰写；本文件目前是框架。

% 纲要内容开始
% 
% * 同一伴随分次的不同交换群
% * 两种滤过计算同一总同调
% * 退化不意味着自然分裂
% 
% 纲要内容结束

\section{隐藏延拓与比较计算}
\label{g2:sec:1005}

待撰写

\subsection{同一伴随分次的不同交换群}
\label{g2:subsec:1005:01}
待撰写

\subsection{两种滤过计算同一总同调}
\label{g2:subsec:1005:02}
待撰写

\subsection{退化不意味着自然分裂}
\label{g2:subsec:1005:03}
待撰写
% 纲要条目排为编号小节标题，并进入目录与书签；说明段仍保留于纲要注释中。

% !TeX root = ../../Book2.tex
% 纲要标题：### 10.6 非零 d₂ 的双复形模型
% 纲要位置：几何学计划纲要.md，第 1423 行。
% 正文状态：初稿；尚未正式审核。

% 纲要内容开始
% 
% 取第一象限双复形 C^{0,1}=C^{1,0}=C^{1,1}=C^{2,0}=k，令从(1,0)到(1,1)、从(0,1)到(1,1)、从(1,0)到(2,0)的三个微分为恒等，其余为零。逐页计算、提升代表元并直接核验总复形无上同调；符号按本卷反交换约定统一。此节是一段连续计算，不再人为分小节。
% 
% ---
% 
% 纲要内容结束

\section{非零 \texorpdfstring{\(\mathrm{d}_2\)}{d₂} 的双复形模型}
\label{g2:sec:1006}

高阶微分可以出现在只有水平、竖直两种微分的双复形中。关键在于一个横向像虽然在纵向上同调中为零，用来消去它的纵向原像还可能产生新的横向像。下面把这个过程放在四个一维向量空间中，从零页一直算到总上同调。

\ProseParagraphBreak

设 \(k\) 为任意域，双复形的非零项为
\[
	C^{0,1}=ku,\qquad C^{1,0}=kv,\qquad C^{1,1}=kw,\qquad C^{2,0}=kz.
\]
选定各自的基，规定
\begin{equation}\label{g2:eq:nonzero-d2-model}
	\mathrm d_{\mathrm h}u=w,\qquad \mathrm d_{\mathrm v}v=w,\qquad \mathrm d_{\mathrm h}v=z,
\end{equation}
其余微分均为零。三个非零箭头在这些基下都是恒等映射。把上排视为 \(q=1\)、下排视为 \(q=0\)，三列依次为 \(p=0,1,2\)，图为
\[
	\begin{tikzcd}[column sep=large,row sep=large]
		ku \arrow[r,"1"] & kw & 0 \\
		0 & kv \arrow[u,"1"] \arrow[r,"1"'] & kz.
	\end{tikzcd}
\]
任意两个可复合的微分都为零：一旦到达 \(w\) 或 \(z\)，已经没有非零的后续箭头。因此 \(\mathrm d_{\mathrm h}^2=\mathrm d_{\mathrm v}^2=0\) 且 \(\mathrm d_{\mathrm h}\mathrm d_{\mathrm v}+\mathrm d_{\mathrm v}\mathrm d_{\mathrm h}=0\)，这确实是本章约定的反交换双复形。

\ProseParagraphBreak

先把目标算出来，便于之后逐页核对。总复形只有两个非零次数：
\begin{equation}\label{g2:eq:nonzero-d2-total-complex}
	0\longrightarrow ku\oplus kv\xrightarrow{\mathrm d}kw\oplus kz\longrightarrow0,
	\qquad [\mathrm d]=\begin{pmatrix}1&1\\0&1\end{pmatrix}.
\end{equation}
左项在次数 \(1\)，右项在次数 \(2\)，因为 \(\mathrm du=w\)、\(\mathrm dv=w+z\)。该矩阵行列式为 \(1\)，所以微分是同构，\(H^n(\operatorname{Tot}C)=0\) 对所有 \(n\) 成立。谱序列也应当最后把所有类消去。

\ProseParagraphBreak

按列取递减滤过 \(F^p\operatorname{Tot}^nC=\bigoplus_{a\ge p}C^{a,n-a}\)。在两个非零次数，它分别是
\[
	F^0\operatorname{Tot}^1=ku\oplus kv,\quad F^1\operatorname{Tot}^1=kv,\quad F^2\operatorname{Tot}^1=0,
\]
\[
	F^0\operatorname{Tot}^2=F^1\operatorname{Tot}^2=kw\oplus kz,\quad F^2\operatorname{Tot}^2=kz,\quad F^3\operatorname{Tot}^2=0.
\]
这些式子也直接验证了总微分保持滤过。

\ProseParagraphBreak

零页就是原来的四项双分次向量空间，\(\mathrm d_0=\mathrm d_{\mathrm v}\)。中间列的 \(kv\xrightarrow{1}kw\) 是同构，取列上同调后两个位置都消失；其余两列各留下一项。因此
\[
	E_1^{0,1}=k[u]_1,\qquad E_1^{2,0}=k[z]_1,
\]
其他项全为零。第一页微分 \(\mathrm d_1\) 的双次数为 \((1,0)\)，两个非零位置之间没有这种箭头，所以 \(\mathrm d_1=0\)、\(E_2=E_1\)。用代表元看，\(\mathrm d_{\mathrm h}u=w\) 在中间列是纵向边界，故 \(\mathrm d_1[u]_1=0\)。

\ProseParagraphBreak

这个消失允许我们修正 \(u\)。取 \(x_1=-v\)，就有
\[
	\mathrm d_{\mathrm h}u+\mathrm d_{\mathrm v}(-v)=w-w=0.
\]
于是修正后的总上链 \(x=u-v\) 满足
\begin{equation}\label{g2:eq:nonzero-d2-corrected-representative}
	\mathrm d(u-v)=w-(w+z)=-z\in F^2\operatorname{Tot}^2C.
\end{equation}
因此它属于 \(Z_2^{0,1}\)，是第二页中从 \([u]_1\) 延续下来的类。按\cref{g2:prop:spectral-sequence-differential}，
\begin{equation}\label{g2:eq:nonzero-d2-value}
	\mathrm d_2:E_2^{0,1}\longrightarrow E_2^{2,0},\qquad [u-v]_2\longmapsto-[z]_2.
\end{equation}
两边均为一维，\(-1\) 在任意域中都是单位；特征 \(2\) 时它等于 \(1\)，映射仍然非零。因此 \(\mathrm d_2\) 是同构，第三页全部为零，此后各页也全为零。

\ProseParagraphBreak

这项微分的两个作用都能在总复形中看见。源中的 \(u\) 无法改成总闭元：要消掉 \(w\) 分量，修正系数必须是 \(-v\)，但剩下的 \(-z\) 已经没有更深的次数 \(1\) 上链可供抵消，因为 \(F^2\operatorname{Tot}^1=0\)。靶中的 \(z\) 本来是一个总闭元，却满足 \(z=\mathrm d(v-u)\)，所以它是总边界；这个边界的原像位于较浅的滤过层，直到第二页才进入当前层的关系中。一张 \(\mathrm d_2\) 同时删除了不能继续提升的源类和本来就应被识别为边界的靶类。

\ProseParagraphBreak

全部页面可归纳为
\[
	\begin{array}{c|c|c}
		\text{页面}&\text{非零位置}&\text{非零微分}\\ \hline
		E_0&(0,1),(1,0),(1,1),(2,0)&\mathrm d_0(v)=w\\
		E_1&(0,1),(2,0)&\text{无}\\
		E_2&(0,1),(2,0)&\mathrm d_2[u-v]_2=-[z]_2\\
		E_3=E_\infty&\text{无}&\text{无}
	\end{array}
\]
第一象限收敛定理给出的 \(E_\infty=\operatorname{gr}_F H^*(\operatorname{Tot}C)=0\)，与\eqref{g2:eq:nonzero-d2-total-complex}的直接计算吻合。若改按行滤过，上排 \(ku\to kw\) 和下排 \(kv\to kz\) 都已经是同构，于是另一条谱序列在第一页就为零。两种滤过改变的是计算中何时看见这些抵消，总复形及其上同调始终相同。
% 本节沿用纲要的连续叙述安排，不另造小节标题。



\begin{chaptersummary}
待撰写
\end{chaptersummary}

\BookExercises
待撰写

% !TeX root = ../../Book2.tex
% 纲要标题：## 第11章　Serre 谱序列与纤维化计算
% 纲要位置：几何学计划纲要.md，第 1429 行。
% 正文状态：初稿；尚未正式审核。

% 纲要内容开始
% 
% > 本章目标与证明范围：由基底骨架拉回滤过构造 Serre 谱序列，先解释局部系数再做计算。代表定理包括乘法、传递和所选有限生成结论；不扩充为所有谱序列的名录。[R1]
% 
% 纲要内容结束

\chapter{Serre 谱序列与纤维化计算}
\label{g2:ch:11}
\BookChapterNumber{g2:ch:11}

\begin{chapterabstract}
	基底的骨架把纤维化的总空间分成逐层可处理的部分。本章从这一几何滤过构造 Serre 谱序列，以沿路运输解释第二页的局部系数，并证明乘法、边缘映射、传递与比较。Hopf 丛、定向球丛、路径纤维化及酉群的计算将非零微分与环结构联系起来；纤维定向类的传递给出 Euler 类，回路球面的计算产生整数除幂代数。最后利用有限生成群在谱序列中的封闭性质，以及逐次消去高阶同伦群的有限截段，证明单连通 CW 复形的同调逐次有限生成当且仅当其同伦群逐次有限生成。
\end{chapterabstract}

\begin{learninggoals}
	\begin{itemize}
		\item 由基底骨架逆像构造几何滤过，识别相对群、运输作用及第二页，并说明固定总次数的收敛。
		\item 使用总次数的 Leibniz 规则和两条边缘映射，把页面上的微分与投影、限制及延拓障碍对应起来。
		\item 计算 Hopf 丛与球丛，定义 Euler 类并证明相应上同调正合列。
		\item 求出回路球面的加性群与所选杯积环，证明酉群、特殊酉群的整数外代数结构。
		\item 由谱序列证明有限生成性在纤维化中的传播及单连通空间的 Serre 有限生成定理。
	\end{itemize}
\end{learninggoals}

第8章的纤维化长正合列把 \(F\to E\to B\) 的同伦群联在一起。对同调，直接写出同样一条长正合列通常行不通。最简单的乘积已经说明原因：\(S^2\times S^1\) 的三次上同调来自一个二次类与一个一次类相乘，它同时使用基底和纤维的信息。Serre 谱序列先把这类混合信息放在坐标 \((2,1)\)，再检查总空间的粘合是否允许它保留下来。

\ProseParagraphBreak

比较两个纤维化
\[
	S^1\longrightarrow S^2\times S^1\longrightarrow S^2,
	\qquad S^1\longrightarrow S^3\longrightarrow S^2.
\]
它们的基底与纤维相同，因此第二页都有四个整数生成元；总空间的上同调却不同。第一张纤维化中，圆周类能够延拓到整个乘积；第二张中，它的延拓受阻，产生从 \((0,1)\) 到 \((2,0)\) 的非零微分。这是本章贯穿的例子：第二页描述局部材料，高阶微分记录这些材料在整体上拼接时遇到的关系。

\ProseParagraphBreak

本章主要沿用第10章的上同调指标 \(\mathrm d_r:(p,q)\mapsto(p+r,q-r+1)\)。有限生成证明另用同调版，其箭头改为 \((p,q)\mapsto(p-r,q+r-1)\)，并用 \(E^r_{p,q}\) 区分页面。局部系统沿用第5章“先走 \(\alpha\) 再走 \(\beta\)”的运输约定；第8章提供提升及路径纤维化，第9章提供 CW 逼近、Hurewicz 与消去同伦群的胞腔构造。关于谱序列及其几何应用的延伸阅读，可参看\cite[第三章，Spectral Sequences and Applications]{BottTu1982}。

% !TeX root = ../../Book2.tex
% 纲要标题：### 11.1 几何滤过及第二页
% 纲要位置：几何学计划纲要.md，第 1383 行。
% 正文状态：未撰写；本文件目前是框架。

% 纲要内容开始
% 
% * 基底骨架的逆像
% * 纤维局部平凡模型与相对群
% * 运输作用及局部系数
% * 第二页的识别与收敛条件
% 
% 纲要内容结束

\section{几何滤过及第二页}
\label{g2:sec:1101}

待撰写

\subsection{基底骨架的逆像}
\label{g2:subsec:1101:01}
待撰写

\subsection{纤维局部平凡模型与相对群}
\label{g2:subsec:1101:02}
待撰写

\subsection{运输作用及局部系数}
\label{g2:subsec:1101:03}
待撰写

\subsection{第二页的识别与收敛条件}
\label{g2:subsec:1101:04}
待撰写
% 纲要条目排为编号小节标题，并进入目录与书签；说明段仍保留于纲要注释中。

% !TeX root = ../../Book2.tex
% 纲要标题：### 11.2 乘积、传递与自然性
% 纲要位置：几何学计划纲要.md，第 1440 行。
% 正文状态：初稿；尚未正式审核。

% 纲要内容开始
% 
% * 杯积和 Leibniz 规则
% * 边缘映射的几何意义
% * 传递映射的代表元解释
% * 映射纤维化之间的比较
% 
% 纲要内容结束

\section{乘积、传递与自然性}
\label{g2:sec:1102}

第二页给出了许多群，真正的计算还需要它们之间的联系。乘法使一个生成元上的微分决定整片表格；边缘映射把表格接回纤维包含与基底投影；传递则把纤维上的一个类转为基底上的障碍。下面始终使用含幺交换系数环 \(R\)，谈到零行的普通系数时另假设纤维道路连通。

\subsection{杯积和 Leibniz 规则}
\label{g2:subsec:1102:01}

纤维的杯积与沿路运输交换，所以 \(\mathcal H^*(F;R)\) 是分次代数的局部系统。在局部平凡的记法中，\(a\otimes u\) 表示一个基底 \(p\) 次类配上一个纤维 \(q\) 次类。将两个这样的类相乘，必须先把第一个纤维因子越过第二个基底因子。因此第二页的乘法为
\begin{equation}\label{g2:eq:serre-product-sign}
	(a\otimes u)(b\otimes v)=(-1)^{q s}(a\smile b)\otimes(u\smile v),\qquad |a|=p,\ |u|=q,\ |b|=s.
\end{equation}
这是有张量识别时的公式。一般局部系数情形，则先按第5章的道路运输把系数移到同一纤维，再作纤维杯积；同一个 \((-1)^{qs}\) 仍须保留。

\begin{theorem}[Serre 谱序列的乘法]\label{g2:thm:serre-multiplicative}
	上同调 Serre 谱序列从第二页起是按总次数分次交换的代数谱序列。对 \(x\in E_r^{p,q}\)、\(y\in E_r^{s,t}\)，有
	\begin{equation}\label{g2:eq:serre-leibniz}
		\mathrm d_r(xy)=(\mathrm d_rx)y+(-1)^{p+q}x\,\mathrm d_ry.
	\end{equation}
	下一页的乘法由本页乘法诱导；第二页乘法由局部系数杯积和\cref{g2:eq:serre-product-sign}给出。目标上同调的滤过满足 \(F^pH^m(E)F^sH^n(E)\subset F^{p+s}H^{m+n}(E)\)，且极限页作为代数同构于这个环的伴随分次。
\end{theorem}
\begin{proof}
	先把几何滤过落实为一个可作乘法的链模型。用\cref{g2:lem:serre-filtered-cw-model}的相容 CW 逼近 \(X_p\to E_p\)，记 \(X=\bigcup_pX_p\)、\(C_*=C_*^{\mathrm{cell}}(X;\mathbb Z)\)。给每个胞腔记录它首次出现的层数；\(C_*(X_p)\) 即由层数不超过 \(p\) 的胞腔生成。积空间使用胞腔积的弱拓扑，积胞腔的层数是两因子层数之和。第4章的交叉积将其胞腔链识别为 \(C_*\otimes C_*\)，边界为 \(\partial(c\otimes d)=\partial c\otimes d+(-1)^{|c|}c\otimes\partial d\)。

	取基底对角映射的胞腔逼近 \(\Delta_c:B\to B\times B\)。它满足
	\[
		\Delta_c(B^k)\subset(B\times B)^k=\bigcup_{i+j\le k}B^i\times B^j.
	\]
	在 Hurewicz 替换中，将 \(\Delta_B\simeq\Delta_c\) 的同伦从总空间的真实对角映射开始提升，得到 \(\widetilde\Delta:E\to E\times E\)，且 \(\widetilde\Delta(E_k)\subset\bigcup_{i+j\le k}E_i\times E_j\)。在上述 CW 模型上逐层比较并作胞腔逼近，得到链映射
	\[
		D:C_*\longrightarrow C_*\otimes C_*,\qquad
		D(C_*(X_k))\subset\sum_{i+j\le k}C_*(X_i)\otimes C_*(X_j).
	\]
	这里逐层比较的依据是：相容 CW 逼近在每层都是弱同伦等价；对一个待延拓的胞腔，比较映射及比较同伦的障碍是该弱等价的一个相对同伦类，因而为零。按胞腔维数延拓，并在已处理子复形上固定选择，就得到相容比较。对有限个积因子的并也同样逐个积胞腔处理。这说明所用链对角确实保留总层数，而不只是普通上同调上的某个对角。

	令 \(C^*=\operatorname{Hom}(C_*,R)\)，滤过取在 \(C_*(X_{p-1})\) 上为零的上链。定义 \(uv=(u\otimes v)D\)。若 \(u\in F^p\)、\(v\in F^s\)，对层数小于 \(p+s\) 的任一积胞腔，至少有一个因子被相应上链消去，所以 \(uv\in F^{p+s}\)。\(D\) 为链映射又给出严格的上链恒等式
	\[
		\mathrm d(uv)=(\mathrm du)v+(-1)^{|u|}u\,\mathrm dv.
	\]
	若 \(u,v\) 是第 \(r\) 页的代表，即 \(\mathrm du\in F^{p+r}\)、\(\mathrm dv\in F^{s+r}\)，右端属于 \(F^{p+s+r}\)。更换代表时，乘法与上式使新增项落入第10章商公式的分母：多一层的近似闭元仍多一层，上边界乘积则改写为一个上边界加一个更高滤过项。因此乘积良定义，并且直接取上式在目标商中的类便得到\cref{g2:eq:serre-leibniz}。此论证对每个 \(r\) 适用，无需从第一微分的导子性质猜测后续微分。

	还要说明结合律与交换律为何在第二页成立。对角的两种三重迭代、交换两因子的对角，以及不同胞腔逼近，均由真实对角提供同伦。把同伦作相对胞腔逼近时，\(B^k\times I\) 只需进入总骨架 \(k+1\)。故所得链同伦提高层数至多一；对偶的上链同伦 \(K\) 降低滤过至多一。对于第 \(r\ge2\) 页的闭代表，总滤过为 \(P\) 时，误差形如 \(\mathrm dK(w)+K(\mathrm dw)\)。第一项是来自 \(F^{P-1}\subset F^{P-r+1}\) 的允许边界，第二项属于 \(F^{P+r-1}\subset F^{P+1}\)，两者在该页均为零。这证明从第二页起的结合律、单位律和总次数交换律。链对角代表真实对角，所以在总上同调上得到通常杯积。

	最后，在一个 \(p\) 胞腔与一个 \(s\) 胞腔上，用基底相对定向类乘纤维类识别第一页。交换中间的 \(q\) 次纤维因子和 \(s\) 次基底因子给出 \((-1)^{qs}\)，故第二页为所述乘法。对真闭代表同样取乘积，所得滤过和极限页乘积便是目标杯积的伴随分次。
\end{proof}

计算时需要区分两个次数：乘积交叉符号是 \((-1)^{qs}\)，导子公式的符号是 \((-1)^{p+q}\)。后者使用整个类的总次数。

\subsection{边缘映射的几何意义}
\label{g2:subsec:1102:02}

设 \(F\xrightarrow{i}E\xrightarrow{p}B\)，其中 \(B,F\) 道路连通。第二页零行是 \(H^*(B;R)\)，零列是纤维上同调中沿全部闭路运输不变的部分：
\[
	E_2^{n,0}=H^n(B;R),\qquad E_2^{0,n}=H^n(F;R)^{\pi_1(B)}.
\]
零行上的类只能被别的类击中，零列上的类只能向外发出微分。因此前者逐页取商，后者逐页取子群。

\begin{proposition}[边缘映射]\label{g2:prop:serre-edge-maps}
	Serre 谱序列中的两条边缘分别分解投影拉回与纤维限制：
	\[
		H^n(B;R)\twoheadrightarrow E_\infty^{n,0}=F^nH^n(E;R)\hookrightarrow H^n(E;R),
	\]
	\[
		H^n(E;R)\twoheadrightarrow E_\infty^{0,n}\hookrightarrow H^n(F;R)^{\pi_1(B)}\hookrightarrow H^n(F;R).
	\]
	两条复合依次为 \(p^*\) 和 \(i^*\)。特别地，能从总空间限制下来的纤维类必为永久圈。
\end{proposition}
\begin{proof}
	把原纤维化映到恒等纤维化 \(B\to B\)。后者仅有零行，谱序列在第二页退化，诱导的总上同调映射为 \(p^*\)；对骨架对的拉回以及相应正合偶取映射，即得第一条分解。由于总次数 \(n\) 中 \(F^{n+1}=0\)，最右的分级商就是实际子群 \(F^n\)。

	再从基点上的纤维化 \(F\to\{b_0\}\) 映入原纤维化。零骨架上的限制把全局类送到各纤维上的相容类，沿边的相容性正是运输不变性。滤过的定义给出 \(F^1H^n(E)=\ker i^*\)：若基底有多个顶点，沿路径运输表明在一个纤维为零等价于在全部顶点纤维为零。于是 \(H^n(E)/F^1\) 就是限制像，得到第二条分解。
\end{proof}

\begin{corollary}\label{g2:cor:serre-section-bottom-row}
	若 \(p\) 有截面，则所有落入零行的微分均为零。
\end{corollary}
\begin{proof}
	截面 \(s\) 满足 \(s^*p^*=\operatorname{id}\)，所以 \(p^*\) 单射。由第一条边缘分解，零行任何类都不能在逐页取商时被消去。
\end{proof}

这个判据常比直接追踪代表元快捷。它控制的是零行；其他行之间仍可能存在微分。

\subsection{传递映射的代表元解释}
\label{g2:subsec:1102:03}

设 \(q>0\)。一个纤维类先要保持运输不变，才能进入 \(E_2^{0,q}\)；随后还要通过 \(\mathrm d_2,\ldots,\mathrm d_q\) 的检验，才能到达第 \(q+1\) 页。最后可能发出的箭头恰为
\begin{equation}\label{g2:eq:serre-transgression}
	\tau=\mathrm d_{q+1}:E_{q+1}^{0,q}\longrightarrow E_{q+1}^{q+1,0}.
\end{equation}
它称为\term[chuandi]{传递}[transgression]。源是 \(H^q(F;R)\) 的一个子群，目标是 \(H^{q+1}(B;R)\) 的一个商；较早微分可能已在零行中加入关系。符号 \(\tau a\in H^{q+1}(B;R)\) 只有在该商确实等于整个基底群时才直接使用。

\begin{proposition}[传递的延拓解释]\label{g2:prop:serre-transgression}
	类 \(a\in H^q(F;R)^{\pi_1(B)}\) 属于 \(E_{q+1}^{0,q}\)，当且仅当它可延拓为 \(E_q=p^{-1}(B^q)\) 上的 \(q\) 次上同调类。选择这样的延拓后，它在 \((q+1)\) 胞腔上继续延拓的障碍组成一个基底上链；该上链所定义的类，模去较早微分的像，就是 \(\tau a\)。
\end{proposition}
\begin{proof}
	先对限制纤维化 \(E_q\to B^q\) 使用上一命题。零列 \((0,q)\) 的所有可能非零出射在第 \(q\) 页以前已经出现；其余箭头落在基底维数以外。因此这个位置的永久圈恰为 \(E_{q+1}^{0,q}\)，而边缘命题把它识别为 \(H^q(E_q;R)\to H^q(F;R)\) 的像。

	用一个闭上链 \(A\) 表示所选延拓，再将其作为上链延拓到 \(E_{q+1}\)。余边界 \(\mathrm d\widetilde A\) 在 \(E_q\) 上为零，故定义
	\[
		[\mathrm d\widetilde A]\in H^{q+1}(E_{q+1},E_q;R)
		\cong C^{q+1}(B;\mathcal H^0(F;R))=C^{q+1}(B;R).
	\]
	这个等式描述每个新胞腔上的障碍值。再取一次余边界为零，故这些值组成胞腔上圈。更换上链延拓只改动相对边界；更换较低骨架上的修正，则加入前面几页已经允许的边界。第10章代表元公式把这些不确定性恰好除去，并把剩余的 \(\mathrm d\widetilde A\) 识别为 \(\mathrm d_{q+1}[a]\)。
\end{proof}

还可用空间对表达一种十分实用的判定。记 \(\delta:H^q(F;R)\to H^{q+1}(E,F;R)\) 为连接同态。若 \(b\in H^{q+1}(B,b_0;R)\) 满足
\begin{equation}\label{g2:eq:transgression-relative-criterion}
	\delta a=p^*b,
\end{equation}
则 \(a\) 可传递，且 \(\tau a\) 是 \(b\) 在 \(E_{q+1}^{q+1,0}\) 中的像。事实上，在相容胞腔模型中取代表并调整一个相对上边界，就能使 \(\mathrm dA=p^*b\)；右端在基底 \(q\) 骨架上为零，于是 \(A\) 先给出上述延拓，再在下一层留下 \(b\)。连接同态采用“延拓后取余边界”的约定，因此这里没有额外负号。

\ProseParagraphBreak

例如，若 \(\widetilde H^j(F;R)=0\) 对 \(0<j<q\) 成立，则从 \((0,q)\) 出发的较早微分及落入 \((q+1,0)\) 的较早微分都没有非零源或靶。此时所有运输不变的 \(q\) 次类都可传递，目标也是整个 \(H^{q+1}(B;R)\)。球丛的特征类正是由这种情形产生。

\subsection{映射纤维化之间的比较}
\label{g2:subsec:1102:04}

给定交换图
\[
	\begin{tikzcd}
		F\arrow[r]\arrow[d,"f_F"]&E\arrow[r,"p"]\arrow[d,"f_E"]&B\arrow[d,"f_B"]\\
		F'\arrow[r]&E'\arrow[r,"p'"]&B',
	\end{tikzcd}
\]
沿路提升与 \(f_E\) 相容到同伦，因此纤维上同调的拉回给出局部系统映射 \(f_B^*\mathcal H^q(F';R)\to\mathcal H^q(F;R)\)。这才是第二页映射的完整输入；当基本群作用不平凡时，只列两个抽象群之间的同态还不够。

\begin{proposition}[自然性与比较]\label{g2:prop:serre-comparison}
	上述图诱导反向的谱序列映射 \(f_r^*:E_r'{}^{p,q}\to E_r^{p,q}\)，与微分、乘法及传递交换；第二页为局部系数上同调上的诱导映射，极限页为滤过映射 \(f_E^*\) 的伴随分次。如果第二页在所有位置都是同构，则 \(f_E^*:H^*(E';R)\to H^*(E;R)\) 是同构。
\end{proposition}
\begin{proof}
	先将 \(f_B\) 胞腔逼近，并在 Hurewicz 替换中提升这段同伦以相应调整 \(f_E\)。所得映射保持骨架逆像，故给出空间对正合列之间的交换图。第10章的导出正合偶构造于是逐页给出 \(f_r^*\)，且与微分交换。逐胞腔的相对群识别使用纤维运输，所以第一微分及第二页映射恰好是所述局部系数映射。对角映射的自然性使乘积相容；传递本身是一条指定微分，故也相容。

	若 \(f_2^*\) 为同构，取各页上同调归纳得到所有 \(f_r^*\) 为同构。固定总次数 \(n\)，两边的目标滤过只有 \(n+1\) 层。从最深一层开始，将相邻滤过项及其商放入短正合列，逐层用五引理，便得到 \(f_E^*\) 在 \(n\) 次的同构。这里使用的是前节已经证明的有限目标滤过。
\end{proof}

对同调版也有同样的比较，箭头方向与连续映射一致。特别地，若基底均单连通，\(f_B\) 与 \(f_F\) 都诱导整数同调同构，则 \(f_E\) 也诱导整数同调同构：第4章的同调万有系数短正合列先将基底同构提升为任意纤维同调系数下的同构，随后逐页比较即可。这个步骤解释了为什么“基底与纤维分别不变”能够推出总空间的同调不变。
% 纲要条目排为编号小节标题，并进入目录与书签；说明段仍保留于纲要注释中。

% !TeX root = ../../Book2.tex
% 纲要标题：### 11.3 球丛和路径纤维化
% 纲要位置：几何学计划纲要.md，第 1447 行。
% 正文状态：未撰写；本文件目前是框架。

% 纲要内容开始
% 
% * Hopf 丛中的非零微分
% * 球丛与 Euler 类的联系
% * 路径—回路纤维化
% * 悬挂和回路空间的所选计算
% 
% 纲要内容结束

\section{球丛和路径纤维化}
\label{g2:sec:1103}

待撰写

\subsection{Hopf 丛中的非零微分}
\label{g2:subsec:1103:01}
待撰写

\subsection{球丛与 Euler 类的联系}
\label{g2:subsec:1103:02}
待撰写

\subsection{路径—回路纤维化}
\label{g2:subsec:1103:03}
待撰写

\subsection{悬挂和回路空间的所选计算}
\label{g2:subsec:1103:04}
待撰写
% 纲要条目排为编号小节标题，并进入目录与书签；说明段仍保留于纲要注释中。

% !TeX root = ../../Book2.tex
% 纲要标题：### 11.4 经典群与分类空间
% 纲要位置：几何学计划纲要.md，第 1454 行。
% 正文状态：初稿；尚未正式审核。

% 纲要内容开始
% 
% * 酉群的整数上同调环
% * 通用圆周丛与 Euler 类
% * 维数、环结构及扭项的核验
% 
% 纲要内容结束

\section{经典群与分类空间}
\label{g2:sec:1104}

第八章把一个酉矩阵的第一列取出来，得到以低一阶酉群为纤维的球面丛。这张映射适合归纳：底球面只有两个非零上同调群，纤维的环结构又来自上一步。谱序列的乘法将使整页微分由少量生成元决定。随后考察总空间可缩的圆周丛，便能看到同一个纤维生成元怎样在另一种纤维化中传递成基底上的特征类。

\subsection{酉群的整数上同调环}
\label{g2:subsec:1104:01}

由\cref{g2:prop:classical-sphere-bundles}，取第一列给出局部平凡丛
\[
	\operatorname{U}(n-1)\xrightarrow{i_n}\operatorname{U}(n)\xrightarrow{p_n}S^{2n-1},
	\qquad i_n(A)=\operatorname{diag}(1,A).
\]
当 \(n\ge2\) 时，底球面单连通，所以纤维上同调的局部系统为常值。这里完全使用整数系数；保留整数才能同时判断生成元的倍数以及挠项是否出现。

\begin{theorem}[酉群的整数上同调环]\label{g2:thm:unitary-cohomology}
	对每个 \(n\ge1\)，可选齐次类 \(x_{2j-1}\in H^{2j-1}(\operatorname{U}(n);\mathbb Z)\)（\(1\le j\le n\)），使
	\[
		H^*(\operatorname{U}(n);\mathbb Z)\cong\Lambda_{\mathbb Z}(x_1,x_3,\ldots,x_{2n-1}).
	\]
	这些选择可以逐个 \(n\) 相容，使 \(i_n^*x_{2j-1}\) 为低一阶酉群中的同名类（\(j<n\)），而 \(x_{2n-1}=p_n^*\omega_{2n-1}\)，其中 \(\omega_{2n-1}\) 是球面的一个整数上同调生成元。这里的外代数具有关系 \(x_ix_j=-x_jx_i\)、\(x_i^2=0\)；每个生成元至多出现一次的递增单项式构成其整数基。
\end{theorem}
\begin{proof}
	初始情形为 \(\operatorname{U}(1)=S^1\)，其环是 \(\Lambda_{\mathbb Z}(x_1)\)。设已知 \(\operatorname{U}(n-1)\) 的结论。记 \(m=2n-1\)，并将纤维中的生成元暂记为 \(\bar x_1,\ldots,\bar x_{m-2}\)。基底球面的同调自由，故第二页及其乘法为
	\[
		E_2^{*,*}\cong H^*(S^m;\mathbb Z)\otimes\Lambda_{\mathbb Z}(\bar x_1,\ldots,\bar x_{m-2}).
	\]
	只有第 \(0,m\) 两列可能非零。由 \(\mathrm d_r\) 的双次数，唯一可能非零的微分为 \(\mathrm d_m\)，从第零列指向第 \(m\) 列。对每个旧生成元，
	\[
		\mathrm d_m:E_m^{0,2j-1}\longrightarrow E_m^{m,2j-m},\qquad j<n,
	\]
	目标的第二指标至多为 \(-1\)，所以 \(\mathrm d_m\bar x_{2j-1}=0\)。底球面的类也没有非零出微分。\cref{g2:thm:serre-multiplicative}的 Leibniz 规则于是说明，所有生成元的乘积都具有零微分。这些乘积生成整页，故 \(E_2=E_\infty\)。注意，纤维本身可能在很高次数仍有非零上同调；真正使微分全部消失的是低次生成元与导数性共同作用。

	每个 \(E_\infty^{p,q}\) 都是有限秩自由阿贝尔群。\cref{g2:prop:free-graded-splitting}说明 \(H^k(\operatorname{U}(n);\mathbb Z)\) 也自由：其有限滤过的每个商都有可提升的基。边缘映射的描述则给出满射
	\[
		i_n^*:H^*(\operatorname{U}(n);\mathbb Z)\longrightarrow H^*(\operatorname{U}(n-1);\mathbb Z).
	\]
	据此选择 \(x_{2j-1}\) 提升 \(\bar x_{2j-1}\)。在另一条边上，\(p_n^*\omega_m\) 的伴随分次类就是 \(E_\infty^{m,0}\) 的生成元，令它为 \(x_m\)。由于所有 \(x_{2j-1}\) 都是奇次类，杯积的分次交换律给出 \(2x_{2j-1}^2=0\)；总上同调无挠，故每个平方为零。这就构造出从所述外代数到 \(H^*(\operatorname{U}(n);\mathbb Z)\) 的环同态。

	最后核对这些乘积确实构成基，以解决乘法延拓。在固定总次数 \(k\) 上，滤过只有第零列和第 \(m\) 列两个可能非零的商。不含 \(x_m\) 的递增单项式限制到纤维，正好给出第零列的一组基；含 \(x_m\) 的单项式可写成 \(x_m x_I\)，其在第 \(m\) 列的类为 \(\omega_m\bar x_I\)，也正好给出该列的一组基。任意上同调类先减去第一组基的线性组合，就落在第 \(m\) 层，再由第二组表示；任何线性关系先限制到纤维，第一组的系数都为零，再投到第 \(m\) 层，第二组的系数也都为零。因此这张环同态为同构。沿归纳过程选择提升，即得所述相容性。
\end{proof}

这个证明给出的生成元有两种来源：旧生成元从纤维提升，新生成元从底球面拉回。例如 \(\operatorname{U}(2)\) 的基为 \(1,x_1,x_3,x_1x_3\)，分别处于次数 \(0,1,3,4\)。其中 \(x_3=p_2^*\omega_3\)，而 \(x_1\) 可取行列式映射 \(\det:\operatorname{U}(2)\rightarrow S^1\) 拉回的一次生成元；行列式在标准 \(\operatorname{U}(1)\) 纤维上恰为恒等映射。

\begin{corollary}[特殊酉群的整数上同调环]\label{g2:cor:special-unitary-cohomology}
	对 \(n\ge2\)，有
	\[
		H^*(\operatorname{SU}(n);\mathbb Z)\cong\Lambda_{\mathbb Z}(y_3,y_5,\ldots,y_{2n-1}),\qquad |y_j|=j.
	\]
\end{corollary}
\begin{proof}
	\cref{g2:prop:classical-sphere-bundles}给出 \(\operatorname{SU}(2)\cong S^3\) 以及 \(\operatorname{SU}(n-1)\rightarrow\operatorname{SU}(n)\rightarrow S^{2n-1}\)。从 \(n=3\) 开始归纳。纤维的生成元处于次数 \(3,5,\ldots,2n-3\)，所以唯一可能的 \(\mathrm d_{2n-1}\) 在每个生成元上都有负的目标第二指标，因而全为零；导数性使全页退化。全部分次块自由，纤维生成元可以提升，新生成元由球面拉回。无挠性使奇次平方为零，含新生成元及不含新生成元的单项式分别给出两层的基。由上一证明中的逐层基论证，得到所述外代数。
\end{proof}

\subsection{通用圆周丛与 Euler 类}
\label{g2:subsec:1104:02}

在上面的酉群计算中，旧生成元全部存活。现在把圆周放进总空间可缩的纤维化，其一次类就必须被消去。消去它的微分落在基底二次上同调中；自然性将把这个类送到每个具体圆周丛的基底上。

\ProseParagraphBreak

取有限支撑复序列构成的单位球面
\[
	S^\infty=\bigcup_{N\ge1}S^{2N-1}\subset\mathbb C^\infty,
\]
赋予与有限维球面相容的 CW 弱拓扑。这里可从实球面的逐维赤道与半球胞腔结构取并，每个有限子复形都包含于足够大的有限维球面。标量右乘给出自由 \(S^1\) 作用，其商为 \(\mathbb CP^\infty\)。在 \(z_j\ne0\) 的坐标片上，仍可选择第 \(j\) 坐标为正实数的唯一单位代表；有限阶段的这些公式相容，给出局部平凡丛
\begin{equation}\label{g2:eq:universal-circle-bundle}
	S^1\longrightarrow S^\infty\xrightarrow{\pi}\mathbb CP^\infty.
\end{equation}
其中 \(S^\infty\) 可缩。事实上，由\cref{g2:thm:finite-subcomplex}，任一球面映射的像包含在有限阶段；再放进维数大于定义域维数的球面，由\cref{g2:prop:sphere-connectivity}可知它零伦。因此 \(S^\infty\) 道路连通且全部正阶同伦群为零，\cref{g2:thm:whitehead}用于点到这个 CW 复形的映射，便得到可缩性。

\ProseParagraphBreak

基底上的\term[tautological-complex-line-bundle]{典范复线丛}[tautological complex line bundle]为
\[
	\gamma^1=\{(\ell,v):\ell\in\mathbb CP^\infty,\ v\in\ell\}\longrightarrow\mathbb CP^\infty.
\]
它的单位圆周丛就是\eqref{g2:eq:universal-circle-bundle}。复线的定向使其底层实二维丛 \(\gamma^1_{\mathbb R}\) 定向；按\cref{g2:def:spherical-euler-class}，记
\[
	u\coloneqq e(\gamma^1_{\mathbb R})=\mathrm d_2(a),\qquad a\in H^1(S^1;\mathbb Z)
\]
为正向纤维生成元的传递。这个 \(u\) 是 \(H^2(\mathbb CP^\infty;\mathbb Z)\cong\mathbb Z\) 的生成元：\(E_2^{0,1}\) 与 \(E_2^{2,0}\) 之间只有 \(\mathrm d_2\)，而总空间的一次、二次上同调都为零，迫使该微分既单射又满射。它与\cref{g2:thm:complex-projective-ring}选用的二次生成元至多相差符号。

\ProseParagraphBreak

由该定理及乘法结构，第二页和微分完整写成
\[
	E_2^{*,*}=\mathbb Z[u]\otimes\Lambda_{\mathbb Z}(a),\qquad |u|=(2,0),\quad |a|=(0,1),
\]
\[
	\mathrm d_2(u^k)=0,\qquad \mathrm d_2(u^ka)=u^{k+1}\quad(k\ge0).
\]
每个正次数的底类都被上面一行的类消去，只留下单位元，正好对应总空间可缩。在有限维 Hopf 丛中，底环被 \(u^{N}=0\) 截断，因而 \(u^{N-1}a\) 留下总空间 \(S^{2N-1}\) 的顶次类；通用丛把这些有限截断相容地接在一起。

\begin{proposition}[复线丛的拉回模型]\label{g2:prop:line-bundle-pullback-model}
	设 \(B\) 为有限 CW 复形，\(L\rightarrow B\) 为连续复线丛。则存在某个 \(N\) 及连续映射 \(f:B\rightarrow\mathbb CP^{N-1}\subset\mathbb CP^\infty\)，使 \(L\cong f^*\gamma^1\)。把 \(L\) 看成定向实二维丛后，其 Euler 类满足
	\[
		e(L_{\mathbb R})=f^*u.
	\]
\end{proposition}
\begin{proof}
	有限 CW 复形紧且 Hausdorff，故可取有限平凡化覆盖 \(U_1,\ldots,U_N\) 及从属的连续单位分解 \(\rho_1,\ldots,\rho_N\)。这里所需的事实是 \(\rho_j\ge0\)、\(\operatorname{supp}\rho_j\subset U_j\) 及 \(\sum_j\rho_j=1\)；存在性见第一卷第4.4节的连续单位分解定理。局部 Hermitian 内积也可用这些权重相加，得到 \(L\) 上的连续正定内积，因为非零向量的各项平方长度都为正。将局部标架归一化，得到单位截面 \(s_j\)。

	对 \(b\in U_j\) 及单位向量 \(v\in L_b\)，唯一写成 \(v=s_j(b)\lambda_j(v)\)，其中 \(\lambda_j(v)\in S^1\)。令
	\[
		J(v)=\bigl(\sqrt{\rho_1(b)}\lambda_1(v),\ldots,\sqrt{\rho_N(b)}\lambda_N(v)\bigr)\in S^{2N-1},
	\]
	并在 \(b\notin U_j\) 时把相应坐标规定为零。支集条件使这些补零后的坐标连续，而 \(\sum\rho_j=1\) 保证所得向量确实为单位向量。对 \(z\in S^1\)，有 \(J(vz)=J(v)z\)，故 \(f(b)=[J(v)]\) 与同一纤维中的单位向量选择无关，并由局部截面看出 \(f\) 连续。

	将 \(J\) 在每条复线上复线性延拓，得到 \(L_b\) 到 \(\mathbb C^N\) 的等距单射，其像为 \(f(b)\) 所代表的复直线。因此
	\[
		L\longrightarrow f^*\gamma^1,\qquad v\longmapsto(b,J(v))
	\]
	是纤维上的复线性同构；在局部非零标架上，其逆由一个连续非零系数的倒数给出，故为丛同构。限制到单位向量便得到圆周丛的拉回交换图
	\[
		\begin{tikzcd}[column sep=large]
			S(L)\arrow[r,"J"]\arrow[d] & S^\infty\arrow[d,"\pi"]\\
			B\arrow[r,"f"'] & \mathbb CP^\infty.
		\end{tikzcd}
	\]
	纤维映射保持复定向。由\cref{g2:prop:spherical-euler-naturality}，其 Euler 类为 \(f^*u\)。
\end{proof}

这个构造解释了“通用”的含义：一个具体线丛的局部标架经加权后，确定一张到射影空间的映射，再由同一个典范线丛拉回。赋予每个丛一个随拉回自然变化的上同调类，正是\term[tezheng-lei]{特征类}[characteristic class]的基本形式。这里得到的 \(e(L_{\mathbb R})\) 已是这样一个类；它由圆周的一个整数生成元及传递构造而来。不同单位分解所产生的映射都给出同一个 Euler 类，因为它们拉回得到的是同一个丛。

\ProseParagraphBreak

在一般秩中，射影空间将由 Grassmann 空间替代，单位向量将由酉标架替代。\chapref{g2:ch:19}将建立相应分类空间与通用丛，\chapref{g2:ch:21}再系统构造复向量丛的特征类。眼前这张圆周丛已显示出关键计算方式：先在通用模型中求纤维类的传递，再利用拉回自然性把答案送到给定基底。

\subsection{维数、环结构及扭项的核验}
\label{g2:subsec:1104:03}

得到一个环的候选答案后，可以分别检验其次数、乘法及整数信息。对 \(\operatorname{U}(n)\)，上一小节的论证已经解决了后两者；把全部基元的次数汇总，还会产生几项直观的核验。

\begin{corollary}[酉群的 Betti 数与顶次类]\label{g2:cor:unitary-betti-polynomial}
	记 \(b_k\) 为整数上同调的秩。对 \(n\ge1\)，有
	\[
		\sum_k b_k(\operatorname{U}(n))t^k=\prod_{j=1}^n(1+t^{2j-1}),
		\qquad \sum_k b_k(\operatorname{U}(n))=2^n,
	\]
	且最高非零次数为 \(n^2\)，由 \(x_1x_3\cdots x_{2n-1}\) 生成，Euler 示性数为零。对 \(n\ge2\)，特殊酉群的相应多项式为 \(\prod_{j=2}^n(1+t^{2j-1})\)，最高非零次数为 \(n^2-1\)，Euler 示性数也为零。
\end{corollary}
\begin{proof}
	外代数的一组基由每个生成元“选取或不选取”得到，所以次数生成函数为所列乘积；令 \(t=1\) 得到总秩。最高次数是
	\[
		1+3+\cdots+(2n-1)=n^2,
	\]
	且只有选取全部生成元的一个单项式达到该次数。令 \(t=-1\)，每个奇次因子都为零，得到 Euler 示性数。特殊酉群删去一次生成元，故顶次降为 \(n^2-1\)，其余计算相同。
\end{proof}

顶次与矩阵的实参数数目一致。单位酉矩阵附近的无穷小变动由斜 Hermitian 矩阵给出：对角线有 \(n\) 个纯虚参数，每对非对角位置由一个复数决定，共有 \(n(n-1)\) 个实参数，合为 \(n^2\)。行列式为一增加迹为零的一个实条件，得到 \(n^2-1\)。这与上同调中唯一顶次单项式的位置吻合。

\ProseParagraphBreak

例如
\[
	\sum_kb_k(\operatorname{U}(3))t^k=(1+t)(1+t^3)(1+t^5)
	=1+t+t^3+t^4+t^5+t^6+t^8+t^9.
\]
次数为 \(4\) 的类是 \(x_1x_3\)，次数为 \(6\) 的类是 \(x_1x_5\)，次数为 \(8\) 的类是 \(x_3x_5\)。把这些乘积列出来，有助于避免将“只有奇次生成元”误读为“只有奇次上同调”。

\ProseParagraphBreak

对任一素数 \(\ell\)，系数短正合列 \(0\rightarrow\mathbb Z\xrightarrow{\times\ell}\mathbb Z\rightarrow\mathbb F_\ell\rightarrow0\) 给出
\[
	H^k(\operatorname{U}(n);\mathbb Z)\xrightarrow{\times\ell}H^k(\operatorname{U}(n);\mathbb Z)
	\longrightarrow H^k(\operatorname{U}(n);\mathbb F_\ell)
	\longrightarrow H^{k+1}(\operatorname{U}(n);\mathbb Z)\xrightarrow{\times\ell}H^{k+1}(\operatorname{U}(n);\mathbb Z).
\]
末箭头单射，故约化模 \(\ell\) 给出整数上同调张量 \(\mathbb F_\ell\) 到模 \(\ell\) 上同调的同构。杯积与系数映射相容，因而
\[
	H^*(\operatorname{U}(n);\mathbb F_\ell)\cong\Lambda_{\mathbb F_\ell}(\bar x_1,\bar x_3,\ldots,\bar x_{2n-1}).
\]
当 \(\ell=2\) 时，奇次平方为零的理由是这些类来自整数类的约化；单凭特征二下的分次交换律不能得到这个结论。

\begin{example}[旋转群中的整数挠类]\label{g2:exm:rotation-torsion-comparison}
	\cref{g2:lem:quaternion-rotation-cover}把 \(\operatorname{SO}(3)\) 识别为 \(S^3/\{\pm1\}=\mathbb RP^3\)。\cref{g2:prop:projective-integer-homology}给出
	\[
		H_0\cong H_3\cong\mathbb Z,\qquad H_1\cong\mathbb Z/2,\qquad H_2=0.
	\]
	整数万有系数定理据此得到
	\[
		H^k(\operatorname{SO}(3);\mathbb Z)\cong
		\begin{cases}
			\mathbb Z,&k=0,3,\\
			\mathbb Z/2,&k=2,\\
			0,&\text{其他次数}.
		\end{cases}
	\]
	其中二次的 \(\mathbb Z/2\) 是 \(\operatorname{Ext}_{\mathbb Z}^1(H_1,\mathbb Z)\)。有理系数却只剩 \(H^*(\operatorname{SO}(3);\mathbb Q)\cong\Lambda_{\mathbb Q}(z_3)\)。所以一个有理外代数答案仍可能遗漏整数挠项；酉群的整数外代数结论，依靠的是整数谱序列中自由分次块与生成元提升的完整论证。
\end{example}

% 纲要条目排为编号小节标题，并进入目录与书签；说明段仍保留于纲要注释中。

% !TeX root = ../../Book2.tex
% 纲要标题：### 11.5 有限生成与进一步方法
% 纲要位置：几何学计划纲要.md，第 1460 行。
% 正文状态：初稿；尚未正式审核。

% 纲要内容开始
% 
% * Serre 类与有限生成群
% * 同调与同伦群的有限生成性
% * Adams 与 Eilenberg–Moore 谱序列
% 
% ---
% 
% 纲要内容结束

\section{有限生成与进一步方法}
\label{g2:sec:1105}

谱序列也能回答比“这个群到底是什么”更粗、却更普遍的问题。例如，有限 CW 复形的同调由有限矩阵计算，但它的高阶同伦群并没有同样直接的矩阵表示。能否至少证明每个同伦群只需有限多个生成元？这一节用谱序列把有限生成性在纤维、基底与总空间之间传递，再用第9章的胞腔附着将同伦群逐层分离出来。以下未注明的同调均取整数系数。

\subsection{Serre 类与有限生成群}
\label{g2:subsec:1105:01}

一个页面上的群经取核、取像和取商，成为下一页的群；目标群又是有限个极限页项的逐次扩张。因此，适合被谱序列追踪的性质，应当在这几种操作下保持。

\begin{definition}[Serre 类]\label{g2:def:serre-class}
	阿贝尔群的一类 \(\mathcal C\)，若包含零群，且对任意短正合列 \(0\to A\to B\to C\to0\) 都有
	\[
		B\in\mathcal C\quad\Longleftrightarrow\quad A,C\in\mathcal C,
	\]
	就称为一个\term[Serre lei]{Serre 类}[Serre class]。这等价于对取子群、取商及有限次扩张封闭。若一个同态的核与余核都属于 \(\mathcal C\)，便称它为模 \(\mathcal C\) 的同构。
\end{definition}

有限生成阿贝尔群组成一个 Serre 类，记为 \(\mathcal C_{\mathrm{fg}}\)。其子群有限生成可由整数自由模的子模定理及有限生成群分解推出，商群显然由生成元的像生成。对扩张，提升商群的一组有限生成元，再加入子群的一组有限生成元，就得到中间群的有限生成集。有限阿贝尔群也组成一个 Serre 类：短正合列中中间群的阶是两端群阶的乘积。

\ProseParagraphBreak

Serre 类的定义本身没有规定张量积与 Tor 的行为。使用第二页时，还须核对这些系数运算。对 \(\mathcal C_{\mathrm{fg}}\)，将两群写成有限个无限循环群与有限循环群的直和，再用\cref{g2:prop:cyclic-ext-tor}，可知它们的张量积和 \(\operatorname{Tor}_1^{\mathbb Z}\) 仍有限生成。于是当 \(H_i(B)\) 与 \(A\) 都逐次有限生成时，同调万有系数短正合列
\begin{equation}\label{g2:eq:finite-generation-coefficients}
	0\longrightarrow H_p(B)\otimes A\longrightarrow H_p(B;A)
	\longrightarrow\operatorname{Tor}_1^{\mathbb Z}(H_{p-1}(B),A)\longrightarrow0
\end{equation}
说明系数同调也有限生成。下面用这一项具体性质，而不对任意 Serre 类作未经核对的推广。

\subsection{同调与同伦群的有限生成性}
\label{g2:subsec:1105:02}

称一个道路连通空间具有\term[tongdiao-youxianxing]{逐次有限生成同调}[degreewise finitely generated homology]，如果每个 \(H_n(-;\mathbb Z)\) 都有限生成。各次数可以有不同数量的生成元，也允许无限多个次数非零。譬如 \(\mathbb CP^\infty\) 满足此条件，尽管它有无限多个胞腔。

\begin{lemma}[纤维化中的有限生成性]\label{g2:lem:fibration-finite-generation}
	设 \(F\to E\to B\) 是道路连通空间的纤维化，基底为单连通 CW 复形。若 \(F,E,B\) 中任意两个具有逐次有限生成同调，则第三个也有。
\end{lemma}
\begin{proof}
	用同调谱序列
	\[
		E^2_{p,q}=H_p(B;H_q(F))\Longrightarrow H_{p+q}(E),\qquad
		\mathrm d_r:E^r_{p,q}\to E^r_{p-r,q+r-1}.
	\]
	单连通性消去了运输作用，且 \(H_0(F)=H_0(B)=\mathbb Z\)。每个固定位置经过有限多页便稳定，每个目标次数也只有有限层。

	先设 \(B,F\) 的同调逐次有限生成。式\eqref{g2:eq:finite-generation-coefficients}使全部第二页项有限生成。取核、取商保留此性质，故全部极限页项也有限生成；目标群由有限个这样的群扩张而成，结论随之成立。

	再设 \(F,E\) 满足条件，归纳证明 \(H_k(B)\) 有限生成。假定次数小于 \(k\) 的基底群已有此性质，则第二页中所有 \(p<k\) 的项都有限生成。同调零行没有来箭头，所以对 \(2\le r\le k\) 有
	\[
		0\longrightarrow E^{r+1}_{k,0}\longrightarrow E^r_{k,0}
		\xrightarrow{\mathrm d_r}\operatorname{im}\mathrm d_r\longrightarrow0,
	\]
	其中像是 \(E^r_{k-r,r-1}\) 的子群，故有限生成。稳定项 \(E^\infty_{k,0}\) 是 \(H_k(E)\) 的一个分级商，也有限生成。倒着经过有限多页，短正合列遂说明 \(E^2_{k,0}=H_k(B)\) 有限生成。\(k=1\) 时没有这些微分，结论直接来自稳定项；这也给出归纳起点。

	最后设 \(B,E\) 满足条件，归纳证明 \(H_k(F)\) 有限生成。零列没有去箭头，可能进入 \((0,k)\) 的微分只有
	\[
		\mathrm d_r:E^r_{r,k-r+1}\longrightarrow E^r_{0,k},\qquad 2\le r\le k+1.
	\]
	所有源项的纤维次数都小于 \(k\)，故按归纳假设及系数公式有限生成。各页有短正合列
	\[
		0\longrightarrow\operatorname{im}\mathrm d_r\longrightarrow E^r_{0,k}
		\longrightarrow E^{r+1}_{0,k}\longrightarrow0.
	\]
	从 \(H_k(E)\) 的有限生成分级商 \(E^\infty_{0,k}\) 倒推，便得 \(E^2_{0,k}=H_k(F)\) 有限生成。零次群为 \(\mathbb Z\)，提供了这里的归纳起点。
\end{proof}

\begin{corollary}\label{g2:cor:loop-finite-generation}
	若单连通 CW 复形 \(B\) 具有逐次有限生成同调，则 \(\Omega B\) 也有。
\end{corollary}
\begin{proof}
	将引理的第三种情形用于 \(\Omega B\to P_{b_0}B\to B\)。道路空间可缩，回路空间因基底单连通而道路连通。
\end{proof}

要把结论送到同伦群，需要先处理只含一个同伦群的空间。下述证明把所需模型与有限性理由一并写出。

\begin{lemma}[Eilenberg–MacLane 空间的有限生成性]\label{g2:lem:em-space-finite-generation}
	若 \(A\) 是有限生成阿贝尔群，则对每个 \(n\ge1\)，空间 \(K(A,n)\) 具有逐次有限生成整数同调。
\end{lemma}
\begin{proof}
	先说明这里选择哪个 CW 模型不影响结论。对 \(n\ge2\)，用第9章的稀疏 CW 模型把 \(K(A,n)\) 换成除一个顶点外没有低于 \(n\) 维胞腔的模型。其 \(n\) 骨架是球面楔和，且 \(H_n=C_n/\operatorname{im}\partial_{n+1}\)。给定另一个 \(K(A,n)\)，把每个 \(n\) 胞腔的球面送到它所代表的 \(A\) 中元素；Hurewicz 同构保证这样的代表存在。每个 \((n+1)\) 胞腔的边界在 \(A\) 中为零，故可延拓；更高胞腔的延拓障碍落在目标已消失的高阶同伦群中，也全为零。所得映射在唯一非零同伦群上为同构，Whitehead 定理给出同伦等价。对 \(n=1\)，先在一骨架实现基本群的指定同态，二胞腔的关系保证延拓，随后用目标高阶同伦群全零延拓。因而同样得到模型唯一性。非 CW 纤维可先取 CW 逼近；\cref{g2:lem:weak-equivalence-homology}保证所用同调不变。

	现处理 \(n=1\)。将 \(A\) 写成 \(\mathbb Z^r\times G\)，其中 \(G\) 是有限阿贝尔群。环面 \(T^r\) 是 \(K(\mathbb Z^r,1)\)，已有有限 CW 模型。对有限群 \(G\)，构造如下空间：在次数 \(k\) 为每个序列 \((g_0,\ldots,g_k)\) 放置一个 \(k\) 单形，按删去一个顶点的面映射粘合，并把相邻重复顶点的退化单形压到相应低维单形。所得实现记为 \(EG\)。它是 CW 复形，\(G\) 通过同时左乘全部顶点自由作用；商 \(BG\) 每维只有有限多个胞腔，因为每维的序列总数有限。

	\(EG\) 可缩有一个直接的公式。对由 \((g_0,\ldots,g_k)\) 与重心坐标 \((t_0,\ldots,t_k)\) 表示的点，送到由
	\[
		(1,g_0,\ldots,g_k),\qquad (s,(1-s)t_0,\ldots,(1-s)t_k)
	\]
	表示的点，\(0\le s\le1\)。它与面、退化的识别相容；\(s=0\) 是原点，\(s=1\) 是顶点 \(1\)。有限群的自由作用在 Hausdorff 空间上给出覆盖：对一点及其有限多个不同的平移，先取互不相交的小邻域，再有限次相交得到与所有非平凡平移不交的邻域。故 \(EG\to BG\) 是具有可缩总空间的通用覆盖，\(BG\) 是 \(K(G,1)\)。其每维胞腔链为有限秩自由群，所以每个同调群有限生成。乘积 \(T^r\times BG\) 的同伦群逐坐标计算，说明它是 \(K(A,1)\)；Künneth 定理及有限生成群的张量、Tor 性质给出其同调有限生成。

	对 \(n\ge2\) 作归纳。路径纤维化
	\[
		\Omega K(A,n)\longrightarrow P K(A,n)\longrightarrow K(A,n)
	\]
	的纤维由同伦长正合列具有 \(K(A,n-1)\) 的同伦群，取 CW 逼近并用刚证的模型唯一性，便知其同调逐次有限生成。总空间可缩，基底单连通。\cref{g2:lem:fibration-finite-generation}中“纤维与总空间推出基底”的情形完成归纳。
\end{proof}

接着把一个空间的同伦群分成有限层。这里先构造有限生成证明所用的部分；第12章将进一步研究层与层之间的附着信息。

\begin{lemma}[有限同伦截段]\label{g2:lem:finite-homotopy-truncations}
	设 \(X\) 为单连通带基点 CW 复形。对每个 \(n\ge1\)，存在包含 \(X\hookrightarrow X[n]\)，使
	\[
		\pi_i(X[n])\cong\pi_i(X)\ (i\le n),\qquad \pi_i(X[n])=0\ (i>n),
	\]
	且包含在 \(H_i\) 上对 \(i\le n\) 为同构。可以选择延拓恒等映射的映射 \(X[n]\to X[n-1]\)（\(n\ge2\)），其同伦纤维为 \(K(\pi_n(X),n)\) 的弱同伦型；\(X[1]\) 可缩。
\end{lemma}
\begin{proof}
	从 \(X\) 出发，依次用\cref{g2:lem:kill-homotopy-group}消去 \(\pi_{n+1},\pi_{n+2},\ldots\)。所加胞腔维数至少为 \(n+2\)，因此所有 \(i\le n\) 次同伦群保持不变，相对胞腔链在次数小于 \(n+2\) 时为零，故相应同调也保持。对无限并，第9章的紧像有限子复形论证保证固定次数的同伦群达到稳定值。\(n=1\) 时，初始基本群已为零，所得 \(X[1]\) 的全部同伦群均为零，Whitehead 定理使它可缩。

	在 \(X[n]\) 的子复形 \(X\) 上取包含到 \(X[n-1]\) 的映射。一个新增 \(d\) 胞腔满足 \(d\ge n+2\)，延拓所需消去的边界类属于 \(\pi_{d-1}(X[n-1])=0\)，所以可逐胞腔延拓到整个 \(X[n]\)。将所得映射作第8章的路径纤维化替换，记纤维为 \(F_n\)。在次数 \(i<n\) 上，映射诱导由 \(X\) 给出的同构；基底在次数 \(i\ge n\) 的群均为零。长正合列因而给出
	\[
		\pi_i(F_n)=0\ (i\ne n),\qquad \pi_n(F_n)\cong\pi_n(X).
	\]
	两空间单连通还保证纤维道路连通。取纤维的 CW 逼近，即得到指定的 Eilenberg–MacLane 同伦型。
\end{proof}

\begin{theorem}[Serre 有限生成定理]\label{g2:thm:serre-finite-generation}
	对单连通 CW 复形 \(X\)，下列两个条件等价：
	\begin{enumerate}
		\item 每个 \(H_n(X;\mathbb Z)\) 都有限生成；
		\item 每个 \(\pi_n(X)\)（\(n\ge2\)）都有限生成。
	\end{enumerate}
	特别地，单连通有限 CW 复形的全部同伦群都有限生成。
\end{theorem}
\begin{proof}
	先假定全部同伦群有限生成。\(X[1]\) 可缩，从它开始沿
	\[
		K(\pi_j(X),j)\longrightarrow X[j]\longrightarrow X[j-1]
	\]
	逐层归纳。这里的箭头按上一引理作纤维化替换，仍用原名表示同伦等价的总空间。Eilenberg–MacLane 引理使每一纤维的同调逐次有限生成，纤维化引理再把此性质传到 \(X[j]\)。固定 \(n\)，取 \(j\ge n\)，包含给出 \(H_n(X)\cong H_n(X[j])\)，所以 \(H_n(X)\) 有限生成。相同论证还说明：若只知道 \(\pi_i(X)\) 在 \(i\le j\) 有限生成，则 \(X[j]\) 的全部同调已经逐次有限生成。

	反过来，假定全部 \(H_i(X)\) 有限生成，归纳证明 \(\pi_n(X)\) 有限生成。起点 \(n=2\) 是 Hurewicz 同构 \(\pi_2(X)\cong H_2(X)\)。设低于 \(n\) 的同伦群已知有限生成。上一段的有限截段结论保证基底 \(Y=X[n-1]\) 的全部同调有限生成。令 \(Z=X[n]\)，考虑纤维化 \(F_n\to Z\to Y\)。纤维 \((n-1)\) 连通，故
	\[
		H_q(F_n)=0\ (0<q<n),\qquad H_n(F_n)\cong\pi_n(F_n)\cong\pi_n(X).
	\]
	在同调谱序列中，这使前 \(n-1\) 个正行为空。总次数 \(n\) 只可能有 \((n,0)\) 和 \((0,n)\) 两项；唯一可能击中后一项的箭头是
	\[
		\mathrm d_{n+1}:H_{n+1}(Y)=E^{n+1}_{n+1,0}\longrightarrow E^{n+1}_{0,n}=H_n(F_n).
	\]
	源端在前几页未改变，因为较早的去箭头都落入空行，零行又无来箭头。目标项的稳定余核是目标滤过的最浅子群，因而有正合列
	\begin{equation}\label{g2:eq:finite-generation-five-terms}
		H_{n+1}(Y)\xrightarrow{\mathrm d_{n+1}}H_n(F_n)
		\longrightarrow H_n(Z)\longrightarrow H_n(Y)\longrightarrow0.
	\end{equation}
	第一映射的像有限生成，因为 \(H_{n+1}(Y)\) 有限生成；它的余核是 \(H_n(Z)\cong H_n(X)\) 的子群，也有限生成。因此中间群 \(H_n(F_n)\cong\pi_n(X)\) 是两个有限生成群的扩张，仍有限生成，完成归纳。

	若 \(X\) 为有限 CW 复形，其胞腔链只有有限多个有限秩自由群，故所有同调群有限生成，最后的断言随即成立。
\end{proof}

证明的关键出现在式\eqref{g2:eq:finite-generation-five-terms}：未知的 \(\pi_n(X)\) 被放进一个纤维首个同调群，左边控制它被消去的部分，右边控制剩下的部分。两部分都有限生成，就足以控制整个未知群。这里无需算出传递的具体整数矩阵。

\ProseParagraphBreak

例如，对每个 \(m\ge2\) 和 \(n\ge2\)，\(\pi_n(S^m)\) 都是有限生成阿贝尔群。结论没有给出其阶，也没有断言它是有限群：\(\pi_m(S^m)=\mathbb Z\)，且前面已算出 \(\pi_3(S^2)=\mathbb Z\)。逐次有限生成与“所有高阶群都是有限群”是不同的断言。单连通假设也确实承担作用，章末将用 \(S^1\vee S^2\) 展示失去这一假设后的现象。

\subsection{Adams 与 Eilenberg–Moore 谱序列}
\label{g2:subsec:1105:03}

到这里，Serre 方法的输入与输出已经清楚：由基底和纤维的同调、运输作用及几何关系，求总空间的同调。路径纤维化使同一方法也能反求回路空间。若目标变成稳定同伦群，或已知基底与总空间而想系统地反求纤维的环结构，就需要不同的构造。

\ProseParagraphBreak

Adams\footnote{John Frank Adams（1930—1989），英国数学家，发展了计算稳定同伦群的谱序列方法，并在球面同伦与上同调运算方面作出重要贡献。}方法把模素数上同调与作用在它上面的自然运算一起使用。例如，两个分次向量空间相同的上同调，如果运算的作用不同，就会给出不同的代数输入。为此需要三项新结构：描述这些运算及其关系的代数；在该代数的模范畴中构造消解和高阶 Ext；把逐层消解实现为谱的塔，并证明它在相应素数完备化意义下收敛。第4章的 Ext 是理解消解的起点，但只计算整数模之间的一次 Ext，还没有包含这些结构。第13章将先讨论上同调运算，第9章的稳定化则提供稳定同伦的入口。

\ProseParagraphBreak

Eilenberg–Moore 方法面向同伦拉回。对于 \(F\to E\to B\)，纤维可以理解为 \(E\) 与一个点在 \(B\) 上的同伦拉回；相应的代数模型涉及把 \(C^*(E;k)\) 与 \(k\) 在 \(C^*(B;k)\) 上作导出张量积。这里 \(k\) 是域，“导出”表示先用适当的自由模型替换，以保留普通张量积会丢失的关系。它所需的新输入是带乘法的上链模型、条形消解，以及将该代数模型与实际同伦纤维比较的定理。在单连通与适当有限型条件下，第二页可用分次上同调代数上的 Tor 描述；这些比较与收敛条件须随新构造一同证明。

\ProseParagraphBreak

一个小的代数计算能说明为什么会出现这种输入。设 \(A=k[u]\)、\(|u|=2\)，增广把 \(u\) 送到零。\(k\) 作为分次 \(A\) 模有自由消解
\[
	0\longrightarrow A\varepsilon\xrightarrow{\ \varepsilon\mapsto u\ }A\longrightarrow k\longrightarrow0,
\]
其中 \(\varepsilon\) 的内部次数为二、消解次数为一。张量 \(k\) 后微分为零，所得 Tor 只有单位和一个总次数 \(2-1=1\) 的生成元；用 \(\varepsilon^2=0\) 的自由微分分次代数模型，可见其乘法为外代数。这个答案与 \(\Omega\mathbb CP^\infty\simeq S^1\) 的上同调吻合；后一个同伦型由 \(\mathbb CP^\infty=K(\mathbb Z,2)\)、路径长正合列及上面的模型唯一性得到。要对一般纤维化作同样的推断，承担几何与代数之间联系的正是新的比较定理。

\ProseParagraphBreak

本章的下一步则回到空间本身：有限截段 \(X[n]\to X[n-1]\) 已把一个同伦群分离为纤维，但仍未说明这些纤维怎样附着成原空间。第12章的障碍类与 Postnikov 分解将记录这部分数据，使这里用于证明有限性的构造进一步成为描述同伦类型的工具。
% 纲要条目排为编号小节标题，并进入目录与书签；说明段仍保留于纲要注释中。


\begin{chaptersummary}
	基底骨架的逆像给出 Serre 谱序列，第二页为基底取纤维同调或上同调局部系统后的群。沿路运输决定第二页是否能够使用常值系数，固定次数的骨架稳定及相容 CW 模型保证有限目标滤过。上同调页面的乘法满足总次数 Leibniz 规则，边缘分别对应投影拉回与纤维限制；传递则是纤维类延拓到下一基底骨架的障碍。Hopf 丛迫使相应传递为同构，定向球丛的两行给出 Euler 类与上同调正合列。可缩的道路空间迫使回路球面各次的微分为同构，进而决定加性群与整数除幂乘法。酉群的球面纤维序列由生成元次数迫使退化，再以自由群和乘积基解决环扩张。有限生成性经取子群、商和有限扩张保持；将单个同伦群放进 Eilenberg–MacLane 纤维后，这一性质由谱序列从同调传到全部高阶同伦群。
\end{chaptersummary}

\BookExercises

\begin{exercise}\label{g2:ex:serre-invariant-not-permanent}
	比较 \(S^1\to S^2\times S^1\to S^2\) 与 \(S^1\to S^3\to S^2\) 的上同调谱序列。证明两者所有纤维类都保持运输不变，但圆周的一次生成元只有在前一个纤维化中才是永久圈。解释“运输不变”与“能够从总空间限制下来”的区别。
\end{exercise}
\begin{solution}
	两者基底均单连通，所以运输作用平凡。第二页同为 \(\mathbb Z[x]/(x^2)\otimes\Lambda(u)\)，其中 \(x\) 位于 \((2,0)\)，\(u\) 位于 \((0,1)\)。唯一可能非零的微分为 \(\mathrm d_2u=kx\)。

	乘积中，第二投影拉回的圆周类限制到纤维恰为 \(u\)。边缘映射遂说明 \(u\) 是永久圈，\(k=0\)，乘积规则又使所有微分为零。Hopf 丛中，\(H^1(S^3)=H^2(S^3)=0\)，所以 \(\mathrm d_2\) 必须无核且无余核，即 \(k=\pm1\)。因此 \(u\) 在第二页被消去，不能是某个总空间上同调类的限制。运输不变只检查沿基底道路的相容性，使类能够进入第二页；更高骨架上的延拓还需要通过后续微分。
\end{solution}

\begin{exercise}\label{g2:ex:serre-torus-monodromy}
	令 \(f:T^2\to T^2\) 为 \(f([x],[y])=([-x],[-y])\)，记其映射环面为 \(M_f\)。用 Serre 谱序列求 \(H^k(M_f;\mathbb Z)\)，并说明哪一步需要处理加性扩张。
\end{exercise}
\begin{solution}
	\(H^0(T^2;\mathbb Z)=H^2(T^2;\mathbb Z)=\mathbb Z\)，\(H^1(T^2;\mathbb Z)=\mathbb Z^2\)。\(f^*\) 在一次是 \(-I\)，在二次是 \(\det(-I)=1\)，在零次也是恒等。因为 \(f^{-1}=f\)，上同调局部运输与这些作用相同。圆周的局部系数上链微分为作用的逆减恒等；因此第一行和第三行各在 \(p=0,1\) 有 \(\mathbb Z\)，中间一行的 \(p=0\) 项为零，\(p=1\) 项为 \((\mathbb Z/2)^2\)。只有两列，故 \(E_2=E_\infty\)。

	零、一、三次各只有一个非零分次块，分别给出 \(\mathbb Z\)。二次有短正合列
	\[
		0\longrightarrow(\mathbb Z/2)^2\longrightarrow H^2(M_f;\mathbb Z)\longrightarrow\mathbb Z\longrightarrow0.
	\]
	商群 \(\mathbb Z\) 自由，任选其生成元的一个提升就给出加性分裂，因而 \(H^2(M_f;\mathbb Z)\cong\mathbb Z\oplus(\mathbb Z/2)^2\)。其他次数为零。这里分裂由自由商保证；仅凭谱序列退化还不能跳过这一步。
\end{solution}

\begin{exercise}\label{g2:ex:serre-covering-coefficients}
	对 \(m\ge2\)，考虑覆盖 \(p:S^1\to S^1\)、\(p(z)=z^m\)。把它看成纤维有 \(m\) 个点的纤维化，写出第二页，并从局部系数映射解释 \(p^*:H^1(S^1;\mathbb Z)\to H^1(S^1;\mathbb Z)\) 为乘 \(m\)。
\end{exercise}
\begin{solution}
	纤维上同调只有零次非零，为 \(M=\mathbb Z^m\)。绕基底一周循环置换这 \(m\) 个点，因而在 \(M\) 上循环置换标准基，记该作用为 \(P\)；取相反循环置换也不改变以下核与余核。局部上链为 \(M\xrightarrow{P^{-1}-I}M\)，其核是对角子群 \(\mathbb Z(1,\ldots,1)\)，余核是所有标准基互相识别后得到的 \(\mathbb Z\)。因此 \(E_2^{0,0}\cong E_2^{1,0}\cong\mathbb Z\)，其余项为零，没有后续微分和加性扩张。

	投影拉回在系数上把一个常值整数 \(a\) 送到所有纤维点上的同一个值，即 \(a\mapsto(a,\ldots,a)\)。在一次局部上同调的余核中，\((1,\ldots,1)\) 是 \(m\) 个相同生成元之和，所以诱导映射为乘 \(m\)。这个算例同时说明，纤维不连通时，零行系数是 \(\mathcal H^0(F;\mathbb Z)\)，不能径直换成一份常值 \(\mathbb Z\)。
\end{solution}


\begin{exercise}\label{g2:ex:serre-global-fiber-basis}
	设 \(F\to E\xrightarrow{p}B\) 为道路连通空间的纤维化，\(B\) 为连通 CW 复形。假设 \(H^*(F;\mathbb Z)\) 有有限个齐次元素组成的整数基，且存在 \(e_1,\ldots,e_N\in H^*(E;\mathbb Z)\)，使它们在一个纤维上的限制组成这组基。证明
	\[
		\bigoplus_{j=1}^NH^{k-|e_j|}(B;\mathbb Z)\longrightarrow H^k(E;\mathbb Z),\qquad
		(b_j)_j\longmapsto\sum_jp^*b_j\smile e_j
	\]
	对每个 \(k\) 都是同构。说明结论为何依赖于这些实际存在的全局类。
\end{exercise}
\begin{solution}
	总空间上同调类在各纤维上的限制与道路运输相容，所以 \(e_j|_F\) 全部运输不变。它们构成每次纤维群的基，故整个局部系统为常值。有限秩自由系数的上链计算给出
	\[
		E_2^{p,q}=H^p(B;\mathbb Z)\otimes H^q(F;\mathbb Z).
	\]
	由边缘映射，每个 \(e_j|_F\) 都是永久圈。零行上的微分向负纤维次数出发，因而为零；每个第二页元素都是底类与这些纤维基元的线性组合，Leibniz 规则使全部第二微分为零。逐页重复，得到 \(E_2=E_\infty\)。

	给左侧按底类次数的递减滤过，右侧取 Serre 滤过。\(p^*b_j\) 的滤过至少为 \(|b_j|\)，而 \(e_j\) 的零列首项为 \(e_j|_F\)，所以题中映射保留滤过。它在每个伴随分次上正好是所列第二页同构。固定 \(k\)，滤过有限，对相邻滤过项的短正合列归纳就得到目标同构。这里也解决了加性扩张，因为选定的全局 \(e_j\) 已给出具体提升。若只有抽象的纤维基而没有全局提升，Hopf 丛的 \(u\) 就可能被微分消去，结论不成立。
\end{solution}

\begin{exercise}\label{g2:ex:circle-bundle-projective-plane-ring}
	设 \(S^1\to E\xrightarrow{p}\mathbb CP^2\) 为定向圆周丛，且 \(e(p)=k x\)，其中 \(x\) 生成 \(H^2(\mathbb CP^2;\mathbb Z)\)、\(k\ne0\)。求全部整数上同调群及杯积环。特别指出 \(|k|>1\) 时二次挠类的平方。
\end{exercise}
\begin{solution}
	第二页为 \(\mathbb Z[x]/(x^3)\otimes\Lambda(u)\)，其中 \(|x|=(2,0)\)、\(|u|=(0,1)\)。按 Euler 类定义，\(\mathrm d_2u=kx\)，所以
	\[
		\mathrm d_2(xu)=kx^2,\qquad \mathrm d_2(x^2u)=0.
	\]
	前两张微分在整数群上单射，分别留下 \(\mathbb Z/|k|\) 的余核。第三页以后没有可能的微分，因此总上同调为
	\[
		H^j(E;\mathbb Z)\cong
		\begin{cases}
			\mathbb Z,&j=0,5,\\
			\mathbb Z/|k|,&j=2,4,\\
			0,&\text{其余 }j.
		\end{cases}
	\]
	每个次数只有一层非零滤过。令 \(a=p^*x\)，再选 \(v\in H^5(E)\) 对应 \(x^2u\) 的生成元。乘法与伴随分次相容，故 \(a^2\) 生成四次群；当 \(|k|>1\) 时尤其非零。其余正次乘积落在五次以上或含 \(a^3\)，均为零。完整环可写成以 \(|a|=2,|v|=5\) 为生成元的分次交换环，关系为 \(ka=0,a^3=0,av=0,v^2=0\)。\(k=\pm1\) 时 \(a=0\)，只留下零次与五次整数群。
\end{solution}

\begin{exercise}\label{g2:ex:serre-hopf-pullback-degree}
	令 \(p:S^3\to S^2\) 为复 Hopf 丛，选择生成元 \(u\in H^2(S^2;\mathbb Z)\) 使 \(e(p)=u\)。对整数 \(k\ne0\)，取度为 \(k\) 的映射 \(f_k:S^2\to S^2\)，记拉回圆周丛为 \(p_k:E_k\to S^2\)。求 \(E_k\) 的全部整数上同调群与正次杯积，并证明 \(p_k\) 没有截面。再比较 \(k=0\) 时的加性群。
\end{exercise}
\begin{solution}
	自然性给出 \(e(p_k)=f_k^*u=ku\)。球丛正合列中的唯一非零乘 Euler 类映射为 \(H^0(S^2)=\mathbb Z\xrightarrow{\times k}H^2(S^2)=\mathbb Z\)。因此当 \(k\ne0\) 时，
	\[
		H^j(E_k;\mathbb Z)=
		\begin{cases}
			\mathbb Z,&j=0,3,\\
			\mathbb Z/|k|\mathbb Z,&j=2,\\
			0,&\text{其余 }j.
		\end{cases}
	\]
	所有正次类的乘积都落在至少四次，故全为零。由于 \(e(p_k)=ku\ne0\)，截面消失结论排除了连续截面。

	当 \(k=0\) 时，正合列中的乘 Euler 类映射全为零，给出 \(H^j(E_0;\mathbb Z)=\mathbb Z\)（\(j=0,1,2,3\)），其余为零。每个次数至多有一层非零滤过，故这里没有加性延拓问题。该结果与 \(S^2\times S^1\) 的加性上同调相同；也可取常值 \(f_0\)，其拉回直接就是这个乘积丛。
\end{solution}

\begin{exercise}\label{g2:ex:even-sphere-euler-two-torsion}
	设 \(m\ge2\) 为偶数，\(p:E\to B\) 为道路连通 CW 复形上的定向 \(S^m\) 丛。证明 \(2e(p)=0\)。由此说明：若 \(H^{m+1}(B;\mathbb Z)\) 没有 \(2\) 挠，则 \(e(p)=0\)；特别地，定向 \(S^2\) 丛在底空间 \(S^3\) 上不可能具有非零 Euler 类。
\end{exercise}
\begin{solution}
	令 \(u\in E_2^{0,m}\) 为纤维定向生成元。因为纤维没有 \(2m\) 次上同调，\(u^2=0\)。在第 \(m+1\) 页以前没有非零微分，故该等式在 \(E_{m+1}\) 仍成立。\(m\) 为偶数，Leibniz 规则以及 \(u\) 与 \(e(p)\) 的交换性给出
	\[
		0=\mathrm d_{m+1}(u^2)=e(p)u+u e(p)=2e(p)u.
	\]
	乘 \(u\) 把 \(H^{m+1}(B;\mathbb Z)\) 同构地识别为顶行 \(E_{m+1}^{m+1,m}\)，所以 \(2e(p)=0\)。无 \(2\) 挠时这就迫使 \(e(p)=0\)。对于 \(B=S^3\)、\(m=2\)，所需群为 \(H^3(S^3;\mathbb Z)=\mathbb Z\)，故结论适用。这里不能把 \(m\) 改为奇数：奇数时 Leibniz 规则的两项互相抵消，而 Hopf 丛的 Euler 类可以有无穷阶。
\end{solution}

\begin{exercise}\label{g2:ex:divided-power-finite-generation}
	设 \(m\ge2\) 为偶数，记 \(A=H^*(\Omega S^{m+1};\mathbb Z)=\Gamma_{\mathbb Z}(a_1)\)。证明 \(A\) 在每个次数都是有限生成阿贝尔群，但作为 \(\mathbb Z\) 代数不是有限生成的。再证明换成有理系数后，\(H^*(\Omega S^{m+1};\mathbb Q)\) 是一个元素生成的多项式代数。
\end{exercise}
\begin{solution}
	每个次数的群为零或 \(\mathbb Z\)，第一项已成立。若 \(A\) 由有限多个元素作为代数生成，把这些元素写成加性基 \(a_j\) 的有限和，取 \(N\) 使其中出现的下标都不超过 \(N\)。它们生成的子代数包含在由 \(a_0,a_1,\ldots,a_N\) 生成的子代数中。选素数 \(\ell>N\)，再把系数约化到 \(\mathbb F_\ell\)。

	任何由 \(a_1,\ldots,a_N\) 构成、次数为 \(\ell m\) 的单项式，其下标之和为 \(\ell\)，各正下标都严格小于 \(\ell\)。重复应用除幂乘法，它等于
	\[
		a_{j_1}\cdots a_{j_s}=\frac{\ell!}{j_1!\cdots j_s!}\,a_\ell,
		\qquad j_1+\cdots+j_s=\ell,\quad 1\le j_i\le N.
	\]
	分子含素因子 \(\ell\)，每个分母都不含，故整数系数被 \(\ell\) 整除。于是所有这些单项式在 \(A\otimes\mathbb F_\ell\) 中为零，而 \(a_\ell\) 的像非零，矛盾。非齐次候选生成元也没有例外，因为取它们的齐次分量只会扩大生成的子代数。

	整数同调群均自由，系数定理把有理上同调识别为 \(A\otimes\mathbb Q\)，且系数映射保持杯积。等式 \(a_1^j=j!a_j\) 此时允许除以 \(j!\)，说明 \(a_1\) 生成整个环。各 \(a_1^j\) 位于不同次数且非零，没有非零多项式关系，故得到 \(\mathbb Q[a_1]\)。这区分了逐次数的有限生成性与整个分次代数的有限生成性。
\end{solution}


\begin{exercise}\label{g2:ex:cohomology-suspension-products}
	设 \(B\) 为带基点的道路连通 CW 复形，\(\alpha\in\widetilde H^a(B;\mathbb Z)\)、\(\beta\in\widetilde H^b(B;\mathbb Z)\)，其中 \(a,b>0\)。证明上同调悬挂满足 \(\sigma(\alpha\smile\beta)=0\)。用 \(B=\mathbb CP^\infty\) 说明它为何不能在所有次数上成为同构。
\end{exercise}
\begin{solution}
	先说明悬挂空间的任意两个正次类的乘积为零。将悬挂 \(\Sigma Y\) 用稍加延长的上下两个开锥 \(U,V\) 覆盖，两开集各自可缩。任意正次类 \(c,d\) 分别在 \(U,V\) 上限制为零，所以由空间对正合列可选相对提升
	\[
		\widetilde c\in H^*(\Sigma Y,U),\qquad \widetilde d\in H^*(\Sigma Y,V).
	\]
	相对杯积落在 \(H^*(\Sigma Y,U\cup V)=H^*(\Sigma Y,\Sigma Y)=0\)，故绝对乘积 \(cd=0\)。

	对 \(Y=\Omega B\)，求值给出 \(\operatorname{ev}:\Sigma\Omega B\to B\)。杯积自然性使
	\[
		\operatorname{ev}^*(\alpha\beta)=\operatorname{ev}^*\alpha\smile\operatorname{ev}^*\beta=0.
	\]
	再作悬挂次数同构，按正文对 \(\sigma\) 的识别得到 \(\sigma(\alpha\beta)=0\)。在 \(B=\mathbb CP^\infty\) 时，二次生成元 \(u\) 的首次悬挂生成 \(H^1(\Omega B)=\mathbb Z\)，但 \(u^j\ne0\)（\(j\ge2\)）全是正次类的乘积，故 \(\sigma(u^j)=0\)。因此首次非零次数的同构不能不加条件地延伸到高次。
\end{solution}

\begin{exercise}\label{g2:ex:unitary-determinant-product}
	设 \(n\ge2\)。证明
	\[
		\Phi:S^1\times\operatorname{SU}(n)\longrightarrow\operatorname{U}(n),\qquad(z,A)\longmapsto\operatorname{diag}(z,1,\ldots,1)A
	\]
	为同胚，写出其逆，并说明行列式在这个乘积中对应哪张投影。这一同胚一般是否为直积群到酉群的群同态？比较所得上同调环与\cref{g2:thm:unitary-cohomology}。
\end{exercise}
\begin{solution}
	记 \(D(z)=\operatorname{diag}(z,1,\ldots,1)\)。逆映射为
	\[
		U\longmapsto\bigl(\det U,D(\det U)^{-1}U\bigr).
	\]
	第二分量的行列式为一，两张公式连续且互逆，因此 \(\Phi\) 为同胚；\(\det\Phi(z,A)=z\)，所以行列式对应第一投影。在直积群中 \((z,I)\) 与 \((1,A)\) 总交换，但它们的像 \(D(z)\) 与 \(A\) 未必交换。例如在前两个坐标中取 \(A=\left(\begin{smallmatrix}0&1\\-1&0\end{smallmatrix}\right)\)，其他坐标为单位阵；当 \(z\ne1\) 时，两张矩阵不交换。因此 \(\Phi\) 一般不是群同态。

	两因子的整数上同调都自由，Künneth 定理及外积乘法给出
	\[
		H^*(\operatorname{U}(n);\mathbb Z)\cong H^*(S^1;\mathbb Z)\otimes H^*(\operatorname{SU}(n);\mathbb Z)
		\cong\Lambda_{\mathbb Z}(x_1,y_3,\ldots,y_{2n-1}).
	\]
	其中一次类正是行列式拉回的生成元，与正文外代数一致。这个结论比较的是空间的杯积环，并未要求所选同胚保持群乘法。
\end{solution}

\begin{exercise}\label{g2:ex:unitary-four-middle-pairing}
	在 \(\operatorname{U}(4)\) 中选择\cref{g2:thm:unitary-cohomology}的整数生成元，并将 \(v=x_1x_3x_5x_7\) 取作十六次生成元。求 \(H^8(\operatorname{U}(4);\mathbb Z)\) 的一组基，并写出杯积 \(H^8\times H^8\rightarrow H^{16}=\mathbb Zv\) 对应的整数矩阵。
\end{exercise}
\begin{solution}
	在 \(1,3,5,7\) 中选取互异项，和等于 \(8\) 的方式只有 \(1+7\) 与 \(3+5\)。所以一组基为 \(a=x_1x_7\)、\(b=x_3x_5\)。平方 \(a^2,b^2\) 都含某个奇次生成元的平方，故均为零。而
	\[
		ab=x_1x_7x_3x_5=(-1)^2x_1x_3x_5x_7=v.
	\]
	两次交换把 \(x_7\) 移到末尾，因而没有负号。\(|a|=|b|=8\) 还给出 \(ba=ab\)，所以矩阵为 \(\left(\begin{smallmatrix}0&1\\1&0\end{smallmatrix}\right)\)。其行列式为 \(-1\)，与闭定向十六维流形中维整数交形式的单模性相符；若把顶次生成元换成 \(-v\)，矩阵整体变号。
\end{solution}


\begin{exercise}\label{g2:ex:hurewicz-mod-finite-generation}
	设 \(X\) 为单连通 CW 复形，\(n\ge2\)，且 \(\pi_i(X)\) 对 \(2\le i<n\) 都有限生成。证明 Hurewicz 同态 \(h:\pi_n(X)\to H_n(X;\mathbb Z)\) 的核与余核都有限生成。这里不预先假定 \(\pi_n(X)\) 或 \(H_n(X)\) 有限生成。
\end{exercise}
\begin{solution}
	\(n=2\) 时 Hurewicz 同态为同构，核与余核为零。设 \(n>2\)，按正文构造 \(F_n\to X[n]\to X[n-1]\)，记总空间为 \(Z\)、基底为 \(Y\)。低阶同伦群有限生成使 \(Y\) 的全部同调逐次有限生成。\(F_n\) 为 \((n-1)\) 连通，且 \(\pi_n(F_n)\to\pi_n(Z)\) 为同构。Hurewicz 自然性给出交换图
	\[
		\begin{tikzcd}
			\pi_n(F_n)\arrow[r,"\cong"]\arrow[d,"\cong"']&\pi_n(Z)\arrow[d,"h_Z"]\\
			H_n(F_n)\arrow[r]&H_n(Z).
		\end{tikzcd}
	\]
	下箭头的核由\cref{g2:eq:finite-generation-five-terms}是 \(H_{n+1}(Y)\) 的像，余核同构于 \(H_n(Y)\)，两者均有限生成。因此右箭头也有有限生成的核与余核。包含 \(X\to X[n]\) 在 \(\pi_n\) 和 \(H_n\) 上都是同构，并与 Hurewicz 映射交换，所以同一结论适用于 \(h\)。这正是 Hurewicz 映射模 \(\mathcal C_{\mathrm{fg}}\) 为同构的含义。
\end{solution}

\begin{exercise}\label{g2:ex:finite-homology-infinite-pi-two}
	令 \(X=S^1\vee S^2\)，基点取楔和点。构造其万有覆盖，并证明 \(H_k(X;\mathbb Z)\) 对每个 \(k\) 都有限生成，但
	\[
		\pi_2(X)\cong\bigoplus_{j\in\mathbb Z}\mathbb Z.
	\]
	说明这个例子对有限生成定理的单连通条件有何意义。
\end{exercise}
\begin{solution}
	把实直线按整数点分成一胞腔，并在每个整数点 \(j\) 附着一个球面 \(S_j^2\)。所得空间记为 \(\widetilde X\)。将实直线按 \(t\mapsto e^{2\pi\mathrm it}\) 投到 \(S^1\)，每个 \(S_j^2\) 同胚地投到 \(S^2\)，得到覆盖 \(\widetilde X\rightarrow X\)：在楔和点附近，圆周的一段小弧连同球面的一小片，其原像是围绕各整数点的互不相交同胚副本。对其余点，覆盖性质直接由圆周覆盖及球面同胚得到。

	\(\widetilde X\) 的一维骨架为实直线，而每个二胞腔沿常值映射附着，因此其基本群为零，所构造的覆盖是万有覆盖。胞腔链中没有三胞腔，二胞腔边界也为零，所以
	\[
		H_2(\widetilde X;\mathbb Z)=\bigoplus_{j\in\mathbb Z}\mathbb Z[S_j^2].
	\]
	\cref{g2:thm:hurewicz-absolute}用于单连通的 \(\widetilde X\)，得到 \(\pi_2(\widetilde X)\cong H_2(\widetilde X;\mathbb Z)\)。覆盖映射的离散纤维在正阶同伦群中为零，纤维化长正合列因而给出 \(\pi_2(\widetilde X)\cong\pi_2(X)\)，得到题中同构。这个群不是有限生成的，因为有限多个元素只涉及有限多个直和坐标。

	另一方面，\(X\) 有一个零胞腔、一个一胞腔和一个二胞腔，胞腔边界全为零。因此 \(H_0(X)=H_1(X)=H_2(X)=\mathbb Z\)，更高同调为零。\(X\) 虽然本身是有限 CW 复形，\(\pi_2(X)\) 仍不是有限生成阿贝尔群；单连通条件不能在有限生成定理中直接删去。绕圆周一圈把 \(S_j^2\) 移到 \(S_{j+1}^2\)，正显示了基本群作用如何产生无穷多个整数生成元。
\end{solution}

\begin{exercise}\label{g2:ex:finite-fundamental-group-finiteness}
	设 \(X\) 为连通有限 CW 复形，且 \(\pi_1(X)\) 为有限群。证明所有 \(\pi_n(X)\)（\(n\ge2\)）都是有限生成阿贝尔群。
\end{exercise}
\begin{solution}
	令 \(m=|\pi_1(X)|\)。万有覆盖 \(\widetilde X\rightarrow X\) 有 \(m\) 层。将 \(X\) 的胞腔特征映射提升到覆盖空间，得到 \(\widetilde X\) 的 CW 结构；每个底胞腔恰有 \(m\) 个提升胞腔。因此 \(\widetilde X\) 也是有限 CW 复形，其每个整数同调群都有限生成。它单连通，\cref{g2:thm:serre-finite-generation}于是说明 \(\pi_n(\widetilde X)\) 对所有 \(n\ge2\) 都有限生成。离散纤维的长正合列给出 \(\pi_n(\widetilde X)\cong\pi_n(X)\)，完成证明。
\end{solution}

\begin{exercise}\label{g2:ex:finite-type-not-finite-cw}
	设 \(m>1\) 为整数，\(G=\mathbb Z/m\)。利用本节的有限群模型说明 \(K(G,1)\) 的每个整数同调群都有限生成，再证明它不可能具有有限 CW 同伦型。
\end{exercise}
\begin{solution}
	本节构造的 \(BG\) 每维只有有限多个胞腔，因为该维单形由有限群 \(G\) 的有限长度序列决定。因此每个胞腔链群都是有限秩自由阿贝尔群，圈群为其子群，仍有限生成，取商后每个 \(H_k(BG;\mathbb Z)\) 都有限生成。\(BG\) 是 \(K(G,1)\)，同伦不变性便给出前一结论。

	假设某个连通有限 CW 复形 \(Y\) 与 \(K(G,1)\) 同伦等价。其万有覆盖 \(\widetilde Y\) 有 \(m\) 层，每个胞腔提升成 \(m\) 个胞腔，所以 \(\widetilde Y\) 也有限，且胞腔计数给出
	\[
		\chi(\widetilde Y)=m\chi(Y).
	\]
	\(Y\) 的全部高阶同伦群为零，覆盖映射在这些群上诱导同构；又因 \(\widetilde Y\) 单连通，它的全部正阶同伦群都为零。\cref{g2:thm:whitehead}因而说明 \(\widetilde Y\) 可缩，所以 \(\chi(\widetilde Y)=1\)。于是 \(1=m\chi(Y)\)，与 \(m>1\) 及 \(\chi(Y)\in\mathbb Z\) 矛盾。逐次同调有限生成允许非零信息出现于无穷多个次数，也并不保证存在只有有限多个胞腔的模型。
\end{solution}


% !TeX root = ../../Book2.tex
% 纲要标题：## 第12章*　障碍理论与 Postnikov 分解
% 纲要位置：几何学计划纲要.md，第 1468 行。
% 正文状态：初稿；尚未正式审核。

% 纲要内容开始
% 
% > 本章目标与证明范围：完整建立简单空间或明确局部系数条件下的首障碍与逐层延拓，随后构造 Postnikov 阶段。障碍群、具体障碍类和消失后的选择自由分别处理。
% 
% 纲要内容结束

\OptionalChapter{障碍理论与 Postnikov 分解}{g2:ch:12}
\BookChapterNumber{g2:ch:12}

\begin{chapterabstract}
	沿胞腔延拓一张映射或一张截面时，边界留下的球面映射决定下一步能否完成。本章把这些局部条件组成带同伦群系数的障碍余圈，证明余圈条件、改变选择的余边界公式、首障碍的选择无关性及延拓之间的差异分类，并由此证明上同调的 Eilenberg–MacLane 可表性。随后构造单连通空间的 Postnikov 塔，以映射柱和相对 Hurewicz 定理将每一阶段恢复为一张上同调类的同伦纤维。最后将理论用于球丛、向量场和结构群约化，比较首截面障碍与 Euler 类，并完成低维底空间上的存在性与同伦分类计算。
\end{chapterabstract}

\begin{learninggoals}
	\begin{itemize}
		\item 将相对映射延拓和截面延拓写成逐胞腔问题，识别各步实际使用的同伦群与局部系数。
		\item 构造障碍余圈，证明改变上一层选择产生余边界，区分固定当前映射与允许修改后的延拓条件。
		\item 说明首障碍何时与选择无关，并用差异类描述延拓存在后的自由度。
		\item 构造 Postnikov 阶段及其 Eilenberg–MacLane 纤维，使用 \(k\)-不变量恢复阶段并判断提升。
		\item 计算球丛和低维向量丛的截面障碍，说明独立截面如何给出结构群约化。
	\end{itemize}
\end{learninggoals}

在圆周上给定 \(f(z)=z^k\)，什么时候能把它填成一张 \(D^2\to S^1\) 的映射？答案是 \(k=0\)：圆盘中的同心圆可以把边界映射收缩到中心，所以一个非零绕数无法被填掉。若边界值来自某个丛中的单位向量，选取圆盘上的平凡化后，同样会出现一个绕数。于是“填一张映射”和“连续选一个向量”可以由同一种延拓判据处理。

\ProseParagraphBreak

把圆盘换成有许多胞腔的空间，问题多了一层：每张胞腔的边界可能单独绕出非零类，但修改相邻低维胞腔上的选择，也会改变这些类。障碍上链记录当前选择留下的全部边界；障碍上同调类则把这种修改带来的变化除去。这一区别正是“当前填不进去”与“调整上一层以后仍然填不进去”的区别，贯穿本章的存在性证明。

\ProseParagraphBreak

延拓还可以按目标的同伦次数来组织。一个空间的二次、三次同伦群分别是什么，并不能说明二次信息与三次信息怎样接在一起。Postnikov 阶段逐层保留这些群，\(k\)-不变量记录相邻阶段的粘合；本章将构造一族所有同伦群都相同、上同调却不同的空间，让这项额外信息直接出现在计算中。

\ProseParagraphBreak

本章使用第1章的相对 CW 同伦延拓性质、第5章的局部系数、第8章的纤维运输与同伦长正合列、第9章的 Hurewicz 定理和消去同伦群的构造。Postnikov 阶段与 Euler 类的比较还使用第11章的传递。读到相应位置时会重新指出这些结果所需的条件。首障碍与逐层延拓在简单纤维的范围内建立，非平凡的基底运输保留为局部系数；Postnikov 的重构证明在单连通空间中展开。几何应用涉及第一卷向量丛的标架、内积和正交补，所用构造在正文中作简要回顾。

% !TeX root = ../../Book2.tex
% 纲要标题：### 12.1 映射和截面的延拓问题
% 纲要位置：几何学计划纲要.md，第 1472 行。
% 正文状态：初稿；尚未正式审核。

% 纲要内容开始
% 
% * 空间对、骨架与相对条件
% * 单个胞腔上的延拓判据
% * 基本群作用产生的系数系统
% 
% 纲要内容结束

\section{映射和截面的延拓问题}
\label{g2:sec:1201}

给圆周上的每一点指定一个单位向量以后，能否把这个选择连续地填到整个圆盘？如果向量已经绕了一整圈，填充时便会遇到障碍。更高维的情形保留了同一结构：一个新胞腔是圆盘，旧数据已经规定它的边界；是否能够填进去，由边界在目标或纤维中的同伦类决定。先把这些逐胞腔的问题写准确，下一节才能把它们组织为上同调类。

\subsection{空间对、骨架与相对条件}
\label{g2:subsec:1201:01}

设 \((X,A)\) 为 CW 对，给定映射 \(f_A:A\to Y\)。\term[yingshe-yantuo-wenti]{映射延拓问题}[extension problem]要求寻找 \(f:X\to Y\)，满足 \(f|_A=f_A\)。这里边界数据是实际给定的映射；若只要求其同伦类相同，就需要另作说明。记
\begin{equation}\label{g2:eq:relative-obstruction-skeleta}
	X_{-1}=A,\qquad X_r=A\cup X^r\quad(r\ge0).
\end{equation}
这些是相对骨架：即使 \(A\) 含高维胞腔，它也从一开始就全部保留。下标 \(r\) 只控制 \(X\setminus A\) 中已经加入的胞腔；它与上一章表示同伦截段的 \(X[n]\) 含义不同。

\ProseParagraphBreak

更一般地，给定纤维化 \(p:E\to X\) 及 \(A\) 上的截面 \(s_A\)，要求寻找满足
\[
	p\circ s=\operatorname{id}_X,\qquad s|_A=s_A
\]
的截面。映射延拓是它的特例：取 \(E=X\times Y\)，截面必形如 \(s(x)=(x,f(x))\)。而对一般提升问题，若已有 \(g:X\to B\)、纤维化 \(p_0:E_0\to B\) 和 \(g|_A\) 的一个指定提升，则先取拉回
\[
	g^*E_0=\{(x,e):g(x)=p_0(e)\}\longrightarrow X.
\]
截面 \(x\mapsto(x,\widetilde g(x))\) 与满足原边界条件的提升 \(\widetilde g\) 一一对应。于是映射、提升和截面可以用同一套逐骨架方法处理。

\ProseParagraphBreak

从 \(X_r\) 到 \(X_{r+1}\) 时，每个新胞腔的边界已经有值，同维胞腔内部则互不相交。只要对每个特征圆盘找到保持边界的延拓，就能沿原来的附着映射粘合。CW 弱拓扑保证最后所得映射连续：连续性只需在 \(A\) 及每个特征圆盘上核对。对无限 CW 复形，这个事实允许逐维构造相容映射；在各个有限阶段分别找到互不相容的延拓，则还没有得到一张定义在 \(X\) 上的映射。

\begin{example}[圆盘边界上的绕数]\label{g2:exm:boundary-winding-extension}
	令 \(X=D^2\)、\(A=S^1\)，取 \(f_A(z)=z^k:S^1\to S^1\)，其中 \(k\in\mathbb Z\)。若它延拓到 \(D^2\)，把圆盘沿半径收缩到中心，就得到 \(f_A\) 的自由零伦，故绕数必须为零。反过来，\(k=0\) 时常值映射本身便给出延拓。因此延拓恰在 \(k=0\) 时存在。这同时是平凡实二维丛上的单位截面延拓问题：单位向量的边界绕数阻止它无零点地填入圆盘。
\end{example}

还须区分两种要求：把一个已经选好的 \(s_r:X_r\to E\) 原样延到下一层，和只固定 \(X_{r-1}\) 上的数据、允许重新选择 \(r\) 胞腔内的截面。第一种要求逐胞腔的障碍值全部为零；第二种允许用新选择改变这些值，因而将在下一节产生“障碍类为零”的判据。

\ProseParagraphBreak

比较两张延拓的同伦也属于同一问题。若 \(s_0,s_1\) 在 \(A\) 上相同，寻找固定 \(A\) 的截面同伦，就是在
\[
	\bigl(X\times I,\ (X\times\partial I)\cup(A\times I)\bigr)
\]
上延拓截面，其中两端取 \(s_0,s_1\)，侧面保持给定值。每个 \(r\) 胞腔乘区间成为一个 \((r+1)\) 维延拓问题，这将解释差异类的次数为何比存在性障碍低一度。

\subsection{单个胞腔上的延拓判据}
\label{g2:subsec:1201:02}

先看普通映射。设 \(q\ge1\)，\(\chi:D^{q+1}\to X_{q+1}\) 是一个新胞腔的特征映射，\(f_q:X_q\to Y\) 已经给定。复合
\[
	S^q\xrightarrow{\chi|_{S^q}}X_q\xrightarrow{f_q}Y
\]
就是待填圆盘上的边界数据。

\begin{proposition}[单个胞腔的映射延拓]\label{g2:prop:cell-extension-criterion}
	设 \(q\ge1\)，\(Y\) 为道路连通空间。映射 \(u:S^q\to Y\) 能延拓到 \(D^{q+1}\)，当且仅当它自由零伦；选择球面基点及目标基点识别后，这等价于对应的 \(\pi_q(Y)\) 类为零。因此，\(f_q:X_q\to Y\) 能保持已有值延到 \(X_{q+1}\)，当且仅当每个相对 \((q+1)\) 胞腔给出的这个类均为零。
\end{proposition}
\begin{proof}
	若 \(U:D^{q+1}\to Y\) 延拓 \(u\)，则 \(H(z,t)=U((1-t)z)\) 把边界映射收缩到 \(U(0)\)。反过来，给定从 \(u\) 到常值的同伦，将 \(D^{q+1}\) 写为 \(S^q\) 的锥，这个同伦就在锥顶处下降为所需延拓。

	由\cref{g2:cor:free-homotopy-orbits}，自由同伦对应基本群对带基点同伦群的作用轨道。任何作用自同构都保持零元，且只有零元能被送到零元，因此自由零伦等价于任一基点识别下的类为零。最后逐胞腔作上述填充；它们在公共边界均等于 \(f_q\)，所以可粘合为下一相对骨架上的连续映射。
\end{proof}

零胞腔要求目标非空；一胞腔要求它的两个端点像可以用道路相连。这两步是非空性和道路分支的问题。二胞腔开始才出现 \(\pi_1\) 中的边界环路；若这个群非交换，仍可判定单个环路是否为单位元，但尚不能把全部障碍当作普通阿贝尔群系数的上链。

\ProseParagraphBreak

截面的圆盘判据需要把边界数据送到同一纤维。设 \(p:E\to X\) 为 Hurewicz 纤维化，拉回到特征圆盘得到 \(\pi:V=\chi^*E\to D^{q+1}\)。第11章的\cref{g2:lem:fibration-disk-model}给出保持基底的纤维同伦等价
\[
	V\ \underset{P}{\stackrel{Q}{\rightleftarrows}}\ D^{q+1}\times F,
\]
其中 \(F=\pi^{-1}(0)\)。给定的边界截面 \(s_\partial:S^q\to V\) 经 \(Q\) 写成 \(Q s_\partial(z)=(z,u(z))\)。由此得到边界的纤维映射 \(u:S^q\to F\)。

\begin{proposition}[单个胞腔的截面延拓]\label{g2:prop:section-cell-extension}
	设 \(q\ge1\)，\(\pi:V\to D^{q+1}\) 为具有非空道路连通纤维的 Hurewicz 纤维化，\(s_\partial\) 为其边界上的截面。在上述纤维同伦乘积模型中，\(s_\partial\) 可延拓为整个圆盘上的截面，当且仅当 \(u:S^q\to F\) 的类为零。这个零值判据与乘积模型及基点识别的选择无关。
\end{proposition}
\begin{proof}
	若截面 \(s\) 已存在，\(Qs\) 的纤维坐标便延拓 \(u\)，所以其类为零。

	反过来，假定 \(u\) 延拓为 \(U:D^{q+1}\to F\)。复合 \(t(x)=P(x,U(x))\) 是圆盘上的截面；其边界未必等于 \(s_\partial\)，但 \(PQ\) 与恒等映射之间的纤维同伦给出从 \(t|_{S^q}\) 到 \(s_\partial\) 的截面同伦。对 CW 对 \((D^{q+1},S^q)\) 应用\cref{g2:lem:relative-homotopy-lift}，基底同伦取 \(H(x,r)=x\)，初始提升取 \(t\)，边界提升取刚才的同伦。所得提升的终时刻是一个严格延拓 \(s_\partial\) 的截面。这一步把纤维同伦模型给出的边界相等到同伦，修正成实际边界相等。

	判据两端已经等价于与模型无关的截面存在性，因此改变模型不改变它是否为零；基点改变也保持零元。
\end{proof}

纤维丛在圆盘上可用实际平凡化完成同样的论证；保留纤维同伦模型则同时适用于路径纤维化。对一般 Serre 纤维化，相同的圆盘论证可以只提升有限 CW 参数：沿径向收缩提升已给边界，得到中心纤维中的球面映射；若它零伦，就把这张圆盘填充作为初始提升，再反向提升径向同伦，并在边界指定原提升的反向过程。相对提升引理保证最终边界恰为原截面。这给出同一零值判据，而不要求整条纤维作为任意参数空间参与提升。

\begin{example}[连通度给出的自动延拓]\label{g2:exm:connectivity-extension-range}
	设 \(n\ge2\)，\(Y\) 为 \((n-1)\) 连通空间，\((X,A)\) 为 CW 对，\(f_A:A\to Y\) 已给定。选择新顶点的像，并用道路填入新一胞腔；随后 \(\pi_q(Y)=0\)（\(1\le q<n\)）使所有维数不超过 \(n\) 的新胞腔都可填入。因此 \(f_A\) 总可延到 \(X_n\)，首次可能出现的障碍位于 \((n+1)\) 胞腔，取值于 \(\pi_n(Y)\)。对具有 \((n-1)\) 连通非空纤维的纤维化，指定截面的相对延拓具有相同范围。
\end{example}

\subsection{基本群作用产生的系数系统}
\label{g2:subsec:1201:03}

单个胞腔只问一个球面类是否为零，所以任意基点运输都不会改变答案。若要把多个胞腔上的障碍相加，或比较两次选择造成的差，便必须说明这些类处于哪些群、不同群之间怎样识别。第5章的局部系统在这里重新出现，系数从纤维同调换成纤维同伦群。

\begin{definition}[简单空间]\label{g2:def:simple-space}
	道路连通空间 \(Y\) 若满足 \(\pi_1(Y)\) 交换，且它对每个 \(\pi_q(Y)\)（\(q\ge2\)）的基点运输作用都平凡，称为\term[jiandan-kongjian]{简单空间}[simple space]。对整数 \(q\ge1\)，若只要求这些作用在 \(2\le j\le q\) 时平凡，则称为 \(q\) 简单空间。\(q=1\) 时这个条件仅要求基本群交换。
\end{definition}

基本群对自身的作用是共轭，因此定义中的交换性恰好使一次群的基点运输也不受路径选择影响。由\cref{g2:prop:basepoint-action}，两条具有相同端点的道路给出的基点变换相差一个环路作用；对简单空间，这个差别消失。于是其各点的 \(\pi_q\) 可以作唯一的路径无关识别。每个单连通空间都是简单空间；\(S^1\) 也简单，因为其基本群为 \(\mathbb Z\)，高阶同伦群全为零。

\begin{proposition}[纤维同伦群的局部系统]\label{g2:prop:homotopy-local-system}
	设 \(p:E\to X\) 为基底道路连通、纤维非空且简单的 Hurewicz 纤维化。对每个 \(q\ge1\)，纤维同伦群组成阿贝尔群的局部系统 \(\Pi_q(p)\)：在 \(x\in X\) 处取群 \(\pi_q(F_x)\)，沿道路 \(\gamma:x\to y\) 的同态由纤维运输 \(T_\gamma:F_x\to F_y\) 诱导，并用纤维内部的路径识别基点。它只依赖于 \(\gamma\) 的定端点同伦类，且
	\begin{equation}\label{g2:eq:homotopy-local-transport}
		\Pi_q(p)(\gamma*\eta)=\Pi_q(p)(\eta)\Pi_q(p)(\gamma).
	\end{equation}
	若只处理次数至 \(q\)，纤维 \(q\) 简单已足够。
\end{proposition}
\begin{proof}
	第11章的道路提升把 \(T_\gamma\) 确定到自由同伦；反向道路给出同伦逆，定端点道路同伦也给出相应映射同伦。若为 \(F_x,F_y\) 各选一点，\(T_\gamma\) 的基点像未必等于已选终点基点，故还须沿 \(F_y\) 内的一条道路变换基点。简单性说明这样的基点变换与所选道路无关。改变 \(T_\gamma\) 的提升选择时，同伦的基点轨迹也只造成同一类已经消失的环路作用，因而群同态无歧义。

	纤维运输的拼接满足 \(T_{\gamma*\eta}\simeq T_\eta T_\gamma\)，诱导同态遂满足\eqref{g2:eq:homotopy-local-transport}；常值道路诱导恒等映射。这正是\cref{g2:def:local-system}的局部系统公理。\(q\ge2\) 时同伦群本来就交换，\(q=1\) 时交换性由简单性提供。对 Serre 纤维化，可通过第11章的 Hurewicz 替换及纤维弱同伦等价搬回这一系统，或直接对球面参数作上述运输。
\end{proof}

“纤维简单”消除的是同一纤维内的基点歧义，并不要求绕基底一周后的运输恒等。例如取反足映射 \(a:S^2\to S^2\)，其映射环面
\[
	E=(S^2\times I)/((z,1)\sim(a(z),0))\longrightarrow S^1
\]
具有单连通纤维，但基底正向绕行在 \(\pi_2(S^2)=\mathbb Z\) 上作用为乘 \(-1\)。这是因为 Hurewicz 同构把该作用对应到反足映射在二次同调上的度 \(-1\)。因此 \(\Pi_2(p)\) 是反号局部系统，不能以常值整数系数取代。定向 \(S^m\) 丛的情况则相反：其顶次运输保持定向，故 \(\Pi_m(p)\cong\mathbb Z\) 为常值系统；未定向时，它记录的正是纤维定向反转。

\ProseParagraphBreak

普通映射问题中，若 \(Y\) 简单，乘积纤维化 \(X\times Y\to X\) 的运输为恒等，因而系数就是常值群 \(\pi_q(Y)\)。若目标不简单，则已给的低维映射也会进入系数。例如一张定义在 \(X_2\) 上的映射决定
\[
	\rho:\pi_1(X)\longrightarrow\pi_1(Y),
\]
因为 \(X\) 的基本群由二骨架决定。固定 \(\rho\) 后，\(\pi_1(Y)\) 在 \(\pi_q(Y)\) 上的作用可沿 \(\rho\) 拉回，成为 \(X\) 上的局部系统。仅给出抽象群 \(\pi_q(Y)\)，还没有指定这组系数。下一节采用简单纤维而允许基底运输不平凡的条件，既保持阿贝尔上同调的运算，也包括非定向球丛这样的几何情形。

\ProseParagraphBreak

因此，第 \(q+1\) 层的障碍将由相对胞腔上链记录在
\[
	C^{q+1}(X,A;\Pi_q(p))
\]
中。每个胞腔的定向决定球面类的符号，道路运输决定不同胞腔的系数如何相加。下一节要证明这些值满足余圈条件，并确定修改已有延拓时产生的余边界；只有完成这两步，才能把逐胞腔判据提升为上同调中的不变量。映射与提升问题的阅读对照可参看\cite[第4.3节，Obstruction Theory]{Hatcher2002}及\cite[第18章，第5节，Obstruction theory]{May1999}。
% 纲要条目排为编号小节标题，并进入目录与书签；说明段仍保留于纲要注释中。

% !TeX root = ../../Book2.tex
% 纲要标题：### 12.2 障碍余圈及其不变量
% 纲要位置：几何学计划纲要.md，第 1478 行。
% 正文状态：初稿；尚未正式审核。

% 纲要内容开始
% 
% * 障碍上链的构造
% * 余圈条件的证明
% * 改变初始延拓时的余边界
% * 障碍类消失与延拓存在
% * 差异类与上同调可表性
% 
% 纲要内容结束

\section{障碍余圈及其不变量}
\label{g2:sec:1202}

一个边界球面类判断一个胞腔能否填入。对整个空间，胞腔的边界之间还有相互关系，而重新选择已经填好的部分也会改变下一步的障碍。上链记录逐胞腔的值，余圈条件记录这些值之间的关系，除去余边界则记录允许改变上一层后的不变量。本节逐一证明这三件事。

\subsection{障碍上链的构造}
\label{g2:subsec:1202:01}

固定道路连通 CW 对 \((X,A)\)、Hurewicz 纤维化 \(p:E\to X\)，以及 \(A\) 上的指定截面 \(s_A\)。假定纤维非空且简单，记上一节构造的局部系统为 \(\Pi_n(p)\)。对某个 \(n\ge1\)，设已经选定延拓 \(s:X_n=A\cup X^n\to E\)。

\ProseParagraphBreak

给每个相对 \((n+1)\) 胞腔选定向及特征映射 \(\chi_e:D^{n+1}\to X_{n+1}\)。沿圆盘的径向收缩把边界截面送到中心纤维，便得到一个球面映射 \(a_e:S^n\to F_{\chi_e(0)}\)。源球面取圆盘的正向边界；它的同伦类就是这个胞腔上的延拓障碍。

\begin{definition}[障碍上链]\label{g2:def:obstruction-cochain}
	在上述条件下，将每个有向相对 \((n+1)\) 胞腔 \(e\) 赋值为 \([a_e]\in\pi_n(F_{\chi_e(0)})\)，得到\term[zhangai-shanglian]{障碍上链}[obstruction cochain]
	\[
		c_{n+1}(s)\in C^{n+1}(X,A;\Pi_n(p)),\qquad c_{n+1}(s)(e)=[a_e].
	\]
	在 \(A\) 内的胞腔上取零。若通过道路把系数搬到选定的基点纤维，则同时按第5章的局部系数约定记录运输。
\end{definition}

这一赋值包含三项选择，但它们的作用可分别追踪。反转胞腔定向使球面类变号；改变圆盘上的径向提升，只产生纤维内同伦；改变连接基点的道路，则按 \(\Pi_n(p)\) 的运输变换系数。纤维简单性保证同一纤维内的基点道路不再产生歧义。因此得到的是一个确定的局部系数上链，而不是把不同纤维中的元素直接相加。

\ProseParagraphBreak

由\cref{g2:prop:section-cell-extension}，保持 \(s\) 在 \(X_n\) 上的全部值时，延拓到 \(X_{n+1}\) 的充要条件是
\begin{equation}\label{g2:eq:obstruction-cochain-vanishing}
	c_{n+1}(s)=0.
\end{equation}
这是上链逐项为零的条件。它与稍后出现的上同调类为零有不同含义。

\begin{example}[带一个关系的二维复形]\label{g2:exm:obstruction-relation-cochain}
	令 \(X=S^1\cup_m e^2\)，二胞腔沿圆周的 \(m\) 次绕行附着，\(m\ge2\)。给定一次骨架映射 \(f:S^1\to S^1\)，其次数为 \(k\)。把它看成乘积纤维化的截面，则
	\[
		c_2(f)(e^2)=mk\in\pi_1(S^1)=\mathbb Z.
	\]
	因为 \(m\ne0\)，保持这个 \(f\) 延拓的条件是 \(k=0\)。另一方面，胞腔上链的余边界 \(C^1(X;\mathbb Z)\to C^2(X;\mathbb Z)\) 是乘 \(m\)，所以 \([c_2(f)]=0\) 对所有 \(k\) 都成立。允许改变一胞腔上的映射后，确实总能把它换为常值再延拓。这个例子先把两种零值判据的差别显示出来。
\end{example}

\subsection{余圈条件的证明}
\label{g2:subsec:1202:02}

虽然 \(s\) 尚未定义在 \((n+1)\) 胞腔内部，障碍值仍须满足 \((n+2)\) 胞腔给出的关系。关键是：整个特征圆盘可缩，它的拉回纤维化有截面；这个截面未必与已有边界数据相同，但足以证明障碍值在胞腔边界上的总和为零。

\begin{lemma}[障碍余圈]\label{g2:lem:obstruction-cocycle}
	设 \((X,A)\) 是道路连通 CW 对，\(p:E\to X\) 为纤维非空且简单的 Hurewicz 纤维化。对 \(n\ge1\) 及截面 \(s:X_n\to E\)，有
	\[
		\delta c_{n+1}(s)=0\quad\text{于 }C^{n+2}(X,A;\Pi_n(p)).
	\]
\end{lemma}
\begin{proof}
	先把逐胞腔的定义扩大到相对圆盘。给定带基点映射 \(\Phi:(D^{n+1},S^n)\to(X_{n+1},X_n)\)，用 \(s\Phi\) 作为边界提升，沿 \(\Phi\) 给出的径向同伦提升到一个纤维，记所得类为 \(\omega_s[\Phi]\)。相对同伦的提升说明它只依赖相对同伦类；将两个圆盘沿加法方向拼接，再提升同一拼接，说明 \(\omega_s\) 是群同态。若 \(\Phi\) 本身已有一个与指定边界相同的提升，则所得球面可填，\(\omega_s[\Phi]=0\)。

	为了同时处理运输，拉回到 \(X\) 的万有覆盖 \(\widetilde X\)。局部系统在那里成为常值系统；对每个骨架记其全部逆像为 \(\widetilde X_r\)。当 \(r\ge2\) 时，\(X_r\hookrightarrow X\) 在基本群上为同构，所以 \(\widetilde X_r\) 单连通。对 \(n\ge2\)，空间对 \((\widetilde X_{n+1},\widetilde X_n)\) 只附着 \((n+1)\) 胞腔，故为 \(n\) 连通；\cref{g2:thm:hurewicz-relative}给出
	\[
		\pi_{n+1}(\widetilde X_{n+1},\widetilde X_n)\xrightarrow{\cong}
		H_{n+1}(\widetilde X_{n+1},\widetilde X_n).
	\]
	特征圆盘在这张同构下成为相应的有向胞腔。因此 \(\omega_s\) 在相对同调上的表示，恰好是障碍上链的求值，而不是另选的一张同态。

	现在取一个 \((n+2)\) 胞腔的特征映射 \(g:D^{n+2}\to\widetilde X_{n+2}\)，并在边界上选基点。其胞腔边界由以下路径得到：先取 \(g|_{S^{n+1}}\) 的球面类，再把它送入 \(\pi_{n+1}(\widetilde X_{n+1},\widetilde X_n)\)，最后用上述 Hurewicz 同构。这与同调边界一致，因为 Hurewicz 映射与边界及包含相容，见\cref{g2:prop:hurewicz-naturality}。

	把纤维化拉回 \(D^{n+2}\)，可选一个在边界基点处等于 \(s\) 的截面。它给出 \(g|_{S^{n+1}}\) 的整体提升。将这个球面写成 \(D^{n+1}/S^n\) 后，所得相对圆盘在其边界恒取选定基点，提升在该边界也恒等于 \(s\) 的值。因此这一相对类被 \(\omega_s\) 送到零。换言之，\(c_{n+1}(s)\) 在任一 \((n+2)\) 胞腔的胞腔边界上取零。

	\(n=1\) 时，\(\widetilde X_1\) 未必单连通，需要把上述相对同伦到胞腔链的步骤直接说明。对二维相对圆盘，将它经过的二维胞腔表示为有向小圆盘，用细弧连接到边界基点；剩余部分映到 \(\widetilde X_1\)。沿细弧切开后，每个小圆盘贡献其胞腔障碍的正值或负值，剩余部分已有截面，贡献为零。重新粘合时，边上的道路成对反向出现；改变连接弧只把相应纤维环路共轭。由于 \(\pi_1(F)\) 交换，共轭消失，所得总值只依赖各二维胞腔的有向次数。这就证明 \(\omega_s\) 经由 \(H_2(\widetilde X_2,\widetilde X_1)\) 分解。上述小圆盘表示来自相对胞腔逼近及圆盘映射的次数分解，与第3章计算胞腔关联数所用的分解相同。

	对三胞腔的附着球面使用这一分解，其每条内部边出现两次、方向相反；也可沿前一段的圆盘提升，直接看出相应相对类的 \(\omega_s\) 为零。因此 \(n=1\) 同样得到 \(\delta c_2(s)=0\)。最后把覆盖中的计算降回 \(X\)。改变提升胞腔会同时改变其中心纤维的识别，恰好插入 \(\Pi_n(p)\) 的运输；第5章的局部系数胞腔余边界正是这一等变求值。因此结论对原来的局部系数成立。\(A\) 内的胞腔原已有截面，所有相应值为零，故上述论证也保留相对条件。
\end{proof}

由引理可定义
\begin{equation}\label{g2:eq:obstruction-class}
	o_{n+1}(s)=[c_{n+1}(s)]\in H^{n+1}(X,A;\Pi_n(p)).
\end{equation}
这称为当前截面的\term[zhangai-lei]{障碍类}[obstruction class]。上同调群描述障碍可能出现的位置；\(o_{n+1}(s)\) 是当前数据给出的具体元素。一个非零群完全可能包含等于零的障碍类。

\subsection{改变初始延拓时的余边界}
\label{g2:subsec:1202:03}

设 \(s,t\) 都定义在 \(X_n\)，并在 \(X_{n-1}\) 上相同。对一个 \(n\) 胞腔，将 \(t\) 的圆盘取正向、\(s\) 的圆盘取负向，沿共同边界粘成 \(S^n\)，再利用圆盘上的纤维同伦模型送入中心纤维。所得类记为
\[
	D_n(s,t)(e^n)\in\Pi_n(p)_{\chi_e(0)}.
\]
这些值组成\term[chayi-shanglian]{差异上链}[difference cochain] \(D_n(s,t)\in C^n(X,A;\Pi_n(p))\)。它记录把 \(s\) 改成 \(t\) 时，在每个 \(n\) 胞腔中实际加入的球面类。

\ProseParagraphBreak

若 \(s,t\) 只在 \(X_{n-1}\) 上由指定截面同伦 \(h\) 相连，则用 \(h\) 填侧面，将上下两张圆盘联成柱体的整个边界，同样得到 \(D_n(s,t;h)\)。为固定符号，柱体取先胞腔、后区间的乘积定向。若 \(H\) 是两端用 \(s,t\)、侧面用 \(h\) 的部分截面，则
\begin{equation}\label{g2:eq:difference-prism-sign}
	D_n(s,t;h)(e^n)=(-1)^n c_{n+1}(H)(e^n\times I).
\end{equation}
因为 \(\partial(e^n\times I)=\partial e^n\times I+(-1)^n(e^n\times\{1\}-e^n\times\{0\})\)，这一约定恰好使上端为 \(t\)、下端为 \(-s\)。

\begin{proposition}[差异与余边界]\label{g2:prop:obstruction-difference}
	设 \((X,A)\) 为道路连通 CW 对，\(p:E\to X\) 为纤维非空且简单的 Hurewicz 纤维化，\(n\ge1\)。若 \(s,t:X_n\to E\) 是在 \(A\) 上相同的截面，\(h\) 是它们在 \(X_{n-1}\) 上固定 \(A\) 的截面同伦，则
	\begin{equation}\label{g2:eq:obstruction-difference}
		\delta D_n(s,t;h)=c_{n+1}(t)-c_{n+1}(s).
	\end{equation}
	此外，给定 \(s\) 和任意 \(b\in C^n(X,A;\Pi_n(p))\)，都能选到在 \(X_{n-1}\) 上与 \(s\) 相同的 \(t\)，使 \(D_n(s,t)=b\)。
\end{proposition}
\begin{proof}
	令 \(W=X_{n+1}\)，在 \(W\times I\) 上拉回纤维化。部分截面 \(H\) 已在 \(X_n\times\partial I\)、\(X_{n-1}\times I\) 和 \(A\times I\) 上定义，这包含相对柱体的 \(n\) 骨架。由上一引理，其 \((n+1)\) 次障碍在每个 \((n+2)\) 胞腔 \(e^{n+1}\times I\) 的边界上取零。因此
	\[
		0=\sum_{e^n}[e^{n+1}:e^n]c_{n+1}(H)(e^n\times I)
		+(-1)^{n+1}\bigl(c_{n+1}(t)-c_{n+1}(s)\bigr)(e^{n+1}).
	\]
	有局部系数时，求和中同时沿各附着道路运输；两端的系数用区间方向的恒等运输识别。代入式\eqref{g2:eq:difference-prism-sign}并乘 \((-1)^n\)，便得到所需公式。

	为实现指定的 \(b\)，在每个相对 \(n\) 胞腔内部取一个小圆盘。在纤维模型中，把该圆盘的边界压成一点，插入代表 \(b(e^n)\) 的球面映射；原来通向小圆盘中心的值通过一条纤维道路与球面基点连接。等价地，这是保持外边界的挤压映射 \(D^n\to D^n\vee S^n\)，在第一因子沿用原映射，在第二因子加入所选球面类。简单性使连接道路不影响所得类，方向取为 \(t-s\)，于是差异正是 \(b(e^n)\)。对 \(n=1\)，此操作就是在原道路中插入指定环路。

	纤维同伦模型转换回原丛后，若外边界只相等到纤维同伦，按\cref{g2:prop:section-cell-extension}证明中的相对提升把它修正为实际相等。各胞腔内部的修改相互独立，在 \(X_{n-1}\) 及 \(A\) 上全不改变，故由 CW 弱拓扑粘合成 \(t\)。它满足 \(D_n(s,t)=b\)。
\end{proof}

因此，保持 \(X_{n-1}\) 上的数据时，\(c_{n+1}(s)\) 恰好可以改变一个任意余边界。障碍类正是这组修改下保持不变的量。普通映射的对应叙述见\cite[第18章第5节，Obstruction theory]{May1999}；本节的局部运输与柱体公式给出了截面版本的同一机制。

\subsection{障碍类消失与延拓存在}
\label{g2:subsec:1202:04}

现在可以把“修改后填入”的几何要求写成一个上同调判据。

\begin{theorem}[逐层延拓判据]\label{g2:thm:obstruction-extension}
	设 \((X,A)\) 为道路连通 CW 对，\(p:E\to X\) 为纤维非空且简单的 Hurewicz 纤维化，\(n\ge1\)，且 \(s:X_n\to E\) 为截面。则以下两条件等价：
	\begin{enumerate}
		\item \(o_{n+1}(s)=0\in H^{n+1}(X,A;\Pi_n(p))\)；
		\item \(s|_{X_{n-1}}\) 可保持不变地延拓为 \(X_{n+1}\) 上的截面。
	\end{enumerate}
	若要求保持原来的整个 \(s|_{X_n}\)，充要条件是上链 \(c_{n+1}(s)=0\)。
\end{theorem}
\begin{proof}
	若 \(o_{n+1}(s)=0\)，选 \(b\in C^n(X,A;\Pi_n(p))\) 使 \(\delta b=-c_{n+1}(s)\)。差异命题给出与 \(s\) 在 \(X_{n-1}\) 上相同的 \(t\)，满足
	\[
		c_{n+1}(t)=c_{n+1}(s)+\delta b=0.
	\]
	单胞腔判据把 \(t\) 延拓到 \(X_{n+1}\)。反过来，若有这样的延拓，记它在 \(X_n\) 上的限制为 \(t\)，则 \(c_{n+1}(t)=0\)，而差异公式使 \(c_{n+1}(s)=-\delta D_n(s,t)\)，故障碍类为零。最后一句直接是\cref{g2:eq:obstruction-cochain-vanishing}。
\end{proof}

\begin{theorem}[首障碍]\label{g2:thm:primary-obstruction}
	设 \((X,A)\) 为道路连通 CW 对，\(p:E\to X\) 为纤维非空、简单且 \((n-1)\) 连通的 Hurewicz 纤维化，\(n\ge1\)，并给定截面 \(s_A\)。则 \(s_A\) 总可延到 \(X_n\)，且
	\[
		o_{n+1}(p,s_A)\in H^{n+1}(X,A;\Pi_n(p))
	\]
	与所选 \(X_n\) 上的延拓无关。它为零当且仅当 \(s_A\) 可延到 \(X_{n+1}\)。这个\term[shouzhangai]{首障碍}[primary obstruction]对保持给定截面的 CW 对拉回自然。若相对胞腔维数均不超过 \(n+1\)，它为零便等价于全空间上的截面存在。
\end{theorem}
\begin{proof}
	非空性允许选择新顶点的值，道路连通性允许填入新一胞腔；\(\pi_j(F)=0\)（\(0<j<n\)）再使维数至 \(n\) 的胞腔都可以填入。故存在所需初始延拓。

	取两张这样的截面 \(s,t\)。在 \(X_{n-1}\) 上构造固定 \(A\) 的截面同伦，只需处理 \(k\) 胞腔乘区间，\(k\le n-1\)；其边界问题取值于 \(\pi_k(F)\)。当 \(k=0\) 时由道路连通性填入，其余均由低阶同伦群消失填入。于是得到前一命题所需的 \(h\)，差异公式证明两张障碍上圈同属一个上同调类。

	若这个类为零，逐层延拓判据给出延到 \(X_{n+1}\) 的截面；反之，取实际延拓的限制来计算，障碍上链本身为零。对胞腔映射的拉回，在每个相对胞腔上，边界球面复合及运输与原构造相容；将复合映射分解成有向胞腔贡献，就得到局部系数上链上的拉回公式。一般映射先相对胞腔逼近，拉回纤维化在基底同伦下纤维同伦等价，故得到相同的障碍类。这证明自然性。最后，在所述维数范围内 \(X_{n+1}=X\)，所以无需再处理更高胞腔。
\end{proof}

纤维较低同伦群的消失在这里承担两个不同任务：它使初始截面存在，也使更低骨架上的初始选择同伦唯一。缺少第二项时，第 \(n+1\) 次障碍类仍对固定的低层数据有意义，却可能随着更低层的选择改变。

\begin{corollary}[由障碍群消失得到延拓]\label{g2:cor:obstruction-groups-vanish}
	设 \((X,A)\) 为道路连通 CW 对，\(p:E\to X\) 为纤维非空且简单的 Hurewicz 纤维化。若
	\[
		H^{q+1}(X,A;\Pi_q(p))=0\qquad(q\ge1),
	\]
	则每个 \(A\) 上的截面都可延拓到 \(X\)。对有限维相对 CW 对，只需核对实际出现的次数。
\end{corollary}
\begin{proof}
	先延拓到 \(X_1\)。在第 \(q\) 步，障碍群为零使当前障碍类为零，所以可以保持 \(X_{q-1}\) 不变，修改 \(q\) 胞腔上的值，再延拓到 \(X_{q+1}\)。一个固定的 \(k\) 胞腔只会在第 \(k\) 步最后修改一次，以后所有步骤都保留它。取这些最终值，便在每个特征圆盘上得到连续且边界相容的截面；CW 弱拓扑使它们定义整个 \(X\) 上的连续截面。
\end{proof}

\subsection{差异类与上同调可表性}
\label{g2:subsec:1202:05}

截面存在以后，还可以问它是否同伦唯一。这一次要填的是两张截面之间的柱体，因此障碍所在的上同调次数降低一。

\begin{proposition}[同伦的差异类]\label{g2:prop:section-difference-class}
	设 \((X,A)\)、\(p:E\to X\) 满足上一推论的空间与纤维条件，\(s,t\) 为在 \(A\) 上相同的全局截面。给定它们在 \(X_{n-1}\) 上固定 \(A\) 的截面同伦 \(h\)，\(n\ge1\)，则 \(D_n(s,t;h)\) 为余圈。其类
	\[
		\Delta_n(s,t;h)\in H^n(X,A;\Pi_n(p))
	\]
	为零，当且仅当 \(h|_{X_{n-2}}\) 可延拓为 \(s|_{X_n}\) 与 \(t|_{X_n}\) 之间的截面同伦。\(n=1\) 时 \(X_{-1}=A\)，所以只固定 \(A\) 上的同伦。
\end{proposition}
\begin{proof}
	两张全局截面的障碍上链均为零，所以\cref{g2:eq:obstruction-difference}给出 \(\delta D_n=0\)。在空间对
	\[
		\bigl(X\times I,\ X\times\partial I\cup A\times I\bigr)
	\]
	上考虑给定的两端与部分同伦。相对胞腔为 \(e^k\times(0,1)\)，维数为 \(k+1\)。相对链的边界去掉两端后只有 \(\partial e^k\times I\)，所以其局部系数上同调自然识别为 \(H^{j-1}(X,A;\Pi_n(p))\)。这就是柱体相对方向使次数降低一的原因。

	用式\eqref{g2:eq:difference-prism-sign}识别同伦障碍，逐层延拓判据便说明：其类为零，恰好允许修改内部的 \(n\) 维胞腔，即 \(e^{n-1}\times I\) 上的同伦，同时固定更低柱体与全部两端，再填入 \(e^n\times I\)。这正是结论所述的延拓，而不是要求保留原来的整个 \(h|_{X_{n-1}}\)。
\end{proof}

如果纤维 \((n-1)\) 连通，上述 \(h\) 在所需低次中总能找到。其不同选择在更低的柱体上又可同伦，因为相关同伦群仍在消失范围内；差异公式于是使 \(\Delta_n\) 的上同调类与这些低次选择无关。接下来两种情形中，这个差异类已经记录全部自由度。

\begin{theorem}[截面同伦类的参数]\label{g2:thm:section-torsor}
	设 \((X,A)\) 为道路连通 CW 对，\(p:E\to X\) 为具有非空、简单、\((n-1)\) 连通纤维的 Hurewicz 纤维化，\(n\ge1\)，并假设指定截面 \(s_A\) 至少有一个全局延拓。以下任一条件成立时，延拓的相对截面同伦类形成 \(H^n(X,A;\Pi_n(p))\) 上的主齐性集：
	\begin{enumerate}
		\item 相对胞腔维数均不超过 \(n\)；
		\item 纤维的全部正阶同伦群除 \(\pi_n\) 外均为零。
	\end{enumerate}
	选定一个延拓 \(s_0\) 后，\([t]\mapsto\Delta_n(s_0,t)\) 给出与该上同调群的双射。
\end{theorem}
\begin{proof}
	低阶同伦群消失使各截面在 \(X_{n-1}\) 上可由相对同伦相连。选定这样的同伦后，逐胞腔球面的拼接给出
	\[
		\Delta_n(s,t)+\Delta_n(t,u)=\Delta_n(s,u),\qquad
		\Delta_n(s,s)=0.
	\]
	改变辅助低次同伦只加一个余边界：将两次同伦再放到一个参数区间上比较，在维数低于 \(n\) 的地方由纤维连通度填充，剩余柱体边界正是差异公式。因而这些等式确实是上同调中的等式。

	若 \(\Delta_n(s,t)=0\)，上一命题给出 \(X_n\) 上的相对同伦。条件 I 下已经完成全空间；条件 II 下，继续填入更高的 \(e^k\times I\) 所遇系数 \(\pi_k(F)\) 均为零，所以同伦延拓到全部 \(X\)。反过来，全局同伦当然使每一级差异障碍为零。这证明所列映射单射。

	给定任意上同调类，取余圈 \(b\in C^n(X,A;\Pi_n(p))\) 表示它。在 \(s_0\) 的各 \(n\) 胞腔中插入 \(b\) 的球面代表，得到 \(t_n\)，并有 \(D_n(s_0,t_n)=b\)。条件 I 下它已经是全局截面；条件 II 下，由 \(\delta b=0\) 及 \(c_{n+1}(s_0)=0\)，得到 \(c_{n+1}(t_n)=0\)，先延到 \(X_{n+1}\)，再利用更高同伦群消失继续延拓。所得截面给出指定的差异类，证明满射。

	因此上同调群通过加入差异类自由且传递地作用于延拓同伦类。这样的集合称为\term[zhuqixingji]{主齐性集}[torsor]：任意两个元素之间有唯一的群元素作为差，但在选择 \(s_0\) 前没有一个自然的零元。
\end{proof}

\begin{theorem}[上同调的可表性]\label{g2:thm:em-representability}
	设 \(G\) 是阿贝尔群，\(n\ge1\)，\((X,A)\) 为 CW 对。取带基点 CW 模型 \(K(G,n)\)，并令 \(\iota_n\in H^n(K(G,n);G)\) 为在 \(\pi_n\cong G\) 上对应恒等映射的类。则有自然双射
	\[
		[(X,A),(K(G,n),*)]\xrightarrow{\cong}H^n(X,A;G),\qquad [f]\longmapsto f^*\iota_n.
	\]
	左端是把 \(A\) 映到基点的映射在固定 \(A\) 的同伦下的类；无相对条件时得到 \([X,K(G,n)]\cong H^n(X;G)\)。
\end{theorem}
\begin{proof}
	第9章已经构造了 \(K(G,n)\)。\(n\ge2\) 时，Hurewicz 定理给出 \(H_n(K;\mathbb Z)\cong G\) 以及低正次同调消失，故万有系数定理把 \(H^n(K;G)\) 识别为 \(\operatorname{Hom}(G,G)\)。\(n=1\) 时，用一次 Hurewicz 定理及 \(G\) 交换得到相同识别；\(H_0=\mathbb Z\) 不产生 Ext 项。因此 \(\iota_n\) 定义明确。

	对乘积纤维化 \(X\times K(G,n)\to X\)，常值映射是指定边界值的一个延拓。纤维为 \((n-1)\) 连通简单空间，且只剩一个正阶同伦群，所以上一定理的条件 II 给出由 \(H^n(X,A;G)\) 参数化的全部映射同伦类。\(X\) 不连通时逐分支应用即可。

	还要验证所得参数确实是 \(f^*\iota_n\)。先把 \(f\) 在 \(X_{n-1}\) 上同伦为常值，并由同伦延拓性质将这段同伦延到全空间。对每个相对 \(n\) 胞腔，此时 \(f\) 压去边界后给出 \(S^n\to K(G,n)\)。差异上链在该胞腔上的值就是此球面代表的 \(G\) 中元素；而 \(\iota_n\) 对这个球面基本类的求值也正是该元素，因为它对应 \(G\) 上的恒等映射。因此两个胞腔上圈相同，给出所列拉回公式。公式与空间对映射的复合相容，所以双射自然。
\end{proof}

这一结论把一个抽象上同调类变成一张实际映射。下一节正要利用它：把一个阶段的粘合类表示为到 \(K(G,n+1)\) 的映射，再取该映射的同伦纤维，便能把相应的 \(G\) 放回所需次数。关于可表性、Postnikov 阶段与障碍问题之间的关系，可参看\cite[第4.3节，Connections with Cohomology]{Hatcher2002}。
% 纲要条目排为编号小节标题，并进入目录与书签；说明段仍保留于纲要注释中。

% !TeX root = ../../Book2.tex
% 纲要标题：### 12.3 Postnikov 塔的构造
% 纲要位置：几何学计划纲要.md，第 1486 行。
% 正文状态：初稿；尚未正式审核。

% 纲要内容开始
% 
% * 逐次消去高阶同伦群
% * Eilenberg–MacLane 纤维
% * k-不变量和阶段之间的粘合
% * 同伦群相同的不同三型
% 
% 纲要内容结束

\section{Postnikov 塔的构造}
\label{g2:sec:1203}

前两节把映射的构造分成逐个胞腔的延拓。本节从另一端分解问题：保留目标空间的低阶同伦群，逐次舍去高阶同伦群。第11章已用这样的截段证明有限生成性；现在要说明相邻截段怎样接在一起，以及这部分接合信息为何不能从同伦群的清单中读出。

\subsection{逐次消去高阶同伦群}
\label{g2:subsec:1203:01}

设 \(X\) 为单连通带基点 CW 复形。\cref{g2:lem:finite-homotopy-truncations}构造了 \(X\hookrightarrow X[n]\)：依次附着 \((n+2),(n+3),\ldots\) 维胞腔，消去 \(\pi_{n+1},\pi_{n+2},\ldots\)，同时保留 \(\pi_i\)（\(i\le n\)）。注意这与取 \(n\) 骨架的方向相反。取骨架是删去高维胞腔，可能重新产生高阶同伦；这里则添加高维胞腔，为原来不能填充的球面提供填充。

\begin{definition}[Postnikov 塔]\label{g2:def:postnikov-tower}
	设 \(X\) 为道路连通带基点空间。一列带基点映射
	\[
		\cdots\xrightarrow{p_{n+1}}P_nX\xrightarrow{p_n}P_{n-1}X\longrightarrow\cdots\longrightarrow P_1X
	\]
	连同相容映射 \(a_n:X\to P_nX\)，若满足 \(p_na_n=a_{n-1}\)、\((a_n)_*:\pi_i(X)\xrightarrow{\cong}\pi_i(P_nX)\)（\(i\le n\)），以及 \(\pi_i(P_nX)=0\)（\(i>n\)），就称为 \(X\) 的\term[Postnikov ta]{Postnikov 塔}[Postnikov tower]。空间 \(P_nX\) 称为第 \(n\) 阶段；高于 \(n\) 的同伦群全为零的空间也称为一个\term[n xing]{\(n\)-型}[\(n\)-type]。
\end{definition}

Postnikov 塔\footnote{Mikhail Mikhailovich Postnikov（1927—2004），苏联、俄罗斯数学家，主要研究代数拓扑与同伦理论，长期任教于莫斯科大学。}的下标记录保留的最高同伦次数。对单连通空间，可取 \(P_1X=\{*\}\)。第一个有可能非平凡的阶段是 \(P_2X\)，它只保留 \(\pi_2(X)\)。\symidx{\(P_nX\)：空间的第 \(n\) 个 Postnikov 阶段}

\begin{theorem}[Postnikov 塔的存在]\label{g2:thm:postnikov-tower-existence}
	每个单连通带基点 CW 复形 \(X\) 都有 Postnikov 塔，且可使每张 \(p_n\)（\(n\ge2\)）为 Hurewicz 纤维化，各阶段具有 CW 同伦型。
\end{theorem}
\begin{proof}
	先用 \(X[n]\) 构造相容的 CW 阶段。包含 \(X\to X[n-1]\) 可延拓为 \(q_n:X[n]\to X[n-1]\)：在一个新加的 \(d\) 胞腔上，\(d\ge n+2\)，延拓障碍属于 \(\pi_{d-1}(X[n-1])=0\)。逐胞腔作延拓，并用相对胞腔逼近保持 \(X\) 上的原映射，就得到 \(q_n|_X=\operatorname{id}_X\) 后接包含。各 \(X[n]\) 已有定义所需的同伦群。

	再由下往上作路径纤维化替换。取 \(P_1X=\{*\}\)，并令 \(j_1:X[1]\to P_1X\) 为同伦等价。假定已有同伦等价 \(j_{n-1}:X[n-1]\to P_{n-1}X\)。把复合 \(f_n=j_{n-1}q_n\) 替换为
	\[
		P_nX=\{(z,\gamma):z\in X[n],\ \gamma:I\to P_{n-1}X,\ \gamma(0)=f_n(z)\},
		\qquad p_n(z,\gamma)=\gamma(1).
	\]
	\cref{g2:prop:fibration-replacement}说明 \(p_n\) 是 Hurewicz 纤维化，而 \(j_n(z)=(z,\text{常值道路 }f_n(z))\) 是同伦等价。令 \(a_n=j_n|_X\)，便有严格等式
	\[
		p_na_n=j_{n-1}q_n|_X=j_{n-1}|_X=a_{n-1}.
	\]
	所以替换没有破坏塔的相容性，也没有改变任何阶段的同伦群。
\end{proof}

这项构造有许多选择，但阶段作为 \(X\) 的截段具有确定的同伦意义。若 \(Y\) 是带基点 \(n\)-型，则限制给出双射
\begin{equation}\label{g2:eq:postnikov-truncation-universal}
	[X[n],Y]_*\xrightarrow{\cong}[X,Y]_*.
\end{equation}
存在性来自新增胞腔的边界映入 \(Y\) 后总能填充，因为边界维数至少为 \(n+1\)。唯一性则把同一延拓问题用于圆柱：一个相对 \(d\) 胞腔对应的同伦延拓障碍在 \(\pi_d(Y)\) 中，仍为零。无限多个胞腔的延拓按弱拓扑粘合，得到全局映射与同伦。因此两个截段之间可延拓 \(X\) 上的相同映射，得到互逆到同伦的比较；这也说明将 \(X[n]\) 换成同伦等价的 \(P_nX\) 只改变模型。

\ProseParagraphBreak

特别地，\(P_2X\) 是 \(K(\pi_2(X),2)\) 的同伦型；如果 \(X\) 已经是 \(N\)-型，则 \(a_N:X\to P_NX\) 在全部同伦群上为同构，CW 模型之间的 Whitehead 定理便使它成为同伦等价。阶段保留的是整个低阶同伦结构，而不仅是某几个群的抽象副本。

\subsection{Eilenberg–MacLane 纤维}
\label{g2:subsec:1203:02}

相邻阶段只相差一个同伦群，纤维正好把这个差别单独装出来。

\begin{proposition}[阶段之间的纤维]\label{g2:prop:postnikov-stage-fiber}
	设 \(X\) 为单连通带基点 CW 复形，取上一构造的纤维化塔。对 \(n\ge2\)，\(p_n:P_nX\to P_{n-1}X\) 的纤维 \(F_n\) 道路连通，并满足
	\[
		\pi_i(F_n)=0\quad(i\ne n),\qquad \pi_n(F_n)\cong\pi_n(X).
	\]
	因此它的 CW 逼近为 \(K(\pi_n(X),n)\)，阶段纤维序列可记为
	\begin{equation}\label{g2:eq:postnikov-stage-fibration}
		K(\pi_n(X),n)\longrightarrow P_nX\xrightarrow{p_n}P_{n-1}X.
	\end{equation}
\end{proposition}
\begin{proof}
	总空间与基底均单连通，长正合列的末段先给出纤维道路连通。对 \(i<n\)，映射 \((p_n)_*\) 是由两张 \(a_n\) 给出的同构；在 \(i=n\) 时，基底的 \(\pi_n\) 与 \(\pi_{n+1}\) 均为零。因此长正合列给出
	\[
		\pi_i(F_n)=0\ (i<n),\qquad \pi_n(F_n)\xrightarrow{\cong}\pi_n(P_nX)=\pi_n(X).
	\]
	对 \(i>n\)，两阶段对应的同伦群全为零，纤维群也为零。CW 逼近保留这些群；第9章的 Eilenberg–MacLane 定义及第11章对其模型唯一性的证明给出最后的识别。
\end{proof}

式\eqref{g2:eq:postnikov-stage-fibration}并没有说 \(P_nX\) 是基底与纤维的乘积。此前的 Hopf 丛已经表明，同样的基底与纤维可能产生不同的总空间。这里纤维只剩一个同伦群，使粘合信息能够由一个上同调类记录，但仍需要把这个类找出来。

\ProseParagraphBreak

整座纤维化塔还保留了 \(X\) 的弱同伦型。说明这一点有助于理解为什么逐阶研究最终足以恢复全部同伦信息。

\begin{proposition}[从塔恢复弱同伦型]\label{g2:prop:postnikov-limit}
	对\cref{g2:thm:postnikov-tower-existence}中的纤维化塔，逆极限
	\[
		P_\infty X=\{(z_1,z_2,\ldots)\in\textstyle\prod_nP_nX:p_n(z_n)=z_{n-1}\}
	\]
	带子空间拓扑。相容映射给出 \(a:X\to P_\infty X\)，它是弱同伦等价。
\end{proposition}
\begin{proof}
	固定 \(i\ge1\)。同伦群系统 \(\pi_i(P_nX)\) 从 \(n\ge i\) 起稳定为 \(\pi_i(X)\)。先证明
	\[
		\pi_i(P_\infty X)\longrightarrow\varprojlim_n\pi_i(P_nX)
	\]
	满射。给定右边相容的同伦类，逐层选带基点代表 \(f_n:S^i\to P_nX\)。已有 \(f_{n-1}\) 后，\(p_nf_n\) 与 \(f_{n-1}\) 带基点同伦；相对基点提升这段同伦，便能调整 \(f_n\) 使 \(p_nf_n=f_{n-1}\) 严格成立。所有坐标合在一起给出到逆极限的连续映射。

	再设 \(f:S^i\to P_\infty X\) 在每个坐标上都零伦。取足够大的 \(N\)，使 \(\pi_{i+1}(P_nX)\to\pi_{i+1}(P_{n-1}X)\) 在 \(n>N\) 时全为同构。先选 \(f_N\) 的零伦 \(G_N:D^{i+1}\to P_NX\)。假定已选 \(G_{n-1}\)，任选 \(f_n\) 的零伦 \(G_n\)。映射 \(p_nG_n\) 与 \(G_{n-1}\) 在边界相同，将它们沿边界反向粘合，得到 \(P_{n-1}X\) 中的一个 \((i+1)\) 球面类。上述同伦群满射允许从 \(P_nX\) 中提升这个类；在 \(G_n\) 的圆盘内部接入相反的球面代表，即可使两张圆盘映射相对边界同伦。相对边界的同伦提升又将其改成严格等式 \(p_nG_n=G_{n-1}\)。递推所得相容零伦给出 \(f\) 在逆极限中的零伦。前 \(N-1\) 个坐标由投影取得，所以舍去有限个初始阶段不会遗漏条件。

	同样的道路构造处理道路分支：各阶段道路连通，而 \(\pi_1(P_nX)=0\) 使逐层选择的道路可以调整为相容道路。因此逆极限道路连通。综上，\(\pi_i(P_\infty X)\) 与稳定群 \(\pi_i(X)\) 同构，且这些同构由 \(a\) 给出。
\end{proof}

逆极限本身未必是 CW 复形，所以这里结论写为弱同伦等价。取它的 CW 逼近后，Whitehead 定理便恢复 \(X\) 的 CW 同伦型。上面的证明实际使用了各次同伦群最终稳定；一般逆极限不能直接把同伦群逐项取极限。

\subsection{k-不变量和阶段之间的粘合}
\label{g2:subsec:1203:03}

前节的可表性定理把 \(H^{n+1}(B;G)\) 中的类表示为映射 \(B\to K(G,n+1)\)。给出这样的映射 \(\kappa\)，就可以拉回路径纤维化，构造
\begin{equation}\label{g2:eq:k-invariant-path-model}
	E_\kappa=\{(b,\lambda):b\in B,\ \lambda(0)=\kappa(b),\ \lambda(1)=*\}
	\longrightarrow B.
\end{equation}
纤维是 \(\Omega K(G,n+1)\)，其唯一非零同伦群是 \(\pi_n\cong G\)。道路从 \(\kappa(b)\) 指向基点，这正是第8章规定的 \(\operatorname{hofib}(\kappa)\) 模型。它拉回的是终点固定的路径纤维化 \(P^{\mathrm{in}}K\xrightarrow{\operatorname{ev}_0}K\)。

\ProseParagraphBreak

关键问题是：每张阶段纤维化是否都能由这种形式得到？下面用一个实际构造证明，而后再把构造得到的映射识别为障碍类。

\begin{lemma}[单一同伦纤维的粘合映射]\label{g2:lem:principal-em-fibration}
	设 \(q:Z\to B\) 是单连通带基点 CW 复形之间的映射，\(n\ge2\)，其同伦纤维具有 \(K(G,n)\) 的弱同伦型，其中 \(G\) 为阿贝尔群。则存在带基点映射 \(\kappa:B\to K(G,n+1)\) 和覆盖 \(B\) 的弱同伦等价
	\[
		j:Z\longrightarrow E_\kappa,\qquad \operatorname{pr}_B j=q.
	\]
	可以按指定的 \(\pi_n(\operatorname{hofib}q)\cong G\) 选择模型，使纤维比较在这个群上为恒等。
\end{lemma}
\begin{proof}
	先假设 \(q\) 胞腔。一般情形取带基点胞腔逼近 \(q'\) 及同伦 \(H:q\simeq q'\)；为 \(q'\) 构造比较后，在所得道路前接上 \(\kappa H(z,-)\)，并将基底坐标改为 \(q(z)\)，便得到原映射的比较。令 \(M=M_q\) 为映射柱，内含 \(Z=Z\times\{0\}\)，另一端为 \(B\)。\(M\) 形变收缩到 \(B\)，且 \((M,Z)\) 为 CW 对。同伦纤维与相对同伦群的长正合列比较给出
	\[
		\pi_i(M,Z)=0\ (i\ne n+1),\qquad \pi_{n+1}(M,Z)\cong G.
	\]
	具体地，把 \(q\) 作路径纤维化替换后，相对圆盘的基底投影与纤维中的边界互相对应，得到 \(\pi_i(M,Z)\cong\pi_{i-1}(\operatorname{hofib}q)\)（\(i\ge2\)）。这也是第8章纤维化长正合列的空间对形式。

	将 \(Z\) 压成一点，记 \(Q=M/Z\)。相对 Hurewicz 定理用于单连通的 \(M,Z\)，说明 \(H_i(M,Z)=0\)（\(i\le n\)），且 \(H_{n+1}(M,Z)\cong G\)。商的相对胞腔链给出 \(\widetilde H_i(Q)\cong H_i(M,Z)\)。此外 \(Q\) 单连通：把商视作在 \(M\) 上沿 \(Z\) 附着一个锥，van Kampen 将其基本群识别为 \(\pi_1(M)\) 除去 \(\pi_1(Z)\) 的像的正规闭包，两者在这里均为零。依次使用首个非零次数的 Hurewicz 定理，便知 \(Q\) 为 \(n\) 连通，且 \(\pi_{n+1}(Q)\cong G\)。相对与绝对 Hurewicz 映射的自然性还说明，商映射诱导
	\[
		\pi_{n+1}(M,Z)\xrightarrow{\cong}\pi_{n+1}(Q).
	\]

	现在从 \(Q\) 开始附着维数至少为 \(n+3\) 的胞腔，逐次消去更高同伦群，得到 \(K=K(G,n+1)\) 及 \(v:Q\to K\)。复合 \(\bar\kappa:M\to Q\xrightarrow{v}K\) 在 \(Z\) 上恒取基点，并在全部相对同伦群上给出同构
	\[
		\pi_i(M,Z)\xrightarrow{\cong}\pi_i(K,*).
	\]
	将 \(\bar\kappa\) 限制到映射柱的另一端 \(B\)，定义 \(\kappa:B\to K\)。对 \(z\in Z\)，反向走过映射柱线段，得到道路 \(\lambda_z(t)=\bar\kappa([z,1-t])\)，其起点为 \(\kappa(q(z))\)、终点为 \(*\)。因此
	\[
		j(z)=(q(z),\lambda_z)
	\]
	是一张所需的连续映射。

	最后验证弱等价。比较空间对 \((M,Z)\) 的长正合列与 \(E_\kappa\to B\to K\) 的同伦纤维序列，得到
	\[
		\begin{tikzcd}[column sep=small]
			\pi_{i+1}(B)\arrow[r]\arrow[d,equal]&\pi_{i+1}(M,Z)\arrow[r]\arrow[d,"\cong"]&\pi_i(Z)\arrow[r]\arrow[d,"j_*"]&\pi_i(B)\arrow[r]\arrow[d,equal]&\pi_i(M,Z)\arrow[d,"\cong"]\\
			\pi_{i+1}(B)\arrow[r]&\pi_{i+1}(K)\arrow[r]&\pi_i(E_\kappa)\arrow[r]&\pi_i(B)\arrow[r]&\pi_i(K).
		\end{tikzcd}
	\]
	交换性来自 \(\bar\kappa\) 及刚写出的零伦 \(\lambda_z\)。两侧相对群的竖箭头已为同构，故长正合列给出所有 \(j_*\) 为同构；低阶项还说明两空间均单连通。所诱导的纤维比较在 \(\pi_n\) 上是某个 \(G\) 的自同构。改变 \(v\) 在唯一非零同伦群上的识别，即后接实现其逆的 \(K(G,n+1)\) 自同伦等价，就可使这张纤维比较为指定的恒等。这样的自同伦等价可按第11章 Eilenberg–MacLane 模型比较中的逐胞腔构造取得。
\end{proof}

构造中两次处理有不同作用：商 \(M/Z\) 把“在总空间上已经为零”的信息变成一个带基点对象；随后附着高维胞腔，只保留所需的 \(G\)。映射柱线段把这份相对信息重新变成 \(\kappa q\) 的一个具体零伦，所以能恢复总空间到同伦纤维的映射。

\ProseParagraphBreak

为了从这个构造读出上同调类，给 \(K(G,n+1)\) 选普遍类
\[
	\iota_{n+1}\in H^{n+1}(K(G,n+1);G)\cong\operatorname{Hom}(G,G),\qquad
	\iota_{n+1}\longleftrightarrow\operatorname{id}_G.
\]
这个同构来自 Hurewicz 与万有系数定理，低于 \(n+1\) 的约化同调均为零。终点固定的路径空间可缩，所以空间对的连接同态 \(\delta_{\mathrm{in}}\) 给出同构。用
\[
	\sigma_{\mathrm{in}}\iota_{n+1}=\delta_{\mathrm{in}}^{-1}\operatorname{ev}_0^*\iota_{n+1}
\]
规定 \(H^n(\Omega K;G)\) 中的恒等类，即按这个类识别纤维的 \(\pi_n\cong G\)。第11章的相对传递判据直接给出
\[
	\mathrm d_{n+1}(\operatorname{id}_G)=\iota_{n+1}.
\]
这里仅使用加性谱序列；其构造对任意阿贝尔系数成立，不要求 \(G\) 自带环乘法。这个识别也明确记录了道路方向：反向道路映射 \(r:P^{\mathrm{in}}K\to P_*K\) 覆盖 \(K\) 上的恒等，故 \(r|_{\Omega K}^*(\sigma\iota_{n+1})=\sigma_{\mathrm{in}}\iota_{n+1}\)。在同一环路空间中，反向环路在 \(\pi_n(\Omega K)\) 上取负，因此入向与出向的两种悬挂类相差负号。引理中允许调整的纤维识别正是按这里的约定选择。

\begin{definition}[阶段的 \(k\)-不变量]\label{g2:def:postnikov-k-invariant}
	设 \(X\) 为单连通带基点 CW 复形，\(n\ge2\)，记 \(G=\pi_n(X)\)。在阶段纤维化\eqref{g2:eq:postnikov-stage-fibration}中，用 \(\pi_n(F_n)\cong G\) 将纤维类 \(\iota_F\in H^n(F_n;G)\) 识别为 \(\operatorname{id}_G\)。定义
	\begin{equation}\label{g2:eq:postnikov-k-transgression}
		k_{n+1}(X)=\mathrm d_{n+1}\iota_F\in H^{n+1}(P_{n-1}X;G),
	\end{equation}
	称为这一阶段的\term[k bubianliang]{\(k\)-不变量}[\(k\)-invariant]。下标 \(n+1\) 在本节记录上同调次数。
\end{definition}

这个传递在整个目标群中取值。因为 \(F_n\) 为 \((n-1)\) 连通，纤维上同调的第 \(1,\ldots,n-1\) 行全为零，没有更早微分改变 \((0,n)\) 与 \((n+1,0)\) 两项。基底单连通又使 \(G\) 的局部系统为常值。

\begin{theorem}[由 \(k\)-不变量恢复阶段]\label{g2:thm:postnikov-k-reconstruction}
	设 \(X\) 为单连通带基点 CW 复形。对 \(n\ge2\)，令 \(G=\pi_n(X)\)，并用\cref{g2:thm:em-representability}将 \(k_{n+1}(X)\) 表示为映射
	\[
		\kappa_{n+1}:P_{n-1}X\longrightarrow K(G,n+1).
	\]
	则 \(P_nX\) 与 \(\operatorname{hofib}(\kappa_{n+1})\) 具有相同的弱同伦型，比较保持到 \(P_{n-1}X\) 的映射及指定的纤维同伦群识别。在 CW 模型上，这给出同伦等价。特别地，若 \(k_{n+1}(X)=0\)，这一阶段分裂为 \(P_{n-1}X\times K(G,n)\) 的同伦型。
\end{theorem}
\begin{proof}
	先在 \(X[n]\to X[n-1]\) 的 CW 模型上应用\cref{g2:lem:principal-em-fibration}，再经阶段同伦等价搬到 \(P_nX\to P_{n-1}X\)。得到某张 \(\kappa\) 及纤维上诱导恒等的弱等价。比较此纤维化与普遍路径纤维化，传递的自然性给出
	\[
		\mathrm d_{n+1}\iota_F=\kappa^*\iota_{n+1}.
	\]
	左边正是定义的 \(k_{n+1}(X)\)。所以构造出的 \(\kappa\) 表示所需类，可表性定理说明它与任意选定的 \(\kappa_{n+1}\) 同伦。

	还需检查更换代表映射不改变总空间。若 \(H:B\times I\to K\) 是 \(\kappa_0\) 到 \(\kappa_1\) 的同伦，则在式\eqref{g2:eq:k-invariant-path-model}中的道路 \(\lambda\) 前接上反向道路 \(t\mapsto H(b,1-t)\)，得到覆盖 \(B\) 的映射 \(E_{\kappa_0}\to E_{\kappa_1}\)。用原向同伦给出逆；往返道路可保持端点收缩，故两复合与恒等映射纤维同伦。因而同伦类相同的代表给出相同的纤维同伦型。

	当类为零时，取常值代表 \(\kappa=0\)。此时式\eqref{g2:eq:k-invariant-path-model}直接成为 \(B\times\Omega K(G,n+1)\)，其 CW 同伦型即 \(B\times K(G,n)\)。
\end{proof}

可以把这一构造画成拉回方块：
\[
	\begin{tikzcd}[column sep=large,row sep=large]
		P_nX\arrow[r]\arrow[d,"p_n"']&P^{\mathrm{in}}K(G,n+1)\arrow[d,"\operatorname{ev}_0"]\\
		P_{n-1}X\arrow[r,"\kappa_{n+1}"']&K(G,n+1).
	\end{tikzcd}
\]
在本图中用同伦等价的模型替换 \(P_nX\)，上边的路径空间可缩，纤维均为 \(K(G,n)\)。底边的一张映射便决定了这一层怎样附着。

\ProseParagraphBreak

\(k_{n+1}\) 与阶段纤维化的首截面障碍相联系，但须遵守两种边界的方向。若障碍 \(o_{n+1}\) 按胞腔正向边界上的截面映射取值，则本书的约定给出
\[
	o_{n+1}=-k_{n+1}.
\]
可在单个胞腔上看见这个符号。把纤维化拉回 \(D^{n+1}\)，在圆盘上选一个局部截面，其边界记为 \(l\)；既定骨架截面的边界记为 \(s\)。两者的差 \(s-l\) 经纤维方向的变形代表障碍 \(o_{n+1}\)。因此与 \(\iota_F\) 配对，若 \(A\) 是骨架上延拓纤维恒等类的上圈，就有 \(A(s)-A(l)=o_{n+1}\)。延拓 \(A\) 为下一层的上链后，传递上链在该胞腔上取值 \(A(l)\)；而 \(A(s)\) 是把 \(A\) 沿骨架截面拉回后所得上链的胞腔余边界。模去余边界，便得到 \(\tau\iota_F=-o_{n+1}\)。这与第11章“延拓后取余边界”的传递约定相容。

\ProseParagraphBreak

对给定 \(f:Y\to P_{n-1}X\)，式\eqref{g2:eq:k-invariant-path-model}直接说明：把 \(f\) 提升到 \(P_nX\)，就是给 \(\kappa_{n+1}f\) 选一个零伦；存在条件为
\begin{equation}\label{g2:eq:postnikov-lift-obstruction}
	f^*k_{n+1}(X)=0\in H^{n+1}(Y;G)
\end{equation}
其中 \(Y\) 为 CW 复形。这一次条件是完整条件，因为纤维 \(K(G,n)\) 没有其他非零同伦群。

\ProseParagraphBreak

一个 \(k\)-不变量的具体表达还依赖于对阶段与系数群的识别。若 \(h:B\to B'\) 是基底比较，\(\alpha:G\to G'\) 是系数同构，则对应的相容关系为 \(\alpha_*k=h^*k'\)。系数变化改变纤维的命名，基底比较则改变类所处的空间；恢复阶段时两者都要同时记录。对非单连通空间，逐次消去高阶同伦群仍产生阶段，但 \(\pi_1\) 对 \(\pi_n\) 的作用给出局部系统 \(\Pi_n\)，相应障碍应取在 \(H^{n+1}(P_{n-1}X;\Pi_n)\) 中。作用非平凡时，不能把它替换为上述到普通 \(K(G,n+1)\) 的常系数映射。

\subsection{同伦群相同的不同三型}
\label{g2:subsec:1203:04}

现在给出同伦群完全相同、同伦型却不同的一族空间。取
\[
	B=K(\mathbb Z,2)=\mathbb CP^\infty,\qquad H^*(B;\mathbb Z)=\mathbb Z[u],\quad |u|=2.
\]
对整数 \(m\)，可表性定理给出代表 \(m u^2\in H^4(B;\mathbb Z)\) 的映射
\[
	\kappa_m:B\longrightarrow K(\mathbb Z,4).
\]
令 \(E_m\) 为其同伦纤维的 CW 逼近。等价的路径模型给出纤维化 \(K(\mathbb Z,3)\to E_m\to K(\mathbb Z,2)\)，以下计算在该模型上进行，CW 逼近不改变同伦群与奇异同调。

\begin{example}[相同同伦群的不同三型]\label{g2:exm:postnikov-two-stage-family}
	上述 \(E_m\) 均单连通，并且
	\[
		\pi_i(E_m)\cong
		\begin{cases}
			\mathbb Z,&i=2,3,\\
			0,&i\ne2,3\quad(i\ge1).
		\end{cases}
	\]
	然而
	\[
		H^3(E_m;\mathbb Z)\cong\ker(\mathbb Z\xrightarrow{\times m}\mathbb Z),\qquad
		H^4(E_m;\mathbb Z)\cong\operatorname{coker}(\mathbb Z\xrightarrow{\times m}\mathbb Z).
	\]
	因此不同正整数 \(m\) 给出两两不同的同伦型；\(m=0\) 与 \(m=1\) 也不同。
\end{example}
\begin{proof}
	对 \(E_m\to K(\mathbb Z,2)\to K(\mathbb Z,4)\) 用同伦长正合列。二次群由基底给出；三次群由连接同态 \(\pi_4(K(\mathbb Z,4))\to\pi_3(E_m)\) 给出；其余正阶项为零。这个计算与 \(m\) 无关，因为 \(\kappa_m\) 在全部同伦群上诱导零映射：两端唯一非零的次数不同。

	上同调能看见 \(m\)。纤维 \(K(\mathbb Z,3)\) 可由 \(S^3\) 附着维数至少为五的胞腔构造，故其 \(H_3=\mathbb Z\)、\(H_4=0\)；万有系数定理给出 \(H^3=\mathbb Z\)、\(H^4=0\)，一次、二次也为零。按入向路径纤维化的悬挂类选择三次生成元 \(v\)。在总次数不超过四的范围，Serre 页只有基底的 \(1,u,u^2\) 和 \((0,3)\) 处的 \(v\)。唯一可能的微分为
	\[
		\mathrm d_4:\mathbb Zv\longrightarrow\mathbb Zu^2,\qquad \mathrm d_4v=m u^2,
	\]
	其中等式来自路径模型的自然性。第三次因此留下这张映射的核，第四次留下其余核；每个次数只有一层，故无加性扩张。

	当 \(m=0\) 时，\(E_0\) 为 \(K(\mathbb Z,2)\times K(\mathbb Z,3)\) 的同伦型，三次、四次上同调均为 \(\mathbb Z\)。当 \(m=1\) 时，两群均为零；当 \(m>1\) 时，第四次群为 \(\mathbb Z/m\)。这些群区分了所列空间。
\end{proof}

这个例子也解释 Whitehead 定理中的映射条件。每个 \(E_m\) 都有同样的抽象同伦群，但若一张映射在全部同伦群上诱导同构，它就必须是 CW 模型之间的同伦等价，因而还必须保持刚算出的上同调。上述已由上同调区分的空间之间不存在这样的实现映射；障碍正包含在 \(k_4=m u^2\) 中。

\ProseParagraphBreak

球面 \(S^2\) 给出 \(m=1\) 的一个具体来源。已有 \(\pi_2(S^2)=\pi_3(S^2)=\mathbb Z\)，所以它的第三阶段位于
\[
	K(\mathbb Z,3)\longrightarrow P_3S^2\longrightarrow K(\mathbb Z,2).
\]
\(S^2\to S^2[3]\) 只添加维数至少为五的胞腔。因而 \(H_2=\mathbb Z\)、\(H_3=H_4=0\)；即使五胞腔能对四次同调增加边界关系，原来的四次群已经为零。万有系数定理遂给出 \(H^3(P_3S^2)=H^4(P_3S^2)=0\)。写 \(k_4(S^2)=m u^2\)，上一计算迫使乘 \(m\) 无核且无余核，故 \(m=\pm1\)。改变纤维生成元的符号后可取 \(m=1\)，所以 \(P_3S^2\) 是 \(\operatorname{hofib}(u^2:\mathbb CP^\infty\to K(\mathbb Z,4))\) 的 CW 同伦型。Hopf 类提供三次同伦群，\(u^2\) 则记录这层怎样接在二次阶段上。

\ProseParagraphBreak

对阶段构造、主纤维化与 \(k\)-不变量的进一步阅读，可参看\cite[第4.3节，Postnikov Towers]{Hatcher2002}及\cite[第22章第4节，Postnikov systems]{May1999}。本节采用上同调次数作为 \(k\) 的下标；比较不同文献时，先核对该类所在的 \(H^{n+1}(P_{n-1}X;\pi_nX)\)，再比较下标。
% 纲要条目排为编号小节标题，并进入目录与书签；说明段仍保留于纲要注释中。

% !TeX root = ../../Book2.tex
% 纲要标题：### 12.4 几何中的截面障碍
% 纲要位置：几何学计划纲要.md，第 1493 行。
% 正文状态：初稿；尚未正式审核。

% 纲要内容开始
% 
% * 球丛的首障碍与 Euler 类
% * 向量场及结构群约化
% * 平面丛与曲面上的向量场
% 
% ---
% 
% 纲要内容结束

\section{几何中的截面障碍}
\label{g2:sec:1204}

给每个点连续选一个单位向量，是一个截面问题。沿一个胞腔的边界走一圈，已经选好的向量可能在纤维球面上绕出非零同伦类，这就阻止它填满胞腔。上一章则从纤维定向类的上同调延拓定义了 Euler 类。本节把这两种延拓放到同一个胞腔上比较，再用向量场、平面丛和曲面上的计算说明：障碍类的次数、系数以及消失后的结论都由纤维决定。

\subsection{球丛的首障碍与 Euler 类}
\label{g2:subsec:1204:01}
设 \(p:E\rightarrow B\) 是定向 \(S^m\) 丛，\(m\ge1\)，其中 \(B\) 为道路连通 CW 复形。这里的定向仍指沿道路运输保持纤维的正向顶次类。\cref{g2:prop:sphere-connectivity}给出 \(\pi_j(S^m)=0\) 对 \(1\le j<m\) 成立，再加上纤维道路连通，故可以从零骨架开始逐胞腔选截面，一直得到 \(s:B^m\rightarrow E_m\)，其中 \(E_m=p^{-1}(B^m)\)。第一次可能失败发生在 \((m+1)\) 胞腔：将丛拉回其特征圆盘并平凡化，边界截面成为 \(a:S^m\rightarrow S^m\)。\cref{g2:lem:sphere-degree-classification}说明它零伦当且仅当 \(\deg a=0\)，所以延拓问题成为一个整数是否为零的问题。

\ProseParagraphBreak

源球面取胞腔圆盘的正向边界，目标球面取纤维定向。\cref{g2:thm:primary-obstruction}于是给出与初始截面选择无关的首障碍
\[
	o_{m+1}(p)\in H^{m+1}(B;\mathbb Z).
\]
系数为普通整数，是因为保持纤维定向的运输在 \(\pi_m(S^m)\cong\mathbb Z\) 上作用恒等。\cref{g2:def:spherical-euler-class}所定义的 \(e(p)=\mathrm d_{m+1}u\) 也在这个群里；其中 \(u\) 是纤维的正向上同调生成元。比较二者时需要保留连接同态的符号，下面的局部计算同时给出识别及其符号。

\begin{theorem}[球丛的首障碍]\label{g2:thm:sphere-primary-euler}
	设 \(p:E\rightarrow B\) 为定向 \(S^m\) 丛，\(m\ge1\)，\(B\) 为道路连通 CW 复形。障碍上链按正向胞腔边界上的截面映射度取值，Euler 类按式\eqref{g2:eq:spherical-euler-transgression}定义，则
	\begin{equation}\label{g2:eq:primary-euler-sign}
		o_{m+1}(p)=-e(p).
	\end{equation}
	若 \(\dim B\le m+1\)，则 \(p\) 有连续截面当且仅当 \(e(p)=0\)。对定向圆周丛 \(m=1\)，这一等价对任意维数的 CW 底空间均成立。
\end{theorem}
\begin{proof}
	先比较同一个胞腔上的两种延拓。固定 \(B^m\) 上的截面 \(s\)。由\cref{g2:prop:serre-transgression}，纤维生成元 \(u\) 可延拓为 \(v\in H^m(E_m;\mathbb Z)\)：底空间只有 \(m\) 维，球丛谱序列中唯一可能离开 \((0,m)\) 的微分已落在其维数以外。以
	\[
		v-p^*s^*v
	\]
	替换 \(v\)，不改变它在纤维上的限制，并使 \(s^*v=0\)。减去这一底空间类，是为了把截面所带的贡献从延拓类中消去。

	取一个正向 \((m+1)\) 胞腔的特征映射 \(\chi:D^{m+1}\rightarrow B\)，在拉回丛 \(D^{m+1}\times S^m\) 中使用保持纤维定向的平凡化。令 \(x\in H^m(\partial D^{m+1};\mathbb Z)\) 为正向边界生成元。边界上的乘积有
	\[
		H^m(\partial D^{m+1}\times S^m;\mathbb Z)=\mathbb Z\operatorname{pr}_1^*x\oplus\mathbb Z\operatorname{pr}_2^*u.
	\]
	\(v\) 在纤维上的值为 \(u\)，所以它在这个边界乘积上的限制必为
	\[
		\operatorname{pr}_2^*u+\lambda\operatorname{pr}_1^*x,\qquad \lambda\in\mathbb Z.
	\]
	边界截面是图映射 \(\Gamma_a(z)=(z,a(z))\)。把上式沿 \(\Gamma_a\) 拉回，条件 \(s^*v=0\) 给出
	\[
		0=a^*u+\lambda x=(\deg a+\lambda)x,
	\]
	从而 \(\lambda=-\deg a\)。

	现在使用空间对 \((D^{m+1}\times S^m,\partial D^{m+1}\times S^m)\) 的连接同态。纤维类 \(\operatorname{pr}_2^*u\) 已定义在整个圆盘乘积上，故连接同态将它送到零；\(\operatorname{pr}_1^*x\) 则被送到圆盘的正向相对生成元乘纤维的零次单位元。后一个正号可以直接由上链检查：若 \(\widetilde x\) 延拓边界代表元，则
	\[
		\langle\mathrm d\widetilde x,[D^{m+1}]\rangle=\langle x,\partial[D^{m+1}]\rangle=1.
	\]
	因此 \(v\) 的上同调延拓障碍在该胞腔上的值是 \(\lambda=-\deg a\)。\cref{g2:prop:serre-transgression}把这些值组成的上圈识别为 \(\mathrm d_{m+1}u\)，而截面首障碍的值为 \(\deg a\)。逐胞腔比较后得到式\eqref{g2:eq:primary-euler-sign}。

	若 \(\dim B\le m+1\)，首障碍为零时，\cref{g2:thm:primary-obstruction}允许先修改 \(B^m\) 上的截面，再将它延拓到全部 \((m+1)\) 胞腔；这已经覆盖整个 \(B\)。反之，全局截面使全部胞腔障碍为零。结合比较式便得到所述等价。若 \(m=1\)，首障碍为零后先延拓到 \(B^2\)，以后每一步的系数群 \(\pi_q(S^1)\)（\(q\ge2\)）都为零，因而可以继续延拓。各骨架上的截面相容，由 CW 弱拓扑粘合为全局连续截面。
\end{proof}

式\eqref{g2:eq:primary-euler-sign}中的负号来自已经固定的两种约定：截面障碍取边界映射的正向度，上一章 Euler 类取“延拓后取余边界”的传递。若采用首障碍来规范 Euler 类，其记号对应这里的 \(-e(p)\)。判定截面存在时只使用这个类是否为零，因而两种规范给出相同判定；计算它在定向基本类上的具体整数时则必须保持所选符号。

\ProseParagraphBreak

对 \(m\ge2\)，\(e(p)=0\) 保证可以选到 \(B^{m+1}\) 上的截面。继续向上延拓仍要考察 \(\pi_{m+1}(S^m),\pi_{m+2}(S^m),\ldots\) 中的障碍，下一步落在 \(H^{m+2}(B;\Pi_{m+1}(p))\) 中，并可能依赖已选截面。球面低阶同伦群的消失给出上述维数范围；圆周的全部高阶同伦群消失，则使其首障碍成为任意维数上的完整判据。若球丛不定向，首障碍仍可构造，但系数须换为由纤维定向反转给出作用的整数局部系统；上述普通整数 Euler 类的陈述以定向为条件。

\subsection{向量场及结构群约化}
\label{g2:subsec:1204:02}
设 \(V\rightarrow B\) 为实秩 \(r\) 向量丛，并给定一个连续正定内积。记 \(S(V)\) 为各纤维单位球面组成的球丛，\(\varepsilon_B^k=B\times\mathbb R^k\) 为平凡秩 \(k\) 丛。无零截面 \(v\) 可逐点归一化为 \(v/\|v\|\)，所以寻找无零截面就是寻找 \(S(V)\) 的截面。对光滑流形上的光滑向量丛，第一卷第12章第3节的内积存在定理保证了这样的光滑内积；其假设为流形有限维、Hausdorff 且第二可数，证明是用局部有限的光滑单位分解将局部 Euclid 内积相加。

\ProseParagraphBreak

还可以从标架看这个问题。局部选取正交标架后，过渡矩阵属于 \(\operatorname O(r)\)。若第一根标架向量能够在不同图上保持不变，过渡矩阵便有 \(\operatorname{diag}(1,A)\) 的形式。这提示我们，整体向量场对应于过渡矩阵所能取值的群变小。

\begin{definition}[正交结构群约化]\label{g2:def:orthogonal-reduction}
	设 \(V\rightarrow B\) 是带连续正定内积的实秩 \(r\) 向量丛，\(0\le k\le r\)。若存在局部正交标架，使所有过渡矩阵都落在子群
	\[
		\left\{\begin{pmatrix}I_k&0\\0&A\end{pmatrix}:A\in\operatorname O(r-k)\right\}\subseteq\operatorname O(r),
	\]
	则称 \(V\) 的正交\term[jiegouqun-yuehua]{结构群约化}[reduction of structure group]到 \(\operatorname O(r-k)\)。这里约定 \(\operatorname O(0)\) 为平凡群。对定向丛，使用正向正交标架，并将两个正交群换成相应的特殊正交群，就得到到 \(\operatorname{SO}(r-k)\) 的约化。
\end{definition}

\begin{proposition}[独立截面与约化]\label{g2:prop:sections-orthogonal-reduction}
	设 \(V\rightarrow B\) 为带连续正定内积的实秩 \(r\) 向量丛，\(1\le k\le r\)。以下条件等价：
	\begin{equivalences}
		\item \(V\) 有 \(k\) 个逐点线性无关的连续截面；
		\item \(V\) 有 \(k\) 个逐点正交归一的连续截面；
		\item 存在秩 \(r-k\) 向量丛 \(W\)，使 \(V\cong\varepsilon_B^k\oplus W\)；
		\item 正交结构群可约化到 \(\operatorname O(r-k)\)。
	\end{equivalences}
	对光滑向量丛及光滑内积，同一等价在全部截面、丛同构和局部标架都要求光滑时成立。若 \(V\) 已定向，则上述条件也等价于到 \(\operatorname{SO}(r-k)\) 的约化。
\end{proposition}
\begin{proof}
	对逐点独立的截面 \(v_1,\ldots,v_k\) 作 Gram--Schmidt 正交化：令
	\[
		w_j=v_j-\sum_{i<j}\langle v_j,e_i\rangle e_i,\qquad e_j=\frac{w_j}{\|w_j\|}.
	\]
	独立性保证每个 \(w_j\) 都处处非零，故这些公式保持连续性；在光滑情形也保持光滑性。由此得到 I 推出 II，反向蕴含直接成立。

	给定 II 中的正交截面，令 \(W_b\) 为它们在 \(V_b\) 中的共同正交补。映射
	\[
		P_b(v)=v-\sum_{i=1}^k\langle v,e_i(b)\rangle e_i(b)
	\]
	是到 \(W_b\) 的投影。在某点选取 \(r-k\) 个局部截面，其投影在该点构成 \(W_b\) 的基；相应 Gram 行列式在附近仍非零，因此这些投影给出 \(W\) 的局部标架。这证明 \(W\) 确为连续子丛，光滑时也为光滑子丛。映射
	\[
		\varepsilon_B^k\oplus W\longrightarrow V,\qquad (b,t_1,\ldots,t_k,w)\longmapsto\sum_i t_i e_i(b)+w
	\]
	是丛同构，其逆由内积系数与 \(P_b\) 给出，故 II 推出 III。一个直和分解中平凡因子的标准截面又给出 I。

	在 \(W\) 上局部正交化标架，将它接到 \(e_1,\ldots,e_k\) 后面，便得到 \(V\) 的局部正交标架。不同图上的前 \(k\) 根相同，过渡矩阵因而具有所述块对角形式，得到 IV。反过来，这种过渡矩阵固定前 \(k\) 根标准向量，所以各图中的这些局部截面粘成全局正交截面。定向情形下，当 \(k<r\) 时，由 \(V\) 的定向及有序截面组赋予 \(W\) 相容定向，再使用 \(W\) 的正向正交标架。当 \(k=r\) 时，可在每个连通分支上按需反转最后一根向量，使全标架保持给定定向。于是同样得到特殊正交群的结论。
\end{proof}

这个证明使用的正交补构造也见第一卷第12章第3节；这里把投影写出，是为了说明它同时适用于连续丛，以及约化中的标架如何实际产生。第一卷第12章第1节的“整体标架与平凡性”正是上述 \(k=r\) 的情形。

\ProseParagraphBreak

把一个纤维中全部有序正交 \(k\) 标架放在一起，得到\term[stiefel-cong]{Stiefel 丛}[Stiefel bundle] \(V_k(V)\rightarrow B\)，其纤维是
\[
	V_k(\mathbb R^r)=\{(e_1,\ldots,e_k):\langle e_i,e_j\rangle=\delta_{ij}\}.
\]
局部正交标架将它平凡化为 \(U\times V_k(\mathbb R^r)\)，过渡矩阵同时作用于每根向量，所以此构造与局部选择相容。它的截面就是 \(k\) 个正交单位截面。\(k=1\) 恢复单位球丛；\(k>1\) 时，延拓所用的纤维同伦群变成 Stiefel 空间的同伦群。这将“找多少个独立向量场”统一成了截面问题。\symidx{\(V_k(\mathbb R^r)\)：有序正交 \(k\) 标架的 Stiefel 空间}

\ProseParagraphBreak

在 \(r\) 维光滑流形 \(M\) 上取 \(V=\mathrm TM\)，处处非零向量场等价于单位切丛的截面，也等价于正交结构群从 \(\operatorname O(r)\) 约化到 \(\operatorname O(r-1)\)。这里连续存在性与光滑存在性一致。事实上，给一个连续无零截面 \(v\)，在每点附近取局部光滑截面 \(w_i\)，使 \(\|w_i-v\|<\|v\|/2\)；可以在局部标架中固定该点的系数，再由连续性缩小邻域来做到。用第一卷第12章第3节的光滑单位分解 \(\{\varphi_i\}\) 形成 \(w=\sum_i\varphi_iw_i\)，便有
\[
	\|w-v\|\le\sum_i\varphi_i\|w_i-v\|<\frac{\|v\|}{2},
\]
故 \(w\) 是光滑无零截面。在维数不超过 \(r\) 的 CW 底空间上，\(r\ge2\) 的定向秩 \(r\) 向量丛只有首个截面障碍落在可见维数内，所以 \(e(V)=0\) 正好判定无零截面的存在。

\ProseParagraphBreak

秩一时单位纤维是离散的 \(S^0\)，延拓首先要求两端所选分支能够连起来。例如 Möbius 线丛
\[
	L=([0,1]\times\mathbb R)/\bigl((1,t)\sim(0,-t)\bigr)\longrightarrow S^1
\]
的截面在区间上由连续函数 \(a\) 表示，端点条件为 \(a(1)=-a(0)\)。若 \(a(0)=0\)，已经有零点；若不为零，介值定理仍迫使 \(a\) 在区间内取零。单位纤维的两个点绕基底一圈后交换，正是无法选取连续分支的原因。这个障碍发生在一维，属于离散纤维的分支问题。

\subsection{平面丛与曲面上的向量场}
\label{g2:subsec:1204:03}
二维底空间配圆周纤维时，整个延拓过程只需一个整数上圈：先在顶点和边上选截面，再读取各面边界的绕数。下面先让过渡函数直接给出这个整数，再在切丛上把它拆成两个局部贡献。

\begin{example}[球面上的旋转粘合]\label{g2:exm:clutched-plane-obstruction}
	给 \(S^2\) 一个定向，把它写成两张闭圆盘 \(D_+\cup D_-\)。在公共边界上用 \(t\in\mathbb R/2\pi\mathbb Z\) 表示 \(D_-\) 的正向边界参数。对整数 \(k\)，将两张平凡定向实平面丛按
	\[
		(t,z)_+\sim(t,e^{\mathrm i kt}z)_-,\qquad z\in\mathbb C\cong\mathbb R^2,
	\]
	粘合，得到定向平面丛 \(V_k\rightarrow S^2\)；Euclid 内积因旋转粘合而给出全局内积。

	在 \(D_+\) 上选常值单位截面 \(1\)，它在 \(D_-\) 的边界坐标中变成 \(e^{\mathrm i kt}\)，绕数为 \(k\)。把 \(D_+\) 作为已经固定截面的子复形，填入 \(D_-\) 的相对首障碍就是这个整数。由于 \(D_+\) 可缩，\(H^2(S^2,D_+;\mathbb Z)\rightarrow H^2(S^2;\mathbb Z)\) 为同构；正向 \(D_-\) 对应满足 \(\langle\omega,[S^2]\rangle=1\) 的生成元 \(\omega\)。于是
	\[
		o_2(S(V_k))=k\omega,\qquad e(V_k)=-k\omega.
	\]
	\cref{g2:thm:sphere-primary-euler}说明 \(V_k\) 有无零截面当且仅当 \(k=0\)。这里即使允许修改 \(D_+\) 上的截面，结论也相同：其沿边界的附加绕数为零，因为它已经延拓到 \(D_+\) 内部。\(k=0\) 时粘合为恒等，两个标准基向量直接粘成整体标架。
\end{example}

对闭定向连通曲面 \(\Sigma_g\) 上的定向平面丛 \(V\)，同一个方法把存在性归结为一个整数。\cref{g2:prop:surface-integer-homology,g2:thm:uct-cohomology}给出 \(H^2(\Sigma_g;\mathbb Z)\cong\mathbb Z\)；取正向生成元 \(\omega\)，写 \(e(V)=q\omega\)。于是 \(V\) 有无零截面当且仅当 \(q=0\)。一旦选定一个单位截面 \(s_0\)，平面定向与内积给出逐纤维的正向四分之一转 \(J\)，\((s_0,Js_0)\) 就是全局正向正交标架。因此每个单位截面都唯一写为
\[
	s=f_1s_0+f_2Js_0,\qquad f=(f_1,f_2):\Sigma_g\longrightarrow S^1.
\]
截面的同伦对应这些映射的同伦。\(S^1\) 是 \(K(\mathbb Z,1)\)，故\cref{g2:thm:em-representability}给出
\[
	\{\text{单位截面的同伦类}\}\cong H^1(\Sigma_g;\mathbb Z)\cong\mathbb Z^{2g}.
\]
这里自由同伦与保基点同伦给出相同的分类：选基点 \(b_0\)，将 \(f\) 换成 \(f/f(b_0)\) 就使基点映到 \(1\)，而任意同伦 \(H\) 也可规范为 \(H(b,t)/H(b_0,t)\)。这个坐标依赖于 \(s_0\)：更换参考截面会平移所有坐标。各整数可由曲面标准 \(2g\) 条生成回路上 \(f\) 的绕数读出；曲面二胞腔的附着字是交换子的乘积，其总绕数为零，所以这些整数可以任意指定。

\begin{example}[球面与环面的切向量场]\label{g2:exm:tangent-sphere-torus-obstruction}
	给 \(S^2\subset\mathbb R^3\) 外向定向，令 \(e_3=(0,0,1)\)。将常向量 \(e_3\) 正交投影到切平面，得到光滑切向量场
	\[
		v(x)=e_3-x_3x.
	\]
	因为 \(\langle v(x),x\rangle=x_3-x_3=0\)，它确实切于球面；其零点恰为南北极 \(S=-e_3,N=e_3\)。删去两极附近的互不相交小圆盘，余下的圆环记为 \(A\)。在 \(A\) 上归一化 \(v\)，就有单位切丛的截面。要填回两个圆盘，分别计算这两个边界截面的绕数。

	在北极附近用 \((x_1,x_2)\) 作坐标，\(x_3=\sqrt{1-x_1^2-x_2^2}\)。向量场在坐标切标架中的两个系数是
	\[
		(-x_3x_1,-x_3x_2).
	\]
	沿小圆周把正函数 \(x_3\) 连续变到 \(1\)，该非零系数向量同伦到 \((-x_1,-x_2)\)，其归一化是圆周上的旋转半圈，度为 \(+1\)。南极附近 \(x_3=-\sqrt{1-x_1^2-x_2^2}\)，同样的系数向量同伦到 \((x_1,x_2)\)，故度也为 \(+1\)。南极这张坐标与外向定向相反，但改用正向坐标会同时反转圆盘参数及坐标切标架的定向，两次符号变化相消，所以局部贡献仍为 \(+1\)。

	取以 \(A\) 为子复形、两张正向圆盘为剩余二胞腔的相对 CW 结构。这两个障碍组成 \(H^2(S^2,A;\mathbb Z)\cong\mathbb Z^2\) 中的 \((1,1)\)。忘掉固定在 \(A\) 上的截面，得到绝对首障碍；相对类到绝对类的映射将两个值相加，因为正向基本圈在两张圆盘上的系数都是 \(1\)。因此
	\[
		o_2(S(\mathrm TS^2))=2\omega,\qquad e(\mathrm TS^2)=-2\omega
	\]
	沿用式\eqref{g2:eq:primary-euler-sign}的规范。这里零点附近正向小圆周上的归一化向量场之度，就是该零点的\term[lingdian-jubu-zhibiao]{局部指标}[local index]；两极的指标总和为 \(2\)，对应首障碍 \(o_2\) 的取值。这证明球面不存在处处非零连续切向量场，也解释了为什么改变原向量场的局部形状不能把所有零点一起消去：所得绝对障碍类保持不变。

	环面 \(T^2=\mathbb R^2/(2\pi\mathbb Z)^2\) 则有两个整体向量场 \(\partial/\partial\theta_1\)、\(\partial/\partial\theta_2\)。它们来自 \(\mathbb R^2\) 的常向量场，在格点平移下不变，所以可下降，且处处独立。故 \(\mathrm TT^2\cong\varepsilon_{T^2}^2\)，首障碍及 Euler 类均为零；选择这个标架后，单位向量场的同伦类由两条基本回路上的绕数 \((a,b)\in\mathbb Z^2\) 完全描述。
\end{example}

这些计算分出了三个不同层次：Möbius 线丛在选取单位纤维的分支时失败；定向平面丛的分支问题自动解决，首障碍是面边界的整数绕数；首障碍消失之后，已有截面的不同同伦类还由低一次的上同调记录。截面是否存在与存在之后有多少种同伦选择，因而分别对应本章的障碍类与差类。
% 纲要条目排为编号小节标题，并进入目录与书签；说明段仍保留于纲要注释中。


\begin{chaptersummary}
	单个胞腔上的延拓取决于边界球面映射是否零伦；沿基底运输这些同伦类，便得到障碍所用的局部系数。逐胞腔障碍组成余圈，改变上一层选择的差异上链满足 \(\delta D=c(t)-c(s)\)。因此上链为零判定当前截面能够延拓，上同调类为零则允许先修改上一骨架。纤维在首个非零同伦群以下连通时，首障碍与这些初始选择无关；延拓存在后，在维数或纤维同伦群的相应范围内，差异类使同伦类组成上同调群的主齐性集。以 \(K(G,n)\) 为目标便得到 \(H^n(-;G)\) 的可表性。单连通空间的 Postnikov 塔逐次保留低阶同伦群，阶段间的纤维为 \(K(\pi_nX,n)\)，\(k_{n+1}\) 则决定这一纤维如何接上去。几何中，单位向量是球丛截面，多个独立向量对应正交结构群约化；首截面障碍与传递定义的 Euler 类按本书约定相差负号，消失性给出低维底空间上的完整存在判据。
\end{chaptersummary}

\BookExercises

\begin{exercise}\label{g2:ex:torus-extension-commutator}
	将环面 \(T^2\) 表示为在 \(S^1\vee S^1\) 上沿交换子附着一个二胞腔。给定整数 \(a,b\)，在两个圆周上分别取绕数为 \(a,b\) 的映射到 \(S^1\)，证明它们可以延拓到整个环面，并写出一个显式延拓。若把目标换为 \(S^1\vee S^1\)，让两个圆周分别映到目标的两个标准生成圆周，则能否延拓？比较两个问题的区别。
\end{exercise}
\begin{solution}
	设环面的两条一胞腔为 \(u,v\)，二胞腔沿 \(uvu^{-1}v^{-1}\) 附着。目标 \(S^1\) 的基本群是交换群 \(\mathbb Z\)，所以边界复合的同伦类为 \(a+b-a-b=0\)，单胞腔延拓判据给出延拓。将环面写为 \(S^1\times S^1\)，映射
	\[
		f(z,w)=z^a w^b
	\]
	就是一个显式选择；在 \(S^1\times\{1\}\) 与 \(\{1\}\times S^1\) 上分别恢复给定幂映射。如果原给定映射只与幂映射有相同绕数，先在一骨架上取相应同伦，再用 CW 同伦延拓性质把末端边界换回原给定映射。

	目标 \(S^1\vee S^1\) 的基本群是两个生成元 \(x,y\) 的自由群。二胞腔边界的像为 \(xyx^{-1}y^{-1}\)，这是非空约化字，故不是单位元，不能延拓。两个问题中的边界在一次同调中都为零，因为同调只记录交换化；能否填圆盘却由实际基本群中的类决定。非交换的一次障碍因而不能只用整数绕数或普通一次同调替代。
\end{solution}

\begin{exercise}\label{g2:ex:reflection-circle-sections}
	令
	\[
		E=(S^1\times[0,1])/((z,1)\sim(\bar z,0))\longrightarrow S^1
	\]
	为以复共轭粘合的圆周丛。证明它有截面，并直接计算截面在纤维内同伦下的类数。再算 \(H^1(S^1;\Pi_1(p))\)，比较两个答案。
\end{exercise}
\begin{solution}
	常值纤维坐标 \(1\) 与 \(-1\) 都满足粘合条件，因此给出截面。任一截面可写成 \(t\mapsto[u(t),t]\)，其中 \(u:[0,1]\to S^1\) 满足 \(u(0)=\overline{u(1)}\)。区间可缩，故可取实值辐角提升 \(u(t)=e^{\mathrm i\theta(t)}\)。端点条件等价于
	\[
		\theta(0)+\theta(1)=2\pi k\qquad(k\in\mathbb Z).
	\]
	改变辐角提升会使 \(\theta\) 整体增加 \(2\pi\ell\)，从而把 \(k\) 改为 \(k+2\ell\)。因此 \(k\) 的奇偶性与提升选择无关。

	若有截面同伦，则相应的 \(u(t,s)\) 在方形上可取连续实值辐角提升。\(\bigl(\theta(0,s)+\theta(1,s)\bigr)/(2\pi)\) 是随 \(s\) 连续的整数，故保持不变，说明奇偶性为同伦不变量。反过来，先用线性插值把 \(\theta(t)\) 变为连接原两端值的仿射函数，再让初端值从 \(\theta(0)\) 移到 \(\pi k\)，末端值始终取 \(2\pi k\) 减初端值。这些同伦都保留端点条件，最后得到常值函数 \(\pi k\)。所以每张截面同伦于常值纤维坐标 \((-1)^k\)，且这两类不同，恰有两类。

	复共轭在 \(\pi_1(S^1)=\mathbb Z\) 上作用为乘 \(-1\)，故 \(\Pi_1(p)\) 是反号局部系统。圆周的局部系数胞腔上链为 \(\mathbb Z\xrightarrow{-2}\mathbb Z\)，因此 \(H^1(S^1;\Pi_1(p))=\mathbb Z/2\)。这与两类截面的差异一致；若误用常值系数，就会得到 \(\mathbb Z\)，错误地预测无穷多个差异类。
\end{solution}

\begin{exercise}\label{g2:ex:homology-invisible-extension-obstruction}
	考虑 Hopf 映射 \(h:S^3\to S^2\)。证明它在全部正次整数同调上诱导零映射，却不能延拓到 \(D^4\)。说明这与第9章 Hurewicz 定理的范围如何相容。
\end{exercise}
\begin{solution}
	\(S^3\) 的正次同调只在三次为 \(\mathbb Z\)，但 \(H_3(S^2)=0\)，因此所有正次诱导同调映射都是零。由\cref{g2:thm:hopf-infinite-order}，Hopf 映射在 \(\pi_3(S^2)\) 中具有无穷阶，尤其非零。单胞腔延拓判据遂说明它不能填成 \(D^4\to S^2\)。

	\(S^2\) 是一连通空间，Hurewicz 定理保证首个非零次数的 \(\pi_2(S^2)\to H_2(S^2)\) 为同构，并不把三次同伦群识别为三次同调。这里的障碍正位于下一次的 \(\pi_3(S^2)\)，所以同调看不见它并不违反该定理。延拓问题必须保留边界所在的具体同伦次数。
\end{solution}

\begin{exercise}\label{g2:ex:obstruction-cochain-versus-class}
	设 \(m\ge2\)，\(X=S^1\cup_m e^2\) 为沿度 \(m\) 映射附着二胞腔的空间。在一骨架上给定 \(s_k:S^1\to S^1\)、\(s_k(z)=z^k\)。计算其延拓障碍上链及在 \(H^2(X;\mathbb Z)\) 中的类，并判断保留整个一骨架或仅保留顶点时分别能否延拓。最后把固定一骨架的问题写成相对障碍类。
\end{exercise}
\begin{solution}
	取胞腔定向使二维边界为乘 \(m\)。附着圆周经 \(s_k\) 后的度为 \(mk\)，所以障碍上链是 \(c^2(s_k)=mk\in C^2(X;\mathbb Z)=\mathbb Z\)。绝对胞腔上链中 \(\delta^1:\mathbb Z\to\mathbb Z\) 为乘 \(m\)，故 \(H^2(X;\mathbb Z)=\mathbb Z/m\)，而 \([c^2(s_k)]=0\) 对全部 \(k\) 成立。

	保持 \(s_k\) 不变时，必须有实际的边界度 \(mk=0\)，即 \(k=0\)。若只保留顶点的值，则可将一胞腔上的映射重新选择为常值，再常值延拓到二胞腔；这对每个 \(k\) 都可做到。这里重新选择不要求与原 \(s_k\) 同伦，所以没有声称非零度的圆周映射能够同伦为常值。

	若取 \(A=X^1\)，相对上链只有二维的 \(\mathbb Z\)，没有一次上链可提供余边界。因而 \(H^2(X,A;\mathbb Z)=\mathbb Z\)，相对首障碍是整数 \(mk\) 本身。相对类为零恰好等价于原 \(s_k\) 能够延拓。相对条件改变的是允许调整的部分，也就改变了应使用的障碍群。
\end{solution}

\begin{exercise}\label{g2:ex:maps-to-sphere-top-dimension}
	设 \(X\) 为维数不超过 \(n\) 的 CW 复形，\(n\ge2\)，\(u\) 为 \(H^n(S^n;\mathbb Z)\) 的正向生成元。证明 \([X,S^n]\to H^n(X;\mathbb Z)\)、\([f]\mapsto f^*u\) 是双射。说明为什么删去维数条件后，\(X=S^3,n=2\) 就给出反例。
\end{exercise}
\begin{solution}
	将映射视为平凡 \(S^n\) 丛的截面。纤维 \(S^n\) 是 \((n-1)\) 连通的简单空间，且已有常值截面；\cref{g2:thm:section-torsor}的维数条件于是满足。以常值截面为原点，差异类给出同伦类与 \(H^n(X;\pi_n(S^n))\cong H^n(X;\mathbb Z)\) 的双射。

	还需识别这张双射。先将 \(f\) 在 \(X^{n-1}\) 上同伦到常值，再用同伦延拓性质得到整个 \(X\) 上的映射。每个正向 \(n\) 胞腔压掉边界后成为 \(S^n\)，差异上链在此处的值就是所得球面映射的度。\(f^*u\) 的胞腔代表在同一胞腔上的值也是这个度，所以两张对应相同。

	若取 \(X=S^3,n=2\)，则 \(H^2(S^3;\mathbb Z)=0\)，但常值映射与 Hopf 映射不是同伦的。事实上 \(S^2\) 单连通，自由同伦不再识别不同的带基点三次类，而 \(\pi_3(S^2)=\mathbb Z\)。更高维胞腔开始看到 \(S^2\) 的高阶同伦群，单个二次上同调类已经不能完成分类。
\end{solution}

\begin{exercise}\label{g2:ex:projective-plane-sphere-map}
	令 \(q:\mathbb RP^2\to\mathbb RP^2/\mathbb RP^1\cong S^2\) 为压掉一骨架的商映射。证明 \(q\) 不是零伦，却在全部正次整数同调上诱导零映射。求 \([\mathbb RP^2,S^2]\) 的元素个数，并用胞腔上链解释这个现象。
\end{exercise}
\begin{solution}
	\(\mathbb RP^2\) 的标准胞腔复形在零、一、二维各有一个胞腔，边界 \(\partial_2\) 为乘 \(2\)、\(\partial_1=0\)。因此 \(H_1(\mathbb RP^2;\mathbb Z)=\mathbb Z/2\)、\(H_2(\mathbb RP^2;\mathbb Z)=0\)。球面的一次同调为零，源空间的二次及更高同调为零，故所有正次 \(q_*\) 都为零。

	对偶上链中 \(\delta^1\) 为乘 \(2\)，于是 \(H^2(\mathbb RP^2;\mathbb Z)=\mathbb Z/2\)。适当选取商球面的定向，\(q\) 把唯一二胞腔以系数 \(1\) 映到球面的二胞腔，所以 \(q^*u\) 由上链 \(1\) 表示，正是模 \(2\) 的非零类。零伦映射的正次上同调拉回为零，因此 \(q\) 不是零伦。

	上一题适用于二维复形与 \(S^2\)，故 \([\mathbb RP^2,S^2]\) 恰有两个元素，分别由常值映射与 \(q\) 代表。此处的二次整数上同调来自一次同调的挠部分，即万有系数定理中的 \(\operatorname{Ext}(\mathbb Z/2,\mathbb Z)\)；它不能只通过二次同调映射读出。
\end{solution}

\begin{exercise}\label{g2:ex:disk-extension-choices}
	在 \(\partial D^n\) 上给定到 \(S^n\) 的常值映射，\(n\ge2\)。求全部延拓 \(D^n\to S^n\) 在固定边界的同伦下的分类，并说明为什么没有存在障碍仍可有无穷多种选择。若改选一个度为 \(r\) 的延拓作为参考，差异坐标怎样改变？
\end{exercise}
\begin{solution}
	常值映射本身提供延拓，所以存在性已经解决。每个延拓将边界送到同一点，故唯一地下降为带基点映射
	\[
		D^n/\partial D^n\cong S^n\longrightarrow S^n.
	\]
	固定边界的同伦也恰好下降为带基点同伦。球面度分类定理因此给出整数分类，任意整数都可由一个相应度的球面映射实现，再复合商映射得到延拓。

	从相对障碍看，\(H^{n+1}(D^n,\partial D^n;\mathbb Z)=0\)，因为相对 CW 对只有一个 \(n\) 胞腔；差异群却是 \(H^n(D^n,\partial D^n;\mathbb Z)=\mathbb Z\)。存在障碍与延拓的差异处在相邻两个次数，前者消失不强迫后者消失。

	以常值延拓为参考时，度为 \(k\) 的延拓坐标为 \(k\)。改用度 \(r\) 的延拓作参考，正负粘合给出的差异为 \(k-r\)。因此全部坐标平移了 \(-r\)，而任意两张延拓的差保持不变。
\end{solution}

\begin{exercise}\label{g2:ex:nonzero-obstruction-group-section}
	设 \(n\ge2\)，考察平凡球丛 \(S^{n+1}\times S^n\to S^{n+1}\)。计算其首障碍群与首障碍类，说明“障碍群不为零”不能推出“没有截面”。再解释 \(H^{n+1}(B;\Pi_n)=0\) 在逐层延拓中实际提供的保证。
\end{exercise}
\begin{solution}
	纤维 \(S^n\) 的首个非零同伦群为 \(\pi_n(S^n)=\mathbb Z\)，平凡丛的运输作用恒等。因此首障碍群为
	\[
		H^{n+1}(S^{n+1};\mathbb Z)=\mathbb Z.
	\]
	但任取 \(v_0\in S^n\)，映射 \(x\mapsto(x,v_0)\) 就是一张全局截面。将它限制到各骨架，所有胞腔边界都已有填充，故障碍上链与障碍类均为零。这是非零障碍群中的零元素。

	障碍群为零时，所有落在这个群中的障碍类都自动为零。于是对于任一当前骨架截面，都能在保持更低一层的条件下调整最高一层，再向下一骨架延拓。这个条件保证整批延拓问题在该步可解；针对一张具体截面，只需它对应的那个类为零。若还有更高维胞腔，则仍需继续处理后续次数。
\end{solution}

\begin{exercise}\label{g2:ex:postnikov-two-stage-maps}
	取正文中由 \(m u^2\)、\(\ell u^2\) 定义的三型 \(E_m,E_\ell\)。设 \(f_a:\mathbb CP^\infty\to\mathbb CP^\infty\) 满足 \(f_a^*u=a u\)，其中 \(a\in\mathbb Z\)。证明：存在覆盖 \(f_a\) 的映射 \(E_m\to E_\ell\)，且它在纤维 \(\pi_3\cong\mathbb Z\) 上诱导乘 \(b\)，当且仅当 \(bm=\ell a^2\)。纤维生成元均按正文的传递约定选取。
\end{exercise}
\begin{solution}
	先设所需映射存在。记两纤维的三次上同调生成元为 \(v_m,v_\ell\)，则纤维映射拉回 \(v_\ell\) 为 \(bv_m\)。传递的自然性给出
	\[
		bm u^2=\mathrm d_4(bv_m)=f_a^*(\mathrm d_4v_\ell)=f_a^*(\ell u^2)=\ell a^2u^2,
	\]
	故条件必要。

	反过来，设 \(bm=\ell a^2\)。取 \(h_b:K(\mathbb Z,4)\to K(\mathbb Z,4)\) 在唯一非零同伦群上诱导乘 \(b\)。映射 \(h_b\kappa_m\) 与 \(\kappa_\ell f_a\) 代表同一个四次上同调类，因此可表性定理给出带基点同伦 \(H:h_b\kappa_m\simeq\kappa_\ell f_a\)。在入向路径模型中，令
	\[
		(b_0,\lambda)\longmapsto\bigl(f_a(b_0),\,\overline{H(b_0,-)}*(h_b\lambda)\bigr).
	\]
	这里横线表示反向走过同伦道路，星号表示依次拼接。新道路从 \(\kappa_\ell f_a(b_0)\) 出发，先到 \(h_b\kappa_m(b_0)\)，再到基点，所以公式确实落在 \(E_\ell\) 中并覆盖 \(f_a\)。在基点纤维上，\(H\) 为常值道路，故诱导的环路映射是 \(\Omega h_b\)，在 \(\pi_3\) 上为乘 \(b\)。这证明条件充分。
\end{solution}

\begin{exercise}\label{g2:ex:postnikov-lift-torsor}
	设 \(n\ge2\)、\(G\) 为阿贝尔群，\(E_\kappa\to B\) 为 \(\kappa:B\to K(G,n+1)\) 的入向路径模型。给定 CW 复形上的映射 \(f:Y\to B\)，并假定它至少有一个提升。证明：在始终覆盖这张固定映射 \(f\) 的同伦关系下，提升类组成 \(H^n(Y;G)\) 的主齐性集，即这个群在提升类上自由且传递地作用。说明为何没有天然指定的零元。
\end{exercise}
\begin{solution}
	提升等价于连续选取道路 \(\lambda_y:\kappa f(y)\to*\)。固定一组选择 \(\lambda_y^0\)。任取连续环路族 \(\omega:Y\to\Omega K(G,n+1)\)，拼接 \(\lambda_y^0*\omega_y\) 给出另一个提升。环路族同伦给出覆盖 \(f\) 的提升同伦，因此这是同伦类上的作用。

	对任意另一提升 \(\lambda_y\)，取基点环路 \(\omega_y=\overline{\lambda_y^0}*\lambda_y\)。插入的往返道路 \(\lambda_y^0*\overline{\lambda_y^0}\) 可保持端点收缩，故 \(\lambda_y^0*\omega_y\) 与 \(\lambda_y\) 在覆盖 \(f\) 的条件下同伦。这证明传递性。如果 \(\lambda^0*\omega\) 与 \(\lambda^0*\omega'\) 在覆盖 \(f\) 的条件下同伦，在每条道路前接上 \(\overline{\lambda^0}\) 并收缩往返部分，就得到 \(\omega\simeq\omega'\)，所以作用自由。

	环路空间的 CW 同伦型为 \(K(G,n)\)，可表性定理把 \([Y,\Omega K(G,n+1)]\) 及环路拼接识别为 \(H^n(Y;G)\) 及其加法。固定 \(\lambda^0\) 后，提升类因而可与这个群一一对应；更换初始提升使该对应整体平移。提升问题本身只给出各提升之间的差，没有指定哪一个应当成为零元。这里固定的是整个基底映射；允许基底映射一起移动会产生另一种分类问题。
\end{solution}

\begin{exercise}\label{g2:ex:postnikov-cp-two-sphere-lift}
	设 \(u\) 为 \(H^2(\mathbb CP^2;\mathbb Z)\) 的标准生成元，\(f_a:\mathbb CP^2\to P_2S^2\simeq\mathbb CP^\infty\) 表示 \(a u\)，其中 \(a\in\mathbb Z\)。求哪些 \(f_a\) 可以提升到 \(P_3S^2\)，哪些可以提升到 \(S^2\)。到 \(S^2\) 的提升允许三角形交换到同伦。用结果说明，只知道目标的二次阶段为何不足以解决四维复形上的映射问题。
\end{exercise}
\begin{solution}
	正文已证明，适当选择纤维生成元后，\(k_4(S^2)=u^2\in H^4(\mathbb CP^\infty;\mathbb Z)\)。因此
	\[
		f_a^*k_4(S^2)=a^2u^2\in H^4(\mathbb CP^2;\mathbb Z).
	\]
	这个四次上同调群由 \(u^2\) 自由生成，所以上式为零当且仅当 \(a=0\)。式\eqref{g2:eq:postnikov-lift-obstruction}遂说明，恰好 \(f_0\) 可以提升到第三阶段。

	若存在提升到 \(S^2\) 的映射，继续复合 \(S^2\to P_3S^2\)，就得到第三阶段的提升，因此仍必须 \(a=0\)。当 \(a=0\) 时，可表性定理说明 \(f_0\) 零伦，故取到 \(S^2\) 的常值映射即可得到交换到同伦的提升。

	二次阶段只检测 \(H^2\)，于是每个整数 \(a\) 都产生一张到 \(P_2S^2\) 的映射。四胞腔还会检测下一阶段的粘合类，它通过平方 \(a^2u^2\) 排除了全部非零 \(a\)。同一结论也可由杯积直接看见：\(S^2\) 的二次类平方为零，任何实际来自 \(S^2\) 的拉回类都必须保留这个关系。
\end{solution}

\begin{exercise}\label{g2:ex:annulus-relative-unit-sections}
	设 \(A=[0,1]\times S^1\)，用 \((r,e^{\mathrm i\theta})\) 表示其点。在平凡定向平面丛 \(A\times\mathbb R^2\) 中，要求单位截面在两个边界上分别为 \(e^{\mathrm i k\theta}\) 和 \(e^{\mathrm i\ell\theta}\)，其中 \(k,\ell\in\mathbb Z\)，两条边界都使用同一个 \(\theta\) 方向。
	\begin{enumerate}
		\item 证明这样的单位截面存在当且仅当 \(k=\ell\)。
		\item 当 \(k=\ell\) 时，证明固定两个边界的截面同伦类恰由整数 \(n\) 分类，并可取代表
		\[
			s_n(r,e^{\mathrm i\theta})=e^{\mathrm i(k\theta+2\pi nr)}.
		\]
	\end{enumerate}
\end{exercise}
\begin{solution}
	1. 单位截面可视为映射 \(f:A\rightarrow S^1\)。两个边界包含映射在 \(A\) 中通过改变 \(r\) 同伦，所以其与 \(f\) 的复合有相同绕数，必要条件为 \(k=\ell\)。反之，\(f(r,e^{\mathrm i\theta})=e^{\mathrm i k\theta}\) 就给出一个延拓。

	2. 给一个满足边界条件的 \(f\)，令 \(g(r,e^{\mathrm i\theta})=f(r,e^{\mathrm i\theta})e^{-\mathrm i k\theta}\)。它在两条边界上恒为 \(1\)。圆环的基本群由内边界生成，故 \(g_*:\pi_1(A)\rightarrow\pi_1(S^1)\) 为零。\cref{g2:prop:covering-lift-criterion}用于 \(t\mapsto e^{2\pi\mathrm it}\)，给出唯一的连续提升 \(G:A\rightarrow\mathbb R\)，使 \(G\) 在内边界恒为 \(0\)。外边界上 \(G\) 为整数值连续函数，因而恒为某个 \(n\in\mathbb Z\)。

	将 \(G\) 线性同伦到 \((r,e^{\mathrm i\theta})\mapsto nr\)，内外边界分别一直保持 \(0,n\)。取指数并乘回 \(e^{\mathrm i k\theta}\)，就得到固定边界的 \(f\simeq s_n\)。若 \(f\) 作固定边界的同伦，则相应提升仍可在内边界规范为零；外边界的整数值在同伦参数下连续，故不能变化。因此不同 \(n\) 对应不同同伦类。这个整数也可沿任何一条由内边界到外边界的径向道路读取：它是消去 \(e^{\mathrm i k\theta}\) 后所得闭路的绕数。
\end{solution}

\begin{exercise}\label{g2:ex:double-mobius-frame}
	令 \(L\rightarrow S^1\) 为 Möbius 线丛。构造 \(L\oplus L\) 的一个整体正交标架，说明一个没有无零截面的线丛与自身直和后可以成为平凡丛，并用纤维的连通性解释这与 \(L\) 的截面障碍为何相容。
\end{exercise}
\begin{solution}
	将 \(L\oplus L\) 写为 \(([0,1]\times\mathbb R^2)/((1,z)\sim(0,-z))\)，Euclid 内积可下降到商。在区间上取
	\[
		v_1(t)=(\cos\pi t,\sin\pi t),\qquad v_2(t)=(-\sin\pi t,\cos\pi t).
	\]
	它们处处正交且长为 \(1\)，并满足 \(v_j(1)=-v_j(0)\)，因此恰好满足商中的端点粘合条件，下降为全局正交标架。映射 \((b,a_1,a_2)\mapsto a_1v_1(b)+a_2v_2(b)\) 给出 \(\varepsilon_{S^1}^2\cong L\oplus L\)。

	单个 \(L\) 的单位纤维是两点集，运输交换这两点，故不能沿一胞腔连续连接端点。直和后单位纤维变成道路连通的圆周，半圈旋转正好连接相反端点，而底空间没有二胞腔来产生下一步障碍。构造的 \(v_1,v_2\) 将这个连通性解释落实成了具体标架。
\end{solution}

\begin{exercise}\label{g2:ex:sphere-one-vector-field-parity}
	对 \(n\ge1\)，证明 \(S^n\) 有处处非零连续切向量场当且仅当 \(n\) 为奇数。偶数维时，从假设存在的单位切向量场构造恒等映射到对径映射的同伦；奇数维时，将环境 Euclid 空间看成复向量空间，直接写出一个光滑单位切向量场。
\end{exercise}
\begin{solution}
	设 \(n\) 为偶数，并假定有无零切向量场。归一化得到 \(s:S^n\rightarrow\mathbb R^{n+1}\)，满足 \(\|s(x)\|=1\) 及 \(\langle s(x),x\rangle=0\)。于是
	\[
		H(x,t)=\cos(\pi t)x+\sin(\pi t)s(x)
	\]
	的范数平方为 \(\cos^2(\pi t)+\sin^2(\pi t)=1\)，所以它是球面内的连续同伦，从 \(H(x,0)=x\) 到 \(H(x,1)=-x\)。\cref{g2:prop:degree-basic,g2:prop:reflection-antipodal-degree}却给出两端的度分别为 \(1\) 与 \((-1)^{n+1}=-1\)，矛盾。

	若 \(n=2q+1\)，将 \(\mathbb R^{2q+2}\) 识别为 \(\mathbb C^{q+1}\)，取 \(s(z)=\mathrm iz\)。复数乘 \(\mathrm i\) 保持实内积范数，并满足 \(\langle\mathrm iz,z\rangle=0\)，故 \(s\) 在单位球面上为光滑单位切向量场。由\cref{g2:prop:sections-orthogonal-reduction}，它还给出 \(\mathrm TS^{2q+1}\cong\varepsilon^1\oplus W\) 的一个分解；这里得到的是一条平凡直和因子。
\end{solution}

\begin{exercise}\label{g2:ex:tautological-plane-circle-sections}
	令 \(L\to\mathbb RP^2\) 为典范实线丛，取带自然内积的秩二丛 \(V=\varepsilon^1\oplus L\)。证明 \(S(V)\) 有截面，并计算其截面在纤维内同伦下的类数。提示：在双覆盖 \(S^2\to\mathbb RP^2\) 上，\(V\) 的纤维坐标变换为 \((a,b)\mapsto(a,-b)\)。
\end{exercise}
\begin{solution}
	平凡因子的单位向量给出两张截面 \(s_+=(1,0)\)、\(s_-=(-1,0)\)，所以存在障碍为零。沿 \(\mathbb RP^2\) 的非平凡基本群元素运输时，单位圆周经反射 \((a,b)\mapsto(a,-b)\)，它在 \(\pi_1(S^1)=\mathbb Z\) 上作用为乘 \(-1\)。因此差异类使用反号整数局部系统 \(\mathbb Z_{\mathrm{sgn}}\)。

	在标准提升胞腔基中，\(\mathbb RP^2\) 的一次边界为 \(g-1\)，二次边界为 \(1+g\)，其中 \(g\) 是覆盖变换。代入反号作用，局部系数上链成为
	\[
		\mathbb Z\xrightarrow{-2}\mathbb Z\xrightarrow{0}\mathbb Z.
	\]
	故 \(H^1(\mathbb RP^2;\mathbb Z_{\mathrm{sgn}})=\mathbb Z/2\)。纤维 \(S^1=K(\mathbb Z,1)\) 没有更高同伦群，\cref{g2:thm:section-torsor}遂给出恰有两个截面同伦类。

	还可辨认两张显式代表。限制到 \(\mathbb RP^1\cong S^1\)，典范线丛成为 Möbius 线丛，单位圆周丛正是\cref{g2:ex:reflection-circle-sections}的反射粘合模型。该题的辐角奇偶性将 \(s_+\) 与 \(s_-\) 分开，因此它们在整个底空间上也不能纤维同伦，恰好代表两类。
\end{solution}


% !TeX root = ../../Book2.tex
% 纲要标题：## 第13章*　上同调运算与模2方法
% 纲要位置：几何学计划纲要.md，第 1501 行。
% 正文状态：未撰写；本文件目前是框架。

% 纲要内容开始
% 
% > 本章目标与证明范围：把 Steenrod 运算从障碍章分离。选定模2的链级构造、Cartan 公式和低阶关系为完整模块；一般 Adem 关系完整陈述并注明外引；一般稳定同伦计算不在本章暗中承诺。
% 
% 纲要内容结束

\OptionalChapter{上同调运算与模2方法}{g2:ch:13}
\BookChapterNumber{g2:ch:13}

\begin{chapterabstract}
待撰写
\end{chapterabstract}

\begin{learninggoals}
待撰写
\end{learninggoals}

待撰写

% !TeX root = ../../Book2.tex
% 纲要标题：### 13.1 自然运算及基本例子
% 纲要位置：几何学计划纲要.md，第 1455 行。
% 正文状态：未撰写；本文件目前是框架。

% 纲要内容开始
% 
% * 系数短正合列与 Bockstein
% * 稳定、非稳定与次数
% * 可表性和 Eilenberg–MacLane 空间
% 
% 纲要内容结束

\section{自然运算及基本例子}
\label{g2:sec:1301}

待撰写

\subsection{系数短正合列与 Bockstein}
\label{g2:subsec:1301:01}
待撰写

\subsection{稳定、非稳定与次数}
\label{g2:subsec:1301:02}
待撰写

\subsection{可表性和 Eilenberg–MacLane 空间}
\label{g2:subsec:1301:03}
待撰写
% 纲要条目排为编号小节标题，并进入目录与书签；说明段仍保留于纲要注释中。

% !TeX root = ../../Book2.tex
% 纲要标题：### 13.2 Steenrod 平方的构造
% 纲要位置：几何学计划纲要.md，第 1511 行。
% 正文状态：未撰写；本文件目前是框架。

% 纲要内容开始
% 
% * 链对角的高阶同伦
% * cup-i 乘积及边界恒等式
% * 平方运算与余链选择无关性
% * 自然性、单位与不稳定条件
% 
% 纲要内容结束

\section{Steenrod 平方的构造}
\label{g2:sec:1302}

待撰写

\subsection{链对角的高阶同伦}
\label{g2:subsec:1302:01}
待撰写

\subsection{cup-i 乘积及边界恒等式}
\label{g2:subsec:1302:02}
待撰写

\subsection{平方运算与余链选择无关性}
\label{g2:subsec:1302:03}
待撰写

\subsection{自然性、单位与不稳定条件}
\label{g2:subsec:1302:04}
待撰写
% 纲要条目排为编号小节标题，并进入目录与书签；说明段仍保留于纲要注释中。

% !TeX root = ../../Book2.tex
% 纲要标题：### 13.3 基本关系和可计算后果
% 纲要位置：几何学计划纲要.md，第 1468 行。
% 正文状态：未撰写；本文件目前是框架。

% 纲要内容开始
% 
% * Cartan 公式的链级证明
% * 低阶关系的代表性证明
% * 一般 Adem 关系及其外部证明输入
% * 实射影空间上的 Steenrod 平方
% * 附着映射和 Hopf 不变量的约束
% 
% 纲要内容结束

\section{基本关系和可计算后果}
\label{g2:sec:1303}

待撰写

\subsection{Cartan 公式的链级证明}
\label{g2:subsec:1303:01}
待撰写

\subsection{低阶关系的代表性证明}
\label{g2:subsec:1303:02}
待撰写

\subsection{一般 Adem 关系及其外部证明输入}
\label{g2:subsec:1303:03}
待撰写

\subsection{实射影空间上的 Steenrod 平方}
\label{g2:subsec:1303:04}
待撰写

\subsection{附着映射和 Hopf 不变量的约束}
\label{g2:subsec:1303:05}
待撰写
% 纲要条目排为编号小节标题，并进入目录与书签；说明段仍保留于纲要注释中。

% !TeX root = ../../Book2.tex
% 纲要标题：### 13.4 与示性类的后续比较
% 纲要位置：几何学计划纲要.md，第 1526 行。
% 正文状态：初稿；尚未正式审核。

% 纲要内容开始
% 
% 本节是已建运算的后续应用：先准确预告 Thom 类上的结论，第20—21章完成 Thom 与示性类后再验证；不以这些后置结论反向证明平方运算。
% 
% * Thom 类上的运算
% * Stiefel–Whitney 类的另一解释
% * Wu 公式的所选版本和范围
% 
% ---
% 
% 纲要内容结束

\section{与示性类的后续比较}
\label{g2:sec:1304}

对一张向量丛，每个纤维的单位圆盘都贡献一个最高次相对上同调生成元。若把这些生成元组成整体的 Thom 类，再对它施加 Steenrod 平方，所得信息就能测量向量丛的粘合。本节先写明这项比较所用的对象与结论，再从已经建立的平方运算和 Poincaré 对偶计算 Wu 类。Thom 类的整体存在与同构在第20章证明，与 Stiefel--Whitney 类及切丛的识别在第21章完成。

\subsection{Thom 类上的运算}
\label{g2:subsec:1304:01}

设 \(V\to B\) 为有限 CW 复形上的实秩 \(r\) 向量丛，\(r\ge1\)，并选取连续正定内积。它的单位圆盘丛与单位球丛分别记为
\[
	D(V)=\{v\in V:\|v\|\le1\},\qquad S(V)=\{v\in V:\|v\|=1\}.
\]
将全部单位球面同时压成一个点，得到\term[Thom-kongjian]{Thom 空间}[Thom space]\footnote{René Thom（1923--2002），法国数学家，发展了配边理论，并研究奇点及突变理论。}
\[
	\operatorname{Th}(V)\coloneqq D(V)/S(V).
\]
每个纤维圆盘的边界被压成这个共同基点，因而纤维在商中成为一个带基点 \(S^r\)。这个构造保留了纤维在底空间上如何粘合，同时让纤维的最高次相对类能够看作商空间上的约化类。单位球丛有径向领圈，故空间对的上同调与商的约化上同调相同：
\[
	H^q(D(V),S(V);\mathbb F_2)\cong\widetilde H^q(\operatorname{Th}(V);\mathbb F_2).
\]
下文优先使用左侧，这样底空间类经投影拉回后，可以直接与相对类相乘。

\begin{definition}[模二 Thom 类]\label{g2:def:mod-two-thom-class}
	设 \(V\to B\) 为带连续正定内积的实秩 \(r\) 向量丛，\(r\ge1\)。若
	\[
		U_V\in H^r(D(V),S(V);\mathbb F_2)
	\]
	限制到每个纤维对 \((D(V_b),S(V_b))\cong(D^r,S^{r-1})\) 时，都是其 \(r\) 次上同调的非零元素，就称 \(U_V\) 为 \(V\) 的模二\term[Thom-lei]{Thom 类}[Thom class]。
\end{definition}

模二系数下，一个一维向量空间只有一个非零元素，所以反转纤维定向不改变这个局部要求。这正是任意实向量丛都能具有模二 Thom 类的原因；整数版本则必须保留纤维定向及其运输。

\ProseParagraphBreak

下面先陈述具体结论，完整证明见\cref{g2:thm:thom-isomorphism,g2:prop:thom-naturality-sum}。此处固定有限 CW 底空间，已足以覆盖本节的例子；更一般底空间的控制随 Thom 同构的完整证明一起处理。

\begin{theorem}[模二 Thom 同构]\label{g2:thm:thom-mod-two-preview}
	设 \(V\to B\) 为有限 CW 复形上的实秩 \(r\) 向量丛，\(r\ge1\)，并选取连续正定内积。则存在唯一的模二 Thom 类 \(U_V\)，且投影 \(p:D(V)\to B\) 给出同构
	\begin{equation}\label{g2:eq:thom-mod-two-isomorphism}
		\Phi_V:H^q(B;\mathbb F_2)\xrightarrow{\cong}H^{q+r}(D(V),S(V);\mathbb F_2),\qquad a\longmapsto p^*a\smile U_V.
	\end{equation}
	Thom 类及此同构与向量丛的拉回相容；两个丛的外积 Thom 类是各自 Thom 类的上同调外积。
\end{theorem}

Thom 同构的局部内容是圆盘对的相对上同调只有最高次非零；整体证明须把局部生成元粘合，并证明杯积映射在粘合后仍为同构。这两步由\secref{g2:sec:2002}分别完成。所述版本及其与示性类的关系可对照\cite[第23章第5--6节]{May1999}。

\ProseParagraphBreak

使用上述 Thom 同构，相对平方运算把 \(U_V\) 送到 \(H^{r+i}(D(V),S(V);\mathbb F_2)\)，所以唯一存在一个 \(i\) 次底空间类，使
\begin{equation}\label{g2:eq:thom-square-identity}
	\operatorname{Sq}^i U_V=p^*w_i^{\mathrm{Sq}}(V)\smile U_V,
	\qquad w_i^{\mathrm{Sq}}(V)=\Phi_V^{-1}(\operatorname{Sq}^i U_V).
\end{equation}
上标暂时标明这组类由平方运算产生；它与射影丛构造的 Stiefel--Whitney 类的识别，完整证明见\cref{g2:thm:thom-square-comparison}。\(U_V\) 的次数是 \(r\)，所以 \(\operatorname{Sq}^i U_V\) 的次数为 \(r+i\)，经 \(\Phi_V^{-1}\) 后恰好降回 \(i\)。这个次数检查也说明，式中参与运算的是总空间对上的 \(U_V\)，而 \(w_i^{\mathrm{Sq}}\) 位于底空间。

\ProseParagraphBreak

对平凡丛 \(\varepsilon_B^r=B\times\mathbb R^r\)，Thom 类是 \(1\times u\)，其中 \(u\in H^r(D^r,S^{r-1};\mathbb F_2)\) 为生成元。正次数平方在 \(1\) 上为零，在 \(u\) 上也为零，因为目标相对上同调已高于 \(r\)。相对外积的 Cartan 公式于是给出
\[
	w_0^{\mathrm{Sq}}(\varepsilon_B^r)=1,\qquad w_i^{\mathrm{Sq}}(\varepsilon_B^r)=0\quad(i>0).
\]
一般丛的 \(\operatorname{Sq}^i U_V\) 则可能非零：每张纤维上看不到的高次类，能够通过底空间方向进入总空间对的上同调。这使平方运算能够读出丛的整体信息。

\subsection{Stiefel–Whitney 类的另一解释}
\label{g2:subsec:1304:02}

第21章要构造的\term[Stiefel-Whitney-lei]{Stiefel--Whitney 类}[Stiefel--Whitney class]\footnote{Eduard Stiefel（1909--1978），瑞士数学家，早期研究拓扑学，后从事数值分析与科学计算。Whitney 的简介见本卷 Alexander--Whitney 映射处。}是一组随实向量丛自然变化的模二上同调类。式\eqref{g2:eq:thom-square-identity}给出它们的运算解释。为看清两种构造为何应当一致，先使用 Thom 同构及其相容性，直接推出这些类的基本规则。

\begin{proposition}[Thom 平方类的基本性质]\label{g2:prop:thom-square-properties}
	对有限 CW 复形上的实向量丛，按式\eqref{g2:eq:thom-square-identity}定义 \(w_i^{\mathrm{Sq}}\)。则对秩 \(r\ge1\) 的丛 \(V\)，有
	\[
		w_0^{\mathrm{Sq}}(V)=1,\qquad w_i^{\mathrm{Sq}}(V)=0\quad(i>r).
	\]
	对有限 CW 复形之间的映射 \(f:B'\to B\) 及同一底空间上的两个正秩实向量丛 \(V,W\)，有
	\begin{align*}
		w_i^{\mathrm{Sq}}(f^*V)&=f^*w_i^{\mathrm{Sq}}(V),\\
		w_k^{\mathrm{Sq}}(V\oplus W)&=\sum_{i+j=k}w_i^{\mathrm{Sq}}(V)\smile w_j^{\mathrm{Sq}}(W).
	\end{align*}
	特别地，添加平凡直和因子不改变这些类。
\end{proposition}
\begin{proof}
	\(\operatorname{Sq}^0U_V=U_V\) 及 \(\operatorname{Sq}^iU_V=0\)（\(i>r\)）结合 \(\Phi_V\) 单射，立即给出前两条。

	给拉回丛使用拉回内积，令 \(F:(D(f^*V),S(f^*V))\to(D(V),S(V))\) 为丛映射。其每张纤维映射都是所给线性识别，因此 \(F^*U_V=U_{f^*V}\)。用平方的自然性和 \(pF=fp'\)，得到
	\[
		\operatorname{Sq}^iU_{f^*V}=F^*\operatorname{Sq}^iU_V
		=(p')^*f^*w_i^{\mathrm{Sq}}(V)\smile U_{f^*V}.
	\]
	再用拉回丛的 Thom 同构单射，即得自然性。

	先将两丛放在不同底空间上，考虑外直和丛 \(V\boxplus W\to B\times C\)。使用纤维中的最大范数，其圆盘对可写成
	\[
		\bigl(D(V)\times D(W),\ S(V)\times D(W)\cup D(V)\times S(W)\bigr).
	\]
	最大范数与 Euclid 直和范数之间的径向变换给出相同的圆盘对同伦型。外积 \(U_V\times U_W\) 限制到每张纤维都是最高次相对生成元，所以按 Thom 类的唯一性，它就是外直和的 Thom 类。\cref{g2:thm:steenrod-cartan}的相对外积形式于是给出
	\[
		\operatorname{Sq}^k(U_V\times U_W)
		=\sum_{i+j=k}\operatorname{Sq}^iU_V\times\operatorname{Sq}^jU_W
		=\Phi_{V\boxplus W}\!\left(\sum_{i+j=k}w_i^{\mathrm{Sq}}(V)\times w_j^{\mathrm{Sq}}(W)\right).
	\]
	由 Thom 同构单射可读出外直和公式。同底空间的 \(V\oplus W\) 是沿对角映射 \(B\to B\times B\) 拉回 \(V\boxplus W\)，自然性将外积变成杯积，便得到所列公式。最后取 \(W\) 为平凡丛，使用上一小节算出的类，得到稳定性。
\end{proof}

这些规则分别来自平方的恒等项、不稳定条件、自然性与 Cartan 公式。这里的证明把两套性质逐项对应起来；\cref{g2:thm:thom-square-comparison}进一步证明线丛归一化，再用分裂原理得到唯一性，从而去掉上标 \(\mathrm{Sq}\)。

\begin{example}[Möbius 线丛的 Thom 平方]\label{g2:exm:mobius-thom-square}
	设 \(L\to S^1\) 为 Möbius 线丛。它的圆盘丛是 Möbius 带，球丛是边界圆周。将边界压成一点后，得到实射影平面的同伦型：圆盘丛收缩到中心圆周，而边界到中心的投影是二重覆盖，所以其同伦余纤维是 \(S^1\cup_2 e^2=\mathbb RP^2\)。

	由上一节的射影空间环计算，\(\widetilde H^*(\operatorname{Th}(L);\mathbb F_2)\) 在一次、二次由 \(t,t^2\) 生成，且 \(t^2\ne0\)。Thom 类必须是唯一非零的一次类 \(U_L=t\)，所以
	\[
		\operatorname{Sq}^1U_L=U_L^2=t^2\ne0.
	\]
	因此，若用 Thom 同构把它搬回 \(S^1\)，所得 \(w_1^{\mathrm{Sq}}(L)\) 就是 \(H^1(S^1;\mathbb F_2)\) 的非零元素。这与平凡线丛的零值相区别。第21章将把这个类与纤维定向绕基底一周后是否反转对应起来。
\end{example}

这个例子显示，相对类的平方可以非零，即使底空间本身没有二次上同调。二次类在 Thom 空间中，Thom 同构再将其次数减一；若一开始就把 \(U_L\) 当成底空间上的一次类，就会错误地令它的平方消失。

\subsection{Wu 公式的所选版本和范围}
\label{g2:subsec:1304:03}

还有一种方法把平方运算组织成几何类：在闭流形上先把 \(\operatorname{Sq}^i x\) 与基本类配对，再用 Poincaré 对偶把这个数值泛函表示为杯积。本小节采用的是将切丛类写成 Wu 类之平方的版本。

\begin{definition}[Wu 类]\label{g2:def:wu-class}
	设 \(M\) 为连通闭光滑 \(n\) 维流形，\([M]_2\in H_n(M;\mathbb F_2)\) 为模二基本类。对 \(0\le i\le n\)，满足
	\begin{equation}\label{g2:eq:wu-defining-pairing}
		\langle v_i(M)\smile x,[M]_2\rangle
		=\langle\operatorname{Sq}^i x,[M]_2\rangle
		\qquad\bigl(x\in H^{n-i}(M;\mathbb F_2)\bigr)
	\end{equation}
	的唯一类 \(v_i(M)\in H^i(M;\mathbb F_2)\)，称为第 \(i\) 个\term[Wu-lei]{Wu 类}[Wu class]。令 \(v(M)=\sum_{i=0}^n v_i(M)\)，称为总 Wu 类。
\end{definition}

这些类由吴文俊\footnote{吴文俊（1919--2017），中国数学家，研究拓扑学中的示性类与嵌入问题，并开创数学机械化方法。}引入。它们的存在唯一性已经可以从本卷前面的结果证明。

\begin{proposition}[Wu 类的确定与消失范围]\label{g2:prop:wu-existence-vanishing}
	对连通闭光滑 \(n\) 维流形 \(M\)，式\eqref{g2:eq:wu-defining-pairing}唯一确定 \(v_i(M)\)。此外，\(v_0(M)=1\)，且 \(v_i(M)=0\) 对 \(i>n/2\) 成立。
\end{proposition}
\begin{proof}
	模二系数下每个流形都有规范的局部定向，因此\cref{g2:cor:closed-poincare-duality}可用于 \(M\)，无须另加整数可定向条件。结合域上的万有系数同构，杯积给出完美配对
	\[
		H^i(M;\mathbb F_2)\otimes H^{n-i}(M;\mathbb F_2)\longrightarrow\mathbb F_2,
		\qquad a\otimes x\longmapsto\langle a\smile x,[M]_2\rangle.
	\]
	由于 \(\operatorname{Sq}^i\) 可加，右侧的 \(x\mapsto\langle\operatorname{Sq}^i x,[M]_2\rangle\) 是线性泛函，完美配对恰好将它表示为唯一的 \(v_i\)。\(i=0\) 时 \(\operatorname{Sq}^0x=x\)，所以 \(v_0=1\)。若 \(i>n/2\)，则输入次数 \(n-i<i\)，不稳定条件使 \(\operatorname{Sq}^i x=0\) 对所有输入成立；完美配对便迫使 \(v_i=0\)。
\end{proof}

现在把这些完全由上同调运算确定的类与切丛联系起来。下面的公式在\cref{g2:thm:wu-formula}中由 Thom 平方与对角嵌入完整证明；其准确陈述也见\cite[第23章第3节]{May1999}。

\begin{theorem}[Wu 公式]\label{g2:thm:wu-tangent-preview}
	设 \(M\) 为连通闭光滑 \(n\) 维流形。它的切丛 Stiefel--Whitney 类与式\eqref{g2:eq:wu-defining-pairing}确定的 Wu 类满足
	\begin{equation}\label{g2:eq:wu-tangent-formula}
		w(\mathrm TM)=\operatorname{Sq}(v(M)),\qquad
		w_k(\mathrm TM)=\sum_{i+j=k}\operatorname{Sq}^i v_j(M).
	\end{equation}
	这里 \(\operatorname{Sq}=\sum_{i\ge0}\operatorname{Sq}^i\)，等式按上同调次数逐项解释。
\end{theorem}

这项比较把对角嵌入 \(M\hookrightarrow M\times M\) 的法丛识别为 \(\mathrm TM\)，再将其 Thom 类的平方写为对偶配对。\cref{g2:thm:wu-formula}利用第20章的 Thom 同构完成这部分几何证明。眼下已经可以独立算出 \(v(M)\) 与 \(\operatorname{Sq}(v(M))\)；公式说明这些计算怎样给出切丛类。

\begin{example}[射影平面的 Wu 类]\label{g2:exm:projective-plane-wu-classes}
	对 \(M=\mathbb RP^2\)，上一节给出 \(H^*(M;\mathbb F_2)=\mathbb F_2[t]/(t^3)\)，\(|t|=1\)，且 \(\langle t^2,[M]_2\rangle=1\)。\(\operatorname{Sq}^1t=t^2\)，故定义配对要求 \(\langle v_1t,[M]_2\rangle=1\)，于是 \(v_1=t\)。高于 \(n/2=1\) 的 Wu 类为零，所以
	\[
		v(\mathbb RP^2)=1+t,\qquad
		\operatorname{Sq}(v(\mathbb RP^2))=1+t+t^2.
	\]
	Wu 比较式因此将给出 \(w(\mathrm T\mathbb RP^2)=1+t+t^2\)。\(v_2=0\) 与 \(w_2=t^2\) 同时成立：总平方还把 \(v_1\) 推向更高次数。

	再取 \(M=\mathbb CP^2\)。由\cref{g2:thm:complex-projective-ring}及系数变换，\(H^*(M;\mathbb F_2)=\mathbb F_2[h]/(h^3)\)，\(|h|=2\)，且 \(\langle h^2,[M]_2\rangle=1\)。一次上同调为零，故 \(v_1=0\)；\(\operatorname{Sq}^2h=h^2\) 使 \(v_2=h\)，更高 Wu 类由消失范围为零。因此
	\[
		v(\mathbb CP^2)=1+h,\qquad
		\operatorname{Sq}(v(\mathbb CP^2))=1+h+h^2.
	\]
	其中 \(\operatorname{Sq}^1h=0\) 是因为三次上同调为零。相应的切丛结论为 \(w_2(\mathrm T\mathbb CP^2)=h\)、\(w_4(\mathrm T\mathbb CP^2)=h^2\)。两种射影平面的计算都只需生成元上的平方与顶次配对；整数可定向性不影响这套模二计算。
\end{example}

Wu 类的定义使用闭流形的普通基本类及完美配对。若有边界，Poincaré--Lefschetz 对偶会在一侧引入相对上同调；若非紧，则须选择紧支或局部有限版本。因此式\eqref{g2:eq:wu-tangent-formula}在本节按连通闭光滑流形陈述，不把同一组普通上同调公式直接套到带边界或非紧情形。对一般实向量丛，式\eqref{g2:eq:thom-square-identity}则是不同的陈述：它使用丛的 Thom 类，而没有使用底空间的流形基本类。
% 纲要条目排为编号小节标题，并进入目录与书签；说明段仍保留于纲要注释中。


\begin{chaptersummary}
待撰写
\end{chaptersummary}

\BookExercises
待撰写


% !TeX root = ../Book2.tex
% 纲要标题：# 第三编　横截性、相交与 Morse 理论
% 纲要位置：几何学计划纲要.md，第 1536 行。
% 正文状态：本编导读初稿；尚未正式审核。

% 纲要内容开始
% 
% 纲要内容结束

\BookPart{横截性、相交与 Morse 理论}{g2:part:03}
\BookPartGuide
两条曲线相切时，轻微移动就可能让一个交点变成两个，也可能让它消失；横截相交则能在足够小的扰动下保持局部形状。要把相交次数变成拓扑不变量，首先需要控制这种局部变化，再研究一族交点怎样组成带边界的流形。本编由光滑逼近和横截性开始，把代数拓扑中的映射度、对偶和胞腔边界解释为光滑几何中的计数与分解。

\ProseParagraphBreak

\chapref{g2:ch:14}说明如何扰动映射而保留指定区域的数据。Sard 定理使正则值可以被选出，参数横截性则把一族映射的总体性质转为大多数参数下的性质。这样得到的逆像具有预期维数，横截同伦的逆像还会提供配边。\chapref{g2:ch:15}处理相邻的构造问题：怎样把抽象流形嵌入 Euclid 空间，怎样用法丛描述子流形的整个邻域，以及怎样把子流形的运动延拓为环境运动。固有满淹没的局部平凡性也在这里证明；其核心是保证沿基底道路的全部提升能够走完所需时间。

\ProseParagraphBreak

有了横截性和邻域模型，\chapref{g2:ch:16}便能把正则值原像、子流形交点和向量场零点赋予符号。改变正则值或扰动时，点集连成一维配边；边界的符号关系给出计数的不变性。前面由基本类与杯积定义的配对，在这里得到横截代表的解释。不动点也可作为映射图像与对角线的交点处理，从而把局部指数之和与上同调上的交错迹联系起来。

\ProseParagraphBreak

\chapref{g2:ch:17}用一个光滑实函数逐层观察流形。次水平集在没有临界值的区间内沿梯度流变形；越过一个非退化临界点时，局部二次型的负方向决定附着把手的指标。把手给出胞腔同伦模型，临界点数因而受到同调的约束。\chapref{g2:ch:18}继续记录临界点之间的负梯度轨线，以轨线计数建立 Morse--Smale 同调。这里需要分析轨线断裂、紧性和粘合，才能把一维轨线空间的端点与微分平方为零对应起来。正文先用模二系数完成构造与比较，再讨论整数取向。

\ProseParagraphBreak

本编可从第一卷的光滑流形理论直接开始：\chapref{g2:ch:14}、\chapref{g2:ch:15}需要第一卷第11章的微分与子流形坐标、第12章的单位分解，以及第13章的常微分方程存在性和光滑依赖；Sard 证明还用到 Taylor 估计、零测集与 Fubini 定理。进入\chapref{g2:ch:16}时，需要回顾本卷\chapref{g2:ch:04}、\chapref{g2:ch:06}的乘积、帽积与对偶，映射度的同调与积分比较还分别使用\chapref{g2:ch:03}的映射度和\chapref{g2:ch:07}的 de Rham 同构。Morse 不等式和 Morse 同调的比较使用\chapref{g2:ch:03}的胞腔链。前一编的全部谱序列计算并非这些构造的共同前提。阅读时可以沿用同一张曲面：先选择横截曲线计算相交，再选高度函数观察把手附着，最后比较临界点之间的轨线与胞腔边界。这样能够具体辨认局部微分条件怎样约束整体拓扑。

% !TeX root = ../../Book2.tex
% 纲要标题：## 第14章　光滑逼近、Sard 与横截性
% 纲要位置：几何学计划纲要.md，第 1538 行。
% 正文状态：初稿；尚未正式审核。

% 纲要内容开始
% 
% > 本章目标与证明范围：先建立光滑逼近与函数空间的必要拓扑，再完成 Sard 和参数横截性。一般位置的结论具有明确正则性、紧性、相对保持和边界条件。
% 
% 纲要内容结束

\chapter{光滑逼近、Sard 与横截性}
\label{g2:ch:14}
\BookChapterNumber{g2:ch:14}

\begin{chapterabstract}
	方程的解集何时成为流形，以及不规则的相交能否通过微小改变变得规则，是本章的两个基本问题。先以卷积、单位分解和坐标内替换证明相对光滑逼近，说明紧源与非紧源上导数控制的含义；再以秩分层、覆盖估计和维数归纳证明光滑 Sard 定理。横截性将正则值推广到子流形约束，给出逆像的维数、切空间及法丛。参数定理把总体横截化为几乎处处切片横截，局部扰动进一步给出保持闭集附近数据的横截逼近。最后处理边界条件，并由横截同伦构造逆像间的配边。
\end{chapterabstract}

\begin{learninggoals}
	\begin{itemize}
		\item 用局部卷积与截断构造光滑逼近，说明相对条件如何在拼接中保持。
		\item 区分 \(C^0\)、\(C^1\) 及更高阶控制，辨认非紧源上的弱拓扑、强拓扑与紧支改动。
		\item 重建光滑 Sard 定理的证明，并说明临界点集与临界值集的差别及有限可微性的作用。
		\item 将横截条件化为法方向上的满射，计算横截逆像的切空间、维数及法丛。
		\item 从约束流形到参数空间的投影证明参数横截定理，区分稠密、剩余与几乎处处。
		\item 构造相对横截扰动，核对边界限制映射，并用紧性将同伦逆像变成配边。
	\end{itemize}
\end{learninggoals}

考虑平面中的方程 \(y=x^2\) 与直线 \(y=a\)。当 \(a>0\) 时，它们交于两点，每个交点的切线方向都不同；当 \(a<0\) 时没有交点；在 \(a=0\) 时，两点合并成一个相切点。交点数在这个特殊时刻改变，但把参数也加进去，全部解组成的集合 \(\{(a,x):a=x^2\}\) 始终是一条光滑曲线。微分拓扑经常利用这一转换：单个对象可能出现退化，放进一个有足够参数的光滑族后，可以研究哪些切片仍然规则，以及规则切片之间如何联系。

\ProseParagraphBreak

第一卷第11章的正则水平集定理已经说明：当方程组的微分满射时，解集具有预期维数。现在还要回答两个问题。第一，正则值是否足够多，以至于总能作很小的调整选到它？Sard 定理\footnote{Arthur Sard（1909—1980），美国数学家，研究临界值、线性逼近与样条。}给出“临界值为零测集”的答案。第二，约束若是一条曲线或一张曲面，而非一个目标点，什么条件代替满射？所需条件是映射微分的像与约束子流形的切空间共同张成整个目标切空间，即横截性。

\ProseParagraphBreak

本章使用第一卷第11章的微分、逆函数定理和子流形坐标，第12章的局部有限单位分解与光滑截断；具体调用时会恢复条件。分析方面使用紧集上的一致连续性、Taylor 估计以及 Lebesgue 零测集和 Fubini 定理，所需含义在证明处交代。本章的 Sard 与横截性主干针对光滑映射，光滑逼近同时控制任意指定的有限阶导数；最优有限阶 Sard 条件在第14.2节准确陈述并注明原论文的证明位置。有关光滑逼近和参数横截性的参考见\cite[第6章]{Lee2012Smooth}。

% !TeX root = ../../Book2.tex
% 纲要标题：### 14.1 光滑逼近及函数空间
% 纲要位置：几何学计划纲要.md，第 1542 行。
% 正文状态：未撰写；本文件目前是框架。

% 纲要内容开始
% 
% * 局部卷积和单位分解
% * 保持闭子集附近数据的相对逼近
% * 紧源上的 Cʳ 拓扑
% * 非紧源上支集与拓扑的选择
% 
% 纲要内容结束

\section{光滑逼近及函数空间}
\label{g2:sec:1401}

待撰写

\subsection{局部卷积和单位分解}
\label{g2:subsec:1401:01}
待撰写

\subsection{保持闭子集附近数据的相对逼近}
\label{g2:subsec:1401:02}
待撰写

\subsection{\texorpdfstring{紧源上的 \(C^r\) 拓扑}{紧源上的 Cʳ 拓扑}}
\label{g2:subsec:1401:03}
待撰写

\subsection{非紧源上支集与拓扑的选择}
\label{g2:subsec:1401:04}
待撰写
% 纲要条目排为编号小节标题，并进入目录与书签；说明段仍保留于纲要注释中。

% !TeX root = ../../Book2.tex
% 纲要标题：### 14.2 Sard 定理的证明
% 纲要位置：几何学计划纲要.md，第 1549 行。
% 正文状态：初稿；尚未正式审核。

% 纲要内容开始
% 
% * 临界集合的秩分层
% * 秩为零部分的覆盖估计
% * 维数归纳与 Fubini 步骤
% * 正则性阶数及反例范围
% 
% 纲要内容结束

\section{Sard 定理的证明}
\label{g2:sec:1402}

方程 \(f(x)=y\) 的解集能否成为光滑子流形，取决于解上的微分是否满射。单独检查一个指定的 \(y\) 可能很难；Sard 定理给出整体信息：对光滑映射，坏的目标值只占一个零测集。证明必须控制临界点的\emph{像}，因为源中的临界点本身可能很多，甚至充满整个源流形。

\subsection{临界集合的秩分层}
\label{g2:subsec:1402:01}

先明确这里的“坏”。设 \(f:M^m\to N^n\) 是 \(C^1\) 映射，\(N\) 无边界。若 \(\mathrm{d}f_x:\mathrm{T}_xM\to\mathrm{T}_{f(x)}N\) 不满射，称 \(x\) 为\term{临界点}[critical point]；否则称为\term{正则点}[regular point]。临界点的像称为\term{临界值}[critical value]。若 \(f^{-1}(y)\) 中每一点都正则，称 \(y\) 为\term{正则值}[regular value]。特别地，不在像中的点也是正则值；一个值只要有一个临界原像，就已经是临界值。

\ProseParagraphBreak

在 \(n\) 维流形 \(N\) 中，集合 \(E\) 称为\term{零测集}[set of measure zero]，若它在每张坐标图中的像都有 \(n\) 维 Lebesgue 测度零。Euclid 空间中的含义是：给定任意 \(\varepsilon>0\)，能用可数个长方体覆盖它，并使体积总和小于 \(\varepsilon\)。这个条件与坐标无关。在较小的紧坐标块上，光滑坐标变换的导数有界，因而是 Lipschitz 映射；一个半径为 \(r\) 的小球的像落在半径至多为 \(Lr\) 的球中，覆盖总体积至多乘上固定常数。再用可数个这样的坐标块覆盖，便知零测性保持。这里用的是同维坐标变化，无须先在 \(N\) 上选择体积形式。

\begin{example}\label{g2:exm:sard-critical-values}
	映射 \(f:\mathbb R^2\to\mathbb R^2\)，\(f(x,y)=(x,y^2)\)，有
	\[
		\mathrm{D}f_{(x,y)}=\begin{pmatrix}1&0\\0&2y\end{pmatrix}.
	\]
	临界点集是 \(y=0\)，临界值集是目标中的横轴。第二坐标为正的值有两个正则原像，第二坐标为负的值没有原像；这两类值都正则。对于常值映射 \(M\to\mathbb R\)，整个 \(M\) 都由临界点组成，临界值却只有一个。
\end{example}

在坐标域 \(U\subset\mathbb R^m\) 内写 \(C=\{x:\operatorname{rank}\mathrm{D}f_x<n\}\)，并按秩分为
\[
	C^{(q)}=\{x:\operatorname{rank}\mathrm{D}f_x=q\},\qquad 0\le q<n.
\]
这些集合未必都是子流形。分层在这里首先是一种分类：当 \(q\ge1\) 时，某个非零的偏导数让我们把一个输出坐标直接改作输入坐标，进而把问题降低一维。

\begin{lemma}\label{g2:lem:sard-rank-reduction}
	设 \(f:U\to\mathbb R^n\) 光滑，\(U\subset\mathbb R^m\) 开，且 \(\operatorname{rank}\mathrm{D}f_a\ge1\)。重排坐标后，在 \(a\) 的某个邻域上存在光滑源坐标 \((t,z)\in J\times V\subset\mathbb R\times\mathbb R^{m-1}\)，使
	\[
		f(t,z)=(t,G(t,z)),\qquad
		\operatorname{rank}\mathrm{D}f_{(t,z)}=1+\operatorname{rank}\mathrm{D}_zG(t,z).
	\]
\end{lemma}
\begin{proof}
	选择非零偏导数并重排坐标，使 \(\partial f^1/\partial x^1(a)\ne0\)。令
	\[
		\Phi(x)=\bigl(f^1(x),x^2,\ldots,x^m\bigr).
	\]
	其导数行列式正是这个非零偏导数。第一卷第11章第3节的逆函数定理保证 \(\Phi\) 在充分小的邻域上可光滑反解；缩小像邻域，可取乘积盒 \(J\times V\)。在新坐标中，第一分量为 \(t\)，其余分量记作 \(G\)，于是
	\[
		\mathrm{D}f=\begin{pmatrix}1&0\\\partial_tG&\mathrm{D}_zG\end{pmatrix}.
	\]
	利用第一行消去第一列下方的元素，矩阵的秩便等于 \(1+\operatorname{rank}\mathrm{D}_zG\)。
\end{proof}

这里没有要求邻域内的秩恒定；非零偏导数只负责选出一个可用的坐标。真正需要继续处理的，正是其余方向上仍会变化的秩。

\subsection{秩为零部分的覆盖估计}
\label{g2:subsec:1402:02}

若 \(\mathrm{D}f_a=0\)，上面的坐标办法暂时失效。此时 Taylor 展开提供另一种控制：若若干阶导数都为零，一个边长为 \(\delta\) 的源立方体的像就比 \(\delta\) 小得多。源中小立方体的数量约为 \(\delta^{-m}\)，而每个像的 \(n\) 维体积约为 \(\delta^{n(k+1)}\)；当后一个指数更大时，总体积趋于零。

\begin{lemma}\label{g2:lem:sard-flat-cover}
	设 \(U\subset\mathbb R^m\) 开，\(m,n\ge1\)，\(f:U\to\mathbb R^n\) 为 \(C^{k+1}\) 映射，\(k\ge1\)。令
	\[
		A_k=\{x\in U:\partial^\alpha f^i(x)=0\text{，对所有 }1\le|\alpha|\le k,\ 1\le i\le n\}.
	\]
	若 \(n(k+1)>m\)，则 \(f(A_k)\) 为零测集。
\end{lemma}
\begin{proof}
	用可数个包含于 \(U\) 的闭立方体覆盖 \(U\)，只需在其中一个闭立方体 \(Q\) 上证明结论。\(Q\) 上各个 \(k+1\) 阶偏导数有共同的上界。将边长为 \(R\) 的 \(Q\) 均分成 \(s^m\) 个边长 \(R/s\) 的小立方体。对每个与 \(A_k\) 相交的小立方体 \(Q_j\)，选取 \(a_j\in A_k\cap Q_j\)。Taylor 公式的低阶项全部为零，故存在与 \(s,j\) 无关的常数 \(B\)，使
	\[
		|f(x)-f(a_j)|\le B|x-a_j|^{k+1}
		\le B\bigl(\sqrt m R/s\bigr)^{k+1}\qquad(x\in Q_j).
	\]
	因此 \(f(A_k\cap Q)\) 被至多 \(s^m\) 个半径为 \(B(\sqrt m R/s)^{k+1}\) 的球覆盖。若 \(\omega_n\) 表示单位 \(n\) 维球的体积，覆盖总体积至多为
	\[
		s^m\omega_n B^n(\sqrt m R/s)^{n(k+1)}
		=C_Qs^{m-n(k+1)}\longrightarrow0.
	\]
	零测性由此成立；可数并仍为零测集，故得到整个 \(A_k\) 的结论。
\end{proof}

这里的多重指标 \(\alpha=(\alpha_1,\ldots,\alpha_m)\) 记录各坐标求导次数，\(|\alpha|=\alpha_1+\cdots+\alpha_m\)。上式只需普通的有限阶 Taylor 余项估计，不要求函数等于自己的 Taylor 级数。对于光滑函数，在一点的全部导数都为零也不能推出其在邻域内恒定；例如 \(e^{-1/x^2}\) 在 \(x=0\) 补值为零后的函数就具有这种现象。覆盖估计恰好能处理这些极其平坦的点。

\subsection{维数归纳与 Fubini 步骤}
\label{g2:subsec:1402:03}

剩下的工作是把前面两种控制接起来。正秩部分沿一个输出坐标切片；零秩但尚未高度平坦的部分，则落在某个偏导数的正则零集上。两种情形都把源的维数减少一。

\ProseParagraphBreak

切片时使用如下 Fubini 结论：若 \(E\subset\mathbb R\times\mathbb R^{n-1}\) 紧，\(n\ge2\)，且每个截面 \(E_t=\{w:(t,w)\in E\}\) 都为 \((n-1)\) 维零测集，则 \(E\) 为 \(n\) 维零测集。因为紧集是 Borel 可测集，对其示性函数使用非负函数的 Fubini--Tonelli 定理，得到
\[
	\operatorname{vol}_n(E)=\int_{\mathbb R}\operatorname{vol}_{n-1}(E_t)\,\mathrm{d}t=0.
\]
这一紧集切片结论也有只用有限覆盖的证明，见\cite[第126页，引理6.2]{Lee2012Smooth}。下面每次都先限制到紧块，使所处理的像确实可测。

\begin{theorem}[Sard 定理]\label{g2:thm:sard}
	设 \(M^m\) 是有限维、Hausdorff、第二可数的光滑流形，允许有边界；\(N^n\) 是无边界光滑流形，具有相同的分离性与可数性条件。若 \(f:M\to N\) 光滑，则它的临界值集在 \(N\) 中为零测集。因此正则值在 \(N\) 中稠密。
\end{theorem}
\begin{proof}
	先证明无边界情形，并对源维数 \(m\) 归纳。当 \(n=0\) 时，所有微分都满射，临界值集为空。当 \(m=0<n\) 时，第二可数的零维流形是可数离散集，像也可数，故为零测集。以下设 \(m,n\ge1\)，并假设所有源维数小于 \(m\) 的光滑映射均已满足结论。

	取可数的源、目标坐标图，并将源图限制在各目标图的原像内。坐标变化保持微分的秩和目标零测性，故只需讨论 \(f:U\subset\mathbb R^m\to\mathbb R^n\)。除临界集 \(C\) 外，记 \(A_j\) 为前一引理中一至 \(j\) 阶偏导数全为零的集合。它们在 \(U\) 内闭，并满足
	\[
		C\supset A_1\supset A_2\supset\cdots.
	\]
	取整数 \(k\ge1\) 使 \(n(k+1)>m\)。我们依次处理
	\[
		C=(C\setminus A_1)\cup\bigcup_{j=1}^{k-1}(A_j\setminus A_{j+1})\cup A_k.
	\]

	对 \(C\setminus A_1\)，若 \(n=1\)，此集合为空：一维目标的微分不满射就只能为零。若 \(n\ge2\)，\cref{g2:lem:sard-rank-reduction}在每个这样的点附近给出 \(f(t,z)=(t,G(t,z))\)。令 \(G_t(z)=G(t,z)\)。秩公式说明 \((t,z)\) 是 \(f\) 的临界点，当且仅当 \(z\) 是 \(G_t:V\to\mathbb R^{n-1}\) 的临界点。对源坐标盒内任意紧块 \(K\)，集合 \(E=f(C\cap K)\) 紧，而它在目标超平面第一坐标等于 \(t\) 上的截面包含于 \(G_t\) 的临界值集。\(G_t\) 的源维数是 \(m-1\)，由归纳假设，每个这样的截面为零测集；Fubini 结论遂给出 \(E\) 零测。用可数紧块穷尽各坐标盒，再用可数个坐标盒覆盖 \(C\setminus A_1\)，得到 \(f(C\setminus A_1)\) 零测。

	对 \(A_j\setminus A_{j+1}\)，取一点 \(a\)。某个分量的某个 \(j+1\) 阶偏导数在 \(a\) 非零，所以能选一个 \(j\) 阶偏导数 \(h=\partial^\alpha f^i\)，使 \(\mathrm{D}h_a\ne0\)。同时 \(h\) 在 \(A_j\) 上为零。第一卷第11章第3节的隐函数定理说明，在 \(a\) 的充分小邻域 \(V_a\) 内，\(H_a=\{h=0\}\) 是光滑超曲面：选取 \(h\) 的一个非零偏导数，即可把相应坐标解成其余 \(m-1\) 个坐标的光滑函数。于是 \(A_j\cap V_a\subset H_a\)。在这些点上，\(\mathrm{d}f=0\)，故限制映射 \(f|_{H_a}\) 的微分也为零，特别地不满射。归纳假设应用于源维数为 \(m-1\) 的 \(f|_{H_a}\)，说明 \(f(A_j\cap V_a)\) 零测。用可数个这样的邻域覆盖 \(A_j\setminus A_{j+1}\)，便得到这一整部分的像零测。

	最后，\cref{g2:lem:sard-flat-cover}直接给出 \(f(A_k)\) 零测。三部分合起来完成归纳。

	若 \(M\) 有边界，在每张边界坐标图中，光滑映射的坐标表达按定义能在每一点附近延拓到半空间之外的开邻域。原映射在边界点的微分就是这个局部延拓的微分；所以原映射的临界点也是延拓映射的临界点。对可数个局部延拓使用刚证的无边界结论，仍得到临界值集零测。最后，当 \(n\ge1\) 时，任意非空目标开集在某张坐标图中包含一个小球，具有正测度，不能全部由临界值组成；\(n=0\) 时所有值都正则。因此正则值稠密。
\end{proof}

这条证明的分工可由下表记忆。它按照点上的导数选择工具，而没有把整片临界集误当成一个光滑子流形。

\begin{center}
	\begin{tabularx}{\textwidth}{@{}lXX@{}}
		\toprule
		临界点所在部分 & 可用的信息 & 控制临界值的方法 \\
		\midrule
		\(C\setminus A_1\) & 某个一阶偏导数非零 & 一个输出坐标作参数，降维后用 Fubini \\
		\(A_j\setminus A_{j+1}\) & 某个高阶偏导数有正则零集 & 限制到包含该部分的超曲面，降低源维数 \\
		\(A_k\) & 低阶 Taylor 项全部消失 & 小立方体数量与像球体积的覆盖估计 \\
		\bottomrule
	\end{tabularx}
\end{center}

\begin{corollary}\label{g2:cor:sard-lower-dimension}
	设 \(f:M^m\to N^n\) 是上述流形间的光滑映射，且 \(m<n\)。则 \(f(M)\) 是 \(N\) 中的零测集；每个正则值的原像均为空。
\end{corollary}
\begin{proof}
	从 \(m\) 维空间到 \(n\) 维空间的线性映射不可能满射，因此每一点都是临界点。像恰为临界值集，应用\cref{g2:thm:sard}即可。
\end{proof}

例如光滑曲线在平面中的像具有二维测度零，即使曲线在非紧参数区间上的像可能稠密，仍不能因此获得正面积。“稠密”说的是闭包，“零测”说的是体积，二者衡量不同的性质。

\subsection{正则性阶数及反例范围}
\label{g2:subsec:1402:04}

光滑性在上述归纳中使任意阶偏导数都继续光滑。对于有限阶可微映射，\(h=\partial^\alpha f^i\) 的可微阶数会随 \(|\alpha|\) 增大而降低，隐函数定理所得超曲面的正则性也随之降低。因此，那个归纳过程不能直接给出最优的有限阶条件。

\ProseParagraphBreak

有限阶 Sard 定理的准确形式是：对 \(C^r\) 映射 \(f:M^m\to N^n\)，若整数 \(r\ge1\) 满足
\[
	r>\max\{m-n,0\},
\]
则临界值集为零测集。这里 \(M,N\) 可取光滑流形，映射只要求 \(C^r\)，源的边界按局部延拓处理。有限阶版本的证明参见\cite[第885--889页，定理4.1、7.2]{Sard1942}；其中高阶平坦集合的细分估计负责避免上述正则性损失。本章后续参数定理使用已完整证明的光滑版本。

\ProseParagraphBreak

这个阶数条件有可见的几何意义。对平面上的函数，\(C^2\) 足够；单有 \(C^1\) 时，临界值甚至可以覆盖整个区间。下面给出一个具体构造，使“更多可微性”所排除的现象看得见。

\begin{example}\label{g2:exm:sard-c1-counterexample}
	存在 \(C^1\) 函数 \(F:\mathbb R^2\to\mathbb R\)，其临界值集包含 \([0,2]\)。

	从 \([0,1]\) 开始，每一步保留每个区间两端各占其长度 \(\lambda=2/5\) 的闭子区间，得到 Cantor 集 \(E\)。第 \(k\) 步删去 \(2^{k-1}\) 个开间隙，每个长度为 \((1-2\lambda)\lambda^{k-1}\)。固定一个在 \((0,1)\) 上为正、端点处平坦、积分为 \(1\) 的光滑函数 \(\rho\)，并在每个第 \(k\) 步间隙 \(I=(a,b)\) 上放置
	\[
		h(x)=\frac{3^{-k}}{b-a}\rho\left(\frac{x-a}{b-a}\right).
	\]
	在 \(E\) 及 \([0,1]\) 外令 \(h=0\)。间隙上这项的积分为 \(3^{-k}\)，上确界至多为
	\[
		\frac{\|\rho\|_\infty}{3(1-2\lambda)}
		\left(\frac1{3\lambda}\right)^{k-1}\longrightarrow0.
	\]
	这保证 \(h\) 连续：在某个固定间隙的端点附近，连续性来自 \(\rho\) 在端点趋于零；当间隙的级数趋于无穷时，则由上述一致上界趋于零。有限多个低级间隙与余下间隙分别处理，即得到每个 \(E\) 中点处的连续性。

	令 \(g(x)=\int_0^xh(t)\,\mathrm{d}t\)，则 \(g\in C^1(\mathbb R)\)，且 \(g'=0\) 于 \(E\)。全部间隙的积分之和为 \(\sum_{k\ge1}2^{k-1}3^{-k}=1\)。一个第 \(k\) 步保留下来的基本区间，其内各级间隙的积分之和为
	\[
		\sum_{j=k+1}^{\infty}2^{j-k-1}3^{-j}=3^{-k}.
	\]
	所以 \(g\) 把保留区间依次映到标准三分 Cantor 构造中的保留区间，且把每个间隙映到相应的中间三分之一。于是 \(g(E)=C\)，其中 \(C\) 是标准三分 Cantor 集。

	现在取 \(F(x,y)=g(x)+g(y)\)。在 \(E\times E\) 上，\(\mathrm{d}F=(g'(x),g'(y))=0\)，而 \(F(E\times E)=C+C=[0,2]\)。最后这个等式可这样验证：\([-1,1]\) 中每个数都能写成 \(\sum_{j\ge1}2\varepsilon_j3^{-j}\)，其中 \(\varepsilon_j\in\{-1,0,1\}\)，因为第一位所给的三个闭区间覆盖 \([-1,1]\)，选定第一位后再缩放余项即可递归选择。把正、负位分别放到两个只含 \(0,2\) 的三进制展开里，得到 \(C-C=[-1,1]\)。再由 \(1-C=C\)，得到 \(C+C=1+(C-C)=[0,2]\)。
\end{example}

因此，零测结论涉及映射的正则性，而不只是处处存在一个连续的微分。在一般维数中，整数阶阈值的最优性及相应构造见\cite[第890页，第8节]{Sard1942}；对当前 \(m=2,n=1\) 的情形，上例已经完整展示 \(C^1\) 与 \(C^2\) 的差别。另一个边界是源和目标维数：当 \(m<n\) 时，\(C^1\) 已足以让整个像零测：在一个紧坐标立方体上，导数有界使映射 Lipschitz；分割所得的 \(s^m\) 个小立方体各自的像都被半径至多 \(B/s\) 的球包含，故总体积至多为 \(Cs^{m-n}\to0\)。再取可数坐标覆盖即可。

% !TeX root = ../../Book2.tex
% 纲要标题：### 14.3 横截逆像和参数定理
% 纲要位置：几何学计划纲要.md，第 1556 行。
% 正文状态：初稿；尚未正式审核。

% 纲要内容开始
% 
% * 线性横截条件与切空间
% * 逆像的维数及法丛
% * 由总体映射到几乎处处参数
% * 稠密性、剩余性和泛型的区别
% 
% 纲要内容结束

\section{横截逆像和参数定理}
\label{g2:sec:1403}

正则值定理处理方程 \(f(x)=y\)：微分能够独立改变全部输出时，解集有预期维数。现在把一个点 \(y\) 换成子流形 \(S\subset N\)，研究条件 \(f(x)\in S\)。沿着 \(S\) 的输出变化不会破坏这个条件，真正需要控制的是离开 \(S\) 的方向。横截性正是把“全部输出方向”改成“全部法方向”后的满秩条件。

\subsection{线性横截条件与切空间}
\label{g2:subsec:1403:01}

设 \(S\subset N\) 是无边界光滑流形 \(N\) 的嵌入子流形。在 \(y\in S\) 处，用包含映射的微分把 \(\mathrm{T}_yS\) 视为 \(\mathrm{T}_yN\) 的子空间。商空间 \(\mathrm{T}_yN/\mathrm{T}_yS\) 记录法方向：两个向量只要相差一个切于 \(S\) 的向量，就代表同一个法方向。这里还没有选择内积，因而“法”暂时不表示垂直。

\begin{definition}[横截性]\label{g2:def:transversality}
	设 \(M,N\) 为无边界光滑流形，\(S\subset N\) 为无边界嵌入子流形，\(f:M\to N\) 光滑。若对每个 \(x\in f^{-1}(S)\) 都有
	\begin{equation}\label{g2:eq:transverse-linear-condition}
		\mathrm{d}f_x(\mathrm{T}_xM)+\mathrm{T}_{f(x)}S=\mathrm{T}_{f(x)}N,
	\end{equation}
	则称 \(f\) 与 \(S\) \term[hengjiexing]{横截}[transverse]，记作 \(f\pitchfork S\)。若只在 \(x\in K\cap f^{-1}(S)\) 处要求此等式，就说 \(f\) 在 \(K\subset M\) 上与 \(S\) 横截。
\end{definition}

横截条件等价于商映射
\[
	\overline{\mathrm{d}f_x}:\mathrm{T}_xM\longrightarrow\mathrm{T}_{f(x)}N/\mathrm{T}_{f(x)}S,
	\qquad v\longmapsto[\mathrm{d}f_x(v)]
\]
满射。因而只需补足 \(S\) 缺少的方向；\(\mathrm{d}f_x\) 本身可以不满射。当 \(S=\{y\}\) 时，\(\mathrm{T}_yS=0\)，这个条件恰好说 \(y\) 是正则值。淹没与每个嵌入子流形都横截；若 \(f(M)\cap S=\varnothing\)，横截条件也成立，因为没有需要检查的点。记号中的横截并不要求相交。

\begin{example}\label{g2:exm:transverse-curves}
	取 \(S=\mathbb R\times\{0\}\subset\mathbb R^2\)。映射 \(f(t)=(t,t)\) 在唯一交点处给出斜向量 \((1,1)\)，它与 \(S\) 的切向量 \((1,0)\) 张成平面，所以 \(f\pitchfork S\)。映射 \(g(t)=(t,t^2)\) 在零处只给出 \((1,0)\)，与 \(S\) 相切，故不横截。映射 \(h(t)=(t,t^2+1)\) 的像避开 \(S\)，也与 \(S\) 横截。\cref{g2:fig:transversality-curves}将这三种情形并排画出。
\end{example}

\begin{figure}[htbp]
	\centering
	\input{Figures/Chapter14-Transversality}
	\caption[横截、相切与空交]{横轴表示子流形 \(S\)，曲线表示映射的像。左图在交点处补足法方向；中图的切线仍落在 \(S\) 内；右图没有交点，也满足横截条件。}
	\label{g2:fig:transversality-curves}
\end{figure}

若 \(A,B\subset N\) 都是嵌入子流形，对包含映射 \(i_A:A\hookrightarrow N\) 使用上述定义，就得到两子流形的横截条件
\[
	\mathrm{T}_yA+\mathrm{T}_yB=\mathrm{T}_yN\qquad(y\in A\cap B).
\]
它对 \(A,B\) 对称。两条空间曲线在 \(\mathbb R^3\) 中若有交点，两个一维切空间之和至多二维，所以不可能横截；相交曲线在图上看起来“交叉”，还不足以判断横截性，环境维数也参与条件。

\subsection{逆像的维数及法丛}
\label{g2:subsec:1403:02}

第一卷第11章第4节的嵌入子流形坐标切片说明：若 \(S\) 的余维为 \(c\)，那么在每个 \(y\in S\) 附近，可以找一个淹没 \(q:V\to\mathbb R^c\)，使 \(S\cap V=q^{-1}(0)\)。具体说，选取把 \(S\) 写成最后 \(c\) 个坐标为零的适配坐标，再令 \(q\) 取出这 \(c\) 个坐标。于是 \(f(x)\in S\) 局部变成 \(q(f(x))=0\)。

\begin{theorem}[横截逆像]\label{g2:thm:transverse-preimage}
	设 \(M^m,N^n\) 为无边界光滑流形，\(S\subset N\) 为余维 \(c\) 的无边界嵌入子流形，\(f:M\to N\) 光滑且 \(f\pitchfork S\)。则非空的 \(W=f^{-1}(S)\) 是 \(M\) 的 \((m-c)\) 维嵌入子流形，并且
	\begin{equation}\label{g2:eq:transverse-preimage-tangent}
		\mathrm{T}_xW=\{v\in\mathrm{T}_xM:\mathrm{d}f_x(v)\in\mathrm{T}_{f(x)}S\}.
	\end{equation}
	若 \(c>m\)，则 \(W=\varnothing\)。若 \(S\) 在 \(N\) 中闭，则 \(W\) 在 \(M\) 中闭。
\end{theorem}
\begin{proof}
	在 \(y=f(x)\in S\) 附近取上述 \(q\)。因为 \(\ker\mathrm{d}q_y=\mathrm{T}_yS\)，横截条件等价于 \(\mathrm{d}(q\circ f)_x\) 满射。为看清这一步，任取 \(a\in\mathbb R^c\)，先用 \(q\) 为淹没找到 \(u\in\mathrm{T}_yN\) 满足 \(\mathrm{d}q_y(u)=a\)，再由横截性把 \(u\) 写成 \(\mathrm{d}f_x(v)+w\)，其中 \(w\in\mathrm{T}_yS\)。于是 \(a=\mathrm{d}(q\circ f)_x(v)\)。反过来，若这个复合满射，对任意 \(u\in\mathrm{T}_yN\) 选择 \(v\) 满足 \(\mathrm{d}(q\circ f)_x(v)=\mathrm{d}q_y(u)\)，便有 \(u-\mathrm{d}f_x(v)\in\mathrm{T}_yS\)。

	因此 \(0\) 是 \(q\circ f:f^{-1}(V)\to\mathbb R^c\) 的正则值。第一卷第11章第4节的正则水平集定理给出局部的 \((m-c)\) 维坐标切片，切空间为 \(\ker\mathrm{d}(q\circ f)_x\)，即\eqref{g2:eq:transverse-preimage-tangent}。这些局部切片都取 \(M\) 的子空间拓扑，且坐标转换来自环境坐标的限制，故拼成所述嵌入子流形。若 \(c>m\)，从 \(m\) 维空间到 \(c\) 维空间不可能满射，所以不存在原像点。最后，\(S\) 闭时由 \(f\) 连续得到 \(W\) 闭。
\end{proof}

例如，两张在 \(N^n\) 中横截相交的子流形 \(A^a,B^b\)，其交非空时维数为 \(a+b-n\)，且在交点处
\[
	\mathrm{T}_y(A\cap B)=\mathrm{T}_yA\cap\mathrm{T}_yB.
\]
这是将定理用于 \(i_A:A\hookrightarrow N\) 的直接结果。横截条件保证方程提供了独立约束；一个余维 \(c\) 的目标子流形，恰好在源中施加 \(c\) 个独立方程。

\ProseParagraphBreak

切空间公式还可以记录随点变化的法方向。对嵌入子流形 \(S\subset N\)，定义其\term[facong]{法丛}[normal bundle]为商向量丛
\[
	\nu_N S\coloneqq\mathrm{T}N|_S/\mathrm{T}S.
\]
在适配坐标中，这个商丛由最后 \(c\) 个坐标方向给出局部标架；换坐标诱导光滑可逆的商空间变换，因而确实定义了秩 \(c\) 的光滑向量丛。若选取 Riemann 度量，正交投影将此商丛识别为 \((\mathrm{T}S)^\perp\)，但商丛的定义本身不依赖度量。

\begin{proposition}[逆像的法丛]\label{g2:prop:transverse-normal-bundle}
	在\cref{g2:thm:transverse-preimage}的条件下，令 \(\bar f=f|_W:W\to S\)。微分诱导规范的光滑向量丛同构
	\[
		\nu_M W\longrightarrow\bar f^{\,*}\nu_N S,
		\qquad[v]\longmapsto[\mathrm{d}f(v)].
	\]
	等价地，有短正合列
	\[
		0\longrightarrow\mathrm{T}W\longrightarrow\mathrm{T}M|_W
		\xrightarrow{\ \overline{\mathrm{d}f}\ }\bar f^{\,*}\nu_N S\longrightarrow0.
	\]
\end{proposition}
\begin{proof}
	由\eqref{g2:eq:transverse-preimage-tangent}，\(\overline{\mathrm{d}f_x}\) 的核恰为 \(\mathrm{T}_xW\)，横截性又保证它满射。线性代数的商空间同构定理逐点给出所述同构。它们还须组成光滑丛映射：在源与目标的适配坐标中，微分由光滑矩阵表示，取法方向就是取相应的商矩阵，因而随点光滑。取商后的矩阵为可逆的 \(c\) 阶矩阵，其逆由伴随矩阵除以非零行列式给出，亦光滑。整个定义只使用微分与自然商映射，所以与坐标无关。
\end{proof}

由此可见，横截逆像除了继承预期的维数，也继承目标子流形的法方向。例如，若 \(S\) 的法丛已经给定定向，就可以通过这个同构给 \(W\) 的法丛定向。后面的相交数构造将使用这项信息。

\subsection{由总体映射到几乎处处参数}
\label{g2:subsec:1403:03}

一个指定映射可能不横截。给它增加几个可自由调节的参数后，参数方向有机会补足原来缺少的法方向。例如
\[
	F:\mathbb R\times\mathbb R\longrightarrow\mathbb R,
	\qquad F(a,t)=t^2-a.
\]
固定 \(a=0\) 时，\(F_0(t)=t^2\) 与 \(\{0\}\) 不横截；总体映射却处处是淹没，因为 \(\partial F/\partial a=-1\)。下面的定理把这种可调节性转化为关于单个参数的结论。

\begin{theorem}[参数横截性]\label{g2:thm:parametric-transversality}
	设 \(P,M,N\) 为无边界光滑流形，\(S\subset N\) 为无边界嵌入子流形，\(F:P\times M\to N\) 光滑。若 \(F\pitchfork S\)，则对几乎处处的 \(a\in P\)，切片映射 \(F_a:M\to N\)，\(F_a(x)=F(a,x)\)，满足 \(F_a\pitchfork S\)。这里的例外集在 \(P\) 的每张坐标图内都为 Lebesgue 零测集；特别地，好参数在 \(P\) 中稠密。
\end{theorem}
\begin{proof}
	先把所有满足约束的参数与源点放在一起。由\cref{g2:thm:transverse-preimage}，
	\[
		Z=F^{-1}(S)\subset P\times M
	\]
	是光滑嵌入子流形。考虑投影 \(\pi:Z\to P\)，\(\pi(a,x)=a\)。我们证明：\(a\) 是 \(\pi\) 的正则值，当且仅当 \(F_a\pitchfork S\)。这样就能将横截性问题交给\cref{g2:thm:sard}。

	固定 \((a,x)\in Z\)，令 \(y=F(a,x)\)、\(Q=\mathrm{T}_yN/\mathrm{T}_yS\)。把总体微分送到这个商空间，分别记参数方向与源方向的部分为
	\[
		B:\mathrm{T}_aP\longrightarrow Q,\qquad
		A:\mathrm{T}_xM\longrightarrow Q.
	\]
	于是总体微分为 \((u,v)\mapsto Bu+Av\)。总体横截性表示 \(\operatorname{im}B+\operatorname{im}A=Q\)，而切空间公式给出
	\[
		\mathrm{T}_{(a,x)}Z=\{(u,v):Bu+Av=0\}.
	\]
	投影微分 \(\mathrm{d}\pi_{(a,x)}\) 满射，等价于每个 \(u\) 都能配一个 \(v\)，使 \(Bu+Av=0\)，即 \(\operatorname{im}B\subset\operatorname{im}A\)。结合总体的满射性，这恰好等价于 \(\operatorname{im}A=Q\)，也就是 \(F_a\) 在 \(x\) 处与 \(S\) 横截。对该参数下的所有 \(x\) 同时检查，就得到所需等价；若纤维为空，两边的条件都成立。

	Sard 定理说明 \(\pi\) 的临界值集为零测集，故其补集中的所有参数都横截。任意非空坐标开集都有正测度，不可能整个包含在例外集中，所以好参数稠密。零维参数空间中，Sard 的结论表示没有坏参数，此时上述线性证明也直接给出每个切片横截。
\end{proof}

上述证明使用的不是 \(F\) 的正则值，而是约束流形 \(Z\) 到参数空间的投影的正则值。正则值条件保证：改变参数以后，可以同时改变源点来补偿约束的变化。这正是把总体可调节性传到单个切片的步骤。横截逆像和参数定理也可参看\cite[第6章，第143--146页]{Lee2012Smooth}。

\begin{example}\label{g2:exm:parametric-parabola}
	回到 \(F(a,t)=t^2-a\)，约束流形 \(Z\) 是参数平面中的抛物线 \(a=t^2\)，投影为 \(\pi(t)=t^2\)。它只有一个临界值 \(a=0\)。当 \(a>0\) 时，两个零点 \(\pm\sqrt a\) 处导数非零；当 \(a<0\) 时，没有零点；两类参数都满足横截性。\(a=0\) 是两个零点合并的时刻。参数定理说明坏参数集很小，却并不声称所有参数都好，也不保证不同好参数的原像具有相同的连通分支数。
\end{example}

\subsection{稠密性、剩余性和泛型的区别}
\label{g2:subsec:1403:04}

参数定理使用有限维参数空间的测度；“映射经过任意小扰动就能横截”则使用函数空间的拓扑。要把这两种表述衔接起来，先需要一个稳定性事实：已经在紧集上满足的横截条件，可以承受足够小的一阶扰动。

\begin{proposition}[紧集上的横截开放性]\label{g2:prop:transversality-open-on-compact}
	设 \(M,N\) 为无边界光滑流形，\(S\subset N\) 为闭的无边界嵌入子流形，\(K\subset M\) 紧。若光滑映射 \(f:M\to N\) 在 \(K\) 上与 \(S\) 横截，则 \(f\) 有一个弱 \(C^1\) 邻域，其中每个光滑映射在 \(K\) 上仍与 \(S\) 横截。特别地，当 \(M\) 紧时，全部 \(C^1\) 横截映射构成 \(C^1(M,N)\) 中的开集。
\end{proposition}
\begin{proof}
	对每个 \(x\in K\)，区分两种情形。若 \(f(x)\notin S\)，因 \(S\) 闭，可取 \(x\) 的相对紧邻域 \(U_x\)，使 \(f(\overline U_x)\) 落在 \(N\setminus S\) 中；充分小的 \(C^0\) 扰动在较小闭邻域上仍避开 \(S\)。这一控制可用覆盖该紧像的有限张坐标图表达。

	若 \(f(x)\in S\)，取适配坐标及局部定义函数 \(q:V\to\mathbb R^c\)。横截性使 \(\mathrm{D}(q\circ f)\) 的某个 \(c\) 阶子式在 \(x\) 非零。缩小相对紧的 \(U_x\)，可使 \(f(\overline U_x)\subset V\)，且这个子式的绝对值在 \(\overline U_x\) 上有正下界。对 \(g\) 在此紧集上的值与一阶导数施加足够小的坐标误差，就能保持 \(g(\overline U_x)\subset V\) 并使同一个子式仍非零。因此 \(q\circ g\) 在这里为淹没，尤其在 \(g^{-1}(S)\cap U_x\) 上满足横截条件。

	这些邻域覆盖 \(K\)。取有限子覆盖，再同时施加相应的有限个紧集 \(C^1\) 条件，便得到所需弱 \(C^1\) 邻域。取 \(K=M\) 得最后结论；该证明只涉及一阶导数，因而也适用于 \(C^1\) 映射。
\end{proof}

闭性在这里用来阻止新的交点从 \(S\) 的遗漏极限点附近出现。例如 \(S=\{1,1/2,1/3,\ldots\}\subset\mathbb R\) 是非闭的零维嵌入子流形。从单点流形到 \(\mathbb R\) 的映射 \(f(*)=0\) 避开 \(S\)，因而横截；任意接近它的某个常值映射 \(f_j(*)=1/j\) 却不横截。这个例子即使源紧，开放性仍然失效。

\ProseParagraphBreak

源的紧性与所选拓扑也须匹配。映射 \(f(x)=x\) 从 \(\mathbb R\) 到自身与 \(\{0\}\) 横截。取光滑函数 \(\rho\)，使其在零附近等于 \(1\)，在 \([-1,1]\) 外为零，令
\[
	g_j(x)=x\bigl(1-\rho(x-j)\bigr).
\]
在每个固定紧集上，\(g_j\) 最终与 \(f\) 完全相同，所以 \(g_j\to f\) 于弱 \(C^\infty\) 拓扑；但 \(g_j\) 在 \(j\) 附近恒为零，均不横截。紧集上的开放性无法一次控制逃往无穷远处的交点。

\begin{definition}[稠密与剩余]\label{g2:def:generic-residual}
	在拓扑空间 \(X\) 中，若 \(A\) 与每个非空开集相交，则称 \(A\) \term[choumi]{稠密}[dense]。若 \(A\) 包含一个可数个开稠密集的交，则称 \(A\) 为\term[shengyuji]{剩余集}[residual set]。一个性质在剩余集上成立时，通常称其为\term[fanxing]{泛型}[generic]性质，使用时须同时说明空间及其拓扑。
\end{definition}

稠密性保证任意小的邻域里有好对象，剩余性则记录可数个开稠密要求能够同时成立。在 Baire 空间中，每个可数开稠密交都稠密。例如 \(\mathbb R\) 具有这个性质：给定非空开区间及开稠密集 \(U_1,U_2,\ldots\)，依次在前一个区间内部与 \(U_j\) 的交中选非空闭区间，令长度趋零；实数的完备性给出一个属于所有区间的点。这个嵌套区间论证说明了完备性如何使可数条件同时实现。一般拓扑空间未必有此性质：在 \(\mathbb Q\) 的相对拓扑中，枚举其点为 \(q_1,q_2,\ldots\)，每个 \(\mathbb Q\setminus\{q_j\}\) 都开且稠密，交却为空。

\ProseParagraphBreak

对光滑映射空间，本章使用的具体结论是：若 \(M,N\) 无边界，\(S\subset N\) 是闭的无边界嵌入子流形，则
\[
	\mathcal T_S=\{f\in C^\infty(M,N):f\pitchfork S\}
\]
在弱 \(C^\infty\) 拓扑及\secref{g2:sec:1401}定义的强 \(C^\infty\) 拓扑中都稠密且为剩余集。稠密性由下一节的\cref{g2:thm:relative-transversality}证明：其中取固定闭集为空，便能在任意给定的强邻域内得到横截映射。剩余性的推导则很短。取紧集穷尽 \(K_1\subset K_2\subset\cdots\)，\(\bigcup_jK_j=M\)，令
\[
	\mathcal U_j=\{f:f\text{ 在 }K_j\text{ 上与 }S\text{ 横截}\}.
\]
由\cref{g2:prop:transversality-open-on-compact}，每个 \(\mathcal U_j\) 在这两种拓扑中都开；全局横截映射的稠密性又使每个 \(\mathcal U_j\) 稠密。最后 \(\mathcal T_S=\bigcap_j\mathcal U_j\)。这里稠密性有单独的扰动证明，不需要先对整个流形映射空间证明一个 Baire 定理。

\ProseParagraphBreak

“几乎处处”与“泛型”也不是同一个量词。将有理数枚举为 \(q_i\)，对每个 \(k\) 取以这些点为中心的开区间，使总长度小于 \(2^{-k}\)，其并 \(U_k\) 开且稠密。于是 \(G=\bigcap_kU_k\) 是稠密剩余集，却有测度零，因为它包含在每个总长度小于 \(2^{-k}\) 的 \(U_k\) 中。相应地，\(\mathbb R\setminus G\) 为满测集，却是可数个闭且无内点的集合之并。函数空间中的泛型结论不附带一个未指定的概率分布；参数横截定理的几乎处处结论则明确使用有限维参数坐标中的 Lebesgue 测度。

% !TeX root = ../../Book2.tex
% 纲要标题：### 14.4 相对与带边界横截性
% 纲要位置：几何学计划纲要.md，第 1563 行。
% 正文状态：初稿；尚未正式审核。

% 纲要内容开始
% 
% * 局部扰动和有限覆盖
% * 保持指定区域不变
% * 边界限制映射的横截条件
% * 同伦族与配边的生成
% 
% ---
% 
% 纲要内容结束

\section{相对与带边界横截性}
\label{g2:sec:1404}

参数定理要求先有一个总体横截的映射族。对任意给定映射，怎样造出这样的族？在目标坐标中加一个小向量，就能补足缺少的方向；在源端乘截断函数，又能把改变限制在指定区域。将这个局部操作逐块实施，便得到保持已给数据的一般位置结果。边界问题还需检查沿边界本身是否有足够的方向，最后才能把一族逆像拼成配边。

\subsection{局部扰动和有限覆盖}
\label{g2:subsec:1404:01}

\begin{lemma}[坐标内的横截扰动]\label{g2:lem:local-transverse-perturbation}
	设 \(M,N\) 是无边界光滑流形，\(S\subset N\) 是嵌入子流形，\(f:M\to N\) 光滑。设 \(K\subset U\subset M\)，其中 \(K\) 紧，\(U\) 开，\(f(U)\) 落在一张目标坐标图内。则可在 \(U\) 内选择紧支扰动，使所得光滑映射 \(g\) 在 \(K\) 上横截于 \(S\)，并在 \(M\setminus U\) 上等于 \(f\)。扰动可以在任意预先指定的有限阶 \(C^r\) 意义下任意小。
\end{lemma}
\begin{proof}
	记目标坐标为 \(y:V\to\mathbb R^n\)。第一卷第12章第3节的光滑截断结论保证存在 \(\rho:M\to[0,1]\)，在 \(K\) 的某邻域上为 \(1\)，且支集紧含于 \(U\)。取充分小的参数球 \(B\subset\mathbb R^n\)，定义
	\[
		F(a,x)=\begin{cases}
			y^{-1}\bigl(y(f(x))+\rho(x)a\bigr),&x\in U,\\
			f(x),&x\notin\operatorname{supp}\rho.
		\end{cases}
	\]
	两个表达在交叠处一致；\(f(\operatorname{supp}\rho)\) 紧含于 \(V\)，故可统一缩小 \(B\)，使平移后的点仍在坐标像内。在 \(\rho>0\) 的区域，关于参数的微分在坐标中为 \(\rho(x)I_n\)，所以总体映射 \(F\) 是淹没。\cref{g2:thm:parametric-transversality}应用于 \(B\times\{\rho>0\}\)，说明几乎每个任意小的 \(a\) 都使 \(F_a\) 在这里横截于 \(S\)，尤其在 \(K\) 上横截。

	令 \(g=F_a\)。在固定紧支集上，\(F_a\) 的各阶坐标导数连续依赖于 \(a\)，并在 \(a=0\) 时等于 \(f\) 的相应导数。因此让 \(a\) 足够小，即可满足任意给定的有限阶误差界。对非紧源上的强拓扑，紧支集只与局部有限控制族中的有限多个紧块相交，同一结论仍适用。
\end{proof}

这个构造说明参数为什么应在目标坐标中取值：源端原本可能缺少的法向方向，由 \(a\) 直接补上。截断函数的过渡区不必横截；下一步会用覆盖把所有尚未处理的点纳入某个 \(\rho=1\) 的区域。

\ProseParagraphBreak

先设 \(M\) 紧，\(S\) 闭。取有限个紧块 \(K_i\subset U_i\)，其内部覆盖 \(M\)，每个 \(f(U_i)\) 落在一张目标图中。依次应用引理，使第 \(i\) 步在 \(K_i\) 上横截。已经处理的 \(K_1\cup\cdots\cup K_{i-1}\) 是紧集；\cref{g2:prop:transversality-open-on-compact}保证，只要本步足够小，就保留那里的横截性。再把总误差分配成有限份，最后得到任意接近 \(f\) 的处处横截映射。这里保持的是先前取得的满秩条件，先前的函数值可以继续改变。

\subsection{保持指定区域不变}
\label{g2:subsec:1404:02}

若映射在某个区域已经用于拼接，那里常常必须逐点不变。例如同伦的两个端点已经给定，不能在“取得一般位置”时顺便更换端点。相对横截性把这种要求直接列入结论。

\begin{theorem}[相对横截逼近]\label{g2:thm:relative-transversality}
	设 \(M,N\) 是有限维、Hausdorff、第二可数的无边界光滑流形，\(S\subset N\) 是闭嵌入子流形，\(A\subset M\) 闭。若光滑映射 \(f:M\to N\) 在 \(A\) 的一个开邻域上横截于 \(S\)，则对每个有限整数 \(r\ge1\) 及 \(f\) 的每个强 \(C^r\) 邻域 \(\mathcal U\)，存在光滑映射 \(g\in\mathcal U\)，使 \(g\pitchfork S\)，并使 \(g=f\) 在 \(A\) 的某个开邻域上成立。\(g\) 还可选为与 \(f\) 光滑同伦，且同伦在这个较小邻域上不变。\(A=\varnothing\) 时即得到普通横截逼近。
\end{theorem}
\begin{proof}
	设 \(f\) 在开集 \(U_0\supset A\) 上横截。由第一卷第4章第3—4节的正规性与覆盖收缩，选开集 \(W,V\)，使
	\[
		A\subset W\subset\overline W\subset V\subset\overline V\subset U_0.
	\]
	可先缩小 \(\mathcal U\)，使其中每个映射在 \(\overline V\) 上横截。具体做法是在 \(\overline V\) 附近取局部有限的紧坐标块，在每块上用紧集横截开放性给出正的 \(C^1\) 容差；强拓扑允许同时要求这些逐块不等式。因而接下来只需处理 \(M\setminus V\)。

	对这个闭集取可数、在整个 \(M\) 上局部有限的相对紧开集族 \(U_i\)，每个 \(\overline{U_i}\) 避开 \(\overline W\)，且 \(f(\overline{U_i})\) 含于一个目标坐标图。取紧块 \(K_i\subset U_i\)，使其内部覆盖 \(M\setminus V\)。这里使用的是第一卷第12章第3节构造单位分解时的局部有限收缩；尤其局部有限性在 \(W\) 的边缘也成立。用逐块 \(C^0\) 容差保证后续映射仍留在所选目标图内。

	从 \(f_0=f\) 开始，在第 \(i\) 步应用\cref{g2:lem:local-transverse-perturbation}，得到 \(f_i\)，使其在 \(K_i\) 上横截，改动支集紧含于 \(U_i\)。同时让改动足够小，以保留 \(K_1,\ldots,K_{i-1}\) 上的横截性，并满足原邻域的误差要求。下面说明无穷多步中如何保留这些要求，避免只在每个有限阶段正确。

	在固定的源、目标坐标块中用各阶导数的上确界测量差。若原邻域在第 \(\ell\) 块给出容差 \(\varepsilon_\ell\)，令第 \(i\) 步在它碰到的各块上耗用不到 \(2^{-i-1}\varepsilon_\ell\)。第 \(i\) 步完成后，紧集开放性给出 \(\delta_i>0\)：以后只要在覆盖 \(K_i\) 的有限坐标块上与 \(f_i\) 的 \(C^1\) 差小于 \(\delta_i\)，便仍横截。对每个 \(j>i\)，额外要求第 \(j\) 次改动在这些块上的大小小于 \(2^{-(j-i+1)}\delta_i\)。总和小于 \(\delta_i\)。一次改动有紧支，所遇原控制块有限，先前的 \(K_i\) 也只有有限多个；所以每一步只是对参数添加有限多个严格上界，可以同时满足。

	每一点的某邻域只遇到有限多个 \(U_i\)，因此 \(f_i\) 在该邻域最终不再改变，定义出光滑映射 \(g\)。上述误差和保证 \(g\in\mathcal U\) 并在每个 \(K_i\) 上横截；在 \(\overline V\) 上的横截性由最初缩小的邻域保证。这些集合覆盖 \(M\)，故 \(g\pitchfork S\)。所有改动支集避开 \(\overline W\)，于是 \(g=f\) 于 \(W\)。

	最后，在各次坐标平移中把选定参数 \(a_i\) 换成 \(t a_i\)，按相同次序组合，\(0\le t\le1\)。各步可一并选得对所有 \(t\) 满足前述坐标内的误差预算，所以组合有定义。源端的局部有限性保证每点附近只有有限次组合，从而给出联合光滑的 \(H(t,x)\)，满足 \(H(0,x)=f(x)\)、\(H(1,x)=g(x)\)，且在 \(W\) 上恒等于 \(f\)。这个同伦本身的每张切片不必横截。
\end{proof}

\begin{corollary}[余维较大的集合可以避开]\label{g2:cor:avoid-high-codimension}
	设 \(M,N\) 是无边界光滑流形，\(S\subset N\) 是闭嵌入子流形，且 \(\dim M<\operatorname{codim}_N S\)。每个光滑映射 \(M\to N\) 都可任意光滑逼近为像避开 \(S\) 的映射，并保持同伦类；若原映射在闭集 \(A\) 的某邻域上避开 \(S\)，可在较小邻域内保持它不变。
\end{corollary}
\begin{proof}
	相对横截逼近给出 \(g\pitchfork S\)。若 \(g(x)\in S\)，则 \(\mathrm{T}_xM\to\mathrm{T}_{g(x)}N/\mathrm{T}_{g(x)}S\) 必须满射，与维数不等式矛盾。因此 \(g^{-1}(S)=\varnothing\)。任意 \(C^\infty\) 邻域只需满足构成它的一组有限阶基本条件，选足够大的有限 \(r\) 即可。
\end{proof}

例如三维空间中的一条闭路可以作任意小的扰动避开给定点；对一张曲面也成立，因为点的余维为3。若障碍物是一张曲面，余维只有1，上述维数论证便不再给出避开结论。相交数将在后续章节解释何时交点反而不能全部消除。

\subsection{边界限制映射的横截条件}
\label{g2:subsec:1404:03}

\begin{theorem}[带边界的横截逆像]\label{g2:thm:boundary-transverse-preimage}
	设 \(M^m\) 是光滑带边界流形，\(N\) 是无边界光滑流形，\(S\subset N\) 是无边界、余维为 \(c\) 的嵌入子流形。若光滑映射 \(f:M\to N\) 以及限制 \(f|_{\partial M}\) 都横截于 \(S\)，则 \(Z=f^{-1}(S)\) 若非空，就是 \(m-c\) 维带边界嵌入子流形，且
	\[
		\partial Z=Z\cap\partial M,\qquad
		\mathrm{T}_xZ=(\mathrm{d}f_x)^{-1}(\mathrm{T}_{f(x)}S).
	\]
	在 \(x\in\partial Z\) 处还有
	\[
		\mathrm{T}_x(\partial Z)=\mathrm{T}_xZ\cap\mathrm{T}_x(\partial M),\qquad
		\mathrm{T}_xZ+\mathrm{T}_x(\partial M)=\mathrm{T}_xM.
	\]
\end{theorem}
\begin{proof}
	内部点由\cref{g2:thm:transverse-preimage}处理。边界点处取 \(S\) 的局部定义函数 \(q:V\to\mathbb R^c\)，即 \(S\cap V=q^{-1}(0)\)，且 \(\ker\mathrm{d}q=\mathrm{T}S\) 沿 \(S\) 成立。令 \(h=q\circ f\)。边界横截条件等价于 \(\mathrm{d}(h|_{\partial M})\) 满射，这正是第一卷第11章第4节“带边界的正则水平集”所需的条件。

	为看清边界为何保留下来，在半空间坐标 \((x^1,\ldots,x^m)\)、\(x^m\ge0\) 中，满射性使前 \(m-1\) 个方向中的某 \(c\) 列给出可逆子式。重排这些坐标后，以
	\[
		\bigl(h^1(x),\ldots,h^c(x),x^{c+1},\ldots,x^{m-1},x^m\bigr)
	\]
	作新坐标。逆函数定理适用于坐标函数的局部光滑延拓，且最后一个坐标仍为 \(x^m\)，故局部逆保持半空间。\(Z\) 在这些坐标中就是前 \(c\) 项为零的半空间切片。切片边界、切空间及最后一个向内方向分别给出陈述中的各个等式。包含映射在这些图中为坐标嵌入，图册与子空间拓扑相容。
\end{proof}

\begin{example}\label{g2:exm:boundary-transversality-failure}
	在上半平面 \(M=\{(x,y):y\ge0\}\) 上令 \(f(x,y)=y+x^2\)，取 \(S=\{0\}\subset\mathbb R\)。由于 \(\partial_y f=1\)，\(f\) 处处为淹没；然而 \(f^{-1}(0)=\{(0,0)\}\)，并不是预期的一维带边界流形。缺少的条件是 \(f|_{\partial M}(x)=x^2\) 在零点处的横截性。改成 \(g(x,y)=x\)，逆像为向上的半直线，才具有一维内部及一个边界点。完整切空间含向内法向，而边界限制只能使用沿边界的方向，两项检查因此有不同作用。
\end{example}

相对逼近也有相同的边界版本：在\cref{g2:thm:relative-transversality}中允许源 \(M\) 带边界，要求 \(f\) 与 \(f|_{\partial M}\) 在固定区域都横截，则可使最终 \(g\) 与 \(g|_{\partial M}\) 同时横截。证明中的坐标扰动在半空间上仍有意义；在 \(\rho>0\) 处，总体映射限制到 \(B\times\operatorname{Int}M\) 和 \(B\times\partial M\) 都因参数方向而成为淹没。对这两张无边界源的映射分别应用参数定理，避开两个零测异常集的并，即可同时满足条件。边界上的满射还蕴含完整切空间上的满射。紧块开放性分别应用于内部的完整导数和边界切向导数，此后的有限覆盖或局部有限构造完全沿用上述误差选择。

\subsection{同伦族与配边的生成}
\label{g2:subsec:1404:04}

一条同伦使每个点的像随时间移动；若整条同伦横截于 \(S\)，所有相交点的运动轨迹就组成高一维的流形。它的边界记录起点和终点的逆像。这是由一般位置构造几何不变量的基本方式。

\begin{definition}[配边]\label{g2:def:unoriented-cobordism}
	设 \(Z_0,Z_1\) 是紧无边界 \(k\) 维光滑流形。若存在紧 \((k+1)\) 维光滑带边界流形 \(W\)，以及微分同胚 \(\partial W\cong Z_0\amalg Z_1\)，就称 \(Z_0,Z_1\)\term[peibian]{配边}[cobordant]。此处暂不指定定向；有向版本还需规定边界的诱导定向。
\end{definition}

\begin{theorem}[同伦的横截逆像]\label{g2:thm:transverse-homotopy-cobordism}
	设 \(M^m\) 是紧无边界光滑流形，\(N\) 是无边界光滑流形，\(S\subset N\) 是闭、余维 \(c\le m\) 的嵌入子流形。若两个光滑映射 \(f_0,f_1:M\to N\) 均横截于 \(S\)，且连续同伦，则存在光滑同伦 \(H:[0,1]\times M\to N\)，在两个端点的时间邻域分别恒等于 \(f_0,f_1\)，并且 \(H\pitchfork S\)。于是
	\[
		W=H^{-1}(S),\qquad
		\partial W=\bigl(\{0\}\times f_0^{-1}(S)\bigr)
		\amalg\bigl(\{1\}\times f_1^{-1}(S)\bigr)
	\]
	给出两个逆像之间的配边。
\end{theorem}
\begin{proof}
	从一条连续同伦 \(F\) 出发，取光滑非降函数 \(\lambda:[0,1]\to[0,1]\)，在0附近等于0、在1附近等于1。把 \(F\) 换成 \(F(\lambda(t),x)\)，所得同伦在两端时间邻域中等于光滑的 \(f_0,f_1\)。\cref{g2:thm:relative-smooth-approximation}适用于这个带边界源，给出保持较小端点邻域的光滑同伦 \(F'\)。

	在这些端点邻域内，\(F'\) 以及它的边界限制都横截于 \(S\)，因为 \(f_0,f_1\) 已横截。应用刚证明的带边界相对横截逼近，固定再小一些的端点邻域，得到所需的 \(H\)。\cref{g2:thm:boundary-transverse-preimage}给出 \(W\) 的流形结构及边界公式。最后，\(S\) 闭，所以 \(W\) 是紧空间 \([0,1]\times M\) 的闭子集，因而紧；这一点正是配边定义中的紧性要求。
\end{proof}

\begin{example}\label{g2:exm:zeros-appear-in-pairs}
	在 \(S^1\) 上考虑 \(f_a(e^{\mathrm i\theta})=\cos\theta-a\)。当 \(|a|<1\) 时，零点有两个且都是正则零点；当 \(|a|>1\) 时没有零点。沿 \(a\) 从0变到2的总体函数，其参数导数等于 \(-1\)，故总体逆像是光滑一维流形，虽然 \(a=1\) 的单张切片有临界零点。两个零点的轨迹在该临界时刻接合并消失。这说明总体横截足以构造配边，却不要求每个时刻的切片都横截。
\end{example}

一般地，当 \(m=c\) 时，横截逆像是有限点集。上述配边是紧一维流形，其边界点数为偶数。可以直接用有限图解释这件事：取有限条闭坐标弧，其相对内部覆盖 \(W\)，遇到边界点时使用半区间图。把所有弧端点和 \(\partial W\) 放入有限集 \(P\)。任一坐标弧删去 \(P\) 后的各开段，在 \(W\setminus P\) 中既开又闭，故其连通分支就是有限条这样的开段。这把 \(W\) 写成有限图：内部顶点的度数为2，边界顶点的度数为1。由各边恰有两个端点，全部顶点度数之和为边数的两倍，故度数为1的顶点有偶数个。因此 \(f_0^{-1}(S)\) 与 \(f_1^{-1}(S)\) 的点数模2相同。后续相交数理论将在加入定向后，把这个奇偶计数细化为带符号计数。


\begin{chaptersummary}
	局部卷积提供平滑操作，截断与局部有限覆盖将它拼成保持闭集附近数据的整体逼近。导数的控制阶数决定哪些性质稳定；非紧源上还要区分逐紧集控制与处处给定容差。光滑 Sard 定理把临界集分为可降维的部分和高阶平坦的部分，分别以归纳、Fubini 和覆盖体积估计处理。横截性要求微分在约束的法方向上满射，因此逆像具有预期维数，法丛则是目标法丛的拉回。参数横截定理将切片的横截性等价为约束流形投影的正则值条件；局部坐标扰动使这一结论可以用于任意映射，并保持指定区域。源带边界时，边界限制映射须另外横截，逆像的边界才恰是端点或边界数据的逆像。对于紧源与闭约束，横截同伦的逆像给出配边，进而保持零维逆像点数的奇偶性。
\end{chaptersummary}

\BookExercises

\begin{exercise}\label{g2:ex:cr-convergence-orders}
	对整数 \(k\ge0\)，在 \(S^1\) 上定义 \(h_n(e^{\mathrm{i}\theta})=n^{-k-1}\sin(n\theta)\)。判断 \(h_n\) 在哪些有限阶 \(C^r\) 拓扑中趋于零，并说明为何它不在 \(C^\infty\) 拓扑中趋于零。
\end{exercise}
\begin{solution}
	对第 \(j\) 阶角变量导数，其最大绝对值为 \(n^{j-k-1}\)。故当 \(r\le k\) 时，所有 \(j\le r\) 的导数一致趋于零，得到 \(C^r\) 收敛。第 \(k+1\) 阶导数的最大绝对值恒为 \(1\)，所以 \(r\ge k+1\) 时不收敛。圆周可用有限个角坐标块覆盖，且换角坐标只差常数，以上计算确实检验了圆周的 \(C^r\) 拓扑。\(C^\infty\) 收敛要求每个有限阶都收敛，第 \(k+1\) 阶已经排除了它。
\end{solution}

\begin{exercise}\label{g2:ex:sphere-valued-approximation}
	设 \(M\) 为紧光滑流形，\(f:M\to S^q\subset\mathbb R^{q+1}\) 连续，\(q\ge1\)。若光滑 \(u:M\to\mathbb R^{q+1}\) 满足 \(\sup_M|u-f|<1/2\)，证明 \(g=u/|u|\) 为光滑映射 \(M\to S^q\)，且 \(g\) 与 \(f\) 同伦。给出 \(\sup_M|g-f|\) 的误差上界，并说明若 \(u=f\) 于一个区域，同伦在该区域有何性质。
\end{exercise}
\begin{solution}
	因为 \(|f|=1\)，有 \(|u|\ge1-|u-f|>1/2\)，所以归一化有定义且光滑。又
	\[
		|g-f|\le|g-u|+|u-f|=\bigl|1-|u|\bigr|+|u-f|\le2|u-f|,
	\]
	故 \(\sup_M|g-f|\le2\sup_M|u-f|\)。令 \(v_t=(1-t)f+tu\)，则 \(|v_t-f|=t|u-f|<1/2\)，于是 \(v_t\ne0\)，而 \(H(x,t)=v_t(x)/|v_t(x)|\) 是从 \(f\) 到 \(g\) 的连续同伦。若 \(u=f\) 于某区域，那里 \(v_t=f\) 且 \(|f|=1\)，故整个同伦固定不动。
\end{solution}

\begin{exercise}\label{g2:ex:uniform-not-strong-convergence}
	取非负 \(b\in C^\infty_c(\mathbb R)\)，使 \(b(0)=1\)、\(\operatorname{supp}b\subset[-1/4,1/4]\)，并令 \(u_n(x)=n^{-1}b(x-n)\)。证明 \(u_n\) 在每个有限阶的全局一致导数范数下趋于零，但不在强 \(C^0\) 拓扑中趋于零。
\end{exercise}
\begin{solution}
	对固定 \(r\)，有
	\[
		\max_{0\le j\le r}\sup_{x\in\mathbb R}|u_n^{(j)}(x)|
		=\frac1n\max_{0\le j\le r}\sup_{x\in\mathbb R}|b^{(j)}(x)|\longrightarrow0.
	\]
	强拓扑允许位置相关的误差。在局部有限紧块 \(K_m=[m-1/2,m+1/2]\)、\(m\in\mathbb Z\) 上同时要求 \(\sup_{K_m}|v|<e^{-|m|}/2\)，得到零映射的一个强 \(C^0\) 邻域。然而 \(u_n(n)=1/n>e^{-n}/2\)，故没有任何 \(u_n\) 属于这个邻域。这也表明非紧源的强拓扑比全局一致拓扑更细。
\end{solution}

\begin{exercise}\label{g2:ex:critical-collapse}
	设 \(f:\mathbb R^2\to\mathbb R^2\)，\(f(x,y)=(x,xy)\)。求临界点集、临界值集及所有正则值的原像。说明临界点集的维数与临界值集的维数为何不同。
\end{exercise}
\begin{solution}
	导数矩阵为 \(\begin{pmatrix}1&0\\y&x\end{pmatrix}\)，行列式为 \(x\)，故临界点集为直线 \(\{0\}\times\mathbb R\)。这条直线全部被映到 \((0,0)\)，所以临界值集只有 \(\{(0,0)\}\)。对于 \(a\ne0\)，值 \((a,b)\) 的原像为唯一的 \((a,b/a)\)，它正则；对于 \(a=0,b\ne0\)，原像为空，也正则。唯一的临界值 \((0,0)\) 的原像是整条临界直线。这里一个一维临界点集被压成零维的临界值集，Sard 定理控制的是后者在目标中的测度。
\end{solution}

\begin{exercise}\label{g2:ex:regular-values-not-open}
	令 \(f:\mathbb R\to\mathbb R\)，\(f(x)=e^{-x}\sin x\)。证明 \(0\) 是正则值，但正则值集不是开集。再证明：若 \(M,N\) 为无边界光滑流形、\(f:M\to N\) 光滑且固有，则正则值集是开集。这里固有指紧集的原像紧。
\end{exercise}
\begin{solution}
	\(f'(x)=e^{-x}(\cos x-\sin x)\)，故临界点恰为 \(x_k=\pi/4+k\pi\)，\(k\in\mathbb Z\)，对应临界值为 \((-1)^ke^{-x_k}/\sqrt2\)。这些值均不为零，并在 \(k\to+\infty\) 时趋于零。另一方面，\(f^{-1}(0)=\{k\pi:k\in\mathbb Z\}\)，各点导数为 \((-1)^ke^{-k\pi}\ne0\)，故 \(0\) 正则，而它的任意邻域都含临界值。

	一般情形下，临界点集 \(C\subset M\) 闭，因为满秩矩阵构成开集。设临界值序列 \(y_j=f(x_j)\) 收敛到 \(y\)，其中 \(x_j\in C\)。取 \(y\) 的一个相对紧坐标邻域；去掉有限项后，\(y_j\) 和 \(y\) 位于某个紧坐标块 \(K\) 内。固有性给出 \(f^{-1}(K)\) 紧，故 \(x_j\) 有子列收敛到 \(x\in C\)，且连续性给出 \(f(x)=y\)。于是临界值集闭，其补集即正则值集开。首个例子不固有：任意包含 \(0\) 的小闭区间的原像都包含充分大的正实数，因而不紧。
\end{solution}

\begin{exercise}\label{g2:ex:c1-small-dimensional-image}
	设 \(U\subset\mathbb R^m\) 开，\(f:U\to\mathbb R^n\) 为 \(C^1\) 映射，且 \(m<n\)。不使用有限阶 Sard 定理，直接证明 \(f(U)\) 为 \(n\) 维零测集，并推广到 \(C^1\) 流形映射。
\end{exercise}
\begin{solution}
	先在一个闭立方体 \(Q\subset U\) 上工作。导数有界给出 \(Q\) 上的 Lipschitz 常数 \(L\)。将边长为 \(R\) 的 \(Q\) 分成 \(s^m\) 个小立方体，每个像均包含在半径为 \(L\sqrt mR/s\) 的球中，故 \(f(Q)\) 的覆盖总体积至多为
	\[
		s^m\omega_n(L\sqrt mR/s)^n=C_Qs^{m-n}\longrightarrow0.
	\]
	用可数个这样的闭立方体覆盖 \(U\)，即得结论。对于流形，取可数坐标图并将源图限制到目标图的原像，使坐标表达成为上述 \(C^1\) 映射。坐标图之间的 \(C^1\) 变换在紧块上 Lipschitz，保持同维零测集，再取可数并即可。源有边界时，在半立方体上使用局部 \(C^1\) 延拓，得到同样估计。
\end{solution}

\begin{exercise}\label{g2:ex:transversality-saddle-graph}
	设 \(f:\mathbb R^2\to\mathbb R^3\)，\(f(x,y)=(x,y,x^2-y^2)\)。分别对竖直轴 \(L=\{(0,0,z):z\in\mathbb R\}\) 及水平面 \(P_a=\{z=a\}\) 判断 \(f\) 的横截性，并写出原像的维数及切空间。
\end{exercise}
\begin{solution}
	与 \(L\) 相交只发生在 \((x,y)=(0,0)\)。此处 \(\mathrm{d}f\) 的像为水平的 \(xy\) 平面，与 \(\mathrm{T}_0L\) 的竖直方向共同张成 \(\mathbb R^3\)，故 \(f\pitchfork L\)。原像为单点，维数 \(2-2=0\)，切空间为零。

	\(P_a\) 的法方向由第三坐标给出，复合函数为 \(h(x,y)=x^2-y^2\)，微分为 \(\mathrm{d}h_{(x,y)}(u,v)=2xu-2yv\)。其唯一临界点为原点，临界值为零。因此 \(a\ne0\) 时 \(f\pitchfork P_a\)，原像为一维双曲线，切空间为 \(\{(u,v):xu-yv=0\}\)。\(a=0\) 时原像为两条直线 \(y=x\) 与 \(y=-x\) 的并，原点处不横截，该交叉在原点附近不是一维流形；离开原点后仍有上述一维切空间。
\end{solution}

\begin{exercise}\label{g2:ex:parametric-sufficient-not-necessary}
	设 \(F:\mathbb R^2\to\mathbb R^2\)，\(F(a,t)=(t,t^3-at)\)，目标子流形为横轴 \(S\)。求出哪些 \(a\) 使 \(F_a\pitchfork S\)，并检查总体映射 \(F\) 是否与 \(S\) 横截。这说明参数定理的总体假设是何种条件？
\end{exercise}
\begin{solution}
	切片的交点满足 \(t(t^2-a)=0\)，法方向导数为 \(3t^2-a\)。当 \(a<0\) 时只有 \(t=0\) 一个交点，导数为 \(-a\ne0\)；当 \(a>0\) 时有 \(0,\pm\sqrt a\) 三个交点，对应导数为 \(-a,2a,2a\)，均不为零。当 \(a=0\) 时只有零点，导数也为零。因此好参数恰为 \(\mathbb R\setminus\{0\}\)。

	总体映射的法分量为 \(H(a,t)=t^3-at\)，其微分为 \(-t\,\mathrm{d}a+(3t^2-a)\,\mathrm{d}t\)，在 \((0,0)\) 处为零，且此点映入 \(S\)。故 \(F\) 不横截。几乎处处参数横截仍然成立，说明总体横截性是参数定理给出的充分条件，并非其结论成立的必要条件。
\end{solution}

\begin{exercise}\label{g2:ex:mobius-regular-zero}
	把 Möbius 线丛的总空间写成 \(E=(\mathbb R\times\mathbb R)/((t,v)\sim(t+1,-v))\)，零截面记为 \(S\)。证明不存在光滑函数 \(h:E\to\mathbb R\)，同时满足 \(h^{-1}(0)=S\) 与 \(0\) 为正则值。说明此现象与法丛有什么关系。
\end{exercise}
\begin{solution}
	若这样的 \(h\) 存在，令 \(H(t,v)=h([t,v])\)。则 \(H(t,0)=0\)，故 \(\partial H/\partial t(t,0)=0\)。正则值条件要求 \(b(t)=\partial H/\partial v(t,0)\) 处处非零。商关系给出 \(H(t+1,-v)=H(t,v)\)，对 \(v\) 求导得 \(b(t+1)=-b(t)\)。连续函数 \(b\) 在任何 \([t,t+1]\) 的两端异号，由介值定理必有零点，矛盾。

	从法丛看，若 \(h\pitchfork\{0\}\)，\cref{g2:prop:transverse-normal-bundle}会给出 \(\nu_E S\cong S\times\mathbb R\)。但零截面的法方向是纤维方向，沿底圆周绕一周发生符号翻转，正是 Möbius 线丛本身。上述 \(b(t)\) 的符号论证也直接证明了这个法线丛不能平凡化。
\end{solution}

\begin{exercise}\label{g2:ex:finite-simultaneous-transversality}
	设 \(M\) 为紧无边界光滑流形，\(N\) 为无边界光滑流形，\(S_1,\ldots,S_k\subset N\) 为闭的无边界嵌入子流形。在 \(C^\infty(M,N)\) 的 \(C^\infty\) 拓扑中，证明同时与所有 \(S_i\) 横截的映射构成开稠密集。能否直接由各次横截性推出映射也与所有非空交集 \(S_i\cap S_j\) 横截？
\end{exercise}
\begin{solution}
	由\cref{g2:prop:transversality-open-on-compact}及\cref{g2:thm:relative-transversality}，每个集合 \(\mathcal T_i=\{f:f\pitchfork S_i\}\) 都开且稠密。其有限交开。给定任意非空开集 \(U\)，先与 \(\mathcal T_1\) 相交，所得仍为非空开集；再依次与其余 \(\mathcal T_i\) 相交，每次由稠密性保留非空性。故有限交也稠密。

	最后的推断不成立，即使交集本身是子流形也如此。取 \(M=S^1\)、\(N=\mathbb R^2\)，\(f(\theta)=(\sin\theta,\sin\theta)\)，令 \(S_1,S_2\) 为两条坐标轴。两类交点都在 \(\theta=0,\pi\)，导数为 \((\cos\theta,\cos\theta)\)，与每条轴分别张成平面，所以两次横截性均成立。但是 \(S_1\cap S_2=\{0\}\)，从一维切空间到二维切空间的微分不可能满射，故 \(f\) 不与该交集横截。同时满足各个子流形的横截条件，并不等于对其交集施加了独立的全部法向约束。
\end{solution}

\begin{exercise}\label{g2:ex:boundary-parabola-levels}
	在上半平面 \(M=\{(x,y):y\ge0\}\) 上令 \(f_a(x,y)=y+x^2-a\)。分别判断 \(f_a\) 与 \(f_a|_{\partial M}\) 何时横截于 \(\{0\}\)，求 \(f_a^{-1}(0)\) 及其边界，并解释 \(a=0\) 时何处失去预期维数。
\end{exercise}
\begin{solution}
	完整微分为 \(2x\,\mathrm{d}x+\mathrm{d}y\)，处处满射，故所有 \(a\) 都使 \(f_a\pitchfork\{0\}\)。边界限制为 \(x^2-a\)：若 \(a>0\)，零点为 \(\pm\sqrt a\)，导数分别为 \(\pm2\sqrt a\)，因而横截；若 \(a<0\)，没有零点，也横截；若 \(a=0\)，零点处导数为零，不横截。

	当 \(a>0\) 时，逆像是弧 \(\{(x,a-x^2):-\sqrt a\le x\le\sqrt a\}\)，以 \(x\) 为参数得到一维带边界流形，边界是 \((\pm\sqrt a,0)\)。当 \(a<0\) 时逆像为空。当 \(a=0\) 时逆像只有原点，维数为零，而无边界正则值公式会给出 \(2-1=1\)。此时约束将唯一的向内方向消耗掉，边界切向方向却无法补偿，正是附加边界条件失效的位置。
\end{solution}

\begin{exercise}\label{g2:ex:relative-value-versus-neighborhood}
	令 \(f(x)=x^2\)、\(A=\{0\}\subset\mathbb R\)、\(S=\{0\}\subset\mathbb R\)。证明不存在与 \(f\) 在 \(A\) 的某邻域相同且横截于 \(S\) 的光滑映射。若只要求 \(g(0)=f(0)\)，构造在每个固定紧区间上任意 \(C^r\) 接近 \(f\) 且横截于 \(S\) 的光滑映射，\(r<\infty\)。
\end{exercise}
\begin{solution}
	若 \(g=f\) 于零点的某邻域，则 \(g(0)=0\)、\(g'(0)=0\)，所以 \(g\) 在这个原像点不横截。只固定点值时，取 \(g_\varepsilon(x)=x^2+\varepsilon x\)，\(\varepsilon\ne0\)。它有两个零点 \(0,-\varepsilon\)，其导数分别为 \(\varepsilon,-\varepsilon\)，均非零，故处处横截于 \(S\)。在 \([-R,R]\) 上，差的零阶上界为 \(|\varepsilon|R\)，一阶上界为 \(|\varepsilon|\)，更高阶差为零，故所有有限阶都随 \(\varepsilon\to0\) 逼近。保持邻域内的映射会同时固定导数，只固定闭集上的值则未必如此。
\end{solution}

\begin{exercise}\label{g2:ex:circle-regular-zero-parity}
	设 \(h:S^1\to\mathbb R\) 光滑，且 \(0\) 为正则值。用横截同伦的逆像证明 \(h^{-1}(0)\) 有偶数个点。再对每个整数 \(k\ge0\) 给出恰有 \(2k\) 个正则零点的例子。
\end{exercise}
\begin{solution}
	正则零点孤立，而零点集是紧圆周的闭子集，故它为有限点集。\(h\) 与常值函数 \(1\) 通过 \((1-t)h+t\) 连续同伦，两端映射都横截于 \(\{0\}\)。\cref{g2:thm:transverse-homotopy-cobordism}给出一个紧一维流形，其边界为 \(h^{-1}(0)\) 与空集的不交并。正文的有限图度数论证说明其边界点数为偶数，故结论成立。

	当 \(k=0\) 时取常值函数 \(1\)。当 \(k\ge1\) 时取 \(h(e^{\mathrm i\theta})=\sin(k\theta)\)。零点为 \(\theta=j\pi/k\)、\(0\le j<2k\)，各点导数为 \(k(-1)^j\ne0\)，恰得 \(2k\) 个正则零点。
\end{solution}

\begin{exercise}\label{g2:ex:zeros-escape-to-infinity}
	取 \(M=\mathbb R\)，考虑 \(H:[0,1]\times M\to\mathbb R\)、\(H(t,x)=e^x-t\)。证明总体映射及其边界限制都横截于 \(\{0\}\)，求 \(H^{-1}(0)\)。解释为何两端的正则零点数奇偶性不同，以及这为何不违背本章的配边结论。
\end{exercise}
\begin{solution}
	\(\mathrm{d}H=-\mathrm{d}t+e^x\,\mathrm{d}x\) 处处满射。\(t=0\) 时函数 \(e^x\) 没有零点，横截条件自动成立；\(t=1\) 时只有 \(x=0\) 一个零点，导数为 \(1\)，也横截。总体逆像为
	\[
		W=\{(t,\log t):0<t\le1\},
	\]
	它是与 \((0,1]\) 微分同胚的一维带边界流形，边界只有 \((1,0)\)。当 \(t\downarrow0\) 时，\(x=\log t\to-\infty\)，所以 \(W\) 非紧，缺失的端点逃向源的无穷远处。两端分别有0个和1个零点；配边结论要求源紧，正是为了保证逆像紧并排除这种现象。
\end{solution}

% !TeX root = ../../Book2.tex
% 纲要标题：## 第15章　嵌入、邻域与光滑族的局部平凡性
% 纲要位置：几何学计划纲要.md，第 1572 行。
% 正文状态：初稿；尚未正式审核。

% 纲要内容开始
% 
% > 本章目标与证明范围：将嵌入、管状邻域、同位延拓与 Ehresmann 按所需局部模型连接。弱嵌入维数界与强 Whitney 界分别安排证明，不能用一次泛型投影跨过消去双点的步骤。
% 
% 纲要内容结束

\chapter{嵌入、邻域与光滑族的局部平凡性}
\label{g2:ch:15}
\BookChapterNumber{g2:ch:15}

\begin{chapterabstract}
	嵌入把抽象流形放进 Euclid 空间，管状邻域把子流形附近的几何化为法丛，领状邻域则把边界附近化为乘积。本章先由坐标和泛型投影建立弱 Whitney 界，再补充双点消去与低维构造，得到强嵌入界。随后用逆函数定理、变半径控制和流构造邻域模型，并把子流形的同位延拓为环境运动。最后证明固有满淹没局部平凡：水平提升给出纤维间的搬运，固有性保证提升存在到所需时刻。边界条件和复环面族说明，局部平凡性的结论取决于所处的光滑或复解析范畴。
\end{chapterabstract}

\begin{learninggoals}
	\begin{itemize}
		\item 从有限坐标与光滑截断构造初始嵌入，计算投影坏方向的维数，并解释非紧源为何还需固有性控制。
		\item 辨认一般位置中的横截双点，说明 Whitney 盘、标架及局部移动如何参与双点消去，并分清低维情形。
		\item 用法向局部指数和变半径缩小构造管状邻域，理解法丛正规化在选择无关性中的作用。
		\item 从向内横截向量场构造领圈，利用有符号的领圈参数建立双化与边界粘合。
		\item 把同位的速度延拓为时间依赖向量场，用共同紧支集获得环境同位，并识别相对条件和非紧限制。
		\item 由水平提升证明 Ehresmann 定理，核对边界推广条件，并区分光滑局部平凡与复解析局部平凡。
	\end{itemize}
\end{learninggoals}

一条光滑曲线在每个点附近都像直线，一张曲面在每个点附近都像平面。这种局部坐标并没有立即给出整条曲线或整张曲面的统一表达。例如，沿一个圆周逐点选择短法线段，必须保证来自不同基点的线段不会相撞，才能用“圆周上的位置加法向位移”描述整个邻域；若圆周随时间运动，还要把不同位置上的速度组织成一个环境向量场。本章研究的正是这些从逐点局部模型到整体构造的步骤。

\ProseParagraphBreak

\secref{g2:sec:1501}先把流形放进高维空间，再寻找保留单射性与微分秩的投影方向。达到 \(2n+1\) 维时，维数比较足以避开所有割线方向；再降一维，一般位置会留下孤立双点。\secref{g2:sec:1502}专门处理这一层障碍。之后转向一个已经嵌入的子流形：\secref{g2:sec:1503}将其邻域与法丛联系起来，\secref{g2:sec:1504}给出边界的乘积邻域。这些邻域构造又使\secref{g2:sec:1505}能够延拓局部规定的运动。

\ProseParagraphBreak

最后考虑一个光滑参数族 \(p:E\to B\)。淹没条件说明每个点附近都有投影坐标，却未保证整条纤维能同时装进同一个乘积图。Ehresmann\footnote{Charles Ehresmann（1905—1979），法国数学家，研究微分拓扑、纤维丛与范畴论。}定理用固有性补上这一环：紧参数范围内的所有提升轨道留在一个紧集里，因而能把一条纤维完整地搬到相邻纤维。读到这里时，应特别留意“每一点的解短时间存在”与“全部初值都能走完指定时间”之间的差别。

\ProseParagraphBreak

本章默认流形有限维、Hausdorff、第二可数，未特别说明边界时均指无边界流形。前章的 Sard 定理、参数横截性与相对扰动用于选取一般位置；第一卷第11—13章的逆函数定理、单位分解、常微分方程光滑依赖与逃逸判据用于邻域和流。用到这些结论时会恢复具体条件。弱嵌入与流的参考见\cite[第6、9章]{Lee2012Smooth}；强 Whitney 界的双点构造参见\cite{Whitney1944SelfIntersections}。

% !TeX root = ../../Book2.tex
% 纲要标题：### 15.1 高维初始嵌入与泛型投影
% 纲要位置：几何学计划纲要.md，第 1576 行。
% 正文状态：初稿；尚未正式审核。

% 纲要内容开始
% 
% * 由局部坐标构造整体嵌入
% * 投影方向的坏集合
% * 弱 Whitney 维数界
% * 非紧源的穷竭及固有性
% 
% 纲要内容结束

\section{高维初始嵌入与泛型投影}
\label{g2:sec:1501}

一张坐标图把流形的一小块放进 Euclid 空间，但不同图中的坐标不能直接首尾相接。先让每张图负责几个独立的坐标分量，便能把局部信息全部保存下来；得到一个维数可能很高的嵌入后，再寻找不会压坏切向量、也不会合并不同点的投影方向。紧流形上的证明由这两步完成。非紧流形还须控制无穷远，否则单射浸入的像仍可能发生拓扑上的聚集。

\subsection{由局部坐标构造整体嵌入}
\label{g2:subsec:1501:01}

先回忆第一卷第11章第4节的区别：浸入要求微分单射；嵌入还要求映射是到子空间像的同胚。对紧源，连续单射到 Hausdorff 空间自动是到像的同胚，因此只需构造单射浸入。

\begin{proposition}\label{g2:prop:finite-coordinate-embedding}
	每个紧光滑 \(n\) 维流形 \(M\)，允许有边界，都能光滑嵌入某个有限维 Euclid 空间。
\end{proposition}
\begin{proof}
	由第一卷第12章第3节的相对紧坐标细化与截断构造，取有限张图 \((U_i,x_i)\)，以及较小开集 \(V_i\) 和光滑函数 \(\rho_i:M\to[0,1]\)，使 \(V_i\) 覆盖 \(M\)，\(\rho_i=1\) 于 \(\overline V_i\) 的邻域，\(\operatorname{supp}\rho_i\Subset U_i\)。把 \(\rho_i x_i\) 在 \(U_i\) 外补零，定义
	\[
		F(p)=\bigl(\rho_1(p),\rho_1(p)x_1(p),\ldots,
		\rho_s(p),\rho_s(p)x_s(p)\bigr)\in\mathbb R^{s(n+1)}.
	\]
	补零光滑，因为每个支集离相应图域的边缘留有邻域。

	若 \(F(p)=F(q)\)，选 \(i\) 使 \(p\in V_i\)。则 \(\rho_i(p)=\rho_i(q)=1\)，尤其 \(q\in U_i\)，而对应的向量块给出 \(x_i(p)=x_i(q)\)，故 \(p=q\)。在 \(p\in V_i\) 附近，第 \(i\) 个向量块就是 \(x_i\)，其微分单射，故 \(\mathrm{d}F_p\) 单射。最后，紧性保证连续逆，因而 \(F\) 为嵌入。边界图中的完整坐标微分同样可逆，所以证明也包括带边界情形。
\end{proof}

保留额外的标量 \(\rho_i\) 有实际作用：它告诉我们哪个坐标块正在工作，并允许从乘积 \(\rho_i x_i\) 恢复坐标。若只把加权坐标相加到同一个 \(\mathbb R^n\) 中，来自不同图的信息可能相互抵消，也可能把不同点混在一起。

\subsection{投影方向的坏集合}
\label{g2:subsec:1501:02}

设已有单射浸入 \(F:M^n\to\mathbb R^q\)。对单位向量 \(v\)，记 \(\pi_v:\mathbb R^q\to v^\perp\) 为正交投影。投影可能以两种方式失败：某条非零切向量被压成零，或者连接两个像点的割线被压成零。这两种情况都能把“坏方向”写成一个光滑映射的像。

\ProseParagraphBreak

在 \(M\) 上选一个辅助丛度量，记单位切丛为 \(U\mathrm{T}M\)。第一卷第12章第3节的丛度量存在定理保证这种选择；单位球在每个局部标架中是光滑超曲面，故 \(\dim U\mathrm{T}M=2n-1\)。考虑
\[
	\tau:U\mathrm{T}M\longrightarrow S^{q-1},\qquad
	\tau(x,u)=\frac{\mathrm{d}F_x(u)}{|\mathrm{d}F_x(u)|},
\]
以及去掉对角线的点对空间上的映射
\[
	\sigma:(M\times M)\setminus\Delta_M\longrightarrow S^{q-1},\qquad
	\sigma(x,y)=\frac{F(x)-F(y)}{|F(x)-F(y)|}.
\]
分母分别由浸入性和单射性保证非零。\(v\) 属于 \(\tau\) 的像，恰好意味着某个非零切向量落入 \(\ker\pi_v\)；\(v\) 属于 \(\sigma\) 的像，恰好意味着某对不同点被投影合并。由于 \(u\) 与 \(-u\)、\((x,y)\) 与 \((y,x)\) 都出现，这两个像已同时包含正、反方向。

\begin{proposition}\label{g2:prop:good-projection-directions}
	设 \(F:M^n\to\mathbb R^q\) 是光滑单射浸入，\(n\ge1\)，\(M\) 允许有边界。若 \(q>2n\)，则几乎每个单位方向 \(v\) 使 \(\pi_v\circ F\) 为浸入；若 \(q>2n+1\)，则几乎每个 \(v\) 同时使它单射。好方向在单位球中稠密。
\end{proposition}
\begin{proof}
	当 \(q>2n\) 时，\(\dim U\mathrm{T}M=2n-1<q-1\)，由\cref{g2:cor:sard-lower-dimension}，\(\tau\) 的像零测。当 \(q>2n+1\) 时，点对空间维数 \(2n<q-1\)，\(\sigma\) 的像也零测；两者的并仍零测。任一非空球面开集有正测度，所以补集稠密。

	若源有边界，将单位切丛分为底点在内部与边界的两部分，维数分别为 \(2n-1\)、\(2n-2\)。点对空间则分为内部乘内部、内部乘边界、边界乘内部、边界乘边界四个无边界流形，维数分别至多 \(2n,2n-1,2n-1,2n-2\)。在每一部分应用同一低维像结论，再取有限并即可。\(n=0\) 时没有切向条件；点对空间可数，割线方向的像也可数。
\end{proof}

\subsection{弱 Whitney 维数界}
\label{g2:subsec:1501:03}

\begin{theorem}[弱 Whitney 嵌入定理]\label{g2:thm:weak-whitney-compact}
	设 \(M\) 是紧光滑 \(n\) 维流形，允许有边界。则 \(M\) 可光滑嵌入 \(\mathbb R^{2n+1}\)。若 \(n\ge1\)，还存在光滑浸入 \(M\to\mathbb R^{2n}\)。
\end{theorem}
\begin{proof}
	若 \(n=0\)，紧零维流形是有限点集，直接选直线中的相同数目互异点即可。以下设 \(n\ge1\)。先由\cref{g2:prop:finite-coordinate-embedding}得到某个 \(F:M\to\mathbb R^q\)。若 \(q\le2n+1\)，补零坐标即可；否则用\cref{g2:prop:good-projection-directions}选好方向，将维数降一，所得仍是单射浸入。紧性使它仍为嵌入，所以可以重复，直到环境维数为 \(2n+1\)。从这里再投影一次，切向坏方向的维数仍小于目标方向球的维数，因此能保留浸入性，得到 \(\mathbb R^{2n}\) 中的浸入。
\end{proof}

最后一次投影有一个明确变化：割线映射的源与方向球同为 \(2n\) 维，零测论证已经不能排除其像覆盖一块开区域。于是微分仍可单射，不同点却可能投到一起。下一节的强界正要处理这些双点。

\begin{example}\label{g2:exm:circle-bad-projection}
	单位圆 \(S^1\subset\mathbb R^2\) 沿竖直方向投到横轴，成为函数 \((x,y)\mapsto x\)。圆的最高点和最低点有相同的像，而最左点和最右点的切线竖直，在那里投影微分为零。对圆周投到直线，转动方向只会转动这些坏点；维数条件不足时，没有理由期待避开全部坏方向。
\end{example}

\subsection{非紧源的穷竭及固有性}
\label{g2:subsec:1501:04}

现在令 \(M\) 为无边界光滑流形，仍假设有限维、Hausdorff 且第二可数。单射浸入未必是嵌入；第一卷第11章第4节的固有单射浸入命题给出可用的判据：若紧集的原像都紧，则到像的逆连续。下面会在构造过程中一直保留这个条件。

\begin{lemma}\label{g2:lem:smooth-proper-exhaustion}
	每个这样的 \(M\) 都有非负光滑函数 \(h:M\to[0,\infty)\)，使每个次水平集 \(h^{-1}([0,c])\) 紧。
\end{lemma}
\begin{proof}
	取第一卷第12章第3节构造中的紧穷竭 \(K_j\subset\operatorname{int}K_{j+1}\)，\(\bigcup_jK_j=M\)。以光滑截断选 \(\chi_j:M\to[0,1]\)，在 \(K_j\) 附近等于1，支集包含于 \(\operatorname{int}K_{j+1}\)。定义
	\[
		h=\sum_{j\ge1}(1-\chi_j).
	\]
	任一点附近属于某个 \(K_m\) 的内部，当 \(j\ge m\) 时对应项恒为零，所以该和局部有限且光滑。若 \(x\notin K_{m+1}\)，则前 \(m\) 项均为1，故 \(h(x)\ge m\)。对给定 \(c\)，选整数 \(m>c\)，次水平集便是 \(K_{m+1}\) 的闭子集，因此紧。
\end{proof}

为了得到有限维的初始嵌入，把 \(h\) 的相邻壳层分配到两个坐标块。由\cref{g2:thm:sard}，对每个整数 \(j\ge0\) 可选正则值
\[
	j<a_j<b_j<j+1.
\]
令 \(D_0=h^{-1}([0,1])\)、\(E_0=h^{-1}([0,a_1])\)，并对 \(j\ge1\) 令
\[
	D_j=h^{-1}([j,j+1]),\qquad E_j=h^{-1}([b_{j-1},a_{j+1}]).
\]
这些 \(E_j\) 是紧带边界流形，\(D_j\subset\operatorname{Int}E_j\)，且同为偶数或同为奇数的不同 \(E_j\) 两两不交。正则值条件在边界附近把 \(h\) 变成一个坐标，因此 \(E_j\) 的边界确为相应水平集，不会出现尖角。

\ProseParagraphBreak

由紧源情形，取嵌入 \(f_j:E_j\to\mathbb R^{2n+1}\)，以及 \(\beta_j\) 在 \(D_j\) 附近等于1、支集包含于 \(\operatorname{Int}E_j\) 的光滑截断。补零后定义
\[
	F(x)=\left(\sum_{j\text{ 偶}}\beta_j(x)f_j(x),\quad
	\sum_{j\text{ 奇}}\beta_j(x)f_j(x),\quad h(x)\right)
	\in\mathbb R^{4n+3}.
\]
每个坐标块在一点附近至多出现一项，且壳层族局部有限，所以 \(F\) 光滑。若 \(F(x)=F(y)\)，先由最后一坐标知 \(h(x)=h(y)\)；选 \(j\) 使 \(x\in D_j\)，便有 \(y\in D_j\)。相应奇偶块给出 \(f_j(x)=f_j(y)\)，从而 \(x=y\)。在 \(x\in D_j\) 附近，这一块就是 \(f_j\)，其微分单射，所以 \(F\) 为浸入。最后坐标 \(h\) 保证 \(F\) 固有：紧目标集的原像是某个紧次水平集的闭子集。因此 \(F\) 已经是有限维空间中的固有嵌入。

\ProseParagraphBreak

普通投影未必保留固有性。为控制它，先把前 \(4n+2\) 个坐标用微分同胚
\[
	z\longmapsto\frac{z}{\sqrt{1+|z|^2}}
\]
压进单位球，保留最后坐标 \(h\)。所得固有嵌入的像包含于某条直线 \(L\) 周围的有界半径管中。若再选投影方向 \(v\notin L\)，这一“管轴”仍留下非零方向，不能全部被压掉。

\begin{lemma}\label{g2:lem:proper-projection-tube}
	设 \(F:M\to\mathbb R^q\) 为固有映射，其像包含于直线 \(L\) 周围的半径 \(R\) 管中。若单位向量 \(v\notin L\)，则 \(\pi_v\circ F\) 固有，且其像仍包含于直线 \(\pi_v(L)\) 周围某个有界半径管中。
\end{lemma}
\begin{proof}
	写 \(z=s+b\)，其中 \(s\in L\)、\(|b|\le R\)。因为 \(v\notin L\)，\(\pi_v|_L\) 单射，故存在 \(c>0\)，使 \(|\pi_v(s)|\ge c|s|\)。若 \(|\pi_v(z)|\le B\)，则
	\[
		c|s|\le|\pi_v(z)|+|\pi_v(b)|\le B+R,
	\]
	从而 \(|z|\le(B+R)/c+R\)。因此紧目标集的原像闭且包含于 \(F\) 的一个紧集原像中，因而紧。另一方面，\(\pi_v(z)\) 到 \(\pi_v(L)\) 的距离不超过 \(|\pi_v(b)|\le R\)，得到管的控制。
\end{proof}

\begin{theorem}[弱 Whitney 嵌入定理的固有形式]\label{g2:thm:weak-whitney-proper}
	每个有限维、Hausdorff、第二可数的无边界光滑 \(n\) 维流形都能固有地光滑嵌入 \(\mathbb R^{2n+1}\)。若 \(n\ge1\)，还存在到 \(\mathbb R^{2n}\) 的固有光滑浸入。
\end{theorem}
\begin{proof}
	若 \(n=0\)，第二可数的零维流形是至多可数的离散集，将各点依次送到不同非负整数便是固有嵌入。以下设 \(n\ge1\)，使用上面构造的有界管内固有嵌入。当环境维数 \(q>2n+1\) 时，先避开\cref{g2:prop:good-projection-directions}的零测坏方向，再避开管轴的两个单位方向。这样得到的投影仍单射、浸入且固有，因而仍是嵌入；\cref{g2:lem:proper-projection-tube}还保留了下一步所需的管。迭代到 \(q=2n+1\) 即得第一项。若 \(n\ge1\)，再选一次避开切向坏方向和管轴的投影，保留浸入性与固有性，得到第二项。
\end{proof}

以上壳层与管轴构造也见\cite[第132--135页，引理6.13--6.14及定理6.15]{Lee2012Smooth}。紧性和固有性在证明中承担了不同工作：前者把有限张坐标图拼成初始嵌入，后者阻止非紧源在投影后从无穷远重新聚到有限位置。关于带边界非紧流形，可在\secref{g2:sec:1504}构造双化以后，将双化流形的固有嵌入限制到原来的闭半边。

% !TeX root = ../../Book2.tex
% 纲要标题：### 15.2 强 Whitney 界的附加论证
% 纲要位置：几何学计划纲要.md，第 1583 行。
% 正文状态：未撰写；本文件目前是框架。

% 纲要内容开始
% 
% * 一般位置的自交模型
% * 消去双点所需的局部构造
% * 高维和低维特例的分开处理
% * 不能由一次泛型投影得到最优界
% 
% 纲要内容结束

\section{强 Whitney 界的附加论证}
\label{g2:sec:1502}

待撰写

\subsection{一般位置的自交模型}
\label{g2:subsec:1502:01}
待撰写

\subsection{消去双点所需的局部构造}
\label{g2:subsec:1502:02}
待撰写

\subsection{高维和低维特例的分开处理}
\label{g2:subsec:1502:03}
待撰写

\subsection{不能由一次泛型投影得到最优界}
\label{g2:subsec:1502:04}
待撰写
% 纲要条目排为编号小节标题，并进入目录与书签；说明段仍保留于纲要注释中。

% !TeX root = ../../Book2.tex
% 纲要标题：### 15.3 法丛与管状邻域
% 纲要位置：几何学计划纲要.md，第 1590 行。
% 正文状态：初稿；尚未正式审核。

% 纲要内容开始
% 
% * 法向指数映射
% * 零截面附近的逆函数论证
% * 非紧子流形的变半径控制
% * 管状邻域及其选择无关性
% 
% 纲要内容结束

\section{法丛与管状邻域}
\label{g2:sec:1503}

在子流形的一点附近，适配坐标把它写成 \(\mathbb R^k\times\{0\}\)。沿整个子流形移动时，这些横向坐标可能发生扭转，所以邻域通常不能写成一个固定乘积。法丛恰好记录横向方向怎样拼接。管状邻域定理说明：只要把每条纤维取得足够短，这个线性模型就能实际描述子流形的一个开邻域。

\subsection{法向指数映射}
\label{g2:subsec:1503:01}

设 \(S\subset N\) 是无边界流形 \(N\) 的无边界嵌入子流形。\secref{g2:sec:1403}已定义内在的法丛
\[
	\nu_N S=\mathrm{T}N|_S/\mathrm{T}S.
\]
一个法向量首先是模去切向量后的等价类。由第一卷第12章第3节的丛度量存在及正交补命题，在 \(N\) 上选光滑 Riemann 度量 \(g\)，便能把每个等价类 \(v\in\nu_NS\) 唯一地表示为正交向量 \(\iota_g(v)\in(\mathrm{T}S)^\perp\)。正交投影的局部矩阵由光滑正定矩阵的逆给出，所以 \(\iota_g\) 随基点光滑。

\ProseParagraphBreak

还要把这个向量变成附近的一个实际点。为此，在 \(\mathrm{T}N\) 上选择一个光滑联络 \(\nabla\)；第一卷第15章第1节用局部平凡联络的单位分解加权证明了其存在。这里只需要联络给出的短时曲线，度量负责选出法向代表。若坐标中的联络系数为 \(\Gamma^a_{bc}\)，初值为 \((x,v)\) 的曲线满足
\[
	\ddot z^a+\Gamma^a_{bc}(z)\dot z^b\dot z^c=0,
	\qquad z(0)=x,\quad\dot z(0)=v.
\]
把 \((z,\dot z)\) 当作一阶方程的未知量，第一卷第13章第1节的存在唯一性与初值光滑依赖定理适用。\(v=0\) 时解是常曲线，因此对 \((x,0)\) 的充分小邻域，解都能存在到时间1。令
\[
	\exp_x^\nabla(v)=z_{x,v}(1).
\]
它在 \(\mathrm{T}N\) 的零截面某邻域上光滑，称为该联络的局部指数映射。方程由 \(\nabla_{\dot z}\dot z=0\) 给出，与坐标选择相容。

\ProseParagraphBreak

在常解 \((x,0)\) 处，对初速度作一阶变化，二次速度项的一阶变化为零，故变分方程为 \(J''=0\)、\(J(0)=0\)、\(J'(0)=w\)。于是 \(J(1)=w\)，即
\[
	\mathrm{d}_v\exp_x^\nabla\big|_{v=0}=\operatorname{id}_{\mathrm{T}_xN}.
\]
同时 \(\exp_x^\nabla(0)=x\)，所以在基点方向上的微分也是恒等。定义法向映射
\[
	\Phi(v)=\exp_x^\nabla(\iota_gv),\qquad v\in\nu_xS.
\]
选择任意联络已足够；此处没有要求测地线在任意长时间存在，也没有使用最短路径性质。

\subsection{零截面附近的逆函数论证}
\label{g2:subsec:1503:02}

零截面处的切空间有规范分解
\[
	\mathrm{T}_{0_x}(\nu_NS)\cong\mathrm{T}_xS\oplus\nu_xS.
\]
第一项来自零截面的微分，第二项是纤维中的速度。在局部丛坐标中，这就是基底坐标与纤维坐标的分解；零纤维坐标使过渡函数中对基点求导的附加项消失，因此分解与平凡化无关。

\begin{proposition}\label{g2:prop:normal-exponential-local}
	对上面构造的 \(\Phi\)，有 \(\Phi(0_x)=x\)，且
	\[
		\mathrm{d}\Phi_{0_x}(u,v)=u+\iota_g(v).
	\]
	故 \(\Phi\) 在每个 \(0_x\) 的某邻域内为局部微分同胚。
\end{proposition}
\begin{proof}
	基点方向和纤维方向的微分已分别算出。由于 \(\mathrm{T}_xN=\mathrm{T}_xS\oplus(\mathrm{T}_xS)^\perp\)，右边是线性同构。第一卷第11章第3节的逆函数定理遂给出局部逆。
\end{proof}

这个结论还没有把整个零截面附近的不同纤维分开。即使每一点附近都有逆，来自相距很远的两个基点的短向量仍可能被送到同一点。紧子流形可用紧性统一选择半径；非紧情形需要让半径随位置缩小。

\begin{example}\label{g2:exm:shrinking-circle-tubes}
	在平面中取圆周的不交并
	\[
		S=\coprod_{j\ge1}\{(3j,0)+j^{-1}(\cos\theta,\sin\theta):\theta\in\mathbb R\}.
	\]
	这是闭的非紧嵌入一维子流形。采用 Euclid 度量和直线指数时，若统一法向半径为 \(\varepsilon>0\)，选 \(j\) 使 \(j^{-1}<\varepsilon\)，则第 \(j\) 个圆的所有向心法向量都可到达圆心，法向映射不再单射。各圆上的半径改为小于 \(1/(2j)\)，则这些问题消失。局部逆的存在并不保证统一的正半径。
\end{example}

\subsection{非紧子流形的变半径控制}
\label{g2:subsec:1503:03}

以下缩小论证既用于法向指数，也会用于比较不同的管状映射。保留紧参数族的版本，可以同时处理一段同位，而不用对无穷多个参数分别选邻域。

\begin{lemma}\label{g2:lem:zero-section-shrinking}
	设 \(S\subset N\) 是闭的无边界嵌入子流形，\(E\to S\) 为有丛度量的光滑向量丛，\(\dim E=\dim N\)。设 \(\Phi_t\) 是依赖于 \(t\in[0,1]\) 的联合光滑映射，在所有零向量附近有定义，并满足
	\[
		\Phi_t(0_x)=x,\qquad \mathrm{d}(\Phi_t)_{0_x}\text{ 为同构}.
	\]
	则存在连续函数 \(\varepsilon:S\to(0,\infty)\)，使每个 \(\Phi_t\) 在
	\[
		E_\varepsilon=\{v\in E:|v|<\varepsilon(\pi(v))\}
	\]
	上都是到 \(N\) 中开集的微分同胚；各局部逆对 \(t\) 也光滑。只有一个映射时，视为常参数族即可。
\end{lemma}
\begin{proof}
	由\cref{g2:lem:smooth-proper-exhaustion}，选非负光滑固有函数 \(h:N\to[0,\infty)\)，记 \(K_j=S\cap\{h\le j\}\)。闭性保证 \(K_j\) 紧。对每个 \(j\)，存在 \(r_j>0\)，使所有 \(t\) 下、基点在 \(K_j\) 上且长度小于 \(r_j\) 的向量都在定义域内，微分可逆，而且
	\[
		|h(\Phi_t(v))-h(\pi(v))|<1/3.
	\]
	这是对紧集 \([0,1]\times\{0_x:x\in K_j\}\) 使用有限覆盖和连续性。还可要求同一个 \(t\) 下，这些向量之间没有重合的像。若做不到，便有 \(t_i\)、不同的 \(v_i,w_i\)，基点都在 \(K_j\)，长度都趋零，且 \(\Phi_{t_i}(v_i)=\Phi_{t_i}(w_i)\)。取子列使 \(t_i\to t\)、\(\pi(v_i)\to x\)、\(\pi(w_i)\to y\)。极限等式给出 \(x=y\)。但映射 \((t,v)\mapsto(t,\Phi_t(v))\) 在 \((t,0_x)\) 的微分可逆，故其局部单射性排除了这样的序列。缩小后可令 \(r_j\) 单调不增。

	选连续正函数 \(\varepsilon(x)\)，使当 \(j-1\le h(x)\le j\) 时有 \(\varepsilon(x)<r_{j+3}\)。例如先在整数点选单调递减且足够小的正数，作分段线性插值，再与 \(h|_S\) 复合即可。必要时用\cref{g2:thm:relative-smooth-approximation}将这个正函数换成更小的光滑正函数。

	设 \(v,w\in E_\varepsilon\)，且 \(\Phi_t(v)=\Phi_t(w)\)。上述高度估计给出 \(|h(\pi(v))-h(\pi(w))|<2/3\)。若 \(j-1\le h(\pi(v))\le j\)，则两个基点都属于 \(K_{j+1}\)，而两个向量的长度都小于 \(r_{j+1}\)：第二个基点至多比第一个提前或推后一层，半径选择中预留的三层已涵盖这一差别。\(K_{j+1}\) 上的单射性于是给出 \(v=w\)。故每个 \(\Phi_t\) 在整个 \(E_\varepsilon\) 上单射。

	它的微分处处可逆，因此是开映射；单射的局部微分同胚便是到其开像的微分同胚。参数逆的光滑性同样来自 \((t,v)\mapsto(t,\Phi_t(v))\) 的局部逆，各局部逆由唯一性相容。
\end{proof}

上面的紧集论证还允许法丛随参数变化。具体地，设 \(K\) 紧且无边界，\(i_t:K\to N\) 为光滑嵌入族，\(E\to[0,1]\times K\) 是光滑向量丛。若一个零截面附近的光滑映射
\[
	\mathcal F(t,x,v)=(t,F(t,x,v))
\]
满足 \(F(t,x,0)=i_t(x)\)，且 \(\mathrm{d}\mathcal F\) 在零截面处可逆，则它在某个共同的小圆盘丛上仍是到相对开像的微分同胚。事实上，带时间的零截面像由紧空间 \([0,1]\times K\) 嵌入 \([0,1]\times N\)；若任意小半径仍有两个向量像重合，取趋零向量及基点的收敛子列，极限像相等便迫使时间和基点都相等，再由该点的局部逆排除重合。零截面紧同时给出定义域与可逆微分的统一半径。此论证只用丛的局部平凡化，不要求把所有 \(E_{t,x}\) 预先识别成一个固定丛。在参数端点，先局部延拓系数；由于第一坐标始终是 \(t\)，逆函数定理所得的局部逆也保持时间，因而可限制回半空间。这一版本将在同位延拓中保持每个时刻的子流形。

若 \(S\) 原来没有在 \(N\) 中闭，先取所有适配坐标域的并 \(N_0\)。它是包含 \(S\) 的开集，且 \(S\) 在 \(N_0\) 中闭：任何 \(z\in N_0\) 都位于某张适配图内，图中的坐标切片是相对闭集，因而一个不在切片上的点不可能是 \(S\) 的聚点。在 \(N_0\) 中应用引理即可。由此，管状邻域存在性不要求原嵌入固有，但所需半径可以趋于零。

\subsection{管状邻域及其选择无关性}
\label{g2:subsec:1503:04}

\begin{definition}\label{g2:def:tubular-neighborhood}
	设 \(S\subset N\) 为无边界嵌入子流形。一张\term[guanzhuanglinyu]{管状邻域映射}[tubular neighborhood map]是从 \(\nu_NS\) 的零截面某邻域 \(V\) 到 \(N\) 的开集 \(U\supset S\) 的微分同胚 \(\tau:V\to U\)，满足 \(\tau(0_x)=x\)，并且纤维方向的微分在取法商后为恒等：
	\[
		\nu_xS\longrightarrow\mathrm{T}_xN\longrightarrow\nu_xS,
		\qquad v\longmapsto\mathrm{d}\tau_{0_x}(0,v)\longmapsto v.
	\]
	称 \(U\) 为 \(S\) 的\term[guanzhuanglinyu]{管状邻域}[tubular neighborhood]。可缩小 \(V\)，使它为上一引理中的变半径圆盘丛。
\end{definition}

\begin{theorem}[管状邻域定理]\label{g2:thm:tubular-neighborhood}
	每个无边界光滑流形中的无边界嵌入子流形都有管状邻域。子流形紧时可取常数法向半径，一般情形可取正的变半径。
\end{theorem}
\begin{proof}
	度量与联络给出的法向指数 \(\Phi\) 满足\cref{g2:prop:normal-exponential-local}，故\cref{g2:lem:zero-section-shrinking}在适当环境开集中给出 \(\tau=\Phi|_{E_\varepsilon}\)。其零截面值正确，纤维微分是 \(\iota_g\)，取商后为恒等，因此满足定义。若 \(S\) 紧，连续正函数 \(\varepsilon\) 有正的下界，进一步取更小常数即可。
\end{proof}

例如 \(S^k\subset\mathbb R^{k+1}\) 的外法向线丛平凡，\(\tau(x,t)=(1+t)x\)、\(|t|<1/2\)，给出一条环形邻域。在 Möbius 带总空间中，零截面的法线丛仍是 Möbius 线丛；其邻域是扭转的带，而非圆周与区间的乘积。管状模型保留的是整个法丛的拼接。

\ProseParagraphBreak

在 \(U=\tau(E_\varepsilon)\) 上，令 \(r=\pi\circ\tau^{-1}\)，得到光滑淹没 \(r:U\to S\)，且 \(r|_S=\operatorname{id}\)。纤维缩放给出
\[
	R_s(\tau(v))=\tau((1-s)v),\qquad0\le s\le1,
\]
因此 \(U\) 光滑形变收缩到 \(S\)。变半径条件只依赖基点，故整条线段 \((1-s)v\) 都留在定义域内。这些结论常用来把邻域中的问题送回子流形。\cref{g2:fig:tubular-retraction}画出圆周的情形。

\begin{figure}[htbp]
	\centering
	\input{Figures/Chapter15-TubularNeighborhood}
	\caption[法向纤维与管状邻域收缩]{圆周附近的短法向纤维填成环形邻域。投影 \(r\) 将同一纤维上的点送回基点 \(x\)；纤维缩放则给出从恒等映射到 \(r\) 的形变。}
	\label{g2:fig:tubular-retraction}
\end{figure}


\begin{proposition}\label{g2:prop:tubular-uniqueness}
	设 \(\tau_0,\tau_1\) 是同一个无边界嵌入子流形 \(S\subset N\) 的两张管状邻域映射，均按上述定义正规化法方向微分。存在法丛零截面的一个共同邻域及联合光滑族 \(\tau_t\)、\(0\le t\le1\)，使每张 \(\tau_t\) 都是管状邻域映射，端点为原映射的限制，并且整族在零截面上不变。
\end{proposition}
\begin{proof}
	先缩小到 \(\tau_1(v)\) 落在 \(\tau_0\) 的像内，令 \(F=\tau_0^{-1}\circ\tau_1\)。它是法丛零截面附近的微分同胚，固定零截面，且诱导的法方向微分为恒等。记纤维缩放为 \(\delta_t(v)=tv\)。对 \(t>0\)，定义
	\[
		F_t=\delta_{1/t}\circ F\circ\delta_t.
	\]
	在局部丛坐标中写 \(F(x,v)=(a(x,v),b(x,v))\)。零截面与正规化条件给出
	\[
		a(x,0)=x,\qquad b(x,0)=0,\qquad \partial_vb(x,0)=I.
	\]
	于是
	\[
		F_t(x,v)=\left(a(x,tv),\frac{b(x,tv)}t\right),\qquad
		\frac{b(x,tv)}t=\int_0^1\partial_vb(x,stv)v\,\mathrm{d}s.
	\]
	积分表达说明这族光滑延到 \(t=0\)，且 \(F_0(x,v)=(x,v)\)。这些表达来自整体纤维缩放，故不同平凡化中的延拓相容。零截面处的微分形如
	\[
		\mathrm{d}(F_t)_{0_x}(u,v)=(u+tB_xv,v),
	\]
	始终可逆。

	令 \(\tau_t=\tau_0\circ F_t\)。紧参数区间保证它们在零截面的某共同邻域都有定义；对这一族使用\cref{g2:lem:zero-section-shrinking}，取得共同的变半径邻域，使各张都为到开像的微分同胚。它们固定零截面，法方向微分都是恒等，且 \(t=0,1\) 时分别为 \(\tau_0,\tau_1\)。
\end{proof}

这就是这里的选择无关性：缩小邻域后，度量、联络和半径所给出的不同管状坐标可以通过固定子流形的光滑同位相连。它并未指定唯一的一张坐标图。若一个邻域参数化在法商上诱导的是丛自同构 \(A\)，先把纤维变量换成 \(A^{-1}v\)，便得到定义中的正规化；把这项识别写清楚，才能比较相同法丛上的两张映射。Euclid 空间中的法向加法与变半径构造还可参看\cite[第138--141页，定理6.24及命题6.25]{Lee2012Smooth}。

% !TeX root = ../../Book2.tex
% 纲要标题：### 15.4 边界的领状邻域
% 纲要位置：几何学计划纲要.md，第 1597 行。
% 正文状态：初稿；尚未正式审核。

% 纲要内容开始
% 
% * 向内横截向量场与局部流
% * 边界邻域的整体拼接
% * 双化和相对粘合中的应用
% 
% 纲要内容结束

\section{边界的领状邻域}
\label{g2:sec:1504}

在半空间图中，边界附近看起来像“边界坐标加一个非负坐标”。各张图的非负坐标并不相同，因而还不能把所有边界点统一向内推一段距离。本节用一个向内向量场把这些局部图连接起来：每条轨道从边界出发，时间成为共同的横向坐标。由此得到的领状邻域可用于双化、粘合和保持边界数据的构造。

\subsection{向内横截向量场与局部流}
\label{g2:subsec:1504:01}

设 \(M\) 为有限维、Hausdorff、第二可数的光滑带边界流形。在边界坐标 \((y,r)\)、\(r\ge0\) 中，向量 \(v\in\mathrm{T}_xM\) 严格指向内部的条件是 \(\mathrm{d}r_x(v)>0\)。第一卷第11章的边界坐标变换说明，另一边界定义坐标沿边界的微分是 \(\mathrm{d}r\) 的正倍数，所以这一条件不依赖所选图。它也说明同一点的向内向量组成凸集，正的加权平均仍向内。

\begin{lemma}\label{g2:lem:boundary-defining-function}
	每个光滑带边界流形 \(M\) 都有光滑函数 \(r:M\to[0,\infty)\)，满足 \(r^{-1}(0)=\partial M\) 且 \(\mathrm{d}r\ne0\) 于 \(\partial M\)。在边界的某个开邻域上，可选择光滑向量场 \(X\)，使 \(Xr=1\)。
\end{lemma}
\begin{proof}
	取由边界坐标图和内点坐标图组成的局部有限覆盖，并取从属光滑单位分解 \(\chi_i\)。边界图上记最后一个坐标为 \(r_i\)，内点图上令 \(r_i=1\)，定义 \(r=\sum_i\chi_i r_i\)，各项在支集外补零。每个内点的所有非零项都为正，故 \(r>0\)；在边界点，来自内点图的项在该点附近消失，其余 \(r_i\) 为零，所以 \(r=0\)。边界上
	\[
		\mathrm{d}r=\sum_i\chi_i\,\mathrm{d}r_i.
	\]
	各个 \(\mathrm{d}r_i\) 是同一个向内法向余向量的正倍数，且权重之和为1，故和非零。

	在各边界图中取 \(\partial_{r_i}\)，用单位分解加权拼成边界附近的光滑向量场 \(Y\)。其边界值严格向内，因此 \(Yr>0\) 在边界的某开邻域 \(U\) 上成立。令 \(X=Y/(Yr)\) 于 \(U\)，便有 \(Xr=1\)。这里使用的是第一卷第12章第3节已经证明的局部有限单位分解；局部有限性保证求和及逐项求导有意义。
\end{proof}

\(r\) 称为一个\term[bianjiedingyihanshu]{边界定义函数}[boundary defining function]。条件 \(\mathrm{d}r\ne0\) 保证它真正测量横向变化；仅有同一零点集的函数 \(r^2\) 不具有这个性质。归一化 \(Xr=1\) 的作用也很具体：沿从边界出发的积分曲线，\(r\) 的值恰好就是已经经过的时间。

\ProseParagraphBreak

在边界处不能直接要求 \(X\) 生成保持 \(M\) 的双侧流。向内的速度在负时间会走到半空间外。我们在边界坐标中把光滑分量局部延拓到 Euclid 开集，应用第一卷第13章第1节的初值存在、唯一和光滑依赖定理，再保留非负时间。因 \(r(0)=0\)、\(r'=1\)，所得小段正时间曲线位于内部。不同局部延拓在半空间内给出同一个方程，解的唯一性使这些正时间解一致。记其为
\[
	F(x,t)=\Phi_t(x),\qquad x\in\partial M,\quad 0\le t<\varepsilon_x.
\]
在 \(t=0\) 处，\(\mathrm{d}F\) 把边界切向量保持为原向量，并把 \(\partial_t\) 送到 \(X_x\)。这个向量不切于边界，因此上述导数为同构，局部逆函数论证给出边界附近的半空间坐标。

\subsection{边界邻域的整体拼接}
\label{g2:subsec:1504:02}

\begin{definition}\label{g2:def:collar}
	设 \(M\) 为光滑带边界流形。若光滑嵌入 \(c:\partial M\times[0,1)\to M\) 满足 \(c(x,0)=x\)，且像是 \(\partial M\) 在 \(M\) 中的开邻域，就称它为一个\term[lingzhuanglinyu]{领状邻域}[collar neighborhood]的参数化，简称领圈。
\end{definition}

\begin{theorem}[领状邻域定理]\label{g2:thm:collar-neighborhood}
	每个有限维、Hausdorff、第二可数的光滑带边界流形都存在领圈。边界可以非紧；领圈的像可取在任何预先给定的边界开邻域内。
\end{theorem}
\begin{proof}
	在前面构造的开集 \(U\) 上使用 \(Xr=1\) 的向量场，再把 \(U\) 与指定的边界邻域相交。由局部存在与光滑依赖，每个边界点有一个边界邻域 \(V_i\) 及正数 \(a_i\)，使从 \(V_i\) 出发的解在 \(0\le t<a_i\) 上存在并留在 \(U\)。可取 \(\{V_i\}\) 局部有限。设 \(\{\psi_i\}\) 为从属的单位分解，令
	\[
		\varepsilon(x)=\frac12\sum_i\psi_i(x)a_i>0.
	\]
	对每个 \(x\)，在 \(\psi_i(x)>0\) 的有限多个指标中选 \(a_i\) 最大者，则 \(\varepsilon(x)<a_i\) 且 \(x\in V_i\)。故 \(F\) 在
	\[
		D=\{(x,t):x\in\partial M,\ 0\le t<\varepsilon(x)\}
	\]
	上有定义，光滑且满足 \(r(F(x,t))=t\)。

	若 \(F(x,t)=F(y,s)\)，对两边取 \(r\) 得 \(t=s\)。两条轨道从相同时间、相同位置反向求解，在返回边界前都位于 \(M\)，唯一性给出 \(x=y\)。所以 \(F\) 在整个 \(D\) 上单射，这一步同时排除了来自遥远边界点的相撞。

	在正时间，流的微分把 \(\mathrm{T}_x\partial M\) 与 \(X_x\) 组成的基送到终点，并保持生成方向 \(X\)，故 \(\mathrm{d}F\) 仍可逆。也可注意 \(r\circ F=t\)，再在每张 \(t\) 切片上用流的局部逆。零时间处已由半空间局部逆得到同样结论。于是 \(F\) 是单射局部微分同胚，其局部逆在交叠处一致，拼成到开像的光滑整体逆。

	最后令 \(c(x,u)=F(x,u\varepsilon(x))\)，\(0\le u<1\)。\((x,u)\mapsto(x,u\varepsilon(x))\) 是到 \(D\) 的微分同胚，因而 \(c\) 是所需领圈。
\end{proof}

非紧边界上，轨道允许走的时间随点变化；最后的重参数化把这个变化吸收到 \(\varepsilon(x)\) 中。统一参数区间 \([0,1)\) 因而不表示每个点都能沿最初向量场运动同样长的时间。相关流与领圈构造可参看\cite[第217--223页，定理9.20、9.24、9.25]{Lee2012Smooth}。

\begin{example}\label{g2:exm:disk-collar}
	闭单位球 \(D^n\) 的一个领圈为 \(c(x,t)=(1-t/2)x\)，其中 \(x\in S^{n-1}\)、\(0\le t<1\)。其像是 \(\{z:1/2<|z|\le1\}\)，在 \(D^n\) 中开。边界点被逐渐推入内部，角方向保持不变；半径 \(1/2\) 没有包含在像中，对应参数端点 \(t=1\) 未取到。
\end{example}

\subsection{双化和相对粘合中的应用}
\label{g2:subsec:1504:03}

取 \(M\) 的两份副本 \(M_+,M_-\)，将对应的边界点识别，得到\term[shuanghua]{双化}[double]
\[
	\operatorname{Dbl}(M)=(M_+\amalg M_-)/\{x_+\sim x_-:x\in\partial M\}.
\]
给定领圈 \(c\) 后，用一个有符号的参数把接缝附近写成
\[
	\Theta:\partial M\times(-1,1)\longrightarrow\operatorname{Dbl}(M),\qquad
	\Theta(x,s)=\begin{cases}
		[c(x,s)]_+,&s\ge0,\\
		[c(x,-s)]_-,&s\le0.
	\end{cases}
\]
两条表达在零处一致。它是到接缝开邻域的同胚：在两半上分别有领圈的逆，它们的第一分量在边界一致，第二分量分别取正号、负号，因而由粘合定义得到连续逆。

\begin{proposition}\label{g2:prop:smooth-boundary-gluing}
	设 \(M_1,M_2\) 为同维光滑带边界流形，\(\varphi:\partial M_1\to\partial M_2\) 为微分同胚。选定两侧领圈后，商空间 \(M_1\cup_\varphi M_2\) 具有光滑无边界流形结构，使两份 \(M_i\) 都以闭的带边界子流形嵌入其中。特别地，\(\operatorname{Dbl}(M)\) 是无边界光滑流形。只沿边界分支的并作同样粘合时，其余边界保留。
\end{proposition}
\begin{proof}
	以 \(\partial M_1\) 为共同边界，在负半边用 \(c_2(\varphi(x),-s)\) 代替上面的表达。两侧内点保留原图；接缝附近取 \(\partial M_1\) 的坐标与 \(s\) 构成乘积图。两张接缝图的变换为 \((y,s)\mapsto(h(y),s)\)，所以光滑。接缝图与内点图的交落在 \(s>0\) 或 \(s<0\)，过渡由领圈、原流形坐标及 \(\varphi\) 的光滑性保证。

	商空间还须满足分离性与可数性。每个接缝点有上述乘积邻域，两侧内点也有原来的开邻域。两个不同接缝点用边界上的不交开集及较小领圈分离；内点与接缝点用避开该内点的小领圈分离；内点之间使用原流形的 Hausdorff 性或两侧互不相交的内部。这证明 Hausdorff 性。内部图及接缝的乘积图均可取可数族，故整体第二可数。各侧在这些坐标中为闭半空间，其补集是另一侧内部，所以是所述闭嵌入。
\end{proof}

粘合光滑对象时，领圈还告诉我们应匹配哪些数据。若两边的函数在领圈内都等于同一个边界函数 \(a(x)\)，它们就粘成光滑函数，因为接缝表达为与 \(s\) 无关的 \(a(x)\)。只要求边界值相等则不够：两侧都取领圈坐标 \(u\)，接起来变成 \(|s|\)，在零处不可微。对一般函数，需要各阶法向导数按 \(u=-s\) 的符号相容；先在领圈内调整为乘积形式，是一种直接保证相容的做法。\cite[第224--225页，定理9.29]{Lee2012Smooth}给出了相同的流形粘合构造。

% !TeX root = ../../Book2.tex
% 纲要标题：### 15.5 同位延拓
% 纲要位置：几何学计划纲要.md，第 1603 行。
% 正文状态：初稿；尚未正式审核。

% 纲要内容开始
% 
% * 同位的速度场
% * 管状邻域延拓与截断
% * 环境流及支集控制
% * 相对同位和非紧限制
% 
% 纲要内容结束

\section{同位延拓}
\label{g2:sec:1505}

把一个圆在平面中平移，可以同时平移整个平面；若只想移动圆附近的区域，还可以把平移速度截断。一般的嵌入族也提出同样的问题：已经规定子流形上的每个点怎样运动，能否让整个环境流形一起运动，并使远处保持原样？困难在于不同时间的速度给在不同位置上。先把这些速度写成随时间变化的向量场，再延拓到环境中，流的唯一性就能保证它准确实现原来的运动。

\subsection{同位的速度场}
\label{g2:subsec:1505:01}

\begin{definition}\label{g2:def:smooth-isotopy}
	设 \(K,N\) 为光滑流形。光滑映射 \(i:[0,1]\times K\to N\) 若对每个 \(t\) 都使 \(i_t=i(t,-)\) 成为嵌入，就称为一个光滑\term[tongwei]{同位}[isotopy]。若光滑映射 \(\Phi:[0,1]\times N\to N\) 的每个 \(\Phi_t\) 都是微分同胚，且 \(\Phi_0=\operatorname{id}_N\)，则称其为\term[huanjingtongwei]{环境同位}[ambient isotopy]。
\end{definition}

同位指定的是一族带参数的嵌入，而不只是一族集合。即使 \(i_t(K)\) 始终是同一个子集，其上的点也可能沿子流形移动。例如 \(i_t(e^{\mathrm i\theta})=e^{\mathrm i(\theta+t)}\) 保持圆周的像不变，却让每个标记点转动。环境同位若延拓这个族，应满足 \(\Phi_t\circ i_0=i_t\)，而非仅把两个像集合相互送到。

\ProseParagraphBreak

对光滑同位定义
\[
	v(t,i_t(x))=\frac{\partial}{\partial t}i(t,x)\in\mathrm{T}_{i_t(x)}N.
\]
因为 \(i_t\) 单射，右端唯一确定左端；这个向量属于环境切空间，未必切于 \(i_t(K)\)。为理解其光滑性，考虑带时间标记的映射
\[
	\widehat i:[0,1]\times K\longrightarrow[0,1]\times N,
	\qquad(t,x)\longmapsto(t,i_t(x)).
\]
其微分为 \((a,u)\mapsto(a,a\partial_ti+\mathrm{d}i_t(u))\)，所以单射。若 \(K\) 紧，\(\widehat i\) 又是从紧空间到 Hausdorff 空间的连续单射，因而是嵌入。于是在其像上，\(v\) 是向量丛 \(\operatorname{pr}_N^*\mathrm{T}N\) 的光滑截面。这正是需要延拓的数据：时间分量保持为零，空间分量等于指定速度。

\subsection{管状邻域延拓与截断}
\label{g2:subsec:1505:02}

先取紧无边界的 \(K\)。在 \(N\) 上选度量，令 \(E_{t,x}\) 为 \(\mathrm{d}i_t(\mathrm{T}_xK)\) 的正交补。这些空间在 \([0,1]\times K\) 上组成光滑向量丛：在局部标架中，正交投影由光滑的 Gram 矩阵及其逆给出。用\secref{g2:sec:1503}的法向指数构造
\[
	\mathcal E(t,x,u)=\bigl(t,\exp_{i_t(x)}u\bigr),\qquad u\in E_{t,x}.
\]
在零截面上它的微分可逆，且第一分量仍为 \(t\)。\cref{g2:lem:zero-section-shrinking}的紧参数版本给出零截面邻域，使 \(\mathcal E\) 成为到 \([0,1]\times N\) 中轨迹邻域的微分同胚。其逆的基点分量定义保持时间的光滑收缩
\[
	R(t,y)=(t,i_t(x)).
\]
这里的法向坐标只用于把空间点拉回同一时刻的子流形。

\ProseParagraphBreak

有了 \(R\)，向量场仍不能直接写成 \(v\circ R\)：这个向量位于 \(\mathrm{T}_{i_t(x)}N\)，而需要的值应位于 \(\mathrm{T}_yN\)。可以用局部标架明确解决这个问题。取覆盖轨迹的有限个向量丛平凡化，记环境切丛的局部标架为 \(e_{\alpha j}\)，并在轨迹上取从属单位分解 \(\lambda_\alpha\)。将
\[
	v=\sum_j v_\alpha^j e_{\alpha j}
\]
中的系数 \(v_\alpha^j\) 及权重 \(\lambda_\alpha\) 与 \(R\) 复合，再在 \((t,y)\) 处使用标架 \(e_{\alpha j}(y)\)。缩小邻域，使涉及的基点和 \(y\) 都落在该平凡化内，就得到
\[
	\widetilde X(t,y)=\sum_{\alpha,j}\lambda_\alpha(R(t,y))\,
	v_\alpha^j(R(t,y))\,e_{\alpha j}(y).
\]
在原轨迹上 \(R\) 是恒等，权重之和为1，故 \(\widetilde X=v\)。各项按其支集补零，得到光滑延拓。

\ProseParagraphBreak

设 \(\Omega\subset N\) 是包含整个运动轨迹 \(i([0,1]\times K)\) 的开集。轨迹紧，因此可在上述时空邻域中选一个光滑截断 \(\chi\)，于轨迹附近等于1，支集含于 \([0,1]\times\Omega\)，且支集的空间投影是紧集。第一卷第12章的光滑截断定理适用于这个带边界的时空区域。令 \(X_t(y)=\chi(t,y)\widetilde X(t,y)\)，在其余地方补零，则
\[
	X_t(i_t(x))=\partial_ti(t,x),\qquad
	\bigcup_{0\le t\le1}\operatorname{supp}X_t\subset L\Subset\Omega.
\]
一个共同的紧集 \(L\) 控制所有时刻的支集，是下一步得到整个时间区间上环境流的关键。

\subsection{环境流及支集控制}
\label{g2:subsec:1505:03}

\begin{lemma}\label{g2:lem:uniformly-supported-evolution}
	设 \(N\) 为无边界光滑流形，\(X_t\) 是 \(0\le t\le1\) 上光滑的时间依赖向量场。若存在紧集 \(L\subset N\)，使 \(X_t=0\) 于 \(N\setminus L\)，对所有 \(t\) 同时成立，则其演化映射 \(\Phi_{t,s}:N\to N\) 对全部 \(s,t\in[0,1]\) 有定义，光滑且为微分同胚，逆为 \(\Phi_{s,t}\)。它在 \(N\setminus L\) 上为恒等。
\end{lemma}
\begin{proof}
	第一卷第13章第1节“演化映射”给出局部存在、联合光滑性与逆映射公式。这里只需保证解能走到指定的时间。\(N\setminus L\) 中的点对全部时间都是零点，所以从这些点出发的解恒定。若一条从 \(L\) 出发的解进入 \(N\setminus L\)，从进入处反向使用唯一性，它在此前也应是同一条常值解，矛盾。因此这条解始终留在 \(L\)。

	若最大存在区间有一个位于 \([0,1]\) 内的有限端点，时空轨道就留在紧集 \([0,1]\times L\) 中。把 \(X_t\) 写成增广向量场 \(\partial_t+X_t\)，对这个紧集用有限个局部初值邻域取得统一的正存在时间，便能从靠近端点的时刻重新启动并越过它，与最大性矛盾。这正是第一卷第13章第2节逃逸判据的证明；端点0、1处先将光滑系数作局部时间延拓即可。故演化在整个指定时间区间存在，原有的逆映射与光滑依赖结论同时成立。
\end{proof}

\begin{theorem}[同位延拓定理]\label{g2:thm:isotopy-extension}
	设 \(K\) 为紧无边界光滑流形，\(N\) 为无边界光滑流形，\(i:[0,1]\times K\to N\) 为光滑同位。对任意包含 \(i([0,1]\times K)\) 的开集 \(\Omega\subset N\)，存在环境同位 \(\Phi_t\)，使
	\[
		\Phi_t\circ i_0=i_t,
	\]
	且所有 \(\Phi_t\) 在同一个紧集 \(L\Subset\Omega\) 外均为恒等。
\end{theorem}
\begin{proof}
	用上一小节的管状延拓与截断构造 \(X_t\)，再用\cref{g2:lem:uniformly-supported-evolution}取 \(\Phi_t=\Phi_{t,0}\)。对任意 \(x\in K\)，曲线 \(t\mapsto i_t(x)\) 满足
	\[
		\frac{\mathrm{d}}{\mathrm{d}t}i_t(x)=X_t(i_t(x)),\qquad i_{t=0}(x)=i_0(x).
	\]
	它与 \(t\mapsto\Phi_t(i_0(x))\) 是同一个初值问题的解，故由唯一性相等。共同紧支集给出所述支集控制。
\end{proof}

\begin{example}\label{g2:exm:translated-sphere-isotopy}
	设 \(a\in\mathbb R^{n+1}\)，\(i_t(x)=x+ta\)，其中 \(x\in S^n\)。整个运动轨迹紧，沿它的速度恒为 \(a\)。选择在轨迹邻域为1的紧支函数 \(\chi\)，以 \(X(y)=\chi(y)a\) 为环境向量场。每条指定轨道始终在 \(\chi=1\) 的区域，故环境流在球面上恰好实现 \(x\mapsto x+ta\)，在远处保持恒等。环境的其他点如何移动由截断决定，因此延拓存在而不唯一。
\end{example}

同位由嵌入组成这个要求不能放松成一般同伦。例如圆周上的常值映射与某些度为零的映射可以连续变形，但常值映射无法由环境微分同胚作用于一个嵌入得到；环境运动始终保留嵌入的单射性与局部秩。

\subsection{相对同位和非紧限制}
\label{g2:subsec:1505:04}

先看需要固定一个环境闭集 \(B\subset N\) 的情形。若存在开集 \(V\supset B\)，使轨迹进入 \(V\) 时的指定速度全部为零，即
\[
	\partial_ti(t,x)=0\qquad\text{当 }i_t(x)\in V,
\]
那么可让延拓后的环境同位在 \(B\) 的某邻域恒为恒等。事实上，速度非零的轨迹点的闭包是轨迹中的紧集，且避开 \(V\)；取一个在该紧集附近为1、支集避开 \(B\) 的截断函数，将前面的延拓场再乘以它。它在速度非零处保持原值，在其余轨迹点仍为零，所以仍延拓原同位，而支集与 \(B\) 分离。若仅要求同位在 \(K\) 的某个闭子集 \(A\) 上逐点不变，则原构造已经保证 \(X_t(i_0(x))=0\) 对 \(x\in A\) 成立，环境流因而逐点固定 \(i_0(A)\)；这并不自动固定其整个邻域。

\ProseParagraphBreak

紧源还可以换成受控的非紧源。一个足够条件是：\(i_0:K\to N\) 为固有嵌入，并存在固定紧集 \(C\subset K\)，使 \(i_t(x)=i_0(x)\) 对所有 \(x\notin C\) 和所有 \(t\) 成立。在这种情况下，\(\widehat i\) 是固有嵌入：紧目标集的原像落在 \([0,1]\times C\) 与 \([0,1]\times i_0^{-1}(L)\) 的并中，其中 \(L\) 为该目标集的空间投影。速度截面在轨迹上的支集紧，可在其附近按同样的局部标架构造延拓，并在其余轨迹上保留零值，再截断为有共同紧支集的场。流的论证不变，得到环境同位。

\begin{example}\label{g2:exm:noncompact-isotopy-no-extension}
	令 \(K=N=\mathbb R\)，定义
	\[
		i_t(x)=\begin{cases}
			\dfrac{\arctan(tx)}{t},&t>0,\\
			x,&t=0.
		\end{cases}
	\]
	积分表达 \(i_t(x)=\int_0^x(1+t^2u^2)^{-1}\,\mathrm{d}u\) 说明它在 \((t,x)\) 上光滑，且 \(\partial_x i_t>0\)，故每个 \(i_t\) 都是嵌入。当 \(t>0\) 时，其像是有界开区间 \((-\pi/(2t),\pi/(2t))\)；而 \(i_0\) 是整个实线的恒等映射。如果存在环境同位延拓它，\(\Phi_t\circ i_0=i_t\) 就迫使 \(\Phi_t=i_t\)，与 \(\Phi_t\) 必须满射到 \(\mathbb R\) 矛盾。非紧源上，各个切片是嵌入本身并不足够。
\end{example}

源 \(K\) 若紧但有边界，可以把速度截面的延拓改在半空间坐标中进行：先将各局部光滑分量延过边界，再用时空单位分解拼接，保留沿轨迹的原值。局部嵌入的半空间标准形由其延拓的浸入标准形给出，因此这一步适用于带边界的轨迹；截断和流的证明随后相同。环境 \(N\) 若有边界，所需延拓场必须沿 \(\partial N\) 切向，才能产生保持 \(N\) 的双侧运动。一个直接适用的相对版本是原同位在环境边界的某邻域内恒定，并把延拓场的支集选在内部；第一卷第13章第2节的边界切向流结论也说明了这一条件的作用。

% !TeX root = ../../Book2.tex
% 纲要标题：### 15.6 固有淹没的局部平凡性
% 纲要位置：几何学计划纲要.md，第 1610 行。
% 正文状态：初稿；尚未正式审核。

% 纲要内容开始
% 
% * 水平分布和局部提升方程
% * 固有性保证的统一存在区间
% * 径向提升与局部平凡化
% * Ehresmann 定理
% * Ehresmann 定理的带边界版本
% * 光滑与全纯局部平凡性
% 
% ---
% 
% 纲要内容结束

\section{固有淹没的局部平凡性}
\label{g2:sec:1506}

淹没定理说明，满秩映射在每个源点附近都像一个坐标投影。然而，一条纤维可能穿过许多张坐标图：把这些局部投影拼在一起，能否得到整个纤维附近的乘积？本节用沿基底路径移动纤维的方法回答这个问题。固有性保证移动过程中没有点在有限时间内逃向无穷远，因而同一套构造能作用于整条纤维。

\subsection{水平分布和局部提升方程}
\label{g2:subsec:1506:01}

先设 \(p:E\to B\) 是无边界光滑流形之间的淹没。记 \(\dim E=m\)、\(\dim B=b\)。每一点的核
\[
	V_e=\ker\mathrm{d}p_e\subseteq\mathrm{T}_eE
\]
有固定维数 \(m-b\)，由满秩矩阵的局部消元可知，它们组成 \(\mathrm{T}E\) 的光滑子丛 \(V\)，称为\term{垂直切丛}[vertical tangent bundle]。若 \(e\in p^{-1}(y)\)，则 \(V_e=\mathrm{T}_e(p^{-1}(y))\)：垂直方向就是留在当前纤维内的运动方向。

在 \(E\) 上选一个辅助 Riemann 度量，令 \(H=V^\perp\)。第一卷第12章已经证明，光滑向量丛可取光滑度量，光滑子丛的正交补仍是光滑子丛。因此
\[
	\mathrm{T}E=V\oplus H,
	\qquad \mathrm{d}p_e|_{H_e}:H_e\longrightarrow\mathrm{T}_{p(e)}B
\]
是线性同构。它的逆随 \(e\) 光滑变化：在局部标架中，这就是一个可逆光滑矩阵的逆。记这个逆为
\begin{equation}\label{g2:eq:horizontal-right-inverse}
	L_e:\mathrm{T}_{p(e)}B\longrightarrow\mathrm{T}_eE,
	\qquad \mathrm{d}p_e\circ L_e=\operatorname{id}.
\end{equation}
像丛 \(H\) 称为一个\term{水平分布}[horizontal distribution]。给定基底的一个速度，\(L_e\) 选出总空间中与它对应的速度。这里不要求水平分布满足 Frobenius 可积条件；提升一条曲线只需求解常微分方程，不需要存在处处切于 \(H\) 的高维子流形。

\ProseParagraphBreak

设 \(\psi:U_0\to\mathbb R^b\) 是基底坐标图，\(u\in\mathbb R^b\)。在 \(p^{-1}(U_0)\) 上定义依赖参数 \(u\) 的向量场
\begin{equation}\label{g2:eq:radial-horizontal-field}
	X_u(e)=L_e\bigl(\mathrm{d}(\psi^{-1})_{\psi(p(e))}u\bigr).
\end{equation}
若 \(e(t)\) 是它的积分曲线，链式法则给出
\[
	\frac{\mathrm{d}}{\mathrm{d}t}\psi(p(e(t)))
	=\mathrm{d}\psi_{p(e(t))}\mathrm{d}p_{e(t)}X_u(e(t))=u.
\]
所以基底坐标沿直线运动：
\begin{equation}\label{g2:eq:horizontal-base-line}
	\psi(p(e(t)))=\psi(p(e(0)))+tu.
\end{equation}
这条等式同时说明提升方程的含义与求解方向：我们先在基底指定直线，再用 \(L\) 确定总空间中的运动。

对一条光滑路径 \(c(t)\subset U_0\)，将上述常向量 \(u\) 换成 \(\frac{\mathrm{d}}{\mathrm{d}t}(\psi\circ c)(t)\)，便得到含时向量场。如果初始点位于 \(p^{-1}(c(0))\)，同一个计算说明解始终投影到 \(c(t)\)。第一卷第13章第1节的局部存在唯一性和初值、参数光滑依赖定理适用：系数关于 \((t,e)\) 光滑，而含时方程也可通过增加时间坐标写成自治方程。下面需要补上的，是这些局部解能存在多久。

\subsection{固有性保证的统一存在区间}
\label{g2:subsec:1506:02}

回忆，\(p\) 固有是指每个紧集 \(K\subset B\) 的原像 \(p^{-1}(K)\) 紧。因而，只要基底路径始终留在同一个紧集内，它的任何水平提升也都被限制在一个紧集内。这与第一卷第13章第2节的逃逸判据正好配合：光滑向量场的最大积分曲线若在有限时刻不能继续，就必定离开定义域中的每个紧集。

\begin{lemma}\label{g2:lem:proper-horizontal-lift}
	设 \(p:E\to B\) 是固有淹没，\(c:[0,1]\to U_0\) 是光滑路径，且 \(c([0,1])\) 包含在 \(U_0\) 的某个紧子集中。对每个 \(e_0\in p^{-1}(c(0))\)，由选定水平分布确定的提升都在整个 \([0,1]\) 上存在且唯一。
\end{lemma}
\begin{proof}
	先由局部存在唯一性得到经过 \(e_0\) 的提升。设它的向前最大存在时刻为 \(T\leq1\)。在 \(t<T\) 时，它的投影等于 \(c(t)\)，所以提升位于紧集 \(p^{-1}(c([0,1]))\subset p^{-1}(U_0)\) 内。把时间作为附加坐标后，积分曲线位于紧集 \([0,1]\times p^{-1}(c([0,1]))\) 内。路径及向量场可在端点附近光滑延拓；有限个局部存在邻域给出这个紧集附近统一的短时存在区间。因此解可以越过 \(T\) 继续，与 \(T\) 是不可延拓的端点矛盾。若只需解到 \(T=1\)，同一论证保证端点值存在。唯一性来自局部唯一性。
\end{proof}

这里的“统一”不要求预先给整条非紧纤维估计一个全局速度界。固有性已经使初始纤维和路径上方的全部区域紧，从而可用有限个局部方程控制所有初始点。对于依赖光滑参数的路径，还可把参数加入自治方程，令其导数为零；当参数限制在紧集内时，同样的有限覆盖论证给出共同的存在区间。随后，初值与参数光滑依赖定理便保证提升同时光滑依赖于起点和参数。关于这两个常微分方程工具，可对照 \cite[第9章，第215--217页]{Lee2012Smooth}。

\ProseParagraphBreak

固有性在这里有实际作用。例如
\[
	p:\mathbb R^2\setminus\{(0,0)\}\longrightarrow\mathbb R,
	\qquad p(x,y)=x
\]
处处为淹没，但 \(p^{-1}(0)\cong\mathbb R\setminus\{0\}\) 有两个连通分支，\(p^{-1}(x)\cong\mathbb R\) 对 \(x\ne0\) 只有一个连通分支。因此 \(0\) 附近不存在以整条纤维为因子的局部乘积。满秩条件没有阻止纤维的一部分从缺失点附近消失；这个映射也确实不固有。

\subsection{径向提升与局部平凡化}
\label{g2:subsec:1506:03}

固定 \(y_0\in B\)，取坐标图 \(\psi:U_0\to\mathbb R^b\)，使 \(\psi(y_0)=0\)。缩小后可取 \(r>0\)，使闭球 \(\overline{B(0,2r)}\) 包含在 \(\psi(U_0)\) 中。令
\[
	U=\psi^{-1}(B(0,r)),\qquad F=p^{-1}(y_0).
\]
对 \(y\in U\)，选择从 \(y_0\) 到 \(y\) 的坐标径向路径
\[
	c_y(t)=\psi^{-1}(t\psi(y)),\qquad 0\leq t\leq1.
\]
每一条路径都位于紧集 \(K=\psi^{-1}(\overline{B(0,r)})\) 内，而且随 \((t,y)\) 光滑变化。

记 \(\varphi_u^t(e)\) 为式 \eqref{g2:eq:radial-horizontal-field} 中向量场 \(X_u\) 的局部流。由 \cref{g2:lem:proper-horizontal-lift}，对任意 \((z,y)\in F\times U\)，表达式
\begin{equation}\label{g2:eq:ehresmann-trivialization}
	\Phi(z,y)=\varphi_{\psi(y)}^1(z)
\end{equation}
都有定义。式 \eqref{g2:eq:horizontal-base-line} 说明
\[
	p(\Phi(z,y))=y.
\]
因此 \(\Phi\) 已经具有局部平凡化所需的投影关系。它把纤维 \(F\) 中每一点沿同一条基底路径搬运到 \(p^{-1}(y)\)，而不是分别在源空间中挑选互不相干的坐标图。

为了从另一条纤维回到 \(F\)，对 \(e\in p^{-1}(U)\) 固定 \(u=\psi(p(e))\)，沿 \(X_u\) 反向流动一个单位时间。它在基底中的坐标是 \((1+t)u\)，\(-1\leq t\leq0\)，仍留在 \(K\) 内，故同样存在。于是候选逆映射为
\begin{equation}\label{g2:eq:ehresmann-inverse}
	\Psi(e)=\bigl(\varphi_{\psi(p(e))}^{-1}(e),p(e)\bigr).
\end{equation}
这里上标 \(-1\) 表示时间 \(-1\)，而参数 \(\psi(p(e))\) 在这次反向流动中保持不变。若在途中不断按当前位置重新设定参数，方程就变了，也不再是原路径的逆向提升。

\subsection{Ehresmann 定理}
\label{g2:subsec:1506:04}

\begin{theorem}[Ehresmann 定理]\label{g2:thm:ehresmann}
	设 \(E,B\) 是无边界光滑流形，\(p:E\to B\) 是固有的满淹没。则每个 \(y_0\in B\) 都有邻域 \(U\)，使得存在微分同胚
	\[
		\Phi:F\times U\longrightarrow p^{-1}(U),\qquad F=p^{-1}(y_0),
	\]
	满足 \(p\circ\Phi=\operatorname{pr}_U\)。特别地，\(p\) 是以紧流形为纤维的光滑纤维丛。
\end{theorem}
\begin{proof}
	淹没定理给出 \(F\) 的光滑子流形结构，固有性给出其紧性。按照前两小节选取水平分布和坐标径向路径，并定义式 \eqref{g2:eq:ehresmann-trivialization} 与 \eqref{g2:eq:ehresmann-inverse}。

	先验证光滑性。向量场 \(X_u(e)\) 关于 \((e,u)\) 光滑，所以流映射 \((t,e,u)\mapsto\varphi_u^t(e)\) 在其存在域内光滑。它的存在域是开集，这是初值与参数光滑依赖定理的一部分。上一小节已证明，定义 \(\Phi\) 所需的每个 \((1,z,\psi(y))\) 和定义 \(\Psi\) 所需的每个 \((-1,e,\psi(p(e)))\) 都在这个存在域内。因此限制到这些初值并复合光滑映射后，\(\Phi\) 与 \(\Psi\) 都光滑。此处分别使用了两个事实：固有性保证解确实走到指定时刻，参数依赖定理保证到达点光滑变化。

	再验证互逆。若 \(u=\psi(y)\)，则 \(p(\varphi_u^1(z))=y\)，故逆向运动使用的仍是同一个向量场 \(X_u\)。流的唯一性给出
	\[
		\varphi_u^{-1}(\varphi_u^1(z))=z,
		\qquad \varphi_u^1(\varphi_u^{-1}(e))=e.
	\]
	于是 \(\Psi\Phi=\operatorname{id}_{F\times U}\)、\(\Phi\Psi=\operatorname{id}_{p^{-1}(U)}\)。投影关系已经验证，定理得证。
\end{proof}

若基底连通，则它道路连通。沿连接两个基点的路径取有限个局部平凡化，逐次比较相邻纤维，就知所有纤维彼此微分同胚。这个结论仍允许全局的粘合发生扭转。例如 \(S^1\to S^1\)、\(z\mapsto z^2\) 是固有满淹没，纤维为两个点，局部有上述乘积；总空间 \(S^1\) 连通，而 \(\{1,2\}\times S^1\) 不连通，因此全局并不是这个乘积。

所选度量、水平分布和基底径向坐标都影响具体的 \(\Phi\)。更换选择后仍得到合法的局部平凡化；两种平凡化之间的坐标变换在每条纤维上是一个微分同胚，并随基点光滑变化。定理保证局部乘积的存在，不指定一个天然的纤维识别。

\subsection{Ehresmann 定理的带边界版本}
\label{g2:subsec:1506:05}

有边界时，水平提升还必须尊重边界。先讨论基底无边界、总空间可能有边界的情形。

\begin{proposition}\label{g2:prop:ehresmann-source-boundary}
	设 \(E\) 是有边界光滑流形，\(B\) 无边界，\(p:E\to B\) 固有、满射且处处为淹没。若限制 \(p|_{\partial E}:\partial E\to B\) 也是淹没，则每个 \(y_0\in B\) 附近都有保持边界的局部平凡化
	\[
		(p^{-1}(U),p^{-1}(U)\cap\partial E)
		\cong(F,\partial F)\times U,
		\qquad \partial F=F\cap\partial E.
	\]
\end{proposition}
\begin{proof}
	关键是构造满足 \eqref{g2:eq:horizontal-right-inverse} 的 \(L\)，并使边界点的 \(L_e\) 取值于 \(\mathrm{T}_e\partial E\)。在边界点附近，\(p|_{\partial E}\) 满秩，故可将基底坐标 \(q^1,\ldots,q^b\) 的拉回补成边界坐标，再加一个边界定义函数 \(r\geq0\)。这些函数给出适配坐标 \((q,z,r)\)，其中 \(p\) 就是 \(q\) 投影。令基底坐标方向的局部提升保持 \((z,r)\) 不变，便得到沿边界切向的局部右逆。内点附近按普通淹没坐标选择局部右逆。

	取从属于这些邻域的光滑单位分解 \(\{\rho_i\}\)，把局部右逆 \(L_i\) 拼成 \(L=\sum_i\rho_iL_i\)。由于 \(\sum_i\rho_i=1\)，有 \(\mathrm{d}p\circ L=\operatorname{id}\)；在边界点，每个有贡献的 \(L_i\) 都取值于边界切空间，因此 \(L\) 也如此。其像与 \(\ker\mathrm{d}p\) 互补。

	现在按前述公式构造 \(X_u\)。它沿 \(\partial E\) 切向，其流保持边界，也保持内部。第一卷第13章第2节已说明这一点：在半空间坐标中，切向条件使法向方程写成 \(\dot r=a(q,z,r)r\)，所以 \(r=0\) 与 \(r>0\) 都在存在区间内保持。紧集上的统一短时存在和逃逸判据也适用于这样的向量场。故固有性保证正、反向提升均存在。原来的 \(\Phi\) 与 \(\Psi\) 保持边界且互逆，给出结论。适配坐标同时说明纤维是有边界子流形，边界恰为 \(F\cap\partial E\)。
\end{proof}

不能直接对任意度量取 \((\ker\mathrm{d}p)^\perp\)，就宣称所得水平分布沿边界切向：正交性并不自动带来这个条件。上面的局部右逆加单位分解，正是为解决这一点而作。

\ProseParagraphBreak

边界上的满秩条件也不能删除。考虑光滑带边界曲面
\[
	E=\{(x,y)\in\mathbb R^2:|y|\leq2,\ x^2+y^2\geq1\},
	\qquad p(x,y)=x.
\]
它的边界由两条水平直线和单位圆组成，彼此不相交，因而没有角点。映射 \(p\) 固有、满射，而且 \(\mathrm{d}p=\mathrm{d}x\) 在二维切空间上处处满射。然而，\(|x|<1\) 时纤维为两段闭区间，\(|x|>1\) 时纤维为一段闭区间，所以在 \(x=\pm1\) 附近不可能局部平凡。问题恰发生在圆周上的 \((\pm1,0)\)：此处 \(p|_{\partial E}\) 的微分为零。

\ProseParagraphBreak

如果基底也有边界，需说明总空间的哪一部分对应基底边界。一个适用于无角点流形的清楚版本如下。

\begin{proposition}\label{g2:prop:ehresmann-base-boundary}
	设 \(E,B\) 是有边界光滑流形，\(p:E\to B\) 是固有满淹没，并满足
	\[
		p^{-1}(\partial B)=\partial E.
	\]
	则 \(p\) 在基底的半空间邻域上也有光滑局部平凡化，纤维是紧的无边界流形。
\end{proposition}
\begin{proof}
	先看 \(e\in\partial E\)。取源和目标的边界坐标，分别记法向坐标为 \(r\geq0\)、\(s\geq0\)。因为 \(s\circ p\) 在 \(r=0\) 上为零，可写成
	\[
		s(p(q,r))=r\,a(q,r).
	\]
	其边界切向导数为零。满秩条件迫使 \(a(q,0)\ne0\)，而内点映到内点迫使 \(a(q,0)>0\)。因此 \(p\) 在法向上为同向的非退化映射；同时，\(\mathrm{d}p\) 将边界切空间满射到基底的边界切空间。于是 \(p|_{\partial E}:\partial E\to\partial B\) 是固有满淹没。

	固定 \(y_0\in\partial B\)。对这个无边界流形之间的限制映射应用 \cref{g2:thm:ehresmann}，得到某个边界邻域 \(V\) 上的平凡化，纤维 \(F=p^{-1}(y_0)\) 紧且无边界。把 \(V\) 缩小，使其闭包位于一张边界坐标图中，再取基底的半空间邻域 \(V\times[0,\varepsilon)\)，记其法向坐标为 \(s\)。在该邻域上选择 \(\mathrm{d}p\) 的光滑右逆，提升 \(\partial/\partial s\) 得到向量场 \(Y\)。上面的法向计算说明，\(Y\) 在 \(\partial E\) 上严格指向内侧。

	从边界平凡化中的点出发，沿 \(Y\) 流动时间 \(s\)。基底的 \(V\) 坐标保持不变，法向坐标恰好由 \(0\) 增为 \(s\)。紧性给出统一的短时存在，缩小 \(\varepsilon\) 后定义出
	\[
		F\times V\times[0,\varepsilon)\longrightarrow p^{-1}(V\times[0,\varepsilon)).
	\]
	任意右侧点的逆构造是沿 \(Y\) 反向流动，直到法向坐标降为零，再使用边界平凡化的逆。在到达零之前，轨道保持在内部；其基底轨道是紧线段，固有性防止在途中逃逸。到达边界时，向量场在局部半空间坐标中可光滑延拓到开集，故流与反向到达映射一直光滑到边界。唯一性证明这两个构造互逆。基底内点附近则直接应用无边界版本。
\end{proof}

例如闭圆盘上的投影 \(\overline D^2\to[-1,1]\)、\((x,y)\mapsto x\) 固有、满射，在二维切空间上处处满秩，但它不满足 \(p^{-1}(\partial B)=\partial E\)。端点上方的纤维为单点，内部纤维为线段，确实不能局部平凡。若还希望纤维本身有边界，那么它与有边界基底的乘积通常会出现角点；此时需另行规定各边界面的对应关系，上述无角点版本不涉及这种更一般的结构。

\subsection{光滑与全纯局部平凡性}
\label{g2:subsec:1506:06}

一个复流形同时是实光滑流形。全纯映射的坐标表达满足复可微条件，而一般的实光滑映射没有这个约束。因此，固有全纯淹没可以应用 \cref{g2:thm:ehresmann}，但所得平凡化首先是实光滑的。下面用复环面说明，全纯平凡化确实可能不存在。

取上半平面 \(\{\tau\in\mathbb C:\operatorname{Im}\tau>0\}\) 中的一个小开圆盘 \(U\)，定义格子与复环面
\[
	\Lambda_\tau=\mathbb Z+\mathbb Z\tau,
	\qquad E_\tau=\mathbb C/\Lambda_\tau.
\]
把这些环面组成一个族，可令 \(\mathbb Z^2\) 作用于 \(\mathbb C\times U\)：
\[
	(m,n)\cdot(z,\tau)=(z+m+n\tau,\tau),
\]
并取商空间 \(E=(\mathbb C\times U)/\mathbb Z^2\)。这个作用自由且适当不连续。后一条可以直接检查：在任意紧集上，\(\operatorname{Im}\tau\) 有正下界，而 \(z\) 有界；若平移前后的点都在该紧集中，先由虚部估计得到整数 \(n\) 有界，再由实部估计得到 \(m\) 有界，故只有有限个群元素能使紧集与其像相交。由于作用由全纯变换给出，局部商坐标使 \(E\) 成为复流形，映射
\[
	p:E\longrightarrow U,\qquad [z,\tau]\longmapsto\tau
\]
是全纯淹没，纤维正是 \(E_\tau\)。

这个族甚至有一个显式的实光滑全局平凡化：
\begin{equation}\label{g2:eq:complex-torus-smooth-family}
	(\mathbb R^2/\mathbb Z^2)\times U\longrightarrow E,
	\qquad ([s,t],\tau)\longmapsto[s+t\tau,\tau].
\end{equation}
它的逆在代表元上由
\[
	t=\frac{\operatorname{Im}z}{\operatorname{Im}\tau},
	\qquad s=\operatorname{Re}z-t\operatorname{Re}\tau
\]
给出；更换格子代表元只改变 \((s,t)\) 的整数部分。因此它是良定义的实光滑微分同胚。纤维 \(\mathbb R^2/\mathbb Z^2\) 紧，故 \(p\) 固有。

\ProseParagraphBreak

这些实光滑同胚的纤维是否也彼此双全纯？这里“双全纯”是指映射及其逆都全纯。下面的判别将把这个问题归结为格子的变化。

\begin{lemma}\label{g2:lem:complex-torus-lattices}
	若 \(\mathbb C/\Lambda\) 与 \(\mathbb C/\Lambda'\) 双全纯，则存在 \(a\in\mathbb C\setminus\{0\}\)，使 \(a\Lambda=\Lambda'\)。反过来，这个条件也给出双全纯同构。
\end{lemma}
\begin{proof}
	商映射 \(\mathbb C\to\mathbb C/\Lambda\) 是万有覆盖。给定双全纯同构 \(f\)，把它与前一个覆盖映射复合。由于 \(\mathbb C\) 单连通，覆盖映射的提升定理给出提升 \(F:\mathbb C\to\mathbb C\)；覆盖映射局部双全纯，所以 \(F\) 也是全纯的。这里用到 \cref{g2:prop:covering-lift-criterion}：\(\mathbb C\) 道路连通、局部道路连通，且基本群为零，因此满足提升所需的子群条件。

	对每个 \(\lambda\in\Lambda\)，差 \(F(z+\lambda)-F(z)\) 取值于离散集 \(\Lambda'\)，又随 \(z\) 连续，故为常量。求导得 \(F'(z+\lambda)=F'(z)\)。于是 \(F'\) 在一个紧基本平行四边形上有界，周期性又使它在全平面上有界。由有界整函数为常数的 Liouville 定理，\footnote{Joseph Liouville（1809--1882），法国数学家。}\(F'=a\)，从而 \(F(z)=az+c\)。这里使用的复分析结论可由 Cauchy 导数估计迅速推出。\footnote{Augustin-Louis Cauchy（1789--1857），法国数学家，研究极限与收敛的严格表述，建立了复函数积分理论的基本结果。}若整函数 \(h\) 满足 \(|h|\leq M\)，则在任意半径 \(R\) 的圆盘中心有 \(|h'|\leq M/R\)，令 \(R\to\infty\) 即得 \(h'=0\)。

	因为 \(f\) 为局部同构，\(a\ne0\)；又因 \(F(z+\lambda)-F(z)=a\lambda\in\Lambda'\)，有 \(a\Lambda\subseteq\Lambda'\)。若某个 \(\mu\in\Lambda'\) 不属于 \(a\Lambda\)，则 \(z\) 与 \(z+\mu/a\) 在 \(\mathbb C/\Lambda\) 中不同，却被 \(f\) 映到同一点，与单射性矛盾。故 \(a\Lambda=\Lambda'\)。反过来，\(z\mapsto az\) 在商空间上直接诱导所需同构。
\end{proof}

对 \(\Lambda_\tau=\mathbb Z+\mathbb Z\tau\)，上述格子条件等价于
\begin{equation}\label{g2:eq:torus-modular-equivalence}
	\sigma=\frac{A\tau+B}{C\tau+D},
	\qquad
	\begin{pmatrix}A&B\\C&D\end{pmatrix}\in\mathrm{SL}(2,\mathbb Z).
\end{equation}
理由是，\((1,\sigma)\) 与 \((a,a\tau)\) 是同一个格子的两组正定向整数基，换基矩阵的行列式为 \(1\)；将两个基的关系相除便得到此式。反过来，若有此式，则 \((C\tau+D,A\tau+B)\) 是 \(\Lambda_\tau\) 的整数基，乘以 \((C\tau+D)^{-1}\) 后得到 \(\Lambda_\sigma\)。

固定一个 \(\tau_0\)，整数矩阵只有可数多个，所以与 \(E_{\tau_0}\) 双全纯的参数至多可数。任意开圆盘却有不可数多个参数，故不存在 \(\tau_0\) 的邻域，使该族中的所有纤维都与 \(E_{\tau_0}\) 双全纯。全纯局部平凡化会给出这样的纤维同构，因此不可能存在。

式 \eqref{g2:eq:complex-torus-smooth-family} 将差别呈现得很具体：所有纤维的实光滑结构相同，格子的复形状却随 \(\tau\) 改变。Ehresmann 证明中的辅助度量、水平补空间和单位分解都是实光滑构造，它们正好提供前一种意义的平凡化。


\begin{chaptersummary}
	坐标与截断给出初始嵌入，切向和割线坏方向的维数估计给出弱 Whitney 界；非紧源还要保持固有性，才能由单射浸入得到嵌入。强界要求继续处理一般位置留下的双点，高维 Whitney 移动、低维模型和非紧源构造分别承担不同任务。法向局部指数在零截面处的微分可逆，变半径控制把这些局部逆拼成管状邻域；边界的向内向量场则给出领圈，为双化与粘合提供横向坐标。同位延拓把嵌入族的速度扩展到环境，再由共同紧支集保证环境流存在。Ehresmann 定理使用另一种存在性控制：固有性把紧基底路径的提升限制在紧集内，径向提升及其反向过程于是给出光滑局部平凡化。源有边界时还须控制边界限制的秩；即使已获光滑平凡化，复结构也可以随参数变化。
\end{chaptersummary}

\BookExercises

\begin{exercise}\label{g2:ex:projection-loses-properness}
	设 \(F:\mathbb R\to\mathbb R^2\)、\(F(t)=(\arctan t,t)\)，令 \(p:\mathbb R^2\to\mathbb R\) 为第一坐标投影。证明 \(F\) 是固有嵌入，而 \(p\circ F\) 是非固有的嵌入。将 \(F(\mathbb R)\) 包含在竖直轴周围一个有界半径管中，并说明本例中\cref{g2:lem:proper-projection-tube}的哪个假设不成立。
\end{exercise}
\begin{solution}
	\(F'(t)=(1/(1+t^2),1)\ne0\)，第二坐标使 \(F\) 单射，且它在像上的连续逆就是第二坐标限制，故为嵌入。若 \(K\subset\mathbb R^2\) 紧，则 \(F^{-1}(K)\) 闭，并被 \(K\) 的第二坐标界限制在某个有界区间内，故紧。因此 \(F\) 固有。

	\(p\circ F=\arctan\) 的导数恒正，它是 \(\mathbb R\) 到 \((-\pi/2,\pi/2)\) 的微分同胚，所以作为到 \(\mathbb R\) 的映射仍为嵌入。但紧区间 \([-\pi/2,\pi/2]\) 的原像为整个 \(\mathbb R\)，故不固有。由于 \(|\arctan t|<\pi/2\)，原曲线处于竖直轴周围半径 \(\pi/2\) 的管中；本次投影沿竖直方向，恰好压掉了管轴，因此引理所要求的投影方向不平行管轴的条件失效。
\end{solution}

\begin{exercise}\label{g2:ex:proper-immersion-finite-branches}
	设 \(f:M\to N\) 是无边界光滑流形间的固有浸入。证明每个纤维 \(f^{-1}(y)\) 有限。若 \(f^{-1}(y)=\{x_1,\ldots,x_s\}\)，证明存在 \(y\) 的开邻域 \(V\) 以及两两不交的 \(x_i\) 的开邻域 \(U_i\)，使 \(f^{-1}(V)\subset\bigcup_iU_i\)，且每个 \(f|_{U_i}\) 为嵌入。这说明像在 \(y\) 附近有多少条局部分支？
\end{exercise}
\begin{solution}
	常秩定理说明浸入在每个源点附近都是嵌入，故纤维是离散集。固有性又保证单点的原像紧。若紧离散纤维无限，其逐点孤立邻域构成没有有限子覆盖的开覆盖，矛盾，故纤维有限。

	在有限多个 \(x_i\) 附近取两两不交的开集 \(U_i\)，使各限制为嵌入。若不存在所需 \(V\)，可选 \(y_j\to y\) 及 \(z_j\notin\bigcup_iU_i\)，满足 \(f(z_j)=y_j\)。集合 \(\{y\}\cup\{y_j:j\ge1\}\) 紧，其固有原像也紧，因此 \(z_j\) 有子列收敛到 \(z\)。连续性给出 \(f(z)=y\)，故 \(z=x_i\) 对某个 \(i\) 成立；但 \(U_i\) 开，子列最终应落入 \(U_i\)，矛盾。于是所需 \(V\) 存在。像在 \(V\) 中由 \(s\) 条嵌入分支组成；不同分支可能相交，也可能局部重合，因此这不等于它们有 \(s\) 个不同的像。
\end{solution}

\begin{exercise}\label{g2:ex:projective-rank-model}
	设 \(m\ge2\)，\(V_m\) 为实对称无迹 \((m+1)\times(m+1)\) 矩阵空间，\(A_m=\operatorname{diag}(1,\ldots,1,-m)\)，\(P_m:V_m\to V_m/\mathbb RA_m\) 为商映射。证明
	\[
		[x]\longmapsto P_m\left(xx^{\mathsf T}-\frac1{m+1}I\right),\qquad |x|=1,
	\]
	给出 \(\mathbb RP^m\) 到维数为 \((m+1)(m+2)/2-2\) 的 Euclid 空间的光滑嵌入。判断哪些 \(m\ge2\) 使这一维数等于强 Whitney 界 \(2m\)，并解释同一个公式在 \(m=1\) 时为何失效。
\end{exercise}
\begin{solution}
	对称矩阵空间的维数是 \((m+1)(m+2)/2\)，无迹条件及取一维商各减一。公式在 \(x\) 和 \(-x\) 处相同，故在射影空间上良定义。若两点像相等，则 \(xx^{\mathsf T}-yy^{\mathsf T}=\lambda A_m\)。左端秩至多 \(2\)，右端在 \(\lambda\ne0\) 时秩为 \(m+1\ge3\)，故 \(\lambda=0\)。于是两秩一正交投影相同，像直线相同，得到 \([x]=[y]\)。

	对 \(v\perp x\)，微分的代表为 \(xv^{\mathsf T}+vx^{\mathsf T}\)。它仍有秩至多 \(2\)，故若在商中为零，就必须为零；把此矩阵作用于 \(x\)，得到 \(v=0\)。所以映射是单射浸入。射影空间紧，因而它是嵌入。

	等式 \((m+1)(m+2)/2-2=2m\) 化为 \((m-2)(m+1)=0\)，在所给范围只有 \(m=2\)。当 \(m=1\) 时，取 \(x=(1,0)\)、\(y=(0,1)\)，便有 \(xx^{\mathsf T}-yy^{\mathsf T}=A_1\)，两个不同射影点在商中重合；秩比较失效的原因是 \(A_1\) 本身只有秩 \(2\)。
\end{solution}

\begin{exercise}\label{g2:ex:projection-many-double-points}
	对正整数 \(r\)，设
	\[
		\gamma_r(t)=(\sin t,\sin(2rt),\cos t),\qquad t\in\mathbb R/2\pi\mathbb Z.
	\]
	证明 \(\gamma_r\) 是光滑嵌入，而其前两个坐标组成的平面投影是恰有 \(2r-1\) 个横截双点的浸入，且没有三重点。
\end{exercise}
\begin{solution}
	第一、第三坐标组成标准圆周参数，所以 \(\gamma_r\) 单射且微分不为零；源紧，故是嵌入。投影的微分为 \((\cos t,2r\cos(2rt))\)。若 \(\cos t=0\)，则 \(t=\pi/2\) 或 \(3\pi/2\) 模 \(2\pi\)，此时 \(\cos(2rt)=(-1)^r\ne0\)，故投影始终为浸入。

	若两个不同参数 \(t,u\) 的投影相同，第一坐标给出 \(u=\pi-t\) 模 \(2\pi\)，而第二坐标变成 \(\sin(2rt)=-\sin(2rt)\)。所以 \(t=k\pi/(2r)\)，其中 \(0\le k<4r\)。排除对合 \(t\mapsto\pi-t\) 的两个不动点 \(\pi/2,3\pi/2\)，其余 \(4r-2\) 个参数两两成对，故共有 \(2r-1\) 个双点。一个给定的正弦值在圆周上至多有两个原像，所以不会出现三重点。

	在配对的 \(t,u=\pi-t\) 处，两条切向量分别为
	\[
		(\cos t,2r\cos(2rt)),\qquad(-\cos t,2r\cos(2rt)).
	\]
	其行列式为 \(4r\cos t\cos(2rt)\)。候选参数处 \(\cos(2rt)=\pm1\)，而 \(\cos t=0\) 的两个参数已经排除，故所有双点都横截。
\end{solution}

\begin{exercise}\label{g2:ex:normal-trivial-defining-map}
	设 \(S\subset N\) 是余维 \(c\) 的无边界嵌入子流形。证明：法丛 \(\nu_NS\) 平凡，当且仅当存在 \(S\) 的开邻域 \(U\subset N\) 和光滑淹没 \(q:U\to\mathbb R^c\)，使 \(q^{-1}(0)=S\)。这里的方程只要求在一个邻域内成立。
\end{exercise}
\begin{solution}
	若存在这样的 \(q\)，正则水平集的切空间公式给出 \(\ker\mathrm{d}q_x=\mathrm{T}_xS\)。因此 \(\mathrm{d}q_x\) 在法商上诱导线性同构 \(\nu_xS\to\mathbb R^c\)。该同构光滑依赖 \(x\)，于是 \([v]\mapsto(x,\mathrm{d}q_xv)\) 是 \(\nu_NS\cong S\times\mathbb R^c\) 的丛平凡化。

	反过来，取平凡化 \(A:\nu_NS\to S\times\mathbb R^c\) 和管状映射 \(\tau:E_\varepsilon\to U\)。令
	\[
		q=\operatorname{pr}_2\circ A\circ\tau^{-1}:U\to\mathbb R^c.
	\]
	\(A(E_\varepsilon)\) 是乘积流形中的开集，而第二坐标投影在该开集上的微分满射，所以 \(q\) 是淹没。它的零点对应法丛零截面，正好是 \(S\)。证明同时说明：一个一般子流形不能用同一组独立方程沿全体基点描述，障碍可能来自法丛的拼接。
\end{solution}

\begin{exercise}\label{g2:ex:diagonal-tubular-model}
	设 \(M\) 为无边界光滑流形，\(\Delta_M\subset M\times M\) 为对角线。用 \([(u,v)]\mapsto v-u\) 证明 \(\nu_{M\times M}\Delta_M\cong\mathrm{T}M\)。选择 \(\mathrm{T}M\) 上一个联络，证明
	\[
		\Psi(x,w)=(x,\exp_x^\nabla w)
	\]
	在 \(\mathrm{T}M\) 的零截面某个变半径邻域上给出对角线的管状邻域映射。交换两个因子的映射在对角线的法丛上诱导什么作用？
\end{exercise}
\begin{solution}
	在 \((x,x)\) 处，\(\mathrm{T}_{(x,x)}(M\times M)=\mathrm{T}_xM\oplus\mathrm{T}_xM\)，而对角线的切空间为 \(\{(u,u)\}\)。线性映射 \((u,v)\mapsto v-u\) 满射，核恰为对角切空间，因此诱导法商同构；它由切向量的自然运算定义，随 \(x\) 光滑，给出所需丛同构。

	\(\Psi(x,0)=(x,x)\)。在零截面处，由局部指数的微分公式得
	\[
		\mathrm{d}\Psi_{(x,0)}(u,w)=(u,u+w),
	\]
	这是线性同构，且其法商分量为 \((u+w)-u=w\)。对角线在 Hausdorff 乘积中闭，故\cref{g2:lem:zero-section-shrinking}直接给出一个变半径邻域，使 \(\Psi\) 成为到开像的微分同胚。以上法商计算又保证正规化条件，所以它是管状映射。

	交换映射把 \((u,v)\) 变成 \((v,u)\)，法商中的差由 \(v-u\) 变为 \(u-v\)，因此所诱导的作用是 \(\mathrm{T}M\) 各纤维上的负恒等。
\end{solution}

\begin{exercise}\label{g2:ex:tubular-normalization-sign}
	对 \(S=\{0\}\subset\mathbb R\)，考虑零点附近的两张邻域参数化 \(\tau_+(v)=v\) 和 \(\tau_-(v)=-v\)。证明不存在共同零点邻域上的联合光滑族 \(\tau_t\)，使各 \(\tau_t\) 都是局部微分同胚、固定零点，且端点分别为这两张参数化。解释为何这不违反\cref{g2:prop:tubular-uniqueness}，并求缩放公式 \(\delta_{1/t}\tau_-\delta_t\) 在 \(t\to0\) 时的极限。
\end{exercise}
\begin{solution}
	若这样的族存在，函数 \(a(t)=\mathrm{d}(\tau_t)_0\) 连续且处处非零，因为每张映射在零点都是局部微分同胚。然而 \(a(0)=1\)、\(a(1)=-1\)，介值定理迫使某个 \(a(t)=0\)，矛盾。

	这里法丛就是 \(\mathbb R\)，\(\tau_-\) 在法商上诱导 \(-\operatorname{id}\)，没有满足管状映射定义中的恒等正规化条件，故唯一性命题不适用。对 \(t>0\)，直接计算
	\[
		\delta_{1/t}\circ\tau_-\circ\delta_t(v)=\frac{-tv}{t}=-v,
	\]
	极限仍为负恒等，并非恒等。若先把纤维识别乘以 \(-1\)，则参数化变成 \(v\mapsto v\)，这时两者才具有同一正规化。
\end{solution}

\begin{exercise}\label{g2:ex:boundary-functions-ratio}
	设 \(r_0,r_1:M\to[0,\infty)\) 都是光滑带边界流形 \(M\) 的边界定义函数，即零点集均为 \(\partial M\)，且边界上的微分均非零。证明在边界的某个开邻域中存在正光滑函数 \(a\)，使 \(r_1=ar_0\)，并求二者在边界上的微分关系。
\end{exercise}
\begin{solution}
	按\cref{g2:lem:boundary-defining-function}的归一化步骤，在边界附近取 \(Xr_0=1\) 的向内向量场。领圈证明中的流参数化 \(F(x,t)\) 满足 \(r_0(F(x,t))=t\)，其中 \(t\) 可在一个随 \(x\) 变化的正区间内取值。令 \(b(x,t)=r_1(F(x,t))\)。因为 \(b(x,0)=0\)，微积分基本定理给出
	\[
		b(x,t)=t\int_0^1\partial_tb(x,ut)\,\mathrm{d}u.
	\]
	积分定义了包括 \(t=0\) 在内的光滑函数。\(r_1\) 在内部为正且边界微分非零，所以 \(\partial_tb(x,0)>0\)；缩小邻域后，该积分仍正。通过 \(F\) 的逆把它搬回 \(M\)，便得到所需 \(a\)。对 \(r_1=ar_0\) 求微分，并用边界上的 \(r_0=0\)，得到 \(\mathrm{d}r_1=a\,\mathrm{d}r_0\)。正因这个比例为正，由任一边界定义函数判断的向内方向都相同。
\end{solution}

\begin{exercise}\label{g2:ex:interior-homotopy-equivalence}
	设 \(M\) 为光滑带边界流形，\(j:\operatorname{Int}M\hookrightarrow M\) 为包含映射。用领圈证明 \(j\) 是同伦等价；允许 \(M\) 及其边界非紧。
\end{exercise}
\begin{solution}
	取领圈 \(c:\partial M\times[0,1)\to M\)。选择非负光滑函数 \(\chi:[0,1)\to[0,1]\)，使它在零附近为1，在 \([1/2,1)\) 上为零。再选 \(\varepsilon>0\)，使 \(\varepsilon\sup|\chi'|<1/2\)。对 \(0\le s\le1\) 定义
	\[
		f_s(t)=t+s\varepsilon\chi(t).
	\]
	有 \(f_s'(t)>1/2\)，且 \(f_s(1/2)=1/2\)，所以 \(f_s\) 将 \([0,1)\) 映入自身，并在 \([1/2,1)\) 上恒等。令 \(H_s(c(x,t))=c(x,f_s(t))\)，在领圈外令 \(H_s\) 为恒等。两种表达在领圈参数接近1的区域一致，故拼成光滑同伦，且 \(H_0=\operatorname{id}_M\)。

	由于 \(f_1(0)=\varepsilon>0\)，\(H_1\) 的像完全落在内部，记所得映射为 \(r:M\to\operatorname{Int}M\)。于是 \(j\circ r\simeq\operatorname{id}_M\)。若起点在内部，则领圈参数 \(t>0\)，整个过程中 \(f_s(t)\ge t>0\)，故限制同伦也留在内部，给出 \(r\circ j\simeq\operatorname{id}_{\operatorname{Int}M}\)。领圈已把非紧边界统一参数化为 \([0,1)\)，构造不要求边界紧。
\end{solution}

\begin{exercise}\label{g2:ex:isotopy-complement-normal}
	设 \(K\) 为紧无边界光滑流形，\(N\) 为无边界光滑流形，\(i_t:K\to N\) 为光滑同位。证明 \(N\setminus i_0(K)\) 与 \(N\setminus i_1(K)\) 微分同胚；再证明经 \(i_0,i_1\) 拉回到 \(K\) 的两个法丛同构。说明同位延拓如何同时给出这两个结论。
\end{exercise}
\begin{solution}
	由\cref{g2:thm:isotopy-extension}，存在环境同位 \(\Phi_t\)，使 \(\Phi_t i_0=i_t\)。\(\Phi_1\) 是 \(N\) 的微分同胚，并把 \(i_0(K)\) 恰好送到 \(i_1(K)\)，因此限制到两个补集仍为微分同胚，逆是 \(\Phi_1^{-1}\) 的限制。

	对 \(\Phi_1 i_0=i_1\) 求微分，得
	\[
		\mathrm{d}\Phi_1\circ\mathrm{d}i_0=\mathrm{d}i_1.
	\]
	所以在每个 \(x\in K\) 上，\(\mathrm{d}\Phi_1\) 把 \(\mathrm{d}i_0(\mathrm{T}_xK)\) 同构地映到 \(\mathrm{d}i_1(\mathrm{T}_xK)\)，并在商空间上诱导同构
	\[
		\frac{\mathrm{T}_{i_0(x)}N}{\mathrm{d}i_0(\mathrm{T}_xK)}\longrightarrow
		\frac{\mathrm{T}_{i_1(x)}N}{\mathrm{d}i_1(\mathrm{T}_xK)},\qquad
		[v]\longmapsto[\mathrm{d}\Phi_1(v)].
	\]
	该映射及其逆光滑依赖 \(x\)，故是法丛同构。这里用法商表达最直接；若以度量正交补代表法向量，则须在应用 \(\mathrm{d}\Phi_1\) 后再投影到新的正交补，因为一般微分同胚并不保持所选度量。
\end{solution}

\begin{exercise}\label{g2:ex:moving-compact-support-escape}
	取 \(\chi\in C_c^\infty(\mathbb R)\)，使 \(\chi=1\) 于零的某邻域，\(\operatorname{supp}\chi\subset[-1,1]\)。在实线上定义时间依赖向量场
	\[
		X_t(x)=\begin{cases}
			(1-t)^{-2}\chi\bigl(x-(1-t)^{-1}\bigr)\,\partial_x,&0\le t<1,\\
			0,&t=1.
		\end{cases}
	\]
	证明 \((t,x)\mapsto X_t(x)\) 在 \([0,1]\times\mathbb R\) 上光滑，每个固定时刻的向量场都有紧支集，但从 \(x(0)=1\) 出发的解不能延续到 \(t=1\)。它说明\cref{g2:lem:uniformly-supported-evolution}中的哪个条件不能只逐时刻要求？
\end{exercise}
\begin{solution}
	在 \(t<1\) 处表达显然光滑。给定有限点 \(x_0\)，当 \(t\) 充分接近1且 \(x\) 位于 \(x_0\) 的一个固定有界邻域时，\((1-t)^{-1}\) 与该邻域的距离大于1，故 \(\chi\) 恒为零。于是向量场在 \((1,x_0)\) 的某邻域内恒为零，包括所有阶导数，证明联合光滑性。

	对每个 \(t<1\)，支集包含在紧区间 \([(1-t)^{-1}-1,(1-t)^{-1}+1]\) 内；\(t=1\) 时支集为空。但曲线
	\[
		x(t)=(1-t)^{-1},\qquad0\le t<1,
	\]
	满足 \(x'(t)=(1-t)^{-2}\) 且 \(\chi(x(t)-(1-t)^{-1})=1\)，因此正是所给初值问题的解。它在 \(t\to1\) 时趋于无穷，无法作为实值连续曲线延续。所有时刻的支集并没有包含在一个共同紧集中；逐时刻紧支集允许有效区域逃向无穷，因而不足以保证非自治方程在整个时间区间存在。
\end{solution}

\begin{exercise}\label{g2:ex:compact-fibers-not-proper}
	令 \(E=(S^1\times\mathbb R)\sqcup(S^1\times(0,\infty))\)，令 \(p:E\to\mathbb R\) 在每个分支上都等于第二坐标投影。证明 \(p\) 是满淹没，每条纤维都紧，但 \(p\) 不固有，并且在 \(0\) 附近不局部平凡。
\end{exercise}
\begin{solution}
	两个分支都是无边界二维流形，投影的微分都满射；第一个分支已保证满射性。若 \(t\leq0\)，纤维是一份 \(S^1\)；若 \(t>0\)，纤维是两份 \(S^1\) 的不交并，所以每条纤维均紧。固定 \(z_0\in S^1\)，在第二个分支中取点列 \((z_0,1/j)\)。它位于 \(p^{-1}([0,1])\) 内，却在 \(E\) 中没有收敛子列：第二分支不包含参数零的点，而两个分支都是开闭集，点列也不能收敛到第一分支。因此 \(p^{-1}([0,1])\) 不紧，\(p\) 不固有。\(0\) 的任意邻域都含正参数，纤维的连通分支数由一变二，故不存在局部平凡化。这说明逐条纤维紧不能代替紧集原像紧。
\end{solution}

\begin{exercise}\label{g2:ex:circle-proper-submersion}
	设 \(k\geq2\)，考虑 \(p:S^1\to S^1\)、\(p(z)=z^k\)。证明它是固有满淹没，写出小圆弧上以 \(k\) 个点为纤维的局部平凡化，并证明它没有全局平凡化。
\end{exercise}
\begin{solution}
	源空间紧，故 \(p\) 固有；每个单位复数都有 \(k\) 个互异的单位根，故满射。用角坐标表示为 \(\theta\mapsto k\theta\)，其导数非零，因此是淹没。设圆弧 \(U\) 上选定光滑幅角 \(\alpha:U\to\mathbb R\)，则
	\[
		\{0,1,\ldots,k-1\}\times U\longrightarrow p^{-1}(U),
		\qquad (j,w)\longmapsto\exp\!\left(\frac{\mathrm{i}(\alpha(w)+2\pi j)}{k}\right)
	\]
	是局部平凡化：其像覆盖恰好 \(k\) 个互不相交的开圆弧，在每一份上与 \(p\) 互逆。若存在全局平凡化，则 \(S^1\) 同胚于 \(k\) 份 \(S^1\) 的不交并；前者连通，后者有 \(k\) 个连通分支，矛盾。
\end{solution}

\begin{exercise}\label{g2:ex:proper-submersion-real-line}
	设 \(p:E\to\mathbb R\) 是无边界非空光滑流形上的固有满淹没。证明 \(p\) 有全局平凡化 \(E\cong p^{-1}(0)\times\mathbb R\)。再证明非空紧无边界流形上不存在取值于 \(\mathbb R\) 的处处淹没光滑函数。
\end{exercise}
\begin{solution}
	选取 \(\mathrm{d}p\) 的光滑右逆，令 \(X\) 为常向量 \(1\) 的提升，则 \(\mathrm{d}p(X)=1\)。若 \(e(t)\) 是其积分曲线，则 \(p(e(t))=p(e(0))+t\)。任意有限时间区间中的基底轨道都包含在某个紧区间 \([a,b]\) 内，因此提升留在紧集 \(p^{-1}([a,b])\) 中。逃逸判据说明 \(X\) 完备；这里每个有限时间区间可以使用不同的紧集，不需要 \(E\) 本身紧。

	记其全局流为 \(\varphi^t\)，令 \(F=p^{-1}(0)\)。则
	\[
		\Phi:F\times\mathbb R\longrightarrow E,\quad \Phi(z,t)=\varphi^t(z),
		\qquad \Psi(e)=(\varphi^{-p(e)}(e),p(e))
	\]
	光滑且互逆，并满足 \(p\Phi(z,t)=t\)，这就是所求全局平凡化。最后，紧流形上的实值连续函数取得最大值；在无边界流形的极大点处，光滑函数的微分为零，因而不能处处为淹没。
\end{solution}

\begin{exercise}\label{g2:ex:horizontal-holonomy-product}
	在 \(E=\mathbb R^2/\mathbb Z^2\) 上用 \(([\theta],[\phi])\) 表示坐标，令 \(p:E\to\mathbb R/\mathbb Z\) 为第一坐标投影。给定 \(\alpha\in\mathbb R\)，以 \(X=\partial_\theta+\alpha\partial_\phi\) 张成水平分布。计算沿基底闭路 \(c(t)=[t]\)、\(0\leq t\leq1\) 的水平搬运，并说明这个计算如何区分丛的全局平凡性与所选水平分布的性质。
\end{exercise}
\begin{solution}
	向量场 \(X\) 在整数平移下不变，故在商空间上良定义，且 \(\mathrm{d}p(X)=\partial_\theta\)。经过 \(([0],[\phi_0])\) 的水平提升为
	\[
		e(t)=([t],[\phi_0+\alpha t]).
	\]
	因此绕基底一周后，初始纤维上的搬运映射为 \([\phi]\mapsto[\phi+\alpha]\)。它在 \(\alpha\notin\mathbb Z\) 时不是恒等映射；但 \(p\) 本身就是两个圆的乘积投影，始终全局平凡。故沿闭路搬运的结果还取决于所选水平分布。非平凡的这一搬运映射不能单独证明丛不平凡；这里改取 \(\alpha=0\) 就得到恒等搬运。
\end{solution}

% !TeX root = ../../Book2.tex
% 纲要标题：## 第16章　映射度、相交与局部指标
% 纲要位置：几何学计划纲要.md，第 1621 行。
% 正文状态：初稿；尚未正式审核。

% 纲要内容开始
% 
% > 本章目标与证明范围：用横截性把代数对偶解释为几何计数，再以一维配边证明不变性。所有整数计数有明确取向规则，非横截情形通过扰动而非朴素点数处理。
% 
% 纲要内容结束

\chapter{映射度、相交与局部指标}
\label{g2:ch:16}
\BookChapterNumber{g2:ch:16}

\begin{chapterabstract}
	正则值的带符号原像数、子流形的相交数、向量场零点指标和不动点指数，都把局部方向信息汇成整数。本章比较映射度的正则值、同调及积分定义，说明固有性如何控制非紧情形；用横截扰动和有向一维配边构造相交数，并与帽积配对比较。局部球面度给出向量场指标及其可加性，截面同伦则保持闭流形上的总指标。最后把不动点看作图像与对角线的交点，计算对角类，得到局部指数之和等于上同调交错迹的 Lefschetz 公式。
\end{chapterabstract}

\begin{learninggoals}
	\begin{itemize}
		\item 由正则值的微分符号计算映射度，并证明它与基本类及顶次形式积分的计算一致。
		\item 核对非紧映射的固有性、形式的紧支集及同伦的整体固有性，判断度公式何时适用。
		\item 按法向优先约定确定相交符号，用一维配边证明扰动不变性，并把交点计数写成帽积配对。
		\item 由小球边界的方向度计算孤立零点指标，证明坐标无关性、非退化公式及扰动后的可加性。
		\item 用截面横截同伦比较闭流形上不同向量场的总指标，并区分局部指标与零点个数。
		\item 从图像、对角类和杯积符号推导 Lefschetz 公式，由非零 Lefschetz 数判定不动点的存在。
	\end{itemize}
\end{learninggoals}

方程 \(x^3-x=0\) 有三个实根 \(-1,0,1\)，函数在这些点的导数分别为 \(2,-1,2\)。若只数根，答案是3；若同时记录穿过目标值的方向，答案则是 \(1-1+1=1\)。把右端换成充分大的正数，方程只有一个实根，导数为正，带符号的答案仍然是1。原像点会出现或消失，方向计数却保留了一个稳定的整数。

\ProseParagraphBreak

同样的现象也出现在曲线相交、向量场零点与不动点问题中。两条曲线相切后可能产生两个交点；一个退化零点可能裂成多个非退化零点；一个自映射的不动点可以随着映射改变而移动。把这些现象放进一参数族，满足横截条件的全部解组成一维流形，其边界记录初、末状态。方向规定使区间的两个端点贡献相反符号，这便是计数保持不变的几何原因。

\ProseParagraphBreak

\secref{g2:sec:1601}先把这种计数与第3章的局部度、第6章的基本类和第7章的积分比较连接起来。\secref{g2:sec:1602}将目标点换成子流形，比较相交符号与法向方向，再证明几何交数等于相应的对偶配对。\secref{g2:sec:1603}考察向量场在小球边界上的朝向，得到能容纳退化零点的指标。\secref{g2:sec:1604}则把方程 \(f(x)=x\) 写成图像与对角线相交，从局部的 \(\det(I-\mathrm{d}f_x)\) 走到上同调上的交错迹。

\ProseParagraphBreak

本章使用第4章的外积与杯积符号、第6章的帽积和 Poincaré 对偶，以及第14章的相对横截性和第15章的管状邻域。关键调用处会恢复定向、维数、紧性及支集条件：它们决定交集是否有限、符号能否连续选择，以及一族解是否会逃到无穷远。正则值与积分度的比较可参阅\cite[第457--459页，定理17.35及命题17.36]{Lee2012Smooth}；有限复形及其回缩空间的 Lefschetz 定理可参阅\cite[第2.C节，The Lefschetz Fixed Point Theorem]{Hatcher2002}。

% !TeX root = ../../Book2.tex
% 纲要标题：### 16.1 映射度的等价定义
% 纲要位置：几何学计划纲要.md，第 1625 行。
% 正文状态：初稿；尚未正式审核。

% 纲要内容开始
% 
% * 正则值处的局部符号
% * 同调定义与积分定义的比较
% * 固有映射的映射度
% * 复合公式及典型计算
% 
% 纲要内容结束

\section{映射度的等价定义}
\label{g2:sec:1601}

圆周映射 \(z\mapsto z^k\) 的度既可以从顶维同调读出，也可以在一个目标点处数原像，还可以由角形式的积分求得。这三个计算测量的是同一件事：源流形按方向把目标覆盖了多少次。对一般流形，临界点会使覆盖片出现、合并或消失，非紧源还允许原像逃到无穷远。本节先解释每个正则原像的符号，再证明三种度的一致性，并把无穷远处所需的条件写清楚。除另有说明，流形均有限维、Hausdorff、第二可数。

\subsection{正则值处的局部符号}
\label{g2:subsec:1601:01}

设 \(f:M^n\to N^n\) 为定向无边界光滑流形间的光滑映射，\(n\ge1\)。若 \(y\) 是正则值，则每个 \(x\in f^{-1}(y)\) 处的微分 \(\mathrm{d}f_x\) 是同构。在源、目标各取一组正向基，定义
\[
	\varepsilon_x(f)=\operatorname{sgn}\det(\mathrm{d}f_x)\in\{1,-1\}.
\]
更换正向基只会让行列式乘以正数，因此该符号与基无关。它记录这一小片是保持方向还是反转方向，不能用行列式的绝对值替代。

\begin{proposition}\label{g2:prop:regular-local-degree-sign}
	设 \(f:M^n\to N^n\) 为定向无边界光滑流形间的光滑映射，\(n\ge1\)，且 \(\mathrm{d}f_x\) 可逆。按定向局部同调生成元定义的局部度满足
	\[
		\deg_x(f)=\varepsilon_x(f).
	\]
	若 \(f\) 固有，则每个正则值的原像有限。
\end{proposition}
\begin{proof}
	\cref{g2:def:local-degree}中的空间对定义可直接用于定向流形：由坐标切除，\(H_n(M,M\setminus\{x\};\mathbb Z)\cong\mathbb Z\)，其正生成元由定向指定。在正向坐标中把 \(x,f(x)\) 移到原点，写 \(f(z)=Az+r(z)\)，其中 \(A=\mathrm{d}f_x\) 可逆、\(|r(z)|/|z|\to0\)。取 \(c>0\) 使 \(|Az|\ge c|z|\)，再把坐标球缩小，使 \(|r(z)|<c|z|/2\)。于是
	\[
		|Az+s r(z)|\ge c|z|/2>0\qquad(z\ne0,\ 0\le s\le1).
	\]
	这给出穿孔空间对上的同伦，所以局部度等于线性映射 \(A\) 的局部度。

	对可逆矩阵作 Gram--Schmidt 分解，上三角部分可在保持对角元为正时连续变为恒等；正交部分则用平面旋转化为恒等，或化为一次坐标反射，分别对应 \(\det A>0\) 与 \(\det A<0\)。这些变化始终保持可逆，故仍是穿孔空间对同伦。恒等映射保持局部生成元，反射使其变号：通过相对球与边界球的连接同构，这正是\cref{g2:prop:reflection-antipodal-degree}已证的反射度。于是局部度为 \(\operatorname{sgn}\det A\)。

	第一卷第11章第3节的逆函数定理又说明正则原像的每个点都孤立。若 \(f\) 固有，纤维 \(f^{-1}(y)\) 是紧集；紧离散集有限，得到最后一项。
\end{proof}

对固有 \(f\) 及其正则值 \(y\)，可以先记带符号计数
\[
	d_y(f)=\sum_{x\in f^{-1}(y)}\varepsilon_x(f),
\]
空原像对应空和0。\cref{g2:thm:sard}保证存在正则值。为什么在变形时要数正负号，而不只数点数？一条有向的一维逆像会把新生或消失的两个端点赋予相反符号。

\begin{proposition}\label{g2:prop:oriented-regular-fiber-cobordism}
	设 \(M^n\) 为紧定向无边界光滑流形，\(N^n\) 为定向无边界光滑流形，\(n\ge1\)。若光滑映射 \(f_0,f_1:M\to N\) 同伦，且 \(y\in N\) 同时为两者的正则值，则 \(d_y(f_0)=d_y(f_1)\)。
\end{proposition}
\begin{proof}
	\cref{g2:thm:transverse-homotopy-cobordism}给出光滑同伦 \(H:M\times[0,1]\to N\)，在两端的时间邻域分别等于 \(f_0,f_1\)，并横截于 \(\{y\}\)。因 \(M\) 紧，\(W=H^{-1}(y)\) 是紧一维带边界流形，边界为两个端点纤维。

	这里给 \(M\times[0,1]\) 取 \(M\) 的方向在前、时间方向在后的乘积定向。给 \(W\) 的取向规定如下：先取 \(\mathrm{T}_yN\) 的正向基，选它们经 \(\mathrm{d}H\) 的任意提升，再接一条 \(\mathrm{T}W\) 的正向基，所得应是环境的正向基。这就是商方向在前、核方向在后的约定；改变提升只添加核方向，不改变取向。

	在端点附近，\(H(x,t)=f_j(x)\)，故 \(\partial_t\) 张成核。提升目标正向基到 \(\mathrm{T}_xM\) 后，其相对 \(M\) 的符号为 \(\varepsilon_x(f_j)\)，所以 \(W\) 的正向切向量是 \(\varepsilon_x(f_j)\partial_t\)。按照外法向优先的边界定向，\(t=1\) 端点的符号是 \(\varepsilon_x(f_1)\)，\(t=0\) 端点的符号是 \(-\varepsilon_x(f_0)\)。第一卷第14章第4节的 Stokes 定理对 \(W\) 上常函数1给出
	\[
		0=\int_W\mathrm{d}1=\int_{\partial W}1
		=d_y(f_1)-d_y(f_0).
	\]
	因而两个带符号计数相等。
\end{proof}

\subsection{同调定义与积分定义的比较}
\label{g2:subsec:1601:02}

第3章对球面定义的度可以原样推广到连通闭定向流形。这里“闭”仍表示紧且无边界。\cref{g2:thm:fundamental-class}给出由定向确定的正生成元 \([M]\in H_n(M;\mathbb Z)\) 与 \([N]\in H_n(N;\mathbb Z)\)。对连续映射 \(f:M\to N\)，唯一满足
\begin{equation}\label{g2:eq:manifold-homological-degree}
	f_*[M]=(\deg f)[N]
\end{equation}
的整数称为 \(f\) 的映射度。这个定义已说明同伦不变性；还须证明正则值计数和积分确实产生这一整数。

\ProseParagraphBreak

比较积分时，一个容易被略过的环节是：在基本类上对积分上链求值，为什么等于通常的流形积分？第7章给出了形式与奇异上同调的积分比较，下面用坐标盒中的基本类补足这一识别，也同时准备非紧情形所需的支集版本。

\begin{lemma}\label{g2:lem:integration-fundamental-pairing}
	设 \(M^n\) 为定向无边界光滑流形，\(n\ge1\)，\(\omega\in\Omega_c^n(M)\)。以 \([M]_{\mathrm{lf}}\) 表示与该定向对应的局部有限基本类，则
	\[
		\bigl\langle\mathcal I_{c,M}[\omega],[M]_{\mathrm{lf}}\bigr\rangle
		=\int_M\omega.
	\]
	左端的求值先在包含 \(\operatorname{supp}\omega\) 的紧集上相对化；若 \(M\) 闭，它就是通常的上同调与基本类配对。
\end{lemma}
\begin{proof}
	\cref{g2:thm:de-rham-support}把紧支形式的积分解释为紧支上同调类。若支集包含于紧集 \(K\)，这个上链在完全落入 \(M\setminus K\) 的单形上为零，因此可对 \([M]_K\) 求值；\cref{g2:prop:lf-fundamental-class}说明 \([M]_{\mathrm{lf}}\) 的相对化恰为 \([M]_K\)。

	先设 \(\omega\) 的支集紧含于一张正向坐标盒的内部。在该图域内取一个内部包含支集的闭盒，将它按正向仿射单形剖分，所得有限光滑链 \(c\) 的边界避开支集。在支集的每一点，\(c\) 的局部同调值都是正生成元；由\cref{g2:lem:compact-localization}的唯一性，\(c\) 表示该支集上的基本类。于是积分上链的定义及有限可加性给出
	\[
		\langle I\omega,c\rangle=\sum_{\sigma\text{ 为盒中单形}}\int_\sigma\omega
		=\int_M\omega.
	\]
	对于一般紧支 \(\omega\)，用有限张较小正向坐标盒覆盖支集，并以从属光滑单位分解写成有限和 \(\omega=\sum_\alpha\omega_\alpha\)，每项都紧支于相应盒内。对每项应用刚才的计算，再用配对和积分的线性性相加即得结论。
\end{proof}

\begin{theorem}[映射度比较定理]\label{g2:thm:degree-comparison}
	设 \(M,N\) 为连通闭定向光滑 \(n\) 维流形，\(n\ge1\)，\(f:M\to N\) 光滑。则对每个正则值 \(y\in N\)，
	\begin{equation}\label{g2:eq:degree-regular-value}
		\deg f=\sum_{x\in f^{-1}(y)}\operatorname{sgn}\det(\mathrm{d}f_x).
	\end{equation}
	对每个光滑顶次形式 \(\omega\in\Omega^n(N)\)，还有
	\begin{equation}\label{g2:eq:degree-integral}
		\int_M f^*\omega=(\deg f)\int_N\omega.
	\end{equation}
	因此正则值计数与正则值的选择无关；任取 \(\int_N\eta=1\) 的顶次形式，积分 \(\int_Mf^*\eta\) 都等于上述同调度。
\end{theorem}
\begin{proof}
	紧源使 \(f\) 固有，故 \(F=f^{-1}(y)\) 有限。先比较基本类在这些点附近的像。切除给出
	\[
		H_n(M,M\setminus F;\mathbb Z)
		\cong\bigoplus_{x\in F}H_n(M,M\setminus\{x\};\mathbb Z),
	\]
	并把 \([M]\) 的限制送到各局部正生成元组成的向量。与第3章的局部度求和相同，空间对映射给出交换图
	\[
		\begin{tikzcd}[column sep=large,row sep=large]
			H_n(M;\mathbb Z) \arrow[r] \arrow[d,"f_*"'] & \displaystyle\bigoplus_{x\in F}H_n(M,M\setminus\{x\};\mathbb Z) \arrow[d,"\sum f_{*,x}"] \\
			H_n(N;\mathbb Z) \arrow[r] & H_n(N,N\setminus\{y\};\mathbb Z).
		\end{tikzcd}
	\]
	沿左、下两边，\([M]\) 变为 \((\deg f)[N]_y\)；沿上、右两边，\cref{g2:prop:regular-local-degree-sign}给出 \(\sum_{x\in F}\varepsilon_x(f)[N]_y\)。右侧每个局部箭头先经切除取到只含这一个原像点的小邻域，再由 \(f\) 的限制诱导。比较系数即得\eqref{g2:eq:degree-regular-value}。若 \(F\) 为空，上方相对群为零，同一论证给出度为零。

	所有顶次形式都闭。由\cref{g2:lem:integration-fundamental-pairing,g2:prop:integration-naturality}及上同调求值的自然性，
	\begin{align*}
		\int_Mf^*\omega
		&=\bigl\langle\mathcal I_M[f^*\omega],[M]\bigr\rangle
		=\bigl\langle f^*\mathcal I_N[\omega],[M]\bigr\rangle\\
		&=\bigl\langle\mathcal I_N[\omega],f_*[M]\bigr\rangle
		=(\deg f)\int_N\omega.
	\end{align*}
	最后，在一张正向坐标图内取非负、非零的紧支顶次形式，再除以其正积分，就得到 \(\int_N\eta=1\) 的形式。因此积分的倍数唯一，且已由基本类证明为整数。
\end{proof}

这个证明补全了第一卷第14章第5节预先陈述的次数积分公式。维数0时，连通闭流形只有一个点，基本类记录该点的定向符号，积分就是带符号求值，同一公式直接成立。积分表达适用于所有顶次形式，包括积分为零的形式；只有把度写成两个积分之比时，才要求分母非零。局部微分给出的正负号决定同调中的整数，而积分把这个整数读成拉回后的总量。正则值与积分的比较也见\cite[第457--459页，定理17.35及命题17.36]{Lee2012Smooth}。

\subsection{固有映射的映射度}
\label{g2:subsec:1601:03}

连通非紧流形的普通顶次同调为零，见\cref{g2:prop:noncompact-top-homology}。因此把\eqref{g2:eq:manifold-homological-degree}中的普通基本类照搬过去，会失去要计数的信息。第6章的局部有限链允许沿整个流形延伸，只要求每个紧集遇到有限项；其基本类仍在每一点具有正的局部值。

\begin{definition}\label{g2:def:proper-manifold-degree}
	设 \(M,N\) 为连通定向无边界光滑 \(n\) 维流形，\(n\ge1\)，\(f:M\to N\) 为固有连续映射。由
	\[
		f_*[M]_{\mathrm{lf}}=(\deg f)[N]_{\mathrm{lf}}
	\]
	确定的整数称为 \(f\) 的映射度。这里使用整数系数的局部有限同调。
\end{definition}

这个定义的两个依据需要同时保留。\cref{g2:thm:poincare-duality-locally-finite}在次数0给出 \(H_n^{\mathrm{lf}}(N;\mathbb Z)\cong H^0(N;\mathbb Z)=\mathbb Z\)，基本类对应1。另一方面，\cref{g2:prop:support-functoriality}中的固有性使推前有定义：目标紧集 \(K\) 的逆像紧，故原局部有限链只有有限条单形可能在推前后遇到 \(K\)。若 \(M,N\) 都闭，这一定义就是前面的普通同调度。

\begin{theorem}[固有映射的度公式]\label{g2:thm:proper-degree-formulas}
	设 \(f:M^n\to N^n\) 为连通定向无边界光滑流形间的固有光滑映射，\(n\ge1\)。对任一正则值 \(y\) 及任一紧支顶次形式 \(\omega\in\Omega_c^n(N)\)，有
	\[
		\deg f=\sum_{x\in f^{-1}(y)}\operatorname{sgn}\det(\mathrm{d}f_x),
		\qquad \int_Mf^*\omega=(\deg f)\int_N\omega.
	\]
	特别地，非零度强迫 \(f\) 满射。
\end{theorem}
\begin{proof}
	正则纤维紧且离散，因而有限。将局部有限基本类相对化到 \(f^{-1}(y)\) 与 \(\{y\}\)，得到与\cref{g2:thm:degree-comparison}相同的交换图；各局部值仍是正生成元，故得到符号求和公式。

	固有性保证 \(\operatorname{supp}f^*\omega\subset f^{-1}(\operatorname{supp}\omega)\) 紧。由\cref{g2:thm:de-rham-support}的固有拉回相容性，\(\mathcal I_{c,M}[f^*\omega]=f^*\mathcal I_{c,N}[\omega]\)。紧支上链与局部有限链的求值满足
	\[
		\langle f^*a,c\rangle=\langle a,f_*c\rangle;
	\]
	两边实际都只涉及支集附近的有限项，逐单形就是同一项。再用\cref{g2:lem:integration-fundamental-pairing}和度的定义，得到积分公式。若 \(y\) 不在像内，它自动是正则值，空和公式便给出 \(\deg f=0\)。
\end{proof}

同伦也须保留无穷远处的控制。设 \(H:M\times[0,1]\to N\) 是连续同伦；称它为固有同伦，是指这个整体映射固有。它使两个端点在局部有限同调上的推前相同。为验证这一点，对局部有限链逐项使用\cref{g2:thm:singular-homotopy-invariance}的棱柱算子。若某个输出棱柱遇到目标紧集 \(K\)，原单形必遇到紧集
\[
	\operatorname{pr}_M\bigl(H^{-1}(K)\bigr).
\]
这样的原单形只有有限条，每条棱柱又只剖成有限个单形，所以棱柱算子仍局部有限，原来的等式 \(\partial P+P\partial=H_{1\#}-H_{0\#}\) 继续成立。因此固有同伦保持映射度。

\begin{example}\label{g2:exm:proper-slices-not-proper-homotopy}
	对 \(0\le t\le1\) 令 \(H_t:\mathbb R\to\mathbb R\)、\(H_t(x)=x+t x^2\)。每个固定时刻的映射都固有：\(t=0\) 时为恒等，\(t>0\) 时二次项使函数在两端趋于正无穷。值0对每张切片都是正则值。\(t=0\) 时只有原像0、符号为正，故度为1；\(t>0\) 时原像为 \(0,-1/t\)，导数分别为 \(1,-1\)，故度为0。

	整体同伦并不固有，因为 \(H^{-1}(0)\) 包含 \((-1/t,t)\)，当 \(t\downarrow0\) 时这些点逃向无穷远。两条相反符号的分支中有一条从无穷远进入，正好说明“每张切片固有”不能代替“整个同伦固有”。
\end{example}

连映射本身都不固有时，正则值计数也可能随位置改变。例如 \(f:\mathbb R\to\mathbb R\)、\(f(x)=e^x\)，每个正数有一个符号为正的原像，每个负数没有原像。其像未覆盖零点附近的负半轴，而趋向0的原像跑向负无穷。积分虽对某些拉回仍收敛，也不能据此指定统一的度：\(\int_{\mathbb R}f^*(a(y)\,\mathrm{d}y)=\int_0^\infty a(y)\,\mathrm{d}y\)，随紧支函数 \(a\) 在正、负半轴间的分布变化，无法等于一个固定整数乘 \(\int_{\mathbb R}a(y)\,\mathrm{d}y\)。

\subsection{复合公式及典型计算}
\label{g2:subsec:1601:04}

\begin{proposition}\label{g2:prop:manifold-degree-laws}
	设 \(M,N,P\) 为连通定向无边界光滑 \(n\) 维流形，\(n\ge1\)，\(f:M\to N\)、\(g:N\to P\) 固有连续。则
	\[
		\deg(g\circ f)=(\deg g)(\deg f),\qquad\deg(\operatorname{id}_M)=1.
	\]
	保定向微分同胚的度为1，反定向微分同胚的度为 \(-1\)。反转源或目标其中一端的定向使度变号，同时反转两端则不变。
\end{proposition}
\begin{proof}
	固有映射的复合仍固有，局部有限推前满足复合律，所以
	\[
		(gf)_*[M]_{\mathrm{lf}}
		=g_*\bigl((\deg f)[N]_{\mathrm{lf}}\bigr)
		=(\deg f)(\deg g)[P]_{\mathrm{lf}}.
	\]
	恒等映射保持基本类。微分同胚是固有映射，每个目标点只有一个正则原像，符号求和公式给出 \(\pm1\)。定向反转使对应基本类乘以 \(-1\)，代入定义即得最后一项。
\end{proof}

若 \(p:M\to N\) 是连通定向流形间的 \(r\) 层光滑覆盖，且局部微分保持定向，则 \(\deg p=r\)；若处处反向，则 \(\deg p=-r\)。有限层覆盖在流形间固有：目标紧集可由有限个闭包位于均匀覆盖邻域的小坐标块覆盖，每块的全部原像是有限个紧块。局部符号在连通源上恒定，故上面两个结论直接由每个纤维的 \(r\) 项求和得到。圆周 \(z\mapsto z^k\) 的度 \(k\) 是这一计算的一维例子。

\begin{example}\label{g2:exm:torus-integer-matrix-degree}
	令 \(T^n=\mathbb R^n/\mathbb Z^n\)，取标准定向。对整数矩阵 \(A\in\operatorname{Mat}_{n\times n}(\mathbb Z)\)，映射
	\[
		f_A([x])=[Ax]
	\]
	良定义。形式 \(\eta=\mathrm{d}x^1\wedge\cdots\wedge\mathrm{d}x^n\) 在整数平移下不变，因而下降到环面，且 \(\int_{T^n}\eta=1\)。直接计算 \(f_A^*\eta=(\det A)\eta\)，所以
	\[
		\deg f_A=\det A.
	\]
	例如 \(A=\begin{pmatrix}2&1\\0&-3\end{pmatrix}\) 给出度 \(-6\)。若 \(\det A\ne0\)，每个点都是正则点，各原像点的符号均为 \(\operatorname{sgn}\det A\)，所以符号求和公式说明每个纤维有 \(|\det A|\) 个点；若 \(\det A=0\)，积分拉回为零，度也为零。行列式同时包含覆盖片数与方向，两者在这一模型中恰好合为一个整数。
\end{example}

\begin{example}\label{g2:exm:antipodal-differential-degree}
	给 \(S^n\subset\mathbb R^{n+1}\) 取外法向优先定向，对径映射 \(a(x)=-x\) 的度为 \((-1)^{n+1}\)，与第3章的同调计算相同。若 \((x,v_1,\ldots,v_n)\) 是环境中的正向基，映射后的外法向是 \(-x\)，切向量为 \(-v_i\)。整个 \(n+1\) 元组乘以 \(-1\)，故其符号为 \((-1)^{n+1}\)。这就是 \(\mathrm{d}a_x\) 相对于源、目标球面定向的符号。直接把它看作“\(n\) 个切向量都变负”会漏掉目标点的外法向也已改变这一事实。
\end{example}

度不等于未经加权的覆盖次数。例如 \(p(x)=x^3-x\) 是实线到实线的固有映射；在正则值0上有三个原像 \(-1,0,1\)，导数符号为 \(+,-,+\)，因此度为1。在充分大的正则值上只有一个原像，符号仍为正。临界值两侧原像数变化，代数和却保持不变。

% !TeX root = ../../Book2.tex
% 纲要标题：### 16.2 相交数的构造
% 纲要位置：几何学计划纲要.md，第 1632 行。
% 正文状态：初稿；尚未正式审核。

% 纲要内容开始
% 
% * 互补维子流形的线性取向
% * 横截扰动与零维交集
% * 一参数族给出的配边
% * 交数与帽积配对的比较
% 
% 纲要内容结束

\section{相交数的构造}
\label{g2:sec:1602}

环面上的一条经线和一条纬线相交一次；把其中一条弯曲，可能额外出现两个交点，却仍不能把两条曲线彻底分开。记录每次穿过的方向后，新出现的两个交点符号相反，计数因而保持不变。第6章用杯积与基本类定义了代数配对，本节将证明，这个配对正是几何中的带符号交点数。先明确局部符号，再用横截扰动与一维配边处理不横截的情形，最后在管状邻域中核对帽积。

\subsection{互补维子流形的线性取向}
\label{g2:subsec:1602:01}

设 \(M\) 是已定向、无边界的光滑 \(n\) 维流形，\(A^p,B^q\subset M\) 是已定向、无边界的嵌入子流形，\(p+q=n\)。在横截交点 \(z\)，两个切空间互补，因此把它们的基接在一起，就能与环境定向比较。

\begin{definition}\label{g2:def:local-intersection-sign}
	在上述条件下，取 \(\mathrm{T}_zA\) 的正向基 \((a_1,\ldots,a_p)\) 和 \(\mathrm{T}_zB\) 的正向基 \((b_1,\ldots,b_q)\)。若
	\[
		(a_1,\ldots,a_p,b_1,\ldots,b_q)
	\]
	是 \(\mathrm{T}_zM\) 的正向基，定义局部相交符号 \(\varepsilon_z(A,B)=1\)；否则定义为 \(-1\)。若 \(A\pitchfork B\) 且 \(A\cap B\) 紧，定义其\term[xiangjiaoshu]{相交数}[intersection number]为
	\[
		I_M(A,B)=\sum_{z\in A\cap B}\varepsilon_z(A,B).
	\]
\end{definition}

更换两片的正向基只会乘上两个正行列式，所以符号与这些选择无关。横截逆像定理给出 \(\dim(A\cap B)=p+q-n=0\)，紧的零维流形是有限点集，因而上式是有限和。若 \(A\) 紧且 \(B\) 在 \(M\) 中闭，交集自动紧。维数为零的定向子流形按带符号的点理解，取向比较也包含这个符号。

\ProseParagraphBreak

为了把交点计数用于一般映射，需要把两片的比较改写为一个线性同构的符号。给法丛 \(\nu_MB=\mathrm{T}M|_B/\mathrm{T}B\) 取定向，使
\begin{equation}\label{g2:eq:normal-first-orientation}
	(\nu_MB,\mathrm{T}B)=\mathrm{T}M.
\end{equation}
这里法方向排在切方向之前：把一个法向基任意提升到环境切空间，再接上 \(B\) 的正向基，所得必须是环境的正向基。提升的改变只增加切向分量，不改变行列式。这个约定称为法向优先。

\begin{proposition}\label{g2:prop:intersection-linear-signs}
	在互补维横截交点处，商投影诱导的同构
	\[
		\mathrm{T}_zA\longrightarrow\nu_zB
	\]
	的定向符号等于 \(\varepsilon_z(A,B)\)。交换两片、反转一片定向及反转环境定向分别给出
	\[
		\varepsilon_z(B,A)=(-1)^{pq}\varepsilon_z(A,B),\qquad
		\varepsilon_z(-A,B)=-\varepsilon_z(A,B),\qquad
		\varepsilon_z^{\,-M}(A,B)=-\varepsilon_z^{\,M}(A,B).
	\]
	反转 \(B\) 的定向同样使符号反转。
\end{proposition}
\begin{proof}
	横截性及维数相等保证所列商映射为同构。相对于一个法向优先的环境基，把 \(A\) 的基分解为法向分量与切向分量；接上 \(B\) 的基后，矩阵为分块三角形，下方切向分量不影响行列式，符号就是法向分量矩阵的符号。交换两组基需要作 \(pq\) 次向量交换，得到 \((-1)^{pq}\)。反转任一基的定向或反转作为比较对象的环境定向，均改变一次符号。
\end{proof}

若 \(X^p\) 是已定向流形，\(f:X\to M\) 横截于 \(B\)，则在 \(x\in f^{-1}(B)\) 处定义
\[
	\varepsilon_x(f,B)=\operatorname{sgn}\bigl(\mathrm{T}_xX\xrightarrow{\mathrm{d}f_x}\mathrm{T}_{f(x)}M\longrightarrow\nu_{f(x)}B\bigr).
\]
当逆像紧时，令 \(I(f,B)=\sum_x\varepsilon_x(f,B)\)。这按源中的原像计数；不同源点即使映到同一个环境点，也分别贡献。嵌入 \(f=i_A\) 的情形正好恢复前面的定义。定义整数符号所需的是 \(X\) 的定向与 \(B\) 法丛的定向；环境及 \(B\) 的定向通过\eqref{g2:eq:normal-first-orientation}提供后者。若不选择这些定向，仍可取交点数模 \(2\)。

\begin{example}\label{g2:exm:torus-basic-intersection}
	在标准定向 \(\mathrm{d}x\wedge\mathrm{d}y\) 的环面 \(\mathbb R^2/\mathbb Z^2\) 中，令 \(a(t)=(t,0)\)、\(b(t)=(0,t)\)，参数均模 \(1\)。它们仅在原点相交，基 \((\partial_x,\partial_y)\) 正向，所以 \(I(a,b)=1\)、\(I(b,a)=-1\)。\(b\) 的正向法线是 \(\partial_x\)，而 \(a\) 的正向法线是 \(-\partial_y\)：后一个负号来自法向必须排在 \(\partial_x\) 之前。
\end{example}

\subsection{横截扰动与零维交集}
\label{g2:subsec:1602:02}

两片可能相切，甚至整段重合；此时直接数共同点不能得到一个稳定的数。横截性允许我们将其中一片略微移动。若用任意映射扰动，还需确认所得映射仍为嵌入。下面直接用环境微分同胚作扰动，整个过程都保持嵌入性质。

\begin{proposition}\label{g2:prop:transverse-embedded-perturbation}
	设 \(M\) 为无边界光滑流形，\(A\subset M\) 为紧无边界嵌入子流形，\(B\subset M\) 为闭的无边界嵌入子流形。对 \(A\) 的任意给定开邻域 \(\Omega\)，存在支撑于 \(\Omega\) 的、任意小的光滑环境同位 \(\Phi_t\)，使 \(\Phi_1(A)\pitchfork B\)。若两片维数互补，终点交集为有限点集。
\end{proposition}
\begin{proof}
	用有限张相对紧坐标图覆盖 \(A\)，其闭包包含于 \(\Omega\)。将各图的坐标向量场乘以紧支截断，再在图外补零，得到有限个紧支光滑场 \(V_1,\ldots,V_N\)，其值在 \(A\) 的每一点张成环境切空间。由\cref{g2:lem:uniformly-supported-evolution}，这些场的短时流 \(\varphi_j^s\) 都是环境微分同胚，并在同一个紧集外恒等。置
	\[
		\Psi(a,z)=\varphi_N^{a_N}\circ\cdots\circ\varphi_1^{a_1}(z),
		\qquad F(a,x)=\Psi(a,x),\quad x\in A.
	\]
	在 \(a=0\) 处，参数导数的像包含所有 \(V_j(x)\)，所以 \(F\) 的微分满射。\(A\) 紧，故可统一缩小参数球，使 \(F\) 在这个球乘 \(A\) 上均为淹没。\cref{g2:thm:parametric-transversality}于是给出任意小参数 \(a\)，使 \(F_a\pitchfork B\)。

	令 \(\Phi_t(z)=\Psi(ta,z)\)。每个时刻都是微分同胚，\(\Phi_0=\operatorname{id}\)，所以限制到 \(A\) 始终为嵌入。参数趋零时，这族映射及其各阶导数在共同紧支集上趋于恒等族，给出所需的小扰动。最后，\(\Phi_1(A)\) 紧，\(B\) 闭，故零维交集紧而有限。
\end{proof}

定向随 \(\Phi_t\) 推送给移动后的 \(A\)。由于环境微分同胚从恒等映射连续变化，它始终保持环境定向。对一般的 \(f:X\to M\)，其中 \(X\) 紧，\cref{g2:thm:relative-transversality}也给出与 \(f\) 同伦的横截扰动。下一小节将证明不同扰动所得的带符号和相同，之后便可将这个和定义为未必横截的 \(A,B\) 或 \(f,B\) 的相交数。

\begin{example}\label{g2:exm:intersection-birth-signs}
	在定向平面中，令 \(B\) 是沿正 \(x\) 方向的横轴，\(A_s\) 是按递增 \(x\) 定向的曲线 \(y=x^2-s\)。\(s<0\) 时局部没有交点；\(s>0\) 时交点为 \(x=\pm\sqrt{s}\)。两片切向量给出
	\[
		\det\begin{pmatrix}1&1\\2x&0\end{pmatrix}=-2x,
	\]
	故左交点的符号为 \(+1\)，右交点为 \(-1\)。两个新交点的总贡献为零。\(s=0\) 的相切切片没有局部横截符号，但整个参数族的交点轨迹 \(s=x^2\) 仍是一条光滑曲线。
\end{example}

\begin{figure}[htbp]
	\centering
	\input{Figures/Chapter16-IntersectionPair}
	\caption{相切附近的交点生成。曲线按横坐标递增定向，横轴向右定向；出现的两个交点分别贡献 \(+1\) 和 \(-1\)。}
	\label{g2:fig:intersection-pair}
\end{figure}

\subsection{一参数族给出的配边}
\label{g2:subsec:1602:03}

第14章已经将同伦横截化，得到连接端点逆像的紧流形。现在给这个配边定向，就能把模二不变性加强为整数不变性。这里需要核对时间方向与法方向的排列，而不只是指出端点成对出现。

\begin{theorem}\label{g2:thm:intersection-homotopy-invariance}
	设 \(M\) 是已定向、无边界的光滑 \(n\) 维流形，\(B\subset M\) 是已定向、闭的无边界嵌入 \(q\) 维子流形，\(X\) 是紧、已定向、无边界的光滑 \(p\) 维流形，\(p+q=n\)。若光滑映射 \(f_0,f_1:X\to M\) 横截于 \(B\) 且连续同伦，则
	\[
		I(f_0,B)=I(f_1,B).
	\]
\end{theorem}
\begin{proof}
	由\cref{g2:thm:transverse-homotopy-cobordism}，可选横截同伦 \(H:X\times[0,1]\to M\)，在两个端点的时间邻域中分别等于 \(f_0,f_1\)。其逆像 \(W=H^{-1}(B)\) 是紧一维带边界流形，边界是两端逆像的不交并。这里 \(X\) 紧、\(B\) 闭保证 \(W\) 紧。

	给 \(X\times[0,1]\) 取先 \(\mathrm{T}X\)、后 \(\partial_t\) 的乘积定向。商微分给出短正合列
	\[
		0\longrightarrow\mathrm{T}W\longrightarrow\mathrm{T}(X\times[0,1])|_W
		\xrightarrow{\ \overline{\mathrm{d}H}\ }H^*\nu B\longrightarrow0.
	\]
	规定 \(W\) 的正向，使法方向排在 \(W\) 的切方向之前，即 \((H^*\nu B,\mathrm{T}W)\) 给出源乘积的定向。更具体地，先列法基的提升，最后列 \(W\) 的正向切向量。

	在端点邻域，\(H\) 与时间无关。若 \(x\in f_i^{-1}(B)\)，记其局部符号为 \(\varepsilon_x\)。法基在 \(\mathrm{T}_xX\) 中的提升相对于 \(X\) 定向具有符号 \(\varepsilon_x\)，故上述规则使 \(W\) 的正向切向量为 \(\varepsilon_x\partial_t\)。采用外法向优先的边界定向，\(t=1\) 端点的符号是 \(\varepsilon_x\)，\(t=0\) 端点的符号是 \(-\varepsilon_x\)。因此
	\[
		\sum_{z\in\partial W}\varepsilon_z=I(f_1,B)-I(f_0,B).
	\]
	紧一维有向流形的边界符号之和为零：按第14章证明中的有限坐标弧分割，把各小弧沿定向排列，每条小弧贡献“终点减起点”，内部端点两两抵消，剩下的就是边界端点。于是右端为零。
\end{proof}

由此可对任意光滑 \(f:X\to M\) 选一个横截扰动 \(f'\)，定义 \(I(f,B)=I(f',B)\)。两个这样的扰动均与 \(f\) 同伦，定理保证良定义。对紧嵌入 \(A\) 使用上一小节的环境扰动，便定义 \(I_M(A,B)\)。其交换及反向规则仍成立：先把两片置于横截位置，再应用\cref{g2:prop:intersection-linear-signs}；更换所移动的一片也不影响结果，因为保持定向的环境微分同胚满足 \(I(A,\Phi(B))=I(\Phi^{-1}(A),B)\)，而 \(\Phi^{-1}|_A\) 与包含映射同伦。

\ProseParagraphBreak

紧性可以用“整个交点轨迹紧”替代，但不能只要求每张切片有有限个交点。例如 \(f_t:\mathbb R\to\mathbb R\)、\(f_t(x)=tx-1\)，以正向原点作为 \(B\)。\(t=0\) 时无交点，\(t>0\) 时恰有一个符号为正的交点 \(x=1/t\)。当 \(t\downarrow0\) 时它逃向无穷，轨迹不是紧配边，所以整数计数发生变化。这个例子说明，在非紧源上应用配边证明时，必须控制整条同伦的逆像。

\subsection{交数与帽积配对的比较}
\label{g2:subsec:1602:04}

以下设环境 \(M\) 闭且已定向，以便使用整数基本类 \([M]\)。\cref{g2:cor:closed-poincare-duality}给出同构
\[
	D_M:H^k(M;\mathbb Z)\longrightarrow H_{n-k}(M;\mathbb Z),\qquad
	D_M(\alpha)=\alpha\frown[M].
\]
对闭已定向嵌入子流形 \(i_B:B^q\hookrightarrow M\)，记
\[
	\mathrm{PD}_M(B)=D_M^{-1}\bigl((i_B)_*[B]\bigr)\in H^{n-q}(M;\mathbb Z).
\]
这里的 \(\mathrm{PD}\) 是 Poincaré 对偶类；它由帽积定义。要证明它检测穿过 \(B\) 的次数，需在 \(B\) 附近找到一个具体相对类，并核对其帽积正好留下 \(B\) 的基本类。

\begin{lemma}\label{g2:lem:normal-relative-class}
	设 \(M\) 是闭、已定向的光滑 \(n\) 维流形，\(B\subset M\) 是闭、已定向的嵌入 \(q\) 维子流形，\(p=n-q\)。法丛采用\eqref{g2:eq:normal-first-orientation}的定向。存在唯一的类
	\[
		u_B\in H^p(M,M\setminus B;\mathbb Z),
	\]
	使它在每个小法向 \(p\) 维盘上的限制，在正向局部基本类上取值 \(1\)。若 \(\overline u_B\) 是它在绝对上同调中的像，则
	\begin{equation}\label{g2:eq:normal-class-cap-fundamental}
		\overline u_B\frown[M]=(i_B)_*[B],\qquad
		\overline u_B=\mathrm{PD}_M(B).
	\end{equation}
\end{lemma}
\begin{proof}
	由\cref{g2:thm:tubular-neighborhood}，可将 \(B\) 的一个开邻域 \(U\) 识别为法丛的一个开圆盘丛；沿纤维径向重参数化，它与整个法丛 \(E\) 微分同胚，并保持零截面及纤维定向。\cref{g2:prop:relative-cohomology-exactness}的切除把 \(H^p(M,M\setminus B)\) 识别为 \(H^p(E,E\setminus B)\)。因此先在法丛中构造相对类。

	对任一法丛可平凡化的开集 \(V\subset B\)，选保持法向定向的平凡化，得到空间对
	\[
		(E|_V,E|_V\setminus V)\cong
		(\mathbb R^p\times V,(\mathbb R^p\setminus\{0\})\times V).
	\]
	相对 Eilenberg--Zilber 比较给出
	\begin{equation}\label{g2:eq:normal-low-degree-local}
		H^j(E|_V,E|_V\setminus V;\mathbb Z)\cong H^{j-p}(V;\mathbb Z).
	\end{equation}
	恢复这一计算的理由：\cref{g2:prop:cross-cup-evaluation}证明中的相对链比较将左侧链模型写成 \(C_*(\mathbb R^p,\mathbb R^p\setminus\{0\})\otimes C_*(V)\)。第一因子的同调只在 \(p\) 次为 \(\mathbb Z\)；由\cref{g2:lem:free-complex-splitting}，其圈边界序列可分裂为一个 \(\mathbb Z\) 加若干两项可缩复形，故它链同伦等价于仅在 \(p\) 次有 \(\mathbb Z\) 的复形。张量后再取 \(\operatorname{Hom}(-,\mathbb Z)\)，便得到上同调的次数平移。特别地，\(j<p\) 时左侧为零，\(j=p\) 时它由每个分支上的纤维取值唯一确定；全部取值为 \(1\) 的类就是正向纤维生成元与 \(1\in H^0(V)\) 的外积。

	这些局部类可以唯一粘合。取覆盖紧流形 \(B\) 的有限个丛平凡化开集，逐个取并。对两个开集 \(V,W\)，相对 Mayer--Vietoris 列包含
	\[
		H^{p-1}(E|_{V\cap W},E|_{V\cap W}\setminus(V\cap W))
		\longrightarrow H^p(E|_{V\cup W},E|_{V\cup W}\setminus(V\cup W))
	\]
	及其后的“限制到 \(V,W\)，再取交集上的差”两条箭头。这条列由相对小链模型的短正合列取对偶得到，和第6章使用的相对 Mayer--Vietoris 构造相同。归纳时，交集位于新加入的平凡化开集中，故\eqref{g2:eq:normal-low-degree-local}适用。其 \(p-1\) 次群为零，且两侧的 \(p\) 次局部类在交集上都以每根正向纤维的取值 \(1\) 为特征，所以差为零。正合性给出延拓的存在与唯一性。相同正合列也保持低于 \(p\) 次的消失，完成归纳。于是得到 \(u_B\)。改变管状邻域或平凡化后，纤维的局部取值不变，唯一性保证所取的是同一相对类。

	接着计算帽积。把 \(u_B\) 切除到 \(U\)，记作 \(u\)。\cref{g2:prop:relative-cap}给出
	\[
		z_U=u\frown[U]_B\in H_q(U;\mathbb Z),
	\]
	其中 \([U]_B\in H_n(U,U\setminus B;\mathbb Z)\) 是第6章的紧集基本类。令 \(r:U\to B\) 为管状收缩。要证明 \(r_*z_U=[B]\)，依\cref{g2:lem:compact-localization}，只需在每个 \(b\in B\) 核对其顶次局部值。

	在 \(b\) 附近选法向优先的乘积坐标 \((z,y)\in\mathbb R^p\times\mathbb R^q\)，其中 \(B\) 为 \(z=0\)。令 \(v_b\in H^q(B,B\setminus\{b\};\mathbb Z)\) 为在 \(B\) 的正向局部基本类上取值 \(1\) 的类。相对帽积的自然性及\eqref{g2:eq:cap-evaluation}在小链模型中的同一计算给出
	\begin{equation}\label{g2:eq:normal-cap-local-test}
		\left\langle v_b,\operatorname{loc}_b(r_*z_U)\right\rangle
		=\left\langle u\smile r^*v_b,[U]_{\{b\}}\right\rangle=1.
	\end{equation}
	说明这里的支集对应：\(u\) 的相对支集为 \(B\)，\(r^*v_b\) 的相对支集为法纤维 \(r^{-1}(b)\)，两者相交于 \(b\)。所以其相对杯积在 \(H^n(U,U\setminus\{b\})\) 中；局部化的基本类正是 \([U]_{\{b\}}\)。最后一个等号是\cref{g2:prop:local-cross-normalization}：先法向、后基底的两个正向局部生成元之积，在环境正向基本类上取值 \(1\)。某个因子为零维时，该计算退为点上乘法，结论相同。

	因此 \(r_*z_U\) 在每点的局部值都与 \([B]\) 相同，基本类的唯一性给出等式。将紧支相对类映入 \(M\)，帽积自然性又给出
	\[
		\overline u_B\frown[M]=(i_U)_*z_U
		=(i_B)_*r_*z_U=(i_B)_*[B].
	\]
	中间等式使用管状邻域沿纤维缩放到零截面的形变收缩。最后由 \(D_M\) 的单射性，\(\overline u_B=\mathrm{PD}_M(B)\)。
\end{proof}

这个相对类把“穿过 \(B\)”变成一个可以拉回的代数条件。法向小盘上的归一化决定每次横穿的贡献，上一证明中的帽积计算则保证它与第6章已经采用的对偶约定一致。

\begin{theorem}[几何相交与对偶配对]\label{g2:thm:geometric-intersection-pairing}
	设 \(M\) 是闭、已定向的光滑 \(n\) 维流形，\(B\subset M\) 是闭、已定向的嵌入 \(q\) 维子流形，\(X\) 是闭、已定向的光滑 \(p\) 维流形，\(p+q=n\)。若光滑映射 \(f:X\to M\) 横截于 \(B\)，且法丛采用法向优先定向，则
	\[
		\sum_{x\in f^{-1}(B)}\varepsilon_x(f,B)
		=\left\langle f^*\mathrm{PD}_M(B),[X]\right\rangle.
	\]
	对任意光滑 \(f\)，左侧采用横截扰动所定义的 \(I(f,B)\) 后，等式仍成立。特别地，若 \(A^p\subset M\) 也是闭已定向嵌入子流形，记 \([A]\)、\([B]\) 为其基本类在 \(M\) 中的推前，则
	\begin{equation}\label{g2:eq:intersection-cup-pairing}
		I_M(A,B)=\langle\mathrm{PD}_M(B),[A]\rangle
		=\langle\mathrm{PD}_M(A)\smile\mathrm{PD}_M(B),[M]\rangle.
	\end{equation}
\end{theorem}
\begin{proof}
	令 \(Z=f^{-1}(B)\)。它是有限点集。上一引理的相对类拉回为
	\[
		f^*u_B\in H^p(X,X\setminus Z;\mathbb Z).
	\]
	基本类的局部化与相对配对自然性给出
	\[
		\left\langle f^*\mathrm{PD}_M(B),[X]\right\rangle
		=\langle f^*u_B,[X]_Z\rangle
		=\sum_{x\in Z}\left\langle(f^*u_B)_x,[X]_x\right\rangle.
	\]
	最后一步在 \(Z\) 的互不相交小坐标邻域中切除；基本类限制到每点都为该点的正向局部类，所以有限直和中的每项系数都是 \(1\)。

	对 \(x\in Z\)，取 \(f(x)\) 附近的法向坐标 \(z:M\supset U\to\mathbb R^p\)。法向部分 \(z\circ f\) 在 \(x\) 的微分是定义 \(\varepsilon_x(f,B)\) 的线性同构。由逆函数定理，它在小邻域内为局部微分同胚，按其行列式符号将 \([X]_x\) 送到正向法盘基本类或其相反数。\(u_B\) 在这个正向法盘上取值 \(1\)，故每项恰等于 \(\varepsilon_x(f,B)\)。这证明横截情形。

	一般 \(f\) 的横截扰动 \(f'\) 与 \(f\) 同伦，所以它们诱导相同的上同调拉回；左侧按此前的不变性定义，等式随之成立。取 \(f=i_A\)，得到\eqref{g2:eq:intersection-cup-pairing}的第一个等式。第二个等式由帽积顺序决定：
	\[
		\begin{aligned}
		\langle\mathrm{PD}_M(B),[A]\rangle
		&=\langle\mathrm{PD}_M(B),\mathrm{PD}_M(A)\frown[M]\rangle\\
		&=\langle\mathrm{PD}_M(A)\smile\mathrm{PD}_M(B),[M]\rangle,
		\end{aligned}
	\]
	最后一行正是\eqref{g2:eq:cap-evaluation}。因此此处杯积先放 \(A\) 的对偶类，再放 \(B\) 的对偶类。
\end{proof}

\begin{corollary}\label{g2:cor:intersection-homological-invariance}
	闭已定向光滑 \(n\) 维流形中的互补维闭已定向嵌入子流形，其相交数只依赖整数同调类。它是这些类的双线性配对；若任一类为挠类，则相交数为零。交换两类乘以 \((-1)^{pq}\)，反转任一子流形或环境的定向均使相交数反号。
\end{corollary}
\begin{proof}
	\(D_M^{-1}\)、杯积及基本类取值均线性，所以\eqref{g2:eq:intersection-cup-pairing}只依赖同调类且双线性。若 \(r[A]=0\)，则 \(rI_M(A,B)=0\)，而整数群无挠，故 \(I_M(A,B)=0\)，另一因子同理。对偶类次数分别为 \(q,p\)，分次交换律给出 \((-1)^{pq}\)。子流形反向使其基本类及对偶类反号；环境反向使 \([M]\) 与 \(D_M\) 都反号，因此两个对偶类各反号，最终的基本类取值再反号，总变化仍为负号。
\end{proof}

\begin{example}\label{g2:exm:torus-poincare-sign-check}
	回到环面中的 \(a,b\)。令 \(\xi,\eta\in H^1(T^2;\mathbb Z)\) 分别从第一、第二圆周拉回正向生成元，则 \(\langle\xi\smile\eta,[T^2]\rangle=1\)。由帽积取值公式，\(D(\xi)=[b]\)、\(D(\eta)=-[a]\)，故
	\[
		\mathrm{PD}(a)=-\eta,\qquad \mathrm{PD}(b)=\xi,\qquad
		\langle(-\eta)\smile\xi,[T^2]\rangle=1.
	\]
	这与两条正向曲线的实际交点符号相同，也说明法向优先约定中的负号如何与杯积交换符号配合。
\end{example}

式\eqref{g2:eq:intersection-cup-pairing}也对任意互补维同调类定义代数相交配对；当类已有光滑子流形代表时，定理将它解释为横截交点数。以 \(\mathbb F_2\) 为系数时，同样的局部类、粘合与帽积证明给出模二相交配对：局部生成元没有正负之分，因而不需要环境及两片的整数定向。

% !TeX root = ../../Book2.tex
% 纲要标题：### 16.3 向量场零点指标
% 纲要位置：几何学计划纲要.md，第 1589 行。
% 正文状态：未撰写；本文件目前是框架。

% 纲要内容开始
% 
% * 局部球面映射度
% * 非退化零点的行列式符号
% * 局部扰动与指标可加性
% * Euler 类的后续解释
% 
% 纲要内容结束

\section{向量场零点指标}
\label{g2:sec:1603}

待撰写

\subsection{局部球面映射度}
\label{g2:subsec:1603:01}
待撰写

\subsection{非退化零点的行列式符号}
\label{g2:subsec:1603:02}
待撰写

\subsection{局部扰动与指标可加性}
\label{g2:subsec:1603:03}
待撰写

\subsection{Euler 类的后续解释}
\label{g2:subsec:1603:04}
待撰写
% 纲要条目排为编号小节标题，并进入目录与书签；说明段仍保留于纲要注释中。

% !TeX root = ../../Book2.tex
% 纲要标题：### 16.4 不动点与 Lefschetz 公式
% 纲要位置：几何学计划纲要.md，第 1646 行。
% 正文状态：未撰写；本文件目前是框架。

% 纲要内容开始
% 
% * 图像和对角线
% * 横截不动点的局部指数
% * 对角类和上同调迹
% * Lefschetz 数非零的存在推论
% 
% ---
% 
% 纲要内容结束

\section{不动点与 Lefschetz 公式}
\label{g2:sec:1604}

待撰写

\subsection{图像和对角线}
\label{g2:subsec:1604:01}
待撰写

\subsection{横截不动点的局部指数}
\label{g2:subsec:1604:02}
待撰写

\subsection{对角类和上同调迹}
\label{g2:subsec:1604:03}
待撰写

\subsection{Lefschetz 数非零的存在推论}
\label{g2:subsec:1604:04}
待撰写
% 纲要条目排为编号小节标题，并进入目录与书签；说明段仍保留于纲要注释中。


\begin{chaptersummary}
	映射度将每个正则原像的微分符号相加，所得整数既是推前基本类的倍数，也是拉回顶次形式的积分倍数。非紧源上，局部有限基本类和紧支形式承担相同的比较，而固有映射与整体固有同伦保证支集和无穷远处的控制。相交数用法向优先定向比较源切空间与目标法空间；横截扰动给出有限计数，有向一维配边保证它不依赖扰动，局部法向类进一步把它识别为帽积配对。

	\ProseParagraphBreak

	向量场的孤立零点由小球边界的方向度赋予指标，非退化时为线性化行列式的符号。保持边界无零的扰动只能重新分配内部指标，不能改变总和；闭流形上的截面同伦也保持总指标，即使流形不可定向。对于不动点，图像与对角线把同一局部度变成 \(\operatorname{sgn}\det(I-\mathrm{d}f_x)\)，对角类的杯积计算再把局部指数和化为 Lefschetz 数。非零总量迫使相应原像、交点或不动点存在，而总量为零仍可能有彼此抵消的局部贡献。
\end{chaptersummary}

\BookExercises

\begin{exercise}\label{g2:ex:degree-sublinear-perturbation}
	设 \(A\in\operatorname{GL}(n,\mathbb R)\)，\(r:\mathbb R^n\to\mathbb R^n\) 光滑，且 \(|r(x)|/|x|\to0\) 当 \(|x|\to\infty\)。证明 \(F(x)=Ax+r(x)\) 固有，计算其度，并证明 \(F\) 满射。若目标正则值的全部原像点处均有 \(\det\mathrm{D}F>0\)，这些原像点的数量有何限制？
\end{exercise}
\begin{solution}
	选 \(c>0\) 使 \(|Ax|\ge c|x|\)。在某个大球之外有 \(|r(x)|\le c|x|/2\)，所以同伦 \(H(x,t)=Ax+t r(x)\) 对所有 \(0\le t\le1\) 同时满足
	\[
		|H(x,t)|\ge c|x|/2.
	\]
	目标紧集有界，因而其在 \(H\) 下的原像包含于一个固定闭球与 \([0,1]\) 的乘积；该原像又闭，所以紧。因此 \(H\) 是整体固有同伦，特别 \(F\) 固有，且
	\[
		\deg F=\deg A=\operatorname{sgn}\det A\ne0.
	\]
	非零度给出满射性。若某正则值的全部原像符号为正，符号和就是原像点数。因原像非空，该数必须为1，同时迫使 \(\det A>0\)。若 \(\det A<0\)，任何正则值的原像都必须至少包含一个反向点。
\end{solution}

\begin{exercise}\label{g2:ex:proper-local-diffeomorphism-degree}
	设 \(f:M^n\to N^n\) 为连通定向无边界光滑流形间的固有光滑映射，\(n\ge1\)，且 \(\mathrm{d}f\) 处处可逆。证明 \(f\) 满射，所有纤维有同一个有限正基数 \(r\)，并有 \(|\deg f|=r\)。若某个纤维只有一个点，证明 \(f\) 为微分同胚。将此结论用于上一题：若 \(\det\mathrm{D}F\) 处处非零，可以得到什么加强？
\end{exercise}
\begin{solution}
	微分的定向符号是 \(M\) 上连续取值于 \(\{1,-1\}\) 的函数，因连通性而恒等于某个 \(\varepsilon\)。任取 \(x_0\in M\)，\(y_0=f(x_0)\) 是正则值，紧离散纤维有有限个 \(r_0\ge1\) 点。因此 \(\deg f=\varepsilon r_0\ne0\)，从而 \(f\) 满射。所有目标点都是正则值，再次使用求和公式，得到
	\[
		\deg f=\varepsilon\#f^{-1}(y),
	\]
	所以每个纤维都有 \(r_0=|\deg f|\) 个点。若某个纤维单点，则全部纤维单点，\(f\) 双射。它的局部逆由逆函数定理光滑；各局部逆在交叠处因单射性一致，故拼成光滑整体逆。

	上一题已给出 \(|\deg F|=1\)，因此若 \(\mathrm{D}F\) 处处可逆，\(F:\mathbb R^n\to\mathbb R^n\) 就是整体微分同胚。这里无穷远处的增长控制与局部微分可逆分别保证固有性与局部逆，两者合起来才得到整体逆。
\end{solution}

\begin{exercise}\label{g2:ex:real-polynomial-degree}
	设 \(p(x)=a_mx^m+\cdots+a_0\) 是非常数实多项式，\(a_m\ne0\)。把实线取标准定向，按 \(m\) 的奇偶性及 \(a_m\) 的符号计算 \(\deg p\)。由此证明：对任意正则值 \(y\)，方程 \(p(x)=y\) 的实根数与 \(m\) 模2同余。证明中须核对所用同伦的整体固有性。
\end{exercise}
\begin{solution}
	令 \(q(x)=p(x)-a_mx^m\)，取同伦 \(H(x,t)=a_mx^m+tq(x)\)。因 \(\deg q<m\)，充分大 \(|x|\) 时有 \(|q(x)|\le|a_m||x|^m/2\)，故
	\[
		|H(x,t)|\ge |a_m||x|^m/2
	\]
	对所有 \(t\in[0,1]\) 同时成立。与上一题同理，\(H\) 整体固有，因而 \(\deg p=\deg(x\mapsto a_mx^m)\)。

	若 \(m\) 为奇数，任取非零目标值，有唯一实原像，导数 \(ma_mx^{m-1}\) 的符号为 \(\operatorname{sgn}a_m\)，故度为该符号。若 \(m\) 为偶数，取与 \(a_m\) 异号的目标值，原像为空，故度为0。等价地，同号非零目标值的两个原像导数符号相反，仍给出0。

	正则值的所有实原像都是单根，且有限。模2后每个 \(\pm1\) 都等于1，所以
	\[
		\#p^{-1}(y)\equiv\deg p\equiv m\pmod2.
	\]
	这既给出奇次多项式至少有一个实根，也说明偶次情形的正则实根总数为偶数。
\end{solution}

\begin{exercise}\label{g2:ex:boundary-degree-balance}
	设 \(W^{n+1}\) 为紧定向光滑带边界流形，\(n\ge1\)，边界的连通分支为 \(M_1,\ldots,M_s\)，均取外法向优先定向。设 \(N^n\) 为连通闭定向光滑流形，\(F:W\to N\) 光滑，记 \(f_j=F|_{M_j}\)。证明
	\[
		\sum_{j=1}^s\deg f_j=0.
	\]
	把结论用于圆环 \(W=\{z\in\mathbb C:1\le|z|\le2\}\)：若把内、外两个边界圆都按逆时针方向参数化，它们到 \(S^1\) 的限制映射度是什么关系？
\end{exercise}
\begin{solution}
	取 \(N\) 上积分为1的光滑顶次形式 \(\eta\)。因为 \(\mathrm{d}\eta=0\)，第一卷的 Stokes 定理和\cref{g2:thm:degree-comparison}给出
	\[
		\sum_j\deg f_j
		=\sum_j\int_{M_j}f_j^*\eta
		=\int_{\partial W}F^*\eta
		=\int_W\mathrm{d}(F^*\eta)=0.
	\]
	关键是所有边界分支使用诱导定向。对平面标准定向的圆环，外边界的诱导方向为逆时针，内边界为顺时针。因此若两个参数圆都取逆时针方向，内边界度要在上述和式中另乘 \(-1\)。公式成为 \(\deg f_{\mathrm{out}}-\deg f_{\mathrm{in}}=0\)，即二者的度相等。单边界情形的延拓障碍可对照\cite[第459页，定理17.38]{Lee2012Smooth}。
\end{solution}

% 第16.2节习题候选，题目及解答供并入章末。
\begin{exercise}\label{g2:ex:torus-slope-intersection}
	给 \(T^2=\mathbb R^2/\mathbb Z^2\) 取标准定向。设 \((r,s),(u,v)\in\mathbb Z^2\) 均为非零本原向量，令 \(A(t)=(rt,st)\)、\(B(t)=(ut,vt)\)，参数 \(t\) 模 \(1\)，并以递增参数定向。证明这两条曲线为嵌入圆周，且 \(I(A,B)=rv-su\)。当 \(rv-su\ne0\) 时求实际交点数；当 \(rv-su=0\) 时给出不交的微小扰动。
\end{exercise}
\begin{solution}
	若 \(A(t)=A(t')\)，则 \(r(t-t'),s(t-t')\in\mathbb Z\)。由 Bézout 等式及 \(\gcd(r,s)=1\)，得 \(t-t'\in\mathbb Z\)，故 \(A\) 单射；其导数不为零，源紧，所以为嵌入，\(B\) 同理。记标准坐标圆周类为 \(a,b\)，则 \([A]=ra+sb\)、\([B]=ua+vb\)。正文已算出 \(I(a,b)=1\)、\(I(b,a)=-1\)；平移一个坐标圆周得到与原圆周不交的同调代表，所以 \(I(a,a)=I(b,b)=0\)。双线性给出 \(I(A,B)=rv-su\)。

	还可直接数交点。取整数矩阵 \(P\in\operatorname{SL}(2,\mathbb Z)\)，使 \(P(r,s)^{\mathsf T}=(1,0)^{\mathsf T}\)；它由 Bézout 等式构造，并诱导保持定向的环面微分同胚。记 \(P(u,v)^{\mathsf T}=(u',v')^{\mathsf T}\)，则 \(v'=rv-su\)。\(B\) 与水平圆周相交恰要求 \(v't\in\mathbb Z\)。若 \(v'\ne0\)，模 \(1\) 恰有 \(|v'|\) 个参数；\(B\) 嵌入，给出相同数量的不同交点。各点的有序切向量行列式均为 \(v'\)，所以全部具有符号 \(\operatorname{sgn}(v')\)。若 \(v'=0\)，本原性使 \(u'=\pm1\)，两曲线的像都是水平圆周；将其中一条平移到 \(y=\varepsilon\ne0\) 即不相交，再用 \(P^{-1}\) 送回即可。
\end{solution}

\begin{exercise}\label{g2:ex:product-sphere-self-intersection}
	给 \(S^2\times S^2\) 取先第一因子、后第二因子的乘积定向。令 \(A=S^2\times\{y_0\}\)、\(B=\{x_0\}\times S^2\)，分别取因子定向。计算它们的全部相交配对。若闭已定向嵌入曲面 \(\Sigma\subset S^2\times S^2\) 满足 \([\Sigma]=m[A]+n[B]\)，求 \(I(\Sigma,\Sigma)\)，并说明 \(mn\ne0\) 时为何任何与 \(\Sigma\) 同调的嵌入曲面都不能与它不交。
\end{exercise}
\begin{solution}
	两片在 \((x_0,y_0)\) 唯一相交，先列 \(A\) 的正向基、再列 \(B\) 的正向基，正好是乘积正向基，所以 \(I(A,B)=1\)。两片维数均为 \(2\)，交换符号为 \((-1)^4=1\)，故 \(I(B,A)=1\)。把 \(y_0\) 移到另一点，得到与 \(A\) 同调且不交的纤维，故 \(I(A,A)=0\)；同理 \(I(B,B)=0\)。于是交矩阵为
	\[
		\begin{pmatrix}0&1\\1&0\end{pmatrix},\qquad
		I(\Sigma,\Sigma)=2mn.
	\]
	如果 \(\Sigma'\) 与 \(\Sigma\) 同调且二者不交，则 \(I(\Sigma,\Sigma')=0\)。同调不变性却使该数等于 \(2mn\)，所以 \(mn\ne0\) 时不可能。这里的自交数按两份代表作横截扰动计算，并非把 \(\Sigma\cap\Sigma\) 当成有限点集。
\end{solution}

\begin{exercise}\label{g2:ex:intersection-boundary-sign}
	设 \(M^n\) 是已定向无边界光滑流形，\(C^{p+1}\subset M\) 是紧已定向带边界嵌入子流形，\(A=\partial C\) 采用外法向优先的边界定向，\(B^q\subset M\) 是已定向闭嵌入无边界子流形，\(p+q=n\)。假设 \(C\) 及 \(A\) 都与 \(B\) 横截。用一维交集 \(W=C\cap B\) 证明 \(I(A,B)=0\)，并计算 \(\partial W\) 的点符号与 \(\varepsilon_z(A,B)\) 的关系，其中 \(W\) 按 \((\nu B,\mathrm{T}W)=\mathrm{T}C\) 定向。
\end{exercise}
\begin{solution}
	带边界横截逆像定理给出紧一维流形 \(W\)，并有 \(\partial W=A\cap B\)。设 \(z\) 为边界交点，\(a_1,\ldots,a_p\) 是 \(A\) 的正向基，\(o\) 是 \(C\) 的向外方向，则 \((o,a_1,\ldots,a_p)\) 给出 \(C\) 的正向。因为 \(A\pitchfork B\)，法基可在 \(\mathrm{T}A\) 内提升；其相对于 \((a_1,\ldots,a_p)\) 的符号为 \(\varepsilon_z(A,B)\)。

	把 \(o\) 从这些法基提升之后移到之前，需要交换 \(p\) 次。因而将 \(W\) 的正向与它的向外方向比较，得到
	\[
		\varepsilon_z(\partial W)=(-1)^p\varepsilon_z(A,B).
	\]
	向外向量为使其切于 \(W\) 而添加的切向分量，不改变这一行列式符号。紧一维有向流形的边界符号之和为零，故
	\[
		0=\sum_{z\in\partial W}\varepsilon_z(\partial W)=(-1)^pI(A,B),
	\]
	即 \(I(A,B)=0\)。
\end{solution}

\begin{exercise}\label{g2:ex:projective-lines-mod-two-intersection}
	在 \(\mathbb RP^2\) 中取两条不同的射影直线 \(L_0,L_1\)。证明它们横截相交于一点，模二相交数为 \(1\)，因而与它们分别模二同调的两个嵌入圆周不能不交。解释为何不能在本例中直接使用本节的整数有向相交配对。
\end{exercise}
\begin{solution}
	射影直线来自 \(\mathbb R^3\) 中的二维线性子空间。两个不同二维子空间的交恰为一维，所以对应的射影直线只有一个交点。作线性坐标变换后，可设它们分别为 \(z_2=0\) 和 \(z_1=0\)，交点为 \([1:0:0]\)。在仿射图 \(z_0\ne0\) 中，两条直线是两个坐标轴，故横截。模二计数因此为 \(1\)。

	本节模二版本的几何相交与对偶配对说明该数只依赖模二同调类。若两条相应代表不交，计数就为 \(0\)，与 \(1\) 矛盾。\(\mathbb RP^2\) 不可定向，且直线的法线丛也是不可定向的：把一条直线提升为 \(S^2\) 上的赤道，沿赤道的一个固定向北法向量在对径识别下被微分送到相反方向，绕投影圆周一周后法向反转。因此本例不能连续选择整数局部符号所需的法向定向，而模二计数不受反号影响。
\end{solution}

\begin{exercise}\label{g2:ex:mixed-monomial-vector-index}
	在标准平面切标架中，把向量场的系数写成复数。设它在原点附近满足
	\[
		F(z)=z^p\overline z^{\,q}+r(z),\qquad r(z)=o(|z|^{p+q}),
	\]
	其中 \(p,q\ge0\)、\(p+q\ge1\)，且 \(F\) 光滑。证明原点是孤立零点，指标为 \(p-q\)。在充分小的闭圆盘内，将该场作一个在边界上充分小的扰动，假设新场全部零点非退化。证明新零点总数 \(N\) 满足
	\[
		N\ge|p-q|,\qquad N\equiv p-q\pmod2.
	\]
\end{exercise}
\begin{solution}
	令 \(d=p+q\)。充分小的 \(|z|>0\) 时，\(|r(z)|<|z|^d/2\)，而 \(|z^p\overline z^{\,q}|=|z|^d\)，所以 \(|F(z)|>|z|^d/2>0\)。同时 \(F(0)=0\)，故原点孤立。对 \(0\le s\le1\)，\(z^p\overline z^{\,q}+s r(z)\) 在这个穿孔邻域内都非零。边界归一化因此把 \(F\) 同伦到主项；在 \(z=\varepsilon e^{\mathrm i\theta}\) 上，主项的单位方向为 \(e^{\mathrm i(p-q)\theta}\)。圆周映射度于是给出 \(\operatorname{ind}_0(F)=p-q\)。

	固定上述小圆盘 \(D\)，记 \(c=\min_{\partial D}|F|>0\)。若扰动 \(G\) 满足 \(\sup_{\partial D}|G-F|<c\)，则 \((1-s)F+sG\) 在边界始终非零。新零点集闭且不碰边界，非退化性又使它离散，故有限。记正、负指标零点数为 \(N_+,N_-\)。由\cref{g2:thm:local-index-additivity,g2:prop:nondegenerate-vector-index}，
	\[
		N_+-N_-=p-q,\qquad N=N_++N_-.
	\]
	因此 \(N\ge|p-q|\)，且 \(N-(p-q)=2N_-\)，得到奇偶约束。当 \(p=q\) 时原零点仍孤立，但指标为零，结论只要求扰动后的非退化零点成偶数出现，允许全部消失。
\end{solution}

\begin{exercise}\label{g2:ex:product-vector-field-index}
	设 \(X\) 是 \(m\) 维无边界光滑流形 \(M\) 上的向量场，\(Y\) 是 \(n\) 维无边界光滑流形 \(N\) 上的向量场，\(m,n\ge1\)。若 \(a,b\) 分别是其孤立零点，定义乘积场 \(Z(x,y)=(X(x),Y(y))\)。用局部相对同调外积的自然性证明
	\[
		\operatorname{ind}_{(a,b)}Z=\operatorname{ind}_aX\operatorname{ind}_bY.
	\]
	对 \(\mathbb R^3=\mathbb C\times\mathbb R\) 上系数为 \((z^k,-t)\) 的场，\(k\ge1\)，计算原点指标。
\end{exercise}
\begin{solution}
	取把 \(a,b\) 移到零点的坐标，系数分别为 \(F:U\to\mathbb R^m\)、\(G:V\to\mathbb R^n\)，并缩小 \(U,V\)，使每个映射仅在原点为零。乘积系数 \(F\times G\) 在 \(U\times V\) 中也只在 \((0,0)\) 为零，所以乘积零点孤立。

	记正向局部同调生成元为 \(\mu_m,\mu_n\)。第4章相对外积的自然性适用于穿孔开集，因为
	\[
		(U\setminus\{0\})\times V\ \cup\ U\times(V\setminus\{0\})
		=(U\times V)\setminus\{(0,0)\}.
	\]
	\cref{g2:prop:local-cross-normalization}证明中的定向乘积链说明，先列 \(U\) 坐标、再列 \(V\) 坐标时，\(\mu_m\times\mu_n=\mu_{m+n}\)，没有额外符号。于是
	\begin{align*}
		(F\times G)_*(\mu_m\times\mu_n)
		&=F_*\mu_m\times G_*\mu_n\\
		&=(\deg_0F)(\deg_0G)(\mu_m\times\mu_n).
	\end{align*}
	\cref{g2:prop:vector-index-well-defined}已将系数映射的局部同调度识别为向量场指标，因此得到乘法公式。此论证同时适用于退化零点，不要求微分可逆。

	在 \(\mathbb C\) 因子，\(z^k\) 的原点指标为 \(k\)；在实线因子，\(-t\) 的指标为 \(-1\)。采用坐标顺序 \((\operatorname{Re}z,\operatorname{Im}z,t)\)，乘积指标为 \(-k\)。
\end{solution}

\begin{exercise}\label{g2:ex:punctured-rotating-frame-index}
	在穿孔平面圆盘中写 \(z=r e^{\mathrm i\theta}\)，对整数 \(k\) 取由标准标架转动 \(k\theta\) 得到的正向标架
	\[
		e_1=(\cos k\theta,\sin k\theta),\qquad
		e_2=(-\sin k\theta,\cos k\theta).
	\]
	求径向场 \(X(x,y)=(x,y)\) 在这组标架中的系数沿小圆周的绕数。证明当 \(k\ne0\) 时，这组标架不能作为处处非退化的连续标架延过原点，并说明所得系数绕数为什么不违背局部指标的坐标无关性。
\end{exercise}
\begin{solution}
	把平面向量写成复数时，标架变换就是乘以 \(e^{\mathrm ik\theta}\)。若 \(c\) 表示径向场在新标架中的复系数，则 \(e^{\mathrm ik\theta}c=r e^{\mathrm i\theta}\)，所以
	\[
		c=r e^{\mathrm i(1-k)\theta}.
	\]
	沿小圆周归一化后，其绕数为 \(1-k\)。

	若该标架能非退化地连续延过原点，第一列向量 \(e_1\) 就在整个小闭圆盘上连续且非零。将它归一化，可得从闭圆盘到 \(S^1\) 的连续映射，其边界限制为 \(e^{\mathrm i\theta}\mapsto e^{\mathrm ik\theta}\)。把边界圆沿径向缩到圆心，得到这一映射与常值映射的同伦，故边界度必须为0。但它的度为 \(k\)，所以 \(k\ne0\) 时不可能延拓。

	向量场指标要求使用包含零点的坐标图及其坐标切标架。这里 \(k\ne0\) 的标架只存在于穿孔圆盘，零点处缺少这份坐标识别，故不能用其系数绕数代替原点指标。标准标架中的径向场指标仍为1；差出的 \(k\) 正是沿圆周旋转标架自身所贡献的绕数。
\end{solution}

\begin{exercise}\label{g2:ex:odd-dimensional-total-index}
	设 \(X\) 是无边界光滑 \(n\) 维流形上的向量场，\(n\ge1\)，\(p\) 是孤立零点。证明
	\[
		\operatorname{ind}_p(-X)=(-1)^n\operatorname{ind}_pX.
	\]
	再仅用本章的总指标不变性，证明奇数维闭流形上任何只有孤立零点的光滑向量场，其总指标为零；因此它不可能恰有一个零点且该零点非退化。若所有零点都非退化，正、负指标零点数有何关系？
\end{exercise}
\begin{solution}
	在零点附近写 \(X\) 的系数为 \(F\)。负场的单位方向映射为 \(a\circ\widehat F\)，其中 \(a:S^{n-1}\to S^{n-1}\)、\(a(u)=-u\)。\cref{g2:prop:reflection-antipodal-degree}给出 \(\deg a=(-1)^n\)，复合公式遂得所求恒等式。\(n=1\) 时 \(a\) 交换 \(S^0\) 的两点，在约化零次同调上乘以 \(-1\)，同样成立。

	现设 \(M\) 闭且 \(n\) 为奇数。零点集有限，且 \(X,-X\) 有相同零点。记 \(J=\sum_p\operatorname{ind}_pX\)。\cref{g2:thm:total-vector-index-invariance}使 \(-X\) 的总指标也为 \(J\)，而逐点的公式又使它等于 \(-J\)。所以 \(2J=0\)；\(J\) 为整数，故 \(J=0\)。这里不需要整体定向，也没有使用指标总和与 Euler 示性数的比较。

	单个非退化零点的指标为 \(\pm1\)，无法构成总和0，因此题述情形不可能。若所有零点都非退化，设其正、负指标数为 \(N_+,N_-\)，则 \(N_+-N_-=J=0\)，所以两类数量相同，零点总数为偶数。
\end{solution}

\begin{exercise}\label{g2:ex:torus-affine-lefschetz}
	设 \(A\in M_2(\mathbb Z)\)、\(b\in\mathbb R^2\)，令 \(f:\mathbb R^2/\mathbb Z^2\to\mathbb R^2/\mathbb Z^2\) 为 \(f([x])=[Ax+b]\)。计算 \(L(f)\)。若 \(\det(I-A)\ne0\)，证明它恰有 \(|\det(I-A)|\) 个不动点，并计算每点指数。
\end{exercise}
\begin{solution}
	平移 \(b\) 可通过 \(tb\) 同伦移去。一次同调上的矩阵为 \(A\)，一次上同调上的矩阵为 \(A^{\mathsf T}\)；零次作用为恒等，二次作用为乘 \(\det A\)。因此
	\[
		L(f)=1-\operatorname{tr}A+\det A=\det(I-A).
	\]
	记 \(B=I-A\)。不动点方程是 \(Bx\equiv b\pmod{\mathbb Z^2}\)。若 \(\det B\ne0\)，先取 \(x_0=B^{-1}b\)，全部解在环面上形成陪集 \([x_0]+\ker\overline B\)。核为 \(B^{-1}\mathbb Z^2/\mathbb Z^2\)，乘以 \(B\) 将其同构到 \(\mathbb Z^2/B\mathbb Z^2\)。整数 Smith 标准形将后者写成两个有限循环群的直和，其阶为两个非零对角因子的绝对值之积，即 \(|\det B|\)。每点微分均为 \(A\)，故均非退化，指数都等于 \(\operatorname{sgn}\det B\)。指数和为 \(|\det B|\operatorname{sgn}\det B=\det B\)，与 \(L(f)\) 一致。
\end{solution}

\begin{exercise}\label{g2:ex:even-complex-projective-fixed-point}
	设 \(r\geq1\)。利用 \(H^*(\mathbb CP^{2r};\mathbb Z)=\mathbb Z[u]/(u^{2r+1})\)、\(|u|=2\)，证明每个连续映射 \(f:\mathbb CP^{2r}\to\mathbb CP^{2r}\) 都有不动点。\cite[第2.C节，The Lefschetz Fixed Point Theorem]{Hatcher2002}
\end{exercise}
\begin{solution}
	二次整数上同调由 \(u\) 生成，故 \(f^*u=du\)，其中 \(d\in\mathbb Z\)。拉回保持杯积，所以 \(f^*(u^k)=d^ku^k\)。只有偶次上同调非零，故
	\[
		L(f)=1+d+d^2+\cdots+d^{2r}.
	\]
	若 \(d\geq0\)，这个和至少为 \(1\)。若 \(d\leq-1\)，则
	\[
		1+d+\cdots+d^{2r}=\frac{d^{2r+1}-1}{d-1}>0,
	\]
	因为分子、分母同为负数。于是对所有整数 \(d\) 都有 \(L(f)>0\)，由 Lefschetz 不动点定理得到结论。证明使用了杯积将二次作用与所有高偶次作用联系起来，单独知道各次群的秩还不足以得到这些幂次关系。
\end{solution}

\begin{exercise}\label{g2:ex:lefschetz-product}
	设 \(M,N\) 是闭、已定向光滑流形，\(f:M\to M\)、\(g:N\to N\) 连续。证明 \(L(f\times g)=L(f)L(g)\)。如果两映射都光滑且全部不动点非退化，再用局部指数给出同一等式的几何证明。
\end{exercise}
\begin{solution}
	有限维上同调的 Künneth 同构将 \(H^k(M\times N;\mathbb R)\) 写为 \(\bigoplus_{p+q=k}H^p(M;\mathbb R)\otimes H^q(N;\mathbb R)\)。由外积自然性，\((f\times g)^*\) 在对应直和项上为 \(f^*\otimes g^*\)。若两矩阵为 \(A,B\)，则张量基中对角元为 \(A_{ii}B_{jj}\)，所以 \(\operatorname{tr}(A\otimes B)=\operatorname{tr}A\operatorname{tr}B\)。因此
	\[
		L(f\times g)=\sum_{p,q}(-1)^{p+q}\operatorname{tr}(f^*|_{H^p(M)})\operatorname{tr}(g^*|_{H^q(N)})=L(f)L(g).
	\]
	在非退化情形，不动点集为两个有限不动点集的乘积。微分是块对角矩阵，故
	\[
		\operatorname{ind}(f\times g,(x,y))
		=\operatorname{sgn}\det(I-\mathrm{d}f_x)\operatorname{sgn}\det(I-\mathrm{d}g_y).
	\]
	分别对 \(x,y\) 求和，再用 Lefschetz 公式，便得同一等式。
\end{solution}

\begin{exercise}\label{g2:ex:degenerate-circle-fixed-indices}
	固定 \(0<\varepsilon<1/4\)，在 \(S^1=\mathbb R/\mathbb Z\) 上定义 \(f([t])=[t+\varepsilon\sin^2(2\pi t)]\)。求全部不动点及其指数。再考虑 \(f_a([t])=[t+\varepsilon\sin^2(2\pi t)-a]\)：当 \(0<a<\varepsilon\) 时求不动点个数和指数，并说明当 \(-\varepsilon<a<0\) 时发生了什么。
\end{exercise}
\begin{solution}
	\(f\) 的位移增量在 \([0,\varepsilon]\subset[0,1)\) 内，所以不动点恰为 \([0]\)、\([1/2]\)。各点附近的 \(t-f(t)=-\varepsilon\sin^2(2\pi t)\) 在两侧都为负，故由一维局部指数公式，两点指数都为零。所有 \(f_a\) 都同伦于恒等映射，所以 \(L(f_a)=0\)。

	若 \(0<a<\varepsilon\)，位移增量位于 \((-1,1)\)，不动点方程恰为 \(\sin^2(2\pi t)=a/\varepsilon\)。一个周期内有四个解，每个原不动点两侧各有一个。局部位移为 \(h_a(t)=a-\varepsilon\sin^2(2\pi t)\)，其导数
	\[
		h_a'(t)=-4\pi\varepsilon\sin(2\pi t)\cos(2\pi t)
	\]
	在四个解处均非零，符号依次交替；围绕每个原不动点的两个指数分别为 \(+1\)、\(-1\)，故各小邻域内的总指数仍为零。若 \(-\varepsilon<a<0\)，增量 \(\varepsilon\sin^2(2\pi t)-a\) 严格位于 \((0,2\varepsilon)\subset(0,1)\)，没有不动点。原来的两个零指数不动点因此都可在小扰动下消失。
\end{solution}

% !TeX root = ../../Book2.tex
% 纲要标题：## 第17章　Morse 函数、把手与整体拓扑
% 纲要位置：几何学计划纲要.md，第 1655 行。
% 正文状态：未撰写；本文件目前是框架。

% 纲要内容开始
% 
% > 本章目标与证明范围：从非退化临界点的局部二次型，经过次水平集的真实粘合，建立 Morse 不等式与 Poincaré–Hopf。这里不把动态轨线计数当作已经建立的 Morse 同调。
% 
% 纲要内容结束

\chapter{Morse 函数、把手与整体拓扑}
\label{g2:ch:17}
\BookChapterNumber{g2:ch:17}

\begin{chapterabstract}
待撰写
\end{chapterabstract}

\begin{learninggoals}
待撰写
\end{learninggoals}

待撰写

% !TeX root = ../../Book2.tex
% 纲要标题：### 17.1 临界点及 Morse 引理
% 纲要位置：几何学计划纲要.md，第 1659 行。
% 正文状态：初稿；尚未正式审核。

% 纲要内容开始
% 
% * Hessian 的内在定义
% * 非退化标准形的逐变量消去
% * Morse 函数的存在与扰动
% * 相对及带边界数据
% 
% 纲要内容结束

\section{临界点及 Morse 引理}
\label{g2:sec:1701}

把山地浸入逐渐上涨的水中，水面经过普通高度时，岸线只是移动；经过谷底或山口时，水域可能多出一个分支，或把原有分支连接起来。对流形上的函数，决定这种变化的首先是临界点附近的形状。二阶导数给出最容易计算的信息，Morse 引理则说明：在非退化临界点附近，换一套光滑坐标就能把函数精确化成二次型。

\subsection{Hessian 的内在定义}
\label{g2:subsec:1701:01}

设 \(f:M\to\mathbb R\) 光滑。临界点 \(p\) 满足 \(\mathrm df_p=0\)，所以 \(f\) 沿任何方向的一阶变化都为零。要比较不同方向的二阶变化，不能只把某张坐标图里的二阶偏导数矩阵当成天然对象：换坐标时通常会出现额外的一阶导数项。临界点条件恰好使这些项消失。

\begin{proposition}\label{g2:prop:critical-hessian-intrinsic}
	设 \(M\) 为无边界光滑流形，\(f\in C^\infty(M)\)，\(p\) 为 \(f\) 的临界点。对 \(v,w\in\mathrm T_pM\)，任取局部光滑向量场 \(V,W\)，使 \(V_p=v\)、\(W_p=w\)，则
	\[
		\operatorname{Hess}_pf(v,w)\coloneqq V(Wf)(p)
	\]
	与延拓的选择无关，并给出 \(\mathrm T_pM\) 上的对称双线性型。在坐标 \(x^1,\ldots,x^n\) 中，它的矩阵为 \((\partial_i\partial_jf(p))\)。
\end{proposition}
\begin{proof}
	写 \(V=\sum_iV^i\partial_i\)、\(W=\sum_jW^j\partial_j\)。乘积法则给出
	\[
		V(Wf)(p)=\sum_{i,j}V^i(p)\partial_iW^j(p)\partial_jf(p)
		+\sum_{i,j}V^i(p)W^j(p)\partial_i\partial_jf(p).
	\]
	第一项因 \(\mathrm df_p=0\) 为零，第二项只取决于 \(v,w\)，并且对它们双线性。混合偏导数相等又给出对称性。定义左端用向量场作用于函数写成，已经不依赖坐标；右端给出所述矩阵。
\end{proof}

这个双线性型称为 \(f\) 在临界点的\term[Hessianxing]{Hessian 型}[Hessian]。若坐标变换写作 \(x=x(y)\)，则在临界点
\[
	H_y=J^{\mathsf T}H_xJ,\qquad J=\left(\frac{\partial x^i}{\partial y^j}(p)\right).
\]
这是合同变换。因此矩阵的秩及正、负惯性指数有坐标无关的意义；矩阵的具体特征值则随坐标缩放改变。在非临界点，\(f(x)=x\) 经过局部坐标变换 \(y=x+x^2\)，其在原点的二阶导数从 \(0\) 变成 \(-2\)。这说明临界点条件在上述构造中的作用。

\begin{definition}\label{g2:def:morse-critical-point}
	设 \(f\) 为光滑 \(n\) 维流形上的实值光滑函数，\(p\) 为内部临界点。若 \(\operatorname{Hess}_pf\) 非退化，就称 \(p\) 为\term[feituohualinjiedian]{非退化临界点}[nondegenerate critical point]。该型的负惯性指数称为 \(p\) 的\term[Morsezhibiao]{Morse 指标}[Morse index]，记作 \(\lambda(p)\)。无边界流形上的光滑函数若所有临界点都非退化，就称为\term[Morsehanshu]{Morse 函数}[Morse function]。
\end{definition}

指标 \(\lambda(p)\) 等于 Hessian 在其上负定的子空间的最大维数，满足 \(0\le\lambda(p)\le n\)。它计数的是独立的下降方向。\(x^2+y^2\)、\(-x^2+y^2\)、\(-x^2-y^2\) 在原点的指标依次为 \(0,1,2\)。其中鞍点同时有上升和下降方向，不能用“极大或极小”两类描述。

\ProseParagraphBreak

定义允许不同临界点有相同的临界值。例如 \(\cos(2\theta)\) 是圆上的 Morse 函数，却有两个同高的极大点和两个同高的极小点。以后逐个跨越临界值时，会另外使用使临界值互异的扰动。

\subsection{非退化标准形的逐变量消去}
\label{g2:subsec:1701:02}

Taylor 展开把函数写成二次项加高阶余项；对次水平集的粘合来说，精确的标准形更方便，因为可以直接沿坐标方向构造变形。下面的证明将高阶项吸收到坐标中。关键操作仍是配平方，但平方前的系数现在是光滑函数，必须同时检查所得变量确实构成坐标。

\begin{theorem}[Morse 引理]\label{g2:thm:morse-lemma}
	设 \(M\) 为无边界光滑 \(n\) 维流形，\(f:M\to\mathbb R\) 光滑，\(p\) 为指标 \(\lambda\) 的非退化临界点。则存在以 \(p\) 为原点的光滑坐标 \((y_1,\ldots,y_n)\)，使
	\begin{equation}\label{g2:eq:morse-normal-form}
		f=f(p)-y_1^2-\cdots-y_\lambda^2+y_{\lambda+1}^2+\cdots+y_n^2.
	\end{equation}
\end{theorem}
\begin{proof}
	先在以 \(p\) 为原点的小凸坐标域内计算，并减去常数 \(f(p)\)。沿线段 \(t\mapsto tx\) 两次使用微积分基本定理，利用 \(f(0)=0\)、\(\mathrm df_0=0\)，得到
	\[
		f(x)=\sum_{i,j=1}^n x_ix_j a_{ij}(x),\qquad
		a_{ij}(x)=\int_0^1(1-t)\partial_i\partial_jf(tx)\,\mathrm dt.
	\]
	矩阵 \(A(x)=(a_{ij}(x))\) 光滑、对称，且 \(A(0)=H/2\)。先作线性坐标变换，把 \(A(0)\) 化成对角矩阵；非退化性保证每个对角元非零。

	现在消去第一变量与其余变量的混合项。缩小邻域，使 \(a_{11}\) 不变号，记其符号为 \(\varepsilon_1\)。配平方得到
	\[
		f(x)=a_{11}(x)\left(x_1+\sum_{j>1}\frac{a_{1j}(x)}{a_{11}(x)}x_j\right)^2
		+\sum_{i,j>1}x_ix_j\left(a_{ij}(x)-\frac{a_{i1}(x)a_{1j}(x)}{a_{11}(x)}\right).
	\]
	令
	\[
		y_1=\sqrt{|a_{11}(x)|}\left(x_1+\sum_{j>1}\frac{a_{1j}(x)}{a_{11}(x)}x_j\right),
		\qquad y_j=x_j\quad(j>1).
	\]
	由于 \(a_{1j}(0)=0\) 对 \(j>1\) 成立，此变换在零处的微分是对角矩阵，第一对角元为 \(\sqrt{|a_{11}(0)|}\)，其余为 \(1\)。第一卷第11章第3节的逆函数定理因而给出局部光滑逆。换成这套坐标后，函数成为
	\[
		\varepsilon_1y_1^2+\sum_{i,j>1}y_iy_j b_{ij}(y),
	\]
	其中 \(B(0)\) 仍是非退化对角矩阵。这里 \(b_{ij}\) 可以依赖 \(y_1\)，并不影响下一步：保持 \(y_1\) 不动，对其余变量配平方即可。

	更一般地，已分离出前 \(k-1\) 个平方后，余项具有 \(\sum_{i,j\ge k}z_iz_j c_{ij}(z)\) 的形式，且剩余系数矩阵在零处非退化、对角。保持前 \(k-1\) 个坐标和后 \(n-k\) 个坐标不动，只把第 \(k\) 个坐标改成
	\[
		\sqrt{|c_{kk}(z)|}\left(z_k+\sum_{j>k}\frac{c_{kj}(z)}{c_{kk}(z)}z_j\right).
	\]
	同一微分检验保证这仍是坐标变换，并分离出一个符号为 \(\operatorname{sgn}c_{kk}(0)\) 的平方。有限次之后得到 \(f=\sum_i\varepsilon_i y_i^2\)。这套坐标中的 Hessian 为 \(2\operatorname{diag}(\varepsilon_1,\ldots,\varepsilon_n)\)，合同不变性说明负号恰有 \(\lambda\) 个。重排坐标并加回 \(f(p)\)，即得公式。
\end{proof}

在标准形中，微分为 \(2\sum_i\varepsilon_i y_i\,\mathrm dy_i\)，所以原点是邻域内唯一的临界点。非退化临界点因此孤立；紧流形上的 Morse 函数只有有限个临界点，因为临界点集是闭的，而这些隔离邻域给出它的一个开覆盖，紧性提供有限子覆盖。

\begin{example}\label{g2:exm:morse-coordinate-completion}
	考虑 \(f(x,y)=x^2+2xy+2y^2+y^3\)。原点的 Hessian 为 \(\left(\begin{smallmatrix}2&2\\2&4\end{smallmatrix}\right)\)，正定。直接写成
	\[
		f=(x+y)^2+y^2(1+y)
	\]
	后，取 \(u=x+y\)、\(v=y\sqrt{1+y}\)，便得到 \(f=u^2+v^2\)。变换在原点的行列式为 \(1\)，所以 \((u,v)\) 确为局部坐标。余项 \(y^3\) 进入了坐标 \(v\)，无须令它等于零或用近似舍去。
\end{example}

退化临界点可能有不同的局部行为。\(x^4+y^2\) 在原点是严格极小值，但 Hessian 有核；任何局部坐标变换都保持 Hessian 的秩，故它不能变成非退化二次型。非退化性给出了能被有限个正、负平方完全描述的那类临界点。

\subsection{Morse 函数的存在与扰动}
\label{g2:subsec:1701:03}

局部标准形很有用，还需要知道流形上能否找到这样的函数。把 \(\mathrm df\) 看成余切丛的一个截面，临界点就是它与零截面的交点。在临界点，截面沿纤维方向的导数正是 Hessian，所以非退化性等价于这个交点横截。扰动时须仍保持截面是某个函数的微分；直接对函数加入局部线性项就能做到这一点。

\begin{theorem}\label{g2:thm:morse-approximation}
	设 \(M\) 为有限维、Hausdorff、第二可数的无边界光滑流形，\(r\ge2\) 为有限整数。每个 \(f\in C^\infty(M)\) 的任意强 \(C^r\) 邻域中都有光滑 Morse 函数。若 \(M\) 紧，则 Morse 函数还构成 \(C^\infty(M)\) 中关于 \(C^2\) 拓扑的开集。
\end{theorem}
\begin{proof}
	先说明局部扰动。在坐标域 \(U\) 中取紧集 \(K\)，并选紧支于 \(U\) 的光滑函数 \(\rho\)，使其在 \(K\) 的一个邻域 \(V\) 上等于 \(1\)。令
	\[
		f_a=f+\rho\sum_{j=1}^n a_jx^j,\qquad a\in\mathbb R^n,
	\]
	并在 \(U\) 外仍取 \(f_a=f\)。在 \(V\) 上，微分的坐标分量为
	\[
		G(a,x)=(\partial_1f(x)+a_1,\ldots,\partial_nf(x)+a_n).
	\]
	\(G\) 关于参数的微分为恒等矩阵，所以它是淹没。\cref{g2:thm:parametric-transversality}说明：任意小参数中都有 \(a\)，使 \(G_a\) 横截于 \(\{0\}\)。在 \(G_a=0\) 处，\(\mathrm D_xG_a\) 就是 \(f_a\) 的 Hessian，横截等价于它可逆。因此 \(f_a\) 在 \(K\) 上只有非退化临界点，且扰动在任意指定有限阶意义下随 \(a\to0\) 趋于零。

	保持已经得到的非退化性需要二阶控制。对任意紧集 \(L\)，条件“\(f\) 在 \(L\) 上的全部临界点非退化”等价于微分截面 \(x\mapsto(x,\mathrm df_x)\) 在 \(L\) 上横截于余切丛的零截面。后者是闭嵌入子流形。由\cref{g2:prop:transversality-open-on-compact}，该条件对截面的 \(C^1\) 小扰动开放，也就对函数的 \(C^2\) 小扰动开放。坐标上，这分别控制非临界区域的一阶导数不变成零，以及临界点邻域内 Hessian 的可逆性。

	若 \(M\) 紧，取有限个上述紧集 \(K_i\)，使其内部覆盖 \(M\)，依次作局部扰动。在每一步把参数取得足够小，就能保留先前紧集上的非退化性，并使总 \(C^r\) 误差小于给定界。最后的函数在所有 \(K_i\) 上满足条件，因而为 Morse 函数。刚才的紧集开放性取 \(L=M\)，也证明了紧情形的开放结论。

	一般的 \(M\) 取可数、局部有限的相对紧坐标域 \(U_i\)，以及 \(K_i\subset U_i\)，使 \(\operatorname{Int}K_i\) 覆盖 \(M\)。每步扰动紧支于 \(U_i\)。对原强邻域在第 \(\ell\) 个控制块上的容差 \(\varepsilon_\ell\)，要求第 \(i\) 步在那里耗用小于 \(2^{-i-1}\varepsilon_\ell\)。第 \(i\) 步完成后，在 \(K_i\) 上选一个保持非退化性的 \(C^2\) 容差 \(\delta_i>0\)；对后来第 \(j\) 步，另要求它在覆盖 \(K_i\) 的坐标块上的 \(C^2\) 改变量小于 \(2^{-(j-i+1)}\delta_i\)。每步只遇到有限多个原控制块和有限多个先前要求，故可以同时满足。局部有限性又使每点附近只有有限次实际改变，所以极限函数光滑。误差级数保证它仍在原邻域内，并保留每个 \(K_i\) 上的非退化性。这是\cref{g2:thm:relative-transversality}证明中的逐块选择，此处每一步都通过函数的局部线性项实施。
\end{proof}

\begin{proposition}\label{g2:prop:distinct-morse-critical-values}
	设 \(M\) 为紧无边界光滑流形，\(f:M\to\mathbb R\) 为 Morse 函数。任意有限阶 \(C^r\) 邻域内都存在 Morse 函数 \(g\)，与 \(f\) 有完全相同的临界点和指标，且不同临界点的 \(g\) 值互异。
\end{proposition}
\begin{proof}
	临界点为有限个 \(p_1,\ldots,p_N\)。取两两不交的坐标邻域 \(U_i\)，以及紧支于 \(U_i\)、在 \(p_i\) 的较小邻域 \(V_i\) 上为 \(1\) 的截断函数 \(\rho_i\)。令 \(g=f+\sum_i c_i\rho_i\)。在每个 \(V_i\) 上，\(g\) 与 \(f\) 相差常数，故临界点和 Hessian 完全相同。在紧集 \(M\setminus\bigcup_iV_i\) 上，\(\mathrm df\) 处处不为零；相对于任一度量，它的范数有正下界。令全部 \(c_i\) 足够小，便使这里的 \(\mathrm dg\) 仍不为零。最后避开有限个等式 \(f(p_i)+c_i=f(p_j)+c_j\)，可使所有临界值互异。参数还可任意小，所以满足所给 \(C^r\) 要求。
\end{proof}

这种调整改变临界高度，却不改变局部的下降方向数。它使相邻两个正则水平之间可以只放一个临界点，正适合下一节逐次描述次水平集的变化。

\subsection{相对及带边界数据}
\label{g2:subsec:1701:04}

在粘合问题中，一部分函数已经规定好，例如边界附近要把高度直接当成领圈参数。此时扰动应避开那一部分。能保持哪些数据，取决于固定区域里原有临界点是否已经非退化。

\begin{proposition}\label{g2:prop:relative-morse-approximation}
	设 \(M\) 为有限维、Hausdorff、第二可数的无边界光滑流形，\(A\subset M\) 闭，\(f\in C^\infty(M)\) 在 \(A\) 的某开邻域上只有非退化临界点。则对任意有限整数 \(r\ge2\)，可在 \(f\) 的任意强 \(C^r\) 邻域内找到 Morse 函数 \(g\)，使 \(g=f\) 在 \(A\) 的某开邻域上成立。
\end{proposition}
\begin{proof}
	在已满足条件的开集 \(U\) 中选 \(A\subset W\subset\overline W\subset V\subset\overline V\subset U\)。先用逐紧块的 \(C^2\) 开放性，缩小给定强邻域，使其中的函数在 \(\overline V\) 上都保持非退化性。再对 \(M\setminus V\) 取局部有限的紧坐标块及支集域，使支集均避开 \(\overline W\)。按\cref{g2:thm:morse-approximation}的局部线性扰动及误差分配逐块处理，得到在整个 \(M\) 上为 Morse 的 \(g\)。全部改动支集避开 \(\overline W\)，所以 \(g=f\) 于 \(W\)。
\end{proof}

带边界时，本章所用数据是边界为正则水平，临界点都在内部。下述构造同时适用于只有外边界的紧流形和两端高度已指定的配边。

\begin{proposition}\label{g2:prop:boundary-morse-function}
	设 \(M\) 为紧光滑带边界 \(n\) 维流形，\(n\ge1\)，并把边界分成互不相交的开闭子集 \(\partial M=\partial_-M\amalg\partial_+M\)，允许其中一部分为空。给定 \(a<b\)，存在光滑函数 \(f:M\to[a,b]\) 及两侧的领圈，使
	\[
		f(x,t)=a+t\quad\text{在 }\partial_-M\text{ 的较小领圈内},\qquad
		f(x,t)=b-t\quad\text{在 }\partial_+M\text{ 的较小领圈内}.
	\]
	这里 \(t\ge0\) 指向内部。函数在内部取值于 \((a,b)\)，所有临界点非退化且临界值互异。可选择在领圈内为 \(\mathrm dt^2+g_{\partial M}\) 的 Riemann 度量；此时 \(\nabla f\) 在 \(\partial_-M\) 严格向内，在 \(\partial_+M\) 严格向外。
\end{proposition}
\begin{proof}
	由\cref{g2:thm:collar-neighborhood}选择两侧互不相交的领圈，并把宽度取得小于 \((b-a)/4\)。在更小的领圈内分别规定 \(f_0=a+t\)、\(f_0=b-t\)，在领圈以外取常数 \((a+b)/2\)。用只依赖 \(t\) 的截断函数把这些表达作凸组合，得到光滑 \(f_0\)。在内部出现的各个参与值都严格介于 \(a,b\)，所以 \(a<f_0<b\)；在固定的较小领圈内，\(\mathrm df_0=\pm\mathrm dt\)，没有临界点。

	对 \(\operatorname{Int}M\) 应用\cref{g2:prop:relative-morse-approximation}，固定一条更窄领圈中的全部函数值。所有实际改动可支撑在远离边界的一个紧集内；在那里 \(f_0\) 与 \(a,b\) 都有正距离。因此同时要求 \(C^0\) 误差足够小，就能保持 \(a<f<b\)。所得函数与 \(f_0\) 在边界邻域一致，故光滑延拓到 \(M\)。临界点避开固定领圈，处在紧集内，因而只有有限个。按\cref{g2:prop:distinct-morse-critical-values}的截断常数构造，在内部各临界点附近作更小的扰动，即可使临界值互异而不改变边界数据及值域。

	在领圈上先取乘积度量 \(\mathrm dt^2+g_{\partial M}\)，再用单位分解延拓为全局正定度量，并保持较小领圈上的度量不变。由梯度定义 \(g(\nabla f,\cdot)=\mathrm df\)，两端分别有 \(\nabla f=\partial_t\) 和 \(\nabla f=-\partial_t\)。因为 \(t\) 向内增加，这给出所需指向。
\end{proof}

特别地，取 \(\partial_-M=\varnothing\)，便得到边界为最高水平、梯度处处外向的 Morse 函数。它将用于带边界的 Poincaré--Hopf 公式。若事先要求在某邻域逐点保留一个退化临界点，Morse 逼近就无法满足；相对命题中的条件准确描述了固定数据应先具备的性质。

% !TeX root = ../../Book2.tex
% 纲要标题：### 17.2 次水平集的变化
% 纲要位置：几何学计划纲要.md，第 1666 行。
% 正文状态：初稿；尚未正式审核。

% 纲要内容开始
% 
% * 无临界区间的梯度变形
% * 跨越单个临界值的局部模型
% * 把手附着映射与指标
% * 从把手到胞腔同伦类型
% 
% 纲要内容结束

\section{次水平集的变化}
\label{g2:sec:1702}

把一个函数的高度从低向高扫描，得到一族不断增大的区域。Morse 引理描述了临界点附近的高度，却还没有说明这个小区域怎样接到已经出现的部分。本节把两件事连接起来：无临界点处用流搬运等高面，临界点附近则直接放入一个有厚度的圆盘。它的厚度保留光滑流形结构，中心圆盘负责改变同伦型。

\subsection{无临界区间的梯度变形}
\label{g2:subsec:1702:01}
记 \(M^a=f^{-1}(( -\infty,a])\)，称为 \(f\) 的\term[cishuipingji]{次水平集}[sublevel set]。若 \(a\) 是正则值且 \(f^{-1}(a)\) 位于 \(M\) 内部，则正则水平集定理给出 \(M^a\) 的新边界 \(f^{-1}(a)\)。在局部坐标中，\(f-a\) 本身就是横向坐标，因而次水平条件成为一个半空间条件。

\begin{proposition}\label{g2:prop:morse-regular-band}
	设 \(M\) 为光滑流形，\(f:M\to\mathbb R\) 光滑，\(a<b\)，且 \(f^{-1}([a,b])\) 是位于 \(M\) 内部的紧集，其中没有临界点。则
	\[
		f^{-1}([a,b])\cong f^{-1}(a)\times[a,b]
	\]
	的微分同胚可取为使 \(f\) 对应第二投影；\(M^a\hookrightarrow M^b\) 是强形变收缩的包含映射，\(M^a\) 与 \(M^b\) 也微分同胚。
\end{proposition}
\begin{proof}
	选光滑 Riemann 度量。由于 \(\mathrm{d}f\ne0\)，向量场
	\[
		Y=\frac{\operatorname{grad}f}{|\operatorname{grad}f|^2}
	\]
	在闭带的某邻域有定义，且 \(Yf=1\)。用紧支截断将它延拓到 \(M\) 上，保持闭带附近的表达。第一卷第13章第1节的光滑流定理及第2节的紧支完备性说明，这个紧支向量场的流 \(\Phi_t\) 对全部时间存在；轨道留在闭带期间满足 \(f(\Phi_t(x))=f(x)+t\)。于是
	\[
		(z,t)\longmapsto\Phi_{t-a}(z),\qquad
		x\longmapsto\bigl(\Phi_{a-f(x)}(x),f(x)\bigr)
	\]
	互为光滑逆。这也解释了紧性的作用：所需时间内的轨道不会在空间中逃走。

	在 \(M^b\) 上令
	\[
		R_s(x)=
		\begin{cases}
			x,&f(x)\le a,\\
			\Phi_{-s(f(x)-a)}(x),&a\le f(x)\le b,
		\end{cases}
		\qquad 0\le s\le1.
	\]
	两式在 \(f=a\) 处一致，给出连续同伦，且逐点固定 \(M^a\)。当 \(s=1\) 时，整个上方闭带落到 \(f^{-1}(a)\)，故得到强形变收缩。

	为得到微分同胚，取 \(a\) 以下一个足够薄的无临界领带；它由闭带的紧性及 \(\mathrm{d}f\ne0\) 延伸得到。在乘积坐标中，用一个光滑严格递增的区间微分同胚把 \([a-\delta,b]\) 压到 \([a-\delta,a]\)，并使它在左端附近为恒等。与更低处的恒等映射拼接，即得 \(M^b\cong M^a\)。这里的压缩移动了边界；前一个强形变收缩则只要求连续，并固定整个旧区域。
\end{proof}

\begin{example}
	在 \(M=\mathbb R^2\setminus\{(0,0)\}\) 上取 \(f(x,y)=x\)。它处处无临界点，但 \(f^{-1}(0)\) 有两个连通分支，其他水平集都是一条直线。\(M^{-1}\) 可缩，而 \(M^1\) 绕缺失的原点有一个不可缩圆周。闭带不紧，正是上面命题无法使用的条件：穿过原点的水平搬运会离开 \(M\)。
\end{example}

\subsection{跨越单个临界值的局部模型}
\label{g2:subsec:1702:02}
现在设临界点 \(p\) 的指标为 \(\lambda\)，写 \(\mu=n-\lambda\)。由\cref{g2:thm:morse-lemma}，可在 \(p\) 附近选坐标 \((u,v)\in\mathbb R^\lambda\times\mathbb R^\mu\)，使
\begin{equation}\label{g2:eq:morse-local-handle-model}
	f(u,v)=c-|u|^2+|v|^2,\qquad c=f(p).
\end{equation}
负方向 \(u\) 在离开临界点时降低高度，正方向 \(v\) 提供横向厚度。选充分小的 \(\varepsilon>0\) 和 \(r>0\)，满足 \(r^2<\varepsilon\)，并使下面的闭集全落在坐标域中：
\[
	H=\bigl\{(u,v):|v|\le r,\ |u|^2\le\varepsilon+|v|^2\bigr\}.
\]
它有一个明确的乘积参数化
\begin{equation}\label{g2:eq:morse-handle-parametrization}
	\Phi:D^\lambda\times D^\mu\longrightarrow H,
	\qquad
	\Phi(x,w)=\bigl(\sqrt{\varepsilon+r^2|w|^2}\,x,rw\bigr).
\end{equation}
缩放因子严格为正，故这是带角流形之间的微分同胚。代入高度可见
\[
	f\circ\Phi
	=c+r^2|w|^2-(\varepsilon+r^2|w|^2)|x|^2.
\]
因此 \(H\subset M^{c+\varepsilon}\)，且
\begin{equation}\label{g2:eq:morse-handle-attaching-face}
	H\cap M^{c-\varepsilon}
	=\Phi(S^{\lambda-1}\times D^\mu)
	\subset f^{-1}(c-\varepsilon).
\end{equation}
新块只沿第一因子的边界贴到旧区域上；它的其余内部点都在旧水平以上。\(\lambda=0\) 时约定 \(S^{-1}=\varnothing\)，新块与旧区域不相交；\(\mu=0\) 时新块是整个 \(n\) 维圆盘，其全部边界都被附着。

\begin{figure}[htbp]
	\centering
	\input{Figures/Chapter17-LocalHandle}
	\caption[临界点附近的把手]{二维指标一模型 \(f-c=-u^2+v^2\)。阴影块 \(H\) 沿左右两段粗曲线贴到旧次水平集；中心线是核心圆盘。下降向量场沿 \(u\) 方向向外、沿 \(v\) 方向向内，将剩余闭带送入旧区域与新块的并。}
	\label{g2:fig:morse-local-handle}
\end{figure}

局部二次型还提供了控制剩余空间的向量场。在这个坐标域内选 Euclid 度量，则负梯度为
\[
	X=(2u,-2v),\qquad Xf=-4(|u|^2+|v|^2).
\]
在自由面 \(|v|=r\) 上，有 \(X(|v|^2-r^2)=-4r^2<0\)，所以流从外面进入 \(H\)；在旧水平 \(f=c-\varepsilon\) 上，它也严格指向旧次水平集。除原点外，高度一直下降。这两个明确的不等式将局部模型与整体流联系起来。

\subsection{把手附着映射与指标}
\label{g2:subsec:1702:03}
\begin{definition}\label{g2:def:morse-handle}
	设 \(W\) 为光滑 \(n\) 维带边界流形，\(0\le\lambda\le n\)。一个\term[bashou]{指标 \(\lambda\) 的把手}[handle of index \(\lambda\)]是乘积 \(D^\lambda\times D^{n-\lambda}\)，其附着面为 \(S^{\lambda-1}\times D^{n-\lambda}\)。给定光滑嵌入
	\[
		\varphi:S^{\lambda-1}\times D^{n-\lambda}\longrightarrow\partial W,
	\]
	将附着面上的点 \((x,w)\) 与 \(\varphi(x,w)\) 识别，再将接缝的角磨圆，称为向 \(W\) 附着这个把手。\(D^\lambda\times\{0\}\) 称为\term[hexinyuanpan]{核心圆盘}[core disk]，\(\{0\}\times D^{n-\lambda}\) 称为\term[gonghexinyuanpan]{共核心圆盘}[cocore disk]。
\end{definition}

附着数据包含整个有厚度的球面邻域。仅指定中心球面 \(S^{\lambda-1}\times\{0\}\) 的像，还没有说明横向圆盘怎样放置；后续手术中的法标架正来自这份数据。

\begin{theorem}[Morse 把手附着定理]\label{g2:thm:morse-handle-attachment}
	设 \(M\) 为光滑 \(n\) 维流形，\(f:M\to\mathbb R\) 光滑，\(a<b\) 为正则值，且紧闭带 \(f^{-1}([a,b])\) 位于 \(M\) 内部。若其中恰有一个临界点 \(p\)，该点非退化且指标为 \(\lambda\)，则 \(M^b\) 微分同胚于向 \(M^a\) 附着一个指标 \(\lambda\) 的把手并磨圆所得的流形。附着映射可由式\eqref{g2:eq:morse-handle-parametrization}的附着面沿无临界带下降到 \(f^{-1}(a)\) 得到。
\end{theorem}
\begin{proof}
	令 \(c=f(p)\)。先取充分小的 \(\varepsilon>0\)，使 \(a<c-\varepsilon<c+\varepsilon<b\)，并能使用上面的坐标块。两侧闭带 \([a,c-\varepsilon]\)、\([c+\varepsilon,b]\) 都没有临界点，\cref{g2:prop:morse-regular-band}将它们识别为领带。因此只需处理 \(a=c-\varepsilon\)、\(b=c+\varepsilon\)；最后把附着映射沿下方的流搬回原来的水平即可。

	记 \(A=M^{c-\varepsilon}\)，取式\eqref{g2:eq:morse-handle-parametrization}中的 \(H\)，令 \(K=A\cup H\)。式\eqref{g2:eq:morse-handle-attaching-face}说明，\(K\) 确实是沿嵌入 \(\Phi|_{S^{\lambda-1}\times D^\mu}\) 作出的把手附着商：两块的交恰为附着面，且 \(H\) 紧，故这个商到它们在 \(M\) 中的并是同胚。坐标参数化给出了新块的光滑结构。

	将局部 Euclid 度量与外部任意度量用单位分解拼接，使其在 \(H\) 的一个邻域内仍为 Euclid 度量。令 \(X=-\operatorname{grad}f\)。旧边界上 \(Xf<0\)，新块的自由面上 \(X(|v|^2-r^2)<0\)，所以 \(X\) 在 \(K\) 的每个边界面都朝向内部。它的唯一零点 \(p\) 已在 \(K\) 内部。

	先把交界的角变成光滑边界。若 \(0<\lambda<n\)，角附近可用两个独立的横向函数
	\[
		s=f-c+\varepsilon,\qquad t=|v|^2-r^2
	\]
	描述 \(K=\{s\le0\}\cup\{t\le0\}\)。两者的微分在 \(s=t=0\) 处线性无关：此处 \(u,v\) 均非零，而 \(\mathrm{d}t\) 没有 \(u\) 分量。取光滑偶函数 \(\eta_\delta\)，满足 \(\eta_\delta(z)\ge|z|\)、\(|\eta_\delta'(z)|\le1\)，并在 \(|z|\ge\delta\) 时等于 \(|z|\)。例如将绝对值与支集很小的偶非负平滑核卷积即可。把角附近的边界换成
	\[
		m_\delta(s,t)=\frac{s+t-\eta_\delta(s-t)}2=0,
	\]
	并取 \(m_\delta\le0\) 一侧。它包含原来的 \(K\)，且新增边界只发生在 \(|s|,|t|\) 都很小的地方，因此与角邻域外的原边界精确拼接。两偏导非负且之和为一；由于 \(Xs<0\)、\(Xt<0\)，有 \(Xm_\delta<0\)，所以新边界光滑，\(X\) 仍严格朝内。记磨圆后的区域为 \(K_\delta\)。取 \(\delta\) 小，可使 \(K_\delta\setminus A\subset\{f<c+\varepsilon\}\)。指标为 \(0\) 或 \(n\) 时没有这种角，直接令 \(K_\delta=K\)。

	剩下的紧区域 \(M^{c+\varepsilon}\setminus\operatorname{int}_M K_\delta\) 没有临界点。其上 \(|\operatorname{grad}f|^2\) 有严格正的下界。因此从新增边界 \(B=\partial K_\delta\cap f^{-1}([c-\varepsilon,c+\varepsilon])\) 沿 \(-X\) 向外运动，必在有限时间到达 \(f=c+\varepsilon\)：若一直不到达，高度将以一个固定正速度增加而始终不超过该值，矛盾。反向轨道同样在有限时间到达 \(B\)。光滑依赖与隐函数定理说明首次到达时间 \(\tau(z)>0\) 随 \(z\in B\) 光滑变化，因为两端的向量场都横截边界。流遂给出
	\[
		M^{c+\varepsilon}\setminus\operatorname{int}_M K_\delta
		\cong\{(z,t):z\in B,\ 0\le t\le\tau(z)\}.
	\]
	没有外部边界的连通分支已完全包含于 \(K_\delta\)，在此乘积中不出现。

	这段外部乘积可以吸收到 \(K_\delta\) 的一个内部领圈中。具体说，在每条领线的区间 \([-h,\tau(z)]\) 上，选依赖 \(z\) 光滑、在左端附近恒等的严格递增微分同胚到 \([-h,0]\)；\(B\) 紧，故可统一取 \(h>0\)。这样的区间映射可由一个处处为正、左端等于一的光滑密度积分而成，再使总积分等于 \(h\)。与领圈外的恒等映射拼接，得到 \(M^{c+\varepsilon}\cong K_\delta\)。至此既确定了附着面，也处理了剩余空间和接缝光滑结构。
\end{proof}

上面证明还给出一个相对结论，后面计算同调时会用到：未磨圆的 \(K=A\cup H\) 是 \(M^{c+\varepsilon}\) 的强形变收缩像。沿 \(X\) 流到第一次进入 \(K\) 即停下，进入时间在各光滑面附近由横截性连续变化；角附近它是到两个边界面的时间的较小者，仍连续。剩余紧区域中 \(Xf\) 有严格负的上界，保证每条轨道都在有限时间进入。将这段时间乘以同伦参数，就得到固定 \(K\)、从而固定 \(A\) 的收缩。这个结论不需要先磨圆。

\ProseParagraphBreak

在曲面上，指标零把手是一个新圆盘；指标一把手是沿两条短边接上的带；指标二把手沿整个圆周填入一个圆盘。中间那条带可能连接两个原有分支，也可能把同一个分支的两个部位连起来。局部指标决定新增块的形状，附着位置则决定全局拓扑怎样改变。

\subsection{从把手到胞腔同伦类型}
\label{g2:subsec:1702:04}
把手中 \(D^{n-\lambda}\) 的厚度可以连续压掉，但在压缩时必须固定已经贴在旧区域上的整个附着面。直接将 \((x,w)\) 送到 \((x,(1-s)w)\) 会移动附着面，无法与旧区域上的恒等映射拼接。下面给出保持粘合位置的收缩。

\begin{lemma}\label{g2:lem:handle-core-retraction}
	乘积圆盘 \(D^\lambda\times D^\mu\) 强形变收缩到
	\[
		Q=(D^\lambda\times\{0\})\cup(S^{\lambda-1}\times D^\mu).
	\]
	因此沿 \(\varphi\) 附着把手所得的空间强形变收缩到沿 \(\varphi(x,0)\) 附着核心 \(D^\lambda\) 所得的空间，且收缩固定旧空间。
\end{lemma}
\begin{proof}
	给出一条先增大 \(|x|\)、同时减小 \(|w|\)，直到碰到 \(|x|=1\) 或 \(w=0\) 的路径。置
	\[
		T(x,w)=\min\{|w|,-\log|x|\},\qquad
		T(0,w)=|w|.
	\]
	当 \(w\ne0\) 时，令
	\[
		F_s(x,w)=\left(e^{sT(x,w)}x,
		\left(1-\frac{sT(x,w)}{|w|}\right)w\right),
		\qquad 0\le s\le1;
	\]
	当 \(w=0\) 时取 \(F_s(x,0)=(x,0)\)。因为 \(T\le-\log|x|\)、\(T\le|w|\)，路径始终留在乘积圆盘内。\(x\to0\) 时 \(-\log|x|\to\infty\)，所以 \(T=|w|\) 在其附近成立；\(w\to0\) 时 \(T\to0\)，第二分量的范数至多为 \(|w|\)。这些估计保证整个公式连续。\(s=1\) 时至少达到 \(|x|=1\) 或 \(w=0\) 之一；原本就在 \(Q\) 上的点有 \(T=0\)，始终不动。\(\lambda=0\) 时这就是圆盘到中心的径向收缩，\(\mu=0\) 时全部点已在核心上。最后在旧空间上取恒等映射，因附着面逐点不动，公式经附着商下降为所需的强形变收缩。
\end{proof}

由把手附着定理后的相对收缩及此引理，单临界点跨越给出
\begin{equation}\label{g2:eq:morse-cell-attachment}
	M^{c+\varepsilon}\simeq
	M^{c-\varepsilon}\cup_\alpha D^\lambda,
	\qquad
	\alpha(x)=\Phi(x,0),\quad x\in S^{\lambda-1},
\end{equation}
且同伦等价固定旧次水平集。这一步把负二次项的个数变成了实际附着胞腔的维数。

\begin{corollary}\label{g2:cor:morse-cell-type}
	设 \(M\) 为紧光滑 \(n\) 维流形。若 \(M\) 无边界，取任意 Morse 函数 \(f\)；若有边界，设 \(f=b-t\) 于一个向内参数为 \(t\) 的领圈，\(f<b\) 于内部，且全部内部临界点非退化。记指标 \(k\) 的临界点数为 \(m_k\)。则 \(M\) 具有有限把手分解，并同伦等价于一个恰有 \(m_k\) 个 \(k\) 胞腔的有限 CW 复形。
\end{corollary}
\begin{proof}
	非退化临界点孤立；紧性和边界附近无临界点保证它们只有有限多个。由\cref{g2:prop:distinct-morse-critical-values}的小扰动，可分开临界值且保持临界点指标，并在带边界情形保持领圈上的函数。取低于最小临界值的水平，此时次水平集为空。在相邻临界值之间取正则水平，逐次应用\cref{g2:thm:morse-handle-attachment}，便得到所述有限把手分解。带边界时最后只剩一个无临界的边界领带，吸收这段领带不再增加把手。

	逐次使用式\eqref{g2:eq:morse-cell-attachment}，每次只增加一个同指标的胞腔。这里附着次序由临界值排列，未必按指标排列，所以还须说明能得到通常的 CW 结构。假设前一阶段已用有限 CW 复形 \(L\) 表示，把新附着映射先送入 \(L\)，再用\cref{g2:thm:cellular-approximation}将 \(S^{\lambda-1}\to L\) 同伦到 \(L^{\lambda-1}\)。\cref{g2:prop:homotopic-attaching-maps}说明这不改变附着后的同伦型，并保持旧空间。将旧空间换成同伦等价的 \(L\) 时，同样可把同伦逆与两段逆同伦放在附着球面的映射柱上，再用圆盘对的同伦延拓性质延入圆盘；\cref{g2:thm:cofibration-pushout}保证映射柱可压回普通附着，故所得空间仍同伦等价。于是每一步都保留已有胞腔，只向正确的低维骨架增加一个 \(\lambda\) 胞腔。有限归纳得到所需 CW 复形，并保持各维胞腔数。
\end{proof}

这两个滤过承担不同任务：实际次水平集按高度增长，CW 骨架按维数增长。前者保留临界点穿越的几何过程，后者适合使用第3章的胞腔链复形。把二者通过相对同伦等价联系起来，才能将临界点个数与同调维数比较。
% 纲要条目排为编号小节标题，并进入目录与书签；说明段仍保留于纲要注释中。

% !TeX root = ../../Book2.tex
% 纲要标题：### 17.3 Morse 不等式
% 纲要位置：几何学计划纲要.md，第 1673 行。
% 正文状态：未撰写；本文件目前是框架。

% 纲要内容开始
% 
% * 相对同调的临界贡献
% * 滤过链群与弱不等式
% * 强不等式和 Euler 等式
% * 完美 Morse 函数及反例
% 
% 纲要内容结束

\section{Morse 不等式}
\label{g2:sec:1703}

待撰写

\subsection{相对同调的临界贡献}
\label{g2:subsec:1703:01}
待撰写

\subsection{滤过链群与弱不等式}
\label{g2:subsec:1703:02}
待撰写

\subsection{强不等式和 Euler 等式}
\label{g2:subsec:1703:03}
待撰写

\subsection{完美 Morse 函数及反例}
\label{g2:subsec:1703:04}
待撰写
% 纲要条目排为编号小节标题，并进入目录与书签；说明段仍保留于纲要注释中。

% !TeX root = ../../Book2.tex
% 纲要标题：### 17.4 Poincaré–Hopf 的证明
% 纲要位置：几何学计划纲要.md，第 1680 行。
% 正文状态：未撰写；本文件目前是框架。

% 纲要内容开始
% 
% * 梯度零点指标与 Morse 指标
% * 指标和的扰动不变性
% * 任意孤立零点向量场的公式
% * 边界版本及指向条件
% 
% 纲要内容结束

\section{Poincaré–Hopf 的证明}
\label{g2:sec:1704}

待撰写

\subsection{梯度零点指标与 Morse 指标}
\label{g2:subsec:1704:01}
待撰写

\subsection{指标和的扰动不变性}
\label{g2:subsec:1704:02}
待撰写

\subsection{任意孤立零点向量场的公式}
\label{g2:subsec:1704:03}
待撰写

\subsection{边界版本及指向条件}
\label{g2:subsec:1704:04}
待撰写
% 纲要条目排为编号小节标题，并进入目录与书签；说明段仍保留于纲要注释中。

% !TeX root = ../../Book2.tex
% 纲要标题：### 17.5 把手方法的深入应用
% 纲要位置：几何学计划纲要.md，第 1637 行。
% 正文状态：未撰写；本文件目前是框架。

% 纲要内容开始
% 
% * 曲面与射影空间的分解
% * 把手消去的局部条件
% * 配边与手术的接口
% 
% ---
% 
% 纲要内容结束

\section{把手方法的深入应用}
\label{g2:sec:1705}

待撰写

\subsection{曲面与射影空间的分解}
\label{g2:subsec:1705:01}
待撰写

\subsection{把手消去的局部条件}
\label{g2:subsec:1705:02}
待撰写

\subsection{配边与手术的接口}
\label{g2:subsec:1705:03}
待撰写
% 纲要条目排为编号小节标题，并进入目录与书签；说明段仍保留于纲要注释中。


\begin{chaptersummary}
待撰写
\end{chaptersummary}

\BookExercises
待撰写

% !TeX root = ../../Book2.tex
% 纲要标题：## 第18章*　Morse–Smale 同调：轨线、紧化与粘合
% 纲要位置：几何学计划纲要.md，第 1695 行。
% 正文状态：未撰写；本文件目前是框架。

% 纲要内容开始
% 
% > 本章目标与证明范围：闭光滑流形、Morse–Smale 数据、F₂ 系数下完成整个同调构造及与奇异同调比较。整数取向另列进阶节，不以“端点成对出现”代替紧化为流形的证明。
% 
% 纲要内容结束

\OptionalChapter{Morse–Smale 同调：轨线、紧化与粘合}{g2:ch:18}
\BookChapterNumber{g2:ch:18}

\begin{chapterabstract}
待撰写
\end{chapterabstract}

\begin{learninggoals}
待撰写
\end{learninggoals}

待撰写

% !TeX root = ../../Book2.tex
% 纲要标题：### 18.1 梯度流的局部动力学
% 纲要位置：几何学计划纲要.md，第 1699 行。
% 正文状态：初稿；尚未正式审核。

% 纲要内容开始
% 
% * 双曲临界点和稳定流形定理
% * 稳定、不稳定流形的维数
% * Morse–Smale 扰动
% * 轨线空间及平移群商
% 
% 纲要内容结束

\section{梯度流的局部动力学}
\label{g2:sec:1801}

上一章的胞腔边界由把手的附着映射定义。若要用下降轨线直接计算边界，首先必须回答：从一个临界点出发的轨线有多少自由度，能够流向另一个临界点的条件又消去多少自由度？稳定与不稳定流形给出这两组参数；横截性保证相交后的维数等于预期值。最后还要删去沿同一条轨线移动时间原点的一个参数。

\subsection{双曲临界点和稳定流形定理}
\label{g2:subsec:1801:01}

在 Euclid 空间中，线性方程 \(\dot z=Az\) 的解为 \(e^{tA}z(0)\)。特征值实部为负的方向在正时间收缩，实部为正的方向在负时间收缩。若没有实部为零的特征值，两类方向便有严格的指数间隔。

\begin{definition}\label{g2:def:hyperbolic-zero}
	设光滑向量场 \(X\) 在 \(p\) 处为零。若线性化 \(\mathrm D_pX:\mathrm T_pM\to\mathrm T_pM\) 的复特征值都不在虚轴上，就称 \(p\) 为\term[shuangqulingdian]{双曲零点}[hyperbolic zero]。实部为负、为正的广义特征空间分别记为 \(E^s_p,E^u_p\)，称为稳定与不稳定子空间。
\end{definition}

线性化在零点处按共轭变换，与第16章使用的 \(\mathrm D_pX\) 是同一对象。对负梯度 \(X=-\operatorname{grad}_g f\)，在 Morse 临界点处有
\[
	\mathrm D_pX=-G^{-1}H,
\]
其中 \(G\) 为正定度量矩阵，\(H\) 为 Hessian 矩阵。它与对称矩阵 \(-G^{-1/2}HG^{-1/2}\) 相似，所以全部特征值实且非零。若临界点的 Morse 指标为 \(\lambda\)，正特征值有 \(\lambda\) 个，负特征值有 \(n-\lambda\) 个。

\begin{theorem}[局部稳定流形定理]\label{g2:thm:hyperbolic-stable-manifold}
	设 \(X\) 是光滑流形上的光滑向量场，\(p\) 为双曲零点。则在 \(p\) 的邻域内存在光滑嵌入圆盘 \(W^s_{\mathrm{loc}}(p)\)、\(W^u_{\mathrm{loc}}(p)\)，分别切于 \(E^s_p,E^u_p\)。充分靠近 \(p\)、在所有正时间内留在该邻域的轨线，恰从局部稳定流形出发，并指数趋于 \(p\)；负时间的对应结论对局部不稳定流形成立。这里的“局部”允许缩小圆盘及用于控制轨线的邻域。
\end{theorem}
\begin{proof}
	给出稳定流形的构造；将时间反向便得到另一部分。取以 \(p\) 为零的坐标，写
	\[
		X(z)=Az+R(z),\qquad R(0)=0,\qquad \mathrm DR(0)=0.
	\]
	记 \(P_s,P_u\) 为沿不变分解 \(E^s\oplus E^u\) 的投影。由矩阵指数的 Jordan 形式，可选 \(a>0,C\ge1\)，使
	\[
		\lVert e^{tA}P_s\rVert\le Ce^{-at},\qquad
		\lVert e^{-tA}P_u\rVert\le Ce^{-at}\quad(t\ge0).
	\]
	多项式因子可吸收到稍慢的指数中。将 \(R\) 乘以小球上的截断函数，可在整个坐标空间得到全局 Lipschitz 函数，其 Lipschitz 常数 \(L\) 任意小，且在更小球内仍是原来的 \(R\)。这里 \(\mathrm DR(0)=0\) 同时控制 \(R\) 的导数和截断产生的项。

	固定 \(0<\beta<a\)，考虑连续路径的完备赋范空间
	\[
		\mathcal B_\beta=
		\left\{z:[0,\infty)\to E^s\oplus E^u:
		\lVert z\rVert_\beta=\sup_{t\ge0}e^{\beta t}|z(t)|<\infty\right\}.
	\]
	给定稳定初值 \(\xi\in E^s\)，把寻找轨线改写成寻找下式的不动点：
	\begin{equation}\label{g2:eq:stable-integral-equation}
		\begin{aligned}
			(\mathcal T_\xi z)(t)
			={}&e^{tA}\xi+\int_0^t e^{(t-s)A}P_sR(z(s))\,\mathrm ds\\
			&-\int_t^\infty e^{(t-s)A}P_uR(z(s))\,\mathrm ds.
		\end{aligned}
	\end{equation}
	第一项把稳定初值向前推进；最后一项从未来施加“没有增长分量”的条件。这正是决定不稳定初值所需的额外约束。两项积分分别满足
	\[
		\lVert\mathcal T_\xi z-\mathcal T_\xi w\rVert_\beta
		\le CL\left(\frac1{a-\beta}+\frac1{a+\beta}\right)
		\lVert z-w\rVert_\beta.
	\]
	例如第一项积分的加权估计为
	\[
		e^{\beta t}\int_0^t Ce^{-a(t-s)}L e^{-\beta s}\,\mathrm ds
		\le \frac{CL}{a-\beta}.
	\]
	选 \(L\) 使总常数 \(\kappa<1\)，逐次迭代 \(\mathcal T_\xi\) 即得唯一不动点 \(z_\xi\)，并有
	\[
		\lVert z_\xi\rVert_\beta\le\frac{C}{1-\kappa}|\xi|.
	\]
	取 \(\xi\) 足够小，这条路径始终处在原向量场未被截断的球中。对积分式求导可得 \(\dot z_\xi=X(z_\xi)\)，而 \(P_sz_\xi(0)=\xi\)。于是初值集合是图像
	\[
		\xi\longmapsto \xi+h(\xi),\qquad
		h(\xi)=-\int_0^\infty e^{-sA}P_uR(z_\xi(s))\,\mathrm ds.
	\]

	还需说明图像光滑，而不只连续。截断后 \(R\) 的各阶导数有界；在上述加权范数中，对积分逐次求导的估计仍由指数积分控制。\(\mathrm D_z\mathcal T_\xi\) 的算子范数小于 \(\kappa\)，因此
	\[
		(I-\mathrm D_z\mathcal T_\xi)^{-1}
		=\sum_{j=0}^\infty(\mathrm D_z\mathcal T_\xi)^j
	\]
	一致收敛。对不动点等式作差商，再移项乘上这个逆，便得 \(z_\xi\) 对 \(\xi\) 的一阶导数；迭代同一步骤得到所有阶导数及其连续性。故 \(h\) 光滑。由于 \(\mathrm DR(0)=0\)，在零路径处求导得到 \(\mathrm Dh(0)=0\)，图像切于 \(E^s\)。

	最后验证刻画的唯一性。若原方程的一条正向轨线始终留在足够小球内，它是有界路径。变参数公式在不稳定分量上从 \(t\) 积分至 \(T\)，再令 \(T\to\infty\)；因 \(e^{(t-T)A}P_uz(T)\to0\)，得到\eqref{g2:eq:stable-integral-equation}。同一算子在普通有界连续路径范数（即 \(\beta=0\)）中也为收缩，可预先将 \(L\) 取得同时满足这一要求。因此该有界路径必等于已构造的 \(z_\xi\)，从而指数衰减并从图像出发。流的唯一性保证图像在轨线仍处于局部邻域期间不变。这完成了稳定流形的构造与刻画。
\end{proof}

非线性稳定流形一般不是线性子空间。例如 \(\dot u=u+v^2,\ \dot v=-v\) 的稳定流形为 \(u=-v^2/3\)，它在零处切于 \(v\) 轴；后面的习题将直接解出这一图像。线性化确定切方向，积分方程则决定曲面本身怎样弯曲。

\subsection{稳定、不稳定流形的维数}
\label{g2:subsec:1801:02}

本章在闭流形上进行计数。紧性保证光滑向量场的流 \(\Phi_t\) 对所有 \(t\in\mathbb R\) 存在，这里使用第一卷第13章第2节的紧流形完备性。对零点 \(p\)，定义
\[
	W^s(p)=\{x:\lim_{t\to+\infty}\Phi_t(x)=p\},\qquad
	W^u(p)=\{x:\lim_{t\to-\infty}\Phi_t(x)=p\}.
\]
它们分别称为\term[wendingliuxing]{稳定流形}[stable manifold]与\term[buwendingliuxing]{不稳定流形}[unstable manifold]。\symidx{\(W^s(p),W^u(p)\)：稳定与不稳定流形}

\begin{proposition}\label{g2:prop:morse-stable-unstable}
	设 \(M^n\) 为闭光滑流形，\(f:M\to\mathbb R\) 为 Morse 函数，\(X=-\operatorname{grad}_g f\)。对指标为 \(\lambda(p)\) 的临界点 \(p\)，\(W^s(p)\)、\(W^u(p)\) 具有自然的单射浸入流形结构，维数分别为 \(n-\lambda(p)\)、\(\lambda(p)\)。在 \(p\) 处，它们的切空间分别是 \(\mathrm D_pX\) 的稳定与不稳定子空间。
\end{proposition}
\begin{proof}
	上一小节已经证明 \(p\) 双曲，并计算了两种线性子空间的维数。局部稳定图像与其各个向后流像
	\[
		\Phi_{-t}(W^s_{\mathrm{loc}}(p)),\qquad t\ge0,
	\]
	覆盖 \(W^s(p)\)：趋于 \(p\) 的轨线在充分大时间后留在局部邻域。局部图像上的稳定坐标满足 \(\dot\xi=A_s\xi+O(|\xi|^2)\)。正定二次型
	\[
		Q(\xi)=\int_0^\infty |e^{tA_s}\xi|^2\,\mathrm dt
	\]
	沿该方程的导数为 \(-|\xi|^2+O(|\xi|^3)\)，所以足够小的 \(Q\) 子水平椭球给出正向不变的局部圆盘。因而可只取整数时间 \(t=0,1,2,\ldots\)，得到递增的圆盘族。相邻圆盘的重叠转换由光滑流给出，其微分可逆；因此这些图拼成同维光滑流形，包含到 \(M\) 的映射为单射浸入。对反向流作相同构造得到 \(W^u(p)\)。切空间与维数来自局部稳定流形定理。
\end{proof}

为使断裂点附近的粘合能够写成具体方程，以下选择在 Morse 坐标中具有精确线性表达的下降场。先在\cref{g2:thm:morse-approximation}给出的 Morse 函数上，用\cref{g2:prop:distinct-morse-critical-values}分开临界值。

\begin{definition}\label{g2:def:normalized-morse-data}
	设 \(M\) 为闭光滑流形，\(f\) 为临界值互异的 Morse 函数。若光滑向量场 \(X\) 满足 \(Xf<0\) 于非临界点，并且在每个临界点的某张 Morse 坐标中有
	\begin{equation}\label{g2:eq:normalized-negative-gradient}
		f=f(p)-|u|^2+|v|^2,\qquad
		X(u,v)=(2u,-2v),
		\quad u\in\mathbb R^{\lambda(p)},\ v\in\mathbb R^{n-\lambda(p)},
	\end{equation}
	就把 \((f,X)\) 称为本章的正规化 Morse 数据。
\end{definition}

这种数据总能选到。取各临界点互不相交的 Morse 坐标邻域，在其中规定 Euclid 度量，再用单位分解与外部度量拼接，保持更小邻域内的度量不变。所得 \(X=-\operatorname{grad}_g f\) 满足\eqref{g2:eq:normalized-negative-gradient}。反过来，定义中的每个 \(X\) 也确实是某个光滑度量的负梯度：在非临界区域令 \(\alpha=-\mathrm df\)，则 \(\alpha(X)>0\)，并有分解
\[
	\mathrm TM=\mathbb RX\oplus\ker\alpha.
\]
在 \(\ker\alpha\) 上任选正定内积 \(h\)，规定
\[
	g(aX+v,bX+w)=ab\,\alpha(X)+h(v,w).
\]
于是 \(g(X,\cdot)=\alpha=-\mathrm df\)。临界点附近的 Euclid 度量满足同一等式，用单位分解拼接时此线性等式仍保持，得到所需全局光滑度量。

\ProseParagraphBreak

对正规化数据，局部稳定与不稳定圆盘分别为 \(u=0\)、\(v=0\)，轨线的表达为
\[
	u(t)=e^{2t}u(0),\qquad v(t)=e^{-2t}v(0).
\]
这项选择将高阶项放到临界点以外，后面可直接用停留时间 \(T\) 与参数 \(e^{-2T}\) 描述经过临界邻域的长轨线。第18.4节将通过延拓映射比较不同的数据选择。

\ProseParagraphBreak

此时全局稳定流形还可明确参数化为 \(\mathbb R^{n-\lambda(p)}\)。取局部稳定球面 \(|v|=\varepsilon\)，对 \(v\ne0\) 定义
\[
	\iota_s(v)=\Phi_{-\frac12\log(|v|/\varepsilon)}
	\left(0,\varepsilon\frac{v}{|v|}\right),\qquad \iota_s(0)=p.
\]
当 \(|v|\le\varepsilon\) 时，这就是局部坐标点 \((0,v)\)；当 \(|v|>\varepsilon\) 时，则沿完整流向后延拓。每条非恒定稳定轨线恰好经过局部球面一次，故它给出全局一一参数化，逆映射由光滑的到达时间给出。零维情形直接为一个点。反向时间同样参数化 \(W^u(p)\cong\mathbb R^{\lambda(p)}\)。这里采用的是前述自然浸入流形结构；特别地，这两个流形都可缩，后面可以在它们上面连续选择方向。

\subsection{Morse–Smale 扰动}
\label{g2:subsec:1801:03}

从 \(p\) 出发并流向 \(q\) 的非恒定轨线，其每个点都在 \(W^u(p)\cap W^s(q)\) 中。两个流形若相切，交集可能有多余维数，计数也会对扰动不稳定。因此需要第14章的横截条件。

\begin{definition}\label{g2:def:morse-smale-data}
	闭光滑流形上的正规化 Morse 数据 \((f,X)\) 满足下列横截性时，称它满足\term[Morse-Smale tiaojian]{Morse--Smale 条件}[Morse--Smale condition]：
	\[
		W^u(p)\pitchfork W^s(q)\qquad
		\text{对全部临界点 }p,q,
	\]
	条件在相交点处的含义是两切空间之和为整个 \(\mathrm T_xM\)；对浸入流形用其上述自然流形结构判断。
\end{definition}

\begin{theorem}[Morse--Smale 扰动]\label{g2:thm:morse-smale-existence}
	设 \((f,X_0)\) 是闭光滑流形上的正规化 Morse 数据，并在各临界点选定更小的、两两不交的闭坐标邻域。则 \(X_0\) 的每个 \(C^\infty\) 邻域中都有光滑向量场 \(X\)，使 \((f,X)\) 满足 Morse--Smale 条件，且 \(X=X_0\) 于这些闭邻域的某个固定邻域。函数 \(f\) 及全部临界点均保持不变。
\end{theorem}
\begin{proof}
	把问题化为有限时间的流映射横截性，这样每次扰动只需有限个参数。对每个临界点 \(p\)，在固定线性邻域内部取局部出球面
	\[
		S_p^-=\{(u,0):|u|=\varepsilon_p\},
	\]
	维数为 \(\lambda(p)-1\)；对每个 \(q\) 取入球面
	\[
		S_q^+=\{(0,v):|v|=\varepsilon_q\},
	\]
	维数为 \(n-\lambda(q)-1\)。空维球面取空集。任何从 \(p\) 到 \(q\) 的非恒定轨线，都恰好经过这两个球面一次。考察
	\[
		E_X:S_p^-\times(0,\infty)\longrightarrow M,\qquad
		E_X(a,T)=\Phi_T^X(a).
	\]
	条件 \(E_X\pitchfork S_q^+\) 等价于 \(W^u(p)\pitchfork W^s(q)\)：\(E_X\) 的微分包含不稳定方向及流方向，而在入球面上
	\(\mathrm T W^s(q)=\mathrm T S_q^+\oplus\mathbb RX\)。
	同一流方向已经在 \(E_X\) 的像内，故两种切空间满秩条件相同。\(p=q\) 的情形在临界点由正、负子空间互补自动成立，且严格下降排除了非恒定同宿轨线。

	先只要求在紧集 \(K_N=S_p^-\times[1/N,N]\) 上横截，\(N\) 为正整数。若 \(E_{X_0}(a,T)\in S_q^+\)，该连接段在两临界邻域之间必经过一个允许扰动的小流盒。对支于该流盒的场 \(Y\)，终点关于 \(X_\theta=X_0+\theta Y\) 的一阶变化为
	\begin{equation}\label{g2:eq:flow-endpoint-variation}
		\left.\frac{\partial}{\partial\theta}\right|_0
		\Phi_T^{X_\theta}(a)
		=\int_0^T
		\mathrm d\Phi_{T-s}^{X_0}\bigl(Y(\Phi_s^{X_0}(a))\bigr)\,\mathrm ds.
	\end{equation}
	这由第一卷第13章第1节的参数光滑依赖及其变分方程得到：对流方程求参数导数，再用基本解矩阵积分。选一小段时间上的光滑截断函数，其积分为一；对任意目标向量 \(w\in\mathrm T_{\Phi_T(a)}M\)，沿这段轨线规定
	\[
		Y(\Phi_s(a))=\rho(s)\,
		(\mathrm d\Phi_{T-s})^{-1}w.
	\]
	用流盒坐标横向延拓并加截断，便得到允许的光滑场，使\eqref{g2:eq:flow-endpoint-variation}恰等于 \(w\)。因 \(f\) 严格下降，流盒可选在一个小高度带中，该轨线不会再次经过其支集。因此有限个这样的场就能在此约束点补足全部终点方向。

	约束点集 \(K_N\cap E_{X_0}^{-1}(S_q^+)\) 紧。取有限覆盖，收集上述有限个场 \(Y_1,\ldots,Y_L\)，得到
	\[
		(a,T,\theta)\longmapsto
		\Phi_T^{X_0+\sum_i\theta_iY_i}(a).
	\]
	在足够小的参数球内，它在 \(K_N\) 附近与 \(S_q^+\) 横截。原来避开 \(S_q^+\) 的其余紧部分仍避开它，因 \(S_q^+\) 闭。由\cref{g2:thm:parametric-transversality}，任意小参数中都有使切片在 \(K_N\) 上横截的选择。充分小的场扰动还保持 \(Xf<0\)：固定临界邻域之外是紧集，原来的 \(-X_0f\) 在那里有正下界。临界邻域内的场完全未改。

	现在同时实现全部条件。临界点对只有有限多个，而 \(N=1,2,\ldots\)，故可把这些紧集横截要求排成可数序列。每项要求由刚才的构造在允许场空间中稠密；由有限时间流的一阶连续依赖及\cref{g2:prop:transversality-open-on-compact}，它又是开条件。逐项选光滑场 \(X_j\)，使第 \(j\) 项及先前各项成立，并把第 \(j\) 次改动的 \(C^j\) 范数取得小于 \(2^{-j}\) 倍的指定容差。每次再为已满足的有限个条件选择保持它们的开邻域，分配后续改动的总误差小于这些邻域的容差。还把所有改动的总 \(C^0\) 误差限制在保持下降条件的界内，并满足最初给定的 \(C^\infty\) 邻域中的有限阶控制。

	于是 \(X_j\) 在每个 \(C^r\) 范数中收敛；在有限坐标覆盖上逐阶取极限，得到光滑 \(X\)。全部紧集横截条件由各自预留的余量保持，固定邻域内的表达也未变。任意有限时间的连接轨线属于某个 \(K_N\)，所以所有不稳定与稳定流形都横截。这便得到所需的 \(X\)。
\end{proof}

这里的扰动改变临界点之间的轨线走向，同时完整保留临界点附近的线性模型。这两种控制将在后面同时使用：横截性给粘合方程的可逆线性部分，局部线性模型给小参数的精确表达。

\subsection{轨线空间及平移群商}
\label{g2:subsec:1801:04}

对不同临界点 \(p,q\)，设 \(\widetilde{\mathcal M}(p,q)\) 为满足
\[
	\dot\gamma(t)=X(\gamma(t)),\qquad
	\lim_{t\to-\infty}\gamma(t)=p,\qquad
	\lim_{t\to+\infty}\gamma(t)=q
\]
的非恒定参数化轨线。由流的唯一性，取值 \(\gamma\mapsto\gamma(0)\) 将其识别为 \(W^u(p)\cap W^s(q)\)。时间平移
\[
	(\tau\cdot\gamma)(t)=\gamma(t+\tau)
\]
只改变零时刻，不改变下降路径。将这一冗余除去，得到\term[guixianmolikongjian]{轨线模空间}[moduli space of trajectories]
\[
	\mathcal M(p,q)=\widetilde{\mathcal M}(p,q)/\mathbb R.
\]
本章的计数都在这个商空间中进行。\symidx{\(\mathcal M(p,q)\)：无参数下降轨线空间}

\begin{proposition}\label{g2:prop:trajectory-quotient}
	设 \((f,X)\) 是闭光滑 \(n\) 维流形上的正规化 Morse--Smale 数据，\(p\ne q\) 为临界点。若 \(\mathcal M(p,q)\ne\varnothing\)，则 \(f(p)>f(q)\)、\(\lambda(p)>\lambda(q)\)，而 \(\mathcal M(p,q)\) 是维数
	\[
		\dim\mathcal M(p,q)=\lambda(p)-\lambda(q)-1
	\]
	的光滑流形。任取 \(f(q)<a<f(p)\) 中的正则值，取唯一交点给出微分同胚
	\begin{equation}\label{g2:eq:trajectory-level-section}
		\mathcal M(p,q)\cong
		W^u(p)\cap W^s(q)\cap f^{-1}(a).
	\end{equation}
	时间平移作用自由且固有。
\end{proposition}
\begin{proof}
	非恒定轨线上 \(\frac{\mathrm d}{\mathrm dt}f(\gamma(t))=Xf<0\)，所以两端临界值严格下降。\cref{g2:thm:transverse-preimage}的交集维数公式给出
	\[
		\dim\widetilde{\mathcal M}(p,q)
		=\lambda(p)+(n-\lambda(q))-n
		=\lambda(p)-\lambda(q).
	\]
	非零流向量 \(X\) 同时切于这两张流形，故交集的非空维数至少为一，得到指标严格下降。沿交集，\(\mathrm df(X)<0\)，所以 \(f=a\) 再切去一维。

	轨线上 \(f\) 连续且严格递减，两端值分别为 \(f(p),f(q)\)，故每条轨线恰有一个点在水平 \(a\) 上。记这个截面为 \(S\)。由到达方程 \(f(\Phi_t(x))=a\) 及其时间导数 \(Xf\ne0\)，隐函数定理给出到达时间对 \(x\) 的光滑依赖。因此
	\[
		S\times\mathbb R\longrightarrow\widetilde{\mathcal M}(p,q),
		\qquad(z,t)\longmapsto\Phi_t(z)
	\]
	为微分同胚；逆映射把一个点送回其唯一截面点，并记录所用时间。于是平移作用对应第二因子上的普通平移，既自由又固有：若初、末两点在两个指定紧集中，它们的时间坐标之差有界，截面坐标也在紧集中，作用映射的相应原像为紧集。取商就得到 \(S\) 的光滑结构及\eqref{g2:eq:trajectory-level-section}。换另一正则水平时，过渡映射是流的首次到达映射，仍光滑，故这个结构不依赖所选水平。
\end{proof}

对于不同临界点，模空间非空的必要条件是指标至少下降一。指标相差一时，商空间为零维；相差二时，它为一维。零维目前只说明每条轨线在附近孤立，是否总共有有限条还要使用下一节的断裂紧性。

\begin{example}\label{g2:exm:sphere-meridian-moduli}
	标准 \(S^n\) 的高度函数只有最高点 \(p\) 与最低点 \(q\)。其负梯度轨线是由北向南的子午线；除去时间平移后，每条子午线由赤道上的唯一交点确定。因此
	\[
		\mathcal M(p,q)\cong S^{n-1}.
	\]
	对 \(n=1\)，它是两条半圆轨线组成的两点集；对 \(n=2\)，它是一圈轨线。将度量在极点附近改成正规化度量，并保持径向方向，可保留这组无参数轨线。这个例子也显示，若不除去平移，即使只有一条几何路径，仍会出现一整条 \(\mathbb R\) 的参数。
\end{example}

% !TeX root = ../../Book2.tex
% 纲要标题：### 18.2 轨线序列和断裂紧化
% 纲要位置：几何学计划纲要.md，第 1706 行。
% 正文状态：初稿；尚未正式审核。

% 纲要内容开始
% 
% * 能量恒等式与临界值下降
% * 选定时间平移后的局部收敛
% * 有限次断裂和能量控制
% * 零维紧性及计数有限性
% 
% 纲要内容结束

\section{轨线序列和断裂紧化}
\label{g2:sec:1802}

轨线可以越来越靠近一个中间临界点，在那里停留越来越久，再继续下降。每一条轨线都从 \(p\) 到 \(q\)，极限却可能要由“从 \(p\) 到 \(r\)”和“从 \(r\) 到 \(q\)”两条轨线合起来描述。要把轨线数用作链复形的系数，首先必须知道这种退化能发生多少次，以及是否还有别的轨线逃失方式。本节利用下降的函数值为轨线选定截点，把紧性问题化成有限个紧水平集中的收敛问题。

\subsection{能量恒等式与临界值下降}
\label{g2:subsec:1802:01}

固定闭光滑流形 \(M\)、Morse 函数 \(f\) 和光滑 Riemann 度量 \(g\)，令 \(X=-\operatorname{grad}_g f\)。前节所选的正规化数据满足这些条件。闭性保证 \(X\) 的流对所有实时间存在：第一卷第13章第2节的紧支完备性定理适用，因为此时整个流形就是一个紧支集。

\begin{proposition}\label{g2:prop:morse-energy-limits}
	设 \(M\) 为闭光滑流形，\(f:M\to\mathbb R\) 为 Morse 函数，\(g\) 为光滑 Riemann 度量。任何非恒定解
	\[
		\dot\gamma(t)=-\operatorname{grad}_g f(\gamma(t)),\qquad t\in\mathbb R,
	\]
	都存在临界点 \(p,q\)，使 \(\gamma(t)\to p\) 当 \(t\to-\infty\)，\(\gamma(t)\to q\) 当 \(t\to+\infty\)。而且 \(f(p)>f(q)\)，并有
	\begin{equation}\label{g2:eq:morse-energy-identity}
		E(\gamma)\coloneqq\int_{-\infty}^{+\infty}|\dot\gamma(t)|_g^2\,\mathrm{d}t
		=f(p)-f(q).
	\end{equation}
\end{proposition}
\begin{proof}
	梯度的定义给出
	\[
		\frac{\mathrm{d}}{\mathrm{d}t}f(\gamma(t))
		=\mathrm{d}f_{\gamma(t)}(\dot\gamma(t))
		=-|\operatorname{grad}_g f(\gamma(t))|_g^2.
	\]
	积分于有限区间 \([a,b]\)，得到
	\[
		\int_a^b|\dot\gamma|_g^2\,\mathrm{d}t
		=f(\gamma(a))-f(\gamma(b)).
	\]
	\(f\) 在紧流形上有界，因此两端的函数值极限存在，且正、负半轴上的能量均有限。还要证明轨线本身收敛，而不只证明其函数值收敛。

	考察正向极限集
	\[
		\omega(\gamma)=\bigcap_{T\ge0}\overline{\{\gamma(t):t\ge T\}}.
	\]
	其中每个集合都是非空紧连通集，且随 \(T\) 增大而递减；因此交也是非空紧连通集。若 \(x\in\omega(\gamma)\) 不是临界点，就可取 \(x\) 的邻域 \(V\)，以及稍大的邻域 \(U\)，使 \(|\operatorname{grad}f|\ge\delta>0\) 于 \(U\)。流的连续性和 \(\overline V\subset U\) 的紧性给出统一的 \(h>0\)：从 \(V\) 出发的轨线在时间 \([0,h]\) 内留在 \(U\)。轨线无限次进入 \(V\)；选取间隔大于 \(h\) 的进入时刻，每次在随后长度 \(h\) 的区间耗费至少 \(h\delta^2\) 的能量，与正向能量有限矛盾。故 \(\omega(\gamma)\) 全由临界点组成。

	Morse 临界点孤立，在紧流形上只有有限个。有限集的非空连通子集只有一个点，因此 \(\omega(\gamma)=\{q\}\)。若轨线不趋于 \(q\)，可取一个离 \(q\) 有正距离的点列，并由紧性取得另一个正向聚点，矛盾。负向极限同理为某个临界点 \(p\)。令有限区间的两端趋于无穷，即得能量恒等式。非恒定轨线在任何有限时刻均不经过临界点，否则由常微分方程解的唯一性，它将恒定；故其能量严格为正，\(f(p)>f(q)\)。
\end{proof}

能量在本章就是下降掉的函数值。它控制的是整个过程的总变化量，并不控制需要多长时间：在临界点附近，梯度很小，轨线可以用很长时间下降很少的高度。这正是断裂极限的来源。

\begin{example}
	在局部标准模型 \(f=c-u^2+v^2\) 中，下降方程为 \(\dot u=2u\)、\(\dot v=-2v\)，所以
	\[
		u(t)=e^{2t}u(0),\qquad v(t)=e^{-2t}v(0).
	\]
	从 \(v=\varepsilon\) 的截面进入，且 \(u(0)=\eta>0\)，到达 \(u=\varepsilon\) 所需的时间为
	\[
		T=\frac12\log\frac{\varepsilon}{\eta}.
	\]
	当 \(\eta\to0\) 时，\(T\to\infty\)；前半段趋于沿 \(v\) 轴进入原点的轨线，后半段经重新选定时间零点后趋于沿 \(u\) 轴离开原点的轨线。原点是两个极限段的共同端点。
\end{example}

\subsection{选定时间平移后的局部收敛}
\label{g2:subsec:1802:02}

非参数轨线没有指定的时间零点。若任意选择参数，同一条轨线作越来越大的平移，便可能在每个有限时间区间都趋于一个恒定临界点。为保留一段真正的运动，应先让轨线在一个固定正则水平上经过时间零点。

\begin{lemma}\label{g2:lem:anchored-trajectory-convergence}
	设 \(M\) 闭，\(X=-\operatorname{grad}_g f\)，其中 \(f\) 为 Morse 函数。设 \(a\) 为正则值，\(\gamma_\nu\) 为满足 \(f(\gamma_\nu(0))=a\) 的下降轨线。则可取子列及一条非恒定下降轨线 \(\gamma\)，使
	\[
		\gamma_\nu\longrightarrow\gamma
		\quad\text{于每个有界时间区间，按 }C^\infty\text{ 收敛}.
	\]
	若所有 \(\gamma_\nu\) 都连接同一对临界点 \(p,q\)，则极限轨线的两个端点 \(p',q'\) 满足
	\[
		f(p)\ge f(p')>a>f(q')\ge f(q).
	\]
\end{lemma}
\begin{proof}
	水平集 \(f^{-1}(a)\) 紧，故可使 \(x_\nu=\gamma_\nu(0)\to x\in f^{-1}(a)\)。设 \(\Phi_t\) 为 \(X\) 的流，则 \(\gamma_\nu(t)=\Phi_t(x_\nu)\)。第一卷第13章第1节的光滑依赖定理说明 \((t,x)\mapsto\Phi_t(x)\) 光滑。限制到任意紧时间区间后，所有阶导数都连续，因而 \(\Phi_t(x_\nu)\to\Phi_t(x)\) 按 \(C^\infty\) 收敛。\(a\) 为正则值，故 \(X(x)\ne0\)，极限轨线非恒定。

	若原轨线连接 \(p,q\)，对每个固定 \(t\) 都有 \(f(q)\le f(\gamma_\nu(t))\le f(p)\)。取极限，再让 \(t\to\pm\infty\)，并用\cref{g2:prop:morse-energy-limits}，便得到所述端点不等式。
\end{proof}

这项引理保留了穿过指定水平的那一段轨线。它的端点可能已经改变，例如原来的 \(p\to q\) 轨线在极限中变成了 \(p'\to q'\)。其余下降过程要到别的水平上重新选定时间零点才能看见。

\ProseParagraphBreak

下面记录一个连接相邻截点的判断。取正则值 \(a>b\)，设 \(\gamma_\nu(0)=x_\nu\in f^{-1}(a)\)，并记其经过 \(f^{-1}(b)\) 的时间为 \(\tau_\nu>0\)，对应点为 \(y_\nu\)。若 \(x_\nu\to x\)、\(y_\nu\to y\)，则：

\begin{itemize}
	\item 若 \(\tau_\nu\) 有界，取子列使其趋于 \(\tau\)，便有 \(y=\Phi_\tau(x)\)，两截点属于同一条极限轨线。
	\item 若 \(\tau_\nu\to\infty\)，则从 \(x\) 出发的正向轨线始终留在 \(f\ge b\)，而从 \(y\) 出发的负向轨线始终留在 \(f\le a\)。若 \([b,a]\) 内恰有一个临界点 \(r\)，两条轨线相应的极限端点就都只能是 \(r\)。
\end{itemize}

第二点的“始终”可逐个固定时间检验：对任意 \(T>0\)，当 \(\tau_\nu>T\) 时，\(f(\Phi_T(x_\nu))\ge b\)，取极限即可；负向论证相同。随后由上一命题，两端必须收敛到带中的临界点。正则端水平不含临界点，带中又只有 \(r\)，因此两段确实在同一个 \(r\) 接合。

\subsection{有限次断裂和能量控制}
\label{g2:subsec:1802:03}

设 \(p\ne q\) 为临界点，且 \(f(p)>f(q)\)。一个从 \(p\) 到 \(q\) 的\term[duanlie guixian]{断裂轨线}[broken trajectory]是有限序列
\[
	\boldsymbol\gamma=(\gamma_1,\ldots,\gamma_\ell),\qquad
	\gamma_i\in\mathcal M(p_{i-1},p_i),\qquad
	p_0=p,\quad p_\ell=q,
\]
其中各段均为非恒定非参数轨线。它的临界值严格下降：
\[
	f(p_0)>f(p_1)>\cdots>f(p_\ell).
\]
两条相接轨线各自拥有时间平移的自由，故断裂轨线记录的是各段的非参数类，没有额外的共同时间参数。未断裂轨线对应 \(\ell=1\)。

\begin{theorem}[断裂轨线紧性]\label{g2:thm:broken-trajectory-compactness}
	设 \(M\) 为闭光滑流形，\(f\) 为临界值互异的 Morse 函数，\(X=-\operatorname{grad}_g f\)。设 \(p,q\) 为满足 \(f(p)>f(q)\) 的临界点。从 \(p\) 到 \(q\) 的下降轨线序列都存在一个有限断裂极限：可取子列、临界点链 \(p=p_0,\ldots,p_\ell=q\) 及时间平移 \(s_{\nu,i}\)，使
	\[
		\gamma_\nu(\,\cdot+s_{\nu,i})\longrightarrow\gamma_i
		\quad\text{于紧时间区间按 }C^\infty\text{ 收敛},
		\qquad
		\gamma_i\in\mathcal M(p_{i-1},p_i).
	\]
	相邻两段的平移满足 \(s_{\nu,i+1}-s_{\nu,i}\to+\infty\)，而能量完整分解为
	\begin{equation}\label{g2:eq:broken-energy-sum}
		E(\gamma_\nu)=f(p)-f(q)=\sum_{i=1}^{\ell}E(\gamma_i).
	\end{equation}
	全部断裂轨线组成一个紧 Hausdorff 空间；其拓扑可由下述有限组正则水平上的交点确定。
\end{theorem}
\begin{proof}
	列出 \(f(q)\) 与 \(f(p)\) 之间的全部临界值，包括端点：
	\[
		c_0=f(p)>c_1>\cdots>c_m=f(q).
	\]
	每个区间 \((c_{j+1},c_j)\) 选一个正则值 \(a_j\)，令 \(L_j=f^{-1}(a_j)\)。每条 \(p\to q\) 轨线恰好与每个 \(L_j\) 相交一次，记交点为 \(x_{\nu,j}\)。有限乘积 \(\prod_{j=0}^{m-1}L_j\) 紧；取共同子列，使所有 \(x_{\nu,j}\to x_j\)。再对有限个相邻水平间的通过时间取子列，使每个时间或者趋于有限值，或者趋于正无穷。

	在 \(x_{\nu,j}\) 处把时间平移为零，由\cref{g2:lem:anchored-trajectory-convergence}得到穿过 \(x_j\) 的完整极限轨线。如果两个相邻水平间的通过时间有限，上面关于相邻截点的判断说明，这两条极限轨线只是同一轨线的不同参数选择，应合并为一段。如果通过时间趋于无穷，两水平之间恰有临界值 \(c_{j+1}\)，且因临界值互异只有一个临界点；同一判断说明，上段的正端点与下段的负端点就是这个点。最上段的负端点只能为 \(p\)：它的函数值介于 \(a_0\) 与 \(f(p)\) 之间，而这一区域仅有临界点 \(p\)。最下段同理以 \(q\) 结束。于是全部非重复段接成所需的临界点链，而不只是互不相关的局部极限。

	为每一段选它穿过的一个 \(L_j\)，以原轨线到达该水平的时刻为 \(s_{\nu,i}\)。跨过相邻两段接点时，通过时间已经趋于无穷，故相邻平移之差也趋于无穷。段数至多为 \(m\)。若将全部不同临界值之间的最小正差记为 \(\Delta\)，还可写出
	\[
		\ell\Delta\le f(p)-f(q).
	\]
	每段的能量由其两端临界值之差给出，有限求和逐项抵消，便得到\eqref{g2:eq:broken-energy-sum}。这条等式说明未被任何局部极限段捕获的能量为零；长时间停留本身没有另带一份正能量。

	最后给出整个断裂空间的拓扑并核实其紧性。任一断裂链都在每个 \(L_j\) 上恰有一个交点，故有评价映射
	\[
		\operatorname{ev}:\overline{\mathcal M}(p,q)
		\longrightarrow\prod_{j=0}^{m-1}L_j.
	\]
	这里 \(\overline{\mathcal M}(p,q)\) 暂表示所有有限断裂链的集合。该映射为单射：一条非恒定轨线由其任意一个点唯一决定，而每段连接两个不同临界值，必穿过所选水平中的至少一个；因而全部截点确定各段及它们的连接次序。用这个单射赋予断裂空间子空间拓扑。

	其像是闭集。事实上，给定一列断裂链及其截点的极限，可能出现的临界点链只有有限种；取子列使链的临界点序列固定。分别对每个非恒定段应用刚才的有限水平论证，使它进一步分裂成有限条轨线，再按原有共同端点串接。得到的仍是从 \(p\) 到 \(q\) 的有限断裂链，且其截点恰为给定极限。因此评价像序列闭。有限乘积是紧度量空间，序列闭即闭，所以断裂空间紧且 Hausdorff。

	在未断裂部分，这个拓扑与正则水平截面的通常拓扑一致：任取一处截点，附近的非断裂轨线到其他固定正则水平的到达时间，由流的光滑依赖和横截性随该截点连续变化；反过来乘积收敛当然蕴含该截点收敛。更换所选正则水平也不改变拓扑，因为同一相邻临界值间没有临界点，流给出两个水平的光滑对应。
\end{proof}

若数据还满足 Morse--Smale 条件，每个非恒定段都满足
\[
	\lambda(p_{i-1})-\lambda(p_i)\ge1,
\]
因为\cref{g2:prop:trajectory-quotient}给出的轨线空间维数为两端指标差减一，非空光滑流形不能有负维数。因此
\[
	\ell\le\lambda(p)-\lambda(q).
\]
能量的临界值下降给出有限性，横截性还把可出现的段数限制在指标差以内。这两种控制分别来自函数的下降和稳定、不稳定流形的相交。

\ProseParagraphBreak

\(\overline{\mathcal M}(p,q)\) 的记号记录了允许断裂后的紧空间。下一节将证明，指标差为二时，每个断裂点都有一条由真实未断轨线形成的半开邻域。因此在链复形所需的零维、一维情形，这个空间确实是原轨线空间的紧化，并具有确定的边界结构。

\begin{lemma}\label{g2:lem:broken-local-crossing}
	设 \((f,g)\) 为闭流形上的正规化 Morse--Smale 数据，未断轨线 \(\gamma_\nu\in\mathcal M(p,q)\) 在上述拓扑中趋于一断裂链 \((\alpha,\beta)\in\mathcal M(p,r)\times\mathcal M(r,q)\)。在 \(r\) 的正规化坐标中选充分小的闭箱
	\[
		B_\varepsilon=\{(u,v):|u|\le\varepsilon,\ |v|\le\varepsilon\},
	\]
	使 \(\alpha\) 沿稳定轴在 \((0,v_0)\) 进入 \(|v|=\varepsilon\)，\(\beta\) 沿不稳定轴在 \((u_0,0)\) 离开 \(|u|=\varepsilon\)。则对充分大的 \(\nu\)，\(\gamma_\nu\) 的相应进入点、离开点趋于 \((0,v_0)\)、\((u_0,0)\)，且这次通过箱子的时间 \(T_\nu\to+\infty\)。
\end{lemma}
\begin{proof}
	在 \(\alpha\) 所穿过的一个正则水平上固定时间零点。评价拓扑的收敛和\cref{g2:lem:anchored-trajectory-convergence}使原轨线在这个锚点附近的任意固定有限时间区间上趋于 \(\alpha\)。从锚点到 \((0,v_0)\) 的时间有限；到达截面处有
	\[
		X(|v|^2-\varepsilon^2)=-4\varepsilon^2\ne0.
	\]
	因此隐函数定理和流的光滑依赖使附近轨线的截面交点连续变化，得到进入点收敛。沿 \(\beta\) 从另一个水平截点向后追踪，同理得到离开点收敛。

	箱内方程为 \(u(t)=e^{2t}u(0)\)、\(v(t)=e^{-2t}v(0)\)。从进入面出发，\(|u|\) 增加而 \(|v|\) 减少，故只会从离开面离开。两次求得的交点属于同一次通过：进入面的函数值至少为 \(f(r)\)，离开面的函数值至多为 \(f(r)\)，而非恒定轨线严格下降，离开后不能再进入。设进入点的第一分量为 \(u_\nu\)。它不能为零，否则该轨线将一直趋于 \(r\) 而无法继续到 \(q\)；于是
	\[
		T_\nu=\frac12\log\frac{\varepsilon}{|u_\nu|}\longrightarrow+\infty,
	\]
	因为 \(u_\nu\to0\)。
\end{proof}

\begin{lemma}\label{g2:lem:continuation-broken-compactness}
	设 \(M\) 为闭光滑流形，\(K\) 为紧光滑参数空间，\((f_{t,\sigma},g_{t,\sigma})\) 为光滑族，其中 \(\sigma\in K\)。假设存在 \(T>0\)，使族在 \(t\le-T\) 和 \(t\ge T\) 分别等于固定的 Morse 数据 \((f_-,g_-)\)、\((f_+,g_+)\)，两端 Morse 函数的临界值分别互异。设 \(\sigma_\nu\to\sigma\)，且
	\[
		\dot\eta_\nu=-\operatorname{grad}_{g_{t,\sigma_\nu}}f_{t,\sigma_\nu}(\eta_\nu),
		\qquad
		\eta_\nu(-\infty)=p,\quad\eta_\nu(+\infty)=q,
	\]
	其中 \(p,q\) 是固定的两端临界点。则可取子列，得到一条参数为 \(\sigma\) 的中心延拓轨线 \(\eta\)，其两端为 \(r_-,r_+\)，以及从 \(p\) 到 \(r_-\) 的有限 \(f_-\) 断裂链和从 \(r_+\) 到 \(q\) 的有限 \(f_+\) 断裂链。允许任一端的链为空。中心轨线按紧时间区间 \(C^\infty\) 收敛，端部各段经时间平移后按相同意义收敛。能量有统一上界，并等于极限各段能量之和。
\end{lemma}
\begin{proof}
	取子列使 \(\eta_\nu(0)\) 在紧流形中收敛。带参数常微分方程的光滑依赖给出一条定义在整个实轴的极限解 \(\eta\)，以及紧时间区间上的 \(C^\infty\) 收敛。此处定义时间零点由外加的族 \(f_{t,\sigma}\) 固定，不能再将中心轨线任意平移。两端方程已经自治，故\cref{g2:prop:morse-energy-limits}在两条半轴上给出临界极限 \(r_-,r_+\)。

	在负端取所有位于 \(f_-(r_-)\) 与 \(f_-(p)\) 之间的正则水平，重复紧性定理的有限截点论证。过渡时间有限的截点合并为一段，发散的截点之间出现唯一的临界点。最后一段接到 \(r_-\)：若极限中心负尾非恒定，就再取一个位于其负端极限值以下、且由该尾穿过的正则水平，中心轨线的局部收敛确定这一截点；若负尾恒定于 \(r_-\)，则以收敛的终点 \(\eta_\nu(-T)\to r_-\) 代替最后的截点。若前一截点到此终点的时间有界，流的唯一性将迫使一条非恒定轨线在有限时间到达临界点，矛盾；所以最后的时间趋于无穷，端点确为 \(r_-\)。由此得到所述负端链。正端用相同论证，时间方向相反。两端各只有有限个临界值，故断裂次数有限。

	沿非自治解求导得到
	\begin{equation}\label{g2:eq:continuation-energy-identity}
		E(\eta_\nu)
		=f_-(p)-f_+(q)
		+\int_{-T}^{T}\partial_t f_{t,\sigma_\nu}(\eta_\nu(t))\,\mathrm{d}t.
	\end{equation}
	这里 \(E(\eta_\nu)=\int_{\mathbb R}|\dot\eta_\nu|_{g_{t,\sigma_\nu}}^2\,\mathrm{d}t\)，等式由有限区间积分再取两端极限得到。紧性使积分项的绝对值不超过
	\[
		2T\max_{[-T,T]\times K\times M}|\partial_t f_{t,\sigma}|.
	\]
	中心区间上的收敛使该积分项趋于 \(\eta\) 的同类积分。两端各段能量分别相加为 \(f_-(p)-f_-(r_-)\)、\(f_+(r_+)-f_+(q)\)；中心能量由同一公式计算。三部分之和正好是原能量的极限，故没有遗漏。
\end{proof}

这一证明还有一个用于延拓复合的形式。若时间族由两个固定长度的过渡区间组成，中间在一套固定 Morse 数据上停留长度 \(R\)，当 \(R\to\infty\) 时，分别在两个过渡区间的中心选定时间零点，可取得两条中心延拓轨线。中间自治段用有限正则水平作相同提取，得到连接两条中心轨线的有限梯度断裂链；两端仍可有上述断裂。能量恒等式中的时间导数只支撑于两个固定过渡区间，故上界与 \(R\) 无关。这个结论控制可能的极限；这些极限附近的局部坐标与粘合将在后面构造延拓映射时使用。

\subsection{零维紧性及计数有限性}
\label{g2:subsec:1802:04}

现在可以回答链群微分中的第一个实际问题：指标相差一的临界点之间，究竟是否只有有限条非参数轨线？

\begin{corollary}\label{g2:cor:zero-dimensional-trajectories-finite}
	设 \(M\) 为闭光滑流形，\((f,g)\) 为 Morse--Smale 数据，\(f\) 的临界值互异。若 \(\lambda(p)-\lambda(q)=1\)，则 \(\mathcal M(p,q)\) 是有限集。
\end{corollary}
\begin{proof}
	\cref{g2:prop:trajectory-quotient}的横截交维数公式使 \(\mathcal M(p,q)\) 为零维光滑流形，因而每一点都是孤立点。\cref{g2:thm:broken-trajectory-compactness}说明任何序列都有有限断裂极限。但若断成至少两段，每段指标严格下降至少一，便有 \(\lambda(p)-\lambda(q)\ge2\)，与假设矛盾。所以所有极限仍在 \(\mathcal M(p,q)\) 内，原空间已经紧。紧离散空间有限：其各个单点邻域组成开覆盖，有限子覆盖即给出有限点集。
\end{proof}

于是可定义
\[
	n_2(p,q)\coloneqq\#\mathcal M(p,q)\pmod2\in\mathbb F_2.
\]
有限性必须先于这个计数；零维只说明局部孤立，不能排除无穷多个点向一个未加入的断裂极限聚集。指标差为一时，正是严格的指标下降排除了这种聚集。

\ProseParagraphBreak

当 \(\lambda(p)-\lambda(q)=2\) 时，最多出现两段，所有可能的新增点都属于
\[
	\bigsqcup_{\lambda(r)=\lambda(q)+1}
	\mathcal M(p,r)\times\mathcal M(r,q).
\]
每个因子已经由上述推论知道有限，故候选断裂点也有限。不过，“只有这些候选极限”仍未说明每个候选点附近是一条半开弧，也没有排除多条弧以同一个方式聚集。下一节的粘合与唯一性将准确补上这项局部几何事实。

% !TeX root = ../../Book2.tex
% 纲要标题：### 18.3 粘合与低维边界结构
% 纲要位置：几何学计划纲要.md，第 1663 行。
% 正文状态：未撰写；本文件目前是框架。

% 纲要内容开始
% 
% * 断裂轨线的局部拼接
% * 线性化和横截性给出的逆
% * 非线性修正与唯一性
% * 一维紧化的边界识别
% 
% 纲要内容结束

\section{粘合与低维边界结构}
\label{g2:sec:1803}

待撰写

\subsection{断裂轨线的局部拼接}
\label{g2:subsec:1803:01}
待撰写

\subsection{线性化和横截性给出的逆}
\label{g2:subsec:1803:02}
待撰写

\subsection{非线性修正与唯一性}
\label{g2:subsec:1803:03}
待撰写

\subsection{一维紧化的边界识别}
\label{g2:subsec:1803:04}
待撰写
% 纲要条目排为编号小节标题，并进入目录与书签；说明段仍保留于纲要注释中。

% !TeX root = ../../Book2.tex
% 纲要标题：### 18.4 复形、不变性与经典同调
% 纲要位置：几何学计划纲要.md，第 1670 行。
% 正文状态：未撰写；本文件目前是框架。

% 纲要内容开始
% 
% * 微分和边界平方为零
% * 延拓映射与参数族链同伦
% * 由把手滤过构造比较映射
% * 同调同构及对选择的独立性
% 
% 纲要内容结束

\section{复形、不变性与经典同调}
\label{g2:sec:1804}

待撰写

\subsection{微分和边界平方为零}
\label{g2:subsec:1804:01}
待撰写

\subsection{延拓映射与参数族链同伦}
\label{g2:subsec:1804:02}
待撰写

\subsection{由把手滤过构造比较映射}
\label{g2:subsec:1804:03}
待撰写

\subsection{同调同构及对选择的独立性}
\label{g2:subsec:1804:04}
待撰写
% 纲要条目排为编号小节标题，并进入目录与书签；说明段仍保留于纲要注释中。

% !TeX root = ../../Book2.tex
% 纲要标题：### 18.5 整数系数与推广边界
% 纲要位置：几何学计划纲要.md，第 1677 行。
% 正文状态：未撰写；本文件目前是框架。

% 纲要内容开始
% 
% 在已完成 F₂ 链复形之后，用不稳定流形取向构造相容的轨线取向，核对断裂边界符号，再说明 Morse–Bott、非紧和等变情形分别增加何种紧性或横截条件；不统一打包为一个“推广定理”。
% 
% ---
% 
% 纲要内容结束

\section{整数系数与推广边界}
\label{g2:sec:1805}

待撰写
% 本节沿用纲要的连续叙述安排，不另造小节标题。


\begin{chaptersummary}
待撰写
\end{chaptersummary}

\BookExercises
待撰写


% !TeX root = ../Book2.tex
% 纲要标题：# 第四编　向量丛、示性类与曲率代表
% 纲要位置：几何学计划纲要.md，第 1733 行。
% 正文状态：本编导读初稿；尚未正式审核。

% 纲要内容开始
% 
% 纲要内容结束

\BookPart{向量丛、示性类与曲率代表}{g2:part:04}
\BookPartGuide
向量丛在每个小邻域上都是乘积，整体上却可能没有一组处处连续的标架。要比较两张丛，可以直接比较过渡函数，也可以把它们送到同一个参数空间，或计算由它们产生的上同调类。若丛带有联络，还能通过曲率得到微分形式。本编依次建立这些描述之间的联系，使拓扑分类、截面零点和曲率积分可以在同一张丛上相互对照。

\ProseParagraphBreak

\chapref{g2:ch:19}从局部平凡化与拉回出发，证明向量丛随底空间映射的同伦保持同构类型。将各局部标架用单位分解组合，就能把纤维嵌入一个固定线性空间，得到以 Grassmann 空间为目标的分类映射。分类映射的同伦唯一性说明这种表示如何摆脱局部选择。本章同时区分固定秩的分类、添加平凡丛后的稳定分类，以及带联络时需要保留的额外数据。

\ProseParagraphBreak

接下来的构造从每根纤维的方向开始。\chapref{g2:ch:20}把盘与边界球面的相对生成元粘成 Thom 类，再证明与它作杯积给出次数平移的同构。沿零截面拉回产生 Euler 类，把它放回盘球丛对的长正合列则得到 Gysin 序列。截面的横截零点使这些上同调类具有几何代表；局部系数保留绕行时的方向变化。Euler 类所记录的首障碍，与更高维胞腔上可能继续出现的障碍，在具体例子中分别考察。

\ProseParagraphBreak

\chapref{g2:ch:21}继续构造不同次数的示性类。射影丛的上同调关系给出实丛的 Stiefel--Whitney 类和复丛的陈类，实丛的复化再产生 Pontryagin 类。分裂原理把计算化到线丛，同时保证所得等式能回到原来的底空间。Thom 平方及 Wu 公式连接这些类与前面的上同调运算；复 Grassmann 流形的 Schubert 计算则把示性类落实为子空间满足相交条件的计数。

\ProseParagraphBreak

\chapref{g2:ch:22}转向光滑丛及联络。不变多项式作用于曲率后得到闭形式，改变联络时，转移形式给出两份代表之间的恰当差。与拓扑示性类比较以后，曲率的积分才具有确定的归一化和拓扑含义。陈--Simons 形式进一步记录联络之间的差异；三维模型把规范变换、整数模糊和边界条件写成可以核对的公式。

\ProseParagraphBreak

阅读\chapref{g2:ch:19}需要第一卷的向量丛、连续单位分解，以及本卷\chapref{g2:ch:01}的 CW 构造；丛的分类可以从这里单独学习。\chapref{g2:ch:20}、\chapref{g2:ch:21}需要本卷\chapref{g2:ch:04}至\chapref{g2:ch:06}的上同调乘积、局部系数和对偶，零点解释另用\chapref{g2:ch:14}至\chapref{g2:ch:16}的横截性与相交。与首障碍、平方运算有关的比较分别回到\chapref{g2:ch:12}、\chapref{g2:ch:13}。\chapref{g2:ch:22}则需第一卷第18章的联络和曲率、本卷\chapref{g2:ch:07}的 de Rham 比较及本编已建的示性类。计算时宜同时注明工作系数：模二类、整数类和曲率的实上同调类回答不同的问题，其中的整数挠信息应由拓扑构造继续保存。

% !TeX root = ../../Book2.tex
% 纲要标题：## 第19章　向量丛及分类空间
% 纲要位置：几何学计划纲要.md，第 1735 行。
% 正文状态：未撰写；本文件目前是框架。

% 纲要内容开始
% 
% > 本章目标与证明范围：在仿紧、可数可平凡化或 CW 等明确条件下建立分类定理。区分有限维 Grassmann 近似、稳定分类和含联络的几何分类。
% 
% 纲要内容结束

\chapter{向量丛及分类空间}
\label{g2:ch:19}
\BookChapterNumber{g2:ch:19}

\begin{chapterabstract}
待撰写
\end{chapterabstract}

\begin{learninggoals}
待撰写
\end{learninggoals}

待撰写

% !TeX root = ../../Book2.tex
% 纲要标题：### 19.1 丛的运算与同伦不变性
% 纲要位置：几何学计划纲要.md，第 1739 行。
% 正文状态：初稿；尚未正式审核。

% 纲要内容开始
% 
% * 过渡函数、拉回和丛同构
% * 区间上的平凡化延拓
% * 同伦拉回丛的同构
% * 球面向量丛的粘合分类
% 
% 纲要内容结束

\section{丛的运算与同伦不变性}
\label{g2:sec:1901}

把圆周上的每一点换成一条直线，既可以得到平凡线丛，也可以得到 Möbius 线丛。纤维维数相同，区别发生在绕行一周后的粘合。分类向量丛首先要确定允许怎样改变局部坐标，以及底空间中的连续变形会不会改变所得的丛。本节把第一卷的光滑丛语言转到连续范畴，再建立分类映射将要使用的同伦不变性。

\subsection{过渡函数、拉回和丛同构}
\label{g2:subsec:1901:01}

以下 \(\mathbb K\) 表示 \(\mathbb R\) 或 \(\mathbb C\)。第一卷第12章第1节中，向量丛的局部平凡化把每根纤维线性识别为同一个向量空间。本章允许一般拓扑基空间，要求平凡化及换架连续；当底空间和过渡函数都光滑时，便恢复那里的光滑向量丛。

\begin{definition}\label{g2:def:topological-vector-bundle}
	设 \(X\) 为拓扑空间。一个秩为 \(n\) 的 \(\mathbb K\)-向量丛是连续映射 \(\pi:E\to X\)，其每根纤维 \(E_x=\pi^{-1}(x)\) 是 \(n\) 维 \(\mathbb K\)-向量空间，并且每点附近存在开集 \(U\) 及同胚
	\[
		\phi_U:E|_U\longrightarrow U\times\mathbb K^n,
	\]
	使第一分量为 \(\pi\)，第二分量在各纤维上线性可逆。这种同胚称为局部平凡化。覆盖恒等映射、逐纤维线性的连续映射称为丛同态；若它是逐纤维可逆的，则称为丛同构。
\end{definition}

逐纤维可逆已经保证逆映射连续：局部表示是连续矩阵 \(A(x)\in\operatorname{GL}_n(\mathbb K)\)，而逆矩阵的各元素由行列式和余子式连续给出。因此，同构要求的不只是逐点存在某个线性同构，而是这些同构随基点连续变化。记同构类的集合为 \(\operatorname{Vect}_n^{\mathbb K}(X)\)，平凡秩 \(n\) 丛记为 \(\varepsilon_X^n=X\times\mathbb K^n\)。

\ProseParagraphBreak

在覆盖 \(\{U_i\}\) 上选平凡化 \(\phi_i\)，坐标列向量满足 \(v_i=g_{ij}(x)v_j\)，其中
\[
	g_{ij}:U_i\cap U_j\longrightarrow\operatorname{GL}_n(\mathbb K),
	\qquad g_{ij}g_{jk}=g_{ik},\quad g_{ii}=I.
\]
这些连续矩阵就是\term[guodu hanshu]{过渡函数}[transition functions]。反过来，满足上述余循环等式的矩阵把各 \(U_i\times\mathbb K^n\) 粘合起来：在重叠处将第 \(j\) 张图的 \((x,v)\) 与第 \(i\) 张图的 \((x,g_{ij}(x)v)\) 识别。余循环等式恰好保证识别的传递性。取商拓扑后，各张图的饱和开集给出局部乘积，因而所得对象确实是向量丛。

\begin{proposition}\label{g2:prop:bundle-transition-isomorphism}
	设 \(E,F\) 是同一基空间上的有限秩 \(\mathbb K\)-向量丛，在共同开覆盖上的过渡函数分别为 \(g_{ij},h_{ij}\)。丛同态 \(E\to F\) 等价于连续矩阵族 \(A_i\)，满足
	\[
		A_i g_{ij}=h_{ij}A_j.
	\]
	同秩时，它是丛同构当且仅当各 \(A_i(x)\) 可逆。特别地，将局部坐标改为 \(v'_i=b_i(x)v_i\)，过渡函数变为 \(g'_{ij}=b_i g_{ij}b_j^{-1}\)。
\end{proposition}
\begin{proof}
	局部同态写成 \((x,v_i)\mapsto(x,A_i(x)v_i)\)。同一个向量也可用第 \(j\) 张图表示；两种输出相同恰好要求 \(A_ig_{ij}v_j=h_{ij}A_jv_j\) 对所有 \(v_j\) 成立。这就是所列等式，它保证局部连续映射粘成整体同态。逆矩阵连续说明同构判据，最后的换架公式由 \(v'_i=b_ig_{ij}b_j^{-1}v'_j\) 得到。
\end{proof}

丛的运算由纤维运算和过渡函数共同确定。例如 \(E\oplus F\)、\(E\otimes F\) 和 \(E^*\) 的过渡矩阵依次为
\[
	\begin{pmatrix}g_{ij}&0\\0&h_{ij}\end{pmatrix},
	\qquad g_{ij}\otimes h_{ij},
	\qquad(g_{ij}^{-1})^{\mathsf T}.
\]
这里复丛的 \(E^*\) 取复线性对偶，所以出现转置而不是共轭转置。矩阵乘法、张量积及逆转置保留余循环等式，故这些纤维运算都得到向量丛；同样有 \(\operatorname{Hom}(E,F)=E^*\otimes F\) 与外幂 \(\bigwedge^kE\)。

\ProseParagraphBreak

对连续映射 \(f:Y\to X\)，定义\term[lahui cong]{拉回丛}[pullback bundle]
\[
	f^*E=\{(y,e)\in Y\times E:f(y)=\pi(e)\},\qquad (y,e)\longmapsto y.
\]
其局部平凡化来自 \(E\) 在 \(U_i\) 上的平凡化，过渡函数为 \(g_{ij}\circ f\)。下面的方块既交换，也刻画了这个构造的泛性质：
\[
	\begin{tikzcd}
		f^*E\arrow[r]\arrow[d]&E\arrow[d,"\pi"]\\
		Y\arrow[r,"f"']&X.
	\end{tikzcd}
\]
具体说，若 \(q:F\to Y\) 是另一向量丛，\(\Phi:F\to E\) 是覆盖 \(f\) 的逐纤维线性连续映射，则唯一的分解为 \(v\mapsto(q(v),\Phi(v))\in f^*E\)。连续性来自积与子空间拓扑。由此得到自然同构 \((f\circ h)^*E\cong h^*(f^*E)\)，以及拉回与直和、张量积、对偶交换；逐点写下对应的有序组或张量即可核对它们与过渡函数相容。

\subsection{区间上的平凡化延拓}
\label{g2:subsec:1901:02}

为比较同伦的两端，先考察 \(E\to X\times I\)，其中 \(I=[0,1]\)。局部坐标能把相邻时间的纤维对应起来，但一张坐标通常覆盖不了整个区间。区间的紧性提供有限次接续；接缝处需要调整坐标，使两段的对应真正一致。

\begin{lemma}\label{g2:lem:bundle-cylinder-local-trivialization}
	设 \(E\to X\times I\) 为有限秩实或复向量丛。对每个 \(x\in X\)，存在开邻域 \(U\)，使 \(E|_{U\times I}\) 平凡。若已给定 \(E|_{U\times\{0\}}\) 的平凡化，可以在缩小 \(U\) 后将其延拓到整条柱体。
\end{lemma}
\begin{proof}
	沿紧集 \(\{x\}\times I\) 取有限个乘积邻域上的平凡化。取这有限个基空间邻域的交，并对时间作细分
	\[
		0=t_0<t_1<\cdots<t_m=1,
	\]
	使每个 \(U\times[t_{j-1},t_j]\) 都落在其中一张平凡化内。这样的细分存在，因为有限个时间开区间覆盖紧区间，具有正的 Lebesgue 数。

	设前一段已使用坐标 \(\psi_{j-1}\)，后一段原来的坐标为 \(\phi_j\)。在接缝 \(t=t_{j-1}\) 上，二者相差一个连续可逆矩阵 \(B_j(y)\)：\(\psi_{j-1}(y,t_{j-1})=B_j(y)\phi_j(y,t_{j-1})\)。将后一段坐标改为 \(\psi_j(y,t)=B_j(y)\phi_j(y,t)\)，就在接缝处相同。有限个闭时间段的粘合引理使这些坐标拼成一个连续平凡化；逆坐标也按同样方式粘合，因此确为同胚。在第一段先乘一个只依赖 \(y\) 的可逆矩阵，就能匹配任意指定的起点平凡化，再依次接续。
\end{proof}

固定这样的柱体平凡化，保持坐标列向量不变，便得到局部运输
\[
	T_U(y;s,t):E_{(y,s)}\longrightarrow E_{(y,t)},
	\qquad 0\le s,t\le1.
\]
它连续依赖 \((y,s,t)\)，并满足 \(T_U(y;t,t)=I\) 和 \(T_U(y;t,r)T_U(y;s,t)=T_U(y;s,r)\)。不同柱体平凡化所得运输可以不同。下一步用单位分解决定每次使用哪一张图及使用多长的时间，使这些局部运输有次序地接成整体同构。

\subsection{同伦拉回丛的同构}
\label{g2:subsec:1901:03}

此处使用第一卷第4章第4节的连续单位分解定理：仿紧 Hausdorff 空间的任意开覆盖都有从属连续单位分解 \(\{\rho_i\}\)，支集族局部有限，\(\operatorname{supp}\rho_i\subset U_i\)，并且 \(\sum_i\rho_i=1\)。局部有限性使每一点附近的求和实际只有有限项。下面用这些函数分配时间，而不把可逆矩阵作平均，因为可逆矩阵的平均未必仍可逆。

\begin{theorem}[丛的同伦不变性]\label{g2:thm:bundle-homotopy-invariance}
	设 \(X\) 为仿紧 Hausdorff 空间，\(E\to X\times I\) 为有限秩实或复向量丛，\(E_0=E|_{X\times\{0\}}\)。则存在覆盖 \(X\times I\) 恒等映射的丛同构
	\[
		\operatorname{pr}_X^*E_0\cong E,
	\]
	它在 \(t=0\) 上为恒等。因此 \(E_0\cong E_1\)。若 \(f_0,f_1:X\to Y\) 连续同伦，\(V\to Y\) 为有限秩向量丛，则 \(f_0^*V\cong f_1^*V\)。
\end{theorem}
\begin{proof}
	由\cref{g2:lem:bundle-cylinder-local-trivialization}取开覆盖 \(\{U_i\}_{i\in A}\)，使各 \(U_i\times I\) 上都有平凡化及运输 \(T_i\)。取从属的局部有限单位分解 \(\rho_i\)，再给指标集选一个良序。良序只用来固定有限个实际运输的先后，最终同构不要求与这项选择无关。

	令
	\[
		a_i(x)=\sum_{j<i}\rho_j(x),\qquad b_i(x)=a_i(x)+\rho_i(x).
	\]
	这些函数连续，因为每点附近只有有限个支集。对目标时间 \(t\in I\)，第 \(i\) 次运输从高度 \(t a_i(x)\) 到 \(t b_i(x)\)。当 \(x\in U_i\) 时，使用
	\[
		L_i(x,t)=T_i\bigl(x;t a_i(x),t b_i(x)\bigr).
	\]
	当 \(x\notin\operatorname{supp}\rho_i\) 时，两高度相同，令 \(L_i\) 为该纤维上的恒等映射。两种规定在交叠处一致；支集包含于 \(U_i\)，所以它们合成两个图拉回丛之间的连续同构。这个支集条件同时保证，在需要使用局部运输的地方，所选坐标始终有定义。

	现在固定 \(x\)。只有有限个 \(\rho_i(x)\) 非零，按既定顺序复合相应 \(L_i(x,t)\)，得到
	\[
		L(x,t):E_{(x,0)}\longrightarrow E_{(x,t)}.
	\]
	这些时间段首尾相接，因为累积和从零增加到一；零权重给出的段只是恒等。要证明 \(L\) 连续，取 \(x\) 的一个只遇到有限多个支集的邻域 \(W\)。在 \(W\times I\) 上，将这有限个指标全部按顺序列入复合，包括可能权重为零的项，便得到一个固定的有限连续复合。它与逐点定义一致，所以 \(L\) 连续。反向复合这些运输给出连续逆，故 \(L\) 是所需丛同构。\(t=0\) 时每一段都是恒等。

	若 \(H:X\times I\to Y\) 是从 \(f_0\) 到 \(f_1\) 的同伦，将前一结论用于 \(H^*V\)。它的两端分别是 \(f_0^*V\)、\(f_1^*V\)，故得到最后的断言。
\end{proof}

这项构造处理了非紧底空间：不同点可以需要不同数量、不同位置的时间坐标图；局部有限性保证附近的搬运仍由有限个连续步骤给出。同构依赖平凡化及单位分解，定理确定的是同构类。

\begin{corollary}\label{g2:cor:contractible-base-trivial-bundle}
	仿紧 Hausdorff 可缩空间上的每个有限秩实或复向量丛都平凡。更一般地，若 \(f:X\to Y\) 是仿紧 Hausdorff 空间之间的同伦等价，则拉回给出 \(\operatorname{Vect}_n^{\mathbb K}(Y)\to\operatorname{Vect}_n^{\mathbb K}(X)\) 的双射。
\end{corollary}
\begin{proof}
	可缩时，恒等映射与某个常值映射 \(c_{x_0}\) 同伦，而 \(c_{x_0}^*E=X\times E_{x_0}\)，故由定理得到平凡性。一般地，若 \(g\) 为 \(f\) 的同伦逆，则 \(g^*f^*=(f\circ g)^*\) 与 \(f^*g^*=(g\circ f)^*\) 在同构类上分别为恒等，故互为逆映射。
\end{proof}

例如，圆环 \(S^1\times I\) 上的向量丛都来自核心圆周；从内圆周限制和从外圆周限制所得的丛同构。这样的同构并不规定绕圆周的平行移动，后者还需要联络。

\begin{corollary}\label{g2:cor:principal-bundle-homotopy}
	设 \(X\) 为仿紧 Hausdorff 空间，\(G\) 为拓扑群，\(P\to X\times I\) 为局部平凡的右主 \(G\) 丛。则 \(P\cong\operatorname{pr}_X^*(P|_{X\times\{0\}})\)，且同构可在零端取为恒等。因而同伦映射拉回主 \(G\) 丛得到同构主丛。
\end{corollary}
\begin{proof}
	右主丛的局部模型是 \(U\times G\)，右作用为第二分量的右乘；换平凡化在第二分量上左乘一个连续的群值函数。因而柱体平凡化仍可逐段调整并粘合，保持第二坐标的局部运输是右 \(G\) 等变的。上一定理只复合这些运输，没有使用纤维加法。以相同的时间分配与局部有限复合构造同构，它自动右 \(G\) 等变。
\end{proof}

\subsection{球面向量丛的粘合分类}
\label{g2:subsec:1901:04}

把 \(S^m\)、\(m\ge1\) 分成两个闭半球 \(D_+^m,D_-^m\)，交集为赤道 \(S^{m-1}\)。半球可缩，其上的向量丛平凡；全部粘合信息因而集中到赤道上的一族可逆线性变换。

\begin{definition}\label{g2:def:clutching-bundle}
	对连续映射 \(g:S^{m-1}\to\operatorname{GL}_n(\mathbb K)\)，将两个平凡丛 \(D_+^m\times\mathbb K^n\)、\(D_-^m\times\mathbb K^n\) 按
	\[
		(x,v)_+\sim(x,g(x)v)_-\qquad(x\in S^{m-1})
	\]
	粘合，所得丛记为 \(E_g\to S^m\)。映射 \(g\) 称为这个构造的\term[clutching hanshu]{粘合函数}[clutching function]。
\end{definition}

赤道附近的局部平凡性可以直接看见。取领圈 \(S^{m-1}\times(-\varepsilon,\varepsilon)\)，在正半侧保持原坐标，在负半侧乘 \(g(x)^{-1}\)；两侧在赤道上给出相同的坐标。因而商空间在每个赤道点附近仍为乘积，远离赤道则沿用半球内的平凡化。

\begin{theorem}[球面上的粘合分类]\label{g2:thm:clutching-classification}
	设 \(m,n\ge1\)。每个 \(S^m\) 上的秩 \(n\) 实或复向量丛同构于某个 \(E_g\)。两粘合函数 \(g,h\) 给出同构丛，当且仅当存在连续映射
	\[
		A_\pm:D_\pm^m\longrightarrow\operatorname{GL}_n(\mathbb K)
	\]
	使赤道上有 \(h=A_-gA_+^{-1}\)。同伦的粘合函数给出同构丛。

	复丛的同构类因而与 \([S^{m-1},\operatorname{GL}_n(\mathbb C)]\) 双射；若只考虑带指定取向的实丛及保向同构，则相应地与 \([S^{m-1},\operatorname{GL}_n^+(\mathbb R)]\) 双射。
\end{theorem}
\begin{proof}
	由\cref{g2:cor:contractible-base-trivial-bundle}，给定丛在两个半球上都可平凡化；沿赤道比较这两套坐标就得到 \(g\)，原丛由同样的识别关系恢复。若有覆盖球面恒等映射的同构 \(E_g\to E_h\)，在两半球内分别表示为 \(A_+,A_-\)。赤道上的同一点先粘合再映射或先映射再粘合，给出
	\[
		A_-(x)g(x)v=h(x)A_+(x)v.
	\]
	这就是所述关系。反过来，满足关系的两半球同构在接缝处相容，因而粘成整体同构。

	若 \(g_t\) 是粘合函数间的同伦，就在赤道乘区间上用 \(g_t\) 粘合两个 \(D_\pm^m\times I\times\mathbb K^n\)。刚才的领圈坐标同时依赖 \(t\)，所以得到 \(S^m\times I\) 上的向量丛，其两端为 \(E_{g_0},E_{g_1}\)。\cref{g2:thm:bundle-homotopy-invariance}说明两端同构。

	最后，\(\operatorname{GL}_n(\mathbb C)\) 与 \(\operatorname{GL}_n^+(\mathbb R)\) 都道路连通。这里可用极分解先将可逆矩阵连续变为酉矩阵或正向正交矩阵，再由酉矩阵的特征角、实正交矩阵的二维旋转块把它连到单位矩阵。\(A_\pm\) 可沿半球收缩同伦到一个常值，而这个常值再连到单位元。因此 \(h=A_-gA_+^{-1}\) 蕴含 \(h\simeq g\)。带取向时先在两半球选保向标架，所有矩阵都落在正行列式群中。于是上述满射在相应同伦类上也单射。
\end{proof}

常用计算还可以把群缩小。极分解给出形变收缩
\[
	A\longmapsto A(A^*A)^{-s/2},\qquad 0\le s\le1,
\]
复数情形终点在 \(\operatorname U(n)\)，实数情形终点在 \(\operatorname O(n)\)，正行列式分支则在 \(\operatorname{SO}(n)\)。正定矩阵的实幂由正交或酉对角化定义，并连续依赖矩阵；在任一固定紧谱区间上用多项式一致逼近相应正实函数，也能直接验证这种连续性。因此上述球面分类可以分别使用 \(\operatorname U(n)\) 和 \(\operatorname{SO}(n)\)。对 \(m\ge2\)，归一化映射 \(g(x)\mapsto g(x_0)^{-1}g(x)\) 将自由同伦类识别为这些连通群的 \(\pi_{m-1}\)：给定自由同伦 \(H(x,t)\)，乘上 \(H(x_0,t)^{-1}\) 就得到定基点同伦。

\ProseParagraphBreak

未指定取向的实丛需要保留换架关系中的两个分支。当 \(m\ge2\) 时，赤道连通，\(\det g\) 符号恒定；必要时在一个半球改用反向标架，可使 \(g\) 落在正行列式群中。因此每个这样的实丛都可定向。遗忘取向后，还须将 \(g\) 与 \(JgJ^{-1}\) 识别，其中 \(J\) 是一个固定反射矩阵。确实，两半球同时反向换架改变了总取向；反之，在 \(h,g\) 均正行列式的等式 \(h=A_-gA_+^{-1}\) 中，\(A_-\) 与 \(A_+\) 的行列式符号必须相同，故两者可同时同伦到 \(I\) 或同时同伦到 \(J\)。普通实丛的分类便是相应 \(\pi_{m-1}\operatorname{SO}(n)\) 对这个反射共轭作用的轨道集。

\ProseParagraphBreak

\(m=1\) 时，赤道 \(S^0\) 有两个点，不能用连通性将二者的行列式符号合并。沿两个闭弧平凡化后，可以把一个接缝矩阵化为单位矩阵，另一个接缝留下一个矩阵 \(B\)。\(\det B>0\) 时实丛平凡，\(\det B<0\) 时同构于 Möbius 线丛与 \(n-1\) 条平凡线的直和：把 \(B\) 分别沿其所在连通分支变到 \(I\) 或 \(\operatorname{diag}(-1,1,\ldots,1)\) 即得。复矩阵群只有一个连通分支，所以圆周上的复向量丛全部平凡。这个一维例子显示，换架的连通分支本身就是分类数据。

% !TeX root = ../../Book2.tex
% 纲要标题：### 19.2 通用丛和 Grassmann 模型
% 纲要位置：几何学计划纲要.md，第 1746 行。
% 正文状态：初稿；尚未正式审核。

% 纲要内容开始
% 
% * 有限维 Grassmann 典范丛
% * 由截面构造分类映射
% * 维数增大与相容嵌入
% * 无限 Grassmann 空间的拓扑
% 
% 纲要内容结束

\section{通用丛和 Grassmann 模型}
\label{g2:sec:1902}

若能把一张秩 \(n\) 向量丛的每根纤维都线性嵌入同一个 \(\mathbb K^N\)，那么每个基点就确定一个 \(n\) 维子空间。原来的丛因而可以用一张到“所有 \(n\) 维子空间组成的空间”的映射来描述。这里 \(\mathbb K\) 表示 \(\mathbb R\) 或 \(\mathbb C\)，秩始终按 \(\mathbb K\) 计算。本节把这一直观构造写成局部坐标，并说明当不同丛需要不同的 \(N\) 时，如何把这些模型放在同一个空间中。

\subsection{有限维 Grassmann 典范丛}
\label{g2:subsec:1902:01}
第1章已用正交投影矩阵定义 Grassmann 空间，并在\cref{g2:thm:grassmann-schubert-cw}中构造其胞腔分解。这里记为
\[
	\operatorname{Gr}_n(\mathbb K^N)=\{V\subseteq\mathbb K^N:\dim_{\mathbb K}V=n\},\qquad 1\le n\le N.
\]
拓扑由 \(V\mapsto P_V\) 给出，其中 \(P_V\) 是标准内积下的正交投影。矩阵条件 \(P=P^*\)、\(P^2=P\)、\(\operatorname{tr}P=n\) 是闭条件，且 \(\|P\|\le1\)，所以该空间紧且 Hausdorff。现在需要一种能同时描述子空间及其内部向量的坐标。

\ProseParagraphBreak

固定正交分解 \(\mathbb K^N=A\oplus A^\perp\)，其中 \(\dim A=n\)。令 \(U_A\) 为投影 \(\pi_A|_V:V\to A\) 可逆的那些 \(V\)。这样的子空间唯一写成线性映射 \(L:A\to A^\perp\) 的图：
\[
	V=\operatorname{graph}L=\{a+L(a):a\in A\}.
\]
设 \(B_L\) 表示矩阵 \(\binom{I}{L}\)，则
\[
	P_V=B_L(I+L^*L)^{-1}B_L^*.
\]
反过来，若按 \(A\oplus A^\perp\) 把 \(P_V\) 分块，则 \(L=P_{21}P_{11}^{-1}\)。这些公式证明图参数与投影矩阵相互连续，且在有限维矩阵变量中光滑。可逆条件是开条件，故 \(U_A\) 开。任意子空间 \(V\) 都属于 \(U_V\)，这些图坐标覆盖整个空间，坐标变换也是由矩阵乘法及可逆矩阵求逆组成。因此实 Grassmann 空间的实维数为 \(n(N-n)\)，复 Grassmann 空间的复维数为 \(n(N-n)\)。

\begin{definition}\label{g2:def:finite-tautological-bundle}
	对 \(\mathbb K=\mathbb R\) 或 \(\mathbb C\) 及 \(1\le n\le N\)，定义\term[dianfan-xiangliangcong]{典范向量丛}[tautological vector bundle]
	\[
		\gamma_n^N=\{(V,v)\in\operatorname{Gr}_n(\mathbb K^N)\times\mathbb K^N:v\in V\},\qquad \pi(V,v)=V,
	\]
	总空间取积空间的子空间拓扑。其在 \(V\) 上的纤维就是 \(V\) 本身。
\end{definition}

它确实是秩 \(n\) 向量丛，而不只是附带线性纤维的映射。在 \(U_A\) 上，局部平凡化及逆为
\[
	(V,v)\longmapsto(V,\pi_Av),\qquad
	(V,a)\longmapsto(V,a+L_Va).
\]
二者连续且逐纤维线性。特别地，选定 \(A\) 的一组基 \(a_1,\ldots,a_n\) 后，\(a_j+L_Va_j\) 就给出 \(U_A\) 上的一组连续局部标架。\(n=1\) 时，基底分别是 \(\mathbb RP^{N-1}\)、\(\mathbb CP^{N-1}\)，所得典范线丛正是前面射影空间中的模型。

\ProseParagraphBreak

一张映射 \(c:X\to\operatorname{Gr}_n(\mathbb K^N)\) 的拉回丛可以直接写成
\[
	c^*\gamma_n^N=\{(x,v)\in X\times\mathbb K^N:v\in c(x)\}.
\]
因此有限维典范丛统一描述了平凡丛 \(X\times\mathbb K^N\) 的所有秩 \(n\) 子丛。这里的 \(N\) 控制可容纳的子丛；下一步证明紧基底上的每张有限秩丛都能放入某个这样的平凡丛。

\subsection{由截面构造分类映射}
\label{g2:subsec:1902:02}
一张非平凡丛通常没有一组处处独立的全局标架，但可以有更多全局截面，在每一点共同张成纤维。多出来的截面允许我们先使用局部标架，再把每个标架在其定义域以外补成零。单位分解恰好提供所需权重。

\begin{proposition}\label{g2:prop:bundle-embedding-classifier}
	设 \(X\) 为仿紧 Hausdorff 空间，\(E\to X\) 为秩 \(n\) 的实或复向量丛，并有有限平凡化覆盖 \(U_1,\ldots,U_m\)。则存在逐纤维线性的丛嵌入
	\[
		J:E\longrightarrow X\times\mathbb K^{mn}
	\]
	和连续映射 \(c_J:X\to\operatorname{Gr}_n(\mathbb K^{mn})\)，满足 \(c_J(x)=J(E_x)\) 及 \(E\cong c_J^*\gamma_n^{mn}\)。特别地，紧 Hausdorff 基底上的每张有限秩向量丛都有这样的有限维模型。
\end{proposition}
\begin{proof}
	取从属连续单位分解 \(\rho_1,\ldots,\rho_m\)，满足 \(\operatorname{supp}\rho_i\subseteq U_i\)。所用的是第一卷第4.4节的连续单位分解定理，其条件正是仿紧与 Hausdorff。将各局部内积加权求和，得到 \(E\) 上的连续正定内积。对每个 \(U_i\) 的局部标架作逐向量正交化，得到正交单位标架 \(e_{i1},\ldots,e_{in}\)；归一化的分母始终为正，故标架连续。

	令 \(s_{ia}=\sqrt{\rho_i}\,e_{ia}\) 于 \(U_i\)，在其余地方补零。支集闭且包含于 \(U_i\)，所以每个不属于 \(U_i\) 的点都有一个邻域使该截面恒为零，补零连续。它们定义满射
	\[
		A_x:\mathbb K^{mn}\longrightarrow E_x,\qquad
		(z_{ia})\longmapsto\sum_{i,a}z_{ia}s_{ia}(x).
	\]
	至少一个 \(\rho_i(x)>0\)，对应的 \(n\) 个截面已张成 \(E_x\)，所以确实满射。令 \(\tau_i(v)\in\mathbb K^n\) 为 \(v\in E_x\) 在第 \(i\) 组正交标架下的坐标。伴随映射给出
	\begin{equation}\label{g2:eq:weighted-bundle-embedding}
		J_x(v)=A_x^*(v)
		=\bigl(\sqrt{\rho_1(x)}\,\tau_1(v),\ldots,
		\sqrt{\rho_m(x)}\,\tau_m(v)\bigr).
	\end{equation}
	不在 \(U_i\) 时相应坐标仍取零。正交性与 \(\sum_i\rho_i=1\) 给出
	\[
		\|J_x(v)\|^2=\sum_i\rho_i(x)\|v\|^2=\|v\|^2.
	\]
	故 \(J_x\) 是等距单射，且 \(A_xJ_x=I_{E_x}\)。其像的正交投影为 \(P_x=J_xA_x\)，矩阵元连续依赖于 \(x\)，所以 \(c_J(x)=\operatorname{im}P_x\) 连续。

	映射 \(v\mapsto(x,J_xv)\) 是 \(E\) 到 \(c_J^*\gamma_n^{mn}\) 的逐纤维线性同构，逆为 \((x,z)\mapsto A_xz\)，故是丛同构。这给出拉回图
	\[
		\begin{tikzcd}[column sep=large]
			E\arrow[r,"{v\mapsto(c_J(x),J_xv)}"]\arrow[d]&\gamma_n^{mn}\arrow[d]\\
			X\arrow[r,"c_J"']&\operatorname{Gr}_n(\mathbb K^{mn}).
		\end{tikzcd}
	\]
	最后，紧 Hausdorff 空间仿紧，且任意局部平凡化覆盖有有限子覆盖，所以满足已用条件。
\end{proof}

这个构造也解释了映射为何取值于 Grassmann 空间：局部标架经加权提供了多个坐标，真正与原纤维对应的是这些坐标共同张成的 \(n\) 维像空间。\term[fenlei-yingshe]{分类映射}[classifying map]是给定一个拉回同构 \(E\cong c^*\gamma\) 的这种映射；不同权重与标架可能给出不同的 \(c\)。下一节将证明，在统一的无限模型中，它们具有同一个同伦类。

\subsection{维数增大与相容嵌入}
\label{g2:subsec:1902:03}
增大环境维数时，把 \(\mathbb K^N\) 嵌入 \(\mathbb K^{N+1}\) 的前 \(N\) 个坐标，便得到闭嵌入
\[
	i_N:\operatorname{Gr}_n(\mathbb K^N)\longrightarrow\operatorname{Gr}_n(\mathbb K^{N+1}),\qquad V\longmapsto V\oplus0.
\]
投影矩阵只是在右下角添零，所以它连续；紧源到 Hausdorff 目标的单射给出嵌入。典范丛相容地满足 \(i_N^*\gamma_n^{N+1}\cong\gamma_n^N\)，同构在纤维上就是 \(v\mapsto(v,0)\)。因此把一个分类映射移入更大的环境，不会改变它表示的秩 \(n\) 丛。

\ProseParagraphBreak

另一个有用构造是正交补丛。在 \(\operatorname{Gr}_n(\mathbb K^N)\) 上令
\[
	(\gamma_n^N)^\perp=\{(V,w):w\in V^\perp\}.
\]
它由连续投影矩阵 \(I-P_V\) 描述，秩为 \(N-n\)，图坐标同样给出局部平凡化。逐纤维相加得到
\begin{equation}\label{g2:eq:tautological-orthogonal-sum}
	\gamma_n^N\oplus(\gamma_n^N)^\perp
	\cong\operatorname{Gr}_n(\mathbb K^N)\times\mathbb K^N.
\end{equation}
拉回这个同构，便知满足\cref{g2:prop:bundle-embedding-classifier}条件的 \(E\) 有某张补丛 \(F\)，使 \(E\oplus F\) 平凡。补丛本身可以非平凡。

\begin{example}\label{g2:exm:sphere-tangent-stable-triviality}
	把单位球面 \(S^m\) 嵌入 \(\mathbb R^{m+1}\)，其切空间为 \(x^\perp\)。于是 \(x\mapsto x^\perp\) 将 \(\mathrm TS^m\) 表示为 \(\gamma_m^{m+1}\) 的拉回。正交补直线由处处非零的法向量 \(x\) 张成，故为平凡线丛。具体同构为
	\[
		\mathrm TS^m\oplus\varepsilon^1\longrightarrow\varepsilon^{m+1},
		\qquad(x,v,t)\longmapsto(x,v+tx),
	\]
	其中 \(\varepsilon^r=S^m\times\mathbb R^r\)。逆映射将 \(z\) 分解为 \(z-\langle z,x\rangle x\) 与 \(\langle z,x\rangle\)。

	然而 \(\mathrm TS^2\) 不平凡。若有连续非零切向量场，将它归一化为 \(v(x)\)，则
	\[
		H(x,t)=\cos(\pi t)x+\sin(\pi t)v(x)
	\]
	始终为单位向量，会给出恒等映射与对径映射之间的同伦。两者在 \(S^2\) 上的度分别为 \(1\)、\(-1\)，与\cref{g2:prop:degree-basic,g2:prop:reflection-antipodal-degree}矛盾。因此加一条平凡线后得到的同构，并不保证原秩下也有同构。
\end{example}

这里有两种不同的增大：\(N\) 增大只改变装载纤维的环境，\(E\mapsto E\oplus\varepsilon^r\) 则实际增加纤维的秩。前者用于固定秩的分类空间，后者产生稳定分类。例子中的切丛说明，稳定操作可能遗失原秩下的信息。

\subsection{无限 Grassmann 空间的拓扑}
\label{g2:subsec:1902:04}
把 \(\mathbb K^\infty\) 定义为有限支撑序列构成的向量空间，即 \(\bigcup_N\mathbb K^N\)。它的每个有限维子空间都包含于某个 \(\mathbb K^N\)：只须对一组有限基向量的支撑取并。因而所有 \(n\) 维子空间可以由前面的有限模型穷尽。

\begin{definition}\label{g2:def:infinite-grassmann}
	对固定 \(n\ge1\)，定义
	\[
		G_n=\operatorname{Gr}_n(\mathbb K^\infty)
		=\bigcup_{N\ge n}\operatorname{Gr}_n(\mathbb K^N).
	\]
	赋予它\term[zhijixian-ruotuopu]{直极限弱拓扑}[weak direct-limit topology]：集合 \(U\subseteq G_n\) 开，当且仅当每个 \(U\cap\operatorname{Gr}_n(\mathbb K^N)\) 都在该有限模型中开。同样，令 \(\gamma_n=\bigcup_N\gamma_n^N\)，总空间也取直极限弱拓扑。
\end{definition}

这些弱拓扑的作用是把“由各有限阶段定义的映射”粘合成连续映射。有限阶段的原有拓扑保留不变：一个有限阶段中的闭集，在所有更大的阶段中仍闭，所以在整个直极限中闭。投影矩阵的各个坐标函数在所有阶段连续，故在 \(G_n\) 上连续，并能分离不同子空间，因此 \(G_n\) 是 Hausdorff 空间。第1章的 Schubert 胞腔与上述嵌入相容：在 \(\mathbb K^N\) 中的胞腔恰好是最大枢轴指标不超过 \(N\) 的那些胞腔，闭包关系保持不变。所以这些有限子复形取并也给出 \(G_n\) 的 CW 结构，其弱拓扑就是这里规定的拓扑。

\begin{lemma}\label{g2:lem:grassmann-compact-stage}
	固定 \(n\ge1\) 及 \(\mathbb K=\mathbb R\) 或 \(\mathbb C\)。在上述弱拓扑下，\(\operatorname{Gr}_n(\mathbb K^\infty)\) 的每个紧子集都包含于某个有限阶段 \(\operatorname{Gr}_n(\mathbb K^N)\)。因此紧空间上的映射及其同伦都可在足够大的有限阶段中进行。
\end{lemma}
\begin{proof}
	若紧集 \(K\) 不包含于任一阶段，可递归选择点 \(x_j\in K\) 与严格递增的整数 \(N_j\)，使第 \(N_j\) 阶段包含全部先前选过的点，而 \(x_j\) 不在该阶段。于是点集 \(A=\{x_j:j\ge1\}\) 与每个有限阶段都只有有限交。对任意子集 \(A'\subseteq A\)，该交仍有限，因有限阶段 Hausdorff而闭。弱拓扑遂说明 \(A'\) 在 \(G_n\) 中闭。特别地，\(A\) 是 \(K\) 的闭子集，且其每个子集都闭，所以 \(A\) 是紧的无限离散空间。这与单点开覆盖没有有限子覆盖矛盾。

	连续映射的紧像满足第一项。若定义域 \(X\) 紧，则 \(X\times I\) 也紧，将第一项用于同伦映射即可。
\end{proof}

判断一张从弱直极限出发的同伦是否连续，还需要同时处理时间参数。下面把这一点明确下来，后面的坐标移动将反复使用它。

\begin{lemma}\label{g2:lem:grassmann-parameter-continuity}
	设 \(Y=\bigcup_NY_N\) 取由各 \(Y_N\) 规定的直极限拓扑，\(L\) 为局部紧 Hausdorff 空间。映射 \(F:Y\times L\to Z\) 连续，当且仅当各限制 \(F|_{Y_N\times L}\) 连续。特别地，这适用于 \(Y=G_n\)、\(L=I\)，以及局部丛坐标中的 \(L=\mathbb K^n\)。
\end{lemma}
\begin{proof}
	必要性直接来自限制映射。反过来，给连续映射集 \(C(L,Z)\) 取紧开拓扑，其子基本开集为
	\[
		[C,O]=\{u:u(C)\subseteq O\},\qquad C\subseteq L\ \text{紧},\quad O\subseteq Z\ \text{开}.
	\]
	每个连续映射 \(F_N:Y_N\times L\to Z\) 的伴随映射 \(y\mapsto F_N(y,\cdot)\) 连续：若 \(F_N(\{y\}\times C)\subseteq O\)，紧性允许把 \(C\) 上的有限个乘积邻域合并，得到 \(y\) 的一个邻域，使同样的包含对该邻域内全部点成立。由 \(Y\) 的终拓扑性质，这些伴随映射合成连续映射 \(Y\to C(L,Z)\)。

	局部紧 Hausdorff 条件保证求值 \(C(L,Z)\times L\to Z\) 连续。具体地，若 \(u(t)\in O\)，可取 \(t\) 的开邻域 \(V\)，使 \(\overline V\) 紧且包含于 \(u^{-1}(O)\)；于是 \([\overline V,O]\times V\) 的求值落在 \(O\) 内。将求值与前述伴随映射复合，便得到 \(F\) 的连续性。
\end{proof}

现在验证无限典范丛仍然局部平凡。对任意有限坐标集 \(J\subset\mathbb N\)、\(|J|=n\)，记 \(A_J\) 为相应的坐标 \(n\) 平面，\(U_J\subseteq G_n\) 为 \(\pi_{A_J}|_V\) 可逆的子空间集合。它与各有限阶段的交都是开集，所以 \(U_J\) 开。任一 \(n\) 个线性无关向量构成的矩阵都有一个非零 \(n\) 阶子式，故这些 \(U_J\) 覆盖 \(G_n\)。有限阶段的平凡化
\[
	(V,v)\longmapsto(V,\pi_{A_J}v)
\]
相容地取并，给出 \(\gamma_n|_{U_J}\to U_J\times A_J\)。它的连续性由总空间的弱拓扑得到，逆映射在每个有限阶段乘 \(A_J\) 上连续，再由\cref{g2:lem:grassmann-parameter-continuity}得到整体连续。这里 \(U_J\) 本身取其与各阶段交所给的弱拓扑；这是开子空间上的终拓扑。于是 \(\gamma_n\) 是一张真正的秩 \(n\) 向量丛，且具有可数平凡化覆盖。

\begin{lemma}\label{g2:lem:universal-numerable-cover}
	无限典范丛 \(\gamma_n\to G_n\) 的可数坐标平凡化覆盖 \(\{U_J\}\) 有从属的局部有限连续单位分解。此外，在每个有限阶段中，只有有限个权重可能非零。
\end{lemma}
\begin{proof}
	枚举全部 \(n\) 元坐标集为 \(J_1,J_2,\ldots\)。对 \(V\in G_n\)，置
	\[
		f_j(V)=\det\bigl(P_V[J_j,J_j]\bigr).
	\]
	这是投影矩阵相应主子式。在一个包含 \(V\) 的有限阶段选正交单位标架矩阵 \(Q\)，则 \(P_V=QQ^*\)，所以 \(f_j(V)=|\det Q_{J_j}|^2\ge0\)。对 \(\det(Q^*Q)\) 作矩阵乘积的子式展开，重复行指标的项由行列式的交替性抵消，剩下每个 \(n\) 元坐标集的两个共轭子式之积，得到
	\[
		\sum_j f_j(V)=\det(Q^*Q)=1.
	\]
	在第 \(N\) 阶段，只有 \(J_j\subseteq\{1,\ldots,N\}\) 时该函数才可能非零；这样的 \(j\) 只有有限个。并且 \(f_j>0\) 的集合正是 \(U_{J_j}\)。

	原函数族的支集未必局部有限，因此先削去每一项的一小部分。令
	\[
		a_j=2^{-j-1},\qquad g_j(V)=\max\{0,f_j(V)-a_j\}.
	\]
	有 \(\operatorname{supp}g_j\subseteq\{f_j\ge a_j\}\subseteq U_{J_j}\)。对固定 \(V\)，取 \(j_0\) 包含所有在 \(V\) 处非零的 \(f_j\)。集合
	\[
		C=\bigcup_{j>j_0}\{f_j\ge a_j\}
	\]
	与每个有限阶段的交只是有限个闭集的并，故 \(C\) 在弱拓扑下闭；且 \(V\notin C\)。开邻域 \(G_n\setminus C\) 不遇到任何 \(j>j_0\) 的支集，证明了局部有限性。

	又因 \(\sum_j a_j=1/2\)，有 \(\sum_jg_j(V)\ge1/2>0\)。局部有限性使这个和连续。于是
	\[
		\rho_j(V)=\frac{g_j(V)}{\sum_i g_i(V)}
	\]
	就是所需单位分解。在每个有限阶段只有有限个 \(f_j\) 非零，相同结论也适用于 \(\rho_j\)。
\end{proof}

最后把有限嵌入构造用于可数局部有限数据。设 \(E\to X\) 有可数平凡化覆盖及从属局部有限单位分解，则式\eqref{g2:eq:weighted-bundle-embedding}按 \(i=1,2,\ldots\) 排列坐标，给出 \(E\to X\times\mathbb K^\infty\) 的逐纤维单射。在每个基点的某个邻域内，只有固定有限个权重可能非零，所以其分类映射局部落在一个固定有限阶段；前面的矩阵公式因此直接保证它连续。若 \(X\) 仿紧 Hausdorff 且有可数平凡化覆盖，先取局部有限细化与从属单位分解，再把属于同一原覆盖成员的权重相加，即得到这样的可数局部有限数据。局部有限支集的并仍闭，故分组后支集仍包含于对应平凡化域。这样，非紧基底也得到无限典范丛的拉回表示。

% !TeX root = ../../Book2.tex
% 纲要标题：### 19.3 分类定理的证明
% 纲要位置：几何学计划纲要.md，第 1753 行。
% 正文状态：未撰写；本文件目前是框架。

% 纲要内容开始
% 
% * 丛到分类映射的构造
% * 同伦类与选择无关性
% * 同伦映射给出同构丛
% * 适用基空间与数值型限制
% 
% 纲要内容结束

\section{分类定理的证明}
\label{g2:sec:1903}

待撰写

\subsection{丛到分类映射的构造}
\label{g2:subsec:1903:01}
待撰写

\subsection{同伦类与选择无关性}
\label{g2:subsec:1903:02}
待撰写

\subsection{同伦映射给出同构丛}
\label{g2:subsec:1903:03}
待撰写

\subsection{适用基空间与数值型限制}
\label{g2:subsec:1903:04}
待撰写
% 纲要条目排为编号小节标题，并进入目录与书签；说明段仍保留于纲要注释中。

% !TeX root = ../../Book2.tex
% 纲要标题：### 19.4 主丛、结构群约化与例子
% 纲要位置：几何学计划纲要.md，第 1760 行。
% 正文状态：未撰写；本文件目前是框架。

% 纲要内容开始
% 
% * 主丛分类和 BG
% * 约化与齐性纤维丛的截面
% * 实线丛、复线丛和低秩球面丛
% * 联络不是纯拓扑分类数据
% 
% ---
% 
% 纲要内容结束

\section{主丛、结构群约化与例子}
\label{g2:sec:1904}

待撰写

\subsection{主丛分类和 BG}
\label{g2:subsec:1904:01}
待撰写

\subsection{约化与齐性纤维丛的截面}
\label{g2:subsec:1904:02}
待撰写

\subsection{实线丛、复线丛和低秩球面丛}
\label{g2:subsec:1904:03}
待撰写

\subsection{联络不是纯拓扑分类数据}
\label{g2:subsec:1904:04}
待撰写
% 纲要条目排为编号小节标题，并进入目录与书签；说明段仍保留于纲要注释中。


\begin{chaptersummary}
待撰写
\end{chaptersummary}

\BookExercises
待撰写

% !TeX root = ../../Book2.tex
% 纲要标题：## 第20章　Thom 同构、Euler 类与 Gysin 序列
% 纲要位置：几何学计划纲要.md，第 1769 行。
% 正文状态：初稿；尚未正式审核。

% 纲要内容开始
% 
% > 本章目标与证明范围：从局部纤维定向构造 Thom 类，在相对上同调中证明同构。随后以零截面与球丛解释 Euler 类和 Gysin，不把这些对象分别当作独立定义表。
% 
% 纲要内容结束

\chapter{Thom 同构、Euler 类与 Gysin 序列}
\label{g2:ch:20}
\BookChapterNumber{g2:ch:20}

\begin{chapterabstract}
	纤维的局部方向可以组成总空间上的相对上同调类。本章从盘球丛、Thom 空间及定向局部系统出发，证明任意 CW 基空间上带局部系数的 Thom 同构，说明有限覆盖粘合与无限胞腔的次数控制如何配合。零截面拉回得到 Euler 类，横截截面的零点流形及孤立零点指标给出它的几何解释，并与前文的球丛传递规范比较。最后由盘球丛对的长正合列推导 Gysin 序列，计算 Hopf 丛和圆丛，再用微分形式建立沿纤维积分及其符号约定。
\end{chapterabstract}

\begin{learninggoals}
	\begin{itemize}
		\item 比较盘球丛对、穿孔总空间对与 Thom 商空间，辨认非紧基底及扭系数时各模型的适用范围。
		\item 用纤维的局部同调描述方向输运，解释模二系数和定向局部系统各自如何处理方向反转。
		\item 完整证明局部乘积与整体 Thom 同构，说明无限 CW 基空间为何由有限个维数层控制固定次数。
		\item 从零截面构造 Euler 类，证明它与首障碍、横截零点类及局部指标和的关系。
		\item 用一个具体三维丛辨认 Euler 类消失与非零截面存在之间的区别。
		\item 从长正合列推导 Gysin 序列，计算球丛上同调，并将其积分映射与微分形式的纤维积分对应。
	\end{itemize}
\end{learninggoals}

球面上的向量场为什么不能总是避开零点？第12章把问题放在逐胞腔延拓中，用边界球面的映射度记录首个障碍；第16—17章则把孤立零点附近的映射度相加，得到 Euler 示性数。两种计算都出现同一种局部整数：绕着一个零点走一圈，向量方向绕纤维球面多少次。本章将这个整数放进一张向量丛的相对上同调，从而把延拓、零点和整体上同调联系起来。

\ProseParagraphBreak

先看平凡丛 \(B\times\mathbb R^r\)。每根纤维都带有圆盘对 \((D^r,S^{r-1})\) 的正生成元 \(u_r\)，于是底空间上的类 \(a\) 可以变成相对类 \(a\times u_r\)，次数升高 \(r\)。若丛有扭转，局部正生成元可能在交叠处变号；选择全局方向，或让系数也按同一符号运输，便能把它们拼成 Thom 类。Thom 同构断言，这个升次数构造已经给出盘球丛对的全部上同调。

\ProseParagraphBreak

随后只需对同一个 Thom 类作两种操作。沿零截面拉回得到 Euler 类，它记录一个截面的零点；把相对上同调放回盘球丛对的长正合列，得到 Gysin\footnote{Werner Gysin（1915--1998），瑞士数学家，研究流形纤维化的同调，提出后来以其名字命名的映射与正合列。}序列，其中一支箭头正是乘 Euler 类。这样，三项构造共享同一个局部生成元，而每个映射都能追溯到一个具体的空间对或截面。

\ProseParagraphBreak

本章主要对象是 CW 复形上的有限秩实向量丛；光滑零点与纤维积分部分另说明紧性、横截性及定向条件。第5章的局部系数、第6章的对偶与帽积、第7章的 de Rham 比较、第12章的首障碍及第19章的丛同伦不变性会在用到时逐项回引。Thom 构造与平方运算的联系可参阅\cite[第23章第5--6节]{May1999}，Euler 类与 Thom 同构的经典组织可参阅\cite[第9节，Oriented Bundles and the Euler Class；第10节，The Thom Isomorphism Theorem]{MilnorStasheff1974}。第21章将继续比较由 Thom 平方产生的类与其他示性类构造。

% !TeX root = ../../Book2.tex
% 纲要标题：### 20.1 Thom 空间及定向
% 纲要位置：几何学计划纲要.md，第 1773 行。
% 正文状态：初稿；尚未正式审核。

% 纲要内容开始
% 
% * 盘丛、球丛与相对模型
% * 局部纤维生成元
% * 系数环及扭定向系统
% 
% 纲要内容结束

\section{Thom 空间及定向}
\label{g2:sec:2001}

一根实 \(r\) 维纤维本身可缩，它的普通上同调看不出维数；把零点以外的部分作为相对项后，便出现一个 \(r\) 维生成元。本章把这个局部生成元沿整张向量丛组织起来。首先需要说明采用哪一个相对空间，以及纤维坐标反转方向时生成元如何变化。

\subsection{盘丛、球丛与相对模型}
\label{g2:subsec:2001:01}
本章设 \(B\) 为 CW 复形，\(p:E\to B\) 为秩 \(r\ge1\) 的实向量丛，零截面记为 \(s_0\)。\cref{g2:lem:cw-paracompact-partitions}保证 \(B\) 仿紧；将局部 Euclid 内积按从属单位分解求和，就得到连续纤维内积。这是\secref{g2:sec:1902}构造丛嵌入时已经使用的方法，局部有限性使它同样适用于无限覆盖。以下固定这样一个内积，并记 \(\|v\|\) 为相应范数。

\begin{definition}\label{g2:def:disk-sphere-thom}
	向量丛 \(E\) 的\term[pancong]{盘丛}[disk bundle]、\term[qiucong]{球丛}[sphere bundle]分别为
	\[
		D(E)=\{v\in E:\|v\|\le1\},\qquad S(E)=\{v\in E:\|v\|=1\}.
	\]
	把整个 \(S(E)\) 压成同一个基点 \(*\)，所得带基点空间
	\[
		\operatorname{Th}(E)=D(E)/S(E)
	\]
	称为 \(E\) 的\term[thom-kongjian]{Thom 空间}[Thom space]。另记 \(E^\times=E\setminus s_0(B)\)。
\end{definition}

在局部正交标架中，盘丛和球丛分别是 \(U\times D^r\) 与 \(U\times S^{r-1}\)。Thom 商把所有这些纤维的边界共同送到 \(*\)；各纤维内部的点仍然保留。所以它同时记住纤维的维数及纤维之间的粘合方式。对基底只有一点的情形，这一构造就是 \(D^r/S^{r-1}\cong S^r\)。

\begin{proposition}\label{g2:prop:thom-relative-models}
	对上述向量丛，\((D(E),S(E))\) 是 NDR 对。包含映射
	\[
		(D(E),S(E))\longrightarrow(E,E^\times)
	\]
	在同调和上同调中诱导同构，系数可以取从 \(B\) 拉回的局部系统。对常系数环 \(R\)，商映射还诱导自然同构
	\begin{equation}\label{g2:eq:thom-relative-quotient}
		\widetilde H^k(\operatorname{Th}(E);R)\xrightarrow{\ \cong\ }H^k(D(E),S(E);R).
	\end{equation}
	约化同调也有相同的商空间比较。
\end{proposition}
\begin{proof}
	先给出控制球面附近的径向同伦。对 \(v\in D(E)\)，令 \(\rho=\|v\|\)，并置
	\[
		u(v)=\min\{1,2(1-\rho)\},\qquad
		H(v,t)=\left(1-t+\frac{t}{\max\{\rho,1/2\}}\right)v.
	\]
	分母处处为正，故包括零截面在内该式连续。其范数始终不超过一；当 \(\rho=1\) 时同伦固定 \(v\)，当 \(u(v)<1\) 时有 \(\rho>1/2\)，从而 \(\|H(v,1)\|=1\)。又 \(u^{-1}(0)=S(E)\)。这逐项满足\cref{g2:def:NDR-pair}的条件，所以球丛包含是\cref{g2:thm:NDR-cofibration}所述的闭余纤维化。

	径向映射 \(v\mapsto v/\max\{1,\|v\|\}\) 给出 \(E\) 到 \(D(E)\) 的强形变收缩；线性插值保持每根纤维，且固定盘丛。同样，\(v\mapsto v/\|v\|\) 给出 \(E^\times\) 到 \(S(E)\) 的强形变收缩。因此包含映射在两个空间和两个相对子空间的同调、上同调上分别诱导同构。把这些映射放入对的长正合列，五引理给出所需的相对同构。所有同伦均保持 \(p(v)\) 不变，故拉回系数的识别也是固定的；局部系数情形使用\cref{g2:thm:twisted-homotopy-invariance}的同伦不变性。

	最后核对商空间比较。取球面附近的开领圈
	\[
		N=\{v\in D(E):\|v\|>1/2\}.
	\]
	上面的同伦限制为 \(N\) 到 \(S(E)\) 的强形变收缩，且在商上把 \(N/S(E)\) 收缩到 \(*\)。诱导同伦连续，是因为商映射与局部紧区间 \(I\) 的乘积仍为商映射，见\cref{g2:lem:quotient-interval}。于是长正合列给出
	\[
		\begin{aligned}
			H^k(D(E),N;R)&\xrightarrow{\ \cong\ }H^k(D(E),S(E);R),\\
			H^k(\operatorname{Th}(E),N/S(E);R)&\cong\widetilde H^k(\operatorname{Th}(E);R).
		\end{aligned}
	\]
	商映射在 \(D(E)\setminus S(E)\) 上是到 \(\operatorname{Th}(E)\setminus\{*\}\) 的同胚。分别从上述两个对中切除 \(S(E)\) 和 \(\{*\}\)，所得对因而同胚。切除条件成立：这两个闭集分别包含于开集 \(N\) 和 \(N/S(E)\) 内。由\cref{g2:prop:relative-cohomology-exactness}的上同调切除，两个对的上同调相同。合并这些同构便得到式\eqref{g2:eq:thom-relative-quotient}，而且合成映射正是商映射诱导的映射。同调的证明逐步相同，所用切除是\cref{g2:thm:singular-excision}。
\end{proof}

盘球丛模型适合使用长正合列，\((E,E^\times)\) 适合讨论零截面附近的几何，而 Thom 空间把相对问题写成一个带基点空间的约化上同调。以上论证说明这些模型为何给出同一份常系数信息。局部系数将在相对模型上保留：从 \(B\) 拉回的局部系统一般不能在压缩整个 \(S(E)\) 后成为商空间上的局部系统。

\ProseParagraphBreak

内积的选择不会改变 Thom 空间的同胚类型。若两组范数为 \(\|\cdot\|_0,\|\cdot\|_1\)，则
\[
	v\longmapsto
	\begin{cases}
		\dfrac{\|v\|_0}{\|v\|_1}v,&v\ne0,\\
		0,&v=0
	\end{cases}
\]
把第一组盘球丛送到第二组盘球丛，交换下标便得逆。在任何局部标架中，两组正定内积在每个基点附近具有统一的范数比较常数，所以该映射在零截面也连续。它因此诱导 Thom 空间的同胚。若 \(B\) 紧，则 \(D(E)\) 紧，且 \(v\mapsto v/(1-\|v\|)\) 将开盘丛同胚到 \(E\)，这时 \(\operatorname{Th}(E)\) 正是总空间 \(E\) 的一点紧化。非紧基底时使用上述盘球丛商定义；\cref{g2:ex:noncompact-thom-not-compactification}将具体比较两种构造。

\subsection{局部纤维生成元}
\label{g2:subsec:2001:02}
固定 \(b\in B\)，记 \(D_b=D(E)\cap E_b\)、\(S_b=S(E)\cap E_b\)。相对球的计算给出
\begin{equation}\label{g2:eq:fiber-orientation-group}
	H_j(D_b,S_b;\mathbb Z)\cong
	\begin{cases}
		\mathbb Z,&j=r,\\
		0,&j\ne r.
	\end{cases}
\end{equation}
这也可由 \(D_b/S_b\cong S^r\) 及\cref{g2:thm:sphere-homology}得到。选一组有序线性坐标，就选出顶维相对基本类 \(o_b=[D_b,S_b]\)；反转坐标方向时它变成 \(-o_b\)。由上一小节的相对比较，它也确定 \(H_r(E_b,E_b\setminus\{0\};\mathbb Z)\) 的一个生成元。

\ProseParagraphBreak

对于任意 Abel 群 \(A\)，上同调万有系数定理给出自然识别
\begin{equation}\label{g2:eq:fiber-top-cohomology-hom}
	H^r(D_b,S_b;A)\cong\operatorname{Hom}_{\mathbb Z}\bigl(H_r(D_b,S_b;\mathbb Z),A\bigr),
\end{equation}
其余次数的相对上同调为零。这里顶维相对同调是自由群，且其余相对同调消失，所以万有系数定理中的 \(\operatorname{Ext}\) 项为零，见式\eqref{g2:eq:uct-cohomology}。当 \(A=\mathbb Z\) 时，与 \(o_b\) 对偶的生成元 \(u_b\) 由
\[
	\langle u_b,o_b\rangle=1
\]
刻画。方向改变使 \(o_b\) 和 \(u_b\) 同时变号。这一配对将是把局部方向转成 Thom 类的标准化条件。

\ProseParagraphBreak

球面取向采用外法向在前的边界约定。长正合列的边界映射于是满足 \(\partial o_b=[S_b]\)。秩一时这个公式读成
\[
	\partial[-1,1]=[1]-[-1]\in\widetilde H_0(S^0;\mathbb Z).
\]
若球面的约化顶维上同调类 \(\omega_b\) 在 \([S_b]\) 上取值一，则对的连接映射满足 \(\delta\omega_b=u_b\)：在链上有 \(\langle\delta\omega_b,o_b\rangle=\langle\omega_b,\partial o_b\rangle=1\)。这个约定固定了后面球丛长正合列中纤维积分的正负号。

\begin{definition}\label{g2:def:bundle-orientation-system}
	向量丛 \(E\) 的\term[xiangliangcong-dingxiang-jubuxitong]{定向局部系统}[orientation local system]记为 \(\mathcal O_E\)，其在 \(b\) 处的群为
	\[
		\mathcal O_{E,b}=H_r(E_b,E_b\setminus\{0\};\mathbb Z).
	\]
	局部平凡化给出这些群之间的识别，沿路径依次使用局部识别便得到输运。若可为每个 \(b\) 选取生成元 \(o_b\in\mathcal O_{E,b}\)，使它们在局部平凡化中相容，则称 \(E\) \term[xiangliangcong-kedingxiang]{可定向}[orientable]；这样一组选择称为 \(E\) 的一个定向。
\end{definition}

这一定向描述每根向量空间纤维的方向，基底本身不需要是流形，也不需要已有定向。它与\chapref{g2:ch:05}中流形的定向局部系统使用同一种局部同调思想；此处取的是纤维在零点的局部同调。

\begin{proposition}\label{g2:prop:bundle-orientation-transport}
	若两组局部坐标的过渡函数为 \(g_{\alpha\beta}:U_\alpha\cap U_\beta\to\operatorname{GL}_r(\mathbb R)\)，则 \(\mathcal O_E\) 的相应过渡函数为
	\[
		\operatorname{sgn}\det g_{\alpha\beta}\in\{1,-1\}.
	\]
	这些函数定义秩一的整数局部系统，其路径输运只依赖端点固定的路径同伦类。在每个路径连通分支上，\(E\) 可定向当且仅当所有闭路在 \(\mathcal O_E\) 上的输运都是恒等。
\end{proposition}
\begin{proof}
	先考察一个可逆线性映射 \(A:\mathbb R^r\to\mathbb R^r\)。极分解把它写成正交矩阵与正定矩阵的乘积；正定矩阵可沿 \((1-t)P+tI\) 连到恒等。行列式为正的正交矩阵由二维平面旋转组成，可以把各旋转角连续缩到零。因此 \(\det A>0\) 时，\(A\) 与恒等映射通过可逆线性映射同伦，在相对顶维同调上作用为 \(+1\)。若 \(\det A<0\)，先复合一次坐标反射便化为前一种情形；反射在相对基本类上作用为 \(-1\)，与\cref{g2:prop:reflection-antipodal-degree}的球面度计算一致。故所求作用正是 \(\operatorname{sgn}\det A\)。

	行列式连续且不为零，所以其符号在交集上局部常值。过渡函数的乘法关系经行列式符号仍成立，因而确实得到局部系统。用这些局部常值识别沿路径输运，细分路径或改变局部平凡化均由过渡关系保证不改变合成；对路径同伦，把参数方形细分到每个小方形落入一个平凡化域，小方形内的输运相消，得到端点固定的同伦不变性。这也就是\secref{g2:sec:0501}中局部系统的路径描述。

	若已有相容生成元，任何闭路都把它送回自身，故输运恒等。反过来，在一个分支中固定 \(b_0\) 及一个生成元 \(o_{b_0}\)，沿从 \(b_0\) 到 \(b\) 的路径输运它。两条路径之差是一条闭路，所以假设保证所得 \(o_b\) 与路径无关；局部平凡化中的输运恒定，又保证这些生成元相容。CW 复形局部道路连通，因此上述逐分支的构造给出所需定向。
\end{proof}

\begin{example}\label{g2:exm:complex-bundle-real-orientation}
	每张复秩 \(m\) 向量丛的底层实秩 \(2m\) 丛都有自然定向。对一组复基 \(e_1,\ldots,e_m\)，取实有序基
	\[
		(e_1,\mathrm i e_1,e_2,\mathrm i e_2,\ldots,e_m,\mathrm i e_m).
	\]
	若复坐标变换矩阵为 \(A\)，其底层实矩阵的行列式为 \(\det_{\mathbb R}A=|\det_{\mathbb C}A|^2>0\)。这个公式可在复化后看出：实矩阵的复化分解为 \(A\) 与 \(\overline A\) 两块，行列式为 \(\det A\cdot\overline{\det A}\)。因此更换复标架不会改变上述实方向，局部方向自动粘合。
\end{example}

\subsection{系数环及扭定向系统}
\label{g2:subsec:2001:03}
取交换有幺环 \(R\)，并令
\[
	\mathcal O_{E,R}=\mathcal O_E\otimes_{\mathbb Z}R.
\]
这是秩一自由 \(R\) 模的局部系统，过渡函数仍为 \(\pm1\)，现将它们视作 \(R\) 中的单位。一个 \(R\) 定向是 \(\mathcal O_{E,R}\) 与常局部系统 \(\underline R\) 的同构，等价地，是一组相容的纤维 \(R\) 模基元。整数定向经系数变换给出 \(R\) 定向；若 \(2=0\) 于 \(R\)，则两个整数方向都给出同一个模 \(R\) 基元，故每张实向量丛都有规范的 \(R\) 定向。特别地，模二系数完全消除了方向反转造成的符号。

\begin{example}\label{g2:exm:mobius-orientation-system}
	设 \(L\to S^1\) 为 Möbius 线丛，其矩形模型沿两端按 \((0,t)\sim(1,-t)\) 粘合。沿基圆一周，纤维的相对基本类变成其相反数，所以 \(\mathcal O_L\) 的单值作用为 \(-1\)，它没有整数定向。若一个类 \(u\in H^1(D(L),S(L);\mathbb Z)\) 在每根纤维上都是生成元，则沿基圆输运后其限制必须仍为自身，同时又要变号。这要求一个整数生成元等于其相反数，矛盾。用 \(\mathbb F_2\) 系数时两个符号相同，方向障碍消失。\secref{g2:sec:1304}已经利用这一模二 Thom 类研究了平方运算，见\cref{g2:exm:mobius-thom-square}。
\end{example}

为了对所有实丛统一处理整数系数，可以让系数也随纤维方向一起变化。固定 \(b\)，式\eqref{g2:eq:fiber-top-cohomology-hom}给出规范识别
\begin{equation}\label{g2:eq:twisted-fiber-endomorphism}
	H^r(D_b,S_b;\mathcal O_{E,R,b})
	\cong\operatorname{Hom}_{\mathbb Z}(\mathcal O_{E,b},\mathcal O_{E,R,b})
	\cong\operatorname{End}_R(\mathcal O_{E,R,b}).
\end{equation}
第二个同构把群同态唯一延拓为 \(R\) 线性映射。右端有一个不依赖任何基元的元素：恒等映射。整数系数时它就是 \(\operatorname{id}_{\mathcal O_{E,b}}\)；一般系数时，在中间一项对应 \(m\mapsto m\otimes1\)。

\begin{definition}\label{g2:def:twisted-thom-class}
	类
	\[
		U_E\in H^r(D(E),S(E);p^*\mathcal O_{E,R})
	\]
	若在每个纤维上的限制都对应式\eqref{g2:eq:twisted-fiber-endomorphism}中的恒等映射，就称为 \(E\) 的\term[niu-thom-lei]{扭 Thom 类}[twisted Thom class]。需要区分系数环时记为 \(U_{E,R}\)。若选定 \(E\) 的 \(R\) 定向 \(\mathcal O_{E,R}\cong\underline R\)，由它得到的常系数类记为 \(u_E\)，称为这个定向的 Thom 类；其纤维限制在选定的相对基本类上取值一。
\end{definition}

这一标准化把两次变号配成了方向无关的条件。具体地，选一个整数方向 \(o_b\)，并令 \(u_b(o_b)=1\)，则局部扭类写成 \(u_b\otimes(o_b\otimes1)\)。更换方向时两个因子都变号，张量保持不变；在式\eqref{g2:eq:twisted-fiber-endomorphism}中，这正是同一个恒等映射。在 Möbius 例中，它也沿闭路保持不变。局部条件能否确定一个全局上同调类，需要进一步的粘合证明；下一节的\cref{g2:thm:thom-isomorphism}将证明这样的类存在、唯一，并说明乘以它为何给出次数平移。

\begin{proposition}\label{g2:prop:bundle-orientation-square}
	存在规范的局部系统同构
	\begin{equation}\label{g2:eq:bundle-orientation-square}
		\mathcal O_{E,R}\otimes_R\mathcal O_{E,R}\xrightarrow{\ \cong\ }\underline R.
	\end{equation}
	在任一整数方向 \(o_b\) 下，它将 \((o_b\otimes a)\otimes(o_b\otimes c)\) 送到 \(ac\)。
\end{proposition}
\begin{proof}
	每个纤维都是秩一自由 \(R\) 模，所以该规则给出一个同构。把 \(o_b\) 换成 \(-o_b\) 时，同一元素的两个系数都变成相反数，乘积不变，故映射与方向选择无关。沿路径输运时，两个定向因子具有同一符号，其乘积为 \((\pm1)^2=1\)，因而这些纤维同构与输运相容，给出局部系统同构。
\end{proof}

这与\cref{g2:prop:orientation-square}中流形定向系统的平方配对是同一种构造。它解释了后面公式中为什么会出现一个或两个定向因子：乘一个扭 Thom 类会添上 \(\mathcal O_{E,R}\)，两个相同的定向因子则可按式\eqref{g2:eq:bundle-orientation-square}规范消去。此后的杯积按“底空间类在前、Thom 类在后”排列，记作 \(p^*a\smile U_E\)。这一顺序同时确定次数交换可能带来的符号。
% 纲要条目排为编号小节标题，并进入目录与书签；说明段仍保留于纲要注释中。

% !TeX root = ../../Book2.tex
% 纲要标题：### 20.2 Thom 同构的证明
% 纲要位置：几何学计划纲要.md，第 1729 行。
% 正文状态：未撰写；本文件目前是框架。

% 纲要内容开始
% 
% * 局部乘积上的杯积同构
% * Mayer–Vietoris 粘合
% * 一般基空间的有限维或穷竭控制
% * 自然性与直和相容
% 
% 纲要内容结束

\section{Thom 同构的证明}
\label{g2:sec:2002}

待撰写

\subsection{局部乘积上的杯积同构}
\label{g2:subsec:2002:01}
待撰写

\subsection{Mayer–Vietoris 粘合}
\label{g2:subsec:2002:02}
待撰写

\subsection{一般基空间的有限维或穷竭控制}
\label{g2:subsec:2002:03}
待撰写

\subsection{自然性与直和相容}
\label{g2:subsec:2002:04}
待撰写
% 纲要条目排为编号小节标题，并进入目录与书签；说明段仍保留于纲要注释中。

% !TeX root = ../../Book2.tex
% 纲要标题：### 20.3 Euler 类与零点
% 纲要位置：几何学计划纲要.md，第 1736 行。
% 正文状态：未撰写；本文件目前是框架。

% 纲要内容开始
% 
% * 零截面拉回
% * 横截截面的零点类
% * 局部指标与 Euler 数
% * 零 Euler 类不足以消除所有高阶障碍
% 
% 纲要内容结束

\section{Euler 类与零点}
\label{g2:sec:2003}

待撰写

\subsection{零截面拉回}
\label{g2:subsec:2003:01}
待撰写

\subsection{横截截面的零点类}
\label{g2:subsec:2003:02}
待撰写

\subsection{局部指标与 Euler 数}
\label{g2:subsec:2003:03}
待撰写

\subsection{零 Euler 类不足以消除所有高阶障碍}
\label{g2:subsec:2003:04}
待撰写
% 纲要条目排为编号小节标题，并进入目录与书签；说明段仍保留于纲要注释中。

% !TeX root = ../../Book2.tex
% 纲要标题：### 20.4 Gysin 序列及积分沿纤维
% 纲要位置：几何学计划纲要.md，第 1793 行。
% 正文状态：未撰写；本文件目前是框架。

% 纲要内容开始
% 
% * 从盘球丛对的长正合列推导
% * 乘 Euler 类的映射
% * Hopf 丛及球丛计算
% * 定向纤维积分的上同调模型
% 
% ---
% 
% 纲要内容结束

\section{Gysin 序列及积分沿纤维}
\label{g2:sec:2004}

待撰写

\subsection{从盘球丛对的长正合列推导}
\label{g2:subsec:2004:01}
待撰写

\subsection{乘 Euler 类的映射}
\label{g2:subsec:2004:02}
待撰写

\subsection{Hopf 丛及球丛计算}
\label{g2:subsec:2004:03}
待撰写

\subsection{定向纤维积分的上同调模型}
\label{g2:subsec:2004:04}
待撰写
% 纲要条目排为编号小节标题，并进入目录与书签；说明段仍保留于纲要注释中。


\begin{chaptersummary}
	盘球丛对把每根纤维的方向记录为一个最高次相对生成元；定向局部系统同时保存这些生成元沿道路的符号变化。局部乘积给出次数平移，有限覆盖通过 Mayer--Vietoris 粘合，一般 CW 基空间则由相对胞腔复形控制。所得 Thom 类唯一，并与拉回、直和及系数变换相容；任意局部系数的上同调都可由与它作杯积移入盘球丛对。

	\ProseParagraphBreak

	零截面拉回 Thom 类得到 \(e_{\mathrm{Th}}(E)\)。它等于球丛的首障碍，也等于横截零点流形的对偶类；秩与底流形维数相等时，求值就是孤立零点的指标和。按前文已固定的传递符号，它与球丛的 \(e(S(E))\) 相差一个负号。Euler 类消失只消除首障碍，高维胞腔还可能阻止非零截面存在。

	\ProseParagraphBreak

	盘球丛对的长正合列经 Thom 同构改写为 Gysin 序列：忘去相对条件对应乘 Euler 类，球丛上的类经连接映射再作 Thom 逆同构，得到降低纤维维数的积分映射。Hopf 丛与二球面上的圆丛展示了它的计算作用。按底空间微分在前、纤维微分在后的顺序，形式积分满足无附加符号的投影公式，并与规范化后的拓扑积分一致；整数序列还保留实微分形式看不到的扭信息。
\end{chaptersummary}

\BookExercises

\begin{exercise}\label{g2:ex:noncompact-thom-not-compactification}
	取 \(\mathbb R\) 上的平凡实线丛 \(E=\mathbb R\times\mathbb R\)。直接从盘球丛模型证明 \(\operatorname{Th}(E)\simeq S^1\)，并计算其约化整数上同调。再与总空间 \(E\) 的一点紧化比较，说明非紧基底时两种构造确实可能不同。
\end{exercise}
\begin{solution}
	盘丛和球丛分别是 \(\mathbb R\times[-1,1]\) 与 \(\mathbb R\times\{-1,1\}\)。同伦
	\[
		F_s(x,t)=((1-s)x,t),\qquad 0\le s\le1,
	\]
	始终把球丛映入球丛，固定中间纤维 \(\{0\}\times[-1,1]\)，所以下降为 Thom 空间上的同伦。下降后的同伦连续，由\cref{g2:lem:quotient-interval}的积商性质保证。终映射的像为中间纤维的商 \([-1,1]/\{-1,1\}\cong S^1\)，因而得到一个强形变收缩。于是
	\[
		\widetilde H^k(\operatorname{Th}(E);\mathbb Z)\cong
		\begin{cases}
			\mathbb Z,&k=1,\\
			0,&k\ne1.
		\end{cases}
	\]
	总空间则同胚于 \(\mathbb R^2\)，其一点紧化为 \(S^2\)，约化整数上同调的非零项在二次。因此它与 \(\operatorname{Th}(E)\) 连同伦等价都不成立。盘球丛商只将纤维的单位边界共同压成一点；总空间的一点紧化还把沿非紧基底逃向无穷远的方向纳入无穷远点的邻域条件。
\end{solution}

\begin{exercise}\label{g2:ex:doubled-bundle-orientation}
	设 \(E\to B\) 为实秩 \(r\) 向量丛。对任意局部实标架 \(e_1,\ldots,e_r\)，用有序标架
	\[
		((e_1,0),\ldots,(e_r,0),(0,e_1),\ldots,(0,e_r))
	\]
	给 \(E\oplus E\) 定向。证明该定向不依赖局部标架，所以即使 \(E\) 不可定向，\(E\oplus E\) 仍有这一规范定向。再用 \((v,w)\mapsto v+\mathrm i w\) 识别它与复化丛 \(E\otimes_{\mathbb R}\mathbb C\) 的底层实丛，比较题给定向与复结构的自然定向。
\end{exercise}
\begin{solution}
	若原标架的变换矩阵是 \(A\in\operatorname{GL}_r(\mathbb R)\)，则直和中这组有序标架的变换矩阵为 \(\operatorname{diag}(A,A)\)，行列式是 \((\det A)^2>0\)。所以所有局部方向相容，且更换任何局部标架都不改变定向。这一论证不要求 \(\det A\) 本身为正。

	在题给识别下，直和定向对应实有序基
	\[
		(e_1,\ldots,e_r,\mathrm i e_1,\ldots,\mathrm i e_r),
	\]
	而复结构的自然定向对应交错排列 \((e_1,\mathrm i e_1,\ldots,e_r,\mathrm i e_r)\)。把前一个排列变为后一个时，\(\mathrm i e_j\) 需要跨过 \(r-j\) 个实基向量，总交换次数为
	\[
		(r-1)+(r-2)+\cdots+1=\frac{r(r-1)}2.
	\]
	故两个定向的比较符号为 \((-1)^{r(r-1)/2}\)：\(r\equiv0,1\pmod4\) 时相同，\(r\equiv2,3\pmod4\) 时相反。两种定向都与标架无关，但所采用的有序基排列不同。
\end{solution}

\begin{exercise}\label{g2:ex:orientation-over-z-four}
	令 \(L\to S^1\) 为 Möbius 线丛。分别对 \(R=\mathbb Z\)、\(\mathbb Z/4\)、\(\mathbb F_2\)，计算 \(H^0(S^1;\mathcal O_{L,R})\)，并判断 \(L\) 是否 \(R\) 可定向。说明存在非零平行截面为何还不足以给出 \(R\) 定向。
\end{exercise}
\begin{solution}
	这三个定向系统沿正向圆周的单值作用都是乘以 \(-1\)。\cref{g2:prop:circle-local-coefficients}给出圆周局部系数零次上同调为单值作用的不变量，因此
	\[
		H^0(S^1;\mathcal O_{L,R})=\ker(-1-1:R\longrightarrow R)=\{a\in R:2a=0\}.
	\]
	在三种系数中依次得到 \(0\)、\(\{0,2\}\subset\mathbb Z/4\)、\(\mathbb F_2\)。一个 \(R\) 定向要求平行截面的值为秩一自由 \(R\) 模的基元，选定一个局部基后，这等价于该值是 \(R\) 中的单位。

	整数系数没有非零平行截面，所以不可定向。\(\mathbb Z/4\) 系数虽有非零平行截面 \(2\)，但它不是单位，所生成的子模只有 \(\{0,2\}\)，不能成为纤维模的一组基；故仍不可定向。\(\mathbb F_2\) 系数的平行截面 \(1\) 是单位，因而给出规范定向。第二个例子说明，在带有零因子的系数环上，“非零”与“生成整个秩一自由模”是两个不同的条件。
\end{solution}

\begin{exercise}\label{g2:ex:mobius-thom-cohomology}
	设 \(L\to S^1\) 为 Möbius 线丛。分别计算 \(H^k(D(L),S(L);\mathbb Z)\)、\(H^k(D(L),S(L);p^*\mathcal O_L)\) 与 \(H^k(D(L),S(L);\mathbb F_2)\)。再用盘球丛对的长正合列独立验证第一组结果，并解释其二次整数类为什么不能充当 Thom 类。
\end{exercise}
\begin{solution}
	由式\eqref{g2:eq:twisted-thom-isomorphism}及 \(\mathcal O_L\otimes\mathcal O_L\cong\mathbb Z\)，普通整数相对上同调对应 \(H^{k-1}(S^1;\mathcal O_L)\)，扭系数相对上同调对应 \(H^{k-1}(S^1;\mathbb Z)\)。定向系统的圆周胞腔上链为
	\[
		\mathbb Z\xrightarrow{\,-2\,}\mathbb Z,
	\]
	所以它的零次上同调为零、一次为 \(\mathbb Z/2\)。于是
	\[
	\begin{array}{c|ccc}
		&k=1&k=2&\text{其余 }k\\\hline
		H^k(D(L),S(L);\mathbb Z)&0&\mathbb Z/2&0\\
		H^k(D(L),S(L);p^*\mathcal O_L)&\mathbb Z&\mathbb Z&0\\
		H^k(D(L),S(L);\mathbb F_2)&\mathbb F_2&\mathbb F_2&0
	\end{array}
	\]
	最后一行使用模二方向系统平凡及圆周的模二上同调。

	独立验证时，\(D(L)\) 收缩到基圆，\(S(L)\) 也是一个圆，而其投影绕基圆两次。故限制 \(H^1(D(L);\mathbb Z)\to H^1(S(L);\mathbb Z)\) 为乘二；两者的零次限制为恒等。长正合列于是给出
	\[
		0\longrightarrow H^1(D(L),S(L);\mathbb Z)\longrightarrow\mathbb Z
		\xrightarrow{\,2\,}\mathbb Z\longrightarrow H^2(D(L),S(L);\mathbb Z)\longrightarrow0,
	\]
	以及其他次数消失，与上表相符。秩一丛的 Thom 类须是一次类，并在每根区间对上取生成元；这里普通整数一次相对群已经为零。二次的扭元素限制到纤维后为零，因为纤维对的二次上同调为零，所以它不能替代 Thom 类。
\end{solution}

\begin{exercise}\label{g2:ex:infinite-thom-cochains}
	令 \(B=\bigvee_{j\ge1}S^1_j\) 具有 CW 弱拓扑，\(E=B\times\mathbb R^2\) 取标准定向。计算 \(\widetilde H^k(\operatorname{Th}(E);\mathbb Z)\)。找出一个三次类，使它在每个圆周上限制后都非零，并说明用有限子楔和的上同调按补零映射取直极限会漏掉什么。
\end{exercise}
\begin{solution}
	\(B\) 有一个零胞腔及可数多个一胞腔，胞腔边界全为零。因此
	\[
		H^0(B;\mathbb Z)=\mathbb Z,\qquad H^1(B;\mathbb Z)=\prod_{j\ge1}\mathbb Z,
	\]
	更高次为零。Thom 同构与常系数商比较给出
	\[
		\widetilde H^k(\operatorname{Th}(E);\mathbb Z)=
		\begin{cases}
			\mathbb Z,&k=2,\\
			\prod_{j\ge1}\mathbb Z,&k=3,\\
			0,&\text{其他次数}.
		\end{cases}
	\]
	取一次胞腔上链 \(a=(1,1,\ldots)\)。它是上循环且不为余边界，\(\Phi_E(a)\) 在限制到任意 \(S^1_j\) 上的丛时都是三次群的生成元，这由 Thom 自然性可知。

	前 \(N\) 个圆周的上同调为 \(\mathbb Z^N\)。若按题给补零映射取直极限，得到的是 \(\bigoplus_{j\ge1}\mathbb Z\)，只含有限支撑序列；它不含 \((1,1,\ldots)\)。实际由子空间包含诱导的上同调映射是反向的限制，而非补零映射。胞腔计算使用全部胞腔系数的直积，因而保留了这个类，也说明固定次数的有限骨架控制不等于有限胞腔支撑。
\end{solution}

\begin{exercise}\label{g2:ex:thom-space-product-euler}
	设 \(E\to B\) 为已定向实秩 \(r\ge1\) 丛，取系数环 \(R\)，记 \(e=e_{\mathrm{Th}}(E)\)。证明 Thom 同构一般只是 \(H^*(B;R)\) 模的次数平移同构；其乘法满足
	\[
		\Phi_E(a)\smile\Phi_E(b)=(-1)^{r\deg b}\Phi_E(a\smile b\smile e).
	\]
	特别计算 \(u_E^2\)，并说明为什么平凡正秩丛的 Thom 空间约化上同调中任意两个类的乘积都为零。
\end{exercise}
\begin{solution}
	记 \(j^*:H^*(D(E),S(E);R)\to H^*(D(E);R)\) 为忘记相对条件的映射。零截面与投影互为同伦逆，而 Euler 类定义给出 \(s_0^*j^*u_E=e\)，故
	\[
		j^*u_E=p^*e.
	\]
	相对杯积与绝对—相对杯积相容：两个相对上链都在 \(S(E)\) 上为零，将其中一个视为绝对上链，不改变它们的杯积。因此
	\[
		u_E\smile u_E=j^*u_E\smile u_E=p^*e\smile u_E=\Phi_E(e).
	\]
	对一般齐次类，把次数 \(r\) 的 \(u_E\) 越过 \(p^*b\)，得到
	\[
		(p^*a\smile u_E)\smile(p^*b\smile u_E)
		=(-1)^{r\deg b}p^*(a\smile b\smile e)\smile u_E,
	\]
	即题给公式。非齐次类按齐次分量展开即可。

	平凡正秩丛有常值非零截面，由\cref{g2:prop:euler-basic-properties}得 \(e=0\)，所以上式恒为零。经商比较，这就是 Thom 空间约化上同调的乘法；它不意味着底空间本身的上同调乘法为零。Thom 同构保留底空间类作用在相对类上的模结构，约化环的内部乘法则额外记住 Euler 类。
\end{solution}

\begin{exercise}\label{g2:ex:odd-rank-euler-torsion}
	设 \(E\to B\) 是奇数秩的定向实向量丛。证明 \(2e_{\mathrm{Th}}(E)=0\)。据此说明：若 \(H^{\operatorname{rank}E}(B;\mathbb Z)\) 没有非零的二阶扭元，则 Euler 类为零；但这一点仍不保证存在非零截面。再用这个结论与\cref{g2:cor:tangent-euler-characteristic}证明闭定向奇维流形的 Euler 示性数为零。
\end{exercise}
\begin{solution}
	记 \(E^{\mathrm{op}}\) 为反转纤维定向后的同一个丛。逐纤维映射 \(-I:E\to E^{\mathrm{op}}\) 保持定向：它相对于原定向的行列式符号为 \((-1)^{\operatorname{rank}E}=-1\)，而目标又反转了一次定向。Euler 类对保持定向的丛同构不变，因此
	\[
		e_{\mathrm{Th}}(E)=e_{\mathrm{Th}}(E^{\mathrm{op}})=-e_{\mathrm{Th}}(E).
	\]
	于是 \(2e_{\mathrm{Th}}(E)=0\)。若相应上同调群没有非零二阶扭元，就有 \(e_{\mathrm{Th}}(E)=0\)。\cref{g2:ex:zero-euler-no-section}中的 \(S^4\) 上定向三维丛满足这个条件，却没有非零截面，故消失条件不充分。

	若 \(M\) 是闭定向奇维流形，则整数求值满足 \(2\langle e_{\mathrm{Th}}(\mathrm TM),[M]\rangle=0\)，而整数群没有非零二阶扭元，所以这个求值为零。由\cref{g2:cor:tangent-euler-characteristic}，它恰是 \(\chi(M)\)。
\end{solution}

\begin{exercise}\label{g2:ex:product-tangent-euler}
	设 \(M,N\) 是闭定向光滑流形，维数分别为正整数 \(m,n\)。乘积按 \(M\) 在前、\(N\) 在后定向。用切丛分解和 Euler 类证明
	\[
		e_{\mathrm{Th}}(\mathrm T(M\times N))
		=\pi_M^*e_{\mathrm{Th}}(\mathrm TM)\smile\pi_N^*e_{\mathrm{Th}}(\mathrm TN),
		\qquad\chi(M\times N)=\chi(M)\chi(N).
	\]
	进一步，若 \(X,Y\) 分别是 \(M,N\) 上仅有非退化零点的向量场，计算乘积向量场 \((X,Y)\) 在 \((x,y)\) 处的指标。
\end{exercise}
\begin{solution}
	乘积坐标给出保持所选定向的丛同构
	\[
		\mathrm T(M\times N)\cong\pi_M^*\mathrm TM\oplus\pi_N^*\mathrm TN.
	\]
	由拉回自然性和直和公式，得到所求的第一式。乘积基本类为 \([M]\times[N]\)，而按同一因子次序的外积满足
	\[
		\bigl\langle\pi_M^*a\smile\pi_N^*b,[M]\times[N]\bigr\rangle
		=\langle a,[M]\rangle\langle b,[N]\rangle
	\]
	对 \(a\in H^m(M;\mathbb Z)\)、\(b\in H^n(N;\mathbb Z)\) 成立。这里没有交换两个因子，故没有附加符号。将两个 Euler 类代入，并在三个流形上使用\cref{g2:cor:tangent-euler-characteristic}，便得第二式。

	乘积向量场的零点恰为两个零点集的笛卡儿积。在正向乘积坐标下，其线性化为分块对角矩阵
	\[
		\begin{pmatrix}\mathrm DX_x&0\\0&\mathrm DY_y\end{pmatrix}.
	\]
	因此每个零点非退化，且
	\[
		\operatorname{ind}_{(x,y)}(X,Y)
		=\operatorname{sgn}(\det\mathrm DX_x\det\mathrm DY_y)
		=\operatorname{ind}_x(X)\operatorname{ind}_y(Y).
	\]
	对所有零点求和，得到两个指标和的乘积，也从零点计数直接解释了 Euler 数的乘法公式。
\end{solution}

\begin{exercise}\label{g2:ex:pullback-tangent-sphere-euler}
	设 \(f:S^2\to S^2\) 是次数为 \(d\in\mathbb Z\) 的光滑映射，给 \(E=f^*\mathrm TS^2\) 取拉回定向。计算 \(e_{\mathrm{Th}}(E)\)，判定 \(E\) 何时存在非零截面。若 \(s\) 是与零截面横截的截面，证明它至少有 \(2|d|\) 个零点，并且零点总数为偶数。
\end{exercise}
\begin{solution}
	令 \(\omega\in H^2(S^2;\mathbb Z)\) 在正基本类上取值一。由自然性、映射度的定义和\cref{g2:ex:euler-tangent-spheres-surfaces}，
	\[
		e_{\mathrm{Th}}(E)=f^*(2\omega)=2d\omega.
	\]
	若有非零截面，Euler 类必须为零，故 \(d=0\)。反过来，若 \(d=0\)，Euler 类为零，而底空间维数等于丛的秩；\cref{g2:thm:euler-primary-comparison}与\cref{g2:thm:sphere-primary-euler}说明首障碍消失后即可得到全空间上的单位截面。因此非零截面存在当且仅当 \(d=0\)。这里先得到连续截面；若要求光滑，可将其充分接近地光滑逼近，紧性保证逼近后仍非零。

	设横截零点中指标为 \(+1,-1\) 的个数分别为 \(N_+,N_-\)。\cref{g2:thm:euler-index-evaluation}给出 \(N_+-N_-=2d\)，故
	\[
		N_++N_-\ge|N_+-N_-|=2|d|,
		\qquad N_++N_-=2d+2N_-.
	\]
	第一式给出下界，第二式说明总数为偶数。\(d=0\) 时下界允许没有零点，也允许某个截面有符号相消的零点对。
\end{solution}

\begin{exercise}\label{g2:ex:disk-relative-euler-index}
	给闭圆盘 \(D^2\subset\mathbb C\) 取标准定向，在平凡定向二维丛上指定一个非零边界截面 \(v:S^1\to\mathbb C^*\)，令 \(k=\deg(v/|v|)\)。证明 \(v\) 可以延拓为处处非零的连续截面当且仅当 \(k=0\)。若一个光滑延拓 \(F:D^2\to\mathbb C\) 的零点全部孤立，证明其指标和为 \(k\)。当 \(v(z)=z^k\) 时，对任意整数 \(k\) 写出达到这个指标和的具体光滑延拓，并解释为何这不与平凡丛的绝对 Euler 类为零矛盾。
\end{exercise}
\begin{solution}
	若有非零延拓，将其除以长度便得到 \(D^2\to S^1\) 的延拓，故边界映射零伦，次数为零。反之，\(k=0\) 时 \(v/|v|:S^1\to S^1\) 零伦。用这个零伦沿半径填入圆盘，可得连续映射 \(g:D^2\to S^1\)，其边界值为 \(v/|v|\)。边界上的正函数 \(|v|\) 也有正延拓，例如取
	\[
		h(re^{\mathrm i\theta})=(1-r)+r|v(e^{\mathrm i\theta})|,
		\qquad h(0)=1.
	\]
	于是 \(hg\) 是所求非零延拓。

	设光滑延拓的零点为 \(z_1,\ldots,z_\ell\)，在每个零点周围取小闭圆盘 \(D_i\)，互不相交且不碰外边界。在挖去这些小圆盘内部的区域上，\(F/|F|\) 是到 \(S^1\) 的连续映射。该区域正向边界的同调类为外边界减去各小圆盘的正向边界；边界总类在区域的一维同调中为零。映入 \(H_1(S^1;\mathbb Z)\) 后便得
	\[
		k-\sum_{i=1}^\ell\deg\bigl((F/|F|)|_{\partial D_i}\bigr)=0.
	\]
	每个内边界次数按小圆盘本身的正向计算，正是对应零点的指标，故总和为 \(k\)。

	对 \(v(z)=z^k\)，若 \(k>0\)，取 \(F(z)=z^k\)；若 \(k<0\)，取 \(F(z)=\overline z^{\,-k}\)；若 \(k=0\)，取 \(F(z)=1\)。前两种延拓只有原点一个零点，其小圆周映射次数分别为 \(k\)，后一种没有零点。绝对 Euler 类属于 \(H^2(D^2;\mathbb Z)=0\)，只描述是否能自由选择一个非零截面；本题要求保持指定边界值，障碍属于 \(H^2(D^2,S^1;\mathbb Z)\cong\mathbb Z\)，恰由边界绕数 \(k\) 给出。
\end{solution}

\begin{exercise}\label{g2:ex:gysin-product-normalization}
	设 \(B\) 满足本章 Thom 同构的条件，\(r\ge2\)，并取平凡定向丛 \(E=B\times\mathbb R^r\)。令 \(u\) 是球面 \(S^{r-1}\) 的正向上同调生成元，仍用 \(u\) 表示它在 \(B\times S^{r-1}\) 上的拉回。证明
	\[
		H^k(B\times S^{r-1};R)=\pi^*H^k(B;R)\oplus\bigl(\pi^*H^{k-r+1}(B;R)\smile u\bigr),
	\]
	并计算两部分的沿纤维积分。写出上同调环的乘法约定，解释 Gysin 序列在这个例子中如何分裂。
\end{exercise}
\begin{solution}
	球面只有零次及 \(r-1\) 次上同调，且都是自由秩一 \(R\) 模。第4章的乘积比较因此给出所述直和；也可在球面的一份零胞腔、一份 \(r-1\) 胞腔的链模型上直接得到同一分解。正向生成元满足 \(u^2=0\)，因为它来自球面，而 \(H^{2r-2}(S^{r-1};R)=0\)。若 \(a\in H^\ell(B;R)\)，乘积中还有
	\[
		u\smile\pi^*a=(-1)^{(r-1)\ell}\pi^*a\smile u.
	\]
	这两条关系连同基底的乘法，给出分次张量积环 \(H^*(B;R)\otimes H^*(S^{r-1};R)\)。

	由正合性，\(\pi_!\pi^*a=0\)；由纤维归一化及投影公式，\(\pi_!(\pi^*b\smile u)=b\)。故
	\[
		\pi_!(\pi^*a+\pi^*b\smile u)=b.
	\]
	平凡丛有非零截面，所以 Euler 类为零，Gysin 序列分解成短正合列
	\[
		0\longrightarrow H^k(B;R)\xrightarrow{\pi^*}H^k(B\times S^{r-1};R)\xrightarrow{\pi_!}H^{k-r+1}(B;R)\longrightarrow0.
	\]
	映射 \(b\mapsto\pi^*b\smile u\) 是右逆，给出所述分裂。这个例子也检验了符号：基底类写在 \(u\) 前面时，积分结果就是 \(b\)。
\end{solution}

\begin{exercise}\label{g2:ex:gysin-sphere-product-base}
	在 \(B=S^2\times S^2\) 上，取两因子的正向生成元 \(x,y\)，于是
	\[
		H^*(B;\mathbb Z)=\mathbb Z[x,y]/(x^2,y^2),\qquad\deg x=\deg y=2.
	\]
	设定向实二平面丛 \(E\to B\) 满足 \(e_{\mathrm{Th}}(E)=ax+by\)，其中 \((a,b)\ne(0,0)\)，令 \(d=\gcd(|a|,|b|)>0\)。计算 \(H^*(S(E);\mathbb Z)\) 的各次群，并用 \(X=\pi^*x\)、\(Y=\pi^*y\) 以及适当的三次、五次生成元描述乘法。
\end{exercise}
\begin{solution}
	所有非零的 Euler 乘法由下面两个整数矩阵给出：
	\[
		H^0(B)=\mathbb Z\longrightarrow H^2(B)=\mathbb Z^2,\quad 1\longmapsto(a,b),
	\]
	\[
		H^2(B)=\mathbb Z^2\longrightarrow H^4(B)=\mathbb Z,\quad (s,t)\longmapsto bs+at.
	\]
	第一个映射单射，其余核为 \(\mathbb Z\oplus\mathbb Z/d\)，可由 Bézout 等式把列向量 \((a,b)\) 通过整数可逆行变换化成 \((d,0)\) 得到。第二个映射的像为 \(d\mathbb Z\)，核由 \((a/d,-b/d)\) 生成。基底的奇数次上同调为零，故 Gysin 序列在各次直接给出
	\[
		H^j(S(E);\mathbb Z)=
		\begin{cases}
			\mathbb Z,&j=0,3,5,\\
			\mathbb Z\oplus\mathbb Z/d,&j=2,\\
			\mathbb Z/d,&j=4,\\
			0,&\text{其余次数}.
		\end{cases}
	\]
	这里 \(\mathbb Z/1=0\)。二次群由 \(X,Y\) 生成，唯一加法关系为 \(aX+bY=0\)。四次群由 \(XY=\pi^*(xy)\) 生成，阶恰为 \(d\)，并有 \(X^2=Y^2=0\)。

	积分把三次群同构到上述第二个矩阵的核，把五次群同构到 \(H^4(B)\)。因此可选择唯一的生成元 \(z,w\)，使
	\[
		\pi_!z=\frac adx-\frac bdy,\qquad \pi_!w=xy.
	\]
	投影公式给出
	\[
		\pi_!(Xz)=x\smile\left(\frac adx-\frac bdy\right)=-\frac bdxy,\qquad
		\pi_!(Yz)=\frac adxy.
	\]
	由于五次积分是同构，得到 \(Xz=-(b/d)w\)、\(Yz=(a/d)w\)。其余正次数乘积由这些关系、分次交换律及五次以上群为零决定。这也展示了为什么只列出加法群还不够：同样的次数中，乘法仍保留 Euler 类两个分量的具体数值。
\end{solution}

\begin{exercise}\label{g2:ex:gysin-fiber-degree}
	设 \(E,E'\to B\) 为同秩 \(r\ge2\) 的定向实向量丛，\(B\) 连通且满足 Thom 同构的条件。设 \(F:S(E)\to S(E')\) 覆盖 \(\operatorname{id}_B\)，并在每条正向球面纤维上具有同一个整数映射度 \(m\)。证明
	\[
		e_{\mathrm{Th}}(E')=m\,e_{\mathrm{Th}}(E),\qquad
		\pi_!F^*b=m\,\pi'_!b.
	\]
	特别说明：纤维反向的同构会怎样改变 Euler 类和沿纤维积分。
\end{exercise}
\begin{solution}
	对 \(0<t\le1\)、\(v\in S(E)\)，定义 \(\overline F(tv)=tF(v)\)，并把各零向量送到相应的零向量。这是从 \((D(E),S(E))\) 到 \((D(E'),S(E'))\) 的连续映射，且覆盖基底恒等。零截面附近的连续性来自 \(\|\overline F(tv)\|=t\)；球面方向无论怎样变化，像的长度都趋于零。

	每条盘纤维上的相对映射度等于边界球面上的映射度 \(m\)，这是盘球对长正合列的自然性。故 \(\overline F^*u_{E'}\) 限制到每条纤维都是 \(m\) 倍正向生成元。由 Thom 同构在零次的唯一性，
	\[
		\overline F^*u_{E'}=m\,u_E.
	\]
	沿零截面拉回，得到 \(e_{\mathrm{Th}}(E')=m\,e_{\mathrm{Th}}(E)\)。对任意基底类 \(c\)，又有
	\[
		\overline F^*\Phi_{E'}(c)=p^*c\smile\overline F^*u_{E'}=m\Phi_E(c).
	\]
	空间对连接映射的自然性给出 \(\delta_EF^*b=\overline F^*\delta_{E'}b\)。两丛秩相同，积分定义中的次数符号也相同，因此代入上述等式便有 \(\pi_!F^*b=m\pi'_!b\)。

	纤维反向同构对应 \(m=-1\)，所以两边按所选方向计算的 Euler 类互为负数，积分也相差负号。若只反转同一丛的定向而保持底层映射为恒等，这正是 Thom 类、Euler 类和纤维积分同时反号的规则。
\end{solution}

\begin{exercise}\label{g2:ex:angular-form-thom}
	设 \(E\to B\) 为光滑定向实二平面丛，\(\pi:S(E)\to B\) 为其单位圆丛。给定一形式 \(\theta\in\Omega^1(S(E))\) 及闭二形式 \(\omega\in\Omega^2(B)\)，满足
	\[
		\int_{S(E)_x}\theta=1,\qquad \mathrm d\theta=\pi^*\omega.
	\]
	在 \(D(E)\setminus s_0(B)\) 上记径向投影为 \(\rho(v)=v/\|v\|\)。取光滑径向函数 \(\chi\)，使其在零截面附近为零，在球丛附近为一。证明
	\[
		U=\mathrm d(\chi\rho^*\theta)-p^*\omega
	\]
	给出实 Thom 类的相对形式代表，并求 \(e_{\mathrm{Th}}(E)\) 的实系数像。最后证明：若这个实 Euler 类非零，则球丛上不存在积分为一的闭一形式。
\end{exercise}
\begin{solution}
	因为 \(\chi\) 在零截面邻域恒为零，\(\chi\rho^*\theta\) 可通过补零光滑延拓到整个盘丛。闭性 \(\mathrm dU=0\) 来自 \(\mathrm d^2=0\) 及 \(\mathrm d\omega=0\)。在球丛上，\(\chi=1\)，所以
	\[
		i^*U=\mathrm d\theta-\pi^*\omega=0.
	\]
	因此 \(U\) 是相对闭形式。限制到单条盘纤维时，\(p^*\omega\) 为零；由外法向优先的 Stokes 公式，
	\[
		\int_{D(E)_x}U=\int_{S(E)_x}\theta=1.
	\]
	每条纤维上都是正向相对生成元，故 Thom 类的唯一性给出 \([U]=u_E\) 的实系数像。

	在零截面附近，\(U=-p^*\omega\)，于是
	\[
		e_{\mathrm{Th}}(E)_{\mathbb R}=[s_0^*U]=-[\omega].
	\]
	这个负号由相对代表在边界必须为零决定，与本节 Gysin 的定向约定一致。若另有闭一形式 \(\eta\) 满足每条纤维积分为一，则实系数 Gysin 序列给出 \(\pi_![\eta]=1\)，从而 \(e_{\mathrm{Th}}(E)_{\mathbb R}=0\)。这与假设矛盾，故所述闭形式不存在。
\end{solution}

% !TeX root = ../../Book2.tex
% 纲要标题：## 第21章　Stiefel–Whitney、陈与 Pontryagin 类
% 纲要位置：几何学计划纲要.md，第 1802 行。
% 正文状态：初稿；尚未正式审核。

% 纲要内容开始
% 
% > 本章目标与证明范围：分别完成整数类与模2类的构造，再证明分裂原理和计算法。拓扑扭信息在此保留，不能用下一章的曲率形式取代。
% 
% 纲要内容结束

\chapter{Stiefel–Whitney、陈与 Pontryagin 类}
\label{g2:ch:21}
\BookChapterNumber{g2:ch:21}

\begin{chapterabstract}
	一张向量丛在不同维数上的扭转，可以由底空间上不同次数的上同调类记录。本章从实、复线丛出发，用射影丛的上同调关系构造 Stiefel--Whitney 类与陈类，再由复化定义 Pontryagin 类。Leray--Hirsch 定理和旗丛提供分裂原理，使线丛计算能够回到原来的向量丛。随后证明 Whitney 公式、顶类与 Euler 类的比较、Thom 平方解释和 Wu 公式，计算射影空间切丛与结构群提升障碍。最后在复 Grassmann 流形上把示性类与 Schubert 条件联系起来，完成四维复空间中二维子空间的上同调环及四条直线的相交计数。
\end{chapterabstract}

\begin{learninggoals}
	\begin{itemize}
		\item 将实线丛的方向输运、复线丛的 Euler 类和过渡函数联系起来，固定第一示性类的符号。
		\item 证明 Leray--Hirsch 定理、射影丛公式和分裂原理，说明形式根计算为何能决定原底空间上的等式。
		\item 构造 Stiefel--Whitney 类、陈类与 Pontryagin 类，掌握自然性、直和及复化公式，并辨认整数二阶扭信息。
		\item 完成 Thom 平方与 Wu 类的比较，用切丛计算检验向量场、几乎复结构和 spin 提升的条件。
		\item 在 Grassmann 流形上解释 Schubert 类及特殊类，计算 \(\operatorname{Gr}_2(\mathbb C^4)\) 的乘法并求解四条直线的相交问题。
	\end{itemize}
\end{learninggoals}

上一章的 Euler 类记录一张丛的截面能否避开零点。它的次数等于丛的实秩，因此只提供一个特定次数的信息：线丛的方向会沿闭路改变，较高秩丛又可能在二维或四维子空间上发生扭转，这些现象需要不同次数的类。示性类把一张丛组织成一组可拉回的上同调数据；同一个定义可以随映射限制到曲线、曲面或更高维的检验空间。

\ProseParagraphBreak

先考虑复秩二丛 \(E=L_1\oplus L_2\)。如果两条线的第一类分别为 \(x_1,x_2\)，它们的和与乘积就是最自然的对称数据：交换两条线不改变它们。但是一般丛并没有给定的线丛分解。我们会将 \(E\) 拉回到旗丛，在那里得到实际分解，并证明上同调拉回单射。于是在线丛上算出的等式能够返回原基空间。读者需要区分这两件事：分解发生在新的基空间，而等式由单射保证仍能检测原来的类。

\ProseParagraphBreak

为了避免先假定乘法公式再用它定义示性类，本章从射影丛的上同调关系读出各阶系数，再证明这些系数具有所需性质。实线丛使用模二一次类，复线丛使用整数二次类，两种构造具有相同的计算形状，系数却保留不同的信息。Pontryagin 类由实丛的复化得到；它的整数直和公式中还可能出现被二消去的修正项。我们将写出这个修正，并在实射影空间上给出非零例子。

\ProseParagraphBreak

本章使用第19章的丛分类、第20章的 Thom 与 Euler 构造，以及第13章的平方运算；所需条件会在具体证明中回引。分裂、平方与三类示性类的组织可参阅\cite[第23章第6--7节]{May1999}。最后一节的 Grassmann 计算把此前的上同调环、对偶和横截相交放在同一个空间里使用，也为后续代数几何中的交叉计算准备一个完整例子。

% !TeX root = ../../Book2.tex
% 纲要标题：### 21.1 线丛与通用生成元
% 纲要位置：几何学计划纲要.md，第 1806 行。
% 正文状态：未撰写；本文件目前是框架。

% 纲要内容开始
% 
% * 实线丛和第一 Stiefel–Whitney 类
% * 复线丛与第一 Chern 类
% * 过渡函数、分类空间和指数序列的比较
% 
% 纲要内容结束

\section{线丛与通用生成元}
\label{g2:sec:2101}

待撰写

\subsection{实线丛和第一 Stiefel–Whitney 类}
\label{g2:subsec:2101:01}
待撰写

\subsection{复线丛与第一 Chern 类}
\label{g2:subsec:2101:02}
待撰写

\subsection{过渡函数、分类空间和指数序列的比较}
\label{g2:subsec:2101:03}
待撰写
% 纲要条目排为编号小节标题，并进入目录与书签；说明段仍保留于纲要注释中。

% !TeX root = ../../Book2.tex
% 纲要标题：### 21.2 射影丛与分裂原理
% 纲要位置：几何学计划纲要.md，第 1812 行。
% 正文状态：未撰写；本文件目前是框架。

% 纲要内容开始
% 
% * Leray–Hirsch 的证明
% * 射影丛上同调
% * 旗丛拉回的单射性
% * 以形式根进行对称多项式计算
% 
% 纲要内容结束

\section{射影丛与分裂原理}
\label{g2:sec:2102}

待撰写

\subsection{Leray–Hirsch 的证明}
\label{g2:subsec:2102:01}
待撰写

\subsection{射影丛上同调}
\label{g2:subsec:2102:02}
待撰写

\subsection{旗丛拉回的单射性}
\label{g2:subsec:2102:03}
待撰写

\subsection{以形式根进行对称多项式计算}
\label{g2:subsec:2102:04}
待撰写
% 纲要条目排为编号小节标题，并进入目录与书签；说明段仍保留于纲要注释中。

% !TeX root = ../../Book2.tex
% 纲要标题：### 21.3 一般示性类的构造
% 纲要位置：几何学计划纲要.md，第 1819 行。
% 正文状态：初稿；尚未正式审核。

% 纲要内容开始
% 
% * 陈类的存在唯一性
% * Stiefel–Whitney 类及实射影丛关系
% * Whitney 和公式与 Wu 公式
% * 对偶、张量和复化
% * Pontryagin 类及其性质
% 
% 纲要内容结束

\section{一般示性类的构造}
\label{g2:sec:2103}

射影丛公式把一张秩 \(r\) 丛的粘合信息写入一个首一多项式关系，旗丛又把这个多项式变成线性因子的乘积。现在用关系的系数定义一般示性类，再利用旗丛拉回的单射性，把线丛上的计算送回原基空间。这一顺序同时解决存在性、选择无关性与计算问题。

以下主要以 CW 复形为基，有限秩丛均带纤维内积。证明中的射影丛及旗丛使用上一节的弱 CW 逼近比较。Thom 类也可沿同一比较运输：盘丛、球丛均为局部平凡丛，拉回到 CW 弱逼近后，总空间映射由纤维化同伦长正合列为弱等价；绝对上同调比较及对的长正合列给出相对上同调同构。因此第20章的 Thom、Euler 及其自然性可用于这些实际拉回丛。

\subsection{陈类的存在唯一性}
\label{g2:subsec:2103:01}
设 \(E\to B\) 为复秩 \(r\ge1\) 丛。在 \(p:\mathbb P(E)\to B\) 上记典范线丛为 \(\lambda\)，置 \(h=c_1(\lambda^*)\)。\cref{g2:thm:projective-bundle}说明 \(1,h,\ldots,h^{r-1}\) 是 \(H^*(B;\mathbb Z)\) 模基。因此 \(h^r\) 可以唯一写成这些幂的线性组合；其各项的总次数必须同为 \(2r\)。

\begin{definition}\label{g2:def:chern-classes-projective}
	复秩 \(r\) 丛 \(E\to B\) 的\term[chen-lei]{陈类}[Chern classes]为唯一满足
	\begin{equation}\label{g2:eq:chern-projective-relation}
		h^r+p^*c_1(E)h^{r-1}+p^*c_2(E)h^{r-2}+\cdots+p^*c_r(E)=0
	\end{equation}
	的类 \(c_i(E)\in H^{2i}(B;\mathbb Z)\)。置 \(c_0(E)=1\)，\(i>r\) 时置 \(c_i(E)=0\)，并称 \(c(E)=1+c_1(E)+\cdots+c_r(E)\) 为总陈类。
\end{definition}

当 \(r=1\) 时，\(\mathbb P(E)=B\)、\(\lambda=E\)，关系变成 \(-c_1(E)+c_1(E)=0\)，故新定义延续上一节的线丛规范。正号排列来自 \(h=c_1(\lambda^*)\)：如果改用 \(c_1(\lambda)=-h\)，关系的系数就要交错变号。

\begin{theorem}\label{g2:thm:characteristic-existence-uniqueness}
	CW 复形上的复有限秩向量丛具有上述自然陈类。若 \(q:Y\to B\) 是使 \(q^*E=L_1\oplus\cdots\oplus L_r\) 的旗丛，记 \(x_j=c_1(L_j)\)，则
	\begin{equation}\label{g2:eq:chern-root-formula}
		q^*c_i(E)=\sum_{j_1<\cdots<j_i}x_{j_1}\cdots x_{j_i},\qquad
		q^*c(E)=\prod_{j=1}^r(1+x_j).
	\end{equation}
	自然性及这个分裂公式唯一确定全部 \(c_i\)。
\end{theorem}
\begin{proof}
	射影丛模基已经保证存在唯一的关系系数。若 \(f:A\to B\)，则拉回方块中的 \(\mathbb P(f^*E)\to\mathbb P(E)\) 将 \(\lambda\) 拉回为新的典范线丛，故也将 \(h\) 拉回为新的 \(h\)。拉回式\eqref{g2:eq:chern-projective-relation}并用新射影丛的模基唯一性，得到 \(c_i(f^*E)=f^*c_i(E)\)。

	为确定关系的系数，先设 \(E=L_1\oplus\cdots\oplus L_r\) 已分裂。在 \(\mathbb P(E)\) 上，丛
	\[
		\operatorname{Hom}_{\mathbb C}(\lambda,p^*E)
		=\bigoplus_{j=1}^r(\lambda^*\otimes p^*L_j)
	\]
	有一个处处非零截面：每点的直线包含映射 \(\lambda_\ell\hookrightarrow E_{p(\ell)}\)。\cref{g2:prop:euler-basic-properties}说明它的底层定向实丛 Euler 类为零；同一命题的直和公式及线丛的 \(c_1=e_{\mathrm{Th}}\) 给出
	\[
		0=\prod_{j=1}^r c_1(\lambda^*\otimes p^*L_j)
		=\prod_{j=1}^r(h+p^*x_j).
	\]
	这里使用的是已证明的 Euler 类乘法，并未预先假定一般陈类的乘法。展开右端，与唯一的首一射影关系比较，得到 \(c_i(E)\) 是 \(x_j\) 的第 \(i\) 个初等对称式。

	一般情形拉回到旗丛，用自然性及刚才的分裂计算得到式\eqref{g2:eq:chern-root-formula}。由\cref{g2:thm:splitting-principle}，\(q^*\) 单射，所以这些等式唯一确定原基空间上的类；任何其他自然赋类若满足同样的分裂公式，拉回后相同，原来也必相同。
\end{proof}

例如秩二时，形式根 \(x_1,x_2\) 给出 \(c_1=x_1+x_2\)、\(c_2=x_1x_2\)，射影丛关系为 \(h^2+c_1h+c_2=0\)。根是旗空间上的实际一次线丛类，未必分别下降到底空间；能下降的这些对称式已经由射影关系独立构造。这是“形式根”计算可以可靠使用的依据。

\subsection{Stiefel–Whitney 类及实射影丛关系}
\label{g2:subsec:2103:02}
对实秩 \(r\) 丛，射影纤维是 \(\mathbb RP^{r-1}\)，用模二系数后它同样有一列由一次类的幂组成的基。这里不需要给原丛或射影纤维选整数方向。

\begin{definition}\label{g2:def:stiefel-whitney-projective}
	设 \(E\to B\) 为实秩 \(r\ge1\) 向量丛，\(p:\mathbb P(E)\to B\) 为实射影丛，\(\lambda\) 为典范实线丛，\(t=w_1(\lambda)\)。唯一满足
	\begin{equation}\label{g2:eq:sw-projective-relation}
		t^r+p^*w_1(E)t^{r-1}+p^*w_2(E)t^{r-2}+\cdots+p^*w_r(E)=0
	\end{equation}
	的类 \(w_i(E)\in H^i(B;\mathbb F_2)\)，称为 \(E\) 的\term[Stiefel-Whitney-lei]{Stiefel--Whitney 类}[Stiefel--Whitney classes]。置 \(w_0=1\)、\(w_i=0\)（\(i>r\)），并记 \(w(E)=\sum_iw_i(E)\)。
\end{definition}

\begin{theorem}\label{g2:thm:stiefel-whitney-construction}
	上述类对 CW 基空间上的实丛自然。若旗丛拉回使 \(q^*E=L_1\oplus\cdots\oplus L_r\)，并记 \(a_j=w_1(L_j)\)，则
	\begin{equation}\label{g2:eq:sw-root-formula}
		q^*w(E)=\prod_{j=1}^r(1+a_j).
	\end{equation}
	自然性及该分裂公式唯一确定这些类；秩一时与\cref{g2:def:first-stiefel-whitney-line}一致。
\end{theorem}
\begin{proof}
	射影丛的模二模基保证关系系数存在唯一，射影拉回方块保证自然性。若 \(E\) 已分裂，\(\operatorname{Hom}_{\mathbb R}(\lambda,p^*E)\) 仍有处处非零的包含截面。它分裂成 \(\lambda^*\otimes p^*L_j\) 的直和。由式\eqref{g2:eq:line-mod-two-euler}及模二 Euler 直和公式，得到
	\[
		0=e_2(\operatorname{Hom}(\lambda,p^*E))
		=\prod_{j=1}^r(t+p^*a_j).
	\]
	与式\eqref{g2:eq:sw-projective-relation}比较唯一系数，即得分裂公式。一般丛再拉回实旗丛；\(q^*\) 在模二上同调上单射，所以相同论证给出唯一性。秩一的关系为 \(t+w_1(E)=0\)，正好恢复线丛类。
\end{proof}

\begin{proposition}\label{g2:prop:first-classes-determinants}
	对 CW 基空间上的实丛 \(E\) 与复丛 \(F\)，有
	\[
		w_1(E)=w_1(\det E),\qquad c_1(F)=c_1(\det F).
	\]
	实丛可定向当且仅当 \(w_1(E)=0\)。
\end{proposition}
\begin{proof}
	在相应旗丛上，行列式线丛是各线因子的张量积，因此其第一类是各根之和；这也是两个总类的一次项。拉回单射使等式在底空间成立。实丛方向的输运为过渡矩阵的行列式符号，恰好是 \(\det E\) 的方向输运，所以可定向性由\cref{g2:prop:real-line-first-class}判定。
\end{proof}

\begin{theorem}\label{g2:thm:top-class-euler}
	若 \(E\to B\) 为 CW 基空间上的实秩 \(r\) 丛，则 \(w_r(E)=e_2(E)\)。若 \(E\) 整数可定向，则 \(w_r(E)=\rho_2 e_{\mathrm{Th}}(E)\)。若 \(F\to B\) 为复秩 \(r\) 丛，则
	\[
		c_r(F)=e_{\mathrm{Th}}(F_{\mathbb R}),
	\]
	右侧采用复自然定向。
\end{theorem}
\begin{proof}
	先在旗丛上分裂成线丛。实情形中，模二 Euler 类是各线的 \(w_1\) 之积，恰为式\eqref{g2:eq:sw-root-formula}的最高项。复情形中，实丛的直和顺序按各复线的正向实二平面排列，正是复自然定向；Euler 直和公式使最高类为 \(x_1\cdots x_r\)，恰为 \(c_r\)。再由拉回单射得到原空间等式。整数定向丛的 Thom 类约化模二后是规范模二 Thom 类，零截面拉回与系数约化交换，得到中间断言。
\end{proof}

\subsection{Whitney 和公式与 Wu 公式}
\label{g2:subsec:2103:03}
两张丛直和后，在线丛分裂模型中只是把两组根放在一起。因此直和公式来自一个简单的乘法，但要让它在原基空间成立，需要同时分裂两张丛并保留拉回单射。

\begin{theorem}[Whitney 和公式]\label{g2:thm:whitney-sum}
	设 \(B\) 为 CW 复形。对其上的复有限秩丛 \(E,F\)，以及实有限秩丛 \(V,W\)，分别有
	\begin{equation}\label{g2:eq:whitney-total-formulas}
		c(E\oplus F)=c(E)c(F),\qquad w(V\oplus W)=w(V)w(W).
	\end{equation}
	两组类都在添上平凡直和因子时不变。自然性、线丛归一化、秩以上的消失及式\eqref{g2:eq:whitney-total-formulas}也唯一刻画各组类。
\end{theorem}
\begin{proof}
	先取 \(E\) 的旗丛，再取拉回 \(F\) 的旗丛。复合拉回仍单射，且两丛同时分裂，设根分别为 \(x_i,y_j\)。直和的根就是它们的合并，所以
	\[
		q^*c(E\oplus F)=\prod_i(1+x_i)\prod_j(1+y_j)
		=q^*(c(E)c(F)).
	\]
	单射性给出复情形。实旗丛及一次模二根给出完全相同的乘法论证。平凡丛分裂成第一类均为零的线丛，其总类为 \(1\)，从而得到稳定性。任何另一组满足所列性质的类，在旗丛上由线丛归一化与直和公式确定为同一个根乘积，故由单射性与本构造一致。
\end{proof}

短正合列 \(0\to E\to G\to F\to0\) 也满足这个公式，因为纤维内积给出 \(G\cong E\oplus E^\perp\)，而 \(E^\perp\to F\) 是丛同构。这里需要的是拓扑分裂，若原丛带有全纯等额外结构，不要求分裂保持那些结构。

\ProseParagraphBreak

现在完成第13章留下的 Thom 平方比较。记实秩 \(r\) 丛的模二 Thom 类为 \(U\)，用第20章同构定义 \(w_i^{\mathrm{Sq}}\) 使
\[
	\operatorname{Sq}^iU=p^*w_i^{\mathrm{Sq}}(E)\smile U.
\]
\cref{g2:prop:thom-square-properties}已经用平方的自然性与 Cartan 公式证明这组类的自然性、直和公式和秩以上消失。第20章提供的 Thom 同构及纤维积直和公式使同一论证适用于本节的基空间，余下只需核对线丛归一化。

\begin{theorem}\label{g2:thm:thom-square-comparison}
	对 CW 复形上的实有限秩丛 \(E\)，射影丛构造的 Stiefel--Whitney 类满足
	\begin{equation}\label{g2:eq:thom-sw-identity-proved}
		\operatorname{Sq}(U)=p^*w(E)\smile U.
	\end{equation}
	因此 \(w_i^{\mathrm{Sq}}(E)=w_i(E)\)。
\end{theorem}
\begin{proof}
	设 \(L\) 为实线丛，\(U\) 是它的一次模二 Thom 类。平方的最高项性质给出 \(\operatorname{Sq}^1U=U\smile U\)。记 \(j^*\) 为忘去相对条件，则零截面与盘丛投影互为同伦逆，故
	\[
		j^*U=p^*e_2(L).
	\]
	相对杯积与绝对—相对杯积相容，所以
	\[
		U\smile U=j^*U\smile U=p^*e_2(L)\smile U.
	\]
	Thom 同构的单射性及式\eqref{g2:eq:line-mod-two-euler}说明 \(w_1^{\mathrm{Sq}}(L)=e_2(L)=w_1(L)\)。其余正次项由秩一消失。于是这组平方类满足刚才唯一性定理的全部条件，必与射影丛构造相同。
\end{proof}

这一步把作用于总空间相对类的平方运算，识别为基空间上的示性类。对流形的对角嵌入应用它，便能把切丛类与第13章已经定义的 Wu 类联系起来。

\begin{theorem}[Wu 公式]\label{g2:thm:wu-formula}
	设 \(M\) 为连通闭光滑 \(n\) 维流形，\(v(M)\) 为\cref{g2:def:wu-class}按模二对偶定义的总 Wu 类。则
	\[
		w(\mathrm TM)=\operatorname{Sq}(v(M)).
	\]
	这个等式不要求 \(M\) 整数可定向。
\end{theorem}
\begin{proof}
	全部系数取 \(\mathbb F_2\)。先说明证明的目标：把对角线的 Thom 类平方化，一边读出法丛的 \(w\)，另一边用对偶基读出 Wu 类的平方。记 \(H=H^*(M;\mathbb F_2)\)，并以 \(\epsilon(a)\) 表示 \(a\) 的 \(n\) 次分量在 \([M]_2\) 上的求值。闭流形有有限 CW 同伦型，故 \(H\) 有限维；Poincaré 对偶给出齐次基 \(b_\alpha\) 及其对偶基 \(b^\alpha\)，满足
	\[
		\epsilon(b_\alpha b^\beta)=\delta_\alpha^\beta.
	\]
	对角嵌入 \(\Delta:M\hookrightarrow M\times M\) 的法丛同构于 \(\mathrm TM\)：商线性映射 \((v,w)\mapsto w-v\) 的核就是对角切空间。由管状邻域的 Thom 类定义推前 \(\Delta_!a\)，即先在邻域内取 \(a\) 与法丛 Thom 类的乘积，再经切除和忘去相对条件得到 \(H^{n+\deg a}(M\times M;\mathbb F_2)\) 中的类。第20章法向局部对偶的同一计算给出
	\[
		\langle\Delta_!a\smile z,[M\times M]_2\rangle
		=\langle a\smile\Delta^*z,[M]_2\rangle.
	\]
	令 \(D=\Delta_!1\)。由域系数 Künneth 同构和上式的完美配对，
	\begin{equation}\label{g2:eq:diagonal-dual-basis}
		D=\sum_\alpha b_\alpha\times b^\alpha.
	\end{equation}
	例如将右边与 \(x\times y\) 相乘后求值，得到 \(\sum_\alpha\epsilon(b_\alpha x)\epsilon(b^\alpha y)=\epsilon(xy)\)，恰与对角拉回相同，故这确是 \(D\)。

	自然性使平方与相对类到绝对类的映射交换；\cref{g2:thm:thom-square-comparison}于是给出
	\[
		\operatorname{Sq}(D)=\Delta_!w(\mathrm TM)
		=D\smile\operatorname{pr}_1^*w(\mathrm TM).
	\]
	最后一个等式是邻域内的杯积公式：\(\operatorname{pr}_1\) 限制在管状邻域后，与收缩到对角线再投影同伦，所以其拉回正是 Thom 同构中的基空间因子。

	在 \(H\otimes H\) 上施加线性映射 \(1\otimes\epsilon\)。对右边，式\eqref{g2:eq:diagonal-dual-basis}给出 \((1\otimes\epsilon)D=1\)，所以结果是 \(w(\mathrm TM)\)。对左边，Cartan 公式和外积相容性给出
	\[
		(1\otimes\epsilon)\operatorname{Sq}(D)
		=\sum_\alpha\operatorname{Sq}(b_\alpha)\,\epsilon(\operatorname{Sq}(b^\alpha)).
	\]
	Wu 类的定义逐次数合起来正是 \(\epsilon(\operatorname{Sq}(x))=\epsilon(v(M)x)\)。因此上式等于
	\[
		\operatorname{Sq}\left(\sum_\alpha b_\alpha\,\epsilon(v(M)b^\alpha)\right)
		=\operatorname{Sq}(v(M)),
	\]
	最后一步是对偶基展开。两边比较即得所求公式，全程只使用模二基本类。
\end{proof}

这就证明了\cref{g2:thm:wu-tangent-preview}的预告。低次时，\(w_1=v_1\)、\(w_2=v_2+\operatorname{Sq}^1v_1\)；后者额外的一项说明 Wu 类本身与切丛类一般不同。第13章中 \(\mathbb RP^2\) 与 \(\mathbb CP^2\) 的计算现在分别给出实际切丛的总类，而不再只是运算产生的候选表达。

\subsection{对偶、张量和复化}
\label{g2:subsec:2103:04}
对偶、张量与复化改变线丛的根，因此可以用同一个分裂过程计算。进行比较时须保留系数：复陈类位于整数上同调，实 Stiefel--Whitney 类位于模二上同调，二者之间使用约化映射 \(\rho_2\)。

\begin{proposition}\label{g2:prop:chern-operations}
	设 \(E,F\) 为 CW 基空间上的复丛，秩分别为 \(r,s\)，形式陈根为 \(x_i,y_j\)。则
	\[
		c_k(E^*)=c_k(\overline E)=(-1)^k c_k(E),\qquad
		c(E\otimes F)=\prod_{i,j}(1+x_i+y_j),
	\]
	第二式在同时分裂空间中计算，再由单射性理解于底空间。特别地，
	\[
		c_1(E\otimes F)=s\,c_1(E)+r\,c_1(F).
	\]
	若 \(L\) 为复线丛、\(\ell=c_1(L)\)，则
	\begin{equation}\label{g2:eq:chern-line-twist}
		c_k(E\otimes L)=\sum_{j=0}^k\binom{r-j}{k-j}c_j(E)\ell^{k-j}.
	\end{equation}
\end{proposition}
\begin{proof}
	在同时分裂空间，\(E^*\)、\(\overline E\) 的线因子第一类为 \(-x_i\)，因为线丛对偶与共轭都变号；第 \(k\) 个初等对称式遂乘以 \((-1)^k\)。张量积实际分裂为全部 \(L_i\otimes K_j\)，每个因子的第一类为 \(x_i+y_j\)，给出乘积公式及其一次项。

	对线丛扭转，把第 \(k\) 个初等对称式写为 \(e_k(x_1+\ell,\ldots,x_r+\ell)\)。从一项的 \(k\) 个因子中选 \(j\) 个贡献 \(x\)，其余贡献 \(\ell\)；固定那 \(j\) 个根后，还可从其余 \(r-j\) 个位置选出 \(k-j\) 个，故系数为 \(\binom{r-j}{k-j}\)。这证明式\eqref{g2:eq:chern-line-twist}。各等式都是对称式，且拉回单射，因此在底空间成立。
\end{proof}

实丛的对偶由内积与自身同构，所以 \(w(E^*)=w(E)\)。若两实丛分别分裂为根 \(a_i,b_j\) 的实线丛，则其张量积满足 \(w(E\otimes F)=\prod_{i,j}(1+a_i+b_j)\)。这些公式的乘积次数不同，不能把总陈类的整数系数式直接改名为总 Stiefel--Whitney 类的式子。

\begin{proposition}\label{g2:prop:real-complex-characteristic}
	对 CW 基空间上的复秩 \(r\) 丛 \(F\)，有
	\begin{equation}\label{g2:eq:complex-real-sw}
		w_{2i}(F_{\mathbb R})=\rho_2c_i(F),\qquad w_{2i+1}(F_{\mathbb R})=0.
	\end{equation}
	对实丛 \(E\) 的复化 \(E_{\mathbb C}=E\otimes_{\mathbb R}\mathbb C\)，有
	\begin{equation}\label{g2:eq:complexification-chern-mod-two}
		\rho_2c_i(E_{\mathbb C})=w_i(E)^2,\qquad 2c_{2j+1}(E_{\mathbb C})=0.
	\end{equation}
\end{proposition}
\begin{proof}
	在复旗丛上分裂 \(F\)。该拉回不仅在整数上同调单射，在模二上同调也单射，因为复射影纤维的幂基约化后仍是一组基。每条复线的底层实二平面定向，故 \(w_1=0\)，而\cref{g2:thm:top-class-euler}给出 \(w_2=\rho_2c_1\)。对各线因子应用 Whitney 公式，得 \(w(F_{\mathbb R})=\prod_i(1+\rho_2x_i)\)，比较次数即得第一组公式。

	映射 \(v+\mathrm i w\mapsto(v,w)\) 给出实丛同构 \((E_{\mathbb C})_{\mathbb R}\cong E\oplus E\)，于是
	\[
		\sum_i\rho_2c_i(E_{\mathbb C})=w(E)^2=\sum_iw_i(E)^2
	\]
	在模二系数下成立；最后一步的混合项成对相消。比较次数得到第二组的第一个式子。复化丛有规范的复线性同构 \(E_{\mathbb C}\cong\overline{E_{\mathbb C}}\)，由共轭的陈公式得 \(c_i=(-1)^ic_i\)，故奇下标项被 \(2\) 消去。
\end{proof}

例如实线丛复化后虽然只有一个陈类，这个类也未必为零：式\eqref{g2:eq:line-complexification-bockstein}精确把它写成方向类的整数 Bockstein。其模二约化为 \(\operatorname{Sq}^1w_1=w_1^2\)，与式\eqref{g2:eq:complexification-chern-mod-two}一致。复化保留了某些实方向信息，但把它从一次模二群移入二次整数扭群。

\subsection{Pontryagin 类及其性质}
\label{g2:subsec:2103:05}
实丛复化以后得到整数陈类。奇下标类可能只剩二阶扭信息，偶下标类则给出位于 \(4i\) 次的 Pontryagin 类。定义中的交错符号使定向实二平面的第一 Pontryagin 类成为 Euler 类的平方。

\begin{definition}\label{g2:def:pontryagin-classes}
	实有限秩丛 \(E\to B\) 的\term[Pontryagin-lei]{Pontryagin 类}[Pontryagin classes]定义为
	\begin{equation}\label{g2:eq:pontryagin-definition}
		p_i(E)=(-1)^i c_{2i}(E_{\mathbb C})\in H^{4i}(B;\mathbb Z),\qquad p(E)=1+p_1(E)+p_2(E)+\cdots.
	\end{equation}
	若 \(2i>\operatorname{rank}_{\mathbb R}E\)，则该项为零。
\end{definition}

Pontryagin\footnote{\OriginalName{Лев Семёнович Понтрягин}（Lev Semenovich Pontryagin，1908--1988），苏联数学家，研究代数拓扑、拓扑群与最优控制理论。}类对实丛定义，不要求它可定向。

\begin{proposition}\label{g2:prop:pontryagin-properties}
	对 CW 基空间上的实丛，Pontryagin 类自然且在加平凡直和因子后不变，并满足
	\[
		\rho_2p_i(E)=w_{2i}(E)^2.
	\]
	若 \(E\) 是定向实秩 \(2r\) 丛，则 \(p_r(E)=e_{\mathrm{Th}}(E)^2\)。若 \(F\) 为复丛，形式陈根为 \(x_j\)，则
	\begin{equation}\label{g2:eq:pontryagin-complex-underlying}
		p(F_{\mathbb R})=\prod_j(1+x_j^2),\qquad
		p_1(F_{\mathbb R})=c_1(F)^2-2c_2(F).
	\end{equation}
\end{proposition}
\begin{proof}
	自然性、稳定性来自复化与陈类的相应性质。模二约化时符号消失，再用式\eqref{g2:eq:complexification-chern-mod-two}即可。

	对定向秩 \(2r\) 丛 \(E\)，实同构 \((E_{\mathbb C})_{\mathbb R}\cong E\oplus E\) 所比较的两种方向须明确。若 \(e_1,\ldots,e_{2r}\) 为正向实基，复自然实方向是交错排列 \((e_1,\mathrm i e_1,\ldots,e_{2r},\mathrm i e_{2r})\)，直和方向则先放全部 \(e_j\)，再放全部 \(\mathrm i e_j\)。它们相差 \((-1)^{(2r)(2r-1)/2}=(-1)^r\)。最高陈类等于复方向 Euler 类，因而
	\[
		c_{2r}(E_{\mathbb C})=(-1)^re_{\mathrm{Th}}(E\oplus E)=(-1)^re_{\mathrm{Th}}(E)^2.
	\]
	代入定义便得 \(p_r=e_{\mathrm{Th}}^2\)。

	最后，复线性同构 \((F_{\mathbb R})_{\mathbb C}\cong F\oplus\overline F\) 可在每根纤维由复结构的 \(\pm\mathrm i\) 特征子空间给出。因此其陈总类在分裂空间为
	\[
		\prod_j(1+x_j)(1-x_j)=\prod_j(1-x_j^2).
	\]
	取偶下标项再乘 \((-1)^i\)，得到式\eqref{g2:eq:pontryagin-complex-underlying}。一次项利用 \(\sum_jx_j^2=(\sum_jx_j)^2-2\sum_{i<j}x_ix_j\) 化成所列公式。
\end{proof}

整数 Pontryagin 类的直和关系还要保留奇陈类的乘积。直接从陈类的 Whitney 公式分开偶、奇下标，可得
\begin{equation}\label{g2:eq:pontryagin-sum-torsion}
	\begin{split}
		p_k(E\oplus F)={}&\sum_{i+j=k}p_i(E)p_j(F)\\
		&+(-1)^k\!\sum_{\substack{a+b=2k\\a,b\text{ 为奇数}}}c_a(E_{\mathbb C})c_b(F_{\mathbb C}).
	\end{split}
\end{equation}
第二行被 \(2\) 消去，但未必等于零。因此
\[
	p(E\oplus F)=p(E)p(F)\quad\text{于 }H^*(B;\mathbb Z[1/2])
\]
总是成立；在整数上同调中，若相关群没有二阶扭，或者相关奇陈乘积全部消失，也可使用同一公式。特别是两个丛都来自复丛的底层实丛时，上面根的成对表达使这些奇项确实为零。

\begin{example}\label{g2:exm:pontryagin-whitney-torsion}
	令 \(L\to\mathbb RP^4\) 为典范实线丛，\(a=w_1(L)\)，\(t=c_1(L_{\mathbb C})=\widetilde\beta_2a\)。由第13章 \(\rho_2\widetilde\beta_2=\operatorname{Sq}^1\) 及最高平方性质，\(\rho_2t=a^2\)，所以 \(\rho_2(t^2)=a^4\ne0\)，从而 \(t^2\ne0\)。每条实线丛的 Pontryagin 类在正次全为零，但
	\[
		p_1(L\oplus L)=-c_2(L_{\mathbb C}\oplus L_{\mathbb C})=-t^2\ne0.
	\]
	这里整个差正是式\eqref{g2:eq:pontryagin-sum-torsion}的奇下标乘积。换成实或有理系数后 \(t\) 消失，便看不到这个例子中的区别。
\end{example}

至此，三个系列的次数与系数可以并列为
\[
\begin{array}{c|c|c|c}
	\text{丛的类型}&\text{类}&\text{所在群}&\text{构造所用数据}\\\hline
	\text{实秩 }r&w_i&H^i(B;\mathbb F_2)&\text{实射影丛关系}\\
	\text{复秩 }r&c_i&H^{2i}(B;\mathbb Z)&\text{复射影丛关系}\\
	\text{实秩 }r&p_i&H^{4i}(B;\mathbb Z)&(-1)^ic_{2i}(E_{\mathbb C})
\end{array}
\]
下一节将这些类作用于实际切丛和结构群约化；下一章再在实、复系数中用曲率给出代表。整数类及其扭项已经在拓扑构造中确定，曲率比较须与这里的规范一致。
% 纲要条目排为编号小节标题，并进入目录与书签；说明段仍保留于纲要注释中。

% !TeX root = ../../Book2.tex
% 纲要标题：### 21.4 几何计算与约束
% 纲要位置：几何学计划纲要.md，第 1827 行。
% 正文状态：未撰写；本文件目前是框架。

% 纲要内容开始
% 
% * 射影空间切丛
% * 球丛和低维切丛
% * 特征数与结构群约化
% * spin 条件与更高障碍的区别
% 
% 纲要内容结束

\section{几何计算与约束}
\label{g2:sec:2104}

待撰写

\subsection{射影空间切丛}
\label{g2:subsec:2104:01}
待撰写

\subsection{球丛和低维切丛}
\label{g2:subsec:2104:02}
待撰写

\subsection{特征数与结构群约化}
\label{g2:subsec:2104:03}
待撰写

\subsection{spin 条件与更高障碍的区别}
\label{g2:subsec:2104:04}
待撰写
% 纲要条目排为编号小节标题，并进入目录与书签；说明段仍保留于纲要注释中。

% !TeX root = ../../Book2.tex
% 纲要标题：### 21.5 Grassmann 流形的胞腔、Schubert 类与相交计算
% 纲要位置：几何学计划纲要.md，第 1834 行。
% 正文状态：初稿；尚未正式审核。

% 纲要内容开始
% 
% * 固定完全旗与 Schubert 胞腔的秩条件
% * 行阶梯形坐标、偶维胞腔与整数上同调基
% * 通用子丛、商丛和特殊 Schubert 类
% * 利用旗的移动及相交配对计算低阶乘法
% * $\operatorname{Gr}(2,4)$ 的杯积计算与四条直线问题
% 
%   在复数域模型上计算特殊除子类的四次幂，并由一般位置保证几何计数；完整的方程消去、重数与退化核验放在 V.27.5。本节不以一般 Schubert 簇光滑为前提，也不追加完整旗簇表示论。
% 
% * 与 V.27.5 中 Chow 环及 Plücker 方程计算的比较
% 
% ---
% 
% 纲要内容结束

\section{Grassmann 流形的胞腔、Schubert 类与相交计算}
\label{g2:sec:2105}

示性类已经能把向量丛的运算变成上同调计算。现在反过来问：这些上同调类是否有具体的几何代表，它们的杯积又在数什么？复 Grassmann 流形提供了一个可以同时看见两面的模型。一个子空间与固定旗相交的维数条件给出 Schubert 簇；通用商丛的陈类恰好代表其中最简单的一类条件。最后，\(\operatorname{Gr}(2,4)\) 上一个四次杯积将回答三维射影空间中四条直线的公共截线问题。

以下记 \(G=\operatorname{Gr}(k,n)=\operatorname{Gr}_k(\mathbb C^n)\)，其中 \(1\le k<n\)，并置 \(q=n-k\)、\(N=kq\)。它是复维数 \(N\) 的紧光滑流形，始终采用复结构给出的实定向及整数上同调。第1章的\cref{g2:thm:grassmann-schubert-cw}已经证明胞腔的特征映射及闭包关系；这里恢复坐标含义，再用它们构造同调类并计算乘法。有关胞腔、相交与通用丛的经典叙述，可参阅\cite[第1章第5节，The Cell Decomposition、The Schubert Calculus、Universal Bundles]{GriffithsHarris1994}。

\subsection{固定完全旗与 Schubert 胞腔的秩条件}
\label{g2:subsec:2105:01}
固定一条\term[wanquanqi]{完全旗}[complete flag]
\[
	0=F_0\subset F_1\subset\cdots\subset F_n=\mathbb C^n,
	\qquad\dim_{\mathbb C}F_i=i.
\]
选取与它相容的基 \(e_1,\ldots,e_n\)，使 \(F_i=\operatorname{span}(e_1,\ldots,e_i)\)。对 \(V\in G\)，整数 \(\dim(V\cap F_i)\) 从零增到 \(k\)，每步最多增加一，所以恰有 \(k\) 个跳跃位置
\[
	1\le d_1<d_2<\cdots<d_k\le n.
\]
第1章用这些递增指标编号胞腔。为了让指标直接记录余维，现在改用 \(\lambda_i=q+i-d_i\)。于是 \(q\ge\lambda_1\ge\cdots\ge\lambda_k\ge0\)。这样的整数列称为放在 \(k\times q\) 矩形中的一个分拆，记其大小为 \(|\lambda|=\sum_i\lambda_i\)。它可画成左对齐的方格行，行长依次为 \(\lambda_1,\ldots,\lambda_k\)；方格总数将成为复余维。

\begin{definition}\label{g2:def:schubert-partition-locus}
	给定上述旗与分拆，令 \(d_i=q+i-\lambda_i\)。具有跳跃位置恰为 \(d_1,\ldots,d_k\) 的子空间组成 Schubert 胞腔 \(C_\lambda(F_\bullet)\)。其闭包
	\begin{equation}\label{g2:eq:schubert-rank-conditions}
		\Omega_\lambda(F_\bullet)=\{V\in G:\dim(V\cap F_{d_i})\ge i,\ 1\le i\le k\}
	\end{equation}
	称为相对于该旗的\term[Schubert-cu]{Schubert 簇}[Schubert variety]。
\end{definition}

只写 \(\dim(V\cap F_{d_i})=i\) 还不足以定义开胞腔：还要排除跳跃提前发生。因此等价的精确条件是同时要求
\[
	\dim(V\cap F_{d_i})=i,\qquad\dim(V\cap F_{d_i-1})=i-1.
\]
例如 \(k=1\) 时，\(C_{(j)}=\mathbb P(F_{n-j})\setminus\mathbb P(F_{n-j-1})\)，而 \(\Omega_{(j)}=\mathbb P(F_{n-j})\) 是复射影空间中的余维 \(j\) 线性子空间。闭包中的不等式允许跳跃提前发生；它们比开胞腔的等式更适合描述极限。

\ProseParagraphBreak

这些条件也确实是矩阵的秩条件。若 \(V\) 由一个满秩 \(n\times k\) 矩阵 \(A\) 的列张成，则
\[
	\dim(V\cap F_d)=k-\operatorname{rank}(\mathbb C^k\xrightarrow{A}\mathbb C^n\longrightarrow\mathbb C^n/F_d).
\]
因而 \(\dim(V\cap F_d)\ge i\) 等价于最后 \(n-d\) 行组成的矩阵秩不超过 \(k-i\)，可由全部 \((k-i+1)\) 阶子式为零表达。这里条件与所选 \(V\) 的基无关，因为更换基只在右侧乘一个可逆矩阵。它说明 Schubert 簇在 Grassmann 坐标中由多项式方程给出，也说明闭集中的奇点为何可能出现：一个秩条件的不同阶退化会在同一闭包中相遇。

\subsection{行阶梯形坐标、偶维胞腔与整数上同调基}
\label{g2:subsec:2105:02}
在 \(C_\lambda\) 内，从右向左消元得到唯一的基
\begin{equation}\label{g2:eq:schubert-echelon-coordinates}
	v_i=e_{d_i}+\sum_{\substack{a<d_i\\a\notin\{d_1,\ldots,d_{i-1}\}}}z_{ia}e_a,
	\qquad 1\le i\le k.
\end{equation}
第 \(i\) 个向量有 \(d_i-i=q-\lambda_i\) 个自由复系数，因而
\[
	C_\lambda\cong\mathbb C^{N-|\lambda|},\qquad
	\Omega_\lambda=\coprod_{\mu_i\ge\lambda_i\ \text{对所有 }i}C_\mu.
\]
胞腔的实维数是 \(2(N-|\lambda|)\)。例如在 \(\operatorname{Gr}(2,4)\) 中，\(\lambda=(1,0)\) 对应 \((d_1,d_2)=(2,4)\)，可以把一组基作为矩阵的行写成
\[
	\begin{pmatrix}z_{11}&1&0&0\\z_{21}&0&z_{23}&1\end{pmatrix}.
\]
三个自由复数给出六维实胞腔；分拆只有一个方格，故它在八维实 Grassmann 流形中余维二。

\begin{proposition}\label{g2:prop:schubert-integral-basis}
	复 Grassmann 流形 \(G=\operatorname{Gr}(k,n)\) 的整数同调与上同调没有扭元，奇数次群均为零。每个 \(\Omega_\lambda\) 有一个由其开胞腔复定向确定的类
	\[
		[\Omega_\lambda]\in H_{2(N-|\lambda|)}(\Omega_\lambda;\mathbb Z).
	\]
	把这些类推入 \(G\)，得到各次同调的基；其 Poincar\'e 对偶
	\begin{equation}\label{g2:eq:schubert-class-definition}
		\sigma_\lambda=\mathrm{PD}_G[\Omega_\lambda]\in H^{2|\lambda|}(G;\mathbb Z)
	\end{equation}
	称为\term[Schubert-lei]{Schubert 类}[Schubert class]，并构成整数上同调的加法基。这些类与所选完全旗无关。
\end{proposition}
\begin{proof}
	由\cref{g2:thm:grassmann-schubert-cw}，上述集合给出有限 CW 结构，且全部胞腔的实维数为偶数。胞腔微分降低一次，所以其源或靶必为零，全部微分为零。于是同调在每个偶数维自由生成于该维胞腔，在奇数维为零。万有系数定理中的 Ext 项也为零，得到对应的上同调结论。

	\(\Omega_\lambda\) 是一个子复形，唯一的最高维胞腔为 \(C_\lambda\)，其余胞腔的实维数至少低二。故其最高维同调为 \(\mathbb Z\)，到最高维胞腔相对同调的映射为同构。选择在 \(C_\lambda\) 上给出复定向生成元的唯一类，便定义了 \([\Omega_\lambda]\)。包含诱导的胞腔链映射将这个生成元送到 \(G\) 中同一胞腔的生成元，所以所有这类类组成同调基。这里的构造只用胞腔复形，并未要求整个闭包是流形。

	\(G\) 是闭定向流形，\cref{g2:cor:closed-poincare-duality}将这个同调基送到上同调基。任意两条完全旗由某个 \(g\in\operatorname{GL}_n(\mathbb C)\) 联系。该群道路连通：极分解先把正定因子收缩到恒等，酉因子再逐个缩小其特征值的辐角。因此可取从恒等到 \(g\) 的道路，给出 \(G\) 上恒等映射与 \(V\mapsto gV\) 的同伦。后一个映射保持各开胞腔的复定向，并把对应的 Schubert 簇相互送到，所以它们给出相同同调类及对偶类。
\end{proof}

几何相交时，还可以用一个实际的闭光滑流形来代表这个类。令 \(\widetilde\Omega_\lambda\) 由所有旗
\[
	0=W_0\subset W_1\subset\cdots\subset W_k,
	\qquad\dim W_i=i,\quad W_i\subset F_{d_i}
\]
组成，并令 \(\rho_\lambda(W_\bullet)=W_k\)。已选 \(W_{i-1}\) 后，选择 \(W_i\) 等价于在 \(F_{d_i}/W_{i-1}\) 中选一条线，因而每一步都是复射影丛，纤维为 \(\mathbb{CP}^{d_i-i}\)。这说明 \(\widetilde\Omega_\lambda\) 是紧光滑复流形，复维数为 \(\sum_i(d_i-i)=N-|\lambda|\)。

\ProseParagraphBreak

映射 \(\rho_\lambda\) 的像恰为 \(\Omega_\lambda\)：给定满足秩条件的 \(V\)，可在嵌套的 \(V\cap F_{d_i}\) 中逐次选择 \(W_i\)。若 \(V\in C_\lambda\)，则维数迫使 \(W_i=V\cap F_{d_i}\)，选择唯一；在这个开集上，常秩核空间随矩阵坐标光滑且复解析地变化。因此 \(\rho_\lambda\) 在开胞腔上是保持复定向的同构，而
\begin{equation}\label{g2:eq:schubert-resolution-class}
	(\rho_\lambda)_*[\widetilde\Omega_\lambda]=[\Omega_\lambda].
\end{equation}
最后一个等式可限制到最高维胞腔核对：两边都给出其正生成元，上一命题的唯一性使它们相等。这个旗模型将闭包低维处的多重选择留在纤维中，同时为相交理论提供一个光滑源。

\subsection{通用子丛、商丛和特殊 Schubert 类}
\label{g2:subsec:2105:03}
在 \(G\) 上有通用子丛 \(S\) 与商丛 \(Q\)：它们在 \(V\) 处的纤维分别为 \(V\) 和 \(\mathbb C^n/V\)。这是第19章典范丛的有限维模型，构成短正合列
\begin{equation}\label{g2:eq:grassmann-tautological-sequence}
	0\longrightarrow S\longrightarrow\underline{\mathbb C}^{\,n}\longrightarrow Q\longrightarrow0.
\end{equation}
选择 Hermitian 内积后可把 \(Q\) 识别为 \(S^\perp\)，所以 Whitney 和公式给出 \(c(S)c(Q)=1\)。这个等式已经把两张丛的陈类联系起来；以下识别它们所记录的几何条件。

\begin{theorem}\label{g2:thm:special-schubert-chern}
	在 \(G=\operatorname{Gr}(k,n)\) 上，令 \(q=n-k\)。对 \(1\le j\le q\)，有
	\begin{equation}\label{g2:eq:special-schubert-quotient}
		\sigma_j:=\sigma_{(j,0,\ldots,0)}=c_j(Q).
	\end{equation}
	特别地，\(\sigma_1=c_1(Q)=c_1(S^*)=c_1(\det S^*)\)。
\end{theorem}
\begin{proof}
	取 \(A=F_{q+1-j}\)，记 \(a=\dim A=q+1-j\)。对于分拆 \((j,0,\ldots,0)\)，除第一个条件外，其余条件由维数不等式自动成立。因此 \(\Omega_j=\{V:V\cap A\ne0\}\)。为把这个非零交条件写成一个横截零集，在 \(G\times\mathbb P(A)\) 上记 \(\ell\) 为从第二因子拉回的重言线丛，考虑秩 \(q\) 丛 \(\ell^*\otimes Q\) 的截面
	\[
		s(V,L):L\hookrightarrow\mathbb C^n\longrightarrow\mathbb C^n/V.
	\]
	其零集为 \(Z=\{(V,L):L\subset A\cap V\}\)。固定 \(L\) 后，\(V/L\) 可任取 \(\mathbb C^n/L\) 中的 \((k-1)\) 维子空间，所以 \(Z\) 是 \(\mathbb P(A)\) 上的光滑 Grassmann 丛。

	还可直接核对截面横截：在 \(L\subset V\) 处，\(G\) 的切空间由图坐标识别为 \(\operatorname{Hom}_{\mathbb C}(V,\mathbb C^n/V)\)。截面在这些方向的微分，除一个无关紧要的负号外，就是限制到 \(L\) 的映射
	\[
		\operatorname{Hom}(V,\mathbb C^n/V)\longrightarrow\operatorname{Hom}(L,\mathbb C^n/V),
	\]
	它满射，因为 \(L\) 上的线性映射可沿任意补空间延拓。故 \(s\) 横截于零截面。由\cref{g2:thm:top-class-euler}以及\cref{g2:thm:euler-transverse-zero-duality}，\(Z\) 的对偶类是 \(c_q(\ell^*\otimes Q)\)，定向就是复定向。

	令 \(h=c_1(\ell^*)\)。分裂原理把 \(Q\) 的形式根记为 \(x_1,\ldots,x_q\)，从而
	\begin{equation}\label{g2:eq:schubert-incidence-top-chern}
		c_q(\ell^*\otimes Q)=\prod_{b=1}^q(h+x_b)
		=\sum_{i=0}^q h^{q-i}c_i(Q).
	\end{equation}
	这是在旗丛拉回后验证的对称多项式恒等式；拉回单射使它在原空间成立。沿 \(\mathbb P(A)=\mathbb{CP}^{a-1}\) 求值时，只有 \(h^{a-1}=h^{q-j}\) 留下系数一，所以式\eqref{g2:eq:schubert-incidence-top-chern}给出 \(c_j(Q)\)。

	把这个系数提取写成同调推前，便可直接使用已证的乘积求值公式。记投影为 \(\pi:G\times\mathbb P(A)\to G\)。对任意 \(\beta\in H^{2(N-j)}(G;\mathbb Z)\)，顶陈类与零点类的比较给出
	\[
	\begin{aligned}
		\langle\beta,\pi_*[Z]\rangle
		&=\langle\pi^*\beta,[Z]\rangle\\
		&=\left\langle\pi^*\beta\smile\sum_i h^{q-i}c_i(Q),[G\times\mathbb{CP}^{a-1}]\right\rangle\\
		&=\langle\beta\smile c_j(Q),[G]\rangle.
	\end{aligned}
	\]
	全部次数均为偶数，交换因子没有符号。另一方面，\(\pi|_Z\) 的像为 \(\Omega_j\)，在其开胞腔上 \(V\cap A\) 恰为一维，故该处投影是保持复定向的同构。最高维胞腔的归一化给出 \(\pi_*[Z]=[\Omega_j]\)。\(G\) 的整数上同调自由，Poincar\'e 配对非退化，因此上述全部求值确定 \(\sigma_j=c_j(Q)\)。最后对式\eqref{g2:eq:grassmann-tautological-sequence}取一次陈类，并用对偶与行列式的公式，得到 \(\sigma_1\) 的其余表达。
\end{proof}

这个证明的几何步骤是先把“至少有一条交线”改成“连这条线也一起选出”。增加 \(\mathbb P(A)\) 参数后，退化条件成为一张向量丛的横截截面零集；随后沿这些参数求值，消去辅助选择，得到所求陈类。

\subsection{利用旗的移动及相交配对计算低阶乘法}
\label{g2:subsec:2105:04}
同一个 Schubert 类可以由许多不同旗给出的簇代表。移动旗的用途，是让相交发生在光滑开胞腔中并且横截；这时每个复交点贡献 \(+1\)，杯积便成为实际计数。

\begin{proposition}\label{g2:prop:schubert-moving-flags}
	设分拆 \(\lambda^{(1)},\ldots,\lambda^{(m)}\) 满足 \(\sum_b|\lambda^{(b)}|=N\)。可以选择处于一般位置的完全旗 \(F_\bullet^{(b)}\)，使各 \(\Omega_{\lambda^{(b)}}(F_\bullet^{(b)})\) 的公共交集是有限个点，全部位于各自开胞腔中且横截，并且
	\[
		\#\bigcap_{b=1}^m\Omega_{\lambda^{(b)}}(F_\bullet^{(b)})
		=\left\langle\prod_{b=1}^m\sigma_{\lambda^{(b)}},[G]\right\rangle.
	\]
	这里一般位置可取为旗参数空间中的一个开稠密集合；交集允许为空。
\end{proposition}
\begin{proof}
	群 \(\operatorname{GL}_n(\mathbb C)\) 在 \(G\) 上传递作用。作用映射 \((g,V)\mapsto gV\) 是淹没：在图坐标 \(\mathrm T_VG=\operatorname{Hom}(V,\mathbb C^n/V)\) 中，矩阵的一阶变化可以任意指定这张线性映射。因此将一个固定光滑胞腔随 \(g\) 移动，\cref{g2:thm:parametric-transversality}保证对一般参数，它横截于任意另一个固定光滑子流形。

	每个 Schubert 簇只有有限个胞腔。依次移动这些旗，同时对全部胞腔组合应用参数横截性，就使任意组合的相交都横截。若所有胞腔都取开胞腔，预期复维数为 \(N-\sum_b|\lambda^{(b)}|=0\)。若至少一个取闭包中的低维胞腔，其复维数至少再低一，预期维数变成负数，横截性迫使交集为空。所以全部交点落在开胞腔中。簇的闭性与 \(G\) 的紧性使总交集紧，而零维横截交集离散，故只有有限点。

	用式\eqref{g2:eq:schubert-resolution-class}的光滑旗模型代表各同调类。其乘积映射到 \(G^m\)，公共交点对应于这个映射对对角子流形的逆像；由于交点都在开胞腔中，那里各旗模型与 Schubert 簇相同，逆像仍然横截。将\cref{g2:thm:geometric-intersection-pairing}用于对角子流形，再用帽积的自然性及各个基本类的推前，所得求值就是所列杯积在 \([G]\) 上的值。这也把带奇点闭包的计数归结为闭光滑源与光滑对角的相交。

	在交点附近，所有切空间和法向商空间都是复向量空间，横截比较映射为复线性同构，其实行列式为复行列式绝对值的平方，严格为正。因此每个点的符号为 \(+1\)。最后，光滑旗模型紧，横截性在参数的小变动下保持；同时交点避开各闭包的低维部分，这一避开条件也由紧性保持。故所选好参数构成开集，参数横截性给出的稠密性则使它成为开稠密集。
\end{proof}

为了从这些计数读出一个类的系数，需要知道哪两个 Schubert 类相互对偶。定义矩形中的补分拆
\[
	\lambda^c=(q-\lambda_k,q-\lambda_{k-1},\ldots,q-\lambda_1).
\]
它把矩形中余下的方格旋转半周再左对齐，满足 \(|\lambda|+|\lambda^c|=N\)。

\begin{proposition}\label{g2:prop:schubert-complement-pairing}
	若 \(|\lambda|+|\mu|=N\)，则
	\begin{equation}\label{g2:eq:schubert-complement-pairing}
		\langle\sigma_\lambda\smile\sigma_\mu,[G]\rangle
		=\begin{cases}1,&\mu=\lambda^c,\\0,&\mu\ne\lambda^c.\end{cases}
	\end{equation}
\end{proposition}
\begin{proof}
	采用相反的两条坐标旗
	\[
		F_d=\operatorname{span}(e_1,\ldots,e_d),\qquad
		F'_d=\operatorname{span}(e_{n-d+1},\ldots,e_n).
	\]
	若某个 \(i\) 满足 \(\lambda_i+\mu_{k+1-i}>q\)，则相应两个旗子空间的维数和不超过 \(n\)，且它们在相反坐标旗中不相交。然而若 \(V\) 同时满足两组秩条件，它在这两个不相交子空间中分别至少有 \(i\) 维和 \(k+1-i\) 维，总计至少 \(k+1\) 维，矛盾。因此交集非空要求全部 \(\lambda_i+\mu_{k+1-i}\le q\)。将这些不等式求和，左边为 \(N=kq\)，所以每一项都必须取等号，即 \(\mu=\lambda^c\)。

	在补分拆情形，记 \(d_i=q+i-\lambda_i\)。上面的两子空间分别是 \(F_{d_i}\) 与 \(\operatorname{span}(e_{d_i},\ldots,e_n)\)，交集恰为 \(\mathbb Ce_{d_i}\)。由于 \(V\) 在二者中的维数之和至少为 \(k+1\)，它必须包含 \(e_{d_i}\)。对全部 \(i\) 使用此论证，唯一可能的交点便是
	\[
		V_0=\operatorname{span}(e_{d_1},\ldots,e_{d_k}).
	\]
	它确实满足两组条件，并位于两个开胞腔中。

	最后核对横截性。取补空间由非枢轴坐标向量张成，图坐标中的切向量为 \(A:V_0\to\mathbb C^n/V_0\)。由式\eqref{g2:eq:schubert-echelon-coordinates}，第一片的切方向是 \(A(e_{d_i})\) 在非枢轴行 \(a<d_i\) 的系数；相反旗的第二片恰好允许非枢轴行 \(a>d_i\) 的系数。这两组参数不交且合起来覆盖全部图坐标，所以切空间互补。复定向使唯一交点的符号为正，由前一命题的相交解释得到式\eqref{g2:eq:schubert-complement-pairing}。
\end{proof}

因此，若某个齐次类写成 \(\alpha=\sum_\lambda a_\lambda\sigma_\lambda\)，则
\[
	a_\lambda=\langle\alpha\smile\sigma_{\lambda^c},[G]\rangle.
\]
一个乘积的系数就能通过增加一项互补几何条件来求出。下面在最小的非射影空间例子中把这些条件逐项算完。

\subsection{\texorpdfstring{\(\operatorname{Gr}(2,4)\) 的杯积计算与四条直线问题}{Gr(2,4) 的杯积计算与四条直线问题}}
\label{g2:subsec:2105:05}
把 \(\operatorname{Gr}(2,4)\) 的点 \(V\) 看成 \(\mathbb{CP}^3\) 中的射影直线 \(\mathbb P(V)\)。固定射影旗“点 \(P\) 位于直线 \(L\)，而 \(L\) 位于平面 \(H\)”后，六个分拆的含义如下。
\begin{center}
	\begin{tabular}{ccl}
		\hline
		分拆 & 复余维 & 所参数化的射影直线 \\
		\hline
		\(\varnothing\) & \(0\) & 全部直线 \\
		\((1)\) & \(1\) & 与 \(L\) 相交 \\
		\((2)\) & \(2\) & 经过 \(P\) \\
		\((1,1)\) & \(2\) & 包含于 \(H\) \\
		\((2,1)\) & \(3\) & 经过 \(P\) 且包含于 \(H\) \\
		\((2,2)\) & \(4\) & 等于给定直线 \(L\) \\
		\hline
	\end{tabular}
\end{center}
其中 \(\Omega_2\) 与 \(\Omega_{11}\) 都同构于 \(\mathbb{CP}^2\)，\(\Omega_{21}\) 是一束直线组成的 \(\mathbb{CP}^1\)，而 \(\Omega_{22}\) 是一点。\(\Omega_1\) 的闭包却不必光滑，因而前面的光滑代表与移动旗论证在此已有实际用途。为简化式子，省略下标分拆中的逗号。

\begin{proposition}\label{g2:prop:gr24-schubert-products}
	\(G=\operatorname{Gr}(2,4)\) 的非零正次乘积由以下等式及交换律决定：
	\begin{equation}\label{g2:eq:gr24-schubert-products}
	\begin{gathered}
		\sigma_1^2=\sigma_2+\sigma_{11},\qquad
		\sigma_1\sigma_2=\sigma_1\sigma_{11}=\sigma_{21},\\
		\sigma_1\sigma_{21}=\sigma_{22},\qquad
		\sigma_2^2=\sigma_{11}^2=\sigma_{22},\qquad
		\sigma_2\sigma_{11}=0.
	\end{gathered}
	\end{equation}
	特别地，\(\sigma_1^3=2\sigma_{21}\)、\(\sigma_1^4=2\sigma_{22}\)，且 \(\langle\sigma_{22},[G]\rangle=1\)。
\end{proposition}
\begin{proof}
	先求余维互补的乘积。经过两个不同点的射影直线唯一，所以 \(\sigma_2^2=\sigma_{22}\)。同时位于两个不同平面的直线是它们的交线，所以 \(\sigma_{11}^2=\sigma_{22}\)。点若不在给定平面中，就没有既经过该点又包含于该平面的直线，故 \(\sigma_2\sigma_{11}=0\)。这也与补分拆配对相符：\((2)\) 与 \((1,1)\) 各自为自身的补。

	在平面 \(H\) 中经过点 \(P\) 的直线组成一束。另取一条不包含于 \(H\) 的直线 \(L\)，并使交点 \(Q=L\cap H\) 不等于 \(P\)。这束直线中恰有连接 \(P,Q\) 的一条与 \(L\) 相交，所以 \(\sigma_1\sigma_{21}=\sigma_{22}\)。一般位置保证交点横截，复定向使系数为一。

	六次群只由 \(\sigma_{21}\) 生成。若 \(P\notin L\)，经过 \(P\) 且与 \(L\) 相交的直线组成平面 \(\langle P,L\rangle\) 中的直线束，正是某条旗的 \(\Omega_{21}\)。在一般位置下相交横截于光滑部分，因而 \(\sigma_1\sigma_2=\sigma_{21}\)。同理，平面 \(H\) 内与一条横穿 \(H\) 的直线相交，等价于经过 \(L\cap H\)，所以 \(\sigma_1\sigma_{11}=\sigma_{21}\)。也可把这两个等式分别与 \(\sigma_1\) 配对，由刚才的唯一交点计数及非退化性确定系数。

	最后写 \(\sigma_1^2=a\sigma_2+b\sigma_{11}\)。分别乘 \(\sigma_2\) 与 \(\sigma_{11}\) 并在 \([G]\) 上求值，前面已经算得的乘积给出 \(a=b=1\)。其余两式由继续乘 \(\sigma_1\) 得到。次数超过八的类为零，因此没有其他正次乘积需要决定。
\end{proof}

这里的两种余维二条件不能混为一类：经过一个点与位于一个平面都给出一份 \(\mathbb{CP}^2\)，但它们在 \(G\) 中有不同同调类。它们各自的自交为一，彼此的交数为零，正是四次群秩为二的几何表现。

\ProseParagraphBreak

也可以将乘法表写成陈类的环表示。记 \(h=\sigma_1\)、\(s=\sigma_{11}\)。取一个非零线性函数 \(\alpha:\mathbb C^4\to\mathbb C\)，它限制到各 \(V\) 上给出 \(S^*\) 的截面。零点集是 \(V\subset\ker\alpha\)，即 \(\Omega_{11}\)。在这个零集处，图坐标微分为 \(A\mapsto\alpha\circ A\)，满射到 \(V^*\)，所以截面横截；于是
\[
	s=c_2(S^*),\qquad c(S)=1-h+s,
	\qquad c(Q)=1+h+(h^2-s).
\]
其中 \(c_2(Q)=\sigma_2=h^2-s\)，与特殊 Schubert 类公式一致。乘法表等价于
\begin{equation}\label{g2:eq:gr24-cohomology-ring}
	H^*(\operatorname{Gr}(2,4);\mathbb Z)
	\cong\mathbb Z[h,s]/(h^3-2hs,\ s^2-h^2s),
	\qquad\deg h=2,\quad\deg s=4.
\end{equation}
为核对没有漏掉整数关系，先用两条关系推出 \(h^3s=0\)、\(s^3=0\)、\(h^5=0\)：例如 \(h^3s=2hs^2=2h^3s\)，相减即得第一式。故商环由 \(1,h,h^2,s,hs,s^2\) 加法生成；这些元素的像分别为 \(1,\sigma_1,\sigma_2+\sigma_{11},\sigma_{11},\sigma_{21},\sigma_{22}\)，构成同调所给自由基的另一组整数基，所以映射既满又单。

\begin{theorem}\label{g2:thm:four-lines-two-transversals}
	在 \(\mathbb{CP}^3\) 中，对处于一般位置的四条射影直线，恰有两条射影直线同时与它们相交；两条解不同，每条解的局部相交重数为一。
\end{theorem}
\begin{proof}
	每条给定直线产生一份 \(\Omega_1\) 条件。由\cref{g2:prop:schubert-moving-flags}，一般位置下四份条件横截相交于各自开胞腔，故都是孤立的光滑交点，局部符号为 \(+1\)。总数于是为
	\[
		\langle\sigma_1^4,[G]\rangle
		=2\langle\sigma_{22},[G]\rangle=2.
	\]
	横截性的局部方程具有可逆复 Jacobian 矩阵，故每个零点都是单零点，其局部重数为一。
\end{proof}

这个数字还可以在一个直观的二次曲面模型中看出。先取三条两两不交的直线。前两条来自分解 \(\mathbb C^4=A\oplus B\) 的两个二维因子；第三条由于与它们都不交，可以写成某个同构 \(T:A\to B\) 的图。借助 \(T\) 识别 \(B\cong A\)，三条线便为
\[
	L_1=\mathbb P(A\oplus0),\quad L_2=\mathbb P(0\oplus A),\quad
	L_3=\{[a:a]:[a]\in\mathbb P(A)\}.
\]
一条同时遇到前两条线的直线必连接 \([a:0]\) 与 \([0:b]\)。它再与第三条相交，当且仅当 \(a,b\) 成比例。因此三条线的公共截线恰为
\[
	\ell_{[a]}=\{[sa:ta]:[s:t]\in\mathbb{CP}^1\},\qquad[a]\in\mathbb P(A).
\]
这些直线扫出二次曲面
\begin{equation}\label{g2:eq:three-lines-ruled-quadric}
	Q=\{[x_1:x_2:y_1:y_2]:x_1y_2-x_2y_1=0\}
	\cong\mathbb{CP}^1\times\mathbb{CP}^1.
\end{equation}
这个同构就是 \(([a],[s:t])\mapsto[sa:ta]\)：秩一的二乘二矩阵可分解为列向量乘行向量，两个因子分别只差非零倍数，所以参数唯一。非零秩一矩阵处方程的梯度不为零，故曲面光滑。固定 \([s:t]\) 得到包含原三条线的一个线族，固定 \([a]\) 得到公共截线的另一个线族。

\ProseParagraphBreak

一般的第四条线不包含于 \(Q\)，其参数代入二次方程后得到一个有两个不同根的齐次二次式。因而它与 \(Q\) 相交于两个点，每点唯一确定一条 \(\ell_{[a]}\)。这两条截线不同：若两个点属于同一条 \(\ell_{[a]}\)，则第四条线就是这条截线，因而包含于 \(Q\)，与选择矛盾。这给出了两条解的具体几何位置。若第四条线与 \(Q\) 相切，两个交点可能合并；若第四条线本身属于适当线族，还可能出现整族解。一般位置在这里排除的正是这些退化。

\subsection{与 V.27.5 中 Chow 环及 Plücker 方程计算的比较}
\label{g2:subsec:2105:06}
Grassmann 流形也有一个具体的射影方程模型。定义 Pl\"ucker 映射\footnote{Julius Pl\"ucker（1801--1868），德国数学家，研究解析几何、直线几何及数学物理。}
\[
	\iota:G\longrightarrow\mathbb P(\textstyle\bigwedge^k\mathbb C^n),
	\qquad\operatorname{span}(v_1,\ldots,v_k)\longmapsto[v_1\wedge\cdots\wedge v_k].
\]
换基只将外积乘以行列式，所以射影点不变。非零可分解外积 \(\xi=v_1\wedge\cdots\wedge v_k\) 又能恢复原子空间为 \(\{v:v\wedge\xi=0\}\)，故映射单射。在一个非零 \(k\) 阶子式对应的图域中，图矩阵的各项都是其他子式与这一子式之比，所以逆坐标光滑且微分单射。\(G\) 紧，因而 \(\iota\) 是到其像的光滑复嵌入。

\ProseParagraphBreak

其拉回的射影重言线在 \(V\) 处正是 \(\bigwedge^kV\)，所以
\[
	\iota^*\mathcal O(1)=\det S^*,\qquad\iota^*c_1(\mathcal O(1))=\sigma_1.
\]
这里 \(\mathcal O(1)\) 表示复射影空间的超平面线丛。若固定一个 \(q\) 维子空间 \(A\)，投影 \(\mathbb C^n\to\mathbb C^n/A\) 的行列式在 \(\bigwedge^k\mathbb C^n\) 上给出线性函数。它在 \(V\) 的 Pl\"ucker 点处为零，恰好当 \(V\cap A\ne0\)，所以特殊除子 \(\Omega_1\) 确为这一嵌入下的超平面截面。杯积中的每个 \(\sigma_1\)，因此也可以理解为加入一条这样的线性方程。

\ProseParagraphBreak

对 \(\operatorname{Gr}(2,4)\)，令 \(p_{ij}\) 为外积的六个坐标。由 \(\xi\wedge\xi=0\) 得到
\begin{equation}\label{g2:eq:gr24-plucker-quadric}
	p_{12}p_{34}-p_{13}p_{24}+p_{14}p_{23}=0.
\end{equation}
反过来，若某个坐标如 \(p_{12}\ne0\)，把它归一化为一，便可取矩阵
\[
	\begin{pmatrix}
		1&0&-p_{23}&-p_{24}\\
		0&1&p_{13}&p_{14}
	\end{pmatrix}.
\]
它的前五个相应子式是指定坐标，最后一个为 \(p_{13}p_{24}-p_{14}p_{23}=p_{34}\)。其他非零坐标图只需重新排列基，所以式\eqref{g2:eq:gr24-plucker-quadric}正好刻画 Pl\"ucker 像。因而这个 Grassmann 流形是 \(\mathbb{CP}^5\) 中的光滑二次超曲面，\(\int_G\sigma_1^4=2\) 与其二次方程相呼应。

第五卷第27章第5节将沿这个同一模型组织 Chow 环与方程计算：子簇先按有理等价形成代数循环类，再通过复几何的循环类映射与这里的整数上同调比较，并处理具体方程的消去、局部重数和退化参数。当前已经建立的是胞腔基本类、特殊陈类与横截相交计数；一般位置下每个点的局部贡献已由复横截性确定为一。转到代数模型后，重数还要在点合并或交集不横截时继续记录总计数，两种描述因此共享同一算例而承担不同的计算任务。
% 纲要条目排为编号小节标题，并进入目录与书签；说明段仍保留于纲要注释中。


\begin{chaptersummary}
	实线丛的第一 Stiefel--Whitney 类记录方向沿闭路的变化，复线丛的第一陈类是其底层定向实平面丛的 Euler 类。对偶重言复线丛在复射影直线上的第一类取值为 \(1\)，这固定了本章的符号。过渡函数的局部辐角把第一陈类表示为整数二上闭链，实线丛复化则给出整数 Bockstein 类。

	\ProseParagraphBreak

	Leray--Hirsch 定理将整体纤维基转成总空间的自由上同调模。射影丛关系的系数定义各阶陈类与 Stiefel--Whitney 类；迭代射影化所得的旗丛使原丛分裂为线丛，且上同调拉回单射。自然性、Whitney 公式及根的对称计算因而成立。顶陈类等于复定向 Euler 类，顶 Stiefel--Whitney 类等于模二 Euler 类；Thom 平方得到同一组实示性类，对角嵌入又给出 \(w(\mathrm TM)=\operatorname{Sq}(v(M))\)。

	\ProseParagraphBreak

	复化定义整数 Pontryagin 类。它们稳定、自然，但一般直和的整数公式需要保留奇次陈类乘积造成的二阶扭修正。射影空间的切丛稳定同构使三类计算具体化；球面的稳定平凡性则说明这些稳定类与依赖原秩的 Euler 类各有作用。示性类还能给出附加结构的条件：\(w_1\) 检测方向，\(c_1\) 检测行列式平凡性，而已定向丛的 spin 提升恰由 \(w_2\) 控制。

	\ProseParagraphBreak

	复 Grassmann 流形的 Schubert 胞腔给出整数加性基，几何入射条件的对偶类提供乘法解释。特殊 Schubert 类等于商丛的陈类，因而线性子空间之间的相交问题可以转成示性类运算。在 \(\operatorname{Gr}_2(\mathbb C^4)\) 中，四次超平面型 Schubert 条件的乘积求值为 \(2\)，对应三维复射影空间中同时与四条一般位置直线相交的两条直线。
\end{chaptersummary}

\BookExercises
\begin{exercise}\label{g2:ex:line-sections-positive-negative-indices}
	取整数 \(p,q\ge0\)，\(p+q=d>0\)。在 \(\mathbb CP^1\) 上，齐次单项式 \(z_0^pz_1^q\) 给出 \(\mathcal O(d)\) 的一个截面。找出全部零点，计算局部指标，并判断这些零点何时横截。再构造 \(\mathcal O(-d)\) 的光滑截面，使零点仍位于同样的位置，且非零局部指标全部为负。
\end{exercise}
\begin{solution}
	\(\mathcal O(d)=(\gamma^*)^{\otimes d}\)。在每条复直线 \(\ell\subset\mathbb C^2\) 上，齐次多项式限制为一个齐次 \(d\) 次函数，即 \((\ell^*)^{\otimes d}\) 的元素，所以题给式确实定义一个整体截面。

	当 \(p>0\) 时，\([0:1]\) 是零点。在附近以 \(w=z_0/z_1\) 作坐标，以 \((w,1)\) 及其对偶的张量幂作标架，截面系数为 \(w^p\)。在小圆周上归一化后，它是 \(e^{\mathrm i\theta}\mapsto e^{\mathrm i p\theta}\)，故局部指标为 \(p\)。同理，当 \(q>0\) 时，\([1:0]\) 是零点，以 \(v=z_1/z_0\) 作坐标，局部系数为 \(v^q\)，指标为 \(q\)。指数为零的那一端不是零点，此外没有零点。因此指标和是 \(p+q=d\)，与 \(c_1(\mathcal O(d))=dh\) 的求值相同。

	局部映射 \(w\mapsto w^p\) 在零点的实微分可逆恰好当 \(p=1\)；若 \(p>1\)，微分为零。另一端同理。所以所有实际零点横截恰好当 \(p,q\) 中每个非零指数都为一。指标和计入重数，而非只数零点位置。

	将此截面作纤维复共轭，得到 \(\overline{\mathcal O(d)}\) 的光滑截面。Hermitian 内积给出复线性同构 \(\overline{\mathcal O(d)}\cong\mathcal O(d)^*=\mathcal O(-d)\)。在共轭标架中的局部系数分别为 \(\overline w^p\)、\(\overline v^q\)；其归一化圆周映射的度分别为 \(-p\)、\(-q\)。随后施加的丛同构在各纤维中保持复方向，故不改变这些指标。所得截面的指标和为 \(-d\)，零点位置仍与原截面相同。
\end{solution}

\begin{exercise}\label{g2:ex:circle-bundle-torsion-line}
	在 \(\mathbb CP^1\) 上取 \(\mathcal O(m)\)，其中 \(m\ge2\)，令 \(Y_m=S(\mathcal O(m))\) 为其单位圆丛，投影为 \(\pi\)。对 \(Y_m\) 上的线丛 \(L=\pi^*\mathcal O(1)\)，计算 \(c_1(L)\) 的阶，并确定哪些整数 \(k\) 使 \(L^{\otimes k}\) 平凡。说明 \(c_1(L)\) 换成实系数后为何为零，并直接写出 \(L^{\otimes m}\) 的非零截面。
\end{exercise}
\begin{solution}
	\(Y_m\) 是紧光滑流形，第17章的 Morse 胞腔模型使它具有有限 CW 同伦型；结合丛的同伦不变性，可以在它上面使用线丛的上同调分类。记 \(h=c_1(\mathcal O(1))\)。给底层实二维丛取复方向，其 Euler 类为 \(mh\)。\cref{g2:thm:thom-gysin}的相应一段为
	\[
		H^0(\mathbb CP^1;\mathbb Z)=\mathbb Z
		\xrightarrow{\ \times m\ }H^2(\mathbb CP^1;\mathbb Z)=\mathbb Z
		\xrightarrow{\pi^*}H^2(Y_m;\mathbb Z)\longrightarrow0.
	\]
	于是 \(H^2(Y_m;\mathbb Z)\cong\mathbb Z/m\)，\(c_1(L)=\pi^*h\) 是其生成元，恰有阶 \(m\)。线丛的张量分类给出
	\[
		L^{\otimes k}\text{ 平凡}\quad\Longleftrightarrow\quad k\,c_1(L)=0
		\quad\Longleftrightarrow\quad m\mid k.
	\]
	负整数 \(k\) 按对偶解释，结论不变。换成实系数后，乘 \(m\) 是 \(\mathbb R\) 上的同构，同一段序列给出 \(H^2(Y_m;\mathbb R)=0\)，所以该非零整数扭类的实像为零。

	还可以直接看出第 \(m\) 个张量幂平凡。由定义，\(L^{\otimes m}=\pi^*\mathcal O(m)\)。一个点 \(y\in Y_m\) 本身就是 \(\mathcal O(m)_{\pi(y)}\) 中的单位向量，因此
	\[
		s(y)=(y,y)\in\bigl(\pi^*\mathcal O(m)\bigr)_y
	\]
	是处处非零的连续截面。它给出完整线标架，从而直接产生这个张量幂的平凡化。
\end{solution}

\begin{exercise}\label{g2:ex:klein-line-complexification}
	给 Klein 瓶 \(K\) 取一个零胞腔、两个有向一胞腔 \(a,b\) 及一个按字 \(aba^{-1}b\) 附着的二胞腔。令 \(\alpha,\beta\in H^1(K;\mathbb F_2)\) 分别在 \(a,b\) 上取值为 \((1,0)\)、\((0,1)\)。对 \(\varepsilon,\eta\in\mathbb F_2\)，记实线丛 \(L_{\varepsilon,\eta}\) 满足 \(w_1(L_{\varepsilon,\eta})=\varepsilon\alpha+\eta\beta\)。计算这些实线丛复化后的第一陈类，判断复化何时平凡，并说明四个不同实线丛会得到多少种复线丛。
\end{exercise}
\begin{solution}
	附着字中 \(a\) 的代数次数为零，\(b\) 的代数次数为二，所以整数胞腔边界为
	\[
		\partial_2(1)=2b,\qquad\partial_1=0.
	\]
	因此整数胞腔上链的非零余边界是
	\[
		\mathrm d^1:\mathbb Z^2\longrightarrow\mathbb Z,\qquad(u,v)\longmapsto2v,
	\]
	从而 \(H^2(K;\mathbb Z)=\mathbb Z/2\)。记二胞腔对偶上链所代表的生成元为 \(z\)。模二余边界全为零，故 \(\alpha,\beta\) 确实构成一次群的一组基。

	按整数 Bockstein 的除二公式，将 \(\alpha\) 提升为整数上链 \((1,0)\)，其余边界为零，所以 \(\widetilde\beta_2(\alpha)=0\)。将 \(\beta\) 提升为 \((0,1)\)，余边界为 \(2\)，除以二后得到二次上链 \(1\)，所以 \(\widetilde\beta_2(\beta)=z\)。由\cref{g2:prop:real-line-complexification-bockstein}，
	\[
		c_1(L_{\varepsilon,\eta}\otimes_{\mathbb R}\mathbb C)=\eta z.
	\]
	因此恰在 \(\eta=0\) 时复化平凡。原来的四个实线丛由四个不同的 \(w_1\) 区分；复化后只有两种同构类，\(\varepsilon\) 的差别消失，\(\eta\) 的差别保留为整数二阶扭类。这个计算在同一底空间上同时展示了实线丛非平凡性被复化消去和被复化保留的两种情况。
\end{solution}

\begin{exercise}\label{g2:ex:leray-hirsch-orientation-torus}
	令 \(a:S^2\to S^2\) 为对径映射，并取映射环面
	\[
		Y=(S^2\times[0,1])/\bigl((x,1)\sim(a(x),0)\bigr),\qquad p:Y\to S^1.
	\]
	判断纤维的整数二次生成元能否延拓为 \(Y\) 上的整数类，计算 \(H^*(Y;\mathbb Z)\)。再证明模二系数下存在满足 Leray--Hirsch 条件的整体纤维基，并计算模二上同调环。
\end{exercise}
\begin{solution}
	对径映射在 \(H^2(S^2;\mathbb Z)\) 上作用为 \(-1\)，因为\cref{g2:prop:reflection-antipodal-degree}给出它的度为 \((-1)^3=-1\)。若整数生成元来自整体类，沿基圆一周运输后必须保持不变，这与生成元变号矛盾。

	用基圆的零胞腔及一胞腔计算局部系数，\(q=0\) 行是常值 \(\mathbb Z\)，\(q=2\) 行的上链微分是 \(-2:\mathbb Z\to\mathbb Z\)。因此 Serre 谱序列的非零第二页项只有
	\[
		E_2^{0,0}=\mathbb Z,\qquad E_2^{1,0}=\mathbb Z,\qquad E_2^{1,2}=\mathbb Z/2.
	\]
	底空间只有两个可能的列，全部高阶微分为零；每个总次数也至多有一个非零项，故没有群延拓选择。于是
	\[
		H^k(Y;\mathbb Z)=
		\begin{cases}
			\mathbb Z,&k=0,1,\\
			\mathbb Z/2,&k=3,\\
			0,&\text{其他次数}.
		\end{cases}
	\]
	特别地，\(H^2(Y;\mathbb Z)=0\) 也直接排除了整数纤维生成元的延拓。

	模二系数下，\(-1=1\)，两行都是常值 \(\mathbb F_2\)。同一谱序列仍无非零微分，且边缘映射
	\[
		H^2(Y;\mathbb F_2)\longrightarrow H^2(S^2;\mathbb F_2)
	\]
	为同构。选 \(u\) 映到纤维生成元，则 \(1,u\) 满足 Leray--Hirsch 条件。令 \(x=p^*\eta\)，其中 \(\eta\) 是基圆的一次非零类。于是 \(1,x,u,xu\) 为加性基。\(x^2=p^*(\eta^2)=0\)，而 \(u^2=0\)，因为四次群为零；\(xu\ne0\) 由 Leray--Hirsch 的唯一表达保证。因此
	\[
		H^*(Y;\mathbb F_2)\cong\mathbb F_2[x,u]/(x^2,u^2),\qquad |x|=1,\quad |u|=2.
	\]
	同一纤维丛在不同系数下是否具有整体纤维基，取决于运输作用是否保留所需生成元。
\end{solution}

\begin{exercise}\label{g2:ex:projectivization-line-twist}
	设 \(E\to B\) 为复秩 \(r\ge1\) 丛，\(L\to B\) 为复线丛。构造覆盖底空间恒等映射的自然同构 \(\mathbb P(E)\cong\mathbb P(E\otimes L)\)，并比较两张重言线丛及其对偶的一阶类。若 \(r=2\)，将两条射影丛关系的系数分别记为 \(a_1,a_2\) 和 \(a'_1,a'_2\)。用 \(\ell=c_1(L)\) 表示 \(a'_1,a'_2\)。
\end{exercise}
\begin{solution}
	在每个 \(b\in B\) 上，将一条线 \(V\subset E_b\) 送到 \(V\otimes L_b\subset E_b\otimes L_b\)。局部选取 \(L\) 的非零标架后，这是通常的线子空间识别；更换标架只乘一个非零复数，不改变线，所以各局部映射粘合。张量 \(L_b^*\) 给出连续逆，因此得到自然同构 \(f:\mathbb P(E)\to\mathbb P(E\otimes L)\)。按每个点的纤维定义，有
	\[
		f^*\lambda_{E\otimes L}\cong\lambda_E\otimes p^*L.
	\]
	令 \(h=c_1(\lambda_E^*)\)、\(h'=f^*c_1(\lambda_{E\otimes L}^*)\)。线丛的张量及对偶公式给出
	\[
		h'=h-p^*\ell.
	\]
	所以射影丛作为空间相同，其用于记录原丛的重言线类仍可能不同。

	当 \(r=2\) 时，原关系为 \(h^2+p^*a_1h+p^*a_2=0\)。代入 \(h=h'+p^*\ell\)，整理得
	\[
		(h')^2+p^*(a_1+2\ell)h'+p^*(a_2+a_1\ell+\ell^2)=0.
	\]
	由射影关系系数的唯一性，
	\[
		a'_1=a_1+2\ell,\qquad a'_2=a_2+a_1\ell+\ell^2.
	\]
	这也说明原向量丛的信息不仅在射影总空间中，还在其重言线丛及投影中。
\end{solution}

\begin{exercise}\label{g2:ex:exterior-square-rank-three-roots}
	设 \(E\to B\) 为复秩三向量丛，记本章构造的陈类为 \(c_1,c_2,c_3\)。用分裂原理计算 \(\bigwedge^2E\) 的三个陈类，说明计算中的每个整数系数如何出现。
\end{exercise}
\begin{solution}
	在使 \(E\) 分裂的旗空间上写 \(q^*E=L_1\oplus L_2\oplus L_3\)，令 \(x_i=c_1(L_i)\)。外积的双线性给出实际丛分解
	\[
		\bigwedge^2(q^*E)\cong(L_1\otimes L_2)\oplus(L_1\otimes L_3)\oplus(L_2\otimes L_3).
	\]
	线丛张量的一阶类可加，所以新丛的三个形式根是
	\[
		y_1=x_1+x_2,\qquad y_2=x_1+x_3,\qquad y_3=x_2+x_3.
	\]
	它们的一次和为 \(2(x_1+x_2+x_3)\)，因为每个 \(x_i\) 出现两次。两两乘积之和为
	\[
	\begin{aligned}
		y_1y_2+y_1y_3+y_2y_3
		&=x_1^2+x_2^2+x_3^2+3(x_1x_2+x_1x_3+x_2x_3)\\
		&=(x_1+x_2+x_3)^2+(x_1x_2+x_1x_3+x_2x_3).
	\end{aligned}
	\]
	每个平方出现一次，每个混合项出现三次，所以最后恰剩一个二阶初等对称项。三个根的乘积为
	\[
		(x_1+x_2)(x_1+x_3)(x_2+x_3)
		=(x_1+x_2+x_3)(x_1x_2+x_1x_3+x_2x_3)-x_1x_2x_3.
	\]
	右边第一项中三重乘积出现三次，减去一次后留下左边的两次；其余六个三次单项式各出现一次。把初等对称多项式识别为 \(q^*c_i\)，再由 \(q^*\) 单射返回底空间，得到
	\[
		c_1\!\left(\bigwedge^2E\right)=2c_1,\qquad
		c_2\!\left(\bigwedge^2E\right)=c_1^2+c_2,\qquad
		c_3\!\left(\bigwedge^2E\right)=c_1c_2-c_3.
	\]
\end{solution}

\begin{exercise}\label{g2:ex:endomorphism-chern-low-classes}
	设 \(E\to B\) 为 CW 复形上的复秩 \(r\ge1\) 向量丛。计算复向量丛 \(\operatorname{End}(E)=E^*\otimes E\) 的第一、第二陈类，用 \(c_1(E),c_2(E)\) 表示结果。再说明将 \(\operatorname{End}(E)\) 换成无迹自同态丛 \(\operatorname{End}_0(E)\) 时，所有正次陈类为何不变。
\end{exercise}
\begin{solution}
	将 \(E\) 拉回到分裂空间，记其形式根为 \(x_1,\ldots,x_r\)。由\cref{g2:prop:chern-operations}，\(\operatorname{End}(E)\) 的根为全部 \(x_j-x_i\)。当 \(i=j\) 时根为零；其余根可按 \((i,j)\) 与 \((j,i)\) 配对，因此
	\[
		c(\operatorname{End}(E))
		=\prod_{i<j}\bigl(1+(x_j-x_i)\bigr)\bigl(1-(x_j-x_i)\bigr)
		=\prod_{i<j}\bigl(1-(x_j-x_i)^2\bigr)
	\]
	成立于分裂空间。每个因子的第一个正次项都是四次，所以 \(c_1(\operatorname{End}(E))=0\)。四次项给出
	\[
	\begin{aligned}
		c_2(\operatorname{End}(E))
		&=-\sum_{i<j}(x_j-x_i)^2\\
		&=-(r-1)\sum_i x_i^2+2\sum_{i<j}x_ix_j\\
		&=-(r-1)\bigl(c_1(E)^2-2c_2(E)\bigr)+2c_2(E)\\
		&=2r\,c_2(E)-(r-1)c_1(E)^2.
	\end{aligned}
	\]
	第二行中，每个 \(x_i^2\) 在含有下标 \(i\) 的 \(r-1\) 对中出现；混合项来自平方展开的 \(-2x_ix_j\)。旗丛拉回单射使这些等式返回原基空间。配对计算始终使用整数多项式恒等式，没有对上同调类除以二，故包含整数扭信息的基空间也适用。

	迹映射在换架共轭下不变，故 \(\operatorname{End}_0(E)=\ker\operatorname{tr}\) 是复秩 \(r^2-1\) 子丛。映射
	\[
		A\longmapsto\left(\frac{\operatorname{tr}A}{r},\ A-\frac{\operatorname{tr}A}{r}I\right)
	\]
	给出 \(\operatorname{End}(E)\cong\underline{\mathbb C}\oplus\operatorname{End}_0(E)\)，逆映射为 \((z,A_0)\mapsto zI+A_0\)。这里的除以 \(r\) 发生在复纤维的线性代数中。平凡线因子的总陈类为一，所以 Whitney 公式说明两张丛的全部正次陈类相同。
\end{solution}

\begin{exercise}\label{g2:ex:three-real-lines-projective-product}
	设 \(m,n\ge1\)，取 \(B=\mathbb RP^m\times\mathbb RP^n\)。记从两因子拉回的典范实线丛为 \(L_a,L_b\)，令 \(a=w_1(L_a)\)、\(b=w_1(L_b)\)，并置
	\[
		E=L_a\oplus L_b\oplus(L_a\otimes L_b).
	\]
	计算 \(w_1(E),w_2(E),w_3(E)\)，判断 \(E\) 是否可定向、是否平凡，并确定哪些 \((m,n)\) 使它存在处处非零截面。
\end{exercise}
\begin{solution}
	射影空间模二环及域系数 Künneth 公式给出
	\[
		H^*(B;\mathbb F_2)=\mathbb F_2[a,b]/(a^{m+1},b^{n+1}).
	\]
	第三条线的第一类为 \(a+b\)，所以 Whitney 公式得到
	\[
	\begin{aligned}
		w(E)&=(1+a)(1+b)(1+a+b),\\
		w_1(E)&=0,\\
		w_2(E)&=a^2+ab+b^2,\\
		w_3(E)&=a^2b+ab^2.
	\end{aligned}
	\]
	由\cref{g2:prop:first-classes-determinants}，\(w_1=0\) 说明 \(E\) 可定向。由于 \(m,n\ge1\)，单项式 \(ab\) 在上述截断多项式环中始终非零，并且与尚未消失的 \(a^2,b^2\) 线性无关。因此 \(w_2(E)\ne0\)，\(E\) 对所有给定 \((m,n)\) 都不平凡。

	当至少一个 \(m,n\) 大于一时，\(a^2b\)、\(ab^2\) 中至少有一个非零；若两项均非零，它们是不同的基单项式，不能相消。所以此时 \(w_3(E)\ne0\)。\cref{g2:thm:top-class-euler}把它识别为模二 Euler 类，而非零截面会使该类为零，故这些情形没有处处非零截面。

	剩下 \(m=n=1\) 时，底空间是二维环面，\(w_3(E)=0\)。给 \(E\) 选内积，单位球丛的纤维为 \(S^2\)。先在零骨架选单位向量；纤维道路连通，使该选择能沿每条一胞腔延拓。每张二维胞腔上的丛平凡，边界截面对应一条 \(S^1\to S^2\) 的映射，而 \(\pi_1(S^2)=0\)，故还能延入圆盘。底空间没有更高胞腔，得到整体单位截面。因此恰在 \((m,n)=(1,1)\) 时存在非零截面。此时 \(w_2(E)=ab\ne0\)，所以一条非零截面仍不足以形成整个三维丛的标架。
\end{solution}

\begin{exercise}\label{g2:ex:real-projective-wu-formula-all-dimensions}
	设 \(n\ge1\)，令 \(a\in H^1(\mathbb RP^n;\mathbb F_2)\) 为非零类。直接由 Wu 类的求值定义与射影空间上的平方公式证明
	\[
		v_i(\mathbb RP^n)=\binom{n-i}{i}a^i\qquad(0\le i\le n),
	\]
	其中二项式系数取模二，且下标大于非负上标时约定为零。说明 \(i>n/2\) 时为何消失，并分别计算 \(n=4,5\) 时的总 Wu 类及由 Wu 公式得到的切丛总 Stiefel--Whitney 类。
\end{exercise}
\begin{solution}
	对 \(0\le i\le n\)，\(H^i(\mathbb RP^n;\mathbb F_2)\) 由 \(a^i\) 生成，因而可写 \(v_i=\lambda_i a^i\)。模二基本类满足 \(\langle a^n,[\mathbb RP^n]_2\rangle=1\)，所以把 Wu 类的定义用于 \(a^{n-i}\)，得到
	\[
	\begin{aligned}
		\lambda_i
		&=\langle v_i a^{n-i},[\mathbb RP^n]_2\rangle\\
		&=\langle\operatorname{Sq}^i(a^{n-i}),[\mathbb RP^n]_2\rangle
		=\binom{n-i}{i}.
	\end{aligned}
	\]
	最后一步使用式\eqref{g2:eq:projective-steenrod-formula}。若 \(i>n/2\)，则 \(i>n-i\)，平方次数超过输入类的次数，不稳定条件使它为零；这也与二项式系数的约定一致。\(i>n\) 时所在上同调群本来为零。

	当 \(n=4\) 时，可能的正次项只有一次和二次，系数分别为 \(\binom31=1\)、\(\binom22=1\)（模二），故
	\[
		v(\mathbb RP^4)=1+a+a^2.
	\]
	由平方公式，\(\operatorname{Sq}(a)=a+a^2\)，\(\operatorname{Sq}(a^2)=a^2+a^4\)，中间的 \(\operatorname{Sq}^1(a^2)=2a^3=0\)。因此\cref{g2:thm:wu-formula}给出
	\[
		w(\mathrm T\mathbb RP^4)=\operatorname{Sq}(1+a+a^2)=1+a+a^4.
	\]
	当 \(n=5\) 时，\(\binom41=0\)、\(\binom32=1\)（模二），所以
	\[
		v(\mathbb RP^5)=1+a^2,\qquad
		w(\mathrm T\mathbb RP^5)=1+a^2+a^4.
	\]
	两种情形中 Wu 类的高次部分都先消失，随后对低次 Wu 类施加平方，仍能产生切丛较高次的非零类。两组结果也分别等于截断后的 \((1+a)^5\)、\((1+a)^6\)，与射影空间切丛的直接计算一致。
\end{solution}

\begin{exercise}\label{g2:ex:projective-stiefel-binary}
	对 \(n\ge1\)，证明 \(\mathbb RP^n\) 的全部正次 Stiefel--Whitney 类为零，当且仅当 \(n+1\) 是二的幂。再对 \(\mathbb RP^3\)、\(\mathbb RP^5\)、\(\mathbb RP^7\) 分别判断可定向性与 spin 结构的存在性，并计算存在时 spin 结构的个数。
\end{exercise}
\begin{solution}
	令 \(m=n+1\)，在 \(\mathbb F_2[a]/(a^m)\) 中有 \(w(\mathrm T\mathbb RP^n)=(1+a)^m\)。将 \(m\) 写成二进制 \(m=\sum_{j\in J}2^j\)。因为特征二下 \((1+a)^{2^j}=1+a^{2^j}\)，所以
	\[
		(1+a)^m=\prod_{j\in J}(1+a^{2^j}).
	\]
	若 \(m=2^k\)，右端为 \(1+a^m=1\)。反过来，若 \(J\) 至少有两个元素，令 \(j_0=\min J\)。右端 \(a^{2^{j_0}}\) 的系数为一：只能从第 \(j_0\) 个因子选择正次项，其他正次项的次数都更大。而 \(2^{j_0}<m\)，这一项在商环中不为零。因此总类不等于一，得到所求判据。

	三个维数都是奇数，所以相应实射影空间都可定向。\(w_2\) 的系数依次是 \(\binom42=6\)、\(\binom62=15\)、\(\binom82=28\)，故 \(\mathbb RP^3\)、\(\mathbb RP^7\) 有 spin 结构，\(\mathbb RP^5\) 没有。前两者的一次模二上同调均为 \(\mathbb F_2\)，所以各有两份 spin 结构。正次 Stiefel--Whitney 类全零是示性类层面的结论；判定实际切丛是否平凡，还需研究完整的标架延拓。
\end{solution}

\begin{exercise}\label{g2:ex:product-projective-chern-numbers}
	给 \(M=\mathbb CP^1\times\mathbb CP^2\) 乘积复结构。记两个因子的正生成元拉回为 \(x,y\)，所以 \(x^2=0\)、\(y^3=0\)，且 \(\langle xy^2,[M]\rangle=1\)。求 \(c_1,c_2,c_3\)、\(p_1(\mathrm TM)\)，计算三个陈数 \(\langle c_1^3,[M]\rangle\)、\(\langle c_1c_2,[M]\rangle\)、\(\langle c_3,[M]\rangle\)，并判断 \(M\) 是否有 spin 结构。
\end{exercise}
\begin{solution}
	乘积切丛分裂为两个因子的切丛拉回，所以
	\[
		c(\mathrm TM)=(1+x)^2(1+y)^3
		=(1+2x)(1+3y+3y^2).
	\]
	按次数收集得到
	\[
		c_1=2x+3y,\qquad c_2=6xy+3y^2,\qquad c_3=6xy^2.
	\]
	因此
	\[
		p_1=c_1^2-2c_2=(12xy+9y^2)-(12xy+6y^2)=3y^2.
	\]
	只有 \(xy^2\) 项在顶次求值中有贡献，于是
	\[
		c_1^3=54xy^2,\qquad c_1c_2=24xy^2,\qquad c_3=6xy^2.
	\]
	三个数分别为 \(54,24,6\)。最后 \(w_2(\mathrm TM)=\rho_2c_1=\overline y\ne0\)，其非零性可通过限制到第二个因子的射影直线检验。因此 \(M\) 没有 spin 结构。
\end{solution}

\begin{exercise}\label{g2:ex:two-spin-lifts-circle}
	设 \(E=S^1\times\mathbb R^3\) 带标准方向与度量。用 \(\operatorname{Spin}(3)=S^3\) 及四元数覆盖 \(\rho:S^3\to\operatorname{SO}(3)\)，对 \(\epsilon\in\{1,-1\}\) 定义
	\[
		Q_\epsilon=(\mathbb R\times S^3)/\!\sim,
		\qquad(t+1,q)\sim(t,\epsilon q).
	\]
	证明 \(\eta_\epsilon([t,q])=(e^{2\pi\mathrm i t},\rho(q))\) 给出 \(E\) 的 spin 结构。证明二者作为覆盖固定标架丛的 spin 结构不同构，却作为抽象主 \(S^3\) 丛都平凡，并说明二者为何穷尽所有 spin 结构。
\end{exercise}
\begin{solution}
	\(\epsilon\) 位于 \(S^3\) 的中心，因此右乘作用 \([t,q]g=[t,qg]\) 有定义，短弧上的实参数提升给出主丛局部平凡化。又 \(\rho(\epsilon q)=\rho(q)\)，故 \(\eta_\epsilon\) 尊重商关系，且逐纤维为 \(\rho\)、满足所需等变性。

	取固定标架丛中的截面 \(s(z)=(z,I_3)\)。其在 \(Q_\epsilon\) 中的原像是
	\[
		(\mathbb R\times\{1,-1\})/\bigl((t+1,a)\sim(t,\epsilon a)\bigr).
	\]
	\(\epsilon=1\) 时，这是两个互不相交的圆；\(\epsilon=-1\) 时，两张局部覆盖沿一周互换，得到一个连通的二重圆周覆盖。任何与 \(\eta\) 交换的同构都须把这两份原像同构地对应，连通分支数不同使它不可能存在。

	作为抽象主丛，\(Q_1\) 已经是乘积。对 \(Q_{-1}\)，取单位纯虚四元数 \(u\)，令 \(h(t)=\cos(\pi t)+u\sin(\pi t)\)，则 \(h(t+1)=-h(t)\)。映射
	\[
		[t,q]\longmapsto(e^{2\pi\mathrm i t},h(t)q)
	\]
	尊重接缝、保持右作用，并在局部具有连续逆，所以给出主丛平凡化。它改变了到固定 \(\operatorname{SO}(3)\) 标架丛的映射，因而不与前面不同构的结论冲突。由\cref{g2:thm:spin-obstruction}，所有 spin 结构同构类受 \(H^1(S^1;\mathbb F_2)=\mathbb F_2\) 自由传递作用，恰有两份，故上述构造已经穷尽。
\end{solution}

\begin{exercise}\label{g2:ex:smooth-projective-surface-numbers}
	设 \(F(z_0,z_1,z_2,z_3)\) 为次数 \(d\ge1\) 的复齐次多项式，其射影零点集 \(X\subset\mathbb CP^3\) 是光滑复曲面；等价地，所定义的齐次方程沿零点集的法向微分处处满射。记 \(h=c_1(\lambda^*)\)，并用同一字母表示其在 \(X\) 上的限制。证明
	\[
		c_1(\mathrm TX)=(4-d)h,\qquad
		c_2(\mathrm TX)=(d^2-4d+6)h^2,\qquad
		\langle h^2,[X]\rangle=d.
	\]
	据此计算 \(X\) 的两个陈数及第一 Pontryagin 数，并证明 \(X\) 有 spin 结构当且仅当 \(d\) 为偶数。
\end{exercise}
\begin{solution}
	在一条直线 \(\ell\subset\mathbb C^4\) 上，\(F\) 的限制是 \(\ell\) 上的次数 \(d\) 齐次函数，所以是 \((\ell^*)^{\otimes d}\) 的元素。因此 \(F\) 给出线丛 \((\lambda^*)^{\otimes d}\) 的光滑截面，零集正是 \(X\)。题设满秩条件说明该截面横截于零，并将复法丛识别为 \((\lambda^*)^{\otimes d}|_X\)。选 Hermitian 度量分裂切丛的短正合列，得到复丛同构
	\[
		\mathrm T\mathbb CP^3|_X\cong\mathrm TX\oplus(\lambda^*)^{\otimes d}|_X.
	\]
	线丛张量公式给出后一项的第一陈类为 \(dh\)，因此
	\[
		c(\mathrm TX)=\frac{(1+h)^4}{1+dh}
		=(1+4h+6h^2)(1-dh+d^2h^2)
		=1+(4-d)h+(d^2-4d+6)h^2,
	\]
	这里只保留实次数不超过四的项。

	由\cref{g2:thm:euler-transverse-zero-duality}，\(X\) 的对偶类是这条复线丛的 Euler 类，即 \(dh\)。复法向实维数为二，交换法向与切向的方向不会产生符号，所以该零点方向与 \(X\) 的复方向一致。于是
	\[
		\langle h^2,[X]\rangle=\langle h^2\smile dh,[\mathbb CP^3]\rangle=d.
	\]
	从而
	\[
	\begin{aligned}
		\langle c_1^2,[X]\rangle&=d(4-d)^2,\\
		\langle c_2,[X]\rangle&=d(d^2-4d+6),\\
		\langle p_1(\mathrm TX),[X]\rangle&=\langle c_1^2-2c_2,[X]\rangle=d(4-d^2).
	\end{aligned}
	\]
	最后 \(w_2(\mathrm TX)=(4-d)\overline h\)。\(d\) 偶数时它为零。\(d\) 奇数时，\(\langle\overline h^{\,2},[X]_2\rangle=1\)，故 \(\overline h\ne0\)，于是 \(w_2\ne0\)。由 spin 提升准则得到所求充要条件。比如 \(d=4\) 时两种陈数为 \(0,24\)，第一 Pontryagin 数为 \(-48\)；这些数均由上述切丛与零点类计算得到。
\end{solution}

\begin{exercise}\label{g2:ex:gr25-cells-pairing}
	列出 \(\operatorname{Gr}(2,5)\) 的各次整数 Betti 数与 Euler 示性数。给出四次、八次上同调的 Schubert 基，并写出两组基之间的相交配对矩阵。
\end{exercise}
\begin{solution}
	此时 \(k=2\)、\(q=3\)、\(N=6\)。矩形中的分拆按大小排列为
	\[
	\begin{array}{c|l}
		|\lambda|&\lambda\\\hline
		0&(0,0)\\
		1&(1,0)\\
		2&(2,0),(1,1)\\
		3&(3,0),(2,1)\\
		4&(3,1),(2,2)\\
		5&(3,2)\\
		6&(3,3).
	\end{array}
	\]
	故从零次到十二次的偶数 Betti 数依次为 \(1,1,2,2,2,1,1\)，全部奇数 Betti 数为零。Euler 示性数为这些数之和，即 \(10\)。这里每个分拆贡献一个偶维胞腔，所以也可直接数得 \(\binom52=10\)。

	四次上同调取基 \((\sigma_2,\sigma_{11})\)，八次取基 \((\sigma_{31},\sigma_{22})\)。补分拆公式为 \(\lambda^c=(3-\lambda_2,3-\lambda_1)\)，所以 \((2,0)^c=(3,1)\)、\((1,1)^c=(2,2)\)。按上述次序，相交配对矩阵为
	\[
		\begin{pmatrix}1&0\\0&1\end{pmatrix}.
	\]
	这里第二组基与第一组基所在次数不同，因为 Grassmann 流形实维数为十二，互补次数是四与八。
\end{solution}

\begin{exercise}\label{g2:ex:gr24-tangent-chern-numbers}
	在 \(G=\operatorname{Gr}(2,4)\) 上令 \(h=\sigma_1\)、\(s=\sigma_{11}\)。证明复切丛同构于 \(S^*\otimes Q\)，计算 \(c_1(\mathrm TG)\)、\(c_2(\mathrm TG)\)、\(c_4(\mathrm TG)\)，并将顶陈数与胞腔计算得到的 Euler 示性数比较。
\end{exercise}
\begin{solution}
	在 \(V\) 附近，另一个二维子空间可写成从 \(V\) 到补空间的线性映射的图。对这些图在零映射处求微分，得到 \(\mathrm T_VG=\operatorname{Hom}_{\mathbb C}(V,\mathbb C^4/V)\)。更换补空间不改变商值的一阶映射，所以这些识别粘成 \(\mathrm TG\cong S^*\otimes Q\)。

	在共同分裂空间上记 \(S^*\) 的根为 \(x_1,x_2\)，\(Q\) 的根为 \(y_1,y_2\)。由正文计算，
	\[
		x_1+x_2=y_1+y_2=h,\qquad x_1x_2=s,\qquad y_1y_2=h^2-s.
	\]
	切丛四个根为 \(x_i+y_j\)。它们之和为 \(4h\)，所以 \(c_1(\mathrm TG)=4h\)。平方和为
	\[
	\begin{aligned}
		\sum_{i,j}(x_i+y_j)^2
		&=2(h^2-2s)+2\bigl(h^2-2(h^2-s)\bigr)+2h^2\\
		&=2h^2.
	\end{aligned}
	\]
	由根的对称多项式恒等式 \(2c_2=c_1^2-\sum_{i,j}(x_i+y_j)^2\)，得到 \(2c_2=14h^2\)。整数上同调没有二扭，故 \(c_2(\mathrm TG)=7h^2\)。

	顶类是四根之积。先对 \(y_j\) 求积，并利用 \(x_i^2=hx_i-s\)，有
	\[
	\begin{aligned}
		c_4(\mathrm TG)
		&=\prod_{i=1}^2\bigl(x_i^2+hx_i+h^2-s\bigr)\\
		&=\prod_{i=1}^2(2hx_i+h^2-2s)\\
		&=3h^4-4h^2s+4s^2=3h^4.
	\end{aligned}
	\]
	最后一步使用 \(s^2=h^2s\)。共同旗丛拉回单射，因此这些根计算给出原空间上的等式。由 \(\int_Gh^4=2\)，顶陈数为 \(6\)。\(G\) 恰有六个偶维胞腔，Euler 示性数也为 \(6\)；两者相等是复切丛顶陈类、实 Euler 类与\cref{g2:cor:tangent-euler-characteristic}的具体核验。
\end{solution}

\begin{exercise}\label{g2:ex:schubert-divisor-singularity}
	令 \(A\subset\mathbb C^4\) 为固定二维子空间，\(\Omega_1=\{V\in\operatorname{Gr}(2,4):V\cap A\ne0\}\)。证明 \(\Omega_1\) 在点 \(V=A\) 处不光滑，并在该点附近给出它的方程。再说明正文中的关联空间 \(Z=\{(V,L):L\subset A\cap V\}\) 在这个点上的投影纤维是什么。
\end{exercise}
\begin{solution}
	选择补空间 \(B\)，把 \(A\) 附近的 \(V\) 写成线性映射 \(T:A\to B\) 的图。则 \(V\cap A\) 对应 \(\ker T\)，所以局部方程为
	\[
		\det T=t_{11}t_{22}-t_{12}t_{21}=0
	\]
	于 \(\operatorname{Hom}(A,B)\cong\mathbb C^4\) 中。在 \(T\ne0\) 且行列式为零时，\(T\) 的秩为一，行列式的至少一个偏导数非零，故该处为复三维光滑超曲面。

	在 \(T=0\) 处，集合是所有秩至多一矩阵组成的锥。若它在那里是复三维光滑子流形，则其切空间必须包含每条曲线 \(t\mapsto tE_{ij}\) 在零点的速度。四个矩阵单位 \(E_{ij}\) 都秩一，且线性张成全部 \(\mathbb C^4\)，这与切空间只有复维数三矛盾。因此原点确实是奇点，而不仅是这个方程的一阶导数恰好同时消失。

	投影到 \(V=A\) 时，条件为任意直线 \(L\subset A\)，所以纤维是 \(\mathbb P(A)\cong\mathbb{CP}^1\)。在其余 \(\Omega_1\) 的点上，\(V\cap A\) 一维，纤维只有一个点。关联空间把奇点处不唯一的交线选择保留下来；正文通过其基本类推前定义同一个 Schubert 类，因此闭包有奇点并不妨碍类的存在。
\end{solution}

\begin{exercise}\label{g2:ex:real-four-lines-count}
	在 \(\mathbb C^4=\mathbb C^2\oplus\mathbb C^2\) 中取射影直线
	\[
		L_1=\mathbb P(\mathbb C^2\oplus0),\quad L_2=\mathbb P(0\oplus\mathbb C^2),\quad L_3=\mathbb P\{(a,a):a\in\mathbb C^2\}.
	\]
	第四条直线取 \(L_T=\mathbb P\{(a,Ta):a\in\mathbb C^2\}\)。分别令
	\[
		T_+=\begin{pmatrix}0&2\\3&0\end{pmatrix},\qquad
		T_-=\begin{pmatrix}0&-1\\2&0\end{pmatrix}.
	\]
	证明每组四条线两两不交，并求所有复公共截线。两组数据都由实矩阵给出，它们各有多少条实公共截线？
\end{exercise}
\begin{solution}
	前面三条直线两两不交。图直线与 \(L_2\) 不交是因为第一分量为零会迫使 \(a=0\)；与 \(L_1\) 不交等价于 \(T\) 可逆，与 \(L_3\) 不交等价于 \(T-I\) 可逆。两个所给矩阵分别满足
	\[
		\det T_+=-6,\quad\det(T_+-I)=-5,\qquad
		\det T_-=2,\quad\det(T_--I)=3,
	\]
	所以四条线两两不交。

	前三条线的公共截线由正文给出，为 \(\ell_{[a]}=\{[sa:ta]:[s:t]\in\mathbb{CP}^1\}\)。它与 \(L_T\) 相交，恰好当 \(Ta\) 与 \(a\) 成比例，即 \([a]\) 是 \(T\) 的特征线。对 \(T_+\)，特征值为 \(\pm\sqrt6\)，可取特征向量 \(a_\pm=(2,\pm\sqrt6)\)。所以两条公共截线为 \(\ell_{[a_+]}\)、\(\ell_{[a_-]}\)，均由实向量张成，是两条不同的实射影直线。

	对 \(T_-\)，特征值为 \(\pm\mathrm i\sqrt2\)，对应特征向量可取 \((1,\mp\mathrm i\sqrt2)\)。仍有两条不同的复公共截线，它们互为复共轭，且都不由实二维子空间定义。因为实矩阵的实特征线必有实特征值，而这里没有实特征值，所以没有实公共截线。两个例子说明复交数 \(2\) 可以对应两条实解，也可以对应一对非实共轭解，不能直接把复计数当作实计数。
\end{solution}

% !TeX root = ../../Book2.tex
% 纲要标题：## 第22章　Chern–Weil 与次级示性形式
% 纲要位置：几何学计划纲要.md，第 1848 行。
% 正文状态：未撰写；本文件目前是框架。

% 纲要内容开始
% 
% > 本章目标与证明范围：在实或复系数上完整比较拓扑类与曲率代表。Chern–Simons 的转移构造给出明确选择依赖；Euler 形式为第三卷高维 Gauss–Bonnet 提供规范一致的输入。
% 
% 纲要内容结束

\chapter{Chern–Weil 与次级示性形式}
\label{g2:ch:22}
\BookChapterNumber{g2:ch:22}

\begin{chapterabstract}
待撰写
\end{chapterabstract}

\begin{learninggoals}
待撰写
\end{learninggoals}

待撰写

% !TeX root = ../../Book2.tex
% 纲要标题：### 22.1 不变多项式与曲率
% 纲要位置：几何学计划纲要.md，第 1852 行。
% 正文状态：未撰写；本文件目前是框架。

% 纲要内容开始
% 
% * 对称不变多项式及极化
% * 丛值外乘与迹
% * Bianchi 恒等式给出的闭性
% 
% 纲要内容结束

\section{不变多项式与曲率}
\label{g2:sec:2201}

待撰写

\subsection{对称不变多项式及极化}
\label{g2:subsec:2201:01}
待撰写

\subsection{丛值外乘与迹}
\label{g2:subsec:2201:02}
待撰写

\subsection{Bianchi 恒等式给出的闭性}
\label{g2:subsec:2201:03}
待撰写
% 纲要条目排为编号小节标题，并进入目录与书签；说明段仍保留于纲要注释中。

% !TeX root = ../../Book2.tex
% 纲要标题：### 22.2 转移形式与联络无关性
% 纲要位置：几何学计划纲要.md，第 1858 行。
% 正文状态：初稿；尚未正式审核。

% 纲要内容开始
% 
% * 联络路径及参数积分
% * 陈–Weil 转移公式
% * 自然性和陈–Weil 同态
% * 规范变换下的变化
% 
% 纲要内容结束

\section{转移形式与联络无关性}
\label{g2:sec:2202}

同一张丛可以选择许多联络，它们的曲率通常也不同。上一节证明了不变多项式作用于曲率后得到闭形式；要让这个形式代表丛本身的类，还须解释改变联络为何只增加一个恰当形式。办法是让联络随一个参数连续变化，把曲率代表的变化率写成外微分，再沿参数积分。所得的原始形式还保留了两条联络之间的比较信息，将成为次级示性形式的来源。

\subsection{联络路径及参数积分}
\label{g2:subsec:2202:01}
固定光滑主右矩阵 Lie 群丛 \(Q\to M\)，其结构群为 \(G\)，Lie 代数为 \(\mathfrak g\)。本节的流形有限维、Hausdorff、第二可数；参数族在闭区间的端点光滑，意思是局部可以光滑延拓到稍大的区间。第一卷第18.2节“主联络的仿射空间”命题说明，两条联络之差属于 \(\Omega^1(M;\operatorname{ad}Q)\)。恢复它在局部的含义：若联络矩阵为 \(A_0,A_1\)，则换架时
\[
	A_j'=g^{-1}A_jg+g^{-1}\mathrm dg\quad(j=0,1),
	\qquad a'=g^{-1}ag,\quad a=A_1-A_0.
\]
相减后微分项消去，故 \(a\) 是一张整体伴随丛值一形式。反过来，将任意这样的 \(a\) 加到一条联络上，仍满足联络的换架规律。因此
\begin{equation}\label{g2:eq:affine-connection-path}
	\nabla_t=\nabla_0+t a,\qquad A_t=A_0+t a,\qquad 0\le t\le1
\end{equation}
给出连接两条联络的路径。对酉联络或正交联络，差分别取值于斜 Hermitian 或反对称矩阵，这条路径也保持相应相容性。

\begin{proposition}\label{g2:prop:curvature-path-variation}
	设 \(Q\to M\) 为光滑主右矩阵群丛，\(\nabla_t\) 是其光滑联络路径，局部矩阵与曲率分别为 \(A_t\)、\(F_t=\mathrm dA_t+A_t\wedge A_t\)。记 \(\dot A_t=\partial_tA_t\)，并以 \(\mathrm D_t\) 表示伴随丛值形式的协变外微分。则 \(\dot A_t\) 是整体伴随丛值一形式，而且
	\begin{equation}\label{g2:eq:curvature-path-variation}
		\partial_tF_t=\mathrm D_t\dot A_t.
	\end{equation}
	若 \(P:\mathfrak g^k\to\mathbb K\) 是对称 \(\operatorname{Ad}(G)\) 不变 \(k\) 线性式，\(k\ge1\)、\(\mathbb K=\mathbb R\) 或 \(\mathbb C\)，则
	\begin{equation}\label{g2:eq:chern-weil-form-variation}
		\partial_tP(F_t,\ldots,F_t)
		=k\,\mathrm dP(\dot A_t,F_t,\ldots,F_t),
	\end{equation}
	右侧有 \(k-1\) 个曲率因子。
\end{proposition}
\begin{proof}
	取不随参数变化的局部截面。过渡函数 \(g\) 与 \(t\) 无关，对换架公式求导即得 \(\dot A_t'=g^{-1}\dot A_tg\)，所以它们确实粘成伴随丛值形式。逐项求导给出
	\[
		\partial_tF_t=\mathrm d\dot A_t+\dot A_t\wedge A_t+A_t\wedge\dot A_t.
	\]
	一形式的协变外微分为 \(\mathrm D_t\dot A_t=\mathrm d\dot A_t+A_t\wedge\dot A_t+\dot A_t\wedge A_t\)，这证明第一式。两个乘积项都是加号，是因为被微分的形式次数为一。

	每个 \(F_t\) 的次数为二，交换两个曲率因子不产生符号。因此 \(P\) 的对称性给出
	\[
		\partial_tP(F_t,\ldots,F_t)
		=kP(\mathrm D_t\dot A_t,F_t,\ldots,F_t).
	\]
	另一方面，\cref{g2:lem:invariant-form-differential}将普通外微分改写为各因子的协变外微分。作用于第一个因子得到上式右边除去 \(k\) 的部分；作用于其余因子的项都含有 \(\mathrm D_tF_t\)，由\cref{g2:prop:cw-bianchi}为零。故
	\[
		\mathrm dP(\dot A_t,F_t,\ldots,F_t)
		=P(\mathrm D_t\dot A_t,F_t,\ldots,F_t),
	\]
	完成第二式的证明。
\end{proof}

沿参数积分是对形式系数作通常的一元积分。更准确地，对光滑族 \(\alpha_t\in\Omega^r(M;\mathbb K)\)，定义
\[
	\left(\int_0^1\alpha_t\,\mathrm dt\right)_x(v_1,\ldots,v_r)
	=\int_0^1(\alpha_t)_x(v_1,\ldots,v_r)\,\mathrm dt.
\]
紧参数区间上的光滑依赖允许在局部坐标中逐项求导，因此积分仍为光滑 \(r\) 形式，并且与 \(\mathrm d_M\) 交换。这个积分不改变底空间形式的次数。

\ProseParagraphBreak

也可以把上述计算放到 \([0,1]\times M\) 上。令拉回丛在参数方向没有联络分量，则它的曲率为
\[
	\widehat F=F_t+\mathrm dt\wedge\dot A_t.
\]
我们将参数微分写在前面。由于 \((\mathrm dt)^2=0\)，有
\begin{equation}\label{g2:eq:chern-weil-cylinder-form}
	P(\widehat F,\ldots,\widehat F)
	=P(F_t,\ldots,F_t)+\mathrm dt\wedge kP(\dot A_t,F_t,\ldots,F_t).
\end{equation}
总空间上的闭性在 \(\mathrm dt\) 分量中恰好给出式\eqref{g2:eq:chern-weil-form-variation}。这样，联络变化率成为参数方向的一份曲率，而转移形式将由含有这份曲率的项积分得到。

\subsection{陈–Weil 转移公式}
\label{g2:subsec:2202:02}
\begin{definition}\label{g2:def:chern-weil-transgression}
	设 \(\Gamma=(\nabla_t)_{0\le t\le1}\) 是固定光滑主 \(G\) 丛上的光滑联络路径，\(P\) 是次数 \(k\ge1\) 的对称不变多项式。形式
	\begin{equation}\label{g2:eq:chern-weil-transgression}
		T_P(\Gamma)=k\int_0^1P(\dot A_t,F_t,\ldots,F_t)\,\mathrm dt
		\in\Omega^{2k-1}(M;\mathbb K)
	\end{equation}
	称为这条路径的\term[zhuanyixingshi]{转移形式}[transgression form]。对式\eqref{g2:eq:affine-connection-path}的仿射路径，简记为 \(T_P(\nabla_0,\nabla_1)\)。
\end{definition}

这里每个被积项都是整体标量形式：\(\dot A_t\) 和 \(F_t\) 都按共轭变换，而 \(P\) 不变。其次数为 \(1+2(k-1)=2k-1\)，比曲率代表低一次。

\begin{theorem}[陈--Weil 转移公式]\label{g2:thm:chern-weil-transgression}
	设 \(Q\to M\) 为光滑主右矩阵群丛，\(P\) 为次数 \(k\ge1\) 的实或复对称不变多项式，\(\Gamma=(\nabla_t)\) 为连接 \(\nabla_0,\nabla_1\) 的光滑联络路径。则
	\begin{equation}\label{g2:eq:chern-weil-transgression-differential}
		\mathrm dT_P(\Gamma)=P(F_1,\ldots,F_1)-P(F_0,\ldots,F_0).
	\end{equation}
	因此闭形式 \(P(F_\nabla,\ldots,F_\nabla)\) 的 de Rham 上同调类与联络无关。
\end{theorem}
\begin{proof}
	参数积分与外微分交换，故\cref{g2:prop:curvature-path-variation}给出
	\[
	\begin{aligned}
		\mathrm dT_P(\Gamma)
		&=\int_0^1 k\,\mathrm dP(\dot A_t,F_t,\ldots,F_t)\,\mathrm dt\\
		&=\int_0^1\partial_tP(F_t,\ldots,F_t)\,\mathrm dt\\
		&=P(F_1,\ldots,F_1)-P(F_0,\ldots,F_0).
	\end{aligned}
	\]
	最后一步是逐个形式系数使用微积分基本定理。任意两条联络可由仿射路径连接，右侧总为恰当形式；结合\cref{g2:thm:chern-weil-closed}的闭性，得到上同调类与联络无关。
\end{proof}

例如复线丛上取 \(P(X)=\mathrm iX/(2\pi)\)。由于 \(k=1\)，转移形式就是 \(\mathrm i(\nabla_1-\nabla_0)/(2\pi)\)，曲率之差是它的外微分。高次数中的额外曲率因子，使同一思路能比较所有陈--Weil 形式。

\ProseParagraphBreak

转移形式本身仍可能随连接路径改变。其可控程度需要一个二参数计算；仅知道两份转移形式的外微分相同，只能推出差闭，还不能推出差恰当。

\begin{proposition}\label{g2:prop:transgression-path-homotopy}
	设 \(Q\to M\) 为光滑主右矩阵群丛，\(P\) 为次数 \(k\ge1\) 的对称不变多项式。若两条光滑联络路径 \(\Gamma_0,\Gamma_1\) 具有相同起点与终点，则 \(T_P(\Gamma_1)-T_P(\Gamma_0)\) 为恰当形式。特别地，对三条联络 \(\nabla_0,\nabla_1,\nabla_2\)，
	\begin{equation}\label{g2:eq:transgression-triangle-mod-exact}
		T_P(\nabla_0,\nabla_2)
		\equiv T_P(\nabla_0,\nabla_1)+T_P(\nabla_1,\nabla_2)
		\pmod{\mathrm d\Omega^{2k-2}(M;\mathbb K)}.
	\end{equation}
	仿射转移还严格满足 \(T_P(\nabla_1,\nabla_0)=-T_P(\nabla_0,\nabla_1)\)。
\end{proposition}
\begin{proof}
	由于联络空间是仿射空间，可用
	\[
		\nabla_{s,t}=(1-s)\nabla_t^{\,0}+s\nabla_t^{\,1}
	\]
	连接两条路径，而且在 \(t=0,1\) 处不随 \(s\) 变化。记局部矩阵为 \(A_{s,t}\)，令 \(a=\partial_tA\)、\(b=\partial_sA\)、\(F=\mathrm d_MA+A\wedge A\)。在方形与 \(M\) 的乘积上取无参数分量的联络，其曲率为
	\[
		\widehat F=F+\mathrm ds\wedge b+\mathrm dt\wedge a.
	\]
	展开 \(P(\widehat F^k)\)，写成
	\[
		\omega+\mathrm ds\wedge U+\mathrm dt\wedge V
		+\mathrm ds\wedge\mathrm dt\wedge W,
	\]
	其中
	\[
	\begin{aligned}
		\omega&=P(F^k),&
		U&=kP(b,F^{k-1}),\\
		V&=kP(a,F^{k-1}),&
		W&=k(k-1)P(a,b,F^{k-2}).
	\end{aligned}
	\]
	\(k=1\) 时约定 \(W=0\)。这里 \(F^j\) 表示连续放入 \(j\) 个曲率因子。核对最后一项的符号：\((\mathrm ds\wedge b)\wedge(\mathrm dt\wedge a)=-\mathrm ds\wedge\mathrm dt\wedge b\wedge a\)，再交换两个一形式 \(a,b\) 又产生一个负号，故得到上述 \(W\)。

	总外微分为零。比较 \(\mathrm ds\wedge\mathrm dt\) 的系数，得到
	\[
		\partial_sV-\partial_tU+\mathrm d_MW=0.
	\]
	在方形上积分，并用固定端点条件 \(b(s,0)=b(s,1)=0\)，有
	\[
	\begin{aligned}
		T_P(\Gamma_1)-T_P(\Gamma_0)
		&=\int_0^1\int_0^1\partial_sV\,\mathrm dt\,\mathrm ds\\
		&=\int_0^1\bigl(U(s,1)-U(s,0)\bigr)\,\mathrm ds
		-\mathrm d_M\int_0^1\int_0^1W\,\mathrm dt\,\mathrm ds\\
		&=-\mathrm d_M\int_0^1\int_0^1W\,\mathrm dt\,\mathrm ds.
	\end{aligned}
	\]
	这给出了差的显式原始形式。

	为证明三角式，先沿 \(\nabla_0\) 到 \(\nabla_1\) 的路径走，再沿 \(\nabla_1\) 到 \(\nabla_2\) 走。作保向参数变换不会改变式\eqref{g2:eq:chern-weil-transgression}，因为 \(\dot A_t\,\mathrm dt\) 按一元换元公式变换。可将两段在端点附近重新参数化为常值，再光滑拼接；积分是两段积分之和。将这条路径与 \(\nabla_0\) 到 \(\nabla_2\) 的仿射路径比较，刚才的二参数计算给出式\eqref{g2:eq:transgression-triangle-mod-exact}。最后，反向走仿射路径使速度换号，换元后直接得到严格的反对称性。
\end{proof}

因此，两条固定联络之间有一个与路径无关的元素
\[
	[T_P(\nabla_0,\nabla_1)]
	\in\Omega^{2k-1}(M;\mathbb K)/\mathrm d\Omega^{2k-2}(M;\mathbb K).
\]
这是形式模去恰当形式的商；只有当两端的曲率代表相等时，转移形式才闭，才能进一步称为 de Rham 上同调类。

\subsection{自然性和陈–Weil 同态}
\label{g2:subsec:2202:03}
将不变多项式的次数加倍，恰好对应曲率形式的次数。记
\[
	I^*(G;\mathbb K)=\bigoplus_{k\ge0}I^k(G;\mathbb K)
\]
为对称不变多项式代数，其中 \(I^k\) 的元素是次数 \(k\) 的齐次多项式，并与其极化后的对称 \(k\) 线性式对应。乘积取普通多项式乘积；极化给出的组合系数已包含在乘积定义中。

\begin{theorem}[陈--Weil 同态]\label{g2:thm:chern-weil-homomorphism}
	设 \(G\) 为矩阵 Lie 群，\(Q\to M\) 为有限维、Hausdorff、第二可数光滑流形上的光滑主右 \(G\) 丛，\(\mathbb K=\mathbb R\) 或 \(\mathbb C\)。对任意光滑联络 \(\nabla\)，令
	\begin{equation}\label{g2:eq:chern-weil-homomorphism}
		\operatorname{cw}_Q(P)=[P(F_\nabla,\ldots,F_\nabla)].
	\end{equation}
	这定义了与联络无关的保单位代数同态
	\[
		\operatorname{cw}_Q:I^*(G;\mathbb K)\longrightarrow
		H_{\mathrm{dR}}^{\mathrm{even}}(M;\mathbb K),
		\qquad I^k\longrightarrow H_{\mathrm{dR}}^{2k}.
	\]
	对光滑映射 \(f:N\to M\)，它满足
	\begin{equation}\label{g2:eq:chern-weil-pullback}
		\operatorname{cw}_{f^*Q}(P)=f^*\operatorname{cw}_Q(P).
	\end{equation}
	对结构群同态 \(\rho:H\to G\) 与主 \(H\) 丛 \(R\to M\)，若 \(Q=R\times_\rho G\)，则
	\[
		\operatorname{cw}_{Q}(P)=\operatorname{cw}_{R}(\rho^*P),
		\qquad
		(\rho^*P)(X_1,\ldots,X_k)=P(\rho_*X_1,\ldots,\rho_*X_k).
	\]
\end{theorem}
\begin{proof}
	光滑主丛存在联络：局部平凡丛先取零联络，用从属光滑单位分解在总空间上加权；权重之和为一保持竖直再现条件，底空间上的权重保持等变性。这是第一卷第18.2节“主联络的存在”定理的构造。闭性由\cref{g2:thm:chern-weil-closed}保证，联络无关性由\cref{g2:thm:chern-weil-transgression}保证，故式\eqref{g2:eq:chern-weil-homomorphism}有定义。

	固定一条联络。对多项式相加、乘标量，代入曲率的运算相应相加、乘标量。对次数分别为 \(k,\ell\) 的 \(P_1,P_2\)，在极化后的乘积中，每项都把 \(k+\ell\) 个二次曲率因子分给两组。因这些因子均为偶次，重排不产生符号；极化的平均系数又使全部相同项的系数和为一。因此
	\[
		(P_1P_2)(F^{k+\ell})=P_1(F^k)\wedge P_2(F^\ell).
	\]
	常数多项式 \(1\) 给出常数零形式 \(1\)，从而得到保单位代数同态。

	拉回联络在局部为 \(\mathrm d+f^*A\)。因为拉回与外微分、外乘交换，其曲率是 \(f^*F\)，所以
	\[
		P(F_{f^*\nabla}^k)=f^*P(F_\nabla^k).
	\]
	在源流形上若另选一条联络，转移公式使所得上同调类不变，于是得到自然性。

	最后，沿 \(\rho\) 扩张结构群后的局部联络矩阵为 \(\rho_*A\)。Lie 代数同态保持括号，故相应曲率为 \(\rho_*F\)。将它代入 \(P\) 得到 \((\rho^*P)(F^k)\)，证明最后一式。
\end{proof}

自然性也在转移形式的层次成立。若拉回整条路径，则
\begin{equation}\label{g2:eq:transgression-pullback}
	T_P(f^*\Gamma)=f^*T_P(\Gamma),
\end{equation}
因为速度、曲率和参数积分都与拉回交换。这是一条形式等式，比上同调类自然性更精细。下一节将选取确定的行列式及 Pfaffian 多项式，把式\eqref{g2:eq:chern-weil-homomorphism}与上一章的拓扑类逐一比较。

\subsection{规范变换下的变化}
\label{g2:subsec:2202:04}
规范变换是覆盖底空间恒等映射的主丛自同构。沿用第一卷第18.4节“联络的规范变换”命题，在固定局部截面中写成
\[
	A^g=g^{-1}Ag+g^{-1}\mathrm dg,\qquad F^g=g^{-1}Fg.
\]
这里的局部群值函数满足伴随换架规律，因而表示同一个整体规范变换；平凡丛中它就是一张整体映射 \(g:M\to G\)。不变多项式立即给出 \(P((F^g)^k)=P(F^k)\)，是代表形式的严格相等。

\begin{proposition}\label{g2:prop:transgression-gauge-naturality}
	设 \(Q\to M\) 为光滑主右矩阵群丛，\(P\) 为次数 \(k\ge1\) 的对称不变多项式，\(g\) 为固定规范变换。对任意联络路径 \(\Gamma\)，同时变换路径上的所有联络，有
	\[
		T_P(\Gamma^g)=T_P(\Gamma).
	\]
	因此仿射转移满足
	\begin{equation}\label{g2:eq:transgression-gauge-invariant-pair}
		T_P(\nabla_0^g,\nabla_1^g)=T_P(\nabla_0,\nabla_1).
	\end{equation}
	只比较一条联络及其变换时，\(T_P(\nabla,\nabla^g)\) 是闭形式，其上同调类与 \(\nabla\) 无关。
\end{proposition}
\begin{proof}
	\(g\) 不随路径参数改变，所以 \(\partial_tA_t^g=g^{-1}\dot A_tg\)，而 \(F_t^g=g^{-1}F_tg\)。将这些式子代入 \(P\)，共轭由不变性消去，被积形式逐项不变。规范变换还把仿射路径送到变换后两端之间的仿射路径，因为附加项 \(g^{-1}\mathrm dg\) 的系数为一，故得到式\eqref{g2:eq:transgression-gauge-invariant-pair}。

	转移公式给出
	\[
		\mathrm dT_P(\nabla,\nabla^g)=P((F^g)^k)-P(F^k)=0.
	\]
	再取另一联络 \(\nabla'\)。用\cref{g2:prop:transgression-path-homotopy}沿四个端点依次比较，模去恰当形式有
	\[
	\begin{aligned}
		T_P(\nabla',(\nabla')^g)
		&\equiv T_P(\nabla',\nabla)+T_P(\nabla,\nabla^g)
		+T_P(\nabla^g,(\nabla')^g)\\
		&=-T_P(\nabla,\nabla')+T_P(\nabla,\nabla^g)
		+T_P(\nabla,\nabla')\\
		&=T_P(\nabla,\nabla^g).
	\end{aligned}
	\]
	两个外侧项由同时规范不变性及反对称性抵消，所以闭转移类与所选联络无关。
\end{proof}

\begin{example}\label{g2:exm:circle-gauge-transgression}
	取 \(S^1=\mathbb R/\mathbb Z\) 上的平凡酉线丛，\(\nabla=\mathrm d\)，并令 \(g(x)=e^{2\pi\mathrm i m x}\)，其中 \(m\in\mathbb Z\)。角度 \(x\) 本身不是圆周上的实函数，但 \(\mathrm dx\) 是整体一形式，且正向积分为 \(1\)。于是
	\[
		A^g=g^{-1}\mathrm dg=2\pi\mathrm i m\,\mathrm dx.
	\]
	对 \(P(X)=\mathrm iX/(2\pi)\)，有
	\[
		T_P(\nabla,\nabla^g)=\frac{\mathrm i}{2\pi}(A^g-0)
		=-m\,\mathrm dx,\qquad
		\int_{S^1}T_P(\nabla,\nabla^g)=-m.
	\]
	\(m\ne0\) 时，这个闭形式不是恰当形式，因为闭曲线上恰当一形式的积分为零。两端曲率都为零，但规范变换的绕数仍留在转移形式中。
\end{example}

因此，曲率代表的规范不变性与转移形式的选择依赖各自有准确的含义：同时改变一对联络，转移形式保持相等；固定一端而用规范变换改变另一端，可以留下非零闭类。后面的三维例子将把同样的现象写成整数改变的积分，从而得到模整数的不变量。
% 纲要条目排为编号小节标题，并进入目录与书签；说明段仍保留于纲要注释中。

% !TeX root = ../../Book2.tex
% 纲要标题：### 22.3 拓扑类的比较
% 纲要位置：几何学计划纲要.md，第 1815 行。
% 正文状态：未撰写；本文件目前是框架。

% 纲要内容开始
% 
% * 线丛上的归一化
% * 分裂原理及一般 Chern 类
% * Pontryagin 形式与 Euler 的 Pfaffian
% * 曲率形式看不到的扭信息
% 
% 纲要内容结束

\section{拓扑类的比较}
\label{g2:sec:2203}

待撰写

\subsection{线丛上的归一化}
\label{g2:subsec:2203:01}
待撰写

\subsection{分裂原理及一般 Chern 类}
\label{g2:subsec:2203:02}
待撰写

\subsection{Pontryagin 形式与 Euler 的 Pfaffian}
\label{g2:subsec:2203:03}
待撰写

\subsection{曲率形式看不到的扭信息}
\label{g2:subsec:2203:04}
待撰写
% 纲要条目排为编号小节标题，并进入目录与书签；说明段仍保留于纲要注释中。

% !TeX root = ../../Book2.tex
% 纲要标题：### 22.4 陈–Simons 与几何应用
% 纲要位置：几何学计划纲要.md，第 1872 行。
% 正文状态：初稿；尚未正式审核。

% 纲要内容开始
% 
% * 陈–Simons 形式
% * 陈–Simons 积分与规范变换
% * 边界、延拓和相对不变量
% * 陈–Simons 变分与陈特征
% 
% ---
% 
% 纲要内容结束

\section{陈–Simons 与几何应用}
\label{g2:sec:2204}

两个联络给出同一个示性上同调类，二者的曲率形式却未必相等。\secref{g2:sec:2202}的转移公式把这个差写成一个外微分；现在保留那个次数低一的形式，便得到能随联络变化的几何量。在三维，二次不变多项式的四次曲率形式因维数消失，三次转移形式仍可以积分。问题随之变成：这个积分依赖哪些选择，改变标架时为何会出现整数？

\subsection{陈–Simons 形式}
\label{g2:subsec:2204:01}
设 \(P\) 是次数为 \(k\ge1\) 的对称不变多项式，\(\nabla_0,\nabla_1\) 是同一张光滑丛上的两个联络。沿直线路径 \(\nabla_t=\nabla_0+t a\)，其中 \(a=\nabla_1-\nabla_0\)，定义
\begin{equation}\label{g2:eq:relative-chern-simons-form}
	\operatorname{CS}_P(\nabla_0,\nabla_1)
	\coloneqq k\int_0^1P(a,F_t,\ldots,F_t)\,\mathrm dt.
\end{equation}
这就是\cref{g2:thm:chern-weil-transgression}中的转移形式，称为相对于 \(\nabla_0\) 的\term[chen-Simons-xingshi]{陈--Simons 形式}[Chern--Simons form]\footnote{James Harris Simons（1938--2024），美国数学家，研究微分几何，并与陈省身共同发展次级示性不变量。陈省身的简介见第21章。}。它是底空间上的整体 \((2k-1)\) 次形式，并满足
\[
	\mathrm d\operatorname{CS}_P(\nabla_0,\nabla_1)=P(F_1^k)-P(F_0^k).
\]
若右边为零，则它给出 de Rham 上同调类。这个条件可能来自两个曲率代表相等，也可能来自底空间维数小于 \(2k\)。

\ProseParagraphBreak

这里的参照联络是定义的一部分。\cref{g2:prop:transgression-path-homotopy}说明，保持两端点而更换联络路径只改变一个恰当形式；对三个联络，还有
\begin{equation}\label{g2:eq:chern-simons-reference-change}
	\operatorname{CS}_P(\nabla_0,\nabla_2)
	\equiv\operatorname{CS}_P(\nabla_0,\nabla_1)+\operatorname{CS}_P(\nabla_1,\nabla_2)
	\pmod{\mathrm d\Omega^{2k-2}}.
\end{equation}
因此换参照联络时，变化由另一份转移形式决定。若原丛已有一个全局标架，可以取该标架下的平凡联络 \(\mathrm d\) 为参照；换一个随点变化的标架，也就换了这个参照联络。这正是下面规范变换公式的来源。次级形式的原始构造及选择依赖可参阅\cite[第3节]{ChernSimons1974}。

\begin{proposition}\label{g2:prop:su-two-chern-simons-form}
	设 \(A\in\Omega^1(M;\mathfrak{su}(2))\) 为光滑流形 \(M\) 上平凡 \(\operatorname{SU}(2)\) 丛在一个固定标架中的联络矩阵，\(F=\mathrm dA+A\wedge A\)。定义
	\begin{equation}\label{g2:eq:su-two-chern-simons-density}
		\operatorname{CS}(A)=\frac1{8\pi^2}\operatorname{tr}\left(A\wedge\mathrm dA+\frac23A\wedge A\wedge A\right).
	\end{equation}
	则 \(\operatorname{CS}(A)=\operatorname{CS}_{c_2}(\mathrm d,\mathrm d+A)\)，且
	\[
		\mathrm d\operatorname{CS}(A)=\frac1{8\pi^2}\operatorname{tr}(F\wedge F)=c_2(\nabla).
	\]
\end{proposition}
\begin{proof}
	\(\mathfrak{su}(2)\) 由无迹反 Hermitian 矩阵组成，故 \(\operatorname{tr}F=0\)。本章采用 \(c(\nabla)=\det(I+\mathrm iF/(2\pi))\)，其二次项为
	\[
		c_2(\nabla)=\frac1{8\pi^2}\left(\operatorname{tr}(F^2)-(\operatorname{tr}F)^2\right)=\frac1{8\pi^2}\operatorname{tr}(F^2).
	\]
	相应极化为 \(P(X,Y)=\operatorname{tr}(XY)/(8\pi^2)\)。从零联络到 \(A\) 的路径是 \(A_t=tA\)，曲率为 \(F_t=t\mathrm dA+t^2A^2\)。代入式\eqref{g2:eq:relative-chern-simons-form}，得到
	\[
		2\int_0^1\frac1{8\pi^2}\operatorname{tr}\bigl(A(t\mathrm dA+t^2A^2)\bigr)\,\mathrm dt
		=\frac1{8\pi^2}\operatorname{tr}\left(A\mathrm dA+\frac23A^3\right).
	\]
	转移恒等式给出它的外微分。这里的乘法均含外乘；迹的分次循环律保证该计算与矩阵的非交换乘法相容。
\end{proof}

若 \(\dim M=3\)，四形式 \(c_2(\nabla)\) 因维数恒为零，因而 \(\operatorname{CS}(A)\) 闭。其积分却可以随 \(A\) 连续变化：四次拓扑类在这个维数不提供数值，三次次级形式还记得联络。

\subsection{陈–Simons 积分与规范变换}
\label{g2:subsec:2204:02}
现在固定一个连通闭定向光滑三维流形 \(M\)。在平凡 \(\operatorname{SU}(2)\) 丛的一张标架中，令
\[
	S(A)=\int_M\operatorname{CS}(A)\in\mathbb R.
\]
为了解它在换标架后的变化，先在群本身上计算一个规范化三次形式。

\begin{lemma}\label{g2:lem:su-two-integral-three-form}
	将 \(\operatorname{SU}(2)\) 与单位三维球面识别为
	\[
		q=(x_0,x_1,x_2,x_3)\longmapsto
		\begin{pmatrix}x_0+\mathrm i x_1&x_2+\mathrm i x_3\\-x_2+\mathrm i x_3&x_0-\mathrm i x_1\end{pmatrix},
		\qquad\sum_{j=0}^3x_j^2=1.
	\]
	给 \(S^3\) 使用 \(\mathbb R^4\) 中单位球的外法向优先边界定向，记由左平移定义的 Lie 代数值一形式为 \(\vartheta=q^{-1}\mathrm dq\)。则
	\begin{equation}\label{g2:eq:su-two-normalized-volume}
		\omega=-\frac1{24\pi^2}\operatorname{tr}(\vartheta^3),\qquad \int_{S^3}\omega=1.
	\end{equation}
	因此对任意光滑映射 \(g:M\to\operatorname{SU}(2)\)，有 \(\int_Mg^*\omega=\deg g\in\mathbb Z\)。
\end{lemma}
\begin{proof}
	在单位元处，与坐标 \(x_1,x_2,x_3\) 对应的切向量为
	\[
		I=\begin{pmatrix}\mathrm i&0\\0&-\mathrm i\end{pmatrix},\qquad
		J=\begin{pmatrix}0&1\\-1&0\end{pmatrix},\qquad
		K=\begin{pmatrix}0&\mathrm i\\\mathrm i&0\end{pmatrix}.
	\]
	它们满足 \(I^2=J^2=K^2=-1\)、\(IJ=K\)、\(JK=I\)、\(KI=J\)，逆序乘积变号。三次外乘在 \((I,J,K)\) 上取值是六个交错矩阵乘积之和；每个交错乘积都等于 \(-1\)，所以
	\[
		\operatorname{tr}(\vartheta^3)(I,J,K)=6\operatorname{tr}(-1)=-12.
	\]
	在点 \((1,0,0,0)\)，这组切基为球面的正向标准正交基。左乘单位四元数保持 \(\mathbb R^4\) 的 Euclid 范数，因而保持球面圆度量；它随单位四元数连续变化，且恒等左乘保向，所以全部左乘保向。\(\vartheta\) 与圆体积形式均左不变，刚才在单位元的计算遂给出
	\[
		-\frac1{24\pi^2}\operatorname{tr}(\vartheta^3)=\frac1{2\pi^2}\,\mathrm{vol}_{S^3}.
	\]
	为核对常数，球坐标 \((\cos\chi,\sin\chi\,u)\)，\(u\in S^2\)、\(0\le\chi\le\pi\)，给出体积
	\[
		\operatorname{vol}(S^3)=\operatorname{area}(S^2)\int_0^\pi\sin^2\chi\,\mathrm d\chi=4\pi\cdot\frac\pi2=2\pi^2.
	\]
	故 \(\omega\) 的积分为一。最后使用\cref{g2:thm:degree-comparison}：该定理在闭连通同维定向流形之间将归一化顶次形式的拉回积分识别为基本类所定义的整数映射度。这里两端恰为三维，故所有条件满足。
\end{proof}

\begin{theorem}\label{g2:thm:chern-simons-gauge-integer}
	设 \(M\) 为连通闭定向光滑三维流形，\(A\) 为平凡 \(\operatorname{SU}(2)\) 丛的一张标架下的光滑联络矩阵。对任意光滑规范变换 \(g:M\to\operatorname{SU}(2)\)，令
	\[
		A^g=g^{-1}Ag+g^{-1}\mathrm dg,\qquad \theta=g^{-1}\mathrm dg,\qquad B=g^{-1}Ag.
	\]
	则有形式恒等式
	\begin{equation}\label{g2:eq:chern-simons-gauge-formula}
		\operatorname{CS}(A^g)-\operatorname{CS}(A)
		=-\frac1{8\pi^2}\mathrm d\operatorname{tr}(\theta\wedge B)+g^*\omega,
	\end{equation}
	并且
	\begin{equation}\label{g2:eq:chern-simons-integer-shift}
		S(A^g)-S(A)=\deg g.
	\end{equation}
	于是 \(S(A)\bmod\mathbb Z\in\mathbb R/\mathbb Z\) 与全局标架无关，并在规范变换下不变。
\end{theorem}
\begin{proof}
	记 \(R=g^{-1}(\mathrm dA)g\)。微分乘积法则和 \(\mathrm d(g^{-1})=-g^{-1}(\mathrm dg)g^{-1}\) 给出
	\[
		\mathrm d\theta=-\theta^2,\qquad \mathrm dB=R-\theta B-B\theta.
	\]
	为把所有符号看清，先省略公因子 \(1/(8\pi^2)\)。迹中三个一次形式作循环置换时符号为正，因此展开 \(B+\theta\) 的三次项后，三个含两个 \(B\) 的项相同，三个含两个 \(\theta\) 的项也相同。于是
	\[
	\begin{aligned}
		&\operatorname{tr}\left((B+\theta)\mathrm d(B+\theta)+\frac23(B+\theta)^3\right)
		-\operatorname{tr}\left(A\mathrm dA+\frac23A^3\right)\\
		&\qquad=\operatorname{tr}(\theta R-B\theta^2)-\frac13\operatorname{tr}(\theta^3).
	\end{aligned}
	\]
	另一方面，
	\[
		\mathrm d\operatorname{tr}(\theta B)
		=\operatorname{tr}(-\theta^2B-\theta(R-\theta B-B\theta))
		=-\operatorname{tr}(\theta R)+\operatorname{tr}(B\theta^2).
	\]
	代回即得式\eqref{g2:eq:chern-simons-gauge-formula}，因为 \(-\operatorname{tr}(\theta^3)/(24\pi^2)=g^*\omega\)。在闭流形上，恰当三次形式的积分由 Stokes 定理为零；\cref{g2:lem:su-two-integral-three-form}将剩下的一项化为 \(\deg g\)。因此变换前后的实数相差整数，模整数的值不变。
\end{proof}

取 \(M=S^3=\operatorname{SU}(2)\)、\(A=0\) 和 \(g(q)=q\)，便得到
\[
	A^g=\vartheta,\qquad F_{A^g}=0,\qquad S(0)=0,\qquad S(\vartheta)=1.
\]
两个平坦联络在规范变换下等价，但在各自矩阵表达中得到的实数不同。整数模糊由换标架的同伦类产生；把数值取到 \(\mathbb R/\mathbb Z\) 后，这两份表达才一致。

\subsection{边界、延拓和相对不变量}
\label{g2:subsec:2204:03}
有边界时，式\eqref{g2:eq:chern-simons-gauge-formula}中的恰当项会留下边界积分。另一方面，在四维内部积分 \(c_2(\nabla)\)，也自然要求在三维边界记录一个陈--Simons 修正。下面把这两件事组织为相对不变量。所有边界都取外法向优先定向，所用 Stokes 定理是第一卷第14.4节的版本：对紧有向光滑流形及光滑至边界的形式，\(\int_W\mathrm d\eta=\int_{\partial W}\eta\)。

\begin{proposition}\label{g2:prop:relative-chern-simons-integer}
	设 \(W\) 为紧定向光滑四维流形，\(M=\partial W\)，\(E\to W\) 为具有 \(\operatorname{SU}(2)\) 结构的复秩二光滑丛，并给边界丛选一张 \(\operatorname{SU}(2)\) 标架 \(s\)。对任意相容光滑联络 \(\nabla\)，记其边界矩阵为 \(A_s\)。则
	\begin{equation}\label{g2:eq:relative-second-chern-number}
		N(W,E,s)=\int_W c_2(\nabla)-\int_M\operatorname{CS}(A_s)
	\end{equation}
	是一个只依赖 \((W,E,s)\) 的整数。若 \(s'=sg\)，其中 \(g:M\to\operatorname{SU}(2)\)，则
	\[
		N(W,E,s')=N(W,E,s)-\sum_{M_\alpha\subset M}\deg(g|_{M_\alpha}).
	\]
\end{proposition}
\begin{proof}
	\cref{g2:thm:collar-neighborhood}在紧光滑边界附近给出乘积领圈。先把边界标架延到这条领圈：沿领圈的法向区间作平行移动，便把边界上的丛识别为领圈上的乘积。于是 \(s\) 在领圈上确定平凡联络。联络构成仿射空间，用只依赖领参数的光滑截断将给定联络与这个平凡联络插值，可得到一个联络 \(\nabla^0\)，使它在较小领圈上恰为 \(\mathrm d\)。

	对 \(\nabla^0\) 与 \(\nabla\) 使用转移公式，再用 Stokes 定理，得到
	\[
	\begin{aligned}
		\int_Wc_2(\nabla)-\int_Wc_2(\nabla^0)
		&=\int_M\operatorname{CS}_{c_2}(\nabla^0|_M,\nabla|_M)\\
		&=\int_M\operatorname{CS}(A_s).
	\end{aligned}
	\]
	因此式\eqref{g2:eq:relative-second-chern-number}等于 \(\int_Wc_2(\nabla^0)\)。将 \(W\) 与反向副本 \(-W\) 沿边界粘合为闭定向四维流形 \(X\)。在第二份 \(-W\) 上取平凡复秩二丛，并用 \(s\) 与第一份 \(E\) 粘合，得到 \(E_X\)。两边联络在领圈标架中都为零，故光滑粘合为 \(\nabla_X\)。于是
	\[
		\int_Wc_2(\nabla^0)=\int_Xc_2(\nabla_X)
		=\langle c_2(E_X),[X]\rangle\in\mathbb Z.
	\]
	最后一个等式是\cref{g2:thm:chern-weil-chern-comparison}，其中整数陈类先映到实上同调，再由积分与基本类配对。闭合后的丛 \(E_X\) 由边界粘合数据确定；改变领圈识别只给出同构的丛，因为两种向内延拓可沿领参数收缩回同一边界标架。因此右侧与联络及辅助选择无关。

	改用 \(s'=sg\) 时，内部曲率积分不变，边界矩阵由 \(A_s\) 变为 \(A_s^g\)。对每个边界分支应用式\eqref{g2:eq:chern-simons-integer-shift}，便得到所列变化公式。边界方向已进入各分支映射度的定义。
\end{proof}

特别地，如果边界上的联络有一份四维延拓，那么
\begin{equation}\label{g2:eq:chern-simons-extension-mod-integer}
	S(A_s)\equiv\int_Wc_2(\nabla)\pmod{\mathbb Z}.
\end{equation}
两份延拓的内部积分相差整数。也可把两份填充沿共同边界粘合，第二份反向，差便等于闭四维流形上的陈数。若需先使联络在领圈成为边界联络的乘积，可用保持边界值的插值；转移形式在边界的限制为零，所以此调整不改变四维积分。这个论证说明，整性来自闭合后的整数拓扑类，而边界联络仍决定模整数的连续数值。

\ProseParagraphBreak

相对上同调也直接记录同一个修正。对包含 \(j:M\hookrightarrow W\)，采用复形
\[
	\mathcal C^r=\Omega^r(W)\oplus\Omega^{r-1}(M),\qquad
	\mathrm D(\alpha,\beta)=(\mathrm d\alpha,j^*\alpha-\mathrm d\beta).
\]
有 \(\mathrm D^2=0\)。这个复形与边界限制为零的相对形式复形计算同一上同调：边界形式可沿领圈延拓并用截断补零，因而限制映射 \(\Omega^*(W)\to\Omega^*(M)\) 逐次满射；由核的短正合列与上述复形的长正合列比较，恒等映射在外侧给出同构，中间也为同构。结合\cref{g2:thm:de-rham-support}，它就是 \(H^*(W,M;\mathbb R)\) 的一个模型。

\ProseParagraphBreak

成对形式 \((c_2(\nabla),\operatorname{CS}(A_s))\) 是闭元，因为边界上 \(\mathrm d\operatorname{CS}(A_s)=j^*c_2(\nabla)\)。其相对求值为
\[
	\langle(\alpha,\beta),[W,M]\rangle=\int_W\alpha-\int_M\beta.
\]
对一个余边界 \(\mathrm D(\eta,\xi)\)，右端由两次 Stokes 公式为零，所以只依赖相对类。式\eqref{g2:eq:relative-second-chern-number}正是这个求值；边界的平凡化说明了为什么这里需要相对数据。

\subsection{陈–Simons 变分与陈特征}
\label{g2:subsec:2204:04}
三维陈--Simons 积分首先给出一个具体的变分问题。将所有 \(\mathfrak{su}(2)\) 值一次形式看成平凡丛的联络空间，允许它们作任意光滑变化。

\begin{proposition}\label{g2:prop:chern-simons-variation}
	设 \(M\) 为闭定向光滑三维流形，\(A\) 为平凡 \(\operatorname{SU}(2)\) 丛的联络矩阵，\(a\in\Omega^1(M;\mathfrak{su}(2))\)。则
	\begin{equation}\label{g2:eq:chern-simons-variation}
		\left.\frac{\mathrm d}{\mathrm dt}\right|_{t=0}S(A+ta)
		=\frac1{4\pi^2}\int_M\operatorname{tr}(a\wedge F_A).
	\end{equation}
	因此 \(S\) 的临界联络恰为平坦联络。
\end{proposition}
\begin{proof}
	对式\eqref{g2:eq:su-two-chern-simons-density}求导，用迹的分次循环律，得到
	\[
	\begin{aligned}
		\delta\operatorname{tr}\left(A\mathrm dA+\frac23A^3\right)
		&=\operatorname{tr}(a\mathrm dA+A\mathrm da+2aA^2)\\
		&=2\operatorname{tr}(aF_A)-\mathrm d\operatorname{tr}(Aa).
	\end{aligned}
	\]
	闭流形上的最后一项积分为零，故得到变分公式。若 \(F_A=0\)，所有变分都为零。反过来，给 \(M\) 临时选一个 Riemann 度量，并对 \(\mathfrak{su}(2)\) 取正定内积 \(\langle X,Y\rangle=-\operatorname{tr}(XY)\)。对实二次形式定义一次形式 \(*\gamma\)，使 \(\beta\wedge *\gamma=\langle\beta,\gamma\rangle\mathrm{vol}_M\) 对所有二次形式 \(\beta\) 成立；在正向正交余标架中，它将 \(e^1\wedge e^2\) 送到 \(e^3\)，其余两项循环处理，故逐点存在且唯一。对 Lie 代数值形式逐个系数应用它。若处处驻定，可取 \(a=*F_A\)。因为一次与二次形式交换时符号为正，
	\[
		0=\int_M\operatorname{tr}((*F_A)\wedge F_A)=-\int_M|F_A|^2\,\mathrm{vol}_M.
	\]
	非负连续函数的积分为零迫使 \(F_A=0\)。辅助度量只用于检验变分泛函非退化，原来的作用量不含这个选择。
\end{proof}

实数作用量在规范变换下可增加整数，因此对整数 \(\ell\)，相位
\[
	\exp\bigl(2\pi\mathrm i\ell S(A)\bigr)
\]
在规范变换下不变。整数条件可以直接在 \(S^3\) 上检验：恒等映射产生一次单位平移，要求相位不变就有 \(e^{2\pi\mathrm i\ell}=1\)。如果流形有边界，变分中还保留 \(-\int_{\partial M}\operatorname{tr}(Aa)/(8\pi^2)\)；固定边界联络使 \(a\) 的切向限制为零，从而消去这一项。边界条件因此是变分问题实际数据的一部分。

\ProseParagraphBreak

指标理论会使用另一种方便相加的曲率组合。对复向量丛的联络定义有限次数意义下的形式
\begin{equation}\label{g2:eq:chern-character-curvature-form}
	\operatorname{ch}(\nabla)=\operatorname{tr}\exp\left(\frac{\mathrm iF}{2\pi}\right)
	=\operatorname{rank}E+\sum_{j\ge1}\frac1{j!}\left(\frac{\mathrm i}{2\pi}\right)^j\operatorname{tr}(F^j).
\end{equation}
在给定有限维底流形上，超过维数的形式自动为零，所以这里只需有限和。每一项来自不变多项式，因而闭，且其上同调类与联络无关。它称为\term[chen-tezheng-xingshi]{陈特征形式}[Chern character form]。

\ProseParagraphBreak

直和曲率是分块对角矩阵，所以 \(\operatorname{ch}(\nabla^E\oplus\nabla^F)=\operatorname{ch}(\nabla^E)+\operatorname{ch}(\nabla^F)\)。张量丛曲率为 \(F_E\otimes I+I\otimes F_F\)，两个偶次矩阵形式可交换，指数因而分解，迹也分解，得到张量乘法公式。在分裂空间上，若陈根为 \(x_j\)，则其上同调类为 \(\sum_j e^{x_j}\)。分裂拉回单射使这个表达返回原空间；前三项为
\[
	\operatorname{ch}(E)=\operatorname{rank}E+c_1(E)+\frac12\bigl(c_1(E)^2-2c_2(E)\bigr)+\cdots
\]
的实或复系数像。特别在刚才的 \(\operatorname{SU}(2)\) 模型中，四次陈特征分量等于 \(-c_2\)，其转移形式就是 \(-\operatorname{CS}(A)\)。

\ProseParagraphBreak

第24章将把这种可加表达放进拓扑 K 理论。另一方面，解析指标是适当 Fredholm 算子的 \(\dim\ker D-\dim\operatorname{coker}D\)；把它与示性形式积分联系起来，需要构造算子、证明 Fredholm 性并建立相应指标定理，第三卷第17章将给出所选版本及其证明范围。本章已经提供其中的曲率表达、拓扑归一化和改变联络时的边界项。对 Euler 形式，\cref{g2:thm:chern-weil-euler-comparison}与上一章的切丛 Euler 求值直接给出闭定向偶维流形上的 Euler 示性数积分；一般椭圆算子的指标则还需要相应的解析输入。
% 纲要条目排为编号小节标题，并进入目录与书签；说明段仍保留于纲要注释中。


\begin{chaptersummary}
待撰写
\end{chaptersummary}

\BookExercises
待撰写


% !TeX root = ../Book2.tex
% 纲要标题：# 第五编　配边、稳定化与 K 理论
% 纲要位置：几何学计划纲要.md，第 1831 行。
% 正文状态：未撰写；本文件目前是框架。

% 纲要内容开始
% 
% 纲要内容结束

\BookPart{配边、稳定化与 K 理论}{g2:part:05}
\BookPartGuide
待撰写

% !TeX root = ../../Book2.tex
% 纲要标题：## 第23章*　配边、Pontryagin–Thom 与光滑结构
% 纲要位置：几何学计划纲要.md，第 1883 行。
% 正文状态：初稿；尚未正式审核。

% 纲要内容开始
% 
% > 本章目标与证明范围：以完整的带标架 Pontryagin–Thom 对应作为代表深结果，而不只画坍缩示意图。一般手术分类和 h-配边的完整证明仍外引，维数和范畴条件逐条声明。
% 
% 纲要内容结束

\OptionalChapter{配边、Pontryagin–Thom 与光滑结构}{g2:ch:23}
\BookChapterNumber{g2:ch:23}

\begin{chapterabstract}
	闭流形可以作为高一维紧流形的边界出现，配边由此把流形组成可计算的群。本章先证明 Stiefel--Whitney 数与 Pontryagin 数的配边不变性，再完整建立带标架 Pontryagin--Thom 对应：管状邻域坍缩与正则逆像互逆，增加环境坐标对应球面映射的悬挂。随后把手术解释为配边中的把手附着，说明基本群、中维交叉与 Whitney 技巧怎样进入分类。最后比较光滑、PL 与拓扑结构，用异光滑球面和四维 Euclid 空间的精确定理展示范畴差异。
\end{chapterabstract}

\begin{learninggoals}
	\begin{itemize}
		\item 区分无定向、定向及带稳定标架的配边，利用领圈构造群运算并核对边界符号。
		\item 用相对基本类证明示性数的配边不变性，并用射影空间检验流形能否成为边界。
		\item 完成坍缩映射与正则逆像的双向构造，分别处理管状邻域、标架及环境嵌入的选择。
		\item 把增加法方向与悬挂对应起来，证明 \(\Omega_n^{\mathrm{fr}}\cong\pi_n^{\mathrm S}\)。
		\item 读出手术迹的把手指标及其基本群效应，识别高维配边定理与四维光滑问题的适用条件。
	\end{itemize}
\end{learninggoals}

在球面上把一条纬线逐渐推到北极，纬线会缩成一点，随后消失。若要求每个时刻都是同一张圆周的嵌入，这个过程在消失时失败；若把所有时刻放在一起，轨迹却是一张光滑圆盘。这提示我们改变比较对象的方法：允许中间截面的拓扑发生变化，要求全部轨迹组成一张高一维的光滑流形。配边正是这一想法的整体形式。

\ProseParagraphBreak

前两章已经把切丛变成上同调类，再通过基本类求值得到整数或模二数。为什么某些示性数能阻止一个流形成为边界？在配边中，边界的稳定切丛来自内部，而整个边界的基本类在内部同调于零。这两条事实便给出不变性的证明。再多保留一组稳定标架，流形的信息还能转换成球面之间的映射；第9章引入的稳定同伦群于是获得了直接的几何解释。

\ProseParagraphBreak

本章的主要证明是带标架 Pontryagin--Thom 对应。所用横截性、管状邻域、领圈和固有淹没分别回到第14、15章；稳定化使用第9章第5节的悬挂定理。后两节从第17章的把手附着转入手术及分类问题。\(h\)-配边、手术障碍、结构提升和异光滑结构的深层结论将在陈述处明确引用正式文献，维数、边界及范畴均写入假设，供读者在具体例子中判断能否使用。

% !TeX root = ../../Book2.tex
% 纲要标题：### 23.1 配边及其几何构造
% 纲要位置：几何学计划纲要.md，第 1887 行。
% 正文状态：初稿；尚未正式审核。

% 纲要内容开始
% 
% * 定向、无定向和带标架配边
% * 粘合、逆元与群结构
% * 示性数的配边不变性
% 
% 纲要内容结束

\section{配边及其几何构造}
\label{g2:sec:2301}

球面是球的边界，环面是实心环面的边界。这两张曲面的拓扑不同，却都能由一张高一维的紧流形填满。若把两个闭流形一起看成某张紧流形的全部边界，就得到一种允许拓扑改变的比较关系。需要同时决定的是：填充过程中应保留哪些几何资料？只要求光滑结构、再要求方向、进一步指定切丛的稳定标架，会得到不同的配边关系。

\subsection{定向、无定向和带标架配边}
\label{g2:subsec:2301:01}
本节的闭流形均指紧且无边界的光滑流形，可以不连通。配边自身允许有边界，但也要求紧；例如用无限延伸的圆柱填充一个端面，不能作为这里的配边。

\begin{definition}\label{g2:def:unoriented-oriented-bordism}
	设 \(M_0,M_1\) 为闭光滑 \(n\) 维流形。若存在紧光滑 \(n+1\) 维流形 \(W\)，使
	\[
		\partial W\cong M_0\amalg M_1,
	\]
	就称 \(M_0\) 与 \(M_1\)\term[peibian]{配边}[bordant]；此时没有指定方向，亦称无定向配边。若 \(M_0,M_1,W\) 均定向，并且上述识别保向地写成
	\[
		\partial W\cong(-M_0)\amalg M_1,
	\]
	就称二者定向配边。这里 \(-M_0\) 表示反转方向，边界方向按外法向优先约定确定。把 \(M_0\) 称为入边界、\(M_1\) 称为出边界。若 \(M\) 与空流形配边，就称 \(M\) 为相应意义下的边界。
\end{definition}

外法向优先的具体含义是：若 \(\nu\) 指向 \(W\) 外部，则 \((v_1,\ldots,v_n)\) 在边界上为正向，恰好当 \((\nu,v_1,\ldots,v_n)\) 在 \(W\) 上为正向。区间的两个端点因而具有相反符号；一般的圆柱也同样有一端带负号。这与\cref{g2:prop:boundary-fundamental-class}中连接同态满足 \(\partial[W,\partial W]=[\partial W]\) 的约定一致。

例如 \(S^n=\partial D^{n+1}\) 是无定向边界，也是取诱导方向后的定向边界。\(\partial(D^2\times S^1)=S^1\times S^1\)，所以二维环面也为边界。将填充球面与环面的三维流形取不交并，就能得到球面到环面的定向配边。配边只记录两端是否能共同出现，\(W\) 内部没有要求每个中间截面都与两端微分同胚。

现在给切向量规定一组随点光滑变化的坐标。实际切丛可能没有全局标架，但加入若干平凡方向后可能出现标架；球面的径向补充就是例子。用 \(\varepsilon_M^r=M\times\mathbb R^r\) 表示带标准基的实平凡丛，底空间清楚时省略下标。

\begin{definition}\label{g2:def:stable-framed-bordism}
	闭光滑 \(n\) 维流形 \(M\) 上的一个\term[wendingqiebiaojia]{稳定切标架}[stable tangent framing]是丛同构
	\[
		\varphi:\mathrm TM\oplus\varepsilon^r\xrightarrow{\ \cong\ }\varepsilon^{n+r},
	\]
	其中 \(r\ge0\)。把 \(\varphi\) 换成 \(\varphi\oplus\operatorname{id}_{\varepsilon^s}\) 称为稳定化；两个稳定切标架在稳定到同一秩后若可经丛同构同伦相连，就视为同一稳定标架。带稳定切标架的流形之间的微分同胚，也要求保持此稳定同伦类。

	紧光滑带边界流形 \(W^{n+1}\) 的稳定切标架诱导边界标架，采用下面的外法向约定。若 \(\partial W\) 的标架分别等于 \(-\varphi_0\) 与 \(\varphi_1\)，就称 \((M_0,\varphi_0)\) 与 \((M_1,\varphi_1)\)\term[daibiaojiapeibian]{带标架配边}[framed bordant]。其中 \(-\varphi\) 是先稳定化，再将一个平凡方向的坐标反号所得的标架。
\end{definition}

把稳定方向的标准正向与目标空间的标准正向固定下来，\(\varphi\) 就确定了 \(\mathrm TM\) 的方向。但是方向只区分基变换行列式的正负，标架还指定了整组坐标如何随点变化；即使方向相同，两组标架之间也可能有无法消去的扭转。忘掉标架得到定向，忘掉定向得到无定向流形。边界上需要延拓的资料随之逐层减少。

下面把边界标架的符号写清楚，以便随后把带标架流形变成球面之间的映射。设
\[
	\Phi:\mathrm TW\oplus\varepsilon^r\longrightarrow\varepsilon^{n+1+r},
\]
取沿边界的向外横截向量场 \(\nu\)，令 \(v\in\mathrm T\partial W\)、\(a\in\mathbb R^r\)、\(t\in\mathbb R\)。将新增的稳定坐标放在最后，定义
\begin{equation}\label{g2:eq:boundary-tangent-framing}
	(\partial\Phi)(v,a,t)=\Phi\bigl(v+(-1)^{n+r}t\nu,a\bigr).
\end{equation}
从 \((v_1,\ldots,v_n,a_1,\ldots,a_r,\nu)\) 把 \(\nu\) 移到最前，经过 \(n+r\) 次交换，故式中的符号恰好保证诱导方向是外法向优先。公式中多出的一维正是边界在 \(W\) 内的法线，它解释了为何这里自然使用稳定标架。

这个定义不依赖所取向外横截向量场：这些向量场的正的法向分量可作线性插值，过程中始终横截，故所得标架同伦。若先把 \(\Phi\) 稳定化一次，式\eqref{g2:eq:boundary-tangent-framing}中的符号会改变，但新增的两个平凡坐标同时改变排列；两种边界标架相差末两坐标上的旋转 \((b,t)\mapsto(-t,b)\)，它经旋转角从 \(0\) 连续变到 \(\pi/2\) 与恒等相连。因此定义与稳定次数相容。类似地，反射哪一个平凡坐标也不影响逆标架的稳定同伦类：在两个平凡方向组成的平面内转动反射轴即可互相连接。

球面的径向同构
\[
	\mathrm TS^n\oplus\varepsilon^1\longrightarrow\varepsilon^{n+1},\qquad(v,t)\longmapsto v+tx
\]
在\cref{g2:eq:sphere-stable-tangent-characteristic}已经出现。按本节的边界约定，球内标准切标架诱导的是 \((v,t)\mapsto v+(-1)^ntx\)。这给出了每个球面上一组确实能延拓到球内的稳定标架。讨论带标架配边时，球面与它所选的标架一起构成对象。

\subsection{粘合、逆元与群结构}
\label{g2:subsec:2301:02}
若 \(W_{01}\) 从 \(M_0\) 配边到 \(M_1\)，\(W_{12}\) 从 \(M_1\) 配边到 \(M_2\)，应能沿共同的 \(M_1\) 把二者接起来。\cref{g2:thm:collar-neighborhood}提供了所需的光滑接缝：每个边界分支附近有 \(M_1\times[0,\varepsilon)\) 的坐标。将两侧的非负参数分别写成 \(u\) 和 \(-u\)，接缝邻域就成为 \(M_1\times(-\varepsilon,\varepsilon)\)。这在商空间上给出光滑流形结构，其剩余边界为 \((-M_0)\amalg M_2\)。定向情形在接缝处一侧为出边界、一侧为入边界，恰好使两个方向相接。

带标架情形还要把标架接起来。在领圈内利用乘积分解识别各切丛，标架就成为一族可逆矩阵；将领圈参数在靠近边界处平滑地压到零，便可把这族标架改成与参数无关。如果给定的边界标架只在稳定后同伦，先稳定到共同秩，再把该同伦放进一段圆柱，使两侧在接缝处逐点一致。对入边界反转外法向时，式\eqref{g2:eq:boundary-tangent-framing}的最后一坐标反号，正与入边界上的 \(-\varphi\) 相抵消。因此拼接后的 \(\mathrm TW\) 上得到同一组稳定标架。

圆柱同时说明了为何反转方向或反射标架应当给出逆元。若 \(\varphi:\mathrm TM\oplus\varepsilon^r\to\varepsilon^{n+r}\)，在 \(M\times[0,1]\) 上定义
\begin{equation}\label{g2:eq:framed-cylinder}
	\Psi(v+s\partial_u,a)=\bigl(\varphi(v,a),(-1)^{n+r}s\bigr).
\end{equation}
其稳定次数为 \(r\)。在上端 \(u=1\)，外法向为 \(\partial_u\)，边界公式给出 \((\varphi(v,a),t)\)；在下端 \(u=0\)，外法向为 \(-\partial_u\)，得到 \((\varphi(v,a),-t)\)。因此它的边界是 \((M,\varphi)\amalg(M,-\varphi)\)。忘掉标架，便恢复 \(\partial(M\times[0,1])=(-M)\amalg M\) 的定向圆柱。

\begin{proposition}\label{g2:prop:bordism-group}
	对每个整数 \(n\ge0\)，闭光滑 \(n\) 维流形的无定向配边、定向配边、带稳定切标架的配边都是等价关系。不交并诱导交换群，分别记作
	\[
		\Omega_n^{\mathrm O},\qquad\Omega_n^{\mathrm{SO}},\qquad\Omega_n^{\mathrm{fr}}.
	\]
	单位元是空流形的类。定向群中 \([M]\) 的逆元为 \([-M]\)，带标架群中 \([M,\varphi]\) 的逆元为 \([M,-\varphi]\)；无定向群中每个元素都是自身的逆元。忘掉标架、忘掉定向分别给出群同态
	\[
		\Omega_n^{\mathrm{fr}}\longrightarrow\Omega_n^{\mathrm{SO}}\longrightarrow\Omega_n^{\mathrm O}.
	\]
\end{proposition}
\begin{proof}
	圆柱给出自反性。交换入边界与出边界，定向情形同时反转 \(W\) 的方向，带标架情形同时反射 \(W\) 的一个稳定坐标，便得到反向配边；无定向情形只须交换边界标识。这证明对称性。上述领圈拼接证明传递性。

	若分别给出 \(M\) 到 \(M'\)、\(N\) 到 \(N'\) 的配边，取它们的不交并就得到 \(M\amalg N\) 到 \(M'\amalg N'\) 的配边，故不交并在等价类上良定义。不交并的结合与交换微分同胚保持各分支所带的结构，给出结合律与交换律，空流形给出单位元。圆柱的全部边界是 \((-M)\amalg M\)，或 \((M,-\varphi)\amalg(M,\varphi)\)，所以这些不交并代表零。无定向时同一个圆柱的边界是两份 \(M\)，故 \(2[M]=0\)。最后，保留较少资料仍是配边，而且与不交并相容，因此两个忘却映射都是群同态。
\end{proof}

\begin{example}\label{g2:exm:zero-dimensional-bordism}
	零维闭流形是有限个点。定向时每个点有一个符号，定义其代数点数为正点数减去负点数。若它是某个紧定向一维流形 \(W\) 的边界，则把常值上同调类 \(1\in H^0(W;\mathbb Z)\) 与边界基本类配对，有
	\[
		\langle1,[\partial W]\rangle
		=\langle1,i_*\partial[W,\partial W]\rangle=0,
	\]
	因为对的长正合列中 \(i_*\partial=0\)。所以代数点数在定向配边下不变。反之，代数点数为零时，可以将每个负点与一个正点配成一对，用一个定向区间连接；这些区间的不交并填满所有点。因此
	\[
		\Omega_0^{\mathrm{SO}}\cong\mathbb Z.
	\]
	无定向时改用 \(\mathbb F_2\) 系数，同一论证说明边界的点数必为偶数；偶数个点又可两两用区间配对，故 \(\Omega_0^{\mathrm O}\cong\mathbb F_2\)。这两个计算中，代数点数与奇偶数分别是完整的不变量。
\end{example}

\subsection{示性数的配边不变性}
\label{g2:subsec:2301:03}
第21章的示性类依赖切丛，第6章的基本类则允许对顶次上同调类求值。二者结合所得的 Stiefel--Whitney 数与 Pontryagin 数已在\cref{g2:def:characteristic-numbers}定义。配边提供了一条检验这些数的办法：若流形确是边界，它的稳定切丛来自填充流形，因而有关示性类也能延拓；能延拓的顶次类在整个边界上的求值为零。

\begin{theorem}\label{g2:thm:characteristic-number-bordism}
	设 \(M_0,M_1\) 为闭光滑 \(n\) 维流形。若二者无定向配边，则每个 Stiefel--Whitney 数相同：对正整数 \(i_1,\ldots,i_s\) 满足 \(\sum i_j=n\)，有
	\[
		\left\langle\prod_jw_{i_j}(\mathrm TM_0),[M_0]_2\right\rangle
		=\left\langle\prod_jw_{i_j}(\mathrm TM_1),[M_1]_2\right\rangle.
	\]
	若二者定向配边且 \(n=4k\)，则每个整数 Pontryagin 数相同：对 \(\sum i_j=k\)，有
	\[
		\left\langle\prod_jp_{i_j}(\mathrm TM_0),[M_0]\right\rangle
		=\left\langle\prod_jp_{i_j}(\mathrm TM_1),[M_1]\right\rangle.
	\]
	因此相应边界的上述各数均为零。每个 Stiefel--Whitney 数定义同态 \(\Omega_n^{\mathrm O}\to\mathbb F_2\)，每个 Pontryagin 数定义同态 \(\Omega_{4k}^{\mathrm{SO}}\to\mathbb Z\)。
\end{theorem}
\begin{proof}
	先证明有关切丛关系，再把求值转移到配边内。设 \(W\) 为相应的紧光滑配边，记 \(A=\partial W\)、\(i:A\hookrightarrow W\)。向外横截向量场给出实际丛同构
	\[
		\mathrm TA\oplus\varepsilon^1\cong i^*\mathrm TW.
	\]
	Stiefel--Whitney 类的 Whitney 公式及自然性（\cref{g2:thm:whitney-sum}）于是给出 \(w_j(\mathrm TA)=i^*w_j(\mathrm TW)\)。对于 Pontryagin 类，复化这一稳定同构，再用 \(p_j(E)=(-1)^jc_{2j}(E_{\mathbb C})\)，同样得到整数等式 \(p_j(\mathrm TA)=i^*p_j(\mathrm TW)\)。这里加的是平凡丛，其总陈类为 \(1\)，所以整数系数的稳定性直接成立。

	无定向情形令 \(R=\mathbb F_2\)，定向情形令 \(R=\mathbb Z\)。相对基本类与边界基本类满足
	\[
		[W,A]_R\in H_{n+1}(W,A;R),\qquad
		\partial[W,A]_R=[A]_R,
	\]
	这是\cref{g2:prop:relative-fundamental-class,g2:prop:boundary-fundamental-class}；模二系数时每个流形都具有该基本类。对的长正合列包含
	\[
		H_{n+1}(W,A;R)\xrightarrow{\partial}H_n(A;R)\xrightarrow{i_*}H_n(W;R),
	\]
	故 \(i_*[A]_R=0\)。对任意 \(\alpha\in H^n(W;R)\)，上同调与同调配对的自然性给出
	\begin{equation}\label{g2:eq:extendible-class-boundary-zero}
		\langle i^*\alpha,[A]_R\rangle=\langle\alpha,i_*[A]_R\rangle=0.
	\end{equation}
	取 \(\alpha=\prod_jw_{i_j}(\mathrm TW)\) 或 \(\alpha=\prod_jp_{i_j}(\mathrm TW)\)，前述切丛关系说明其边界限制恰是所需的示性类乘积。模二时 \([A]_2=[M_0]_2+[M_1]_2\)，而定向整数情形 \([A]=-[M_0]+[M_1]\)，所以零求值分别给出定理中的两条等式。

	在不交并上，基本类是各分支基本类之和，示性类逐分支限制为原来的示性类，故每个求值具有可加性。这也证明它们在配边群上诱导所述同态。
\end{proof}

这个证明的核心是式\eqref{g2:eq:extendible-class-boundary-zero}：全体边界在 \(W\) 内同调于零。它直接保留了模二与整数信息。Euler 数则还依赖实际切丛的秩，不能由 \(\mathrm TA\oplus\varepsilon^1\) 恢复；例如 \(S^2\) 是 \(D^3\) 的边界，却有 \(\chi(S^2)=2\)。

\begin{example}\label{g2:exm:projective-bordism-obstructions}
	在 \(\mathbb RP^2\) 上，令 \(a\in H^1(\mathbb RP^2;\mathbb F_2)\) 为非零类。式\eqref{g2:eq:real-projective-tangent-classes}给出
	\[
		w(\mathrm T\mathbb RP^2)=(1+a)^3=1+a+a^2.
	\]
	由于 \(\langle a^2,[\mathbb RP^2]_2\rangle=1\)，其 \(w_1^2\) 数与 \(w_2\) 数均为 \(1\)。因此 \(\mathbb RP^2\) 不能成为任何紧光滑三维流形的全部边界。

	给 \(\mathbb CP^2\) 取复定向，并令 \(h=c_1(\lambda^*)\) 满足 \(\langle h^2,[\mathbb CP^2]\rangle=1\)。由式\eqref{g2:eq:complex-projective-pontryagin}，
	\[
		p_1(\mathrm T\mathbb CP^2)=3h^2,\qquad
		\langle p_1,[\mathbb CP^2]\rangle=3.
	\]
	所以 \(\mathbb CP^2\) 不是定向边界。更进一步，\(m[\mathbb CP^2]\) 的此项 Pontryagin 数为 \(3m\)，当 \(m\ne0\) 时不会消失，故 \([\mathbb CP^2]\) 在 \(\Omega_4^{\mathrm{SO}}\) 中具有无限阶。
\end{example}

如果 \(M\) 带稳定切标架，\(\mathrm TM\oplus\varepsilon^r\) 平凡便使所有正次 Stiefel--Whitney 类与 Pontryagin 类为零。因而这些数可以先检验一张流形有没有稳定标架；但在已有标架后，它们不能分辨标架本身的扭转。下一节将用球面映射记录这部分资料。

陈数所需的结构还包括复结构。若 \(M^{2m}\) 的 \(\mathrm TM\oplus\varepsilon^{2r}\) 已选复结构，并且配边 \(W^{2m+1}\) 的稳定切丛 \(\mathrm TW\oplus\varepsilon^{2r+1}\) 也有复结构，按外法向约定在边界上诱导给定的稳定复结构，就可以讨论稳定复配边。此时边界陈类是 \(W\) 上陈类的限制，式\eqref{g2:eq:extendible-class-boundary-zero}的整数证明原样适用，得到陈数的配边不变性。普通定向配边只要求方向能延拓，故若要用陈数作障碍，还须指定这项复结构的延拓条件。
% 纲要条目排为编号小节标题，并进入目录与书签；说明段仍保留于纲要注释中。

% !TeX root = ../../Book2.tex
% 纲要标题：### 23.2 Pontryagin–Thom 对应
% 纲要位置：几何学计划纲要.md，第 1893 行。
% 正文状态：未撰写；本文件目前是框架。

% 纲要内容开始
% 
% * 嵌入和法丛标架
% * 管状邻域坍缩映射
% * 横截逆像构造反向映射
% * 同伦与配边的双向对应
% * 稳定维数与选择无关性
% 
%   将嵌入维数的增大与 II.9.5 的悬挂映射对应；分别证明管状邻域选择、标架同伦及环境嵌入选择不会改变稳定类。
% 
% 纲要内容结束

\section{Pontryagin–Thom 对应}
\label{g2:sec:2302}

待撰写

\subsection{嵌入和法丛标架}
\label{g2:subsec:2302:01}
待撰写

\subsection{管状邻域坍缩映射}
\label{g2:subsec:2302:02}
待撰写

\subsection{横截逆像构造反向映射}
\label{g2:subsec:2302:03}
待撰写

\subsection{同伦与配边的双向对应}
\label{g2:subsec:2302:04}
待撰写

\subsection{稳定维数与选择无关性}
\label{g2:subsec:2302:05}
待撰写
% 纲要条目排为编号小节标题，并进入目录与书签；说明段仍保留于纲要注释中。

% !TeX root = ../../Book2.tex
% 纲要标题：### 23.3 手术与把手
% 纲要位置：几何学计划纲要.md，第 1903 行。
% 正文状态：初稿；尚未正式审核。

% 纲要内容开始
% 
% * 局部手术和配边的把手分解
% * 基本群与中维障碍
% * Whitney 技巧
% * h-配边、s-配边与手术障碍
% 
% 纲要内容结束

\section{手术与把手}
\label{g2:sec:2303}

配边把两个流形放在同一个高一维流形的边界上。若沿高度切开这个配边，中间截面会怎样变化？第17章的把手附着已经给出答案所需的局部模型。本节从这个模型出发，说明怎样有目的地改变基本群，以及为什么同样的办法走到中维时会遇到相交问题。

\subsection{局部手术和配边的把手分解}
\label{g2:subsec:2303:01}
设 \(M\) 是闭光滑 \(n\) 维流形，\(0\le p\le n-1\)。要在一个 \(p\) 维球面附近动手，首先需要的不只是嵌入 \(S^p\hookrightarrow M\)，而是一个乘积邻域
\[
	\varphi:S^p\times D^{n-p}\hookrightarrow M.
\]
它等价于给中心球面选择一组法标架，再应用管状邻域定理。记 \(C=M\setminus\varphi(S^p\times\operatorname{int}D^{n-p})\)。两块
\[
	S^p\times D^{n-p},\qquad D^{p+1}\times S^{n-p-1}
\]
具有同一个边界 \(S^p\times S^{n-p-1}\)，因而可以相互替换。

\begin{definition}\label{g2:def:surgery}
	沿 \(\varphi\) 的\term[shoushu]{\(p\) 维手术}[surgery]是流形
	\[
		M_\varphi=C\cup_{\varphi|_{S^p\times S^{n-p-1}}}(D^{p+1}\times S^{n-p-1}),
	\]
	粘合处使用领圈并磨圆。把 \(D^{p+1}\times D^{n-p}\) 沿 \(S^p\times D^{n-p}\) 接到 \(M\times\{1\}\)，所得 \((n+1)\) 维配边称为这次手术的\term[shoushuji]{手术迹}[trace of surgery]，记为 \(W_\varphi\)。
\end{definition}

\begin{proposition}\label{g2:prop:surgery-trace}
	\(W_\varphi\) 从 \(M\) 出发只有一个指标 \(p+1\) 的把手，另一端边界是 \(M_\varphi\)。反向看，它从 \(M_\varphi\) 出发只有一个指标 \(n-p\) 的把手。若 \(M\) 定向，按定向兼容的附着方式构造，则 \(\partial W_\varphi=(-M)\amalg M_\varphi\)。
\end{proposition}
\begin{proof}
	乘积圆盘的边界由两个面组成：
	\[
		\partial(D^{p+1}\times D^{n-p})
		=(S^p\times D^{n-p})\cup(D^{p+1}\times S^{n-p-1}).
	\]
	第一个面附着后成为内部接缝，第二个面与 \(C\) 合成新边界。这正是定义中的替换。对照\cref{g2:def:morse-handle}，核心圆盘是 \(D^{p+1}\times\{0\}\)，故指标为 \(p+1\)。倒转配边方向，交换两个圆盘因子，原来的自由面成为附着面，新的核心为 \(D^{n-p}\)，指标成为 \(n-p\)。最后按外法向优先的边界定向约定，乘积的下端带相反定向，上端带原定向，得到所述边界公式。
\end{proof}

三个维数分别是：接受手术的流形维数 \(n\)、被取走的中心球面维数 \(p\)、手术迹的把手指标 \(p+1\)。这里把“\(p\) 维手术”按球面维数命名；其他文献还可能按把手指标命名。

\begin{proposition}\label{g2:prop:surgery-euler}
	一次 \(p\) 维手术满足
	\[
		\chi(M_\varphi)-\chi(M)=(-1)^{n-p-1}-(-1)^p.
	\]
	因此奇数维时 Euler 示性数不变，偶数维时改变量为 \(-2(-1)^p\)。
\end{proposition}
\begin{proof}
	利用领圈把粘合分解加厚成开覆盖，Mayer--Vietoris 序列给出 Euler 示性数的加法公式。两次分解具有相同的 \(C\) 和交界 \(S^p\times S^{n-p-1}\)，相减后只剩
	\[
		\chi(D^{p+1}\times S^{n-p-1})-\chi(S^p\times D^{n-p})
		=1+(-1)^{n-p-1}-1-(-1)^p.
	\]
	球面的 Euler 示性数来自它的零胞腔和顶维胞腔，圆盘可缩；故公式也包括 \(S^0\) 出现的端点情形。
\end{proof}

\begin{example}\label{g2:exm:surface-surgery}
	在定向曲面上，\(p=0\) 的手术取走两个小圆盘，再用圆柱 \(D^1\times S^1\) 接起两个边界圆。如果两个圆盘位于同一连通曲面，就增加一个把手，\(\chi\) 减少 \(2\)；如果它们分别位于两张曲面，就把两张曲面连成其连通和。反向的 \(p=1\) 手术切去一个圆柱邻域，再用两个圆盘封口，\(\chi\) 增加 \(2\)。沿不分离圆周切开会减少一个把手，沿分离圆周切开则产生两个分支。相同的局部替换可以有不同的全局效果，区别来自附着球面的位置。
\end{example}

任意紧光滑配边 \((W;M_0,M_1)\) 都可以这样逐步扫描。\Cref{g2:prop:boundary-morse-function}给出在两端领圈内等于高度参数、内部只有有限个非退化临界点且临界值互异的函数。逐次应用\cref{g2:thm:morse-handle-attachment}，每跨过一个指标 \(\lambda\) 的临界点，截面便经历一次 \((\lambda-1)\) 维手术。指标 \(0\) 和 \(n+1\) 分别使一个 \(S^n\) 分支诞生或消失，需要另计。这样得到的是有限把手分解；把手的位置和法标架保留了仅凭配边群元素看不到的信息。

\subsection{基本群与中维障碍}
\label{g2:subsec:2303:02}
先考察手术能否消去一个指定环路。直觉上，沿这个环路接上圆盘便会让它可缩；手术还必须同时补上横向厚度，保证结果仍是流形。手术迹将这两件事分开处理。

\begin{proposition}\label{g2:prop:surgery-fundamental-group}
	设 \(M\) 连通，\(1\le p\le n-3\)，并沿 \(\varphi:S^p\times D^{n-p}\hookrightarrow M\) 做手术。则 \(M_\varphi\) 连通，且
	\[
		\pi_1(M_\varphi)\cong
		\begin{cases}
			\pi_1(M)/\langle\!\langle[\alpha]\rangle\!\rangle,&p=1,\\
			\pi_1(M),&p\ge2.
		\end{cases}
	\]
	这里 \(\alpha=\varphi|_{S^1\times\{0\}}\)，用一条道路把它接到固定基点；\(\langle\!\langle[\alpha]\rangle\!\rangle\) 表示它生成的正规子群。更换这条道路只会共轭 \([\alpha]\)，不改变该正规子群。
\end{proposition}
\begin{proof}
	由\cref{g2:prop:surgery-trace,g2:lem:handle-core-retraction}，从 \(M\) 到 \(W_\varphi\) 的同伦效果是附着一个 \((p+1)\) 胞腔。对于 \(X\cup_\beta D^k\)，\(k\ge2\)，取圆盘内部及圆盘中心的补集作为开覆盖：前者可缩，后者收缩到 \(X\)，交集收缩到 \(S^{k-1}\)。\Cref{g2:thm:van-kampen}于是给出
	\[
		\pi_1(X\cup_\beta D^k)=
		\begin{cases}
			\pi_1(X)/\langle\!\langle\beta_*\pi_1(S^1)\rangle\!\rangle,&k=2,\\
			\pi_1(X),&k\ge3.
		\end{cases}
	\]
	第二行用到了 \(S^{k-1}\) 单连通。因而从 \(M\) 出发计算 \(\pi_1(W_\varphi)\)，正好得到命题右侧。

	反向的核心胞腔维数是 \(n-p\ge3\)。它的附着球面连通，不能把原有的不同连通分支接在一起；而 \(W_\varphi\) 从连通的 \(M\) 附着一个把手得到，故本身连通。因此 \(M_\varphi\) 也连通。对反向附着再用上式，\(\pi_1(M_\varphi)\to\pi_1(W_\varphi)\) 是同构。把两端的计算结合起来即得结论。
\end{proof}

维数限制在证明中有明确位置：\(n-p\ge3\) 保证反向附着不改变基本群。如果反向把手的指标降到 \(2\)，它也能消去关系；降到 \(1\)，则可能新增环路或连接分支。三维流形沿圆周的手术恰好落在前一种情形，因此不能套用上面的商群公式。

\ProseParagraphBreak

在高维定向流形中，有限组基本群生成元可用互不相交的嵌入圆周表示。每个圆周的法丛是定向的，而 \(S^1\) 上的定向实向量丛平凡，因而可以选择法标架。当 \(n\ge4\) 时，逐次作 \(1\) 维手术便能消去这些生成元。这说明存在同一有向配边类中的单连通代表；这一步改变了原流形的同伦型。

\ProseParagraphBreak

想继续消去更高同伦群，就要用嵌入的 \(S^p\) 代表待消去元素，并为它选择法标架。横截性提供一个有用的维数判断：两个 \(p\) 维代表在 \(n\) 维流形中的预期相交维数是 \(2p-n\)。在 \(2p<n\) 时，有希望通过一般位置使它们不相交；在 \(n=2p\) 时，交点成为孤立点，不能再靠一个小扰动全部移走。这是中维在手术理论中反复出现的原因。

\ProseParagraphBreak

设 \(M^{2q}\) 闭、连通、定向。其自由中维同调上的相交配对可写成
\[
	\lambda(x,y)=\bigl\langle\operatorname{PD}(x)\smile\operatorname{PD}(y),[M]\bigr\rangle,
	\qquad x,y\in H_q(M;\mathbb Z)/\text{挠子群}.
\]
这里 \(\operatorname{PD}(x)\in H^q(M;\mathbb Z)\) 是 Poincar\'e 对偶类，模挠意义下的取值不依赖代表；杯积的分次交换律给出 \(\lambda(y,x)=(-1)^q\lambda(x,y)\)。因而 \(q\) 偶数时得到对称形式，\(q\) 奇数时得到斜对称形式。若两个类由互不相交的球面代表，则交数为零。若一个球面具有平凡法丛，沿其中一个非零法向量场略作推移，就得到与自身不相交的同调代表，所以它的自交数也必须为零。法标架与相交数据都是能否实施手术的实际条件。

\begin{example}\label{g2:exm:surgery-hyperbolic-pair}
	在 \(S^q\times S^q\) 中，取 \(a=[S^q\times\{y\}]\)、\(b=[\{x\}\times S^q]\)，并选择乘积定向，则
	\[
		\bigl(\lambda(a_i,a_j)\bigr)_{(a_1,a_2)=(a,b)}
		=\begin{pmatrix}0&1\\(-1)^q&0\end{pmatrix}.
	\]
	两个因子球面各自都有平凡法丛，却具有非零交数，所以不可能同时由互不相交的代表实现。手术中的代数整理需要选择合适的半组生成元并保留与之对偶的数据；仅逐项检查每个生成元的自交数还不够。
\end{example}

基本群非平凡时，还需记录交点的道路信息。对两个带基点道路的、基本群像为零的球面代表，在交点 \(z\) 沿第一个球面到达 \(z\)，再沿第二个球面返回基点，得到 \(g_z\in\pi=\pi_1(M)\)。把交点符号 \(\varepsilon_z\) 一同记下，形成群环中的和
\[
	\lambda=\sum_z\varepsilon_z g_z\in\mathbb Z[\pi].
\]
固定基点道路后，这些标签可比较；改变选择会引入相应的群作用。取增广 \(\mathbb Z[\pi]\to\mathbb Z\)、把每个群元素送到 \(1\)，才得到普通整数交数。因此两个符号相反但群标签不同的交点，虽然在整数计数中相消，却未必能用下述圆盘消去。中维手术障碍保留群环配对及与自交、法标架有关的二次数据，原因就在这里。

\subsection{Whitney 技巧}
\label{g2:subsec:2303:03}
设定向子流形 \(A^a,B^b\subset X^{a+b}\) 横截相交，\(u,v\) 是符号相反的两个交点。在 \(A\) 中选一条连接 \(u,v\) 的嵌入弧，在 \(B\) 中再选一条反向连接的嵌入弧，使弧内部避开其他交点。它们合成一条带两个角的闭曲线，称为 Whitney 圆周。若能用一个圆盘填住它，就可以尝试把 \(A\) 的这一小片沿圆盘推过 \(B\)，使两次穿越一起消失。

\ProseParagraphBreak

这个圆盘还必须光滑嵌入，内部与 \(A\cup B\) 不相交，边界按上述两条弧落在两张片上。沿两条边界弧由片内方向给出的法向数据也须能延拓到整个圆盘，形成与两片相容的乘积邻域；这称为圆盘具有 Whitney 标架。只有闭曲线可缩，尚未检查这两件事。下面把已有这种局部乘积模型时的消点操作写出。

\begin{lemma}\label{g2:lem:local-whitney-move}
	设上述两个交点具有一个带 Whitney 标架的圆盘，其内部避开 \(A\cup B\)，且其小邻域内没有其他片。则存在支撑在该邻域中的光滑同痕，把 \(A\) 移成 \(A'\)，使
	\[
		A'\cap B=(A\cap B)\setminus\{u,v\}.
	\]
	这里同痕作用于 \(A\)，比较相交时保持 \(B\) 为原来的子流形。
\end{lemma}
\begin{proof}
	标架把圆盘的小邻域识别为一个平面邻域 \(U\) 乘上 \(\mathbb R^{a-1}\times\mathbb R^{b-1}\) 中的小圆盘。在这个识别下，两张片分别写成
	\[
		C_A\times\mathbb R^{a-1}\times\{0\},\qquad
		C_B\times\{0\}\times\mathbb R^{b-1},
	\]
	其中平面弧 \(C_A,C_B\) 围成待消去的二角形。把 \(C_A\) 推过这个二角形，可得到支撑在 \(U\) 内的平面同痕 \(G_t\)，使 \(G_1(C_A)\) 与 \(C_B\) 不相交。

	这个平面操作可以直接实现：把 \(C_B\) 拉直为横轴，把相关的 \(C_A\) 写成 \(y=h(x)\)，其中两个横截零点之间 \(h<0\)，两端外侧 \(h>0\)。选紧支光滑函数 \(k\ge0\)，使 \(h+k>0\) 于相关区间，令 \(h_t=h+tk\)。每条图像都嵌入；取沿图像等于 \(k(x)\partial_y\)、离开图像附近便截断为零的时变向量场，它的流将 \(h_0\) 送到 \(h_t\)，给出所需 \(G_t\)。二角形的两边嵌入且内部没有其他片，允许先作上述平面坐标调整。

	取横向变量上的光滑截断 \(\rho(x,y)\)，在 \((0,0)\) 附近为 \(1\)，靠近横向邻域边界为 \(0\)，定义
	\[
		H_t(z,x,y)=\bigl(G_{t\rho(x,y)}(z),x,y\bigr).
	\]
	每个固定 \((x,y)\) 的平面映射都是微分同胚，逆映射也光滑依赖参数，故 \(H_t\) 是同痕，并可与邻域外的恒等映射拼接。移动后的第一张片仍有 \(y=0\)，第二张片仍有 \(x=0\)，所以二者若相交，只能发生在 \(x=y=0\)。此处 \(\rho=1\)，相交等价于 \(G_1(C_A)\cap C_B\)，已经为空。邻域外没有变化，故恰好消去指定的两个交点。
\end{proof}

上述引理把几何操作与圆盘的存在问题分开了。以下是高维中常用的一组充分条件。

\begin{theorem}[Whitney 技巧]\label{g2:thm:whitney-trick}
	设 \(X^{a+b}\) 是单连通定向光滑流形，\(A^a,B^b\) 是其中连通、定向、横截相交的闭嵌入子流形。若 \(a,b\ge3\)，则任意两个符号相反的交点可以通过 \(A\) 的光滑同痕消去，而不改变其余交点。
\end{theorem}

这里引用完整的 Whitney 圆盘与标架构造，见 \cite[定理6.6及引理6.7]{Milnor1965HCobordism}；群环标签版本见 \cite[第7.3节]{Ranicki2002Surgery}。其维数机制如下：单连通性使 Whitney 圆周能够由圆盘填充；对填充圆盘作相对一般位置调整，可能出现的自交维数是 \(4-(a+b)<0\)，与 \(A\)、\(B\) 内部的预期相交维数分别是 \(2-b<0\)、\(2-a<0\)。这些不等式使圆盘可以嵌入且内部避开两片。剩下的边界法向扭转由标架构造处理，随后应用\cref{g2:lem:local-whitney-move}。标架这一步是引用定理的一部分，不能只由“圆盘可缩”替代。

\ProseParagraphBreak

在四维中若 \(a=b=2\)，刚才三个预期维数都变成 \(0\)。填充圆盘本身可以有双点，也可以穿过原来的两张曲面；试图消去一对交点可能又产生新的交点。单连通性仍解决边界圆周可缩的问题，却不能保证存在内部干净且带所需标架的嵌入圆盘。已经给定这种圆盘时，局部消点引理仍然成立；失去的是由一般位置自动得到圆盘的步骤。下一节的四维光滑现象与这道障碍密切相关。

\subsection{h-配边、s-配边与手术障碍}
\label{g2:subsec:2303:04}
配边允许中间截面改变许多次。若两端的包含映射都已经是同伦等价，这些把手在同伦论意义下没有留下新信息，能否把它们全部取消？

\begin{definition}\label{g2:def:h-cobordism}
	紧光滑配边 \((W;M_0,M_1)\) 称为\term[hpeibian]{\(h\)-配边}[h-cobordism]，如果两个包含 \(M_i\hookrightarrow W\) 都是同伦等价。称它为乘积配边，是指存在微分同胚
	\[
		F:M_0\times[0,1]\longrightarrow W,
		\qquad F(x,0)=x,\quad F(M_0\times\{1\})=M_1.
	\]
\end{definition}

\begin{theorem}[\(h\)-配边定理]\label{g2:thm:h-cobordism}
	设 \((W;M_0,M_1)\) 为光滑 \(h\)-配边，\(M_0,M_1\) 是闭连通流形，\(\dim W\ge6\)。如果 \(M_0\) 单连通，则 \(W\) 是乘积配边，特别地 \(M_0\) 与 \(M_1\) 微分同胚。
\end{theorem}

这是 Smale 的定理；完整证明可见 \cite[第6--9节，定理9.1]{Milnor1965HCobordism}。包含为同伦等价保证 \(W\) 和 \(M_1\) 也单连通。证明先把相对把手链复形中的代数消去实现为把手滑动，再用 Whitney 技巧把代数交数化成恰好一次的几何相交，最后取消相邻指标的把手。这个过程所用的高维把手整理是本节的外部输入。定理给出的是固定下端 \(M_0\) 的乘积结构；任意预先指定的上端参数化不一定也逐点固定。

\ProseParagraphBreak

非单连通时，相对万有覆盖的链复形是 \(\mathbb Z[\pi]\) 上的带基自由链复形，其中 \(\pi=\pi_1(M_0)\)。同伦等价使它链可缩，但把它约去所需的基变换未必能由初等变换及把基元素乘以 \(\pm g\)、\(g\in\pi\) 实现。这份差别由\term[whiteheadniuzi]{Whitehead 扭子}[Whitehead torsion]记录：
\[
	\tau(W,M_0)\in\operatorname{Wh}(\pi),\qquad
	\operatorname{Wh}(\pi)=K_1(\mathbb Z[\pi])/\langle[\pm g]:g\in\pi\rangle.
\]
这里 \(K_1(R)\) 是把可逆矩阵稳定化后再模去初等矩阵所得的群；利用链收缩构造奇、偶维链群之间的基变换，取其稳定类，便定义上述扭子。改变胞腔提升和定向会产生 \(\pm g\)，所以要再除掉这些类。构造及其不变性详见 \cite[第8.1--8.2节]{Ranicki2002Surgery}。

\begin{theorem}[\(s\)-配边定理]\label{g2:thm:s-cobordism}
	设 \((W;M_0,M_1)\) 为光滑 \(h\)-配边，两端闭且连通，\(\dim W\ge6\)。则它是固定 \(M_0\) 的乘积配边，当且仅当 \(\tau(W,M_0)=0\)。
\end{theorem}

这一版本见 \cite[第8.2节]{Ranicki2002Surgery}；扭子为零的 \(h\)-配边称为 \(s\)-配边。由于 \(\operatorname{Wh}(\{1\})=0\)，单连通版本是它的特例。使用定理时，首先要有两个包含映射的同伦等价，再检验扭子；两端同调群同构尚未给出这些条件。

\ProseParagraphBreak

最后说明手术怎样把一个“接近同伦等价”的映射改成同伦等价。设 \(M^n,N^n\) 闭、连通、定向。一个\term[duyifayingse]{度一法映射}[degree-one normal map]包含度为 \(1\) 的映射 \(f:M\to N\)、\(N\) 上的稳定实向量丛 \(\xi\)，以及稳定丛同构 \(b:\nu_M\cong f^*\xi\)。这里 \(\nu_M\) 是稳定法丛，同构也可表示为选定 \(\mathrm{T}M\oplus f^*\xi\) 的稳定平凡化。\(\xi\) 是数据的一部分，不预设为 \(\nu_N\)。称两份这样的数据\term[fapeibian]{法配边}[normally bordant]，是指存在配边上的映射及稳定法丛数据延拓它们，目标丛保持为 \(\xi\)。

\begin{theorem}[手术障碍定理]\label{g2:thm:simply-connected-surgery-obstruction}
	设 \(N^n\) 是单连通闭定向光滑流形，\(n\ge5\)，\((f,b):M^n\to N^n\) 是度一法映射。存在只依赖其法配边类的障碍
	\[
		\theta(f,b)\in L_n(\mathbb Z),
	\]
	它为零当且仅当 \((f,b)\) 法配边于一个同伦等价 \(f':M'\to N\)。这里
	\[
		L_n(\mathbb Z)\cong
		\begin{cases}
			\mathbb Z,&n\equiv0\pmod4,\\
			\mathbb Z/2,&n\equiv2\pmod4,\\
			0,&n\text{ 为奇数}.
		\end{cases}
	\]
\end{theorem}

这里引用手术障碍定理及单连通 \(L\) 群计算 \cite[定理3.2、13A.1]{Wall1999Surgery}。原定理表述为单同伦等价；此处基本群平凡，Whitehead 群为零，因而与同伦等价一致。\(L_n(\mathbb Z)\) 是二次手术障碍群：在偶数维，它保留中维相交形式及二次细化模去可通过标准手术消去的数据后的类。\(4k\) 维的整数由相交形式的符号差除以 \(8\) 给出，\(4k+2\) 维的二值障碍由相应中维二次形式的二值不变量给出；这些代数构造和上表的计算也包含在所引结果中。

\ProseParagraphBreak

例如，当目标维数为 \(5\) 或 \(7\) 时，一旦已经给出度一法映射，其障碍群为零，所以可以在法配边中得到同伦等价。输出是新的 \(M'\) 及 \(M'\simeq N\)；从同伦等价走到微分同胚还需进一步信息。对固定同伦类型的光滑流形进行分类，还须研究法数据之间的差别、手术群的作用和 \(s\)-配边关系。把这几个层次分清，下一节同一拓扑类型上不同光滑结构的例子才不会显得与手术理论矛盾。
% 纲要条目排为编号小节标题，并进入目录与书签；说明段仍保留于纲要注释中。

% !TeX root = ../../Book2.tex
% 纲要标题：### 23.4 光滑、PL 与拓扑类别
% 纲要位置：几何学计划纲要.md，第 1860 行。
% 正文状态：未撰写；本文件目前是框架。

% 纲要内容开始
% 
% * 结构的遗忘及提升问题
% * 低维和高维的不同现象
% * 典型异光滑结构的定理与限制
% 
% ---
% 
% 纲要内容结束

\section{光滑、PL 与拓扑类别}
\label{g2:sec:2304}

待撰写

\subsection{结构的遗忘及提升问题}
\label{g2:subsec:2304:01}
待撰写

\subsection{低维和高维的不同现象}
\label{g2:subsec:2304:02}
待撰写

\subsection{典型异光滑结构的定理与限制}
\label{g2:subsec:2304:03}
待撰写
% 纲要条目排为编号小节标题，并进入目录与书签；说明段仍保留于纲要注释中。


\begin{chaptersummary}
	配边把闭流形放到高一维紧流形的边界上比较。不交并给出群运算，圆柱给出逆元，领圈使配边能够拼接。边界稳定切丛来自内部，因此 Stiefel--Whitney 数和整数 Pontryagin 数在相应配边下不变；全体边界在内部同调于零是求值消失的关键。Euler 数还依赖实际秩，不能直接套用这一稳定性论证。

	\ProseParagraphBreak

	法标架把管状邻域的纤维统一成 \(\mathbb R^k\)，坍缩由此产生 \(S^{n+k}\to S^k\)。正则逆像通过微分取回子流形及其法标架。相对横截性把同伦变成配边，保持领圈的管状邻域又把配边变成同伦，两种构造因而互逆。增加一个环境坐标等于增加一个法标架向量，对应映射的悬挂；消去嵌入的选择后得到 \(\Omega_n^{\mathrm{fr}}\cong\pi_n^{\mathrm S}\)。

	\ProseParagraphBreak

	手术把 \(S^p\times D^{n-p}\) 换为 \(D^{p+1}\times S^{n-p-1}\)，其配边迹附着一个指标 \(p+1\) 的把手。低指标手术改变基本群，中维手术受到相交形式及标架数据的约束。高维的一般位置与 Whitney 消交支持 \(h\)-配边和手术理论；四维中圆盘与圆盘仍会相交，相关光滑问题呈现另一种结构。光滑、PL、拓扑三种等价关系须分别辨认，异光滑球面与异光滑 \(\mathbb R^4\) 正是同胚类型不足以确定光滑类型的具体实例。
\end{chaptersummary}

\BookExercises
% 第二卷第23章习题候选：对应23.1；由章入口统一收录。
\begin{exercise}\label{g2:ex:euler-boundary-formula}
	设 \(W\) 为紧光滑 \(d\) 维流形，\(d\ge1\)，不要求定向。用模二系数的 Poincaré--Lefschetz 对偶证明
	\[
		\chi(\partial W)=\bigl(1-(-1)^d\bigr)\chi(W).
	\]
	由此说明，Euler 示性数为奇数的闭光滑流形不能是无定向边界，并解释这与 \(S^2=\partial D^3\) 相容的原因。
\end{exercise}
\begin{solution}
	以下同调均取 \(\mathbb F_2\) 系数。先说明 Euler 交错和是有限的：把 \(W\) 双化为闭光滑流形 \(\operatorname{Dbl}(W)\)，由\cref{g2:cor:morse-cell-type}，双化具有有限 CW 同伦型。把两份副本都映回原来的 \(W\)，给出一个连续折叠映射，其在第一份 \(W\) 上是恒等，所以 \(H_k(W)\) 是双化的 \(H_k\) 的直和因子。边界本身是闭流形，其同调也有限；对的长正合列再保证相对同调有限。故这些 Euler 示性数均可由同调维数的有限交错和计算。

	\cref{g2:thm:poincare-lefschetz}在模二系数下给出
	\[
		\dim H_k(W,\partial W)=\dim H^{d-k}(W)=\dim H_{d-k}(W).
	\]
	因此 \(\chi(W,\partial W)=(-1)^d\chi(W)\)。对的长正合列中各项维数的交错和为零，于是
	\[
		\chi(W)=\chi(\partial W)+\chi(W,\partial W)
		=\chi(\partial W)+(-1)^d\chi(W),
	\]
	即得所求等式。若 \(d\) 偶，则边界的 Euler 示性数为零；若 \(d\) 奇，则它为 \(2\chi(W)\)。两种情形均为偶数，故奇数 Euler 示性数阻止流形成为无定向边界。对于 \(D^3\)，右端是 \(2\chi(D^3)=2\)，正好等于 \(\chi(S^2)\)；边界的 Euler 示性数可以非零。
\end{solution}

\begin{exercise}\label{g2:ex:projective-product-bordism-numbers}
	令 \(X=\mathbb RP^2\times\mathbb RP^2\)、\(Y=\mathbb CP^2\)。计算这两张实四维流形的全部 Stiefel--Whitney 数，即对应于
	\[
		w_1^4,\quad w_1^2w_2,\quad w_1w_3,\quad w_2^2,\quad w_4
	\]
	的五项求值。它们能否区分 \(X\) 与 \(Y\) 的无定向配边类？说明仅凭\cref{g2:thm:characteristic-number-bordism}能得到的结论。
\end{exercise}
\begin{solution}
	设 \(a,b\) 是从两因子拉回的一次生成元。模二上同调环为
	\[
		H^*(X;\mathbb F_2)=\mathbb F_2[a,b]/(a^3,b^3),\qquad
		\langle a^2b^2,[X]_2\rangle=1.
	\]
	积的切丛是两因子切丛拉回之和，因此 Whitney 公式给出
	\[
		w(\mathrm TX)=(1+a+a^2)(1+b+b^2).
	\]
	逐次展开，得到
	\[
		w_1=a+b,\quad w_2=a^2+ab+b^2,\quad
		w_3=a^2b+ab^2,\quad w_4=a^2b^2.
	\]
	于是 \(w_1^4=a^4+b^4=0\)，\(w_1^2w_2=(a^2+b^2)(a^2+ab+b^2)=2a^2b^2=0\)，且 \(w_1w_3=2a^2b^2=0\)。另一方面，\(w_2^2=a^2b^2=w_4\)，二者求值都为 \(1\)。

	对于 \(Y\)，复切丛给出 \(w_1=w_3=0\)、\(w_2=\overline h\)、\(w_4=\overline h^{\,2}\)，且 \(\langle\overline h^{\,2},[Y]_2\rangle=1\)。所以全部答案为
	\[
	\begin{array}{c|ccccc}
		& w_1^4&w_1^2w_2&w_1w_3&w_2^2&w_4\\ \hline
		X&0&0&0&1&1\\
		Y&0&0&0&1&1
	\end{array}
	\]
	这些数不能区分二者，且都证明 \(X\)、\(Y\) 各自不是无定向边界。在 \(X\amalg Y\) 上五个数相加后全部为零，所以这些数不再给出边界障碍。\cref{g2:thm:characteristic-number-bordism}证明的是配边推出各数相等；要由此表反过来断言二者配边，还需要这个必要条件的充分性，或另行构造具体配边。
\end{solution}

\begin{exercise}\label{g2:ex:connected-sum-bordism}
	设 \(M,N\) 为连通闭定向光滑 \(n\) 维流形，\(n\ge2\)。构造 \(M\amalg N\) 到定向连通和 \(M\mathbin{\#}N\) 的配边，并推出每个 Pontryagin 数在定向连通和下可加。判断 \(\mathbb CP^2\mathbin{\#}\mathbb CP^2\) 和 \(\mathbb CP^2\mathbin{\#}(-\mathbb CP^2)\) 是否为定向边界，并分别给出理由。
\end{exercise}
\begin{solution}
	从 \((M\amalg N)\times[0,1]\) 出发，在两个出边界上各选一张嵌入的 \(n\) 维闭球，沿这两张球附着一只 \(n+1\) 维的一把手 \([-1,1]\times D^n\)：它的两个端盘分别接到两因子，按反转边界方向的方式粘合。原来的上端边界被去掉两张球的内部，再接上把手侧面 \([-1,1]\times S^{n-1}\)。这正是定向连通和的颈部，故新的上端为 \(M\mathbin{\#}N\)。在球边缘的领圈坐标中，接缝是二维象限乘 \(S^{n-1}\)；把象限的折角换成一条与两侧相切的光滑弧，就得到光滑边界，并在该小邻域以外保持原构造。所得紧定向流形 \(V\) 满足
	\[
		\partial V=(-M)\amalg(-N)\amalg(M\mathbin{\#}N).
	\]
	因此 \([M\mathbin{\#}N]=[M]+[N]\) 于 \(\Omega_n^{\mathrm{SO}}\)。若 \(n=4k\)，每个 Pontryagin 数是这个群上的同态，故在连通和下等于两因子数值之和。

	\(\mathbb CP^2\mathbin{\#}\mathbb CP^2\) 的 \(p_1\) 数为 \(3+3=6\)，所以不是定向边界。\(\mathbb CP^2\mathbin{\#}(-\mathbb CP^2)\) 则有配边类 \([\mathbb CP^2]+[-\mathbb CP^2]=0\)，因而确是定向边界。这里边界存在性来自上述配边与反向圆柱的粘合；其 Pontryagin 数 \(3-3=0\) 与此相符。
\end{solution}

\begin{exercise}\label{g2:ex:pt-change-target-frame}
	设 \(f:S^{n+k}\to S^k\) 光滑，在基点附近恒取基点，且 \(0\) 为正则值。给定 \(A\in\operatorname{GL}(k,\mathbb R)\)，将 \(v\mapsto Av\) 延拓为一点紧化的自映射 \(\widehat A:S^k\to S^k\)。求 \(\widehat A\circ f\) 在零处的法标架，并分别说明 \(\det A>0\)、\(\det A<0\) 对应的稳定配边类。
\end{exercise}
\begin{solution}
	零点集合仍为 \(M=f^{-1}(0)\)。在法方向上，链式法则给出 \(\mathrm d(\widehat A\circ f)=A\,\mathrm df\)，故新标架为
	\[
		b_{\widehat A\circ f}=b_fA^{-1}.
	\]
	当 \(\det A>0\) 时，\(A^{-1}\) 可在 \(\operatorname{GL}^+(k,\mathbb R)\) 中连接到恒等，因此新旧标架同伦。若 \(\det A<0\)，可将 \(A^{-1}\) 连到反射 \(R=\operatorname{diag}(-1,1,\ldots,1)\)，于是稳定切标架变为 \(\phi_{b_f}(v,Ra)\)，即在一个平凡方向上反射。按逆元约定，这给出原稳定带标架配边类的负元。\(\widehat A\) 在无穷远处只需连续；复合映射在基点邻域恒定，在零点邻域光滑，若需要全局光滑代表，可保持这些区域作相对逼近。
\end{solution}

\begin{exercise}\label{g2:ex:cp-line-normal-framing-obstruction}
	把 \(\mathbb CP^1\) 作为标准射影直线嵌入 \(\mathbb CP^2\)。求其复法线丛的第一陈类，判断它能否成为某张光滑映射 \(\mathbb CP^2\to S^2\) 的正则值逆像。进一步判断增加平凡实方向能否使这张实法丛平凡。
\end{exercise}
\begin{solution}
	记 \(h\in H^2(\mathbb CP^1;\mathbb Z)\) 为正生成元。切丛限制分解为 \(\mathrm T\mathbb CP^2|_{\mathbb CP^1}\cong\mathrm T\mathbb CP^1\oplus\nu\)。由式\eqref{g2:eq:complex-projective-tangent-classes}，两边的第一陈类满足
	\[
		3h=2h+c_1(\nu),\qquad c_1(\nu)=h.
	\]
	法丛为复线丛，其底层定向实平面丛的 Euler 类就是 \(h\)，故不是平凡丛。如果它是 \(f^{-1}(y)\) 且 \(y\) 为正则值，微分会给出 \(\nu\cong\mathbb CP^1\times\mathrm T_yS^2\)，与此矛盾。

	稳定问题须另用稳定类判断。第21章的复丛实化比较给出 \(w_2(\nu_{\mathbb R})=c_1(\nu)\bmod2\ne0\)，而增加平凡实丛保持 \(w_2\)。所以它即使稳定后也不能平凡。Euler 类非零本身不足以完成最后一步，例如 \(\mathrm TS^2\) 的 Euler 类非零，却有 \(\mathrm TS^2\oplus\varepsilon^1\cong\varepsilon^3\)。
\end{solution}

\begin{exercise}\label{g2:ex:compact-level-framed-boundary}
	设 \(f:\mathbb R^{n+1}\to\mathbb R\) 光滑，\(0\) 为正则值，且在某个球外 \(f>0\)。证明 \(M=f^{-1}(0)\)，配上法向量 \(\nabla f\)，代表零元 \(\Omega_n^{\mathrm{fr}}\)。说明紧性假设在此有什么作用。
\end{exercise}
\begin{solution}
	子水平集 \(W=\{f\le0\}\) 闭且有界，因而紧。正则值定理给出 \(\partial W=M\)，外法向可取 \(\nu=\nabla f\)。环境标准切标架限制到 \(W\)，诱导的边界标架为
	\[
		(v,t)\longmapsto v+(-1)^nt\nabla f.
	\]
	因此这一稳定切标架代表零。法标架 \(b=\nabla f\) 所给的稳定切标架为 \((v,t)\mapsto v+t\nabla f\)：\(n\) 偶时它与上述标架相同，\(n\) 奇时它是上述类的逆元，仍为零。空水平集也满足结论。

	球外的正性使所取填充 \(W\) 紧，从而成为定义中的配边。若只知道 \(M\) 紧，负子水平集仍可能无界，例如 \(f(x)=1-|x|^2\)；此时不能直接把 \(\{f\le0\}\) 称作紧配边，须另选填充并核对标架。
\end{solution}

\begin{exercise}\label{g2:ex:standard-sphere-surgery}
	设 \(0\le p\le n-2\)，把 \(S^n\) 写成磨圆后的 \(\partial(D^{p+1}\times D^{n-p})\)。沿其中的 \(S^p\times D^{n-p}\) 使用乘积标架作手术，证明结果微分同胚于 \(S^{p+1}\times S^{n-p-1}\)，并核对 Euler 示性数的改变量。
\end{exercise}
\begin{solution}
	乘积圆盘磨圆后是 \(D^{n+1}\)，故其边界为标准 \(S^n\)。取走第一个面 \(S^p\times\operatorname{int}D^{n-p}\)，余下的是第二个面 \(D^{p+1}\times S^{n-p-1}\)。手术再接上一个这样的面，粘合映射在边界两个因子上都是恒等，于是结果是
	\[
		(D^{p+1}\cup_{S^p}D^{p+1})\times S^{n-p-1}
		\cong S^{p+1}\times S^{n-p-1}.
	\]
	这里两个圆盘的双倍采用通常的光滑领圈，磨圆得到标准球面。Euler 示性数之差为
	\[
		(1+(-1)^{p+1})(1+(-1)^{n-p-1})-(1+(-1)^n)
		=(-1)^{n-p-1}-(-1)^p,
	\]
	与\cref{g2:prop:surgery-euler}一致。特别地，\(S^4\) 沿标准圆周的手术给出 \(S^2\times S^2\)，其 Euler 示性数由 \(2\) 增至 \(4\)。
\end{solution}

\begin{exercise}\label{g2:ex:kill-fundamental-group-by-surgery}
	证明每个连通闭定向光滑 \(n\) 维流形 \(M\)，\(n\ge4\)，都与某个单连通闭定向光滑 \(n\) 维流形有向配边。请说明有限性、圆周嵌入、法标架及所得基本群四个环节。
\end{exercise}
\begin{solution}
	紧流形具有有限 CW 同伦型，故 \(\pi_1(M)\) 有有限生成集 \(g_1,\ldots,g_r\)。将代表这些元素的环路作光滑逼近，再对有限个圆周的并使用一般位置。当 \(n\ge4\) 时，圆周自身或两圆周之间的双点预期维数为 \(2-n<0\)，而微分为零的情形也可由一般位置避开。因此可用互不相交的嵌入圆周代表这些自由同伦类；基点道路的改变只产生共轭，不影响它们生成的正规子群。

	定向 \(M\) 和各圆周后，法丛获得定向。每个定向秩 \(n-1\) 实向量丛在圆周上平凡：沿割开圆周后的区间选标架，两端的差为 \(\operatorname{GL}^+(n-1,\mathbb R)\) 中一个矩阵；这个群道路连通，沿区间用一条道路修正标架，即能使两端吻合。管状邻域定理遂给出互不相交的 \(S^1\times D^{n-1}\) 邻域。

	对这些圆周逐个作定向兼容的 \(1\) 维手术。由于 \(1\le n-3\)，\cref{g2:prop:surgery-fundamental-group}适用，最终基本群为
	\[
		\pi_1(M)/\langle\!\langle g_1,\ldots,g_r\rangle\!\rangle=\{1\}.
	\]
	每步保持连通，并由其定向手术迹提供有向配边。把有限个迹沿共同边界粘合，就得到从 \(M\) 到最终单连通流形的有向配边。
\end{solution}

\begin{exercise}\label{g2:ex:surgery-trace-not-h-cobordism}
	设 \(M^n\) 闭且连通，\(n\ge4\)，在 \(M\) 上作一次 \(1\) 维手术，迹记为 \(W\)。证明 \(H_2(W,M;\mathbb Z)\cong\mathbb Z\)，从而 \(W\) 不是 \(h\)-配边。若附着圆周在 \(M\) 中可缩，解释为什么两端基本群同构仍不能改变这个结论。
\end{exercise}
\begin{solution}
	由\cref{g2:lem:handle-core-retraction}，流形对 \((W,M)\) 相对同伦等价于 \((M\cup_\alpha D^2,M)\)。相对胞腔链复形仅在维数 \(2\) 有一个 \(\mathbb Z\)，其他维数为零，所以其微分全为零，得到
	\[
		H_j(W,M;\mathbb Z)=
		\begin{cases}\mathbb Z,&j=2,\\0,&j\ne2.\end{cases}
	\]
	如果 \(M\hookrightarrow W\) 是同伦等价，所有相对同调群都应为零，矛盾。因此这个迹不是 \(h\)-配边。

	附着圆周可缩时，它生成的正规子群平凡，\cref{g2:prop:surgery-fundamental-group}确实给出 \(\pi_1(M_\varphi)\cong\pi_1(M)\)。然而基本群只检测一维环路，而新附着的核心是二维胞腔，相对 \(H_2\) 仍然记录它。故两端基本群相同不能代替 \(h\)-配边所要求的包含映射同伦等价。
\end{solution}

\begin{exercise}\label{g2:ex:trivial-whitehead-group}
	利用整数矩阵的初等变换证明 \(\operatorname{Wh}(\{1\})=0\)，并据此从 \(s\)-配边定理推出单连通的 \(h\)-配边定理。
\end{exercise}
\begin{solution}
	此时群环是 \(\mathbb Z\)，Whitehead 群为 \(K_1(\mathbb Z)/\langle[-1]\rangle\)。任一可逆整数矩阵 \(A\) 的行列式为 \(\pm1\)。其第一列各整数的最大公因数为 \(1\)，否则该公因数也整除行列式。使用整数行加法、行交换和变号，Euclid 算法可把第一列变为 \((1,0,\ldots,0)^{\mathsf T}\)；随后列加法消去第一行的其余元素。对剩余方阵归纳，得到单位矩阵。

	行加法和列加法对应初等矩阵。行交换可由初等矩阵及一个对角 \(-1\) 实现，变号本身也是对角 \(-1\)。故 \(K_1(\mathbb Z)\) 中每个类都由 \([-1]\) 生成，再除掉它便得到零群。

	现在设 \((W;M_0,M_1)\) 是两端闭连通、\(\dim W\ge6\) 的光滑 \(h\)-配边，且 \(M_0\) 单连通。其扭子只能落在零群 \(\operatorname{Wh}(\{1\})\)，所以必为零。\Cref{g2:thm:s-cobordism}给出固定 \(M_0\) 的微分同胚 \(M_0\times[0,1]\cong W\)，即所需 \(h\)-配边结论。
\end{solution}

\begin{exercise}\label{g2:ex:atlas-change-not-exotic}
	在实线的底层拓扑空间上，分别取由全局坐标 \(x\mapsto x\) 和 \(x\mapsto x^3\) 生成的最大光滑图册 \(\mathcal A\)、\(\mathcal B\)。证明这两个最大图册不同，但两个光滑流形微分同胚。进一步判断底空间的恒等映射在两个方向上是否光滑。
\end{exercise}
\begin{solution}
	坐标变换 \(x\mapsto x^3\) 光滑，但其逆映射 \(y\mapsto y^{1/3}\) 在原点不可微。因此两张坐标不属于同一个光滑图册，所生成的最大图册不同。

	设 \(h\colon(\mathbb R,\mathcal B)\rightarrow(\mathbb R,\mathcal A)\) 为底空间映射 \(h(x)=x^3\)。源端使用坐标 \(\beta(x)=x^3\)，靶端使用标准坐标 \(\alpha(x)=x\)，则其坐标表达为 \(\alpha h\beta^{-1}(y)=y\)。逆映射在相反的坐标中也表达为恒等映射，故 \(h\) 是微分同胚。

	底空间的恒等映射 \(\operatorname{id}\colon(\mathbb R,\mathcal A)\rightarrow(\mathbb R,\mathcal B)\) 的坐标表达是 \(x\mapsto x^3\)，所以光滑；反方向的坐标表达是立方根，所以不光滑。这说明恒等映射失去微分同胚性质，并不能排除其他重新识别底空间的微分同胚。
\end{solution}

\begin{exercise}\label{g2:ex:homotopy-seven-sphere-classes}
	设 \(\Sigma\) 是与 \(S^7\) 同胚的闭光滑流形，并任选一个定向。证明它的全部正次数 Stiefel--Whitney 类、Pontryagin 类以及切丛 Euler 类均为零。说明为什么这些类不能区分 \(M_{2,-1}\) 与标准 \(S^7\)。
\end{exercise}
\begin{solution}
	同胚给出 \(H^q(\Sigma;\mathbb Z)=H^q(S^7;\mathbb Z)\) 及模二版本，因此 \(H^q(\Sigma;\mathbb Z/2)=0\) 对 \(1\le q\le6\) 成立。于是 \(w_q(\mathrm T\Sigma)=0\) 对这些 \(q\) 成立。\(q>7\) 时，Stiefel--Whitney 类由秩条件为零。

	剩下 \(w_7\)。由\cref{g2:cor:tangent-euler-characteristic}，
	\[
		\langle e_{\mathrm{Th}}(\mathrm T\Sigma),[\Sigma]\rangle
		=\chi(\Sigma)=1+(-1)^7=0.
	\]
	\(H^7(\Sigma;\mathbb Z)\cong\mathbb Z\)，且在定向基本类上的求值是同构，故 \(e_{\mathrm{Th}}(\mathrm T\Sigma)=0\)。再由\cref{g2:thm:top-class-euler}，\(w_7\) 是该类的模二约化，因而也为零。

	Pontryagin 类 \(p_i\) 位于 \(H^{4i}(\Sigma;\mathbb Z)\)。\(i=1\) 时目标群为零，\(i\ge2\) 时次数超过维数，仍为零。这些消失结论同时适用于标准 \(S^7\) 和\cref{g2:thm:milnor-exotic-seven-sphere}中的 \(M_{2,-1}\)，但后两者不微分同胚。因而这组示性类在此没有足够的信息完成光滑分类。
\end{solution}

\begin{exercise}\label{g2:ex:exotic-r-four-proper-coordinate}
	设 \(R\) 是异光滑四维欧氏空间。证明不存在固有的单射光滑浸入 \(F\colon R\rightarrow\mathbb R^4\)。这里固有指任意紧集的原像紧。
\end{exercise}
\begin{solution}
	若这样的 \(F\) 存在，源、靶维数相同，浸入条件说明每点的微分都是线性同构。逆函数定理给出局部光滑逆，所以 \(F(R)\) 是开集。

	固有性又保证 \(F(R)\) 是闭集。具体地，设 \(y_j=F(x_j)\) 收敛于 \(y\)。集合 \(C=\{y\}\cup\{y_j:j\ge1\}\) 紧，故 \(F^{-1}(C)\) 紧。取 \(x_j\) 的收敛子列，极限记作 \(x\)；连续性给出 \(F(x)=y\)，所以 \(y\in F(R)\)。欧氏空间可度量，逐列闭性即闭性。

	\(F(R)\) 非空且开闭，而 \(\mathbb R^4\) 连通，故 \(F\) 满射。结合单射性，它是双射；各点的局部光滑逆在交叠处由单射性一致，组成整体光滑逆。因此 \(F\) 是微分同胚，与 \(R\) 异光滑矛盾。
\end{solution}

\begin{exercise}\label{g2:ex:exotic-r-four-one-point}
	设 \(R\) 是异光滑四维欧氏空间，\(R^+=R\cup\{\infty\}\) 是其单点紧化。
	\begin{enumerate}
		\item 证明 \(R^+\) 在拓扑上同胚于 \(S^4\)。若 \(R\) 的光滑结构能延拓成 \(R^+\) 上的光滑流形结构，证明所得闭流形不与标准 \(S^4\) 微分同胚。
		\item 假设进一步存在紧集 \(K\subset R\)，没有任何光滑嵌入的 \(S^3\subset R\) 能把 \(K\) 与无穷远端分开。证明在这种情况下，\(R\) 的光滑结构根本不能延拓过 \(\infty\)。
	\end{enumerate}
\end{exercise}
\begin{solution}
	1. 取同胚 \(h\colon R\rightarrow\mathbb R^4\)。同胚及其逆都将紧集映为紧集，因而令 \(h(\infty)=\infty\)，就得到单点紧化之间的同胚 \(R^+\rightarrow(\mathbb R^4)^+\cong S^4\)。若某个延拓后的光滑流形与标准 \(S^4\) 微分同胚，限制这个微分同胚到 \(R=R^+\setminus\{\infty\}\)，得到 \(R\cong_{\mathrm{DIFF}}S^4\setminus\{p\}\)。绕球面的一个旋转把 \(p\) 移到北极，再作球极投影，就给出右端与标准 \(\mathbb R^4\) 的微分同胚，产生矛盾。

	2. 若有延拓，\(\infty\) 有光滑坐标邻域。由于 \(K\) 紧且不含 \(\infty\)，可在这个邻域内取一个以 \(\infty\) 为中心、与 \(K\) 不交的小闭坐标四维球 \(B\)。其边界是 \(R\) 内的光滑嵌入 \(S^3\)。\(K\) 位于 \(R^+\setminus B\)，而 \(\infty\) 位于 \(\operatorname{int}B\)；从 \(K\) 到无穷远的任何连续路径都要穿过 \(\partial B\)。这与假设的不可分隔性质矛盾。
\end{solution}

% !TeX root = ../../Book2.tex
% 纲要标题：## 第24章*　拓扑 K 理论、Bott 周期性与指标的拓扑侧
% 纲要位置：几何学计划纲要.md，第 1923 行。
% 正文状态：初稿；尚未正式审核。

% 纲要内容开始
% 
% > 本章目标与证明范围：以有限 CW 空间上的复 K 理论为完整主模块，基本乘积定理的证明可在紧 Hausdorff 参数空间上进行。Bott 采用多项式粘接的具体证明路线；实 KO 的八周期和一般指标定理作为准确外部输入，不以复数特例混称全理论。
% 
% 纲要内容结束

\OptionalChapter{拓扑 K 理论、Bott 周期性与指标的拓扑侧}{g2:ch:24}
\BookChapterNumber{g2:ch:24}

\begin{chapterabstract}
	向量丛的直和可以做加法，却通常不能直接做减法。群完备化把它变成虚丛的群，张量积进一步给出环结构。本章建立相对群、悬挂和长正合列，并用多项式粘接完整证明复 Bott 周期性：先把圆周上的连续粘接函数逼近为 Laurent 多项式，再把高次矩阵线性化，最后由单位圆内外的谱子丛得到基本乘积定理及其逆映射。陈特征把有限 CW 空间的有理 K 群与偶奇上同调联系起来；椭圆主符号则给出余切丛上的紧支 K 类，通向解析指标和拓扑指标的比较。
\end{chapterabstract}

\begin{learninggoals}
	\begin{itemize}
		\item 从向量丛的直和构造 \(K^0\)，区分实际同构、稳定同构与虚丛相等，并计算低维球面的例子。
		\item 用映射锥和三元组理解相对 K 群，构造悬挂、连接映射及长正合列。
		\item 完成粘接函数的多项式逼近、线性化及连续谱分解，构造 \(K^0(X\times S^2)\) 分解的两个方向。
		\item 推出复 K 理论的二周期、球面系数及有限 CW 算例，辨认实 KO 八周期使用的不同输入。
		\item 证明陈特征的乘法性与有理同构，理解椭圆符号、紧支类和指标推前各自记录的信息。
	\end{itemize}
\end{learninggoals}

第19章把球面上的向量丛描述为沿赤道粘接两块平凡丛，第21章又用陈类测量复丛的扭曲。现在考虑一个更具体的问题：\(S^2\) 上的 Hopf 线丛 \(H\) 与平凡线丛有什么“差”？两者的秩相同，却有不同的第一陈类。若能把 \([H]-[\underline{\mathbb C}]\) 当作一个可运算的对象，就能把秩与扭曲分开，并在直和、张量积和拉回中追踪这份差异。这正是 K 理论引入虚丛的出发点。

\ProseParagraphBreak

更令人惊讶的是，把任意紧参数空间 \(X\) 乘上一张二球面，并不会产生难以控制的全新类型。令 \(\beta=[H]-1\)，本章将证明每个虚丛都能唯一写成 \(a+b\beta\)，其中 \(a,b\in K^0(X)\)。\(a\) 是在球面一点上的限制，\(b\) 则保留粘接函数绕单位圆变化时的整族信息。对单个矩阵而言，单位圆内特征值的个数是一个整数；对连续矩阵族而言，相应广义特征空间拼成一条向量丛。这一步把线性代数与全局拓扑联系起来，也是证明中最值得看清的构造。

\ProseParagraphBreak

主要准备是第19章的丛粘接、同伦不变性和平凡丛补空间，第1章的映射锥及同伦延拓性质，以及第21章的陈类和分裂原理。长证明中用到的复矩阵谱投影会在原处由 Jordan 块和圆周积分展开。最后一节把本章所得用于椭圆算子的主符号；解析 Fredholm 性及一般指标等式在使用处给出条件和正式文献，实 KO 的八周期也在相应陈述后注明来源。

% !TeX root = ../../Book2.tex
% 纲要标题：### 24.1 向量丛的群完备化
% 纲要位置：几何学计划纲要.md，第 1927 行。
% 正文状态：初稿；尚未正式审核。

% 纲要内容开始
% 
% * K⁰、约化 K⁰ 与相对群
% * 稳定同构和虚丛
% * 粘合构造与球面 K 群
% * 环结构和外积
% 
% 纲要内容结束

\section{向量丛的群完备化}
\label{g2:sec:2401}

第19章把向量丛分类为到 Grassmann 空间的映射，却还没有把全部丛组织成一个可作减法的对象。直和是现成的加法，但一张正秩丛不可能通过直和得到零丛，因此通常没有逆元。本节把“比较两张丛的差”作为新的对象：增加同一张辅助丛不会改变差，稳定同构由此成为等式，张量积则成为乘法。后面正合列与周期性所计算的，正是这套稳定的丛信息。

\subsection{\texorpdfstring{\(K^0\)、约化 \(K^0\) 与相对群}{K⁰、约化 K⁰ 与相对群}}
\label{g2:subsec:2401:01}
本节的底空间先取紧 Hausdorff 空间 \(X\)，丛取有限秩连续复向量丛。允许秩随连通分支变化；秩函数局部常值，紧性保证它只取有限多个值。记 \(\operatorname{Vect}_{\mathbb C}(X)\) 为这些丛的同构类所成集合，包含零丛。直和使它成为交换幺半群：运算满足结合律、交换律并有零元，但尚未保证逆元，也不能预先作消去。

若把 \((E,F)\) 理解为形式差 \(E-F\)，那么 \((E\oplus G,F\oplus G)\) 应表示同一个差。不过只允许这种直接增添，还不足以让所有形式差自然组成群；下面的关系把所需稳定比较一次写全。

\begin{definition}\label{g2:def:topological-k-zero}
	对紧 Hausdorff 空间 \(X\)，在复向量丛对上规定
	\[
		(E,F)\sim(E',F')
		\quad\Longleftrightarrow\quad
		E\oplus F'\oplus G\cong E'\oplus F\oplus G
	\]
	对某张复向量丛 \(G\to X\) 成立。所得等价类集合记作 \(K^0(X)\)，称为 \(X\) 的复\term[tuopu-K-qun]{拓扑 K 群}[topological K-group]。丛对的类写作 \([E]-[F]\)，其中 \([E]\) 表示 \((E,0)\) 的类。运算为
	\[
		([E]-[F])+([E']-[F'])=[E\oplus E']-[F\oplus F'].
	\]
	这样的形式差称为\term[xucong]{虚丛}[virtual bundle]。\symidx{\(K^0(X)\)：复向量丛的群完备化}
\end{definition}

\begin{proposition}\label{g2:prop:k-group-completion}
	上述关系为等价关系，\(K^0(X)\) 是交换群，其零元为 \([0]\)，\([E]-[F]\) 的逆元为 \([F]-[E]\)。映射
	\[
		j:\operatorname{Vect}_{\mathbb C}(X)\longrightarrow K^0(X),\qquad E\longmapsto[E]
	\]
	具有如下泛性质：对任意阿贝尔群 \(A\) 及保持直和、零元的映射 \(a:\operatorname{Vect}_{\mathbb C}(X)\to A\)，存在唯一群同态 \(\overline a:K^0(X)\to A\) 满足 \(\overline a\circ j=a\)。
\end{proposition}
\begin{proof}
	自反性可取 \(G=0\)，对称性由交换等式两侧得到。若比较 \((E,F)\)、\((E',F')\) 使用辅助丛 \(G\)，比较 \((E',F')\)、\((E'',F'')\) 使用 \(G'\)，则先后代入两份同构可得
	\[
		E\oplus F''\oplus(F'\oplus G\oplus G')
		\cong E''\oplus F\oplus(F'\oplus G\oplus G').
	\]
	这证明传递性，并显示为什么关系中保留辅助丛。对两份稳定同构取直和，就得到它们之和的稳定同构，所以加法良定义。结合与交换来自丛直和的相应同构，而 \((E\oplus F,F\oplus E)\sim(0,0)\)，故所列逆元成立。

	定义 \(\overline a([E]-[F])=a(E)-a(F)\)。若两对等价，对定义中的同构应用 \(a\)，在群 \(A\) 中消去 \(a(G)\)，便得到相同的差。因此它良定义且可加。每个元素都是两个 \(j\) 的值之差，所以这一要求也确定了唯一的 \(\overline a\)。相应交换图为
	\[
		\begin{tikzcd}[column sep=large]
			\operatorname{Vect}_{\mathbb C}(X)\arrow[r,"j"]\arrow[dr,"a"']&K^0(X)\arrow[d,dashed,"\exists!\,\overline a"]\\
			&A.
		\end{tikzcd}
	\]
\end{proof}

这个构造称为幺半群的\term[qun-wanbeihua]{群完备化}[group completion]。它把所有可加的丛不变量统一到同一个群上。例如，单点上的复向量空间只由维数决定，故 \(K^0(\{*\})\cong\mathbb Z\)。对一般 \(X\)，逐点取秩给出
\[
	\operatorname{rk}:K^0(X)\longrightarrow C(X,\mathbb Z),\qquad
	\operatorname{rk}([E]-[F])(x)=\dim E_x-\dim F_x,
\]
其中 \(C(X,\mathbb Z)\) 是连续整数值函数群。这个整数差称为虚秩。连通时它是单个整数；不连通时各分支可以不同。

\begin{definition}\label{g2:def:reduced-relative-k}
	对带基点的紧 Hausdorff 空间 \((X,x_0)\)，定义\term[yuehua-K-qun]{约化 K 群}[reduced K-group]
	\[
		\widetilde K^0(X)=\ker\bigl(i^*:K^0(X)\longrightarrow K^0(\{x_0\})=\mathbb Z\bigr),
	\]
	其中 \(i\) 为基点包含，\(i^*\) 取纤维的虚维数。对有限 CW 对 \((X,A)\)，定义\term[xiangdui-K-qun]{相对 K 群}[relative K-group]
	\[
		K^0(X,A)=\widetilde K^0(X/A).
	\]
	若 \(A\ne\varnothing\)，商空间以压缩后的 \(A\) 为基点；若 \(A=\varnothing\)，约定此处的 \(X/A\) 为 \(X_+=X\sqcup\{*\}\)。\symidx{\(\widetilde K^0(X)\)：基点纤维虚秩为零的 K 群}\symidx{\(K^0(X,A)\)：有限 CW 对的相对 K 群}
\end{definition}

记 \(\varepsilon^r=X\times\mathbb C^r\)，并以 \(1=[\varepsilon^1]\) 表示平凡线丛的类。基点取秩有分裂 \(m\mapsto m1\)，故
\begin{equation}\label{g2:eq:k-rank-splitting}
	K^0(X)\cong\mathbb Z\oplus\widetilde K^0(X),\qquad
	x\longmapsto\bigl(\operatorname{rk}_{x_0}x,\ x-(\operatorname{rk}_{x_0}x)1\bigr).
\end{equation}
当 \(X\) 不连通时，约化只消去所选基点处的虚秩，其他分支上的秩仍可非零。例如 \(K^0(S^0)\cong\mathbb Z^2\)，选其中一点为基点后有 \(\widetilde K^0(S^0)\cong\mathbb Z\)。又因为 \(K^0(X_+)\cong K^0(X)\oplus\mathbb Z\)，空子空间约定给出 \(K^0(X,\varnothing)=K^0(X)\)，而 \(K^0(X,X)=0\)。下一节将把商空间定义改写为 \(A\) 上带有指定同构的数据，从而构造相对映射和连接同态。

\subsection{稳定同构和虚丛}
\label{g2:subsec:2401:02}
形式差的两个成员可以简化为“一张实际丛减去一个平凡丛”。其几何依据是第19章的补丛构造：由\cref{g2:prop:bundle-embedding-classifier}，紧基底上的有限秩丛可嵌入某个 \(\varepsilon^N\)，正交补就是所需补丛。局部常秩情形按有限多个开闭秩层分别构造，再把环境维数增到同一个 \(N\)，所以同样成立。

\begin{proposition}\label{g2:prop:k-stable-equality}
	设 \(X\) 紧且 Hausdorff。每个 \(x\in K^0(X)\) 都可写成 \([E]-[\varepsilon^N]\)。对两张复向量丛 \(E,F\to X\)，有
	\[
		[E]=[F]\text{ 于 }K^0(X)
		\quad\Longleftrightarrow\quad
		E\oplus\varepsilon^N\cong F\oplus\varepsilon^N
	\]
	对某个 \(N\ge0\) 成立。因此群中的实际丛等式恰好是第19章所说的稳定同构。
\end{proposition}
\begin{proof}
	给定 \([E]-[F]\)，取补丛 \(F'\) 使 \(F\oplus F'\cong\varepsilon^N\)。在形式差的两项中加上 \(F'\)，便得
	\[
		[E]-[F]=[E\oplus F']-[\varepsilon^N].
	\]
	若 \([E]=[F]\)，定义给出某个 \(G\) 使 \(E\oplus G\cong F\oplus G\)。再取 \(G'\) 使 \(G\oplus G'\cong\varepsilon^N\)，在两边加 \(G'\) 即得稳定同构。反过来，存在这样的稳定同构时，在等价关系中直接取 \(G=\varepsilon^N\) 即可。
\end{proof}

特别地，在带基点的空间上，每个 \(\widetilde K^0(X)\) 的类可写成 \([E]-[\varepsilon^N]\)，且 \(\operatorname{rank}E_{x_0}=N\)。若 \(X\) 连通，\(E\) 在所有点的秩都为 \(N\)。对连通 \(X\)，两项 \([E]-\operatorname{rank}E\) 与 \([F]-\operatorname{rank}F\) 相同，正是说两丛补加适当数量的平凡线以后同构；补加数量可不同，是因为已经减掉各自的秩。

\ProseParagraphBreak

虚丛的“负号”来自群运算。例如 \(-[E]=[0]-[E]\) 的虚秩是负数，而对偶丛 \(E^*\) 与 \(E\) 的实际秩相同。因而对偶通常不是直和的逆元。另一方面，若一个短正合列
\[
	0\longrightarrow E'\longrightarrow E\longrightarrow E''\longrightarrow0
\]
由 \(X\) 上有限秩复向量丛组成，那么 Hermitian 度量使 \(E'\) 的正交补映到 \(E''\) 的每根纤维上成为同构，故 \(E\cong E'\oplus E''\)，从而 \([E]=[E']+[E'']\)。这里讨论连续复向量丛；所取正交补是连续的，论证并没有要求它保持额外的全纯结构。

\begin{proposition}\label{g2:prop:k-pullback-homotopy}
	连续映射 \(f:X\to Y\) 在紧 Hausdorff 空间之间诱导群同态
	\[
		f^*:K^0(Y)\longrightarrow K^0(X),\qquad
		f^*([E]-[F])=[f^*E]-[f^*F].
	\]
	它满足 \((g\circ f)^*=f^*g^*\)、\(\operatorname{id}^*=\operatorname{id}\)，同伦的映射诱导相同同态。带基点映射还保持约化群，有限 CW 对的映射通过商空间诱导相对群同态。
\end{proposition}
\begin{proof}
	拉回保持丛同构与直和，所以保持定义中的等价关系和群运算。复合拉回与逐次拉回之间的纤维识别给出函子等式。若 \(f_0\simeq f_1\)，把任一 \(Y\) 上的丛沿同伦拉回到 \(X\times I\)，\cref{g2:thm:bundle-homotopy-invariance}说明两端限制同构，故它们在 \(K^0\) 中的类相同。基点纤维与商空间映射的自然性分别给出最后两项。
\end{proof}

可缩的非空紧 Hausdorff 空间于是有 \(K^0(X)\cong\mathbb Z\)，且任一基点下 \(\widetilde K^0(X)=0\)。具体地，收缩把每张丛同构到一个常值纤维丛；从此例开始，丛的稳定计算与空间的同伦类型联系起来。

\subsection{粘合构造与球面 K 群}
\label{g2:subsec:2401:03}
球面可分为两个可缩半球，丛在各半球都能平凡化，所以差异只留在赤道的粘合函数中。\cref{g2:thm:clutching-classification}已经证明，对 \(m\ge2\)，秩 \(r\) 复丛的同构类由
\[
	\pi_{m-1}(\operatorname{U}(r))
\]
分类。这里先用极分解把 \(\operatorname{GL}_r(\mathbb C)\) 缩到 \(\operatorname{U}(r)\)，再把赤道基点处的矩阵乘回单位矩阵，把自由同伦化为基点同伦。若 \(g\) 是粘合函数，则加一条平凡线把它变为 \(\operatorname{diag}(g,1)\)。

\begin{proposition}\label{g2:prop:k-clutching-spheres}
	对 \(m\ge2\)，粘合构造给出自然群同构
	\[
		\varinjlim_r\pi_{m-1}(\operatorname{U}(r))
		\xrightarrow{\ \cong\ }\widetilde K^0(S^m),
		\qquad[g]\longmapsto[E_g]-r.
	\]
	直接极限沿 \(g\mapsto\operatorname{diag}(g,1)\) 取；两项代表相同，是指它们在某个共同的有限秩阶段同伦。
\end{proposition}
\begin{proof}
	上一命题把约化类写成 \([E]-r\)，球面连通保证 \(\operatorname{rank}E=r\)，故粘合分类给出满射。若两粘合丛给出相同约化类，\cref{g2:prop:k-stable-equality}说明补加平凡线后两丛同构。再用固定秩的粘合分类，就得到稳定后的粘合函数同伦。这证明单射。

	直和对应分块矩阵 \(\operatorname{diag}(g,h)\)。把两块先稳定到同阶，令
	\[
		R_t=\begin{pmatrix}\cos t\,I&-\sin t\,I\\\sin t\,I&\cos t\,I\end{pmatrix},\qquad
		0\le t\le\pi/2.
	\]
	则 \(\operatorname{diag}(g,I)R_t\operatorname{diag}(I,h)R_t^{-1}\) 从 \(\operatorname{diag}(g,h)\) 连到 \(\operatorname{diag}(gh,I)\)，且在赤道基点始终为单位矩阵。拓扑群的球面映射中，逐点乘法与同伦群的拼接相同：把代表写在立方体上，两种运算具有同一单位元并满足逐坐标的交换律
	\[
		(f_1*f_2)(g_1*g_2)=(f_1g_1)*(f_2g_2),
	\]
	其中 \(*\) 表示沿同一个坐标拼接。将某两项取常值单位映射，就在同伦类上得到逐点乘积等于拼接。因此分块和表示同伦群的和，所构造双射为群同构。
\end{proof}

低维计算已经可以完成。第19章说明圆周上的复丛全部平凡，故 \(K^0(S^1)\cong\mathbb Z\)、\(\widetilde K^0(S^1)=0\)。对 \(S^2\)，\cref{g2:prop:classical-low-homotopy}给出
\[
	\det_*:\pi_1(\operatorname{U}(r))\xrightarrow{\ \cong\ }\pi_1(S^1)=\mathbb Z,
\]
且增加单位矩阵块保持这个同构。因此 \(\widetilde K^0(S^2)\cong\mathbb Z\)，分类整数就是 \(\det g\) 的绕数。由\cref{g2:prop:first-classes-determinants,g2:prop:chern-transition-comparison}，它也等于 \(\langle c_1(E_g),[S^2]\rangle\)，只须保持同一组赤道方向与坐标约定。

下面固定与第21章一致的正生成元。将 \(S^2\) 识别为复定向的 \(\mathbb CP^1\)，取
\[
	z=z_1/z_0,\qquad w=z_0/z_1=1/z,\qquad H=\gamma^*.
\]
对典范线丛 \(\gamma\)，局部标架为 \(e_z=(1,z)\)、\(e_w=(w,1)=z^{-1}e_z\)；对偶标架因而满足 \(e_w^*=z e_z^*\)。沿用第19章的关系 \((x,v)_+\sim(x,g(x)v)_-\)，取 \(D_+\) 为 \(w\) 圆盘、\(D_-\) 为 \(z\) 圆盘，并以 \(z=e^{\mathrm i\theta}\)、\(\theta\) 增大参数化赤道，则 \(g(z)=z\) 的粘合丛正是 \(H\)。第21章已经核定 \(c_1(H)=h\)、\(\langle h,[\mathbb CP^1]\rangle=1\)，所以
\begin{equation}\label{g2:eq:k-two-sphere-additive}
	\widetilde K^0(S^2)=\mathbb Z\,([H]-1),\qquad
	K^0(S^2)=\mathbb Z\,1\oplus\mathbb Z\,([H]-1).
\end{equation}
秩 \(r\ge1\)、第一陈数为 \(d\) 的丛满足 \([E]=r+d([H]-1)\)。事实上它还同构于 \(H^{\otimes d}\oplus\varepsilon^{r-1}\)：这张丛的粘合矩阵 \(\operatorname{diag}(z^d,1,\ldots,1)\) 与 \(g\) 有相同行列式绕数，故在 \(\operatorname{U}(r)\) 中同伦。负数 \(d\) 按对偶张量幂理解。

\begin{example}\label{g2:exm:k-sphere-tangent-line}
	\(\mathbb CP^1\) 的复切线丛满足 \(c_1(\mathrm T\mathbb CP^1)=2h\)，见式\eqref{g2:eq:complex-projective-tangent-classes}，因此
	\[
		[\mathrm T\mathbb CP^1]-1=2([H]-1)\ne0.
	\]
	球面的实切丛加一条实平凡线可平凡化，与这个结论并不冲突：这里用的是复一维切丛，稳定操作只允许加入复平凡方向。实、复理论的稳定信息不同，后面讨论周期性时也须保留基域。
\end{example}

\subsection{环结构和外积}
\label{g2:subsec:2401:04}
直和给出加法，张量积给出另一种运算。它在实际丛上满足 \((E\oplus F)\otimes G\cong(E\otimes G)\oplus(F\otimes G)\)，所以可以把张量积按分配律延到虚丛。

\begin{proposition}\label{g2:prop:k-ring-structure}
	对紧 Hausdorff 空间 \(X\)，下式把 \(K^0(X)\) 变为含幺交换环：
	\[
		([E]-[F])([G]-[J])
		=[E\otimes G\oplus F\otimes J]-[E\otimes J\oplus F\otimes G].
	\]
	单位元为 \(1=[\varepsilon^1]\)。连续映射的拉回是含幺环同态；基点取秩也是含幺环同态，因而 \(\widetilde K^0(X)\) 是 \(K^0(X)\) 的理想。
\end{proposition}
\begin{proof}
	固定一张实际丛 \(G\)，映射 \(E\mapsto[E\otimes G]\) 对直和可加，故由群完备化的泛性质唯一延拓为 \(K^0(X)\to K^0(X)\) 的同态。该同态又对 \(G\) 的直和可加，因此对第二个变量再次用泛性质，得到唯一的双线性乘法。展开两份差，就得到所列公式，也同时证明它与代表选择无关。

	实际丛的张量积有结合、交换及单位同构，这些等式对形式差双线性延拓，故得到环公理。拉回保持直和、张量积及平凡线丛，因而保持乘法和单位。纤维维数满足 \(\dim(E_x\otimes G_x)=\dim E_x\dim G_x\)，所以取秩也乘法相容，其核便为理想。
\end{proof}

\begin{definition}\label{g2:def:k-external-product}
	设 \(X,Y\) 为紧 Hausdorff 空间，\(p_X,p_Y\) 为积空间的投影。对 \(x\in K^0(X)\)、\(y\in K^0(Y)\)，定义\term[K-lilun-waiji]{K 理论外积}[external product in K-theory]
	\[
		x\boxtimes y=p_X^*x\,p_Y^*y\in K^0(X\times Y).
	\]
	对实际丛，它由 \(E\boxtimes F=p_X^*E\otimes p_Y^*F\) 表示。\symidx{\(x\boxtimes y\)：K 理论外积}
\end{definition}

外积可加于每个变量，且对连续映射 \(f:X'\to X\)、\(g:Y'\to Y\)，有
\begin{equation}\label{g2:eq:k-external-naturality}
	(f\times g)^*(x\boxtimes y)=f^*x\boxtimes g^*y.
\end{equation}
这是因为 \(p_X\circ(f\times g)=f\circ p_{X'}\)，另一投影亦然，再使用拉回的乘法相容性。若 \(X=Y\)，沿对角映射 \(\Delta:X\to X\times X\) 拉回便得到原来的乘法：\(\Delta^*(x\boxtimes y)=xy\)。交换两因子的映射则把 \(x\boxtimes y\) 送到 \(y\boxtimes x\)；当前两项均为零次，尚无分次符号。

\ProseParagraphBreak

若 \(X,Y\) 有基点，\(x\in\widetilde K^0(X)\)、\(y\in\widetilde K^0(Y)\)，那么 \(x\boxtimes y\) 在 \(X\times\{y_0\}\) 与 \(\{x_0\}\times Y\) 上的限制均为零。这为把外积看作压缩积 \(X\wedge Y\) 上的类提供了必要条件；下一节将通过相对正合性构造相应的约化外积，并证明其与这里的绝对外积相容。

在 \(S^2\) 的例子中，我们已得到 \(1\) 与 \([H]-1\) 组成的加法基。环结构还要求计算 \(([H]-1)^2\)，而积空间 \(X\times S^2\) 上的丛更需要知道两项系数是否都能从 \(K^0(X)\) 取得。这两件事将在\secref{g2:sec:2403}中由多项式粘合构造一起解决。
% 纲要条目排为编号小节标题，并进入目录与书签；说明段仍保留于纲要注释中。

% !TeX root = ../../Book2.tex
% 纲要标题：### 24.2 正合性和悬挂
% 纲要位置：几何学计划纲要.md，第 1934 行。
% 正文状态：未撰写；本文件目前是框架。

% 纲要内容开始
% 
% * 映射锥与相对 K 群
% * 悬挂定义负次数
% * 长正合列和切除
% * 乘积结构的自然性
% 
% 纲要内容结束

\section{正合性和悬挂}
\label{g2:sec:2402}

待撰写

\subsection{映射锥与相对 K 群}
\label{g2:subsec:2402:01}
待撰写

\subsection{悬挂定义负次数}
\label{g2:subsec:2402:02}
待撰写

\subsection{长正合列和切除}
\label{g2:subsec:2402:03}
待撰写

\subsection{乘积结构的自然性}
\label{g2:subsec:2402:04}
待撰写
% 纲要条目排为编号小节标题，并进入目录与书签；说明段仍保留于纲要注释中。

% !TeX root = ../../Book2.tex
% 纲要标题：### 24.3 复 Bott 周期性的多项式粘接证明
% 纲要位置：几何学计划纲要.md，第 1941 行。
% 正文状态：未撰写；本文件目前是框架。

% 纲要内容开始
% 
% 主模块固定有限 CW 空间上的复 K 理论，采用多项式粘接—基本乘积定理的证明路线；相对版本和胞腔粘合逐项说明。[R8]
% 
% * $\mathbb{CP}^1$ 上的 Hopf 线丛与 Bott 元
% 
%   取 $\beta=[H]-1\in\widetilde K^0(S^2)$，其中 $H$ 为所选 Hopf 线丛，并核验 $(H-1)^2=0$。
% 
% * $X\times S^2$ 上向量丛的粘接函数模型
% * Laurent 多项式逼近与粘接同伦的保持
% * 增添平凡因子后的多项式粘接线性化
% * 单位圆内外的谱分解与连续子丛
% * $K^0(X\times S^2)$ 的基本乘积定理及逆映射
% 
%   构造并证明 $K^0(X)\otimes_{\mathbb Z}\mathbb Z[H]/(H-1)^2\longrightarrow K^0(X\times S^2)$ 为同构。满射由多项式粘接约化得到，单射由逆映射及其同伦无关性得到；不能用未证的 K 理论 Künneth 公式替代。
% 
% * 相对群、约化外积与二次悬挂同构
% 
%   由积空间到楔和及压缩积的相对正合列推出 $\beta\boxtimes-:\widetilde K^0(X)\xrightarrow{\sim}\widetilde K^0(\Sigma^2X)$，并核验所用自然性。
% 
% * 复周期性、球面系数及有限 CW 空间的计算
% 
% 纲要内容结束

\section{复 Bott 周期性的多项式粘接证明}
\label{g2:sec:2403}

待撰写

\subsection{\texorpdfstring{\(\mathbb{CP}^1\) 上的 Hopf 线丛与 Bott 元}{CP¹ 上的 Hopf 线丛与 Bott 元}}
\label{g2:subsec:2403:01}
待撰写

\subsection{\texorpdfstring{\(X\times S^2\) 上向量丛的粘接函数模型}{X×S² 上向量丛的粘接函数模型}}
\label{g2:subsec:2403:02}
待撰写

\subsection{Laurent 多项式逼近与粘接同伦的保持}
\label{g2:subsec:2403:03}
待撰写

\subsection{增添平凡因子后的多项式粘接线性化}
\label{g2:subsec:2403:04}
待撰写

\subsection{单位圆内外的谱分解与连续子丛}
\label{g2:subsec:2403:05}
待撰写

\subsection{\texorpdfstring{\(K^0(X\times S^2)\) 的基本乘积定理及逆映射}{K⁰(X×S²) 的基本乘积定理及逆映射}}
\label{g2:subsec:2403:06}
待撰写

\subsection{相对群、约化外积与二次悬挂同构}
\label{g2:subsec:2403:07}
待撰写

\subsection{复周期性、球面系数及有限 CW 空间的计算}
\label{g2:subsec:2403:08}
待撰写
% 纲要条目排为编号小节标题，并进入目录与书签；说明段仍保留于纲要注释中。

% !TeX root = ../../Book2.tex
% 纲要标题：### 24.4 陈特征与指标问题
% 纲要位置：几何学计划纲要.md，第 1963 行。
% 正文状态：初稿；一般指标定理采用正文所列的准确外部输入。

% 纲要内容开始
% 
% * 有理陈特征和乘法
% * 陈特征的有理同构
% * 椭圆符号的紧支 K 类
% * 解析指标、拓扑指标与指标定理
% 
% ---
% 
% 纲要内容结束

\section{陈特征与指标问题}
\label{g2:sec:2404}

一张虚丛可以有许多不同的代表，而上同调类便于加减和计算。第21章的陈类已经把丛变成了上同调数据，不过总陈类把直和变成乘法，尚不能直接成为 \(K^0\) 到上同调环的环同态。本节先换一组适合加法与张量积的示性类，随后解释：为什么一个微分方程的最高阶部分也会确定一张虚丛，以及这份拓扑数据如何给出解空间之间的维数差。

\subsection{有理陈特征和乘法}
\label{g2:subsec:2404:01}
设 \(L\) 是复线丛，记 \(x=c_1(L)\)。线丛张量积满足 \(c_1(L\otimes L')=x+x'\)，所以指数式 \(e^x\) 正好把张量积变成乘法。若一张秩 \(r\) 的丛分裂为线丛之和，就把这些指数相加。分裂原理允许我们将这个想法用于一般复向量丛。

\begin{definition}\label{g2:def:chern-character}
	设 \(X\) 为有限 CW 空间，\(E\to X\) 为秩 \(r\) 的复向量丛。在\cref{g2:thm:splitting-principle}的分裂空间上记其陈根为 \(x_1,\ldots,x_r\)。定义\term[chen tezheng]{陈特征}[Chern character]
	\[
		\operatorname{ch}(E)\coloneqq\sum_{a=1}^r e^{x_a}
		=\sum_{j\ge0}\operatorname{ch}_j(E),\qquad
		\operatorname{ch}_0(E)=r,\quad
		\operatorname{ch}_j(E)=\frac1{j!}\sum_{a=1}^r x_a^j\in H^{2j}(X;\mathbb Q).
	\]
	这里各个对称多项式先写成 \(c_1(E),\ldots,c_r(E)\) 的多项式，再视为 \(X\) 上的类。若 \(X\) 不连通，零次项逐个连通分支取秩。
\end{definition}

这个定义中没有无穷级数的收敛问题：\(H^q(X;\mathbb Q)=0\) 当 \(q>\dim X\)，所以只有有限项可能非零。为何各项确实降到 \(X\)，也可以直接从 Newton 恒等式看出。令 \(p_j=\sum_a x_a^j\)、\(c_0=1\)，则
\[
	p_j-c_1p_{j-1}+c_2p_{j-2}-\cdots+(-1)^{j-1}c_{j-1}p_1+(-1)^j j c_j=0,
\]
其中 \(c_j=0\) 当 \(j>r\)。因此 \(p_j\) 是陈类的整数系数多项式，除以 \(j!\) 后给出唯一的有理类。前几项为
\begin{equation}\label{g2:eq:chern-character-low-terms}
	\operatorname{ch}(E)=r+c_1+\frac{c_1^2-2c_2}{2}
	+\frac{c_1^3-3c_1c_2+3c_3}{6}+\cdots.
\end{equation}
这些分母说明，有理系数从定义开始就有实际作用。例如在 \(\mathbb CP^2\) 上取 \(H=\gamma^*\)、\(h=c_1(H)\)，便有 \(\operatorname{ch}(H)=1+h+h^2/2\)。

\begin{proposition}\label{g2:prop:chern-character-ring}
	对有限 CW 空间 \(X\)，公式
	\[
		\operatorname{ch}([E]-[F])=\operatorname{ch}(E)-\operatorname{ch}(F)
	\]
	给出自然的含幺环同态
	\[
		\operatorname{ch}:K^0(X)\longrightarrow H^{\mathrm{even}}(X;\mathbb Q)
		\coloneqq\bigoplus_{j\ge0}H^{2j}(X;\mathbb Q).
	\]
	它与外积相容，且 \(\operatorname{ch}_j(E^*)=(-1)^j\operatorname{ch}_j(E)\)。
\end{proposition}
\begin{proof}
	先在同一个分裂空间上把 \(E,F\) 都分裂。\Cref{g2:thm:splitting-principle}保证其拉回在有理上同调中单射，所以在那里验证的等式能够降回 \(X\)。设根分别为 \(x_a,y_b\)。直和的根是两组根的并，故
	\[
		\operatorname{ch}(E\oplus F)=\sum_a e^{x_a}+\sum_b e^{y_b}.
	\]
	张量积则分裂为所有 \(L_a\otimes M_b\) 的直和，其根是 \(x_a+y_b\)，于是
	\[
		\operatorname{ch}(E\otimes F)
		=\sum_{a,b}e^{x_a+y_b}
		=\Bigl(\sum_a e^{x_a}\Bigr)\Bigl(\sum_b e^{y_b}\Bigr).
	\]
	这里根的次数为 \(2\)，杯积交换，有限截断的指数因此满足通常的乘法恒等式。平凡线丛的根为零，特征为 \(1\)。加法性使 \(\operatorname{ch}\) 经群完备化唯一延拓；若两对丛在稳定同构后相等，其陈特征之差也相等，因此虚丛的表示无关。上述乘法式随之延拓到虚丛。

	陈类与拉回相容，定义中的多项式也与拉回相容，故所得同态自然。对 \(X\times Y\) 上两个投影的拉回应用乘法式，就得到外积公式。最后，\(E^*\) 的根为 \(-x_a\)，其 \(j\) 次幂和是原来的 \((-1)^j\) 倍，给出对偶公式。
\end{proof}

对秩零类，零次分量消失。特别地，本章的 Bott 元满足
\begin{equation}\label{g2:eq:chern-bott-positive}
	\operatorname{ch}(\beta)=\operatorname{ch}([H]-1)=e^h-1=h
	\quad\text{于 }\widetilde H^*(S^2;\mathbb Q),
\end{equation}
因为 \(h^2=0\)。这既固定了接下来悬挂比较的正方向，也说明 Bott 元在上同调中对应二次生成元。

\subsection{陈特征的有理同构}
\label{g2:subsec:2404:02}
负一次 \(K\) 群由悬挂定义，因而其陈特征也应该由悬挂得到。记
\[
	s:H^q(X,A;\mathbb Q)\xrightarrow{\cong}
	\widetilde H^{q+1}(\Sigma(X/A);\mathbb Q)
\]
为与\secref{g2:sec:2402}锥坐标方向一致的上同调悬挂同构。将相对 \(K\) 类视为商空间上的约化类，可定义相对的偶陈特征。对 \(u\in K^{-1}(X,A)=\widetilde K^0(\Sigma(X/A))\)，定义
\[
	\operatorname{ch}^{-1}(u)\coloneqq s^{-1}\operatorname{ch}(u)
	\in H^{\mathrm{odd}}(X,A;\mathbb Q)
	\coloneqq\bigoplus_{j\ge0}H^{2j+1}(X,A;\mathbb Q).
\]
悬挂后的约化零次上同调为零，所以右端恰由正的奇次数组成。绝对情形采用 \(X/\varnothing=X_+\)，与本章的未约化 \(K^{-1}(X)\) 一致。

\begin{theorem}\label{g2:thm:rational-chern-character}
	对有限 CW 对 \((X,A)\)，陈特征给出自然同构
	\[
	\begin{aligned}
		K^0(X,A)\otimes_{\mathbb Z}\mathbb Q&\xrightarrow{\cong}H^{\mathrm{even}}(X,A;\mathbb Q),\\
		K^{-1}(X,A)\otimes_{\mathbb Z}\mathbb Q&\xrightarrow{\cong}H^{\mathrm{odd}}(X,A;\mathbb Q).
	\end{aligned}
	\]
	绝对情形的第一行是含幺环同构。
\end{theorem}
\begin{proof}
	证明按胞腔数归纳。需要先把球面上的比较与连接映射的相容性说明清楚，否则不能仅凭两边都具有长正合列就应用五引理。

	由\cref{g2:thm:complex-bott-periodicity,g2:prop:k-sphere-coefficients}，偶维球面 \(S^{2m}\) 的约化零次 \(K\) 群由 \(m\) 个 Bott 元的约化外积生成。式\eqref{g2:eq:chern-bott-positive}及外积公式把这个生成元送到 \(h_1\boxtimes\cdots\boxtimes h_m\)，即压缩积 \(S^2\wedge\cdots\wedge S^2\cong S^{2m}\) 的有向上同调生成元。另一种奇偶性的球面群为零。奇维球面悬挂一次成为偶维球面；按刚才的定义，负一次陈特征同样把相应生成元送到奇次上同调生成元。零维情形则由秩给出。故两种比较在每个球面及其有限楔和上都是有理同构。

	再看连接映射。\Cref{g2:thm:k-long-exact}中的连接由映射锥的商映射 \(X\cup_A CA\to\Sigma A\) 拉回给出；上同调的连接也是同一商映射的拉回，再经悬挂识别。陈特征与拉回相容，而我们用这同一个悬挂方向定义 \(\operatorname{ch}^{-1}\)，所以由负次数群组成的正合列与上同调列之间各个方块都交换。当用 Bott 周期性将次数折回偶、奇两类时，式\eqref{g2:eq:chern-bott-positive}说明乘 \(\beta\) 正对应两次上同调悬挂，因而这些方块仍然交换。这一步只使用已经证明的 Bott 周期性，并未用待证的有理同构来推断周期性。

	设 \(X^n\) 是 \(n\) 维骨架。相对商 \(X^n/X^{n-1}\) 是有限个 \(S^n\) 的楔和，故相对群上的比较已经是同构。对两个正合列取连续五项
	\[
		K^{j-1}(X^{n-1})\longrightarrow K^j(X^n,X^{n-1})
		\longrightarrow K^j(X^n)\longrightarrow K^j(X^{n-1})
		\longrightarrow K^{j+1}(X^n,X^{n-1}),
	\]
	其中次数按二周期解释。将此列张量 \(\mathbb Q\)，正合性保持：把每个元素写成分数后清除有限个分母，就把核像问题化回原来的整数正合列。其下方放置相同奇偶性的有理上同调列。由骨架归纳及球面计算，中间一项以外的四个竖箭头均为同构，五引理给出中间一项的同构。由有限零骨架出发，经过有限次便得到绝对结论。对任意有限 CW 对再取其长正合列，应用同样的五引理，得到相对结论。乘法相容性已由\cref{g2:prop:chern-character-ring}证明。
\end{proof}

因此，对有限 CW 空间，复 \(K\) 理论与普通上同调的差异包含两部分：有理化前的扭元，以及有理上同调中哪些类实际来自整数 \(K\) 类的整性条件。\Cref{g2:exm:k-moore-torsion}中的 Moore 空间给出第一种现象：其约化 \(K^0\) 群是 \(\mathbb Z/q\)，而正次数有理上同调全为零，陈特征看不到这些非零扭元。上面的 \(\mathbb CP^2\) 算例则说明，来自实际线丛的特征也可以带有分数系数。

\subsection{椭圆符号的紧支 K 类}
\label{g2:subsec:2404:03}
设 \(M\) 是闭光滑流形，\(E,F\to M\) 是有限秩复光滑向量丛，\(P:C^\infty(M,E)\to C^\infty(M,F)\) 是 \(m\ge1\) 阶线性微分算子。在局部坐标和丛标架下写成
\[
	P=\sum_{|\alpha|\le m}a_\alpha(x)\partial_x^\alpha.
\]
固定 Fourier 约定，以 \(\partial_{x_j}\) 对应 \(i\xi_j\)。其\term[zhufuhao]{主符号}[principal symbol]为
\[
	\sigma_m(P)(x,\xi)=\sum_{|\alpha|=m}a_\alpha(x)(i\xi)^\alpha:
	E_x\longrightarrow F_x.
\]
换坐标时，最高阶导数按余切变量作张量变换；导数落在换坐标或换标架的系数上时只贡献低阶项。因此这个公式给出 \(\mathrm T^*M\) 上由 \(\pi^*E\) 到 \(\pi^*F\) 的整体丛映射。对 \(t>0\) 有 \(\sigma_m(P)(x,t\xi)=t^m\sigma_m(P)(x,\xi)\)。

\begin{definition}\label{g2:def:elliptic-symbol-k-class}
	若对每个 \(x\in M\) 及非零 \(\xi\in\mathrm T_x^*M\)，\(\sigma_m(P)(x,\xi)\) 都是复线性同构，则称 \(P\) 为\term[tuoyuan suanzi]{椭圆算子}[elliptic operator]。取一个 Riemann 度量，以 \(D(\mathrm T^*M)\)、\(S(\mathrm T^*M)\) 表示单位闭余球丛及其边界。定义其符号类为
	\[
		[\sigma(P)]\coloneqq
		[\pi^*E,\pi^*F,\sigma_m(P)|_{S(\mathrm T^*M)}]
		\in K^0\bigl(D(\mathrm T^*M),S(\mathrm T^*M)\bigr).
	\]
\end{definition}

这里使用的正是\cref{g2:prop:k-triple-model}的三元组模型：两张丛在球丛内部不必同构，但在边界上由可逆主符号给出同构。将向量的长度径向压入单位球后，余切丛的一点紧化同胚于 \(D(\mathrm T^*M)/S(\mathrm T^*M)\)。因而这个群也写成
\[
	K_c^0(\mathrm T^*M)\coloneqq\widetilde K^0\bigl((\mathrm T^*M)^+\bigr),
\]
称为\term[jinzhi K lilun]{紧支 \(K\) 理论}[compactly supported K-theory]。下标 \(c\) 表达的是相对平凡化数据：在紧集以外，两个丛由指定同构相抵消。丛本身并不需要在零截面外消失。

\ProseParagraphBreak

还有一个有用的等价描述。在 \(\mathrm T^*M\) 上考虑两项复形
\[
	0\longrightarrow\pi^*E\xrightarrow{\ \sigma_m(P)\ }\pi^*F\longrightarrow0.
\]
它在零截面外正合，零截面又因 \(M\) 闭而紧，所以它代表一个紧支类。一般有限复形在紧集外正合，也可给出这样的类：取 Hermitian 度量，将偶次项与奇次项分别相加，以 \(\mathrm d+\mathrm d^*\) 从偶次和映到奇次和。正合处的拉普拉斯算子 \(\mathrm d\mathrm d^*+\mathrm d^*\mathrm d\) 可逆，所以这个奇偶映射可逆。两项复形时便恢复上述三元组。此模型及其乘积的详细构造可见\cite[第2节，第489--494页]{AtiyahSinger1968I}。

\begin{proposition}\label{g2:prop:symbol-class-invariance}
	闭光滑流形上复向量丛间椭圆微分算子的符号类不依赖定义单位余球时所取的度量。连续变形的可逆主符号给出相同的 \(K_c^0(\mathrm T^*M)\) 类；特别地，增加低于主阶数的项不改变该类。直和算子的符号类是原来两类之和。
\end{proposition}
\begin{proof}
	两个度量之间以正定二次型的凸组合连接。沿每条非零余切射线作径向重新缩放，便把各个单位余球丛连续识别；主符号的正齐次性说明重新缩放只乘一个正实数。这个因子可在正数中连接到 \(1\)，因此三元组类不变。主符号的同伦正是三元组允许的边界同构同伦。低阶项不进入 \(\sigma_m(P)\)，直和则给出主符号的块对角和，与三元组的加法一致。
\end{proof}

在圆周 \(M=\mathbb R/2\pi\mathbb Z\) 上，\(P=\mathrm d/\mathrm d\theta\) 作用于复值函数，主符号是 \(i\xi\)。单位余球有两个分支，边界同构分别是常数 \(i\) 与 \(-i\)。在这两个分支上各取一条从该常数到 \(1\) 的复非零道路，就将边界同构同伦为恒等；恒等同构延拓到整个单位余球丛，所以
\begin{equation}\label{g2:eq:circle-derivative-symbol-zero}
	[\sigma(\mathrm d/\mathrm d\theta)]=0\quad\text{于 }K_c^0(\mathrm T^*S^1).
\end{equation}
这是一次相对 \(K\) 类的计算，尚未数微分方程的解。

\subsection{解析指标、拓扑指标与指标定理}
\label{g2:subsec:2404:04}
现在从同一算子 \(P\) 出发，分别构造一个分析整数与一个拓扑整数。分析一侧必须先保证两个被相减的维数有限。所用的椭圆分析结论是：闭光滑 \(M\) 上的 \(m\) 阶椭圆算子，在 Sobolev 完备化后给出 Fredholm 算子
\[
	P:H^s(M,E)\longrightarrow H^{s-m}(M,F),
\]
其像闭，核与余核有限维，核及伴随核中的弱解均光滑，且指标与 \(s\) 无关。\(H^s\) 在整数 \(s\ge0\) 时可以理解为直到 \(s\) 阶弱导数平方可积的截面空间；完备化使“闭像”成为可用的分析性质。这里调用\cite[第6节，式(6.1)，第517--519页]{AtiyahSinger1968I}中的结论，Sobolev 截面与 Fredholm 分析分别安排在第三卷第11.2节、第13.2节，指标的几何应用见第三卷第17.4节的安排。

\begin{definition}\label{g2:def:analytic-index}
	设 \(M\) 闭且光滑，\(P:C^\infty(M,E)\to C^\infty(M,F)\) 为复向量丛间的椭圆微分算子。它的\term[jiexi zhibiao]{解析指标}[analytic index]为
	\[
		\operatorname{ind}_{\mathrm a}(P)\coloneqq
		\dim_{\mathbb C}\ker P-\dim_{\mathbb C}\operatorname{coker}P.
	\]
	这里余核可以在上述 Sobolev 空间中取商；由椭圆正则性，也可用光滑截面取商。选度量后，余核与形式伴随 \(P^*\) 的核自然对应。
\end{definition}

拓扑一侧要把余切丛上的紧支类送到一点上的 \(K^0(\mathrm{pt})=\mathbb Z\)。与普通上同调的 Gysin 映射相似，这需要法丛上的 Thom 类。复向量丛恰好具有一个可直接构造的 \(K\) 理论 Thom 类。

\begin{proposition}\label{g2:prop:complex-k-thom}
	设 \(V\to X\) 是有限 CW 空间上的秩 \(r\) 复向量丛。以 \(p:V\to X\) 为投影，\(v\in p^*V\) 为总空间上的典范截面。外代数复形
	\[
		0\longrightarrow p^*\Lambda^0V
		\xrightarrow{v\wedge}p^*\Lambda^1V
		\xrightarrow{v\wedge}\cdots
		\xrightarrow{v\wedge}p^*\Lambda^rV\longrightarrow0
	\]
	定义一个类 \(\lambda_V\in K_c^0(V)\)。乘这个类给出同构
	\[
		\Phi_V:K^j(X)\xrightarrow{\cong}K_c^j(V),\qquad
		a\longmapsto p^*a\cdot\lambda_V,\qquad j=0,-1.
	\]
	其相对版本也成立；若 \(X\) 是闭光滑流形上一张有限秩实光滑向量丛的总空间，\(V\to X\) 为复光滑向量丛，同样得到紧支版本 \(K_c^j(X)\cong K_c^j(V)\)。
\end{proposition}
\begin{proof}
	取 Hermitian 度量。记 \(\varepsilon(v)\) 为楔乘 \(v\)，\(\iota(v)=\varepsilon(v)^*\) 为其伴随；在每个纤维上
	\[
		\varepsilon(v)\iota(v)+\iota(v)\varepsilon(v)=|v|^2\operatorname{id}.
	\]
	所以 \(v\ne0\) 时，\(\iota(v)/|v|^2\) 是这个复形的收缩同伦，复形正合。等价地，\(\varepsilon(v)+\iota(v)\) 从偶次外幂和映到奇次外幂和，在零截面外可逆。由此得到紧支三元组。度量间的道路给出同伦，外代数构造与丛同构相容，故此类整体良定义。

	在一维复纤维上，复形是 \(\mathbb C\xrightarrow{z}\mathbb C\)。按\cref{g2:prop:k-triple-model}的粘合方向，原点盘的标架在无穷远盘中等于 \(z\) 倍标架。这是典范线丛 \(\gamma\to\mathbb CP^1\) 的过渡式；因此
	\begin{equation}\label{g2:eq:k-koszul-bott-sign}
		\lambda_{\mathbb C}=[\gamma]-1=[H]^{-1}-1=-\beta.
	\end{equation}
	最后一个等式使用 \(\beta^2=0\)。这也可以从本章粘接约定核对：正向 \(H\) 将无穷远盘的坐标乘 \(z\) 送入原点盘，而上述三元组将原点盘的坐标乘 \(z\) 送入无穷远盘，故二者互为逆线丛。

	对直和有带分次的外代数同构 \(\Lambda^*(V\oplus W)\cong\Lambda^*V\widehat\otimes\Lambda^*W\)，微分对应张量总微分。因此 \(\lambda_{V\oplus W}=\lambda_V\boxtimes\lambda_W\)（在共同底空间上取相应拉回）。在平凡秩 \(r\) 丛上，乘 \(\lambda_V\) 就是乘 \((-1)^r\beta^{\boxtimes r}\)，由复 Bott 周期性是同构。

	一般 \(X\) 按有限骨架过滤。将 \(V\) 限制到一个胞腔的特征圆盘上，它平凡；对于一对相邻骨架，相对 Thom 空间的商是这些圆盘上相对 Thom 空间的有限楔和，因而相对映射正是刚才证明的平凡丛版本。Thom 类由整体外代数复形给出，限制到骨架和边界时保持同一复形，故各个长正合列方块交换。从零骨架开始应用五引理，得到结论。对 \((X,A)\) 用相对骨架作同样的归纳，证明相对版本。

	最后，若底空间 \(X\) 是闭光滑流形上一张有限秩实光滑向量丛的总空间，以半径递增的闭球丛穷尽它。每个紧集都包含在某个球丛的内部；紧支 \(K\) 类因而来自相应球丛及边界的相对群。球丛与边界具有有限 CW 对的同伦型，上一段的相对同构可用，并与扩大半径相容。取这些群的直极限，便得到所述紧支版本。此时外代数复形本身只在纤维方向具有紧支；与底空间的紧支类相乘后，所得复形在总空间中才具有紧支。
\end{proof}

这里采用的是原始指标论文的外代数 \(V\) 约定，见\cite[第2节，第493--494页]{AtiyahSinger1968I}。式\eqref{g2:eq:k-koszul-bott-sign}说明它与本章正陈根的 Bott 元相差一个负号；下述推前和末端整数识别始终使用同一套 \(\lambda\) 约定。

\ProseParagraphBreak

取光滑嵌入 \(i:M\hookrightarrow\mathbb R^N\)，记其实法丛为 \(\nu\)。度量把 \(\mathrm T^*M\) 识别为 \(\mathrm TM\)，而切映射给出嵌入 \(\mathrm Ti:\mathrm TM\hookrightarrow\mathrm T\mathbb R^N\)。其法丛同构于从 \(M\) 拉回的 \(\nu\oplus\nu\)：第一份来自底空间的法位移，第二份来自切向量的法分量。这个分解可由一个联络写出，改变联络只产生同伦。按
\[
	(a,b)\longmapsto a+ib
\]
给两份法方向赋予复结构，就得到 \(p^*(\nu\otimes_{\mathbb R}\mathbb C)\)。它的复秩为 \(N-\dim M\)。因此此处的复 Thom 类来自两份相同实法丛的组合。

\ProseParagraphBreak

先沿这个复法丛作 Thom 映射，再将管状邻域中的紧支类延为环境中零，得到
\[
	(\mathrm Ti)_!:K_c^0(\mathrm TM)\longrightarrow
	K_c^0(\mathrm T\mathbb R^N).
\]
“延为零”由一点紧化之间的压缩映射实现：环境中管状邻域以外的部分送到无穷远点，故其方向是 \(K\) 类的拉回。将 \((x,v)\in\mathrm T\mathbb R^N\) 写成 \(x+iv\in\mathbb C^N\)，前一命题及 Bott 周期性给出
\[
	K_c^0(\mathrm T\mathbb R^N)=K_c^0(\mathbb C^N)
	=\mathbb Z\,\lambda_{\mathbb C^N},\qquad
	\lambda_{\mathbb C^N}=\lambda_{\mathbb C}^{\boxtimes N}=(-1)^N\beta^{\boxtimes N}.
\]
令 \(b_N\) 把 \(\lambda_{\mathbb C^N}\) 送到 \(1\)。于是\term[tuopu zhibiao]{拓扑指标}[topological index]是复合
\begin{equation}\label{g2:eq:topological-index-pushforward}
	\operatorname{ind}_{\mathrm t}:K_c^0(\mathrm T^*M)
	\xrightarrow{\ \cong\ }K_c^0(\mathrm TM)
	\xrightarrow{\ (\mathrm Ti)_!\ }K_c^0(\mathbb C^N)
	\xrightarrow{\ b_N\ }\mathbb Z.
\end{equation}
这里既固定了法丛中“底法方向在前、纤维法方向在后”的次序，也固定了末端生成元，因而不会因混用 Bott 元的符号改变指标。

\ProseParagraphBreak

这个整数与嵌入和辅助度量无关。度量的变化给出一族识别，管状邻域的变化给出同伦的压缩映射；Thom 类在这些同伦下自然。增加一个平凡环境实方向时，推前额外乘 \(\lambda_{\mathbb C}\)，而末端 \(b_{N+1}\) 正好消去这一因子，故整数不变。对两个嵌入 \(i_0:M\to\mathbb R^{N_0}\)、\(i_1:M\to\mathbb R^{N_1}\)，在乘积环境中用
\[
	x\longmapsto\bigl(\cos t\,i_0(x),\sin t\,i_1(x)\bigr),
	\qquad 0\le t\le\pi/2,
\]
连接其稳定化。每个中间映射至少有一个非零系数的嵌入分量，因而仍是嵌入；沿这族嵌入的法丛及管状邻域给出同伦的推前。末端交换环境实坐标只在复化后作复线性变换，保持 \(\lambda\) 类。上述构造与独立性的完整拓扑处理见\cite[第3节，第497--499页]{AtiyahSinger1968I}。

\begin{theorem}[Atiyah--Singer 指标定理]\label{g2:thm:atiyah-singer-topological-index}
	设 \(M\) 为闭光滑流形，\(E,F\to M\) 为有限秩复光滑向量丛，\(P:C^\infty(M,E)\to C^\infty(M,F)\) 为椭圆微分算子。则
	\[
		\operatorname{ind}_{\mathrm a}(P)
		=\operatorname{ind}_{\mathrm t}([\sigma(P)]),
	\]
	其中右端使用式\eqref{g2:eq:topological-index-pushforward}的复法丛推前及 \(\lambda_{\mathbb C^N}\mapsto1\) 的归一化。
\end{theorem}

这个结论称为\term[Atiyah-Singer-zhibiao-dingli]{Atiyah--Singer 指标定理}[Atiyah--Singer index theorem]。\footnote{Michael Francis Atiyah（1929--2019），英国数学家，参与建立拓扑 \(K\) 理论，并研究几何与数学物理之间的联系。Isadore Manuel Singer（1924--2021），美国数学家，研究微分几何与泛函分析。}这里引用\cite[第6节，定理(6.7)，第520页]{AtiyahSinger1968I}的结论，取其等变陈述中的作用群为平凡群。完整证明还要建立椭圆算子的解析指标在这些拓扑推前下的相容性；原文第8、9节完成这一环节。解析指标从有限维核与余核得到整数，拓扑指标从符号的紧支类及法丛得到整数，定理把这两种构造联系起来。

\begin{example}\label{g2:exm:circle-derivative-index}
	对圆周上的 \(P=\mathrm d/\mathrm d\theta\)，周期解 \(Pf=0\) 恰是常函数，核为一维。方程 \(Pf=g\) 有周期光滑解，当且仅当
	\[
		\int_0^{2\pi}g(\theta)\,\mathrm d\theta=0.
	\]
	必要性来自对导数积分；充分性由 \(f(\theta)=\int_0^\theta g(t)\,\mathrm dt\) 给出，零积分保证周期性，且 \(g\) 的周期性保证各阶导数在接缝吻合。因此余核由平均值给出，也是一维，解析指标为 \(1-1=0\)。另一方面，式\eqref{g2:eq:circle-derivative-symbol-zero}已算出符号类为零，拓扑指标随之为零。这个例子中，核与余核都非零，二者的差却只由符号的拓扑类决定。
\end{example}
% 纲要条目排为编号小节标题，并进入目录与书签；说明段仍保留于纲要注释中。


\begin{chaptersummary}
	K 群以虚丛保留直和的加法关系，以张量积保留乘法。紧空间上的补丛把虚丛写成 \([E]-[\underline{\mathbb C}^N]\)，相对类则同时记录子空间上的一个稳定同构。映射锥使这份边界数据成为一个约化类，悬挂与余纤维列由此给出长正合列和切除。

	\ProseParagraphBreak

	复 Bott 周期性的证明把连续粘接函数依次化为 Laurent 多项式和一次矩阵。谱投影将一次矩阵分成圆内与圆外两类子丛：圆内部分贡献 \(H\)，圆外部分贡献常丛。对多项式次数、负幂及同伦选择的比较保证逆映射良定义，从而得到 \(K^0(X\times S^2)=K^0(X)\oplus K^0(X)\beta\) 以及 \(\beta^2=0\)。取约化部分后，二次悬挂保持复 K 群；球面系数、六项正合列和逐胞腔计算随之得到。

	\ProseParagraphBreak

	陈特征把直和与张量积分别变成上同调的加法与乘法，并在有限 CW 空间上给出有理同构。它揭示有理信息，也解释扭元为何可能不被示性类看见。椭圆算子的可逆主符号给出余切丛上的紧支 K 类；复 Thom 类和嵌入把它推到一个整数。解析指标由核与余核定义，拓扑指标由这个推前定义，Atiyah--Singer 指标定理断言二者相等。
\end{chaptersummary}

\BookExercises

\begin{exercise}\label{g2:ex:k-disconnected-reduction}
	设 \(X\) 是 \(m\ge2\) 个点组成的离散空间，取第一个点为基点。求 \(K^0(X)\) 的环结构及 \(\widetilde K^0(X)\)。再求“在每一个点虚秩都为零”的类组成的子群，比较它与约化群。
\end{exercise}
\begin{solution}
	一张丛由每个点上的复向量空间决定，同构类由 \((r_1,\ldots,r_m)\in\mathbb N^m\) 唯一确定。直和逐坐标相加，所以群完备化为 \(\mathbb Z^m\)；张量积逐纤维进行，乘法也逐坐标进行：
	\[
		(a_1,\ldots,a_m)(b_1,\ldots,b_m)=(a_1b_1,\ldots,a_mb_m).
	\]
	单位元为 \((1,\ldots,1)\)。基点取秩是第一坐标投影，故 \(\widetilde K^0(X)=\{0\}\times\mathbb Z^{m-1}\)。在每点虚秩均为零则要求全部坐标都为零，所得子群为零。约化群只消去所选基点处的秩；例如 \((0,1,0,\ldots,0)\) 是约化群中的非零元。
\end{solution}

\begin{exercise}\label{g2:ex:hopf-dual-sum-trivialization}
	设 \(H\to S^2\) 是本章的 Hopf 线丛。用粘合矩阵给出 \(H\oplus H^*\cong\varepsilon^2\) 的具体同伦证明，并计算 \([H^*]-1\) 与 \([H]-1\) 的关系。解释为什么这个关系不意味着 \([H^*]=-[H]\)。
\end{exercise}
\begin{solution}
	按本章赤道约定，直和的粘合矩阵为 \(\operatorname{diag}(z,z^{-1})\)，\(|z|=1\)。下式在整个参数区间取值于 \(\operatorname{SU}(2)\)：
	\[
		G_t(z)=\begin{pmatrix}\cos t\,z&-\sin t\\\sin t&\cos t\,\overline z\end{pmatrix},
		\qquad0\le t\le\pi/2.
	\]
	两列正交且为单位向量，行列式为 \(\cos^2t+\sin^2t=1\)。起点是所给粘合函数，终点为常矩阵 \(J=\left(\begin{smallmatrix}0&-1\\1&0\end{smallmatrix}\right)\)。再用实旋转矩阵把 \(J\) 连到 \(I\)，便得到粘合函数的零伦。\cref{g2:thm:clutching-classification}于是给出实际秩二丛的平凡化类型。

	在 \(K^0(S^2)\) 中，这说明 \([H]+[H^*]=2\)，所以 \([H^*]-1=-([H]-1)\)。反之，\([H^*]\) 的虚秩为 \(1\)，\(-[H]\) 的虚秩为 \(-1\)，二者不相等。负元关系成立的是已经减去秩的两项。
\end{solution}

\begin{exercise}\label{g2:ex:k-pullback-two-sphere-degree}
	设 \(f:S^2\to S^2\) 的映射度为整数 \(d\)，记 \(\eta=[H]-1\)。求 \(f^*:K^0(S^2)\to K^0(S^2)\) 在加法基 \(1,\eta\) 下的矩阵，并分别求其核与余核。最后说明 \(f^*H\) 与哪张线丛同构。
\end{exercise}
\begin{solution}
	拉回保持平凡线丛，所以 \(f^*1=1\)。由第一陈类的自然性，\(c_1(f^*H)=f^*h=dh\)，故式\eqref{g2:eq:k-two-sphere-additive}给出 \(f^*\eta=d\eta\)。因此所求矩阵为
	\[
		\begin{pmatrix}1&0\\0&d\end{pmatrix}.
	\]
	当 \(d\ne0\) 时，核为零，余核为 \(\mathbb Z/|d|\mathbb Z\)；当 \(d=0\) 时，核为 \(\mathbb Z\eta\)，余核为 \(\mathbb Z\)。秩一复丛由第一陈类分类，所以 \(f^*H\cong H^{\otimes d}\)，负张量幂按对偶理解。这个同构由实际线丛分类得到，比只写出它在 \(K^0\) 中的等式更精确。
\end{solution}

\begin{exercise}\label{g2:ex:k-external-product-visible-chern}
	在 \(S^2\times S^2\) 上，记 \(L=p_1^*H\)、\(M=p_2^*H\)。证明 \(([L]-1)([M]-1)\ne0\)，但该类在 \(S^2\times\{*\}\) 和 \(\{*\}\times S^2\) 上的限制均为零。可直接使用第21章的陈类和公式，无需使用 K 理论乘积定理。
\end{exercise}
\begin{solution}
	令 \(a=p_1^*h\)、\(b=p_2^*h\)。普通整数上同调的 K\"unneth 公式给出 \(a^2=b^2=0\)，而 \(ab\) 在乘积基本类上的求值为 \(1\)，故非零。展开虚丛乘积为
	\[
		([L]-1)([M]-1)=[(L\otimes M)\oplus\varepsilon^1]-[L\oplus M].
	\]
	若它为零，两张秩二丛应当稳定同构，故具有相同陈类。可是
	\[
		c((L\otimes M)\oplus\varepsilon^1)=1+a+b,
		\qquad c(L\oplus M)=(1+a)(1+b)=1+a+b+ab.
	\]
	二次陈类之差为非零的 \(ab\)，矛盾。在第一张因子球面上，\(M\) 限制为平凡线丛；在第二张因子球面上，\(L\) 限制为平凡线丛。因此乘积在两者上均为零。这个类记录两因子共同参与的信息，分别限制到各因子会失去它。
\end{solution}

\begin{exercise}\label{g2:ex:k-retraction-torus}
	设有限 CW 对 \((X,A)\) 的包含有连续收缩映射 \(r:X\to A\)，即 \(r|_A=\operatorname{id}_A\)。证明对每个 \(q\ge0\)，
	\[
		K^{-q}(X)\cong K^{-q}(X,A)\oplus K^{-q}(A).
	\]
	将此用于 \(X=T^2=S^1\times S^1\)、\(A=S^1\times\{1\}\)，计算两者的零次、负一次相对群及 \(K^0(T^2),K^{-1}(T^2)\)。
\end{exercise}
\begin{solution}
	由 \(i^*r^*=\operatorname{id}\)，限制 \(i^*\) 在每个次数上满射。长正合列中进入相对群的连接映射因前一次限制满射而为零，所以
	\[
		0\longrightarrow K^{-q}(X,A)\xrightarrow{j}K^{-q}(X)
		\xrightarrow{i^*}K^{-q}(A)\longrightarrow0
	\]
	短正合，并以 \(r^*\) 为截面。明确的分解把 \(u\) 写成 \((u-r^*i^*u)+r^*i^*u\)，其中第一项在 \(j\) 的像内且有唯一相对原像。

	对环面，取向第一因子的投影作 \(r\)。标准胞腔分解有两条边 \(a,b\) 和一张沿 \(aba^{-1}b^{-1}\) 附着的二胞腔。压掉 \(A\) 就压掉 \(a\)，附着字变为 \(bb^{-1}\)，它零伦。由\cref{g2:prop:homotopic-attaching-maps}，
	\[
		T^2/A\simeq S^1\vee S^2.
	\]
	本章的球面计算给出 \(\widetilde K^0(S^1)=0\)、\(\widetilde K^0(S^2)=\mathbb Z\)，以及 \(\widetilde K^{-1}(S^1)=\mathbb Z\)、\(\widetilde K^{-1}(S^2)=0\)。因此
	\[
		K^0(T^2,A)\cong\mathbb Z,\qquad K^{-1}(T^2,A)\cong\mathbb Z.
	\]
	又 \(K^0(A)=\mathbb Z\)、\(K^{-1}(A)=\mathbb Z\)，代入刚证明的分裂，得到 \(K^0(T^2)\cong\mathbb Z^2\)、\(K^{-1}(T^2)\cong\mathbb Z^2\)。这是加性群的计算，乘法还须保留两个方向及其分次符号。
\end{solution}

\begin{exercise}\label{g2:ex:k-mapping-torus-degree}
	设 \(f:S^2\to S^2\) 为保持基点的度数 \(d\in\mathbb Z\) 的连续映射，定义映射环面
	\[
		T_f=(S^2\times[0,1])/\bigl((x,1)\sim(f(x),0)\bigr).
	\]
	用两段圆柱及相对正合列计算 \(K^0(T_f)\)、\(K^{-1}(T_f)\)，特别求出 \(d=-1\) 时的结果。说明计算是否要求 \(T_f\to S^1\) 是纤维丛。
\end{exercise}
\begin{solution}
	给 \(S^2\) 取一张零胞腔和一张二胞腔，保持基点的 \(f\) 就是胞腔映射，故 \(T_f\) 有有限 CW 结构。记 \(U\) 为参数区间 \([0,1/3]\cup[2/3,1]\) 的像，\(V\) 为 \([1/4,3/4]\) 的像。细分圆柱后，这些集合形成子复形覆盖。\(U\) 由接缝处的映射柱及一段领圈组成，形变收缩到 \(t=0\) 的目标球面；\(V\) 也形变收缩到一张球面。交集有两条圆柱，故
	\[
		U\simeq S^2,\qquad V\simeq S^2,\qquad U\cap V\simeq S^2\sqcup S^2.
	\]
	按靠近 \(t=0\)、靠近 \(t=1\) 的次序排列交集分支，进入 \(U\) 的两张映射在这些识别下分别为恒等与 \(f\)，进入 \(V\) 的两张映射均为恒等。

	这里使用的 Mayer--Vietoris 序列来自\cref{g2:thm:k-long-exact}：对 \((T_f,U)\) 与 \((V,U\cap V)\) 写长正合列，利用切除识别相对项；两条限制映射作差，逐项正合性便给出
	\[
		\cdots\longrightarrow K^i(T_f)\longrightarrow K^i(U)\oplus K^i(V)
		\xrightarrow{D_i}K^i(U\cap V)\longrightarrow K^{i+1}(T_f)\longrightarrow\cdots.
	\]
	次数按本章的二周期约定理解，\(K^1=K^{-1}\)。例如在中间项，差为零表示两份限制相同；先用一对的正合性把它们的差提升到相对项，再通过切除移到另一对，即得到整体原像。对连接项同样使用相对列的核像等式，得到相邻两处的正合性。这正是两段粘合所需的序列。

	在上述三个同伦模型中，差映射为
	\[
		D_i(a,b)=(a-b,f^*a-b).
	\]
	把第二个输出减去第一个输出，并以 \((a,a-b)\) 作为输入坐标，就将这张映射化为一份恒等映射与 \(f^*-1\) 的直和。因此其核、余核分别等于 \(1-f^*\) 的核、余核。由于 \(K^{-1}(S^2)=0\)，正合列给出
	\[
		K^0(T_f)\cong\ker(1-f^*:K^0(S^2)\to K^0(S^2)),\qquad
		K^{-1}(T_f)\cong\operatorname{coker}(1-f^*).
	\]
	以 \(1,\beta\) 为 \(K^0(S^2)=\mathbb Z\oplus\mathbb Z\beta\) 的基，\(f^*1=1\)。又 \(c_1:\widetilde K^0(S^2)\to H^2(S^2;\mathbb Z)\) 为同构，而度数 \(d\) 的映射把二次上同调乘以 \(d\)，故 \(f^*\beta=d\beta\)。所以
	\[
	\begin{aligned}
		K^0(T_f)&\cong\mathbb Z\oplus\ker(1-d:\mathbb Z\to\mathbb Z),\\
		K^{-1}(T_f)&\cong\mathbb Z\oplus\operatorname{coker}(1-d:\mathbb Z\to\mathbb Z).
	\end{aligned}
	\]
	当 \(d=-1\) 时，结果为 \(K^0(T_f)=\mathbb Z\)、\(K^{-1}(T_f)\cong\mathbb Z\oplus\mathbb Z/2\)。当 \(d=1\) 时，两群均为 \(\mathbb Z^2\)；其他次数时，前者为 \(\mathbb Z\)，后者为 \(\mathbb Z\oplus\mathbb Z/|1-d|\)。整个计算只用了映射柱、同伦型与粘合正合性，未要求 \(f\) 为同胚，也未要求 \(T_f\to S^1\) 是纤维丛。
\end{solution}

\begin{exercise}\label{g2:ex:k-suspension-square-zero}
	设 \(Y\) 是带零胞腔基点的有限 CW 空间。证明 \(K^0(\Sigma Y)\) 中的约化理想满足
	\[
		\widetilde K^0(\Sigma Y)\cdot\widetilde K^0(\Sigma Y)=0.
	\]
	解释这为什么并不要求 \(\widetilde K^0(\Sigma Y)\) 本身为零。
\end{exercise}
\begin{solution}
	乘法是外积沿对角映射的拉回。对约化类，它只取决于约化对角
	\[
		\overline\Delta:\Sigma Y\longrightarrow\Sigma Y\wedge\Sigma Y.
	\]
	将悬挂点写为 \([t,y]\)，\(0\le t\le1\)。在乘积 \(\Sigma Y\times\Sigma Y\) 中，令
	\[
		a(t)=\min(2t,1),\qquad b(t)=\max(2t-1,0).
	\]
	映射 \([t,y]\mapsto([a(t),y],[b(t),y])\) 落在楔和中：前半段第二坐标为基点，后半段第一坐标为基点。把两个参数函数分别沿直线从 \(t\) 变到 \(a(t),b(t)\)，就得到从普通对角到此映射的同伦；端点及 \(y\) 为基点时始终取基点，所以同伦在悬挂商上良定义。再投影到压缩积，便将 \(\overline\Delta\) 零伦化。

	因此对任意约化类 \(u,v\)，有 \(uv=\overline\Delta^*(u\boxtimes v)=0\)。但单个约化类可以非零，例如取 \(Y=S^1\)，则 \(\Sigma Y\cong S^2\)，其 Bott 元生成 \(\widetilde K^0(S^2)\cong\mathbb Z\)。这里消失的是两个约化类的乘积，而不是加性群。
\end{solution}

\begin{exercise}\label{g2:ex:k-family-spectral-bundle}
	在第一份 \(S^2=\mathbb CP^1\) 上，把 \(H=\gamma^*\) 实现为平凡丛 \(\underline{\mathbb C}^2\) 的一个线子丛，记相应正交投影为 \(P_x\)。令
	\[
		C_x=2(I-P_x),\qquad f_x(z)=zI+C_x\quad(|z|=1),
	\]
	并由 \(f\) 在 \(S^2\times S^2\) 上粘接一个秩二丛 \(V\)。若 \(\beta_1,\beta_2\) 是两因子上 Bott 元的拉回，求 \([V]\)。证明：虽然每个固定 \(x\) 上的限制 \(V|_{\{x\}\times S^2}\) 都同构于 \(\underline{\mathbb C}\oplus H\)，\(V\) 仍不与 \(W=\underline{\mathbb C}\oplus\operatorname{pr}_2^*H\) 稳定同构。
\end{exercise}
\begin{solution}
	取 \(\mathbb C^2\) 的标准 Hermitian 内积。\(\gamma^\perp\) 的第一陈类为 \(-c_1(\gamma)=c_1(H)\)，由球面上线丛的分类，它同构于 \(H\)，故所要求的线子丛确实存在。\(C\) 在 \(H\) 上为零，在其正交补 \(H^\perp\) 上为乘以2；\(zI+C\) 的两个特征值分别为 \(z\) 与 \(z+2\)，在单位圆上均不为零。

	在 \(H\) 上，粘接函数 \(z\) 给出 \(H\boxtimes H\)；在 \(H^\perp\) 上，\(2+sz\)（\(0\le s\le1\)）始终非零，把粘接函数变为常数2，再改换一个半球上的平凡化即可去掉这一常数。于是
	\[
		V\cong\operatorname{pr}_1^*(H^\perp)\oplus(H\boxtimes H).
	\]
	利用 \([H^\perp]=2-[H]=1-\beta_1\)，得到
	\[
		[V]=(1-\beta_1)+(1+\beta_1)(1+\beta_2)
		=2+\beta_2+\beta_1\beta_2.
	\]
	固定 \(x\) 时，两条来自第一因子的线都成为固定的一维向量空间，所以限制确为 \(\underline{\mathbb C}\oplus H\)。但是 \([W]=2+\beta_2\)，两者之差为 \(\beta_1\beta_2\)。由\cref{g2:thm:k-basic-product}，\(K^0(S^2\times S^2)=K^0(S^2)\oplus K^0(S^2)\beta_2\)，而 \(\beta_1\ne0\)，故这个差非零。稳定同构会给出相同的 K 类，因此不可能成立。这个例子说明，逐点计算绕数还没有记录单位圆内谱空间怎样随参数组成子丛。
\end{solution}

\begin{exercise}\label{g2:ex:k-polynomial-winding}
	设 \(q(z)=z^2-3z+1\)，\(m\) 是非负整数。计算单位圆上粘接函数 \(z^{-m}q(z)\) 所定义的线丛的 K 类。写出正文给出的 \(q\) 的一次矩阵线性化，并求其单位圆内谱子丛的秩。说明多项式次数为何不足以决定该秩。
\end{exercise}
\begin{solution}
	两个根是 \(r_-=(3-\sqrt5)/2\) 与 \(r_+=(3+\sqrt5)/2\)，满足 \(0<r_-<1<r_+\)。在单位圆上，\(z-sr_-\) 始终不为零，把 \(z-r_-\) 同伦到 \(z\)；\(sz-r_+\) 也始终不为零，把 \(z-r_+\) 同伦到常数 \(-r_+\)。因此 \(q(z)\) 的粘接同伦于 \(-r_+z\)。去掉非零常数后，\(z^{-m}q(z)\) 的线丛就是 \(H^{1-m}\)，所以
	\[
		[C(\mathbb C,z^{-m}q)]=1+(1-m)\beta.
	\]
	这里用了 \(\beta^2=0\)，因而整数幂均满足 \([H]^j=1+j\beta\)。

	按\cref{g2:eq:k-linearization-matrix}，可以取
	\[
		L_2(q)(z)=
		\begin{pmatrix}
			1&-z&0\\
			0&1&-z\\
			1&-3&1
		\end{pmatrix}.
	\]
	线性化公式给出 \([C(\mathbb C^3,L_2(q))]=2+[H]=3+\beta\)。另一方面，谱分解给出 \(3+\operatorname{rank}S(L_2(q))\,\beta\)。由 \(\widetilde K^0(S^2)=\mathbb Z\beta\)，比较系数即得 \(\operatorname{rank}S(L_2(q))=1\)。\(q\) 虽是二次多项式，却只有一个根位于单位圆内；圆外因子可同伦为非零常数，不贡献 Bott 元。比如 \(z^2\) 同样是二次多项式，所给类却为 \(1+2\beta\)，其对应谱秩为2。
\end{solution}

\begin{exercise}\label{g2:ex:chern-endomorphism-rank-two}
	设 \(E\to X\) 是有限 CW 空间上的秩二复向量丛，\(c_1(E)=0\)，且 \(\dim X\le8\)。计算 \(\operatorname{ch}(\operatorname{End}E)\)，并说明其中为何没有次数为 \(2\) 和 \(6\) 的分量。
\end{exercise}
\begin{solution}
	在分裂空间上设 \(E\) 的根为 \(x,y\)。由 \(c_1(E)=0\) 得 \(y=-x\)，从而 \(c_2(E)=xy=-x^2\)。因为 \(\operatorname{End}E=E^*\otimes E\)，它的四个根是 \(0,0,2x,-2x\)。于是
	\[
		\operatorname{ch}(\operatorname{End}E)=2+e^{2x}+e^{-2x}
		=4+4x^2+\frac43x^4
		=4-4c_2(E)+\frac43c_2(E)^2.
	\]
	超过次数 \(8\) 的项因空间维数而为零。次数 \(2\) 和 \(6\) 对应根的一次和三次幂和；成对的 \(2x,-2x\) 在奇数次幂中抵消。也可用 \(\operatorname{End}E\cong(\operatorname{End}E)^*\) 及陈特征的对偶公式，看出所有 \(\operatorname{ch}_{2j+1}\) 在有理系数下消失。分裂空间的上同调拉回单射，故等式在 \(X\) 上成立。
\end{solution}

\begin{exercise}\label{g2:ex:chern-lens-torsion-line}
	设 \(q\ge2\)，\(X=M_q=S^1\cup_{\deg q}D^2\)。取复线丛 \(L\to X\)，使 \(c_1(L)\) 为 \(H^2(X;\mathbb Z)\cong\mathbb Z/q\) 的一个生成元。
	\begin{enumerate}
		\item 证明 \(L^{\otimes q}\) 平凡，但 \(L\) 不稳定平凡。
		\item 证明 \(u=[L]-1\in\widetilde K^0(X)\) 的加法阶恰为 \(q\)，而 \(\operatorname{ch}(u)=0\)。区分第一问的张量幂与这里的加法倍数。
	\end{enumerate}
\end{exercise}
\begin{solution}
	1. \Cref{g2:prop:line-bundle-cohomology-classification}给出复线丛的上同调分类，第21章将其二次分类类识别为第一陈类；线丛张量积的第一陈类相加。因此 \(c_1(L^{\otimes q})=q c_1(L)=0\)，故 \(L^{\otimes q}\) 平凡。若 \(L\oplus\varepsilon^r\cong\varepsilon^{r+1}\)，取行列式线丛便得到 \(L\cong\varepsilon^1\)，这与 \(c_1(L)\ne0\) 矛盾。

	2. 若 \(nu=0\)，稳定等式意味着 \(L^{\oplus n}\) 与秩 \(n\) 平凡丛稳定同构（先取 \(n>0\)）。对这个直和取第一陈类，得到 \(n c_1(L)=0\)，所以 \(q\mid n\)。另一方面，\cref{g2:exm:k-moore-torsion}已经计算 \(\widetilde K^0(M_q)\cong\mathbb Z/q\)，所以 \(qu=0\)。二者合起来说明 \(u\) 的阶正好是 \(q\)。

	\(c_1(L)\) 是扭元，进入有理上同调后为零。因此 \(\operatorname{ch}(L)=e^{c_1(L)}=1\)，\(\operatorname{ch}(u)=0\)。第一问给出乘法关系 \([L]^q=1\)，第二问给出加法关系 \(q([L]-1)=0\)；后一个关系在本题中来自 \(K\) 群的计算，不能只把前一个等式中的幂换成倍数。
\end{solution}

\begin{exercise}\label{g2:ex:circle-derivative-zeroth-order}
	在复值周期函数上考虑 \(P_a=\mathrm d/\mathrm d\theta+a\)，其中 \(a\in\mathbb C\) 为常数，圆周周期为 \(2\pi\)。求其核与余核的维数，证明指标恒为零，并解释当 \(a\) 经过 \(i\mathbb Z\) 时，解空间维数的跳跃为何不违背指标的稳定性。
\end{exercise}
\begin{solution}
	齐次方程的解为 \(f(\theta)=Ce^{-a\theta}\)。它具有周期 \(2\pi\) 当且仅当 \(e^{-2\pi a}=1\)，即 \(a\in i\mathbb Z\)。所以核在这些参数处为一维，在其他参数处为零。

	对非齐次方程，积分因子给出
	\[
		f(\theta)=e^{-a\theta}\left(C+\int_0^\theta e^{at}g(t)\,\mathrm dt\right).
	\]
	周期条件等价于
	\[
		(e^{2\pi a}-1)C=\int_0^{2\pi}e^{at}g(t)\,\mathrm dt.
	\]
	若 \(a\notin i\mathbb Z\)，对任意 \(g\) 都唯一确定 \(C\)，故余核为零。若 \(a\in i\mathbb Z\)，可解条件是右端为零；这是一条非零的线性条件，例如取 \(g(t)=e^{-at}\) 时积分为 \(2\pi\)。所以余核为一维。由微分方程反复求导可知，在函数值匹配后各阶导数也周期匹配，所得解确为光滑周期函数。

	因此两种情况下核、余核的维数分别为 \((0,0)\) 与 \((1,1)\)，指标始终为零。参数经过 \(i\mathbb Z\) 时两者同时跳跃，差保持不变。所有 \(P_a\) 的主符号都是 \(i\xi\)，由式\eqref{g2:eq:circle-derivative-symbol-zero}，符号类恒为零，也给出同一拓扑指标。
\end{solution}

\begin{exercise}\label{g2:ex:adjoint-symbol-inverse}
	设 \(M\) 为闭光滑流形，\(P:C^\infty(M,E)\to C^\infty(M,F)\) 为复向量丛间的椭圆微分算子。选 Riemann 度量及 Hermitian 丛度量，以积分恒等式 \(\langle Pu,v\rangle_{L^2}=\langle u,P^*v\rangle_{L^2}\) 定义形式伴随。证明
	\[
		[\sigma(P^*)]=-[\sigma(P)]\quad\text{于 }K_c^0(\mathrm T^*M),
		\qquad\operatorname{ind}_{\mathrm a}(P^*)=-\operatorname{ind}_{\mathrm a}(P).
	\]
	分别从三元组和核、余核说明这两个负号的来源。
\end{exercise}
\begin{solution}
	分部积分把每个最高阶偏导数变为其负号，系数同时取 Hermitian 伴随；导数落在系数或体积密度上产生的都是低阶项。在采用 \(i\xi\) 的主符号约定下，这正给出 \(\sigma_m(P^*)=\sigma_m(P)^*\)。记余球边界上的同构为 \(\alpha:\pi^*E\to\pi^*F\)，则伴随的三元组为 \((\pi^*F,\pi^*E,\alpha^*)\)。

	原三元组的负元为 \((\pi^*F,\pi^*E,\alpha^{-1})\)。在边界上，\(\alpha^*\alpha\) 是 \(\pi^*E\) 的正定自同构，故
	\[
		\alpha_t=\bigl((1-t)I+t\alpha^*\alpha\bigr)\alpha^{-1},\qquad0\le t\le1,
	\]
	始终可逆，并连接 \(\alpha^{-1}\) 与 \(\alpha^*\)。所以两份三元组同伦，得到符号类的负号。

	分析一侧，闭像与椭圆正则性给出 \(\operatorname{coker}P\cong\ker P^*\) 以及 \(\operatorname{coker}P^*\cong\ker P\)。因而
	\[
		\operatorname{ind}_{\mathrm a}(P^*)
		=\dim\ker P^*-\dim\ker P
		=-\operatorname{ind}_{\mathrm a}(P).
	\]
	这里负号来自交换核与余核；拓扑一侧则来自交换三元组的两张丛并反转边界同构。指标定理将这两种说明联系起来。
\end{solution}


% !TeX root = ../Book2.tex
% 纲要标题：# 第二卷附录
% 纲要位置：几何学计划纲要.md，第 1972 行。
% 正文状态：未撰写；本文件目前是框架。

% 纲要内容开始
% 
% 纲要内容结束

\appendix
\renewcommand{\thechapter}{\Alph{chapter}}
\ctexset{chapter={number=\Alph{chapter}}}
\BookAppendixPart{g2:part:appendices}

% !TeX root = ../../Book2.tex
% 纲要标题：## 附录 A　结、链环与低维拓扑
% 纲要位置：几何学计划纲要.md，第 1974 行。
% 正文状态：未撰写；本文件目前是框架。
% 补充范围与先修：三维分解采用紧、连通、可定向三维流形的明确版本，边界环面情形另核条件；使用第一卷基本群、覆盖和本附录结补空间、Dehn 填充。素分解、JSJ 与几何化的深层定理准确外引；几何模型接第三卷空间形式及几何群论附录，Ricci 流输入见第三卷第27章，不承担一般三维分类证明。

% 纲要内容开始
% 
% > 不承担一般三维分类；从两个实际可计算不变量进入更深结构。
% 
% > 补充写作范围与外部输入：补充紧、连通、可定向三维流形的素分解、JSJ 与几何化图景，边界环面情形明确条件；深层定理准确外引。依赖第一卷基本群、覆盖及本附录的结补空间和 Dehn 填充；几何模型接第三卷空间形式及几何群论附录，Ricci 流输入见第三卷第27章，不承担一般三维分类证明。
% 
% * **结图与补空间**
%   * 环境同位及 Reidemeister 操作
%   * Wirtinger 展示及结群
%   * 链环数的相交定义
% 
% * **Seifert 曲面与 Alexander 不变量**
%   * Seifert 曲面的构造
%   * 无限循环覆盖和 Alexander 模
%   * Seifert 矩阵的多项式计算
%   * 三叶结和八字结的核验
% 
% * **三维手术**
%   * Dehn 填充和斜率
%   * 手术图及边界变化
%   * 三维分类所需的额外定理
% 
% * **三维分解与几何化**
%   * 连通和与素分解
%   * 不可约性与不可压缩环面
%   * JSJ 分解的选定陈述
%   * 八种几何与几何化定理
%   * 结补空间、Dehn 填充与几何变化
% 
% 
% ---
% 
% 纲要内容结束

\chapter{结、链环与低维拓扑}
\label{g2:app:A}
\BookChapterNumber{g2:app:A}

待撰写

\section{结图与补空间}
\label{g2:sec:A01}
待撰写

\subsection{环境同位及 Reidemeister 操作}
\label{g2:subsec:A01:01}
待撰写

\subsection{Wirtinger 展示及结群}
\label{g2:subsec:A01:02}
待撰写

\subsection{链环数的相交定义}
\label{g2:subsec:A01:03}
待撰写

\section{Seifert 曲面与 Alexander 不变量}
\label{g2:sec:A02}
待撰写

\subsection{Seifert 曲面的构造}
\label{g2:subsec:A02:01}
待撰写

\subsection{无限循环覆盖和 Alexander 模}
\label{g2:subsec:A02:02}
待撰写

\subsection{Seifert 矩阵的多项式计算}
\label{g2:subsec:A02:03}
待撰写

\subsection{三叶结和八字结的核验}
\label{g2:subsec:A02:04}
待撰写

\section{三维手术}
\label{g2:sec:A03}
待撰写

\subsection{Dehn 填充和斜率}
\label{g2:subsec:A03:01}
待撰写

\subsection{手术图及边界变化}
\label{g2:subsec:A03:02}
待撰写

\subsection{三维分类所需的额外定理}
\label{g2:subsec:A03:03}
待撰写

\section{三维分解与几何化}
\label{g2:sec:A04}
待撰写

\subsection{连通和与素分解}
\label{g2:subsec:A04:01}
待撰写

\subsection{不可约性与不可压缩环面}
\label{g2:subsec:A04:02}
待撰写

\subsection{JSJ 分解的选定陈述}
\label{g2:subsec:A04:03}
待撰写

\subsection{八种几何与几何化定理}
\label{g2:subsec:A04:04}
待撰写

\subsection{结补空间、Dehn 填充与几何变化}
\label{g2:subsec:A04:05}
待撰写

% !TeX root = ../../Book2.tex
% 纲要标题：## 附录 B　分层空间、局部拓扑与交叉同调
% 纲要位置：几何学计划纲要.md，第 1941 行。
% 正文状态：未撰写；本文件目前是框架。

% 纲要内容开始
% 
% > 基础部分按“分层—控制—局部拓扑—允许链”深入。Whitney 和 Thom 型条件各有定义；一般层论对偶和中间延拓不冒充本附录已完成的全套理论。
% 
% * **分层的几何条件**
%   * 局部有限分层与边界条件
%   * Whitney (a)、(b) 的切空间极限
%   * 横截切片与链接
%   * 代数、解析和光滑分层的区别
% 
% * **控制数据与局部平凡性**
%   * 管状控制数据及相容性
%   * 分层向量场和积分
%   * 第一同位引理的选定固有版本
%   * 缺少控制或固有性时的失败
% 
% * **锥、链接与允许链**
%   * 局部锥模型与普通对偶的失败
%   * 允许函数（perversity）的维数限制与允许链
%   * 锥公式的代表性计算
%   * 交叉同调的对偶定理：精确陈述及证明输入
% 
% * **层方法与消失圈**
%   * 可构造层和局部系统
%   * 交叉链层与中间延拓的动机
%   * 平滑族与退化族的附近圈
%   * 结点退化的消失圈计算
% 
% ---
% 
% 纲要内容结束

\chapter{分层空间、局部拓扑与交叉同调}
\label{g2:app:B}
\BookChapterNumber{g2:app:B}

待撰写

\section{分层的几何条件}
\label{g2:sec:B01}
待撰写

\subsection{局部有限分层与边界条件}
\label{g2:subsec:B01:01}
待撰写

\subsection{Whitney (a)、(b) 的切空间极限}
\label{g2:subsec:B01:02}
待撰写

\subsection{横截切片与链接}
\label{g2:subsec:B01:03}
待撰写

\subsection{代数、解析和光滑分层的区别}
\label{g2:subsec:B01:04}
待撰写

\section{控制数据与局部平凡性}
\label{g2:sec:B02}
待撰写

\subsection{管状控制数据及相容性}
\label{g2:subsec:B02:01}
待撰写

\subsection{分层向量场和积分}
\label{g2:subsec:B02:02}
待撰写

\subsection{第一同位引理的选定固有版本}
\label{g2:subsec:B02:03}
待撰写

\subsection{缺少控制或固有性时的失败}
\label{g2:subsec:B02:04}
待撰写

\section{锥、链接与允许链}
\label{g2:sec:B03}
待撰写

\subsection{局部锥模型与普通对偶的失败}
\label{g2:subsec:B03:01}
待撰写

\subsection{允许函数（perversity）的维数限制与允许链}
\label{g2:subsec:B03:02}
待撰写

\subsection{锥公式的代表性计算}
\label{g2:subsec:B03:03}
待撰写

\subsection{交叉同调的对偶定理：精确陈述及证明输入}
\label{g2:subsec:B03:04}
待撰写

\section{层方法与消失圈}
\label{g2:sec:B04}
待撰写

\subsection{可构造层和局部系统}
\label{g2:subsec:B04:01}
待撰写

\subsection{交叉链层与中间延拓的动机}
\label{g2:subsec:B04:02}
待撰写

\subsection{平滑族与退化族的附近圈}
\label{g2:subsec:B04:03}
待撰写

\subsection{结点退化的消失圈计算}
\label{g2:subsec:B04:04}
待撰写

% !TeX root = ../../Book2.tex
% 纲要标题：## 附录 C　等变拓扑、稳定结构与计算
% 纲要位置：几何学计划纲要.md，第 2036 行。
% 正文状态：初稿；正文、证明、例子及习题解答已写，尚未正式审核。
% 补充范围与先修：等变 de Rham 与 Cartan 模型固定紧 Lie 群的光滑作用及实系数；与 Borel 构造的比较明确所用有限维逼近和不变性条件。环面局部化另列紧性、定向、固定点分支及系数局部化条件；具体积分公式的主要证明和辛几何应用放第四卷矩映射与局部化，不与稳定同伦、持续同调支线混用。

% 纲要内容开始
% 
% > 保持三条支线独立，不用一个共同定义掩盖三种不同的稳定化。
% 
% > 补充写作范围与外部输入：等变 de Rham 与 Cartan 模型固定紧 Lie 群光滑作用及实系数，说明与 Borel 构造比较的条件；环面局部化另列系数、紧性、定向和固定点分支条件。积分公式的主要证明及辛几何应用放第四卷矩映射与局部化，不与稳定同伦、持续同调支线混用。
% 
% * **等变上同调**
%   * Borel 构造及纤维序列
%   * 固定点和自由作用的对照
%   * 圆群等变球面算例
%   * 等变 de Rham 复形与 Cartan 模型
%   * 紧 Lie 群作用下的比较条件
%   * 环面局部化所需的系数与紧性条件
% 
% * **谱与稳定同伦**
%   * 悬挂稳定化和谱
%   * 稳定同伦群及系数谱
%   * 广义同调理论的进一步输入
% 
% * **持续同调的代数结构**
%   * 有限滤过复形及诱导映射
%   * 域系数条形码的分解
%   * 边界矩阵约化
%   * 稳定性需要的距离与有限性条件
% 
% ---
% 
% 纲要内容结束

\chapter{等变拓扑、稳定结构与计算}
\label{g2:app:C}
\BookChapterNumber{g2:app:C}

\begin{chapterabstract}
	群作用、悬挂与滤过给拓扑不变量提出了三类进一步的问题。等变上同调通过 Borel 构造保存稳定子信息，圆群作用的球面计算将它化为射影丛关系；Cartan 模型以微分形式实现同一理论，局部化说明固定点数据何时足以恢复整体。谱将反复悬挂后的同伦数据组织起来，球谱、普通同调谱与周期 K 理论谱展示不同的系数行为。持续同调则保存滤过中同调类的出现与消失，以区间分解和边界矩阵约化得到条形码，并用交错与瓶颈距离表达扰动稳定性。
\end{chapterabstract}

\begin{learninggoals}
	\begin{itemize}
		\item 从关联丛构造 \(X_G\)，比较自由轨道与固定点，计算圆群等变球面的环及固定点限制。
		\item 写出 Cartan 复形的分次、群不变性与微分，核验平方为零，说明紧群比较及环面局部化的适用条件。
		\item 从悬挂过渡映射定义谱的稳定同伦群，用球谱、普通同调谱和 K 理论谱辨认系数及次数。
		\item 将有限滤过转为向量空间及线性映射，分解为区间模，并从秩函数读出区间重数。
		\item 手工约化边界矩阵，解释出生与死亡配对，使用交错和瓶颈距离估计过滤值扰动的影响。
	\end{itemize}
\end{learninggoals}

当圆群旋转一个球面时，纬圈都成为轨道，南北两极却保持不动。若只计算轨道空间的上同调，这种区别几乎全部消失。另一方面，将一个带基点空间反复悬挂，有些原本依赖维数的同伦群会进入稳定范围。再看一个逐次加入顶点、边和面的复形：最终同调能够说明剩下几个洞，却不能告诉我们这些洞何时出现、何时被填掉。本附录分别为这三个问题建立合适的对象。

\ProseParagraphBreak

三节可以按需要分别阅读。等变部分使用第7章 de Rham 理论、第19章主丛分类、第20—21章 Thom 类与射影丛公式；稳定同伦部分使用第8—12章的悬挂、同伦群与表示空间，K 理论例子另接第24章；持续同调部分主要使用第2章链复形及有限维线性代数。每一条支线的“稳定”都有自己的含义，下面的对应关系有助于在阅读中保持对象和任务的一致。

\begin{table}[htbp]
	\centering
	\caption{三条支线的对象与比较方式}
	\label{g2:tab:appendix-c-three-directions}
	\begin{tabularx}{\textwidth}{@{}lXX@{}}
		\toprule
		支线 & 加入的数据 & 所研究的比较 \\
		\midrule
		等变上同调 & 群作用、稳定子与通用主丛 & 与普通商、固定点及微分形式模型的比较 \\
		谱与稳定同伦 & 空间序列及悬挂结构映射 & 反复悬挂后的同伦类及广义同调 \\
		持续同调 & 同一滤过中各次包含诱导的映射 & 区间分解以及过滤值扰动后的条形码距离 \\
		\bottomrule
	\end{tabularx}
\end{table}

本附录完整展开代表模型和有限代数计算。紧群的等变 de Rham 比较、紧光滑环面作用的局部化、稳定同伦范畴的基础构造以及条形码匹配定理，在实际使用处列明所选版本和证明输入。等变积分的主要证明及辛约化应用接第四卷第12.6节；持续同调与 Floer 作用滤过的联系接第四卷附录 E。

\section{等变上同调}
\label{g2:sec:C01}

一点上的平凡圆群作用与圆周上的自由旋转，轨道空间都是一点。前者的点被整个圆群固定，后者的点只有单位元保持不动，这个差别在取商时消失了。等变上同调用一个同时保存空间与群作用的构造代替普通商空间。圆群作用于二球面的例子会把这种信息写成一条明确的环关系。

\subsection{Borel 构造及纤维序列}
\label{g2:subsec:C01:01}

取紧矩阵 Lie 群 \(G\)，并取第19章构造的通用主右丛 \(EG\to BG\)。\cref{g2:prop:universal-matrix-principal-bundle}已经证明：\(EG\) 可缩，右 \(G\) 作用自由，而且轨道映射局部平凡。可缩性使这一辅助空间不增加普通同伦信息；自由性则让对角作用能够取一个保留群数据的商。

\begin{definition}\label{g2:def:borel-construction}
	设 \(X\) 是左 \(G\) 空间。其\term[Borel gouzao]{Borel 构造}[Borel construction]，也称同伦商，定义为
	\[
		X_G=EG\times_GX=(EG\times X)/\!\sim,
		\qquad(eg,x)\sim(e,gx).
	\]
	等价地，商掉的是自由右作用 \((e,x)g=(eg,g^{-1}x)\)。对含幺交换环 \(R\)，定义\term[dengbianshangtongdiao]{等变上同调}[equivariant cohomology]
	\[
		H_G^*(X;R)=H^*(X_G;R).
	\]
	右端是普通奇异上同调，乘法为杯积。\symidx{\(X_G=EG\times_GX\)：Borel 构造}\symidx{\(H_G^*(X;R)\)：等变上同调}
\end{definition}

这个构造沿用第19章的关联丛关系，左右作用的次序因而已有固定含义。名称中的 Borel 指 Armand Borel。\footnote{Armand Borel（1923--2003），瑞士数学家，研究代数拓扑与线性代数群。}投影
\[
	p:X_G\longrightarrow BG,\qquad[e,x]\longmapsto[e]
\]
把 \(X_G\) 组织成纤维为 \(X\) 的关联丛：若 \(s:V\to EG\) 是局部截面，便有丛图
\[
	V\times X\longrightarrow p^{-1}(V),\qquad(b,x)\longmapsto[s(b),x].
\]
若 \(e=s(b)g\)，其逆把 \([e,x]\) 送到 \((b,gx)\)。这也直接核验了丛图与商关系相容。

\ProseParagraphBreak

由\cref{g2:thm:bundle-serre}，上述纤维丛是 Serre 纤维化。选定 \(e_0\in EG\) 后，\(x\mapsto[e_0,x]\) 把 \(X\) 识别为 \([e_0]\) 上的纤维，故有纤维序列
\begin{equation}\label{g2:eq:borel-fiber-sequence}
	X\longrightarrow X_G\xrightarrow{p}BG.
\end{equation}
这里“纤维序列”说明 \(X\) 也代表 \(p\) 的同伦纤维；它给出第8章的同伦长正合列，并可使用第11章的纤维化谱序列。它并没有把 \(X_G\) 写成 \(X\times BG\)，作用的粘合数据仍在过渡函数里。

\begin{proposition}\label{g2:prop:borel-model-independence}
	若 \(EG\) 与 \(E'G\) 都是可缩、自由且轨道映射局部平凡的通用主 \(G\) 空间，则它们所给的 \(H_G^*(X;R)\) 自然同构。连续等变映射 \(a:X\to Y\) 诱导上同调环同态 \(a_G^*:H_G^*(Y;R)\to H_G^*(X;R)\)，等变同伦的映射诱导相同同态。
\end{proposition}
\begin{proof}
	比较空间取 \((EG\times E'G\times X)/G\)，其中右作用为
	\[
		(e,e',x)g=(eg,e'g,g^{-1}x).
	\]
	忘记 \(e'\) 或 \(e\) 分别得到映到两种 Borel 构造的纤维丛，纤维分别为 \(E'G\) 与 \(EG\)。例如在第一种 Borel 构造上，这张丛来自将主丛 \(EG\to BG\) 沿 \(p\) 拉回后，以 \(E'G\) 为关联纤维；局部丛图给出所需平凡化。两个纤维都可缩，纤维化同伦长正合列说明两个投影在全部同伦群与道路分支上诱导同构。依\cref{g2:lem:weak-equivalence-homology}，弱同伦等价诱导奇异同调同构，再由上同调万有系数定理得到任意 \(R\) 系数上同调同构；这些同构保持杯积。

	等变映射给出 \([e,x]\mapsto[e,a(x)]\)，等变同伦则以同一公式下降为 Borel 构造间的同伦。投影比较与这些映射交换，因而同构自然；对三个辅助空间用三重乘积比较，又保证不同次比较彼此相容。
\end{proof}

投影的拉回 \(p^*:H^*(BG;R)\to H_G^*(X;R)\) 使等变上同调成为 \(H^*(BG;R)\) 上的代数。纤维包含则给出
\[
	j^*:H_G^*(X;R)\longrightarrow H^*(X;R),
\]
称为忘却映射。它把辅助的群参数忘掉。因 \(p\circ j\) 是常映射，来自 \(H^{>0}(BG;R)\) 的类都映为零；但是忘却映射未必满射，因此一般不能只把群参数置零就得到全部普通上同调。

\subsection{固定点和自由作用的对照}
\label{g2:subsec:C01:02}

若作用平凡，商关系只作用于 \(EG\) 这一因子，故直接得到
\[
	X_G=BG\times X.
\]
例如 \(H_G^*(\mathrm{pt};R)=H^*(BG;R)\)。对域 \(k\) 和有限 CW 复形 \(X\)，Künneth 公式还给出
\[
	H_G^*(X;k)=H^*(BG;k)\otimes_kH^*(X;k)
\]
作为分次代数的同构；乘法含通常的分次交换符号。有限性保证每个总次数只需处理有限维纤维同调，普通上同调的张量表达因而适用。

\begin{proposition}\label{g2:prop:free-action-borel-quotient}
	设紧 Lie 群 \(G\) 在 \(X\) 上自由左作用，且 \(X\to X/G\) 是局部平凡主丛。则
	\[
		q:X_G\longrightarrow X/G,\qquad[e,x]\longmapsto Gx
	\]
	是纤维为 \(EG\) 的局部平凡纤维丛，因而诱导自然环同构
	\[
		q^*:H^*(X/G;R)\xrightarrow{\ \cong\ }H_G^*(X;R).
	\]
\end{proposition}
\begin{proof}
	在 \(X/G\) 的开集 \(V\) 上取截面 \(s\)。每个 \(x\) 唯一写成 \(g s(b)\)，其中 \(b=Gx\)。于是
	\[
		q^{-1}(V)\longrightarrow V\times EG,\qquad
		[e,g s(b)]\longmapsto(b,eg)
	\]
	良定义，逆为 \((b,e)\mapsto[e,s(b)]\)。因此 \(q\) 是所述纤维丛。\(EG\) 可缩，故由同伦长正合列及弱等价的同调判据，\(q\) 在上同调上诱导同构；自然性和环结构均来自连续映射的拉回。
\end{proof}

对于圆群 \(T=S^1\)，可取 \(ET=S^\infty\)、\(BT=\mathbb CP^\infty\)，右作用是单位复数的标量乘法。若 \(T\) 在一点上平凡作用，等变上同调是 \(H^*(BT;\mathbb Z)\cong\mathbb Z[u]\)，其中 \(|u|=2\)。若 \(T\) 在自身上以左乘作用，映射 \([e,z]\mapsto ez\) 把 \(T_T\) 同胚到 \(ET\)，所以其等变上同调只有零次的 \(\mathbb Z\)。这两个作用的轨道空间都是一点，等变上同调却清楚地区分了它们。后一例的忘却映射不能得到 \(H^1(S^1;\mathbb Z)\) 的生成元。

\ProseParagraphBreak

固定点则产生相反的现象。若 \(x\in X\) 被整个 \(G\) 固定，映射
\[
	s_x:BG\longrightarrow X_G,\qquad[e]\longmapsto[e,x]
\]
满足 \(p\circ s_x=\operatorname{id}\)，所以 \(p^*\) 单射。更一般地，令 \(X^G=\{x:gx=x\text{ 对所有 }g\in G\}\)，其包含诱导
\[
	H_G^*(X;R)\longrightarrow H_G^*(X^G;R)
	=H^*(BG\times X^G;R).
\]
固定点上的这些限制可以用于检测整体等变类。

\begin{proposition}\label{g2:prop:borel-homogeneous-orbit}
	设 \(H\subset G\) 为闭子群，\(G\) 在 \(G/H\) 上左乘。则
	\[
		(G/H)_G\cong EG/H,
		\qquad H_G^*(G/H;R)\cong H^*(BH;R).
	\]
	在这个识别下，\(H^*(BG;R)\) 的作用来自自然投影 \(EG/H\to EG/G\) 的拉回。
\end{proposition}
\begin{proof}
	定义 \([e,gH]\mapsto egH\)。关系 \([ea,gH]=[e,agH]\) 的两边都映到 \(eagH\)，故映射良定义；其逆为 \(eH\mapsto[e,H]\)，两者在商拓扑中连续且互逆。\(EG\) 作为 \(H\) 空间仍自由、可缩，\(EG\to EG/H\) 由闭子群商的局部截面给出主 \(H\) 丛，因此可以作 \(EH\to BH\) 的模型。最后的代数结构由同一个投影公式得到。
\end{proof}

这个公式把每条轨道上留下的信息解释为稳定子 \(H\) 的分类空间。当 \(H=G\) 时得到一个固定点的 \(BG\)，当 \(H=\{1\}\) 时得到自由轨道的可缩空间。介于二者之间的稳定子也会留下相应信息。用普通轨道空间替换 Borel 构造，会把这些不同的分类空间全部压成一点。

\subsection{圆群等变球面算例}
\label{g2:subsec:C01:03}

取 \(T=S^1\) 在 \(S^2=\mathbb CP^1\) 上的作用
\[
	z\cdot[z_0:z_1]=[z_0:zz_1].
\]
它的固定点为 \(N=[1:0]\)、\(S=[0:1]\)，其余轨道都是纬圈，轨道空间为闭区间。为了计算等变上同调，先把这张球面上的作用转为 \(BT\) 上的一张射影丛。

\ProseParagraphBreak

记 \(\mathbb C_1\) 为权 \(1\) 的表示，即 \(z\) 作用为乘以 \(z\)，并固定
\begin{equation}\label{g2:eq:equivariant-circle-generator}
	L=ET\times_T\mathbb C_1,\qquad u=c_1(L)\in H^2(BT;\mathbb Z).
\end{equation}
\symidx{\(u=c_1(ET\times_T\mathbb C_1)\)：圆群的二次系数生成元}
在 \(ET=S^\infty\) 的右标量模型中，\([e,w]\mapsto we\) 将 \(L\) 识别为 \(\mathbb CP^\infty\) 的重言线丛 \(\gamma\)。因此若第21章的正生成元记作 \(h=c_1(\gamma^*)\)，本节的 \(u=-h\)。给出这项符号选择，是为了同时读出表示权、固定点限制与后面的 Cartan 变量。

\begin{theorem}\label{g2:thm:equivariant-sphere-ring}
	对上述圆群作用及式\eqref{g2:eq:equivariant-circle-generator}的生成元，存在二次类 \(x\)，使
	\begin{equation}\label{g2:eq:equivariant-sphere-ring}
		H_T^*(S^2;\mathbb Z)\cong\mathbb Z[u,x]/\bigl(x(x+u)\bigr),
		\qquad |u|=|x|=2.
	\end{equation}
	作为 \(\mathbb Z[u]\) 模，它以 \(1,x\) 为基。忘却映射将 \(u\) 送到零、将 \(x\) 送到 \(H^2(\mathbb CP^1;\mathbb Z)\) 的复定向正生成元。固定点限制为
	\begin{equation}\label{g2:eq:sphere-fixed-restriction}
		i^*(a(u)+x b(u))=\bigl(a(u),a(u)-u b(u)\bigr).
	\end{equation}
\end{theorem}
\begin{proof}
	球面是表示 \(\mathbb C_0\oplus\mathbb C_1\) 的线子空间射影空间。取关联丛与取线子空间相容，故
	\[
		(S^2)_T\cong\mathbb P(\underline{\mathbb C}\oplus L).
	\]
	具体地，\([e,\ell]\) 对应于关联纤维中由 \([e,v]\)、\(v\in\ell\) 组成的线。更换代表 \(e\) 时，\(\ell\) 同时按群作用变化，所以这条线保持不变；局部平凡化给出连续的逆。

	令 \(\lambda\) 为这张射影丛上的重言线丛，取 \(x=c_1(\lambda^*)\)。\cref{g2:thm:projective-bundle}给出模基 \(1,x\)，\cref{g2:def:chern-classes-projective}则将射影丛关系的系数识别为 \(E=\underline{\mathbb C}\oplus L\) 的陈类。由直和公式，\(c_1(E)=u\)、\(c_2(E)=0\)，于是
	\[
		x^2+ux=0.
	\]
	\(H^*(BT;\mathbb Z)=\mathbb Z[u]\) 来自\cref{g2:thm:complex-projective-ring}在无限射影空间上的计算，所以得到所述环。纤维上的 \(\lambda\) 是通常的重言线丛，因此 \(x\) 限制为正生成元。

	固定点 \(N\) 对应于 \(E\) 的平凡线，\(S\) 对应于 \(L\)。在这两个截面上，\(\lambda^*\) 分别限制为平凡线和 \(L^*\)，故
	\[
		x|_N=0,\qquad x|_S=-u.
	\]
	\(u\) 从底空间拉回，在两个截面上仍为 \(u\)。代入模基的唯一表达即得式\eqref{g2:eq:sphere-fixed-restriction}。
\end{proof}

限制到固定点究竟丢失多少信息，现在可以逐项计算。若 \(i^*(a+xb)=0\)，第一分量给 \(a=0\)，第二分量给 \(ub=0\)。\(\mathbb Z[u]\) 没有零因子，故 \(b=0\)，所以此例的固定点限制单射。它的像是
\begin{equation}\label{g2:eq:sphere-congruence-image}
	\{(P,Q)\in\mathbb Z[u]\oplus\mathbb Z[u]:P-Q\text{ 能被 }u\text{ 整除}\}.
\end{equation}
若差能被 \(u\) 整除，取 \(a=P\)、\(b=(P-Q)/u\) 就得到原像，证明这一描述也充分。两个固定点的类因而还要满足一条粘合关系。

\ProseParagraphBreak

将 \(u\) 变为可逆元后，每一对 Laurent 多项式都满足上述可解性条件，故
\[
	H_T^*(S^2;\mathbb Z)[u^{-1}]
	\cong\mathbb Z[u,u^{-1}]\oplus\mathbb Z[u,u^{-1}].
\]
右侧两份系数环由两个固定点分别贡献。对应于 \((1,0)\)、\((0,1)\) 的类是
\[
	e_N=\frac{x+u}{u},\qquad e_S=-\frac{x}{u}.
\]
利用 \(x(x+u)=0\) 可直接验证 \(e_N+e_S=1\)、\(e_Ne_S=0\)、\(e_N^2=e_N\)、\(e_S^2=e_S\)。这是局部化之后整体环分成固定点两部分的一个完整计算。

\begin{figure}[htbp]
	\centering
	% !TeX root = ../Book2.tex
\begin{tikzpicture}[>=stealth,every node/.style={font=\small}]
	\shade[ball color=blue!10] (0,0) circle (1.45);
	\draw[thick] (0,0) circle (1.45);
	\draw[dashed,gray] (-1.45,0) arc (180:0:1.45 and .4);
	\draw[blue!70!black,->] (1.45,0) arc (0:-165:1.45 and .4);
	\fill (0,1.45) circle (2pt) node[above] {\(N\)};
	\fill (0,-1.45) circle (2pt) node[below] {\(S\)};
	\draw[->] (1.85,0)--(3.15,0) node[midway,above] {取轨道};
	\draw[very thick] (3.65,-1.45)--(3.65,1.45);
	\fill (3.65,1.45) circle (2pt) node[above] {\(N\)};
	\fill (3.65,-1.45) circle (2pt) node[below] {\(S\)};
	\node[right] at (4.35,.65) {\(x|_N=0\)};
	\node[right] at (4.35,-.65) {\(x|_S=-u\)};
	\node at (2.8,-2.1) {\(i^*(a+xb)=(a,a-ub)\)};
\end{tikzpicture}

	\caption{圆群沿纬圈旋转二球面。普通商把每条纬圈压为区间中的一点；固定点限制则给出两份 \(\mathbb Z[u]\)，在局部化之前还须满足 \(P-Q\in u\mathbb Z[u]\) 的关系。}
	\label{g2:fig:equivariant-sphere}
\end{figure}

若圆群在同一个球面上平凡作用，得到的却是 \(\mathbb Z[u,x]/(x^2)\)。两种作用的普通空间相同，区别写在 \(ux\) 这一项里；取商得到的区间又没有这种二次信息。在后面的微分形式模型中，关系中的 \(u\) 将成为记录无穷小群作用的多项式变量，而 \(x\) 会有一个可显式写出的等变闭代表。

\subsection{等变 de Rham 复形与 Cartan 模型}
\label{g2:subsec:C01:04}

Borel 构造给出了等变上同调的拓扑定义。若群与流形都是光滑的，还可以寻找一种由微分形式组成的计算模型。不过，单取不变微分形式并不够：一点上的微分形式只有常数，而 \(H_{S^1}^*(\mathrm{pt};\mathbb R)=\mathbb R[u]\) 在每个非负偶数次都有类。模型必须同时容纳流形上的形式与群方向的多项式。

\ProseParagraphBreak
设 \(G\) 是 Lie 群，\(\mathfrak g\) 是其 Lie 代数，\(M\) 是有限维 Hausdorff、第二可数光滑无边流形，并给定光滑左作用 \(a_g\colon M\to M\)。本节固定基本向量场的符号
\begin{equation}\label{g2:eq:cartan-fundamental-field}
	\xi_M(x)=\left.\frac{\mathrm d}{\mathrm dt}\right|_{t=0}\exp(t\xi)\cdot x,
	\qquad \xi\in\mathfrak g.
\end{equation}
把 \(S(\mathfrak g^*)\otimes\Omega(M)\) 中的元素看成以 \(\xi\in\mathfrak g\) 为变量、取值于微分形式的多项式 \(\alpha(\xi)\)。其中一个线性多项式的次数规定为 \(2\)，一个普通 \(b\) 次形式的次数仍为 \(b\)。这样规定的原因将在微分中显现：插入 \(\xi_M\) 降低一个形式次数，同时增加一个多项式次数，合计提高一次。

\begin{definition}\label{g2:def:cartan-complex}
	\term[Cartan moxing]{Cartan 模型}[Cartan model]的次数 \(k\) 部分为
	\[
		\Omega_G^k(M)=\bigoplus_{2a+b=k}\bigl(S^a(\mathfrak g^*)\otimes\Omega^b(M)\bigr)^G.
	\]
	这里上标 \(G\) 要求多项式满足
	\begin{equation}\label{g2:eq:cartan-equivariant-polynomial}
		a_g^*\alpha(\operatorname{Ad}_g\xi)=\alpha(\xi)
		\qquad(g\in G,\ \xi\in\mathfrak g).
	\end{equation}
	在这个空间上定义\term[dengbian weifen]{等变微分}[equivariant differential]
	\begin{equation}\label{g2:eq:cartan-differential}
		(\mathrm d_G\alpha)(\xi)=\mathrm d\bigl(\alpha(\xi)\bigr)-\iota_{\xi_M}\alpha(\xi).
	\end{equation}
\end{definition}

这里的 Cartan 指 Henri Cartan。\footnote{Henri Cartan（1904--2008），法国数学家，在多复变函数、层论和代数拓扑中有重要贡献。等变上同调的 Cartan 模型来自他的工作。}第一卷使用的 Cartan 公式则属于微分形式的运算规则。为便于在这里使用，回忆第一卷第14.1节已经证明
\[
	\mathcal L_X=\mathrm d\iota_X+\iota_X\mathrm d,
\]
其中 \(\mathcal L_X\) 由向量场正向流的拉回求导定义。这个符号与式\eqref{g2:eq:cartan-fundamental-field}一致。

\begin{proposition}\label{g2:prop:cartan-is-complex}
	\(\mathrm d_G\) 保持式\eqref{g2:eq:cartan-equivariant-polynomial}的条件，是次数 \(1\) 的分次导子，并且在 \(\Omega_G(M)\) 上满足 \(\mathrm d_G^2=0\)。因而 \((\Omega_G(M),\mathrm d_G)\) 是分次交换微分代数。
\end{proposition}
\begin{proof}
	基本向量场满足 \((a_g)_*\xi_M=(\operatorname{Ad}_g\xi)_M\)，所以外微分和插入运算都与式\eqref{g2:eq:cartan-equivariant-polynomial}相容。若选择 \(\mathfrak g\) 的基 \(e_1,\ldots,e_r\)，相应线性坐标为 \(u_1,\ldots,u_r\)，则
	\[
		\mathrm d_G=\mathrm d-\sum_{j=1}^r u_j\iota_{(e_j)_M}.
	\]
	每一项的总次数都是 \(1\)。外微分与插入都是按形式次数取符号的导子，而多项式次数为偶数，故该符号也就是按总次数取的符号。

	固定 \(\xi\) 后，\(\mathrm d^2=0\)、\(\iota_{\xi_M}^2=0\) 与 Cartan 公式给出
	\[
		(\mathrm d_G^2\alpha)(\xi)=-\mathcal L_{\xi_M}\alpha(\xi).
	\]
	在等变性条件中取 \(g=\exp(t\xi)\)。由于 \(\operatorname{Ad}_{\exp(t\xi)}\xi=\xi\)，得到 \(a_{\exp(t\xi)}^*\alpha(\xi)=\alpha(\xi)\)。对 \(t\) 求导便得右端为零。
\end{proof}

这也解释了为何定义要取 \(G\) 不变部分：在任意多项式值形式上，\(\mathrm d_G^2\) 通常只是 \(-\mathcal L_{\xi_M}\)，并不为零。还应区分不变性与水平性。一个普通形式 \(\beta\) 称为\term[jiben xingshi]{基本形式}[basic form]，是指它既满足 \(a_g^*\beta=\beta\)，又满足所有 \(\iota_{\xi_M}\beta=0\)。自由适当作用下，基本形式恰是商流形上的形式的拉回。Cartan 模型要求的是式\eqref{g2:eq:cartan-equivariant-polynomial}，并未逐项要求水平；插入项正用来记录沿轨道方向的变化。

\begin{example}\label{g2:ex:cartan-point-circle}
	在一点上，\(\Omega_G(\mathrm{pt})=S(\mathfrak g^*)^G\)，微分为零。特别地，令 \(T=S^1=\mathbb R/\mathbb Z\)，取满足 \(\exp(t\xi)=e^{2\pi it}\) 的生成元，并以 \(u(\xi)=1\) 规定线性坐标，则模型为 \(\mathbb R[u]\)，\(\lvert u\rvert=2\)。这里的 \(u\) 将在下一小节与前文权 \(1\) 关联线的一阶陈类相认。

	若 \(T\) 左平移作用于自身，令 \(\theta\in\mathbb R/\mathbb Z\) 为角坐标。虽然 \(\theta\) 不是全局实值函数，\(\mathrm d\theta\) 却是全局形式；不变形式由 \(1\) 与 \(\mathrm d\theta\) 张成，且
	\[
		\mathrm d_T(\mathrm d\theta)=-u.
	\]
	复形为 \(\mathbb R[u]\otimes\Lambda(\mathrm d\theta)\)：所有正次 \(u^j\) 都是边界，奇次没有非零闭元，所以其上同调只有零次的 \(\mathbb R\)。这与自由作用的商是一点相符。
\end{example}

\begin{example}\label{g2:ex:cartan-rotating-sphere}
	继续取式\eqref{g2:eq:sphere-fixed-restriction}中的圆群作用
	\[
		z\cdot[z_0:z_1]=[z_0:zz_1].
	\]
	在两极之外设
	\[
		s=\frac{|z_1|^2}{|z_0|^2+|z_1|^2},\qquad
		\theta=\frac{\arg(z_1/z_0)}{2\pi},\qquad
		\omega=\mathrm ds\wedge\mathrm d\theta.
	\]
	\(s\) 从北极 \(N=[1:0]\) 的 \(0\) 增至南极 \(S=[0:1]\) 的 \(1\)。虽然角坐标在两极失效，\(\omega\) 却光滑延拓到整个球面：在 \(w=z_1/z_0=x+iy\) 坐标中，
	\[
		\omega=\frac{1}{\pi(1+|w|^2)^2}\,\mathrm dx\wedge\mathrm dy,
	\]
	改用 \(1/w\) 后仍是同型表达式。它按复定向满足 \(\int_{\mathbb CP^1}\omega=1\)。基本向量场为 \(V=\partial_\theta\)，因而
	\[
		\iota_V\omega=-\mathrm ds=\mathrm d\mu,\qquad \mu=-s.
	\]
	于是总次数为 \(2\) 的多项式值形式
	\begin{equation}\label{g2:eq:cartan-sphere-form}
		\widetilde\omega=\omega+u\mu=\omega-su
	\end{equation}
	满足 \(\mathrm d_T\widetilde\omega=u(\mathrm d\mu-\iota_V\omega)=0\)。其固定点限制是 \(0\) 与 \(-u\)。这里零次形式 \(\mu\) 恰好抵消了 \(\omega\) 沿轨道方向的插入；这就是等变闭形式比普通闭形式多出一个多项式项的具体原因。
\end{example}

\subsection{紧 Lie 群作用下的比较条件}
\label{g2:subsec:C01:05}

Cartan 复形可以按定义对任意 Lie 群写下。它何时计算 Borel 上同调，则需要比较定理。以下等变 de Rham 定理的紧连通版本取自 Meinrenken\cite[\S6, Theorem 6.1]{Meinrenken2006EquivariantCartan}；Atiyah--Bott\cite[\S4, Theorem 4.13]{AtiyahBott1984MomentMap}给出了等价的 Weil 模型版本。两者的基本场符号须与式\eqref{g2:eq:cartan-fundamental-field}配套，下面用圆群连接把本书的归一化具体写出。

\begin{theorem}[等变 de Rham 定理]\label{g2:thm:equivariant-de-rham}
	设紧 Lie 群 \(G\) 光滑作用于有限维 Hausdorff、第二可数光滑无边流形 \(M\)。则有自然的分次环同构
	\[
		H^*(\Omega_G(M),\mathrm d_G)\cong H_G^*(M;\mathbb R).
	\]
	该同构与等变光滑映射的拉回相容，也与来自一点的系数环作用相容。这里不要求 \(M\) 紧、作用自由或 \(M\) 定向。
\end{theorem}

引用文献的紧连通版本可以推出这里的全部紧群版本。事实上，令 \(G^\circ\) 为单位连通分支，\(\Gamma=G/G^\circ\) 为有限群。Lie 代数相同，并且
\[
	\Omega_G(M)=\Omega_{G^\circ}(M)^\Gamma.
\]
对有限群取不变量在实向量空间上是正合的：用 \(|\Gamma|^{-1}\sum_{\gamma\in\Gamma}\gamma\) 平均即可把不变的上同调类选成不变代表，并把不变边界选成不变原像。因此左端的上同调是 \(H^*(\Omega_{G^\circ}(M))^\Gamma\)。拓扑一侧有正则有限覆盖
\[
	EG\times_{G^\circ}M\longrightarrow EG\times_G M
\]
及覆盖变换群 \(\Gamma\)。有限覆盖的转移满足“拉回后求和”与“先转移再拉回”的关系，除以 \(|\Gamma|\) 后给出
\[
	H_G^*(M;\mathbb R)\cong H_{G^\circ}^*(M;\mathbb R)^\Gamma.
\]
连通版本的比较同构具有自然性，因而与这两个 \(\Gamma\) 作用相容，遂得所述结论。对非连通群只保留 Lie 代数的不变性会漏掉这个有限群作用；定义中取整个 \(G\) 的不变量正为此所需。

\ProseParagraphBreak
比较定理并不是直接把 \(EG\times_G M\) 当成有限维流形来应用\cref{g2:thm:de-rham}。其有限维实现可以沿第19.4节的构造来理解。把紧群嵌入某个 \(\mathrm U(r)\)，取复 Stiefel 流形 \(E_N=V_r(\mathbb C^N)\)。其右 \(G\) 作用自由，且随 \(N\) 增大连通性趋于无穷；忘去标架的最后一列所得球面纤维序列给出了这个连通性增长。于是
\[
	M_N=E_N\times_G M
\]
是有限维光滑流形，对每个固定次数，充分大的 \(N\) 已保留该次数的 Borel 上同调。这里始终是先固定所要计算的次数，再取足够大的有限阶段。

\ProseParagraphBreak
在主右丛 \(E_N\to E_N/G\) 上选连接，将多项式变量代入曲率，再将从 \(M\) 拉回的形式作水平投影，得到落在 \(M_N\) 上的形式。紧性可以提供不变度量与连接；这一构造是第22章陈--Weil 方法在等变复形中的延伸。比较定理还要证明该映射在稳定次数内诱导同构，所用的 Weil 代数是连接变量与曲率变量组成的无上同调障碍的代数模型。仅仅对 \(M\) 上的形式作群平均，不能产生一点上的多项式环；连接模型及其与 Borel 构造的比较正是这里引用的理论内容。Meinrenken\cite[\S\S5--7]{Meinrenken2006EquivariantCartan}分别给出 Weil 模型、Cartan 模型及连接所诱导的映射。

\ProseParagraphBreak
圆群情形足以明确一个容易混淆的负号。设 \(E\to B\) 为主右 \(T\) 丛，\(a\) 为实值连接形式，在周期为 \(1\) 的右作用生成场 \(V_E\) 上满足 \(a(V_E)=1\)。在 \(E\times M\) 上，构成 Borel 型商的右作用是
\[
	(e,x)\cdot z=(ez,z^{-1}x),
\]
其基本场为 \(V_E-V_M\)。若 \(\alpha\) 是 \(M\) 上的 \(T\) 不变形式，省略拉回符号，则其水平投影为
\[
	\operatorname{hor}(\alpha)=\alpha+a\wedge\iota_{V_M}\alpha.
\]
插入 \(V_E-V_M\) 后两项抵消，故该形式基本，能够下降到 \(E\times_T M\)。再由 Cartan 公式及 \(\mathcal L_{V_M}\alpha=0\)，直接算得
\[
	\mathrm d\operatorname{hor}(\alpha)=\operatorname{hor}(\mathrm d\alpha)+\mathrm da\wedge\iota_{V_M}\alpha.
\]
因此连接给出的复形映射是
\begin{equation}\label{g2:eq:circle-cartan-map}
	\mathcal C_a\left(\sum_j u^j\alpha_j\right)
	=\sum_j(-\mathrm da)^j\wedge\operatorname{hor}(\alpha_j).
\end{equation}
上式满足 \(\mathrm d\mathcal C_a=\mathcal C_a\mathrm d_T\)：\(\mathrm da\) 闭，且 \(-u\iota_{V_M}\) 在代入后恰好产生 \(+\mathrm da\wedge\iota_{V_M}\)。这一核对同时决定了多项式变量的符号。

\ProseParagraphBreak
具体说，权 \(1\) 关联线 \(E\times_T\mathbb C\) 的局部酉连接形式为 \(2\pi i a\)，曲率为 \(2\pi i\,\mathrm da\)。依照\cref{g2:prop:chern-weil-line-normalization}的约定，
\begin{equation}\label{g2:eq:equivariant-u-curvature}
	c_1(E\times_T\mathbb C)_{\mathbb R}
	=\left[\frac{i}{2\pi}(2\pi i\,\mathrm da)\right]
	=[-\mathrm da].
\end{equation}
所以式\eqref{g2:eq:circle-cartan-map}确实把 \(u\) 送到本节前半定义的权 \(1\) 关联线的一阶陈类。对于 \(S^{2N+1}\) 上的右标量作用，该关联线是重言线 \(\gamma\)，故 \(u=c_1(\gamma)=-h\)。\(u\) 既有明确的拓扑含义，也有确定的曲率代表，不能在切换基本场方向时只改微分中的一个符号。

\ProseParagraphBreak
现在把比较定理用于\cref{g2:ex:cartan-rotating-sphere}。等变闭形式 \(\omega-su\) 在两极的限制为 \(0,-u\)，而\cref{g2:thm:equivariant-sphere-ring}中的 \(x\) 具有相同限制。式\eqref{g2:eq:sphere-fixed-restriction}说明实系数下该限制映射单射，故
\[
	[\omega-su]=x\in H_T^2(\mathbb CP^1;\mathbb R).
\]
由此得到 \(x\) 的微分形式代表；它忘掉群变量后的普通 de Rham 类是正生成元 \([\omega]\)，固定点处则留下表示的权。

\begin{example}\label{g2:ex:noncompact-group-cartan-failure}
	取加法群 \(G=\mathbb R\) 作用于一点。其 Lie 代数为一维，伴随作用平凡，Cartan 复形仍是零微分的 \(\mathbb R[u]\)。但 \(\mathbb R\) 作为拓扑群可缩，其通用丛可以取 \(\mathbb R\to\mathrm{pt}\)，故 \(B\mathbb R\) 可取为一点，Borel 上同调只有零次。可见比较定理中的群紧性不能直接删掉。另一个界限来自系数：微分形式产生实上同调，\cref{g2:ex:finite-isotropy-orbit}中 \(BC_m\) 的整数扭类不会出现在 Cartan 模型里。
\end{example}

\subsection{环面局部化所需的系数与紧性条件}
\label{g2:subsec:C01:06}

球面例子的限制映射把一个全局等变类送到两极上的多项式。其像仍满足 \(P-Q\) 可被 \(u\) 整除这一条件；允许 \(u\) 作分母以后，两极的数据才可以独立指定。\term[dengbian jubuhua]{等变局部化}[equivariant localization]把这个现象推广到环面作用。

\ProseParagraphBreak
设 \(T=(S^1)^r\)，并固定系数域 \(\mathbf k=\mathbb Q\) 或 \(\mathbb R\)。对每个圆因子都采用前文权 \(1\) 关联线的约定，于是
\[
	R_T=H_T^*(\mathrm{pt};\mathbf k)=\mathbf k[u_1,\ldots,u_r],\qquad |u_j|=2.
\]
一个特征标 \(\chi\colon T\to S^1\) 写成 \(\chi(z_1,\ldots,z_r)=z_1^{m_1}\cdots z_r^{m_r}\)，其关联线的一阶陈类为
\[
	\lambda_\chi=m_1u_1+\cdots+m_ru_r.
\]
取由所有非零特征标的 \(\lambda_\chi\) 生成的乘法集 \(S\subset R_T\)。因为采用特征零域，非零特征标总给出非零线性多项式。

\ProseParagraphBreak
对一个 \(R_T\) 模 \(A\)，\(S^{-1}A\) 中的元素写作 \(a/s\)。尤其有
\[
	a/1=0\quad\Longleftrightarrow\quad\text{存在 }s\in S\text{ 使 }sa=0.
\]
所以局部化抹去的是被某个非零权的乘积消灭的部分。它也保持模的正合列：分母可以相乘来统一，分子落入核的条件就能在乘上另一个分母后化为原正合列中的条件。以下证明正要使用这两个性质。

\begin{theorem}[环面局部化定理]\label{g2:thm:torus-localization}
	设 \(X\) 是有限 \(T\)-CW 复形，或是带光滑 \(T\) 作用的紧光滑无边流形。令 \(F=X^T\) 为固定点集，\(i\colon F\hookrightarrow X\) 为包含。则限制诱导同构
	\[
		S^{-1}i^*\colon S^{-1}H_T^*(X;\mathbf k)
		\xrightarrow{\;\cong\;}S^{-1}H_T^*(F;\mathbf k),
		\qquad \mathbf k=\mathbb Q\text{ 或 }\mathbb R.
	\]
	这一同构不要求 \(X\) 定向。
\end{theorem}

先证明有限 \(T\)-CW 情形。所谓 \(T\)-胞腔是 \(T/H\times D^n\)，群只作用在第一因子；有限 \(T\)-CW 复形由有限个这样的胞腔沿 \(T/H\times S^{n-1}\) 等变粘合而成。\(H=T\) 的胞腔组成固定点子复形 \(F\)。其余胞腔的稳定子 \(H\) 是 \(T\) 的真闭子群。

\begin{proof}[有限 \(T\)-CW 情形的证明]
	对真闭子群 \(H<T\)，商 \(T/H\) 是正维环面。取该商的一个非零特征标并拉回到 \(T\)，得到在 \(H\) 上恒为 \(1\) 的非零特征标 \(\chi\)。于是 \(\lambda_\chi\) 在 \(H^*(BH;\mathbf k)\) 上的限制为零，因为相应的一维 \(H\) 表示平凡。

	\cref{g2:prop:borel-homogeneous-orbit}给出 \((T/H)_T\simeq BH\)，故单个相对胞腔满足
	\[
		H_T^q(T/H\times D^n,T/H\times S^{n-1};\mathbf k)
		\cong H^{q-n}(BH;\mathbf k).
	\]
	该同构来自圆盘对的相对上同调，并保持 \(R_T\) 模结构。右端被 \(\lambda_\chi\in S\) 消灭，因此局部化后为零。

	以 \(F\) 为起点，逐层加入其余的有限个等变胞腔。相对胞腔的计算、切除及三元组长正合列说明：若前一步的相对群局部化为零，加入下一层以后仍为零。有限归纳得到 \(S^{-1}H_T^*(X,F;\mathbf k)=0\)。最后将对 \((X,F)\) 的长正合列局部化，就得到所述同构。
\end{proof}

紧光滑流形版本采用 Atiyah--Bott\cite[\S3, Theorem 3.5]{AtiyahBott1984MomentMap}的局部化定理。其证明利用紧群作用的局部切片与紧性：非固定轨道有真稳定子，局部模型的等变上同调由相应权消灭；有限个局部模型使这些权能组成一个有限乘积。该文以复系数表述支撑集；同一特征标消灭及有限粘合的论证适用于 \(\mathbb Q\) 和 \(\mathbb R\)。上面的胞腔证明把这个有限性机制直接显示出来，也给出了无需光滑结构的版本。

\ProseParagraphBreak
在 \(F\) 上作用平凡，故 \(F_T=BT\times F\)，并且
\[
	H_T^*(F;\mathbf k)=R_T\otimes_{\mathbf k}H^*(F;\mathbf k).
\]
局部化于是把全空间的计算归结到固定分支及其普通上同调。不过，局部化以前可能存在被权消灭的类，也可能像球面那样存在固定点数据之间的整除条件；它们都是原来的等变上同调所保留的信息。\cref{g2:ex:weighted-sphere-coefficients}进一步说明：在正特征系数下，有限稳定子的阶可能变成零，非零权也可能变成零，不能沿用这里的乘法集与结论。

\ProseParagraphBreak
紧性在定理中保证上述有限性。如果已有有限 \(T\)-CW 结构，就可直接用第一种假设。对一般非紧流形，局部的消灭元未必能合并为一个全局消灭元；\cref{g2:ex:localization-noncompact-orbits}给出由无穷多个不同稳定子造成的反例。这与比较定理的条件不同：\cref{g2:thm:equivariant-de-rham}不要求 \(M\) 紧，但这里的局部化结论需要适当的全局有限性。

\ProseParagraphBreak
若还要把局部化用于积分，则应加入定向。设 \(M\) 紧、无边且定向，\(T\) 光滑作用于 \(M\)。对每个固定分支 \(F_a\)，给其法丛 \(\nu_a\) 选取 \(T\) 不变定向，并以 \(\mathrm TM|_{F_a}=\mathrm TF_a\oplus\nu_a\) 规定 \(F_a\) 的定向。法方向没有零权；等变 Euler 类 \(e_T(\nu_a)\) 的最高多项式部分是带定向符号的非零权之积。其余部分含正次普通上同调类，因有限维而幂零，所以 \(e_T(\nu_a)\) 在局部化后可逆。此时等变积分满足
\begin{equation}\label{g2:eq:equivariant-localized-integral}
	\int_M\alpha=\sum_a\int_{F_a}\frac{i_a^*\alpha}{e_T(\nu_a)}.
\end{equation}
这是 Atiyah--Bott\cite[\S3, formula (3.8)]{AtiyahBott1984MomentMap}的局部化积分公式。它还使用等变 Thom 类和推前；第四卷第12.6节将展开所选情形的证明与应用。

\ProseParagraphBreak
对本节的 \(\mathbb CP^1\)，北极切表示的权为 \(1\)，南极为 \(-1\)，复定向给出 \(e_T(\mathrm T_NM)=u\)、\(e_T(\mathrm T_SM)=-u\)。代入式\eqref{g2:eq:cartan-sphere-form}，右端成为
\[
	\frac{0}{u}+\frac{-u}{-u}=1=\int_{\mathbb CP^1}\omega.
\]
这个计算同时用到了多项式变量的归一化、两极的权与普通面积形式的定向；三者在同一约定下相符，公式才给出正确的积分。

\section{谱与稳定同伦}
\label{g2:sec:C02}

第9章把 \(\pi_{n+q}(S^n)\) 沿悬挂排列，发现固定差值 \(q\) 后，足够靠后的群相同。第24章则把一张向量丛所给的 \(K\) 类悬挂两次，并用 Bott 元比较悬挂前后的群。两种构造都要求同时保存一列空间及其比较映射。谱把这样的序列组织成一个稳定对象，使悬挂成为可逆操作，也让不同同调理论能够使用相同的正合列语言。

\subsection{悬挂稳定化和谱}
\label{g2:subsec:C02:01}

以下空间带有基点，CW 模型的基点取为零胞腔。这样基点包含具有同伦延拓性质，约化锥、约化悬挂和压缩积可以直接使用前文的商空间构造。回忆\cref{g2:def:reduced-suspension,g2:prop:suspension-loop-adjunction}：
\[
	\Sigma X=S^1\wedge X,\qquad
	[\Sigma X,Y]_*\cong[X,\Omega Y]_*.
\]
若 \(X\) 为 \((r-1)\) 连通、\(r\ge2\)，\cref{g2:thm:freudenthal}说明 \(\pi_i(X)\to\pi_{i+1}(\Sigma X)\) 在 \(i\le2r-2\) 时为同构，在 \(i=2r-1\) 时为满射。悬挂将空间的连通度也提高一，因此固定差值后最终进入同构范围。

\begin{definition}\label{g2:def:sequential-prespectrum}
	一个顺序\term[yupu]{预谱}[prespectrum]由带基点空间 \(E_n\)、\(n\ge0\)，以及带基点的结构映射
	\[
		\sigma_n:\Sigma E_n\longrightarrow E_{n+1}
	\]
	组成。计算时各 \(E_n\) 取良基点的 CW 模型。预谱映射 \(f:E\to E'\) 是一列带基点映射 \(f_n:E_n\to E'_n\)，满足
	\[
		\sigma'_n\circ\Sigma f_n=f_{n+1}\circ\sigma_n.
	\]
	若每个结构映射的伴随 \(\widetilde\sigma_n:E_n\to\Omega E_{n+1}\) 都是弱同伦等价，则称它为 \(\Omega\)-预谱。
\end{definition}

结构映射不是装饰：它规定前一层的类怎样进入后一层。即使每个 \(E_n\) 都是同一个指定空间，选择不同的 \(\sigma_n\) 仍可能得到不同的稳定群，\cref{g2:ex:spectrum-degree-two-system}会给出这样的计算。要求 \(\Omega\)-预谱的伴随映射为弱同伦等价，则意味着每层已经能通过下一层的回路空间恢复自身。

\begin{example}\label{g2:exm:suspension-prespectrum}
	带基点 CW 复形 \(X\) 给出悬挂预谱
	\[
		E_n=\Sigma^nX,\qquad
		\sigma_n:\Sigma(\Sigma^nX)\xrightarrow{\cong}\Sigma^{n+1}X.
	\]
	取 \(X=S^0\) 时得到球预谱 \(E_n=S^n\)。这些结构映射是空间同胚，但其伴随 \(S^n\to\Omega S^{n+1}\) 只在 Freudenthal 范围内控制同伦群；“结构映射是同胚”与“伴随映射是弱同伦等价”是不同要求。
\end{example}

由 \(\sigma_n\) 和悬挂同态形成
\begin{equation}\label{g2:eq:prespectrum-transition}
	\pi_{q+n}(E_n)\xrightarrow{\ \Sigma\ }
	\pi_{q+n+1}(\Sigma E_n)
	\xrightarrow{(\sigma_n)_*}\pi_{q+n+1}(E_{n+1}).
\end{equation}
固定整数 \(q\)，只取 \(n\) 足够大使 \(q+n\ge2\)，便得到一列阿贝尔群。

\begin{definition}\label{g2:def:prespectrum-homotopy-groups}
	预谱 \(E\) 的第 \(q\) 个稳定同伦群定义为
	\[
		\pi_q(E)=\varinjlim_n\pi_{q+n}(E_n),\qquad q\in\mathbb Z,
	\]
	过渡同态取式\eqref{g2:eq:prespectrum-transition}。一张预谱映射若在所有 \(\pi_q\) 上诱导同构，就称为稳定等价。\symidx{\(\pi_q(E)\)：预谱或谱的稳定同伦群}
\end{definition}

一个极限类由某层中的球面映射表示；两个代表相等，是指它们经过有限次过渡后成为同一个同伦类。删去有限多个起始项不改变极限，所以负的 \(q\) 也合法：例如 \(q=-3\) 时，从 \(n\ge5\) 开始，取的仍是普通的非负高阶同伦群。这里没有定义负维球面空间的同伦群。

\ProseParagraphBreak

若 \(E\) 是 \(\Omega\)-预谱，式\eqref{g2:eq:prespectrum-transition}也可写为伴随映射诱导的
\[
	\pi_{q+n}(E_n)\xrightarrow{(\widetilde\sigma_n)_*}
	\pi_{q+n}(\Omega E_{n+1})\cong\pi_{q+n+1}(E_{n+1}),
\]
故每一步都是同构。一般预谱则要通过谱化把这些稳定信息放进适当的 \(\Omega\)-模型。这里采用的\term[pu]{谱}[spectrum]语言，指预谱经这项处理后在稳定同伦范畴中确定的对象；不同具体模型用稳定等价比较。谱化的存在及该范畴所需结构将在\cref{g2:subsec:C02:03}列明。对悬挂预谱所得的谱，记为 \(\Sigma^\infty X\)，称为\term[xuangua pu]{悬挂谱}[suspension spectrum]；\(\mathbb S=\Sigma^\infty S^0\) 称为\term[qiu pu]{球谱}[sphere spectrum]。\symidx{\(\Sigma^\infty X\)：空间的悬挂谱}\symidx{\(\mathbb S\)：球谱}

\begin{proposition}\label{g2:prop:suspension-spectrum-stable-range}
	设 \(X\) 是 \((r-1)\) 连通的带基点 CW 复形，\(r\ge2\)，基点为零胞腔。对整数 \(q\)，若 \(n\ge0\)、\(q+n\ge2\) 且 \(n\ge q-2r+2\)，则自然映射
	\[
		\pi_{q+n}(\Sigma^nX)\longrightarrow\pi_q(\Sigma^\infty X)
	\]
	为同构。
\end{proposition}
\begin{proof}
	\(\Sigma^nX\) 为 \((r+n-1)\) 连通。可由第9章双锥模型逐次得到这一点：悬挂的单连通性来自两个锥的 van Kampen 计算，约化同调悬挂同构把低次消失向上移一位，Hurewicz 定理便排除低于所需次数的首个非零同伦群。于是 Freudenthal 定理用于 \(\Sigma^nX\) 时，式\eqref{g2:eq:prespectrum-transition}在
	\[
		q+n\le2(r+n)-2
	\]
	时为同构。这个不等式正是 \(n\ge q-2r+2\)。将 \(n\) 增加一后它仍成立，所以从此项开始所有过渡同态都是同构，直接极限即为这一项。
\end{proof}

例如 \(X=S^3\)、\(q=7\) 时，可取 \(r=3\)，上式保证从 \(n=3\) 开始，\(\pi_{10}(S^6)\) 已等于 \(\pi_7(\Sigma^\infty S^3)\)。这给出的是何时停止继续悬挂的标准；具体群的数值还要计算。稳定化的便利在于，进入这个范围以后，提高两边维数不会再改变问题。

\subsection{稳定同伦群及系数谱}
\label{g2:subsec:C02:02}

一个谱同时拥有所有整数次数的稳定群。若把谱看作一种同调理论的系数对象，那么 \(\pi_q(E)\) 就是该理论在一点上的第 \(q\) 次同调。不同系数谱使同一空间的胞腔被不同的系数组检测。先从三个熟悉的来源看这些群。

\begin{proposition}\label{g2:prop:sphere-spectrum-low-coefficients}
	球谱满足
	\[
		\pi_q(\mathbb S)=\pi_q^{\mathrm S},\qquad
		\pi_0(\mathbb S)\cong\mathbb Z,\qquad
		\pi_q(\mathbb S)=0\quad(q<0).
	\]
	在零次同构中，各层球面的恒等映射表示整数 \(1\)。
\end{proposition}
\begin{proof}
	球预谱的过渡同态就是悬挂，所以第一式正是\cref{g2:def:stable-sphere-homotopy}的定义。\(q=0\) 时，\(n\ge2\) 的群 \(\pi_n(S^n)\) 由度数识别为 \(\mathbb Z\)，悬挂保持恒等映射及度数，因此极限为 \(\mathbb Z\)。若 \(q<0\)，取 \(n\) 大到 \(q+n\ge2\)，则 \(q+n<n\)。球面 \(S^n\) 的 \((n-1)\) 连通性使每一项都为零，极限也为零。
\end{proof}

这也解释了为何谱的零次群不是空间的道路分支集：\(\pi_0(\mathbb S)\) 来自任意高维球面之间的等维映射。负次群的消失则是球谱的一项性质；一般谱可以有非零负次群。

\ProseParagraphBreak

第二个来源是 Eilenberg--MacLane 空间。给定阿贝尔群 \(A\)，第9章构造了 \(K(A,n)\)，而第12章的\cref{g2:thm:em-representability}证明
\[
	[X,K(A,n)]_*\cong\widetilde H^n(X;A)\qquad(n\ge1).
\]
选定各模型及其唯一非零同伦群与 \(A\) 的识别，可以取弱等价
\[
	K(A,n)\longrightarrow\Omega K(A,n+1)
\]
在该群上实现恒等。这个映射可按\cref{g2:thm:eilenberg-maclane-existence}的胞腔模型构造：先对目标回路空间取 CW 模型，在 \(n\) 维骨架上选取 \(A\) 的生成元的代表，在 \(n+1\) 维利用关系延拓；更高维延拓所遇到的目标同伦群都为零。所得映射在唯一非零同伦群上为恒等，因而是弱同伦等价。取 \(K(A,0)=A\) 为带零基点的离散空间，零层的映射则选取每个回路分支的代表。将这些映射取伴随，得到 \(\Omega\)-预谱，其谱记为 \(HA\)，称为\term[Eilenberg MacLane pu]{Eilenberg--MacLane 谱}[Eilenberg--MacLane spectrum]。\symidx{\(HA\)：阿贝尔群 \(A\) 的 Eilenberg--MacLane 谱}

\begin{proposition}\label{g2:prop:em-spectrum-coefficients}
	对任意阿贝尔群 \(A\)，
	\[
		\pi_q(HA)=\begin{cases}A,&q=0,\\0,&q\ne0.\end{cases}
	\]
\end{proposition}
\begin{proof}
	取 \(n\) 足够大，使 \(n\ge2\)、\(q+n\ge2\)。\(K(A,n)\) 的同伦群只在次数 \(n\) 非零，所以 \(\pi_{q+n}(K(A,n))\) 只有在 \(q=0\) 时为 \(A\)。此时已选的伴随映射在 \(A\) 上为恒等，故过渡同态也是恒等。取直接极限即得结论。
\end{proof}

由 \(HA\) 表示的同调与上同调正是普通的 \(A\) 系数理论，构造与比较见\cite[第22章第1—2节]{May1999}。系数谱的作用是把“在每个次数用哪个表示空间”与这些表示空间之间的悬挂关系同时保存起来，而不是另外添加一组独立的群。

\ProseParagraphBreak

第三个来源是复 \(K\) 理论。第19章的向量丛分类及取虚秩给出表示空间 \(\mathbb Z\times BU\)；奇次使用稳定酉群 \(U=\bigcup_N\mathrm U(N)\)。Bott 周期性的表示空间形式把
\[
	\mathbb Z\times BU,\ U,\ \mathbb Z\times BU,\ U,\ldots
\]
组织成 \(\Omega\)-预谱，其谱记为 \(KU\)。把所有有限 CW 空间上的自然周期比较实现为这些表示空间之间的弱等价，属于表示空间理论的输入，见\cite[第24章第2节及第25章第7节]{May1999}。在有限 CW 空间上，所得理论与第24章从向量丛定义的复 \(K\) 理论相同。谱系数与上同调系数的次数相反：
\begin{equation}\label{g2:eq:ku-spectrum-coefficients}
	\pi_q(KU)=K^{-q}(\mathrm{pt})
	\cong\begin{cases}\mathbb Z,&q\text{ 为偶数},\\0,&q\text{ 为奇数}.
	\end{cases}
\end{equation}
对 \(q\ge0\)，这些值由\cref{g2:prop:k-sphere-coefficients}给出；将二周期性向两个方向延拓后，公式对所有整数 \(q\) 成立。特别地，\(\pi_{-2}(KU)=\mathbb Z\)，所以 \(KU\) 不是负次群全为零的谱。

\ProseParagraphBreak

实 \(K\) 理论同样由谱 \(KO\) 表示，其系数为 \(\pi_q(KO)=KO^{-q}(\mathrm{pt})\)。沿用第24章已外引的实 Bott 周期性，\(q\bmod8=0,1,\ldots,7\) 时的值依次为
\[
	\mathbb Z,\ \mathbb Z/2,\ \mathbb Z/2,\ 0,\ \mathbb Z,\ 0,\ 0,\ 0.
\]
该表示空间与八周期输入见\cite[第II章第3节；第IV章第3节]{Weibel2013KBook}。实、复的系数已经在一维与二维处分开，因而不能仅把 \(KU\) 的周期数改成八就得到实理论。

\begin{center}
	\begin{tabular}{@{}l l l@{}}
		\toprule
		系数谱 & \(\pi_q\) 的特征 & 对应几何信息 \\
		\midrule
		\(\mathbb S\) & 零次为 \(\mathbb Z\)，负次为零 & 球面映射的稳定类 \\
		\(HA\) & 只在零次为 \(A\) & 普通 \(A\) 系数同调与上同调 \\
		\(KU\) & 偶次 \(\mathbb Z\)，奇次零，二周期 & 复向量丛的稳定差 \\
		\(KO\) & 八周期，并含 \(\mathbb Z/2\) & 实向量丛的稳定差 \\
		\bottomrule
	\end{tabular}
\end{center}

这些例子中有些空间或系数群重复出现，但比较映射与周期结构各有来源。正是这些结构决定了它们在空间附着胞腔时如何响应。

\subsection{广义同调理论的进一步输入}
\label{g2:subsec:C02:03}

在空间层面，映射锥主要控制同调，而同伦群还有非线性现象。把悬挂变成可逆操作以后，映射锥与同伦纤维只相差一次悬挂，因而同一构造也能产生稳定同伦群的长正合列。要准确使用这句话，需要的不只是定义一个直接极限，还需要一个有正确映射、压缩积和余纤维序列的范畴。

\ProseParagraphBreak

本节采用稳定同伦范畴 \(\mathrm{SH}\) 的下列结构。它由顺序预谱通过 CW 替换、谱化并使稳定等价可逆得到；\([E,F]_{\mathrm{SH}}\) 表示其中的映射群。悬挂 \(\Sigma\) 是自等价，所以可以取任意整数次 \(\Sigma^q\)，并有
\[
	\pi_q(E)=[\Sigma^q\mathbb S,E]_{\mathrm{SH}}.
\]
其中 \(q<0\) 的 \(\Sigma^q\mathbb S\) 是球谱的退悬挂，属于谱范畴。每张谱映射都可完成为余纤维序列
\[
	E\longrightarrow F\longrightarrow C\longrightarrow\Sigma E.
\]
对它取映射群会产生长正合列。谱的压缩积采用先作适当 CW 替换后的导出压缩积，记为 \(\wedge^{\mathbb L}\)；它保持余纤维序列与楔和，\(\mathbb S\) 是其单位。球谱在此范畴中是紧对象，即从它的任意整数悬挂出发的映射与任意楔和可交换为直和。

\ProseParagraphBreak

上述结构是稳定同伦理论的基础输入；\cite[第25章第6—7节]{May1999}说明其对象、谱化与所需性质，并给出进一步文献。这里使用的是这些已列出的性质。空间逐层的普通压缩积不能直接替代导出压缩积，任意预谱的逐层同伦类也不能直接替代 \([E,F]_{\mathrm{SH}}\)：替换和稳定等价的可逆化负责使它们真正只依赖稳定对象。

\begin{definition}\label{g2:def:spectrum-generalized-theories}
	固定谱 \(E\)。对带基点 CW 复形 \(X\)，定义
	\begin{align}
		\widetilde E_q(X)&=\pi_q(E\wedge^{\mathbb L}\Sigma^\infty X),\label{g2:eq:spectrum-homology}\\
		\widetilde E^{\,q}(X)&=[\Sigma^\infty X,\Sigma^q E]_{\mathrm{SH}}.\label{g2:eq:spectrum-cohomology}
	\end{align}
	对无指定基点的空间，令 \(E_q(X)=\widetilde E_q(X_+)\)、\(E^q(X)=\widetilde E^{\,q}(X_+)\)。对 CW 对 \((X,A)\)，相对理论取商 \(X/A\)，并约定 \(X/\varnothing=X_+\)。它们称为 \(E\) 表示的\term[guangyitongdiao]{广义同调}[generalized homology]与广义上同调理论。
\end{definition}

上同调在目标上移位，同调在取映射的球谱上移位，因此一点上的次数恰为
\[
	E_q(\mathrm{pt})=\pi_q(E),\qquad
	E^q(\mathrm{pt})=\pi_{-q}(E).
\]
由这个等式可以随时恢复式\eqref{g2:eq:ku-spectrum-coefficients}的负号。对球面还得到
\begin{equation}\label{g2:eq:generalized-sphere-groups}
	\widetilde E_q(S^r)=\pi_{q-r}(E),\qquad
	\widetilde E^{\,q}(S^r)=\pi_{r-q}(E),\qquad r\ge0.
\end{equation}
例如 \(\widetilde{KU}_0(S^2)=\pi_{-2}(KU)=\mathbb Z\)，而普通约化零次同调 \(\widetilde H_0(S^2;\mathbb Z)=0\)。广义理论保留同伦不变、切除与正合性，但不要求一点上只有零次系数。

\begin{proposition}\label{g2:prop:spectrum-exact-wedge}
	采用上述稳定范畴结构，对任意谱 \(E\)，式\eqref{g2:eq:spectrum-homology}与式\eqref{g2:eq:spectrum-cohomology}分别给出带基点 CW 复形上的约化同调与上同调理论。对于 CW 子复形包含 \(A\hookrightarrow X\)，有长正合列
	\[
		\cdots\longrightarrow\widetilde E_q(A)\longrightarrow\widetilde E_q(X)
		\longrightarrow\widetilde E_q(X/A)\longrightarrow\widetilde E_{q-1}(A)\longrightarrow\cdots
	\]
	及
	\[
		\cdots\longrightarrow\widetilde E^{\,q}(X/A)\longrightarrow\widetilde E^{\,q}(X)
		\longrightarrow\widetilde E^{\,q}(A)\longrightarrow\widetilde E^{\,q+1}(X/A)\longrightarrow\cdots.
	\]
	对任意一族带基点 CW 复形 \(X_\alpha\)，
	\[
		\widetilde E_q\left(\bigvee_\alpha X_\alpha\right)
		\cong\bigoplus_\alpha\widetilde E_q(X_\alpha),\qquad
		\widetilde E^{\,q}\left(\bigvee_\alpha X_\alpha\right)
		\cong\prod_\alpha\widetilde E^{\,q}(X_\alpha).
	\]
\end{proposition}
\begin{proof}
	子复形包含是余纤维化，故其映射锥与 \(X/A\) 同伦等价。悬挂谱把相应的空间余纤维序列送为稳定余纤维序列。对这条序列先与 \(E\) 作导出压缩积，再取 \(\pi_q\)，由映射群的正合性得到第一条长列；直接以 \(\Sigma^qE\) 为目标取逆变映射群，得到第二条。最后一个空间箭头 \(X/A\to\Sigma A\) 决定两列的连接同态，所以这些列也对空间对映射自然。

	在稳定范畴中，\(\Sigma^\infty(\Sigma X)\cong\Sigma\Sigma^\infty X\)。将其代入定义并把一次悬挂移到映射源端或目标端，就得到
	\[
		\widetilde E_q(\Sigma X)\cong\widetilde E_{q-1}(X),\qquad
		\widetilde E^{\,q}(\Sigma X)\cong\widetilde E^{\,q-1}(X).
	\]
	悬挂谱及导出压缩积保持楔和，球谱的紧性又使 \(\pi_q\) 将楔和变为直和，于是得到同调的楔和公式。上同调的公式来自楔和的余积性质：一张从楔和到固定目标的稳定映射，等价于逐个给出各分支上的映射，因此得到乘积。最后，CW 空间上的同伦等价变为稳定范畴中的同构，给出同伦不变性；CW 对的切除来自相应商空间的同伦等价。
\end{proof}

这个证明指出了公理来自哪些结构，而这些结构本身需要谱化和稳定范畴的构造。对普通同调，我们还能用第2—4章的链复形直接完成相同的正合计算；对一般谱，适当的映射锥与导出压缩积承担了相应职责。

\ProseParagraphBreak

反方向也有表示定理：定义在全部带基点 CW 复形及其同伦类上的约化上同调理论，若满足正合性、悬挂同构及任意楔和变乘积的公理，就可由一个 \(\Omega\)-谱表示。这是 Brown\footnote{Edgar Henry Brown, Jr.（1926—2021），美国数学家，主要研究同伦论与广义上同调。}可表性定理的上同调理论版本，参见\cite[第22章第2节]{May1999}。其中“全部 CW 复形”和“任意楔和”是输入的一部分；仅在有限复形上给出一些群和正合列，还需要证明相容延拓，才能直接调用这个版本。

\ProseParagraphBreak

谱的系数表提供胞腔计算的起点，却未给出全部计算。若 \(X^r/X^{r-1}\) 是若干 \(r\) 球面的楔和，式\eqref{g2:eq:generalized-sphere-groups}确定相对群；长正合列中的连接同态仍由胞腔附着映射决定。第11章的谱序列也可以组织这样的逐骨架计算。这里的“谱”是保存悬挂结构的几何对象，“谱序列”则是逐页计算群的代数装置，两者在这种计算中相遇，但各有自己的定义。

\section{持续同调的代数结构}
\label{g2:sec:C03}

同一个空间随阈值扩大，会先出现新的连通分支，随后把分支连接起来；一个圈也可能先形成，随后被加入的面填满。逐个计算同调群只能告诉我们每一步有多少个圈。要追踪哪些圈从前一步保留下来，还需要包含映射诱导的同态。持续同调把这些同态一起保存，域系数则使整套数据能够分解为一组有出生时刻和死亡时刻的区间。

\subsection{有限滤过复形及诱导映射}
\label{g2:subsec:C03:01}

固定一个有限单纯复形 \(K\) 和一个域 \(F\)。例如，\(F=\mathbb F_2\) 时边界系数只有 \(0,1\)，方向符号也随之消失；使用 \(\mathbb Q\) 时仍保留通常的带符号边界。以下使用普通同调，因此一个非空连通复形的零维同调是 \(F\)。

\begin{definition}\label{g2:def:finite-filtration-persistence}
	给 \(K\) 的每个单形指定实数 \(f(\sigma)\)，并要求
	\[
		\tau\subseteq\sigma\quad\Longrightarrow\quad f(\tau)\leq f(\sigma).
	\]
	称 \(f\) 为\term[lvguohanshu]{滤过函数}[filtration function]，令
	\[
		K_t=\{\sigma\in K:f(\sigma)\leq t\},\qquad t\in\mathbb R.
	\]
	于是每个 \(K_t\) 都是子复形，且 \(s\leq t\) 时 \(K_s\subseteq K_t\)。固定 \(q\geq0\)，记
	\[
		V_t=H_q(K_t;F),\qquad v_s^t:H_q(K_s;F)\longrightarrow H_q(K_t;F)
	\]
	为包含映射诱导的同态。这些向量空间与同态组成一个\term[chixumo]{持续模}[persistence module]：它们满足 \(v_t^t=\mathrm{id}\) 及 \(v_t^u v_s^t=v_s^u\)。
\end{definition}

上述相容性来自\cref{g2:def:chain-map}的链映射：包含把循环送到循环，把边界送到边界，所以能在同调商空间上定义映射。链的包含是单射，同调上的映射却可能有核：一个原来填不掉的循环，在较大的复形中可能成为边界。

\begin{definition}\label{g2:def:persistent-rank}
	对于 \(s\leq t\)，称
	\[
		H_q^{s,t}(K;F)=\operatorname{im}v_s^t,\qquad
		\beta_q(s,t)=\dim_F\operatorname{im}v_s^t
	\]
	分别为\term[chixutongdiao]{持续同调}[persistent homology]及其\term[zhifubianliang]{秩不变量}[rank invariant]。
	\symidx{\(H_q^{s,t}(K;F)\)：从参数 \(s\) 保留到 \(t\) 的持续同调}
	\symidx{\(\beta_q(s,t)\)：持续同调的维数}
\end{definition}

\(\beta_q(s,t)\) 数的是在 \(s\) 时已经存在、到 \(t\) 时仍能辨认的独立同调类。刚刚在 \(t\) 出生的类一般不属于这个像。只有取 \(s=t\)，它才退化为普通 Betti 数。

\ProseParagraphBreak

设滤过函数的不同取值为 \(a_0<\cdots<a_N\)。在每段 \([a_i,a_{i+1})\) 上，复形与诱导映射保持不变；在 \(a_0\) 之前为空，在 \(a_N\) 之后恒为 \(K\)。因此，全部实参数数据由有限向量空间图
\[
	V_{a_0}\longrightarrow V_{a_1}\longrightarrow\cdots\longrightarrow V_{a_N}
\]
决定。称这种只有有限多个变化时刻、各向量空间有限维且在最早时刻之前为零的持续模为本节的\term[youxiangouzaoxingchixumo]{有限构造型持续模}[finite constructible persistence module]。变化时刻本身采用变化后的值，这一约定将对应左闭右开的区间。

\begin{example}\label{g2:exm:persistence-triangle}
	取一个填满的三角形，顶点为 \(v_0,v_1,v_2\)，边记为 \(e_{01},e_{12},e_{02}\)，二维单形记为 \(\tau\)。令三个顶点的滤过值都为 \(0\)，再令
	\[
		f(e_{01})=1,\qquad f(e_{12})=2,\qquad f(e_{02})=3,\qquad f(\tau)=4.
	\]
	阈值为 \(0\) 时有三个连通分支，阈值为 \(1,2\) 时依次合并为两个和一个分支；第三条边在 \(3\) 时闭合出一个圈，面在 \(4\) 时把圈填掉。
	\ProseParagraphBreak
	更具体地，沿用\cref{g2:exm:triangle-chain-complex}的方向约定，
	\[
		\mathrm d e_{ij}=v_j-v_i,\qquad z=e_{01}+e_{12}-e_{02},\qquad \mathrm d\tau=z.
	\]
	在 \(3\leq t<4\) 时，\(z\) 是非零循环，复形没有二维链，所以 \([z]\neq0\)。在 \(t\geq4\) 时，同一条链成为边界。这给出一维同调的一段寿命 \([3,4)\)。零维同调的三段寿命则为 \([0,1)\)、\([0,2)\)、\([0,\infty)\)。其中无限长的一段表示最终连通复形留下的零维类。
\end{example}

\begin{figure}[htbp]
	\centering
	% !TeX root = ../Book2.tex
\begin{tikzpicture}[>=stealth,every node/.style={font=\small},x=1cm,y=1cm]
	\foreach \k in {0,1,2,3,4} {
		\begin{scope}[xshift={2.35*\k cm}]
			\ifnum\k=4 \fill[blue!12] (0,0)--(1.2,0)--(.6,1)--cycle;\fi
			\ifnum\k>0 \draw[thick] (0,0)--(1.2,0);\fi
			\ifnum\k>1 \draw[thick] (1.2,0)--(.6,1);\fi
			\ifnum\k>2 \draw[thick,blue!70!black] (.6,1)--(0,0);\fi
			\foreach \p in {(0,0),(1.2,0),(.6,1)} \fill \p circle (1.8pt);
			\node at (.6,-.4) {\(t=\k\)};
		\end{scope}
	}
	\node[left] at (-.2,-1.3) {\(H_0\)};
	\node[left] at (-.2,-3.1) {\(H_1\)};
	\draw[gray,->] (0,-3.65)--(10.9,-3.65) node[right] {\(t\)};
	\foreach \k in {0,1,2,3,4} {
		\draw[gray] ({2.35*\k},-3.65)--({2.35*\k},-3.73) node[below] {\(\k\)};
	}
	\draw[very thick,->] (0,-1.2)--(10.7,-1.2);
	\fill (0,-1.2) circle (2pt);
	\draw[very thick] (0,-1.8)--(2.35,-1.8);
	\fill (0,-1.8) circle (2pt);
	\draw[fill=white,thick] (2.35,-1.8) circle (2pt);
	\draw[very thick] (0,-2.4)--(4.7,-2.4);
	\fill (0,-2.4) circle (2pt);
	\draw[fill=white,thick] (4.7,-2.4) circle (2pt);
	\draw[very thick,blue!70!black] (7.05,-3.1)--(9.4,-3.1);
	\fill[blue!70!black] (7.05,-3.1) circle (2pt);
	\draw[fill=white,thick,blue!70!black] (9.4,-3.1) circle (2pt);
\end{tikzpicture}

	\caption{三角形滤过及其条形码。实心端点表示出生时刻计入区间，空心端点表示死亡时刻已不属于区间。}
	\label{g2:fig:persistence-triangle}
\end{figure}

\subsection{域系数条形码的分解}
\label{g2:subsec:C03:02}

“一个类从出生保持到死亡”可以写成最简单的持续模。给定 \(b<d\leq\infty\)，定义\term[qujianmo]{区间模}[interval module]
\[
	I[b,d)_t=\begin{cases}F,&b\leq t<d,\\0,&\text{其余情形},\end{cases}
\]
两个参数都在区间内时，结构映射为恒等映射，其余结构映射为零。它只携带一个在指定时间段存在的独立方向。下述定理说明，有限构造型持续模正是有限多个这样的方向的直和。

\begin{theorem}[区间分解]\label{g2:thm:finite-persistence-decomposition}
	域 \(F\) 上的每个有限构造型持续模都同构于有限多个区间模 \(I[b,d)\) 的直和，其中 \(b\) 是变化时刻，\(d\) 是较晚的变化时刻或 \(\infty\)。区间及其重数由持续模唯一决定。
\end{theorem}

\begin{proof}
	先证明有限图 \(V_0\to\cdots\to V_N\) 的分解。为便于统一处理最后仍存活的向量，暂时添加 \(V_{N+1}=0\)；这一项只是证明中的记号。用 \(v_i^j\) 表示从第 \(i\) 项到第 \(j\) 项的复合同态。
	\ProseParagraphBreak
	若全部空间为零，结论成立。否则，取最早的非零项 \(V_b\)，选取 \(0\neq x_b\in V_b\)，并令 \(d>b\) 是 \(v_b^d(x_b)=0\) 的最小指标。在 \(b\leq i<d\) 时置 \(x_i=v_b^i(x_b)\)。这些向量非零，并沿箭头依次映到下一项，最后映到零，因而张成一个区间子模。要把它分裂为直和因子，还须选出相容的补空间。
	\ProseParagraphBreak
	选线性泛函 \(\lambda_{d-1}:V_{d-1}\to F\)，使 \(\lambda_{d-1}(x_{d-1})=1\)，并定义
	\[
		\lambda_i=\lambda_{d-1}\circ v_i^{d-1},\qquad b\leq i<d.
	\]
	于是 \(\lambda_i(x_i)=1\)，且 \(\lambda_{i+1}v_i^{i+1}=\lambda_i\)。取 \(W_i=\ker\lambda_i\) 于上述区间内，区间外取 \(W_i=V_i\)。内部箭头保持这些核；入口处 \(V_{b-1}=0\)，出口处目标补空间就是整个 \(V_d\)。故 \(W_i\) 组成子模，且
	\[
		V\cong I[b,d)\oplus W.
	\]
	这里的区间暂以整数指标理解。\(\sum_i\dim W_i<\sum_i\dim V_i\)，反复应用上述步骤便得到有限直和分解。
	\ProseParagraphBreak
	现在证明唯一性。令 \(r(i,j)=\operatorname{rank}v_i^j\)。一个区间 \([b,d)\) 对 \(r(i,j)\) 贡献 \(1\)，当且仅当 \(b\leq i\leq j<d\)。因此，若 \(m_{b,d}\) 为该区间的重数，则
	\[
		r(i,j)=\sum_{b\leq i,\,d>j}m_{b,d}.
	\]
	先取关于 \(i\) 的差，选出恰在 \(b\) 出生的区间；再取关于 \(j\) 的差，选出恰在 \(d\) 死亡的区间。由此得到
	\begin{equation}\label{g2:eq:persistence-multiplicity}
		m_{b,d}=r(b,d-1)-r(b-1,d-1)-r(b,d)+r(b-1,d),
	\end{equation}
	其中 \(r(-1,j)=0\)，添加的零项使 \(r(i,N+1)=0\)。所有重数因而都由原来的映射决定。
	\ProseParagraphBreak
	最后把指标 \(i\) 换回变化时刻 \(a_i\)。死亡指标 \(d\leq N\) 对应实区间 \([a_b,a_d)\)；死亡指标 \(N+1\) 表示在最后一项仍存活，应对应 \([a_b,\infty)\)。在相邻变化时刻之间使用恒等映射延拓，便得到实参数的分解与唯一性。
\end{proof}

\begin{definition}\label{g2:def:persistence-barcode}
	\cref{g2:thm:finite-persistence-decomposition}中带重数的区间集合称为持续模的\term[tiaoxingma]{条形码}[barcode]，记为 \(\mathcal B(V)\)。把每个有限区间 \([b,d)\) 画成平面上的点 \((b,d)\)，并另行保留 \((b,\infty)\)，得到它的\term[chixutu]{持续图}[persistence diagram]。
	\symidx{\(\mathcal B(V)\)：持续模 \(V\) 的条形码}
\end{definition}

分解中选出的向量与补空间可以改变，但区间重数不会改变。特别地，\(\dim V_t\) 等于覆盖时刻 \(t\) 的条数，而 \(\operatorname{rank}v_s^t\) 等于同时覆盖 \(s,t\) 的条数。只给每个时刻的维数会丢失后一项信息，\cref{g2:ex:persistence-same-betti}给出一个只有三个时刻的例子。

\ProseParagraphBreak

也可以把离散持续模写成分次 \(F[u]\)-模 \(\bigoplus_{i\geq0}V_i\)，规定 \(u\cdot x=v_i^{i+1}(x)\) 对 \(x\in V_i\) 成立，并在最后一项以后恒等延拓。一次乘 \(u\) 就是向后移动一步。区间 \([b,d)\) 对应一个在次数 \(b\) 出生、满足 \(u^{d-b}x=0\) 的生成元；无限区间对应没有这种关系的自由生成元。\(F[u]\) 上的分次模结构定理由此给出另一条分解证明，见\cite[第3节]{ZomorodianCarlsson2005Persistence}。上面的直接证明把所需的补空间选择和唯一性公式都留在线性代数之内。

\subsection{边界矩阵约化}
\label{g2:subsec:C03:03}

若在每个阈值重新求一次同调，再比较这些商空间，计算中会反复处理同一批边界。按单形进入的顺序约化边界矩阵，可以同时得到全部出生和死亡。关键是换基必须保持每一步已经出现的链空间：后来的单形可以减去早先的单形，早先的单形却不能掺入尚未出现的单形。

\ProseParagraphBreak

将所有单形排成 \(\sigma_1,\ldots,\sigma_m\)，使滤过值非递减；相同滤过值时，面排在其真上方面之前。例如可先按维数排序，再任意打破剩余的并列。给每个单形定向，把各维边界统一记在一个 \(m\times m\) 矩阵 \(D\) 中：
\[
	\mathrm d\sigma_j=\sum_iD_{ij}\sigma_i.
\]
当 \(D_{ij}\neq0\) 时有 \(i<j\) 及 \(\dim\sigma_i=\dim\sigma_j-1\)。\cref{g2:def:chain-complex}中的 \(\mathrm d^2=0\) 在这个总链空间上成为
\[
	D^2=0.
\]
矩阵的第 \(j\) 列记为 \(D_j\)。对非零列 \(c\)，令 \(\operatorname{low}(c)\) 表示其最大非零行号；零列不定义 \(\operatorname{low}\)。

\begin{definition}\label{g2:def:persistence-column-reduction}
	\term[bianjiejuzhenyuehua]{边界矩阵约化}[boundary matrix reduction]按 \(j=1,\ldots,m\) 执行以下步骤，并同时记录换基矩阵 \(V\)：
	\begin{enumerate}
		\item 初始化 \(R_j=D_j\)、\(V_j=\mathbf e_j\)，其中 \(\mathbf e_j\) 是第 \(j\) 个坐标向量。
		\item 若 \(R_j\neq0\)，且已有 \(k<j\) 满足 \(\operatorname{low}(R_k)=\operatorname{low}(R_j)=i\)，则作
		\[
			c=\frac{(R_j)_i}{(R_k)_i},\qquad R_j\leftarrow R_j-cR_k,\qquad V_j\leftarrow V_j-cV_k,
		\]
		并对更新后的列重复这一步。
		\item 当 \(R_j=0\)，或它的最大非零行号与所有已处理非零列都不同时，保留该列并转向下一列。
	\end{enumerate}
\end{definition}

每次消元都严格降低当前的最大非零行号，因此循环会终止。只有同维单形的边界列才会有相同的非零行号，所以换基保持链的维数。最终有 \(R=DV\)，其中 \(V\) 是对角元全为 \(1\) 的上三角矩阵；每个前缀的链空间都被 \(V\) 保持。\(R\) 的非零列具有互不相同的最大非零行号，因而线性无关。若采用稠密存储，至多 \(m\) 列、每列至多 \(m\) 次消元、每次至多 \(m\) 次域运算，给出 \(O(m^3)\) 的运算上界。

\begin{lemma}\label{g2:lem:persistence-pivot-birth}
	若 \(\operatorname{low}(R_j)=i\)，则 \(R_i=0\)。此外，\(\dim\sigma_j=\dim\sigma_i+1\)，而同一指标 \(i\) 至多与一个 \(j\) 这样配对。
\end{lemma}

\begin{proof}
	因 \(DR_j=D^2V_j=0\)，并且 \(R_j\) 在第 \(i\) 行的系数非零、在更高行的系数为零，故
	\[
		D_i=-\frac{1}{(R_j)_i}\sum_{\ell<i}(R_j)_\ell D_\ell.
	\]
	因此 \(D_i\) 属于前 \(i-1\) 列的张成空间。前缀内的列操作不改变此空间，而已经约化的非零列彼此线性无关，故第 \(i\) 列约化后只能为零。维数关系由边界降低一维得出，配对唯一性则来自非零列的最大非零行号互不相同。
\end{proof}

\begin{theorem}\label{g2:thm:persistence-reduction-barcode}
	对上述约化结果，每一对 \(\operatorname{low}(R_j)=i\) 给出 \(q=\dim\sigma_i\) 维的区间
	\[
		[\,f(\sigma_i),f(\sigma_j)\,),
	\]
	每一个既满足 \(R_i=0\)、又没有作为任何非零列最大非零行号出现的指标 \(i\)，给出 \(\dim\sigma_i\) 维的区间 \([f(\sigma_i),\infty)\)。删去端点相同的空区间后，所得正是各维持续同调的条形码。
\end{theorem}

\begin{proof}
	把矩阵列视为原来单形基下的链。对于一对 \((i,j)\)，令 \(c=(R_j)_i\)，并取
	\[
		z_i=c^{-1}R_j,\qquad w_j=c^{-1}V_j.
	\]
	于是 \(\mathrm d w_j=z_i\)，且 \(\mathrm d z_i=0\)。链 \(z_i\) 的最大指标为 \(i\)，该项系数为 \(1\)；链 \(w_j\) 的最大指标为 \(j\)，该项系数为 \(c^{-1}\neq0\)。对于没有被配对的零列指标 \(i\)，取 \(z_i=V_i\)，则 \(\mathrm d z_i=R_i=0\)，最大指标仍为 \(i\)。
	\ProseParagraphBreak
	由\cref{g2:lem:persistence-pivot-birth}，配对的左端不会又成为某一对的右端，每个指标恰被上述一个向量占据。这些向量相对于原单形基构成上三角且对角元非零的换基，所以在每个前缀内仍是一组基。边界在这组基上只有 \(w_j\mapsto z_i\mapsto0\) 以及未配对的 \(z_i\mapsto0\)。
	\ProseParagraphBreak
	因此，整个滤过链复形分裂为若干基本块。在一个配对块中，\(z_i\) 出现之前没有链；\(z_i\) 出现后有一个 \(q\) 维同调类；\(w_j\) 出现时，该类成为边界，且二者组成的块同调为零。未配对块则在 \(z_i\) 出现以后始终贡献一个同调类。取同调便得到所述区间。
	\ProseParagraphBreak
	若 \(f(\sigma_i)=f(\sigma_j)\)，两条链在同一真实参数同时出现，该块始终没有同调，所以它贡献空区间。相同滤过值内的排序可能改变具体配对；实际实参数持续模保持不变，\cref{g2:thm:finite-persistence-decomposition}保证删去空区间后的条形码保持不变。
\end{proof}

这个证明还给出了区间的循环代表 \(z_i\)。配对只用 \(R\) 即可读取，要恢复这些代表的填充链 \(w_j\)，则保留同步计算的 \(V\)。注意 \(R=DV\) 只改变了列基；证明所用的零关系是 \(DR=0\)。完整的链换基在上一证明中构造。列消元与分次模方法的关系可参见\cite[第4节]{ZomorodianCarlsson2005Persistence}。

\begin{example}\label{g2:exm:persistence-triangle-reduction}
	对\cref{g2:exm:persistence-triangle}，取顺序
	\[
		v_0,v_1,v_2,e_{01},e_{12},e_{02},\tau.
	\]
	其边界矩阵是
	\[
		D=\begin{pmatrix}
			0&0&0&-1&0&-1&0\\
			0&0&0&1&-1&0&0\\
			0&0&0&0&1&1&0\\
			0&0&0&0&0&0&1\\
			0&0&0&0&0&0&1\\
			0&0&0&0&0&0&-1\\
			0&0&0&0&0&0&0
		\end{pmatrix}.
	\]
	前三列为零；第四列最大非零行为 \(2\)，第五列为 \(3\)。第六列先减去第五列，再减去第四列：
	\[
		D_6-D_5-D_4=(v_2-v_0)-(v_2-v_1)-(v_1-v_0)=0.
	\]
	所以 \(R_6=0\)，同步记录的循环是 \(V_6=e_{02}-e_{12}-e_{01}=-z\)。第七列的最大非零行为 \(6\)，此前没有相同的最大非零行，故保持不变。最终配对为
	\[
		(2,4),\qquad(3,5),\qquad(6,7),
	\]
	未配对零列只有第 \(1\) 列。因此零维条形码为 \(\{[0,1),[0,2),[0,\infty)\}\)，一维条形码为 \(\{[3,4)\}\)。相应的死亡链可分别取 \(e_{01},e_{12},-\tau\)，其边界为 \(v_1-v_0,v_2-v_1,-z\)。这直接解释了两次分支合并及一次圈被填满。
\end{example}

\subsection{稳定性需要的距离与有限性条件}
\label{g2:subsec:C03:04}

给滤过值加一点误差，出生和死亡时刻通常也会移动。一个非常短的区间还可能完全消失。因此，比较两张条形码时，既要允许端点移动，也要允许以与长度相称的代价删除短区间。持续图中的对角线正好记录后一种操作。

\begin{definition}\label{g2:def:barcode-bottleneck}
	设 \(\mathcal B,\mathcal C\) 是两张有限条形码，其中每个区间为 \([b,d)\)，允许 \(d=\infty\)。一个\term[bufenpipei]{部分匹配}[partial matching]在两边各选一些区间，并将它们一一配对，重复区间按不同的副本计数。两个有限区间的配对代价为
	\[
		c([b,d),[b',d'))=\max\{|b-b'|,|d-d'|\}.
	\]
	两个无限区间的配对代价为 \(|b-b'|\)，有限区间与无限区间的配对代价为 \(\infty\)。未匹配的有限区间代价为 \((d-b)/2\)，未匹配的无限区间代价为 \(\infty\)。一个部分匹配的代价取上述所有代价的最大值；全空时为零。\term[pingjingjuli]{瓶颈距离}[bottleneck distance]定义为
	\[
		d_{\mathrm B}(\mathcal B,\mathcal C)=\inf_{\text{部分匹配}}\{\text{匹配代价}\}.
	\]
	\symidx{\(d_{\mathrm B}\)：条形码的瓶颈距离}
\end{definition}

对于有限条形码，部分匹配只有有限种，因此上式的下确界能够达到。未匹配区间的代价来自
\[
	\inf_{x\in\mathbb R}\max\{|b-x|,|d-x|\}=\frac{d-b}{2},
\]
最小值在 \(x=(b+d)/2\) 取得。故可以把一个未匹配区间理解为持续图的点 \((b,d)\) 匹配到对角线 \(b=d\)，对角线允许无限重数。无限长区间必须与无限长区间配对，才能有有限距离；最终同调的维数因而也受到控制。

\begin{definition}\label{g2:def:persistence-interleaving}
	给定持续模 \((V,v)\)、\((W,w)\) 和 \(\varepsilon\geq0\)，一个\term[jiaocuo]{交错}[interleaving]由线性映射
	\[
		\alpha_t:V_t\longrightarrow W_{t+\varepsilon},\qquad
		\beta_t:W_t\longrightarrow V_{t+\varepsilon}
	\]
	组成，要求它们与结构映射相容，即对 \(s\leq t\) 有
	\[
		w_{s+\varepsilon}^{t+\varepsilon}\alpha_s=\alpha_t v_s^t,
		\qquad v_{s+\varepsilon}^{t+\varepsilon}\beta_s=\beta_t w_s^t,
	\]
	并且往返一次等于前移 \(2\varepsilon\)：
	\[
		\beta_{t+\varepsilon}\alpha_t=v_t^{t+2\varepsilon},\qquad
		\alpha_{t+\varepsilon}\beta_t=w_t^{t+2\varepsilon}.
	\]
	存在这些映射时称二者 \(\varepsilon\)-交错。\term[jiaocuojuli]{交错距离}[interleaving distance]是所有可行 \(\varepsilon\) 的下确界，记为 \(d_{\mathrm I}(V,W)\)，无可行值时取 \(\infty\)。
	\symidx{\(d_{\mathrm I}\)：持续模的交错距离}
\end{definition}

交错把“两个滤过相差至多一个时间误差”写成了线性映射。往返关系保证来回比较之后，得到的仍是原来同调类向后传递的结果。

\begin{proposition}\label{g2:prop:filtration-sup-interleaving}
	设 \(f,g\) 是同一个有限单纯复形 \(K\) 上的两个滤过函数。若
	\[
		\|f-g\|_\infty=\max_{\sigma\in K}|f(\sigma)-g(\sigma)|\leq\varepsilon,
	\]
	则各维域系数持续同调 \(H_q(K^f_\bullet;F)\) 与 \(H_q(K^g_\bullet;F)\) 是 \(\varepsilon\)-交错的。
\end{proposition}

\begin{proof}
	若 \(f(\sigma)\leq t\)，则 \(g(\sigma)\leq t+\varepsilon\)。交换 \(f,g\) 也成立，所以对每个 \(t\) 都有
	\[
		K^f_t\subseteq K^g_{t+\varepsilon}\subseteq K^f_{t+2\varepsilon},\qquad
		K^g_t\subseteq K^f_{t+\varepsilon}\subseteq K^g_{t+2\varepsilon}.
	\]
	第一步包含诱导 \(\alpha_t,\beta_t\)。所有箭头都是同一复形中子复形的包含，故它们随 \(t\) 相容；两次包含的复合正是原滤过前移 \(2\varepsilon\) 的包含。取同调即得交错的全部等式。
\end{proof}

从条形码匹配也能构造交错。下面先证明这个方向，以看清短区间代价中的因子 \(1/2\) 是怎样出现的。

\begin{proposition}\label{g2:prop:matching-implies-interleaving}
	对有限构造型持续模，有
	\[
		d_{\mathrm I}(V,W)\leq d_{\mathrm B}(\mathcal B(V),\mathcal B(W)).
	\]
\end{proposition}

\begin{proof}
	取一个代价不超过 \(\varepsilon\) 的部分匹配。对于配对区间 \(I[b,d)\)、\(I[b',d')\)，定义 \(\alpha_t\) 在源与平移后的目标都为 \(F\) 时为恒等映射，其余时为零；反向 \(\beta_t\) 同样定义。端点估计给出 \(b'-\varepsilon\leq b\) 和 \(d'-\varepsilon\leq d\)，所以平移后的目标不会比源更晚开始或更晚结束。这正保证重叠处的恒等映射与结构映射相容。
	\ProseParagraphBreak
	如果 \(t,t+2\varepsilon\in[b,d)\)，则端点估计又给出 \(t+\varepsilon\in[b',d')\)，故往返映射为恒等；其余情况源或最终目标为零，往返映射与结构映射都为零。无限区间的检查相同，只需使用出生端点的不等式。
	\ProseParagraphBreak
	对于未匹配的区间，\(d-b\leq2\varepsilon\)，故它从任意 \(t\) 到 \(t+2\varepsilon\) 的结构映射均为零。它与零模通过零映射组成 \(\varepsilon\)-交错。把各个区间的交错取直和，得到 \(V,W\) 的 \(\varepsilon\)-交错。对匹配代价取下确界便得结论。
\end{proof}

反过来，从两组相容线性映射恢复一一匹配，还需要一个匹配定理。单独比较各时刻的维数不足以完成这一步，因为匹配要同时控制所有区间的两端。

\begin{theorem}[持续模等距定理]\label{g2:thm:persistence-isometry}
	对于同一域上两个有限构造型持续模 \(V,W\)，有
	\[
		d_{\mathrm B}(\mathcal B(V),\mathcal B(W))=d_{\mathrm I}(V,W).
	\]
\end{theorem}

这里引用\cite[第5章，The Isometry Theorem]{ChazalDeSilvaGlisseOudot2016}的等距定理。其代数稳定性方向把交错产生的秩约束转化为区间匹配；与\cref{g2:prop:matching-implies-interleaving}合起来得到上式。书中的定理覆盖更一般的持续模，本节的有限维、有限变化时刻以及左闭右开约定保证可直接应用。特别地，如果存在一个 \(\varepsilon\)-交错，则瓶颈距离至多为 \(\varepsilon\)；有限条形码的最优部分匹配存在，所以能实际选出代价至多为 \(\varepsilon\) 的匹配。

\begin{corollary}\label{g2:cor:finite-filtration-stability}
	同一个有限单纯复形上的滤过函数 \(f,g\) 满足
	\[
		d_{\mathrm B}\bigl(\mathcal B(H_q(K^f_\bullet;F)),\mathcal B(H_q(K^g_\bullet;F))\bigr)
		\leq\|f-g\|_\infty.
	\]
\end{corollary}

\begin{proof}
	\cref{g2:prop:filtration-sup-interleaving}给出以 \(\|f-g\|_\infty\) 为参数的交错，\cref{g2:thm:persistence-isometry}把它转化为同样界限的瓶颈距离。
\end{proof}

例如，三角形中若所有滤过值的误差至多为 \(\varepsilon<1/2\)，原来的一维区间 \([3,4)\) 长度大于 \(2\varepsilon\)，便不能以允许的代价匹配到对角线，必须有另一个一维区间与它配对，且出生和死亡时刻分别相差至多 \(\varepsilon\)。长度不超过 \(2\varepsilon\) 的区间则可能消失。这一结论把测量误差与可追踪的同调特征长度联系起来。

\ProseParagraphBreak

本节的假设各有具体作用。域系数保证线性泛函、补空间及除法消元都可使用；有限单形数保证算法终止，且条形码只有有限条；实参数上在变化时刻采用变化后的值，使区间统一为左闭右开；距离中允许匹配到对角线，则容纳了短寿命特征的消失。把整数系数直接代入区间分解会遗漏映射中的整除信息，\cref{g2:ex:persistence-integer-obstruction}将这一点化为一个乘 \(2\) 的例子。

\begin{chaptersummary}
	\begin{itemize}
		\item Borel 构造 \(X_G=EG\times_GX\) 是 \(BG\) 上纤维为 \(X\) 的丛。自由主作用给出普通商的上同调，固定点保留 \(H^*(BG)\)。圆群等变球面满足 \(x(x+u)=0\)，固定点限制的像由整除关系控制。
		\item Cartan 模型将多项式元置于偶数次数，以外微分和内乘组成等变微分。紧性、不变性及系数约定保证它与 Borel 模型的比较；环面局部化通过倒置非零权消去非固定轨道的贡献。
		\item 谱的结构映射组织悬挂过渡，稳定同伦群由相应直极限给出。球谱、普通同调谱和 K 理论谱具有不同的系数；广义同调的正合性与表示需要所选稳定范畴的基础定理。
		\item 有限滤过的同调及诱导映射组成持续模。域系数下的区间分解保存各同调类的生存范围，区间重数由秩数据决定；只有每一步的 Betti 数不足以恢复这些信息。
		\item 边界矩阵约化通过保持过滤次序的列变换读出生死配对。过滤值的统一扰动界先给交错，再由匹配定理得到瓶颈距离估计，半开区间、零长区间和无限条均按同一约定处理。
	\end{itemize}
\end{chaptersummary}

\BookExercises

\begin{exercise}\label{g2:ex:equivariant-hopf-sphere}
	圆群 \(T\) 在 \(S^{2m+1}\subset\mathbb C^{m+1}\) 上按标量左乘，\(m\ge0\)。用自由作用商和球丛 Gysin 序列两种方法计算 \(H_T^*(S^{2m+1};\mathbb Z)\)。在商 \(\mathbb CP^m\) 上取 \(h=c_1(\gamma^*)\)，核对来自 \(BT\) 的 \(u\) 对应于 \(h\) 还是 \(-h\)。
\end{exercise}
\begin{solution}
	作用自由，Hopf 轨道映射是主丛，所以\cref{g2:prop:free-action-borel-quotient}给出 \(H_T^*(S^{2m+1};\mathbb Z)=H^*(\mathbb CP^m;\mathbb Z)=\mathbb Z[h]/(h^{m+1})\)。要比较 \(u\)，将左主丛改写为右主丛 \(Q\)，令 \(x\cdot z=z^{-1}x\)。其权 \(1\) 关联线是 \(\gamma^*\)：若 \(\ell_x\) 是 \(\mathbb Cx\) 上满足 \(\ell_x(x)=1\) 的复线性泛函，则
	\[
		[x,w]\longmapsto w\ell_x
	\]
	良定义，因为 \(\ell_{z^{-1}x}=z\ell_x\)。在 Borel 构造到 \(BT\) 和到 \(\mathbb CP^m\) 的两张投影下，所拉回的主右丛同构，故 \(u\) 对应 \(c_1(\gamma^*)=h\)。这里右作用的转换解释了为何与 \(BT\) 上 \(u=c_1(\gamma)\) 的写法不矛盾。

	另一方面，Borel 构造是 \(BT\) 上复秩 \(m+1\) 丛 \(L^{\oplus(m+1)}\) 的单位球丛。其底层定向实丛的 Euler 类为顶陈类 \(u^{m+1}\)。\cref{g2:thm:thom-gysin}中乘 \(u^{m+1}\) 的映射在 \(\mathbb Z[u]\) 上单射，因此总空间的上同调由其余核给出，且由底空间拉回生成。所得环为
	\[
		H_T^*(S^{2m+1};\mathbb Z)=\mathbb Z[u]/(u^{m+1}),
	\]
	与第一种方法及 \(u\leftrightarrow h\) 的对应一致。
\end{solution}

\begin{exercise}\label{g2:ex:weighted-sphere-coefficients}
	设 \(m\ge2\)，圆群按 \(z\cdot[z_0:z_1]=[z_0:z^m z_1]\) 作用于 \(\mathbb CP^1\)。计算整系数等变上同调环及固定点限制的像。若素数 \(p\mid m\)，改用 \(\mathbb F_p\) 系数后，证明固定点限制即使倒置 \(u\) 仍不是单射。比较有理系数下的情形。
\end{exercise}
\begin{solution}
	权 \(m\) 的关联线是 \(L^{\otimes m}\)，第一陈类为 \(mu\)。Borel 构造因而是 \(\mathbb P(\underline{\mathbb C}\oplus L^{\otimes m})\)。取重言对偶的一阶类 \(x\)，射影丛公式给出
	\[
		H_T^*(\mathbb CP^1;\mathbb Z)=\mathbb Z[u,x]/\bigl(x(x+mu)\bigr),
	\]
	并以 \(1,x\) 为 \(\mathbb Z[u]\) 模基。两个固定点限制为 \(x|_N=0\)、\(x|_S=-mu\)，所以像为 \(P-Q\in mu\mathbb Z[u]\) 的多项式对。由于 \(mu\) 非零且不是零因子，整系数限制仍单射，但只倒置 \(u\) 后其像仍要求差属于 \(m\mathbb Z[u,u^{-1}]\)。

	改用 \(\mathbb F_p\) 后，整数上同调无挠，万有系数定理或射影丛基直接给出 \(\mathbb F_p[u,x]/(x^2)\)。其中非零的 \(x\) 限制为 \((0,0)\)，它又是自由模基的一项，故倒置 \(u\) 后仍非零。若采用有理系数，\(m\) 已可逆，继续倒置 \(u\) 便可解任意 \(P-Q=mu b\)，因此限制变为同构。有限稳定子中的 \(m\) 扭信息决定了这三种系数下的差别。
\end{solution}

\begin{exercise}\label{g2:ex:finite-isotropy-orbit}
	设 \(C_m\subset T=S^1\) 是 \(m\) 次单位根群，\(m\ge2\)。计算轨道 \(T/C_m\) 的整系数与有理系数等变上同调，并说明为什么其轨道空间的一点不能恢复整系数答案。
\end{exercise}
\begin{solution}
	由\cref{g2:prop:borel-homogeneous-orbit}，所求为 \(H^*(BC_m)\)。将 \(T/C_m\) 识别为单位圆时，可用 \(zC_m\mapsto z^m\)，所以 \(T\) 作用成为权 \(m\) 的旋转。其 Borel 构造是 \(L^{\otimes m}\) 的单位圆丛，Euler 类为 \(mu\)。Gysin 序列中的乘法 \(\mathbb Z[u]\xrightarrow{mu}\mathbb Z[u]\) 单射，故奇次群为零，零次为 \(\mathbb Z\)，每个正偶次为 \(\mathbb Z/m\)。这些类均由底空间拉回的 \(u\) 的幂生成，得到
	\[
		H_T^*(T/C_m;\mathbb Z)\cong\mathbb Z[u]/(mu).
	\]
	有理化后 \(m\) 可逆，关系强迫 \(u=0\)，所以只剩零次的 \(\mathbb Q\)。普通轨道空间是一点，其整数上同调没有这些正次扭类；扭类记录的是被普通取商压去的有限稳定子。
\end{solution}

\begin{exercise}\label{g2:ex:cartan-free-torus}
	令 \(T=S^1\) 作用于 \(M=(\mathbb R/\mathbb Z)^2\)，作用为 \(t\cdot(\theta,\varphi)=(\theta+t,\varphi)\)。直接计算 Cartan 复形的上同调环，并与商空间的 de Rham 上同调比较。指出复形中 \(\mathrm d\theta\) 与 \(\mathrm d\varphi\) 的不同作用。
\end{exercise}
\begin{solution}
	沿 \(\theta\) 平移不变的形式，其系数只依赖 \(\varphi\)。因此整个复形是分次张量积
	\[
		\bigl(\mathbb R[u]\otimes\Lambda(\mathrm d\theta),\ \mathrm d_T(\mathrm d\theta)=-u\bigr)
		\otimes\bigl(\Omega(S^1_\varphi),\mathrm d\bigr).
	\]
	第一个因子到 \(\mathbb R\) 的映射把 \(u\) 与 \(\mathrm d\theta\) 都送到零。它是链同伦等价：在多项式基上令 \(h(1)=0\)、\(h(u^j)=-u^{j-1}\mathrm d\theta\)（\(j\ge1\)），并令 \(h(u^j\mathrm d\theta)=0\)，则 \(\mathrm d_T h+h\mathrm d_T\) 等于恒等映射减去常数投影。张量上第二个复形后仍给出链同伦，故上同调环等于
	\[
		H^*(S^1_\varphi;\mathbb R)=\Lambda_{\mathbb R}([\mathrm d\varphi]).
	\]
	这正是自由作用商 \(M/T=S^1_\varphi\) 的上同调。形式 \(\mathrm d\varphi\) 水平且闭，下降为商上的生成元；\(\mathrm d\theta\) 沿轨道不水平，它在等变微分下产生 \(-u\)，使群方向的正次多项式成为边界。
\end{solution}

\begin{exercise}\label{g2:ex:cartan-sphere-lifts}
	采用\cref{g2:ex:cartan-rotating-sphere}的 \(s,\omega,V,u,x\)。求所有满足 \(\mathrm d\mu=\iota_V\omega\) 的光滑函数 \(\mu\)，并在 \(H_T^2(\mathbb CP^1;\mathbb R)\) 中识别 \([\omega+u\mu]\)。其中哪些类来自整数等变上同调？说明在北极规定 \(\mu(N)=0\) 起了什么作用。
\end{exercise}
\begin{solution}
	因为 \(\iota_V\omega=-\mathrm ds\)，方程等价于 \(\mathrm d(\mu+s)=0\)。球面连通，所以全部解为 \(\mu=-s+c\)，其中 \(c\in\mathbb R\)。于是
	\[
		[\omega+u\mu]=[\omega-su]+cu=x+cu.
	\]
	其两极限制分别为 \(cu\) 与 \((c-1)u\)。由\cref{g2:thm:equivariant-sphere-ring}，整数二次群以 \(u,x\) 为自由基，实系数群由它扩张系数得到。因此 \(x+cu\) 是整数类的像，当且仅当 \(c\in\mathbb Z\)。这也能直接从北极限制看出必要性，再用整数基看出充分性。

	条件 \(\mu(N)=0\) 选定 \(c=0\)，从所有普通类为 \([\omega]\) 的等变提升中选出 \(x\)。普通面积形式只能确定微分 \(\mathrm d\mu\)，不能确定加法常数；固定一个点上的值才完成这项归一化。若 \(c\) 为整数，给重言对偶线的等变结构再张量权 \(c\) 的平凡线，其一阶类就是 \(x+cu\)。
\end{solution}

\begin{exercise}\label{g2:ex:localization-noncompact-orbits}
	令 \(T=(S^1)^2\)，并令 \(H_n=\ker\chi_n\)，其中 \(\chi_n(z_1,z_2)=z_1z_2^n\)，\(n=1,2,\ldots\)。在光滑非紧流形 \(M=\coprod_{n\ge1}T/H_n\) 上取逐分支的左作用。证明 \(M^T=\varnothing\)，但对 \(\mathbf k=\mathbb Q\) 或 \(\mathbb R\)，\(S^{-1}H_T^*(M;\mathbf k)\ne0\)，其中 \(S\) 仍由所有非零特征标类生成。解释有限胞腔证明中的哪一步不能用于这个例子。
\end{exercise}
\begin{solution}
	每个 \(H_n\) 都是真一维子环面，故每个轨道均无固定点。空间 \(M\) 是可数个圆的不交并，具有第二可数的一维光滑结构；它不紧。由齐性轨道的 Borel 计算与不交并的上同调，逐次数有
	\[
		H_T^*(M;\mathbf k)=\prod_{n\ge1}H^*(BH_n;\mathbf k).
	\]
	这里乘积按每一个上同调次数理解。以 \(t\mapsto(t^{-n},t)\) 参数化 \(H_n\)，并以 \(v\) 表示 \(H^2(BH_n;\mathbf k)\) 的生成元，则限制为 \(u_1\mapsto-nv\)、\(u_2\mapsto v\)。所以该分支的系数环是
	\[
		\mathbf k[u_1,u_2]/(u_1+nu_2).
	\]
	取零次类 \(\mathbf 1=(1,1,\ldots)\)。任意 \(s\in S\) 是有限个非零整系数线性式的乘积。在第 \(n\) 个分支上，一个因子 \(au_1+bu_2\) 限制为 \((b-an)v\)。若 \(a\ne0\)，它至多在一个正整数 \(n\) 上为零；若 \(a=0\)，则因 \(b\ne0\) 而在所有分支上都非零。故有限个因子的乘积在除有限多个以外的分支上均非零，\(s\mathbf 1\ne0\)。

	因此没有 \(s\in S\) 消灭 \(\mathbf 1\)，局部化中的 \(\mathbf 1/1\) 非零。每个轨道的等变上同调单独都被一个权消灭，但无穷多个轨道需要无穷多个不同的权；乘法集中的分母只有有限个因子。有限胞腔证明所用的有限归纳在这里不能替代为对所有分支的一次统一消灭。
\end{solution}

\begin{exercise}\label{g2:ex:spectrum-degree-two-system}
	令预谱 \(E\) 的第 \(n\) 层为 \(S^n\)，每个结构映射 \(\Sigma S^n=S^{n+1}\to S^{n+1}\) 的度数为 \(2\)。计算 \(\pi_0(E)\) 和所有负次稳定群，并判断它是否与球预谱稳定等价。
\end{exercise}
\begin{solution}
	当 \(n\ge2\) 时，\(\pi_n(S^n)\cong\mathbb Z\) 由度数识别。悬挂保持度数，随后与度数 \(2\) 的结构映射复合使度数乘 \(2\)，所以
	\[
		\pi_0(E)=\varinjlim\bigl(\mathbb Z\xrightarrow{\times2}\mathbb Z\xrightarrow{\times2}\cdots\bigr).
	\]
	把第 \(n\) 项的整数 \(a\) 送到 \(a/2^n\in\mathbb Q\)。这个赋值与过渡相容，因为 \(2a/2^{n+1}=a/2^n\)。所得像为所有分母是二次幂的有理数，即 \(\mathbb Z[1/2]\)。若两个分数相等，乘以足够大的共同二次幂后，两个整数代表就在同一个阶段相等，所以该映射也单射。因此 \(\pi_0(E)\cong\mathbb Z[1/2]\)。

	对 \(q<0\)，充分大的 \(n\) 满足 \(2\le q+n<n\)，球面连通性使各项 \(\pi_{q+n}(S^n)=0\)，故 \(\pi_q(E)=0\)。球预谱的零次稳定群是 \(\mathbb Z\)，与 \(\mathbb Z[1/2]\) 不同构：前者的非零元不可能无限次被 \(2\) 整除，后者的每个元素都可被 \(2\) 整除。因此两者不稳定等价，尽管每层的空间完全相同。
\end{solution}

\begin{exercise}\label{g2:ex:generalized-moore-homology}
	设 \(m\ge2\)、\(r\ge2\)，令 \(M=S^r\cup_m e^{r+1}\)，其中附着映射的度数为 \(m\)。对任意谱 \(E\)，从映射锥长正合列证明有自然短正合列
	\[
		0\longrightarrow\pi_{q-r}(E)/m\pi_{q-r}(E)
		\longrightarrow\widetilde E_q(M)
		\longrightarrow\{a\in\pi_{q-r-1}(E):ma=0\}\longrightarrow0.
	\]
	据此分别计算 \(E=H\mathbb Z\) 与 \(E=KU\) 时的全部群。
\end{exercise}
\begin{solution}
	空间余纤维序列为 \(S^r\xrightarrow{m}S^r\to M\to S^{r+1}\)。稳定范畴中的度 \(m\) 球面映射，是恒等映射在阿贝尔映射群中的 \(m\) 倍，所以在任意广义同调上都作用为乘 \(m\)。用球面系数公式，长正合列的相关部分为
	\[
		\pi_{q-r}(E)\xrightarrow{\times m}\pi_{q-r}(E)
		\longrightarrow\widetilde E_q(M)\longrightarrow
		\pi_{q-r-1}(E)\xrightarrow{\times m}\pi_{q-r-1}(E).
	\]
	在两端取余核和核，即得短正合列。一般情形不能仅凭这条列指定自然分裂。

	对 \(H\mathbb Z\)，系数只在零次为 \(\mathbb Z\)，乘 \(m\) 在 \(\mathbb Z\) 上单射，故 \(\widetilde H_q(M;\mathbb Z)\) 只在 \(q=r\) 为 \(\mathbb Z/m\)，其余为零。对 \(KU\)，所有系数均无 \(m\) 挠元，右端核为零；偶次系数为 \(\mathbb Z\)，奇次为零，因此
	\[
		\widetilde{KU}_q(M)=
		\begin{cases}\mathbb Z/m,&q-r\text{ 为偶数},\\0,&q-r\text{ 为奇数}.
		\end{cases}
	\]
	这个结果对所有整数 \(q\) 成立。同一张附着映射在普通理论中只留下一个次数的挠群，在复 \(K\) 同调中则按二周期重复出现。
\end{solution}

\begin{exercise}\label{g2:ex:stable-wedge-direct-sum-product}
	令 \(X=\bigvee_{j=1}^{\infty}S^2\)，带 CW 弱拓扑。计算 \(\widetilde{KU}_0(X)\) 与 \(\widetilde{KU}^{\,0}(X)\)，并说明为何前者是直和、后者是乘积。比较这两个群的基数。
\end{exercise}
\begin{solution}
	每份二球面的两种群分别为 \(\pi_{-2}(KU)=\mathbb Z\) 与 \(\pi_2(KU)=\mathbb Z\)。\cref{g2:prop:spectrum-exact-wedge}于是给出
	\[
		\widetilde{KU}_0(X)=\bigoplus_{j=1}^{\infty}\mathbb Z,
		\qquad\widetilde{KU}^{\,0}(X)=\prod_{j=1}^{\infty}\mathbb Z.
	\]
	同调的直和意味着每个类只含有限多个分支的非零贡献；这来自球谱的紧性。上同调是从整个楔和到表示谱的映射，余积性质允许在每个分支上独立指定一个整数，没有有限支撑要求，因而得到乘积。

	直和是可数集：按有限支撑所在的最大下标取可数并，每个部分包含于某个有限次 \(\mathbb Z^N\)。乘积包含全部二进制序列 \(\{0,1\}^{\mathbb N}\)，由对角线论证是不可数集。因此两个群甚至不可能作为集合双射，更不可能同构。有限楔和时直和与乘积相同，不能由有限例子省略无限楔和公理的区别。
\end{solution}

\begin{exercise}\label{g2:ex:persistence-same-betti}
	考虑两个在时刻 \(0,1,2\) 分别为 \(F,F^2,F\) 的持续模，最早时刻之前为零，最后时刻以后恒等延拓。第一个箭头均为 \(x\mapsto(x,0)\)，第二个箭头分别为 \((a,b)\mapsto a\) 和 \((a,b)\mapsto b\)。计算两张条形码及从时刻 \(0\) 到时刻 \(2\) 的秩，并说明各时刻的 Betti 数为什么不能确定条形码。
\end{exercise}
\begin{solution}
	记 \(F^2\) 的标准基为 \(e_1,e_2\)。第一种情形中，时刻 \(0\) 的生成元依次映为 \(e_1\) 和 \(1\)，始终存活；\(e_2\) 在时刻 \(1\) 出生，下一步映为零。所以条形码是
	\[
		\{[0,\infty),[1,2)\},\qquad r(0,2)=1.
	\]
	第二种情形中，时刻 \(0\) 的生成元映为 \(e_1\)，随后死亡；时刻 \(1\) 新出现的 \(e_2\) 映为末项的 \(1\)，以后始终存活。因此条形码是
	\[
		\{[0,2),[1,\infty)\},\qquad r(0,2)=0.
	\]
	两者各时刻的维数都是 \(1,2,1\)，但第一个持续模把最早的类保留到最后，第二个则以一个新类接替它。条形码记录的映射信息正体现在这个差别上。
\end{solution}

\begin{exercise}\label{g2:ex:persistence-square}
	用对角线 \(ac\) 将四边形 \(abcd\) 分为三角形 \(abc,acd\)。令所有顶点在 \(0\) 出现，四条外边在 \(1\) 出现，对角线在 \(2\) 出现，两个面分别在 \(3,4\) 出现。计算各维条形码，并给出两个一维区间的循环代表及死亡链。若把对角线与面 \(abc\) 的滤过值都改为 \(3\)，条形码怎样变化？
\end{exercise}
\begin{solution}
	把外边定向为 \(e_{ab},e_{bc},e_{cd},e_{ad}\)，对角线定向为 \(e_{ac}\)。时刻 \(1\) 的外圈是
	\[
		z=e_{ab}+e_{bc}+e_{cd}-e_{ad}.
	\]
	这是连通图中唯一的独立圈：图有 \(4\) 条边、\(4\) 个顶点，边界矩阵秩为 \(3\)，故一维循环空间维数为 \(1\)。对角线出现后，新增独立循环
	\[
		y=e_{ab}+e_{bc}-e_{ac}.
	\]
	它与 \(z\) 独立，因为 \(y\) 的对角线系数非零，而 \(z\) 的该系数为零。两个面的边界为
	\[
		\mathrm d[abc]=y,\qquad \mathrm d[acd]=z-y.
	\]
	在时刻 \(3\)，\(y\) 被填掉，\(z\) 仍非零；在时刻 \(4\)，\(z=\mathrm d([abc]+[acd])\) 也被填掉。因此
	\[
		\mathcal B(H_1)=\{[1,4),[2,3)\}.
	\]
	循环代表可取 \(z,y\)，死亡链分别为 \([abc]+[acd]\) 和 \([abc]\)。时刻 \(1\) 四个顶点同时连成一个分支，所以
	\[
		\mathcal B(H_0)=\{[0,\infty),[0,1),[0,1),[0,1)\}.
	\]
	两个二维面的边界独立，故 \(H_2=0\)，更高维也为零。
	\ProseParagraphBreak
	若对角线与第一个面同时在 \(3\) 出现，新增的链 \(y\) 立即成为该面的边界，不在任何参数处贡献同调。矩阵排序会给出空区间 \([3,3)\)，删除后只剩 \(\mathcal B(H_1)=\{[1,4)\}\)；零维条形码不变。
\end{solution}

\begin{exercise}\label{g2:ex:persistence-distance-zero}
	设 \(V\) 是具有有限条形码的持续模。直接从定义证明，\(V\) 到零模的交错距离等于条形码中最长有限区间长度的一半；若存在无限区间，则距离为 \(\infty\)。全空条形码的距离约定为零。再计算 \(\{[0,L)\}\) 与 \(\{[h,L+h)\}\) 的瓶颈距离，其中 \(L>0\)、\(h\geq0\)。
\end{exercise}
\begin{solution}
	与零模之间的两组映射只能是零映射。因此存在 \(\varepsilon\)-交错，当且仅当所有结构映射 \(v_t^{t+2\varepsilon}\) 都为零。区间 \([b,d)\) 满足此条件，当且仅当 \(d-b\leq2\varepsilon\)：若严格长于 \(2\varepsilon\)，取 \(t=b\)，结构映射为恒等；若长度不超过 \(2\varepsilon\)，\(t\) 和 \(t+2\varepsilon\) 不可能同时落在该左闭右开区间内。取有限直和，便得到所述最长长度公式。无限区间在任意有限前移后仍有非零结构映射，所以不存在有限交错参数。
	\ProseParagraphBreak
	对于题中的两张单区间条形码，部分匹配只有两种实质选择：将两条配对，代价为 \(h\)；或两条都匹配到对角线，代价为 \(L/2\)。因此
	\[
		d_{\mathrm B}\bigl(\{[0,L)\},\{[h,L+h)\}\bigr)=\min\{h,L/2\}.
	\]
	当位移很大时，最优匹配会把两条都送到对角线。这说明瓶颈距离同时衡量端点移动与特征消失的代价。
\end{solution}

\begin{exercise}\label{g2:ex:persistence-integer-obstruction}
	把两项持续模 \(\mathbb Z\xrightarrow{\times2}\mathbb Z\) 与 \(\mathbb Z\xrightarrow{\mathrm{id}}\mathbb Z\) 比较，证明它们不同构，虽然各项及箭头的秩相同。再证明前者不能分解为整数区间模的直和，其中整数区间模在区间内取 \(\mathbb Z\)，内部映射取恒等。最后将系数分别换为 \(\mathbb Q\) 和 \(\mathbb F_2\)，计算前者的条形码；两时刻取 \(0,1\)，此前为零、此后恒等延拓。
\end{exercise}
\begin{solution}
	两个箭头的余核分别为 \(\mathbb Z/2\mathbb Z\) 和零，持续模同构会诱导箭头余核的同构，所以它们不同构。它们的两个对象秩都为 \(1\)，箭头的像也都秩为 \(1\)；秩没有记录像 \(2\mathbb Z\) 在目标中的指数。
	\ProseParagraphBreak
	若第一个图能分解为整数区间模，第一项秩为 \(1\) 表明恰有一条区间覆盖时刻 \(0\)。箭头非零要求这条区间也覆盖时刻 \(1\)；末项秩为 \(1\) 又排除了其他区间。于是箭头在某组整数基下应为恒等，其余核为零，与 \(\mathbb Z/2\mathbb Z\) 矛盾。
	\ProseParagraphBreak
	换为 \(\mathbb Q\) 后，乘 \(2\) 是同构，选择末项的基向量为 \(2\) 即把箭头写成恒等，条形码为 \(\{[0,\infty)\}\)。换为 \(\mathbb F_2\) 后，箭头为零，第一个类在 \(1\) 死亡，末项出现另一个类，条形码为 \(\{[0,1),[1,\infty)\}\)。域系数既使区间分解成立，也会影响实际得到的条形码。
\end{solution}

% !TeX root = ../../Book2.tex
% 纲要标题：## 附录 D　叶状结构与光滑映射奇点
% 纲要位置：几何学计划纲要.md，第 2063 行。
% 正文状态：初稿；含正文、证明、例子、习题与完整解答，尚未正式审核。

% 纲要内容开始
% 
% > 以 holonomy 和喷射横截性两条明确工具链展开，不把叶空间和函数奇点视为同一分类问题。
% 
% * **叶、横截面与 holonomy**
%   * 沿叶路径的横向运输
%   * 芽群和共轭不变量
%   * 有限 holonomy 与 Reeb 局部稳定
% 
% * **叶空间与整体动力学**
%   * 环面线性叶状结构
%   * Reeb 分量及非分离叶空间
%   * Reeb 整体稳定与圆周纤维丛
% 
% * **喷射、折叠与尖点**
%   * 喷射丛和接触阶
%   * 秩分层及喷射横截
%   * 折叠与尖点的局部标准形
%   * 稳定性及维数限制
% 
% ---
% 
% 纲要内容结束

\chapter{叶状结构与光滑映射奇点}
\label{g2:app:D}
\BookChapterNumber{g2:app:D}

\begin{chapterabstract}
	叶状结构的局部坐标把流形切成平行叶片，但沿一张叶绕行后，邻近叶片可能发生非平凡的横向位移。横截面返回映射芽记录这种位移；有限 holonomy 的局部模型、环面直线叶与 Reeb 分量展示了它和叶空间拓扑的关系。对于光滑映射，喷射把指定阶数的导数组织成几何对象，横截性由此控制临界集的维数。折叠与尖点的标准形进一步说明纤维怎样合并或分开，以及这些局部形态怎样在扰动后保留。两条支线分别围绕沿叶运输与导数数据展开，最后都需要区分局部模型和整体结论。
\end{chapterabstract}

\begin{learninggoals}
	\begin{itemize}
		\item 构造沿叶路径的横向运输，核对路径复合次序，并辨认完整返回芽与线性返回映射的差别。
		\item 解释有限 holonomy 的商模型，计算附近叶的覆盖次数及叶空间，理解紧叶假设的作用。
		\item 计算环面线性叶状结构与 Reeb 分量的叶、返回映射及商拓扑，说明整体纤维丛结论需要的条件。
		\item 用喷射丛保存有限阶导数，计算矩阵秩层的余维，并通过多项式扰动理解喷射横截性。
		\item 推导折叠与平面尖点的局部标准形，计算临界值曲线和局部纤维，区分局部持续、一般位置与整体稳定。
	\end{itemize}
\end{learninggoals}

在一张叶上选一个点，再在它旁边放一张小横截面。沿叶绕一圈，怎样判断最初邻近的点是否回到了原处？第一卷第13章已经用这个问题引入横截面返回映射芽。本附录将把路径选择、同伦与坐标改变的影响逐一说明，再用这一工具比较两种看似相近的现象：紧叶附近可以是一束封闭的叶，也可以有非紧叶不断绕向它。

\ProseParagraphBreak

另一类问题从映射的临界点出发。曲面投影到平面时，投影边缘常是一条折叠曲线，曲线上还可能出现尖点。只知道微分降秩，尚不能分辨这两种形态；需要再看核方向的二阶、三阶变化。喷射把这些导数放在同一空间中，便于使用第14章的横截工具。标准形则把抽象条件变成可以画图、可以计算原像数目的公式。

\ProseParagraphBreak

叶状结构部分接第一卷第13.4节的 Frobenius 定理及第13.5节的叶空间与返回映射；光滑映射部分接本卷第14章的参数横截性、第15章的淹没与纤维丛，以及第17章的 Morse 引理。正文会在调用处提醒所需条件。有限 holonomy 的一般邻域构造、Reeb 整体稳定、光滑准备定理以及一般位置和整体稳定的指定结果，分别在使用处给出来源；局部运输、有限群模型、代表叶空间、喷射扰动和折叠、尖点的坐标计算在这里展开。

\section{叶、横截面与 holonomy}
\label{g2:sec:D01}

在一张叶状坐标图里，各张叶片看起来都是平行的；沿叶绕行一周后，附近叶片却可能错开。横截面上的返回映射把这种错开变成一个可计算的对象。本节先把返回映射的选择问题处理清楚，再解释有限返回群为何能控制一张紧叶的整个邻域。

\subsection{沿叶路径的横向运输}
\label{g2:subsec:D01:01}

本节取无边界光滑流形 \(M^{p+q}\) 上的光滑正则叶状结构 \(\mathcal F\)，叶维数为 \(p\)、余维数为 \(q\)。第一卷第13.4节的 Frobenius 定理说明，常秩对合分布能给出这样的局部坐标。我们在此使用的是定理已经得到的坐标结论：每点附近有
\[
	\varphi:U\longrightarrow B^p\times B^q,\qquad z\longmapsto(x,y),
\]
其中 \(B^p,B^q\) 是开球，连通的集合 \(y=\text{常数}\) 是\term[yepian]{叶片}[plaque]。在充分小的连通过渡区域，坐标变化具有形式
\[
	(x,y)\longmapsto\bigl(a(x,y),b(y)\bigr).
\]
第二项不依赖 \(x\)，因为它在每张连通叶片上恒定。叶由能够沿相交叶片接续的点组成，并带有这些叶片给出的流形拓扑；这项拓扑未必等于 \(M\) 的子空间拓扑。

\begin{definition}\label{g2:def:foliation-transversal-germ}
	通过 \(x\in M\) 的\term[hengjiemian]{横截面}[transversal]是一个 \(q\) 维局部嵌入子流形 \(T\)，满足
	\[
		\mathrm T_xM=\mathrm T_x\mathcal F\oplus\mathrm T_xT.
	\]
	把 \(T\) 缩小后，这个直和条件在各点成立。若两个定义在 \(x\) 的邻域上的光滑映射在某个更小邻域内相同，就说它们在 \(x\) 处有相同的\term[yingxingshunya]{映射芽}[map germ]。记局部微分同胚 \(h:(T,x)\to(T',x')\) 的芽为 \([h]_x\)。
\end{definition}

在叶状图中，横向坐标投影 \(\pi_y\) 在横截面上的微分可逆，故其限制是局部微分同胚。因此，同一张叶片内两点 \(x,x'\) 处的横截面之间有局部对应
\begin{equation}\label{g2:eq:plaque-transverse-transport}
	(\pi_y|_{T'})^{-1}\circ(\pi_y|_T).
\end{equation}
它把一点移到与它处于同一叶片的另一横截面交点。这里无需给叶选择度量，也无需规定在叶上走多快；保留下来的是叶片的接续关系。

\ProseParagraphBreak

令 \(\alpha:[0,1]\to L\) 是一张叶中的连续路径，端点为 \(x_0,x_m\)，在端点选定横截面 \(T_0,T_m\)。将区间有限细分，使每段路径落在一张叶状图的一张叶片内，在分点处另选小横截面 \(T_i\)。逐段用式\eqref{g2:eq:plaque-transverse-transport}得到 \(h_i:(T_{i-1},x_{i-1})\to(T_i,x_i)\)，再取复合
\[
	h_\alpha=h_m\circ\cdots\circ h_1.
\]
每一步只在足够小的邻域上有定义；有限次缩小起点邻域便使全部复合有定义。我们关心的是最终的芽，而不是某个任意选定的定义域。

\begin{figure}[htbp]
	\centering
	% !TeX root = ../Book2.tex
\begin{tikzpicture}[x=1cm,y=1cm,>=stealth,font=\small]
	\fill[blue!5] (0,-0.6) rectangle (3.1,1.3);
	\fill[blue!5] (2.7,-0.45) rectangle (5.8,1.45);
	\fill[blue!5] (5.4,-0.3) rectangle (8.5,1.6);
	\draw[gray!65] (0,-0.6) rectangle (3.1,1.3);
	\draw[gray!65] (2.7,-0.45) rectangle (5.8,1.45);
	\draw[gray!65] (5.4,-0.3) rectangle (8.5,1.6);
	\draw[thick,->] (0.4,0) .. controls (2.4,0.1) and (5.4,0.3) .. (8.1,0.6);
	\draw[blue!70!black,dashed,->] (0.4,0.55) .. controls (2.4,0.7) and (5.4,0.85) .. (8.1,1.1);
	\draw[thick] (0.6,-0.35)--(0.6,0.95) node[above] {\(T_0\)};
	\draw[thick] (2.9,-0.2)--(2.9,1.1) node[above] {\(T_1\)};
	\draw[thick] (5.6,0)--(5.6,1.3) node[above] {\(T_2\)};
	\draw[thick] (7.9,0.2)--(7.9,1.5) node[above] {\(T_3\)};
	\fill (0.6,0.011) circle (1.5pt) node[below left] {\(x_0\)};
	\fill[blue!70!black] (0.6,0.565) circle (1.5pt) node[left] {\(v\)};
	\fill (7.9,0.578) circle (1.5pt) node[below left] {\(x_3\)};
	\fill[blue!70!black] (7.9,1.081) circle (1.5pt) node[above left] {\(h_\alpha(v)\)};
	\node at (4.2,-0.12) {\(\alpha\)};
	\node[gray!80!black] at (1.8,-0.43) {\(U_1\)};
	\node[gray!80!black] at (4.5,-0.28) {\(U_2\)};
	\node[gray!80!black] at (7.2,-0.13) {\(U_3\)};
\end{tikzpicture}

	\caption[沿叶路径的横向运输]{沿叶路径的横向运输。实线表示所选叶上的路径，虚线表示从横截面上的邻近点出发、按相同图链接续的叶片；每一步只使用足够小的局部邻域。}
	\label{g2:fig:holonomy-transverse-transport}
\end{figure}

\begin{proposition}\label{g2:prop:holonomy-transport}
	固定端点横截面后，\([h_\alpha]_{x_0}\) 与细分、中间横截面及叶状图的选择无关，只依赖于 \(\alpha\) 在叶 \(L\) 内保持端点不动的同伦类。沿用本书先走 \(\alpha\)、再走 \(\beta\) 的路径乘法 \(\alpha*\beta\)，有
	\begin{equation}\label{g2:eq:holonomy-path-composition}
		h_{\alpha*\beta}=h_\beta\circ h_\alpha,\qquad
		h_{\alpha^{-1}}=h_\alpha^{-1},\qquad h_{\text{常路径}}=\operatorname{id},
	\end{equation}
	其中等式均为芽的等式。
\end{proposition}
\begin{proof}
	先在一张图内考虑三张横截面。记 \(p_i=\pi_y|_{T_i}\)，从第一张到第三张的直接对应为 \(p_3^{-1}p_1\)，经过第二张的对应为
	\[
		(p_3^{-1}p_2)(p_2^{-1}p_1)=p_3^{-1}p_1.
	\]
	因此插入分点和中间横截面不改变芽。若改用另一张叶状图，其横向坐标在小重叠区域满足 \(y'=b(y)\)，相应对应变成
	\[
		(b\circ p_2)^{-1}\circ(b\circ p_1)=p_2^{-1}\circ p_1.
	\]
	所以一次局部换图也不改变芽。两种覆盖路径的有限图链有共同细分：在每个分点附近缩小到两张相关图的重叠区域，并用叶内的小邻域使所用路径段都处于同一张叶片。逐段应用上述等式，再消去中间横截面的坐标变换，就得到两种构造的相同芽。有限次比较允许一次性缩小所有所需的定义域。

	再设 \(F:[0,1]^2\to L\) 是保持端点不动的同伦。对 \(L\) 使用叶自身的流形拓扑，故 \(F\) 的每一点都有一个小矩形邻域，其像落在一张叶状图的一张叶片内。紧性与 Lebesgue 数引理给出有限矩形网格，使每个小格的像都有这个性质。在网格顶点的像上选横截面；当多个顶点映到同一点时，可选同一横截面。沿一个小格的边界接续横向对应，复合为恒等芽，因为四条边都可在同一张图中计算，四次横向坐标的变换逐项消去。

	于是沿网格的一条路线，可以把某小格的两条相邻边换成另外两条，端到端的运输芽保持不变。依次越过全部小格，就把同伦下边的路线移到上边；左右两边是端点处的常路径，运输均为恒等。故上下两条路径的运输芽相同。这个论证用到有限个芽的等式，仍可在一个足够小的共同起点邻域内实现。最后，拼接两条路径的图链正好给出式\eqref{g2:eq:holonomy-path-composition}的复合次序；反向走过同一图链给出逆映射，常路径则给出恒等映射。
\end{proof}

证明中特意把同伦放在叶内：邻近叶上的点可以在 \(M\) 中很接近，却没有在同一张叶内连接它们的短路径。横向运输所保存的是指定叶中路径的变形，而不是任意环境同伦。

\subsection{芽群和共轭不变量}
\label{g2:subsec:D01:02}

闭路径的起点与终点相同，因此可以使用同一横截面。不同闭路径的返回芽组成一个群，记录附近叶片绕过这张叶时可能发生的全部局部变化。

\begin{definition}\label{g2:def:leaf-holonomy-group}
	固定 \(x\in L\) 及通过它的横截面 \(T\)。所有基于 \(x\) 的叶内闭路径所给的芽构成\term[hengjiemianfanhuinyingxingshunqun]{横截面返回映射芽群}[holonomy group]
	\[
		H(L,x,T)=\{[h_\gamma]_x:[\gamma]\in\pi_1(L,x)\}
		\subset\operatorname{Diff}_{\mathrm{germ}}(T,x).
	\]
	其中右端是固定 \(x\) 的光滑局部微分同胚芽在通常复合下构成的群。\symidx{\(H(L,x,T)\)：叶的横截面返回映射芽群}
\end{definition}

由于路径乘法先走左边，式\eqref{g2:eq:holonomy-path-composition}说明 \([\gamma]\mapsto h_\gamma\) 是反同态；也可将它看作基本群在横向芽上的右作用。若希望使用通常的左表示，则应写成
\begin{equation}\label{g2:eq:holonomy-left-representation}
	\rho:\pi_1(L,x)\longrightarrow\operatorname{Diff}_{\mathrm{germ}}(T,x),
	\qquad\rho([\gamma])=h_\gamma^{-1}.
\end{equation}
确实，\(\rho([\alpha*\beta])=h_\alpha^{-1}\circ h_\beta^{-1}=\rho([\alpha])\circ\rho([\beta])\)。两种约定得到同一个像群；此后构造群作用时采用 \(\rho\)。

\begin{proposition}\label{g2:prop:holonomy-conjugacy}
	改变横截面或在同一张叶内改变基点，所得返回芽群由横向运输共轭。因此芽群的同构类型、平凡性及有限性只依赖于叶状结构和所选的叶。
\end{proposition}
\begin{proof}
	先在同一点 \(x\) 把 \(T\) 换成 \(T'\)，以同一张叶状图内的叶片对应记作 \(c:(T,x)\to(T',x)\)。计算一条闭路径时，在起点先从 \(T'\) 返回 \(T\)，在终点再从 \(T\) 转到 \(T'\)，故新返回芽为 \(c\circ h_\gamma\circ c^{-1}\)。

	若基点从 \(x\) 移到 \(y\)，选叶内路径 \(\delta:x\to y\)，其运输为 \(h_\delta:(T_x,x)\to(T_y,y)\)。对基于 \(y\) 的闭路径 \(\gamma\)，有
	\[
		h_{\delta*\gamma*\delta^{-1}}
		=h_\delta^{-1}\circ h_\gamma\circ h_\delta.
	\]
	路径 \(\gamma\mapsto\delta*\gamma*\delta^{-1}\) 给出基本群的基点变换，故这个等式恰好给出两返回群的共轭。更换 \(\delta\) 可能改变共轭映射，却不改变“互为共轭”这一结论。
\end{proof}

在横截面上取坐标，使 \(x\) 对应于 \(0\in\mathbb R^q\)。对返回芽求微分得到
\[
	[h]_x\longmapsto\mathrm D_0h\in\operatorname{GL}(q,\mathbb R).
\]
其像称为\term[xianxinghengjiemianfanhuinqun]{线性横截面返回群}[linear holonomy group]。横截面的切空间自然映同构到法空间 \(\nu_x\mathcal F=\mathrm T_xM/\mathrm T_x\mathcal F\)，所以这项线性数据也可直接放在法空间中。换横截面时，链式法则把共轭芽变为共轭线性映射；它因而给出更容易计算的共轭不变量。

\begin{example}\label{g2:exm:parabolic-holonomy}
	在圆柱 \((\mathbb R/\mathbb Z)\times\mathbb R\) 上考虑无零点向量场
	\[
		X=\frac{\partial}{\partial\theta}+r^2\frac{\partial}{\partial r}.
	\]
	中心圆 \(L=\{r=0\}\) 是一张紧叶。取横截面 \(\theta=0\)，沿 \(\theta\) 增加一周的方向跟随邻近积分曲线，横向方程为 \(\mathrm dr/\mathrm d\theta=r^2\)。由初值 \(r(0)=s\) 解得
	\[
		r(\theta)=\frac{s}{1-\theta s},\qquad h(s)=\frac{s}{1-s}.
	\]
	对充分小的 \(s\)，解在 \(0\leq\theta\leq1\) 上存在，故这个公式确实定义返回芽。重复绕行得到
	\[
		h^n(s)=\frac{s}{1-ns},\qquad n\in\mathbb Z.
	\]
	这些芽两两不同，而 \(\mathrm D_0h^n=1\)。因此 \(H(L)\cong\mathbb Z\)，线性返回群却是平凡群。一阶导数没有看到二次项造成的累积偏移。
\end{example}

一个芽要求两个映射在某个邻域内完全一致；它比有限阶 Taylor 数据包含更多信息。\cref{g2:exm:parabolic-holonomy}说明，判断叶的稳定性时只检查线性返回群还不够。下一小节的有限性条件施加在完整的芽群上。

\subsection{有限 holonomy 与 Reeb 局部稳定}
\label{g2:subsec:D01:03}

有限芽群有一项很实用的特点：可以把各元素原先不同的定义域缩成一个共同邻域，再通过换坐标把整个群变成有限线性群。下面的平均公式同时说明这两件事。

\begin{lemma}\label{g2:lem:finite-germ-linearization}
	设 \(H\) 是固定 \(0\in\mathbb R^q\) 的光滑局部微分同胚芽的有限群。则可选取 \(0\) 的邻域及其光滑坐标，使 \(H\) 在该邻域上的代表构成实际群作用，并且在新坐标中作用为一组正交线性变换。微分映射 \(H\to\operatorname{GL}(q,\mathbb R)\) 是单射。
\end{lemma}
\begin{proof}
	对每个 \(g\in H\) 选代表 \(f_g\)，并取 \(f_e=\operatorname{id}\)。群有限，所以有限个等式 \(f_g\circ f_h=f_{gh}\) 都能在足够小的共同邻域内成立，所有需要的复合也都能在那里定义。令 \(A_g=\mathrm D_0f_g\)，则 \(A_{gh}=A_gA_h\)。在共同邻域上定义
	\[
		\Psi(y)=\frac1{|H|}\sum_{g\in H}A_g^{-1}f_g(y).
	\]
	有 \(\Psi(0)=0\) 及 \(\mathrm D_0\Psi=I\)。对每个 \(h\in H\)，把求和指标换成 \(k=gh\)，得到
	\[
		\Psi(f_h(y))
		=\frac1{|H|}\sum_{g\in H}A_g^{-1}f_{gh}(y)
		=A_h\Psi(y).
	\]
	这里先把邻域再缩小，使 \(y\) 和各 \(f_h(y)\) 都落在 \(\Psi\) 的定义域内。逆函数定理使 \(\Psi\) 在一个更小邻域上为微分同胚。

	在 \(\mathbb R^q\) 上取任意内积 \(\langle\ ,\ \rangle_0\)，再令
	\[
		\langle u,v\rangle_H=\frac1{|H|}\sum_{g\in H}\langle A_gu,A_gv\rangle_0.
	\]
	这仍是正定内积，且被全部 \(A_h\) 保持。选此内积下足够小的开球 \(B\)，使 \(\Psi^{-1}(B)\) 包含在上述各等式有效的邻域中。共轭关系说明 \(f_h\) 将 \(\Psi^{-1}(B)\) 映到自身，逆映射是 \(f_{h^{-1}}\)；这样便得到实际的群作用。在新坐标中它就是 \(A_h|_B\)。若 \(A_h=I\)，则 \(f_h\) 在这个邻域内等于恒等映射，故 \(h=e\) 作为芽也成立。
\end{proof}

现在设 \(L\) 是紧叶，\(x\in L\)，其返回芽群为有限群 \(H\)。记式\eqref{g2:eq:holonomy-left-representation}的核为 \(K\)。它是 \(\pi_1(L,x)\) 的正规子群，对应一张连通正规覆盖
\[
	\widehat L\longrightarrow L,\qquad
	\pi_1(L,x)/K\cong H.
\]
称 \(\widehat L\) 为\term[hengjiemianfanhuinfugai]{横截面返回覆盖}[holonomy cover]。这张覆盖的点可以表示为从 \(x\) 出发的叶内路径：同终点的两条路径，在它们拼成的闭路径具有恒等返回芽时识别。覆盖的层数是 \(|H|\)，所以 \(\widehat L\) 也是紧流形。

\ProseParagraphBreak

这里还需固定覆盖变换和横向作用的配合。若 \(g=\rho([\gamma])\in H\)，就在路径代表上定义 \(g[\alpha]=[\gamma*\alpha]\)。核 \(K\) 的正规性保证定义无关选择；直接计算表明这是左覆盖变换作用。横向同时用 \(g=h_\gamma^{-1}\) 作用。因此
\[
	h_{\gamma*\alpha}(h_\gamma^{-1}(v))=h_\alpha(v).
\]
这个等式说明两处作用应当同时使用：更换到同一点的路径代表时，横向起点也作相应变换，最终运输到的点才保持不变。

\begin{proposition}\label{g2:prop:finite-holonomy-model}
	设有限群 \(H\) 在连通紧流形 \(\widehat L\) 上自由光滑左作用，并在 \(\mathbb R^q\) 上忠实正交作用。取以 \(0\) 为心的开球 \(B\)，令
	\begin{equation}\label{g2:eq:finite-holonomy-model}
		E=(\widehat L\times B)/H,\qquad
		g\cdot(z,v)=(gz,gv).
	\end{equation}
	则 \(E\) 是光滑流形，乘积叶片给出叶状结构，中心叶是 \(L=\widehat L/H\)。通过 \([z,v]\) 的叶微分同胚于 \(\widehat L/H_v\)，其中 \(H_v=\{g:gv=v\}\)，其到 \(L\) 的投影是 \([H:H_v]\) 层覆盖。该叶的返回芽群同构于 \(H_v\)，叶空间同胚于 \(B/H\)。
\end{proposition}
\begin{proof}
	对角作用自由，因为任何固定 \((z,v)\) 的元素都固定 \(z\)。有限自由作用的商是流形：在 \(\widehat L\) 中取与其非平凡平移互不相交的小开集 \(V\)，则 \(V\times B\) 到商空间给出局部坐标；不同坐标之间按某个 \(g\) 改变横向坐标，仍保持乘积叶片。这也构造了 \(E\to L\) 的局部丛图。

	固定 \(v\)，映射 \(z\mapsto[z,v]\) 将 \(\widehat L\) 的像给成一张叶。两个点 \((z,v)\)、\((z',v)\) 在商中相同，当且仅当存在 \(g\in H_v\) 使 \(z'=gz\)，故这张叶为 \(\widehat L/H_v\)。它的投影到 \(\widehat L/H\) 是覆盖，纤维大小为 \(|H|/|H_v|\)。每条叶都由这种方式得到；不同的 \(v,w\) 给出同一张叶，当且仅当 \(w\in Hv\)。

	从 \((z,v)\) 出发，把商中一条闭叶内路径提升到 \(\widehat L\times B\)，提升的第二坐标保持为 \(v\)，终点为 \((gz,v)\)，其中 \(g\in H_v\)。转回起点处的横截面时，返回芽是 \(w\mapsto g^{-1}w\) 在 \(v\) 处的芽。反之，\(\widehat L\) 连通，任意 \(g\in H_v\) 都能用从 \(z\) 到 \(gz\) 的路径得到。线性作用忠实；若 \(g\) 在 \(v\) 的某邻域恒等，它就在整个向量空间恒等，因而 \(g=e\)。所以返回群确与 \(H_v\) 同构。

	最后，\([z,v]\mapsto Hv\) 是连续的开映射，其纤维恰为叶。开性可先在乘积空间用投影的开性验证，再通过有限群的开商映射下降。因此它诱导叶空间与 \(B/H\) 的同胚，而不只是集合上的双射。
\end{proof}

\begin{example}\label{g2:exm:mobius-finite-holonomy}
	取 \(\widehat L=\mathbb R/2\mathbb Z\)、\(H=\mathbb Z/2\mathbb Z\)、\(B=(-\varepsilon,\varepsilon)\)，非平凡元素作用为
	\[
		(\theta,s)\longmapsto(\theta+1,-s).
	\]
	式\eqref{g2:eq:finite-holonomy-model}给出开 Möbius 带上的圆叶状结构。中心圆的返回芽为 \(s\mapsto-s\)，返回群是二阶群；非中心圆的稳定子平凡，返回群也平凡，并且到中心圆的投影是二层覆盖。叶空间是 \([0,\varepsilon)\)，商坐标可取 \(|s|\)。每张叶都是圆，却不能把这个叶状邻域写成保持中心叶的 \(S^1\times(-\varepsilon,\varepsilon)\) 乘积，因为乘积的中心叶返回群平凡，而返回群在叶状微分同胚下保持共轭类型。
\end{example}

一般叶状结构中，紧叶附近是否一定有上述模型，是 Reeb 局部稳定定理所回答的问题。名称中的 Reeb 指法国数学家 Georges Reeb。\footnote{Georges Reeb（1920--1993），法国数学家，研究微分拓扑、叶状结构与动力系统。}

\begin{theorem}[Reeb 局部稳定定理]\label{g2:thm:reeb-local-stability}
	设 \(\mathcal F\) 是无边界光滑流形 \(M\) 上的光滑正则叶状结构，\(L\) 是紧叶，完整的横截面返回映射芽群 \(H\) 有限。则 \(L\) 有任意小的开\term[baoheneiyu]{饱和邻域}[saturated neighborhood]，即由完整叶组成的开邻域 \(U\)，使 \((U,\mathcal F|_U)\) 与式\eqref{g2:eq:finite-holonomy-model}的叶状模型微分同胚，并将 \(L\) 对应到中心叶。这里 \(\widehat L\) 是横截面返回覆盖，\(H\) 在横向开球上的作用取自其线性化。特别地，\(U\) 中的每张叶都紧，且具有有限返回芽群。
\end{theorem}
\begin{proof}
	本证明使用如下邻域存在性结果：若一张叶 \(S\) 嵌入且具有有限返回芽群，则 \(S\) 在原流形中的某开邻域，与其模型 \((\widehat S\times\nu_x\mathcal F)/H\) 中零叶的某开邻域叶状微分同胚，并在零叶上为给定识别。这一结果允许 \(S\) 非紧，给出的邻域此时不要求饱和；其一般证明见\cite[第2节，引理1，222--226页]{CrainicMarcut2015ReebStability}。下面完整说明紧性怎样把这项邻域结论变成定理所需的饱和模型。

	紧叶 \(L\) 的包含是从紧流形到 Hausdorff 流形的单射浸入，故为嵌入，可以应用上述结果。记所得叶状微分同胚为 \(F:W\to V\)，其中 \(W\) 是模型零叶的开邻域，\(V\subset M\) 是开邻域。有限芽群由\cref{g2:lem:finite-germ-linearization}线性化，并选好不变内积。\(\widehat L\) 是紧流形，故 \(W\) 在 \(\widehat L\times\nu_x\mathcal F\) 中的原像包含 \(\widehat L\times B\)，其中 \(B\) 是充分小的不变开球。这一步可从 \(\widehat L\times\{0\}\) 的有限乘积邻域子覆盖直接得到。

	令 \(U=F((\widehat L\times B)/H)\)。这是 \(M\) 的开子集。由\cref{g2:prop:finite-holonomy-model}，模型中的每张叶 \(N=\widehat L/H_v\) 都紧。其像包含于原叶状结构的一张叶 \(L'\)；叶状局部坐标说明 \(F|_N:N\to L'\) 对叶自身的拓扑光滑，并且局部微分同胚，所以像在 \(L'\) 中开。这个像又是紧集，而 \(L'\) 是 Hausdorff 流形，故像在 \(L'\) 中闭。叶 \(L'\) 连通，像非空，因此像就是整张 \(L'\)。这就证明 \(U\) 饱和；邻近叶的紧性和有限返回群也由模型命题得到。

	若另给 \(L\) 的任意环境开邻域 \(O\)，在模型中把 \(W\) 换成 \(W\cap F^{-1}(O\cap V)\)，仍是紧零叶的开邻域。重复刚才有限子覆盖的选择，将 \(B\) 缩小便得到 \(U\subset O\)。故这种饱和邻域确实可以任意小。
\end{proof}

\begin{corollary}\label{g2:cor:trivial-holonomy-product}
	若紧叶 \(L\) 的返回芽群平凡，则存在任意小的饱和邻域，叶状微分同胚于 \(L\times B^q\)，各叶对应于 \(L\times\{v\}\)。若余维为 \(1\)，且叶状结构在 \(L\) 附近横向可定向，则有限返回芽群必平凡，因而同样得到这种乘积邻域。
\end{corollary}
\begin{proof}
	当 \(H\) 平凡时，\(K=\pi_1(L,x)\)，相应的覆盖就是 \(\widehat L=L\)，所以模型直接化为乘积。在余维一且横向可定向时，各返回芽保持横向定向，其导数为正数。有限群中的每个导数都有有限乘法阶，正实数中具有有限阶的元素只有 \(1\)。\cref{g2:lem:finite-germ-linearization}的单射性于是给出 \(H=\{e\}\)。
\end{proof}

定理的局部范围是整张紧叶的邻域，因此比单点处的叶状坐标强得多。紧性保证能统一选择横向宽度，并把模型中的紧叶识别为环境中的完整叶；有限芽群保证全部返回可以同时控制。\cref{g2:exm:parabolic-holonomy}虽有紧叶和平凡线性返回群，却有无限的完整芽群，因而不能从它推出上述饱和乘积结论。下一节将直接观察这种局部返回行为如何影响叶空间。

\section{叶空间与整体动力学}
\label{g2:sec:D02}

横截面上的返回映射描述邻近叶怎样绕行；把整条叶压成一点，得到的商空间则记录不同叶能否分开。局部坐标盒总像一叠平行切片，商空间却可能是圆周，也可能连不同点都无法用开集区分。本节从环面上可以显式求解的流出发，再构造一个紧叶与非紧叶共存的实心环面。

\subsection{环面线性叶状结构}
\label{g2:subsec:D02:01}

令 \(T^2=\mathbb R^2/\mathbb Z^2\)，在上面取处处非零的向量场
\[
	X_\alpha=\partial_x+\alpha\partial_y,\qquad \alpha\in\mathbb R.
\]
它的叶由轨道 \(t\mapsto[x_0+t,y_0+\alpha t]\) 给出。这也是闭一形式 \(\mathrm dy-\alpha\mathrm dx\) 的核所确定的叶状结构。第一卷第13.5节已证明有理斜率的叶闭合、无理斜率的叶稠密；这里恢复计算，并把结果与叶空间及 holonomy 联系起来。

\begin{definition}\label{g2:def:foliation-leaf-space}
	设 \(\mathcal F\) 是光滑流形 \(M\) 上的叶状结构。将同一叶内的点识别所得集合 \(M/\mathcal F\)，赋予商映射 \(q:M\to M/\mathcal F\) 所确定的商拓扑，称为\term[yekongjian]{叶空间}[leaf space]。由整条叶组成的子集称为\term[baoheji]{饱和集}[saturated set]；商中的开集恰对应 \(M\) 中的饱和开集。
\end{definition}

\begin{proposition}\label{g2:prop:torus-foliation-quotient}
	对上述 \(T^2\) 上的光滑线性叶状结构，有以下结论。
	\begin{enumerate}
		\item 若 \(\alpha=p/q\)，其中 \(p,q\) 互素且 \(q>0\)，则每条叶为嵌入圆周，最小正流周期为 \(q\)。映射 \([x,y]\mapsto[qy-px]\) 将叶空间微分同胚地识别为圆周，并给出以这些叶为纤维的光滑纤维丛。
		\item 若 \(\alpha\) 无理，则每条叶自身微分同胚于 \(\mathbb R\)，其像在 \(T^2\) 中稠密。叶空间有不止一个点，但其开集只有空集与全空间。
		\item 两种情形中，每条叶的 holonomy 群都平凡。
	\end{enumerate}
\end{proposition}
\begin{proof}
	轨道在时刻 \(t\) 回到初始点，当且仅当 \(t\in\mathbb Z\) 且 \(\alpha t\in\mathbb Z\)。有理情形的返回时间正是 \(q\mathbb Z\)，所以轨道是紧圆周的单射浸入，因而是嵌入。令 \(P([x,y])=[qy-px]\)。其微分为处处非零的 \(q\mathrm dy-p\mathrm dx\)，且沿叶为零。若两点的 \(P\) 值相同，则对实代表有 \(q\Delta y-p\Delta x=k\in\mathbb Z\)。由互素性选整数 \(m,n\) 使 \(qn-pm=k\)，于是 \((\Delta x-m,\Delta y-n)\) 沿方向 \((q,p)\)。这说明它们恰在同一条叶上。

	映射 \(P\) 是满淹没，源紧，故由\cref{g2:thm:ehresmann}成为光滑纤维丛。其局部投影坐标给商空间赋予圆周的光滑结构。也可以先看拓扑：它诱导从紧叶空间到 Hausdorff 圆周的连续双射，因而是同胚。

	无理情形没有非零返回时间，轨道在叶自身的流形结构下以 \(\mathbb R\) 参数化。为恢复稠密性的理由，将 \(0,\alpha,\ldots,N\alpha\) 的小数部分放入 \(N\) 个等长区间；两数之差及其适当正负号给出
	\[
		0<\delta=k\alpha-m<1/N
	\]
	其中 \(k,m\) 为整数。模 \(1\) 的整数倍 \(\alpha\) 因而包含步长小于 \(1/N\) 的一组网点，随 \(N\) 增大可任意逼近圆周上的点。对环面中任意开矩形，先选其中第一坐标 \(u\)，再考察时刻 \(u-x_0+n\)；第二坐标的取值是一个无理旋转轨道，必能进入该矩形。故每条叶稠密。

	每条叶在截面 \(x=0\) 上只交于可数多个点，而该截面不可数，所以不止一条叶。若 \(U\) 是非空饱和开集，每条稠密叶都与 \(U\) 相交，饱和性又使它整条包含于 \(U\)，所以 \(U=T^2\)。这给出第二项的商拓扑。

	截面 \(x=0\) 的首次正返回映射为 \(R_\alpha([y])=[y+\alpha]\)。有理情形中，沿完整闭叶一周对应 \(R_\alpha^q=\mathrm{id}\)，故 holonomy 平凡。无理情形的叶为 \(\mathbb R\)，其全部叶内环路可缩；由\cref{g2:prop:holonomy-transport}的同伦不变性，holonomy 也平凡。
\end{proof}

无理斜率的首次返回映射确实在不断移动截面上的点，但这些运动没有沿叶内环路回到原来的叶内基点。这说明返回轨道的稠密性与环路芽群是两种信息。叶的自身拓扑也必须保留：能选 \(|n_j|\to\infty\) 使 \(n_j\alpha\bmod1\to0\)，于是轨道像在环境中趋回出发点，而叶内参数并未趋回零。因此无理叶到环面的单射浸入不是嵌入。

\ProseParagraphBreak

这个例子还显示一种整体不稳定性。若 \(\alpha_j\to\alpha\)，对应线分布在 \(C^\infty\) 意义下收敛；常系数的各阶导数没有额外误差。可是有理数与无理数都在实数中稠密，任意小的斜率改变都可能把全部圆周叶换成稠密直线叶。微分同胚保持叶的紧性，因而这两类叶状结构不可能通过保持叶的微分同胚等同。

\subsection{Reeb 分量及非分离叶空间}
\label{g2:subsec:D02:02}

现在希望让一张紧叶成为其他叶的极限。取实心环面
\[
	V=D^2\times(\mathbb R/\mathbb Z),\qquad
	D^2=\{(x,y):x^2+y^2\le1\},\qquad \rho=x^2+y^2,
\]
纵向圆坐标记为 \(t\)。角坐标 \(t\) 在实值意义下只局部定义，但 \(\mathrm dt\) 是全局一形式。考虑
\begin{equation}\label{g2:eq:reeb-solid-torus-form}
	\alpha=\mathrm d\rho+(1-\rho)\mathrm dt.
\end{equation}
在中心圆 \(\rho=0\) 上，它等于 \(\mathrm dt\)；在 \(\rho>0\) 处，径向分量 \(\mathrm d\rho\) 非零。因此 \(\alpha\) 处处非零，其核是一张光滑二维分布。

\ProseParagraphBreak

直接计算得到 \(\mathrm d\alpha=-\mathrm d\rho\wedge\mathrm dt\)，从而 \(\alpha\wedge\mathrm d\alpha=0\)。这个等式说明核分布对合：若 \(X,Y\) 都在核中，外微分公式给出 \(\alpha([X,Y])=-\mathrm d\alpha(X,Y)=0\)。最后一个等式也可由取 \(\alpha(Z)=1\) 后在 \((Z,X,Y)\) 上计算 \(\alpha\wedge\mathrm d\alpha\) 得到。第一卷第13.4节的 Frobenius 定理于是适用。在边界附近，将同一公式延伸到 \(\rho>1\) 的小领圈，仍为非零可积形式，所以边界处的叶状结构也由通常无边流形上的局部定理限制而来。

\begin{proposition}\label{g2:prop:reeb-component-leaves}
	式\eqref{g2:eq:reeb-solid-torus-form}在实心环面上确定的叶状结构有一张边界环面叶 \(L_\partial=\partial D^2\times S^1\)。全部内部叶以 \(c\in\mathbb R/\mathbb Z\) 编号，表达为
	\begin{equation}\label{g2:eq:reeb-plane-leaves}
		L_c=\{(z,t):|z|<1,\ t=c+\log(1-|z|^2)\pmod1\}.
	\end{equation}
	每条 \(L_c\) 是嵌入开圆盘，且在 \(V\) 中的闭包为 \(L_c\cup L_\partial\)。沿边界叶的子午圈，单侧 holonomy 为恒等；沿正向纵圈，以向内坐标 \(s=1-\rho\) 表示的返回芽为 \(s\mapsto e s\)。
\end{proposition}
\begin{proof}
	在 \(\rho=1\) 上，\(\alpha=\mathrm d\rho\)，其核恰是边界切空间，所以整个连通边界为一张叶。在内部，令
	\[
		F(z,t)=t-\log(1-|z|^2)\pmod1.
	\]
	则 \(\mathrm dF=\alpha/(1-\rho)\)。每个纤维由式\eqref{g2:eq:reeb-plane-leaves}给出，是开圆盘上的光滑图像；投影到 \(z\) 是其光滑逆，所以它连通且嵌入。由于 \(F\) 是淹没，这些纤维正是内部的最大叶。

	固定边界点 \((e^{i\varphi},t_0)\)。选实代表 \(c,t_0\)，令
	\[
		r_k^2=1-\exp(t_0-c-k).
	\]
	当 \(k\) 充分大时，\(0<r_k<1\)，并且 \((r_ke^{i\varphi},t_0)\in L_c\)，趋向所选边界点。这证明边界全部包含于闭包；内部的其余纤维由连续的 \(F\) 分开，不会增加其他内部极限点。

	沿子午圈只改变极角 \(\varphi\)，\(\rho,t\) 均不变，故横向点原样返回。沿纵圈令 \(t\) 从 \(0\) 增至 \(1\)，保持极角不变；要切于叶，必须满足
	\[
		\frac{\mathrm d\rho}{\mathrm dt}=-(1-\rho),\qquad
		\frac{\mathrm ds}{\mathrm dt}=s.
	\]
	因此 \(s(1)=e s(0)\)。将初始 \(s\) 限于足够小邻域，就能走完整个纵圈。反向纵圈则给出收缩芽 \(s\mapsto e^{-1}s\)。其不同整数次幂互异，故边界叶的 holonomy 群无限。
\end{proof}

这个模型称为一个\term[Reeb fenliang]{Reeb 分量}[Reeb component]。边界叶虽紧，其邻近叶却是向它卷绕的非紧平面；上一节局部稳定中的有限 holonomy 条件在这里恰好不成立。图中的曲线只画二维叶在固定极角截面内的交线，完整的叶还包含极角方向。

\begin{figure}[htbp]
	\centering
	% !TeX root = ../Book2.tex
\begin{tikzpicture}[>=stealth,x=6cm,y=1.25cm,every node/.style={font=\small}]
	\draw[->] (-.02,0)--(1.12,0) node[right] {\(\rho\)};
	\draw[->] (0,-2.8)--(0,.25) node[above] {\(t\)};
	\draw[very thick,blue!70!black] (1,-2.65)--(1,.12);
	\node[above] at (1,.12) {\(\rho=1\)};
	\foreach \c in {.15,.65,1.15} {
		\draw[thick,teal!70!black,->,domain=-.05:-2.55,samples=55,variable=\v] plot ({1-exp(\v-\c)},\v);
	}
	\draw[dashed,gray] (0,-1)--(1,-1);
	\draw[dashed,gray] (0,-2)--(1,-2);
	\node[left] at (0,-1) {\(-1\)};
	\node[left] at (0,-2) {\(-2\)};
	\node[right,align=left] at (1.05,-1.35) {边界叶\(L_\partial\)};
	\node[align=center] at (.3,-2.35) {内部叶的截线};
\end{tikzpicture}

	\caption[Reeb 分量的叶向边界聚集]{Reeb 分量的叶向边界聚集。纵向圆周先展开为实坐标 \(t\)，内部叶的截线满足 \(\rho=1-e^{t-c}\)；右侧直线属于紧边界叶。}
	\label{g2:fig:reeb-component}
\end{figure}

\begin{proposition}\label{g2:prop:reeb-leaf-space}
	上述 Reeb 分量的叶空间非 \(T_1\)，从而非 Hausdorff。更具体地，包含边界叶类的唯一开邻域是整个叶空间。
\end{proposition}
\begin{proof}
	设饱和开集 \(U\subset V\) 包含 \(L_\partial\)。每条 \(L_c\) 都有趋于边界的点列，故它与 \(U\) 相交；饱和性便给出 \(L_c\subset U\)。于是 \(U=V\)。另一方面，\(L_c\) 在 \(V\) 中不闭，商中单点 \(\{[L_c]\}\) 因而不闭：商拓扑下集合闭当且仅当它的逆像闭。这正是否定 \(T_1\) 条件。
\end{proof}

内部叶之间仍能用内部函数 \(F\) 区分，所以内部叶空间是一个圆周。加进边界叶后，新增的这个点却只有全空间这一开邻域。由此可以看见，非分离性来自具体的叶聚集方式，而不是局部叶片坐标失效。

\subsection{Reeb 整体稳定与圆周纤维丛}
\label{g2:subsec:D02:03}

局部 Reeb 稳定给一张紧叶附近的饱和模型。要使整张流形成为纤维丛，还需要保证这些模型覆盖全部叶，并使所得横向参数能够组成一个分离的空间。先在所有叶均已有良好邻域的情形，把这一步证明出来。

\begin{definition}\label{g2:def:transverse-orientation}
	若光滑叶状结构的法丛 \(\nu\mathcal F=\mathrm TM/\mathrm T\mathcal F\) 可定向，就称 \(\mathcal F\) \term[hengxiang kedingxiang]{横向可定向}[transversely orientable]。余维一时，选定横向定向等价于给局部横向坐标选择一致的正方向；相容坐标的横向导数因而为正。
\end{definition}

\begin{theorem}\label{g2:thm:compact-trivial-holonomy-fibration}
	设 \(M\) 为紧、连通、无边界光滑流形，\(\mathcal F\) 为横向可定向的光滑余维一叶状结构。若每条叶都紧且 holonomy 群平凡，则 \(M/\mathcal F\) 是光滑圆周，商映射是以叶为纤维的光滑纤维丛。特别地，所有叶彼此微分同胚。
\end{theorem}
\begin{proof}
	\cref{g2:cor:trivial-holonomy-product}给每条叶 \(L\) 一个饱和邻域
	\[
		\Psi:L\times(-\varepsilon,\varepsilon)\xrightarrow{\;\cong\;}U,
	\]
	其中每个切片是完整的一条叶。故 \(U/\mathcal F\) 是一个开区间。先证明这些局部参数所在的商是 Hausdorff。给两条不同叶 \(L,K\)，若 \(K\subset U\)，它对应某个非零参数；若 \(K\cap U=\varnothing\)，则在 \(U\) 外。两种情形都可选足够小的 \(\delta>0\)，使
	\[
		C=\Psi(L\times[-\delta,\delta])
	\]
	不与 \(K\) 相交。\(L\) 紧，故 \(C\) 是紧集，在 Hausdorff 的 \(M\) 中闭；它还是完整叶的并。因此 \(\Psi(L\times(-\delta,\delta))\) 与 \(M\setminus C\) 是分别包含两条叶的不交饱和开集。这就分开了两个叶类。

	局部横向坐标的过渡光滑：在重叠区域中，两个坐标函数都沿叶常值，取一条小横截面后，它们的关系就是两张横向坐标间的光滑微分同胚。由 \(M\) 紧，从上述饱和邻域取有限子覆盖，得到商空间的有限区间图册，也就得到第二可数性。故叶空间 \(B=M/\mathcal F\) 是紧、连通、Hausdorff 的光滑无边一维流形。给定的横向定向使这些区间图的过渡递增。

	为说明 \(B\) 为什么是圆周，在各区间图中选正向量场，用单位分解拼成处处非零的正向量场。紧性使其流完备；每条轨道因维数为一而是开集，轨道分割与连通性说明只有一条轨道。于是轨道映射 \(\mathbb R\to B\) 是满的局部微分同胚。返回时间组成离散闭子群；若只有零，它就是 \(\mathbb R\) 到 \(B\) 的微分同胚，与紧性矛盾。因此返回时间为某个 \(a\mathbb Z\)、\(a>0\)，得到 \(B\cong\mathbb R/a\mathbb Z\cong S^1\)。

	在每张 \(U\) 中，商映射就是投影 \(L\times(-\varepsilon,\varepsilon)\to(-\varepsilon,\varepsilon)\)，所以它是光滑满淹没。它又因 \(M\) 紧而固有，\cref{g2:thm:ehresmann}给出光滑纤维丛；沿连通基底的路径逐次使用局部平凡化，便知所有纤维微分同胚。
\end{proof}

这里“每条叶都紧”不能换成“存在一张紧叶且它的 holonomy 平凡”。例如在 \(T^2\) 上取
\[
	X=\partial_x+a(y)\partial_y,
\]
其中 \(a\) 是圆周上的非负光滑函数，在闭弧 \([0,1/3]\) 上恒为零，在余下开弧 \((1/3,1)\) 上严格为正，并在两个端点平坦。闭弧内部的每个水平圆都是 holonomy 平凡的紧叶，其邻域为乘积。然而在正值开弧内，方程 \(y'=a(y)\) 的轨道从一端渐近另一端；它不能在有限时间碰到零点，否则常微分方程解的唯一性会被破坏。因此这些叶是非紧的绕行轨道，全部叶并未构成同一类紧纤维。

\ProseParagraphBreak

一个经典的整体结果用更强的单叶假设排除这种混合。

\begin{theorem}[Reeb 整体稳定定理]\label{g2:thm:global-reeb-stability}
	设 \(M\) 为紧、连通、无边界光滑流形，\(\mathcal F\) 为横向可定向的光滑余维一叶状结构。若存在紧叶 \(L\) 且 \(\pi_1(L)\) 有限，则所有叶均微分同胚于 \(L\)，并存在以这些叶为纤维的光滑纤维丛 \(M\to S^1\)。
\end{theorem}

这个版本采用 Guillemin、Miranda、Pires\cite{GuilleminMirandaPires2011Foliations}中关于 Reeb 整体稳定的陈述。有限基本群先使 holonomy 有限，余维一的横向定向又使有限芽群平凡。整体定理还要证明紧叶能在整个连通流形中延续，这一步是此处引用的一般结果；一旦获得全部叶的紧性和平凡 holonomy，\cref{g2:thm:compact-trivial-holonomy-fibration}便解释了圆周纤维丛结论的来由。上面的混合例子中，紧叶都是圆周，其基本群为 \(\mathbb Z\)，恰不满足这个单叶假设。

\ProseParagraphBreak

即使叶空间是圆周，总空间仍可能带有全局扭转。对任意紧连通光滑流形 \(L\) 及其微分同胚 \(\varphi:L\to L\)，沿用式\eqref{g2:eq:mapping-torus-convention}的映射环面关系
\[
	M_\varphi=(L\times[0,1])/\bigl((x,1)\sim(\varphi(x),0)\bigr).
\]
投影到 \([t]\in S^1\) 是光滑纤维丛，其纤维叶的 holonomy 都平凡；沿基底绕一圈搬运整条纤维，却得到 \(\varphi\)。前一种沿叶内环路移动横截面上的邻近点，后一种沿横向基底环路移动纤维中的点。\cref{g2:ex:foliation-mapping-torus}用球面的一个映射说明：叶内芽群全部平凡，仍然可以留下非乘积的整体粘合。

\section{喷射、折叠与尖点}
\label{g2:sec:D03}

曲面向平面投影时，正则点附近可以直接把像坐标用作源坐标；沿着轮廓线，微分却失去一个方向。微分的秩能够找出这条线，但还不能说明附近的像怎样折回，也不能区分轮廓线上的普通点与尖点。为继续观察二阶、三阶变化，我们把一个点处有限阶的 Taylor 数据组成几何空间，再研究映射的这些数据如何穿过指定的秩层。

\subsection{喷射丛和接触阶}
\label{g2:subsec:D03:01}

本节中，\(M^n,N^p\) 都是有限维、Hausdorff、第二可数的光滑无边流形，映射取 \(C^\infty\)。先固定 \(a\in M\) 与 \(b\in N\)。我们关心的是 \(a\) 的某个邻域到 \(N\) 的映射，且要求它把 \(a\) 送到 \(b\)；为讨论点附近的行为，定义域可以随时缩小。

\begin{definition}\label{g2:def:map-jet}
	设非负整数 \(k\)，两张局部光滑映射 \(f,g\) 满足 \(f(a)=g(a)=b\)。若在源、靶坐标中，它们在 \(a\) 处的全部至 \(k\) 阶偏导数相同，就说它们在 \(a\) 处有相同的 \(k\) 阶\term[pen she]{喷射}[jet]，也称有 \(k\) 阶\term[jiechu]{接触}[contact]，记作
	\[
		j_a^k f=j_a^k g.
	\]
	这里包含零阶条件，即两张映射在 \(a\) 处取值相同。固定源点 \(a\) 与靶点 \(b\) 的等价类集合记为 \(J^k(M,N)_{a,b}\)，所有这些集合的不交并记为 \(J^k(M,N)\)。
\end{definition}

这个定义必须与坐标选择无关。若将坐标表达换成 \(\psi\circ f\circ\varphi^{-1}\)，逐次应用链式法则可知，它的一个至 \(k\) 阶偏导数，只用到 \(f\) 的至 \(k\) 阶偏导数，以及坐标变换相应阶数的导数。因此 \(f,g\) 的数据相同，换坐标后的数据也相同；再用逆坐标变换得到反向蕴含。偏导数相等本身又是等价关系，所以喷射确实由映射定义，而不依赖描述它的坐标。

\ProseParagraphBreak

多重指标 \(\alpha=(\alpha_1,\ldots,\alpha_n)\) 记录各方向的求导次数，\(|\alpha|=\sum_i\alpha_i\)、\(\alpha!=\prod_i\alpha_i!\)。在坐标原点取为 \(a\) 的情况下，一个喷射可写为有限 Taylor 多项式
\[
	\sum_{|\alpha|\leq k}\frac{\partial^\alpha f(a)}{\alpha!}h^\alpha.
\]
这里的多项式只表示指定点处的数据；其系数随源点变化时，还须满足来自同一张映射的相容关系。

\begin{proposition}\label{g2:prop:jet-bundle-charts}
	设 \(M^n,N^p\) 为有限维、Hausdorff、第二可数的光滑无边流形，\(k\geq0\)。喷射空间 \(J^k(M,N)\) 自然成为维数
	\[
		n+p\binom{n+k}{k}
	\]
	的光滑流形。源、靶投影
	\[
		\alpha(j_a^k f)=a,\qquad\beta(j_a^k f)=f(a)
	\]
	以及 \((\alpha,\beta):J^k(M,N)\to M\times N\) 都是淹没；最后一张映射是纤维维数为 \(p\bigl(\binom{n+k}{k}-1\bigr)\) 的光滑纤维丛。每张光滑映射 \(f:M\to N\) 给出光滑的\term[pen she yantuo]{喷射延拓}[jet extension]
	\[
		j^k f:M\longrightarrow J^k(M,N),\qquad a\longmapsto j_a^k f.
	\]
\end{proposition}
\begin{proof}
	取源坐标 \(x:U\to\mathbb R^n\) 与靶坐标 \(y:V\to\mathbb R^p\)。源点在 \(U\)、靶点在 \(V\) 的喷射由坐标
	\[
		\left(x(a),y(b),\left(\partial^\mu(y\circ f\circ x^{-1})(x(a))\right)_{1\leq|\mu|\leq k}\right)
	\]
	唯一确定。反过来，给定这些数，可用以 \(x(a)\) 为中心、系数为指定导数除以阶乘的多项式实现；把源邻域取得足够小，其像便仍在 \(y(V)\) 内。因此这张坐标图的像是
	\[
		x(U)\times y(V)\times\mathbb R^{p(\binom{n+k}{k}-1)}.
	\]
	非负整数多重指标 \(|\mu|\leq k\) 的数目为 \(\binom{n+k}{k}\)：添加一个记录 \(k-|\mu|\) 的变量，再数总和为 \(k\) 的 \(n+1\) 元组即可。

	坐标过渡由有限次链式法则给出，是各阶导数与坐标变换导数的光滑组合，逆过渡也具有同样性质。这些图于是定义光滑结构。不同源点或靶点可由底空间分离；相同底点而不同导数的数据可在同一张喷射坐标图中分离，故所得空间 Hausdorff。可数的源、靶图册与 Euclid 纤维的可数基给出第二可数性。三个投影在这些坐标中都为坐标投影，局部乘积还给出纤维丛结构。最后，\(j^k f\) 的坐标分量就是 \(x\)、\(f(x)\) 及其有限阶导数，均光滑。
\end{proof}

\symidx{\(J^k(M,N)\)：映射的 \(k\) 阶喷射丛}\symidx{\(j^k f\)：光滑映射的喷射延拓}
零阶喷射空间是 \(J^0(M,N)=M\times N\)，而 \(j^0f\) 就是图映射 \(a\mapsto(a,f(a))\)。一阶喷射则可规范地写为
\[
	j_a^1f=(a,f(a),\mathrm d f_a),\qquad
	J^1(M,N)_{a,b}=\operatorname{Hom}(\mathrm T_aM,\mathrm T_bN).
\]
所以 \(J^1(M,N)\to M\times N\) 是向量丛。更高阶的数据在一般坐标变换下会混入较低阶项，整个高阶喷射纤维并没有同样的规范线性结构。例如对实值函数换源坐标 \(x=\varphi(z)\)，二阶导数满足
\[
	\mathrm D^2(f\circ\varphi)_z(v,w)
	=\mathrm D^2f_x(\mathrm D\varphi_zv,\mathrm D\varphi_zw)
	+\mathrm Df_x\bigl(\mathrm D^2\varphi_z(v,w)\bigr).
\]
在临界点，最后一项消失，才恢复第17章 Hessian 的坐标无关双线性型。

\begin{example}\label{g2:exm:jet-versus-germ}
	在 \(\mathbb R\) 上，\(f(t)=t^2\) 与 \(g(t)=t^2+t^3\) 在零点有相同的二阶喷射，却没有相同的三阶喷射。函数芽相等则要求在零点的某个邻域内处处相等，比任何固定阶的接触都更强。

	即使所有有限阶喷射都相同，光滑函数芽也未必相同。令
	\[
		h(t)=\begin{cases}e^{-1/t^2},&t\ne0,\\0,&t=0.\end{cases}
	\]
	每次求导在 \(t\ne0\) 处都得到 \(1/t\) 的多项式乘 \(e^{-1/t^2}\)，且这类表达在 \(t\to0\) 时趋于零。逐次延拓可知 \(h\) 光滑、全部导数在零处为零，所以 \(j_0^kh=j_0^k0\) 对所有 \(k\) 成立；然而 \(h\) 在每个穿孔邻域为正。喷射负责有限阶信息，邻域内的映射等价还需进一步的判别。
\end{example}

还有一项相容性值得直接看见。\(j^k f\) 是源投影 \(\alpha\) 的截面，但任意光滑截面未必由映射产生。在 \(J^1(\mathbb R,\mathbb R)\) 的坐标 \((x,y,A)\) 中，截面 \(s(x)=(x,0,1)\) 就不能写成 \(j^1f\)：第二坐标迫使 \(f\equiv0\)，第三坐标却要求 \(f'\equiv1\)。因此，取得一般位置时必须实际扰动 \(f\)，使扰动后的喷射仍来自同一张映射。喷射丛的构造也可参见\cite[第II章第2节，定理2.7]{GolubitskyGuillemin1973}。

\subsection{秩分层及喷射横截}
\label{g2:subsec:D03:02}

一阶喷射保存了微分矩阵，所以先在矩阵空间中理解秩条件。秩不超过 \(r\) 可用若干子式为零表示，但这些方程在较低秩的点可能相互依赖。将“秩恰好等于 \(r\)”单独取出，才得到具有固定维数的光滑层。

\begin{proposition}\label{g2:prop:matrix-rank-strata}
	设 \(0\leq r\leq\min\{n,p\}\)。秩为 \(r\) 的实 \(p\times n\) 矩阵组成 \(\operatorname{Hom}(\mathbb R^n,\mathbb R^p)\) 的光滑嵌入子流形，其余维为
	\[
		(n-r)(p-r).
	\]
	因此
	\[
		\Sigma_r=\{(a,b,A)\in J^1(M,N):\operatorname{rank}A=r\}
	\]
	也是具有这一余维的光滑嵌入子流形。
\end{proposition}
\begin{proof}
	先设 \(r>0\)。在任一秩 \(r\) 的矩阵附近，交换行列后可保证左上方 \(r\times r\) 块 \(A\) 可逆。将矩阵写成
	\[
		\begin{pmatrix}A&B\\ C&D\end{pmatrix},
	\]
	其中 \(D\) 为 \((p-r)\times(n-r)\) 块。可逆的块行、列操作给出
	\[
		\begin{pmatrix}I&0\\-CA^{-1}&I\end{pmatrix}
		\begin{pmatrix}A&B\\ C&D\end{pmatrix}
		\begin{pmatrix}I&-A^{-1}B\\0&I\end{pmatrix}
		=\begin{pmatrix}A&0\\0&D-CA^{-1}B\end{pmatrix}.
	\]
	于是秩恰为 \(r\)，当且仅当 \(D=CA^{-1}B\)。这把秩层写成以 \(A,B,C\) 为自由变量的光滑图；独立方程的数目就是 \(D\) 的元素数 \((p-r)(n-r)\)。所有非零 \(r\) 阶子式对应的图覆盖秩层。\(r=0\) 时，该层只有零矩阵，余维为 \(np\)，也符合公式。满秩情形若一个补块维数为零，方程为空，秩层便是开集。

	在 \(J^1(M,N)\) 的坐标图内，只需将上述矩阵层与源、靶坐标相乘。换坐标把微分矩阵左右乘以可逆矩阵，保持秩，因而这些局部子流形拼成 \(\Sigma_r\)。
\end{proof}

这里的下标 \(r\) 一律表示秩，而非余秩。若 \(j^1f\pitchfork\Sigma_r\)，\cref{g2:thm:transverse-preimage}便说明
\[
	S_r(f)=\{a\in M:\operatorname{rank}\mathrm d f_a=r\}
	=(j^1f)^{-1}(\Sigma_r)
\]
在非空时是维数 \(n-(n-r)(p-r)\) 的子流形。若这个数为负，则该层为空。要让这个结论成为一般位置工具，剩下的问题是：在喷射必须来自映射的约束下，能否仍取得横截性？

\begin{lemma}\label{g2:lem:jet-polynomial-perturbation}
	设 \(M^n,N^p\) 为有限维、Hausdorff、第二可数的光滑无边流形，\(k\geq0\)，\(W\subset J^k(M,N)\) 为光滑嵌入子流形。设 \(K\subset U\subset M\)，其中 \(K\) 紧、\(U\) 为源坐标域，且光滑映射 \(f\) 将 \(U\) 映入一张靶坐标图。则存在任意小的紧支光滑扰动 \(g\)，在 \(M\setminus U\) 上等于 \(f\)，并使 \(j^kg\) 在 \(K\) 上横截于 \(W\)。扰动可在任意预先指定的有限阶 \(C^s\) 意义下任意小。
\end{lemma}
\begin{proof}
	以 \(x\) 表示源坐标、\(y\) 表示靶坐标，选取紧支于 \(U\) 的光滑截断函数 \(\rho\)，使它在 \(K\) 的开邻域 \(U_0\) 上为 \(1\)。令参数 \(a=(a_\mu)_{|\mu|\leq k}\)，每个 \(a_\mu\in\mathbb R^p\)，并记
	\[
		P_a(x)=\sum_{|\mu|\leq k}a_\mu x^\mu,\qquad
		y\circ f_a=y\circ f+\rho P_a
	\]
	于 \(U\)，支集外取 \(f_a=f\)。参数取得足够小时，紧支上的像仍在靶坐标域中，因而得到光滑总体族。

	在 \(U_0\) 上考虑 \(\Phi(a,x)=j_x^k f_a\)。固定 \(x\) 后，多项式的系数可以独立指定其在 \(x\) 处的全部至 \(k\) 阶 Taylor 系数：将多项式从以零为中心改写成以 \(x\) 为中心，给出可逆线性变换。因此参数方向的微分覆盖喷射坐标中全部靶值与导数方向；源点方向又覆盖源坐标。\(\Phi\) 是淹没，从而横截于 \(W\)。

	\cref{g2:thm:parametric-transversality}应用于有限维参数球与 \(U_0\) 的乘积，得到任意接近零的好参数 \(a\)，使 \(j^kf_a\pitchfork W\) 于 \(U_0\)。取 \(g=f_a\)。固定紧支集上的各阶导数连续依赖有限个参数，所以还能满足给定的有限阶误差界。
\end{proof}

普通横截扰动只需加入常向量；现在要独立改变一至 \(k\) 阶导数，就加入至 \(k\) 次的全部多项式项。这是喷射横截性仍可由有限维参数定理得到的原因。

\ProseParagraphBreak

下面使用第14.1节的强 Whitney \(C^\infty\) 拓扑：一个基本邻域控制某个有限阶数的坐标导数，并允许在局部有限个紧坐标块上指定不同的正容差。弱拓扑则每次只控制有限多个紧块；源紧时两者一致。强 Whitney \(C^\infty\) 映射空间具有 Baire 性，即可数个开稠密集的交仍稠密。这里采用\cite[第II章第3节，命题3.3，第44—46页]{GolubitskyGuillemin1973}的函数空间结论；其证明逐步缩小各阶喷射邻域，使局部各阶导数一致收敛到一张光滑映射。这个输入负责把可数个局部一般位置要求同时实现。

\begin{theorem}[Thom 喷射横截性]\label{g2:thm:jet-transversality}
	设 \(M,N\) 为有限维、Hausdorff、第二可数的光滑无边流形，\(k\geq0\)，\(W\subset J^k(M,N)\) 为光滑嵌入子流形。则
	\[
		\mathcal T_W=\{f\in C^\infty(M,N):j^k f\pitchfork W\}
	\]
	在强 Whitney \(C^\infty\) 拓扑中为稠密剩余集。对可数多个这样的喷射条件，满足全部条件的映射仍构成稠密剩余集。若 \(M\) 紧且 \(W\) 闭，\(\mathcal T_W\) 还是开集。
\end{theorem}
\begin{proof}
	将 \(W\) 覆盖为可数个紧块 \(C_\ell\) 的并，使每个块都紧含于一张 \(W\) 的适配图内，且其源、靶投影分别紧含于源图 \(U_\ell\) 与靶图 \(V_\ell\)。这是由 \(W\) 的第二可数性与局部紧性得到的。令 \(\mathcal U_\ell\) 为满足以下条件的映射集合：在所有 \(j_a^kf\in C_\ell\) 的点上，\(j^kf\) 与 \(W\) 横截。

	先证 \(\mathcal U_\ell\) 开。相关源点都属于紧集 \(\alpha(C_\ell)\)。若某个 \(f\in\mathcal U_\ell\) 没有保持此性质的紧集 \(C^{k+1}\) 邻域，可选 \(f_i\) 在该紧集附近至 \(k+1\) 阶趋于 \(f\)，以及坏点 \(a_i\)，使 \(j_{a_i}^kf_i\in C_\ell\)。取子列，喷射在 \(C_\ell\) 中收敛到 \(j_a^kf\)。\(W\) 的切空间随点连续，\(j^kf_i\) 的一阶微分也收敛；法方向微分不满射是闭条件，所以极限在 \(a\) 仍不横截，矛盾。因而一个弱 \(C^{k+1}\) 邻域已经足够，也就得到强拓扑中的开放性。

	再证稠密性。固定任意 \(f\) 与它的强邻域，选源截断函数 \(\rho\)，在 \(\alpha(C_\ell)\) 附近为 \(1\)、紧支于 \(U_\ell\)；选靶截断函数 \(\chi\)，在 \(\beta(C_\ell)\) 附近为 \(1\)、紧支于 \(V_\ell\)。在这两张坐标图的适用区域作扰动
	\[
		y\circ f_a=y\circ f+\rho(x)\chi(f(x))P_a(x),
	\]
	其他地方保持 \(f\) 不变。两个截断函数保证公式光滑拼接；把参数统一取小，还保证像留在靶图内。

	若 \(j_x^kf_a\in C_\ell\)，则 \(x\in\alpha(C_\ell)\) 且 \(f_a(x)\in\beta(C_\ell)\)。后一个紧集在 \(\chi=1\) 的开区域内有正的坐标余量，故参数足够小时，\(f(x)\) 也在该区域内。因此这类点附近两个截断函数都恒为 \(1\)，总体喷射与上一引理一样是淹没。对这一开区域应用参数横截性，取得任意小的参数，使扰动属于 \(\mathcal U_\ell\)。扰动紧支，故只碰到强邻域控制族中的有限多个紧块；所有有限阶误差可以同时取小，所得映射仍在原邻域内。这证明 \(\mathcal U_\ell\) 稠密。

	因为 \(W=\bigcup_\ell C_\ell\)，有 \(\mathcal T_W=\bigcap_\ell\mathcal U_\ell\)。前述 Baire 性给出其稠密性，而这个表达本身给出剩余性。可数多个喷射条件只需再对相应的开稠密族取一次可数交。最后，若 \(M\) 紧且 \(W\) 闭，\cref{g2:prop:transversality-open-on-compact}应用于 \(j^kf\)，给出 \(j^kf\) 的 \(C^1\) 开放性；它由 \(f\) 的 \(C^{k+1}\) 小扰动保持，故 \(\mathcal T_W\) 开。
\end{proof}

这一定理及其紧块证明也见\cite[第II章第4节，定理4.9及推论4.10，第54—56页]{GolubitskyGuillemin1973}。秩层 \(\Sigma_r\) 通常不是闭集，但它们是有限多个嵌入子流形，因而定理仍保证同时横截于所有秩层的映射稠密且为剩余集。非闭层的使用依靠证明中的紧块定位，不能直接套用“整个层闭”的开放性结论。

\begin{example}\label{g2:exm:surface-jet-rank}
	对于两曲面之间的映射，\(n=p=2\)。秩零层的余维为 \(4>2\)，故一阶喷射同时横截于各秩层的映射没有秩零点。秩一层余维为 \(1\)，其原像是一条光滑临界曲线，或为空。

	在秩一矩阵附近，行列式是这个层的局部定义函数。因此令 \(\lambda=\det\mathrm Df\)，此处 \(j^1f\pitchfork\Sigma_1\) 就等价于 \(\mathrm d\lambda\ne0\)。例如
	\[
		f(x,y)=(x,xy+y^3),\qquad\lambda(x,y)=x+3y^2
	\]
	满足这个条件。临界曲线为 \(x=-3y^2\)，在原点的切线是 \(\partial_y\) 方向，而 \(\ker\mathrm d f_0\) 也正是这个方向。相比之下，\((x,y)\mapsto(x,y^2)\) 的临界曲线为 \(y=0\)，微分的核 \(\partial_y\) 横穿它。前一个映射的尖点将于\cref{g2:exm:plane-cusp-model}展开计算；这两种行为已经说明，秩一层的横截性保证临界集光滑，却仍需二阶、三阶条件来辨认临界点的类型。
\end{example}

\subsection{折叠与尖点的局部标准形}
\label{g2:subsec:D03:03}

秩下降指出微分在哪些方向失去信息。接下来要问：沿这些方向，函数最先留下的是二次项还是三次项？第17章的 Morse 引理已经说明，非退化二次项能够吸收全部高阶余项。折叠是这一事实在有多个输出分量时的形式；尖点则出现在一个核方向的二次项消失、三次项仍保留下来的位置。

\ProseParagraphBreak

这里使用的坐标变换同时允许改变源与靶。两个映射芽 \(f:(\mathbb R^n,0)\to(\mathbb R^p,0)\)、\(g:(\mathbb R^n,0)\to(\mathbb R^p,0)\) 称为\term[zuoyoudengjia]{左右等价}[right-left equivalent]，是指存在固定原点的局部微分同胚 \(\phi,\psi\)，使 \(g=\psi\circ f\circ\phi^{-1}\)。也称 \(\mathcal A\)-等价。这样的等价保持局部纤维之间的对应；求标准形，就是在这类变换下尽量简化映射，而等式要求在一个实际邻域上成立。

\ProseParagraphBreak

设 \(f:M^n\to N^p\) 光滑，\(n\ge p\)，并在 \(a\) 处有秩 \(p-1\)。令
\[
	K_a=\ker\mathrm df_a,\qquad Q_a=\mathrm T_{f(a)}N/\operatorname{im}\mathrm df_a.
\]
于是 \(\dim K_a=n-p+1\)、\(\dim Q_a=1\)。在坐标中把二阶导数限制到核方向，再投影到余核，得到对称双线性映射
\begin{equation}\label{g2:eq:intrinsic-kernel-hessian}
	B_a:K_a\times K_a\longrightarrow Q_a,\qquad B_a(v,w)=[\mathrm D^2f_a(v,w)].
\end{equation}
这个映射具有坐标无关的含义。改变源坐标时，二阶链式法则多出的项落在 \(\operatorname{im}\mathrm df_a\) 中；改变靶坐标时，多出的项含 \(\mathrm df_a(v)\) 与 \(\mathrm df_a(w)\)，两者都为零。余下的变化只是 \(K_a\) 与 \(Q_a\) 中的线性换基。因此，可以在任选 \(Q_a\cong\mathbb R\) 后判断 \(B_a\) 是否非退化。

\begin{definition}\label{g2:def:fold-point}
	设 \(f:M^n\to N^p\) 是无边界光滑流形之间的光滑映射，\(n\ge p\)。若 \(\operatorname{rank}\mathrm df_a=p-1\)，且式\eqref{g2:eq:intrinsic-kernel-hessian}的核方向二次型非退化，就称 \(a\) 为\term[zhediedian]{折叠点}[fold point]。
\end{definition}

\begin{theorem}[折叠标准形]\label{g2:thm:fold-normal-form}
	设 \(f:M^n\to N^p\) 光滑，\(n\ge p\)，\(a\) 为折叠点。记 \(k=n-p+1\)。存在以 \(a\)、\(f(a)\) 为原点的坐标，使映射成为
	\begin{equation}\label{g2:eq:fold-normal-form}
		(u_1,\ldots,u_{p-1},z_1,\ldots,z_k)
		\longmapsto\left(u_1,\ldots,u_{p-1},-\sum_{i=1}^{\lambda}z_i^2+\sum_{i=\lambda+1}^{k}z_i^2\right).
	\end{equation}
	反过来，式\eqref{g2:eq:fold-normal-form}中的原点是折叠点。
\end{theorem}
\begin{proof}
	秩为 \(p-1\) 意味着可以选取 \(p-1\) 个输出坐标，其拉回的微分线性无关。把这些拉回补成源坐标，便把 \(f\) 写成
	\[
		f(u,z)=(u,h(u,z)),\qquad u\in\mathbb R^{p-1},\quad z\in\mathbb R^k.
	\]
	在原点有 \(h_z=0\)，而 \(h_{zz}\) 可逆。隐函数定理给出唯一的光滑函数 \(z=a(u)\)，使 \(h_z(u,a(u))=0\)。将 \(z\) 平移为 \(a(u)+w\)，再在靶端把最后坐标减去 \(h(u,a(u))\)，即可假设
	\[
		h(u,0)=0,\qquad h_z(u,0)=0.
	\]
	这里靶坐标变换为 \((u,v)\mapsto(u,v-h(u,a(u)))\)，其逆显式存在。

	对 \(z\) 变量沿线段积分两次，得到
	\[
		h(u,z)=\sum_{i,j=1}^k z_i z_j A_{ij}(u,z),\qquad
		A_{ij}(u,z)=\int_0^1(1-t)\frac{\partial^2h}{\partial z_i\partial z_j}(u,tz)\,\mathrm dt.
	\]
	\(A\) 光滑、对称，且 \(A(0,0)\) 可逆。先对 \(z\) 作一个常系数线性变换，使 \(A(0,0)\) 对角化。随后按\cref{g2:thm:morse-lemma}证明中的配平方步骤，保持全部参数 \(u\) 不变，只改变 \(z\)。第一步例如取
	\[
		w_1=\sqrt{|A_{11}(u,z)|}\left(z_1+\sum_{j>1}\frac{A_{1j}(u,z)}{A_{11}(u,z)}z_j\right),\qquad w_j=z_j\ (j>1).
	\]
	缩小邻域后 \(A_{11}\) 不变号；该变换关于 \(z\) 的微分在原点可逆，故连同未动的 \(u\) 构成局部微分同胚。它分离出一个带固定符号的平方。剩余系数可以依赖 \(u,w_1\)，但仍可保持这些变量不动，继续对其余变量配平方。有限次后得到式\eqref{g2:eq:fold-normal-form}。二次型的惯性在足够小邻域内不变，故负号数固定。反向结论由标准形在核方向上的 Hessian 为对角矩阵 \(2\operatorname{diag}(-I_\lambda,I_{k-\lambda})\) 直接得到。
\end{proof}

这个证明把 Morse 引理的坐标选择做成了随 \(u\) 光滑变化的一族。若不指定一维余核的定向，改变最后一个靶坐标的符号会交换 \(\lambda\) 与 \(k-\lambda\)；因此此时折叠型只确定无序对 \(\{\lambda,k-\lambda\}\)。

\ProseParagraphBreak

在标准形中，临界集是 \(z=0\)，维数为 \(p-1\)，且 \(f\) 在其上的限制是浸入。其切空间由 \(u\) 方向张成，微分的核由 \(z\) 方向张成，二者直和为源切空间。对于平面到平面的映射，模型就是 \((x,y)\mapsto(x,y^2)\)：临界值线的一侧有两个原像，另一侧没有原像，线上有一个合并的原像。这解释了“折叠”所描述的几何变化。

\ProseParagraphBreak

现在把维数固定为 \(n=p=2\)。在秩为 \(1\) 的点附近，同样可取 \(f(x,y)=(x,h(x,y))\)。临界集由 \(h_y=0\) 给出。若 \(h_{yy}\ne0\)，该点是折叠；若 \(h_{yy}=0\)，还需看沿核方向的下一阶变化，以及参数 \(x\) 能否把这种退化展开。

\begin{definition}\label{g2:def:plane-cusp}
	光滑曲面之间的映射在一点处的芽，若左右等价于
	\begin{equation}\label{g2:eq:plane-cusp-normal-form}
		(x,y)\longmapsto(x,y^3+xy),
	\end{equation}
	就称该点为\term[jiandian]{尖点}[cusp]。在一般奇点分类中也称简单尖点；本小节的“尖点”采用这个含义。
\end{definition}

\begin{theorem}[Whitney 尖点标准形]\label{g2:thm:plane-cusp-recognition}
	设 \(h\) 为原点附近的光滑实函数，\(h(0,0)=0\)，且 \(f(x,y)=(x,h(x,y))\)。若在原点满足
	\begin{equation}\label{g2:eq:plane-cusp-criterion}
		h_y=h_{yy}=0,\qquad h_{xy}\ne0,\qquad h_{yyy}\ne0,
	\end{equation}
	则 \(f\) 的原点是尖点。
\end{theorem}

证明使用光滑准备定理的下列有限生成结论：若 \(h(0,y)\) 在 \(y=0\) 恰为三阶零，则任意光滑函数芽 \(r(x,y)\) 都能写成
\begin{equation}\label{g2:eq:cusp-preparation-input}
	r(x,y)=a_0(x,h(x,y))+a_1(x,h(x,y))y+a_2(x,h(x,y))y^2,
\end{equation}
其中 \(a_i\) 为靶平面原点附近的光滑函数芽。这是\cite[第IV章定理3.6、推论3.5及推论3.11]{GolubitskyGuillemin1973}的具体情形，也是该书第VI章定理2.4所用的输入。其有限性可由 \((x,h)=(x,y^3)\) 看出：写 \(h(0,y)=y^3b(y)\)、\(b(0)\ne0\)，再用 \(h(x,y)-h(0,y)\) 可被 \(x\) 整除，便得到上述理想相等；商空间因此以 \(1,y,y^2\) 为基。准备定理把这一商空间中的有限生成提升为实际光滑函数的表达式\eqref{g2:eq:cusp-preparation-input}。下面完整说明这项输入怎样给出坐标标准形。

\begin{proof}
	先在靶端减去 \(h(x,0)\)，使 \(h(x,0)=0\)；这不改变式\eqref{g2:eq:plane-cusp-criterion}，并使 \(h_x(0,0)=0\)。将准备表达式用于 \(r(x,y)=y^3\)，得到
	\[
		y^3=A(x,h)y^2+B(x,h)y+C(x,h).
	\]
	这里以及以下省略 \(h(x,y)\) 的完整变量。令 \(x=0\)，比较关于 \(y\) 的零、一、二阶项，可得 \(A(0,0)=B(0,0)=C(0,0)=0\)。置
	\[
		\eta=y-\frac{A(x,h)}3,\qquad
		P=-B-\frac{A^2}{3},\qquad
		R=C+\frac{AB}{3}+\frac{2A^3}{27}.
	\]
	这些定义把三次式精确改写为
	\begin{equation}\label{g2:eq:cusp-depressed-relation}
		\eta^3+P(x,h)\eta=R(x,h).
	\end{equation}
	需要核验的关键是：\(P,R\) 确实可以作为靶坐标。

	在原点有 \(\eta=0\)、\(\eta_y=1\)。对式\eqref{g2:eq:cusp-depressed-relation}分别作一次 \(x\) 微分、三次 \(y\) 微分，以及一次 \(x\) 和一次 \(y\) 微分。利用 \(h_x=h_y=h_{yy}=0\)，得到
	\[
		R_x=0,\qquad R_v h_{yyy}=6,\qquad P_x=R_v h_{xy},
	\]
	其中 \(v\) 表示 \(P,R\) 的第二个靶变量，全部导数均在原点求值。由假设，\(R_v\ne0\)、\(P_x\ne0\)。因此
	\[
		\det\mathrm D(P,R)_0=P_xR_v\ne0.
	\]
	而源端变换 \((x,y)\mapsto(U,V)=(P(x,h),\eta)\) 的微分矩阵，第一行是 \((P_x,0)\)，第二行的末项是 \(1\)，也可逆。逆函数定理说明两端都是真正的光滑坐标变换。在靶端取 \((U',V')=(P(x,v),R(x,v))\)，式\eqref{g2:eq:cusp-depressed-relation}恰好给出
	\[
		U'=U,\qquad V'=V^3+UV,
	\]
	即所需的尖点标准形。
\end{proof}

条件 \(h_{xy}\ne0\) 还使 \(h_y=0\) 成为一条光滑曲线。沿这条曲线，微分的核由 \(\partial_y\) 张成；在尖点处核恰与临界曲线相切，而 \(h_{yyy}\ne0\) 保证这种相切沿曲线只发生一次简单退化。这与折叠处核横穿临界曲线的情形相对照。

\begin{example}\label{g2:exm:plane-cusp-model}
	对模型 \(f(x,y)=(x,y^3+xy)\)，有
	\[
		\mathrm Df=\begin{pmatrix}1&0\\y&x+3y^2\end{pmatrix}.
	\]
	因此临界曲线与临界值曲线分别参数化为
	\[
		(x,y)=(-3t^2,t),\qquad f(-3t^2,t)=(-3t^2,-2t^3).
	\]
	后者满足 \(4a^3+27b^2=0\)，其中 \((a,b)\) 是靶坐标。\(t\ne0\) 时 \(h_{yy}=6t\ne0\)，都是折叠点；原点满足式\eqref{g2:eq:plane-cusp-criterion}，是唯一的尖点。

	纤维计算也很直接：\(f(x,y)=(a,b)\) 等价于
	\[
		x=a,\qquad y^3+ay-b=0.
	\]
	当 \(a\ge0\) 时，三次函数 \(y^3+ay\) 严格递增，故有唯一实根。当 \(a<0\)，令 \(s=\sqrt{-a/3}\)；极大、极小值分别在 \(-s,s\) 取得，值为 \(2s^3,-2s^3\)。所以 \(|b|<2s^3\) 时有三个不同实根，\(|b|>2s^3\) 时只有一个，\(|b|=2s^3\) 时有两个不同实根，其中一个是二重根。等价地，\(-4a^3-27b^2>0\) 的区域有三个原像，小于零的区域有一个原像。原点的三重根对应尖点本身。这些根在 \((a,b)\to0\) 时都趋于零，故也是映射芽的局部纤维描述。
\end{example}

\begin{figure}[htbp]
	\centering
	% !TeX root = ../Book2.tex
\begin{tikzpicture}[>=stealth,every node/.style={font=\small},x=.85cm,y=.85cm]
	\draw[->,gray] (-3,0)--(1.1,0) node[right] {\(x\)};
	\draw[->,gray] (0,-1.55)--(0,1.55) node[above] {\(y\)};
	\draw[very thick,blue!65!black,domain=-.95:.95,samples=65,variable=\t] plot ({-3*\t*\t},{\t});
	\fill (0,0) circle (2pt);
	\node[below] at (-1.1,-1.65) {源：\((-3t^2,t)\)};
	\draw[->,thick] (1.4,.5)--(2.6,.5) node[midway,above] {\(f\)};
	\begin{scope}[xshift=5.4cm]
		\fill[blue!8] (0,0) plot[domain=0:1,samples=45,variable=\t] ({-3*\t*\t},{2*\t*\t*\t})--(-3,-2) plot[domain=1:0,samples=45,variable=\t] ({-3*\t*\t},{-2*\t*\t*\t})--cycle;
		\draw[->,gray] (-3.2,0)--(1,0) node[right] {\(a\)};
		\draw[->,gray] (0,-2.3)--(0,2.3) node[above] {\(b\)};
		\draw[very thick,blue!65!black,domain=-1:1,samples=80,variable=\t] plot ({-3*\t*\t},{-2*\t*\t*\t});
		\fill (0,0) circle (2pt);
		\node at (-2,.35) {3 个原像};
		\node at (.65,1.05) {1 个原像};
		\node[below] at (-1.2,-2.4) {靶：\((-3t^2,-2t^3)\)};
	\end{scope}
\end{tikzpicture}

	\caption{尖点模型的临界曲线与临界值曲线。临界曲线上除原点外均为折叠点；靶端两支之间有三个局部原像，外部有一个。}
	\label{g2:fig:plane-cusp-curves}
\end{figure}

\subsection{稳定性及维数限制}
\label{g2:subsec:D03:04}

标准形把一张映射在某点附近的行为固定下来。稳定性进一步问：映射受到小扰动时，这种行为是否仍可通过坐标变化恢复？临界点可以移动，临界值也可以移动，所以比较时需要允许两端一起调整。

\begin{definition}\label{g2:def:smooth-map-stability}
	设 \(M,N\) 为无边界光滑流形，\(M\) 紧。若 \(f\in C^\infty(M,N)\) 有一个 Whitney \(C^\infty\) 邻域 \(\mathcal U\)，使每个 \(g\in\mathcal U\) 都满足
	\[
		g=\psi\circ f\circ\phi^{-1}
	\]
	其中 \(\phi\in\operatorname{Diff}(M)\)、\(\psi\in\operatorname{Diff}(N)\)，则称 \(f\) 为\term[wendingyingshe]{稳定映射}[stable mapping]。
\end{definition}

这里的邻近性就是前面采用的光滑映射拓扑；紧源上可以用有限张坐标图中各阶导数的控制描述。定义中的微分同胚是整个 \(M,N\) 上的映射。若只讨论一个奇点在小邻域中的持续存在，并允许其所在点随扰动移动，则是在讨论局部奇点型的稳定性。

\begin{proposition}\label{g2:prop:plane-fold-cusp-persistence}
	平面光滑映射的折叠点在足够小的 \(C^2\) 扰动下，仍在其附近出现同型折叠点；尖点在足够小的 \(C^3\) 扰动下，仍在其附近出现同型尖点。这里可先固定一个足够小的闭坐标邻域，并在更小的内部邻域中寻找扰动后的点。
\end{proposition}
\begin{proof}
	先取 \(f(x,y)=(x,h(x,y))\) 的坐标。对足够接近的 \(g\)，其第一分量仍可与 \(y\) 组成源坐标，所以可把 \(g\) 写成 \((x,\widetilde h(x,y))\)，且 \(\widetilde h\) 在所需阶数上接近 \(h\)。

	折叠情形有 \(h_y(0)=0\)、\(h_{yy}(0)\ne0\)。在 \(x=0\) 的一条短竖直线段上，\(h_y\) 严格单调，且两端异号。足够小的 \(C^2\) 扰动保留这两项性质，故 \(\widetilde h_y\) 有一个零点，且其 \(y\) 导数非零。折叠标准形给出同型局部模型。

	尖点情形考虑 \(G=(h_y,h_{yy})\)。在原点，
	\[
		\mathrm DG_0=\begin{pmatrix}h_{xy}&0\\h_{xyy}&h_{yyy}\end{pmatrix}
	\]
	可逆。一个非退化零点在小的 \(C^1\) 扰动下仍存在：记 \(A=\mathrm DG_0\)，先取闭球 \(B\)，使 \(\|I-A^{-1}\mathrm DG\|<1/3\) 于 \(B\)；再使 \(\widetilde G\) 足够接近 \(G\)，便有 \(\|I-A^{-1}\mathrm D\widetilde G\|<1/2\)，且 \(\|A^{-1}\widetilde G(0)\|\) 小于球半径的一半。映射 \(z\mapsto z-A^{-1}\widetilde G(z)\) 将 \(B\) 映入自身并为压缩，故有唯一不动点，即 \(\widetilde G\) 的零点。所需 \(C^1\) 接近由 \(h\) 的 \(C^3\) 接近保证。这个零点处的 \(\widetilde h_{xy}\) 与 \(\widetilde h_{yyy}\) 仍非零，故尖点识别定理给出相同标准形。
\end{proof}

\begin{theorem}[曲面映射的一般位置]\label{g2:thm:generic-surface-singularities}
	设 \(M,N\) 为有限维、Hausdorff、第二可数的无边界光滑曲面。\(C^\infty(M,N)\) 中存在一个关于 Whitney \(C^\infty\) 拓扑的稠密剩余集，其中每张映射的奇点都为折叠点或尖点。
\end{theorem}

这一结论采用\cite[第VI章定理2.5及第4节]{GolubitskyGuillemin1973}的版本。第一阶秩分层的横截性排除了余维为 \(4\) 的秩零层，并使临界集成为曲线；进一步对二阶喷射的相应层施加横截性，才使核与临界曲线的相切成为上述简单尖点。若 \(M\) 紧，这些映射的临界集是有限个圆的并，尖点只有有限个：临界集闭而光滑；尖点集在其中由相切条件定义为闭集，识别条件又使它离散，紧性便给出有限性。

\ProseParagraphBreak

局部类型还未记录不同临界点映到同一目标位置时的关系。一个具体的全局判据如下。设紧源映射 \(f:M^n\to N^p\)、\(n\ge p\)，只有折叠奇点，令 \(S(f)\) 为其临界集。其限制 \(f|_{S(f)}\) 是浸入。称这个浸入具有\term[zhengchangjiaocha]{正常交叉}[normal crossings]，是指对任意映到同一点 \(b\) 的不同临界点 \(a_1,\ldots,a_r\)，超平面
\[
	H_i=\mathrm df_{a_i}(\mathrm T_{a_i}S(f))\subset\mathrm T_bN
\]
的法向余方向独立；等价地，\(\operatorname{codim}(H_1\cap\cdots\cap H_r)=r\)。特别地，\(r\le p\)。对于曲面到曲面，这要求临界值曲线的自交只有横截双点。

\ProseParagraphBreak

在这个紧源、只有折叠点的范围内，映射 \(f\) 稳定当且仅当 \(f|_{S(f)}\) 有正常交叉。这里调用\cite[第III章定理4.4，结合第1节定理1.5]{GolubitskyGuillemin1973}；证明还使用把无穷小稳定性转为实际坐标变化的定理。当靶为 \(\mathbb R\) 时，折叠就是 Morse 临界点，而正常交叉恰要求所有临界值互异。这使第17章分开临界值的扰动具有了整体稳定性的解释。

\begin{example}\label{g2:exm:morse-equal-values-unstable}
	函数 \(f:S^1\to\mathbb R\)、\(f(\theta)=\cos2\theta\) 的四个临界点都非退化，但临界值只有 \(-1,1\)。在其中一个极大点附近加上一个小的截断常数 \(\varepsilon\)，即可保留四个临界点及其指标，使该极大值变为 \(1+\varepsilon\)。具体扰动可沿用\cref{g2:prop:distinct-morse-critical-values}的构造：截断函数在临界点附近恒为 \(1\)，过渡区原微分有正下界，所以参数足够小时不会产生新临界点。

	扰动后的函数具有三个不同临界值。左右等价会把临界点双射到临界点，并用靶端微分同胚双射其临界值集合，因此保持不同临界值的个数。两函数遂不可能左右等价，而 \(\varepsilon\) 可任意小，故 \(f\) 不稳定。局部的四个 Morse 模型都保留下来了，变化发生在它们的临界值之间。
\end{example}

维数限制也直接进入一般位置结论。曲面映射中的尖点是离散的，因为需要在二维源中同时满足两个独立退化条件；若源维数增加，这些条件的逆像可以成为正维集合，进一步退化也可能出现。一个三维模型为
\begin{equation}\label{g2:eq:swallowtail-model}
	F(a,b,t)=(a,b,t^4+at^2+bt).
\end{equation}
它称为\term[yanweidian]{燕尾点}[swallowtail]的模型。在原点，沿核方向的前一、二、三阶导数都为零，四阶导数不为零；而
\[
	\mathrm d(h_t)_0=\mathrm db,\qquad \mathrm d(h_{tt})_0=2\,\mathrm da,\qquad
	\mathrm d(h_{ttt})_0=24\,\mathrm dt
\]
线性无关。这正是三个退化条件可以横截相交的情形。该三维标准形及其横截识别见\cite[第VII章定理4.1]{GolubitskyGuillemin1973}；\cref{g2:ex:swallowtail-distinction}还用一个有限维代数直接区别它与折叠、尖点。曲面的一般位置结论因而有明确的维数范围；推广到其他源、靶维数时，需要重新检查相应喷射层的余维及局部标准形。

\begin{chaptersummary}
	\begin{itemize}
		\item 沿叶路径的横向运输只依赖端点固定的叶内同伦类。改变横截面或基点使返回芽群共轭；导数只保留线性信息，非平凡返回芽甚至可以与恒等映射具有相同的全部有限阶导数。
		\item 紧叶的有限 holonomy 通过有限覆盖与横向有限群作用给出局部商模型。附近叶的稳定子决定其覆盖次数；返回群平凡时得到饱和的乘积邻域。
		\item 有理和无理环面斜率分别产生圆周叶空间与不可分辨的叶类。Reeb 分量的内部平面卷向紧边界叶，产生非分离商；整体成为圆周上的纤维丛还需要遍及全部叶的紧性与返回条件，或 Reeb 整体稳定中更强的单叶假设。
		\item 喷射保留指定阶数的导数，秩 \(r\) 层的余维为 \((n-r)(p-r)\)。加入带截断的多项式可独立扰动喷射坐标，参数横截性据此给出一般位置；预期维数为负时，相应秩层被避开。
		\item 折叠由核方向的非退化二次型控制，平面尖点还需要三阶项和混合导数。局部模型在小扰动后可以移动并保留，但整体稳定还要控制不同临界点的像如何相遇；较高维数中也会出现新的奇点型。
	\end{itemize}
\end{chaptersummary}

\BookExercises

\begin{exercise}\label{g2:ex:submersion-trivial-holonomy}
	设 \(p:M\to B\) 为光滑淹没，以各纤维的连通分支为叶。证明任意叶的横截面返回映射芽群平凡。更具体地，对叶内路径 \(\alpha:x\to y\) 和端点横截面 \(T_x,T_y\)，把 \(h_\alpha\) 用 \(p\) 在这两张横截面上的限制表示出来。
\end{exercise}
\begin{solution}
	淹没的常秩局部坐标把 \(p\) 写成横向坐标投影，故其纤维确实给出正则叶状结构。横截条件说明 \(\mathrm d(p|_{T_x})_x\) 和 \(\mathrm d(p|_{T_y})_y\) 都是同构；缩小横截面后，两者是到 \(b=p(x)=p(y)\) 附近开集的局部微分同胚。沿每张叶片运动保持 \(p\) 的值，有限次接续后仍有
	\[
		p\circ h_\alpha=p|_{T_x}.
	\]
	因此在足够小的共同邻域上，
	\[
		h_\alpha=(p|_{T_y})^{-1}\circ(p|_{T_x}).
	\]
	对于闭路径使用同一横截面 \(T_x=T_y\)，右端就是恒等芽。所以每条闭路径的返回芽都平凡，无需假设这张叶单连通或紧。公式也说明返回平凡的原因：一个全局横向参数 \(p\) 在整个运输过程中保持不变。
\end{solution}

\begin{exercise}\label{g2:ex:arbitrary-order-parabolic-holonomy}
	设整数 \(k\geq2\)，在 \((\mathbb R/\mathbb Z)\times\mathbb R\) 上考虑向量场 \(X_k=\partial_\theta+r^k\partial_r\)。求中心圆沿 \(\theta\) 正向绕行一周的返回芽及其各次迭代。证明其返回群无限，而返回映射与恒等映射在 \(0\) 处具有相同的 \(k-1\) 阶导数数据。
\end{exercise}
\begin{solution}
	横向初值问题为 \(r'=r^k\)、\(r(0)=s\)。对 \(s\ne0\) 分离变量，得到
	\[
		r(\theta)=s\bigl(1-(k-1)\theta s^{k-1}\bigr)^{-1/(k-1)}.
	\]
	当 \(s\) 足够接近零且 \(\theta\) 落在给定有界区间内时，括号内为正，公式光滑地延拓到 \(s=0\)，并直接满足初值方程。因此
	\[
		h(s)=s\bigl(1-(k-1)s^{k-1}\bigr)^{-1/(k-1)},\qquad
		h^n(s)=s\bigl(1-(k-1)n s^{k-1}\bigr)^{-1/(k-1)}
	\]
	对每个 \(n\in\mathbb Z\) 都成立为芽的等式。负 \(n\) 对应沿反方向行进；每个固定 \(n\) 只需重新缩小 \(s\) 的邻域。

	在 \(u=0\) 处展开 \((1-u)^{-1/(k-1)}\)，得到
	\[
		h^n(s)=s+n s^k+O(s^{2k-1}).
	\]
	所以 \(n\ne m\) 时，两返回芽的 \(k\) 阶系数不同，全部迭代两两不同，返回群为无限循环群。对于 \(n=1\)，\(h(s)-s\) 从 \(k\) 次项开始，故直到 \(k-1\) 阶的导数都与恒等映射相同。任给一个预先固定的有限阶数，只要选更大的 \(k\)，就能让该阶数以内的导数数据完全看不到这个非平凡返回芽。
\end{solution}

\begin{exercise}\label{g2:ex:cyclic-holonomy-disk-model}
	设 \(m\geq2\)、\(\widehat L=\mathbb R/\mathbb Z\)，取复平面中以零为心的开圆盘 \(B\)。令循环群 \(C_m\) 的生成元作用于 \(\widehat L\times B\) 为
	\[
		(\theta,z)\longmapsto\bigl(\theta+1/m,e^{2\pi i/m}z\bigr).
	\]
	在商上取乘积叶状结构。计算中心叶与非中心叶的返回芽群、各叶到中心叶的覆盖次数及叶空间。解释为什么叶空间即使同胚于圆盘，这个叶状邻域也不能与中心圆的乘积叶状邻域等同。
\end{exercise}
\begin{solution}
	群在 \(\widehat L\) 上作用自由，在 \(B\) 上作用忠实，故可用\cref{g2:prop:finite-holonomy-model}。在 \(z=0\) 处稳定子为整个 \(C_m\)，中心叶为 \(\widehat L/C_m\)，其返回芽群为 \(C_m\)，到自身的投影是一层覆盖。若 \(z\ne0\)，等式 \(e^{2\pi i j/m}z=z\) 强迫 \(m\mid j\)，所以稳定子平凡。这些叶微分同胚于 \(\widehat L\)，返回芽群平凡，到中心叶的投影为 \(m\) 层覆盖。全部叶都是圆，但它们在邻域中绕行的次数不同。

	叶空间是 \(B/C_m\)。若 \(B\) 的半径为 \(R\)，映射 \(z\mapsto z^m\) 的像是半径 \(R^m\) 的开圆盘，两个点有相同像当且仅当它们相差一个 \(m\) 次单位根。它还是开映射：在非零点附近由逆函数定理成立，在零点附近则直接由小圆盘的像仍为小圆盘得出。因此诱导的连续双射 \(B/C_m\to\{w:|w|<R^m\}\) 是开映射，从而是同胚。

	这张商圆盘没有记录中心叶的有限返回群。若存在把中心叶对应到中心叶的叶状微分同胚通向乘积邻域，\cref{g2:prop:holonomy-conjugacy}及叶状微分同胚对叶片运输的保持性将使两张中心叶的返回群同构；乘积情形为平凡群，本题却为 \(C_m\)，矛盾。商映射 \(z\mapsto z^m\) 在零点的微分为零，也提醒我们不能把叶空间的这个拓扑坐标直接当作原流形上的淹没横向坐标。
\end{solution}

\begin{exercise}\label{g2:ex:torus-foliation-periods}
	在 \(T^2\) 上取闭一形式 \(\omega_\alpha=\mathrm dy-\alpha\mathrm dx\)。求它在全部闭路上的积分值组成的子群。若 \(\alpha=p/q\) 为既约分数，说明 \(q\omega_\alpha\) 如何给出叶空间上的圆周坐标。若 \(\alpha\) 无理，证明任何沿叶常值的连续实值函数都为常数，并解释这与 \(\omega_\alpha\) 处处非零为何相容。
\end{exercise}
\begin{solution}
	闭路提升到 \(\mathbb R^2\) 后，两个端点之差为某个整数向量 \((m,n)\)，故积分为 \(n-\alpha m\)。反过来每个整数向量均由一条闭路实现，所以周期群为 \(\mathbb Z+\alpha\mathbb Z\)。若 \(\alpha=p/q\)，互素性给出这个群恰为 \(q^{-1}\mathbb Z\)。函数
	\[
		P([x,y])=[qy-px]\in\mathbb R/\mathbb Z
	\]
	良定义，且 \(P^*\mathrm dt=q\omega_\alpha\)。正文已证其纤维恰为连通的叶，所以它就是所需圆周坐标。

	若 \(\alpha\) 无理，正文中的小步长逼近也说明 \(\mathbb Z+\alpha\mathbb Z\) 在 \(\mathbb R\) 中稠密。若 \(u:T^2\to\mathbb R\) 连续且沿叶常值，固定一条稠密叶及其中一点 \(z\)。任意 \(w\in T^2\) 都是这条叶上一列点的极限，连续性给出 \(u(w)=u(z)\)，故 \(u\) 常值。形式 \(\omega_\alpha\) 有局部原函数 \(y-\alpha x\)，但绕闭路后的值会相差周期；非零闭形式并不自动给出一个全局实值原函数。无理时也没有常数倍把全部周期变成整数，因而不能用上述积分取模方式为这些叶构造圆周纤维坐标。
\end{solution}

\begin{exercise}\label{g2:ex:flat-reeb-holonomy}
	令
	\[
		b(s)=\begin{cases}e^{-1/s^2},&s>0,\\0,&s\le0,\end{cases}
	\]
	在实心环面上将式\eqref{g2:eq:reeb-solid-torus-form}改为 \(\alpha_b=\mathrm d\rho+b(1-\rho)\mathrm dt\)。证明仍得到边界为紧叶、内部叶为开圆盘的叶状结构。设边界纵圈的返回芽在向内坐标 \(s=1-\rho\) 中为 \(h\)，证明 \(h\) 非恒等且有无限阶，但它与恒等映射的任意有限阶导数在零点都相同。
\end{exercise}
\begin{solution}
	函数 \(b\) 光滑且在零点平坦；内部 \(b(1-\rho)>0\)。在中心圆上 \(\alpha_b=b(1)\mathrm dt\ne0\)，其余点有非零径向分量，且 \(\alpha_b\wedge\mathrm d\alpha_b=0\)。内部叶由
	\[
		t+F(\rho)=c\pmod1,\qquad
		F(\rho)=\int_0^\rho\frac{\mathrm du}{b(1-u)}
	\]
	给出。\(F\) 从零严格增至无穷；例如靠近1时 \(b(1-u)\le(1-u)^2\)，故积分发散。每条叶都是开圆盘上的图像，并像正文一样聚集到全部边界。边界上形式等于 \(\mathrm d\rho\)，仍是一张环面叶。

	沿纵圈，\(s\) 满足 \(s'=b(s)\)，所以 \(h\) 是这条方程的时间1映射。对小的 \(s>0\)，有 \(h(s)>s\)；其正次迭代在适当小邻域中严格递增，因此没有正次迭代是恒等芽。对 \(s\le0\)，方程平稳，\(h(s)=s\)；这也给出了跨边界领圈中的光滑延拓。

	再估计正侧。对每个整数 \(N\)，存在 \(C_N\) 使 \(0\le b(s)\le C_Ns^N\)。先用 \(b(s)\le Cs^2\) 说明从小的初值 \(s\) 出发，解在时间 \([0,1]\) 内保持在 \([s,2s]\)：若第一次到达 \(2s\) 的时间为 \(t_0\le1\)，则
	\[
		s=\int_0^{t_0}b(s(t))\,\mathrm dt\le4Cs^2,
	\]
	对 \(s<1/(4C)\) 矛盾。于是对任意 \(N\) 都有
	\[
		0<h(s)-s=\int_0^1b(s(t))\,\mathrm dt\le C_N(2s)^N.
	\]
	光滑函数 \(h(s)-s\) 因而在零点平坦：若某阶导数首次非零，Taylor 公式会给出该阶的非零主项，与更高次的上述界矛盾。所以所有有限阶导数均与恒等映射相同，尽管映射芽本身不同。这也说明有限阶数据能够辨认什么，取决于所讨论的对象和附加非退化条件。
\end{solution}

\begin{exercise}\label{g2:ex:foliation-mapping-torus}
	令 \(A:S^2\to S^2\) 为对径映射 \(A(x)=-x\)，考虑映射环面 \(M_A\)。证明其纤维叶状结构横向可定向，全部叶为 \(S^2\)，且 holonomy 群平凡；再证明 \(M_A\) 不可定向，因而不微分同胚于 \(S^2\times S^1\)。说明这里的全局粘合与叶内 holonomy 分别记录什么。
\end{exercise}
\begin{solution}
	映射环面有光滑投影 \(p:M_A\to S^1\)，圆周非零一形式的拉回 \(p^*\mathrm dt\) 给出其法方向的定向。每条纤维都是 \(S^2\)，既紧又单连通；沿纤维内环路的 holonomy 因同伦不变性而平凡。

	对径映射的次数为 \((-1)^3=-1\)：它由球体上的线性映射 \(-I:\mathbb R^3\to\mathbb R^3\) 限制得到，这张线性映射反转三维定向，同时将外法向送到外法向，所以边界球面的定向也反转。若 \(M_A\) 可定向，结合已定向的基底方向，便给每条纤维一个连续选择的定向。把它拉回 \(S^2\times[0,1]\)，沿区间连续变化不可能翻转球面定向；但端点粘合要求这个定向被 \(A\) 保持，与 \(A\) 反定向矛盾。因此 \(M_A\) 不可定向，而乘积 \(S^2\times S^1\) 可定向，两者不微分同胚。

	沿一条纤维内的环路移动邻近叶的交点，所得芽为恒等；沿基底绕一圈运输整条纤维，所得粘合却为 \(A\)。两个过程的道路所在空间和被移动对象不同，所以前者平凡不会消除后者的扭转。
\end{solution}

\begin{exercise}\label{g2:ex:jet-hessian-coordinate}
	设 \(f:(\mathbb R^n,0)\to\mathbb R\) 光滑，\(\varphi:(\mathbb R^n,0)\to(\mathbb R^n,0)\) 为局部微分同胚。推导 \(f\circ\varphi\) 的二阶导数变换式，说明临界点处 Hessian 的非退化性与正负惯性指数为何不依赖坐标。在一维用 \(f(x)=x\)、\(\varphi(t)=t+t^2\) 说明一般点的二阶导数不是这样变换的张量。
\end{exercise}
\begin{solution}
	对 \(\mathrm D(f\circ\varphi)=\mathrm Df_{\varphi}\circ\mathrm D\varphi\) 再求一次导数，乘积法则给出
	\[
		\mathrm D^2(f\circ\varphi)_0(v,w)
		=\mathrm D^2f_0(\mathrm D\varphi_0v,\mathrm D\varphi_0w)
		+\mathrm Df_0\bigl(\mathrm D^2\varphi_0(v,w)\bigr).
	\]
	若 \(\mathrm Df_0=0\)，记 \(A=\mathrm D\varphi_0\)，两组 Hessian 矩阵便满足 \(H'=A^{\mathsf T}HA\)。\(A\) 可逆，所以核的维数保持；将一个正定或负定子空间用 \(A^{-1}\) 拉回，维数也保持。两个方向分别应用这一观察，得到最大正定、最大负定子空间的维数相同，即惯性指数相同。

	一维例子中，原函数 \(f''(0)=0\)，却有 \((f\circ\varphi)''(0)=2\)，差别来自 \(f'(0)\varphi''(0)=2\)。因此在非临界点，遗漏后一项会把同一个函数的坐标计算误认为两种 Hessian。
\end{solution}

\begin{exercise}\label{g2:ex:jet-polynomial-parameters}
	对 \(f:\mathbb R^2\to\mathbb R\)，在零点附近取截断函数恒为 \(1\)，并作扰动
	\[
		f_a(x,y)=f(x,y)+a_0+a_1x+a_2y+a_3x^2+a_4xy+a_5y^2.
	\]
	计算参数对零点二阶喷射各坐标的作用，证明全部二阶喷射数据都能独立改变。再说明若只允许加入常数，为何不能用同样办法改善函数临界点的 Hessian。
\end{exercise}
\begin{solution}
	按 \((f,f_x,f_y,f_{xx},f_{xy},f_{yy})\) 排列二阶喷射数据，扰动增加的向量为
	\[
		(a_0,a_1,a_2,2a_3,a_4,2a_5).
	\]
	这是 \(\mathbb R^6\) 上的可逆线性映射，因此任何指定的六项小变化都能由唯一参数给出。在其他源点处，将二次多项式改写为以该点为中心的 Taylor 多项式，仍得到可逆的系数变换。

	若只加入常数，全部正阶导数都不变，临界点集合及每个点的 Hessian 也都不变。常数参数可以移动函数值，却不能把退化临界点变成非退化临界点。这正是喷射扰动需要包含相应阶数多项式项的原因。
\end{solution}

\begin{exercise}\label{g2:ex:generic-rank-dimensions}
	设 \(M^2,N^p\) 是光滑无边流形，\(f:M\to N\) 的一阶喷射横截于全部秩层。分别在 \(p=2,3,4\) 时列出可能出现的秩及其点集维数。特别说明 \(p=4\) 时能推出浸入，但还缺少什么信息才能判断嵌入。
\end{exercise}
\begin{solution}
	秩 \(r\) 层的预期维数为 \(2-(2-r)(p-r)\)。当 \(r=0\) 时，它为 \(2-2p<0\)，三种情形都无秩零点。秩二是满源秩，点集开放。

	当 \(p=2\) 时，秩一点集维数为 \(2-1=1\)，可为临界曲线或空集。当 \(p=3\) 时，它的维数为 \(2-2=0\)，所以秩一点孤立；若源紧，由于已无秩零点，秩一点集闭，因而有限。当 \(p=4\) 时，预期维数为 \(2-3=-1\)，没有秩一点，故处处秩二，\(f\) 为浸入。

	浸入只控制每个点附近的微分；不同源点仍可能具有同一像。因此还需检查单射性与像的拓扑，例如紧源时单射浸入才由紧性和 Hausdorff 性成为嵌入。单点喷射条件没有同时比较两个不同源点，这类问题需要分别研究多点数据。
\end{solution}

\begin{exercise}\label{g2:ex:explicit-fold-completion}
	考虑 \(f:\mathbb R^3\to\mathbb R^2\)、\(f(x,y,z)=(x,y^2-z^2+xy+2xz)\)。求临界集与临界值集，写出使 \(f\) 化为折叠标准形的源、靶坐标。它与 \(g(x,y,z)=(x,y^2+z^2)\) 在原点是否左右等价？
\end{exercise}
\begin{solution}
	秩下降条件是 \(2y+x=0\)、\(-2z+2x=0\)。所以临界集为 \((x,-x/2,x)\)，临界值集为 \((x,3x^2/4)\)。配平方得到
	\[
		y^2-z^2+xy+2xz=(y+x/2)^2-(z-x)^2+3x^2/4.
	\]
	在源取 \(u=x\)、\(v=y+x/2\)、\(w=z-x\)，在靶取 \(U=a\)、\(V=b-3a^2/4\)。这两张坐标变换都具有显式光滑逆，并使 \(f\) 成为 \((u,v,w)\mapsto(u,v^2-w^2)\)。

	其核方向二次型惯性为 \((1,1)\)；\(g\) 的二次型为正定，惯性为 \((2,0)\)。左右等价只能对核作可逆线性换基，并对一维余核乘一个非零实数，因此保持惯性对，至多交换正、负数目。这两个无序惯性对不同，故不等价。也可以从纤维看出差别：\(g\) 的临界值线一侧没有原像，而不定折叠的二次型可以取任意正、负小值。
\end{solution}

\begin{exercise}\label{g2:ex:translated-cusp}
	对 \(\varepsilon\in\mathbb R\)，令 \(f_\varepsilon(x,y)=(x,y^3+xy+\varepsilon y^2)\)。求其尖点位置，并给出把它精确化为 \((X,Y)\mapsto(X,Y^3+XY)\) 的源、靶坐标。尖点的移动是否表示奇点型消失？
\end{exercise}
\begin{solution}
	记第二分量为 \(h\)。方程 \(h_y=h_{yy}=0\) 为
	\[
		3y^2+x+2\varepsilon y=0,\qquad 6y+2\varepsilon=0.
	\]
	解为 \((x,y)=(\varepsilon^2/3,-\varepsilon/3)\)，其像是 \((\varepsilon^2/3,-\varepsilon^3/27)\)。又因 \(h_{xy}=1\)、\(h_{yyy}=6\)，该点满足尖点识别条件。

	取源坐标
	\[
		X=x-\varepsilon^2/3,\qquad Y=y+\varepsilon/3.
	\]
	代入后，
	\[
		h=Y^3+(x-\varepsilon^2/3)Y+2\varepsilon^3/27-\varepsilon x/3.
	\]
	所以在靶平面取 \(A=a-\varepsilon^2/3\)、\(B=b-2\varepsilon^3/27+\varepsilon a/3\)，映射恰成为 \((X,Y)\mapsto(X,Y^3+XY)\)。这些都是可逆的仿射变换。尖点离开了旧原点，却在新位置保留同一个模型；局部稳定比较允许这样的移动。
\end{solution}

\begin{exercise}\label{g2:ex:degenerate-cubic-unfolding}
	令 \(f_0(x,y)=(x,y^3)\)，并令 \(f_\varepsilon(x,y)=(x,y^3+\varepsilon y)\)。比较 \(\varepsilon>0\)、\(\varepsilon=0\)、\(\varepsilon<0\) 时的临界集及奇点型。由此说明尖点识别条件中的 \(h_{xy}\ne0\) 为什么不能删掉。
\end{exercise}
\begin{solution}
	Jacobi 行列式为 \(3y^2+\varepsilon\)。当 \(\varepsilon>0\) 时处处非零，所以 \(f_\varepsilon\) 处处为局部微分同胚。当 \(\varepsilon<0\) 时，临界集为两条直线
	\[
		y=\pm\sqrt{-\varepsilon/3}.
	\]
	这些点有 \(h_{yy}=6y\ne0\)，故都是折叠点。令 \(s=\sqrt{-\varepsilon/3}\)，两条临界值线分别为 \(b=-2s^3\) 与 \(b=2s^3\)。

	当 \(\varepsilon=0\) 时，临界集为 \(y=0\)，其每个点都满足 \(h_y=h_{yy}=0\)、\(h_{yyy}=6\)，但 \(h_{xy}=0\)。该映射没有尖点模型的参数展开；任意小的正 \(\varepsilon\) 扰动就能把附近全部奇点消去，任意小的负扰动则产生两条折叠曲线。左右等价保持临界集，所以这些变化不能由坐标变换消除。非零三阶项本身因而不够，混合导数条件负责使退化随横向参数发生简单变化。
\end{solution}

\begin{exercise}\label{g2:ex:swallowtail-distinction}
	对光滑映射芽 \(f:(\mathbb R^n,0)\to(\mathbb R^p,0)\)，令 \(\mathcal E_n\) 为原点处光滑实函数芽的环，并定义 \(A_f=\mathcal E_n/(f_1,\ldots,f_p)\)。证明左右等价诱导 \(A_f\) 的环同构。计算三张 \(\mathbb R^3\to\mathbb R^3\) 映射
	\[
		(a,b,t)\mapsto(a,b,t^2),\qquad
		(a,b,t)\mapsto(a,b,t^3+at),\qquad
		(a,b,t)\mapsto(a,b,t^4+at^2+bt)
	\]
	对应代数的实维数，说明燕尾模型不能左右等价于前两个模型。
\end{exercise}
\begin{solution}
	源坐标变换 \(\phi\) 通过函数拉回给出 \(\mathcal E_n\) 的环自同构，因此把分量生成的理想及其商环一起带过去。若 \(\psi\) 是固定原点的靶微分同胚，则沿线段积分可写
	\[
		\psi_i(y)=\sum_j y_j c_{ij}(y),\qquad
		c_{ij}(y)=\int_0^1\frac{\partial\psi_i}{\partial y_j}(sy)\,\mathrm ds.
	\]
	所以每个 \(\psi_i\circ f\) 都属于 \((f_1,\ldots,f_p)\)。对 \(\psi^{-1}\) 作同样计算，得到相反包含，两组分量生成同一个理想。两类坐标变化合起来便证明左右等价不变性。

	题中每个理想都含 \(a,b\)，所以依次成为
	\[
		A_2\cong\mathcal E_1/(t^2),\qquad
		A_3\cong\mathcal E_1/(t^3),\qquad
		A_4\cong\mathcal E_1/(t^4).
	\]
	带积分余项的 Taylor 公式说明每个光滑函数模 \(t^k\) 都可由 \(1,t,\ldots,t^{k-1}\) 表示；若这些单项式的线性组合属于 \((t^k)\)，在零处取前 \(k-1\) 阶导数便知所有系数均为零。因此三个代数的实维数分别为 \(2,3,4\)。环同构保持实向量空间维数，燕尾模型遂与折叠、沿额外参数延拓的尖点模型均不等价。
\end{solution}

% !TeX root = ../../Book2.tex
% 纲要标题：## 附录 E　有理同伦与微分分次代数
% 纲要位置：几何学计划纲要.md，第 2085 行。
% 正文状态：未撰写；本文件目前是框架。
% 证明范围与先修：固定单连通、有限型 CW 空间及其有理模型，使用本卷第4、7、9、12章的乘积、de Rham 比较、同伦群与 Postnikov 语言。PL 多项式微分形式模型在有理数域上构造；光滑流形的 de Rham 代数在实数域上构造，经系数扩张后比较，不能将实模型直接说成有理模型。完成选定球面、射影空间与 Massey 乘积算例；一般 Sullivan 存在、唯一性与有理同伦对应定理准确外引，Kähler 形式性注明第四卷背景及外部输入。不扩展为一般 operad 或高阶范畴课程。

% 纲要内容开始
% 
% > 本附录证明范围与外部输入：固定单连通有限型 CW 空间，依赖本卷乘积、de Rham 比较、同伦群和 Postnikov 语言。严格区分有理数域上的 PL 多项式微分形式模型与实数域上的 de Rham 代数，经系数扩张后比较。完成球面、射影空间与 Massey 乘积代表算例；一般 Sullivan 存在、唯一性及有理同伦对应定理准确外引，Kähler 形式性另注明外部输入。不扩展为一般 operad 或高阶范畴课程。
% 
% * **分次交换微分代数与模型**
%   * 分次交换代数与微分的符号
%   * 微分代数态射与拟同构
%   * 有理多项式微分形式模型
%   * 光滑流形的实 de Rham 代数
%   * 系数扩张与模型比较
% 
% * **最小模型与有理同伦**
%   * 自由分次交换代数与生成元
%   * 最小性、滤过与逐阶构造
%   * 单连通有限型空间的最小模型
%   * 生成元与有理同伦群
%   * Sullivan 主定理及其使用条件
% 
% * **球面与射影空间计算**
%   * 奇维球面的模型
%   * 偶维球面的模型
%   * 复射影空间的模型
%   * 乘积空间与映射的模型计算
% 
% * **Massey 乘积、形式性与几何**
%   * 三重 Massey 乘积与定义条件
%   * 代表选择与不定性
%   * 形式性与上同调代数模型
%   * 可核验的非形式性障碍
%   * Kähler 流形形式性的外部输入
% 
% ---
% 
% 纲要内容结束

\chapter{有理同伦与微分分次代数}
\label{g2:app:E}
\BookChapterNumber{g2:app:E}

待撰写

\section{分次交换微分代数与模型}
\label{g2:sec:E01}
待撰写

\subsection{分次交换代数与微分的符号}
\label{g2:subsec:E01:01}
待撰写

\subsection{微分代数态射与拟同构}
\label{g2:subsec:E01:02}
待撰写

\subsection{有理多项式微分形式模型}
\label{g2:subsec:E01:03}
待撰写

\subsection{光滑流形的实 de Rham 代数}
\label{g2:subsec:E01:04}
待撰写

\subsection{系数扩张与模型比较}
\label{g2:subsec:E01:05}
待撰写

\section{最小模型与有理同伦}
\label{g2:sec:E02}
待撰写

\subsection{自由分次交换代数与生成元}
\label{g2:subsec:E02:01}
待撰写

\subsection{最小性、滤过与逐阶构造}
\label{g2:subsec:E02:02}
待撰写

\subsection{单连通有限型空间的最小模型}
\label{g2:subsec:E02:03}
待撰写

\subsection{生成元与有理同伦群}
\label{g2:subsec:E02:04}
待撰写

\subsection{Sullivan 主定理及其使用条件}
\label{g2:subsec:E02:05}
待撰写

\section{球面与射影空间计算}
\label{g2:sec:E03}
待撰写

\subsection{奇维球面的模型}
\label{g2:subsec:E03:01}
待撰写

\subsection{偶维球面的模型}
\label{g2:subsec:E03:02}
待撰写

\subsection{复射影空间的模型}
\label{g2:subsec:E03:03}
待撰写

\subsection{乘积空间与映射的模型计算}
\label{g2:subsec:E03:04}
待撰写

\section{Massey 乘积、形式性与几何}
\label{g2:sec:E04}
待撰写

\subsection{三重 Massey 乘积与定义条件}
\label{g2:subsec:E04:01}
待撰写

\subsection{代表选择与不定性}
\label{g2:subsec:E04:02}
待撰写

\subsection{形式性与上同调代数模型}
\label{g2:subsec:E04:03}
待撰写

\subsection{可核验的非形式性障碍}
\label{g2:subsec:E04:04}
待撰写

\subsection{Kähler 流形形式性的外部输入}
\label{g2:subsec:E04:05}
待撰写



\backmatter
\cleardoublepage
\pagenumbering{roman}
\bookmarksetup{startatroot}
% !TeX root = ../Book2.tex
\cleardoublepage
% !TeX root = ../../Book2.tex
\chapter*{后记}
\addcontentsline{toc}{chapter}{后记}
待撰写

\cleardoublepage
% !TeX root = ../../Book2.tex
\PrintBookBibliography
% 仅在核读并引用真实资料后补入书目。

\cleardoublepage
% !TeX root = ../../Book2.tex
\PrintBookIndex{chinese}{中文索引}

\cleardoublepage
% !TeX root = ../../Book2.tex
\PrintBookIndex{foreign}{外文索引}

\cleardoublepage
% !TeX root = ../../Book2.tex
\PrintBookIndex{symbols}{符号索引}


\end{document}
